% !TeX program = pdflatex
% Wykład (rozdziały 1--32) + ramkowany kod numeryczny.
% Zadania i audyty znajdują się w oddzielnych plikach TEX/PDF.
% Kompilacja: latexmk -cd -pdf books/pl/partial_differential_equations.tex
\pdfobjcompresslevel=0
\documentclass[11pt,a4paper,oneside]{book}

% Wspólna preambuła skryptu „Równania różniczkowe cząstkowe”.
% Preambula włączona bezpośrednio do jednego pliku.
\usepackage[T1]{fontenc}
\usepackage[utf8]{inputenc}
\usepackage[polish]{babel}
\usepackage{lmodern}
\usepackage{amsmath,amssymb,amsthm,mathtools}
\usepackage{esint}
\usepackage{microtype}
\usepackage[a4paper,margin=2.7cm]{geometry}
\usepackage{enumitem}
\usepackage{xcolor}
\usepackage{graphicx}
\graphicspath{{../../assets/pl/partial_differential_equations/}}
\usepackage{tikz}
\usetikzlibrary{arrows.meta,calc,positioning}
\usepackage[unicode,colorlinks=true,linkcolor=blue!45!black,urlcolor=blue!45!black]{hyperref}

\setlength{\parindent}{0pt}
\setlength{\parskip}{4pt}
\setlength{\emergencystretch}{2em}
\setlist[enumerate]{topsep=4pt,itemsep=2pt,leftmargin=*}
\setlist[itemize]{topsep=4pt,itemsep=2pt,leftmargin=*}
\numberwithin{equation}{chapter}
\theoremstyle{plain}
\newtheorem{twierdzenie}{Twierdzenie}[chapter]
\newtheorem{stwierdzenie}[twierdzenie]{Stwierdzenie}
\newtheorem{lemat}[twierdzenie]{Lemat}
\newtheorem{wniosek}[twierdzenie]{Wniosek}
\theoremstyle{definition}
\newtheorem{definicja}[twierdzenie]{Definicja}
\newtheorem{przyklad}[twierdzenie]{Przykład}
\newtheorem{kontrprzyklad}[twierdzenie]{Kontrprzykład}
\theoremstyle{remark}
\newtheorem{uwaga}[twierdzenie]{Uwaga}

\newcommand{\R}{\mathbb{R}}
\newcommand{\C}{\mathbb{C}}
\newcommand{\N}{\mathbb{N}}
\newcommand{\T}{\mathbb{T}}
\newcommand{\eps}{\varepsilon}
\newcommand{\pd}{\partial}
\DeclareMathOperator{\supp}{supp}
\DeclareMathOperator{\tr}{tr}
\DeclareMathOperator{\diver}{div}

\title{Równania różniczkowe cząstkowe\\\large Tom I. Podstawy teorii PDE i rozwiązań słabych}
\author{}
\date{19 listopada 2026}

% BEGIN GLOBAL CODE LISTINGS (generated from latex/pl/code-listings.tex)
\usepackage{xcolor}
\usepackage{listings}
\definecolor{codebg}{RGB}{248,249,251}
\definecolor{codeborder}{RGB}{180,191,203}
\definecolor{codecomment}{RGB}{61,113,79}
\definecolor{codekeyword}{RGB}{41,78,132}
\lstdefinelanguage{Julia}{
  morekeywords={abstract,baremodule,begin,break,catch,const,continue,do,else,elseif,
    end,export,false,finally,for,function,global,if,import,in,isa,let,local,macro,
    module,mutable,primitive,quote,return,struct,true,try,type,using,where,while,
    missing,nothing,println,print},
  sensitive=true,morecomment=[l]{\#},morecomment=[n]{\#=}{=\#},
  morestring=[b]",
}
\lstdefinestyle{skrypt}{
  backgroundcolor=\color{codebg},frame=single,rulecolor=\color{codeborder},
  framesep=5pt,frameround=tttt,framerule=.5pt,
  basicstyle=\ttfamily\footnotesize,keywordstyle=\color{codekeyword}\bfseries,
  commentstyle=\color{codecomment}\itshape,stringstyle=\color{purple!65!black},
  numbers=left,numberstyle=\tiny\color{gray!75!black},numbersep=8pt,
  stepnumber=1,breaklines=true,breakatwhitespace=false,
  columns=fullflexible,keepspaces=true,showstringspaces=false,tabsize=4,
  linewidth=.90\linewidth,captionpos=b,aboveskip=11pt,belowskip=12pt,
  xleftmargin=12pt,xrightmargin=2pt,
}
\lstdefinestyle{pythonpde}{style=skrypt,language=Python}
\lstdefinestyle{python}{style=skrypt,language=Python}
\lstdefinestyle{julia}{style=skrypt,language=Julia}
\lstset{style=skrypt}
\renewcommand{\lstlistingname}{Kod}
\renewcommand{\lstlistlistingname}{Spis kodów}
% END GLOBAL CODE LISTINGS
\begin{document}
\frontmatter
\hypersetup{pageanchor=false}
\maketitle
\hypersetup{pageanchor=true}
\tableofcontents
\chapter*{Od redakcji}
\addcontentsline{toc}{chapter}{Od redakcji}
Tom I zawiera ciągły wykład z teorii PDE: od podstaw i metody
charakterystyk, przez równania Laplace'a i Poissona oraz
równanie ciepła, po równanie falowe, metody Fouriera, teorię dystrybucji, przestrzenie Sobolewa, metody wariacyjne, teorię spektralną i eliptyczne zagadnienia brzegowe. Zachowano wspólną
numerację i odsyłacze. Zbiór zadań do rozdziałów 1--32
oraz szczegółowy audyt matematyczno-redakcyjny są
przygotowane jako osobne dokumenty, niezależne od wykładu.
W tekście przez \emph{obszar} rozumiemy niepusty otwarty
spójny podzbiór przestrzeni euklidesowej. Określenie
\emph{dziedzina} rezerwujemy dla dziedziny funkcji lub operatora;
angielskie \emph{field} w algebrze oznacza \emph{ciało},
a \emph{vector field} w geometrii — \emph{pole wektorowe}.
W nazwach matematycznych stosujemy poprawną odmianę nazwisk:
równanie Hamiltona--Jacobiego, nierówność Harnacka,
twierdzenie Laxa--Milgrama, oszacowanie Schaudera,
metoda Galerkina. Używamy \emph{iloczynu tensorowego},
a nie sztucznej nazwy odczasownikowej.
Każda całka oznaczona zawiera obszar całkowania i różniczkę
miary, na przykład $\int_\Omega f(x)\,\mathrm d\lambda_d(x)$;
na torusie $\T^d$ używamy miary Lebesgue'a o ustalonej
normalizacji, określanej przy pierwszym użyciu.
Kod numeryczny jest zapisany bezpośrednio w obramowanych
listingach tego skryptu; do kompilacji książki nie potrzeba
osobnych plików Python ani wcześniej wydawanych archiwów.
Aby uruchomić przykład, można przenieść dany listing do pliku
źródłowego w Pythonie i zainstalować biblioteki wskazane w kodzie.
W przykładach wizualizacji korzystamy z Plotly;
interaktywne wykresy można następnie zapisać jako HTML.
Rysunki wiru Taylora--Greena wykonano bezpośrednio w TikZ.

\subsection*{Jak czytać rozdziały w porządku logicznym}
Pierwsze przykłady PDE i charakterystyki można czytać niezależnie,
ale zaawansowana część teorii ma świadomy porządek zależności.
Szeregi Fouriera na torusie przedstawiamy przed teorią dystrybucji;
następnie wprowadzamy dystrybucje, transformatę Fouriera na
$\R^d$, przestrzenie Sobolewa i ich zanurzenia. Dopiero po nich
budujemy metody wariacyjne, operatory nieograniczone oraz
słabą teorię eliptyczną i ewolucyjną.

Nie każde odwołanie do dalszego twierdzenia oznacza cykl w dowodzie:
czasem jest to jawna zapowiedź późniejszego, niezależnie
udowodnionego wyniku. W szczególności wykład o funkcjach
harmonicznych odwołuje się do ich analityczności dowiedzionej
później w tym samym rozdziale. Rozdziały o operatorach
pseudoróżniczkowych korzystają z nieunitarnej konwencji
transformaty Fouriera już zdefiniowanej w części podstawowej.
Jawne odsyłacze i zależności między rozdziałami indeksuje
aplikacja Biblioteki matematycznej, niezależnie od kompilacji skryptu.
\mainmatter

% Rozdział 1.
\chapter[Wprowadzenie do PDE]{Równania różniczkowe cząstkowe: pojęcia i pierwsze zjawiska}
\label{chap:pde-intro}

Równania różniczkowe cząstkowe (ang. \emph{partial differential equations}, PDE)
opisują związki między pochodnymi funkcji wielu zmiennych. Przy ich pomocy
formułuje się prawa przepływu ciepła, rozchodzenia się fal, ruchu płynów,
dynamiki kwantowej i geometrii rozmaitości. Ten sam operator różniczkowy
może prowadzić do zupełnie różnych zagadnień w zależności od tego, na
jakim obszarze pracujemy, jakie narzucamy warunki i w jakiej klasie
szukamy rozwiązań.

Celem rozdziału jest zbudowanie języka potrzebnego w dalszej części
skryptu. Wprowadzimy rząd i typ nieliniowości równania, symbol główny,
powierzchnie charakterystyczne i podstawowe klasy równań. Następnie,
na przykładach, oddzielimy trzy problemy: istnienie rozwiązania,
jego jednoznaczność i stabilną zależność od danych.

\section{Co jest niewiadomą w równaniu PDE?}
\label{sec:pde-unknown}

Niech $n\geq1$, $\Omega\subset\R^n$ będzie niepustym zbiorem otwartym,
a $u\colon\Omega\to\R$ funkcją niewiadomą. Równanie różniczkowe
cząstkowe wiąże wartości $u$ i jej pochodnych cząstkowych. Przykładowo:
\begin{align}
 -\Delta u&=f &&\text{(równanie Poissona)},\label{eq:poisson-intro}\\
 u_t-\kappa\Delta_xu&=0,\quad \kappa>0
 &&\text{(równanie ciepła)},\label{eq:heat-intro}\\
 u_{tt}-c^2\Delta_xu&=0,\quad c>0
 &&\text{(równanie falowe)},\label{eq:wave-intro}\\
 u_t+u\,u_x&=0 &&\text{(równanie Burgersa bez lepkości)}.
 \label{eq:burgers-intro}
\end{align}
W równaniach ciepła i falowym niewiadoma zależy od czasu $t$ i
zmiennych przestrzennych $x=(x_1,\ldots,x_d)$; w podanym równaniu
Burgersa jest jedna zmienna przestrzenna, czyli $d=1$. Zapis $\Delta_x$ oznacza
laplasjan liczony wyłącznie względem zmiennych przestrzennych, tj.
$\Delta_xu=\sum_{j=1}^d\partial_{x_j}^2u$.

\begin{uwaga}[Równanie a zagadnienie z dodatkowymi warunkami]
Jeśli na niepustym zbiorze otwartym istnieje rozwiązanie klasyczne
$u$ równania $-\Delta u=f$, to bez dodatkowych warunków nie jest
ono jedyne: dla każdej stałej $b\in\R$ funkcja $u+b$ również
spełnia to równanie. Dopiero po
określeniu obszaru, warunków brzegowych i klasy funkcji otrzymujemy
konkretne \emph{zagadnienie}. Przykładowo każda funkcja afiniczna
$u(x)=a\cdot x+b$ jest harmoniczna w $\R^n$; nie można zatem mówić
o jednoznaczności bez dalszych założeń.
\end{uwaga}

W pierwszych rozdziałach przez rozwiązanie klasyczne równania rzędu
$m$ rozumiemy funkcję klasy $C^m$ spełniającą równanie punktowo. Dla równań
ewolucyjnych wygodniej jest osobno liczyć pochodne czasowe i
przestrzenne; na przykład $C^{1,2}$ dla kolejności argumentów $(t,x)$ oznacza ciągłość funkcji, jej pierwszej pochodnej czasowej
oraz wszystkich pochodnych przestrzennych rzędu pierwszego i drugiego.
Nie wymaga się tu istnienia pochodnych mieszanych czasowo-przestrzennych. Rozwiązania słabe i dystrybucyjne
wprowadzimy później.

\section{Rząd równania i rodzaje nieliniowości}
\label{sec:pde-types}

Niech $\N_0=\{0,1,2,\ldots\}$ i $n\geq1$.
Dla wielowskaźnika $\alpha=(\alpha_1,\ldots,\alpha_n)\in\N_0^n$
przyjmujemy $|\alpha|=\alpha_1+\cdots+\alpha_n$ i
$\partial^\alpha=\partial_{x_1}^{\alpha_1}\cdots
\partial_{x_n}^{\alpha_n}$. Niech $m\geq1$, a $J^m u(x)$ oznacza
uporządkowaną rodzinę wszystkich wartości $\partial^\alpha u(x)$,
$|\alpha|\leq m$ (tzw. $m$-dżet funkcji w punkcie).

\begin{definicja}[Równanie rzędu $m$]
Niech $n,m\geq1$ będą liczbami całkowitymi,
$\Omega\subset\R^n$ będzie otwarta, a
$N_m=\#\{\alpha\in\N_0^n:|\alpha|\leq m\}$.
Niech $U\subset\Omega\times\R^{N_m}$ będzie zbiorem otwartym
i $F\colon U\to\R$. Dla funkcji $u\in C^m(\Omega,\R)$
takiej, że $(x,J^m u(x))\in U$ dla każdego $x\in\Omega$,
równanie $F(x,J^m u(x))=0$ nazywamy równaniem rzędu $m$,
jeżeli $F$ zależy nietrywialnie od współrzędnych $z_\alpha$
o $|\alpha|=m$, tj. istnieją dwa punkty zbioru $U$ o tym samym
$x$ i tych samych współrzędnych $z_\beta$ dla $|\beta|<m$,
w których $F$ przyjmuje różne wartości. Pojęcie rzędu odnosi się do przedstawionego równania.
W szczególnych punktach część najwyższego rzędu może znikać;
na przykład operator $x\,\partial_x^2+\partial_x$ ma rząd $2$
na $\R$, choć w punkcie $x=0$ jego symbol drugiego rzędu
zeruje się dla każdego kowektora.
\end{definicja}

\begin{definicja}[Liniowość, półliniowość i quasiliniowość]
Niech $\Omega\subset\R^n$ będzie otwarta i $m\geq1$.
\begin{enumerate}
\item Równanie \emph{liniowe} rzędu $m$ ma postać
$Lu=f$, gdzie $Lu=\sum_{|\alpha|\leq m}a_\alpha(x)\partial^\alpha u$,
a współczynniki $a_\alpha$ oraz $f$ nie zależą od $u$.
Jest \emph{jednorodne}, jeśli $f=0$.
\item Równanie jest \emph{półliniowe} (semiliniowe), jeśli ma postać
$\sum_{|\alpha|=m}a_\alpha(x)\partial^\alpha u
+G(x,J^{m-1}u)=0$. Pochodne najwyższego rzędu wchodzą zatem
liniowo, a ich współczynniki zależą tylko od $x$.
\item Równanie jest \emph{quasiliniowe}, jeśli można je zapisać jako
$\sum_{|\alpha|=m}a_\alpha(x,J^{m-1}u)\partial^\alpha u
+G(x,J^{m-1}u)=0$.
\item Równanie nazywamy \emph{w pełni nieliniowym}, gdy jego zależność
od pochodnych najwyższego rzędu nie jest afiniczna.
\end{enumerate}
Terminy te dotyczą przyjętego rzędu i postaci równania; równania
liniowe są szczególnym przypadkiem półliniowych, a półliniowe
szczególnym przypadkiem quasiliniowych. Klasyfikacja odnosi się do
\emph{wyrażenia różniczkowego}, nie wyłącznie do zbioru jego
rozwiązań: równania $\Delta u=0$ i $(\Delta u)^2=0$ mają ten
sam zbiór rozwiązań klasycznych, ale inne postaci algebraiczne.
\end{definicja}

\begin{przyklad}[Rozpoznawanie typu]
Dla niewiadomej dwóch zmiennych, oznaczanych przez $(x,y)$
w pierwszych czterech przykładach, a przez $(t,x)$ w ostatnim:
\begin{enumerate}
\item $u_{xx}+u_{yy}=\sin x$ jest liniowe i niejednorodne.
\item $u_{xx}+u_{yy}+u^3=0$ jest półliniowe.
\item $(1+u_x^2)u_{xx}+u_{yy}=0$ jest quasiliniowe: współczynnik
przy $u_{xx}$ zależy od pochodnej pierwszego rzędu.
\item $u_{xx}u_{yy}-u_{xy}^2=1$ jest w pełni nieliniowe
(równanie typu Monge'a--Ampère'a).
\item $u_t+u\,u_x=0$ jest quasiliniowe pierwszego rzędu.
\end{enumerate}
\end{przyklad}

\begin{stwierdzenie}[Zasada superpozycji]
Niech $V,W$ będą przestrzeniami wektorowymi funkcji nad $\R$,
$L\colon V\to W$ operatorem liniowym, $u_1,u_2\in V$
oraz $f_1,f_2,f\in W$. Jeśli $Lu_1=f_1$ i
$Lu_2=f_2$, to dla dowolnych $a,b\in\R$ zachodzi
$L(au_1+bu_2)=af_1+bf_2$. W szczególności zbiór rozwiązań
równania jednorodnego $Lu=0$ jest podprzestrzenią wektorową
$V$, a zbiór rozwiązań $Lu=f$, jeżeli jest niepusty,
ma postać $u_0+\ker L$.
\end{stwierdzenie}
\begin{proof}
Z liniowości $L$ otrzymujemy $L(au_1+bu_2)=aLu_1+bLu_2
=af_1+bf_2$. Zbiór $\ker L$ zawiera wektor zerowy i jest zamknięty
na kombinacje liniowe, więc jest podprzestrzenią wektorową $V$.
Jeżeli $Lu_0=f$, to dla $u\in V$ mamy $Lu=f$ wtedy i tylko wtedy,
gdy $L(u-u_0)=0$, czyli $u-u_0\in\ker L$. Zbiór rozwiązań
jest zatem dokładnie równy $u_0+\ker L$.
\end{proof}

\begin{uwaga}
Dla równania nieliniowego nie wolno zazwyczaj dodawać rozwiązań.
Rozważmy na przykład równanie $w_x=w^2$ na
$\Omega=(-1,1)^2$, gdzie $w=w(x,y)$.
Funkcje $w_1(x,y)=1/(2-x)$ i $w_2(x,y)=1/(3-x)$
są rozwiązaniami, ponieważ $w_j'=w_j^2$.
Ich suma nie jest jednak rozwiązaniem, gdyż
$(w_1+w_2)'=w_1^2+w_2^2\ne(w_1+w_2)^2$.
Podobnie $2w_1$ nie jest rozwiązaniem: pokazuje to,
że nie działa również zwykłe mnożenie rozwiązań przez skalar.
\end{uwaga}

\section{Dane początkowe i brzegowe}
\label{sec:pde-data}

Niech $\Omega\subset\R^d$ będzie obszarem ograniczonym, tj.
niepustym, spójnym zbiorem otwartym, i załóżmy, że jej brzeg
$\partial\Omega$ jest co najmniej klasy $C^1$ tam, gdzie używamy
wektora normalnego $\nu$ skierowanego na zewnątrz.

\begin{definicja}[Podstawowe warunki brzegowe]
Niech $\Omega\subset\R^d$ będzie obszarem ograniczonym o brzegu
klasy $C^1$, $g,\beta\colon\partial\Omega\to\R$ będą zadanymi
funkcjami, a $\nu$ jednostkowym wektorem normalnym skierowanym
na zewnątrz. Dla rzeczywistej funkcji $u$ ciągłej na
$\overline\Omega$ (w warunkach Neumanna i Robina dodatkowo
klasy $C^1$ aż do brzegu) rozważamy:
\begin{enumerate}
\item \emph{warunek Dirichleta}: $u|_{\partial\Omega}=g$;
\item \emph{warunek Neumanna}: $\partial_\nu u|_{\partial\Omega}=g$,
 gdzie $\partial_\nu u=\nabla u\cdot\nu$;
\item \emph{warunek Robina}: $\partial_\nu u+\beta u=g$
 na $\partial\Omega$, gdzie $\beta$ jest zadaną funkcją.
\end{enumerate}
Warunki te rozumiemy na razie punktowo, dla funkcji o odpowiedniej
regularności aż do brzegu.
\end{definicja}

Dla zagadnień ewolucyjnych na $(0,T)\times\Omega$, gdzie $T>0$,
warunek początkowy zapisujemy na przykład jako $u(0,x)=u_0(x)$.
Równanie falowe drugiego rzędu względem czasu wymaga zwykle także
warunku $u_t(0,x)=u_1(x)$. Jeśli $\Omega$ jest ograniczona,
potrzebne są ponadto odpowiednie warunki na bocznym brzegu
$(0,T)\times\partial\Omega$.

\begin{przyklad}[Zgodność danych na krawędzi czasoprzestrzeni]
Niech $T>0$, $\Omega\subset\R^d$ będzie obszarem ograniczonym,
$u_0\colon\overline\Omega\to\R$ oraz
$g\colon[0,T]\times\partial\Omega\to\R$.
Jeżeli $u\in C([0,T]\times\overline\Omega)$ spełnia
$u(0,x)=u_0(x)$ dla $x\in\overline\Omega$ oraz
$u(t,x)=g(t,x)$ dla $(t,x)\in[0,T]\times\partial\Omega$,
to koniecznie $u_0(x)=g(0,x)$ dla każdego $x\in\partial\Omega$.
Jest to najprostszy \emph{warunek zgodności} danych; wyższa
regularność często narzuca dalsze warunki.
\end{przyklad}

\begin{uwaga}[Zagadnienie Cauchy'ego]
Przez zagadnienie Cauchy'ego rozumiemy poszukiwanie rozwiązania
na podstawie danych zadanych na odpowiedniej hiperpłaszczyźnie
lub hiperpowierzchni. Nie należy utożsamiać go z zagadnieniem
Dirichleta. Na przykład zadanie wartości $u$ i pochodnej normalnej
na części brzegu dla równania Laplace'a jest zagadnieniem Cauchy'ego,
które, jak zobaczymy w~\S\ref{sec:pde-hadamard}, jest niestabilne.
\end{uwaga}

\section{Symbol główny i powierzchnie charakterystyczne}
\label{sec:pde-symbol}

Pojęcie symbolu ujawnia, co z równania pozostaje przy badaniu
oscylacji o bardzo dużej częstości. Niech $m\geq1$,
$\Omega\subset\R^n$ będzie otwarta, a
$P=\sum_{|\alpha|\leq m}a_\alpha(x)\partial^\alpha$
liniowym operatorem różniczkowym o współczynnikach ciągłych
na $\Omega$. Zakładamy, że przynajmniej jeden współczynnik
najwyższego rzędu nie znika identycznie.

\begin{definicja}[Symbol główny]
\emph{Symbolem głównym} operatora
$P=\sum_{|\alpha|\leq m}a_\alpha(x)\partial^\alpha$ rzędu $m$
na otwartym zbiorze $\Omega\subset\R^n$ nazywamy funkcję
$p_m\colon\Omega\times\R^n\to\C$, będącą dla każdego ustalonego $x$
wielomianem jednorodnym stopnia $m$ względem $\xi$ (dopuszczamy
wielomian zerowy w punktach zaniku części głównej), daną wzorem
\begin{equation}\label{eq:principal-symbol}
 p_m(x,\xi)=\sum_{|\alpha|=m}a_\alpha(x)\xi^\alpha.
\end{equation}
W naszej konwencji $\partial^\alpha$ nie zawiera czynników $-i$.
Jeśli ten sam operator zapiszemy w konwencji $D_j=-i\partial_{x_j}$,
to $\partial^\alpha=i^{|\alpha|}D^\alpha$, więc jego symbol główny
w konwencji $D$ jest równy $i^mp_m(x,\xi)$.
\end{definicja}

\begin{stwierdzenie}[Test wysokich częstości]
Niech $\Omega\subset\R^n$ będzie otwarta, $m\geq1$ oraz
$P=\sum_{|\alpha|\leq m}a_\alpha(x)\partial^\alpha$ będzie liniowym
operatorem rzędu $m$ o ciągłych współczynnikach, działającym na
funkcjach zespolonych (w razie potrzeby po zespoleniu operatora
rzeczywistego). Niech $p_m$ będzie jego symbolem głównym
w konwencji~\eqref{eq:principal-symbol}, a $x\in\Omega$, $\xi\in\R^n$. Dla $\lambda>0$ oraz
$v_\lambda(z)=e^{i\lambda z\cdot\xi}$ zachodzi
\[
 e^{-i\lambda x\cdot\xi}Pv_\lambda(x)
 =i^m\lambda^m p_m(x,\xi)+O(\lambda^{m-1})
 \quad\text{gdy }\lambda\to\infty.
\]
Stała w symbolu $O(\lambda^{m-1})$ może zależeć od $x$ i $\xi$.
\end{stwierdzenie}
\begin{proof}
Mamy $\partial^\alpha v_\lambda(x)=
(i\lambda)^{|\alpha|}\xi^\alpha v_\lambda(x)$.
Po podstawieniu do definicji $P$ wyrazy z $|\alpha|=m$
dają dokładnie $i^m\lambda^mp_m(x,\xi)$, a skończona
suma pozostałych wyrazów jest rzędu co najwyżej $\lambda^{m-1}$.
\end{proof}

\begin{definicja}[Eliptyczność]
Liniowy operator skalarny $P$ rzędu $m$ na zbiorze otwartym
$\Omega\subset\R^n$, o symbolu głównym $p_m$ danym przez
\eqref{eq:principal-symbol}, nazywamy \emph{eliptycznym w punkcie}
$x\in\Omega$, jeżeli $p_m(x,\xi)\ne0$ dla każdego
$\xi\in\R^n\setminus\{0\}$. Jest \emph{eliptyczny na $\Omega$},
gdy warunek zachodzi w każdym punkcie.
\end{definicja}

\begin{przyklad}
Dla $P=-\Delta=-\sum_{j=1}^n\partial_{x_j}^2$
przyjęta konwencja daje $p_2(x,\xi)=-|\xi|^2$.
Symbol nie zeruje się poza $\xi=0$, więc laplasjan jest
operatorem eliptycznym. Znak symbolu nie ma znaczenia
dla samej definicji eliptyczności.
\end{przyklad}

\begin{uwaga}[Eliptyczność a jednostajna eliptyczność]
Dla operatora drugiego rzędu o części głównej
$-\sum_{i,j=1}^n a_{ij}(x)\partial_{x_ix_j}$, gdzie
$A(x)=(a_{ij}(x))_{i,j=1}^n$ jest macierzą rzeczywistą i symetryczną,
warunek jednostajnej eliptyczności (w konwencji dolnego
ograniczenia) ma postać
$\xi^{\mathsf T}A(x)\xi\ge\lambda|\xi|^2$
dla wszystkich $x\in\Omega$, $\xi\in\R^n$ i pewnej
\emph{jednej} stałej $\lambda>0$. Samo niezerowanie się
symbolu głównego w poszczególnych punktach nie daje
takiej stałej. Przykładowo operator
$P=-x^2\partial_x^2$ jest eliptyczny na $(0,1)$,
ale nie jest jednostajnie eliptyczny: dla $x\downarrow0$
wartość współczynnika $x^2$ dąży do zera.
Wymóg dolnego oszacowania przez $\lambda|\xi|^2$
dotyczy operatorów drugiego rzędu w postaci dodatnio
określonej; nie jest definicją wszystkich eliptycznych
operatorów dowolnego rzędu.
\end{uwaga}

\begin{definicja}[Powierzchnia charakterystyczna]
Niech $P$ będzie liniowym operatorem skalarnym rzędu $m$ na
$\Omega\subset\R^n$, a $\Sigma\subset\Omega$ hiperpowierzchnią klasy $C^1$.
W otoczeniu otwartym $U\subset\Omega$ rozważanego punktu przyjmijmy
$\Sigma\cap U=\{z\in U:\phi(z)=0\}$, gdzie
$\phi\in C^1(U)$ i $\mathrm d\phi(z)\ne0$ dla $z\in\Sigma\cap U$.
Mówimy, że $\Sigma$ jest \emph{charakterystyczna dla $P$
w punkcie} $x\in\Sigma$, jeżeli
$p_m(x,\mathrm d\phi(x))=0$. W przeciwnym razie jest
\emph{niecharakterystyczna w $x$}. Warunek nie zależy od
wyboru lokalnej funkcji definiującej $\phi$.
\end{definicja}

\begin{proof}[Uzasadnienie niezależności od funkcji definiującej]
Jeśli $\phi$ i $\psi$ są dwiema lokalnymi funkcjami definiującymi
tę samą hiperpowierzchnię klasy $C^1$, to
$T_x\Sigma=\ker\mathrm d\phi(x)=\ker\mathrm d\psi(x)$.
Istotnie, różniczkowanie funkcji definiującej wzdłuż krzywej
na $\Sigma$ daje zero, a zarówno przestrzeń styczna, jak i jądro
niezerowej różniczki mają wymiar $n-1$. Wybierzmy $v\in\R^n$
tak, by $\mathrm d\phi(x)v=1$. Dla dowolnego $w\in\R^n$
wektor $w-\mathrm d\phi(x)w\,v$ należy do wspólnego jądra,
więc $\mathrm d\psi(x)w=\mathrm d\psi(x)v\,\mathrm d\phi(x)w$.
Stąd $\mathrm d\psi(x)=a\,\mathrm d\phi(x)$, gdzie
$a=\mathrm d\psi(x)v\ne0$, ponieważ $\mathrm d\psi(x)\ne0$.
Ponieważ $p_m(x,\cdot)$ jest jednorodny stopnia $m$,
$p_m(x,\mathrm d\psi(x))=a^mp_m(x,\mathrm d\phi(x))$.
Oba wyrażenia zerują się zatem jednocześnie.
\end{proof}

\begin{przyklad}[Stożek charakterystyczny równania falowego]
Dla operatora $P=\partial_t^2-c^2\Delta_x$ w $\R^{1+d}$
symbol główny jest równy
$p_2(t,x;\tau,\xi)=\tau^2-c^2|\xi|^2$.
Kowektory $(\tau,\xi)\ne0$ spełniające
$\tau^2=c^2|\xi|^2$ są charakterystyczne. W przypadku
jednego wymiaru przestrzennego proste $t-x/c=\mathrm{const}$
i $t+x/c=\mathrm{const}$ są charakterystyczne.
Natomiast powierzchnia początkowa $t=0$ jest
niecharakterystyczna, bo $p_2(1,0)=1$.
\end{przyklad}

\begin{uwaga}[Pułapka: symbol równania ciepła]
Dla $P=\partial_t-\kappa\Delta_x$, traktowanego jako zwykły
operator drugiego rzędu we wszystkich zmiennych $(t,x)$,
symbol główny wynosi $-\kappa|\xi|^2$: nie zależy od $\tau$.
W takim sensie powierzchnia $t=0$ jest charakterystyczna
($p_2((t,x),(1,0))=0$),
chociaż zadanie początkowe dla równania ciepła jest naturalne
i stabilne w odpowiednich przestrzeniach. Dlatego równania
paraboliczne klasyfikuje się przy użyciu \emph{ważonego}
porządku pochodnych: jednej pochodnej czasowej przypisuje się
wagę $2$, a przestrzennej wagę $1$. Przy konwencji symbolu
$D_j=-i\partial_{x_j}$ prowadzi to dla operatora ciepła do
symbolu parabolicznego $i\tau+\kappa|\xi|^2$, który dla
rzeczywistych $(\tau,\xi)\ne0$ nie znika. W szczególności
nie wolno bez dodatkowych zastrzeżeń stosować do równania
ciepła intuicji charakterystyk równania falowego.
\end{uwaga}

\section{Klasyfikacja operatorów drugiego rzędu}
\label{sec:pde-classification}

Dla równania drugiego rzędu w dwóch zmiennych rzeczywistych
rozważmy część główną
\begin{equation}\label{eq:2d-principal}
 A(x,y)u_{xx}+2B(x,y)u_{xy}+C(x,y)u_{yy}.
\end{equation}
Współczynniki $A,B,C$ są rzeczywiste i ciągłe, a odpowiadająca
im macierz symetryczna to
$\mathsf A(x,y)=\begin{pmatrix}A&B\\B&C\end{pmatrix}$.

\begin{definicja}[Klasyfikacja w dwóch zmiennych]
Niech $\Omega\subset\R^2$ będzie otwarta, $A,B,C\colon\Omega\to\R$
będą ciągłe, a $\mathsf A=\begin{pmatrix}A&B\\B&C\end{pmatrix}$
będzie macierzą części głównej~\eqref{eq:2d-principal}.
Ustalmy punkt $(x,y)\in\Omega$.
Oznaczmy $\delta(x,y)=B(x,y)^2-A(x,y)C(x,y)$.
Równanie z częścią główną~\eqref{eq:2d-principal} ma w tym punkcie
typ \emph{eliptyczny}, gdy $\delta<0$;
\emph{hiperboliczny}, gdy $\delta>0$;
\emph{paraboliczny w tradycyjnej klasyfikacji dwuzmiennowej},
gdy $\delta=0$ i $\mathsf A\ne0$.
Jeśli $\mathsf A=0$, część drugiego rzędu zanika i tej
klasyfikacji tam nie stosujemy.
\end{definicja}

\begin{stwierdzenie}[Interpretacja geometryczna i niezmienniczość typu]
Niech $\mathsf A=\begin{pmatrix}A&B\\B&C\end{pmatrix}$
będzie rzeczywistą, symetryczną macierzą $2\times2$, różną od zera.
Wtedy:
\begin{enumerate}
\item $B^2-AC<0$ wtedy i tylko wtedy, gdy $\mathsf A$ jest
 dodatnio albo ujemnie określona;
\item $B^2-AC>0$ wtedy i tylko wtedy, gdy $\mathsf A$ ma
 wartości własne przeciwnych znaków;
\item $B^2-AC=0$ wtedy i tylko wtedy, gdy $\mathsf A$ ma rząd $1$.
\end{enumerate}
Pod lokalną zmianą współrzędnych klasy $C^2$ o nieosobliwym
jakobianie liczba dodatnich, ujemnych i zerowych wartości własnych
macierzy części głównej nie zmienia się. W konsekwencji typ
równania jest niezmienniczy względem takich zmian współrzędnych.
\end{stwierdzenie}
\begin{proof}
Macierz $\mathsf A$ ma rzeczywiste wartości własne
$\lambda_1,\lambda_2$, przy czym
$\lambda_1\lambda_2=\det\mathsf A=AC-B^2=-\delta$.
Jeśli $\delta<0$, iloczyn jest dodatni i obie niezerowe
wartości własne mają ten sam znak. Jeśli $\delta>0$,
mają znaki przeciwne. Gdy $\delta=0$, jedna wartość własna
jest zerowa, a druga niezerowa, ponieważ $\mathsf A\ne0$.

Niech $z=\Phi(x)$ będzie lokalną zmianą współrzędnych,
$J=D\Phi(x)$ jej nieosobliwym jakobianem i
$u(x)=v(\Phi(x))$, gdzie $\Phi$ i $v$ są klasy $C^2$.
Piszemy $J_{ki}=\partial_{x_i}\Phi_k(x)$.
Dwukrotne zastosowanie reguły łańcuchowej daje
\[
 \partial_{x_i x_j}u(x)=
 \sum_{k,l=1}^2(\partial_{z_k z_l}v)(\Phi(x))J_{ki}J_{lj}
 +\sum_{k=1}^2(\partial_{z_k}v)(\Phi(x))
       \partial_{x_i x_j}\Phi_k(x).
\]
Jeśli $a_{ij}$ są elementami $\mathsf A$, to w zapisie
$\sum_{k,l=1}^2 b_{kl}\partial_{z_k z_l}v$ otrzymujemy
$b_{kl}=\sum_{i,j=1}^2J_{ki}a_{ij}J_{lj}=(J\mathsf A J^{\mathsf T})_{kl}$.
Macierz $(b_{kl})$ jest symetryczna; po zebraniu pochodnych
mieszanych $\partial_{z_1z_2}v=\partial_{z_2z_1}v$ ich łączny
współczynnik wynosi $2b_{12}$.
Drugi składnik wzoru zawiera jedynie pierwsze pochodne $v$.
Prawo bezwładności
Sylvestera zapewnia równość liczby dodatnich, ujemnych
i zerowych wartości własnych obu macierzy.
\end{proof}

\begin{przyklad}[Trzy podstawowe typy]
\begin{enumerate}
\item $u_{xx}+u_{yy}=0$: $A=C=1$, $B=0$, więc $\delta=-1$;
 typ eliptyczny.
\item $u_{xx}-u_{yy}=0$: $A=1$, $C=-1$, $B=0$, więc
 $\delta=1$; typ hiperboliczny.
\item $u_{xx}+2u_{xy}+u_{yy}=0$: $A=B=C=1$, więc
 $\delta=0$; typ paraboliczny w klasyfikacji dwuzmiennowej.
 Część główna ma postać $(\partial_x+\partial_y)^2$.
\end{enumerate}
\end{przyklad}

\begin{przyklad}[Równanie zmieniające typ]
Rozważmy na $\R^2$ równanie Tricomiego
\begin{equation}\label{eq:tricomi}
 u_{xx}+x\,u_{yy}=0.
\end{equation}
Mamy $A=1$, $B=0$, $C=x$, zatem $\delta=-x$. Równanie jest
eliptyczne w obszarze $x>0$, hiperboliczne w obszarze $x<0$
i typu parabolicznego na prostej $x=0$.
\end{przyklad}

\begin{center}
\begin{tikzpicture}[scale=.9,>=Stealth]
 \fill[gray!10] (-4,-1.25) rectangle (0,1.25);
 \fill[gray!23] (0,-1.25) rectangle (4,1.25);
 \draw[->] (-4.3,0)--(4.45,0) node[right] {$x$};
 \draw[->] (0,-1.5)--(0,1.5);
 \draw[line width=1.1pt] (0,-1.25)--(0,1.25);
 \node[align=center] at (-2,0.65) {hiperboliczny\\$x<0$};
 \node[align=center] at (2,0.65) {eliptyczny\\$x>0$};
 \node[above,align=center] at (0,1.6) {$x=0$: typ paraboliczny};
\end{tikzpicture}
\par\smallskip\small
Typ równania Tricomiego zmienia się na prostej $x=0$.
\end{center}

\begin{uwaga}[Dwa znaczenia słowa „paraboliczny”]
Typ paraboliczny w klasyfikacji wyrażenia
$Au_{xx}+2Bu_{xy}+Cu_{yy}$ dotyczy macierzy części
drugiego rzędu w \emph{dwóch równorzędnych zmiennych}.
Równanie ciepła $u_t-\kappa\Delta_xu=0$ jest natomiast
parabolicznym \emph{równaniem ewolucyjnym} z wyróżnionym
czasem i z innym skalowaniem pochodnych. Terminy są
powiązane historycznie, ale nie należy utożsamiać tych
definicji bez wskazania konwencji.
\end{uwaga}

W większej liczbie zmiennych sama nieokreśloność macierzy
części głównej nie wystarcza do pełnej definicji
hiperboliczności. Dla równania falowego wyróżnia się
kowektor czasowy oraz wymaga odpowiednich własności
rzeczywistych pierwiastków wielomianu symbolu w kierunku
czasowym. Systematyczną teorię rozwinie rozdział o
równaniach hiperbolicznych.

\section{Pierwsze modele: co można udowodnić od razu?}
\label{sec:pde-models}

\subsection{Laplasjan: zasada maksimum i jednoznaczność}

\begin{twierdzenie}[Słaba zasada maksimum]
\label{thm:laplace-max}
Niech $\Omega\subset\R^n$ będzie niepustym, ograniczonym
zbiorem otwartym i niech $u\in C^2(\Omega)\cap
C(\overline\Omega)$ będzie funkcją rzeczywistą. Jeśli
$\Delta u\geq0$ w $\Omega$, to
\begin{equation}\label{eq:maximum-principle}
 \max_{\overline\Omega}u\leq\max_{\partial\Omega}u.
\end{equation}
Jeśli $\Delta u\leq0$, to
$\min_{\overline\Omega}u\geq\min_{\partial\Omega}u$.
W szczególności, gdy $\Delta u=0$, oba wnioski zachodzą.
\end{twierdzenie}
\begin{proof}
Dla $\eps>0$ określmy $u_\eps(x)=u(x)+\eps|x|^2$.
W $\Omega$ zachodzi
$\Delta u_\eps=\Delta u+2n\eps>0$.
Funkcja $u_\eps$ jest ciągła na zwartym zbiorze
$\overline\Omega$, zatem osiąga maksimum. Gdyby osiągała
je w punkcie wewnętrznym, jej hesjan byłby w tym punkcie
ujemnie półokreślony, stąd $\Delta u_\eps\leq0$,
co daje sprzeczność. Maksimum $u_\eps$ leży zatem na brzegu.
Wobec tego dla każdego $x\in\overline\Omega$ mamy
$u(x)+\eps|x|^2\leq
\max_{z\in\partial\Omega}(u(z)+\eps|z|^2)$.
Niech $M=\max_{z\in\overline\Omega}|z|^2<\infty$.
Prawa strona jest nie większa niż
$\max_{\partial\Omega}u+\eps M$, a $\eps|x|^2\geq0$.
Stąd $u(x)\leq\max_{\partial\Omega}u+\eps M$ dla każdego
$\eps>0$. Przechodząc z $\eps$ do zera, otrzymujemy
$u(x)\leq\max_{\partial\Omega}u$; biorąc maksimum po $x$,
dostajemy~\eqref{eq:maximum-principle}.
Dla $\Delta u\leq0$ stosujemy wykazany wynik do $-u$.
\end{proof}

\begin{wniosek}[Jednoznaczność i stabilność zagadnienia Dirichleta]
\label{cor:dirichlet-stability}
Niech $\Omega\subset\R^n$ będzie niepustym ograniczonym
zbiorem otwartym, a $u,v\in C^2(\Omega)\cap
C(\overline\Omega)$ będą funkcjami rzeczywistymi spełniającymi $\Delta u=\Delta v$
w $\Omega$. Wówczas
\[
 \max_{x\in\overline\Omega}|u(x)-v(x)|
 \leq\max_{x\in\partial\Omega}|u(x)-v(x)|.
\]
W szczególności dwa rozwiązania klasyczne zagadnienia
$-\Delta u=f$ z takimi samymi wartościami brzegowymi
muszą być równe.
\end{wniosek}
\begin{proof}
Różnica $w=u-v$ jest harmoniczna. Twierdzenie
\ref{thm:laplace-max}, zastosowane do $w$ i $-w$,
daje dla każdego $x\in\overline\Omega$ nierówności
$w(x)\leq\max_{\partial\Omega}w$ oraz
$-w(x)\leq\max_{\partial\Omega}(-w)$.
Obie prawe strony są nie większe niż
$M=\max_{\partial\Omega}|w|$, więc $|w(x)|\leq M$.
Jeśli wartości brzegowe $u$ i $v$ są równe, to $M=0$
i $u=v$ na $\overline\Omega$.
\end{proof}

Zauważmy, że udowodniliśmy \emph{jednoznaczność i stabilność
względem wartości brzegowych przy ustalonej prawej stronie $f$},
ale nie istnienie rozwiązania Dirichleta dla dowolnego obszaru
i dowolnych danych. Oszacowanie nie dotyczy zmiany prawej
strony równania. Istnienie wymaga osobnej argumentacji.

\subsection{Warunek Neumanna: brak jednoznaczności i warunek zgodności}

\begin{stwierdzenie}[Konieczny warunek zgodności]
Niech $n\geq1$, $\Omega\subset\R^n$ będzie obszarem ograniczonym
o brzegu klasy $C^1$, a $f\in C(\overline\Omega,\C)$ i
$g\in C(\partial\Omega,\C)$. Niech $\nu$ będzie jednostkową
normalną zewnętrzną. Jeśli $u\in C^2(\overline\Omega,\C)$ spełnia
$-\Delta u=f$ w $\Omega$ oraz $\partial_\nu u=g$
na $\partial\Omega$, to
\begin{equation}\label{eq:neumann-compatibility}
 \int_\Omega f(x)\,\mathrm d\lambda_n(x)
 +\int_{\partial\Omega}g(x)\,\mathrm d\sigma(x)=0.
\end{equation}
Tutaj $\lambda_n$ oznacza miarę Lebesgue'a na $\R^n$,
a $\sigma$ miarę powierzchniową na $\partial\Omega$.
Jeśli $u,v\in C^2(\overline\Omega,\C)$ rozwiązują to samo
zagadnienie Neumanna, to funkcja $u-v$ jest stała na $\Omega$.
\end{stwierdzenie}
\begin{proof}
Ponieważ $\nabla u$ jest polem klasy $C^1$ aż do brzegu,
a $\Omega$ jest obszarem ograniczonym o brzegu klasy $C^1$,
możemy zastosować twierdzenie Gaussa do pola $\nabla u$.
Daje ono
$\int_\Omega\Delta u(x)\,\mathrm d\lambda_n(x)
=\int_{\partial\Omega}\partial_\nu u(x)\,\mathrm d\sigma(x)$.
Po podstawieniu $\Delta u=-f$ oraz $\partial_\nu u=g$
i przeniesieniu obu całek na jedną stronę otrzymujemy
\eqref{eq:neumann-compatibility}.

Dla $w=u-v$ zachodzi $\Delta w=0$ i $\partial_\nu w=0$.
Stosując twierdzenie Gaussa do pola $\overline w\nabla w$ klasy $C^1$
(w przypadku zespolonym osobno do części rzeczywistej i urojonej)
i korzystając z równości
$\operatorname{div}(\overline w\nabla w)=|\nabla w|^2+\overline w\Delta w$,
otrzymujemy pierwszą tożsamość Greena:
\[
 \int_\Omega|\nabla w(x)|^2\,\mathrm d\lambda_n(x)
 =\int_{\partial\Omega}\overline{w(x)}\partial_\nu w(x)\,
 \mathrm d\sigma(x)
 -\int_\Omega\overline{w(x)}\Delta w(x)\,\mathrm d\lambda_n(x)=0.
\]
Gdyby $|\nabla w(x_0)|>0$ dla pewnego $x_0\in\Omega$,
to z ciągłości istniałaby kula zawarta w $\Omega$, na której
$|\nabla w|^2$ byłaby ograniczona z dołu przez dodatnią stałą.
Całka byłaby wtedy dodatnia, wbrew otrzymanej równości.
Zatem $\nabla w=0$ w $\Omega$. Na każdej kuli zawartej
w $\Omega$ funkcja $w$ jest stała, co wynika z całkowania
pochodnej wzdłuż odcinków. Jest więc lokalnie stała,
a spójność $\Omega$ implikuje jej stałość na całym obszarze.
Odwrotnie, dodanie do rozwiązania dowolnej stałej zespolonej
nie zmienia ani laplasjanu, ani pochodnej normalnej.
Jeżeli więc rozwiązanie istnieje, zbiór wszystkich rozwiązań
jest dokładnie rodziną $u+\C$.
\end{proof}

\subsection{Transport: rozwiązanie przesuwa się po charakterystykach}

\begin{stwierdzenie}[Transport ze stałą prędkością]
\label{prop:transport}
Niech $c\in\R$ i $u_0\in C^1(\R)$. Funkcja
$u\colon[0,\infty)\times\R\to\R$ określona wzorem
$u(t,x)=u_0(x-ct)$ jest rozwiązaniem klasycznym
$u_t+c\,u_x=0$ z warunkiem $u(0,x)=u_0(x)$.
Jest to jedyne rozwiązanie klasy $C^1$ na
$[0,\infty)\times\R$ z tym warunkiem.
\end{stwierdzenie}
\begin{proof}
Różniczkując, dostajemy $u_t=-c\,u_0'(x-ct)$
i $u_x=u_0'(x-ct)$; równanie jest spełnione.
Jeśli $v$ jest dowolnym innym rozwiązaniem, to dla
ustalonego $\xi\in\R$ funkcja
$h(t)=v(t,\xi+ct)$ spełnia
$h'(t)=v_t(t,\xi+ct)+c v_x(t,\xi+ct)=0$.
Zatem $v(t,\xi+ct)=v(0,\xi)=u_0(\xi)$, czyli
$v(t,x)=u_0(x-ct)$.
\end{proof}

\begin{center}
\begin{tikzpicture}[>=Stealth,scale=.95]
 \draw[->] (-3.5,0)--(4.3,0) node[right] {$x$};
 \draw[->] (-3.15,-.15)--(-3.15,3.1) node[above] {$t$};
 \foreach \x in {-2.4,-1.3,-.2,.9} {
   \draw[blue!55!black,thick] (\x,0)--({\x+1.3},2.6);
 }
 \draw[densely dashed] (-3.15,2.6)--(3.7,2.6);
 \fill[blue!55!black] (-1.3,0) circle (1.8pt);
 \fill[blue!55!black] (0,2.6) circle (1.8pt);
 \node[below] at (-1.3,0) {$\xi$};
 \node[above] at (0,2.6) {$(\xi+ct,t)$};
 \node[anchor=west] at (2.6,1) {$c>0$};
\end{tikzpicture}
\par\smallskip\small
Charakterystyki równania $u_t+c\,u_x=0$ są prostymi $x=\xi+ct$;
wartość $u$ pozostaje na każdej z nich stała (rysunek dla $c>0$).
\end{center}

\begin{przyklad}[Utrata gładkości mimo gładkich danych]
\label{ex:burgers-blowup}
Rozważmy dla $t\geq0$, $x\in\R$ równanie
$u_t+u\,u_x=0$ z warunkiem $u(0,x)=-\sin x$.
Załóżmy, że rzeczywista funkcja $u$ jest rozwiązaniem klasy $C^1$ na
$[0,T)\times\R$, gdzie $T>0$. Dla $\xi\in\R$ połóżmy
$X(t)=\xi-t\sin\xi$ i $h(t)=u(t,X(t))$.
Z reguły łańcuchowej i równania Burgersa wynika
$h'(t)=-(h(t)+\sin\xi)u_x(t,X(t))$.
Funkcja $q(t)=h(t)+\sin\xi$ spełnia zatem liniowe równanie
$q'(t)=-u_x(t,X(t))q(t)$ z $q(0)=0$.
Współczynnik jest ciągły, więc po pomnożeniu przez czynnik
całkujący $\exp(\int_0^t u_x(s,X(s))\,\mathrm d\lambda_1(s))$
otrzymujemy funkcję o zerowej pochodnej i zerowej wartości
początkowej. Stąd $q(t)=0$ na całym $[0,T)$, a więc
$X'(t)=u(t,X(t))$. Otrzymujemy równanie charakterystyki
\begin{equation}\label{eq:burgers-characteristics}
 x=\xi-t\sin\xi,\qquad u(t,x)=-\sin\xi.
\end{equation}
Dla $0\leq t<1$ pochodna odwzorowania
$\xi\mapsto\xi-t\sin\xi$ wynosi $1-t\cos\xi\geq1-t>0$.
Dla ustalonego $t$ odwzorowanie $F_t(\xi)=\xi-t\sin\xi$
jest więc ściśle rosnące, a z $|F_t(\xi)-\xi|\leq t$
wynika, że jego granice przy $\xi\to\pm\infty$ są równe
odpowiednio $\pm\infty$. Z ciągłości wynika surjektywność.
Twierdzenie o funkcji uwikłanej zastosowane do równania
$F_t(\xi)-x=0$ (pochodna względem $\xi$ jest dodatnia)
daje gładką zależność odwrotności $\xi=\xi(t,x)$ od obu zmiennych,
również lokalnie w otoczeniu $t=0$.
Różniczkując $x=\xi-t\sin\xi$, dostajemy
$\xi_x=(1-t\cos\xi)^{-1}$ oraz
$\xi_t=\sin\xi/(1-t\cos\xi)$.
Dla $u=-\sin\xi$ mamy więc $u_t=-\cos\xi\,\xi_t$,
$u_x=-\cos\xi\,\xi_x$ i $u_t+uu_x=0$.
Ponadto $\xi(0,x)=x$, więc spełniony jest warunek początkowy.
W szczególności
\begin{equation}\label{eq:burgers-gradient}
 u_x(t,x)=\frac{-\cos\xi}{1-t\cos\xi}.
\end{equation}
Dla $x=0$ odpowiada $\xi=0$, a więc
$u_x(t,0)=-1/(1-t)\to-\infty$ gdy $t\uparrow1$.
Rozwiązanie pozostaje ograniczone, ale jego pochodna
przestrzenna staje się nieograniczona w skończonym czasie.
Ponieważ dla $t<1$ odwzorowanie $F_t$ jest bijekcją, każda
funkcja klasy $C^1$ spełniająca te same dane musi być równa
skonstruowanemu rozwiązaniu w zakresie $t<1$. Gdyby istniało
rozwiązanie klasy $C^1$ na $[0,T)\times\R$ dla $T>1$,
jego pochodna $u_x$ byłaby ciągła w $(1,0)$, co jest
sprzeczne z powyższą granicą.
To najprostszy mechanizm powstawania osobliwości i
fal uderzeniowych.
\end{przyklad}

\subsection{Równanie ciepła: natychmiastowe rozprzestrzenianie}

\begin{stwierdzenie}[Jądro ciepła na całej przestrzeni]
\label{prop:heat-kernel}
Niech $d\geq1$, $\kappa>0$, $u_0\in C_c^\infty(\R^d)$.
Dla $t>0$, $x\in\R^d$ zdefiniujmy
\begin{equation}\label{eq:heat-kernel}
 \Gamma_\kappa(t,x)=\frac{1}{(4\pi\kappa t)^{d/2}}
 \exp\!\left(-\frac{|x|^2}{4\kappa t}\right),\qquad
 u(t,x)=\int_{\R^d}\Gamma_\kappa(t,x-y)u_0(y)\,
 \mathrm d\lambda_d(y).
\end{equation}
Wówczas $u\in C^\infty((0,\infty)\times\R^d)$,
$u_t-\kappa\Delta_xu=0$ i
$u(t,\cdot)\to u_0$ jednostajnie, gdy $t\downarrow0$.
Jeśli dodatkowo $u_0\geq0$ i $u_0\not\equiv0$, to
$u(t,x)>0$ dla każdego $t>0$ oraz $x\in\R^d$.
\end{stwierdzenie}
\begin{proof}
Dla $t>0$ różniczkowanie wzoru na $\Gamma_\kappa$ daje
\[
 \partial_t\Gamma_\kappa=
 \left(-\frac d{2t}+\frac{|x|^2}{4\kappa t^2}\right)\Gamma_\kappa,
 \qquad
 \partial_{x_j}\Gamma_\kappa=-\frac{x_j}{2\kappa t}\Gamma_\kappa.
\]
Ponowne różniczkowanie względem $x_j$ i sumowanie po $j$
prowadzi do
$\Delta_x\Gamma_\kappa=
(-d/(2\kappa t)+|x|^2/(4\kappa^2t^2))\Gamma_\kappa$.
Stąd $\partial_t\Gamma_\kappa=\kappa\Delta_x\Gamma_\kappa$.
Niech $K$ będzie zwartym podzbiorem $(0,\infty)\times\R^d$
i niech $S=\supp u_0$. Na iloczynie odpowiednio powiększonego
zwartego otoczenia $K$ i $S$ każda ustalona pochodna
$(t,x)\mapsto\Gamma_\kappa(t,x-y)$ jest ograniczona
jednostajnie względem $(t,x,y)$: jest ciągła, a czas jest
oddzielony od zera. Iloczyn tej stałej i $|u_0(y)|$ jest
całkowalną majorantą, więc twierdzenie o różniczkowaniu całki
z parametrem stosuje się do pochodnych dowolnego rzędu.
Wynika stąd gładkość $u$ i
$u_t-\kappa\Delta_xu=\int_{\R^d}(\partial_t-\kappa\Delta_x)
\Gamma_\kappa(t,x-y)u_0(y)\,\mathrm d\lambda_d(y)=0$.

Podstawienie $z=2\sqrt{\kappa t}\,w$ ma jakobian
$(2\sqrt{\kappa t})^d$. Korzystając z twierdzenia Tonellego
dla nieujemnego iloczynu funkcji gaussowskich i wzoru
$\int_{\R}e^{-s^2}\,\mathrm d\lambda_1(s)=\sqrt\pi$, dostajemy
\[
 \int_{\R^d}\Gamma_\kappa(t,z)\,\mathrm d\lambda_d(z)
 =\pi^{-d/2}\int_{\R^d}e^{-|w|^2}\,\mathrm d\lambda_d(w)=1.
\]
Dlatego
$u(t,x)-u_0(x)=\int_{\R^d}\Gamma_\kappa(t,z)
\bigl(u_0(x-z)-u_0(x)\bigr)\,\mathrm d\lambda_d(z)$.
Funkcja $u_0$ jest jednostajnie ciągła. Ustalmy $\eps>0$
i wybierzmy $r>0$ tak, aby
$|u_0(x-z)-u_0(x)|<\eps$ dla $|z|<r$, niezależnie od $x$.
Część całki po $|z|<r$ ma moduł nie większy od $\eps$;
część po $|z|\geq r$ jest ograniczona przez
$2\|u_0\|_\infty\int_{|z|\geq r}\Gamma_\kappa(t,z)
\,\mathrm d\lambda_d(z)$, a ostatnia całka dąży do zera,
gdy $t\downarrow0$. Istotnie, po podstawieniu $z=\sqrt t\,w$
jest ona równa całce funkcji $\Gamma_\kappa(1,w)$ po
$\{|w|\geq r/\sqrt t\}$. Funkcje podcałkowe, przedłużone
przez zero na dopełnienie tego zbioru, zbiegają punktowo
do zera i są zdominowane przez całkowalną funkcję
$\Gamma_\kappa(1,w)$. Twierdzenie o zbieżności zdominowanej
daje żądaną granicę. Otrzymane oszacowanie jest niezależne
od $x$, a $\eps>0$ było dowolne, co dowodzi zbieżności
jednostajnej na $\R^d$.

Jeżeli $u_0\geq0$ i $u_0\not\equiv0$, to z ciągłości
wynika, że $u_0$ jest dodatnia na pewnym niepustym zbiorze
otwartym. Ponieważ $\Gamma_\kappa(t,x-y)>0$ dla wszystkich
$t>0$, $x,y\in\R^d$, całka we wzorze na $u(t,x)$
jest ściśle dodatnia.
\end{proof}

Jest to pierwsza istotna różnica między równaniem ciepła
a równaniem falowym: dla nieujemnych, niezerowych danych $u_0$
o zwartym nośniku, nawet bardzo małym, dla każdego $t>0$
rozwiązanie równania ciepła jest dodatnie w całej $\R^d$.
Nie oznacza to, że zaburzenie ma wszędzie porównywalną
wielkość; daleko od nośnika danych wartości są bardzo małe.

\begin{stwierdzenie}[Energetyczna jednoznaczność równania ciepła]
\label{prop:heat-uniqueness}
Niech $d\geq1$, $\Omega\subset\R^d$ będzie obszarem ograniczonym
o brzegu klasy $C^1$, $T>0$ i $\kappa>0$.
Jeśli $u,v\in C^{1,2}([0,T]\times\overline\Omega)$ są funkcjami rzeczywistymi i
spełniają $w_t-\kappa\Delta_xw=0$ dla $w=u-v$
na $(0,T)\times\Omega$, a ponadto
$u(0,\cdot)=v(0,\cdot)$ w $\Omega$ oraz
$u=v$ na $[0,T]\times\partial\Omega$, to $u=v$
na $[0,T]\times\overline\Omega$.
\end{stwierdzenie}
\begin{proof}
Niech $E(t)=\frac12\int_\Omega w(t,x)^2\,
\mathrm d\lambda_d(x)$. Funkcje $w$ i $w_t$ są ciągłe na zwartym zbiorze
$[0,T]\times\overline\Omega$, więc iloczyn $ww_t$ jest ograniczony.
Ponieważ $\Omega$ ma skończoną miarę, różniczkowanie pod znakiem
całki jest uzasadnione twierdzeniem o zbieżności zdominowanej.
Dla $0<t<T$, z równania $w_t=\kappa\Delta_xw$ i twierdzenia
Gaussa zastosowanego do pola $w\nabla_xw$ otrzymujemy
\begin{align*}
 E'(t)&=\int_\Omega w(t,x)w_t(t,x)\,\mathrm d\lambda_d(x)\\
 &=\kappa\int_{\partial\Omega}w(t,x)\partial_\nu w(t,x)\,
 \mathrm d\sigma(x)
 -\kappa\int_\Omega|\nabla_xw(t,x)|^2\,\mathrm d\lambda_d(x)\\
 &=-\kappa\int_\Omega|\nabla_xw(t,x)|^2\,\mathrm d\lambda_d(x)\leq0.
\end{align*}
Całka brzegowa znika, ponieważ $w=0$ na brzegu.
Energia jest ciągła także dla $t=0,T$. Ponieważ $E(0)=0$
i $E(t)\geq0$, jej monotoniczność daje $E(t)=0$ dla każdego
$t\in[0,T]$. Gdyby dla pewnego $t$ oraz $x\in\Omega$ zachodziło
$w(t,x)\ne0$, ciągłość dawałaby kulę zawartą w $\Omega$,
na której $w(t,\cdot)^2$ byłoby ograniczone z dołu przez dodatnią
stałą; przeczyłoby to $E(t)=0$. Zatem $w=0$ w $\Omega$
dla każdego $t$, a przez ciągłość także na $\overline\Omega$.
\end{proof}

\subsection{Równanie falowe: skończona prędkość propagacji}

\begin{twierdzenie}[Wzór d'Alemberta]
\label{thm:dalembert}
Niech $c>0$, $u_0\in C^2(\R)$ i $u_1\in C^1(\R)$.
Jedynym rozwiązaniem klasy $C^2$ na $[0,\infty)\times\R$
zagadnienia
$u_{tt}-c^2u_{xx}=0$, $u(0,x)=u_0(x)$,
$u_t(0,x)=u_1(x)$ jest
\begin{equation}\label{eq:dalembert}
 u(t,x)=\frac{u_0(x-ct)+u_0(x+ct)}{2}
 +\frac{1}{2c}\int_{x-ct}^{x+ct}u_1(s)\,\mathrm d\lambda_1(s).
\end{equation}
Jeśli $\supp u_0\cup\supp u_1\subset[-R,R]$ dla $R>0$,
to $u(t,x)=0$ dla $|x|>R+ct$.
\end{twierdzenie}
\begin{proof}
Wprowadzamy zmienne $\xi=x+ct$ i $\eta=x-ct$.
Odwzorowanie jest odwracalne, gdyż
$t=(\xi-\eta)/(2c)$ i $x=(\xi+\eta)/2$.
Dla $v(\xi,\eta)=u((\xi-\eta)/(2c),(\xi+\eta)/2)$
reguła łańcuchowa daje $\partial_t=c(\partial_\xi-\partial_\eta)$
i $\partial_x=\partial_\xi+\partial_\eta$ na funkcjach
złożonych. Ponieważ pochodne mieszane funkcji klasy $C^2$
są równe, mamy
$u_{tt}-c^2u_{xx}=-4c^2v_{\xi\eta}$ w odpowiadających sobie
punktach. Z $c>0$ wynika równoważność z równaniem
$v_{\xi\eta}=0$. Obraz półpłaszczyzny $t\geq0$ to
$\{(\xi,\eta):\xi\geq\eta\}$, nie cała płaszczyzna.
Dla każdego ustalonego $\xi$ pochodna $v_\xi(\xi,\eta)$
jest stała względem $\eta<\xi$, ponieważ $v_{\xi\eta}=0$.
Z ciągłości aż do przekątnej jest równa $a(\xi)=v_\xi(\xi,\xi)$.
Funkcja $a$ jest klasy $C^1$, więc istnieje $F\in C^2(\R)$
z $F'=a$. Dla każdego ustalonego $\eta$ funkcja
$v(\xi,\eta)-F(\xi)$ ma zerową pochodną względem $\xi>\eta$,
a z ciągłości jej wartość jest równa
$G(\eta)=v(\eta,\eta)-F(\eta)$. Funkcja $G$ jest klasy $C^2$.
Otrzymujemy zatem $v(\xi,\eta)=F(\xi)+G(\eta)$ dla $\xi\geq\eta$.
Zatem $u(t,x)=F(x+ct)+G(x-ct)$.
Warunki początkowe dają
$F(x)+G(x)=u_0(x)$ i
$cF'(x)-cG'(x)=u_1(x)$.
Różniczkując pierwszą równość i rozwiązując układ względem
$F',G'$, otrzymujemy
$F'=(u_0'+u_1/c)/2$ i
$G'=(u_0'-u_1/c)/2$.
Dla $\xi=x+ct$ i $\eta=x-ct$ korzystamy z równości
$F(\xi)+G(\eta)=F(\eta)+G(\eta)+\int_\eta^\xi F'(s)\,\mathrm d\lambda_1(s)$.
Po podstawieniu $F(\eta)+G(\eta)=u_0(\eta)$ oraz wzoru na $F'$
otrzymujemy
\[
 F(\xi)+G(\eta)=\frac{u_0(\xi)+u_0(\eta)}2
 +\frac1{2c}\int_\eta^\xi u_1(s)\,\mathrm d\lambda_1(s),
\]
czyli~\eqref{eq:dalembert}.
Dla sprawdzenia istnienia niech $H$ będzie funkcją pierwotną
$u_1$, klasy $C^2$. Prawa strona~\eqref{eq:dalembert} jest sumą
funkcji $(u_0/2+H/(2c))(x+ct)$ oraz
$(u_0/2-H/(2c))(x-ct)$, a więc jest klasy $C^2$.
Dla każdej funkcji $h\in C^2(\R)$ mamy
$\partial_t^2h(x\pm ct)=c^2h''(x\pm ct)$ i
$\partial_x^2h(x\pm ct)=h''(x\pm ct)$, więc suma spełnia
równanie falowe. Przy $t=0$ suma jest równa $u_0(x)$,
a jej pochodna czasowa jest równa $u_1(x)$.
Każde rozwiązanie musiało mieć postać~\eqref{eq:dalembert},
co dowodzi również jednoznaczności.

Jeżeli $|x|>R+ct$, to oba punkty $x-ct$, $x+ct$ leżą
po tej samej stronie przedziału $[-R,R]$, a przedział
$[x-ct,x+ct]$ nie przecina nośnika $u_1$.
Wszystkie składniki~\eqref{eq:dalembert} są więc równe zeru.
\end{proof}

\begin{center}
\begin{tikzpicture}[>=Stealth,scale=1]
 \coordinate (P) at (.7,3.0);
 \coordinate (L) at (-1.6,0);
 \coordinate (R) at (3.0,0);
 \fill[gray!14] (L)--(P)--(R)--cycle;
 \draw[->] (-4.0,0)--(4.05,0) node[right] {$x$};
 \draw[->,gray!70!black] (0,-.2)--(0,3.9) node[above,black] {$t$};
 \draw[blue!65!black,thick] (L)--(P)--(R);
 \draw[very thick] (L)--(R);
 \fill (P) circle (2pt);
 \node[above right] at (P) {$(x_0,t_0)$};
 \node[below] at (L) {$x_0-ct_0$};
 \node[below] at (R) {$x_0+ct_0$};
 \node[align=center,fill=gray!14,inner sep=2pt] at (.7,1.0) {obszar\\zależności};
\end{tikzpicture}
\par\smallskip\small
Dla równania falowego w jednym wymiarze wartość w punkcie
$(x_0,t_0)$ zależy wyłącznie od danych początkowych na
$[x_0-ct_0,x_0+ct_0]$.
\end{center}

\begin{stwierdzenie}[Lokalny bilans energii w zwężającym się przedziale]
\label{prop:pde-v31-wave-cone}
Niech $c>0$, $a,b\in\R$, $a<b$, $0<T<(b-a)/(2c)$ oraz
niech rzeczywista funkcja $u\in C^2([0,T]\times\R)$ spełnia $u_{tt}-c^2u_{xx}=0$.
Dla $0\le t\le T$ połóżmy $L(t)=a+ct$, $R(t)=b-ct$ i
\[
 E(t)=\frac12\int_{L(t)}^{R(t)}
 \bigl(u_t(t,x)^2+c^2u_x(t,x)^2\bigr)\,\mathrm d\lambda_1(x).
\]
Wtedy
\begin{equation}\label{eq:pde-v31-wave-cone}
 E'(t)=-\frac c2\bigl(u_t-cu_x\bigr)^2(t,R(t))
       -\frac c2\bigl(u_t+cu_x\bigr)^2(t,L(t))\le0.
\end{equation}
Jeśli ponadto $u(0,x)=u_t(0,x)=0$ dla $x\in[a,b]$,
to $u(t,x)=0$ dla $0\le t\le T$ i $L(t)\le x\le R(t)$.
\end{stwierdzenie}
\begin{proof}
Oznaczmy $e=(u_t^2+c^2u_x^2)/2$. Z równania falowego,
reguły iloczynu i równości pochodnych mieszanych otrzymujemy
$\partial_t e=u_tu_{tt}+c^2u_xu_{xt}
=c^2(u_tu_{xx}+u_xu_{tx})=c^2\partial_x(u_tu_x)$.
Reguła Leibniza dla całki o ruchomych końcach, z
$L'=c$ oraz $R'=-c$, daje
\[
 E'(t)=c^2[u_tu_x]_{L(t)}^{R(t)}
       -ce(t,R(t))-ce(t,L(t)).
\]
Przy prawym końcu wyrażenie $c^2u_tu_x-ce$ jest równe
$-c(u_t-cu_x)^2/2$, a przy lewym końcu wyrażenie
$-c^2u_tu_x-ce$ jest równe $-c(u_t+cu_x)^2/2$.
Otrzymujemy~\eqref{eq:pde-v31-wave-cone} dla $0<t<T$
(i z jednostronnymi pochodnymi na końcach przedziału czasu).
Jeśli dane początkowe są zerowe na $[a,b]$, to również
$u_x(0,x)=0$ na $(a,b)$, a przez ciągłość także w końcach
przedziału. Stąd $E(0)=0$.
Ponieważ $E\ge0$ i $E$ jest nierosnąca, mamy $E(t)=0$.
Ciągłość pochodnych daje $u_t=u_x=0$ w otwartym stożku
$0<t<T$, $L(t)<x<R(t)$. Dla każdego punktu tego stożka
odcinek pionowy dla $0<s\leq t$ pozostaje w otwartym stożku,
gdyż $L(s)\leq L(t)<x<R(t)\leq R(s)$. Całkowanie
$u_t(s,x)=0$ i ciągłość przy $s=0$ dają $u(t,x)=u(0,x)=0$. Ciągłość daje też równość na jego brzegu.
\end{proof}

\begin{uwaga}[Dlaczego bilans lokalny jest ważniejszy od samego wzoru]
Oszacowanie~\eqref{eq:pde-v31-wave-cone} dowodzi skończonej
prędkości propagacji bez użycia jawnego wzoru d'Alemberta.
Ta sama idea, po modyfikacji geometrii i wyrazów brzegowych,
pojawia się przy wielowymiarowym równaniu falowym i układach
hiperbolicznych; tam jawny wzór na ogół nie jest dostępny.
\end{uwaga}

\section{Poprawność w sensie Hadamarda}
\label{sec:pde-hadamard}

Wyrażenie „rozwiązać równanie” nie jest wystarczające,
jeśli nie kontrolujemy wpływu niewielkich perturbacji danych.
Ustalmy przestrzenie danych $X$ i rozwiązań $Y$, wyposażone
w określone topologie, oraz dopuszczalny zbiór danych
$D\subset X$.

\begin{definicja}[Poprawność Hadamarda]
Zagadnienie jest \emph{poprawnie postawione w sensie Hadamarda}
w parze przestrzeni $(X,Y)$ na zbiorze danych $D\subset X$,
jeśli:
\begin{enumerate}
\item dla każdego $f\in D$ istnieje rozwiązanie $u\in Y$;
\item takie rozwiązanie jest jednoznaczne w $Y$;
\item odwzorowanie rozwiązujące $S\colon D\to Y$, $S(f)=u$,
jest ciągłe względem topologii podprzestrzeni na $D$.
\end{enumerate}
W zależności od zagadnienia „dane” mogą obejmować prawą stronę,
warunki początkowe, warunki brzegowe lub ich kombinację.
\end{definicja}

Trzeci warunek jest równie istotny jak pozostałe: dane
pochodzące z pomiarów lub obliczeń numerycznych zawsze
zawierają niewielki błąd.

\begin{kontrprzyklad}[Niestabilne zagadnienie Cauchy'ego dla laplasjanu]
\label{ex:laplace-illposed}
Niech $\Omega=(0,\pi)\times(0,1)$ i rozważmy równanie
$u_{xx}+u_{yy}=0$ w $\Omega$ z danymi na dolnym brzegu
$u(x,0)=0$ i $u_y(x,0)=g(x)$.
Dla każdej liczby naturalnej $k\geq1$ określmy
\begin{equation}\label{eq:hadamard-counterexample}
 \eps_k=e^{-k/2},\qquad
 u_k(x,y)=\frac{\eps_k}{k}\sinh(ky)\sin(kx).
\end{equation}
Mamy $(u_k)_{xx}=-k^2u_k$ oraz $(u_k)_{yy}=k^2u_k$,
więc $u_k$ jest funkcją harmoniczną. Ponadto $\sinh(0)=0$ i
$\cosh(0)=1$ dają $u_k(x,0)=0$ oraz
$\partial_yu_k(x,0)=\eps_k\sin(kx)$.
Dla każdego ustalonego $r\in\N_0$ zachodzi
\[
 \|g_k\|_{C^r([0,\pi])}\leq C_r k^r e^{-k/2}\longrightarrow0,
 \qquad g_k(x)=\eps_k\sin(kx),
\]
przy stałej $C_r$ niezależnej od $k$: każda pochodna rzędu
$j\leq r$ jest ograniczona w module przez $k^j e^{-k/2}$.
Dla ustalonego $r$ wykładnik $r\log k-k/2$ dąży do $-\infty$.
A więc dane dążą do zera nawet w topologii $C^\infty([0,\pi])$,
wyznaczonej przez półnormy $C^r$ dla $r\in\N_0$.
Tymczasem na poziomie $y=3/4$ otrzymujemy
\[
 \|u_k(\cdot,3/4)\|_{L^2(0,\pi)}
 =\sqrt{\frac\pi2}\,\frac{e^{-k/2}}{k}\sinh(3k/4)
 \underset{k\to\infty}{\sim}
 \sqrt{\frac\pi8}\,\frac{e^{k/4}}{k}
 \longrightarrow\infty.
\]
W obliczeniu użyliśmy tożsamości
$\int_0^\pi\sin^2(kx)\,\mathrm d\lambda_1(x)=\pi/2$
oraz $\sinh(3k/4)=e^{3k/4}(1-e^{-3k/2})/2$.
Przyjmijmy przestrzeń danych $X=C^\infty([0,\pi])$,
a za $D\subset X$ dowolny zbiór danych zawierający $0$ i wszystkie $g_k$.
Niech $Y$ będzie przestrzenią dopuszczalnych rozwiązań klasycznej
regularności, zawierającą $0$ i wszystkie $u_k$, w której
odwzorowanie $u\mapsto u(\cdot,3/4)$ do $L^2(0,\pi)$ jest ciągłe.
Gdyby zagadnienie było poprawnie postawione w sensie Hadamarda,
jednoznaczność wymuszałaby $S(0)=0$ oraz $S(g_k)=u_k$.
Ciągłość $S$ implikowałaby wtedy zbieżność przekrojów
$u_k(\cdot,3/4)$ do zera w $L^2(0,\pi)$, sprzeczną z obliczeniem.
Tak więc poprawność w sensie Hadamarda nie zachodzi w tych
przestrzeniach; jeśli rozwiązania są jednoznaczne, zawodzi
ciągła zależność od danych. Prosta $y=0$ jest niecharakterystyczna
dla operatora Laplace'a przedłużonego na otoczenie brzegu,
gdyż jej kowektor normalny $(0,1)$ nie zeruje symbolu.
\end{kontrprzyklad}

\begin{uwaga}[Dlaczego to nie przeczy zasadzie maksimum?]
We wniosku~\ref{cor:dirichlet-stability} podawaliśmy
\emph{wartości funkcji na całym brzegu}.
W kontrprzykładzie~\ref{ex:laplace-illposed} zadaliśmy
\emph{wartość i pochodną normalną tylko na części brzegu}.
To dwa zupełnie różne zagadnienia. W szczególności
zasada maksimum nie oszacowuje rozwiązania w prostokącie
przez dane Cauchy'ego wyłącznie na dolnej krawędzi.
\end{uwaga}

\begin{kontrprzyklad}[Równanie ciepła odwrócone w czasie]
\label{ex:backward-heat}
Niech $\kappa>0$, $T>0$, $\T=\R/(2\pi\mathbb Z)$.
Rozważmy na $[0,T]\times\T$ równanie
$u_t=-\kappa u_{xx}$. Dla każdej liczby całkowitej $k\geq1$
funkcja
\[
 u_k(t,x)=\exp\!\bigl(\kappa k^2(t-T/2)\bigr)\sin(kx)
\]
jest dobrze określona na torusie, ponieważ $k$ jest liczbą całkowitą.
Mamy $(u_k)_t=\kappa k^2u_k$ i $(u_k)_{xx}=-k^2u_k$, zatem
jest to gładkie rozwiązanie. Jej dane początkowe
$u_k(0,x)=e^{-\kappa k^2T/2}\sin(kx)$ dążą do zera
w $C^\infty(\T)$: pochodna rzędu $j$ ma normę supremum
nie większą niż $k^je^{-\kappa k^2T/2}$, która dla ustalonego
$j$ dąży do zera. Na torusie przyjmujemy miarę będącą obrazem
miary Lebesgue'a z $[0,2\pi)$ przez odwzorowanie ilorazowe;
jej masa wynosi $2\pi$. Ponieważ
$\int_0^{2\pi}\sin^2(kx)\,\mathrm d\lambda_1(x)=\pi$,
mamy $\|u_k(T,\cdot)\|_{L^2(\T)}
=\sqrt\pi\,e^{\kappa k^2T/2}\to\infty$.
Dla przestrzeni danych z topologią $C^\infty(\T)$ i przestrzeni
rozwiązań zawierającej $0,u_k$ oraz kontrolującej normę
$L^2(\T)$ w chwili $T$, jednoznaczność i ciągła zależność od
danych nie mogą zachodzić jednocześnie, zgodnie z argumentem
z poprzedniego kontrprzykładu. Zatem problem ewolucji wstecznej
nie jest poprawnie postawiony w sensie Hadamarda;
formalna zmiana znaku przy laplasjanie odwraca mechanizm
wygładzania w gwałtowne wzmacnianie wysokich częstości.
\end{kontrprzyklad}

\section{Skalowanie i związek typu równania z geometrią}
\label{sec:pde-scaling}

Dla równań ewolucyjnych czas i przestrzeń nie zawsze
skalują się w ten sam sposób. Współczynniki $c$ i $\kappa$
mają odpowiednio wymiary długość/czas oraz
długość$^2$/czas. To sugeruje dwa odmienne skalowania.

\begin{stwierdzenie}[Skalowania równań ciepła i falowego]
Niech $\lambda,\kappa,c>0$, $d\geq1$ i niech
$Q\subset\R\times\R^d$ będzie zbiorem otwartym.
\begin{enumerate}
\item Jeśli $u$ jest rozwiązaniem klasycznym
$u_t-\kappa\Delta_xu=0$ na $Q$, to
$u_\lambda(t,x)=u(\lambda^2t,\lambda x)$
rozwiązuje to samo równanie na
$Q_\lambda=\{(t,x):(\lambda^2t,\lambda x)\in Q\}$.
\item Jeśli $u$ jest rozwiązaniem klasycznym
$u_{tt}-c^2\Delta_xu=0$ na $Q$, to
$v_\lambda(t,x)=u(\lambda t,\lambda x)$
rozwiązuje to samo równanie na
$\{(t,x):(\lambda t,\lambda x)\in Q\}$.
\end{enumerate}
\end{stwierdzenie}
\begin{proof}
Dziedziny funkcji przeskalowanych są otwarte jako przeciwobrazy
$Q$ przez ciągłe odwzorowania liniowe. Ponieważ odwzorowania
te osobno skalują czas i zmienne przestrzenne, zachowana
jest wymagana regularność klasyczna. Reguła łańcuchowa daje
$\partial_tu_\lambda=\lambda^2u_t(\lambda^2t,\lambda x)$
i $\Delta_xu_\lambda=\lambda^2\Delta_xu(\lambda^2t,\lambda x)$.
Stąd
$(\partial_t-\kappa\Delta_x)u_\lambda(t,x)
=\lambda^2(u_t-\kappa\Delta_xu)(\lambda^2t,\lambda x)=0$.
Podobnie $\partial_t^2v_\lambda=\lambda^2u_{tt}(\lambda t,\lambda x)$
oraz $\Delta_xv_\lambda=\lambda^2\Delta_xu(\lambda t,\lambda x)$,
więc $(\partial_t^2-c^2\Delta_x)v_\lambda(t,x)
=\lambda^2(u_{tt}-c^2\Delta_xu)(\lambda t,\lambda x)=0$.
\end{proof}

Przy skali przestrzennej $\ell>0$ podstawienie $x=\ell z$,
$t=\tau s$ do równania ciepła daje współczynnik
$\kappa\tau/\ell^2$ przy laplasjanie w zmiennej $z$.
Wybór $\tau=\ell^2/\kappa$ sprowadza go do jedynki.
Dla równania falowego analogiczny współczynnik wynosi
$c^2\tau^2/\ell^2$, więc naturalną skalą czasu jest $\ell/c$.
To różnica strukturalna, nie tylko rachunkowa.

\section{Co będzie potrzebne dalej?}
\label{sec:pde-outlook}

Podsumujmy rozróżnienia, które pozostaną ważne przez cały skrypt:
\begin{enumerate}
\item \emph{Równanie a zagadnienie.} Dwa takie same operatory
mogą prowadzić do różnych problemów przy innych danych
początkowych lub brzegowych.
\item \emph{Istnienie, jednoznaczność, stabilność.} Każde z tych
zagadnień wymaga osobnego rozumowania. Udowodnienie
jednoznaczności nie implikuje istnienia.
\item \emph{Typ operatora.} Eliptyczność, hiperboliczność
i paraboliczność opisują zasadniczo różne mechanizmy.
\item \emph{Klasa rozwiązań.} Rozwiązania klasyczne są tylko
początkiem. Naturalnymi obiektami bywają funkcje Sobolewa,
dystrybucje oraz rozwiązania słabe i lepkościowe.
\item \emph{Rola geometrii.} Symbol główny, charakterystyki,
struktura brzegu i skalowanie decydują o propagacji
i regularności rozwiązań.
\end{enumerate}

W następnym rozdziale rozwiniemy metodę charakterystyk dla
równań pierwszego rzędu: od liniowego transportu przez
quasiliniowe zagadnienia Cauchy'ego po tworzenie się
fal uderzeniowych.

% Rozdział 2.

\chapter[Charakterystyki i PDE pierwszego rzędu]{Równania pierwszego rzędu i metoda charakterystyk}
\label{chap:pde-characteristics}

W tym rozdziale zamienimy równania pierwszego rzędu na układy
równań zwyczajnych. Wersja liniowa metody pozwala rozwiązywać
równania transportu i ciągłości, natomiast wersja quasiliniowa
wyjaśnia powstawanie osobliwości z gładkich danych. Dla równań
w pełni nieliniowych skonstruujemy układ charakterystyczny
Charpita i zastosujemy go do równań Hamiltona--Jacobiego.

W całym rozdziale $d\geq1$, $T>0$ oraz $\lambda_d$ oznacza
miarę Lebesgue'a na $\R^d$. Wszystkie pochodne występujące
w sformułowaniach klasycznych rozumiemy w sensie punktowym.
W szczególności przez $C^1([0,T]\times\R^d)$ rozumiemy
funkcje o pochodnych pierwszego rzędu ciągłych aż do
$t=0,T$ (jednostronnie w tych punktach).

\section{Pochodna wzdłuż krzywej i najprostsze równanie transportu}
\label{sec:pde2-basic}

Niech $u\in C^1((0,T)\times\R^d)$ i niech
$X\colon I\to\R^d$ będzie krzywą klasy $C^1$, gdzie
$I\subset(0,T)$ jest przedziałem. Reguła łańcuchowa daje
\begin{equation}\label{eq:pde2-chain}
 \frac{\mathrm d}{\mathrm dt}u(t,X(t))
 =u_t(t,X(t))+\nabla_xu(t,X(t))\cdot\dot X(t).
\end{equation}
Dobierając $\dot X(t)$ do współczynników równania,
możemy przekształcić PDE w równanie zwyczajne
wzdłuż krzywej $t\mapsto(t,X(t))$. Taką krzywą nazywamy
\emph{charakterystyką} równania transportu.

\begin{stwierdzenie}[Transport ze stałą prędkością i źródłem]
\label{prop:pde2-constant}
Niech $c\in\R^d$, $u_0\in C^1(\R^d)$,
$f\in C^1([0,T]\times\R^d)$. Istnieje dokładnie jedno
rozwiązanie $u\in C^1([0,T]\times\R^d)$ zagadnienia
\[
 u_t+c\cdot\nabla_xu=f(t,x),\qquad u(0,x)=u_0(x).
\]
Jest ono dane wzorem
\begin{equation}\label{eq:pde2-constant-formula}
 u(t,x)=u_0(x-tc)+\int_0^t f(s,x-(t-s)c)\,\mathrm d\lambda_1(s).
\end{equation}
\end{stwierdzenie}
\begin{proof}
Charakterystyka wychodząca z $\xi$ ma postać
$X(t;\xi)=\xi+tc$. Dla każdego rozwiązania klasy $C^1$
reguła \eqref{eq:pde2-chain} pokazuje, że
$\frac{\mathrm d}{\mathrm dt}u(t,\xi+tc)=f(t,\xi+tc)$.
Całkując od $0$ do $t$ i podstawiając $\xi=x-tc$,
otrzymujemy \eqref{eq:pde2-constant-formula}; dowodzi to
jednoznaczności. Odwrotnie, oznaczmy prawą stronę wzoru
\eqref{eq:pde2-constant-formula} przez $u(t,x)$.
Na każdym zwartym otoczeniu parametrów $(t,x)$ funkcja
podcałkowa oraz jej pochodne względem $t$ i $x$ są ciągłe
i ograniczone dla $s\in[0,T]$. Reguła Leibniza dla całki
z parametrem daje więc ciągłe pochodne pierwszego rzędu,
także jednostronnie na końcach przedziału czasu.
Po podstawieniu $x=\xi+tc$ otrzymujemy
$u(t,\xi+tc)=u_0(\xi)+\int_0^t f(s,\xi+sc)\,\mathrm d\lambda_1(s)$.
Różniczkując względem $t$, dostajemy
$u_t(t,\xi+tc)+c\cdot\nabla_xu(t,\xi+tc)=f(t,\xi+tc)$.
Ponieważ dla ustalonego $t$ przesunięcie $\xi\mapsto\xi+tc$
jest bijekcją, równanie zachodzi dla każdego $x\in\R^d$.
Dla $t=0$ całka znika, więc $u(0,x)=u_0(x)$.
\end{proof}

\begin{przyklad}[Tłumienie wzdłuż charakterystyk]
Niech $\gamma>0$. Dla równania
$u_t+c\cdot\nabla u+\gamma u=0$ z danymi $u_0$
otrzymujemy $u(t,x)=e^{-\gamma t}u_0(x-tc)$.
Transport przesuwa profil, podczas gdy czynnik wykładniczy
tłumi jego amplitudę.
\end{przyklad}

\section{Zmienne pole prędkości: przepływ charakterystyk}
\label{sec:pde2-linear}

Rozważmy równanie
\begin{equation}\label{eq:pde2-linear}
 u_t+b(t,x)\cdot\nabla_xu+c(t,x)u=f(t,x)
 \quad\text{na }(0,T)\times\R^d.
\end{equation}
Dla przejrzystości przyjmiemy dość mocne założenia
zapewniające globalny przepływ; później będziemy je osłabiać.

\begin{twierdzenie}[Metoda charakterystyk dla transportu liniowego]
\label{thm:pde2-linear}
Niech $b\colon[0,T]\times\R^d\to\R^d$ będzie klasy $C^1$
i niech istnieje $M>0$ takie, że
$\|D_xb(t,x)\|\leq M$ i $|b(t,x)|\leq M(1+|x|)$
dla wszystkich $(t,x)\in[0,T]\times\R^d$, gdzie
$\|A\|=\sup_{|v|=1}|Av|$ oznacza normę operatorową
indukowaną przez normę euklidesową. Niech $c,f\in C^1([0,T]\times\R^d)$
oraz $u_0\in C^1(\R^d)$. Wówczas istnieje dokładnie jedno
rozwiązanie $u\in C^1([0,T]\times\R^d)$ zagadnienia
\eqref{eq:pde2-linear} z $u(0,x)=u_0(x)$.

Niech $X(t;\xi)$ będzie rozwiązaniem
\begin{equation}\label{eq:pde2-flow-ode}
 \dot X(t;\xi)=b(t,X(t;\xi)),\qquad X(0;\xi)=\xi.
\end{equation}
Dla każdego $t\in[0,T]$ odwzorowanie
$X_t\colon\xi\mapsto X(t;\xi)$ jest dyfeomorfizmem klasy $C^1$
przestrzeni $\R^d$ na siebie. Rozwiązanie jest określone przez
\begin{equation}\label{eq:pde2-solution}
 \begin{aligned}
 u(t,X(t;\xi))={}&e^{-C(t;\xi)}u_0(\xi)\\
 &+\int_0^t e^{-(C(t;\xi)-C(s;\xi))}
 f(s,X(s;\xi))\,\mathrm ds,
 \end{aligned}
\end{equation}
gdzie $C(t;\xi)=\int_0^t c(r,X(r;\xi))\,\mathrm dr$.
\end{twierdzenie}
\begin{proof}
Pole $b$ jest lokalnie lipschitzowskie względem $x$,
a dzięki warunkowi na $D_xb$ jest jednostajnie
lipschitzowskie na całej przestrzeni. Na przedziale istnienia rozwiązania jego równanie
całkowe i warunek wzrostu dają
\[
 1+|X(t;\xi)|\le 1+|\xi|
 +M\int_0^t(1+|X(s;\xi)|)\,\mathrm d\lambda_1(s).
\]
Z nierówności Gronwalla wynika
$1+|X(t;\xi)|\le(1+|\xi|)e^{Mt}$.
Gdyby maksymalny przedział istnienia kończył się przed $T$,
trajektoria pozostawałaby w ustalonej kuli domkniętej.
Pole $b$ jest ograniczone na iloczynie tej kuli i $[0,T]$,
więc $X$ jest lipschitzowska względem czasu i ma granicę
w końcu przedziału istnienia. Lokalne twierdzenie o istnieniu
ODE z tym warunkiem początkowym przedłuża trajektorię,
co przeczy maksymalności. To samo ograniczenie zapewnia
istnienie granicy przy $t=T$; równanie całkowe i ciągłość
pola dają jednostronną pochodną w tym punkcie. Twierdzenie o zależności rozwiązań
ODE od warunków początkowych daje $C^1$-gładkość $X$
względem $(t,\xi)$.

Dla ustalonego $t\in[0,T]$ i $x\in\R^d$ rozważmy
równanie $\dot Y(s)=-b(t-s,Y(s))$, $Y(0)=x$,
dla $s\in[0,t]$. Pole po prawej stronie spełnia te same
ograniczenia wzrostu i pochodnej przestrzennej co $b$,
więc poprzedni argument zapewnia istnienie na całym
$[0,t]$. Krzywa $r\mapsto Y(t-r)$ rozwiązuje pierwotne
równanie na $[0,t]$. Z jednoznaczności wynika
$X(t;Y(t))=x$ oraz, jeśli $x=X(t;\xi)$, to $Y(t)=\xi$.
Odwrotność $X_t^{-1}$ ma zatem postać $x\mapsto Y(t)$
i jest klasy $C^1$ z twierdzenia o zależności od danych.
W szczególności $D_\xi X(t;\xi)$ jest odwracalna.

Rozważmy odwzorowanie $\Phi(t,\xi)=(t,X(t;\xi))$.
Jego macierz pochodnej ma postać
\[
 D\Phi(t,\xi)=
 \begin{pmatrix}1&0\\ \partial_t X(t;\xi)&D_\xi X(t;\xi)\end{pmatrix},
\]
więc jej wyznacznik jest niezerowy.
Wewnątrz przedziału czasu twierdzenie o funkcji odwrotnej
zapewnia lokalną odwrotność klasy $C^1$. Ponieważ każde $X_t$
jest bijekcją, lokalne odwrotności są zgodne i dają wspólną
odwrotność $\Phi^{-1}(t,x)=(t,X_t^{-1}(x))$.
Na końcach przedziału stosujemy lokalnie to samo twierdzenie
do przedłużenia $X$ klasy $C^1$ względem czasu:
dla $t<0$ przyjmujemy $2X(0;\xi)-X(-t;\xi)$, a dla $t>T$
przyjmujemy $2X(T;\xi)-X(2T-t;\xi)$.
Pochodne czasowe i przestrzenne obu przedłużeń zgadzają się
z odpowiednimi jednostronnymi pochodnymi na brzegu.
Macierz pochodnej w punkcie brzegowym pozostaje
nieosobliwa. Otrzymujemy więc wspólną regularność $C^1$
odwrotności aż do końców przedziału czasu.

Dla dowolnego rozwiązania $u$ postawmy
$U(t;\xi)=u(t,X(t;\xi))$. Z \eqref{eq:pde2-chain}
wynika $\dot U+c(t,X(t;\xi))U=f(t,X(t;\xi))$.
Mnożąc przez czynnik całkujący $e^{C(t;\xi)}$,
otrzymujemy \eqref{eq:pde2-solution}.

Odwrotnie, prawa strona \eqref{eq:pde2-solution} wyznacza
funkcję $U\in C^1$; dzięki dyfeomorfizmowi $X_t$
definiuje ona $u(t,x)=U(t;X_t^{-1}(x))$ klasy $C^1$.
Różniczkowanie po charakterystykach daje
\eqref{eq:pde2-linear}. Każde inne rozwiązanie prowadzi
do tej samej funkcji $U$, co kończy dowód jednoznaczności.
\end{proof}

\begin{przyklad}[Transport w polu dylatacji]
Niech $b(x)=x$, $c=f=0$, $u_0\in C^1(\R^d)$.
Wówczas $X(t;\xi)=e^t\xi$ i
$u(t,x)=u_0(e^{-t}x)$. Profil rozszerza się
wraz z przepływem, ale wartości wzdłuż charakterystyk
nie zmieniają się.
\end{przyklad}

\begin{uwaga}[Odwracalność przepływu a odwracalność ewolucji]
W twierdzeniu \ref{thm:pde2-linear} w skończonym czasie
przepływ przestrzenny jest odwracalny. Nie oznacza to,
że każdy operator ewolucyjny PDE jest odwracalny
w dobrze postawionej klasie funkcji: dla równania ciepła
występuje wygładzanie, a dla transportu quasiliniowego
charakterystyki mogą przestać tworzyć układ współrzędnych.
\end{uwaga}

\section{Jak zmieniają się objętość i masa?}
\label{sec:pde2-jacobian}

Transport wartości skalarnej i transport gęstości to
dwa różne zagadnienia. Dla tego samego pola $b$ porównajmy
$u_t+b\cdot\nabla u=0$ oraz
$\rho_t+\diver_x(\rho b)=0$. Drugie równanie nazywamy
\emph{równaniem ciągłości}; w nim dodatkowo występuje
czynnik $(\diver_x b)\rho$.

\begin{stwierdzenie}[Wzór Jacobiego--Liouville'a]
\label{prop:pde2-liouville}
Niech $b$ spełnia założenia twierdzenia
\ref{thm:pde2-linear}, a $X$ będzie przepływem
\eqref{eq:pde2-flow-ode}. Dla $J(t;\xi)=\det D_\xi X(t;\xi)$
zachodzi
\begin{equation}\label{eq:pde2-jacobian}
 \frac{\mathrm d}{\mathrm dt}J(t;\xi)
 =\bigl(\diver_x b\bigr)(t,X(t;\xi))J(t;\xi),
 \qquad J(0;\xi)=1.
\end{equation}
W szczególności
$J(t;\xi)=\exp\!\bigl(\int_0^t
(\diver_x b)(s,X(s;\xi))\,\mathrm ds\bigr)>0$.
\end{stwierdzenie}
\begin{proof}
Ustalmy $\xi\in\R^d$ i oznaczmy
$A(t)=D_\xi X(t;\xi)$. Zapiszmy równanie
\eqref{eq:pde2-flow-ode} w postaci całkowej:
\[
 X(t;\xi)=\xi+\int_0^t b(s,X(s;\xi))\,\mathrm d\lambda_1(s).
\]
Na iloczynie $[0,T]$ i dowolnego zwartego otoczenia $\xi$
funkcja podcałkowa oraz jej pochodna względem $\xi$
są ciągłe, a więc ograniczone. Można zatem różniczkować
pod znakiem całki. Z reguły łańcuchowej wynika
\[
 A(t)=I_d+\int_0^t D_xb(s,X(s;\xi))A(s)\,\mathrm d\lambda_1(s),
\]
gdzie $I_d$ jest macierzą jednostkową rozmiaru $d$.
Prawa strona jest klasy $C^1$ względem $t$, ponieważ
funkcja podcałkowa jest ciągła. Twierdzenie o pochodnej
całki o zmiennej górnej granicy daje
$\dot A(t)=D_xb(t,X(t;\xi))A(t)$ oraz $A(0)=I_d$.
W końcach przedziału pochodne rozumiemy jednostronnie.
Macierz $A(t)$ jest odwracalna na mocy twierdzenia
\ref{thm:pde2-linear}. Zatem wzór na pochodną wyznacznika
daje $\frac{\mathrm d}{\mathrm dt}\det A
=\det A\,\tr(A^{-1}\dot A)$.
Cykliczność śladu i $\dot A=D_xb\,A$ prowadzą do
$\tr(A^{-1}\dot A)=\tr(D_xb)=\diver_xb$.
Oznaczmy $h(t)=(\diver_xb)(t,X(t;\xi))$ oraz
$H(t)=\int_0^t h(s)\,\mathrm d\lambda_1(s)$.
Z równania \eqref{eq:pde2-jacobian} wynika
$\frac{\mathrm d}{\mathrm dt}(e^{-H(t)}J(t;\xi))=0$.
Ponieważ $J(0;\xi)=\det I_d=1$ oraz $H(0)=0$,
otrzymujemy $J(t;\xi)=e^{H(t)}>0$.
\end{proof}

\begin{twierdzenie}[Transport gęstości i zachowanie masy]
\label{thm:pde2-mass}
Niech $b\in C^2([0,T]\times\R^d;\R^d)$ spełnia
warunki wzrostu i lipschitzowskości z twierdzenia
\ref{thm:pde2-linear}. Zakładamy ponadto, że wszystkie
pochodne przestrzenne $b$ rzędu $1$ i $2$ są ograniczone.
Dla dowolnej $\rho_0\in C_c^1(\R^d)$ istnieje dokładnie jedno
rozwiązanie $\rho\in C^1([0,T]\times\R^d)$ równania
\[
 \rho_t+\diver_x(\rho b)=0,\qquad \rho(0,x)=\rho_0(x).
\]
Jest ono zadane zależnością
\begin{equation}\label{eq:pde2-mass-flow}
 \rho(t,X(t;\xi))J(t;\xi)=\rho_0(\xi).
\end{equation}
Ponadto $\int_{\R^d}\rho(t,x)\,\mathrm d\lambda_d(x)
=\int_{\R^d}\rho_0(x)\,\mathrm d\lambda_d(x)$ dla $t\in[0,T]$.
\end{twierdzenie}
\begin{proof}
Ponieważ $\diver_x(\rho b)=b\cdot\nabla\rho+
(\diver_xb)\rho$, twierdzenie \ref{thm:pde2-linear}
można zastosować z $c=\diver_xb$, $f=0$.
Założenie $b\in C^2$ zapewnia, że $\diver_x b\in C^1$,
więc twierdzenie zapewnia istnienie i jednoznaczność
rozwiązania klasy $C^1$. Wzdłuż charakterystyki zachodzi
$\frac{\mathrm d}{\mathrm dt}\rho(t,X(t;\xi))
=-(\diver_x b)(t,X(t;\xi))\rho(t,X(t;\xi))$.
Po zastosowaniu reguły iloczynu i \eqref{eq:pde2-jacobian}
otrzymujemy
\[
 \frac{\mathrm d}{\mathrm dt}
 \bigl(\rho(t,X(t;\xi))J(t;\xi)\bigr)=0.
\]
W chwili $t=0$ iloczyn wynosi $\rho_0(\xi)$,
co dowodzi \eqref{eq:pde2-mass-flow}.
Ponieważ $J(t;\xi)>0$, funkcja $\rho(t,X_t(\xi))$
jest niezerowa wtedy i tylko wtedy, gdy $\rho_0(\xi)\ne0$.
Homeomorfizm $X_t$ zachowuje domknięcia zbiorów,
więc $\supp\rho(t,\cdot)=X_t(\supp\rho_0)$.
Jest to zbiór zwarty jako ciągły obraz zbioru zwartego.
Obie gęstości są ciągłe i mają zwarte nośniki,
a zatem są całkowalne również wtedy, gdy zmieniają znak.
Możemy więc zastosować twierdzenie o zmianie zmiennych
dla dyfeomorfizmu $x=X_t(\xi)$.
Dodatniość $J$ i \eqref{eq:pde2-mass-flow} dają
\[
 \int_{\R^d}\rho(t,x)\,\mathrm d\lambda_d(x)
 =\int_{\R^d}\rho(t,X(t;\xi))J(t;\xi)
 \,\mathrm d\lambda_d(\xi)
 =\int_{\R^d}\rho_0(\xi)\,\mathrm d\lambda_d(\xi).
\]
\end{proof}

\begin{uwaga}[Warunki regularności]
Wzór \eqref{eq:pde2-mass-flow} zachowuje sens również
przy słabszych założeniach. W poprzednim dowodzie
zastosowaliśmy dodatkową regularność $b$, aby pozostać
w klasie rozwiązań klasycznych; teoria dystrybucyjna
pozwoli te wymagania osłabić.
\end{uwaga}

\begin{stwierdzenie}[Kontrola jakobianu i norm przy transporcie]
\label{prop:pde-v31-transport-lp}
Niech $T>0$, $d\ge1$, a $b$ spełnia założenia
twierdzenia~\ref{thm:pde2-linear} z $\|D_xb\|\le M$.
Niech $X_t$ i $J(t;\xi)=\det D_\xi X_t(\xi)$ będą
odpowiednio przepływem i jego jakobianem. Wówczas
\[
 e^{-dMt}\le J(t;\xi)\le e^{dMt}
 \qquad(0\le t\le T,\ \xi\in\R^d).
\]
Jeśli $u_0\in C^1(\R^d)\cap L^p(\R^d)$ dla
$1\le p<\infty$, a $u$ rozwiązuje
$u_t+b\cdot\nabla u=0$, $u(0)=u_0$, to
\begin{equation}\label{eq:pde-v31-transport-lp}
 e^{-dMt/p}\|u_0\|_p\le\|u(t,\cdot)\|_p
 \le e^{dMt/p}\|u_0\|_p.
\end{equation}
Dla ograniczonego $u_0\in C^1(\R^d)$ zachodzi ponadto
$\|u(t,\cdot)\|_\infty=\|u_0\|_\infty$.
Jeśli $\diver_xb=0$, to $J\equiv1$ i wszystkie wymienione
normy $L^p$ są zachowane.
\end{stwierdzenie}
\begin{proof}
Niech $(e_1,\ldots,e_d)$ będzie standardową bazą
ortonormalną $\R^d$. Dla każdej rzeczywistej macierzy
$A$ rozmiaru $d$, z nierówności Cauchy'ego--Schwarza,
\[
 |\tr A|\le\sum_{j=1}^d|\langle Ae_j,e_j\rangle|
 \le\sum_{j=1}^d |Ae_j|\le d\|A\|.
\]
Stosując to oszacowanie do $A=D_xb(t,x)$, otrzymujemy
$|(\diver_xb)(t,x)|\le dM$.
Wzór~\eqref{eq:pde2-jacobian} i jego rozwiązanie dają
pierwsze oszacowanie. Wzdłuż charakterystyk
$u(t,X_t(\xi))=u_0(\xi)$. Ponieważ $X_t$ jest
dyfeomorfizmem, wzór na zmianę zmiennych daje
\[
 \|u(t,\cdot)\|_p^p
 =\int_{\R^d}|u_0(\xi)|^pJ(t;\xi)\,
   \mathrm d\lambda_d(\xi).
\]
Z oszacowań na $J$ wynika
$e^{-dMt}\|u_0\|_p^p\le\|u(t,\cdot)\|_p^p
\le e^{dMt}\|u_0\|_p^p$. W szczególności całka jest
skończona; wyciągnięcie pierwiastka stopnia $p$
daje \eqref{eq:pde-v31-transport-lp}.
Dla ograniczonej funkcji ciągłej supremum modułu jest
równe supremum istotnemu względem miary Lebesgue'a:
jeśli liczba jest mniejsza od supremum, zbiór punktów,
w których moduł ją przekracza, jest niepusty i otwarty,
a więc ma dodatnią miarę. Ponieważ $X_t$ jest bijekcją,
$\sup_x|u(t,x)|=\sup_\xi|u_0(\xi)|$; stąd wynika równość
norm $L^\infty$. W przypadku zerowej dywergencji
równanie na $J$ ma rozwiązanie stałe równe $1$.
\end{proof}

\section{Niecharakterystyczne dane Cauchy'ego na hiperpowierzchni}
\label{sec:pde2-noncharacteristic}

Rozważmy równanie bez wyróżnionego czasu
\begin{equation}\label{eq:pde2-static}
 a(x)\cdot\nabla u(x)+c(x)u(x)=f(x),\qquad x\in\Omega,
\end{equation}
gdzie $\Omega\subset\R^n$ jest otwarta, $n\geq2$.
Gdy dane zadano na hiperpowierzchni $\Sigma$, istotne
jest, czy charakterystyki przecinają ją poprzecznie.

\begin{definicja}[Niecharakterystyczność dla transportu]
Niech $\Omega\subset\R^n$ będzie otwarta, $a\in C^1(\Omega;\R^n)$,
a $\Sigma\subset\Omega$ hiperpowierzchnią klasy $C^2$.
Mówimy, że $\Sigma$ jest niecharakterystyczna w punkcie
$x_0\in\Sigma$ dla operatora $a\cdot\nabla$, jeżeli
$a(x_0)\notin T_{x_0}\Sigma$. Jeśli w pewnym otwartym
otoczeniu $O\subset\Omega$ punktu $x_0$ mamy
$\Sigma\cap O=\{x\in O:\phi(x)=0\}$, gdzie $\phi\in C^2(O)$
i $\mathrm d\phi(x_0)\ne0$, to
$T_{x_0}\Sigma=\ker\mathrm d\phi(x_0)$. Niecharakterystyczność
jest zatem równoważna warunkowi
$a(x_0)\cdot\nabla\phi(x_0)\ne0$.
\end{definicja}

\begin{twierdzenie}[Lokalne zagadnienie Cauchy'ego]
\label{thm:pde2-cauchy}
Niech $\Omega\subset\R^n$ będzie otwarta,
$a\in C^1(\Omega;\R^n)$, $c,f\in C^1(\Omega)$,
$\Sigma\subset\Omega$
hiperpowierzchnią klasy $C^2$, a $x_0\in\Sigma$.
Niech $g\in C^1(\Sigma)$ w otoczeniu $x_0$.
Jeśli $\Sigma$ jest niecharakterystyczna w $x_0$ dla
operatora $a\cdot\nabla$, to istnieje otwarte otoczenie
$U\subset\Omega$ punktu $x_0$, na którym zagadnienie
\eqref{eq:pde2-static} z danymi $u|_{U\cap\Sigma}=g$
ma rozwiązanie $u\in C^1(U)$. Otoczenie $U$ można wybrać
tak, aby rozwiązanie było jedyne na $U$.
W szczególności rozwiązanie jest jednoznaczne jako zarodek
w $x_0$: dwa rozwiązania lokalne są równe na pewnym
wspólnym otoczeniu tego punktu.
\end{twierdzenie}
\begin{proof}
Wybierzmy lokalną parametryzację
$\gamma\colon V\to\Sigma$, gdzie $V\subset\R^{n-1}$
jest otwartym otoczeniem zera, $\gamma$ jest klasy $C^2$,
$\gamma(0)=x_0$, a wektory
$\partial_{y_1}\gamma(0),\ldots,\partial_{y_{n-1}}\gamma(0)$
są liniowo niezależne. Niech $\Phi_s(x)$ oznacza lokalny
przepływ pola $a$, tj. $\partial_s\Phi_s(x)=a(\Phi_s(x))$
i $\Phi_0(x)=x$. Twierdzenie o lokalnym istnieniu
przepływu pola klasy $C^1$ zapewnia wspólny przedział
$(-\delta_0,\delta_0)$ dla punktów początkowych w pewnym
otoczeniu $x_0$ oraz regularność $C^1$ względem czasu
i punktu początkowego. Zmniejszając $V$, możemy przyjąć,
że $\gamma(V)$ leży w tym otoczeniu oraz w dziedzinie danych $g$.
Zdefiniujmy $\Psi(s,y)=\Phi_s(\gamma(y))$ na
$(-\delta_0,\delta_0)\times V$.

Macierz pochodnej $D\Psi(0,0)$ ma kolumny
$a(x_0)$ oraz $\partial_{y_j}\gamma(0)$. Założenie
niecharakterystyczności mówi dokładnie, że tworzą bazę
$\R^n$. Z twierdzenia o funkcji odwrotnej istnieją
otoczenia $W$ punktu $(0,0)$ i $U$ punktu $x_0$,
dla których $\Psi\colon W\to U$ jest dyfeomorfizmem.
Wybierzmy lokalną funkcję definiującą $\phi$ dla $\Sigma$
w otoczeniu $x_0$, jak w poprzedniej definicji.
Ponieważ
$\partial_s(\phi\circ\Psi)(0,0)=\mathrm d\phi(x_0)a(x_0)\ne0$,
po zmniejszeniu otoczenia pochodna ta ma stały niezerowy znak.
Wybierzmy iloczyn $W'=(-\delta,\delta)\times V'\subset W$
zawierający $(0,0)$ tak, aby $\Psi(W')$ leżał w dziedzinie
$\phi$ i aby powyższy warunek znaku zachodził na $W'$.
Dla każdego $y\in V'$ mamy $\phi(\Psi(0,y))=0$.
Ścisła monotoniczność funkcji $s\mapsto\phi(\Psi(s,y))$
pokazuje, że jej jedynym zerem jest $s=0$.
Przyjmując $U=\Psi(W')$, otrzymujemy otwarte otoczenie $x_0$,
na którym $\Psi|_{W'}$ jest dyfeomorfizmem i
$U\cap\Sigma=\Psi(\{0\}\times V')$.
W szczególności cały odcinek od $(0,y)$ do $(s,y)$
leży w $W'$ dla każdego $(s,y)\in W'$.

Niech $v(s,y)=u(\Psi(s,y))$. Równanie
\eqref{eq:pde2-static} sprowadza się do ODE
$\partial_sv+c(\Psi(s,y))v=f(\Psi(s,y))$
z danymi $v(0,y)=g(\gamma(y))$. Dla
$C(s,y)=\int_0^s c(\Psi(r,y))\,\mathrm dr$ rozwiązanie ma
postać
\[
 v(s,y)=e^{-C(s,y)}g(\gamma(y))+
 \int_0^s e^{-(C(s,y)-C(r,y))}
 f(\Psi(r,y))\,\mathrm dr.
\]
Całki dla $s<0$ rozumiemy jako całki z odwróconą
orientacją granic. Na zwartych podzbiorach iloczynu $W'$
współczynniki $c\circ\Psi$, $f\circ\Psi$ oraz ich pierwsze
pochodne są ciągłe i lokalnie ograniczone. Reguła Leibniza
dla całek z parametrem, zastosowana na zwartych otoczeniach
zawierających odpowiednie odcinki całkowania, daje $v\in C^1(W')$.
Różniczkowanie wzoru względem $s$ daje
$\partial_s v+c(\Psi(s,y))v=f(\Psi(s,y))$, a podstawienie
$s=0$ daje $v(0,y)=g(\gamma(y))$.
Złożenie $u=v\circ(\Psi|_{W'})^{-1}$ jest klasy $C^1$.
Ponieważ $\partial_s\Psi=a\circ\Psi$, reguła łańcuchowa
pokazuje, że $u$ spełnia \eqref{eq:pde2-static} na $U$.
Równość $U\cap\Sigma=\Psi(\{0\}\times V')$ zapewnia
spełnienie danych na całym $U\cap\Sigma$.
Każde inne rozwiązanie na $U$ po złożeniu z $\Psi$
spełnia ten sam skalarny problem ODE na $(-\delta,\delta)$
dla każdego $y\in V'$, więc z jednoznaczności ODE jest równe $u$.
Dla dwóch rozwiązań zdefiniowanych na innych otoczeniach
wybieramy mniejszy iloczyn zawarty w przeciwobrazie
ich wspólnej dziedziny przez $\Psi$. Ten sam argument
dowodzi jednoznaczności zarodka.
\end{proof}

\begin{center}
\begin{tikzpicture}[>=Stealth,scale=.92]
 \draw[->] (-3.3,-1.35)--(3.8,-1.35) node[right] {$x_1$};
 \draw[->] (-3.1,-1.55)--(-3.1,2.5) node[above] {$x_2$};
 \draw[very thick] (-2.55,-.65)..controls (-.8,-.8) and (1.3,-.18)..(2.9,-.45);
 \node[below] at (2.6,-.6) {$\Sigma$};
 \foreach \x/\yy in {-1.75/-.64,-.65/-.52,.55/-.37,1.65/-.34} {
   \draw[blue!65!black,thick,->] (\x,\yy)--({\x+.48},{\yy+1.85});
   \fill (\x,\yy) circle (1.5pt);
 }
 \fill[red!70!black] (.55,-.37) circle (2.4pt);
 \node[below] at (.55,-.76) {$x_0$};
 \node[above] at (1.7,1.4) {charakterystyki};
\end{tikzpicture}
\par\smallskip\small
Dane są zadawane na $\Sigma$; przechodzące przez nią
poprzecznie charakterystyki tworzą lokalny układ współrzędnych.
\end{center}

\begin{kontrprzyklad}[Dane na powierzchni charakterystycznej]
\label{ex:pde2-char-data}
Rozważmy $u_x=0$ na $\R^2$, gdzie $u=u(x,y)$,
oraz dane $u(x,0)=g(x)$ na prostej $\Sigma=\{y=0\}$.
Każde rozwiązanie ma postać $u(x,y)=h(y)$, stąd
istnienie wymaga, by $g$ była stała. Gdy $g\equiv c$,
rozwiązań jest nieskończenie wiele: można wybrać
dowolne $h\in C^1(\R)$ z $h(0)=c$. Prosta $y=0$
jest charakterystyczna, bo pole $(1,0)$ leży
w jej przestrzeni stycznej.
\end{kontrprzyklad}

\section{Quasiliniowe równania pierwszego rzędu}
\label{sec:pde2-quasilinear}

Teraz prędkość charakterystyki może zależeć od niewiadomej:
\begin{equation}\label{eq:pde2-quasilinear}
 u_t+a(t,x,u)\cdot\nabla_xu=b(t,x,u),
 \qquad u(0,x)=u_0(x).
\end{equation}
W odróżnieniu od przypadku liniowego nie znamy przepływu
z góry: równocześnie musimy wyznaczać jego położenie
$X$ oraz wartość $U$ rozwiązania na charakterystyce.

\begin{twierdzenie}[Lokalna konstrukcja i jednoznaczność]
\label{thm:pde2-quasilinear}
Niech $a\colon(-T,T)\times\R^d\times\R\to\R^d$
i $b\colon(-T,T)\times\R^d\times\R\to\R$
będą klasy $C^2$, a $u_0\in C^2(\R^d)$.
Dla dowolnego $\xi_0\in\R^d$ istnieją otwarte otoczenie
$V\subset\R^d$ punktu $\xi_0$, liczba $0<\delta<T$
i otwarty zbiór $Q\subset(-\delta,\delta)\times\R^d$
zawierający $(0,\xi_0)$, na którym zagadnienie
\eqref{eq:pde2-quasilinear} ma rozwiązanie $u\in C^1(Q)$
z danymi $u(0,x)=u_0(x)$ dla $(0,x)\in Q$.
Rozwiązanie jest jednoznaczne jako zarodek klasy $C^1$
w otoczeniu $(0,\xi_0)$.

Dokładniej, niech $(X,U)$ rozwiązują układ ODE
\begin{equation}\label{eq:pde2-quasi-chars}
 \begin{aligned}
 \dot X(t;\xi)&=a(t,X(t;\xi),U(t;\xi)),& X(0;\xi)&=\xi,\\
 \dot U(t;\xi)&=b(t,X(t;\xi),U(t;\xi)),& U(0;\xi)&=u_0(\xi).
 \end{aligned}
\end{equation}
Układ ma rozwiązania dla wszystkich $(t,\xi)\in
(-\delta,\delta)\times V$, a po odpowiednim wyborze
$\delta,V$ odwzorowanie
$(t,\xi)\mapsto(t,X(t;\xi))$ jest dyfeomorfizmem
tego iloczynu na $Q$. Rozwiązanie spełnia
$u(t,X(t;\xi))=U(t;\xi)$.
\end{twierdzenie}
\begin{proof}
Prawa strona układu \eqref{eq:pde2-quasi-chars}
jest klasy $C^2$ w otwartym zbiorze
$(-T,T)\times\R^{d+1}$, a odwzorowanie danych
$\xi\mapsto(\xi,u_0(\xi))$ jest klasy $C^2$.
Lokalne twierdzenie o istnieniu ODE i zależności od danych
zapewnia wspólny przedział istnienia wokół zera dla
punktów początkowych z pewnego otoczenia
$(\xi_0,u_0(\xi_0))$. Z ciągłości odwzorowania danych
wynika więc istnienie rozwiązań $(X,U)$ klasy co najmniej
$C^1$ względem $(t,\xi)$ na iloczynie otwartego przedziału
i otoczenia $\xi_0$.
Dla $\Theta(t,\xi)=(t,X(t;\xi))$ mamy
\[
 D\Theta(0,\xi_0)=
 \begin{pmatrix}1&0\\
 a(0,\xi_0,u_0(\xi_0))&I_d
 \end{pmatrix},
 \qquad \det D\Theta(0,\xi_0)=1.
\]
Stosujemy twierdzenie o funkcji odwrotnej. W otoczeniu,
na którym $\Theta$ jest dyfeomorfizmem, wybieramy mniejszy
iloczyn $(-\delta,\delta)\times V$ zawierający $(0,\xi_0)$,
z $\delta<T$, i kładziemy $Q=\Theta((-\delta,\delta)\times V)$.
Zbiór $Q$ jest otwarty, a ponieważ $\Theta$ zachowuje
pierwszą współrzędną, leży w $(-\delta,\delta)\times\R^d$.

Niech $\xi=\Xi(t,x)$ oznacza jej lokalną odwrotność
współrzędnych przestrzennych. Kładziemy
$u(t,x)=U(t;\Xi(t,x))$. Różniczkując tożsamość
$u(t,X(t;\xi))=U(t;\xi)$ po $t$ przy ustalonym $\xi$,
otrzymujemy
\[
 u_t(t,X)+\nabla_xu(t,X)\cdot a(t,X,U)
 =b(t,X,U).
\]
Ponieważ $U=u(t,X)$, jest to \eqref{eq:pde2-quasilinear}.
W chwili $t=0$ mamy $X(0;\xi)=\xi$ i
$U(0;\xi)=u_0(\xi)$, więc warunek początkowy zachodzi.

Jeśli $v\in C^1$ jest innym lokalnym rozwiązaniem z tymi
danymi na otwartym otoczeniu $(0,\xi_0)$, pole
$(t,y)\mapsto a(t,y,v(t,y))$ jest tam klasy $C^1$.
Lokalne twierdzenie o istnieniu ODE daje wspólny przedział
$(-\varepsilon,\varepsilon)$ dla rozwiązań
$\dot Y(t;\xi)=a(t,Y(t;\xi),v(t,Y(t;\xi)))$,
$Y(0;\xi)=\xi$, dla $\xi$ z pewnego otoczenia $V_1\subset V$
punktu $\xi_0$. Wybieramy $\varepsilon\le\delta$ i $V_1$
tak, aby trajektorie pozostawały w dziedzinie $v$,
a $(0,\xi)$ należało do zbioru danych początkowych.
Funkcja $Z(t;\xi)=v(t,Y(t;\xi))$ spełnia z reguły
łańcuchowej i równania dla $v$
$\partial_t Z=b(t,Y,Z)$ oraz $Z(0;\xi)=u_0(\xi)$.
Para $(Y,Z)$ rozwiązuje więc \eqref{eq:pde2-quasi-chars}
z tymi samymi danymi co $(X,U)$. Jednoznaczność ODE daje
$(Y,Z)=(X,U)$ na całym wspólnym iloczynie.
W konsekwencji $v(t,X(t;\xi))=U(t;\xi)=u(t,X(t;\xi))$.
Obraz $(-\varepsilon,\varepsilon)\times V_1$ przez
dyfeomorfizm $\Theta$ jest otwartym otoczeniem $(0,\xi_0)$,
na którym $v=u$.
\end{proof}

\begin{uwaga}[Co może uniemożliwić przedłużenie?]
Lokalna konstrukcja wymaga zarówno istnienia rozwiązań
układu ODE, jak i odwracalności odwzorowania $\xi\mapsto X(t;\xi)$.
W zagadnieniach quasiliniowych nawet przy globalnym
istnieniu charakterystyk może wystąpić
$\det D_\xi X=0$: wtedy klasyczne współrzędne
charakterystyczne zawodzą. W modelach skalarnych jest
to typowy mechanizm utraty regularności rozwiązania.
\end{uwaga}

\section{Fale proste i kryterium wybuchu pochodnej}
\label{sec:pde2-blowup}

Rozważmy w jednym wymiarze równanie autonomiczne
\begin{equation}\label{eq:pde2-simplewave}
 u_t+a(u)u_x=0,\qquad u(0,x)=u_0(x),
\end{equation}
gdzie $a\in C^2(\R)$, $u_0\in C^2(\R)$.
Wzdłuż charakterystyki wartość $u$ jest stała, więc
\begin{equation}\label{eq:pde2-simplewave-map}
 X(t;\xi)=\xi+t\,a(u_0(\xi)),\qquad
 U(t;\xi)=u_0(\xi).
\end{equation}
Różniczkując względem $\xi$, dostajemy
$X_\xi(t;\xi)=1+t\,a'(u_0(\xi))u_0'(\xi)$.

\begin{twierdzenie}[Pierwsza osobliwość gradientu]
\label{thm:pde2-breaking}
Niech $a,u_0\in C^2(\R)$ i załóżmy, że
$a\circ u_0$ jest ograniczona, a funkcja
$q(\xi)=a'(u_0(\xi))u_0'(\xi)$ jest ograniczona z dołu.
Jeżeli $m=\inf_{\xi\in\R}q(\xi)<0$, określmy
$T_*=-1/m$. Wtedy dla każdego $0\leq t<T_*$
odwzorowanie $X_t(\xi)=\xi+t a(u_0(\xi))$
jest ściśle rosnącym dyfeomorfizmem $\R\to\R$ i
\eqref{eq:pde2-simplewave-map} definiuje rozwiązanie
klasy $C^1$ na $[0,T_*)\times\R$. Jego pochodna wynosi
\begin{equation}\label{eq:pde2-gradient-formula}
 u_x(t,X(t;\xi))=
 \frac{u_0'(\xi)}{1+t a'(u_0(\xi))u_0'(\xi)}.
\end{equation}
Jeśli ponadto minimum $m$ jest osiągane w punkcie
$\xi_*$, to $|u_x(t,X(t;\xi_*))|\to\infty$
gdy $t\uparrow T_*$. Nie istnieje zatem przedłużenie
rozwiązania jako funkcji $C^1$ na otoczenie
$(T_*,X(T_*;\xi_*))$.
\end{twierdzenie}
\begin{proof}
Dla $0\leq t<T_*$ mamy
$X_\xi(t;\xi)=1+tq(\xi)\geq1+tm>0$.
Odwzorowanie $X_t$ jest więc ściśle rosnące.
Z ograniczoności $a\circ u_0$ wynika
$X_t(\xi)-\xi=O(1)$ dla $\xi\to\pm\infty$,
więc $X_t(\xi)\to\pm\infty$ przy $\xi\to\pm\infty$.
Z ciągłości wynika surjektywność, a ze ścisłej
monotoniczności injektywność $X_t$.
Odwzorowanie $(t,\xi)\mapsto(t,X(t;\xi))$ ma wyznacznik
macierzy pochodnej równy $1+tq(\xi)>0$.
Twierdzenie o funkcji odwrotnej daje lokalne odwrotności
klasy $C^1$; ich zgodność wynika z bijektywności $X_t$.
Przy $t=0$ stosujemy to twierdzenie do tego samego wzoru
na $X(t;\xi)$, określonego także dla ujemnych $t$.
W ten sposób $\Xi(t,x)=X_t^{-1}(x)$ jest klasy $C^1$
względem obu zmiennych aż do $t=0$.

Definiujemy $u(t,x)=u_0(\Xi(t,x))$. Jest to funkcja klasy
$C^1$, a ponieważ $X_0$ jest identycznością, mamy
$u(0,x)=u_0(x)$. Z równości $u(t,X(t;\xi))=u_0(\xi)$
przez różniczkowanie względem $t$ otrzymujemy
\[
 u_t(t,X(t;\xi))+u_x(t,X(t;\xi))a(u_0(\xi))=0.
\]
Ponieważ $u_0(\xi)=u(t,X(t;\xi))$ i $X_t$ jest surjekcją,
jest to równanie \eqref{eq:pde2-simplewave} w każdym punkcie
dziedziny. Różniczkowanie tej samej tożsamości względem
$\xi$ daje $u_x(t,X(t;\xi))X_\xi(t;\xi)=u_0'(\xi)$.
Dzieląc przez dodatnią pochodną $X_\xi$, uzyskujemy
\eqref{eq:pde2-gradient-formula}.

Jeżeli $q(\xi_*)=m<0$, to
$a'(u_0(\xi_*))u_0'(\xi_*)\ne0$; w szczególności
$u_0'(\xi_*)\ne0$. Mianownik
\eqref{eq:pde2-gradient-formula} przy $\xi_*$
dąży do zera od strony dodatniej, a licznik jest niezerowy.
Ponadto $X(t;\xi_*)\to X(T_*;\xi_*)$, ponieważ wzór
na charakterystykę jest ciągły także w $t=T_*$.
Zatem moduł pochodnej dąży do nieskończoności wzdłuż
krzywej dochodzącej do wskazanego punktu. Ciągłość
pochodnej przestrzennej ewentualnego przedłużenia
w tym punkcie przeczyłaby tej granicy.
\end{proof}

\begin{przyklad}[Burgers: dokładny czas powstania osobliwości]
\label{ex:pde2-burgers}
Dla równania $u_t+u u_x=0$ mamy $a(u)=u$.
Wybierzmy $u_0(x)=-\tanh x$. Wówczas
$q(x)=u_0'(x)=-1/\cosh^2x$ osiąga minimum $-1$
w $x=0$, zaś $a\circ u_0$ jest ograniczona.
Twierdzenie \ref{thm:pde2-breaking} daje $T_*=1$,
$X(t;\xi)=\xi-t\tanh\xi$ oraz
$u_x(t,0)=-1/(1-t)$ dla $t<1$.
Charakterystyki istnieją jako krzywe gładkie także
później, ale nie mogą już wyznaczać globalnego
klasycznego rozwiązania w tym samym schemacie.
\end{przyklad}

\begin{center}
\begin{tikzpicture}[>=Stealth,scale=.86]
 \draw[->] (-3.4,0)--(3.6,0) node[right] {$x$};
 \draw[->] (-3.1,-.12)--(-3.1,3.75) node[above] {$t$};
 \foreach \xi in {-2.4,-1.6,-.8,-.4,.4,.8,1.6,2.4} {
  \pgfmathsetmacro{\dd}{-(exp(2*\xi)-1)/(exp(2*\xi)+1)}
  \draw[blue!65!black,thick] (\xi,0)--({\xi+0.95*\dd},2.95);
 }
 \draw[densely dashed] (-3.1,3.1)--(3.35,3.1);
 \node[anchor=west] at (3.55,3.08) {$t\approx T_*$};
 \node[below] at (0,0) {$0$};
 \node[align=left,anchor=west] at (3.6,1.45) {zbliżanie się\\charakterystyk};
\end{tikzpicture}
\par\smallskip\small
Dla malejących danych Burgersa charakterystyki
zbliżają się do siebie. Rysunek ma charakter schematyczny;
czas i nachylenia przeskalowano niezależnie.
\end{center}

\begin{uwaga}[Eksperyment obliczeniowy]
W rozdziale~\ref{chap:pde-code} zamieszczono kod Python,
który dla przykładu Burgersa z danymi $u_0=-\tanh x$
rysuje charakterystyki i profile rozwiązań dla
$t\in\{0,0.5,0.9,0.99\}$. Program korzysta z odwracania
monotonicznego odwzorowania $\xi\mapsto X(t;\xi)$ przez
interpolację. Porównanie z dokładnym wzorem
\eqref{eq:pde2-gradient-formula} pozwala sprawdzić,
że $u_x(t,0)=-1/(1-t)$, podczas gdy $|u(t,x)|\leq1$.
Nie jest to algorytm kontynuacji rozwiązania przez szok:
dla $t\geq1$ odwzorowanie etykiet nie jest już globalnie
odwracalna w klasie $C^1$.
\end{uwaga}

\begin{uwaga}[Znaczenie czasu $T_*$]
Twierdzenie \ref{thm:pde2-breaking} dotyczy dokładnie
zadanego równania $u_t+a(u)u_x=0$, bez źródła,
i wskazanej klasy danych. Nie jest ogólnym twierdzeniem
o czasie istnienia dowolnych równań quasiliniowych.
Gdy $q(\xi)\geq0$ dla wszystkich $\xi$, charakterystyki
nie przecinają się dla $t\geq0$, o ile
$a\circ u_0$ pozostaje ograniczona; w takim przypadku
wzór daje globalne rozwiązanie klasy $C^1$.
\end{uwaga}

\section{Od przecinających się charakterystyk do fal uderzeniowych}
\label{sec:pde2-shocks}

Równanie $u_t+a(u)u_x=0$ można, gdy $a=f'$,
zapisać w postaci prawa zachowania
\begin{equation}\label{eq:pde2-conservation}
 u_t+\partial_x f(u)=0.
\end{equation}
Po pojawieniu się nieciągłości wzór z pochodnymi
klasycznymi przestaje mieć sens. Potrzebne jest
sformułowanie całkowe, które później zinterpretujemy
w języku dystrybucji.

\begin{definicja}[Rozwiązanie słabe prawa zachowania]
\label{def:pde2-weak}
Niech $f\in C^1(\R)$, $u_0\in L^1_{\mathrm{loc}}(\R)$,
a $u\in L^1_{\mathrm{loc}}([0,T)\times\R)$, przy czym
$f(u)$ jest lokalnie całkowalna. Funkcję $u$ nazywamy
\emph{rozwiązaniem słabym} zagadnienia
\eqref{eq:pde2-conservation} z danymi $u_0$, jeżeli
dla każdej $\varphi\in C_c^\infty([0,T)\times\R)$
zachodzi
\begin{equation}\label{eq:pde2-weak-form}
 \int_0^T\!\int_\R
 \bigl(u(t,x)\varphi_t(t,x)+f(u(t,x))\varphi_x(t,x)\bigr)
 \,\mathrm d\lambda_1(x)\,\mathrm dt
 +\int_\R u_0(x)\varphi(0,x)\,\mathrm d\lambda_1(x)=0.
\end{equation}
Przestrzeń $C_c^\infty([0,T)\times\R)$ oznacza
restrykcje funkcji gładkich w pewnym otoczeniu
$[0,T)\times\R$, których nośniki są zwarte i nie dochodzą
do czasu $T$.
\end{definicja}

\begin{stwierdzenie}[Warunek Rankine'a--Hugoniota]
\label{prop:pde2-rh}
Niech $f\in C^1(\R)$ i $s\in C^1((0,T))$.
Załóżmy, że $u$ jest klasy $C^1$ po obu stronach
krzywej $x=s(t)$, z jednostronnymi przedłużeniami
klasy $C^1$ do krzywej, i posiada ślady
$u_L(t)$ i $u_R(t)$ na tej krzywej oraz spełnia
\eqref{eq:pde2-conservation} klasycznie poza nią.
Wówczas $u$ spełnia równanie w sensie dystrybucji
na $(0,T)\times\R$ (bez warunku początkowego), tj.
$\int_0^T\int_\R (u\varphi_t+f(u)\varphi_x)\,
\mathrm d\lambda_1(x)\,\mathrm d\lambda_1(t)=0$
dla każdej $\varphi\in C_c^\infty((0,T)\times\R)$,
wtedy i tylko wtedy, gdy
\begin{equation}\label{eq:pde2-rh}
 s'(t)\bigl(u_R(t)-u_L(t)\bigr)
 =f(u_R(t))-f(u_L(t))
\end{equation}
dla każdego $t\in(0,T)$, przy ciągłości śladów.
\end{stwierdzenie}
\begin{proof}
Niech $\varphi\in C_c^\infty((0,T)\times\R)$.
Wybierzmy $0<a<b<T$ oraz $R>0$, tak aby nośnik
$\varphi$ leżał w $(a,b)\times(-R,R)$ oraz
$|s(t)|<R$ dla $t\in[a,b]$. Oznaczmy
$D_L=\{(t,x):a<t<b,\ -R<x<s(t)\}$ oraz
$D_R=\{(t,x):a<t<b,\ s(t)<x<R\}$.
Na każdym z tych zbiorów, z równania klasycznego
i reguły iloczynu, mamy
\[
 \partial_t(u\varphi)+\partial_x(f(u)\varphi)
 =u\varphi_t+f(u)\varphi_x.
\]
Całkujemy po obu zbiorach, stosując twierdzenie Gaussa
do pola $(u\varphi,f(u)\varphi)$ w zmiennych $(t,x)$.
Brzegi są kawałkami klasy $C^1$, pole ma jednostronne
przedłużenia klasy $C^1$, a $\varphi$ znika na zewnętrznych
krawędziach prostokąta. Na krzywej $(t,s(t))$
wektor normalny zewnętrzny dla $D_L$ ma kierunek
$(-s'(t),1)$, a dla $D_R$ kierunek przeciwny.
Element długości krzywej wynosi
$\sqrt{1+s'(t)^2}\,\mathrm d\lambda_1(t)$,
co skraca czynnik normalizujący normalnej.
Wkład od lewej strony wynosi więc
$\int_a^b(f(u_L)-s'u_L)\varphi(t,s(t))\,\mathrm d\lambda_1(t)$,
a od prawej strony
$\int_a^b(s'u_R-f(u_R))\varphi(t,s(t))\,\mathrm d\lambda_1(t)$.
Ich suma wynosi
\[
 \int_0^T\bigl(s'(t)[u](t)-[f(u)](t)\bigr)
 \varphi(t,s(t))\,\mathrm dt,
\]
gdzie $[u]=u_R-u_L$ oraz $[f(u)]=f(u_R)-f(u_L)$.
Ślad funkcji testowej na krzywej można zadawać dowolnie
bez zakładania gładkości samej krzywej ponad klasę $C^1$.
Jeśli zachodzi \eqref{eq:pde2-rh}, suma jest zerowa
dla każdej funkcji testowej, co dowodzi wystarczalności.
Dla dowodu konieczności załóżmy, że równanie zachodzi
w sensie dystrybucji. Niech $\eta\in C_c^\infty((0,T))$
i niech $J\subset(0,T)$ będzie zwartym przedziałem
zawierającym jej nośnik. Zwartość $s(J)$
pozwala wybrać $\chi\in C_c^\infty(\R)$, równą $1$
w otoczeniu $s(J)$. Wtedy
$\varphi(t,x)=\eta(t)\chi(x)$ jest klasy $C_c^\infty$
i spełnia $\varphi(t,s(t))=\eta(t)$ dla $t\in J$.
Zatem dla ciągłej funkcji
$b(t)=s'(t)[u](t)-[f(u)](t)$ zachodzi
$\int_0^T b(t)\eta(t)\,\mathrm dt=0$ dla każdego takiego $\eta$.
Gdyby $b(t_0)\ne0$, ciągłość dawałaby przedział, na którym
$b$ ma stały, niezerowy znak; test nieujemny o nośniku w tym
przedziale prowadziłby do sprzeczności. Stąd $b=0$
w każdym punkcie $(0,T)$.
Warunek rozstrzyga wyłącznie poprawność konserwatywną;
do jednoznacznego wyboru skoku potrzebna jest jeszcze
nierówność entropijna.
\end{proof}

\begin{przyklad}[Skok Burgersa]
Dla $f(u)=u^2/2$ i stałych stanów po obu stronach
fali $u_L\ne u_R$ warunek \eqref{eq:pde2-rh} daje
$\sigma=s'=(u_L+u_R)/2$.
Przy $u_L>u_R$ zachodzi $u_L>\sigma>u_R$:
charakterystyki z obu stron wpadają w falę uderzeniową.
Jeżeli $u_L<u_R$, taka nieciągłość spełnia wprawdzie
warunek Rankine'a--Hugoniota, ale nie jest dopuszczalnym
rozwiązaniem entropijnym.
\end{przyklad}

\begin{przyklad}[Fala rozrzedzenia]
\label{ex:pde2-rarefaction}
Dla równania Burgersa z danymi Riemanna
$u_0(x)=u_L$ dla $x<0$ oraz $u_0(x)=u_R$ dla $x>0$,
przy $u_L<u_R$, funkcja
\[
 u(t,x)=
 \begin{cases}
 u_L,&x\leq u_Lt,\\
 x/t,&u_Lt<x<u_Rt,\\
 u_R,&x\geq u_Rt
 \end{cases}
 \qquad(t>0)
\]
jest ciągła na każdym przekroju $t>0$ i spełnia
równanie Burgersa klasycznie w każdym z trzech
otwartych obszarów. Bezpośrednio sprawdzamy, że
w środkowej części $u_t=-x/t^2$, $u_x=1/t$,
zatem $u_t+uu_x=0$. Ciągłość na liniach
$x=u_Lt$ i $x=u_Rt$ sprawia, że oba skoki, wartości $u$
i strumienia $f(u)=u^2/2$, są zerowe.
Dla $\varphi\in C_c^\infty([0,T)\times\R)$ oraz
$0<\varepsilon<T$ stosujemy całkowanie przez części
na trzech częściach zbioru $(\varepsilon,T)\times\R$.
Wkłady na obu liniach wewnętrznych znoszą się,
a przy $t=T$ i na przestrzennych krańcach nośnika
test znika. Na dolnej krawędzi pozostaje
\[
 \int_\varepsilon^T\!\int_\R
 (u\varphi_t+f(u)\varphi_x)\,
 \mathrm d\lambda_1(x)\,\mathrm d\lambda_1(t)
 +\int_\R u(\varepsilon,x)\varphi(\varepsilon,x)\,
 \mathrm d\lambda_1(x)=0.
\]
Niech $M=\max\{|u_L|,|u_R|\}$. Zarówno $u(t,\cdot)$,
jak i $u_0$ są ograniczone w module przez $M$.
Poza przedziałem
$[\min(0,u_Lt),\max(0,u_Rt)]$ funkcje te są równe.
Stąd dla $t>0$
\[
 \int_\R |u(t,x)-u_0(x)|\,\mathrm d\lambda_1(x)
 \le 2Mt\bigl(\max(0,u_R)-\min(0,u_L)\bigr).
\]
W szczególności zachodzi zbieżność do danych początkowych
w $L^1_{\mathrm{loc}}$. W całce brzegowej możemy więc
zastąpić $u(\varepsilon,\cdot)$ przez $u_0$ z błędem
dążącym do zera, a $\varphi(\varepsilon,\cdot)$ dąży
jednostajnie do $\varphi(0,\cdot)$ na wspólnym zwartym
nośniku przestrzennym. W całce podwójnej korzystamy
z twierdzenia o zbieżności zdominowanej, ponieważ
$u$ oraz $f(u)$ są ograniczone, a test i jego pochodne
mają zwarty nośnik. Przejście z $\varepsilon$ do zera
daje dokładnie \eqref{eq:pde2-weak-form}.
To rozwiązanie jest entropijne; precyzyjną definicję
warunków entropijnych podamy w rozdziale o prawach zachowania.
\end{przyklad}

\begin{uwaga}[Dlaczego samo rozwiązanie słabe nie wystarcza?]
Dla danych z przykładu \ref{ex:pde2-rarefaction}
można też skonstruować skok poruszający się z prędkością
$(u_L+u_R)/2$. Spełnia on ten sam całkowy warunek
\eqref{eq:pde2-weak-form}, lecz narusza warunek entropijny.
Dla tego samego zagadnienia mamy więc co najmniej
dwa różne rozwiązania słabe. Selekcja rozwiązania
fizycznego wymaga dodatkowej informacji.
\end{uwaga}

\section{Równania w pełni nieliniowe i układ Charpita}
\label{sec:pde2-charpit}

Dla ogólnego równania pierwszego rzędu
\begin{equation}\label{eq:pde2-general-F}
 F(x,u(x),\nabla u(x))=0,
\end{equation}
gdzie $x\in\Omega\subset\R^n$,
$F\colon\Omega\times\R\times\R^n\to\R$,
rozpatrujemy dodatkowe zmienne $p\in\R^n$
reprezentujące gradient. Pojawia się
\emph{układ charakterystyczny Charpita}:
\begin{equation}\label{eq:pde2-charpit-system}
 \dot X=F_p(X,U,P),\qquad
 \dot U=P\cdot F_p(X,U,P),\qquad
 \dot P=-F_x(X,U,P)-F_u(X,U,P)P.
\end{equation}
Tutaj $F_x,F_p$ są gradientami po odpowiednich
zmiennych, a $F_u$ jest pochodną skalarną.

\begin{lemat}[Zachowanie równania i warunku styczności]
\label{lem:pde2-contact}
Niech $F\in C^2$ i niech $(X,U,P)$ będzie rozwiązaniem
\eqref{eq:pde2-charpit-system} na przedziale $I$.
Wtedy $s\mapsto F(X(s),U(s),P(s))$ jest stała.
Jeśli rozwiązania zależą klasy $C^1$ od parametru
$y\in V\subset\R^{n-1}$ i $F(X,U,P)=0$, to
\begin{equation}\label{eq:pde2-contact-ode}
 \frac{\mathrm d}{\mathrm ds}
 \bigl(U_{y_j}-P\cdot X_{y_j}\bigr)
 =-F_u(X,U,P)\bigl(U_{y_j}-P\cdot X_{y_j}\bigr)
\end{equation}
dla każdego $j=1,\ldots,n-1$.
\end{lemat}
\begin{proof}
Reguła łańcuchowa daje
\[
 \frac{\mathrm d}{\mathrm ds}F(X,U,P)
 =F_x\cdot F_p+F_u(P\cdot F_p)
 +F_p\cdot(-F_x-F_uP)=0.
\]
Różniczkując $R_j=U_{y_j}-P\cdot X_{y_j}$ po $s$,
wykorzystując \eqref{eq:pde2-charpit-system}, otrzymujemy
\[
 \dot R_j=P_{y_j}\cdot F_p+(F_x+F_uP)\cdot X_{y_j}.
\]
Ponieważ $F(X,U,P)=0$ dla każdego $(s,y)$,
mamy także
$F_x\cdot X_{y_j}+F_u U_{y_j}+F_p\cdot P_{y_j}=0$.
Po podstawieniu do wzoru na $\dot R_j$ dostajemy
$\dot R_j=-F_u(U_{y_j}-P\cdot X_{y_j})$, czyli
\eqref{eq:pde2-contact-ode}.
\end{proof}

\begin{twierdzenie}[Lokalna konstrukcja Charpita]
\label{thm:pde2-charpit}
Niech $n\geq2$, $\Omega\subset\R^n$ będzie otwarta,
$F\in C^2(\Omega\times\R\times\R^n)$,
$V\subset\R^{n-1}$ otwarty oraz
$\gamma\in C^2(V;\Omega)$ regularną parametryzacją
hiperpowierzchni. Niech $g\in C^2(V)$ i
$p_0\in C^1(V;\R^n)$ spełniają dla wszystkich $y\in V$
\begin{equation}\label{eq:pde2-charpit-initial}
 F(\gamma(y),g(y),p_0(y))=0,\qquad
 p_0(y)\cdot\partial_{y_j}\gamma(y)=\partial_{y_j}g(y)
 \quad(j=1,\ldots,n-1).
\end{equation}
Jeżeli w $y_0\in V$ wektor
$F_p(\gamma(y_0),g(y_0),p_0(y_0))$
nie należy do $T_{\gamma(y_0)}\Sigma$,
gdzie $\Sigma=\gamma(V)$ lokalnie, to istnieje
otoczenie $U$ punktu $\gamma(y_0)$ i rozwiązanie
$u\in C^2(U)$ równania \eqref{eq:pde2-general-F}
przyjmujące dane $u(\gamma(y))=g(y)$ w pobliżu $y_0$.
Jego gradient na powierzchni danych jest równy $p_0$.
W klasie rozwiązań $C^2$, które realizują te same
dane $g,p_0$, konstrukcja jest lokalnie jednoznaczna.
\end{twierdzenie}
\begin{proof}
Rozwiążmy układ \eqref{eq:pde2-charpit-system}
z danymi $(X,U,P)(0,y)=(\gamma(y),g(y),p_0(y))$.
Zależność rozwiązań ODE od parametrów daje funkcje
klasy $C^1$ względem $(s,y)$ w pobliżu $(0,y_0)$.
Lemat \ref{lem:pde2-contact} i pierwszy warunek
\eqref{eq:pde2-charpit-initial} pokazują, że
$F(X,U,P)=0$; drugi warunek i równanie
\eqref{eq:pde2-contact-ode} dają
$U_{y_j}=P\cdot X_{y_j}$ przez jednoznaczność
skalarnego równania liniowego. Ponadto z samego
układu wynika $U_s=P\cdot X_s$.

W punkcie $(0,y_0)$ kolumny macierzy $D_{(s,y)}X$
są równe $F_p$ i $\partial_{y_j}\gamma$;
warunek poprzeczności zapewnia, że macierz ta jest
odwracalna. Z twierdzenia o funkcji odwrotnej $X$
stanowi lokalny układ współrzędnych w $\Omega$.
Definiujemy $u(X(s,y))=U(s,y)$. Wspomniane wcześniej
równości są równoważne tożsamości
$\mathrm dU=P\cdot\mathrm dX$, zatem
$\nabla u(X(s,y))=P(s,y)$. Ponieważ $P$ i lokalna
odwrotność $X^{-1}$ są klasy $C^1$, gradient $\nabla u$
jest klasy $C^1$ na $U$, a więc $u\in C^2(U)$.
Ponieważ $F(X,U,P)=0$, otrzymujemy równanie
\eqref{eq:pde2-general-F}.

Jeśli $u\in C^2$ jest innym rozwiązaniem z tym samym
gradientem początkowym, to wzdłuż odpowiednich
charakterystyk jego wartości i pochodne spełniają
\eqref{eq:pde2-charpit-system}: wystarczy zróżniczkować
$F(x,u,\nabla u)=0$ i skorzystać z symetrii hesjanu.
Jednoznaczność ODE implikuje zgodność obu rozwiązań
na wspólnym otoczeniu punktu początkowego.
\end{proof}

\begin{przyklad}[Eikonal i wybór gałęzi]
Niech $F(x,u,p)=|p|^2-1$ w $\R^2$.
Na prostej $\Sigma=\{(x,0):x\in\R\}$ zadamy
$u(x,0)=0$. Zgodność gradientu stycznego wymusza
$p_x=0$, a równanie wymusza $p_y=\pm1$.
Są zatem dwie dopuszczalne gałęzie danych gradientu,
które prowadzą do rozwiązań $u(x,y)=y$ i $u(x,y)=-y$.
Oba są niecharakterystyczne dla odpowiedniego wektora
$F_p=2p$. Samo zadanie wartości na prostej nie określa
więc jednoznacznie rozwiązania w pełni nieliniowego.
\end{przyklad}

\section{Równanie Hamiltona--Jacobiego i charakterystyki Hamiltona}
\label{sec:pde2-hj}

Ważną szczególną postacią równania pierwszego rzędu jest
\begin{equation}\label{eq:pde2-hj}
 u_t+H(x,\nabla_xu)=0,
 \qquad u(0,x)=u_0(x).
\end{equation}
W tej sekcji $H\in C^2(\R^d\times\R^d)$ nie zależy
jawnie od czasu ani od wartości $u$.

\begin{stwierdzenie}[Metoda charakterystyk Hamiltona]
\label{prop:pde2-hamilton}
Niech $u_0\in C^2(\R^d)$, $H\in C^2(\R^d\times\R^d)$
i niech $\xi_0\in\R^d$. Rozwiążmy lokalnie układ
\begin{equation}\label{eq:pde2-hamilton-odes}
 \dot X=H_p(X,P),\qquad
 \dot P=-H_x(X,P),\qquad
 \dot U=P\cdot H_p(X,P)-H(X,P),
\end{equation}
z warunkami $X(0;\xi)=\xi$,
$P(0;\xi)=\nabla u_0(\xi)$ i $U(0;\xi)=u_0(\xi)$.
Dla dostatecznie małych $|t|$ i $\xi$ bliskich $\xi_0$
mapa $\xi\mapsto X(t;\xi)$ jest lokalnie odwracalna.
W tych współrzędnych
$u(t,X(t;\xi))=U(t;\xi)$ definiuje rozwiązanie
klasy $C^1$ równania \eqref{eq:pde2-hj}, a ponadto
$\nabla_xu(t,X(t;\xi))=P(t;\xi)$.
\end{stwierdzenie}
\begin{proof}
Istnienie ODE i lokalna odwracalność wynikają jak
w twierdzeniu \ref{thm:pde2-quasilinear}, ponieważ
$D_\xi X(0;\xi)=I$. Ustalmy współrzędną $\xi_j$
i oznaczmy $R_j=U_{\xi_j}-P\cdot X_{\xi_j}$.
Korzystając z \eqref{eq:pde2-hamilton-odes} oraz
reguły łańcuchowej, obliczamy
\begin{align*}
 \dot R_j&=\partial_{\xi_j}(P\cdot H_p-H)
 -\dot P\cdot X_{\xi_j}-P\cdot\partial_{\xi_j}H_p\\
 &=P_{\xi_j}\cdot H_p-H_x\cdot X_{\xi_j}
 -H_p\cdot P_{\xi_j}+H_x\cdot X_{\xi_j}=0.
\end{align*}
Początkowo
$R_j(0,\xi)=\partial_{\xi_j}u_0-\nabla u_0\cdot e_j=0$,
zatem $R_j=0$. Wynika stąd
$\nabla_xu(t,X(t;\xi))=P(t;\xi)$. Z kolei pochodna
$u(t,X(t;\xi))=U(t;\xi)$ po $t$ daje
$u_t+\nabla_xu\cdot H_p=P\cdot H_p-H$, a więc
$u_t+H=0$, jak należało wykazać.
\end{proof}

\begin{przyklad}[Hamiltonian kinetyczny i powstawanie kaustyki]
\label{ex:pde2-hj-blowup}
Niech $d=1$ i $H(p)=p^2/2$. Wtedy
$\dot P=0$, $\dot X=P$ oraz $\dot U=P^2/2$.
W konsekwencji
\[
 X(t;\xi)=\xi+t u_0'(\xi),\qquad
 P(t;\xi)=u_0'(\xi),\qquad
 U(t;\xi)=u_0(\xi)+\frac t2|u_0'(\xi)|^2.
\]
Dla $u_0(x)=-x^2/2$ otrzymujemy
$X(t;\xi)=(1-t)\xi$. Jeśli $0\leq t<1$, to
\[
 u(t,x)=-\frac{x^2}{2(1-t)},\qquad
 u_x(t,x)=-\frac{x}{1-t}.
\]
W chwili $t=1$ mapa $\xi\mapsto X(t;\xi)$ traci
odwracalność, a dla ustalonego $x\ne0$ gradient
rozwiązania staje się nieograniczony, gdy $t\uparrow1$.
Zjawisko utraty odwracalności rzutu
$(X,P)\mapsto X$ nazywa się w tym kontekście
powstawaniem \emph{kaustyki}. Rozwiązania lepkościowe
umożliwią później kontynuację odpowiednich zadań
Hamiltona--Jacobiego poza pierwszą kaustyką.
\end{przyklad}

\section{Podsumowanie: kiedy metoda charakterystyk działa?}
\label{sec:pde2-summary}

Metoda charakterystyk wykorzystuje zależność
\eqref{eq:pde2-chain}, ale jej wynik zależy od struktury
równania:
\begin{enumerate}
\item Dla transportu liniowego z wystarczająco regularnym
polem prędkości charakterystyki tworzą przepływ
dyfeomorfizmów i zapewniają globalne rozwiązanie klasyczne
na skończonym przedziale czasu.
\item Dla równań quasiliniowych lokalne rozwiązanie istnieje,
dopóki układ charakterystyczny daje odwracalną mapę
początkowych etykiet $\xi$ na położenia $x$.
\item Dla praw zachowania po utracie gładkości naturalne
stają się rozwiązania słabe. Sama słaba postać równania
nie wystarcza zwykle do jednoznaczności; potrzebne są
warunki entropijne.
\item Dla równań w pełni nieliniowych trzeba śledzić
również gradient $p=\nabla u$. Warunek poprzeczności
$F_p$ względem powierzchni danych jest kluczowy,
a rozwiązanie może istnieć jedynie lokalnie nawet
przy bardzo gładkich danych.
\end{enumerate}

W rozdziale następnym przejdziemy do równania Laplace'a,
teorii funkcji harmonicznych i zagadnień Dirichleta.


\chapter[Równania Laplace'a i Poissona]{Równania Laplace'a i Poissona: funkcje harmoniczne i potencjał}
\label{chap:pde-laplace}

W tym rozdziale badamy dwa równania na zbiorze otwartym $\Omega\subset\R^n$:
$\Delta u=0$ oraz $-\Delta u=f$. Pierwsze opisuje między innymi
stacjonarną temperaturę i potencjały bez źródeł; drugie uwzględnia
źródła. W przeciwieństwie do równania transportu nie wystarcza
znajomość danych na dowolnej hiperpowierzchni. Naturalnym zadaniem
jest zwykle problem Dirichleta na \emph{całym} brzegu.

W całym rozdziale, o ile nie zaznaczono inaczej,
$n\geq2$ jest liczbą całkowitą, a funkcje skalarne są
rzeczywiste. Przez $\lambda_k$ oznaczamy miarę Lebesgue'a
w $\R^k$ dla każdej liczby całkowitej $k\geq1$, przez
$\sigma$ miarę powierzchniową na rozpatrywanej gładkiej
hiperpowierzchni, a przez $\omega_n=\sigma(S^{n-1})$
pole jednostkowej sfery. Ze wzoru na całkowanie
we współrzędnych biegunowych wynika
$\lambda_n(B_r(0))=\omega_n\int_0^r s^{n-1}\,
\mathrm d\lambda_1(s)=\omega_n r^n/n$.
Znak laplasjanu ustalamy jako $\Delta=\sum_{j=1}^n\partial_{x_j}^2$;
rozwiązanie fundamentalne będziemy normalizować przez $-\Delta\Phi=\delta_0$.

\section{Własność wartości średniej}
\label{sec:pde3-mean}

\begin{definicja}[Funkcja harmoniczna]
Niech $\Omega\subset\R^n$ będzie niepustym zbiorem otwartym.
Funkcję rzeczywistą $u\in C^2(\Omega)$ nazywamy \emph{harmoniczną}
na $\Omega$, jeśli $\Delta u=0$ w każdym punkcie $\Omega$.
Funkcję $u\in C^2(\Omega)$ nazywamy \emph{subharmoniczną},
gdy $\Delta u\geq0$, oraz \emph{superharmoniczną},
gdy $\Delta u\leq0$. Później rozszerzymy te pojęcia
na funkcje bez klasycznych drugich pochodnych.
\end{definicja}

Dla $u\in C(\Omega)$, $x_0\in\Omega$ oraz
$0<r<\operatorname{dist}(x_0,\partial\Omega)$ zapisujemy
średnią sferyczną
\begin{equation}\label{eq:pde3-sphere-mean}
 M_u(x_0,r)=\frac{1}{\omega_n r^{n-1}}
 \int_{\partial B_r(x_0)}u(y)\,\mathrm d\sigma(y)
 =\frac1{\omega_n}\int_{S^{n-1}}u(x_0+r\theta)\,
 \mathrm d\sigma(\theta).
\end{equation}
Druga równość wynika ze zmiany zmiennych
$y=x_0+r\theta$: miara powierzchniowa skaluje się
przez czynnik $r^{n-1}$, gdyż pochodna tego
odwzorowania na każdej przestrzeni stycznej
jest mnożeniem przez $r$.
Jeśli $\Omega=\R^n$, przyjmujemy
$\operatorname{dist}(x_0,\partial\Omega)=+\infty$.

\begin{twierdzenie}[Własność wartości średniej]
\label{thm:pde3-mean}
Niech $n\geq2$, $\Omega\subset\R^n$ będzie otwarta,
$u\in C^2(\Omega)$ i $\Delta u=0$ w $\Omega$.
Dla każdego $x_0\in\Omega$ oraz $r>0$ takich, że
$\overline{B_r(x_0)}\subset\Omega$, zachodzi
\begin{equation}\label{eq:pde3-both-means}
 u(x_0)=\frac{1}{\omega_n r^{n-1}}
 \int_{\partial B_r(x_0)}u(y)\,\mathrm d\sigma(y)
 =\frac{1}{\lambda_n(B_r(x_0))}
 \int_{B_r(x_0)}u(y)\,\mathrm d\lambda_n(y).
\end{equation}
\end{twierdzenie}
\begin{proof}
Ustalmy $x_0\in\Omega$ i oznaczmy
$d_0=\operatorname{dist}(x_0,\partial\Omega)$.
Rozważmy $M(r)=M_u(x_0,r)$ dla $0<r<d_0$.
Dla każdego zwartego przedziału $[a,b]\subset(0,d_0)$
funkcja $(r,\theta)\mapsto u(x_0+r\theta)$ oraz jej
pochodna $\nabla u(x_0+r\theta)\cdot\theta$ są ciągłe
na $[a,b]\times S^{n-1}$, a więc ograniczone.
Ponieważ sfera ma skończoną miarę, można różniczkować
pod znakiem całki. Na $\partial B_r(x_0)$ jednostkowa
normalna zewnętrzna wynosi $\nu(y)=(y-x_0)/r$.
Zmiana zmiennych $y=x_0+r\theta$, a następnie
twierdzenie Gaussa zastosowane do pola $\nabla u$
na kuli dają
\[
 M'(r)=\frac{1}{\omega_n r^{n-1}}
 \int_{\partial B_r(x_0)}\partial_\nu u(y)\,
 \mathrm d\sigma(y)
 =\frac{1}{\omega_n r^{n-1}}
 \int_{B_r(x_0)}\Delta u(y)\,\mathrm d\lambda_n(y)=0.
\]
Zatem $M$ jest stała na $(0,d_0)$. Ponadto
\[
 |M(r)-u(x_0)|\leq
 \sup_{\theta\in S^{n-1}}|u(x_0+r\theta)-u(x_0)|
 \longrightarrow0\quad(r\downarrow0)
\]
z ciągłości $u$ w $x_0$, co daje średnią sferyczną.
Dla $\overline{B_r(x_0)}\subset\Omega$ mamy $r<d_0$.
We współrzędnych biegunowych otrzymujemy
\[
 \begin{aligned}
 \int_{B_r(x_0)}u(y)\,\mathrm d\lambda_n(y)
 &=\omega_n\int_0^r M(s)s^{n-1}\,\mathrm d\lambda_1(s)\\
 &=u(x_0)\frac{\omega_n r^n}{n}
 =u(x_0)\lambda_n(B_r(x_0)).
 \end{aligned}
\]
Po podzieleniu przez dodatnią miarę kuli otrzymujemy
drugą równość w~\eqref{eq:pde3-both-means}.
\end{proof}

\begin{stwierdzenie}[Odwrotność własności średniej]
\label{prop:pde3-converse}
Niech $n\geq2$, $\Omega\subset\R^n$ będzie otwarta,
a $u\in C^2(\Omega)$. Jeśli dla każdego $x\in\Omega$
i wszystkich dostatecznie małych $r>0$ spełniona jest
pierwsza równość w~\eqref{eq:pde3-both-means}, to
$\Delta u=0$ w $\Omega$. Ten sam wniosek zachodzi przy
założeniu własności średniej po kulach.
\end{stwierdzenie}
\begin{proof}
Ustalmy $x\in\Omega$ i wybierzmy $r_0>0$ takie,
że $\overline{B_{r_0}(x)}\subset\Omega$.
Dla $|h|<r_0$ wzór Taylora z resztą całkową daje
\[
 \begin{aligned}
 u(x+h)&=u(x)+\nabla u(x)\cdot h
 +\tfrac12 h^{\mathsf T}D^2u(x)h+E(h),\\
 E(h)&=\int_0^1(1-t)h^{\mathsf T}
 \bigl(D^2u(x+th)-D^2u(x)\bigr)h\,\mathrm d\lambda_1(t).
 \end{aligned}
\]
Z ciągłości hesjanu wynika
$|E(h)|\leq |h|^2\eta(|h|)$, gdzie
$\eta(r)=\frac12\sup_{|z-x|\leq r}\|D^2u(z)-D^2u(x)\|$
dąży do zera przy $r\downarrow0$, a norma macierzy
jest normą operatorową indukowaną przez normę
euklidesową. W szczególności reszta dla $h=r\theta$
jest $o(r^2)$ jednostajnie względem $\theta\in S^{n-1}$.

Odbicie jednej współrzędnej zachowuje miarę sfery,
więc całki z $\theta_i$ oraz z $\theta_i\theta_j$
dla $i\ne j$ znikają. Permutacje współrzędnych dają
równość wszystkich całek z $\theta_i^2$.
Ponieważ $\sum_{i=1}^n\theta_i^2=1$ na sferze,
każda z tych całek wynosi $\omega_n/n$. Stąd
\[
 M_u(x,r)=u(x)+\frac{r^2}{2n}\Delta u(x)+e(r),
 \qquad |e(r)|\leq r^2\eta(r).
\]
Jeśli $M_u(x,r)=u(x)$ dla wszystkich dostatecznie małych
$r$, po odjęciu $u(x)$ i podzieleniu przez $r^2$
przechodzimy do granicy, otrzymując $\Delta u(x)=0$.

Dla średniej po kuli całkujemy ostatni wzór
w zmiennej promieniowej. Otrzymujemy
\[
 \begin{aligned}
 \frac{1}{\lambda_n(B_r(x))}
 \int_{B_r(x)}u(y)\,\mathrm d\lambda_n(y)
 &=\frac{n}{r^n}\int_0^r M_u(x,s)s^{n-1}\,
       \mathrm d\lambda_1(s)\\
 &=u(x)+\frac{r^2}{2(n+2)}\Delta u(x)+E_r .
 \end{aligned}
\]
Ponieważ $\eta$ jest niemalejąca,
$|E_r|\leq \frac{n}{n+2}r^2\eta(r)=o(r^2)$.
Własność średniej po kulach również daje więc
$\Delta u(x)=0$. Punkt $x$ był dowolny.
\end{proof}

\begin{uwaga}[Dlaczego średnia po kulach nie jest przypadkowa?]
Wartość funkcji harmonicznej w środku jest dokładnie
równowagą jej wartości wokół środka. Dla funkcji
subharmonicznej o klasie $C^2$ rachunek z dowodu pokazuje,
że $M_u(x,r)$ jest niemalejąca jako funkcja $r$, więc
$u(x)\leq M_u(x,r)$. Dla funkcji superharmonicznej
nierówność ma kierunek przeciwny.
\end{uwaga}

\begin{twierdzenie}[Własność wartości średniej bez założenia $C^2$]
\label{thm:pde-v31-continuous-mean}
Niech $n\ge2$, $\Omega\subset\R^n$ będzie otwarta,
a $u\in C(\Omega)$. Załóżmy, że dla każdego
$x\in\Omega$ i każdego $r>0$ z
$\overline{B_r(x)}\subset\Omega$ zachodzi
\[
 u(x)=\frac1{\omega_nr^{n-1}}
 \int_{\partial B_r(x)}u(y)\,\mathrm d\sigma(y).
\]
Wówczas $u\in C^\infty(\Omega)$ i $\Delta u=0$ w $\Omega$.
\end{twierdzenie}
\begin{proof}
Ustalmy $x_0\in\Omega$ i wybierzmy $R>0$ takie,
że $\overline{B_{2R}(x_0)}\subset\Omega$.
Wybierzmy radialną, nieujemną funkcję
$\rho\in C_c^\infty(B_1(0))$ o całce $1$,
na przykład odpowiednio unormowaną funkcję postaci
$\rho(z)=c\,\exp(-1/(1-4|z|^2))$ dla $|z|<1/2$
i $\rho(z)=0$ w pozostałych punktach.
Dla $0<\eps<R/2$ połóżmy
$\rho_\eps(z)=\eps^{-n}\rho(z/\eps)$.
Zmiana zmiennych $z=\eps w$ daje
$\int_{\R^n}\rho_\eps\,\mathrm d\lambda_n=1$.
Dla $x\in B_{3R/2}(x_0)$ określmy
\[
 u_\eps(x)=\int_{B_\eps(0)}\rho_\eps(z)u(x-z)\,
 \mathrm d\lambda_n(z).
\]
Argument $x-z$ należy do $B_{2R}(x_0)$,
więc całka jest dobrze określona. Zmieniając zmienne
$y=x-z$ i korzystając z nośnika $\rho_\eps$, mamy
\[
 u_\eps(x)=\int_{B_{2R}(x_0)}\rho_\eps(x-y)u(y)\,
 \mathrm d\lambda_n(y).
\]
Funkcja $u$ jest ograniczona na $\overline{B_{2R}(x_0)}$,
a każda pochodna $\rho_\eps$ jest ograniczona na $\R^n$.
Dla każdej pochodnej względem $x$ istnieje zatem
stała całkowalna na $B_{2R}(x_0)$, dominująca moduł
odpowiedniej funkcji podcałkowej.
Wielokrotne różniczkowanie pod całką daje
$u_\eps\in C^\infty(B_{3R/2}(x_0))$.

Radialność pozwala napisać $\rho_\eps(z)=k_\eps(|z|)$.
Dla $x\in B_{3R/2}(x_0)$ i $0<s<\eps$ mamy
$\overline{B_s(x)}\subset B_{2R}(x_0)$.
Stosując współrzędne biegunowe oraz założoną własność
średniej na sferze o środku $x$, otrzymujemy
\[
 \begin{aligned}
 u_\eps(x)
 &=\int_0^\eps k_\eps(s)s^{n-1}
       \int_{S^{n-1}}u(x-s\theta)\,\mathrm d\sigma(\theta)
       \,\mathrm d\lambda_1(s)\\
 &=u(x)\omega_n\int_0^\eps k_\eps(s)s^{n-1}
       \,\mathrm d\lambda_1(s)
 =u(x).
 \end{aligned}
\]
Zmiana $\theta\mapsto-\theta$ zachowuje miarę sfery,
a ostatnia równość wynika z normalizacji $\rho_\eps$.
Zatem $u$ jest gładka na $B_{3R/2}(x_0)$, bez potrzeby
przechodzenia z $\eps$ do zera. Stwierdzenie
\ref{prop:pde3-converse} daje $\Delta u=0$ w tym zbiorze.
Dowolność $x_0$ kończy dowód.

Warto zauważyć także inną drogę: dla dowolnego (niekoniecznie
radialnego) nieujemnego jądra wygładzającego o całce $1$
i nośniku w $B_\eps(0)$ powyższy splot jest gładki.
Dla $x\in B_R(x_0)$ oraz $0<r<R/2$ twierdzenie Fubiniego
można zastosować do całki po $\partial B_r(x)\times B_\eps(0)$,
ponieważ funkcja podcałkowa jest ograniczona, a obie miary
są skończone. Własność średniej $u$ dla kul o środkach
$x-z$ daje własność średniej $u_\eps$; stąd
$\Delta u_\eps=0$ na $B_R(x_0)$. Jednostajna ciągłość
$u$ na $\overline{B_{2R}(x_0)}$ daje
$u_\eps\to u$ jednostajnie na $\overline{B_R(x_0)}$.
Po udowodnieniu twierdzenia o jądrze Poissona
(\ref{thm:pde3-poisson}) można zastosować jego wzór
do $u_\eps$ i przejść do granicy, korzystając z dodatniości
i normalizacji jądra. Otrzymuje się w ten sposób
alternatywny dowód harmoniczności i gładkości $u$.
\end{proof}

\begin{uwaga}
Kluczowe jest, że własność średniej zakładamy dla
\emph{wszystkich} kul o domknięciu zawartym w $\Omega$.
Nie należy bez dodatkowego argumentu zastępować tego
warunku wymaganiem dotyczącym jednego promienia
zależnego od środka kuli.
\end{uwaga}

\section{Mocna zasada maksimum, jednoznaczność i twierdzenie Liouville'a}
\label{sec:pde3-max}

Słabą zasadę maksimum udowodniliśmy już w twierdzeniu
\ref{thm:laplace-max}. Własność wartości średniej daje
twierdzenie silniejsze, które nie wymaga ograniczoności obszaru.

\begin{twierdzenie}[Mocna zasada maksimum dla funkcji harmonicznych]
\label{thm:pde3-strong-max}
Niech $\Omega\subset\R^n$ będzie obszarem, $n\geq2$,
a $u\in C^2(\Omega)$ funkcją harmoniczną. Jeśli
$u$ osiąga maksimum lokalne albo minimum lokalne w pewnym
punkcie $x_0\in\Omega$, to jest stała na $\Omega$.
\end{twierdzenie}
\begin{proof}
Niech $u$ osiąga maksimum lokalne w $x_0$ i oznaczmy
$M=u(x_0)$. Wybierzmy $r_0>0$ tak, aby
$\overline{B_{r_0}(x_0)}\subset\Omega$ oraz $u\leq M$
na $B_{r_0}(x_0)$. Dla każdego $0<r<r_0$ własność
wartości średniej daje
\[
 \int_{\partial B_r(x_0)}(M-u(y))\,\mathrm d\sigma(y)=0.
\]
Funkcja podcałkowa jest ciągła i nieujemna.
Gdyby była dodatnia w pewnym punkcie sfery,
byłaby ograniczona z dołu przez dodatnią stałą
na pewnym niepustym względnie otwartym fragmencie
sfery. Każdy taki fragment ma dodatnią miarę
powierzchniową, co przeczyłoby zerowaniu się całki.
Stąd $u=M$ na każdej sferze $\partial B_r(x_0)$
dla $0<r<r_0$, a zatem na całej kuli $B_{r_0}(x_0)$.

Jeśli $M$ jest maksimum globalnym, można zakończyć
dowód bez analityczności: zbiór $\{x\in\Omega:u(x)=M\}$
jest niepusty i względnie domknięty z ciągłości $u$.
Powyższy argument zastosowany w dowolnym jego punkcie
pokazuje, że jest również otwarty. Spójność $\Omega$
wymusza jego równość z $\Omega$.

Dla maksimum jedynie lokalnego nie możemy zakładać
nierówności $u\leq M$ na całym obszarze.
Korzystamy wówczas z twierdzenia o analityczności
funkcji harmonicznych (twierdzenie~\ref{thm:pde3-analytic}).
Niech $w=u-M$ i niech $\mathbb N_0=\{0,1,2,\ldots\}$.
Rozważmy zbiór
\[
 A=\{x\in\Omega:D^\alpha w(x)=0
       \text{ dla każdego }\alpha\in\mathbb N_0^n\},
\]
gdzie $D^\alpha=\partial_{x_1}^{\alpha_1}\cdots
\partial_{x_n}^{\alpha_n}$ i $D^0w=w$.
Zbiór $A$ zawiera $B_{r_0}(x_0)$ i jest względnie
domknięty jako część wspólna zbiorów zerowych
ciągłych pochodnych $D^\alpha w$.
Jeżeli $x\in A$, to rzeczywista analityczność $w$
oznacza, że na pewnej kuli o środku $x$ funkcja $w$
jest sumą swojego szeregu Taylora. Wszystkie jego
współczynniki $D^\alpha w(x)/\alpha!$ są zerowe
(tutaj $\alpha!=\prod_{j=1}^n\alpha_j!$).
Stąd $w$ znika na tej kuli, a wraz z nią wszystkie
jej pochodne. Kula ta zawiera się więc w $A$.
Wynika stąd, że $A$ jest również otwarty.
Ze spójności $\Omega$ otrzymujemy $A=\Omega$,
a w szczególności $w=0$ na $\Omega$.

Nie ma tu błędnego koła: dowód analityczności korzysta
z twierdzenia o jądrze Poissona, którego jednoznaczność
wynika ze słabej zasady maksimum udowodnionej w rozdziale
pierwszym, a konstrukcja i normalizacja jądra korzystają
z funkcji Greena i wzoru reprezentacyjnego Greena.
Żaden z tych argumentów nie wymaga obecnego twierdzenia.
Przypadek minimum sprowadza się do maksimum funkcji $-u$.
\end{proof}

\begin{stwierdzenie}[Oszacowanie pochodnych w środku kuli]
\label{prop:pde3-gradient}
Niech $n\geq2$, $R>0$, $x_0\in\R^n$,
a $u\in C^2(B_R(x_0))\cap C(\overline{B_R(x_0)})$
będzie harmoniczna. Wówczas dla każdego
$j=1,\dots,n$ zachodzi
\begin{equation}\label{eq:pde3-derivative-bound}
 |\partial_{x_j}u(x_0)|\leq \frac nR
 \sup_{y\in B_R(x_0)}|u(y)|.
\end{equation}
\end{stwierdzenie}
\begin{proof}
Ustalmy $0<\rho<R$ i połóżmy $\delta=(R-\rho)/2$.
Jeśli $x\in B_\delta(x_0)$, to dla każdego
$z\in\overline{B_\rho(0)}$ mamy
$|x+z-x_0|\leq|x-x_0|+|z|<\delta+\rho<R$.
Możemy więc zastosować własność wartości średniej
po kuli o ustalonym promieniu $\rho$ i ruchomym środku $x$:
\[
 u(x)=\frac{1}{\lambda_n(B_\rho(0))}
 \int_{B_\rho(0)}u(x+z)\,\mathrm d\lambda_n(z).
\]
Funkcja podcałkowa oraz jej pochodne względem $x$
są ciągłe i ograniczone dla
$(x,z)\in\overline{B_\delta(x_0)}\times
\overline{B_\rho(0)}$, ponieważ wszystkie argumenty $x+z$
leżą w zwartej kuli o promieniu $\delta+\rho<R$.
Różniczkując względem $x_j$ pod znakiem całki
w punkcie $x_0$ i podstawiając $y=x_0+z$, otrzymujemy
\[
 \partial_{x_j}u(x_0)
 =\frac{1}{\lambda_n(B_\rho(0))}
 \int_{B_\rho(x_0)}\partial_{x_j}u(y)\,
 \mathrm d\lambda_n(y).
\]
Nie potrzebujemy przy tym zakładać harmoniczności
pochodnych $u$, która wymagałaby dodatkowej regularności.
Dla $e_j$ oznaczającego $j$-ty wektor standardowej
bazy $\R^n$ pole $y\mapsto u(y)e_j$ jest klasy $C^1$
na otoczeniu domkniętej kuli, a jego dywergencją jest
$\partial_{x_j}u$. Twierdzenie Gaussa daje zatem
\[
 \partial_{x_j}u(x_0)=\frac{1}{\lambda_n(B_\rho(0))}
 \int_{\partial B_\rho(x_0)}u(y)\nu_j(y)\,
 \mathrm d\sigma(y),
\]
gdzie $\nu_j(y)$ jest $j$-tą współrzędną jednostkowej
normalnej zewnętrznej $\nu(y)=(y-x_0)/\rho$.
Oznaczmy $K=\sup_{y\in B_R(x_0)}|u(y)|<\infty$,
co wynika z ciągłości $u$ na domkniętej kuli.
Korzystając z $|\nu_j|\leq1$, otrzymujemy
\[
 |\partial_{x_j}u(x_0)|
 \leq K\frac{\omega_n\rho^{n-1}}{\omega_n\rho^n/n}
 =\frac{nK}{\rho}.
\]
Lewa strona i $K$ nie zależą od $\rho$. Przechodząc
z $\rho$ do $R$ od dołu, dostajemy
\eqref{eq:pde3-derivative-bound}.
\end{proof}

\begin{wniosek}[Twierdzenie Liouville'a]
\label{cor:pde3-liouville}
Niech $n\geq2$ i $u\in C^2(\R^n)$ będzie harmoniczna
oraz ograniczona. Wtedy $u$ jest stała.
\end{wniosek}
\begin{proof}
Niech $K=\sup_{x\in\R^n}|u(x)|<\infty$.
Dla każdego $x_0\in\R^n$, $R>0$ oraz $j=1,\ldots,n$
możemy zastosować \eqref{eq:pde3-derivative-bound},
gdyż $u$ jest harmoniczna na $B_R(x_0)$
i ciągła na jej domknięciu. Otrzymujemy
$|\partial_{x_j}u(x_0)|\leq nK/R$.
Przechodząc z $R$ do nieskończoności przy ustalonym $x_0$,
dostajemy $\partial_{x_j}u(x_0)=0$ dla każdego $j$.
Ponieważ $x_0$ był dowolny, mamy $\nabla u=0$ na $\R^n$.
Dla dowolnych $x,y\in\R^n$ reguła łańcuchowa
i twierdzenie podstawowe rachunku całkowego dają
\[
 u(y)-u(x)=\int_0^1
 \nabla u(x+t(y-x))\cdot(y-x)\,\mathrm d\lambda_1(t)=0.
\]
Zatem $u$ jest stała.
\end{proof}

\begin{kontrprzyklad}[Dlaczego ograniczoność jest ważna?]
Funkcja $u(x)=x_1$ jest harmoniczna na $\R^n$,
lecz nie jest stała. Twierdzenie Liouville'a nie
obowiązuje bez ograniczoności (ani analogicznej
kontroli wzrostu).
\end{kontrprzyklad}

\section{Rozwiązanie fundamentalne laplasjanu}
\label{sec:pde3-fundamental}

Rozwiązaniem fundamentalnym operatora nazywa się funkcję lub dystrybucję,
która reaguje na punktowe źródło. Wprowadzimy ją najpierw
bez formalnej teorii dystrybucji: zapis
$-\Delta\Phi=\delta_0$ będzie oznaczał precyzyjną
tożsamość całkową, gdzie $\delta_0$ jest funkcjonałem
$\varphi\mapsto\varphi(0)$ na funkcjach testowych.
Tutaj $C_c^\infty(\R^n)$ oznacza przestrzeń funkcji
nieskończenie różniczkowalnych o zwartym nośniku, a
$L^1_{\mathrm{loc}}(\R^n)$ klasę funkcji całkowalnych
na każdym zwartym podzbiorze $\R^n$.

\begin{twierdzenie}[Rozwiązanie fundamentalne $-\Delta$]
\label{thm:pde3-phi}
Niech $n\geq2$ i $\omega_n=\sigma(S^{n-1})$. Dla
$x\in\R^n\setminus\{0\}$ połóżmy
\begin{equation}\label{eq:pde3-fundamental}
 \Phi_n(x)=
 \begin{cases}
 \displaystyle\frac{1}{(n-2)\omega_n}|x|^{2-n},&n\geq3,\\[4pt]
 \displaystyle-\frac1{2\pi}\log|x|,&n=2.
 \end{cases}
\end{equation}
Wartość $\Phi_n(0)$ ustalamy dowolnie, na przykład
$\Phi_n(0)=0$; nie wpływa ona na poniższe całki.
Wówczas $\Phi_n\in L^1_{\mathrm{loc}}(\R^n)$,
$\Delta\Phi_n=0$ na $\R^n\setminus\{0\}$ oraz
\begin{equation}\label{eq:pde3-distdelta}
 \int_{\R^n}\Phi_n(x)(-\Delta\varphi(x))\,
 \mathrm d\lambda_n(x)=\varphi(0)
 \quad\text{dla każdej }\varphi\in C_c^\infty(\R^n).
\end{equation}
\end{twierdzenie}
\begin{proof}
Niech $h\in C^2((0,\infty))$ i $w(x)=h(|x|)$
dla $x\ne0$. Oznaczmy $r=|x|$. Ponieważ
$\partial_{x_j}r=x_j/r$, reguła łańcuchowa daje
\[
 \partial_{x_j}w=h'(r)\frac{x_j}{r},\qquad
 \partial_{x_j}^2w
 =h''(r)\frac{x_j^2}{r^2}
 +h'(r)\left(\frac1r-\frac{x_j^2}{r^3}\right).
\]
Sumując względem $j$ i korzystając z
$\sum_{j=1}^n x_j^2=r^2$, otrzymujemy
$\Delta w=h''(r)+(n-1)h'(r)/r$.
Równanie $\Delta w=0$ jest więc równoważne
$(r^{n-1}h'(r))'=0$, skąd $h'(r)=Cr^{1-n}$.
Całkowanie daje $h(r)=A+Cr^{2-n}/(2-n)$ dla $n\geq3$
oraz $h(r)=A+C\log r$ dla $n=2$.

Aby otrzymać jednostkowe punktowe źródło dla operatora
$-\Delta$, wybieramy $C=-1/\omega_n$; wtedy strumień
$-\nabla w$ przez każdą sferę o środku w zerze wynosi
$-h'(r)\omega_n r^{n-1}=1$.
Przy $A=0$ i $\omega_2=2\pi$ otrzymujemy dokładnie
\eqref{eq:pde3-fundamental}. Oznaczmy przez $h_n(r)$
profil radialny tej funkcji, tak że
$\Phi_n(x)=h_n(|x|)$ dla $x\ne0$. Mamy
$h_n'(r)=-1/(\omega_n r^{n-1})$ i
$\nabla\Phi_n(x)=-x/(\omega_n |x|^n)$.
Wykażemy poniżej, że wybrany strumień daje również
żądaną tożsamość całkową.

Funkcja $\Phi_n$ jest gładka poza zerem.
Dla $n\geq3$ i $a>0$ współrzędne biegunowe dają
\[
 \int_{B_a(0)}|\Phi_n(x)|\,\mathrm d\lambda_n(x)
 =\frac1{n-2}\int_0^a r\,\mathrm d\lambda_1(r)
 =\frac{a^2}{2(n-2)}.
\]
Dla $n=2$ oraz $0<a\leq1$ analogicznie
\[
 \int_{B_a(0)}|\Phi_2(x)|\,\mathrm d\lambda_2(x)
 =\int_0^a(-\log r)r\,\mathrm d\lambda_1(r)
 =\frac{a^2}{4}-\frac{a^2}{2}\log a.
\]
W ostatnim rachunku używamy
$\lim_{r\downarrow0}r^2\log r=0$.
Obie całki są skończone i dążą do zera wraz z $a$.
W połączeniu z gładkością poza zerem daje to
$\Phi_n\in L^1_{\mathrm{loc}}(\R^n)$.

Niech $\varphi\in C_c^\infty(\R^n)$ i wybierzmy $R>0$
tak, aby jej nośnik był zawarty w $B_R(0)$.
Dla $0<\eps<R$ rozważmy pierścień
$A_{\eps,R}=B_R(0)\setminus\overline{B_\eps(0)}$.
Funkcje $\Phi_n$ i $\varphi$ są gładkie na otoczeniu
jego domknięcia. Stosujemy twierdzenie Gaussa do pola
$\varphi\nabla\Phi_n-\Phi_n\nabla\varphi$,
którego dywergencja jest równa
$\varphi\Delta\Phi_n-\Phi_n\Delta\varphi
=-\Phi_n\Delta\varphi$ na pierścieniu.
Jest to druga tożsamość Greena, wyprowadzona tu
bez odwoływania się do dalszej części rozdziału.
Na zewnętrznym brzegu $\varphi$ i jej pochodne
znikają, gdyż $\varphi=0$ w jego otoczeniu.
Na wewnętrznym brzegu jednostkowa normalna zewnętrzna
dla pierścienia wynosi $\nu(y)=-y/\eps$,
więc $\partial_\nu\Phi_n(y)=-h_n'(\eps)
=1/(\omega_n\eps^{n-1})$.
Otrzymujemy
\[
 \int_{A_{\eps,R}}\Phi_n(y)(-\Delta\varphi(y))\,
 \mathrm d\lambda_n(y)
 =\int_{\partial B_\eps(0)}
 \bigl(\varphi(y)\,\partial_\nu\Phi_n(y)
       -\Phi_n(y)\,\partial_\nu\varphi(y)\bigr)\,\mathrm d\sigma(y).
\]
Pierwszy składnik całki brzegowej jest średnią
sferyczną funkcji $\varphi$:
\[
 \left|\frac1{\omega_n\eps^{n-1}}
 \int_{\partial B_\eps(0)}\varphi(y)\,\mathrm d\sigma(y)
 -\varphi(0)\right|
 \leq\sup_{|y|=\eps}|\varphi(y)-\varphi(0)|
 \longrightarrow0.
\]
Dla drugiego składnika, z $K=\sup_{\R^n}|\nabla\varphi|$,
mamy oszacowanie
\[
 \left|\int_{\partial B_\eps(0)}
 \Phi_n(y)\partial_\nu\varphi(y)\,\mathrm d\sigma(y)\right|
 \leq K\omega_n\eps^{n-1}|h_n(\eps)|.
\]
Prawa strona wynosi $K\eps/(n-2)$ dla $n\geq3$,
a dla $n=2$ jest równa $K\eps|\log\eps|$.
W obu przypadkach dąży do zera.
Wreszcie całka po $\R^n$ różni się od całki po
$A_{\eps,R}$ tylko całką po $B_\eps(0)$ (z dokładnością
do zbiorów miary zero), ponieważ $\Delta\varphi=0$
poza $B_R(0)$. Moduł tej różnicy jest nie większy niż
$\sup_{\R^n}|\Delta\varphi|
\int_{B_\eps(0)}|\Phi_n(y)|\,\mathrm d\lambda_n(y)$,
co dąży do zera na mocy wcześniejszych obliczeń.
Przechodząc do granicy otrzymujemy
\eqref{eq:pde3-distdelta}.
\end{proof}

\begin{uwaga}[Wymiar jeden]
Dla $n=1$ odpowiednikiem rozwiązania fundamentalnego
jest $\Phi_1(x)=-|x|/2$. Istotnie, dla dowolnej
$\varphi\in C_c^\infty(\R)$ całkowanie przez części
na obu półprostych daje
\[
 \begin{aligned}
 \int_{-\infty}^0(-x)\varphi''(x)\,\mathrm d\lambda_1(x)
 &=\varphi(0),\\
 \int_0^\infty x\varphi''(x)\,\mathrm d\lambda_1(x)
 &=\varphi(0).
 \end{aligned}
\]
Wyrazy brzegowe w nieskończoności znikają dzięki
zwartemu nośnikowi, a iloczyn $x\varphi'(x)$ znika
w zerze. Stąd $\int_\R \Phi_1(-\varphi'')\,
\mathrm d\lambda_1=\varphi(0)$, czyli
$-\Phi_1''=\delta_0$ w sensie dystrybucji.
Przypadek jednowymiarowy jest w istocie teorią ODE,
dlatego w dalszych konstrukcjach potencjału przyjmujemy $n\geq2$.
\end{uwaga}

\begin{twierdzenie}[Potencjał Newtona i równanie Poissona]
\label{thm:pde3-newton}
Niech $n\geq2$ i $f\in C_c^\infty(\R^n)$. Określmy
\begin{equation}\label{eq:pde3-newton}
 (\Phi_n*f)(x)=\int_{\R^n}\Phi_n(x-y)f(y)\,
 \mathrm d\lambda_n(y).
\end{equation}
Wówczas $u=\Phi_n*f\in C^\infty(\R^n)$ i
$-\Delta u=f$ na $\R^n$.
\end{twierdzenie}
\begin{proof}
Ustalmy $x_0\in\R^n$ i zwartą kulę
$K=\overline{B_1(x_0)}$. Niech $L=\supp f$;
dla $f=0$ teza jest oczywista, więc zakładamy $L\ne\emptyset$.
Zbiór $S=K-L=\{x-y:x\in K,\ y\in L\}$ jest zwarty
jako ciągły obraz zbioru zwartego $K\times L$.
Zmiana zmiennych $z=x-y$ ma moduł wyznacznika równy $1$.
Ponieważ $f$ jest ograniczona, a $\Phi_n$ całkowalna
na $S$, dla każdego $x\in K$ całka definiująca $u$ jest
bezwzględnie zbieżna i
\[
 u(x)=\int_{\R^n}\Phi_n(z)f(x-z)\,\mathrm d\lambda_n(z).
\]
Dla każdego wielowskaźnika $\alpha\in\mathbb N_0^n$
pochodna $D^\alpha f$ ma nośnik zawarty w $L$,
a $M_\alpha=\sup_{\R^n}|D^\alpha f|<\infty$.
Wobec tego dla $x\in K$ mamy
\[
 |\Phi_n(z)D^\alpha f(x-z)|
 \leq M_\alpha|\Phi_n(z)|\chi_S(z).
\]
Prawa strona jest całkowalna względem $\lambda_n$.
Na wnętrzu $K$ można więc wielokrotnie różniczkować
pod całką: na każdym kroku ilorazy różnicowe
są zdominowane przez powyższe oszacowanie dla
odpowiedniej pochodnej następnego rzędu, z twierdzenia
o wartości średniej, a ciągłość otrzymanych całek
wynika z twierdzenia o zbieżności zdominowanej.
Dostajemy
\[
 D^\alpha u(x)=\int_{\R^n}
 \Phi_n(z)D^\alpha f(x-z)\,\mathrm d\lambda_n(z)
 \qquad(x\in B_1(x_0)).
\]
Ponieważ $x_0$ był dowolny, $u\in C^\infty(\R^n)$.

Dla ustalonego $x\in\R^n$ funkcja
$\varphi_x(z)=f(x-z)$ należy do $C_c^\infty(\R^n)$.
Reguła łańcuchowa daje
$\partial_{z_j}\varphi_x(z)=-\partial_{x_j}f(x-z)$,
a przy drugim różniczkowaniu znak minus znika.
Zatem $\Delta_z\varphi_x(z)=(\Delta f)(x-z)$.
Stosując \eqref{eq:pde3-distdelta} do $\varphi_x$
i korzystając z uzasadnionego różniczkowania splotu,
otrzymujemy
\[
 -\Delta u(x)=\int_{\R^n}
 \Phi_n(z)(-\Delta_z\varphi_x(z))\,\mathrm d\lambda_n(z)
 =\varphi_x(0)=f(x).
\]
Równanie zachodzi punktowo na całej przestrzeni.
\end{proof}

\begin{uwaga}[Rozwiązanie szczególne a niejednoznaczność]
Jeżeli $u,v\in C^2(\R^n)$ spełniają
$-\Delta u=-\Delta v=f$, to ich różnica jest harmoniczna.
Odwrotnie, dodanie dowolnej funkcji harmonicznej
na $\R^n$ do rozwiązania zachowuje równanie.
Zatem dla $f\in C_c^\infty(\R^n)$ wszystkie rozwiązania
klasyczne mają postać $\Phi_n*f+h$, gdzie
$h\in C^2(\R^n)$ i $\Delta h=0$.
Przykładowo $\Phi_n*f+x_1$ jest innym rozwiązaniem.
Do jednoznaczności potrzebne są dodatkowe warunki
eliminujące swobodę dodawania funkcji harmonicznych.
\end{uwaga}

\section{Funkcje Greena i tożsamości Greena}
\label{sec:pde3-green}

\begin{stwierdzenie}[Dwie tożsamości Greena]
\label{prop:pde3-green-identities}
Niech $\Omega\subset\R^n$ będzie obszarem ograniczonym
o brzegu klasy $C^1$ i niech $u,v\in C^2(\overline\Omega)$
będą funkcjami rzeczywistymi. Przez $C^2(\overline\Omega)$
rozumiemy tutaj funkcje klasy $C^2$ w $\Omega$, których
wartości i pochodne do rzędu drugiego przedłużają się
ciągle na $\overline\Omega$. Niech $\nu$ oznacza
jednostkową normalną zewnętrzną do $\partial\Omega$,
a $\partial_\nu v=\nabla v\cdot\nu$. Wtedy
\begin{align}
 \int_\Omega u(x)\Delta v(x)\,\mathrm d\lambda_n(x)
 &=-\int_\Omega\nabla u(x)\cdot\nabla v(x)\,\mathrm d\lambda_n(x)
 +\int_{\partial\Omega}u(x)\,\partial_\nu v(x)\,\mathrm d\sigma(x),
 \label{eq:pde3-green1}\\
 \int_\Omega(u(x)\Delta v(x)-v(x)\Delta u(x))\,\mathrm d\lambda_n(x)
 &=\int_{\partial\Omega}(u(x)\partial_\nu v(x)-v(x)\partial_\nu u(x))
 \,\mathrm d\sigma(x).\label{eq:pde3-green2}
\end{align}
\end{stwierdzenie}
\begin{proof}
Pole wektorowe $u\nabla v$ jest klasy $C^1$
aż do brzegu, a jego dywergencja wynosi
$\diver(u\nabla v)=\nabla u\cdot\nabla v+u\Delta v$.
Twierdzenie Gaussa daje
\[
 \int_\Omega(\nabla u\cdot\nabla v+u\Delta v)\,
 \mathrm d\lambda_n
 =\int_{\partial\Omega}u\,\partial_\nu v\,\mathrm d\sigma.
\]
Wszystkie całki są skończone, ponieważ funkcje
podcałkowe są ciągłe na odpowiednich zbiorach
zwartych, a ograniczony obszar o brzegu klasy $C^1$
ma skończoną objętość i skończoną miarę brzegu.
Przeniesienie składnika z gradientami daje
\eqref{eq:pde3-green1}. Zapisując tę samą tożsamość
dla pary $(v,u)$ i odejmując ją od pierwszej,
korzystamy z równości $\nabla u\cdot\nabla v
=\nabla v\cdot\nabla u$. Otrzymujemy
\eqref{eq:pde3-green2}.
\end{proof}

\begin{definicja}[Funkcja Greena dla $-\Delta$]
\label{def:pde3-green}
Niech $\Omega\subset\R^n$ będzie obszarem ograniczonym,
$n\geq2$. Dla ustalonego $x\in\Omega$ funkcją Greena
(z warunkiem Dirichleta) nazywamy funkcję
$G_\Omega(x,\cdot)$ taką, że:
\begin{enumerate}
\item jest lokalnie całkowalna w $\Omega$ i klasy
$C^2$ na $\Omega\setminus\{x\}$;
\item $-\Delta_yG_\Omega(x,y)=\delta_x$ w sensie
dystrybucji, to znaczy
\[
 \int_\Omega G_\Omega(x,y)(-\Delta\varphi(y))\,
 \mathrm d\lambda_n(y)=\varphi(x)
 \quad(\varphi\in C_c^\infty(\Omega));
\]
\item funkcja $G_\Omega(x,\cdot)$ ma ciągłe przedłużenie
na $\overline\Omega\setminus\{x\}$, którego wartości
na $\partial\Omega$ są równe zeru.
\end{enumerate}
Wartość w punkcie $y=x$ nie ma znaczenia dla tożsamości
całkowej. Drugi warunek wraz z pierwszym wymusza
$\Delta_yG_\Omega(x,y)=0$ dla $y\ne x$: dla funkcji
testowych o nośniku z dala od $x$ dwukrotne całkowanie
przez części daje zerową całkę z ciągłej funkcji
$\Delta_yG_\Omega$ pomnożonej przez dowolną taką funkcję
testową. Gdyby ta ciągła funkcja była niezerowa w pewnym
punkcie, test nieujemny o nośniku w otoczeniu, na którym
ma ona stały niezerowy znak, dawałby sprzeczność.
W zastosowaniach do całek brzegowych wymagamy również
odpowiedniej gładkości funkcji Greena względem $y$ aż do brzegu.
\end{definicja}

\begin{twierdzenie}[Wzór reprezentacyjny Greena]
\label{thm:pde3-representation}
Niech $n\geq2$, $\Omega\subset\R^n$ będzie obszarem ograniczonym o brzegu klasy $C^2$, $f\in C(\overline\Omega)$
i $u\in C^2(\overline\Omega)$ spełnia $-\Delta u=f$
w $\Omega$. Załóżmy, że dla każdego $x\in\Omega$
istnieje funkcja Greena $G_\Omega(x,y)=\Phi_n(x-y)-h_x(y)$,
gdzie $h_x$ jest harmoniczna i klasy $C^2$ na pewnym
otoczeniu $\overline\Omega$ (w zmiennej $y$).
Wówczas dla $x\in\Omega$
\begin{equation}\label{eq:pde3-green-representation}
 u(x)=\int_\Omega G_\Omega(x,y)f(y)\,\mathrm d\lambda_n(y)
 -\int_{\partial\Omega}u(y)\partial_{\nu_y}G_\Omega(x,y)
 \,\mathrm d\sigma(y).
\end{equation}
\end{twierdzenie}
\begin{proof}
Ustalmy $x\in\Omega$ i wybierzmy
$0<\eps<\operatorname{dist}(x,\partial\Omega)$.
Dla skrócenia zapisu połóżmy $G(y)=G_\Omega(x,y)$.
Funkcje $u$ i $G$ są klasy $C^2$ aż do brzegu
zbioru $\Omega_\eps=\Omega\setminus\overline{B_\eps(x)}$;
jego brzeg składa się z $\partial\Omega$ i $\partial B_\eps(x)$.
Druga tożsamość Greena dla $u$ i $G_\Omega(x,\cdot)$
na $\Omega_\eps=\Omega\setminus\overline{B_\eps(x)}$
daje, ponieważ $\Delta G_\Omega=0$ poza $x$,
\[
 \int_{\Omega_\eps}G_\Omega(x,y)f(y)\,\mathrm d\lambda_n(y)
 =\int_{\partial\Omega_\eps}
 \bigl(u(y)\partial_{\nu_y}G_\Omega(x,y)
 -G_\Omega(x,y)\partial_{\nu_y}u(y)\bigr)
 \,\mathrm d\sigma(y).
\]
Na $\partial\Omega$ normalna jest normalną zewnętrzną
do $\Omega$ i $G=0$, więc pozostaje całka
$\int_{\partial\Omega}u(y)\partial_{\nu_y}G(y)\,
\mathrm d\sigma(y)$.
Na $\partial B_\eps(x)$ normalna zewnętrzna
do $\Omega_\eps$ wynosi $\nu_y=-(y-x)/\eps$.
Z radialności $\Phi_n$ i rachunku w twierdzeniu
\ref{thm:pde3-phi} mamy
$\partial_{\nu_y}\Phi_n(x-y)=1/(\omega_n\eps^{n-1})$.
Wkład wewnętrznej sfery jest więc równy
\[
 \begin{aligned}
 I_\eps={}&\frac1{\omega_n\eps^{n-1}}
 \int_{\partial B_\eps(x)}u(y)\,\mathrm d\sigma(y)\\
 &-\int_{\partial B_\eps(x)}
 \Phi_n(x-y)\partial_{\nu_y}u(y)\,\mathrm d\sigma(y)\\
 &+\int_{\partial B_\eps(x)}
 (h_x(y)\partial_{\nu_y}u(y)
 -u(y)\partial_{\nu_y}h_x(y))\,\mathrm d\sigma(y).
 \end{aligned}
\]
Pierwszy składnik dąży do $u(x)$ z ciągłości $u$.
Wybierzmy $a>0$ z $\overline{B_a(x)}\subset\Omega$
i ograniczmy się do $\eps<a$.
Moduł drugiego składnika jest nie większy niż
$\sup_{\overline{B_a(x)}}|\nabla u|\,
\omega_n\eps^{n-1}|h_n(\eps)|$, gdzie $h_n$
oznacza profil radialny rozwiązania fundamentalnego.
Oszacowanie to dąży do zera, jak wykazano
w twierdzeniu~\ref{thm:pde3-phi}.
Moduł trzeciego składnika jest ograniczony przez
$C_x\omega_n\eps^{n-1}$, gdzie można przyjąć
$C_x=\sup_{\overline{B_a(x)}}(|h_x||\nabla u|
+|u||\nabla h_x|)<\infty$.
Ponieważ $n\geq2$, ten składnik również znika.
Zatem $I_\eps\to u(x)$.

Pozostaje przejście do granicy w całce objętościowej.
Funkcja $f$ jest ograniczona, a $G$ jest całkowalna
w $\Omega$: w otoczeniu $x$ wynika to z lokalnej
całkowalności $\Phi_n$ i ograniczoności $h_x$,
poza tym otoczeniem z ciągłości na zbiorze zwartym.
W szczególności moduł pominiętej całki po $B_\eps(x)$
jest nie większy niż
\[
 \sup_{\overline\Omega}|f|
 \left(\int_{B_\eps(0)}|\Phi_n(z)|\,\mathrm d\lambda_n(z)
 +\sup_{\overline{B_a(x)}}|h_x|\,\lambda_n(B_\eps(0))\right)
 \longrightarrow0.
\]
Ostatecznie otrzymujemy
\[
 \int_\Omega G_\Omega(x,y)f(y)\,\mathrm d\lambda_n(y)
 =\int_{\partial\Omega}u(y)\partial_{\nu_y}G_\Omega(x,y)
 \,\mathrm d\sigma(y)+u(x).
\]
Przeniesienie całki brzegowej daje
\eqref{eq:pde3-green-representation}.
\end{proof}

\section{Jawne jądro Poissona i zagadnienie Dirichleta na kuli}
\label{sec:pde3-poisson-kernel}

W tym podrozdziale $n\geq2$, $R>0$,
$B_R=B_R(0)$ i $\zeta\in\partial B_R$.
Dla $y\in B_R\setminus\{0\}$ wprowadzamy punkt odwrócony
$y^*=R^2y/|y|^2$, leżący poza kulą. Zachodzi tożsamość
\begin{equation}\label{eq:pde3-kelvin-identity}
 |z-y^*|=\frac R{|y|}|z-y| \qquad (|z|=R).
\end{equation}
Istotnie, $|y^*|=R^2/|y|>R$, a dla $|z|=R$
\[
 \begin{aligned}
 |z-y^*|^2
 &=R^2-2R^2\frac{z\cdot y}{|y|^2}
       +\frac{R^4}{|y|^2}\\
 &=\frac{R^2}{|y|^2}(|z|^2-2z\cdot y+|y|^2).
 \end{aligned}
\]
Pierwiastkując obie strony, otrzymujemy
\eqref{eq:pde3-kelvin-identity}.

\begin{stwierdzenie}[Funkcja Greena dla kuli]
\label{prop:pde3-ball-green}
Dla $n\geq3$ i $x,y\in B_R$, $x\ne y$, $y\ne0$, funkcja
\begin{equation}\label{eq:pde3-ball-green}
 G_R(x,y)=\Phi_n(x-y)-\Bigl(\frac R{|y|}\Bigr)^{n-2}
 \Phi_n(x-y^*)
\end{equation}
jest funkcją Greena dla kuli w zmiennej $x$.
Dla $n=2$ odpowiednia funkcja ma postać
\[
 G_R(x,y)=\frac1{2\pi}\log\frac{|y|\,|x-y^*|}{R\,|x-y|}.
\]
W obu przypadkach formuły przedłużają się dla $y=0$
oraz $x\in B_R\setminus\{0\}$ przez granicę (dla $n\geq3$:
$G_R(x,0)=\Phi_n(x)-1/((n-2)\omega_n R^{n-2})$;
dla $n=2$: $G_R(x,0)=(2\pi)^{-1}\log(R/|x|)$).
Ponadto $G_R(x,y)=G_R(y,x)$ dla wszystkich
$x,y\in B_R$ z $x\ne y$. W szczególności jest to
funkcja Greena również w zmiennej $y$ przy ustalonym $x$.
\end{stwierdzenie}
\begin{proof}
Dla $y\ne0$ i $x\ne y$ wprowadźmy pomocnicze oznaczenia
\[
 Q_R(x,y)=R^2-2x\cdot y+\frac{|x|^2|y|^2}{R^2},
 \qquad H_R(x,y)=h_n\bigl(\sqrt{Q_R(x,y)}\bigr),
\]
gdzie $h_n$ jest profilem radialnym $\Phi_n$.
Bezpośrednie rozwinięcie kwadratu odległości daje
$Q_R(x,y)=|y|^2|x-y^*|^2/R^2$.
Dla $n\geq3$ wynika stąd
$H_R(x,y)=(R/|y|)^{n-2}\Phi_n(x-y^*)$.
Dla $n=2$ mamy
$H_R(x,y)=-(2\pi)^{-1}\log(|y|\,|x-y^*|/R)$.
W obu przypadkach podany wzór przyjmuje zatem postać
$G_R(x,y)=\Phi_n(x-y)-H_R(x,y)$.

Z nierówności Cauchy'ego--Schwarza wynika
$Q_R(x,y)\geq(R-|x||y|/R)^2>0$,
gdy $|x||y|<R^2$. Dlatego $H_R$ jest dobrze określona
i klasy $C^\infty$ na tym otwartym zbiorze par $(x,y)$,
również dla $y=0$ i na przekątnej $x=y$ wewnątrz kuli.
W szczególności $H_R(x,0)=h_n(R)$, co daje obie
podane granice dla ustalonego $x\ne0$.
Jeżeli $0\leq r<R$, to dla $|x|\leq R$ i $|y|\leq r$
zachodzi $Q_R(x,y)\geq(R-r)^2>0$.
Zatem $H_R$ oraz wszystkie jej pochodne są ograniczone
na tym zwartym zbiorze. Własność ta będzie potrzebna
przy całkowaniu funkcji Greena względem bieguna.

Przy ustalonym $y\in B_R\setminus\{0\}$ funkcja
$x\mapsto H_R(x,y)$ jest harmoniczna na otoczeniu
$\overline{B_R}$: dla $n\geq3$ jest wielokrotnością
rozwiązania fundamentalnego z biegunem $y^*$ poza
domkniętą kulą, a dla $n=2$ różni się od
$\Phi_2(x-y^*)$ o stałą niezależną od $x$.
Dla $y=0$ jest stała, więc również harmoniczna.
Funkcja $\Phi_n(x-y)$ jest lokalnie całkowalna
i spełnia $-\Delta_x\Phi_n(x-y)=\delta_y$,
na mocy twierdzenia~\ref{thm:pde3-phi} i zmiany zmiennych.
Odejmowanie $H_R$ zachowuje tę tożsamość.

Dla $|x|=R$ mamy $Q_R(x,y)=|x-y|^2>0$,
więc $H_R(x,y)=\Phi_n(x-y)$ i $G_R(x,y)=0$.
Ciągłość aż do brzegu poza biegunem wynika z jawnych
wzorów. Wszystkie warunki definicji funkcji Greena
są więc spełnione w zmiennej $x$.
Ponieważ $Q_R(x,y)=Q_R(y,x)$, mamy
$H_R(x,y)=H_R(y,x)$, a radialność daje
$\Phi_n(x-y)=\Phi_n(y-x)$. Stąd
$G_R(x,y)=G_R(y,x)$, także wtedy, gdy jeden z punktów
jest zerem. Symetria pozwala stosować tę samą konstrukcję
w zmiennej $y$ przy ustalonym $x$; harmoniczna poprawka
$H_R(x,\cdot)$ jest gładka na otoczeniu domkniętej kuli.
W szczególności spełnione są założenia dotyczące
funkcji Greena we wzorze reprezentacyjnym
\eqref{eq:pde3-green-representation}.
\end{proof}

\begin{center}
\begin{tikzpicture}[scale=.87,>=Stealth]
 \fill[blue!5] (0,0) circle (2);
 \draw[thick] (0,0) circle (2);
 \draw[densely dashed,gray!65] (0,0)--(4.25,0);
 \fill (0,0) circle (1.6pt);
 \node[below left] at (0,0) {$0$};
 \fill[blue!75!black] (1.1,0) circle (2.5pt);
 \node[above] at (1.1,0) {$y$};
 \fill[red!75!black] (3.63,0) circle (2.5pt);
 \node[above] at (3.63,0) {$y^*$};
 \draw[<->,thin] (0,-.3)--(2,-.3);
 \node[below] at (1,-.3) {$R$};
 \node[align=left,anchor=west] at (4.2,1.0) {$|y|\,|y^*|=R^2$\\punkt odwrócony leży poza kulą};
\end{tikzpicture}
\par\smallskip\small
Inwersja w sferze $\partial B_R$: punkt $y$ wewnątrz
odpowiada punktowi $y^*$ na zewnątrz. Pozwala to
zbudować harmoniczną poprawkę rozwiązania fundamentalnego.
\end{center}

\begin{twierdzenie}[Jądro Poissona dla kuli]
\label{thm:pde3-poisson}
Niech $n\geq2$, $R>0$ i $g\in C(\partial B_R)$.
Dla $x\in B_R$ oraz $\zeta\in\partial B_R$ zdefiniujmy
jądro $P_R$ i funkcję $u$ na $B_R$:
\begin{equation}\label{eq:pde3-poisson}
 P_R(x,\zeta)=\frac{R^2-|x|^2}{\omega_n R|x-\zeta|^n},
 \qquad
 u(x)=\int_{\partial B_R}P_R(x,\zeta)g(\zeta)
 \,\mathrm d\sigma(\zeta).
\end{equation}
Wówczas $u\in C^\infty(B_R)\cap C(\overline{B_R})$,
$\Delta u=0$ w $B_R$ i $u|_{\partial B_R}=g$.
Jest to jedyne rozwiązanie z klasy
$C^2(B_R)\cap C(\overline{B_R})$ tego problemu Dirichleta.
Ponadto $P_R(x,\zeta)>0$ oraz
$\int_{\partial B_R}P_R(x,\zeta)\,\mathrm d\sigma(\zeta)=1$.
\end{twierdzenie}
\begin{proof}
Najpierw wykażemy związek z normalną pochodną funkcji
Greena. Ustalmy biegun $y\in B_R$ oraz punkt
$\zeta\in\partial B_R$. Jednostkowa normalna zewnętrzna
do kuli wynosi $\nu_\zeta=\zeta/R$.
Korzystamy z jednolitej postaci poprawki harmonicznej
z dowodu stwierdzenia~\ref{prop:pde3-ball-green}:
\[
 G_R(x,y)=\Phi_n(x-y)-h_n(\sqrt{Q_R(x,y)}).
\]
W otoczeniu $\zeta$ obie funkcje po prawej stronie
są gładkie w zmiennej $x$, ponieważ $\zeta\ne y$
i $Q_R(\zeta,y)=|\zeta-y|^2>0$.
Wzór $h_n'(r)=-1/(\omega_n r^{n-1})$ jest prawdziwy
zarówno dla $n=2$, jak i $n\geq3$.
Ponadto
$\nabla_x Q_R(x,y)=-2y+2|y|^2x/R^2$.
Z reguły łańcuchowej otrzymujemy
\[
 \nabla_x H_R(x,y)=
 -\frac{|y|^2x/R^2-y}{\omega_n Q_R(x,y)^{n/2}},
 \qquad H_R(x,y)=h_n(\sqrt{Q_R(x,y)}).
\]
Wobec $\nabla\Phi_n(z)=-z/(\omega_n|z|^n)$
oraz $Q_R(\zeta,y)=|\zeta-y|^2$ mamy
\[
 \begin{aligned}
 -\nabla_xG_R(\zeta,y)\cdot\nu_\zeta
 &=\frac{(\zeta-y)\cdot\zeta
       -( |y|^2\zeta/R^2-y)\cdot\zeta}
       {\omega_n R|\zeta-y|^n}\\
 &=\frac{R^2-|y|^2}{\omega_n R|\zeta-y|^n}.
 \end{aligned}
\]
Rachunek obejmuje również $n=2$ oraz $y=0$,
bez przechodzenia do granicy w formule inwersyjnej.
Dla bieguna $y=0$ poprawka $H_R(x,0)$ jest stała,
a pochodna normalna samego rozwiązania fundamentalnego
daje $1/(\omega_n R^{n-1})$ po zmianie znaku.
Z symetrii funkcji Greena otrzymujemy
\[
 P_R(x,\zeta)=-\nabla_yG_R(x,\zeta)\cdot\nu_\zeta,
 \qquad x\in B_R,\quad \zeta\in\partial B_R,
\]
gdzie pochodna dotyczy drugiej zmiennej i jest obliczona
przez gładkie przedłużenie w otoczeniu brzegu.

Wykażmy teraz harmoniczność jądra względem $x$.
Ustalmy niepusty zwarty zbiór $K\subset B_R$ i
$\zeta\in\partial B_R$. Istnieje $\rho<R$ taka, że
$K\subset B_\rho$. Wybierzmy $\eta>0$ tak małe, aby
$\eta<R-\rho$ i $\rho(R+\eta)<R^2$.
Dla $x\in B_\rho$ i $y\in B_\eta(\zeta)$ mamy
$|x-y|\geq R-\eta-\rho>0$ i $|x||y|<R^2$.
Jednolity wzór na $G_R$ jest więc gładki na iloczynie
tych otwartych zbiorów. Dla $y\in B_\eta(\zeta)\cap B_R$
mamy $\Delta_xG_R(x,y)=0$ w $B_\rho$.
Różniczkując tę tożsamość względem $y$ w ustalonym
kierunku $\nu_\zeta$ i przechodząc do $y=\zeta$
z wnętrza kuli, dostajemy
\[
 \Delta_xP_R(x,\zeta)=
 -\nu_\zeta\cdot\nabla_y\Delta_xG_R(x,\zeta)=0.
\]
Komutacja pochodnych i przejście graniczne są
uzasadnione gładkością na wskazanym iloczynie.
Nie różniczkujemy tu ograniczenia $G_R(x,\cdot)$
do samej sfery, lecz funkcję określoną także
w kierunku normalnym.

Dla każdego wielowskaźnika $\alpha\in\mathbb N_0^n$
pochodna $D_x^\alpha P_R$ jest ciągła i ograniczona
na $\overline{B_\rho}\times\partial B_R$, ponieważ
$|x-\zeta|\geq R-\rho>0$. Funkcja $g$ jest ograniczona,
a sfera ma skończoną miarę. Zatem różniczkowanie
całki definiującej $u$ względem $x$ dowolną liczbę razy
jest dopuszczalne na $B_\rho$ na mocy twierdzenia
o zbieżności zdominowanej, stosowanego również
do ilorazów różnicowych. W szczególności
$\Delta u(x)=\int_{\partial B_R}\Delta_x P_R(x,\zeta)
g(\zeta)\,\mathrm d\sigma(\zeta)=0$.
Dowolność $\rho<R$ daje $u\in C^\infty(B_R)$
i harmoniczność w całej kuli.

Uzasadnijmy normalizację jądra bez powoływania się na
istnienie szukanego rozwiązania Dirichleta. Stosujemy
wzór reprezentacyjny~\eqref{eq:pde3-green-representation}
do funkcji stałej $1$, dla której $-\Delta1=0$.
Poprawka harmoniczna w funkcji Greena dla kuli
jest gładka na otoczeniu domkniętej kuli przy ustalonym
biegunie $x$, więc założenia wzoru są spełnione.
Wynika stąd $1=\int_{\partial B_R}P_R(x,\zeta)\,
\mathrm d\sigma(\zeta)$ dla każdego $x\in B_R$.
Dodatniość wynika z $R^2-|x|^2>0$, dodatniości
$\omega_n R$ oraz $|x-\zeta|>0$.
W szczególności całka Poissona jest średnią ważoną
danych brzegowych i $|u(x)|\leq\|g\|_\infty$.

Ustalmy $\zeta_0\in\partial B_R$ i $\eps>0$.
Z jednostajnej ciągłości $g$ wybierzmy $\delta>0$
tak, że $|g(\zeta)-g(\zeta_0)|<\eps$, gdy
$|\zeta-\zeta_0|<\delta$. Posługując się dodatniością
i normalizacją $P_R$ dostajemy
\[
 |u(x)-g(\zeta_0)|\leq\eps+
 2\|g\|_\infty
 \int_{\substack{\zeta\in\partial B_R\\
 |\zeta-\zeta_0|\geq\delta}}P_R(x,\zeta)\,
 \mathrm d\sigma(\zeta).
\]
Jeżeli $|x-\zeta_0|<\delta/2$, to na zbiorze
całkowania $|x-\zeta|\geq|\zeta-\zeta_0|
-|x-\zeta_0|\geq\delta/2$.
Zatem
\[
 \int_{|\zeta-\zeta_0|\geq\delta}P_R(x,\zeta)\,
 \mathrm d\sigma(\zeta)
 \leq \frac{(R^2-|x|^2)\omega_n R^{n-1}}
 {\omega_n R(\delta/2)^n}\longrightarrow0
 \quad(x\to\zeta_0,\ x\in B_R),
\]
gdzie całkujemy po odpowiedniej części $\partial B_R$.
Stąd $\limsup_{x\to\zeta_0,\ x\in B_R}
|u(x)-g(\zeta_0)|\leq\eps$. Dowolność $\eps>0$
daje żądaną granicę. Definiując $u=g$ na brzegu,
otrzymujemy funkcję ciągłą na $\overline{B_R}$,
gdyż $g$ jest ciągła na sferze, a granice od wnętrza
zgadzają się z jej wartościami.
Jeżeli $v$ jest innym rozwiązaniem w podanej klasie,
to $u-v$ jest harmoniczna w kuli i znika na brzegu.
Słaba zasada maksimum (twierdzenie~\ref{thm:laplace-max})
zastosowana do $u-v$ oraz $v-u$ daje $u=v$.
\end{proof}

\begin{przyklad}[Równanie Dirichleta na dysku]
Niech $n=2$, $R=1$, a $g(\cos\theta,\sin\theta)=\cos(3\theta)$.
Połóżmy $u(x_1,x_2)=x_1^3-3x_1x_2^2$.
Jest to wielomian, więc jest gładki także w środku dysku.
Mamy $\partial_{x_1}^2u=6x_1$ oraz
$\partial_{x_2}^2u=-6x_1$, skąd $\Delta u=0$.
Tożsamość $\cos(3\theta)=\cos^3\theta
-3\cos\theta\sin^2\theta$ daje
$u(r\cos\theta,r\sin\theta)=r^3\cos(3\theta)$.
Dla $r=1$ otrzymujemy dane $g$.
Z jednoznaczności w twierdzeniu~\ref{thm:pde3-poisson}
funkcja ta jest dokładnie całką Poissona z tych danych.
\end{przyklad}

\begin{stwierdzenie}[Nierówność Harnacka w kuli]
\label{prop:pde3-harnack}
Niech $n\geq2$, $R>0$, $u\in C^2(B_R)\cap C(\overline{B_R})$
będzie harmoniczna i nieujemna. Dla każdego $0\leq r<R$
i $|x|\leq r$ zachodzi
\begin{equation}\label{eq:pde3-harnack}
 \frac{R^{n-2}(R-r)}{(R+r)^{n-1}}u(0)
 \leq u(x)\leq
 \frac{R^{n-2}(R+r)}{(R-r)^{n-1}}u(0).
\end{equation}
\end{stwierdzenie}
\begin{proof}
Ciągłość $u$ na domkniętej kuli i nieujemność we wnętrzu
dają $u|_{\partial B_R}\geq0$. Twierdzenie
\ref{thm:pde3-poisson} i jednoznaczność rozwiązania
zapewniają reprezentację $u$ jako całki Poissona
z własnych wartości brzegowych. Możemy więc porównać
jądro w punkcie $x$ i w punkcie $0$.
Ponieważ $P_R(0,\zeta)=1/(\omega_n R^{n-1})$,
\[
 \frac{P_R(x,\zeta)}{P_R(0,\zeta)}
 =\frac{R^{n-2}(R^2-|x|^2)}{|x-\zeta|^n}.
\]
Używając $R-|x|\leq|x-\zeta|\leq R+|x|$,
otrzymujemy
\[
 \frac{R^{n-2}(R-|x|)}{(R+|x|)^{n-1}}
 \leq\frac{P_R(x,\zeta)}{P_R(0,\zeta)}
 \leq\frac{R^{n-2}(R+|x|)}{(R-|x|)^{n-1}}.
\]
Dla $0\leq s<R$ funkcja
$(R-s)/(R+s)^{n-1}$ jest malejąca, ponieważ jej
dodatni licznik maleje, a dodatni mianownik rośnie.
Analogicznie $(R+s)/(R-s)^{n-1}$ jest rosnąca.
Podstawiając $s=|x|\leq r$, otrzymujemy oszacowania
niezależne od $\zeta$. Mnożymy nierówności przez
nieujemne $P_R(0,\zeta)u(\zeta)$ i całkujemy po
$\partial B_R$. Całka z $P_R(0,\zeta)u(\zeta)$
jest równa $u(0)$, a z $P_R(x,\zeta)u(\zeta)$
jest równa $u(x)$. Otrzymujemy
\eqref{eq:pde3-harnack}, również wtedy, gdy $u(0)=0$,
bez dzielenia przez wartość funkcji.
\end{proof}

\begin{uwaga}[Nierówność Harnacka dla kul wewnętrznych]
Niech $\Omega\subset\R^n$ będzie obszarem, a
$u\in C^2(\Omega)$ nieujemną funkcją harmoniczną.
Nierówność Harnacka obowiązuje po translacji
na każdej kuli o domknięciu zawartym w $\Omega$.
W szczególności dla $\overline{B_{2a}(z)}\subset\Omega$
i $w\in\overline{B_a(z)}$ wstawienie $R=2a$, $r=a$
do \eqref{eq:pde3-harnack} daje
\[
 c_n u(z)\leq u(w)\leq C_n u(z),\qquad
 c_n=\frac{2^{n-2}}{3^{n-1}},\quad C_n=3\cdot2^{n-2}.
\]
Wykażmy, że dla każdego niepustego zwartego
$K\subset\Omega$ istnieje stała $C_K\geq1$,
niezależna od $u$, taka, że
$\sup_K u\leq C_K\inf_K u$.

Ustalmy $z_0\in\Omega$. Dla każdego $z\in K$
istnieje droga ciągła $\gamma\colon[0,1]\to\Omega$
od $z_0$ do $z$, ponieważ obszar w $\R^n$ jest
spójny drogowo. Jej obraz jest zwarty; można więc
wybrać $a>0$ tak, aby
$\overline{B_{2a}(\gamma(t))}\subset\Omega$ dla
wszystkich $t\in[0,1]$.
Z jednostajnej ciągłości drogi istnieje podział
$0=t_0<\cdots<t_m=1$, dla którego
$|\gamma(t_{j+1})-\gamma(t_j)|<a$.
Wzdłuż tych punktów kolejne zastosowania oszacowania
dają $c_n^m u(z_0)\leq u(z)\leq C_n^m u(z_0)$.
Jedno dalsze zastosowanie daje dla $w\in B_a(z)$
nierówności
$c_n^{m+1}u(z_0)\leq u(w)\leq C_n^{m+1}u(z_0)$.

Kule $B_a(z)$ uzyskane w ten sposób pokrywają $K$;
ze zwartości wybieramy skończone podpokrycie.
Niech $N$ będzie największą liczbą $m+1$
wśród wybranych łańcuchów. Ponieważ
$0<c_n\leq1\leq C_n$, dla wszystkich $w\in K$
mamy $c_n^N u(z_0)\leq u(w)\leq C_n^N u(z_0)$.
Ostatecznie $\sup_Ku\leq (C_n/c_n)^N\inf_Ku$.
Konstrukcja stałej zależy wyłącznie od geometrii
wybranych kul i dróg, a nie od funkcji $u$.
Nie zakładamy ciągłego przedłużenia $u$ na
\emph{brzeg całego obszaru}.
\end{uwaga}

\section{Problem Poissona i bariery porównawcze}
\label{sec:pde3-poisson-equation}

\begin{twierdzenie}[Problem Poissona na kuli]
\label{thm:pde3-ball-poisson}
Niech $n\geq2$, $R>0$, $f\in C_c^\infty(B_R)$,
a $g\in C(\partial B_R)$. Istnieje dokładnie jedno
$u\in C^2(B_R)\cap C(\overline{B_R})$ takie, że
$-\Delta u=f$ w $B_R$ i $u=g$ na $\partial B_R$.
Ponadto $u\in C^\infty(B_R)$ i dla każdego $x\in B_R$
ma postać
\begin{equation}\label{eq:pde3-poisson-solution}
 u(x)=\int_{B_R}G_R(x,y)f(y)\,\mathrm d\lambda_n(y)
 +\int_{\partial B_R}P_R(x,\zeta)g(\zeta)
 \,\mathrm d\sigma(\zeta).
\end{equation}
\end{twierdzenie}
\begin{proof}
Niech $\widetilde f\colon\R^n\to\R$ oznacza przedłużenie
$f$ przez zero. Ponieważ $\supp f$ jest zwartym
podzbiorem $B_R$, funkcja $f$ znika w pewnym otoczeniu
$\partial B_R$ po wewnętrznej stronie sfery.
Stąd $\widetilde f\in C_c^\infty(\R^n)$.
Połóżmy $w=\Phi_n*\widetilde f$. Z twierdzenia
\ref{thm:pde3-newton} wynika, że $w\in C^\infty(\R^n)$
oraz $-\Delta w=\widetilde f$ na $\R^n$.

Dane $g-w|_{\partial B_R}$ są ciągłe na sferze.
Twierdzenie~\ref{thm:pde3-poisson} dostarcza funkcji
$h\in C^\infty(B_R)\cap C(\overline{B_R})$ harmonicznej
w kuli i równej $g-w$ na jej brzegu. Funkcja $u=w+h$
spełnia równanie $-\Delta u=f$ i warunek brzegowy $u=g$.
Jeśli $u_1,u_2\in C^2(B_R)\cap C(\overline{B_R})$
są dwoma rozwiązaniami, ich różnica jest harmoniczna
i znika na brzegu. Słaba zasada maksimum zastosowana
do tej różnicy i jej przeciwności daje $u_1=u_2$.

Wykażmy wzór~\eqref{eq:pde3-poisson-solution}.
Korzystamy z poprawki harmonicznej
$H_R(x,y)=h_n(\sqrt{Q_R(x,y)})$ skonstruowanej
w dowodzie stwierdzenia~\ref{prop:pde3-ball-green},
gdzie $Q_R(x,y)=R^2-2x\cdot y+|x|^2|y|^2/R^2$.
Dla $x\ne y$ mamy $G_R(x,y)=\Phi_n(x-y)-H_R(x,y)$.
Wybierzmy $0<r<R$ tak, aby $\supp\widetilde f\subset B_r$
(gdy $f=0$, można wybrać dowolne takie $r$), oraz
$\eta>0$ spełniające $(R+\eta)r<R^2$. Dla
$|x|\leq R+\eta$ i $|y|\leq r$ zachodzi
\[
 Q_R(x,y)\geq
 \left(R-\frac{(R+\eta)r}{R}\right)^2>0.
\]
Poprawka $H_R$ jest więc gładka na otoczeniu zwartego
zbioru $\overline{B_{R+\eta}}\times\overline{B_r}$.
Dla każdego wielowskaźnika $\alpha\in\mathbb N_0^n$
istnieje skończona stała
\[
 M_\alpha=\sup_{\substack{|x|\leq R+\eta\\ |y|\leq r}}
 |D_x^\alpha H_R(x,y)|.
\]
Definiujemy dla $x\in B_{R+\eta}$
\[
 h_f(x)=\int_{B_r}H_R(x,y)f(y)\,\mathrm d\lambda_n(y).
\]
Funkcja $M_\alpha|f(y)|$ jest całkowalna w $B_r$
i dominuje moduł odpowiedniej pochodnej funkcji
podcałkowej. Twierdzenie o zbieżności zdominowanej,
stosowane do ilorazów różnicowych i do samych pochodnych,
pozwala różniczkować pod całką dowolną liczbę razy.
W konsekwencji $h_f\in C^\infty(B_{R+\eta})$ oraz
\[
 D^\alpha h_f(x)=\int_{B_r}D_x^\alpha H_R(x,y)
 f(y)\,\mathrm d\lambda_n(y).
\]
Dla każdego $y\in B_r$ funkcja $H_R(\cdot,y)$
jest harmoniczna w $B_R$, również dla $y=0$.
Stąd $\Delta h_f=0$ w $B_R$.

Dla $x\in B_R$ obie funkcje podcałkowe
$\Phi_n(x-y)f(y)$ i $H_R(x,y)f(y)$ są całkowalne:
pierwsza dzięki lokalnej całkowalności rozwiązania
fundamentalnego i ograniczonemu nośnikowi danych,
druga dzięki powyższemu oszacowaniu z $\alpha=0$.
Punkt $y=x$ ma miarę Lebesgue'a równą zeru, więc
dowolna wartość funkcji Greena w tym punkcie nie
zmienia całki. Zatem funkcja
\[
 v(x)=w(x)-h_f(x)=
 \int_{B_R}G_R(x,y)f(y)\,\mathrm d\lambda_n(y)
\]
jest gładka w kuli i spełnia $-\Delta v=f$.
Ponadto $w$ i $h_f$ są ciągłe na domkniętej kuli.
Jeśli $x\in\partial B_R$, to $|x-y|\geq R-r>0$
dla $y\in\supp f$, a $H_R(x,y)=\Phi_n(x-y)$.
Wynika stąd bezpośrednio, że $w(x)=h_f(x)$,
czyli $v(x)=0$ na brzegu.

Niech $p_g$ będzie całką Poissona z danych $g$.
Funkcja $v+p_g$ jest gładka w kuli, ciągła na jej
domknięciu, spełnia $-\Delta(v+p_g)=f$ i jest równa
$g$ na brzegu. Z wykazanej już jednoznaczności
$u=v+p_g$, co daje~\eqref{eq:pde3-poisson-solution}.
\end{proof}

\begin{stwierdzenie}[Oszacowanie przez barierę kwadratową]
\label{prop:pde3-barrier}
Niech $n\geq2$, $R>0$ i niech $\Omega\subset B_R(0)$
będzie obszarem. Niech $f\in C(\Omega)$ będzie rzeczywista
i ograniczona, a $u\in C^2(\Omega)\cap C(\overline\Omega)$
będzie rzeczywistym rozwiązaniem równania $-\Delta u=f$
w $\Omega$. Wówczas
\begin{equation}\label{eq:pde3-barrier}
 \sup_{x\in\Omega}|u(x)|\leq
 \sup_{x\in\partial\Omega}|u(x)|+
 \frac{R^2}{2n}\sup_{x\in\Omega}|f(x)|.
\end{equation}
\end{stwierdzenie}
\begin{proof}
Oznaczmy $M=\sup_\Omega|f|$ oraz $A=\sup_{\partial\Omega}|u|$.
Liczba $M$ jest skończona z założenia, a $A$ jest
skończona dzięki ciągłości $u$ na zwartym zbiorze
$\overline\Omega$. Połóżmy $b(x)=(R^2-|x|^2)/(2n)$
dla $x\in\R^n$. Ponieważ
$\overline\Omega\subset\overline{B_R(0)}$,
mamy $0\leq b\leq R^2/(2n)$ na $\overline\Omega$.
Dla każdego $j=1,\ldots,n$ zachodzi
$\partial_{x_j}^2b=-1/n$, zatem $-\Delta b=1$.

Funkcja $w_+=u-Mb$ należy do
$C^2(\Omega)\cap C(\overline\Omega)$ i spełnia
$\Delta w_+=-f+M\geq0$. Na brzegu mamy
$w_+\leq u\leq A$, gdyż $M\geq0$ i $b\geq0$.
Ze słabej zasady maksimum, twierdzenie~\ref{thm:laplace-max},
otrzymujemy $w_+\leq A$ w $\Omega$.
Analogicznie $w_-=-u-Mb$ spełnia
$\Delta w_-=f+M\geq0$ oraz $w_-\leq A$ na brzegu.
Stąd $w_-\leq A$ w $\Omega$.
Obie nierówności dają punktowe oszacowanie
\[
 |u(x)|\leq A+Mb(x)
 =A+\frac{R^2-|x|^2}{2n}M
 \qquad(x\in\Omega).
\]
Biorąc supremum po $x\in\Omega$ i używając
$b(x)\leq R^2/(2n)$, dostajemy
\eqref{eq:pde3-barrier}. Nie wymagamy przy tym
gładkości brzegu obszaru.
\end{proof}

\begin{przyklad}[Wymuszenie stałe]
Niech $n\geq2$, $R>0$, $f\equiv1$ w $B_R(0)$,
a dane Dirichleta będą równe $0$.
Wtedy $u(x)=(R^2-|x|^2)/(2n)$ jest wielomianem
spełniającym $-\Delta u=1$ w kuli i $u=0$ na sferze.
Z zasady maksimum jest to jedyne rozwiązanie
w klasie $C^2(B_R)\cap C(\overline{B_R})$.
Mamy $\sup_{B_R}|u|=u(0)=R^2/(2n)$,
$\sup_{\partial B_R}|u|=0$ i $\sup_{B_R}|f|=1$.
Zatem w nierówności~\eqref{eq:pde3-barrier} zachodzi
równość; współczynnika $R^2/(2n)$ przy normie wymuszenia
nie można zmniejszyć w całej rozważanej klasie
obszarów i danych.
Zauważmy, że $f\equiv1$ w kuli nie ma w niej zwartego
nośnika, więc istnienia rozwiązania w tym przykładzie
nie wywodzimy z twierdzenia~\ref{thm:pde3-ball-poisson},
lecz sprawdzamy je bezpośrednio.
\end{przyklad}

\section{Harmoniczność, analityczność i funkcje holomorficzne}
\label{sec:pde3-analytic}

\begin{twierdzenie}[Analityczność funkcji harmonicznych]
\label{thm:pde3-analytic}
Niech $n\geq2$, $\Omega\subset\R^n$ będzie otwarta,
a $u\in C^2(\Omega)$ harmoniczna. Wtedy $u$ jest
rzeczywiście analityczna w $\Omega$, czyli w otoczeniu
każdego punktu jest sumą zbieżnego szeregu potęgowego.
W szczególności $u\in C^\infty(\Omega)$.
\end{twierdzenie}
\begin{proof}
Ustalmy $x_0\in\Omega$ i wybierzmy $R>0$ tak, aby
$\overline{B_R(x_0)}\subset\Omega$.
Z twierdzenia~\ref{thm:pde3-poisson}, zastosowanego
po translacji do kuli o środku $x_0$, i z jednoznaczności
rozwiązania zagadnienia Dirichleta mamy
\[
 u(x)=\int_{\partial B_R(x_0)}
 \frac{R^2-|x-x_0|^2}{\omega_n R|x-\zeta|^n}
 u(\zeta)\,\mathrm d\sigma(\zeta)
 \qquad(x\in B_R(x_0)).
\]
Oznaczmy $h=x-x_0$ i $\theta=(\zeta-x_0)/R\in S^{n-1}$.
Po zmianie zmiennych na sferze otrzymujemy
\[
 u(x_0+h)=\frac{1-|h|^2/R^2}{\omega_n}
 \int_{S^{n-1}}(1+q_\theta(h))^{-n/2}
 u(x_0+R\theta)\,\mathrm d\sigma(\theta),
\]
gdzie $q_\theta(h)=-2\theta\cdot h/R+|h|^2/R^2$.
Dla ustalonego $\theta$ i $|q_\theta(h)|<1$ stosujemy
rozwinięcie dwumianowe
\[
 (1+q_\theta(h))^{-n/2}
 =\sum_{k=0}^\infty a_k q_\theta(h)^k,\qquad
 a_0=1,\quad
 a_k=\frac{(-1)^k}{k!}\prod_{\ell=0}^{k-1}
 \left(\frac n2+\ell\right)\quad(k\geq1).
\]

Aby uzasadnić, że otrzymujemy rzeczywiście szereg
potęgowy w $n$ zmiennych, a nie tylko jednostajnie
zbieżny szereg wielomianów, oszacujemy również
sumę modułów ich współczynników z wagami.
Dla wielomianu $p(h)=\sum_\alpha p_\alpha h^\alpha$
i $\rho>0$ oznaczmy
\[
 L_\rho(p)=\sum_\alpha |p_\alpha|\rho^{|\alpha|},
 \qquad h^\alpha=\prod_{j=1}^n h_j^{\alpha_j},
 \quad |\alpha|=\sum_{j=1}^n\alpha_j,
\]
gdzie $\alpha\in\mathbb N_0^n$.
Z reguły mnożenia wielomianów i nierówności trójkąta
wynika $L_\rho(p_1p_2)\leq L_\rho(p_1)L_\rho(p_2)$.
Ponadto $|p(h)|\leq L_\rho(p)$, gdy $|h_j|\leq\rho$
dla każdego $j$. Wybierzmy $\rho=R/(4n)$.
Ponieważ $|\theta_j|\leq1$, mamy niezależnie od $\theta$
\[
 L_\rho(q_\theta)
 \leq\frac{2n\rho}{R}+\frac{n\rho^2}{R^2}
 =\frac12+\frac1{16n}\leq\frac9{16}=:c<1.
\]
Domknięty sześcian $|h_j|\leq\rho$ zawiera się
w $B_R(0)$, gdyż $|h|\leq\sqrt n\,\rho<R$.
Na tym sześcianie rozwinięcie dwumianowe jest poprawne.
Co więcej,
\[
 \sum_{k=0}^\infty |a_k|L_\rho(q_\theta^k)
 \leq\sum_{k=0}^\infty |a_k|c^k<\infty,
\]
ponieważ iloraz kolejnych wyrazów szeregu po prawej
dąży do $c<1$.

Niech $b_{k,\alpha}(\theta)$ oznacza współczynnik
przy $h^\alpha$ w wielomianie $q_\theta^k$.
Są to ciągłe funkcje $\theta$, a dla ustalonego
$\alpha$ jedynie skończenie wiele $k$ może dać
niezerowy współczynnik, ponieważ każdy jednomian
w $q_\theta^k$ ma stopień co najmniej $k$.
Połóżmy
\[
 d_\alpha=\frac1{\omega_n}
 \sum_{k=0}^{|\alpha|}a_k
 \int_{S^{n-1}} b_{k,\alpha}(\theta)u(x_0+R\theta)
 \,\mathrm d\sigma(\theta).
\]
Dla $M=\sup_{\partial B_R(x_0)}|u|<\infty$
nierówność trójkąta i twierdzenie Tonellego dla
nieujemnych składników dają
\[
 \sum_\alpha |d_\alpha|\rho^{|\alpha|}
 \leq M\sum_{k=0}^\infty |a_k|
 \sup_{\theta\in S^{n-1}}L_\rho(q_\theta^k)
 \leq M\sum_{k=0}^\infty |a_k|c^k<\infty.
\]
To samo oszacowanie uzasadnia całkowanie wyraz po
wyrazie i zmianę kolejności sumowania rozwinięcia
dwumianowego oraz jednomianów. Otrzymujemy
\[
 u(x_0+h)=\left(1-\frac{|h|^2}{R^2}\right)
 \sum_{\alpha\in\mathbb N_0^n}d_\alpha h^\alpha
 \qquad(|h_j|\leq\rho,\ j=1,\ldots,n).
\]
Mnożenie przez skończony wielomian zachowuje
bezwzględną zbieżność z wagami $\rho^{|\alpha|}$
na mocy tej samej nierówności dla współczynników.
W szczególności otrzymujemy szereg potęgowy
przedstawiający $u$ w otwartym otoczeniu $x_0$,
zbieżny bezwzględnie i jednostajnie na wskazanym
domkniętym sześcianie. Na każdym mniejszym sześcianie
można różniczkować ten szereg wyraz po wyrazie
dowolną liczbę razy: czynniki wielomianowe w
$|\alpha|$ pochodzące od różniczkowania są
kontrolowane przez geometryczny czynnik
$(\rho'/\rho)^{|\alpha|}$ dla $0<\rho'<\rho$.
Otrzymane współczynniki są zatem współczynnikami
Taylora funkcji $u$ w $x_0$.
Dowolność $x_0$ daje analityczność i gładkość $u$.
\end{proof}

\begin{stwierdzenie}[Harmoniczna sprzężona w wymiarze dwa]
\label{prop:pde3-conjugate}
Niech $\Omega\subset\R^2$ będzie obszarem jednospójnym,
a $u\in C^2(\Omega)$ rzeczywistą funkcją harmoniczną.
Istnieje rzeczywista funkcja harmoniczna $v\in C^2(\Omega)$,
jednoznaczna z dokładnością do stałej addytywnej, taka,
że $F(x+iy)=u(x,y)+iv(x,y)$ jest holomorficzna w $\Omega$
traktowanej jako podzbiór $\C$.
\end{stwierdzenie}
\begin{proof}
Rozważmy formę różniczkową
$\alpha=-u_y\,\mathrm dx+u_x\,\mathrm dy$ klasy $C^1$.
Korzystając z reguły różniczkowania zewnętrznego
i antysymetrii iloczynu zewnętrznego, obliczamy
\[
 \mathrm d\alpha
 =-(u_{yx}\,\mathrm dx+u_{yy}\,\mathrm dy)
       \wedge\mathrm dx
 +(u_{xx}\,\mathrm dx+u_{xy}\,\mathrm dy)
       \wedge\mathrm dy
 =\Delta u\,\mathrm dx\wedge\mathrm dy=0.
\]
Stosujemy twierdzenie o dokładności zamkniętych form
różniczkowych stopnia $1$ klasy $C^1$ na jednospójnym
obszarze: całka takiej formy wzdłuż drogi kawałkami
klasy $C^1$ zależy jedynie od jej końców.
Ustalmy $(x_*,y_*)\in\Omega$ i $c\in\R$.
Definiujemy $v(x,y)=c+\int_\gamma\alpha$, gdzie
$\gamma$ jest dowolną drogą kawałkami klasy $C^1$
w $\Omega$ od $(x_*,y_*)$ do $(x,y)$.
Jednospójność i zamkniętość formy zapewniają
niezależność definicji od wyboru drogi.

W otoczeniu dowolnego punktu $(x,y)\in\Omega$
możemy do wybranej drogi dołączyć krótki odcinek
równoległy do jednej z osi, zawarty w $\Omega$.
Z addytywności całki krzywoliniowej otrzymujemy
\[
 \begin{aligned}
 v(x+h,y)-v(x,y)&=-\int_0^h u_y(x+t,y)\,
                       \mathrm d\lambda_1(t),\\
 v(x,y+h)-v(x,y)&=\int_0^h u_x(x,y+t)\,
                       \mathrm d\lambda_1(t).
 \end{aligned}
\]
Dla $h<0$ całki rozumiemy z odwróconą orientacją.
Ciągłość pochodnych $u$ daje po podzieleniu przez $h$
i przejściu do granicy
$v_x=-u_y$ i $v_y=u_x$. Obie pochodne są klasy $C^1$,
zatem $v\in C^2(\Omega)$ i $\mathrm dv=\alpha$.
Równość pochodnych mieszanych $u$ daje
$\Delta v=-u_{yx}+u_{xy}=0$.

Równości $u_x=v_y$ i $u_y=-v_x$ są równaniami
Cauchy'ego--Riemanna. Przy regularności $C^1$
funkcji $F$ zapewniają jej holomorficzność.
Jeżeli $v_1,v_2$ są dwiema funkcjami sprzężonymi
do $u$, to z równań Cauchy'ego--Riemanna wynika
$\nabla(v_1-v_2)=0$. Różnica jest lokalnie stała,
a ponieważ $\Omega$ jest spójna, jest stała na
całym obszarze. Odwrotnie, dodanie dowolnej stałej
rzeczywistej do $v$ zachowuje wszystkie warunki.
\end{proof}

\begin{kontrprzyklad}[Przeszkoda topologiczna dla sprzężenia]
Na pierścieniu $\Omega=\{z\in\C:1<|z|<2\}$
połóżmy $u(x,y)=\frac12\log(x^2+y^2)$.
Oznaczając $s=x^2+y^2$, obliczamy
\[
 u_x=\frac{x}{s},\quad u_y=\frac{y}{s},\qquad
 u_{xx}=\frac{y^2-x^2}{s^2},\quad
 u_{yy}=\frac{x^2-y^2}{s^2}.
\]
Stąd $\Delta u=0$ na $\Omega$. Gdyby istniała
globalna funkcja sprzężona $v$ do $u$, równania
Cauchy'ego--Riemanna wymuszałyby
\[
 \mathrm dv=\alpha
 =-\frac{y}{x^2+y^2}\,\mathrm dx
 +\frac{x}{x^2+y^2}\,\mathrm dy.
\]
Niech $a=3/2$ i $\gamma(t)=(a\cos t,a\sin t)$
dla $0\leq t\leq2\pi$. Jest to zamknięta krzywa
w $\Omega$, zorientowana przeciwnie do ruchu
wskazówek zegara. Mamy
\[
 \int_\gamma\alpha
 =\int_0^{2\pi}
 \frac{a^2\sin^2t+a^2\cos^2t}{a^2}
 \,\mathrm d\lambda_1(t)=2\pi.
\]
Tymczasem dla formy dokładnej $\mathrm dv$
całka wzdłuż $\gamma$ byłaby równa
$v(\gamma(2\pi))-v(\gamma(0))=0$.
Sprzeczność dowodzi, że nie istnieje globalna
jednowartościowa funkcja sprzężona do $u$.
Każdy punkt pierścienia ma jednak otoczenie
będące kulą o domknięciu zawartym w $\Omega$;
na takiej kuli istnieje funkcja sprzężona na mocy
stwierdzenia~\ref{prop:pde3-conjugate}.
\end{kontrprzyklad}

\section{Wygładzanie harmoniczne na dysku: rozwiązanie przez szeregi Fouriera}
\label{sec:pde3-fourier-disk}

Jawna postać jądra Poissona ma szczególnie przejrzystą
interpretację na jednostkowym dysku. Niech
$g\in C^2(S^1)$ będzie funkcją rzeczywistą; oznacza to,
że funkcja $G(\theta)=g(e^{i\theta})$ jest klasy $C^2$
i $2\pi$-okresowa na $\R$. Ustalamy normalizację
współczynników Fouriera:
\[
 \begin{aligned}
 a_0&=\frac1{2\pi}\int_0^{2\pi}G(\theta)\,
             \mathrm d\lambda_1(\theta),\\
 a_k&=\frac1\pi\int_0^{2\pi}G(\theta)\cos(k\theta)\,
             \mathrm d\lambda_1(\theta),\\
 b_k&=\frac1\pi\int_0^{2\pi}G(\theta)\sin(k\theta)\,
             \mathrm d\lambda_1(\theta)\qquad(k\geq1).
 \end{aligned}
\]
Dwukrotne całkowanie przez części, z wykorzystaniem
okresowości $G$, $G'$ oraz funkcji trygonometrycznych,
daje
\[
 \begin{aligned}
 a_k&=-\frac1{\pi k^2}\int_0^{2\pi}G''(\theta)
                \cos(k\theta)\,\mathrm d\lambda_1(\theta),\\
 b_k&=-\frac1{\pi k^2}\int_0^{2\pi}G''(\theta)
                \sin(k\theta)\,\mathrm d\lambda_1(\theta).
 \end{aligned}
\]
Stąd $|a_k|+|b_k|\leq4\|G''\|_\infty/k^2$.
W szczególności szereg współczynników bezwzględnych
jest zbieżny. Poniższy argument z jądrem Poissona
pokaże również, że szereg trygonometryczny na brzegu
przedstawia dokładnie funkcję $G$, bez konieczności
przywoływania osobnego twierdzenia o zbieżności
szeregów Fouriera.

\begin{stwierdzenie}[Rozwinięcie harmoniczne na dysku]
\label{prop:pde3-fourier-disk}
Niech $g\in C^2(S^1)$ będzie rzeczywista, a
$a_0,a_k,b_k$ jej współczynnikami Fouriera z powyższą
normalizacją. Jedynym rozwiązaniem zagadnienia
$\Delta u=0$ w $B_1(0)\subset\R^2$ oraz $u|_{S^1}=g$
w klasie $C^2(B_1)\cap C(\overline{B_1})$
jest funkcja
\begin{equation}\label{eq:pde3-fourier-disk}
 u(r\cos\theta,r\sin\theta)=a_0+
 \sum_{k=1}^\infty r^k
 (a_k\cos k\theta+b_k\sin k\theta),\quad 0\leq r<1.
\end{equation}
Szereg jest bezwzględnie i jednostajnie zbieżny
na domkniętym dysku, a wszystkie jego pochodne
współrzędnych kartezjańskich można obliczać wyraz
po wyrazie na zwartych podzbiorach dysku.
\end{stwierdzenie}
\begin{proof}
Dla $0\leq r<1$ i $\psi\in\R$ szereg geometryczny
o ilorazie $re^{i\psi}$ daje
\[
 1+2\sum_{k=1}^\infty r^k\cos(k\psi)
 =\operatorname{Re}\frac{1+re^{i\psi}}{1-re^{i\psi}}
 =\frac{1-r^2}{1-2r\cos\psi+r^2}.
\]
Ostatnią równość otrzymujemy, mnożąc licznik
i mianownik ułamka zespolonego przez
$1-re^{-i\psi}$ i biorąc część rzeczywistą.
Dla ustalonego $r<1$ szereg jest jednostajnie
zbieżny względem $\psi$, gdyż $|\cos(k\psi)|\leq1$.
Dla $x=re^{i\theta}$, $\zeta=e^{it}$ mamy
$|x-\zeta|^2=1-2r\cos(\theta-t)+r^2$.
Ponieważ $\omega_2=2\pi$ oraz parametryzacja
$t\mapsto e^{it}$ ma jednostkową prędkość,
wzór na jądro Poissona z twierdzenia
\ref{thm:pde3-poisson} przyjmuje postać
\[
 P_1(re^{i\theta},e^{it})
 =\frac1{2\pi}\left(1+2\sum_{k=1}^\infty
                       r^k\cos(k(\theta-t))\right).
\]
Możemy całkować szereg wyraz po wyrazie względem
$t\in[0,2\pi]$ po pomnożeniu przez ograniczoną
funkcję $G(t)$. Korzystając z tożsamości
$\cos(k(\theta-t))=\cos(k\theta)\cos(kt)
+\sin(k\theta)\sin(kt)$, otrzymujemy dokładnie
\eqref{eq:pde3-fourier-disk}.
Jest to całka Poissona z danych $g$, a zatem
z przywołanego twierdzenia jest harmoniczna,
ciągła na domkniętym dysku i jedyna w podanej klasie.

Uzasadnijmy również bezpośrednio własności szeregu.
Każdy jego składnik jest wielomianem w zmiennych
$x,y$, ponieważ dla $z=x+iy$ zachodzi
$r^k\cos(k\theta)=\operatorname{Re}z^k$ oraz
$r^k\sin(k\theta)=\operatorname{Im}z^k$.
Dla nieujemnych liczb całkowitych $p,q$ i $m=p+q$
mamy dla $k\geq m$
\[
 \partial_x^p\partial_y^q z^k
 =i^q k(k-1)\cdots(k-m+1)z^{k-m},
\]
przy czym dla $m=0$ iloczyn jest równy $1$,
a dla $k<m$ pochodna jest zerowa.
W szczególności dla $m=2$ suma drugich pochodnych
względem $x$ i $y$ znika, co dowodzi harmoniczności
każdego składnika, także w środku dysku.

Dla $0<\rho<1$, $|z|\leq\rho$ i $k\geq\max(1,m)$
moduł pochodnej rzędu $m$ składnika szeregu
jest nie większy niż
$(|a_k|+|b_k|)k^m\rho^{k-m}$.
Szereg tych majorant jest zbieżny, ponieważ
$\sum_{k\geq\max(1,m)}k^{m-2}\rho^{k-m}<\infty$.
Z kryterium Weierstrassa otrzymujemy jednostajną
zbieżność szeregu każdej pochodnej na
$\overline{B_\rho}$. Twierdzenie o różniczkowaniu
szeregów funkcyjnych, stosowane kolejno do
pochodnych na mniejszych otwartych kulach, daje
regularność $C^\infty$ oraz możliwość obliczenia
laplasjanu wyraz po wyrazie w całym dysku.
Argument obejmuje również punkt $z=0$.

Na domkniętym dysku sam szereg jest jednostajnie
i bezwzględnie zbieżny dzięki oszacowaniu
$\sum_k(|a_k|+|b_k|)<\infty$.
Jego suma jest więc ciągła aż do brzegu.
Dla każdego $\theta$ przechodząc z $r$ do $1$
od dołu, otrzymujemy
\[
 G(\theta)=a_0+\sum_{k=1}^\infty
             (a_k\cos k\theta+b_k\sin k\theta),
\]
ponieważ granica całki Poissona jest równa
$G(\theta)$, a jednostajna zbieżność pozwala
przejść do granicy w szeregu. W szczególności
potwierdziliśmy, że suma szeregu na brzegu
jest dokładnie zadaną funkcją, a nie tylko
funkcją o tych samych formalnych współczynnikach.
\end{proof}

Współczynnik wysokiej częstości $k$ zostaje
na okręgu o promieniu $r$ pomnożony przez $r^k$.
Dla $0<r<1$ mamy $r^k=e^{-k|\log r|}$, co daje
wykładnicze tłumienie wysokich częstości.
Można zatem precyzyjnie opisać niestabilność
odtwarzania danych brzegowych.

Ustalmy $0<r<1$ i niech $T_r$ przypisuje ciągłym
danym $g$ profil $\theta\mapsto u(re^{i\theta})$,
gdzie $u$ jest całką Poissona. Jest to liniowe
odwzorowanie przestrzeni $C(S^1)$ w $C(S^1)$,
w której normą jest supremum modułu.
Jest ono różnowartościowe: jeżeli profil jest
zerowy, słaba zasada maksimum daje $u=0$
na $B_r$, a analityczność i zasada identyczności
z dowodu twierdzenia~\ref{thm:pde3-strong-max}
dają $u=0$ na $B_1$ i w konsekwencji $g=0$.

Dla $g_k(e^{i\theta})=\cos(k\theta)$, $k\geq1$,
mamy $\|g_k\|_\infty=1$ i na mocy
\eqref{eq:pde3-fourier-disk}
$T_rg_k(e^{i\theta})=r^k\cos(k\theta)$.
Zatem $\|T_rg_k\|_\infty=r^k\to0$,
podczas gdy $\|g_k\|_\infty=1$.
Odwrotność $T_r^{-1}$ określona na obrazie $T_r$
nie jest więc ciągła w zerze względem normy
supremum, nawet dla ciągu gładkich danych.
To odpowiednik niestabilności Hadamarda
z rozdziału~\ref{chap:pde-intro}.

\begin{center}
\begin{tikzpicture}[>=Stealth,scale=1]
  \draw[->] (0,0)--(5.6,0) node[right] {$k$};
  \draw[->] (0,0)--(0,3.3) node[above] {$r^k$};
  \draw[densely dotted] (0,2.9)--(5.25,2.9);
  \node[left] at (0,2.9) {$1$};
  \draw[blue!70!black,thick,domain=0:5.1,samples=100]
    plot (\x,{2.9*exp(-.25*\x)});
  \draw[red!70!black,thick,domain=0:5.1,samples=100]
    plot (\x,{2.9*exp(-.9*\x)});
  \node[blue!70!black,anchor=west] at (3.5,1.5) {$r=e^{-1/4}$};
  \node[red!70!black,anchor=west] at (3.05,.70) {$r=e^{-9/10}$};
  \node[below] at (2.8,0) {częstość $k$};
\end{tikzpicture}
\par\smallskip\small
Tłumienie modów Fouriera w rozwiązaniu harmonicznym
na dysku: im dalej od brzegu, tym szybsze zanikanie
wysokich częstości.
\end{center}

\section{Zagadnienia brzegowe na bardziej ogólnych obszarach}
\label{sec:pde3-general-domain}

Jawny wzór Poissona jest szczególnym przywilejem kuli.
Dla ogólnego obszaru funkcja Greena może nie mieć
prostej postaci, a istnienie ciągłego rozwiązania Dirichleta
zależy od geometrii brzegu. Aby rozdzielić twierdzenia
udowodnione od materiału wymagającego późniejszych narzędzi,
formułujemy poniższy wynik jako zapowiedź bez dowodu.

\begin{uwaga}[Zapowiedź twierdzenia o ogólnych obszarach]
\label{prop:pde3-future}
Niech $n\geq2$ i niech $\Omega\subset\R^n$ będzie obszarem ograniczonym o brzegu klasy $C^2$.
Wtedy dla każdego $g\in C(\partial\Omega)$ istnieje
dokładnie jedna funkcja
$u\in C^2(\Omega)\cap C(\overline\Omega)$ spełniająca
$\Delta u=0$ w $\Omega$ i $u|_{\partial\Omega}=g$.
Jest to klasyczny wynik teorii potencjału. Nie traktujemy go
tu jako wyniku udowodnionego: dowód wymaga metody Perrona
i barier brzegowych lub odpowiednich twierdzeń o regularności
eliptycznej, które wprowadzimy w dalszej części skryptu.
\end{uwaga}

Warto odróżnić dwie kwestie. \emph{Jednoznaczność}
na każdym obszarze ograniczonym wynika z zasady maksimum,
niezależnie od regularności jej brzegu. Natomiast
\emph{istnienie} dla wszystkich ciągłych danych brzegowych
nie jest prawdziwe na każdym obszarze: punkt brzegowy
może być nieregularny w sensie teorii potencjału.
To zapowiedź kryterium Wienera i pojęcia pojemności,
które pojawi się w dalszej części skryptu.

\section{Podsumowanie i mapa dalszych związków}
\label{sec:pde3-summary}

Poznaliśmy dwa wzajemnie uzupełniające się spojrzenia
na równanie Laplace'a. Własność wartości średniej,
zasady maksimum, nierówność Harnacka i analityczność
ujawniają \emph{sztywność} funkcji harmonicznych.
Z kolei rozwiązanie fundamentalne, potencjał Newtona,
funkcja Greena i jądro Poissona dają \emph{konstrukcje}
rozwiązań, a nie tylko ich jakościowe własności.

Zasadnicze pytanie pozostaje otwarte w pełnej ogólności:
jak rozwiązywać $-\Delta u=f$ na obszarze, dla której
nie znamy jawnego jądra? Odpowiedź dadzą później
przestrzenie Sobolewa, metoda wariacyjna i teoria
regularności eliptycznej. W następnym rozdziale
porównamy te własności z równaniem ciepła i jego
charakterystycznym mechanizmem wygładzania.

% Rozdział 4: równanie ciepła.

\chapter[Równanie ciepła]{Równanie ciepła: jądro Gaussa, maksimum, wygładzanie i ewolucja}
\label{chap:pde-heat}

W poprzednim rozdziale warunek $\Delta u=0$ wymuszał silną
sztywność funkcji. Teraz badamy równanie, w którym laplasjan
steruje zmianą w czasie:
\begin{equation}\label{eq:pde4-heat}
 \partial_tu(t,x)-\kappa\Delta_xu(t,x)=f(t,x),\qquad
 t>0,\quad x\in\R^n,\quad \kappa>0.
\end{equation}
Gdy $f=0$, równanie opisuje swobodną dyfuzję. Fundamentalne pytania
są odmienne niż dla równania Laplace'a: jak odzyskać rozwiązanie
z profilu początkowego, jakie informacje zostają zachowane i dlaczego
każdy dodatni czas poprawia regularność? Wyprowadzimy rozwiązanie
całego zagadnienia Cauchy'ego, zasady maksimum, formułę Duhamela,
bilanse energii i przykłady pokazujące, że bez określenia klasy
rozwiązań nawet jednoznaczność może zawieść.

W całym rozdziale $n\geq1$ jest liczbą całkowitą,
$\kappa>0$, a funkcje skalarne są rzeczywiste, o ile
nie zaznaczono inaczej. Przez $\lambda_k$ oznaczamy
miarę Lebesgue'a w $\R^k$ dla $k\geq1$.
Dla $1\leq p<\infty$ symbole $L^p(\R^n)$ i $\|\cdot\|_p$
oznaczają przestrzeń klas funkcji mierzalnych
o całkowalnej $p$-tej potędze modułu oraz jej normę.
Przestrzeń $L^\infty(\R^n)$ składa się z klas funkcji
istotnie ograniczonych, a jej normą jest istotne
supremum modułu. Przez $\mathrm{BUC}(\R^n)$ oznaczamy
przestrzeń funkcji ograniczonych i jednostajnie
ciągłych z normą supremum. W dalszych twierdzeniach
\emph{silna ciągłość} oznacza zbieżność w normie danej
przestrzeni, a nie tylko punktową ani słabą. Będziemy korzystać z elementarnych
twierdzeń o całkowaniu: Fubiniego--Tonellego, zbieżności
zdominowanej, nierówności Minkowskiego i ciągłości przesunięć
w $L^p$; w rozdziałach o analizie funkcjonalnej wrócimy
do ich szerszego kontekstu.

\section{Jądro Gaussa: skąd bierze się rozwiązanie?}
\label{sec:pde4-gauss}

\begin{definicja}[Jądro cieplne]
\label{def:pde4-kernel}
Niech $n\geq1$ i $\kappa>0$. Dla $t>0$ oraz $x\in\R^n$
definiujemy
\begin{equation}\label{eq:pde4-kernel}
 \Gamma_\kappa(t,x)=\frac{1}{(4\pi\kappa t)^{n/2}}
 \exp\!\left(-\frac{|x|^2}{4\kappa t}\right).
\end{equation}
Zapis $\Gamma_\kappa(t)*g$ oznacza splot w zmiennej przestrzennej:
$(\Gamma_\kappa(t)*g)(x)=\int_{\R^n}\Gamma_\kappa(t,x-y)g(y)
\,\mathrm d\lambda_n(y)$, jeżeli całka ma sens.
\end{definicja}

\begin{twierdzenie}[Równanie, normalizacja i skalowanie]
\label{thm:pde4-gauss-basic}
Dla każdego $n\geq1$, $\kappa>0$, $t>0$ zachodzi:
$\Gamma_\kappa(t,x)>0$,
$(\partial_t-\kappa\Delta_x)\Gamma_\kappa(t,x)=0$
dla $x\in\R^n$, oraz
\begin{equation}\label{eq:pde4-normalize}
 \int_{\R^n}\Gamma_\kappa(t,x)\,\mathrm d\lambda_n(x)=1.
\end{equation}
Ponadto $\Gamma_\kappa(t,x)=t^{-n/2}\Gamma_\kappa(1,x/\sqrt t)$.
\end{twierdzenie}
\begin{proof}
Ustalmy $n\geq1$, $\kappa>0$ i oznaczmy
$\Gamma=\Gamma_\kappa$. Funkcja $\Gamma$ jest
dodatnia i klasy $C^\infty$ na $(0,\infty)\times\R^n$.
Różniczkując oba czynniki w jej definicji, otrzymujemy
\[
 \partial_t\Gamma(t,x)=
 \left(-\frac n{2t}+\frac{|x|^2}{4\kappa t^2}\right)
 \Gamma(t,x),\qquad
 \partial_{x_j}\Gamma(t,x)=-\frac{x_j}{2\kappa t}\Gamma(t,x).
\]
Ponowne różniczkowanie względem $x_j$ daje
\[
 \partial_{x_j}^2\Gamma(t,x)
 =\left(\frac{x_j^2}{4\kappa^2t^2}
 -\frac1{2\kappa t}\right)\Gamma(t,x).
\]
Sumując względem $j=1,\ldots,n$ i mnożąc przez
$\kappa$, dostajemy $\kappa\Delta_x\Gamma=\partial_t\Gamma$.

Dla normalizacji oznaczmy
$I=\int_\R e^{-z^2}\,\mathrm d\lambda_1(z)$.
Całka jest skończona, ponieważ dla $|z|\geq1$
mamy $e^{-z^2}\leq e^{-|z|}$.
Z twierdzenia Tonellego i zmiany zmiennych
we współrzędnych biegunowych w $\R^2$ wynika
\[
 I^2=\int_{\R^2}e^{-|z|^2}\,\mathrm d\lambda_2(z)
 =2\pi\int_0^\infty e^{-r^2}r\,\mathrm d\lambda_1(r)
 =\pi.
\]
Ponieważ $I>0$, mamy $I=\sqrt\pi$.
Ponowne zastosowanie twierdzenia Tonellego
do iloczynu $n$ funkcji jednowymiarowych daje
$\int_{\R^n}e^{-|z|^2}\,\mathrm d\lambda_n(z)=\pi^{n/2}$.
Podstawienie $x=\sqrt{4\kappa t}\,z$, którego
jakobian ma moduł $(4\kappa t)^{n/2}$, daje
\[
 \int_{\R^n}\Gamma_\kappa(t,x)\,\mathrm d\lambda_n(x)
 =(4\pi\kappa t)^{-n/2}(4\kappa t)^{n/2}\pi^{n/2}
 =1.
\]
Wreszcie wstawienie $x/\sqrt t$ do wzoru na
$\Gamma_\kappa(1,\cdot)$ daje dokładnie
$t^{n/2}\Gamma_\kappa(t,x)$.
\end{proof}

\begin{uwaga}[Skalowanie paraboliczne]
Niech $u$ będzie klasy $C^{1,2}$ na
$(0,\infty)\times\R^n$ (jedna ciągła pochodna
czasowa i ciągłe pochodne przestrzenne do rzędu
drugiego) i niech $\partial_tu-\kappa\Delta_xu=0$.
Dla $\rho>0$ połóżmy $u_\rho(t,x)=u(\rho^2t,\rho x)$.
Reguła łańcuchowa daje
\[
 (\partial_t-\kappa\Delta_x)u_\rho(t,x)
 =\rho^2(\partial_tu-\kappa\Delta_xu)(\rho^2t,\rho x)=0.
\]
Czas skaluje się zatem jak kwadrat odległości,
w odróżnieniu od skali falowej $t\sim|x|$.

Skalowanie sugeruje poszukiwanie rozwiązania o stałej
masie w postaci $v(t,x)=t^{-n/2}F(x/\sqrt t)$.
Dla $z=x/\sqrt t$ rachunek daje
\[
 (\partial_t-\kappa\Delta_x)v(t,x)
 =-t^{-n/2-1}
 \left(\kappa\Delta F(z)+\tfrac12 z\cdot\nabla F(z)
 +\tfrac n2 F(z)\right).
\]
Wyrażenie w nawiasie jest dywergencją pola
$\kappa\nabla F(z)+zF(z)/2$.
Wystarczy więc znaleźć profil spełniający
$\kappa\nabla F(z)+zF(z)/2=0$.
Dla $F\in C^1(\R^n)$ jest to równoważne
$\nabla(e^{|z|^2/(4\kappa)}F(z))=0$, czyli
$F(z)=C e^{-|z|^2/(4\kappa)}$.
Warunek jednostkowej masy daje
$C=(4\pi\kappa)^{-n/2}$.
Otrzymujemy jądro Gaussa. Nie twierdzimy, że każde
rozwiązanie równania dla profilu musi spełniać ten
silniejszy warunek zerowania pola; jest on tutaj
wystarczającą metodą konstrukcji szukanego jądra.
\end{uwaga}

\begin{center}
\begin{tikzpicture}[>=Stealth,x=1.15cm,y=3.8cm]
 \draw[->] (-3.2,0)--(3.55,0) node[right] {$x$};
 \draw[->] (0,0)--(0,1.04) node[above] {$\Gamma_1(t,x)$};
 \draw[blue!70!black,thick,domain=-3:3,samples=110]
   plot (\x,{1/sqrt(4*pi*.12)*exp(-\x*\x/(4*.12))});
 \draw[red!70!black,thick,domain=-3:3,samples=110]
   plot (\x,{1/sqrt(4*pi*.45)*exp(-\x*\x/(4*.45))});
 \draw[green!45!black,thick,domain=-3:3,samples=110]
   plot (\x,{1/sqrt(4*pi*1.1)*exp(-\x*\x/(4*1.1))});
 \node[blue!70!black,anchor=west] at (1.5,.79) {$t=0{,}12$};
 \node[red!70!black,anchor=west] at (1.5,.68) {$t=0{,}45$};
 \node[green!45!black,anchor=west] at (1.5,.57) {$t=1{,}1$};
\end{tikzpicture}
\par\smallskip\small
Jądro cieplne w wymiarze jeden ($\kappa=1$): profil
staje się niższy i szerszy, lecz całkowita masa pozostaje równa $1$.
\end{center}

\section{Własność półgrupy i przybliżona jedność}
\label{sec:pde4-semigroup}

\begin{lemat}[Składanie jąder Gaussa]
\label{lem:pde4-convolution}
Dla $n\geq1$, $\kappa>0$ oraz $s,t>0$ zachodzi
$\Gamma_\kappa(t)*\Gamma_\kappa(s)=\Gamma_\kappa(t+s)$.
\end{lemat}
\begin{proof}
Ustalmy $x\in\R^n$. Całka definiująca splot
jest skończona, ponieważ jedno jądro jest ograniczone,
a drugie ma całkę $1$. Rozwinięcie kwadratów daje
\[
 \frac{|x-y|^2}{t}+\frac{|y|^2}{s}
 =\frac{t+s}{ts}\left|y-\frac{s}{t+s}x\right|^2
 +\frac{|x|^2}{t+s}.
\]
Oznaczmy $a=ts/(t+s)>0$. Po podstawieniu
$z=y-sx/(t+s)$ otrzymujemy
\[
 \begin{aligned}
 (\Gamma_\kappa(t)*\Gamma_\kappa(s))(x)
 &=\frac{e^{-|x|^2/(4\kappa(t+s))}}
 {(4\pi\kappa)^n(ts)^{n/2}}
 \int_{\R^n}e^{-|z|^2/(4\kappa a)}
 \,\mathrm d\lambda_n(z)\\
 &=\frac{(4\pi\kappa a)^{n/2}}
 {(4\pi\kappa)^n(ts)^{n/2}}
 e^{-|x|^2/(4\kappa(t+s))}\\
 &=\frac{e^{-|x|^2/(4\kappa(t+s))}}
 {(4\pi\kappa(t+s))^{n/2}}
 =\Gamma_\kappa(t+s,x).
 \end{aligned}
\]
W drugiej równości wykorzystaliśmy normalizację
jądra w chwili $a$, a w trzeciej tożsamość
$a/(ts)=1/(t+s)$.
\end{proof}

\begin{twierdzenie}[Przybliżona jedność]
\label{thm:pde4-approx-identity}
Niech $n\geq1$, $\kappa>0$. Dla $g\in\mathrm{BUC}(\R^n)$
funkcje $P_tg:=\Gamma_\kappa(t)*g$ spełniają
$\|P_tg-g\|_\infty\to0$ przy $t\downarrow0$.
Dla $1\leq p<\infty$ i $g\in L^p(\R^n)$ zachodzi
$\|P_tg-g\|_p\to0$ przy $t\downarrow0$.
Dla każdej $g\in L^p$ ($1\leq p\leq\infty$),
$\|P_tg\|_p\leq\|g\|_p$ dla $t>0$.
\end{twierdzenie}
\begin{proof}
Dla ustalonego $t>0$ jądro $\Gamma_\kappa(t,\cdot)$
należy do każdego $L^q(\R^n)$, $1\leq q\leq\infty$:
jest ograniczone, całkowalne, a dla $q<\infty$
można oszacować jego $q$-tą potęgę przez
$\|\Gamma_\kappa(t,\cdot)\|_\infty^{q-1}$
razy samo jądro. Dla $g\in L^p(\R^n)$
nierówność Höldera z wykładnikiem sprzężonym $q$
($1/p+1/q=1$) daje więc bezwzględną zbieżność
całki definiującej $P_tg(x)$ dla każdego $x$.
Wynik nie zależy od wyboru reprezentanta klasy $g$,
ponieważ translacja zbioru miary zero ma miarę zero.

Potrzebujemy najpierw oszacowania ogona jądra.
Dla każdego $\delta>0$, korzystając ze skalowania
z twierdzenia~\ref{thm:pde4-gauss-basic}, mamy
\[
 \int_{|y|\geq\delta}\Gamma_\kappa(t,y)\,
 \mathrm d\lambda_n(y)
 =\int_{|z|\geq\delta/\sqrt t}\Gamma_\kappa(1,z)\,
 \mathrm d\lambda_n(z)\longrightarrow0
 \quad(t\downarrow0).
\]
Istotnie, funkcje charakterystyczne tych zbiorów
dążą punktowo do zera, a funkcją dominującą
jest całkowalne jądro $\Gamma_\kappa(1,\cdot)$.

Dla $g\in\mathrm{BUC}(\R^n)$ normalizacja daje
\[
 |P_tg(x)-g(x)|\leq
 \int_{\R^n}\Gamma_\kappa(t,y)|g(x-y)-g(x)|
 \,\mathrm d\lambda_n(y).
\]
Ustalmy $\eps>0$ i wybierzmy $\delta>0$ tak,
aby $|g(x-y)-g(x)|<\eps$ dla $|y|<\delta$
i wszystkich $x\in\R^n$. Dzieląc całkę na
$|y|<\delta$ oraz $|y|\geq\delta$, otrzymujemy
\[
 \|P_tg-g\|_\infty\leq\eps+
 2\|g\|_\infty\int_{|y|\geq\delta}\Gamma_\kappa(t,y)
 \,\mathrm d\lambda_n(y).
\]
Granica górna przy $t\downarrow0$ nie przekracza
$\eps$, a dowolność $\eps$ daje żądaną zbieżność.

Niech teraz $1\leq p<\infty$. Przesunięcie
$\tau_yg=g(\cdot-y)$ jest izometrią $L^p(\R^n)$.
Dla $g\in C_c^\infty(\R^n)$ zbieżność
$\|\tau_yg-g\|_p\to0$ przy $y\to0$ wynika
z jednostajnej ciągłości i wspólnego ograniczonego
nośnika dla $|y|\leq1$.
Dla dowolnego $g\in L^p$ wybieramy
$\phi\in C_c^\infty(\R^n)$ z $\|g-\phi\|_p<\eps$,
korzystając z gęstości tej klasy, i otrzymujemy
$\|\tau_yg-g\|_p\leq2\eps+\|\tau_y\phi-\phi\|_p$.
Zatem przesunięcia są ciągłe w normie $L^p$.

Nierówność całkowa Minkowskiego i normalizacja
jądra dają
\[
 \begin{aligned}
 \|P_tg\|_p&\leq\int_{\R^n}\Gamma_\kappa(t,y)
                  \|\tau_yg\|_p\,\mathrm d\lambda_n(y)=\|g\|_p,\\
 \|P_tg-g\|_p&\leq\int_{\R^n}\Gamma_\kappa(t,y)
                  \|\tau_yg-g\|_p\,\mathrm d\lambda_n(y).
 \end{aligned}
\]
Prawe strony są skończone, ponieważ
$\|\tau_yg-g\|_p\leq2\|g\|_p$.
Dla $\eps>0$ z ciągłości przesunięć wybieramy
$\delta>0$ tak, aby $\|\tau_yg-g\|_p<\eps$
dla $|y|<\delta$. Dzieląc drugą całkę jak poprzednio,
dostajemy granicę górną nie większą niż $\eps$.
Dowolność $\eps$ kończy dowód.

Dla $p=\infty$ nierówność kontrakcji wynika
bezpośrednio z $|g(x-y)|\leq\|g\|_\infty$
dla prawie wszystkich $y$ oraz normalizacji jądra.
\end{proof}

Ustalmy $\kappa>0$ i oznaczmy $P_0g=g$.
Dla $s,t>0$ i $g\in L^p$, $1\leq p\leq\infty$,
podwójna całka używana w $P_tP_sg(x)$ jest
bezwzględnie zbieżna. Istotnie, twierdzenie Tonellego
dla całki z $|g|$ oraz lemat~\ref{lem:pde4-convolution}
dają
\[
 \int_{\R^n}\int_{\R^n}\Gamma_\kappa(t,y)
 \Gamma_\kappa(s,z)|g(x-y-z)|\,
 \mathrm d\lambda_n(z)\,\mathrm d\lambda_n(y)
 =P_{t+s}|g|(x)<\infty.
\]
Można zatem zastosować twierdzenie Fubiniego
i otrzymać $P_tP_sg=P_{t+s}g$. Ta sama równość
zachodzi dla $s=0$ lub $t=0$ z definicji $P_0$.

Jeśli $g\in\mathrm{BUC}(\R^n)$, to
\[
 |P_tg(x+h)-P_tg(x)|
 \leq\sup_{z\in\R^n}|g(z+h)-g(z)|,
\]
co pokazuje, że $P_t$ zachowuje $\mathrm{BUC}$
i jest w niej kontrakcją.
W przestrzeni $X=L^p(\R^n)$ dla $1\leq p<\infty$
albo $X=\mathrm{BUC}(\R^n)$ dla $g\in X$ mamy
\[
 \|P_tg-P_sg\|_X\leq\|P_{t-s}g-g\|_X
 \qquad(0\leq s\leq t).
\]
Wynika to z własności półgrupy i kontrakcji
$P_s$. Twierdzenie~\ref{thm:pde4-approx-identity}
daje zatem ciągłość odwzorowania $t\mapsto P_tg$
w normie $X$ na $[0,\infty)$, również z lewej
strony w każdym punkcie dodatnim.

Na $L^\infty(\R^n)$ silna ciągłość w chwili zero
nie musi zachodzić. Istotnie, niech
$g(x)=\chi_{\{x_1>0\}}(x)$.
Z symetrii jądra względem zmiany znaku pierwszej
współrzędnej mamy $P_tg(0)=1/2$ dla każdego $t>0$.
Ponadto
\[
 P_tg(x)=\int_{\{y_1<x_1\}}\Gamma_\kappa(t,y)\,
 \mathrm d\lambda_n(y),
\]
więc funkcja po lewej jest ciągła w $x$ dzięki
twierdzeniu o zbieżności zdominowanej (hiperpłaszczyzna
$y_1=x_1$ ma miarę Lebesgue'a zero).
Dla każdego $\eta>0$ na pewnej półkuli
$B_\delta(0)\cap\{x_1>0\}$ dodatniej miary
zachodzi $|P_tg(x)-1/2|<\eta$ i $g(x)=1$.
Stąd $\|P_tg-g\|_\infty\geq1/2-\eta$,
a dowolność $\eta$ daje
$\|P_tg-g\|_\infty\geq1/2$ dla każdego $t>0$.
Przestrzeń $\mathrm{BUC}$ jest więc naturalną
przestrzenią danych ograniczonych dla silnie
ciągłej półgrupy cieplnej.

\section{Zagadnienie Cauchy'ego na całej przestrzeni}
\label{sec:pde4-cauchy}

\begin{twierdzenie}[Jawne rozwiązanie z danych początkowych]
\label{thm:pde4-cauchy}
Niech $n\geq1$, $\kappa>0$ i $g\in\mathrm{BUC}(\R^n)$.
Definiujemy dla $t>0$,
\begin{equation}\label{eq:pde4-cauchy}
 u(t,x)=\int_{\R^n}\Gamma_\kappa(t,x-y)g(y)
 \,\mathrm d\lambda_n(y),\qquad u(0,x)=g(x).
\end{equation}
Wówczas $u\in C([0,\infty);\mathrm{BUC}(\R^n))$,
$u\in C^\infty((0,\infty)\times\R^n)$,
$u_t-\kappa\Delta_xu=0$ dla $t>0$ oraz
$\sup_{t\geq0}\|u(t,\cdot)\|_\infty\leq\|g\|_\infty$.
Ponadto $g\geq0$ implikuje $u\geq0$.
Jest to jedyne rozwiązanie z klasy
$C^{1,2}((0,T]\times\R^n)\cap C([0,T]\times\R^n)$,
ograniczone na $[0,T]\times\R^n$, o tych danych,
dla każdego ustalonego $T>0$.
\end{twierdzenie}
\begin{proof}
Ustalmy $\kappa>0$ i zapiszmy $\Gamma=\Gamma_\kappa$.
Dla każdego $t>0$ całka w~\eqref{eq:pde4-cauchy}
jest bezwzględnie zbieżna, gdyż $g$ jest ograniczona,
a jądro ma całkę $1$. Wykażmy gładkość rozwiązania
we wszystkich zmiennych dla dodatnich czasów.

Dla każdej liczby całkowitej $m\geq0$ oraz
wielowskaźnika $\alpha\in\mathbb N_0^n$ pochodna
$\partial_t^mD_z^\alpha\Gamma(t,z)$ jest postaci
$A_{m,\alpha}(t,z)\Gamma(t,z)$, gdzie $A_{m,\alpha}$
jest wielomianem względem $z$ o współczynnikach
gładkich dla $t>0$. Istotnie, zaczynamy od $A_{0,0}=1$.
Różniczkowanie wyrażenia $A\Gamma$ daje odpowiednio
\[
 \begin{aligned}
 \partial_{z_j}(A\Gamma)
 &=\left(\partial_{z_j}A-\frac{z_j}{2\kappa t}A\right)\Gamma,\\
 \partial_t(A\Gamma)
 &=\left(\partial_tA+
 \left(-\frac n{2t}+\frac{|z|^2}{4\kappa t^2}\right)A\right)\Gamma.
 \end{aligned}
\]
Obie operacje zachowują opisaną postać, co dowodzi
jej dla wszystkich pochodnych przez indukcję.

Ustalmy $0<a<b<\infty$ i $L>0$ oraz rozważmy
$(t,x)\in[a,b]\times\overline{B_L(0)}$.
Dla ustalonej pochodnej istnieją stała $C>0$
i liczba całkowita $d\geq0$, niezależne od $t,x,y$,
takie, że $|A_{m,\alpha}(t,x-y)|\leq C(1+|x-y|)^d$.
Wynika to z ograniczoności skończonej liczby
współczynników wielomianu na $[a,b]$.
Ponadto z nierówności
$|y|^2\leq2|x-y|^2+2|x|^2$ dostajemy
$|x-y|^2\geq |y|^2/2-L^2$. W konsekwencji
\[
 e^{-|x-y|^2/(4\kappa t)}
 \leq e^{-|x-y|^2/(4\kappa b)}
 \leq e^{L^2/(4\kappa b)}e^{-|y|^2/(8\kappa b)}.
\]
Ponieważ $(4\pi\kappa t)^{-n/2}\leq(4\pi\kappa a)^{-n/2}$
i $|x-y|\leq L+|y|$, istnieje stała $C'>0$ taka,
że na rozważanym zbiorze
\[
 |\partial_t^mD_x^\alpha\Gamma(t,x-y)g(y)|
 \leq C'\|g\|_\infty(1+|y|)^d
       e^{-|y|^2/(8\kappa b)}.
\]
Prawa strona jest całkowalna w $\R^n$, ponieważ
każdy wielomian wzrostu promieniowego jest
całkowalny po pomnożeniu przez taką funkcję Gaussa.

Na wnętrzu wskazanego cylindra ilorazy różnicowe
są zdominowane przez oszacowanie odpowiedniej
pochodnej następnego rzędu, na mocy twierdzenia
o wartości średniej. Twierdzenie o zbieżności
zdominowanej uzasadnia zarówno różniczkowanie
pod całką, jak i ciągłość otrzymanych pochodnych.
Dowolność $a,b,L$ daje $u\in C^\infty((0,\infty)\times\R^n)$
i dla $t>0$ wzór
\[
 \partial_t^mD_x^\alpha u(t,x)
 =\int_{\R^n}\partial_t^mD_x^\alpha\Gamma(t,x-y)
 g(y)\,\mathrm d\lambda_n(y).
\]
Stąd, stosując twierdzenie~\ref{thm:pde4-gauss-basic},
otrzymujemy
$(\partial_t-\kappa\Delta_x)u(t,x)=0$.

Z poprzedniej sekcji wiemy, że $P_t$ zachowuje
przestrzeń $\mathrm{BUC}$ oraz że $t\mapsto P_tg$
jest ciągłe w jej normie na $[0,\infty)$,
gdzie $P_0g=g$. W szczególności $u(t,\cdot)\to g$
jednostajnie na $\R^n$ przy $t\downarrow0$.
Uzasadnijmy również ciągłość łączną: dla $s\geq0$,
$x_0\in\R^n$ i $(t,x)\to(s,x_0)$, $t\geq0$, mamy
\[
 |u(t,x)-u(s,x_0)|\leq
 \|P_tg-P_sg\|_\infty+
 |P_sg(x)-P_sg(x_0)|\longrightarrow0.
\]
Pierwszy składnik znika dzięki ciągłości normowej,
a drugi dzięki ciągłości funkcji $P_sg$.

Kontrakcyjność daje $\|u(t,\cdot)\|_\infty\leq\|g\|_\infty$
również dla $t=0$. Jeżeli $g\geq0$, to $u\geq0$
dzięki dodatniości jądra. Na każdym przedziale
$[0,T]$ skonstruowana funkcja należy więc do klasy
wymaganej w tezie o jednoznaczności.
Twierdzenie~\ref{thm:pde4-global-uniqueness},
udowodnione poniżej przy użyciu zasady maksimum,
zapewnia jednoznaczność w tej klasie. Jego dowód
nie korzysta z obecnej reprezentacji całkowej.
\end{proof}

\begin{wniosek}[Natychmiastowe rozchodzenie się zaburzenia]
\label{cor:pde4-positive}
Niech $n\geq1$, $\kappa>0$, a
$g\in C_c(\R^n)$ będzie nieujemna i nierówna identycznie
zeru. Wówczas rozwiązanie~\eqref{eq:pde4-cauchy}
ma własność $u(t,x)>0$ dla każdego $t>0$, $x\in\R^n$.
\end{wniosek}
\begin{proof}
Funkcja ciągła o zwartym nośniku jest ograniczona
i jednostajnie ciągła na $\R^n$, więc spełnia
założenia twierdzenia~\ref{thm:pde4-cauchy}.
Ponieważ $g$ jest nieujemna i nie jest identycznie
zerowa, istnieje $y_0\in\R^n$ takie, że $g(y_0)>0$.
Z ciągłości wybieramy $r>0$ i $c=g(y_0)/2>0$
tak, aby $g\geq c$ na $\overline{B_r(y_0)}$.
Dla ustalonego $t>0$ i $x\in\R^n$ funkcja
$y\mapsto\Gamma_\kappa(t,x-y)$ jest ciągła
i ściśle dodatnia, a więc osiąga dodatnie minimum
$m_{t,x}$ na tej zwartej kuli.
Nieujemność $g$ na całej przestrzeni daje
\[
 u(t,x)\geq c\int_{B_r(y_0)}
 \Gamma_\kappa(t,x-y)\,\mathrm d\lambda_n(y)
 \geq c\,m_{t,x}\lambda_n(B_r(y_0))>0.
\]
W szczególności zwarty nośnik nieujemnych,
niezerowych danych początkowych nie pociąga
za sobą ograniczonego nośnika rozwiązania
w żadnej dodatniej chwili. Nie twierdzimy tu,
że dane zmieniające znak dają rozwiązanie dodatnie.
\end{proof}

\begin{stwierdzenie}[Zachowanie masy i tłumienie norm]
\label{prop:pde4-mass}
Niech $n\geq1$, $\kappa>0$, $t>0$.
Dla $g\in L^1(\R^n)$ zachodzi
\begin{equation}\label{eq:pde4-mass}
 \int_{\R^n}(P_tg)(x)\,\mathrm d\lambda_n(x)
 =\int_{\R^n}g(x)\,\mathrm d\lambda_n(x).
\end{equation}
Dla $g\in L^1(\R^n)$ zachodzi również
\begin{equation}\label{eq:pde4-L1Linf}
 \|P_tg\|_\infty\leq (4\pi\kappa t)^{-n/2}\|g\|_1,
 \qquad
 \|P_tg\|_2\leq (8\pi\kappa t)^{-n/4}\|g\|_1.
\end{equation}
\end{stwierdzenie}
\begin{proof}
Niech $g\in L^1(\R^n)$ i ustalmy $t>0$.
Twierdzenie Tonellego dla nieujemnej funkcji
podcałkowej oraz normalizacja jądra dają
\[
 \begin{aligned}
 &\int_{\R^n}\int_{\R^n}\Gamma_\kappa(t,x-y)|g(y)|
       \,\mathrm d\lambda_n(y)\,\mathrm d\lambda_n(x)\\
 &\quad=\int_{\R^n}|g(y)|
   \left(\int_{\R^n}\Gamma_\kappa(t,x-y)\,
                 \mathrm d\lambda_n(x)\right)\mathrm d\lambda_n(y)
 =\|g\|_1.
 \end{aligned}
\]
Wewnętrzna całka ma wartość $1$ dzięki niezmienniczości
miary Lebesgue'a względem przesunięć.
W szczególności podwójna całka z $g$ zamiast
$|g|$ jest bezwzględnie zbieżna. Stosując
twierdzenie Fubiniego otrzymujemy
\[
 \int_{\R^n}P_tg(x)\,\mathrm d\lambda_n(x)
 =\int_{\R^n}g(y)\,\mathrm d\lambda_n(y),
\]
czyli~\eqref{eq:pde4-mass}. Jest to zachowanie
całki z funkcji, również dla danych zmieniających
znak. Jeśli $g\geq0$ prawie wszędzie, dodatniość
jądra daje $P_tg\geq0$ i stąd
$\|P_tg\|_1=\|g\|_1$. Dla dowolnych danych
$|P_tg|\leq P_t|g|$, co daje tylko kontrakcję
w normie $L^1$.

Maksimum jądra jest osiągane w zerze:
$\Gamma_\kappa(t,0)=(4\pi\kappa t)^{-n/2}$.
Zatem dla każdego $x\in\R^n$ zachodzi
\[
 |P_tg(x)|\leq(4\pi\kappa t)^{-n/2}\|g\|_1.
\]
Biorąc istotne supremum po $x$, otrzymujemy
pierwszą nierówność w~\eqref{eq:pde4-L1Linf}.

Dla drugiej obliczamy normę jądra wprost.
Korzystając z całki Gaussa, mamy
\[
 \begin{aligned}
 \|\Gamma_\kappa(t,\cdot)\|_2^2
 &=(4\pi\kappa t)^{-n}
 \int_{\R^n}e^{-|x|^2/(2\kappa t)}\,\mathrm d\lambda_n(x)\\
 &=(4\pi\kappa t)^{-n}(2\pi\kappa t)^{n/2}
 =(8\pi\kappa t)^{-n/2}.
 \end{aligned}
\]
Z niezmienniczości normy względem przesunięć
i całkowej nierówności Minkowskiego wynika
\[
 \begin{aligned}
 \|P_tg\|_2
 &\leq\int_{\R^n}|g(y)|
       \|\Gamma_\kappa(t,\cdot-y)\|_2\,\mathrm d\lambda_n(y)\\
 &= (8\pi\kappa t)^{-n/4}\|g\|_1.
 \end{aligned}
\]
Prawa strona jest skończona dla $g\in L^1$,
co zarazem uzasadnia, że $P_tg\in L^2$.
\end{proof}

\section{Paraboliczna zasada maksimum i jednoznaczność}
\label{sec:pde4-maximum}

Niech $\Omega\subset\R^n$ będzie obszarem ograniczonym,
a $T>0$. Wprowadzamy cylinder $Q_T=(0,T)\times\Omega$,
jego domknięcie $\overline{Q_T}=[0,T]\times\overline\Omega$
oraz \emph{brzeg paraboliczny}
\[
 \partial_pQ_T=(\{0\}\times\overline\Omega)
 \cup([0,T]\times\partial\Omega).
\]
Brzeg paraboliczny obejmuje dolną podstawę i ścianę
boczną cylindra, ale nie obejmuje punktów
$\{T\}\times\Omega$. To właśnie na brzegu parabolicznym
kontrolujemy dane rozwiązania w zagadnieniu postępowym;
na górnej podstawie nie zadaje się niezależnych danych.

Przez $C^{1,2}((0,T]\times\Omega)$ rozumiemy tutaj
funkcje o ciągłej pochodnej czasowej i ciągłych
pochodnych przestrzennych do rzędu drugiego
na wskazanym zbiorze; w chwili $T$ pochodna czasowa
jest rozumiana jako pochodna lewostronna.
Nie zakładamy istnienia pochodnych na brzegu
przestrzennym ani w chwili początkowej.

\begin{twierdzenie}[Słaba paraboliczna zasada maksimum]
\label{thm:pde4-max}
Niech $n\geq1$, $\Omega\subset\R^n$ będzie obszarem
ograniczonym, $T>0$, $\kappa>0$, a
$u\in C([0,T]\times\overline\Omega)\cap
C^{1,2}((0,T]\times\Omega)$ będzie funkcją rzeczywistą
spełniającą
$\partial_tu-\kappa\Delta_xu\leq0$
na $(0,T]\times\Omega$.
Wówczas
\begin{equation}\label{eq:pde4-max}
 \max_{[0,T]\times\overline\Omega}u
 \leq \max_{\partial_pQ_T}u.
\end{equation}
W konsekwencji, jeśli funkcje $u_1,u_2$ z tej samej
klasy spełniają
$(\partial_t-\kappa\Delta_x)u_1
\leq(\partial_t-\kappa\Delta_x)u_2$
na $(0,T]\times\Omega$ oraz $u_1\leq u_2$
na $\partial_pQ_T$, to $u_1\leq u_2$
na całym $\overline{Q_T}$.
W szczególności rozwiązania tego samego równania
z tymi samymi danymi na brzegu parabolicznym
są jednoznaczne w podanej klasie.
\end{twierdzenie}
\begin{proof}
Ustalmy $\eps>0$ i połóżmy $v(t,x)=u(t,x)-\eps t$.
Funkcja $v$ ma tę samą regularność co $u$ oraz spełnia
\[
 \partial_tv-\kappa\Delta_xv
 =\partial_tu-\kappa\Delta_xu-\eps\leq-\eps<0.
\]
Ponieważ $\Omega$ jest ograniczony,
$\overline{Q_T}$ jest zwarty. Funkcja $v$ osiąga
więc na nim maksimum. Przypuśćmy, że pewien
punkt maksimum $(t_0,x_0)$ nie należy do
$\partial_pQ_T$. Wówczas $0<t_0\leq T$ i $x_0\in\Omega$.
Funkcja $x\mapsto v(t_0,x)$ osiąga maksimum
lokalne w punkcie wewnętrznym $x_0$. Stąd
$\nabla_xv(t_0,x_0)=0$, jej hesjan przestrzenny
jest niedodatnio określony i w szczególności
$\Delta_xv(t_0,x_0)\leq0$.

Jeśli $t_0<T$, funkcja $t\mapsto v(t,x_0)$
ma maksimum w punkcie wewnętrznym przedziału
$(0,T)$, a więc $\partial_tv(t_0,x_0)=0$.
Jeśli natomiast $t_0=T$, dla każdego dostatecznie
małego $h>0$ mamy
\[
 \frac{v(T,x_0)-v(T-h,x_0)}{h}\geq0.
\]
Przechodząc z $h$ do zera otrzymujemy
$\partial_tv(T,x_0)\geq0$ w sensie pochodnej
lewostronnej. W obu przypadkach
$(\partial_t-\kappa\Delta_x)v(t_0,x_0)\geq0$,
co przeczy ścisłej nierówności dla $v$.

Każdy punkt maksimum $v$ należy zatem do brzegu
parabolicznego. Na tym brzegu $v\leq u$, gdyż
$t\geq0$, a na całym domkniętym cylindrze
$u=v+\eps t\leq v+\eps T$. Otrzymujemy
\[
 \max_{\overline{Q_T}}u
 \leq\max_{\overline{Q_T}}v+\eps T
 =\max_{\partial_pQ_T}v+\eps T
 \leq\max_{\partial_pQ_T}u+\eps T.
\]
Dowolność $\eps>0$ daje~\eqref{eq:pde4-max}.
Nie użyliśmy żadnego założenia o gładkości
brzegu przestrzennego $\Omega$.

Dla porównania funkcji $u_1,u_2$ stosujemy
wykazaną nierówność do $w=u_1-u_2$.
Z założeń wynika
$(\partial_t-\kappa\Delta_x)w\leq0$
oraz $w\leq0$ na brzegu parabolicznym, więc $w\leq0$
na $\overline{Q_T}$. Jeżeli obie funkcje spełniają
to samo równanie i mają te same dane, stosujemy
porównanie także w odwrotnym kierunku.
\end{proof}

\begin{twierdzenie}[Jednoznaczność w klasie rozwiązań ograniczonych]
\label{thm:pde4-global-uniqueness}
Niech $n\geq1$, $\kappa>0$, $T>0$. Jeśli
$u,v\in C([0,T]\times\R^n)\cap
C^{1,2}((0,T]\times\R^n)$ są rzeczywistymi
funkcjami ograniczonymi na $[0,T]\times\R^n$,
spełniają
\[
 \partial_tu-\kappa\Delta_xu=0,\qquad
 \partial_tv-\kappa\Delta_xv=0
 \quad\text{na }(0,T]\times\R^n,
\]
to zachodzi oszacowanie stabilności względem
danych początkowych
\[
 \sup_{(t,x)\in[0,T]\times\R^n}|u(t,x)-v(t,x)|
 \leq\sup_{x\in\R^n}|u(0,x)-v(0,x)|.
\]
W szczególności, jeśli $u(0,x)=v(0,x)$ dla wszystkich
$x\in\R^n$, to $u=v$ na $[0,T]\times\R^n$.
\end{twierdzenie}
\begin{proof}
Niech $w=u-v$ i oznaczmy
$M=\sup_{[0,T]\times\R^n}|w|<\infty$ oraz
$A=\sup_{\R^n}|w(0,\cdot)|\leq M$.
Funkcja $w$ spełnia jednorodne równanie ciepła.
Nie możemy jednak bezpośrednio zastosować
twierdzenia~\ref{thm:pde4-max} do całej przestrzeni,
gdyż cylinder jest nieograniczony.

Dla $\eps>0$ wprowadzamy funkcję
$b_\eps(t,x)=\eps(|x|^2+2n\kappa t)$.
Ponieważ $\Delta_x|x|^2=2n$, zachodzi
\[
 \partial_tb_\eps=2n\kappa\eps
 =\kappa\Delta_xb_\eps.
\]
Funkcja $b_\eps$ jest więc dokładnym rozwiązaniem
równania ciepła. Jest nieujemna dla $t\geq0$
i rośnie kwadratowo wraz z $|x|$.
Ustalmy dowolny punkt $(t_*,x_*)\in[0,T]\times\R^n$.
Dla danego $\eps>0$ wybierzmy $R>0$ tak, aby
$R>|x_*|$ i $\eps R^2>M$.

Na cylindrze $[0,T]\times\overline{B_R(0)}$
rozważamy $w_\eps=w-A-b_\eps$.
Funkcja ta jest ciągła na domknięciu cylindra,
ma wymaganą regularność wewnątrz i na górnej
podstawie oraz spełnia
$(\partial_t-\kappa\Delta_x)w_\eps=0$.
Na dolnej podstawie mamy
\[
 w_\eps(0,x)=w(0,x)-A-\eps|x|^2\leq0
 \qquad(x\in\overline{B_R(0)}).
\]
Na ścianie bocznej dla $0\leq t\leq T$ i $|x|=R$
otrzymujemy
$w_\eps(t,x)\leq M-A-\eps R^2<0$.
Twierdzenie~\ref{thm:pde4-max} daje więc
$w_\eps\leq0$ na całym domkniętym cylindrze.
W szczególności
\[
 w(t_*,x_*)\leq A+\eps(|x_*|^2+2n\kappa t_*).
\]
Chociaż promień $R$ zależy od $\eps$, powyższe
oszacowanie w ustalonym punkcie zachodzi dla
każdego $\eps>0$. Przechodząc z $\eps$ do zera,
otrzymujemy $w(t_*,x_*)\leq A$.

Stosujemy ten sam argument do $-w$, dla którego
stałe $M$ i $A$ są takie same. Otrzymujemy
$-w(t_*,x_*)\leq A$. Punkt był dowolny, zatem
$|w|\leq A$ na $[0,T]\times\R^n$, co dowodzi
oszacowania. Dla zgodnych danych początkowych
$A=0$, więc $w=0$.
Dowód nie korzysta z reprezentacji rozwiązania
za pomocą splotu z jądrem Gaussa.
\end{proof}

\begin{uwaga}[Dlaczego trzeba podać klasę rozwiązań?]
Ograniczoność na całym $[0,T]\times\R^n$ jest
w powyższym dowodzie istotna: pozwala dla danego
$\eps$ wybrać promień kuli niezależny od czasu,
na której ścianie bocznej dominuje bariera.
Nie należy utożsamiać tego założenia z samą
ograniczonością na zbiorach zwartych, która
wynika już z ciągłości rozwiązania.
Pod koniec rozdziału przedstawimy niezerowe
rozwiązanie jednorodnego równania ciepła z zerowymi
danymi początkowymi, niespełniające globalnego
warunku ograniczoności.
\end{uwaga}

\section{Wygładzanie i analityczność przestrzenna}
\label{sec:pde4-smoothing}

\begin{twierdzenie}[Oszacowania pochodnych]
\label{thm:pde4-smoothing}
Niech $n\geq1$, $\kappa>0$, $g\in L^\infty(\R^n)$,
a $u(t,x)=P_tg(x)$ dla $t>0$.
Dla każdego wielowskaźnika $\alpha\in\mathbb N_0^n$
istnieje stała $C_{\alpha,n}>0$, niezależna od
$g,t,\kappa$, taka że dla każdego $t>0$
\begin{equation}\label{eq:pde4-derivative}
 \|\partial_x^\alpha u(t,\cdot)\|_\infty
 \leq C_{\alpha,n}(\kappa t)^{-|\alpha|/2}\|g\|_\infty.
\end{equation}
Co więcej, dla każdego ustalonego $t>0$ funkcja
$x\mapsto u(t,x)$ przedłuża się do funkcji holomorficznej
$U_t$ na $\C^n$, spełniającej
\[
 |U_t(x+i\eta)|\leq
 e^{|\eta|^2/(4\kappa t)}\|g\|_\infty
 \qquad(x,\eta\in\R^n).
\]
\end{twierdzenie}
\begin{proof}
Oznaczmy $a=\kappa t$ i $K(z)=(4\pi)^{-n/2}e^{-|z|^2/4}$
dla $z\in\R^n$. Wówczas $\Gamma_\kappa(t,x)=a^{-n/2}K(x/\sqrt a)$.
Dla wielowskaźnika $\alpha=(\alpha_1,\ldots,\alpha_n)$
o długości $|\alpha|=\sum_j\alpha_j$ reguła łańcuchowa daje
\[
 \partial_x^\alpha\Gamma_\kappa(t,x)
 =a^{-(n+|\alpha|)/2}(\partial^\alpha K)(x/\sqrt a).
\]
Ponieważ $\partial_{z_j}K=-z_jK/2$, indukcyjnie otrzymujemy
$\partial^\alpha K=Q_\alpha K$, gdzie $Q_0=1$ oraz
$Q_{\alpha+e_j}=\partial_{z_j}Q_\alpha-z_jQ_\alpha/2$;
$e_j$ jest $j$-tym wektorem standardowej bazy.
Wielomian $Q_\alpha$ ma niezerowy składnik jednorodny
najwyższego stopnia $(-1/2)^{|\alpha|}z^\alpha$.
Zatem $C_{\alpha,n}=\|\partial^\alpha K\|_1$ jest
skończone i dodatnie. Podstawienie $x=\sqrt a\,z$
daje $\|\partial_x^\alpha\Gamma_\kappa(t,\cdot)\|_1
=a^{-|\alpha|/2}C_{\alpha,n}$.
Argument dominacji z dowodu twierdzenia
\ref{thm:pde4-cauchy} wymaga tylko istotnej
ograniczoności $g$, nie jej ciągłości: wybieramy
ograniczonego reprezentanta klasy. Wolno więc
różniczkować splot dowolną liczbę razy dla $t>0$.
Z niezmienniczości miary Lebesgue'a względem przesunięć
otrzymujemy
\[
 |\partial_x^\alpha u(t,x)|
 \leq\|g\|_\infty\int_{\R^n}
 |\partial_x^\alpha\Gamma_\kappa(t,x-y)|\,\mathrm d\lambda_n(y)
 =a^{-|\alpha|/2}C_{\alpha,n}\|g\|_\infty.
\]
Biorąc supremum otrzymujemy~\eqref{eq:pde4-derivative}.

Dla ustalonego $t>0$ definiujemy
\[
 U_t(z)=\int_{\R^n}E(z,y)g(y)\,\mathrm d\lambda_n(y),
 \qquad
 E(z,y)=\frac1{(4\pi a)^{n/2}}
 \exp\!\left(-\frac{\sum_{j=1}^n(z_j-y_j)^2}{4a}\right).
\]
W wykładniku występuje suma kwadratów zespolonych,
a nie kwadrat normy hermitowskiej. Dla $z=x+i\eta$
mamy $\operatorname{Re}\sum_j(z_j-y_j)^2=|x-y|^2-|\eta|^2$,
więc $|E(x+i\eta,y)|=e^{|\eta|^2/(4a)}\Gamma_\kappa(t,x-y)$.
Normalizacja jądra daje bezwzględną zbieżność
całki dla każdego $z$ oraz oszacowanie podane w tezie.

Dla każdego $y$ funkcja $E(\cdot,y)$ jest holomorficzna
na $\C^n$. Jeżeli $|\operatorname{Re}z|,|\operatorname{Im}z|\leq L$,
to z $|x-y|^2\geq |y|^2/2-L^2$ wynika
\[
 |E(z,y)|\leq(4\pi a)^{-n/2}e^{L^2/(2a)}e^{-|y|^2/(8a)}.
\]
Ponadto
\[
 \begin{aligned}
 \partial_{z_j}E&=-\frac{z_j-y_j}{2a}E,\\
 \partial_{z_k}\partial_{z_j}E
 &=\left(\frac{(z_j-y_j)(z_k-y_k)}{4a^2}
 -\frac{\delta_{jk}}{2a}\right)E,
 \end{aligned}
\]
gdzie $\delta_{jk}=1$ dla $j=k$ i $0$ w przeciwnym
przypadku. Funkcję $E$ i obie grupy pochodnych,
pomnożone przez $g(y)$, dominuje na wskazanym
zbiorze jedna całkowalna funkcja postaci
$C\|g\|_\infty(1+|y|)^2e^{-|y|^2/(8a)}$.

Ustalmy $z_0\in\C^n$ i wybierzmy $L$ tak duże,
aby odcinki $z_0+sh$, $0\leq s\leq1$, $|h|\leq1$,
mieściły się w rozważanym zbiorze. Wzór Taylora
z resztą całkową dla funkcji $s\mapsto E(z_0+sh,y)$ daje
\[
 \begin{aligned}
 E(z_0+h,y)-E(z_0,y)
 &=\sum_j\partial_{z_j}E(z_0,y)h_j\\
 &\quad+\int_0^1(1-s)\sum_{j,k}
 \partial_{z_j}\partial_{z_k}E(z_0+sh,y)h_jh_k\,
 \mathrm d\lambda_1(s).
 \end{aligned}
\]
Ponieważ $\sum_{j,k}|h_jh_k|\leq n|h|^2$,
moduł reszty pomnożonej przez $g$ jest ograniczony
przez całkowalną majorantę razy $|h|^2$. Po całkowaniu
otrzymujemy
\[
 U_t(z_0+h)=U_t(z_0)+\sum_jh_j
 \int_{\R^n}\partial_{z_j}E(z_0,y)g(y)\,\mathrm d\lambda_n(y)
 +O(|h|^2).
\]
Wyraz liniowy jest zespolono-liniowy względem $h$.
Funkcja $U_t$ jest różniczkowalna w sensie Frécheta
nad $\C$ w każdym punkcie $\C^n$, a zatem holomorficzna.
Jej ograniczenie do $\R^n$ jest równe $P_tg$,
w szczególności jest rzeczywiście analityczne.
\end{proof}

\begin{przyklad}[Wygasanie wysokich częstości]
Niech $\xi\in\R^n$ i $g(x)=\cos(\xi\cdot x)$.
Funkcja $g$ należy do $\mathrm{BUC}(\R^n)$,
gdyż jest ograniczona i ma ograniczony gradient.
Połóżmy $u(t,x)=e^{-\kappa|\xi|^2t}\cos(\xi\cdot x)$.
Mamy $\partial_tu=-\kappa|\xi|^2u$ oraz
$\partial_{x_j}^2u=-\xi_j^2u$, więc
$\partial_tu-\kappa\Delta_xu=0$. Funkcja jest gładka,
spełnia dane początkowe i jest ograniczona przez $1$
na całym $[0,\infty)\times\R^n$. Z twierdzeń
\ref{thm:pde4-cauchy} i~\ref{thm:pde4-global-uniqueness}
wynika $u(t,\cdot)=P_tg$. Dla ustalonego $t>0$
mnożnik $e^{-\kappa t|\xi|^2}$ maleje wykładniczo
względem kwadratu normy częstości przestrzennej.
Dla $\xi=0$ rozwiązanie jest stałe i nie zanika.
\end{przyklad}

\begin{lemat}[Generator półgrupy cieplnej na gładkich danych]
\label{lem:pde-v31-heat-generator}
Niech $n\geq1$, $\kappa>0$ oraz
$h\in C^2(\R^n)$, przy czym $h$ i wszystkie pochodne
przestrzenne rzędu co najwyżej dwa są ograniczone
i jednostajnie ciągłe na $\R^n$.
Wtedy dla $t>0$ zachodzi
\begin{equation}\label{eq:pde-v31-heat-generator}
 P_th-h=\kappa\int_0^t P_s(\Delta h)\,\mathrm d\lambda_1(s),
 \qquad
 \lim_{t\downarrow0}\left\|\frac{P_th-h}{t}
                 -\kappa\Delta h\right\|_\infty=0.
\end{equation}
Całkę rozumiemy punktowo lub, równoważnie, jako
całkę funkcji ciągłej o wartościach w $\mathrm{BUC}(\R^n)$.
\end{lemat}
\begin{proof}
Dla $s>0$ mamy $P_sh(x)=\int_{\R^n}\Gamma_\kappa(s,z)
h(x-z)\,\mathrm d\lambda_n(z)$. Różniczkując pod
całką względem $x$ do rzędu drugiego, korzystamy
z ograniczoności pochodnych $h$ i całkowalności
jądra: ilorazy różnicowe są zdominowane przez
jądro pomnożone przez supremum odpowiedniej pochodnej $h$.
Otrzymujemy $\Delta_xP_sh=P_s(\Delta h)$.
Z twierdzenia~\ref{thm:pde4-cauchy} wynika więc
$\partial_sP_sh=\kappa\Delta_xP_sh=\kappa P_s(\Delta h)$.

Alternatywnie można przenieść laplasjan z jądra
na dane. Ustalmy $s>0$, $x\in\R^n$ i oznaczmy
$K_x(y)=\Gamma_\kappa(s,x-y)$. Mamy
$\Delta_x K_x=\Delta_y K_x$. Całkowanie przez części
na $B_R(0)$ daje
\[
 \int_{B_R}(h\Delta_y K_x-K_x\Delta h)\,\mathrm d\lambda_n
 =\int_{\partial B_R}(h\partial_\nu K_x-K_x\partial_\nu h)
 \,\mathrm d\sigma.
\]
Dla $R>|x|$ moduł całki brzegowej jest nie większy niż
\[
 \begin{aligned}
 &\sigma(S^{n-1})R^{n-1}(4\pi\kappa s)^{-n/2}
 e^{-(R-|x|)^2/(4\kappa s)}\\
 &\qquad{}\cdot\left(\|h\|_\infty\frac{R+|x|}{2\kappa s}
 +\|\nabla h\|_\infty\right).
 \end{aligned}
\]
W wymiarze jeden przyjmujemy $\sigma(S^0)=2$,
a całka brzegowa jest sumą po końcach przedziału.
Oszacowanie dąży do zera przy $R\to\infty$.
Całki objętościowe są bezwzględnie zbieżne,
gdyż $h,\Delta h$ są ograniczone, a jądro
i jego pochodne są całkowalne. Przejście do granicy
daje tę samą tożsamość $\Delta_xP_sh=P_s(\Delta h)$.

Dla $0<\delta<t$ całkujemy punktowo równanie
względem czasu, otrzymując
\[
 P_th(x)-P_\delta h(x)
 =\kappa\int_\delta^t P_s(\Delta h)(x)\,\mathrm d\lambda_1(s).
\]
Ponieważ $\Delta h\in\mathrm{BUC}$, odwzorowanie
$s\mapsto P_s(\Delta h)$ jest ciągłe w normie
na $[0,t]$. Jego całka w tej przestrzeni istnieje
jako granica sum Riemanna; ciągłość funkcjonałów
obliczania wartości w punkcie daje zgodność z całką
punktową. Mamy
$\|\int_0^\delta P_s(\Delta h)\,\mathrm d\lambda_1(s)\|_\infty
\leq\delta\|\Delta h\|_\infty$ oraz
$P_\delta h\to h$ w normie supremum.
Przejście z $\delta$ do zera daje pierwszą równość.

Dzieląc ją przez $t$ i odejmując $\kappa\Delta h$,
otrzymujemy
\[
 \left\|\frac{P_th-h}{t}-\kappa\Delta h\right\|_\infty
 \leq\kappa\sup_{0\leq s\leq t}
 \|P_s(\Delta h)-\Delta h\|_\infty
 \longrightarrow0\quad(t\downarrow0).
\]
Ostatnia granica wynika z silnej ciągłości półgrupy.
Dla tych danych generator półgrupy na $\mathrm{BUC}$
działa zatem jako $\kappa\Delta$; nie charakteryzujemy
tutaj całej jego dziedziny.
\end{proof}

\section{Wymuszenie: zasada Duhamela}
\label{sec:pde4-duhamel}

\begin{twierdzenie}[Wzór Duhamela]
\label{thm:pde4-duhamel}
Niech $T>0$, $n\geq1$, $\kappa>0$, $g\in C_c^\infty(\R^n)$,
a $f\in C^\infty([0,T]\times\R^n)$ ma tę własność, że
$f$ i wszystkie jej pochodne względem $x$ i $t$ są
ograniczone, a nośniki $f(t,\cdot)$ są zawarte
we wspólnym zbiorze zwartym $K\subset\R^n$
dla $0\leq t\leq T$. Pochodne czasowe na końcach
przedziału rozumiemy jednostronnie.
Wtedy funkcja
\begin{equation}\label{eq:pde4-duhamel}
 u(t,x)=P_tg(x)+\int_0^t(P_{t-s}f(s,\cdot))(x)
 \,\mathrm d\lambda_1(s),\quad 0\leq t\leq T,
\end{equation}
przy konwencji $P_0h=h$, jest klasy $C^{1,2}$
na $[0,T]\times\R^n$, spełnia
$\partial_tu-\kappa\Delta_xu=f$ oraz $u(0,\cdot)=g$.
Równanie zachodzi również w chwili $0$ w sensie
prawostronnej pochodnej czasowej oraz w chwili $T$
w sensie pochodnej lewostronnej. Ponadto
\begin{equation}\label{eq:pde4-duhamel-bound}
 \|u(t,\cdot)\|_\infty\leq\|g\|_\infty+
 \int_0^t\|f(s,\cdot)\|_\infty\,\mathrm d\lambda_1(s).
\end{equation}
Jest to jedyne rozwiązanie ograniczone na
$[0,T]\times\R^n$ w klasie
$C([0,T]\times\R^n)\cap C^{1,2}((0,T]\times\R^n)$
z tymi danymi.
\end{twierdzenie}
\begin{proof}
Oznaczmy $X=\mathrm{BUC}(\R^n)$ z normą supremum
oraz $F_\alpha(s)=D_x^\alpha f(s,\cdot)$ dla
$\alpha\in\mathbb N_0^n$; w szczególności $F_0(s)=f(s,\cdot)$.
Dla każdego $\alpha$ funkcja $F_\alpha(s)$ należy do $X$,
ponieważ jest ograniczona i ma ograniczony gradient.
Z ograniczoności pochodnej mieszanej
$\partial_sD_x^\alpha f$ otrzymujemy
\[
 \|F_\alpha(s)-F_\alpha(s')\|_\infty
 \leq |s-s'|\sup_{[0,T]\times\R^n}
                  |\partial_tD_x^\alpha f|.
\]
Odwzorowania $F_\alpha:[0,T]\to X$ są więc ciągłe.
Analogiczne warunki przestrzenne spełnia każda
pochodna $g$.

Najpierw uzasadnijmy ciągłość funkcji podcałkowych
w normie $X$ również przy $s=t$. Dla $r,r'\geq0$
oraz $h,h'\in X$ kontrakcyjność daje
\[
 \|P_rh-P_{r'}h'\|_\infty
 \leq\|h-h'\|_\infty+\|P_rh'-P_{r'}h'\|_\infty.
\]
Drugi składnik dąży do zera dla $r\to r'$
przy ustalonym $h'$, dzięki silnej ciągłości półgrupy.
Wynika stąd łączna ciągłość odwzorowania
$(r,h)\mapsto P_rh$ w $[0,\infty)\times X$.
W szczególności dla każdego $\alpha$ funkcja
\[
 H_\alpha(t,s)=P_{t-s}F_\alpha(s),
 \qquad (t,s)\in D_T:=\{(t,s):0\leq s\leq t\leq T\},
\]
jest ciągła w normie $X$ na zwartym trójkącie $D_T$.
Na przekątnej $s=t$ ma wartość $F_\alpha(t)$.
W konsekwencji jest jednostajnie ciągła i ograniczona
na $D_T$. Całki względem $s$ istnieją w $X$ jako
całki funkcji ciągłych i zgadzają się z całkami
punktowymi po obliczeniu wartości w dowolnym $x$.

Dla $r>0$ przenosimy pochodne przestrzenne splotu
na dane, jak w dowodzie lematu
\ref{lem:pde-v31-heat-generator}. Ograniczoność
odpowiednich pochodnych i całkowalność jądra
uzasadniają ten rachunek przez zbieżność zdominowaną.
Dla $r=0$ ta sama tożsamość wynika z definicji $P_0$.
Zatem
$D_x^\alpha(P_rf(s,\cdot))=P_rF_\alpha(s)$
oraz $D_x^\alpha P_rg=P_rD^\alpha g$.
Przy różniczkowaniu całki w~\eqref{eq:pde4-duhamel}
ilorazy różnicowe są ograniczone niezależnie
od $s$ przez supremum odpowiedniej pochodnej
następnego rzędu funkcji $f$ na całym pasku,
na mocy twierdzenia o wartości średniej
i kontrakcyjności $P_{t-s}$.
Dla każdego $t$ dostajemy
\[
 D_x^\alpha u(t,\cdot)
 =P_tD^\alpha g+\int_0^t H_\alpha(t,s)\,
                       \mathrm d\lambda_1(s).
\]
Odwzorowanie po prawej stronie jest ciągłe
w normie $X$ względem $t$: dla $t,t'$ porównujemy
całki na $[0,\min(t,t')]$ przez jednostajną ciągłość
$H_\alpha$ i na pozostałym przedziale przez jego
długość $|t-t'|$ razy ograniczenie $H_\alpha$.
W szczególności w chwili $0$ pochodne przestrzenne
mają wartości $D^\alpha g$.

Pozostaje różniczkowanie względem czasu. Oznaczmy
$F=F_0$, $G(s)=\Delta_x f(s,\cdot)$ oraz
$V(t)=\int_0^tP_{t-s}F(s)\,\mathrm d\lambda_1(s)$.
Z lematu~\ref{lem:pde-v31-heat-generator}
i własności półgrupy wynika dla $0\leq r_1\leq r_2$
oraz $0\leq s\leq T$ tożsamość w $X$
\[
 P_{r_2}F(s)-P_{r_1}F(s)
 =\kappa\int_{r_1}^{r_2}P_qG(s)\,\mathrm d\lambda_1(q).
\]
Istotnie, stosujemy lemat dla $h=F(s)$,
odejmując tożsamości dla $r_2$ i $r_1$;
przy $r_1=0$ korzystamy z $P_0F(s)=F(s)$.
Nie pojawia się więc osobliwość przy $r_1=0$.

Dla $0\leq t<T$ i $0<\tau\leq T-t$ mamy
\[
 \begin{aligned}
 \frac{V(t+\tau)-V(t)}{\tau}
 &=\frac{\kappa}{\tau}\int_0^t
       \int_{t-s}^{t+\tau-s}P_qG(s)\,
       \mathrm d\lambda_1(q)\,\mathrm d\lambda_1(s)\\
 &\quad+\frac1{\tau}\int_t^{t+\tau}P_{t+\tau-s}F(s)\,
                                  \mathrm d\lambda_1(s).
 \end{aligned}
\]
Odwzorowanie $(q,s)\mapsto P_qG(s)$ jest
ciągłe w normie $X$ na zwartym prostokącie
$[0,T]\times[0,T]$, a zatem jednostajnie ciągłe.
Moduł w normie różnicy między pierwszym składnikiem
a $\kappa\int_0^tP_{t-s}G(s)\,\mathrm d\lambda_1(s)$
jest ograniczony przez $\kappa t$ razy supremum
$\|P_qG(s)-P_{t-s}G(s)\|_\infty$ dla
$0\leq s\leq t$ i $|q-(t-s)|\leq\tau$,
$q\in[0,T]$. Oszacowanie to dąży do zera.
Drugi składnik dąży w $X$ do $F(t)$,
ponieważ dla $s\in[t,t+\tau]$ mamy
$0\leq t+\tau-s\leq\tau$ oraz $|s-t|\leq\tau$.
Łączna ciągłość w $(r,s)$ uzasadnia tę granicę
również dla $t=0$.

Dla $0<t\leq T$ i $0<\tau\leq t$ analogicznie
rozpisujemy iloraz lewostronny:
\[
 \begin{aligned}
 \frac{V(t)-V(t-\tau)}{\tau}
 &=\frac{\kappa}{\tau}\int_0^{t-\tau}
      \int_{t-\tau-s}^{t-s}P_qG(s)\,
       \mathrm d\lambda_1(q)\,\mathrm d\lambda_1(s)\\
 &\quad+\frac1{\tau}\int_{t-\tau}^tP_{t-s}F(s)\,
                                   \mathrm d\lambda_1(s).
 \end{aligned}
\]
Pierwszy składnik ma tę samą granicę: oprócz
jednostajnej ciągłości używamy oszacowania brakującej
całki na $[t-\tau,t]$ przez
$\kappa\tau\sup_{s\in[0,T]}\|G(s)\|_\infty$.
Drugi składnik dąży do $F(t)$.
Otrzymujemy dla wszystkich $t\in[0,T]$
(z odpowiednimi pochodnymi jednostronnymi na końcach)
\[
 V'(t)=F(t)+\kappa\int_0^tP_{t-s}G(s)\,
                     \mathrm d\lambda_1(s)
      =F(t)+\kappa\Delta_xV(t).
\]
Pochodna ta jest ciągła w normie $X$ na $[0,T]$
dzięki wcześniej wykazanej ciągłości całek.
W szczególności $V'(0)=F(0)$.

Lemat~\ref{lem:pde-v31-heat-generator} daje również
$\frac{\mathrm d}{\mathrm dt}P_tg=\kappa P_t\Delta g$
w normie $X$ (dla $t=0$ prawostronnie).
Sumując tożsamości, dostajemy
$\partial_tu-\kappa\Delta_xu=f$ i $u(0,\cdot)=g$,
a w chwili początkowej
$\partial_tu(0,\cdot)=\kappa\Delta g+F(0)$.
Ciągłość w normie $X$ funkcji $u$, jej pochodnej
czasowej oraz pochodnych przestrzennych do rzędu
drugiego daje ich ciągłość łączną w $(t,x)$,
przez tę samą nierówność trójkąta, której użyto
w dowodzie twierdzenia~\ref{thm:pde4-cauchy}.
Stąd wynika deklarowana regularność.

Oszacowanie~\eqref{eq:pde4-duhamel-bound} wynika
z kontrakcyjności $P_r$ oraz nierówności trójkąta
dla całki w $X$. Prawa strona jest ograniczona
na $[0,T]$. Różnica dowolnych dwóch ograniczonych
rozwiązań z tym samym wymuszeniem i danymi
początkowymi spełnia jednorodne równanie ciepła
z zerowymi danymi; twierdzenie
\ref{thm:pde4-global-uniqueness} daje jej zerowanie.
\end{proof}

\begin{przyklad}[Źródło stałe w czasie]
Niech $T>0$, $n\geq1$, $\kappa>0$,
$f(t,x)=h(x)$ dla $h\in C_c^\infty(\R^n)$ oraz
$g=0$. Wzór~\eqref{eq:pde4-duhamel}, po podstawieniu
$r=t-s$, daje
\[
 u(t,x)=\int_0^tP_rh(x)\,\mathrm d\lambda_1(r).
\]
Ciągłość $r\mapsto P_rh$ w normie supremum daje
$\partial_tu=P_th$, także prawostronnie w $0$.
Przeniesienie laplasjanu pod całkę oraz na dane,
uzasadnione ograniczonością drugich pochodnych $h$,
prowadzi do
\[
 \kappa\Delta_xu=\kappa\int_0^tP_r\Delta h\,
                       \mathrm d\lambda_1(r)=P_th-h
\]
na mocy lematu~\ref{lem:pde-v31-heat-generator}.
Zatem $\partial_tu-\kappa\Delta_xu=h$,
$u(0,\cdot)=0$ i $\partial_tu(0,\cdot)=h$.
Ponadto $\|u(t,\cdot)\|_\infty\leq t\|h\|_\infty$.
\end{przyklad}

\section{Dyfuzja w obszarze ograniczonym: energia i widmo}
\label{sec:pde4-bounded}

Na obszarze ograniczonym wzór splotowy na całej
przestrzeni na ogół nie spełnia zadanych warunków
brzegowych. Warunki Dirichleta określają temperaturę
na brzegu, a warunki Neumanna --- jej pochodną
normalną, związaną ze strumieniem ciepła przez brzeg.
W tej sekcji wyprowadzimy podstawowe wnioski
dla jednorodnego warunku Dirichleta
\emph{przy założeniu istnienia dostatecznie gładkiego
rozwiązania}. Ogólną teorię istnienia oprzemy później
na przestrzeniach Sobolewa i półgrupach operatorów.

\begin{twierdzenie}[Bilans energii]
\label{thm:pde4-energy}
Niech $n\geq1$, $\Omega\subset\R^n$ będzie obszarem
ograniczonym o brzegu klasy $C^1$, $T>0$, $\kappa>0$.
Niech $u\in C^{1,2}([0,T]\times\overline\Omega)$
będzie funkcją rzeczywistą, spełniającą
$\partial_tu-\kappa\Delta_xu=0$ na $(0,T)\times\Omega$
oraz $u=0$ na $[0,T]\times\partial\Omega$.
Przez podaną klasę regularności rozumiemy, że funkcja
i wskazane pochodne przedłużają się ciągle na
domknięty cylinder, z jednostronnymi pochodnymi
czasowymi na końcach przedziału.
Wówczas dla $0\leq t\leq T$
\begin{equation}\label{eq:pde4-energy}
 \frac12\int_\Omega |u(t,x)|^2\,\mathrm d\lambda_n(x)
 +\kappa\int_0^t\!\int_\Omega|\nabla_xu(s,x)|^2
 \,\mathrm d\lambda_n(x)\,\mathrm d\lambda_1(s)
 =\frac12\int_\Omega|u(0,x)|^2\,\mathrm d\lambda_n(x).
\end{equation}
W szczególności $t\mapsto\|u(t,\cdot)\|_{L^2(\Omega)}$
jest nierosnąca. W powyższej klasie rozwiązanie
z zadanymi danymi początkowymi i zerowym warunkiem
Dirichleta jest jednoznaczne.
\end{twierdzenie}
\begin{proof}
Połóżmy $E(t)=\frac12\int_\Omega |u(t,x)|^2
\,\mathrm d\lambda_n(x)$. Na zwartym zbiorze
$[0,T]\times\overline\Omega$ funkcje $u$ i $\partial_tu$
są ciągłe i ograniczone. Ilorazy różnicowe
$\frac12u(t,x)^2$ względem czasu są zdominowane
przez $\sup|u|\sup|\partial_tu|$, a $\lambda_n(\Omega)$
jest skończona. Z twierdzenia o zbieżności zdominowanej
wynika zatem $E\in C^1([0,T])$, przy jednostronnych
pochodnych na końcach, oraz
$E'(t)=\int_\Omega u(t,x)\partial_tu(t,x)
\,\mathrm d\lambda_n(x)$.
Dla $0<t<T$ podstawiamy równanie ciepła
i stosujemy pierwszą tożsamość Greena
\eqref{eq:pde3-green1} dla ustalonego $t$:
\[
 \begin{aligned}
 E'(t)
 &=\kappa\int_\Omega u(t,x)\Delta_x u(t,x)
             \,\mathrm d\lambda_n(x)\\
 &=-\kappa\int_\Omega|\nabla_xu(t,x)|^2\,\mathrm d\lambda_n(x)\\
 &\quad+\kappa\int_{\partial\Omega}
 u(t,\zeta)\partial_\nu u(t,\zeta)\,\mathrm d\sigma(\zeta).
 \end{aligned}
\]
Tutaj $\nu$ jest jednostkową normalną zewnętrzną,
a $\sigma$ miarą powierzchniową.
W wymiarze $n=1$ ograniczony obszar jest przedziałem
$(a,b)$; ten sam rachunek wynika bezpośrednio
z jednowymiarowego całkowania przez części,
z wyrazem brzegowym
$\kappa[u(t,x)\partial_xu(t,x)]_{x=a}^{x=b}$.
W obu przypadkach wyraz brzegowy znika
dzięki warunkowi Dirichleta.

Funkcja $t\mapsto\int_\Omega|\nabla_xu(t,x)|^2
\,\mathrm d\lambda_n(x)$ jest ciągła na $[0,T]$,
co wynika z ciągłości i ograniczoności gradientu
na domkniętym cylindrze. Mamy więc
\[
 E'(t)=-\kappa\|\nabla_xu(t,\cdot)\|_{L^2(\Omega)}^2.
\]
Równość przedłuża się przez ciągłość do końców
przedziału. Całkując względem czasu od $0$ do $t$,
otrzymujemy~\eqref{eq:pde4-energy}.
Ponieważ $E'\leq0$, energia i norma rozwiązania
są nierosnące.

Jeżeli $u_1,u_2$ spełniają założenia twierdzenia
i mają te same dane początkowe, to $w=u_1-u_2$
spełnia te same warunki z zerowymi danymi.
Bilans energii daje $E_w(t)=0$ dla każdego $t$,
ponieważ obie całki po lewej stronie
są nieujemne, a prawa strona znika.
Gdyby $w(t,x_0)\ne0$ w pewnym $x_0\in\Omega$,
ciągłość dawałaby kulę wewnętrzną, na której
$|w(t,\cdot)|$ jest ograniczona z dołu przez
dodatnią stałą. Przeczyłoby to zerowej normie
$L^2(\Omega)$. Stąd $w=0$ w cylindrze, a na jego
brzegu przestrzennym również z warunku Dirichleta.
\end{proof}

\begin{lemat}[Elementarna nierówność Poincarégo w obszarze ograniczonym]
\label{lem:pde4-poincare}
Niech $n\geq1$, $L>0$ i $\Omega\subset(-L,L)^n$
będzie obszarem ograniczonym o brzegu klasy $C^1$.
Jeżeli $v\in C^1(\overline\Omega)$ jest rzeczywista
i $v|_{\partial\Omega}=0$, to
\begin{equation}\label{eq:pde4-poincare}
 \int_\Omega |v(x)|^2\,\mathrm d\lambda_n(x)
 \leq4L^2\int_\Omega|\nabla v(x)|^2\,\mathrm d\lambda_n(x).
\end{equation}
\end{lemat}
\begin{proof}
Przedłużamy $v$ przez zero do domkniętego sześcianu
$\overline Q=[-L,L]^n$, otrzymując funkcję $V$.
Jest ona ciągła: w punktach brzegu $\Omega$
wartości funkcji $v$ dążą do zera, a poza
$\overline\Omega$ przedłużenie jest lokalnie zerowe.
Wprowadzamy również mierzalną funkcję $W$ na $Q$,
równą $\partial_{x_1}v$ w $\Omega$ i zeru poza
$\Omega$. Jest ona ograniczona, gdyż pochodna
$v$ jest ciągła na zwartym zbiorze $\overline\Omega$.

Dla $n\geq2$ ustalmy $x'\in(-L,L)^{n-1}$
i rozważmy otwarty podzbiór prostej
$I_{x'}=\{s\in(-L,L):(s,x')\in\Omega\}$.
Jest on sumą co najwyżej przeliczalnej rodziny
parami rozłącznych przedziałów otwartych $(a_j,b_j)$.
Każdy ich koniec odpowiada punktowi $\partial\Omega$,
również gdy jest równy $-L$ lub $L$:
jest granicą punktów $\Omega$, ale nie należy
do $\Omega$. Na każdej takiej składowej funkcja
$s\mapsto v(s,x')$ jest klasy $C^1$, a jej wartości
na obu końcach są zerowe. Z twierdzenia podstawowego
rachunku całkowego, najpierw na podprzedziałach
zwartych w $(a_j,b_j)$, a następnie przez
ciągłość na końcach i ograniczoność pochodnej,
wynika
\[
 \int_{a_j}^{b_j}W(s,x')\,\mathrm d\lambda_1(s)=0,
 \qquad
 \int_{a_j}^{q}W(s,x')\,\mathrm d\lambda_1(s)=V(q,x')
 \quad(a_j<q<b_j).
\]
Dla dowolnego $q\in[-L,L]$ możemy rozłożyć całkę
po $[-L,q]$ na przecięcia z tymi składowymi,
ponieważ $W=0$ poza $I_{x'}$, a całka jest
bezwzględnie zbieżna. Wszystkie całe składowe
dają wkład zerowy; co najwyżej jedna składowa
zawiera $q$ i daje wkład $V(q,x')$.
Jeśli $q\notin I_{x'}$, wszystkie wkłady są zerowe,
zgodnie z definicją $V$. Otrzymujemy w każdym przypadku
\[
 V(q,x')=\int_{-L}^{q}W(s,x')\,\mathrm d\lambda_1(s).
\]
W szczególności funkcja $q\mapsto V(q,x')$ jest
absolutnie ciągła na $[-L,L]$ i ma pochodną
równą $W(\cdot,x')$ prawie wszędzie.
Argument nie wymaga, by przekrój miał skończenie
wiele składowych.

Z nierówności Cauchy'ego--Schwarza dostajemy
dla każdego $q\in[-L,L]$
\[
 |V(q,x')|^2\leq 2L
 \int_{-L}^L|W(s,x')|^2\,\mathrm d\lambda_1(s).
\]
Całkując względem $q$ na przedziale długości $2L$
i następnie względem $x'$ na $(-L,L)^{n-1}$,
na mocy twierdzenia Tonellego otrzymujemy
\[
 \int_Q |V|^2\,\mathrm d\lambda_n
 \leq4L^2\int_Q|W|^2\,\mathrm d\lambda_n.
\]
Dla $n=1$ stosujemy bezpośrednio ten sam argument
na przedziale $[-L,L]$, bez zmiennej $x'$.
Z definicji przedłużeń całki po $Q$ są równe
odpowiednim całkom po $\Omega$, a
$|\partial_{x_1}v|^2\leq|\nabla v|^2$.
To daje~\eqref{eq:pde4-poincare}.
\end{proof}

\begin{wniosek}[Wykładniczy zanik energii przy warunku Dirichleta]
\label{cor:pde4-exponential}
Niech $n\geq1$, $T,L,\kappa>0$ i niech
$\Omega\subset(-L,L)^n$ będzie obszarem ograniczonym
o brzegu klasy $C^1$.
Niech $u\in C^{1,2}([0,T]\times\overline\Omega)$
będzie rzeczywistym rozwiązaniem równania
$\partial_tu-\kappa\Delta_xu=0$ na $(0,T)\times\Omega$
z $u=0$ na $[0,T]\times\partial\Omega$.
Wówczas
\begin{equation}\label{eq:pde4-exponential}
 \|u(t,\cdot)\|_{L^2(\Omega)}\leq
 \exp\!\left(-\frac{\kappa t}{4L^2}\right)
 \|u(0,\cdot)\|_{L^2(\Omega)}\qquad(0\leq t\leq T).
\end{equation}
\end{wniosek}
\begin{proof}
Niech $E(t)=\frac12\|u(t,\cdot)\|_{L^2(\Omega)}^2$.
Dla każdego ustalonego $t$ funkcja $u(t,\cdot)$
spełnia założenia lematu~\ref{lem:pde4-poincare}.
Z bilansu energii i~\eqref{eq:pde4-poincare} wynika
\[
 E'(t)=-\kappa\|\nabla_xu(t,\cdot)\|_{L^2(\Omega)}^2
 \leq-\frac{\kappa}{4L^2}\|u(t,\cdot)\|_{L^2(\Omega)}^2
 =-\frac{\kappa}{2L^2}E(t).
\]
Wobec tego pochodna funkcji
$t\mapsto e^{\kappa t/(2L^2)}E(t)$ jest niedodatnia.
Porównując jej wartości w chwilach $0$ i $t$,
otrzymujemy $E(t)\leq e^{-\kappa t/(2L^2)}E(0)$.
Ponieważ $E\geq0$, mnożymy nierówność przez $2$
i pierwiastkujemy, dostając~\eqref{eq:pde4-exponential}.
Zanik wykładniczy dotyczy zarówno energii, jak
i normy $L^2$, ale wykładnik dla normy jest o połowę
mniejszy niż dla energii. Nie twierdzimy,
że stała otrzymana z elementarnej nierówności
Poincarégo jest optymalna.
\end{proof}

\begin{przyklad}[Mody własne na odcinku]
Niech $\Omega=(0,L)$, gdzie $L>0$, oraz $\kappa>0$.
Dla każdego całkowitego $m\geq1$ określamy
$\phi_m(x)=\sin(m\pi x/L)$ i $\lambda_m=(m\pi/L)^2$.
Mamy $\phi_m''=-\lambda_m\phi_m$ oraz
$\phi_m(0)=\phi_m(L)=0$.
Funkcja
\[
 u_m(t,x)=e^{-\kappa\lambda_m t}\phi_m(x)
\]
spełnia $\partial_tu_m=-\kappa\lambda_m u_m$
i $\partial_x^2u_m=-\lambda_m u_m$, zatem
rozwiązuje równanie ciepła, znika na brzegu
i przyjmuje w chwili $0$ dane $\phi_m$.
Ponieważ $\sin^2\theta=(1-\cos(2\theta))/2$,
mamy $\|\phi_m\|_{L^2(0,L)}^2=L/2$.
Stąd
\[
 \|u_m(t,\cdot)\|_{L^2(0,L)}^2
 =\frac L2 e^{-2\kappa\lambda_m t},\qquad
 \|\partial_xu_m(t,\cdot)\|_{L^2(0,L)}^2
 =\lambda_m\frac L2 e^{-2\kappa\lambda_m t}.
\]
Różniczkując połowę pierwszej równości względem
czasu, otrzymujemy dokładnie bilans energii
z twierdzenia~\ref{thm:pde4-energy}.
W szczególności większe $\lambda_m$ odpowiadają
szybszemu zanikowi amplitudy.

Z liniowości równania dla każdego $N\geq1$
i $c_1,\ldots,c_N\in\R$ funkcja
$\sum_{m=1}^N c_m u_m$ również jest rozwiązaniem,
z danymi początkowymi $\sum_{m=1}^N c_m\phi_m$.
Jego jednoznaczność wynika z bilansu energii.
Ogólne rozwinięcia w szereg funkcji własnych
laplasjanu wymagają ponadto uzasadnienia
zbieżności i możliwości różniczkowania szeregu;
wrócimy do nich po wprowadzeniu teorii Fouriera.
\end{przyklad}

\section{Zachowanie dla dużych czasów i nieodwracalność}
\label{sec:pde4-large-time}

\begin{twierdzenie}[Zanik na całej przestrzeni]
\label{thm:pde4-large-time}
Niech $n\geq1$, $\kappa>0$.
Jeśli $g\in L^1(\R^n)$, to $\|P_tg\|_\infty\to0$
przy $t\to\infty$.
Jeśli $g\in L^2(\R^n)$, to $\|P_tg\|_2\to0$
przy $t\to\infty$.
\end{twierdzenie}
\begin{proof}
Dla $g\in L^1(\R^n)$ pierwsze oszacowanie
w~\eqref{eq:pde4-L1Linf} daje
\[
 0\leq\|P_tg\|_\infty\leq
 (4\pi\kappa t)^{-n/2}\|g\|_1\longrightarrow0,
\]
ponieważ $n\geq1$ i $\kappa>0$.
Jeżeli $g\in L^1(\R^n)\cap L^2(\R^n)$, drugie
oszacowanie w~\eqref{eq:pde4-L1Linf} daje
\[
 0\leq\|P_tg\|_2\leq
 (8\pi\kappa t)^{-n/4}\|g\|_1\longrightarrow0.
\]

Niech teraz $g\in L^2(\R^n)$ będzie dowolna.
Ustalmy $\eps>0$. Z gęstości $C_c^\infty(\R^n)$
w $L^2(\R^n)$ wybieramy $h\in C_c^\infty(\R^n)$
takie, że $\|g-h\|_2<\eps/2$.
Funkcja $h$ należy również do $L^1(\R^n)$,
więc istnieje $T_\eps>0$ takie, że
$\|P_th\|_2<\eps/2$ dla wszystkich $t\geq T_\eps$.
Z liniowości i kontrakcyjności $P_t$ w $L^2$,
wykazanej w twierdzeniu~\ref{thm:pde4-approx-identity},
otrzymujemy dla $t\geq T_\eps$
\[
 \|P_tg\|_2
 \leq\|P_t(g-h)\|_2+\|P_th\|_2
 \leq\|g-h\|_2+\|P_th\|_2<\eps.
\]
Dowolność $\eps$ daje żądaną zbieżność.
Argument wymaga tylko aproksymacji danych w $L^2$,
nie zakłada ich całkowalności w $L^1$.
\end{proof}

Należy odróżnić silny zanik dla każdego ustalonego
$g\in L^2$ od zaniku normy operatorowej półgrupy
na $L^2$. Powyższy dowód nie daje wspólnego tempa
zaniku dla wszystkich danych o normie $1$.
Istotnie, ustalmy niezerową funkcję
$h\in C_c^\infty(\R^n)$ z $\|h\|_2=1$ i dla $R>0$
połóżmy $h_R(x)=R^{-n/2}h(x/R)$.
Zmiana zmiennych $x=Ry$ daje $\|h_R\|_2=1$.
Ze skalowania jądra i podstawienia $y=Rz$ w splocie
wynika
\[
 (P_th_R)(x)=R^{-n/2}(P_{t/R^2}h)(x/R).
\]
Zatem $\|P_th_R\|_2=\|P_{t/R^2}h\|_2\to1$
przy $R\to\infty$ dla każdego ustalonego $t>0$,
na mocy silnej ciągłości w chwili zero.
W połączeniu z kontrakcyjnością daje to
$\|P_t\|_{L^2\to L^2}=1$ dla każdego $t\geq0$.
Na całej przestrzeni nie otrzymujemy więc
jednolitego wykładniczego zaniku analogicznego
do oszacowania~\eqref{eq:pde4-exponential}
na ograniczonym obszarze z warunkiem Dirichleta.

Zanik nie musi zachodzić także dla ogólnych danych
ograniczonych: dla $g\equiv1$ normalizacja jądra
daje $P_tg\equiv1$. Z kolei dla nieujemnego
$g\in L^1(\R^n)$, $g\not=0$ w $L^1$, norma $L^1$
pozostaje stała z własności zachowania masy
\eqref{eq:pde4-mass}. Zanik normy supremum
dla takich danych nie oznacza zaniku całkowitej masy.

\begin{kontrprzyklad}[Wsteczne równanie ciepła jest niestabilne]
\label{ex:pde4-backward}
Niech $\Omega=(0,\pi)$, $\kappa=1$ i $T>0$.
Rozważmy przestrzeń $S$ skończonych kombinacji
liniowych funkcji $\phi_m(x)=\sin(mx)$,
gdzie $m\geq1$ jest liczbą całkowitą.
Z przykładu modów własnych oraz jednoznaczności
w twierdzeniu~\ref{thm:pde4-energy} wynika, że
odwzorowanie ewolucji od chwili $0$ do $T$ na $S$
jest dobrze określone i ma postać
\[
 A_T\left(\sum_{m=1}^N c_m\phi_m\right)
 =\sum_{m=1}^N e^{-m^2T}c_m\phi_m.
\]
Funkcje $\phi_m$ są liniowo niezależne:
z tożsamości
$2\sin(mx)\sin(kx)=\cos((m-k)x)-\cos((m+k)x)$
otrzymujemy dla $m\ne k$
$\int_0^\pi\phi_m(x)\phi_k(x)\,\mathrm d\lambda_1(x)=0$,
a $\|\phi_m\|_{L^2(0,\pi)}^2=\pi/2$.
Ponieważ wszystkie mnożniki $e^{-m^2T}$ są dodatnie,
$A_T:S\to S$ jest liniową bijekcją.

Dla danych początkowych $\phi_m$ mamy
$u_m(t,x)=e^{-m^2t}\sin(mx)$.
Stąd
\[
 \|\phi_m\|_\infty=1,\qquad
 \|A_T\phi_m\|_\infty=e^{-m^2T}\longrightarrow0.
\]
Jeżeli $v_m=A_T\phi_m$, to $v_m\to0$ w normie
supremum, natomiast $A_T^{-1}v_m=\phi_m$
nie dąży do zera w tej normie.
Odwrotne odwzorowanie $A_T^{-1}:S\to S$
jest więc nieciągłe w zerze względem normy supremum.

Również w normie $L^2(0,\pi)$ mamy
\[
 \|v_m\|_2=e^{-m^2T}\sqrt{\pi/2}\longrightarrow0,
 \qquad \|A_T^{-1}v_m\|_2=\sqrt{\pi/2}.
\]
W obu normach iloraz
$\|A_T^{-1}v_m\|/\|v_m\|=e^{m^2T}$ jest
nieograniczony. Nawet ograniczenie do dziedziny $S$
nie pozwala więc uzyskać ciągłego odwracania.
Każde odwrotne odwzorowanie na większej klasie
danych końcowych, zgodne z odwracaniem tych modów,
musiałoby ograniczać się do $A_T^{-1}$ na $S$,
a zatem również nie mogłoby być ciągłe w zerze
w żadnej z tych dwóch norm.

Jest to przykład niestabilności w sensie Hadamarda:
nawet w klasie, w której istnienie i jednoznaczność
odwrócenia są zapewnione, ciągła zależność danych
początkowych od końcowych zawodzi.
Nie wynika stąd, że dwa różne dane z $S$
prowadzą do tego samego profilu końcowego.
Termin nieodwracalność oznacza tu brak stabilnego
odwracania w rozpatrywanych normach, nie brak
algebraicznej odwrotności na obrazie ewolucji.
\end{kontrprzyklad}

\section{Nieoczekiwana niejednoznaczność: przykład Tychonoffa}
\label{sec:pde4-tychonoff}

Powyższa jednoznaczność wymagała ograniczoności
rozwiązania w całym cylindrze czasoprzestrzennym.
Pokażemy teraz, że nie jest to wyłącznie techniczny
warunek dowodu.

\begin{twierdzenie}[Istnienie niezerowego rozwiązania o zerowych danych]
\label{thm:pde4-tychonoff}
Istnieje niezerowa funkcja
$u\in C^\infty([0,\infty)\times\R)$
spełniająca $u_t=u_{xx}$ dla $t>0$, $x\in\R$,
przy czym $u(0,x)=0$ dla wszystkich $x\in\R$.
Funkcja ta nie może być ograniczona na żadnym
cylindrze $[0,T]\times\R$ z $T>0$.
\end{twierdzenie}
\begin{proof}
Niech $h(t)=e^{-1/t^2}$ dla $t>0$ i $h(t)=0$
dla $t\leq0$. Na dodatniej półosi funkcja $h$
jest klasy $C^\infty$. Najpierw definiujemy formalnie
szereg, a następnie uzasadnimy jego zbieżność
i gładkość sumy także przy $t=0$. Kładziemy
\begin{equation}\label{eq:pde4-tychonoff}
 u(t,x)=\sum_{k=0}^\infty\frac{h^{(k)}(t)}{(2k)!}
 x^{2k}\quad(t>0),\qquad u(0,x)=0.
\end{equation}
Uzasadnijmy najpierw oszacowanie pochodnych $h$.
Dla $t>0$ i $z=a+ib$ spełniającego $|z-t|\leq t/4$
mamy $a\geq3t/4$, $|b|\leq t/4$ oraz $|z|\leq5t/4$.
Wobec tego
\[
 \operatorname{Re}\frac1{z^2}
 =\frac{a^2-b^2}{|z|^4}
 \geq\frac{t^2/2}{(5t/4)^4}
 =\frac{128}{625t^2}.
\]
Możemy przyjąć $c=128/625>0$, niezależne od $t$.
Funkcja $z\mapsto e^{-1/z^2}$ jest holomorficzna
na otoczeniu domkniętego dysku o środku $t$
i promieniu $t/4$, gdyż dysk nie zawiera zera.
Na jego brzegu moduł tej funkcji nie przekracza
$e^{-c/t^2}$. Oszacowanie Cauchy'ego na tym okręgu
daje dla każdego całkowitego $k\geq0$
\begin{equation}\label{eq:pde4-tychonoff-est}
 |h^{(k)}(t)|\leq k!\left(\frac4t\right)^k
 e^{-c/t^2},\qquad k\geq0.
\end{equation}
Ponieważ $\binom{2k}{k}\geq1$, mamy $(2k)!\geq(k!)^2$.
Dla każdego $A>0$ i $|x|\leq A$
otrzymujemy
\[
 |u(t,x)|\leq e^{-c/t^2}
 \sum_{k=0}^\infty\frac{(4A^2/t)^k}{k!}
 =\exp\!\left(-\frac c{t^2}+\frac{4A^2}t\right).
\]
Majoranta wykładnicza jest skończona dla każdego
$t>0$, więc szereg definiujący $u$ jest bezwzględnie
zbieżny. Oszacowanie dowodzi także jednostajnej
zbieżności $u(t,\cdot)$ do zera na $[-A,A]$
przy $t\downarrow0$.

Uzasadnijmy dokładniej oszacowanie wszystkich pochodnych.
Dla ustalonych $p,q\in\mathbb N_0$ oraz $2k\geq q$ mamy
\[
 \partial_t^p\partial_x^q
 \left(\frac{h^{(k)}(t)x^{2k}}{(2k)!}\right)
 =\frac{h^{(k+p)}(t)x^{2k-q}}{(2k-q)!}.
\]
Jeżeli $2k<q$, pochodna jest identycznie zerowa.
Wybierzmy $B=\max\{1,A\}$ dla $|x|\leq A$.
Oszacowanie~\eqref{eq:pde4-tychonoff-est} daje
\[
 \left|\frac{h^{(k+p)}(t)x^{2k-q}}{(2k-q)!}\right|
 \le e^{-c/t^2}\left(\frac4t\right)^p
 \frac{(k+p)!}{(2k-q)!}
 \left(\frac{4B^2}{t}\right)^k.
\]
Dla $k\geq1$ i $2k\geq q$ mamy
\[
 \frac{(k+p)!}{(2k-q)!}
 =\frac{(k+p)!}{k!}\frac{k!}{(2k)!}
   \frac{(2k)!}{(2k-q)!}
 \leq\frac{(k+p)^p(2k)^q}{k!}
 \leq C_{p,q}\frac{(1+k)^{p+q}}{k!},
\]
gdzie pusty iloczyn ma wartość $1$.
Dla $k=0$ dopuszczalny jest tylko przypadek $q=0$;
zwiększamy wtedy stałą tak, by obejmowała $p!$.

Dla $m\in\mathbb N_0$ i $z\geq0$ mamy
$\sum_{k\geq0}(1+k)^m z^k/k!=Q_m(z)e^z$,
gdzie $Q_0=1$ oraz
$Q_{m+1}=(1+z)Q_m+zQ_m'$.
Istotnie, stosujemy operator $1+z\partial_z$
do szeregu, co jest dozwolone, gdyż ma on
nieskończony promień zbieżności (z kryterium ilorazowego).
Rekurencja pokazuje, że $Q_m$ ma stopień $m$
i nieujemne współczynniki, więc
$Q_m(z)\leq C_m(1+z)^m$ dla $z\geq0$.
Stąd
cały szereg pochodnych jest ograniczony na
$0<t\le1$, $|x|\le A$ przez
\[
 C_{p,q,A}t^{-p}(1+t^{-1})^{p+q}
 \exp\left(-\frac c{t^2}+\frac{4B^2}{t}\right).
\]
Najpierw uzasadnijmy różniczkowanie dla $t>0$.
Na prostokącie $[a,b]\times[-A,A]$, $0<a<b$,
każdy składnik szeregu pochodnych jest ograniczony
przez
\[
 C_{p,q}(4/a)^p
 \frac{(1+k)^{p+q}}{k!}(4B^2/a)^k .
\]
Szereg tych majorant jest zbieżny z kryterium
ilorazowego. Kryterium Weierstrassa daje więc
jednostajną zbieżność każdego szeregu pochodnych
na takim prostokącie. Twierdzenie o różniczkowaniu
szeregów funkcyjnych, stosowane kolejno na mniejszych
otwartych prostokątach, daje $u\in C^\infty$
dla $t>0$ i identyfikuje każdą jej pochodną
$\partial_t^p\partial_x^q u$ z odpowiednim szeregiem.

Sprawdźmy teraz gładkość w chwili zero.
Dla dostatecznie małych $t>0$, konkretnie
$t\leq c/(8B^2)$, mamy
$-c/t^2+4B^2/t\leq-c/(2t^2)$.
Dla każdego $N\geq0$ i ustalonych $p,q,A$
otrzymujemy z wcześniejszego oszacowania
\[
 \sup_{|x|\leq A}|\partial_t^p\partial_x^q u(t,x)|
 =o(t^N)\qquad(t\downarrow0),
\]
gdyż wykładniczy czynnik $e^{-c/(2t^2)}$
zanika szybciej od każdej potęgi $t$.

Przedłużmy $u$ przez zero na $t<0$ i oznaczmy
dla każdego $p,q$ przez $U_{p,q}$ funkcję równą
$\partial_t^p\partial_x^q u$ dla $t>0$ oraz zeru
dla $t\leq0$. Każda $U_{p,q}$ jest ciągła
na $\R^2$ dzięki lokalnie jednostajnemu oszacowaniu
z $N=0$. W punkcie $(0,x)$ iloraz różnicowy
$U_{p,q}$ względem czasu od lewej jest zerowy,
a od prawej dąży do zera na mocy oszacowania
z $N=1$. Względem $x$ pochodna w chwili zero
jest zerowa, ponieważ $U_{p,q}(0,\cdot)\equiv0$.
Stąd na całej $\R^2$
\[
 \partial_tU_{p,q}=U_{p+1,q},\qquad
 \partial_xU_{p,q}=U_{p,q+1}.
\]
Indukcyjnie wynika, że $U_{0,0}$ jest klasy
$C^\infty(\R^2)$ i wszystkie jej pochodne
znikają na $\{0\}\times\R$.
W szczególności jej ograniczenie do $t\geq0$
ma wymaganą gładkość.

Dla $t>0$ różniczkowanie wyraz po wyrazie daje
\[
 \begin{aligned}
 u_t(t,x)&=\sum_{k=0}^\infty
     \frac{h^{(k+1)}(t)}{(2k)!}x^{2k},\\
 u_{xx}(t,x)&=\sum_{k=1}^\infty
     \frac{h^{(k)}(t)}{(2k-2)!}x^{2k-2}\\
 &=\sum_{j=0}^\infty\frac{h^{(j+1)}(t)}{(2j)!}x^{2j}.
 \end{aligned}
\]
Zatem $u_t=u_{xx}$, również przy $t=0$,
gdzie obie pochodne znikają.
Wreszcie $u(t,0)=h(t)>0$ dla każdego $t>0$,
więc nie jest to rozwiązanie zerowe.
Gdyby było ograniczone na pewnym cylindrze $[0,T]\times\R$,
twierdzenie~\ref{thm:pde4-global-uniqueness}
wymusiłoby $u=0$, co jest sprzeczne z $u(t,0)>0$.
\end{proof}

\begin{uwaga}[Co ten kontrprzykład rzeczywiście oznacza?]
Twierdzenie~\ref{thm:pde4-global-uniqueness}
zapewnia jednoznaczność w klasie rozwiązań
ograniczonych na każdym pasku $[0,T]\times\R^n$.
Istnieją również twierdzenia o jednoznaczności
przy innych warunkach wzrostu, których tutaj
nie formułujemy. Przykład Tychonoffa pokazuje,
że sama gładkość na domkniętym pasku i zerowe
dane początkowe nie wystarczają.
Nie wynika z niego, że każde rozwiązanie
nieograniczone prowadzi do niejednoznaczności.
Ponieważ rozwiązanie i wszystkie jego pochodne
znikają w chwili zero, a $u(t,0)>0$ dla $t>0$,
nie jest ono analityczne względem czasu
w otoczeniu chwili początkowej.
Nie przeczy to twierdzeniu o analityczności
przestrzennej splotu z jądrem Gaussa.
\end{uwaga}

\section{Podsumowanie i związek z kolejnym rozdziałem}
\label{sec:pde4-summary}

Kluczowe mechanizmy równania ciepła są widoczne
w~\eqref{eq:pde4-kernel}: dodatniość jądra daje
zachowanie porządku dla rozwiązań splotowych
i natychmiastową propagację niezerowych danych
nieujemnych o zwartym nośniku. Normalizacja
zapewnia zachowanie całki dla danych z $L^1$,
a skalowanie Gaussa wyjaśnia oszacowania zaniku.
Półgrupa jest silnie ciągła na $L^p$ dla $p<\infty$
($p\geq1$) oraz na $\mathrm{BUC}$, lecz nie na
całym $L^\infty$.

Jawny wzór rozwiązuje jednorodne zagadnienie
Cauchy'ego z danymi w $\mathrm{BUC}$, a wzór
Duhamela obejmuje niejednorodne równanie przy
podanych założeniach gładkości danych i wymuszenia.
Jednoznaczność klasycznych rozwiązań na całej
przestrzeni wymaga określenia klasy rozwiązań:
dla funkcji ograniczonych na każdym pasku
czasowym zapewnia ją zasada maksimum, natomiast bez
kontroli wzrostu przykład Tychonoffa daje
niejednoznaczność.

Dla obszarów ograniczonych wyprowadziliśmy bilans
energii i wykładniczy zanik normy przy jednorodnym
warunku Dirichleta, zakładając istnienie odpowiednio
gładkiego rozwiązania. Nie stanowi to jeszcze
ogólnego twierdzenia o istnieniu.
Na całej przestrzeni silny zanik w $L^2$ nie oznacza
zaniku normy operatorowej półgrupy. Wygładzanie
przestrzenne oraz tłumienie wysokich częstości
prowadzą do niestabilności odtwarzania danych
początkowych; nie oznacza ona braku algebraicznej
odwrotności na przestrzeni skończonych sum modów.
W następnym rozdziale porównamy te fakty z równaniem
falowym, gdzie propagacja ma skończoną prędkość,
a klasyczna energia ma inny charakter.


% Rozdział 5. Równanie falowe.

\chapter[Równanie falowe]{Równanie falowe: charakterystyki, energia i propagacja}
\label{chap:pde-wave}

Równanie ciepła wygładza ograniczone dane początkowe,
a niezerowe dane nieujemne o zwartym nośniku prowadzą
do rozwiązania dodatniego w każdym punkcie przestrzeni
w każdej dodatniej chwili. Równanie falowe
zachowuje się inaczej: zaburzenie nie przekracza określonej
prędkości, a oscylacje nie są samoczynnie tłumione. W tym
rozdziale wyprowadzimy jawne wzory, udowodnimy jednoznaczność
na podstawie energii i rozróżnimy słabą oraz silną zasadę
Huygensa.

Ustalmy całkowite $n\geq1$, $c>0$ i oznaczmy
$\square_c u=\partial_t^2u-c^2\Delta_xu$. Zagadnienie
Cauchy'ego ma postać
\begin{equation}\label{eq:pde5-wave}
 \square_cu=0\quad\text{dla }t>0,\ x\in\R^n,\qquad
 u(0,x)=g(x),\quad u_t(0,x)=h(x).
\end{equation}
Dwa warunki początkowe są potrzebne, ponieważ równanie
jest drugiego rzędu względem czasu. Funkcje skalarne
są w tym rozdziale rzeczywiste, o ile nie zaznaczono
inaczej. Przez $\lambda_k$ oznaczamy miarę Lebesgue'a
w $\R^k$ dla $k\geq1$, a przez $\sigma$ standardową
miarę powierzchniową na rozważanej gładkiej
hiperpowierzchni; w wymiarze jeden na brzegu
przedziału jest to miara licząca.
Przez $C^2$ na domkniętym pasku czasowym rozumiemy
funkcje o ciągłych pochodnych do rzędu drugiego
także na jego końcach, gdzie pochodne czasowe
są jednostronne. Wszystkie
rachunki z całkowaniem i różniczkowaniem wykonujemy
najpierw dla danych dostatecznie gładkich; słabe
rozwiązania wprowadzimy dopiero w kolejnych częściach.

\section{Równanie w jednym wymiarze: rozkład na fale biegnące}
\label{sec:pde5-dalembert}

Rozważmy $u_{tt}-c^2u_{xx}=0$ dla $(t,x)\in(0,\infty)\times\R$.
Operator rozkłada się na czynniki
$\partial_t^2-c^2\partial_x^2
=(\partial_t-c\partial_x)(\partial_t+c\partial_x)$.
Wprowadzamy współrzędne charakterystyczne
$\xi=x-ct$, $\eta=x+ct$. Odwzorowanie
$(t,x)\mapsto(\xi,\eta)$ jest odwracalne, a jego
odwrotność dana jest wzorami
$t=(\eta-\xi)/(2c)$ oraz $x=(\eta+\xi)/2$.
Półpłaszczyźnie $t>0$ odpowiada zbiór
$D=\{(\xi,\eta)\in\R^2:\xi<\eta\}$.
Jeżeli $U(\xi,\eta)=u((\eta-\xi)/(2c),(\eta+\xi)/2)$,
to z reguły łańcuchowej
\[
 \partial_tu=-cU_\xi+cU_\eta,\qquad
 \partial_xu=U_\xi+U_\eta,\qquad
 u_{tt}-c^2u_{xx}=-4c^2U_{\xi\eta}.
\]
W ostatniej równości używamy przemienności
pochodnych mieszanych dla funkcji klasy $C^2$.

Wykażmy globalną, a nie tylko lokalną reprezentację
rozwiązania na $D$. Z $U_{\xi\eta}=0$ wynika, że
dla każdego ustalonego $\xi$ funkcja
$\eta\mapsto U_\xi(\xi,\eta)$ jest stała na spójnym
przedziale $(\xi,\infty)$. Podobnie dla ustalonego
$\eta$ funkcja $\xi\mapsto U_\eta(\xi,\eta)$
jest stała na $(-\infty,\eta)$.
Definiujemy $a(\xi)=U_\xi(\xi,\xi+1)$ oraz
$b(\eta)=U_\eta(\eta-1,\eta)$. Funkcje $a,b$
są klasy $C^1$ na $\R$ i na $D$ zachodzi
$U_\xi(\xi,\eta)=a(\xi)$, $U_\eta(\xi,\eta)=b(\eta)$.
Wybierając funkcje pierwotne $F'=a$ i $G'=b$,
otrzymujemy $\nabla(U-F-G)=0$ na $D$, gdzie
$F=F(\xi)$ i $G=G(\eta)$. Zbiór $D$ jest wypukły,
więc różnica jest stała: jej pochodna wzdłuż
dowolnego odcinka w $D$ znika.
Włączając tę stałą do $F$, dostajemy
\[
 u(t,x)=F(x-ct)+G(x+ct),\qquad F,G\in C^2(\R).
\]
Jeżeli $u$ jest klasy $C^2$ także przy $t=0$,
reprezentacja i wzory na pochodne przedłużają się
do $t=0$ przez ciągłość. Ten argument uzasadnia
reprezentację na całej półpłaszczyźnie, potrzebną
w dowodzie jednoznaczności.

W rozdziale wprowadzającym otrzymaliśmy już wzór d'Alemberta
(twierdzenie~\ref{thm:dalembert}). Wracamy do niego tutaj,
aby powiązać reprezentację z charakterystykami, energią i
obszarem zależności. Powtórzony krótki dowód ma charakter
samodzielnego przypomnienia, nie stanowi odrębnego twierdzenia
o szerszych założeniach.

\begin{twierdzenie}[Wzór d'Alemberta]\label{thm:pde5-dalembert}
Niech $c>0$, $g\in C^2(\R)$ i $h\in C^1(\R)$.
Istnieje dokładnie jedno rozwiązanie
$u\in C^2([0,\infty)\times\R)$ zagadnienia
\eqref{eq:pde5-wave} dla $n=1$. Jest nim
\begin{equation}\label{eq:pde5-dalembert}
 u(t,x)=\frac{g(x-ct)+g(x+ct)}2
 +\frac1{2c}\int_{x-ct}^{x+ct}h(y)\,\mathrm d\lambda_1(y).
\end{equation}
\end{twierdzenie}
\begin{proof}
Najpierw wyprowadzimy konieczną postać rozwiązania.
Na mocy globalnej reprezentacji powyżej każde
rozwiązanie ma postać $F(x-ct)+G(x+ct)$.
Warunki początkowe i ciągłość pochodnych przy $t=0$
dają $F+G=g$ oraz $-cF'+cG'=h$ na $\R$.
Różniczkując pierwszą równość i rozwiązując układ,
dostajemy $F'=(g'-h/c)/2$ oraz $G'=(g'+h/c)/2$.
Niech $H(x)=\int_0^x h(y)\,\mathrm d\lambda_1(y)$,
gdzie dla $x<0$ całkę rozumiemy z przeciwną orientacją.
Ponieważ $h\in C^1(\R)$, mamy $H\in C^2(\R)$
oraz $H'=h$. Zatem
\[
 F(x)=\tfrac12g(x)-\tfrac1{2c}H(x)+A,\qquad
 G(x)=\tfrac12g(x)+\tfrac1{2c}H(x)+B.
\]
Warunek $F+G=g$ wymusza $A+B=0$.
Po wstawieniu argumentów $x-ct$ i $x+ct$
otrzymujemy~\eqref{eq:pde5-dalembert}.
W szczególności każde rozwiązanie musi mieć
tę samą postać, co daje jednoznaczność.

Pozostaje wykazać istnienie. Wzór po prawej stronie
\eqref{eq:pde5-dalembert} definiuje funkcję klasy
$C^2$ (nawet dla wszystkich $t\in\R$, przy
orientowanej całce). Oznaczmy $a=x-ct$, $b=x+ct$.
Z reguły łańcuchowej i twierdzenia podstawowego
rachunku całkowego otrzymujemy
\[
 \begin{aligned}
 u_t&=\tfrac c2(-g'(a)+g'(b))+\tfrac12(h(a)+h(b)),\\
 u_x&=\tfrac12(g'(a)+g'(b))+\tfrac1{2c}(h(b)-h(a)),\\
 u_{tt}&=\tfrac{c^2}2(g''(a)+g''(b))
                   +\tfrac c2(h'(b)-h'(a)),\\
 u_{xx}&=\tfrac12(g''(a)+g''(b))
                   +\tfrac1{2c}(h'(b)-h'(a)).
 \end{aligned}
\]
Stąd $u_{tt}=c^2u_{xx}$. W chwili $t=0$ mamy
$a=b=x$, więc $u(0,x)=g(x)$ oraz $u_t(0,x)=h(x)$.
Wzór istotnie daje rozwiązanie zagadnienia Cauchy'ego.
Dowód nie wymaga ograniczoności ani kontroli wzrostu
danych lub rozwiązania w nieskończoności.
\end{proof}

\begin{uwaga}[Dokładny obszar zależności]
Ze wzoru wynika, że wartość $u(t,x)$ zależy jedynie
od $g$ w dwóch punktach $x-ct,x+ct$ i od $h$ na
odcinku $[x-ct,x+ct]$. W szczególności dane z punktów
poza tym odcinkiem nie wpływają na $u(t,x)$.
Dokładniej, dla $t>0$ zmiana wartości $g$ poza
końcami odcinka nie zmienia $u(t,x)$, natomiast
zmiana $h$ o funkcję gładką, nieujemną, niezerową
i o nośniku zawartym w dowolnym niepustym
podprzedziale $(x-ct,x+ct)$ zmienia tę wartość
o dodatnią liczbę. Nie oznacza to więc, że
sygnał jest odczuwany wyłącznie na czole fali:
człon z $h$ zależy także od danych we wnętrzu
podstawy stożka wstecznego.
\end{uwaga}

\begin{przyklad}[Dwa przesuwające się impulsy]
Jeżeli $h=0$ i $g\in C_c^2(\R)$, to
$u(t,x)=\frac12g(x-ct)+\frac12g(x+ct)$.
Początkowy impuls rozdziela się na dwa profile
biegnące w przeciwnych kierunkach bez zmiany kształtu.
Jeżeli natomiast $g=0$, $h\in C_c^1(\R)$,
$h\geq0$ i $h\not\equiv0$, to $\int_\R h\,\mathrm d\lambda_1>0$:
z ciągłości funkcja jest dodatnia i ograniczona
z dołu przez dodatnią stałą na pewnym przedziale.
Wybierzmy $R>0$ takie, że $\operatorname{supp}h\subset[-R,R]$.
Dla ustalonego $x$ oraz $t\geq(|x|+R)/c$ mamy
$[-R,R]\subset[x-ct,x+ct]$, a więc
$u(t,x)=(2c)^{-1}\int_\R h(y)\,\mathrm d\lambda_1(y)>0$.
Ta druga sytuacja pokazuje, dlaczego silna zasada
Huygensa nie zachodzi w wymiarze jeden.
\end{przyklad}

\begin{center}
\begin{tikzpicture}[scale=.88,>=Stealth]
  \fill[blue!8] (-2.7,0)--(2.7,0)--(0,3.35)--cycle;
  \draw[->] (-4,0)--(4,0) node[right] {$x$};
  \draw[->] (0,-.25)--(0,3.9) node[above] {$t$};
  \draw[blue!65!black,thick] (-2.7,0)--(0,3.35)--(2.7,0);
  \fill (0,3.35) circle (1.8pt) node[above right] {$(t_0,x_0)$};
  \draw[<->,blue!50!black] (-2.7,-.3)--(2.7,-.3)
    node[midway,below] {$[x_0-ct_0,x_0+ct_0]$};
  \node[blue!65!black] at (0,1.05) {stożek wsteczny};
  \node at (3.4,3.2) {$x_0=0$};
\end{tikzpicture}
\par\smallskip\small
Obszar zależności w jednym wymiarze; osie i wierzchołek
narysowano dla $x_0=0$, a szerokość stożka odpowiada $c t_0$.
\end{center}

\section{Energia i jednoznaczność w wielu wymiarach}
\label{sec:pde5-energy}

\begin{definicja}[Energia falowa]
Niech $T>0$, $c>0$ i $u\in C^2([0,T]\times\R^n)$.
Jeżeli dla ustalonego $t\in[0,T]$ pochodne
$u_t(t,\cdot)$ oraz $\nabla_xu(t,\cdot)$ są
całkowalne z kwadratem, definiujemy energię
w tej chwili przez
\begin{equation}\label{eq:pde5-energy}
 E_u(t)=\frac12\int_{\R^n}
 \bigl(|u_t(t,x)|^2+c^2|\nabla_xu(t,x)|^2\bigr)
 \,\mathrm d\lambda_n(x).
\end{equation}
\end{definicja}

\begin{twierdzenie}[Zachowanie energii]\label{thm:pde5-energy}
Niech $c,T>0$, $n\ge1$, a $u\in C^2([0,T]\times\R^n)$
spełnia $\square_cu=0$. Załóżmy, że istnieje
zwarty zbiór $K\subset\R^n$, poza którym
$u_t(t,\cdot)=\nabla_xu(t,\cdot)=0$
dla każdego $t\in[0,T]$. Wówczas $E_u(t)=E_u(0)$
dla $0\leq t\leq T$.
\end{twierdzenie}
\begin{proof}
Wybierzmy $R>0$ tak, aby $K\subset B(0,R)$,
i oznaczmy $B=B(0,R)$.
Dla każdego $t\in[0,T]$ całkę definiującą energię
można ograniczyć do $B$, ponieważ jej funkcja
podcałkowa znika poza $K$.
Wszystkie pochodne występujące w różniczkowaniu
tej funkcji względem czasu są ciągłe i ograniczone
na zwartym zbiorze $[0,T]\times\overline B$.
Z twierdzenia o różniczkowaniu pod całką,
uzasadnionego przez zbieżność zdominowaną na
zbiorze skończonej miary, wynika
\[
 E_u'(t)=\int_B
 \bigl(u_tu_{tt}+c^2\nabla_xu\cdot\nabla_xu_t\bigr)
 \,\mathrm d\lambda_n(x).
\]
W szczególności energia jest klasy $C^1$ względem
czasu, z pochodnymi jednostronnymi na końcach.
Z równości
$\diver(u_t\nabla_xu)=\nabla_xu_t\cdot\nabla_xu+
u_t\Delta_xu$ oraz twierdzenia Gaussa otrzymujemy
\[
 \begin{aligned}
 E_u'(t)
 &=\int_B u_t(u_{tt}-c^2\Delta_xu)
      \,\mathrm d\lambda_n(x)\\
 &\quad+c^2\int_{\partial B}u_t\partial_\nu u
      \,\mathrm d\sigma=0.
 \end{aligned}
\]
Pierwsza całka znika na mocy równania, a druga
dlatego, że obie pochodne $u_t$ i $\nabla_xu$
znikają na $\partial B$.
Dla $n=1$ twierdzenie Gaussa zastępujemy
twierdzeniem podstawowym rachunku całkowego na
przedziale $(-R,R)$; wartości brzegowe strumienia
również znikają. Stąd $E_u(t)=E_u(0)$ dla każdego
$t\in[0,T]$.
Założenie o wspólnym zwartym nośniku jest tutaj
warunkiem wystarczającym do przeprowadzenia dowodu,
nie warunkiem koniecznym zachowania energii.
Sama skończoność energii w poszczególnych chwilach
nie uzasadnia powyższego różniczkowania pod całką.
\end{proof}

Najważniejszy wynik nie wymaga jednak założeń o zachowaniu
rozwiązania przy $|x|\to\infty$. Zamiast całej przestrzeni
rozważymy kurczące się kule.

\begin{twierdzenie}[Lokalna energia i skończona prędkość propagacji]
\label{thm:pde5-cone}
Niech $n\ge1$, $c,T>0$, $x_0\in\R^n$ oraz
$u\in C^2([0,T]\times\R^n)$ spełnia $\square_cu=0$.
Jeżeli $u(0,x)=u_t(0,x)=0$ dla wszystkich
$x\in\overline B(x_0,cT)$, to $u(T,x_0)=0$.
Co więcej, w stożku
$\{(t,x):0\le t<T,\ |x-x_0|<c(T-t)\}$
mamy $u(t,x)=0$.
\end{twierdzenie}
\begin{proof}
Dla $0\le t<T$ połóżmy $r(t)=c(T-t)$,
$B_t=B(x_0,r(t))$ i
$e=\frac12(u_t^2+c^2|\nabla_xu|^2)$. Definiujemy
$E(t)=\int_{B_t}e(t,x)\,\mathrm d\lambda_n(x)$.
Dla $n\geq2$ zapisujemy całkę we współrzędnych
biegunowych o środku $x_0$:
\[
 E(t)=\int_0^{r(t)}\!\int_{S^{n-1}}
 e(t,x_0+\rho\theta)\rho^{n-1}\,
 \mathrm d\sigma(\theta)\,\mathrm d\lambda_1(\rho).
\]
Funkcja $e$ i jej pochodna czasowa są ciągłe
na zwartym zbiorze
$[0,T]\times\overline B(x_0,cT)$.
Wewnętrzną całkę możemy więc różniczkować
względem $t$ pod znakiem całki, a otrzymana
pochodna jest ciągła względem $(t,\rho)$
dla $0\leq\rho\leq cT$. Reguła Leibniza dla całki
o zmiennej górnej granicy daje
\[
 \begin{aligned}
 E'(t)
 &=\int_{B_t}e_t(t,x)\,\mathrm d\lambda_n(x)\\
 &\quad+r'(t)r(t)^{n-1}
     \int_{S^{n-1}}e(t,x_0+r(t)\theta)
       \,\mathrm d\sigma(\theta)\\
 &=\int_{B_t}e_t(t,x)\,\mathrm d\lambda_n(x)
    -c\int_{\partial B_t}e(t,x)\,\mathrm d\sigma(x).
 \end{aligned}
\]
W ostatnim przejściu użyliśmy skalowania miary
powierzchniowej przez czynnik $r(t)^{n-1}$
oraz $r'(t)=-c$.
Dla $n=1$ stosujemy bezpośrednio regułę Leibniza
do całki od $x_0-r(t)$ do $x_0+r(t)$: oba
końce dają składniki $-ce(t,x_0\pm r(t))$.
W obu przypadkach $E$ jest klasy $C^1$ dla
$0\leq t<T$, przy pochodnej prawostronnej w zerze.

Z równania falowego wynika punktowa tożsamość
\[
 e_t=u_tu_{tt}+c^2\nabla_xu\cdot\nabla_xu_t
     =c^2\diver(u_t\nabla_xu).
\]
Pole $u_t\nabla_xu$ jest klasy $C^1$ względem $x$,
więc twierdzenie Gaussa na kuli (w wymiarze jeden
twierdzenie podstawowe rachunku całkowego) daje
\[
\begin{aligned}
 E'(t)&=c^2\int_{\partial B_t}u_t\partial_\nu u
 \,\mathrm d\sigma-c\int_{\partial B_t}e\,\mathrm d\sigma\\
 &=-\frac c2\int_{\partial B_t}
 \bigl((u_t-c\partial_\nu u)^2+c^2|\nabla_{\tan}u|^2\bigr)
 \,\mathrm d\sigma\leq0.
\end{aligned}
\]
Tutaj $\nu$ jest zewnętrzną normalną jednostkową,
$\partial_\nu u=\nabla_xu\cdot\nu$, zaś
$\nabla_{\tan}u=\nabla_xu-(\partial_\nu u)\nu$.
Dla $n=1$ rozumiemy całkę brzegową jako sumę po
dwóch końcach przedziału, a składnik styczny jest
zerowy. Do dopełnienia kwadratu użyliśmy
ortogonalnego rozkładu gradientu na składową
normalną i styczną:
$|\nabla_xu|^2=|\partial_\nu u|^2+|\nabla_{\tan}u|^2$.
Z założeń początkowych $u(0,\cdot)$ jest
identycznie zerowa na otwartej kuli $B_0$,
więc także jej gradient znika na tej kuli.
Wraz z $u_t(0,\cdot)=0$ daje to $E(0)=0$.
Ponieważ $E'\leq0$ oraz $E\geq0$,
mamy $E(t)=0$ dla $0\leq t<T$.
Dla ustalonego $t$ nieujemna, ciągła funkcja
$e(t,\cdot)$ musi zatem znikać w każdym punkcie
$B_t$: gdyby była dodatnia w pewnym punkcie,
byłaby ograniczona z dołu przez dodatnią stałą
na pewnej kuli dodatniej miary, wbrew $E(t)=0$.
W szczególności $u_t=\nabla_xu=0$ w rozważanym
stożku. Jeżeli $x\in B_t$, to $x\in B_s$ dla
wszystkich $0\leq s\leq t$, gdyż promień kul
maleje. Stąd
$u(t,x)=u(0,x)+\int_0^t u_t(s,x)\,
\mathrm d\lambda_1(s)=0$.
Wreszcie $x_0\in B_t$ dla wszystkich $t<T$,
więc z ciągłości $u(T,x_0)=0$.
\end{proof}


\begin{twierdzenie}[Lokalne oszacowanie energii dla wymuszonego równania falowego]
\label{thm:pde5-cone-forcing}
Niech $n\ge1$, $c,T>0$, $x_0\in\R^n$, a
$u\in C^2([0,T]\times\R^n)$ spełnia
$u_{tt}-c^2\Delta_xu=f$ dla ciągłej funkcji
$f:[0,T]\times\R^n\to\R$.
Dla $0\le t<T$ oznaczmy $B_t=B(x_0,c(T-t))$ i
\[
 E(t)=\frac12\int_{B_t}
 \bigl(|u_t(t,x)|^2+c^2|\nabla_xu(t,x)|^2\bigr)
 \,\mathrm d\lambda_n(x).
\]
Wtedy dla każdego $0\le t<T$ zachodzi
\begin{equation}\label{eq:pde5-forced-cone-bound}
 \sqrt{E(t)}\le\sqrt{E(0)}+
 \frac1{\sqrt2}\int_0^t
 \|f(s,\cdot)\|_{L^2(B_s)}\,\mathrm d\lambda_1(s).
\end{equation}
W szczególności jeśli początkowe dane $u,u_t$ znikają
na $\overline B(x_0,cT)$, zaś $f=0$ w stożku
$\{(t,x):0\le t<T,\ x\in B_t\}$, to $u$ znika
w tym stożku oraz w jego wierzchołku $(T,x_0)$.
\end{twierdzenie}
\begin{proof}
Z twierdzenia o pochodnej całki po ruchomej kuli,
przeprowadzonego w dowodzie twierdzenia~\ref{thm:pde5-cone},
oraz z tożsamości
$\partial_t e=c^2\diver(u_t\nabla_xu)+u_tf$
otrzymujemy
\[
 E'(t)=-\frac c2\int_{\partial B_t}
 \bigl((u_t-c\partial_\nu u)^2+
 c^2|\nabla_{\tan}u|^2\bigr)\,\mathrm d\sigma
 +\int_{B_t}u_t(t,x)f(t,x)\,\mathrm d\lambda_n(x).
\]
Pierwszy składnik jest niedodatni. Dla drugiego
z nierówności Cauchy'ego--Schwarza otrzymujemy
\[
 E'(t)\leq
 \|u_t(t,\cdot)\|_{L^2(B_t)}\|f(t,\cdot)\|_{L^2(B_t)}
 \leq\sqrt{2E(t)}\|f(t,\cdot)\|_{L^2(B_t)}.
\]
Norma wymuszenia na ruchomej kuli jest ciągła
względem czasu dla $t<T$: po podstawieniu
$x=x_0+r(t)y$ jej kwadrat jest równy
$r(t)^n\int_{B(0,1)}|f(t,x_0+r(t)y)|^2
\,\mathrm d\lambda_n(y)$, a ciągłość wynika
ze zbieżności zdominowanej na ustalonej kuli,
gdyż $f$ jest ograniczona na domkniętym stożku.
Nie wolno bez uzasadnienia dzielić przez $\sqrt{E(t)}$,
ponieważ energia może znikać. Dla $\delta>0$
z reguły łańcuchowej i powyższej nierówności mamy
\[
 \begin{aligned}
 \frac{\mathrm d}{\mathrm dt}\sqrt{E(t)+\delta}
 &=\frac{E'(t)}{2\sqrt{E(t)+\delta}}\\
 &\leq\frac1{\sqrt2}\sqrt{\frac{E(t)}{E(t)+\delta}}\,
       \|f(t,\cdot)\|_{L^2(B_t)}\\
 &\leq\frac1{\sqrt2}\|f(t,\cdot)\|_{L^2(B_t)}.
 \end{aligned}
\]
Całkując na $[0,t]$, otrzymujemy oszacowanie
$\sqrt{E(t)+\delta}\leq\sqrt{E(0)+\delta}+
2^{-1/2}\int_0^t\|f(s,\cdot)\|_{L^2(B_s)}
\,\mathrm d\lambda_1(s)$.
Prawa strona poza wyrazem początkowym nie zależy
od $\delta$. Przechodząc z $\delta\downarrow0$,
dostajemy~\eqref{eq:pde5-forced-cone-bound}.
Jeżeli dane początkowe znikają na $\overline B_0$,
to również gradient danych $u(0,\cdot)$ znika na
$B_0$, zatem $E(0)=0$. Jeśli ponadto $f=0$
w rozważanym stożku, prawa strona oszacowania
jest zerowa dla każdego $t<T$. Argument z
dowodu twierdzenia~\ref{thm:pde5-cone} daje
punktowe zanikanie pochodnych, a następnie
$u(t,x)=0$ w stożku przez całkowanie po pionowym
odcinku zawartym w stożku. Ciągłość daje
$u(T,x_0)=0$.
\end{proof}

\begin{wniosek}[Jednoznaczność bez ograniczeń wzrostu]
\label{cor:pde5-uniqueness}
Niech $n\ge1$, $c,T>0$ i niech
$u,v\in C^2([0,T]\times\R^n)$ rozwiązują to samo
jednorodne równanie falowe i mają takie same dane
początkowe w całej $\R^n$. Wówczas $u=v$.
\end{wniosek}
\begin{proof}
Do $w=u-v$ stosujemy twierdzenie~\ref{thm:pde5-cone}
dla dowolnego $(T',x_0)$, $0<T'\le T$;
ograniczenie $w$ do $[0,T']\times\R^n$ ma
wymaganą regularność, spełnia równanie i ma
zerowe dane początkowe na $\overline B(x_0,cT')$.
Stąd $w(T',x_0)=0$. Ponieważ $T'$ i $x_0$
były dowolne, a w chwili $0$ dane są równe,
otrzymujemy $w=0$ na całym pasku.

W ten sam sposób stosując
twierdzenie~\ref{thm:pde5-cone-forcing} do różnicy
dwóch rozwiązań niejednorodnych, otrzymujemy
lokalną zasadę zależności od danych:
jeżeli oba profile początkowe są zgodne na
$\overline B(x_0,cT')$ i wymuszenia są zgodne
w stożku wstecznym punktu $(T',x_0)$, to wartości
rozwiązań w tym punkcie są równe.
Nie potrzeba zakładać skończonej energii
globalnej ani ograniczoności rozwiązań.
\end{proof}

\begin{uwaga}[Kontrast z równaniem ciepła]
W rozdziale~\ref{chap:pde-heat} skonstruowaliśmy
niezerowe, gładkie rozwiązanie Tychonoffa o zerowych
danych początkowych, co wymagało porzucenia ograniczeń
wzrostu. Dla klasycznego równania falowego na całej
$\R^n$ taki fenomen nie występuje w klasie $C^2$:
lokalna tożsamość energetyczna zapewnia jednoznaczność
bez jakichkolwiek warunków przy $|x|\to\infty$.
\end{uwaga}

\section{Wzór Kirchhoffa w trzech wymiarach}
\label{sec:pde5-kirchhoff}

Dla $v\in C^2(\R^3)$, $x\in\R^3$ i $r>0$
wprowadzamy średnią sferyczną
\begin{equation}\label{eq:pde5-sphere}
 M_rv(x)=\frac1{4\pi}\int_{S^2}v(x+r\omega)
 \,\mathrm d\sigma_{S^2}(\omega),
\end{equation}
gdzie $S^2=\{\omega\in\R^3:|\omega|=1\}$,
a $\mathrm d\sigma_{S^2}$ jest miarą powierzchniową
o całkowitej masie $4\pi$. Definiujemy $M_0v(x)=v(x)$.

\begin{lemat}[Równanie dla średnich sferycznych]
\label{lem:pde5-darboux}
Jeżeli $v\in C^2(\R^3)$, to dla $r>0$
\begin{equation}\label{eq:pde5-darboux}
 \partial_r^2 M_rv(x)+\frac2r\partial_rM_rv(x)
 =M_r(\Delta v)(x).
\end{equation}
\end{lemat}
\begin{proof}
Dla ustalonego $x$ niech $V(r,\omega)=v(x+r\omega)$.
Rozkład laplasjanu na składową radialną i sferyczną daje
$\Delta v(x+r\omega)=\partial_r^2V+(2/r)\partial_rV
+r^{-2}\Delta_{S^2}V$. Całka po sferze z
$\Delta_{S^2}V$ jest zerowa, ponieważ sfera nie
ma brzegu (można to sprawdzić przez lokalną
formułę dywergencji na sferze). Całkujemy otrzymaną
tożsamość po $S^2$ i dzielimy przez $4\pi$.
\end{proof}

\begin{twierdzenie}[Wzór Kirchhoffa]
\label{thm:pde5-kirchhoff}
Niech $c>0$, $g\in C^3(\R^3)$ oraz
$h\in C^2(\R^3)$. Wówczas istnieje dokładnie
jedno rozwiązanie $u\in C^2([0,\infty)\times\R^3)$
zagadnienia~\eqref{eq:pde5-wave}, dane wzorem
\begin{equation}\label{eq:pde5-kirchhoff}
 u(t,x)=\frac{\mathrm d}{\mathrm dt}
 \bigl(t M_{ct}g(x)\bigr)+t M_{ct}h(x),\qquad t\ge0.
\end{equation}
\end{twierdzenie}
\begin{proof}
Najpierw załóżmy, że $u$ jest rozwiązaniem $C^2$.
Dla ustalonego $x$ kładziemy
$W(t,r)=rM_ru(t,\cdot)(x)$, $r\ge0$. Z lematu
\ref{lem:pde5-darboux} otrzymujemy dla $r>0$
$W_{tt}-c^2W_{rr}=0$, gdyż
$\partial_r^2(rM_ru)=rM_r\Delta u$.
Mamy $W(t,0)=0$,
$W(0,r)=rM_rg(x)$ i $W_t(0,r)=rM_rh(x)$.
Przedłużając dane nieparzyście na $r<0$
i stosując wzór d'Alemberta z twierdzenia
\ref{thm:pde5-dalembert}, otrzymujemy
$W_r(t,0)=\frac{\mathrm d}{\mathrm dr}(rM_rg)|_{r=ct}
 +c^{-1}(rM_rh)|_{r=ct}$.
Ponieważ $W_r(t,0)=u(t,x)$, dostajemy
\eqref{eq:pde5-kirchhoff}.

Dla istnienia zdefiniujmy dla $v\in C^2(\R^3)$
$V_v(t,x)=tM_{ct}v(x)$. Z tożsamości Darboux
otrzymujemy dla $t>0$ bezpośrednio
\[
 (\partial_t^2-c^2\Delta_x)V_v
 =t c^2\left(\partial_r^2M_rv+\frac2r\partial_rM_rv
              -M_r\Delta v\right)\bigg|_{r=ct}=0.
\]
Dla $g\in C^3(\R^3)$ otrzymujemy $\partial_tV_g\in C^2$,
a dla $h\in C^2(\R^3)$ także $V_h\in C^2$.
Aby uzasadnić, że $\square_c(\partial_tV_g)=0$
przy tej skończonej regularności, aproksymujemy $g$
lokalnie w $C^3$ przez funkcje gładkie; skończony
zasięg całek sferycznych daje lokalną zbieżność
$\partial_tV_{g_j}\to\partial_tV_g$ w $C^2$.
Dla danych gładkich można skomutować $\square_c$
z $\partial_t$, a przejście do granicy dowodzi żądanej równości.
Ponadto $M_0v=v$ i $\partial_rM_rv|_{r=0}=0$,
ponieważ średnia z $\omega$ na $S^2$ jest zerowa.
Zatem $u=\partial_tV_g+V_h$ spełnia
$u(0,x)=g(x)$ i $u_t(0,x)=h(x)$.
Jednoznaczność wynika z wniosku~\ref{cor:pde5-uniqueness}.
\end{proof}

\begin{uwaga}[Jakie dane widzi punkt w trzech wymiarach?]
Dla ustalonego $(t,x)$ wartość rozwiązania
zależy od $g,h$ (i pochodnych $g$) na sferze
$\{y\in\R^3:|y-x|=ct\}$. Dane z wnętrza kuli,
ale z dala od tej sfery, nie mają wpływu na
$u(t,x)$. Jest to silna zasada Huygensa dla
równania falowego w wymiarze trzy.
\end{uwaga}

\section{Dwa wymiary i zasada Huygensa}
\label{sec:pde5-huygens}

\begin{twierdzenie}[Wzór Poissona w wymiarze dwa]
\label{thm:pde5-2d}
Niech $c>0$, $g\in C_c^3(\R^2)$, $h\in C_c^2(\R^2)$.
Jedyne rozwiązanie $C^2$ równania falowego w
$\R^2$ z tymi danymi jest dane wzorem
\begin{equation}\label{eq:pde5-2d}
\begin{aligned}
 u(t,x)={}&\frac{\mathrm d}{\mathrm dt}
 \left[\frac1{2\pi c}\int_{|y-x|<ct}
 \frac{g(y)}{\sqrt{c^2t^2-|y-x|^2}}\,
 \mathrm d\lambda_2(y)\right]\\
 &+\frac1{2\pi c}\int_{|y-x|<ct}
 \frac{h(y)}{\sqrt{c^2t^2-|y-x|^2}}\,
 \mathrm d\lambda_2(y),\qquad t>0.
\end{aligned}
\end{equation}
\end{twierdzenie}
\begin{proof}
Rozszerzamy $g,h$ na $\R^3$ tak, aby były niezależne
od trzeciej zmiennej: $\widetilde g(y,z)=g(y)$,
$\widetilde h(y,z)=h(y)$. Wzór Kirchhoffa definiuje
wtedy rozwiązanie trójwymiarowe niezależne od $z$
(jednoznaczność ze wniosku~\ref{cor:pde5-uniqueness}).
Dla $r>0$ parametryzujemy dwie półsfery
$|y-x|^2+z^2=r^2$ przez
$z=\pm\sqrt{r^2-|y-x|^2}$. Element pola na każdej
półsferze wynosi
$r(r^2-|y-x|^2)^{-1/2}\,\mathrm d\lambda_2(y)$.
Dlatego dla $\varphi$ niezależnej od $z$ mamy
\[
 M_r\varphi(x,0)=\frac1{2\pi r}
 \int_{|y-x|<r}\frac{\varphi(y)}{\sqrt{r^2-|y-x|^2}}
 \,\mathrm d\lambda_2(y).
\]
Wstawiając $r=ct$ do wzoru Kirchhoffa,
otrzymujemy dokładnie \eqref{eq:pde5-2d}.
\end{proof}

Zwróćmy uwagę na różnicę geometryczną: w wymiarze
trzy całkowanie odbywa się po \emph{sferze},
a w wymiarze dwa po \emph{całym dysku} z osobliwą,
ale całkowalną wagą. Dla $g=0$ i nieujemnej,
niezerowej funkcji $h\in C_c^2(\R^2)$ druga całka
jest ściśle dodatnia, gdy dysk $B(x,ct)$ obejmuje
nośnik $h$ (przy $t$ odpowiednio dużym).
Oznacza to ogon fali, którego nie ma w wersji
trójwymiarowej.

\begin{uwaga}[Silna i słaba zasada Huygensa]
Słaba zasada oznacza, że wpływ danych początkowych
nie może wyjść poza domknięty stożek $|x-y|\le ct$.
W przedstawionym modelu jest to wniosek z
\ref{thm:pde5-cone}. Silna zasada oznacza zależność
od danych wyłącznie na czole stożka, czyli na
$|x-y|=ct$. Dla stałych współczynników w przestrzeni
Euklidesowej silna zasada zachodzi w wymiarach
nieparzystych $n\ge3$, lecz zawodzi w wymiarach
parzystych i w wymiarze jeden. Uogólnienia
na wszystkie wymiary wymagają dalszych rachunków
z rozwiązaniami fundamentalnymi; na tym etapie
udowodniliśmy przypadki $n=1,2,3$.
\end{uwaga}

\section{Zewnętrzne wymuszenie i wzór Duhamela dla fal}
\label{sec:pde5-forcing}

\begin{twierdzenie}[Wzór dla niejednorodnego równania w jednym wymiarze]
\label{thm:pde5-forcing}
Niech $c,T>0$, $g\in C^2(\R)$, $h\in C^1(\R)$,
$f\in C^1([0,T]\times\R)$ i załóżmy, że wszystkie
te funkcje oraz pochodne występujące we wzorze są
ciągłe (nie wymagamy ograniczoności w nieskończoności).
Wówczas funkcja
\begin{equation}\label{eq:pde5-forcing}
\begin{aligned}
 u(t,x)={}&\frac{g(x-ct)+g(x+ct)}2
 +\frac1{2c}\int_{x-ct}^{x+ct}h(y)\,\mathrm d\lambda_1(y)\\
 &+\frac1{2c}\int_0^t\int_{x-c(t-s)}^{x+c(t-s)}
 f(s,y)\,\mathrm d\lambda_1(y)\,\mathrm ds
\end{aligned}
\end{equation}
rozwiązuje $u_{tt}-c^2u_{xx}=f$, $u(0)=g$,
$u_t(0)=h$ na $[0,T]\times\R$. Jest jedyna
w klasie rozwiązań $C^2$.
\end{twierdzenie}
\begin{proof}
Jednorodna część rozwiązuje problem bez wymuszenia.
Oznaczmy trzeci składnik przez $v(t,x)$.
Różniczkując dwukrotnie z wykorzystaniem
reguły Leibniza otrzymujemy
\[
 v_t(t,x)=\frac12\int_0^t\bigl[
 f(s,x+c(t-s))+f(s,x-c(t-s))\bigr] \,\mathrm ds
\]
oraz
\[
 v_{tt}(t,x)=f(t,x)+\frac c2\int_0^t
 \bigl[f_x(s,x+c(t-s))-f_x(s,x-c(t-s))\bigr]
 \,\mathrm ds.
\]
Jednocześnie
$v_{xx}(t,x)=\frac1{2c}\int_0^t
 \bigl[f_x(s,x+c(t-s))-f_x(s,x-c(t-s))\bigr]
 \,\mathrm ds$.
Stąd $v_{tt}-c^2v_{xx}=f$ i
$v(0,x)=v_t(0,x)=0$. Jednoznaczność wynika
z wniosku~\ref{cor:pde5-uniqueness} zastosowanego
do różnicy rozwiązań.
\end{proof}

\section{Warunki brzegowe: struna na odcinku}
\label{sec:pde5-string}

Na odcinku $\Omega=(0,L)$ rozważmy
$u_{tt}-c^2u_{xx}=0$ z warunkiem
$u(t,0)=u(t,L)=0$. Zagadnienie można rozwiązać
metodą rozdzielenia zmiennych. Dla $m\in\N$
połóżmy $\omega_m=cm\pi/L$; wtedy
\[
 u_m(t,x)=\bigl(a_m\cos(\omega_mt)
 +b_m\sin(\omega_mt)\bigr)\sin(m\pi x/L)
\]
jest rozwiązaniem klasycznym. Jeśli zadane dane
początkowe są skończonymi kombinacjami funkcji
$\sin(m\pi x/L)$, dobieramy
$a_m$ jako współczynniki danych $g$,
a $b_m$ jako współczynniki $h$ podzielone przez
$\omega_m$.

\begin{przyklad}[Brak tłumienia modów]
Dla $g(x)=\sin(m\pi x/L)$ i $h=0$ dostajemy
$u(t,x)=\cos(\omega_mt)\sin(m\pi x/L)$.
Amplituda nie maleje monotonicznie, podczas gdy
odpowiedni mod równania ciepła zanika jak
$e^{-\kappa(m\pi/L)^2t}$. Właśnie dlatego
regularizujący charakter półgrupy cieplnej
nie przenosi się na grupę ewolucji falowej.
\end{przyklad}

\begin{stwierdzenie}[Energia struny]
\label{prop:pde5-stringenergy}
Jeśli $u\in C^2([0,T]\times[0,L])$ spełnia
$u_{tt}=c^2u_{xx}$ dla $0<t<T$, $0<x<L$,
a także $u(t,0)=u(t,L)=0$, to
\[
 E(t)=\frac12\int_0^L(u_t^2+c^2u_x^2)\,
 \mathrm d\lambda_1(x)
\]
jest stała na $[0,T]$.
\end{stwierdzenie}
\begin{proof}
Z warunku Dirichleta i regularności mamy
$u_t(t,0)=u_t(t,L)=0$ dla $0<t<T$. Zatem
$E'(t)=c^2[u_tu_x]_{x=0}^{x=L}=0$.
Ciągłość energii daje równość także w chwilach
$t=0,T$.
\end{proof}

\section{Dlaczego schemat numeryczny musi respektować stożek?}
\label{sec:pde5-numerics}

Na regularnej siatce $x_j=j\Delta x$, $t_m=m\Delta t$
przybliżamy drugie pochodne różnicami centralnymi:
\begin{equation}\label{eq:pde5-leapfrog}
 \frac{U_j^{m+1}-2U_j^m+U_j^{m-1}}{(\Delta t)^2}
 =c^2\frac{U_{j+1}^m-2U_j^m+U_{j-1}^m}{(\Delta x)^2}.
\end{equation}
Wprowadzamy liczbę Couranta $\lambda=c\Delta t/\Delta x$.
Dla modu $U_j^m=e^{i(kj\Delta x-\omega m\Delta t)}$
otrzymujemy relację dyspersji
\begin{equation}\label{eq:pde5-cfl}
 \sin^2\!\left(\frac{\omega\Delta t}2\right)
 =\lambda^2\sin^2\!\left(\frac{k\Delta x}2\right).
\end{equation}
Jeżeli $\lambda\le1$, istnieją rzeczywiste
częstotliwości $\omega$ dla wszystkich częstości
siatkowych. Gdy $\lambda>1$, istnieją częstości,
dla których prawa strona przekracza $1$, co
prowadzi do wykładniczego wzrostu rozwiązań
rekurencji. Jest to konieczny warunek stabilności
metody centralnej; w tym konkretnym jednowymiarowym
schemacie z prawidłowo zadanymi danymi
$\lambda\le1$ jest również klasycznym warunkiem
stabilności w odpowiedniej normie energetycznej
(z dodatkową ostrożnością dla częstości granicznych
przy $\lambda=1$). Pełny dowód numeryczny
pozostawimy do działu metod różnicowych.

\begin{uwaga}[Granica stabilności przy $\lambda=1$]
Dla częstotliwości siatkowej $k\Delta x=\pi$ amplituda
modu spełnia rekurencję $A_{m+1}+2A_m+A_{m-1}=0$.
Jej ogólne rozwiązanie to $(\alpha+\beta m)(-1)^m$:
podwójny pierwiastek $-1$ dopuszcza więc liniowy wzrost
przy dowolnie zadanych dwóch pierwszych warstwach.
Dla prędkości początkowej równej zeru i poprawnego
pierwszego kroku $A_1=-A_0$, skąd $\beta=0$.
Analiza stabilności w przypadku granicznym musi uwzględniać
przestrzeń danych i sposób inicjalizacji schematu;
proste kryterium rzeczywistości $\omega$ nie wystarcza.
\end{uwaga}

\begin{uwaga}[Zasięg numeryczny a zasięg fizyczny]
Schemat \eqref{eq:pde5-leapfrog} korzysta z
sąsiadów $j-1,j,j+1$ w poprzedniej warstwie
czasowej. Jeden krok może więc przenieść informację
o odległość $\Delta x$, podczas gdy fizyczna fala
przebywa $c\Delta t$. Warunek $c\Delta t\le\Delta x$
pozwala fizycznemu stożkowi zmieścić się w
numerycznym obszarze zależności. To tłumaczy
geometrię warunku CFL, ale nie zastępuje
dowodu stabilności.
\end{uwaga}

\section{Podsumowanie}
\label{sec:pde5-summary}

Dla równania falowego rozwiązanie w punkcie jest
wyznaczone przez dane w skończonym obszarze
początkowym; wynika to z kurczącej się energii
stożkowej. Równania ciepła i falowe są pod tym
względem fundamentalnie odmienne. Wzory d'Alemberta,
Kirchhoffa i Poissona pokazują również, że w samym
równaniu falowym silna zasada Huygensa zależy
od parzystości wymiaru. Następny rozdział poświęcimy
systematycznemu wykorzystaniu transformaty i
szeregów Fouriera, łącząc otrzymane tu konstrukcje
z metodami spektralnymi.

% ROZDZIAŁ 6: do wstawienia przed rozdziałem z kodem numerycznym.

\chapter[Metody Fouriera i rozdzielenie zmiennych]{Szeregi Fouriera, rozdzielenie zmiennych i zagadnienia Sturma--Liouville'a}
\label{chap:pde-fourier}

W rozdziałach~\ref{chap:pde-heat} i~\ref{chap:pde-wave} uzyskaliśmy rozwiązania na
całej przestrzeni i w prostych obszarach. Na przedziale lub prostokącie
bardziej naturalne jest korzystanie z \emph{modów własnych} operatora
przestrzennego. Metoda rozdzielenia zmiennych będzie skuteczna tylko wtedy,
gdy uzasadnimy, że funkcje własne tworzą dostatecznie bogatą rodzinę.
Najpierw przeprowadzimy odpowiedni rachunek dla szeregów Fouriera,
a następnie zbadamy jego wersję dla regularnych zagadnień
Sturma--Liouville'a. Formalny rachunek dystrybucji przedstawimy w rozdziale~\ref{chap:pde-distributions},
a transformatę Fouriera na $\R^n$ i dystrybucje temperowane ---
w rozdziale~\ref{chap:pde-fourier-Rd}, po wprowadzeniu
przestrzeni Schwartza. Rozróżnienie szeregów Fouriera na torusie
i transformaty na całej przestrzeni jest istotne także dla normalizacji.

\section{Szereg Fouriera: współczynniki, ortogonalność i zbieżność}
\label{sec:pde6-fourier}

Przez $\mathbb T=\R/(2\pi\mathbb Z)$ rozumiemy okrąg z
współrzędną $x\in[-\pi,\pi]$ i utożsamionymi końcami.
Funkcję z $L^1(-\pi,\pi)$ utożsamiamy z jej
$2\pi$-okresowym przedłużeniem, określonym z dokładnością
do zbioru miary zero; wartości na końcach przedziału
nie wpływają na całki. Dla takiej funkcji definiujemy
$\widehat f(k)=(2\pi)^{-1}\int_{-\pi}^{\pi}f(x)e^{-ikx}\,\mathrm d\lambda_1(x)$,
$k\in\mathbb Z$, oraz sumę częściową
$S_Nf(x)=\sum_{|k|\le N}\widehat f(k)e^{ikx}$.
W tym rozdziale $L^2(\mathbb T)$ oznacza zespoloną przestrzeń Hilberta
z iloczynem skalarnym
$\langle f,g\rangle=(2\pi)^{-1}\int_{-\pi}^{\pi}f(x)\overline{g(x)}\,\mathrm d\lambda_1(x)$.
Pojęcie całki Lebesgue'a i podstawowe własności $L^p$ traktujemy
jako znane; szerszą teorię funkcjonalną rozwiniemy w dalszych tomach.
Zwracamy uwagę na normalizację: w tym rozdziale norma $L^2(\T)$
pochodzi od miary $\mathrm d\lambda_1(x)/(2\pi)$.
W dalszych rozdziałach o PDE na torusie stosujemy zwykle
nieznormalizowaną miarę $\mathrm d\lambda_d(x)$; wzory
na wartości własne i współczynniki Fouriera pozostają zgodne,
ale norma względem miary nieznormalizowanej w wymiarze $d$
jest $(2\pi)^{d/2}$ razy większa od normy względem
miary unormowanej do masy $1$ (na torusie
o okresie $2\pi$ w każdej współrzędnej).
Na odcinku $(0,L)$ używamy natomiast zwykłej
miary Lebesgue'a. Ponieważ miara unormowana na
$\mathbb T$ ma masę $1$, nierówność Cauchy'ego--Schwarza
daje $(2\pi)^{-1}\|f\|_{L^1(-\pi,\pi)}\leq\|f\|_2$,
więc współczynniki Fouriera są określone dla każdego
$f\in L^2(\mathbb T)$.

\begin{lemat}[Ortogonalność i nierówność Bessela]
\label{lem:pde6-bessel}
Niech $f\in L^2(\mathbb T)$ oraz $N\geq0$ będzie całkowite. Wówczas
$\langle e^{ikx},e^{i\ell x}\rangle=\delta_{k\ell}$,
$\|S_Nf\|_2^2=\sum_{|k|\le N}|\widehat f(k)|^2$ i
$\sum_{|k|\le N}|\widehat f(k)|^2\le\|f\|_2^2$.
\end{lemat}
\begin{proof}
Dla $k=\ell$ całka z jedynki względem miary
unormowanej wynosi $1$. Jeśli $k\ne\ell$, to
\[
 \frac1{2\pi}\int_{-\pi}^{\pi}e^{i(k-\ell)x}\,
 \mathrm d\lambda_1(x)
 =\frac{e^{i(k-\ell)\pi}-e^{-i(k-\ell)\pi}}
 {2\pi i(k-\ell)}=0,
\]
ponieważ $k-\ell$ jest całkowite. Otrzymujemy
ortonormalność. Rozwijając kwadrat normy skończonej
sumy $S_Nf$ i korzystając z tej ortonormalności,
dostajemy $\|S_Nf\|_2^2=\sum_{|k|\le N}|\widehat f(k)|^2$.
Ponadto dla $|\ell|\leq N$
\[
 \langle f-S_Nf,e^{i\ell x}\rangle
 =\widehat f(\ell)-\sum_{|k|\leq N}
 \widehat f(k)\delta_{k\ell}=0.
\]
Zatem $S_N$ jest rzutem ortogonalnym na przestrzeń
rozpiętą przez te wykładniki. W szczególności
$\langle f-S_Nf,S_Nf\rangle=0$, skąd
\[
 \|f\|_2^2=\|S_Nf\|_2^2+\|f-S_Nf\|_2^2.
\]
Nieujemność ostatniego składnika daje nierówność Bessela.
\end{proof}

Nie można jednak wywnioskować z ortogonalności, że $S_Nf\to f$.
Do tego potrzebujemy zupełności układu wykładniczego.

\begin{lemat}[Jądro Fejéra]
\label{lem:pde6-fejer}
Dla całkowitego $N\geq0$ niech
\begin{equation}\label{eq:pde6-fejer}
 F_N(t)=\frac1{N+1}\left|\sum_{j=0}^Ne^{ijt}\right|^2
 =\frac1{N+1}\left(\frac{\sin((N+1)t/2)}{\sin(t/2)}\right)^2
\end{equation}
(w punktach $t\in2\pi\mathbb Z$ przyjmujemy
$F_N(t)=N+1$, zgodnie z pierwszą postacią wzoru). Wówczas $F_N\ge0$,
$(2\pi)^{-1}\int_{-\pi}^{\pi}F_N(t)\,\mathrm d\lambda_1(t)=1$ oraz
dla każdego $\delta\in(0,\pi)$
$\int_{\delta\le|t|\le\pi}F_N(t)\,\mathrm d\lambda_1(t)\to0$.
\end{lemat}
\begin{proof}
Nieujemność jest bezpośrednia z definicji.
Dla $t\notin2\pi\mathbb Z$ z sumy geometrycznej
$\sum_{j=0}^N e^{ijt}=(1-e^{i(N+1)t})/(1-e^{it})$
i tożsamości $|1-e^{i\alpha}|=2|\sin(\alpha/2)|$
otrzymujemy drugą postać~\eqref{eq:pde6-fejer}.
Rozwijając kwadrat modułu, mamy
\[
 F_N(t)=\frac1{N+1}\sum_{j,\ell=0}^N e^{i(j-\ell)t}
 =\sum_{|k|\leq N}\left(1-\frac{|k|}{N+1}\right)e^{ikt}.
\]
Dla ustalonego $|k|\leq N$ istnieje dokładnie
$N+1-|k|$ par $(j,\ell)$ z różnicą $j-\ell=k$:
dla $k\geq0$ wybieramy $\ell=0,\ldots,N-k$,
a dla $k<0$ zamieniamy role indeksów.
Wyraz stały jest równy $1$, a pozostałe
wykładniki mają całkę zero na pełnym okresie,
skąd normalizacja. Gdy $\delta\leq|t|\leq\pi$,
mamy $|\sin(t/2)|\geq\sin(\delta/2)>0$.
Dlatego
\[
 \int_{\delta\leq|t|\leq\pi}F_N(t)\,\mathrm d\lambda_1(t)
 \leq\frac{2\pi}{(N+1)\sin^2(\delta/2)}
 \longrightarrow0.
\]
\end{proof}

\begin{twierdzenie}[Fejér i zupełność układu trygonometrycznego]
\label{thm:pde6-fejer}
Niech $f\in C(\mathbb T)$. Średnie Cesàra
$\sigma_N f=(N+1)^{-1}\sum_{j=0}^N S_jf$
zbieżne są jednostajnie do $f$ na $\mathbb T$.
W konsekwencji wielomiany trygonometryczne są gęste
w $C(\mathbb T)$ w normie supremum i w $L^2(\mathbb T)$
w normie $L^2$.
\end{twierdzenie}
\begin{proof}
W średniej z sum częściowych współczynnik
przy $\widehat f(k)e^{ikx}$ wynosi
$(N+1-|k|)/(N+1)$ dla $|k|\leq N$, ponieważ
wyraz ten występuje dla $j=|k|,\ldots,N$.
Z drugiej strony podstawienie $y=x-t$ daje
\[
 \frac1{2\pi}\int_{-\pi}^{\pi}e^{ikt}f(x-t)\,
 \mathrm d\lambda_1(t)=e^{ikx}\widehat f(k).
\]
Po podstawieniu całkujemy po odcinku
$[x-\pi,x+\pi]$, który możemy zastąpić
$[-\pi,\pi]$, gdyż $f(y)e^{-iky}$ jest
$2\pi$-okresowa i całka po pełnym okresie
nie zależy od jego początku (przez podział
odcinka i przesunięcie o wielokrotność okresu).
Rozwinięcie skończone z lematu~\ref{lem:pde6-fejer}
daje więc
\[
 \sigma_Nf(x)=\frac1{2\pi}\int_{-\pi}^{\pi}
 F_N(t)f(x-t)\,\mathrm d\lambda_1(t).
\]
Odejmujemy $f(x)$ i dzielimy całkę na $|t|<\delta$
oraz $\delta\le|t|\le\pi$. Z jednostajnej ciągłości $f$
możemy dobrać $\delta$ tak, aby pierwszy składnik
miał moduł co najwyżej $\varepsilon$ dla wszystkich $x,N$.
Drugi ma moduł nie większy niż
$2\|f\|_\infty(2\pi)^{-1}\int_{|t|\ge\delta}F_N(t)
\,\mathrm d\lambda_1(t)$, co dąży do zera.
Ponieważ $F_N$ ma unormowaną całkę $1$, pierwszy
składnik jest rzeczywiście ograniczony przez
$\varepsilon$, niezależnie od $N$. Zatem dla
dostatecznie dużych $N$ norma supremum różnicy
nie przekracza $2\varepsilon$. Dowodzi to
jednostajnej zbieżności.

Uzasadnijmy również gęstość w $L^2$ z uwzględnieniem
okresowości. Funkcje proste są gęste w $L^2$
na przedziale skończonej miary. Dla mierzalnego
$A\subset(-\pi,\pi)$ oraz $\eta>0$ regularność
miary Lebesgue'a daje zbiór zwarty $K\subset A$
i otwarty $O\supset A$, $O\subset(-\pi,\pi)$,
takie że $\lambda_1(O\setminus K)<\eta$.
Gdy $K$ jest niepusty, wybieramy
$0<r<\operatorname{dist}(K,\R\setminus O)$.
Funkcja
\[
 \varphi(x)=\max\{0,1-\operatorname{dist}(x,K)/r\}
\]
jest ciągła, równa $1$ na $K$, ma zwarty nośnik
zawarty w $O$ i wartości w $[0,1]$.
Jeśli $K=\varnothing$, bierzemy $\varphi=0$.
W obu przypadkach
$\|\chi_A-\varphi\|_2^2\leq\eta/(2\pi)$.
Przedłużenie $\varphi$ okresowo jest ciągłe,
bo $\varphi$ znika w otoczeniu końców przedziału.
Aproksymując w ten sposób każdy składnik
funkcji prostej i stosując nierówność trójkąta,
otrzymujemy gęstość $C(\mathbb T)$ w $L^2(\mathbb T)$.
Dla $f\in L^2$ i $\varepsilon>0$ wybieramy
$g\in C(\mathbb T)$ z $\|f-g\|_2<\varepsilon/2$,
a następnie wielomian trygonometryczny $p$
z $\|g-p\|_\infty<\varepsilon/2$.
Ponieważ $\|g-p\|_2\leq\|g-p\|_\infty$,
mamy $\|f-p\|_2<\varepsilon$.
\end{proof}

\begin{twierdzenie}[Zbieżność w $L^2$ i tożsamość Parsevala]
\label{thm:pde6-parseval}
Dla dowolnego $f\in L^2(\mathbb T)$ zachodzą
$S_Nf\to f$ w normie $L^2$ oraz
\begin{equation}\label{eq:pde6-parseval}
 \frac1{2\pi}\int_{-\pi}^{\pi}|f(x)|^2\,\mathrm d\lambda_1(x)
 =\sum_{k\in\mathbb Z}|\widehat f(k)|^2.
\end{equation}
\end{twierdzenie}
\begin{proof}
Dla $\varepsilon>0$ wybieramy wielomian trygonometryczny $p$
z $\|f-p\|_2<\varepsilon$. Dla $N$ nie mniejszych niż stopień
$p$ mamy $S_Np=p$. Ponieważ $f-S_Nf$ jest
prostopadła do przestrzeni, do której należy
$S_Nf-p$, zachodzi
\[
 \|f-p\|_2^2=\|f-S_Nf\|_2^2+\|S_Nf-p\|_2^2.
\]
Zatem $\|f-S_Nf\|_2\leq\|f-p\|_2<\varepsilon$,
co dowodzi zbieżności.
W konsekwencji $\|S_Nf\|_2\to\|f\|_2$,
a z tożsamości $\|S_Nf\|_2^2=
\sum_{|k|\leq N}|\widehat f(k)|^2$
w lemacie~\ref{lem:pde6-bessel} otrzymujemy
\eqref{eq:pde6-parseval}, ponieważ sumy
nieujemnych wyrazów rosną do pełnej sumy po $\mathbb Z$.
\end{proof}

\begin{uwaga}[Jakiego rodzaju zbieżność została udowodniona?]
Sama zbieżność w normie $L^2$ nie pociąga
za sobą zbieżności jednostajnej ani punktowej
całego ciągu. Nie rozstrzygamy tutaj odrębnego
zagadnienia zbieżności prawie wszędzie sum
częściowych Fouriera. Dla funkcji ciągłych jednostajnie zbieżne są
\emph{średnie Fejéra}, niekoniecznie zwykłe sumy $S_N$.
W dalszej części do różniczkowania szeregów użyjemy
gładkich danych, dla których zachodzi dodatkowe szybkie
zanikanie współczynników.
\end{uwaga}

\begin{lemat}[Zanikanie współczynników gładkiej funkcji]
\label{lem:pde6-decay}
Jeśli $f\in C^m(\mathbb T)$, $m\geq1$ całkowite, to dla $k\ne0$
$|\widehat f(k)|\le\|f^{(m)}\|_{L^1(-\pi,\pi)}/(2\pi|k|^m)$.
Jeśli $f\in C^\infty(\mathbb T)$, to szereg Fouriera wraz
ze wszystkimi szeregami otrzymanymi przez różniczkowanie
zbiega jednostajnie.
\end{lemat}
\begin{proof}
Dla $k\ne0$ całkowanie przez części daje
\[
 \widehat f(k)=\frac1{ik}\widehat{f'}(k),
\]
ponieważ $f(\pi)e^{-ik\pi}=f(-\pi)e^{ik\pi}$.
Powtarzając ten rachunek $m$ razy (pochodne
do rzędu $m-1$ mają zgodne wartości na końcach
przedziału), dostajemy
$\widehat f(k)=(ik)^{-m}\widehat{f^{(m)}}(k)$,
co po oszacowaniu modułu całki daje żądaną nierówność.

Dla $f\in C^\infty$ i ustalonego całkowitego
$r\geq0$ wybieramy $m\geq r+2$. Wówczas
\[
 \sum_{k\ne0}|k|^r|\widehat f(k)|
 \leq\frac{\|f^{(m)}\|_{L^1(-\pi,\pi)}}{2\pi}
 \sum_{k\ne0}|k|^{r-m}<\infty.
\]
Kryterium Weierstrassa daje bezwzględną
i jednostajną zbieżność szeregu
$\sum_{k\in\mathbb Z}(ik)^r\widehat f(k)e^{ikx}$,
przy czym dla $r=0$ wyraz $k=0$ oznacza
$\widehat f(0)$, a dla $r>0$ wynosi zero.
Oznaczmy jego ciągłą sumę przez $G_r$.
Z całkowania przez części także dla $k=0$
(jeśli $r>0$, całka z $f^{(r)}$ po okresie
jest zerowa) współczynniki tego szeregu
są dokładnie współczynnikami Fouriera $f^{(r)}$.
Twierdzenie~\ref{thm:pde6-parseval} daje więc
zbieżność tych samych sum częściowych do $f^{(r)}$
w normie $L^2$. Z jednostajnej zbieżności
mamy zarazem zbieżność w $L^2$ do $G_r$.
Jednoznaczność granicy w $L^2$ daje
$G_r=f^{(r)}$ prawie wszędzie, a ponieważ
obie funkcje są ciągłe, równość zachodzi
w każdym punkcie: niezerowa różnica
w jednym punkcie byłaby niezerowa na
pewnym otwartym łuku dodatniej miary.
W szczególności $G_0=f$ i wszystkie
różniczkowania wyraz po wyrazie prowadzą
do właściwych pochodnych.
\end{proof}

\section{Zagadnienie własne laplasjanu na odcinku}
\label{sec:pde6-dirichlet}

Dla $L>0$ rozważmy
\begin{equation}\label{eq:pde6-eigen}
 -X''(x)=\lambda X(x),\quad 0<x<L,\qquad X(0)=X(L)=0.
\end{equation}
\begin{twierdzenie}[Pełne spektrum Dirichleta w jednym wymiarze]
\label{thm:pde6-spectrum}
Dla $L>0$ zagadnienie~\eqref{eq:pde6-eigen} ma
nietrywialne rozwiązanie $X\in C^2([0,L])$
wtedy i tylko wtedy, gdy
$\lambda=\lambda_n=(n\pi/L)^2$ dla pewnego całkowitego $n\ge1$.
Jego przestrzeń rozwiązań jest rozpięta przez
$X_n(x)=\sin(n\pi x/L)$. Ponadto funkcje
$e_n(x)=\sqrt{2/L}\sin(n\pi x/L)$, $n\ge1$,
tworzą bazę ortonormalną w $L^2(0,L)$.
\end{twierdzenie}
\begin{proof}
Mnożąc~\eqref{eq:pde6-eigen} przez $\overline X$
i całkując przez części otrzymujemy
$\lambda\int_0^L|X(x)|^2\,\mathrm d\lambda_1(x)
=\int_0^L|X'(x)|^2\,\mathrm d\lambda_1(x)$.
Ponieważ $X$ nie jest identycznie zerowa, całka z $|X|^2$
jest dodatnia. Gdyby całka z $|X'|^2$ była zerowa,
ciągłość $X'$ wymuszałaby $X'=0$ na całym odcinku.
Wtedy $X$ byłaby stała, a z $X(0)=0$ wynikałoby $X=0$,
sprzecznie z założeniem. Stąd $\lambda$ jest rzeczywista
i dodatnia. Ogólne rozwiązanie ODE to
$A\sin(\sqrt\lambda x)+B\cos(\sqrt\lambda x)$.
Warunek w $0$ daje $B=0$, a nietrywialność wymusza $A\ne0$.
Warunek w $L$ jest więc równoważny
$\sqrt\lambda L=n\pi$ dla pewnego całkowitego $n\ge1$.
Odwrotnie, bezpośrednie różniczkowanie $X_n$ daje
$-X_n''=\lambda_n X_n$ i oba warunki brzegowe.

Dla $m,n\ge1$ tożsamość
$2\sin a\sin b=\cos(a-b)-\cos(a+b)$ daje
\[
 \int_0^L\sin(m\pi x/L)\sin(n\pi x/L)\,
 \mathrm d\lambda_1(x)=\frac L2\delta_{mn}.
\]
Istotnie, całka z $\cos(j\pi x/L)$ znika dla każdego
niezerowego całkowitego $j$, a dla $j=0$ wynosi $L$.
Stąd rodzina $(e_n)$ jest ortonormalna.

Dla zupełności niech $f\in L^2(0,L)$. Przedłużamy ją
nieparzyście do $(-L,L)$, otrzymując $\widetilde f$.
Mamy $\int_{-L}^L|\widetilde f|^2\,\mathrm d\lambda_1
=2\int_0^L|f|^2\,\mathrm d\lambda_1$.
Podstawienie $x=L\theta/\pi$ pozwala zastosować
twierdzenie~\ref{thm:pde6-parseval} do funkcji
$\theta\mapsto\widetilde f(L\theta/\pi)$.
Współczynniki wykładnicze na okresie $2L$ są równe
\[
 c_k=\frac1{2L}\int_{-L}^L\widetilde f(x)
 e^{-ik\pi x/L}\,\mathrm d\lambda_1(x).
\]
Nieparzystość daje $c_0=0$, $c_{-n}=-c_n$ oraz
$c_n=-\frac{i}{L}\int_0^Lf(x)\sin(n\pi x/L)\,
\mathrm d\lambda_1(x)$ dla $n\ge1$.
Nie wymaga to rzeczywistości funkcji $f$.
Grupując wyrazy o indeksach $n,-n$, otrzymujemy
\[
 \sum_{|k|\le N}c_ke^{ik\pi x/L}
 =\sum_{n=1}^N\langle f,e_n\rangle_{L^2(0,L)}e_n(x).
\]
Lewa strona zbiega do $\widetilde f$ w $L^2(-L,L)$;
ograniczenie do $(0,L)$ nie zwiększa normy błędu.
Zatem prawa strona zbiega do $f$ w $L^2(0,L)$,
co dowodzi zupełności. Z ortonormalności i przejścia
do granicy norm sum częściowych wynika również
$\|f\|_{L^2(0,L)}^2=\sum_{n\ge1}|\langle f,e_n\rangle|^2$.
\end{proof}

\begin{wniosek}[Nierówność Poincarégo dla funkcji gładkich]
\label{cor:pde6-poincare}
Jeżeli $v\in C^1([0,L])$, $v(0)=v(L)=0$, to
\begin{equation}\label{eq:pde6-poincare}
 \int_0^L|v(x)|^2\,\mathrm d\lambda_1(x)
 \le \left(\frac L\pi\right)^2
 \int_0^L|v'(x)|^2\,\mathrm d\lambda_1(x).
\end{equation}
Stała $(L/\pi)^2$ jest optymalna.
\end{wniosek}
\begin{proof}
Niech $v_n=\langle v,e_n\rangle_{L^2(0,L)}$.
Oznaczmy $d_n(x)=\sqrt{2/L}\cos(n\pi x/L)$
dla $n\ge1$. Całkowanie przez części daje
\[
 \langle v',d_n\rangle
 = [v d_n]_0^L-\int_0^L v(x)d_n'(x)\,
 \mathrm d\lambda_1(x)
 = (n\pi/L)v_n,
\]
gdyż $v$ znika na obu końcach i $d_n'=-(n\pi/L)e_n$.
Z tożsamości $2\cos a\cos b=\cos(a-b)+\cos(a+b)$
i całek kosinusów obliczonych w poprzednim dowodzie
wynika ortonormalność $(d_n)_{n\ge1}$.
Nierówność Bessela dla tej rodziny (nie potrzebujemy
jej zupełności) i tożsamość Parsevala dla bazy sinusowej dają
\[
 \|v'\|_2^2\ge\sum_{n\ge1}(n\pi/L)^2|v_n|^2
 \ge(\pi/L)^2\sum_{n\ge1}|v_n|^2
 =(\pi/L)^2\|v\|_2^2.
\]
To jest~\eqref{eq:pde6-poincare}. Dla
$v(x)=\sin(\pi x/L)$ mamy $\|v\|_2^2=L/2$ i
$\|v'\|_2^2=(\pi/L)^2 L/2$, więc żadna mniejsza
stała nie może być poprawna dla wszystkich $v$.
\end{proof}

\begin{uwaga}[Po co zachowywać stałe w oszacowaniach?]
Nierówność~\eqref{eq:pde6-poincare} wzmacnia
oszacowanie elementarne uzyskane w rozdziale~\ref{chap:pde-heat}.
Najmniejsza wartość własna $\lambda_1=\pi^2/L^2$
wyznacza dokładną szybkość zaniku pierwszego modu
równania ciepła i najniższą częstość drgań struny.
\end{uwaga}

\section{Rozdzielenie zmiennych: równanie ciepła}
\label{sec:pde6-heat}

Niech $\kappa,L>0$, a $u_t=\kappa u_{xx}$ dla $t>0$,
$x\in(0,L)$, z $u(t,0)=u(t,L)=0$ i $u(0,x)=g(x)$.
Podstawienie $u(t,x)=T(t)X(x)$ daje
$T'X=\kappa TX''$. Dla nietrywialnego iloczynu wybieramy
$t_*$ z $T(t_*)\ne0$ i kładziemy
$\lambda=-T'(t_*)/(\kappa T(t_*))$. Otrzymujemy
$-X''=\lambda X$ dla wszystkich $x$, wraz z zerowymi
warunkami Dirichleta. Wybór $x_*$ z $X(x_*)\ne0$
daje następnie $T'=-\kappa\lambda T$ dla wszystkich $t$.
Nie dzielimy przez funkcje w ich punktach zerowych.
Dla $\lambda=\lambda_n$ otrzymujemy
$T(t)=T(0)e^{-\kappa\lambda_nt}$; stałą $T(0)$ włączamy
do amplitudy modu.


\begin{lemat}[Jądro Gaussa i kontrakcja na odcinku]
\label{lem:pde6-dirichlet-heat}
Niech $L,\kappa>0$ i niech
$g\in C([0,L])$ spełnia $g(0)=g(L)=0$.
Oznaczmy przez $\widetilde g$ jej nieparzyste
$2L$-okresowe przedłużenie na $\R$. Dla $t>0$ połóżmy
\[
 (Q_tg)(x)=\int_{\R}(4\pi\kappa t)^{-1/2}
 e^{-z^2/(4\kappa t)}\widetilde g(x-z)
 \,\mathrm d\lambda_1(z),\qquad x\in[0,L].
\]
Wówczas $Q_tg$ znika dla $x=0,L$,
$\|Q_tg\|_{C([0,L])}\le\|g\|_{C([0,L])}$ oraz
$\|Q_tg-g\|_{C([0,L])}\to0$ dla $t\downarrow0$.
Ponadto współczynnik sinusowy $Q_tg$ odpowiadający
$\sin(n\pi x/L)$ jest równy
$e^{-\kappa(n\pi/L)^2t}b_n$, gdzie
$b_n=(2/L)\int_0^Lg(x)\sin(n\pi x/L)\,\mathrm d\lambda_1(x)$.
\end{lemat}
\begin{proof}
Przedłużenie jest ciągłe i jednostajnie ciągłe na $\R$,
a jego norma supremum jest równa normie $g$.
Oznaczmy jądro Gaussa w definicji przez $G_t(z)$.
Jest ono nieujemne, a podstawienie $z=2\sqrt{\kappa t}\,r$
i całka Gaussa dają $\int_{\R}G_t\,\mathrm d\lambda_1=1$.
Stąd bezwzględna zbieżność całki i kontrakcja.
Ten sam wzór określa funkcję na całej prostej.
Jednostajna ciągłość $\widetilde g$ i normalizacja jądra
dają ciągłość splotu. Okresowość $\widetilde g$ daje
$Q_tg(x+2L)=Q_tg(x)$, a podstawienie $z\mapsto-z$
i parzystość $G_t$ dają $Q_tg(-x)=-Q_tg(x)$.
W szczególności $Q_tg(0)=0$ i
$Q_tg(L)=Q_tg(-L)=-Q_tg(L)$.

Dla $\varepsilon>0$ dobieramy $\delta>0$ z jednostajnej
ciągłości tak, aby dla $|z|<\delta$ różnica
$|\widetilde g(x-z)-\widetilde g(x)|$ była mniejsza
niż $\varepsilon$ dla wszystkich $x$. Wówczas
\[
 \|Q_tg-g\|_\infty\le\varepsilon+2\|g\|_\infty
 \int_{|z|\ge\delta}G_t(z)\,\mathrm d\lambda_1(z).
\]
Podstawienie $z=2\sqrt{\kappa t}\,r$ zamienia ostatnią
całkę na $\pi^{-1/2}\int_{|r|\ge\delta/(2\sqrt{\kappa t})}
e^{-r^2}\,\mathrm d\lambda_1(r)$, która dąży do zera
z całkowalności Gaussa. Ponieważ $\varepsilon$ jest
dowolne, otrzymujemy zbieżność jednostajną.

Niech $K_t(z)=\sum_{j\in\mathbb Z}G_t(z+2jL)$.
Dla $z\in[-L,L]$ i $|j|\ge1$ mamy
$|z+2jL|\ge(2|j|-1)L$, więc szereg zbiega jednostajnie
na okresie dla ustalonego $t>0$ przez kryterium Weierstrassa.
Podział prostej na przedziały długości $2L$ i okresowość
$\widetilde g$ dają
\[
 Q_tg(x)=\int_{-L}^L K_t(z)\widetilde g(x-z)\,
 \mathrm d\lambda_1(z).
\]
Dla współczynników $c_n(f)=(2L)^{-1}
\int_{-L}^L f(x)e^{-in\pi x/L}\,\mathrm d\lambda_1(x)$
całkowanie szeregu i podstawienie na każdym przedziale dają
\[
 2L c_n(K_t)=\int_{\R}G_t(z)e^{-in\pi z/L}\,
 \mathrm d\lambda_1(z).
\]
Czynnik $2L$ wynika z tego, że splot zapisano względem
zwykłej miary Lebesgue'a. Z twierdzenia Fubiniego
(ograniczoność $\widetilde g$ i całkowalność jądra)
oraz podstawienia $y=x-z$ w całce po pełnym okresie
wynika $c_n(Q_tg)=2L c_n(K_t)c_n(\widetilde g)$.

Dla $\xi\in\R$ niech
$M_t(\xi)=\int_{\R}G_t(z)e^{-i\xi z}\,\mathrm d\lambda_1(z)$.
Różniczkowanie pod znakiem całki jest uzasadnione
całkowalnością $|z|G_t(z)$. Ponieważ
$zG_t(z)=-2\kappa tG_t'(z)$, całkowanie przez części
z zanikającym wyrazem brzegowym daje
\[
 M_t'(\xi)=2i\kappa t\int_{\R}G_t'(z)e^{-i\xi z}
 \,\mathrm d\lambda_1(z)=-2\kappa t\xi M_t(\xi).
\]
Warunek $M_t(0)=1$ daje $M_t(\xi)=e^{-\kappa t\xi^2}$.
Stąd $c_n(Q_tg)=e^{-\kappa(n\pi/L)^2t}c_n(\widetilde g)$.
Przejście od współczynników wykładniczych funkcji
nieparzystych do sinusowych, jak w dowodzie
twierdzenia~\ref{thm:pde6-spectrum}, daje tezę również
dla zespolonej funkcji $g$.
\end{proof}

\begin{twierdzenie}[Rozwiązanie przez szereg sinusowy]
\label{thm:pde6-heat}
Niech $\kappa,L>0$ i $g\in C([0,L])$ spełnia
$g(0)=g(L)=0$. Definiujemy
$b_n=(2/L)\int_0^Lg(x)\sin(n\pi x/L)\,\mathrm d\lambda_1(x)$.
Wówczas
\begin{equation}\label{eq:pde6-heat}
 u(t,x)=\sum_{n=1}^\infty b_n e^{-\kappa\lambda_nt}
 \sin(n\pi x/L),\qquad t>0,\;x\in[0,L],
\end{equation}
jest klasycznym rozwiązaniem równania ciepła dla $t>0$,
ciągłym na $[0,\infty)\times[0,L]$ i przyjmuje dane $g$
w sposób jednostajny, gdy $t\downarrow0$. Rozwiązanie to
jest jedyne w klasie funkcji ciągłych na każdym
$[0,T]\times[0,L]$, klasy $C^{1,2}$ we wnętrzu
i spełniających dane i warunki Dirichleta.
\end{twierdzenie}
\begin{proof}
Z ograniczoności $g$ mamy $|b_n|\le2\|g\|_\infty$.
Dla $t\ge\delta>0$ szereg i wszystkie szeregi
z pochodnymi czasowymi i przestrzennymi zbieżne są
jednostajnie, bo dla dowolnych $r,s\ge0$
$\sum_{n\ge1}n^{r+2s}e^{-\kappa(n\pi/L)^2\delta}<\infty$.
Równanie i warunki brzegowe sprawdzamy wyraz po wyrazie.
Dla $t>0$ całkowanie jednostajnie zbieżnego szeregu
i ortogonalność sinusów dają współczynniki sinusowe
$b_ne^{-\kappa\lambda_nt}$. Z lematu~\ref{lem:pde6-dirichlet-heat}
te same współczynniki ma $Q_tg$. Zupełność układu
sinusowego (twierdzenie~\ref{thm:pde6-spectrum})
daje równość w $L^2(0,L)$, a ciągłość obu funkcji
daje równość w każdym punkcie $[0,L]$.
Zatem $u(t,\cdot)=Q_tg$ i
$\|u(t,\cdot)-g\|_\infty\to0$ dla $t\downarrow0$.
Definiujemy $u(0,x)=g(x)$. Ciągłość w $(0,x_0)$
wynika z oszacowania
$|u(t,x)-g(x_0)|\le\|u(t,\cdot)-g\|_\infty+|g(x)-g(x_0)|$.
Dla dodatnich czasów wynika ze zbieżności jednostajnej
na pasach $t\ge\delta$. Nie zakładaliśmy punktowej
zbieżności zwykłych sum częściowych szeregu Fouriera $g$.

Jednoznaczność wynika z parabolicznej zasady maksimum
z rozdziału~\ref{chap:pde-heat}, zastosowanej do różnicy
i jej przeciwieństwa na każdym skończonym przedziale
czasowym. Dla funkcji zespolonych stosujemy ten argument
osobno do części rzeczywistej i urojonej.
\end{proof}

\begin{wniosek}[Ostry wykładnik zaniku w normie $L^2$]
\label{cor:pde6-decay}
Jeżeli $g\in C([0,L])$ i $g(0)=g(L)=0$, to dla
rozwiązania~\eqref{eq:pde6-heat}
\begin{equation}\label{eq:pde6-decay}
 \|u(t,\cdot)\|_{L^2(0,L)}
 \le e^{-\kappa\pi^2t/L^2}\|g\|_{L^2(0,L)},\qquad t\ge0.
\end{equation}
Stała w wykładniku jest optymalna.
\end{wniosek}
\begin{proof}
Dla $t>0$ stosujemy ortogonalność sinusów i Parsevala:
\[
\|u(t)\|_2^2=\frac L2\sum_{n\ge1}|b_n|^2e^{-2\kappa\lambda_nt}
\le e^{-2\kappa\lambda_1t}\|g\|_2^2.
\]
Dla $t=0$ jest to równość. Przy $g(x)=\sin(\pi x/L)$
zachodzi równość dla każdego $t$.
\end{proof}


\begin{stwierdzenie}[Rozwiązanie cieplne dla danych $L^2$]
\label{prop:pde6-L2-semigroup}
Niech $\kappa,L>0$, $g\in L^2(0,L)$ i niech
$g_n=\langle g,e_n\rangle_{L^2(0,L)}$ dla ortonormalnych
funkcji $e_n$ z twierdzenia~\ref{thm:pde6-spectrum}.
Szereg $u(t)=\sum_{n\ge1}e^{-\kappa\lambda_nt}g_ne_n$
wyznacza funkcję $u\in C([0,\infty);L^2(0,L))$,
przy czym $u(0)=g$, a dla każdego $t>0$ otrzymujemy
funkcję gładką względem $x\in[0,L]$, klasyczną
w $t>0$ i spełniającą zerowe warunki Dirichleta.
Zachodzi także $\|u(t)\|_2\le e^{-\kappa\lambda_1t}\|g\|_2$.
\end{stwierdzenie}
\begin{proof}
Z twierdzenia Parsevala mamy $\sum_{n\ge1}|g_n|^2<\infty$.
Dla $t\ge0$ szereg zbiega w $L^2$, a
\[
 \|u(t)-g\|_2^2=\sum_{n\ge1}
 |e^{-\kappa\lambda_nt}-1|^2|g_n|^2\longrightarrow0
 \quad(t\downarrow0)
\]
na mocy twierdzenia o zbieżności zmajoryzowanej
dla szeregów, ponieważ każdy składnik jest ograniczony
przez $|g_n|^2$. Dla dowolnej chwili $t_0\ge0$
stosujemy ten sam argument do
$\sum_{n\ge1}|e^{-\kappa\lambda_nt}-e^{-\kappa\lambda_nt_0}|^2
|g_n|^2$, z tą samą sumowalną majorantą.
Otrzymujemy ciągłość $L^2$ w czasie.
Dla $t\ge\delta>0$ nierówność Cauchy'ego--Schwarza
zapewnia jednostajną zbieżność szeregów pochodnych:
\[
 \sum_{n\ge1}|g_n|n^k e^{-\kappa\lambda_n\delta}
 \le\Bigl(\sum_{n\ge1}|g_n|^2\Bigr)^{1/2}
 \Bigl(\sum_{n\ge1}n^{2k}e^{-2\kappa\lambda_n\delta}\Bigr)^{1/2}
 <\infty.
\]
Dla pochodnej $r$ razy po $x$ i $s$ razy po $t$
stosujemy to oszacowanie z $k=r+2s$ oraz stałym
czynnikiem zależnym od $L,\kappa,r,s$. Kryterium
Weierstrassa daje jednostajną zbieżność wszystkich
szeregów pochodnych, w tym szeregu bez pochodnych.
Jego suma punktowa reprezentuje granicę w $L^2$,
gdyż zbieżność jednostajna na odcinku implikuje
zbieżność w $L^2$. Różniczkowanie wyraz po wyrazie
daje równanie i warunki brzegowe dla $t>0$.
Nie przypisujemy danym $L^2$ wartości punktowych
w chwili $t=0$. Wreszcie z Parsevala wynika
$\|u(t)\|_2^2=\sum_{n\ge1}e^{-2\kappa\lambda_nt}|g_n|^2
\le e^{-2\kappa\lambda_1t}\|g\|_2^2$.
\end{proof}

\section{Rozdzielenie zmiennych: równanie falowe i energia modów}
\label{sec:pde6-wave}

Rozważmy $u_{tt}=c^2u_{xx}$ na $(0,L)$,
z zerowymi warunkami Dirichleta oraz danymi
$u(0,x)=g(x)$, $u_t(0,x)=h(x)$.
Podstawiając $u(t,x)=T_n(t)X_n(x)$, gdzie
$X_n(x)=\sin(n\pi x/L)$ i $\lambda_n=(n\pi/L)^2$,
otrzymujemy $(T_n''+c^2\lambda_nT_n)X_n=0$.
Ponieważ $X_n$ nie jest identycznie zerowa, mamy
$T_n''+c^2\lambda_nT_n=0$ dla wszystkich $t$.
Oznaczmy $\omega_n=c\sqrt{\lambda_n}>0$. Rozwiązanie
tego równania z danymi $T_n(0)=a_n$, $T_n'(0)=b_n$
ma postać $a_n\cos(\omega_nt)+(b_n/\omega_n)\sin(\omega_nt)$,
co sprawdzamy przez różniczkowanie i podstawienie $t=0$.

\begin{twierdzenie}[Struna: reprezentacja modalna]
\label{thm:pde6-wave}
Niech $c,L>0$ i $g,h\in C_c^\infty(0,L)$.
Niech $a_n=(2/L)\int_0^Lg(x)\sin(n\pi x/L)
\,\mathrm d\lambda_1(x)$ i
$b_n=(2/L)\int_0^Lh(x)\sin(n\pi x/L)\,\mathrm d\lambda_1(x)$.
Wówczas funkcja
\begin{equation}\label{eq:pde6-wave}
 u(t,x)=\sum_{n=1}^\infty
 \left[a_n\cos(c\sqrt{\lambda_n}t)
 +\frac{b_n}{c\sqrt{\lambda_n}}
 \sin(c\sqrt{\lambda_n}t)\right]\sin(n\pi x/L)
\end{equation}
jest jedynym rozwiązaniem klasy $C^2(\R\times[0,L])$
powyższego zagadnienia wśród rozwiązań klasycznych.
Jego energia falowa
$E(t)=\frac12\int_0^L(|u_t(t,x)|^2+c^2|u_x(t,x)|^2)
\,\mathrm d\lambda_1(x)$ jest stała.
\end{twierdzenie}
\begin{proof}
Nieparzyste $2L$-okresowe przedłużenia $g,h$
są klasy $C^\infty$, ponieważ obie funkcje
zerują się w otoczeniach końców przedziału.
Po przeskalowaniu okresu lemat~\ref{lem:pde6-decay}
i zależność między współczynnikami wykładniczymi
i sinusowymi z dowodu twierdzenia~\ref{thm:pde6-spectrum}
dają dla każdego całkowitego $m\ge0$ oszacowanie
$|a_n|+|b_n|\le C_m n^{-m}$ z pewną stałą $C_m$.
Niech $A_n(t)=a_n\cos(\omega_nt)
+(b_n/\omega_n)\sin(\omega_nt)$.
Dla całkowitych $r,s\ge0$ pochodna $r$ razy po $x$
i $s$ razy po $t$ wyrazu szeregu jest ograniczona
w module przez
\[
 (n\pi/L)^r\omega_n^s
 \bigl(|a_n|+|b_n|/\omega_n\bigr).
\]
Ponieważ $\omega_n=cn\pi/L$, wybór $m\ge r+s+2$
daje sumowalną majorantę niezależną od $t,x$.
Kryterium Weierstrassa i twierdzenie o różniczkowaniu
szeregów funkcji dają gładkość sumy oraz możliwość
różniczkowania wyraz po wyrazie. Każdy wyraz spełnia
równanie falowe i warunki brzegowe. W chwili $t=0$
szeregi dla $u$ i $u_t$ są odpowiednio szeregami
sinusowymi $g$ i $h$, ponieważ $A_n(0)=a_n$
i $A_n'(0)=b_n$. Z zupełności układu sinusowego
wynika ich zbieżność w $L^2(0,L)$ do $g,h$;
zbieżność jednostajna daje ciągłe sumy, więc równości
z $g,h$ zachodzą w każdym punkcie $[0,L]$.

Dla dowolnego rozwiązania klasy $C^2(\R\times[0,L])$
z zerowymi warunkami Dirichleta różniczkowanie energii
pod znakiem całki i całkowanie przez części dają
\[
 \begin{aligned}
 E'(t)
 &=\operatorname{Re}\int_0^L
 (u_{tt}\overline{u_t}+c^2u_x\overline{u_{xt}})
 \,\mathrm d\lambda_1(x)\\
 &=c^2\operatorname{Re}[u_x\overline{u_t}]_0^L=0.
 \end{aligned}
\]
W drugim kroku użyliśmy równania falowego,
równości pochodnych mieszanych i $u_t(t,0)=u_t(t,L)=0$,
wynikającej z różniczkowania warunków brzegowych po $t$.
Rachunek jest wersją tożsamości energii z rozdziału
~\ref{chap:pde-wave}, poprawną także dla funkcji zespolonych.
Jeżeli $w$ jest różnicą dwóch rozwiązań o tych samych danych,
to $w_t(0,x)=w_x(0,x)=0$, a więc $E_w(0)=0$.
Z zachowania nieujemnej energii i ciągłości pochodnych
wynika $w_t=w_x=0$ dla wszystkich $t,x$. Ponieważ
$w(t,0)=0$, otrzymujemy $w(t,x)=0$.

Możemy również obliczyć energię modalną. Ortogonalność
sinusów oraz kosinusów o dodatnich indeksach daje
\[
 E(t)=\frac L4\sum_{n\ge1}
 \bigl(|A_n'(t)|^2+\omega_n^2|A_n(t)|^2\bigr).
\]
Najpierw stosujemy ortogonalność do sum skończonych,
a potem przechodzimy do granicy dzięki jednostajnej
zbieżności szeregów $u_t,u_x$, a zatem zbieżności w $L^2$.
Rozpisując kwadraty modułów i korzystając z
$\cos^2(\omega_nt)+\sin^2(\omega_nt)=1$, otrzymujemy
$|A_n'|^2+\omega_n^2|A_n|^2
=|b_n|^2+\omega_n^2|a_n|^2$:
wyrazy mieszane z $\operatorname{Re}(a_n\overline{b_n})$
znoszą się. W konsekwencji
\[
 E(t)=\frac L4\sum_{n\ge1}
 (|b_n|^2+\omega_n^2|a_n|^2)=E(0).
\]
Zachowanie energii dotyczy zatem każdego modu osobno,
nie tylko sumy wszystkich modów.
\end{proof}

\begin{przyklad}[Wygładzanie a brak wygładzania]
\label{ex:pde6-smoothing}
Dla równania ciepła amplituda modu $n$ zanika jak
$e^{-\kappa(n\pi/L)^2t}$. Dla równania falowego
współczynnik czasowy modu oscyluje, może chwilowo
znikać, lecz energia nietrywialnego modu pozostaje
dodatnia i stała. Na przykład dla
$g(x)=\sin(n\pi x/L)$, $h=0$ otrzymujemy odpowiednio
$e^{-\kappa\lambda_nt}g(x)$ i
$\cos(c\sqrt{\lambda_n}t)g(x)$.
Chociaż ta funkcja $g$ nie ma zwartego nośnika w $(0,L)$,
oba skończone wzory są rozwiązaniami klasycznymi,
co sprawdzamy przez bezpośrednie różniczkowanie.
Dla równania ciepła w ustalonej chwili $t>0$
mnożnik zanika szybciej niż dowolna potęga $n^{-1}$.
Dla równania falowego w chwili $t=2L/c$ wszystkie
mnożniki $\cos(c\sqrt{\lambda_n}t)$ są równe $1$.
Nie ma więc analogicznego tłumienia wysokich
częstości dla całej ewolucji falowej. Ten kontrast
wyjaśnia spektralnie wygładzanie paraboliczne.
\end{przyklad}

\section{Niejednorodne równania i wzór modalny Duhamela}
\label{sec:pde6-forcing}

Niech $\kappa,L,T>0$. Rozważmy zagadnienie
\[
 \begin{cases}
 u_t-\kappa u_{xx}=f(t,x),&0<t<T,\quad 0<x<L,\\
 u(t,0)=u(t,L)=0,&0\le t\le T,\\
 u(0,x)=g(x),&0\le x\le L.
 \end{cases}
\]
Zakładamy, że $f\in C^\infty([0,T]\times[0,L])$
oraz istnieje zbiór zwarty $K\subset(0,L)$ taki, że
$\operatorname{supp}f(t,\cdot)\subset K$ dla każdego $t\in[0,T]$.
Ponadto niech $g\in C_c^\infty(0,L)$.
Oznaczmy $X_n(x)=\sin(n\pi x/L)$ i $\lambda_n=(n\pi/L)^2$
dla $n\ge1$. Wprowadzamy współczynniki sinusowe
\[
 \begin{aligned}
 f_n(t)&=\frac2L\int_0^L f(t,x)X_n(x)\,\mathrm d\lambda_1(x),\\
 g_n&=\frac2L\int_0^L g(x)X_n(x)\,\mathrm d\lambda_1(x).
 \end{aligned}
\]
Stosujemy więc współczynniki względem funkcji $X_n$,
a nie względem unormowanych funkcji własnych $e_n$.

Najpierw załóżmy, że istnieje rozwiązanie o ciągłych
pochodnych $u_t,u_x,u_{xx}$ na domkniętym prostokącie.
Niech
$b_n(t)=(2/L)\int_0^L u(t,x)X_n(x)\,\mathrm d\lambda_1(x)$.
Ciągłość $u_t$ na zbiorze zwartym pozwala różniczkować
pod znakiem całki. Dwukrotne całkowanie przez części daje
\[
 \begin{aligned}
 \int_0^L u_{xx}X_n\,\mathrm d\lambda_1
 &=[u_xX_n-uX_n']_0^L
   +\int_0^L u X_n''\,\mathrm d\lambda_1\\
 &=-\lambda_n\int_0^L u X_n\,\mathrm d\lambda_1 .
 \end{aligned}
\]
Wyrazy brzegowe znikają, ponieważ zarówno $X_n$,
jak i $u(t,\cdot)$ spełniają zerowe warunki Dirichleta.
Mnożąc równanie przez $(2/L)X_n$ i całkując względem $x$,
otrzymujemy zatem
$b_n'(t)+\kappa\lambda_nb_n(t)=f_n(t)$ i $b_n(0)=g_n$.
Po pomnożeniu przez $e^{\kappa\lambda_nt}$ dostajemy
$(e^{\kappa\lambda_nt}b_n(t))'=e^{\kappa\lambda_nt}f_n(t)$.
Całkowanie od $0$ do $t$ i pomnożenie przez
$e^{-\kappa\lambda_nt}$ prowadzą do wzoru
\begin{equation}\label{eq:pde6-duhamel}
 b_n(t)=e^{-\kappa\lambda_nt}g_n
 +\int_0^t e^{-\kappa\lambda_n(t-s)}f_n(s)
 \,\mathrm d\lambda_1(s).
\end{equation}

Teraz konstruujemy rozwiązanie, nie zakładając jego
istnienia. Definiujemy $b_n$ przez~\eqref{eq:pde6-duhamel}.
Pokażemy, że
$u(t,x)=\sum_{n\ge1}b_n(t)X_n(x)$
jest rozwiązaniem klasy $C^\infty$ na domkniętym
prostokącie (z pochodnymi na brzegu rozumianymi
jako jednostronne lub ich ciągłe przedłużenia).

Dla dowolnych całkowitych $m,j\ge0$ mamy
\[
 |g_n|\le C_m n^{-m},\qquad
 \sup_{t\in[0,T]}|f_n^{(j)}(t)|\le C_{m,j}n^{-m}.
\]
Istotnie, można różniczkować $f_n$ względem czasu
pod znakiem całki. Wszystkie pochodne $f$ i $g$
znikają w ustalonych otoczeniach obu końców odcinka.
Wielokrotne całkowanie przez części względem $x$
nie daje więc wyrazów brzegowych i za każdym razem
wprowadza czynnik $L/(n\pi)$. Normy pochodnych $f$
są ograniczone jednostajnie względem czasu na
domkniętym prostokącie.

Z nieujemności $\kappa\lambda_n$ i $0\le s\le t\le T$
wynika
\[
 \sup_{0\le t\le T}|b_n(t)|
 \le |g_n|+T\sup_{0\le t\le T}|f_n(t)|.
\]
Zatem $b_n$ również zanika szybciej niż dowolna
odwrotna potęga $n$, jednostajnie względem $t$.
Różniczkując~\eqref{eq:pde6-duhamel} według reguły
Leibniza, otrzymujemy $b_n'=f_n-\kappa\lambda_nb_n$.
Dla $j\ge0$ mamy ponadto
$b_n^{(j+1)}=f_n^{(j)}-\kappa\lambda_nb_n^{(j)}$.
Indukcja daje szybkie zanikanie każdej pochodnej
czasowej: jeżeli $b_n^{(j)}$ spełnia oszacowanie
dla każdej potęgi $n^{-m}$, to w wyrazie
$\lambda_nb_n^{(j)}$ stosujemy je z wykładnikiem $m+2$,
ponieważ $\lambda_n=(n\pi/L)^2$.

Dla dowolnych całkowitych $r,j\ge0$ moduł pochodnej
$\partial_x^r\partial_t^j(b_n(t)X_n(x))$ jest
ograniczony przez $(n\pi/L)^r\sup_t|b_n^{(j)}(t)|$.
Wybierając oszacowanie z $m\ge r+2$, otrzymujemy
sumowalną majorantę niezależną od $(t,x)$.
Kryterium Weierstrassa i twierdzenie o różniczkowaniu
szeregów funkcji uzasadniają zbieżność jednostajną
szeregu oraz wszystkich szeregów jego pochodnych.
Zatem
\[
 u_t-\kappa u_{xx}
 =\sum_{n\ge1}(b_n'+\kappa\lambda_nb_n)X_n
 =\sum_{n\ge1}f_n(t)X_n=f(t,\cdot).
\]
Ostatnia równość wynika z zupełności układu sinusowego
(twierdzenie~\ref{thm:pde6-spectrum}) w $L^2(0,L)$
i ciągłości obu stron, gdyż szereg z $f_n(t)$ zbiega
jednostajnie na domkniętym odcinku.
Analogicznie, $b_n(0)=g_n$ i zbieżność szeregu
sinusowego do $g$ w $L^2$ oraz jednostajnie
dają $u(0,x)=g(x)$ dla wszystkich $x\in[0,L]$.
Warunki brzegowe zachodzą wyraz po wyrazie.

Jednoznaczność w klasie rozwiązań ciągłych na domkniętym
prostokącie i klasy $C^{1,2}$ we wnętrzu wynika
z twierdzenia~\ref{thm:pde6-heat}: różnica dwóch rozwiązań
spełnia jednorodne równanie ciepła z zerowymi danymi.
Jeżeli $P_t$ oznacza operator rozwiązania jednorodnego
z sekcji~\ref{sec:pde6-heat}, przy $P_0g=g$, to wzór
modalny można zapisać również jako
\[
 u(t,x)=(P_tg)(x)+
 \int_0^t(P_{t-s}f(s,\cdot))(x)\,\mathrm d\lambda_1(s).
\]
Istotnie, szereg dla $P_{t-s}f(s,\cdot)$ jest jednostajnie
zbieżny dla $0\le s\le t\le T$ i $0\le x\le L$, ponieważ
$e^{-\kappa\lambda_n(t-s)}\le1$ oraz
$\sup_s|f_n(s)|\le C_{2,0}n^{-2}$.
Całkowanie wyraz po wyrazie daje dokładnie
\eqref{eq:pde6-duhamel}. Jest to zasada Duhamela
dla rozważanego zagadnienia Dirichleta.

\section{Regularne zagadnienie Sturma--Liouville'a}
\label{sec:pde6-sturm}

Przy rozdzieleniu zmiennych w równaniach ze
zmiennymi współczynnikami otrzymujemy często operator
$\mathcal L y=-(py')'+qy$. Zdefiniujmy dokładnie
zagadnienie i przestrzeń, w której badamy funkcje własne.

\begin{definicja}[Regularne zagadnienie Dirichleta]
\label{def:pde6-sturm}
Niech $L>0$, $p\in C^1([0,L];\R)$,
$q,w\in C([0,L];\R)$, $p(x)>0$, $w(x)>0$
i $q(x)\ge0$ na $[0,L]$. Zagadnieniem
Sturma--Liouville'a z warunkami Dirichleta nazywamy
\begin{equation}\label{eq:pde6-sturm}
 -(p(x)y'(x))'+q(x)y(x)=\lambda w(x)y(x),
 \quad y(0)=y(L)=0.
\end{equation}
Rozwiązaniem jest $y\in C^2([0,L])$ spełniające
równanie punktowo. Wartością własną nazywamy liczbę
$\lambda\in\mathbb C$, dla której istnieje rozwiązanie
$y\not\equiv0$; funkcje mogą być zespolone. Pracujemy w przestrzeni
$H=L^2((0,L),w(x)\,\mathrm d\lambda_1(x))$ z
iloczynem skalarnym $\langle f,g\rangle_H
=\int_0^Lf(x)\overline{g(x)}w(x)\,\mathrm d\lambda_1(x)$.
Z ciągłości i dodatniości $w$ na zwartym odcinku wynika
$0<\min_{[0,L]}w\le\max_{[0,L]}w<\infty$.
Norma $H$ jest więc równoważna zwykłej normie $L^2$,
a $H$ jest przestrzenią Hilberta.
\end{definicja}

\begin{lemat}[Tożsamość Lagrange'a i rzeczywistość widma]
\label{lem:pde6-lagrange}
Niech $y,z\in C^2([0,L])$. Wtedy
\begin{equation}\label{eq:pde6-lagrange}
 \int_0^L\bigl((\mathcal Ly)(x)\overline{z(x)}
 -y(x)\overline{(\mathcal Lz)(x)}\bigr)\,\mathrm d\lambda_1(x)
 =\left[p\left(y\overline{z'}-y'\overline z\right)\right]_0^L.
\end{equation}
Jeżeli $y\not\equiv0$ jest funkcją własną
z jednorodnymi warunkami Dirichleta, to $\lambda>0$
jest rzeczywista. Funkcje własne dla różnych
wartości własnych są ortogonalne w $H$.
\end{lemat}
\begin{proof}
Całkowanie przez części, przy rzeczywistych $p,q$, daje
\[
 \int_0^L(\mathcal Ly)\overline z\,\mathrm d\lambda_1
 =[-py'\overline z]_0^L
 +\int_0^L(py'\overline{z'}+qy\overline z)\,
 \mathrm d\lambda_1.
\]
Odejmując analogiczną równość dla
$\int_0^L y\overline{\mathcal Lz}\,\mathrm d\lambda_1$,
otrzymujemy~\eqref{eq:pde6-lagrange}.
Dla funkcji własnej $y$ z warunkami Dirichleta
powyższa równość z $z=y$ daje
\[
 \lambda\|y\|_H^2
 =\int_0^L(p(x)|y'(x)|^2+q(x)|y(x)|^2)\,
 \mathrm d\lambda_1(x)>0.
\]
Istotnie, gdyby prawa strona była zerem, dodatniość $p$,
nieujemność $q$ i ciągłość $y'$ wymuszałyby $y'=0$
na całym odcinku, a $y(0)=0$ dawałoby $y=0$.
Ponieważ $\|y\|_H^2>0$, liczba $\lambda$ jest
rzeczywista i dodatnia.
Dla pary funkcji własnych $y,z$ o wartościach
$\lambda,\mu$ stosujemy tożsamość Lagrange'a;
ze znikania brzegów i rzeczywistości liczb własnych
wynika $(\lambda-\mu)\langle y,z\rangle_H=0$.
\end{proof}

\begin{twierdzenie}[Dyskretność i zupełność widma]
\label{thm:pde6-sturm}
Przy założeniach definicji~\ref{def:pde6-sturm}
istnieje rosnący do nieskończoności ciąg
$0<\lambda_1<\lambda_2<\cdots$ i odpowiednie
rzeczywiste funkcje własne $y_n\in C^2([0,L])$,
które po unormowaniu tworzą bazę ortonormalną $H$.
Są to wszystkie wartości własne zagadnienia,
a każda z nich jest prosta.
\end{twierdzenie}
\begin{proof}
Skonstruujemy zwarty operator samosprzężony na $H$,
będący odwrotnością zagadnienia Dirichleta z wagą.
Użyjemy twierdzenia spektralnego dla zwartych
operatorów samosprzężonych w przestrzeni Hilberta:
istnieje baza ortonormalna złożona z wektorów własnych
o niezerowych wartościach własnych oraz bazy jądra;
niezerowe wartości własne mają skończone krotności
i mogą gromadzić się jedynie w zerze.
Jest to jawnie wykorzystywany wynik analizy funkcjonalnej.

\emph{Konstrukcja funkcji Greena.}
Niech $v_1,v_2$ będą rzeczywistymi rozwiązaniami
$\mathcal Lv=0$ o danych
$v_1(0)=0$, $v_1'(0)=1$ oraz
$v_2(L)=0$, $v_2'(L)=-1$.
Istnienie i jednoznaczność wynikają z twierdzenia
dla liniowych układów ODE o ciągłych współczynnikach:
równanie zapisujemy jako układ $v'=z/p$, $z'=qv$,
gdzie $z=pv'$.
Nietrywialne rozwiązanie równania jednorodnego
nie może znikać na obu końcach, ponieważ rachunek
energii z lematu~\ref{lem:pde6-lagrange} przy $\lambda=0$
wymusza wtedy $v=0$.
W szczególności $v_1(L)\ne0$.
Z równań $(pv_1')'=qv_1$, $(pv_2')'=qv_2$ wynika
\[
 [p(v_1'v_2-v_1v_2')]'=qv_1v_2-v_1qv_2=0.
\]
Zatem $C=p(s)(v_1'(s)v_2(s)-v_1(s)v_2'(s))$
jest stałą, a obliczenie w $s=L$ daje
$C=p(L)v_1(L)\ne0$.
Definiujemy
\begin{equation}\label{eq:pde6-green}
 G(x,s)=\frac{v_1(\min\{x,s\})v_2(\max\{x,s\})}{C}.
\end{equation}
Jądro jest rzeczywiste, ciągłe i symetryczne na
$[0,L]^2$ oraz $G(0,s)=G(L,s)=0$.
Dla $s\in(0,L)$ spełnia równanie jednorodne po obu
stronach punktu $x=s$. W tym punkcie zachodzi
\[
 -p(s)(\partial_xG(s^+,s)-\partial_xG(s^-,s))
 =\frac{p(s)(v_1'(s)v_2(s)-v_1(s)v_2'(s))}{C}=1.
\]
Jest to warunek skoku funkcji Greena, wyrażający
$\mathcal L_xG(x,s)=\delta_s$ we wnętrzu odcinka.
Nie potrzebujemy tu ogólnej teorii dystrybucji:
własność odwrotności sprawdzimy bezpośrednio.

\emph{Ograniczoność, zwartość i samosprzężoność.}
Niech $\mathrm d\nu(s)=w(s)\,\mathrm d\lambda_1(s)$
oraz $M=\nu((0,L))>0$. Definiujemy
$Tf(x)=\int_0^LG(x,s)f(s)\,\mathrm d\nu(s)$.
Ponieważ $\|f\|_{L^1(\nu)}\le\sqrt M\|f\|_H$,
całka jest bezwzględnie zbieżna dla każdego $x$,
nie zależy od reprezentanta $f$, a jednostajna
ciągłość $G$ daje ciągłość $Tf$ względem $x$.
Ponadto $\|Tf\|_\infty\le\|G\|_\infty\sqrt M\|f\|_H$,
skąd $\|Tf\|_H\le M\|G\|_\infty\|f\|_H$.

Wybieramy podziały $x_j=jL/N$, $j=0,\ldots,N$,
oraz ciągłe odcinkami liniowe funkcje kapeluszowe
$\theta_j$ przyjmujące wartości $\delta_{ij}$ w węzłach $x_i$.
Są nieujemne, sumują się do $1$, a $\theta_j(x)>0$
implikuje $|x-x_j|\le L/N$.
Interpolant
$G_N(x,s)=\sum_{i,j=0}^NG(x_i,x_j)\theta_i(x)\theta_j(s)$
jest kombinacją wypukłą wartości $G$ w węzłach
odległych od $(x,s)$ o co najwyżej $\sqrt2 L/N$.
Z jednostajnej ciągłości $G$ wynika
$\|G_N-G\|_\infty\to0$.
Operator $T_N$ o jądrze $G_N$ ma obraz zawarty
w przestrzeni rozpiętej przez $\theta_0,\ldots,\theta_N$,
a więc ma skończony rząd. Powyższe oszacowanie
zastosowane do różnicy jąder daje
$\|T_N-T\|_{H\to H}\le M\|G_N-G\|_\infty\to0$.
Korzystamy z faktu, że granica operatorów zwartych
w normie operatorowej jest zwarta, więc $T$ jest zwarty.

Dla $f,h\in H$ całka podwójna z modułu
$G(x,s)f(s)\overline{h(x)}$ względem $\nu\otimes\nu$
jest skończona, bo nie przekracza
$\|G\|_\infty\|f\|_{L^1(\nu)}\|h\|_{L^1(\nu)}$.
Z twierdzenia Fubiniego i rzeczywistej symetrii jądra
otrzymujemy
\[
 \langle Tf,h\rangle_H
 =\int_0^L\int_0^L G(x,s)f(s)\overline{h(x)}
 \,\mathrm d\nu(s)\,\mathrm d\nu(x)
 =\langle f,Th\rangle_H.
\]
Operator jest ograniczony i określony na całym $H$,
a więc jest samosprzężony.

\emph{Regularność i własność odwrotności.}
Dla $f\in H$ funkcja $F=wf$ należy do $L^1(0,L)$.
Wprowadzamy absolutnie ciągłe funkcje
\[
 A(x)=\int_0^x v_1(s)F(s)\,\mathrm d\lambda_1(s),
 \qquad
 B(x)=\int_x^L v_2(s)F(s)\,\mathrm d\lambda_1(s).
\]
Mamy $A'=v_1F$, $B'=-v_2F$ prawie wszędzie.
Z definicji jądra funkcja $y=Tf$ ma postać
$y=(v_2 A+v_1 B)/C$, zatem jest absolutnie ciągła
i znika w $0,L$.
Różniczkowanie prawie wszędzie daje
$y'=(v_2'A+v_1'B)/C$, ponieważ składniki z $A',B'$
znoszą się. Prawa strona jest ciągła, więc z twierdzenia
podstawowego rachunku całkowego dla funkcji absolutnie
ciągłych wynika, że $y$ jest klasy $C^1$ i wzór na $y'$
zachodzi wszędzie (na końcach jednostronnie).
Funkcja $py'=(pv_2'A+pv_1'B)/C$ jest absolutnie ciągła.
Różniczkując ją prawie wszędzie i używając równań
jednorodnych dla $v_1,v_2$, dostajemy
\[
 \begin{aligned}
 (py')'
 &=q(v_2A+v_1B)/C\\
 &\quad +p(v_2'v_1-v_1'v_2)F/C
 =qy-F.
 \end{aligned}
\]
Zatem $\mathcal L(Tf)=wf$ prawie wszędzie.
Jeśli $F$ ma ciągły reprezentant, prawa strona wzoru
na $(py')'$ jest ciągła, więc $py'\in C^1$, a
$p\in C^1$ i $p>0$ dają $y\in C^2$.
Dla dowolnego $f\in H$ nie twierdzimy jeszcze,
że $Tf$ jest klasy $C^2$.
Dla ciągłego wymuszenia $F$ powyższa konstrukcja
z $f=F/w$ daje rozwiązanie klasy $C^2$ równania
$\mathcal Ly=F$; jednoznaczność wynika z rachunku energii
dla równania jednorodnego.

\emph{Dodatniość i przejście do wartości własnych.}
Całkowanie przez części jest dozwolone, ponieważ
$py'$ jest absolutnie ciągła, $y\in C^1$ i $y(0)=y(L)=0$.
Z równania dla $y=Tf$ wynika
\[
 \langle Tf,f\rangle_H
 =\int_0^L y\overline{\mathcal Ly}\,\mathrm d\lambda_1
 =\int_0^L(p|y'|^2+q|y|^2)\,\mathrm d\lambda_1\ge0.
\]
Jeśli energia jest zerowa, ciągłość $y'$ i dodatniość $p$
dają $y'=0$, a warunki brzegowe dają $y=0$.
Wtedy $\mathcal Ly=wf=0$ prawie wszędzie, czyli $f=0$.
Zatem $T$ jest ściśle dodatni i ma trywialne jądro.

Twierdzenie spektralne daje bazę ortonormalną
wektorów własnych $Ty_n=\mu_n y_n$, gdzie $\mu_n>0$.
Możemy wybrać rzeczywiste wektory własne, ponieważ
rzeczywistość jądra daje $T\overline f=\overline{Tf}$.
Istotnie, rzeczywista i urojona część każdego wektora
własnego należą do tej samej przestrzeni własnej;
ortonormalizację można przeprowadzić na jej części
rzeczywistej, która rozpina ją nad $\mathbb C$.
Każda przestrzeń własna odpowiadająca $\mu_n\ne0$
jest skończenie wymiarowa.
Przestrzeń $H$ jest nieskończenie wymiarowa (na przykład
zawiera dowolnie wiele liniowo niezależnych funkcji
charakterystycznych rozłącznych podprzedziałów).
Ponieważ jądro jest trywialne, niezerowych wartości
własnych jest nieskończenie wiele i można je uporządkować
nierosnąco z krotnościami, $\mu_n\downarrow0$.

Wybierając reprezentanta $y_n=\mu_n^{-1}Ty_n$,
otrzymujemy $y_n\in C^1([0,L])$ i warunki Dirichleta.
Równanie $\mathcal L(Ty_n)=wy_n$ daje
$\mathcal Ly_n=\mu_n^{-1}wy_n$ prawie wszędzie.
Wymuszenie $\mu_n^{-1}wy_n$ jest ciągłe, więc
z wykazanej regularności wynika $y_n\in C^2([0,L])$
oraz równanie punktowe. Stąd $\lambda_n=\mu_n^{-1}$
są wartościami własnymi rozważanego zagadnienia.
Odwrotnie, jeśli $\mathcal Ly=\lambda wy$ dla
nietrywialnego $y$ z warunkami Dirichleta, to
$\lambda>0$ z lematu~\ref{lem:pde6-lagrange}.
Funkcje $T(\lambda y)$ i $y$ rozwiązują to samo
zagadnienie z ciągłym wymuszeniem $\lambda wy$;
jego jednoznaczność daje $Ty=\lambda^{-1}y$.
Nie ma więc innych wartości własnych.

Dla ustalonej $\lambda$ rozwiązanie równania
$(py')'=(q-\lambda w)y$ z $y(0)=0$ jest jednoznacznie
wyznaczone przez $y'(0)$, z twierdzenia o liniowych
układach ODE. Gdyby również $y'(0)=0$, rozwiązanie
byłoby zerowe. Dwa nietrywialne rozwiązania z zerowym
warunkiem w $0$ są więc proporcjonalne, co daje
prostotę wartości własnych. Ostatecznie możemy
uporządkować je ściśle rosnąco,
$0<\lambda_1<\lambda_2<\cdots\to\infty$,
a odpowiadające im unormowane funkcje własne
tworzą bazę ortonormalną $H$.
\end{proof}

\begin{uwaga}[Co jest tu nowe?]
Dla stałych współczynników mieliśmy jawne funkcje
$\sin(n\pi x/L)$; przy zmiennych współczynnikach
możemy zwykle nie znać ich w postaci elementarnej,
ale nadal mamy dyskretne wartości własne oraz
zupełny układ modów. To właśnie uzasadnia metodę
rozdzielenia zmiennych również bez jawnych wzorów.
Twierdzenie o oscylacji Sturma, mówiące o liczbie
zer kolejnych funkcji własnych, omówimy osobno,
gdy będziemy dysponować szerszą teorią ODE i
zagadnień spektralnych.
\end{uwaga}

\section{Wyższe wymiary: prostokąt i geometria widma}
\label{sec:pde6-rectangle}

Niech $L_1,L_2>0$ i $\Omega=(0,L_1)\times(0,L_2)$.
Rozważamy klasyczne zagadnienie własne Dirichleta
$-\Delta X=\lambda X$ w $\Omega$, $X|_{\partial\Omega}=0$,
dla $X\in C^2(\overline\Omega)$.
Pojedyncze rozwiązania rozdzielone szukamy w postaci
$X(x,y)=U(x)V(y)$. Jeśli iloczyn jest nietrywialny,
to wybór punktu $y_*$, w którym $V(y_*)\ne0$, w równaniu
$-U''V-UV''=\lambda UV$ daje
$-U''=\alpha U$, gdzie
$\alpha=\lambda+V''(y_*)/V(y_*)$.
Następnie wybór $x_*$ z $U(x_*)\ne0$ daje
$-V''=\beta V$, gdzie $\beta=\lambda-\alpha$.
Warunki brzegowe wymuszają
$U(0)=U(L_1)=V(0)=V(L_2)=0$.
Z twierdzenia~\ref{thm:pde6-spectrum} wynika więc,
że z dokładnością do stałego czynnika otrzymujemy
\begin{equation}\label{eq:pde6-rectangle}
 X_{mn}(x,y)=\sin\left(\frac{m\pi x}{L_1}\right)
 \sin\left(\frac{n\pi y}{L_2}\right),\quad
 \lambda_{mn}=\pi^2\left(\frac{m^2}{L_1^2}
 +\frac{n^2}{L_2^2}\right),\quad m,n\ge1.
\end{equation}
Odwrotnie, każda z tych funkcji spełnia równanie
i wszystkie warunki brzegowe, co wynika bezpośrednio
z różniczkowania obu czynników. Pozostaje wykazać,
że rozdzielenie zmiennych nie pomija żadnych wartości
własnych ani funkcji własnych.

Niech $e_m^{(1)}$ i $e_n^{(2)}$ oznaczają unormowane
funkcje sinusowe na odcinkach $(0,L_1)$ i $(0,L_2)$:
\[
 e_m^{(1)}(x)=\sqrt{2/L_1}\sin(m\pi x/L_1),\qquad
 e_n^{(2)}(y)=\sqrt{2/L_2}\sin(n\pi y/L_2).
\]
Piszemy $E_{mn}(x,y)=e_m^{(1)}(x)e_n^{(2)}(y)$,
czyli $E_{mn}=2(L_1L_2)^{-1/2}X_{mn}$.
Z twierdzenia Fubiniego i ortonormalności baz
jednowymiarowych otrzymujemy
\[
 \langle E_{mn},E_{ij}\rangle_{L^2(\Omega)}
 =\delta_{mi}\delta_{nj}.
\]
Tutaj całkujemy względem zwykłej miary Lebesgue'a
$\lambda_2$ na prostokącie.

Wykażemy zupełność tej rodziny. Niech $f\in L^2(\Omega)$
będzie ortogonalna do wszystkich $E_{mn}$.
Z twierdzenia Fubiniego dla prawie każdego $y$
funkcja $f(\cdot,y)$ należy do $L^2(0,L_1)$.
Dla $m\ge1$ określamy prawie wszędzie
\[
 F_m(y)=\int_0^{L_1} f(x,y)\overline{e_m^{(1)}(x)}
 \,\mathrm d\lambda_1(x).
\]
Z nierówności Cauchy'ego--Schwarza i unormowania
$e_m^{(1)}$ wynika
$|F_m(y)|^2\le\int_0^{L_1}|f(x,y)|^2\,\mathrm d\lambda_1(x)$.
W szczególności $F_m\in L^2(0,L_2)$.
Iloczyn $f\overline{E_{mn}}$ jest całkowalny na
prostokącie, ponownie z nierówności Cauchy'ego--Schwarza.
Z twierdzenia Fubiniego i założonej ortogonalności mamy
\[
 \langle F_m,e_n^{(2)}\rangle_{L^2(0,L_2)}
 =\langle f,E_{mn}\rangle_{L^2(\Omega)}=0
 \quad(n\ge1).
\]
Z zupełności bazy sinusowej na $(0,L_2)$ wynika
$F_m=0$ prawie wszędzie. Ponieważ zbiór indeksów $m$
jest przeliczalny, istnieje jeden zbiór pełnej miary
w $(0,L_2)$, na którym wszystkie $F_m(y)$ znikają
i $f(\cdot,y)\in L^2(0,L_1)$.
Dla każdego takiego $y$ zupełność bazy na $(0,L_1)$
daje $f(\cdot,y)=0$ w $L^2(0,L_1)$.
Ponowne zastosowanie twierdzenia Fubiniego daje
$\|f\|_{L^2(\Omega)}=0$. Dopełnienie ortogonalne
rozpatrywanej rodziny jest więc trywialne, co
dowodzi, że $(E_{mn})_{m,n\ge1}$ jest bazą ortonormalną.

Niech teraz $X\not\equiv0$ będzie dowolną klasyczną
funkcją własną z warunkami Dirichleta.
Dwukrotne całkowanie przez części osobno względem
każdej ze zmiennych, przy ustalonej drugiej zmiennej,
daje
\[
 \langle-\Delta X,E_{mn}\rangle
 =\langle X,-\Delta E_{mn}\rangle.
\]
Wyrazy brzegowe znikają, ponieważ $X$ oraz $E_{mn}$
zerują się na wszystkich bokach prostokąta.
Stąd $(\lambda-\lambda_{mn})\langle X,E_{mn}\rangle=0$
dla wszystkich $m,n$. Co najmniej jeden współczynnik
jest niezerowy z zupełności bazy, więc $\lambda$
jest jedną z liczb~\eqref{eq:pde6-rectangle}.
Dla ustalonego $\lambda$ istnieje tylko skończenie
wiele par $(m,n)$ z $\lambda_{mn}=\lambda$, gdyż
$m\le L_1\sqrt\lambda/\pi$ i $n\le L_2\sqrt\lambda/\pi$.
Rozwinięcie $X$ w bazie zawiera wyłącznie te wyrazy,
a równość w $L^2$ jest równością punktową dzięki
ciągłości. Zatem przestrzeń własna jest dokładnie
rozpięta przez $E_{mn}$ odpowiadające tym parom,
a jej wymiar jest liczbą takich par.
W szczególności w każdym ograniczonym przedziale
liczbowym jest tylko skończenie wiele wartości własnych.

\begin{uwaga}[Geometria prostokąta a krotności]
Dla kwadratu $L_1=L_2=L$ mamy na przykład
$\lambda_{12}=\lambda_{21}=5\pi^2/L^2$.
Są to jedyne dodatnie pary całkowite z $m^2+n^2=5$,
więc ta wartość własna ma krotność dokładnie $2$.
Jeżeli natomiast $L_1^2/L_2^2$ jest niewymierne,
to z równości $\lambda_{mn}=\lambda_{ij}$ wynika
\[
 m^2-i^2=\frac{L_1^2}{L_2^2}(j^2-n^2).
\]
Lewa strona jest całkowita. Gdyby $j^2-n^2\ne0$,
stosunek $L_1^2/L_2^2$ byłby wymierny. Zatem
$j=n$, a następnie $m=i$, bo wszystkie indeksy
są dodatnie. W tym przypadku każda wartość własna
jest prosta.
\end{uwaga}

\section{Dlaczego nie można bezwarunkowo odwrócić ewolucji cieplnej?}
\label{sec:pde6-backward}

Niech $\kappa,L,T>0$ i $H=L^2(0,L)$ z normą względem
miary Lebesgue'a. Rozważmy operator $P_T:H\to H$,
przyporządkowujący danym początkowym $g$ rozwiązanie
równania ciepła z jednorodnymi warunkami Dirichleta
w chwili $T$. Z twierdzenia~\ref{thm:pde6-spectrum}
i stwierdzenia~\ref{prop:pde6-L2-semigroup} wynika
\[
 P_Tg=\sum_{n\ge1} e^{-\kappa\lambda_nT}g_ne_n,\qquad
 g_n=\langle g,e_n\rangle_H,\qquad
 \lambda_n=(n\pi/L)^2.
\]
Szereg zbiega w $H$, a operator jest liniowy i ograniczony:
z tożsamości Parsevala wynika
$\|P_Tg\|_H\le e^{-\kappa\lambda_1T}\|g\|_H$.

Rozważmy zagadnienie odwrotne: dla danych końcowych
$h\in H$ znaleźć $g\in H$ takie, że $P_Tg=h$.
Należy rozróżnić trzy aspekty poprawności zagadnienia
w sensie Hadamarda.

\emph{Jednoznaczność.}
Wszystkie mnożniki $e^{-\kappa\lambda_nT}$ są dodatnie.
Jeżeli $P_Tg=0$, to każdy współczynnik $g_n$ jest zerowy,
a zupełność bazy daje $g=0$. Operator $P_T$ jest zatem
iniektywny, więc dla każdego $h$ z jego obrazu dane
początkowe są wyznaczone jednoznacznie.

\emph{Istnienie.}
Dla $h_n=\langle h,e_n\rangle_H$ warunek $P_Tg=h$
wymusza $g_n=e^{\kappa\lambda_nT}h_n$.
Z Parsevala wynika warunek konieczny
\[
 \sum_{n\ge1} e^{2\kappa\lambda_nT}|h_n|^2<\infty.
\]
Jest on również wystarczający: jeśli zachodzi,
to ortonormalność i zupełność przestrzeni Hilberta
zapewniają zbieżność szeregu
$g=\sum_{n\ge1}e^{\kappa\lambda_nT}h_ne_n$ w $H$.
Wstawienie jego współczynników do wzoru na $P_Tg$
daje $P_Tg=h$. Tym samym
\[
 \operatorname{ran}P_T
 =\left\{h\in H:
 \sum_{n\ge1}e^{2\kappa\lambda_nT}|h_n|^2<\infty\right\}.
\]
Obraz jest gęsty w $H$, ponieważ zawiera każdą skończoną
kombinację liniową funkcji $e_n$.
Nie jest jednak całym $H$. Istotnie, szereg
$h_*=\sum_{n\ge1}e^{-\kappa\lambda_nT}e_n$
zbiega w $H$, ponieważ $\lambda_n$ rośnie jak $n^2$,
ale dla jego współczynników powyższa suma
ważona jest równa $\sum_{n\ge1}1=\infty$.
Zatem $h_*\notin\operatorname{ran}P_T$.
W szczególności obraz $P_T$ nie jest domknięty.

\emph{Stabilność.}
Nawet jeśli ograniczymy dane końcowe do
$\operatorname{ran}P_T$, odwrotność nie jest ciągła
w normie $H$ odziedziczonej z przestrzeni danych.
Niech $h^{(n)}=P_Te_n=e^{-\kappa\lambda_nT}e_n$.
Wówczas
\[
 \|h^{(n)}\|_H=e^{-\kappa\lambda_nT}\longrightarrow0,
 \qquad
 \|P_T^{-1}h^{(n)}\|_H=\|e_n\|_H=1.
\]
Ponieważ $P_T^{-1}0=0$, przeczy to ciągłości
odwrotności w zerze. Przez liniowość odwrotność nie
jest ciągła w żadnym punkcie swojej dziedziny $\operatorname{ran}P_T$.
Dowolnie małe błędy danych końcowych mogą zatem
odpowiadać błędom danych początkowych, które nie
dążą do zera. Jest to spektralna wersja
niestabilności Hadamarda z rozdziału~\ref{chap:pde-heat}.

\section{Podsumowanie i dalsza perspektywa}
\label{sec:pde6-summary}

Dzięki średnim Fejéra uzyskaliśmy zupełność układu
trygonometrycznego i twierdzenie Parsevala.
Na ich podstawie skonstruowaliśmy jawne rozwiązania
ciepła i fal na odcinku oraz ostre oszacowanie zaniku.
Metoda ta przechodzi na zmienne współczynniki przez
regularną teorię Sturma--Liouville'a i na prostokąty
przez iloczyny tensorowe. W kolejnych rozdziałach
rozszerzymy ją na rozwiązania słabe, operatory
samosprzężone, metody wariacyjne i transformatę
Fouriera na $\R^n$.



\chapter[Dystrybucje i słabe pochodne]{Dystrybucje, słabe pochodne i rozwiązania fundamentalne}
\label{chap:pde-distributions}

W poprzednich rozdziałach rozwiązanie klasyczne musiało mieć tyle pochodnych,
ile wymagało równanie. W praktyce taki warunek jest często zbyt silny:
profil fali może mieć skok, dane początkowe mogą być tylko całkowalne,
a rozwiązanie eliptyczne może być najpierw otrzymane przez minimalizację
energii. Potrzebujemy rachunku różniczkowego, w którym można różniczkować
nawet funkcję skokową. Nie chodzi o arbitralne przypisywanie jej wartości
pochodnej w punkcie skoku, lecz o zastąpienie równań punktowych przez
tożsamości całkowe z funkcjami testowymi.

\section{Przestrzeń funkcji testowych}
\label{sec:pde7-tests}

\begin{definicja}[Funkcje testowe i ich zbieżność]
\label{def:pde7-test}
Niech $d\ge1$ oraz $\Omega\subset\R^d$ będzie niepustym zbiorem otwartym.
Przez $\mathcal D(\Omega)=C_c^\infty(\Omega;\C)$ oznaczamy przestrzeń
zespolonych funkcji nieskończenie różniczkowalnych o zwartym nośniku
zawartym w $\Omega$. Nośnikiem funkcji $\varphi$ jest domknięcie
w $\Omega$ zbioru $\{x\in\Omega:\varphi(x)\ne0\}$.
Notacja $K\Subset\Omega$ oznacza, że $K$ jest zwarty i $K\subset\Omega$.
Dla niepustego $K\Subset\Omega$ definiujemy
\[
 \mathcal D_K(\Omega)=
 \{\varphi\in\mathcal D(\Omega):\supp\varphi\subset K\}
\]
oraz półnormy
\[
 p_{K,m}(\varphi)=\max_{|\alpha|\le m}
 \sup_{x\in K}|\partial^\alpha\varphi(x)|,\qquad m\in\N_0.
\]
Tutaj $\alpha=(\alpha_1,\ldots,\alpha_d)\in\N_0^d$,
$|\alpha|=\sum_{i=1}^d\alpha_i$ oraz
$\partial^\alpha=\partial_{x_1}^{\alpha_1}\cdots
\partial_{x_d}^{\alpha_d}$.
Dla $K=\varnothing$ przyjmujemy $\mathcal D_K(\Omega)=\{0\}$
i półnormy równe zeru.

Przestrzeń $\mathcal D_K(\Omega)$ wyposażamy w topologię
wyznaczoną przez półnormy $p_{K,m}$. Na $\mathcal D(\Omega)$
przyjmujemy standardową topologię lokalnie wypukłą granicy
induktywnej: jest to najsilniejsza topologia lokalnie wypukła,
dla której wszystkie włożenia $\mathcal D_K(\Omega)\to
\mathcal D(\Omega)$ są ciągłe. Wyznaczają ją wszystkie półnormy
na $\mathcal D(\Omega)$, których ograniczenia do każdej
$\mathcal D_K(\Omega)$ są ciągłe. Istotnie, topologia
wyznaczona przez tę rodzinę jest lokalnie wypukła,
a każde włożenie jest ciągłe z definicji rodziny.
W każdej innej topologii lokalnie wypukłej o tej
własności półnormy ją wyznaczające mają ciągłe
ograniczenia do wszystkich $\mathcal D_K(\Omega)$,
więc należą do rozważanej rodziny. Otrzymana
topologia jest zatem najsilniejszą taką topologią.

Ciąg $\varphi_j$ zbiega do $\varphi$ w $\mathcal D(\Omega)$
wtedy i tylko wtedy, gdy istnieje jeden zbiór zwarty
$K\Subset\Omega$ zawierający nośnik $\varphi$ oraz nośniki
wszystkich $\varphi_j$ od pewnego indeksu, a ponadto
$p_{K,m}(\varphi_j-\varphi)\to0$ dla każdego $m\in\N_0$.
\end{definicja}

Uzasadnijmy podany opis zbieżności. Wystarczy rozpatrzyć
$\varphi=0$, ponieważ odejmowanie ustalonej funkcji jest
ciągłe, a suma dwóch zbiorów zwartych jest zwarta.
Jeśli wszystkie rozważane nośniki od pewnego indeksu
zawierają się w $K$ i półnormy dążą do zera, to ciąg
zbiega w $\mathcal D_K(\Omega)$, a ciągłość włożenia
daje zbieżność w $\mathcal D(\Omega)$.

Odwrotnie, niech $\varphi_j\to0$ w topologii granicy
induktywnej. Wybieramy rosnący ciąg zbiorów zwartych
$K_\ell\Subset\Omega$, których wnętrza pokrywają $\Omega$
i które zawierają ostatecznie każdy zbiór zwarty w $\Omega$.
Można wziąć
\[
 K_\ell=\{x\in\Omega:|x|\le\ell,\
 \operatorname{dist}(x,\R^d\setminus\Omega)\ge1/\ell\},
 \qquad \ell\ge1,
\]
przyjmując odległość od zbioru pustego równą $+\infty$.
Gdy dopełnienie jest niepuste, jego funkcja odległości
jest ciągła na $\R^d$, więc $K_\ell$ jest domknięty
i ograniczony oraz zawarty w $\Omega$.
Dla $\Omega=\R^d$ otrzymujemy domknięte kule.
Każdy niepusty zwarty podzbiór $\Omega$ jest ograniczony
i ma dodatnią odległość od niepustego dopełnienia
$\Omega$: ciągła funkcja odległości osiąga na nim
dodatnie minimum. Wystarczy więc wybrać $\ell$
z oboma warunkami definiującymi $K_\ell$ spełnionymi
ostro na tym podzbiorze. W szczególności wnętrza
$K_\ell$ pokrywają $\Omega$.
Jeśli nie istnieje jeden wspólny zwarty zbiór dla nośników
od pewnego indeksu, możemy kolejno wybrać rosnące indeksy
$j_\ell$ i punkty $x_\ell\notin K_\ell$ takie, że
$\varphi_{j_\ell}(x_\ell)\ne0$. Istnienie punktu wynika
z tego, że $K_\ell$ jest domknięty i nie zawiera
nośnika $\varphi_{j_\ell}$.
Definiujemy półnormę
\[
 q(\psi)=\sum_{\ell\ge1}
 \frac{|\psi(x_\ell)|}{|\varphi_{j_\ell}(x_\ell)|},
 \qquad \psi\in\mathcal D(\Omega).
\]
Dla każdej ustalonej $\psi$ suma ma tylko skończenie
wiele niezerowych składników: jej nośnik jest zawarty
w pewnym $K_{\ell_0}$, a dla $\ell\ge\ell_0$
punkt $x_\ell$ leży poza tym zbiorem. Na każdej
$\mathcal D_K(\Omega)$ półnorma $q$ jest zatem
ograniczona przez pewną skończoną wielokrotność $p_{K,0}$.
Z definicji topologii granicy induktywnej $q$ jest ciągła.
Jednak $q(\varphi_{j_\ell})\ge1$ dla każdego $\ell$,
co przeczy zbieżności do zera. Istnieje więc wspólny
zwarty zbiór $K$ zawierający nośniki od pewnego indeksu.

Dla każdego $m\in\N_0$ półnorma
\[
 q_m(\psi)=\max_{|\alpha|\le m}
 \sup_{x\in\Omega}|\partial^\alpha\psi(x)|
\]
jest skończona i ciągła w topologii granicy induktywnej:
na $\mathcal D_K(\Omega)$ pokrywa się z $p_{K,m}$,
ponieważ wszystkie pochodne znikają poza nośnikiem
funkcji. Zatem $q_m(\varphi_j)\to0$, a na ustalonym
$K$ otrzymujemy wymagane $p_{K,m}(\varphi_j)\to0$.
Dowodzi to obu implikacji.

Warto odróżnić zwarty nośnik w $\Omega$ od samej
gładkości na domknięciu $\Omega$. Jeśli $\Omega\ne\R^d$,
zwartość nośnika i jego rozłączność z domkniętym
dopełnieniem $\Omega$ zapewniają dodatnią odległość
od tego dopełnienia (dla funkcji zerowej wniosek jest
bezpośredni). Funkcja testowa znika więc w pewnym
otoczeniu brzegu od strony $\Omega$, a przedłużenie
jej przez zero na $\R^d$ jest gładkie. Gdy
$\Omega=\R^d$, przedłużanie nie jest potrzebne.

Istnieje wiele funkcji testowych. Dla przykładu funkcja
\[
 \eta(x)=\begin{cases}
 \exp\!\bigl(-1/(1-|x|^2)\bigr),& |x|<1,\\
 0,& |x|\ge1
 \end{cases}
\]
jest klasy $C_c^\infty(\R^d)$ i jest dodatnia w $B_1(0)$.
Aby to uzasadnić, połóżmy
\[
 \rho(s)=
 \begin{cases}e^{-1/s},&s>0,\\0,&s\le0.\end{cases}
\]
Dla $s>0$ każda pochodna $\rho^{(m)}(s)$ ma postać
$P_m(1/s)e^{-1/s}$, gdzie $P_m$ jest wielomianem:
dla $m=0$ mamy $P_0=1$, a różniczkowanie przy
$z=1/s$ daje $P_{m+1}(z)=z^2(P_m(z)-P_m'(z))$.
Dla dowolnego całkowitego $r\ge0$ zachodzi
$z^r e^{-z}\to0$ przy $z\to+\infty$, na przykład
z oszacowania $e^z\ge z^{r+1}/(r+1)!$.
Każda z powyższych pochodnych, a także jej iloraz
przez $s$, dąży więc do zera przy $s\downarrow0$.
Indukcyjnie ilorazy różnicowe w zerze pokazują,
że $\rho$ jest klasy $C^\infty(\R)$ i wszystkie
jej pochodne w zerze znikają.
Ponieważ $\eta(x)=\rho(1-|x|^2)$, gładkość wynika
z reguły różniczkowania złożenia. Nośnikiem
$\eta$ jest domknięta kula jednostkowa.
Jeśli $\overline{B_r(a)}\subset\Omega$, $r>0$,
to funkcja $x\mapsto\eta((x-a)/r)$ jest funkcją
testową w $\Omega$, dodatnią na $B_r(a)$.

\begin{uwaga}[Dlaczego potrzebny jest wspólny nośnik?]
Dla $d=1$ przesunięcia $\varphi_j(x)=\eta(x-j)$ na $\R$
mają wszystkie pochodne jednostajnie ograniczone
względem $j$, a na każdym ustalonym zbiorze zwartym
funkcje i ich pochodne od pewnego indeksu znikają.
Mimo to nie dążą do zera w $\mathcal D(\R)$, ponieważ
$\supp\varphi_j=[j-1,j+1]$ i żaden zbiór zwarty
nie zawiera tych nośników od pewnego indeksu.
Sama zbieżność jednostajna wszystkich pochodnych
na zbiorach zwartych nie jest zatem zbieżnością
w przestrzeni funkcji testowych.
\end{uwaga}

\section{Dystrybucje jako ciągłe funkcjonały liniowe}
\label{sec:pde7-def}

\begin{definicja}[Dystrybucja]
\label{def:pde7-distribution}
Dla niepustego otwartego $\Omega\subset\R^d$ dystrybucją na $\Omega$
nazywamy ciągły funkcjonał \emph{liniowy} $T:\mathcal D(\Omega)\to\C$.
Zbiór wszystkich takich funkcjonałów oznaczamy przez $\mathcal D'(\Omega)$.
Wartość funkcjonału zapisujemy jako $\langle T,\varphi\rangle$;
nie jest to na ogół iloczyn skalarny, a w szczególności
zapis jest liniowy względem $\varphi$, bez sprzężenia zespolonego.
\end{definicja}

Związek z definicją topologiczną można wyrazić za pomocą lokalnych
oszacowań. Jest to pierwszy ważny przykład tego, że element
przestrzeni sprzężonej do przestrzeni funkcji można opisywać
przez oszacowania na jej półnormach.

\begin{twierdzenie}[Lokalne oszacowanie dystrybucji]
\label{thm:pde7-local-order}
Niech $\Omega\subset\R^d$ będzie otwarta, a $T$ liniowym funkcjonałem
na $\mathcal D(\Omega)$. Wówczas $T\in\mathcal D'(\Omega)$ wtedy i tylko
wtedy, gdy dla każdego $K\Subset\Omega$ istnieją $m_K\in\N_0$
oraz $C_K>0$ takie, że
\begin{equation}\label{eq:pde7-estimate}
 |\langle T,\varphi\rangle|\le C_K p_{K,m_K}(\varphi),
 \qquad \varphi\in\mathcal D_K(\Omega).
\end{equation}
Dla $K\ne\varnothing$ prawa strona oznacza
$C_K\max_{|\alpha|\le m_K}\sup_{x\in K}
|\partial^\alpha\varphi(x)|$, a dla $K=\varnothing$
jest równa zeru zgodnie z definicją półnorm.
\end{twierdzenie}
\begin{proof}
Jeżeli $T$ jest ciągły, to przez ciągłość włożenia
jego ograniczenie do $\mathcal D_K(\Omega)$ jest ciągłe.
Ponieważ półnormy $p_{K,m}$ rosną względem $m$,
każde otoczenie zera zawiera zbiór postaci
$\{\varphi:p_{K,m}(\varphi)<\delta\}$ dla pewnych
$m$ i $\delta>0$ (w skończonym przecięciu warunków
wybieramy największy indeks i najmniejszy promień).
Ciągłość w zerze daje zatem $|T(\varphi)|<1$,
gdy $p_{K,m}(\varphi)<\delta$.
Jeżeli $p_{K,m}(\varphi)>0$, stosujemy tę implikację
do funkcji $\delta\varphi/(2p_{K,m}(\varphi))$.
Z liniowości otrzymujemy
$|T(\varphi)|\le(2/\delta)p_{K,m}(\varphi)$.
Jeżeli $p_{K,m}(\varphi)=0$, to w szczególności
$p_{K,0}(\varphi)=0$; funkcja znika zarówno na $K$,
jak i poza $K$, więc jest zerowa. Dla pustego $K$
oszacowanie również zachodzi, gdyż przestrzeń
$\mathcal D_K(\Omega)$ jest trywialna.

Odwrotnie, jeśli zachodzi~\eqref{eq:pde7-estimate},
to półnorma $\varphi\mapsto|T(\varphi)|$ ma ciągłe
ograniczenie do każdej $\mathcal D_K(\Omega)$.
Z charakterystyki topologii granicy induktywnej
w sekcji~\ref{sec:pde7-tests} jest ona ciągła na
$\mathcal D(\Omega)$. To daje ciągłość $T$.

Równoważnie, funkcjonał liniowy jest dystrybucją,
jeśli zachowuje zbieżność ciągów do zera. Istotnie,
ciągłość implikuje tę własność. Gdyby dla pewnego
$K$ nie zachodziło~\eqref{eq:pde7-estimate}, to dla
każdego $j\ge1$ istniałaby $\psi_j\in\mathcal D_K(\Omega)$
taka, że $|T(\psi_j)|>j p_{K,j}(\psi_j)$.
Wtedy $\varphi_j=\psi_j/T(\psi_j)$ spełnia
$T(\varphi_j)=1$ i $p_{K,j}(\varphi_j)<1/j$.
Dla każdego ustalonego $m$ i $j\ge m$ mamy
$p_{K,m}(\varphi_j)\le p_{K,j}(\varphi_j)$, więc
$\varphi_j\to0$ w $\mathcal D(\Omega)$, sprzecznie
z zachowaniem zbieżności do zera.
\end{proof}

Liczba $m_K$ może zależeć od $K$. Mówimy wtedy o \emph{lokalnie
skończonym rzędzie} dystrybucji. Nie należy bez dodatkowych założeń
wnioskować, że jedna liczba $m$ działa na wszystkich zwartych
podzbiorach obszaru, \emph{nawet gdy jest on ograniczony}.
Istotnie, na $\Omega=(0,1)$ określmy funkcjonał
\[
 \langle T,\varphi\rangle
 =\sum_{k=2}^\infty 2^{-k^2}(-1)^k\varphi^{(k)}(1/k).
\]
W następnej sekcji zapiszemy go jako
$\sum_{k=2}^\infty 2^{-k^2}\partial^k\delta_{1/k}$.
Dla każdego $K\Subset(0,1)$ zbiór
$I_K=\{k\ge2:1/k\in K\}$ jest skończony.
Jeśli $\supp\varphi\subset K$, to wszystkie pochodne
$\varphi$ znikają poza $K$, więc suma zawiera tylko
składniki o indeksach z $I_K$. Jest zatem liniowa
i dla niepustego $I_K$ spełnia oszacowanie
\[
 |T(\varphi)|\le
 \left(\sum_{k\in I_K}2^{-k^2}\right)
 p_{K,\max I_K}(\varphi).
\]
Dla pustego $I_K$ funkcjonał na $\mathcal D_K$ znika.
Z twierdzenia~\ref{thm:pde7-local-order} wynika, że
$T$ jest dystrybucją.

Nie istnieje jednak wspólny skończony rząd $m$.
Ustalmy $k\ge2$, oznaczmy $a=1/k$ i wybierzmy
$r>0$ tak, aby $K=[a-r,a+r]\subset(0,1)$
nie zawierał żadnego innego punktu $1/\ell$.
Wybierzmy $\theta\in C_c^\infty((-1,1))$ z $\theta(0)=1$
(np. $\theta(s)=\eta(2s)/\eta(0)$
z poprzedniej sekcji w wymiarze jeden). Niech
$\psi(s)=s^k\theta(s)/k!$.
Wtedy $\psi^{(k)}(0)=1$ z reguły Leibniza.
Dla $0<\varepsilon<\min\{1,r\}$ połóżmy
\[
 \varphi_\varepsilon(x)=\varepsilon^k
 \psi((x-a)/\varepsilon).
\]
Jej nośnik zawiera się w $K$, a dla każdego $j\ge0$
mamy $\varphi_\varepsilon^{(j)}(x)=
\varepsilon^{k-j}\psi^{(j)}((x-a)/\varepsilon)$.
Stąd dla $m<k$ zachodzi
$p_{K,m}(\varphi_\varepsilon)\to0$ przy
$\varepsilon\downarrow0$, podczas gdy
\[
 \langle T,\varphi_\varepsilon\rangle
 =2^{-k^2}(-1)^k
\]
nie zależy od $\varepsilon$ i nie jest zerem.
Na tym $K$ najmniejszy dopuszczalny rząd wynosi
dokładnie $k$: oszacowanie rzędu $k$ wynika
bezpośrednio z jedynego składnika sumy, a niższe
rzędy wyklucza skonstruowany ciąg (np. $\varepsilon\to0$
wzdłuż dowolnego ciągu dodatniego).
Ponieważ $k$ jest dowolne, ograniczoność zbioru
otwartego nie zapewnia globalnego skończonego rzędu.

\begin{przyklad}[Dystrybucja regularna i miara punktowa]
\label{ex:pde7-regular}
Jeśli $f\in L^1_{\mathrm{loc}}(\Omega)$, to wzór (niezależny od wyboru reprezentanta klasy funkcji $f$)
\begin{equation}\label{eq:pde7-regular}
 \langle T_f,\varphi\rangle
 =\int_\Omega f(x)\varphi(x)\,\mathrm d\lambda_d(x)
\end{equation}
definiuje dystrybucję, ponieważ dla $\supp\varphi\subset K$ zachodzi
$|T_f(\varphi)|\le\|\varphi\|_\infty\int_K|f(x)|\,\mathrm d\lambda_d(x)$.
Oba warunki: liniowość i ciągłość, są więc spełnione.
Dla $a\in\Omega$ wzór $\langle\delta_a,\varphi\rangle=\varphi(a)$
definiuje \emph{deltę Diraca}. Jest to funkcjonał liniowy:
dla $\supp\varphi\subset K$ mamy
$|\varphi(a)|\le p_{K,0}(\varphi)$, gdy $a\in K$,
a gdy $a\notin K$, wartość jest zerowa.
Zatem delta Diraca ma rząd zerowy.
\end{przyklad}

\begin{stwierdzenie}[Iniektywność odwzorowania $f\mapsto T_f$]
\label{prop:pde7-injective}
Jeżeli $f,g\in L^1_{\mathrm{loc}}(\Omega)$ i $T_f=T_g$, to $f=g$
prawie wszędzie w $\Omega$ względem miary Lebesgue'a.
\end{stwierdzenie}
\begin{proof}
Wystarczy rozważyć $h=f-g$. Wybierzmy $a\in\Omega$
i $r>0$ tak, aby $\overline{B_{2r}(a)}\Subset\Omega$.
Ustalmy nieujemną funkcję $\rho\in C_c^\infty(B_1(0))$
o całce $1$ (skalujemy funkcję $\eta$ do kuli
o promieniu $1/2$ i dzielimy przez jej dodatnią,
skończoną całkę). Połóżmy
$\rho_\eps(x)=\eps^{-d}\rho(x/\eps)$ dla $0<\eps<r$.
Dla każdego $x\in B_r(a)$ funkcja
$y\mapsto\rho_\eps(x-y)$ jest testowa i ma nośnik
w $B_{2r}(a)$. Ponieważ $T_h=0$, mamy
$(h*\rho_\eps)(x)=0$ dla wszystkich takich $x$.
Aby przejść do granicy, przedłużmy $h$ po obcięciu do
$B_{2r}(a)$ przez zero i oznaczmy otrzymaną funkcję
przez $H\in L^1(\R^d)$. Z nierówności Minkowskiego
wynika
\[
 \|H*\rho_\eps-H\|_{L^1(\R^d)}
 \le\int_{B_1(0)}\rho(z)
 \|H(\cdot-\eps z)-H\|_{L^1(\R^d)}
 \,\mathrm d\lambda_d(z)\longrightarrow0.
\]
Istotnie, dla $v\in C_c(\R^d)$ i $|b|\le1$
funkcja $v(\cdot-b)-v$ ma nośnik w ustalonym
zbiorze zwartym (suma nośnika $v$ i domkniętej
kuli jednostkowej). Jednostajna ciągłość $v$
daje więc $\|v(\cdot-b)-v\|_1\to0$ dla $b\to0$.
Korzystamy z gęstości $C_c(\R^d)$ w $L^1(\R^d)$.
Uzasadnijmy ją bez odwoływania się do wygładzania.
Funkcje proste o nośniku skończonej miary są gęste
w $L^1$: obcinamy funkcję do ograniczonej kuli,
następnie ograniczamy jej moduł i aproksymujemy
część rzeczywistą oraz urojoną funkcjami prostymi
z dowolnie małym błędem jednostajnym na tej kuli.
Dwa pierwsze błędy w normie $L^1$ dążą do zera
z całkowalności funkcji, a ostatni jest ograniczony
przez błąd jednostajny razy miarę kuli.
Dla mierzalnego zbioru $A$ skończonej miary
i $\tau>0$ regularność miary Lebesgue'a daje
zwarty $K_0\subset A$ i otwarty $O\supset A$
z $\lambda_d(O\setminus K_0)<\tau$.
Jeśli $K_0=\varnothing$, funkcja zerowa
aproksymuje $\chi_A$ z błędem mniejszym niż $\tau$.
W przeciwnym razie wybieramy $r_0>0$ mniejszy
od odległości $K_0$ od $\R^d\setminus O$
(dowolny skończony $r_0>0$, gdy $O=\R^d$).
Funkcja
\[
 w(x)=\max\{0,1-\operatorname{dist}(x,K_0)/r_0\}
\]
jest ciągła, równa $1$ na $K_0$, ma zwarty
nośnik w $O$ i przyjmuje wartości w $[0,1]$.
Zatem $|\chi_A-w|\le\chi_{O\setminus K_0}$,
skąd $\|\chi_A-w\|_1<\tau$.
Stosując tę konstrukcję do składników funkcji
prostej i nierówność trójkąta, uzyskujemy gęstość. Dla dowolnego $\gamma>0$
wybieramy $v$ z $\|H-v\|_1<\gamma$. Niezmienniczość
miary Lebesgue'a względem przesunięć daje
\[
 \|H(\cdot-b)-H\|_1
 \le2\gamma+\|v(\cdot-b)-v\|_1.
\]
Z dowolności $\gamma$ wynika ciągłość przesunięć.
W oszacowaniu splotu stosujemy twierdzenie
o zbieżności zmajoryzowanej z całkowalną
majorantą $2\|H\|_1\rho(z)$.
Samo oszacowanie wynika z podstawienia
$y=x-\eps z$, normalizacji $\rho$ i twierdzenia
Tonellego zastosowanego do modułu różnicy.

Na $B_r(a)$ sploty $H*\rho_\eps$ są zerowe,
ponieważ wszystkie używane punkty $y$ leżą
w $B_{2r}(a)$, gdzie $H=h$ prawie wszędzie.
Dlatego $\|h\|_{L^1(B_r(a))}
\le\|H*\rho_\eps-H\|_{L^1(\R^d)}\to0$.
Zatem $h=0$ prawie wszędzie na tej kuli.
Z drugiej przeliczalności $\R^d$ wybieramy
przeliczalne pokrycie $\Omega$ takimi kulami.
Przeliczalna suma zbiorów miary zero ma miarę zero,
co daje tezę na całym $\Omega$.
Ogólną teorię przybliżonych jedynek rozwiniemy później
w sekcji~\ref{sec:pde7-mollifiers}; nie jest ona
przesłanką niniejszego dowodu.
\end{proof}

\begin{uwaga}[Dlaczego delta nie jest zwykłą funkcją?]
Załóżmy, że $\delta_0=T_f$ dla pewnego
$f\in L^1_{\mathrm{loc}}(\R^d)$. Wybierzmy funkcję
$\psi\in\mathcal D(\R^d)$, $0\le\psi\le1$,
$\psi(0)=1$, $\supp\psi\subset B_1(0)$
(na przykład $\psi(x)=\eta(2x)/\eta(0)$), i połóżmy
$\psi_\varepsilon(x)=\psi(x/\varepsilon)$.
Wtedy $\langle\delta_0,\psi_\varepsilon\rangle=1$,
natomiast dla $0<\varepsilon<1$
\[
 |\langle T_f,\psi_\varepsilon\rangle|
 \le\int_{B_\varepsilon(0)}|f(x)|\,\mathrm d\lambda_d(x)
 \longrightarrow0.
\]
Istotnie, $f$ jest całkowalna na domkniętej kuli
jednostkowej, a miara $B_\varepsilon(0)$ dąży do zera,
więc stosujemy absolutną ciągłość całki Lebesgue'a.
Otrzymujemy sprzeczność.
\end{uwaga}

\section{Pochodne dystrybucyjne i skoki}
\label{sec:pde7-derivatives}

\begin{definicja}[Pochodna dystrybucyjna]
\label{def:pde7-derivative}
Dla $T\in\mathcal D'(\Omega)$ i wielowskaźnika
$\alpha\in\N_0^d$ określamy $\partial^\alpha T\in\mathcal D'(\Omega)$
wzorem
\begin{equation}\label{eq:pde7-distributional-derivative}
 \langle\partial^\alpha T,\varphi\rangle
 =(-1)^{|\alpha|}\langle T,\partial^\alpha\varphi\rangle,
 \qquad\varphi\in\mathcal D(\Omega).
\end{equation}
\end{definicja}

Definicja ma sens: nośnik $\partial^\alpha\varphi$
zawiera się w nośniku $\varphi$, a dla $\varphi\in\mathcal D_K(\Omega)$
mamy
\[
 p_{K,m}(\partial^\alpha\varphi)
 \le p_{K,m+|\alpha|}(\varphi).
\]
Zatem jeśli $T$ spełnia lokalne oszacowanie
~\eqref{eq:pde7-estimate} z $m_K,C_K$, to
\[
 |\langle\partial^\alpha T,\varphi\rangle|
 \le C_K p_{K,m_K+|\alpha|}(\varphi).
\]
Funkcjonał $\partial^\alpha T$ jest liniowy i ciągły
na mocy twierdzenia~\ref{thm:pde7-local-order}.
W szczególności każda dystrybucja ma pochodne
dowolnego rzędu. Dla wielowskaźników $\alpha,\beta$
dwukrotne zastosowanie definicji daje
\[
 \langle\partial^\beta\partial^\alpha T,\varphi\rangle
 =(-1)^{|\alpha|+|\beta|}
 \langle T,\partial^\alpha\partial^\beta\varphi\rangle
 =\langle\partial^{\alpha+\beta}T,\varphi\rangle.
\]
W ostatnim kroku wykorzystaliśmy komutowanie
pochodnych gładkiej funkcji testowej. Pochodne
mieszane dystrybucji zatem również komutują.

\begin{stwierdzenie}[Zgodność z klasycznym różniczkowaniem]
\label{prop:pde7-classical}
Jeśli $u\in C^1(\Omega)$, to $\partial_jT_u=T_{\partial_j u}$.
Jeśli $T\in\mathcal D'(\Omega)$ i $a\in C^\infty(\Omega)$,
to iloczyn $aT$, dany wzorem
$\langle aT,\varphi\rangle=\langle T,a\varphi\rangle$,
jest dystrybucją i zachodzi reguła Leibniza
$\partial_j(aT)=(\partial_ja)T+a\partial_jT$.
\end{stwierdzenie}
\begin{proof}
Dla $u\in C^1(\Omega)$ iloczyn $u\varphi$ ma zwarty
nośnik w $\Omega$ i przedłuża się przez zero do
funkcji klasy $C^1$ na $\R^d$. Z twierdzenia Fubiniego
i podstawowego twierdzenia rachunku całkowego
całka z jej pochodnej względem $x_j$ jest zerowa
(całkujemy najpierw po całej prostej w tej zmiennej).
Reguła Leibniza dla funkcji daje więc
\[
 \langle\partial_jT_u,\varphi\rangle
 =-\int_\Omega u(x)\partial_j\varphi(x)\,\mathrm d\lambda_d(x)
 =\int_\Omega \partial_ju(x)\varphi(x)\,\mathrm d\lambda_d(x).
\]
To dowodzi pierwszej tezy.

Dla $a\in C^\infty(\Omega)$ iloczyn $a\varphi$ ma
nośnik zawarty w nośniku $\varphi$. Z wielowskaźnikowej
reguły Leibniza wynika dla niepustego $K$
\[
 p_{K,m}(a\varphi)\le 2^m p_{K,m}^{\,a}p_{K,m}(\varphi),
 \qquad
 p_{K,m}^{\,a}=\max_{|\gamma|\le m}
 \sup_{x\in K}|\partial^\gamma a(x)|.
\]
Istotnie, dla każdej pochodnej $\partial^\alpha$
sumę współczynników dwumianowych w regule Leibniza
stanowi $2^{|\alpha|}\le2^m$. Wszystkie pochodne
$a$ są ograniczone na $K$. Stosując lokalne oszacowanie
dla $T$ do $a\varphi$, otrzymujemy ciągłość liniowego
funkcjonału $aT$ (dla pustego $K$ jest to trywialne).
Ponadto
\[
 \begin{aligned}
 \langle(\partial_ja)T+a\partial_jT,\varphi\rangle
 &=\langle T,(\partial_ja)\varphi\rangle
   -\langle T,\partial_j(a\varphi)\rangle\\
 &=-\langle T,a\partial_j\varphi\rangle
 =\langle\partial_j(aT),\varphi\rangle.
 \end{aligned}
\]
Jest to żądana reguła Leibniza, poprawna także
dla zespolonych $a$ i $\varphi$.
\end{proof}

\begin{przyklad}[Funkcja Heaviside'a]
\label{ex:pde7-heaviside}
Niech $H(x)=0$ dla $x<0$ i $H(x)=1$ dla $x>0$;
wartość w zerze nie ma znaczenia. Dla $\varphi\in\mathcal D(\R)$
\[
 \langle H',\varphi\rangle
 =-\int_0^\infty\varphi'(x)\,\mathrm d\lambda_1(x)
 =\varphi(0)=\langle\delta_0,\varphi\rangle.
\]
Zatem $H'=\delta_0$ w $\mathcal D'(\R)$.
Utożsamiamy tu funkcję lokalnie całkowalną z jej
dystrybucją regularną. Dla $u(x)=|x|$ całkujemy osobno
na obu półprostych:
\[
 \begin{aligned}
 \langle T_u',\varphi\rangle
 &=-\int_{-\infty}^0(-x)\varphi'(x)\,\mathrm d\lambda_1(x)
   -\int_0^\infty x\varphi'(x)\,\mathrm d\lambda_1(x)\\
 &=-\int_{-\infty}^0\varphi(x)\,\mathrm d\lambda_1(x)
   +\int_0^\infty\varphi(x)\,\mathrm d\lambda_1(x).
 \end{aligned}
\]
Wyrazy brzegowe w zerze znikają dzięki czynnikowi $x$,
a w nieskończoności dzięki zwartemu nośnikowi.
Zatem $T_u'=T_{\operatorname{sgn}}$. Ponieważ
$\operatorname{sgn}=2H-1$ prawie wszędzie i pochodna
dystrybucji stałej jest zerowa, mamy $T_u''=2\delta_0$.
\end{przyklad}

\begin{przyklad}[Pochodna delty]
Bezpośrednio z definicji pochodnej otrzymujemy
\[
 \langle\partial^\alpha\delta_a,\varphi\rangle
 =(-1)^{|\alpha|}\partial^\alpha\varphi(a).
\]
Jest to dystrybucja rzędu co najwyżej $|\alpha|$,
a na każdym zwartym $K$ zawierającym $a$ w swoim
wnętrzu jej rząd wynosi dokładnie $|\alpha|$.
Aby uzasadnić dolne ograniczenie, dla $k=|\alpha|>0$
wybieramy $\theta\in C_c^\infty(B_1(0))$ z $\theta(0)=1$
i kładziemy $\psi(x)=x^\alpha\theta(x)/\alpha!$,
gdzie $x^\alpha=\prod_{i=1}^d x_i^{\alpha_i}$ oraz
$\alpha!=\prod_{i=1}^d\alpha_i!$.
Reguła Leibniza daje $\partial^\alpha\psi(0)=1$.
Dla dostatecznie małego $\varepsilon>0$ funkcje
$\varphi_\varepsilon(x)=\varepsilon^k
\psi((x-a)/\varepsilon)$ należą do $\mathcal D_K(\Omega)$
i spełniają $p_{K,m}(\varphi_\varepsilon)\to0$
dla $m<k$, natomiast
$\langle\partial^\alpha\delta_a,\varphi_\varepsilon\rangle
=(-1)^k$. Wyklucza to oszacowania niższego rzędu.
Dla $k=0$ otrzymujemy deltę rzędu zerowego.

W szczególności $\delta'_0$ nie jest dystrybucją
wyznaczoną przez lokalnie skończoną miarę Radona
(również zespoloną): taka miara $\mu$ spełnia
$|\int\varphi\,\mathrm d\mu|\le|\mu|(K)p_{K,0}(\varphi)$,
gdzie $|\mu|$ oznacza jej wahanie całkowite na $K$.
Dla $\delta'_0$ właśnie wykazaliśmy niemożliwość
oszacowania rzędu zerowego.
\end{przyklad}

\begin{twierdzenie}[Wzór skoku przez hiperpowierzchnię]
\label{thm:pde7-jump}
Niech $\Omega\subset\R^d$ będzie otwarta, a
$S\subset\Omega$ będzie względnie domkniętą,
zorientowaną, zanurzoną hiperpowierzchnią klasy $C^1$
bez brzegu w $\Omega$. Zakładamy, że w każdym
punkcie $S$ istnieje otoczenie, w którym $S$ jest
wykresem funkcji klasy $C^1$ (po obrocie współrzędnych)
i rozdziela to otoczenie na dwie strony $\Omega_-$
i $\Omega_+$. Nazwy stron są zgodne z orientacją:
jednostkowa normalna $\nu$ jest skierowana
od strony minus do strony plus.
Niech $u\in C^1(\Omega\setminus S)$ i niech lokalnie
przy każdym punkcie $S$ ograniczenia $u$ do obu stron
mają przedłużenia klasy $C^1$ na ich domknięcia.
Oznaczmy ślady na $S$ przez $u_-$ i $u_+$.
Wtedy, w sensie dystrybucji,
\begin{equation}\label{eq:pde7-jump}
 \partial_jT_u=T_{(\partial_ju)_{\mathrm{reg}}}
 +(u_+-u_-)\nu_j\delta_S,
\end{equation}
gdzie $(\partial_ju)_{\mathrm{reg}}$ oznacza klasyczną
pochodną poza $S$, a dla ciągłej funkcji $b$ na $S$
\[
 \langle b\delta_S,\varphi\rangle
 =\int_S b(s)\varphi(s)\,\mathrm d\sigma(s),
 \qquad \sigma=\mathcal H^{d-1}\!\restriction_S.
\]
Wartości funkcji objętościowych na $S$ nie wpływają
na ich dystrybucje regularne.
\end{twierdzenie}
\begin{proof}
W lokalnych współrzędnych wykresowych hiperpowierzchnia
ma miarę Lebesgue'a zero z twierdzenia Fubiniego:
każdy przekrój w kierunku osi wykresu jest co najwyżej
jednopunktowy. Przeliczalne pokrycie takimi otoczeniami
daje $\lambda_d(S)=0$. Na każdym zwartym $K\Subset\Omega$
zbiór $K\cap S$ jest zwarty dzięki względnemu domknięciu
$S$. Można go pokryć skończenie wieloma mniejszymi
otoczeniami wykresowymi o ograniczonej podstawie
i ograniczonej pochodnej funkcji wykresu.
Wzór na miarę powierzchni wykresu
$\sqrt{1+|\nabla h|^2}\,\mathrm d\lambda_{d-1}$
daje $\sigma(K\cap S)<\infty$.
W wymiarze $d=1$ stosujemy zamiast tego miarę
liczącą: $K\cap S$ jest skończonym zbiorem punktów.
Z lokalnych przedłużeń klasy $C^1$ i skończonego
pokrycia $K$ wynika lokalna całkowalność $u$
i $(\partial_ju)_{\mathrm{reg}}$. Ponadto
\[
 |\langle b\delta_S,\varphi\rangle|
 \le \sup_{K\cap S}|b|\,\sigma(K\cap S)p_{K,0}(\varphi)
\]
dla $\supp\varphi\subset K$ (gdy przecięcie jest
puste, całka znika). Zatem wyraz powierzchniowy
jest dystrybucją rzędu zerowego.

Ustalmy $\varphi\in\mathcal D(\Omega)$.
Najpierw załóżmy, że jej nośnik jest zawarty
w jednym otoczeniu wykresowym. Możemy wybrać
lokalny cylinder wykresowy o podstawie prostokątnej,
zawierający nośnik w swoim wnętrzu, którego górna
i dolna ściana nie przecinają wykresu.
Wzór Gaussa na obu częściach cylindra stosujemy
do pola $u\varphi\,e_j$, gdzie $e_j$ jest $j$-tym
wektorem standardowej bazy. Obie części mają
brzeg odcinkami klasy $C^1$, a pole ma ciągłe
pochodne aż do strony wykresu.
Na sztucznych ścianach cylindra pole znika,
ponieważ $\varphi$ ma nośnik we wnętrzu.
Normalna zewnętrzna części minus na $S$ to $\nu$,
a części plus to $-\nu$. W konsekwencji
\[
 \begin{aligned}
 -\int_{\Omega_-}u\partial_j\varphi\,\mathrm d\lambda_d
 &=\int_{\Omega_-}(\partial_ju)\varphi\,\mathrm d\lambda_d
   -\int_S u_-\varphi\nu_j\,\mathrm d\sigma,\\
 -\int_{\Omega_+}u\partial_j\varphi\,\mathrm d\lambda_d
 &=\int_{\Omega_+}(\partial_ju)\varphi\,\mathrm d\lambda_d
   +\int_S u_+\varphi\nu_j\,\mathrm d\sigma.
 \end{aligned}
\]
Tutaj całki dotyczą części lokalnego cylindra.
Sumując, otrzymujemy~\eqref{eq:pde7-jump} dla takiego
testu. Dla testu o nośniku rozłącznym z $S$ wzór
wynika z klasycznego całkowania przez części.

W ogólnym przypadku pokrywamy zwarty nośnik
$\varphi$ skończoną rodziną otoczeń wykresowych
i otoczeń rozłącznych z $S$. Korzystamy z gładkiego
podziału jedności podporządkowanego temu pokryciu
w otoczeniu nośnika. Dokładniej, wybieramy skończenie
wiele nieujemnych funkcji testowych $\eta_i$
o nośnikach w tych otoczeniach, których suma $A$
jest dodatnia w otoczeniu nośnika $\varphi$
(istnienie wynika z funkcji na kulach z sekcji
~\ref{sec:pde7-tests} i zwartości nośnika).
Funkcje $\varphi_i=\varphi\eta_i/A$ tam, gdzie $A>0$,
przedłużamy przez zero poza ten zbiór.
Są gładkie, gdyż nośnik $\varphi$ jest zawarty
w otwartym zbiorze $\{A>0\}$, oraz mają zwarte
nośniki w odpowiednich otoczeniach.
Mamy $\varphi=\sum_i\varphi_i$. Stosując wzór
lokalny do każdej $\varphi_i$ i sumując skończenie
wiele równości, otrzymujemy tezę dla $\varphi$.
\end{proof}

\begin{wniosek}[Skok strumienia i równanie ciągłości]
Niech $F=(F_1,\ldots,F_d)$ będzie polem rzeczywistym,
którego każda składowa spełnia założenia
twierdzenia~\ref{thm:pde7-jump}. Wtedy
\[
 \operatorname{div}T_F
 =T_{(\operatorname{div}F)_{\mathrm{reg}}}
 +[F]\cdot\nu\,\delta_S,\qquad [F]=F_+-F_-.
\]
Równanie $\operatorname{div}T_F=0$ jest równoważne
warunkom $\operatorname{div}F=0$ poza $S$ oraz
$[F]\cdot\nu=0$ na $S$.
\end{wniosek}
\begin{proof}
Sumujemy wzór~\eqref{eq:pde7-jump} dla składowych
$F_j$ i pochodnych $\partial_j$.
Jeśli $\operatorname{div}T_F=0$, testy o nośnikach
w $\Omega\setminus S$ dają zerowanie klasycznej
dywergencji prawie wszędzie na tym zbiorze na mocy
stwierdzenia~\ref{prop:pde7-injective}; z ciągłości
jest ona zerowa punktowo. Ponieważ $\lambda_d(S)=0$,
pozostaje $b\delta_S=0$, gdzie $b=[F]\cdot\nu$ jest
ciągłą funkcją rzeczywistą na $S$.
Gdyby $b(s_0)\ne0$, w dostatecznie małym otoczeniu
$s_0$ funkcja $b$ miałaby stały znak i moduł
ograniczony od dołu przez dodatnią stałą.
Wybierając nieujemną funkcję testową o nośniku
w tym otoczeniu, dodatnią w otoczeniu $s_0$,
otrzymalibyśmy niezerową całkę powierzchniową:
każdy niepusty otwarty fragment wykresu ma dodatnią
miarę powierzchniową (dla $d=1$ miarę liczącą).
Sprzeczność dowodzi $b=0$. Odwrotnie, gdy oba
warunki zachodzą, oba składniki wzoru znikają.
Jest to mechanizm stojący za warunkiem
Rankine'a--Hugoniota z rozdziału~\ref{chap:pde-characteristics}.
\end{proof}

\section{Nośnik, lokalność i zbieżność dystrybucji}
\label{sec:pde7-support}

Dla otwartego $U\subset\Omega$ każdą funkcję
$\varphi\in\mathcal D(U)$ przedłużamy przez zero
do $\Omega$. Otrzymujemy funkcję testową w $\Omega$,
ponieważ nośnik jest zwartym podzbiorem $U$.
Ograniczenie dystrybucji $T$ do $U$ definiujemy przez
$\langle T|_U,\varphi\rangle=\langle T,\widetilde\varphi\rangle$,
gdzie $\widetilde\varphi$ jest tym przedłużeniem.
Dla każdego $K\Subset U$ półnormy funkcji
$\varphi\in\mathcal D_K(U)$ i jej przedłużenia na $K$
są równe. Zatem lokalne oszacowanie
~\eqref{eq:pde7-estimate} dla $T$ daje to samo
oszacowanie dla $T|_U$. W szczególności
$T|_U\in\mathcal D'(U)$.

\begin{definicja}[Nośnik dystrybucji]
Dla $T\in\mathcal D'(\Omega)$ nośnik $\supp T$ jest
dopełnieniem w $\Omega$ sumy wszystkich otwartych
zbiorów $U\subset\Omega$ takich, że $T|_U=0$.
Jest więc względnie domknięty w $\Omega$.
\end{definicja}

Uzasadnijmy, że dystrybucja znika również na tej sumie.
Jeżeli $\varphi$ jest funkcją testową na sumie
otwartych zbiorów, na których $T$ znika, to jej
zwarty nośnik ma skończone pokrycie takimi zbiorami.
Rozkładamy $\varphi$ na skończoną sumę funkcji
testowych $\varphi_i$ o nośnikach w odpowiednich
zbiorach pokrycia, stosując konstrukcję z dowodu
twierdzenia~\ref{thm:pde7-jump} (funkcje odcinające
na kulach, a następnie dzielenie przez ich sumę
dodatnią w otoczeniu nośnika).
Każdy składnik spełnia $\langle T,\varphi_i\rangle=0$,
zatem $\langle T,\varphi\rangle=0$.
W konsekwencji $\Omega\setminus\supp T$ jest
największym otwartym zbiorem, na którym $T$ znika.
W szczególności dla testu o nośniku rozłącznym
z $\supp T$ mamy $\langle T,\varphi\rangle=0$.
Ten sam argument pokazuje lokalność równości:
jeżeli $T_1|_{U_i}=T_2|_{U_i}$ na każdym zbiorze
otwartego pokrycia $\Omega$, to $T_1=T_2$.

Dla $f\in L^1_{\mathrm{loc}}(\Omega)$ zachodzi
\[
 \supp T_f=\Omega\setminus
 \bigcup\{U\subset\Omega: U\text{ otwarty},\
 f=0\text{ prawie wszędzie na }U\}.
\]
Istotnie, jeśli $f=0$ prawie wszędzie na $U$,
to $T_f|_U=0$ bezpośrednio z definicji.
Odwrotna implikacja wynika z iniektywności
odwzorowania $f\mapsto T_f$ na $U$
(stwierdzenie~\ref{prop:pde7-injective}).
Prawa strona wzoru to nośnik istotny funkcji $f$;
nie zależy on od reprezentanta klasy prawie wszędzie.
Ponadto na otwartej sumie występującej we wzorze
$f$ znika prawie wszędzie: z drugiej przeliczalności
można wybrać przeliczalne podpokrycie tej sumy.

Jeżeli $T$ znika na $U$, to dla każdego testu na $U$
jego pochodne również są testami na $U$.
Wzór~\eqref{eq:pde7-distributional-derivative} daje
więc $(\partial^\alpha T)|_U=0$. W szczególności
$\supp(\partial^\alpha T)\subset\supp T$.
Podobnie, jeśli $a\in C^\infty(\Omega)$, to
$\supp(aT)\subset\supp a\cap\supp T$:
na dopełnieniu nośnika $T$ funkcjonał znika,
a na dopełnieniu nośnika $a$ sam iloczyn
$a\varphi$ jest zerowy. Stosujemy tu wykazaną
własność znikania na sumie zbiorów otwartych.

Na przykład $\supp\delta_a=\{a\}$, ponieważ
delta znika poza $a$, natomiast w każdym
otoczeniu $a$ istnieje funkcja testowa
o wartości $1$ w $a$. Również
$\supp(\partial^\alpha\delta_a)=\{a\}$:
inkluzja wynika z powyższego, a test postaci
wielomian razy funkcja odcinająca, skonstruowany
w sekcji~\ref{sec:pde7-derivatives}, daje
niezerową wartość funkcjonału w dowolnym
otoczeniu $a$.

\begin{definicja}[Zbieżność dystrybucji]
Mówimy, że ciąg $T_n\in\mathcal D'(\Omega)$ zbiega do
$T\in\mathcal D'(\Omega)$ w sensie dystrybucji, jeśli
\[
 \langle T_n,\varphi\rangle\longrightarrow
 \langle T,\varphi\rangle
 \qquad\text{dla każdej }\varphi\in\mathcal D(\Omega).
\]
Jest to zbieżność słaba* względem pary
$\mathcal D'(\Omega),\mathcal D(\Omega)$.
\end{definicja}

Warto sprawdzić, że granica punktowa funkcjonałów
liniowych w tej definicji rzeczywiście jest
dystrybucją, jeśli wszystkie granice skalarne istnieją.

\begin{stwierdzenie}[Granica punktowa dystrybucji]
Niech $T_n\in\mathcal D'(\Omega)$ i załóżmy, że dla
każdej $\varphi\in\mathcal D(\Omega)$ istnieje skończona
granica $\lim_{n\to\infty}\langle T_n,\varphi\rangle$.
Wówczas wzór
$\langle T,\varphi\rangle=\lim_{n\to\infty}
\langle T_n,\varphi\rangle$ definiuje dystrybucję.
Co więcej, dla każdego $K\Subset\Omega$ istnieją
$m\in\N_0$ i $C>0$ niezależne od $n$ takie, że
\[
 |\langle T_n,\varphi\rangle|\le C p_{K,m}(\varphi)
 \qquad(\varphi\in\mathcal D_K(\Omega),\ n\ge1).
\]
\end{stwierdzenie}
\begin{proof}
Najpierw uzasadnijmy zupełność $\mathcal D_K(\Omega)$
w topologii jej półnorm. Dla niepustego $K$ topologię
tę wyznacza metryka
\[
 d_K(\varphi,\psi)=\sum_{m=0}^\infty 2^{-m-1}
 \min\{1,p_{K,m}(\varphi-\psi)\}.
\]
Nierówność trójkąta wynika z podaddytywności
$\min\{1,t\}$ dla $t\ge0$ i półnorm, a rozdzielanie
punktów z półnormy $p_{K,0}$.
Zbieżność w tej metryce jest równoważna
zbieżności we wszystkich półnormach: w jedną
stronę szacujemy każdy składnik szeregu,
a w drugą najpierw odcinamy jego mały ogon,
następnie szacujemy skończenie wiele składników.
Tak samo warunek Cauchy'ego w metryce jest
równoważny warunkowi Cauchy'ego w każdej półnormie.

Niech $(\psi_\ell)$ będzie takim ciągiem.
Przedłużamy funkcje przez zero do $\R^d$.
Dla każdego wielowskaźnika $\alpha$ ciąg
$\partial^\alpha\psi_\ell$ jest jednostajnie
Cauchy'ego na $\R^d$, ponieważ wszystkie
pochodne znikają poza $K$. Ma więc jednostajną
granicę $g_\alpha$, ciągłą i znikającą poza $K$.
Dla $x\in\R^d$, $t\in\R$ i wektora standardowej
bazy $e_j$ mamy
\[
 \partial^\alpha\psi_\ell(x+te_j)
 -\partial^\alpha\psi_\ell(x)
 =\int_0^t
 \partial^{\alpha+e_j}\psi_\ell(x+se_j)\,
 \mathrm d\lambda_1(s)
\]
(z całką zorientowaną, gdy $t<0$).
Jednostajna zbieżność pozwala przejść do granicy
po obu stronach. Otrzymujemy tę samą tożsamość
dla $g_\alpha$ i $g_{\alpha+e_j}$.
Z ciągłości funkcji podcałkowej wynika
$\partial_jg_\alpha=g_{\alpha+e_j}$.
Indukcyjnie $g_0$ jest gładka, a jej pochodnymi
są $g_\alpha$. Jej nośnik zawiera się w $K$,
więc $g_0|_\Omega\in\mathcal D_K(\Omega)$,
a ciąg $\psi_\ell$ zbiega do tej funkcji
we wszystkich półnormach. Metryka jest zupełna.
Dla pustego $K$ zupełność jest trywialna.

Ustalmy teraz $K\Subset\Omega$.
Z istnienia granic skalarnych wynika punktowa
ograniczoność rodziny funkcjonałów:
dla każdej $\varphi\in\mathcal D_K(\Omega)$
liczby $\langle T_n,\varphi\rangle$ są ograniczone.
Definiujemy domknięte zbiory
\[
 A_j=\{\varphi\in\mathcal D_K(\Omega):
 |\langle T_n,\varphi\rangle|\le j
 \text{ dla wszystkich }n\ge1\},\qquad j\ge1.
\]
Są domknięte dzięki ciągłości każdego $T_n$ na
$\mathcal D_K(\Omega)$, a ich suma pokrywa
całą przestrzeń $\mathcal D_K(\Omega)$.
Z twierdzenia Baire'a dla przestrzeni metrycznej
zupełnej (jawnie wykorzystywanego tutaj wyniku
topologii) pewien $A_j$ ma niepuste wnętrze.
Istnieją więc $\varphi_0\in A_j$, $m\in\N_0$
i $\delta>0$ takie, że
$\varphi_0+\psi\in A_j$ dla $p_{K,m}(\psi)<\delta$.
Z liniowości mamy wówczas
\[
 |\langle T_n,\psi\rangle|
 \le |\langle T_n,\varphi_0+\psi\rangle|
     +|\langle T_n,\varphi_0\rangle|\le2j.
\]
Dla $p_{K,m}(\varphi)>0$ podstawiamy
$\psi=\delta\varphi/(2p_{K,m}(\varphi))$.
Otrzymujemy
$|\langle T_n,\varphi\rangle|
\le(4j/\delta)p_{K,m}(\varphi)$ dla każdego $n$.
Jeśli półnorma jest zerowa, to $\varphi=0$,
więc oszacowanie również zachodzi.
Dla pustego $K$ wybieramy dowolne $C>0$ i $m=0$.

Granica punktowa jest liniowa przez liniowość
każdego $T_n$, a przejście do granicy w powyższym
oszacowaniu daje
$|\langle T,\varphi\rangle|\le C p_{K,m}(\varphi)$.
Ponieważ $K$ było dowolne, twierdzenie
~\ref{thm:pde7-local-order} daje $T\in\mathcal D'(\Omega)$.
\end{proof}

\begin{stwierdzenie}[Różniczkowanie a granica dystrybucyjna]
\label{prop:pde7-derivative-limit}
Niech $\Omega\subset\R^d$ będzie otwarta i
$T_n\to T$ w $\mathcal D'(\Omega)$. Dla każdego
wielowskaźnika $\alpha$ zachodzi
$\partial^\alpha T_n\to\partial^\alpha T$ w $\mathcal D'(\Omega)$.
Dla każdej ustalonej $a\in C^\infty(\Omega)$ zachodzi
również $aT_n\to aT$ w tym sensie.
\end{stwierdzenie}
\begin{proof}
Dla każdego $\varphi\in\mathcal D(\Omega)$ mamy
\[
 \begin{aligned}
 \langle\partial^\alpha T_n,\varphi\rangle
 &=(-1)^{|\alpha|}\langle T_n,\partial^\alpha\varphi\rangle\\
 &\longrightarrow
 (-1)^{|\alpha|}\langle T,\partial^\alpha\varphi\rangle
 =\langle\partial^\alpha T,\varphi\rangle.
 \end{aligned}
\]
Ponieważ $\partial^\alpha\varphi$ jest funkcją testową,
przejście do granicy jest uprawnione.
Analogicznie $a\varphi$ jest funkcją testową, więc
$\langle aT_n,\varphi\rangle=\langle T_n,a\varphi\rangle
\to\langle T,a\varphi\rangle=\langle aT,\varphi\rangle$.
\end{proof}

Jeśli $f_n\to f$ w $L^1_{\mathrm{loc}}(\Omega)$,
to dla każdego testu o nośniku zawartym w $K$
mamy
\[
 |\langle T_{f_n}-T_f,\varphi\rangle|
 \le p_{K,0}(\varphi)
 \int_K|f_n-f|\,\mathrm d\lambda_d\longrightarrow0.
\]
Otrzymujemy zbieżność dystrybucyjną, ale odwrotna
implikacja nie zachodzi, jak pokazuje poniższy przykład.

\begin{przyklad}[Gładkie funkcje zbieżne do delty]
Jeśli $\rho\in\mathcal D(\R^d)$ spełnia $\rho\ge0$ i
$\int_{\R^d}\rho(x)\,\mathrm d\lambda_d(x)=1$, a
$\rho_\varepsilon(x)=\varepsilon^{-d}\rho(x/\varepsilon)$
dla $\varepsilon>0$, to
$T_{\rho_\varepsilon}\to\delta_0$ w $\mathcal D'(\R^d)$
przy $\varepsilon\downarrow0$.
Rzeczywiście, po podstawieniu $x=\varepsilon y$
otrzymujemy
\[
 \begin{aligned}
 |\langle T_{\rho_\varepsilon}-\delta_0,\varphi\rangle|
 &\le \int_{\R^d}\rho(y)
 |\varphi(\varepsilon y)-\varphi(0)|\,\mathrm d\lambda_d(y)\\
 &\le\sup_{y\in\supp\rho}
 |\varphi(\varepsilon y)-\varphi(0)|
 \longrightarrow0.
 \end{aligned}
\]
Ostatnia granica wynika z ciągłości $\varphi$ w zerze
i ograniczoności nośnika $\rho$.
Nie może to być zbieżność do żadnej funkcji
$f\in L^1_{\mathrm{loc}}(\R^d)$ w normach lokalnych $L^1$:
taka zbieżność implikowałaby $T_f=\delta_0$,
co wykluczyliśmy w sekcji~\ref{sec:pde7-def}.
\end{przyklad}

\section{Wygładzanie i przybliżone jedynki}
\label{sec:pde7-mollifiers}

Wybierzmy odtąd $\rho\in C_c^\infty(B_1(0))$,
$\rho\ge0$, o całce równej $1$, i definiujmy
$\rho_\varepsilon(x)=\varepsilon^{-d}\rho(x/\varepsilon)$
dla $\varepsilon>0$. Podstawienie $x=\varepsilon z$
daje $\int_{\R^d}\rho_\varepsilon\,\mathrm d\lambda_d=1$,
a $\supp\rho_\varepsilon\subset B_\varepsilon(0)$.
Dla $u\in L^1_{\mathrm{loc}}(\Omega)$ i punktów
$x\in\Omega_\varepsilon
=\{x\in\Omega:\operatorname{dist}(x,\R^d\setminus\Omega)>\varepsilon\}$
wprowadzamy
$u_\varepsilon(x)=\int_\Omega u(y)\rho_\varepsilon(x-y)
\,\mathrm d\lambda_d(y)$. Dla dystrybucji stosujemy analogiczną
definicję $T_\varepsilon(x)=\langle T(y),\rho_\varepsilon(x-y)\rangle$.
Przyjmujemy odległość od pustego dopełnienia
równą $+\infty$, gdy $\Omega=\R^d$.
Dla $x\in\Omega_\varepsilon$ domknięta kula
$\overline{B_\varepsilon(x)}$ jest zwartym podzbiorem
$\Omega$, więc całka dla $u$ jest bezwzględnie
zbieżna, a argument funkcjonału $T$ jest testem.
W szczególności $(T_u)_\varepsilon=u_\varepsilon$.
Zapis $T(y)$ w nawiasie jest tylko sposobem wskazania zmiennej
funkcji testowej; nie oznacza wartości punktowej dystrybucji.
Splot jest tu \emph{lokalny}: dla danego $x$ musimy mieć
$\operatorname{supp}(\rho_\varepsilon(x-\cdot))\Subset\Omega$.
Nie wolno automatycznie rozciągać takiej definicji na całą
$\R^d$ bez przedłużenia $T$ lub założeń o nośniku.

\begin{twierdzenie}[Wygładzanie dystrybucji]
\label{thm:pde7-mollify}
Niech $\Omega\subset\R^d$ będzie otwarta, $T\in\mathcal D'(\Omega)$
i $\rho_\varepsilon$ jak wyżej. Wówczas
$T_\varepsilon\in C^\infty(\Omega_\varepsilon)$ oraz
\begin{equation}\label{eq:pde7-deriv-mollify}
 \partial^\alpha T_\varepsilon
 =(\partial^\alpha T)*\rho_\varepsilon
 =T*(\partial^\alpha\rho_\varepsilon)
 \quad\text{na }\Omega_\varepsilon.
\end{equation}
Ponadto dla każdej $\varphi\in\mathcal D(\Omega)$
i dostatecznie małego $\varepsilon>0$ zachodzi
\[
 \int_{\Omega_\varepsilon}T_\varepsilon(x)\varphi(x)
 \,\mathrm d\lambda_d(x)
 \longrightarrow\langle T,\varphi\rangle
 \qquad(\varepsilon\downarrow0).
\]
Jest to lokalna zbieżność dystrybucyjna na rosnących
zbiorach $\Omega_\varepsilon$. Gdy $\Omega=\R^d$,
jest to zwykła zbieżność w $\mathcal D'(\R^d)$.
\end{twierdzenie}
\begin{proof}
Ustalmy $\varepsilon>0$ oraz $x_0\in\Omega_\varepsilon$.
Wybierzmy $r>0$ tak, aby
$K=\overline{B_{\varepsilon+r}(x_0)}\Subset\Omega$.
Dla $|x-x_0|<r$ wszystkie testy
$F_x(y)=\rho_\varepsilon(x-y)$ oraz ich pochodne
względem $x$ mają nośniki zawarte w $K$.
Dla dowolnego wielowskaźnika $\alpha$ i wektora
standardowej bazy $e_j$ iloraz różnicowy
\[
 \frac{\partial_x^\alpha F_{x+he_j}
             -\partial_x^\alpha F_x}{h}
\]
zbiega do $\partial_x^{\alpha+e_j}F_x$
we wszystkich półnormach $p_{K,m}$.
Istotnie, po zastosowaniu dowolnej pochodnej
względem $y$ do ilorazu, podstawowe twierdzenie
rachunku całkowego zapisuje go jako średnią
odpowiedniej pochodnej względem $x_j$
na odcinku od $x$ do $x+he_j$.
Jednostajna ciągłość tej pochodnej na rozpatrywanych
zbiorach zwartych daje zbieżność jednostajną w $y$.
Argument bez ilorazu daje również ciągłość
$x\mapsto\partial_x^\alpha F_x$ we wszystkich
półnormach. Lokalna ciągłość $T$ pozwala zatem
indukcyjnie różniczkować działanie funkcjonału:
\[
 \partial_x^\alpha T_\varepsilon(x)
 =\langle T(y),(\partial^\alpha\rho_\varepsilon)(x-y)\rangle.
\]
Wszystkie te pochodne są ciągłe.
Ponieważ
$\partial_y^\alpha\rho_\varepsilon(x-y)
=(-1)^{|\alpha|}(\partial^\alpha\rho_\varepsilon)(x-y)$,
definicja pochodnej dystrybucyjnej daje
\[
 \langle\partial^\alpha T(y),\rho_\varepsilon(x-y)\rangle
 =\langle T(y),(\partial^\alpha\rho_\varepsilon)(x-y)\rangle.
\]
Dwa czynniki $(-1)^{|\alpha|}$ znoszą się.
To dowodzi~\eqref{eq:pde7-deriv-mollify}.

Ustalmy teraz $\varphi\in\mathcal D(\Omega)$,
przedłużoną przez zero na $\R^d$. Dla funkcji
zerowej teza jest oczywista. Dla niezerowej
wybieramy $r>0$ tak, aby
$K=\supp\varphi+\overline{B_{2r}(0)}\Subset\Omega$.
Dla $0<\varepsilon<r$ nośnik $\varphi$
zawiera się w $\Omega_\varepsilon$, a funkcje
$y\mapsto\varphi(x)\rho_\varepsilon(x-y)$
dla wszystkich $x$ mają nośniki w $K$.
Uzasadnijmy zamianę działania $T$ z całką.
Wybieramy domknięty prostokąt w $\R^d$ zawierający
nośnik $\varphi$ w swoim wnętrzu.
Sumy Riemanna całki funkcji
$(x,y)\mapsto\varphi(x)\rho_\varepsilon(x-y)$
względem $x$ zbiegają do całki we wszystkich
półnormach $p_{K,m}$: po każdej pochodnej
względem $y$ błąd jest ograniczony przez objętość
prostokąta razy moduł jednostajnej ciągłości
odpowiedniej pochodnej na iloczynie prostokąta i $K$,
dla średnicy podziału dążącej do zera.
Różniczkowanie samej całki względem $y$
jest dozwolone przez ciągłość i ograniczoność
pochodnych na tym iloczynie.
Zatem granica sum jest funkcją testową z nośnikiem
w $K$. Ciągłość $T$ na $\mathcal D_K(\Omega)$
pozwala przenieść działanie $T$ przez granicę sum.
Po stronie skalarnej są to sumy Riemanna
funkcji $\varphi T_\varepsilon$, przedłużonej przez
zero poza $\Omega_\varepsilon$; jest ona ciągła,
gdyż nośnik $\varphi$ leży wewnątrz tego zbioru.
Otrzymujemy
\[
 \begin{aligned}
 \int_{\Omega_\varepsilon}T_\varepsilon(x)\varphi(x)
 \,\mathrm d\lambda_d(x)
 &=\left\langle T(y),\int_{\R^d}
 \varphi(x)\rho_\varepsilon(x-y)\,\mathrm d\lambda_d(x)
 \right\rangle\\
 &=\langle T,\check\rho_\varepsilon*\varphi\rangle,
 \end{aligned}
\]
gdzie $\check\rho(z)=\rho(-z)$, a
$\check\rho_\varepsilon(z)=\rho_\varepsilon(-z)$.
Dla każdego wielowskaźnika $\alpha$ zachodzi
\[
 \partial^\alpha(\check\rho_\varepsilon*\varphi)(y)
 =\int_{\R^d}\rho(z)\partial^\alpha\varphi(y+\varepsilon z)
 \,\mathrm d\lambda_d(z).
\]
Jednostajna ciągłość $\partial^\alpha\varphi$
na $\R^d$, zwarty nośnik $\rho$ i jej całka $1$
dają zbieżność jednostajną do $\partial^\alpha\varphi$.
Wszystkie nośniki dla $\varepsilon<r$ są zawarte w $K$.
Z charakterystyki zbieżności w $\mathcal D(\Omega)$
wynika $\check\rho_\varepsilon*\varphi\to\varphi$,
a ciągłość $T$ daje tezę.
\end{proof}

\begin{twierdzenie}[Aproksymacja w $L^p_{\mathrm{loc}}$]
\label{thm:pde7-approx-id}
Niech $\Omega\subset\R^d$ będzie otwarta,
$1\le p<\infty$ i $u\in L^p_{\mathrm{loc}}(\Omega)$.
Wówczas $u_\varepsilon\to u$ w $L^p(K)$ dla każdego
$K\Subset\Omega$.
\end{twierdzenie}
\begin{proof}
Wystarczy rozważyć niepusty $K\Subset\Omega$.
Wybieramy $r>0$ tak, aby
$K' = K+\overline{B_{2r}(0)}\Subset\Omega$.
Niech $v=u\chi_{K'}$ na $\Omega$ i $v=0$
poza $\Omega$. Wtedy $v\in L^p(\R^d)$, a dla
$x\in K$, $0<\varepsilon<r$, zachodzi
$u_\varepsilon(x)=(v*\rho_\varepsilon)(x)$,
gdyż w całce występują tylko punkty z $K'$.
Ponadto $u=v$ prawie wszędzie na $K$.
Podstawienie $y=x-\varepsilon z$ i normalizacja
jądra dają różnicę jako całkę przesunięć $v$
minus $v$. Z całkowej nierówności Minkowskiego
otrzymujemy
\[
 \|v*\rho_\varepsilon-v\|_{L^p(\R^d)}
 \le\int_{\R^d}\rho(z)
 \|v(\cdot-\varepsilon z)-v\|_{L^p(\R^d)}
 \,\mathrm d\lambda_d(z).
\]
Uzasadnijmy ciągłość przesunięć bez odwołania
do gęstości funkcji gładkich, która sama może
wynikać z twierdzenia o wygładzaniu.
Funkcje ciągłe o zwartym nośniku są gęste w $L^p$
dla $p<\infty$: funkcje proste o nośniku skończonej
miary są gęste, a regularność miary pozwala
aproksymować ich składniki charakterystyczne
funkcjami ciągłymi o zwartym nośniku.
Dokładniej, dla zbioru $A$ skończonej miary
wybieramy zwarty $K_0\subset A$ i otwarty $O\supset A$
z $\lambda_d(O\setminus K_0)<\gamma$.
Funkcja ciągła o wartościach w $[0,1]$,
równa $1$ na $K_0$ i o zwartym nośniku w $O$,
może być skonstruowana za pomocą odległości
od $K_0$, jak w dowodzie twierdzenia Fejéra
w rozdziale o szeregach Fouriera.
Dla pustego $K_0$ bierzemy funkcję zerową.
Błąd w normie $L^p$ jest co najwyżej $\gamma^{1/p}$.
Nierówność trójkąta daje żądaną gęstość.

Dla $w\in C_c(\R^d)$ różnice $w(\cdot-b)-w$
przy $|b|\le1$ mają wspólny zwarty nośnik,
a ich normy supremum dążą do zera przy $b\to0$.
Stąd również ich normy $L^p$ dążą do zera.
Jeśli $\|v-w\|_p<\gamma$, to izometryczność
przesunięć daje
\[
 \|v(\cdot-b)-v\|_p
 \le 2\gamma+\|w(\cdot-b)-w\|_p.
\]
Z dowolności $\gamma$ wynika ciągłość przesunięć.
Dla $z\in\supp\rho$ mamy $|\varepsilon z|\le\varepsilon$,
więc normy pod całką w oszacowaniu Minkowskiego
dążą do zera jednostajnie względem $z$.
Ich całka dąży do zera, ponieważ $\int\rho=1$.
W konsekwencji
\[
 \|u_\varepsilon-u\|_{L^p(K)}
 \le\|v*\rho_\varepsilon-v\|_{L^p(\R^d)}
 \longrightarrow0.
\]
Dowodzi to tezy.
\end{proof}

\begin{uwaga}[Nie mylić sposobów zbieżności]
Dla $p=\infty$ zbieżność w normie na ogół nie zachodzi.
Na przykład dla funkcji Heaviside'a $H$ na $\R$
i każdej funkcji ciągłej $v$ na $[-1,1]$ mamy
$\|v-H\|_{L^\infty(-1,1)}\ge1/2$.
Istotnie, jeśli ta norma jest równa $M<\infty$,
to przez ciągłość $|v(x)|\le M$ dla wszystkich
$x<0$ w przedziale, a $|v(x)-1|\le M$ dla $x>0$.
Przechodząc do zera z obu stron, otrzymujemy
$1\le |v(0)|+|1-v(0)|\le2M$.
W szczególności nierówność dotyczy każdego
wygładzenia $H_\varepsilon$.
Dla ciągłej funkcji $u$ na $\Omega$ sytuacja
jest inna: na każdym zwartym $K\Subset\Omega$
dla małych $\varepsilon$ mamy
\[
 \sup_{x\in K}|u_\varepsilon(x)-u(x)|
 \le\sup_{\substack{x\in K\\|z|\le1}}
 |u(x-\varepsilon z)-u(x)|\longrightarrow0,
\]
z jednostajnej ciągłości $u$ na zwartym
otoczeniu $K$.
\end{uwaga}

\begin{center}
\begin{tikzpicture}[scale=.86,>=Stealth]
 \draw[->] (-4.3,0)--(4.7,0) node[below] {$x$};
 \draw[->] (0,-.1)--(0,3.5) node[left] {$\rho_\varepsilon(x)$};
 \draw[thick,blue!65!black,domain=-1.998:1.998,samples=101,smooth]
   plot (\x,{1.35*exp(1-1/(1-(\x/2)^2))});
 \draw[thick,red!70!black,domain=-.999:.999,samples=101,smooth]
   plot (\x,{2.7*exp(1-1/(1-(\x)^2))});
 \draw[blue!65!black,densely dashed] (-2,0)--(-2,-.2) (2,0)--(2,-.2);
 \draw[red!70!black,densely dashed] (-1,0)--(-1,-.2) (1,0)--(1,-.2);
 \node[below] at (-2,-.2) {$-2$};
 \node[below] at (-1,-.2) {$-1$};
 \node[below] at (1,-.2) {$1$};
 \node[below] at (2,-.2) {$2$};
 \node[blue!65!black,anchor=west] at (2.05,1.05) {$\varepsilon=2$};
 \node[red!70!black,anchor=west] at (.75,2.65) {$\varepsilon=1$};
 \node[below] at (0,-.1) {$0$};
\end{tikzpicture}
\par\smallskip\small
Jednowymiarowe profile znormalizowanego jądra
$\rho(x)=c\eta(x)$ z sekcji~\ref{sec:pde7-tests},
gdzie $c=(\int_{\R}\eta\,\mathrm d\lambda_1)^{-1}$.
Wspólna umowna skala pionowa: profil dla
$\varepsilon=1$ jest dwukrotnie wyższy i dwukrotnie
węższy od profilu dla $\varepsilon=2$.
Obie funkcje mają całkę $1$ i znikają poza
odpowiednimi przedziałami nośników.
\end{center}

\section{Słabe pochodne i pierwsze przestrzenie Sobolewa}
\label{sec:pde7-weak}

\begin{definicja}[Pochodna słaba]
\label{def:pde7-weak}
Niech $\Omega\subset\R^d$ będzie otwarta,
$u\in L^1_{\mathrm{loc}}(\Omega)$ i $\alpha\in\N_0^d$.
Funkcję $v\in L^1_{\mathrm{loc}}(\Omega)$ nazywamy
\emph{słabą pochodną} $\partial^\alpha u$, jeżeli
\begin{equation}\label{eq:pde7-weak}
 \int_\Omega u(x)\partial^\alpha\varphi(x)\,\mathrm d\lambda_d(x)
 =(-1)^{|\alpha|}\int_\Omega v(x)\varphi(x)\,
 \mathrm d\lambda_d(x)
 \quad (\varphi\in\mathcal D(\Omega)).
\end{equation}
\end{definicja}

Z definicji pochodnej dystrybucyjnej
~\eqref{eq:pde7-distributional-derivative}
wynika, że~\eqref{eq:pde7-weak} jest równoważne
równości $T_v=\partial^\alpha T_u$:
wystarczy pomnożyć obie strony tożsamości całkowej
przez $(-1)^{|\alpha|}$.
Jeśli $v_1,v_2$ spełniają tę tożsamość dla tego samego
$u$ i $\alpha$, to $T_{v_1-v_2}=0$. Ze stwierdzenia
~\ref{prop:pde7-injective} wynika
$v_1=v_2$ prawie wszędzie. Słaba pochodna jest więc
jednoznacznie wyznaczoną klasą funkcji modulo równość
prawie wszędzie. Zmiana reprezentanta $u$ nie wpływa
na tożsamość~\eqref{eq:pde7-weak}.

Jeśli $u\in C^k(\Omega)$, to dla $|\alpha|\le k$
klasyczna pochodna $\partial^\alpha u$ jest słabą
pochodną: stosujemy kolejno stwierdzenie
~\ref{prop:pde7-classical}.
Jeśli natomiast $v$ jest słabą pochodną $\partial^\alpha u$
i $w$ jest słabą pochodną $\partial^\beta v$, to
\[
 T_w=\partial^\beta T_v
 =\partial^\beta\partial^\alpha T_u
 =\partial^{\alpha+\beta}T_u.
\]
Zatem $w$ jest słabą pochodną $\partial^{\alpha+\beta}u$.
Nie wynika stąd istnienie pośrednich słabych pochodnych
z samego istnienia pochodnej wyższego rzędu.

Pochodna dystrybucyjna istnieje zawsze, ale słaba
pochodna w powyższym znaczeniu istnieje tylko wtedy,
gdy odpowiadająca jej dystrybucja jest regularna.
To rozróżnienie jest istotne: pochodna $H$ jest deltą,
która nie pochodzi od funkcji lokalnie całkowalnej.

\begin{definicja}[Przestrzeń $W^{1,p}$]
Niech $\Omega\subset\R^d$ będzie otwarta i $1\le p\le\infty$.
Przez $W^{1,p}(\Omega)$ oznaczamy zbiór tych
$u\in L^p(\Omega)$, dla których wszystkie pochodne
dystrybucyjne $\partial_jT_u$ są reprezentowane przez
funkcje z $L^p(\Omega)$. Przestrzeń wyposażamy w normę
$\|u\|_{W^{1,p}}=\|u\|_{L^p}+\sum_{j=1}^d
\|\partial_ju\|_{L^p}$.
\end{definicja}

Każda funkcja z $L^p(\Omega)$ dla $1\le p\le\infty$
jest lokalnie całkowalna: na zwartym zbiorze $K$
o skończonej mierze stosujemy nierówność Höldera
(dla $p=1$ wystarczy ograniczenie całki, a dla
$p=\infty$ oszacowanie przez $\lambda_d(K)\|u\|_\infty$).
Liniowość pochodnych dystrybucyjnych i ich jednoznaczność
dają liniowość słabych pochodnych na $W^{1,p}(\Omega)$.
Wobec tego jest to przestrzeń wektorowa, a podany
wzór rzeczywiście definiuje normę: jednorodność
i nierówność trójkąta wynikają z odpowiednich
własności norm $L^p$, a zerowanie normy wymusza
$u=0$ prawie wszędzie już z pierwszego składnika.
Pełną teorię przestrzeni Sobolewa, włącznie ze
zupełnością i twierdzeniami o zanurzeniach,
rozwiniemy w rozdziale~\ref{chap:pde-sobolev}.

\begin{uwaga}[Jak ten skrót ma się do dalszych rozdziałów?]
W tym rozdziale przestrzeń $W^{1,p}(\Omega)$ pojawia się
jedynie jako pierwszy przykład funkcji o słabych pochodnych.
Rozdział~\ref{chap:pde-sobolev} zbuduje ogólną teorię
$W^{k,p}(\Omega)$ na obszarach, wraz z zupełnością,
aproksymacją i przedłużaniem. Z kolei
rozdział~\ref{chap:pde-fourier-Rd} wprowadza
przestrzenie $H^s(\R^d)$ definiowane przez transformatę Fouriera.
Dla całkowitych rzędów $m\ge0$ na całej przestrzeni
zachodzi $H^m(\R^d)=W^{m,2}(\R^d)$ jako równość
przestrzeni z równoważnymi normami; uzasadnimy
ją w dalszej teorii. Na zbiorach otwartych należy
najpierw ustalić definicję przestrzeni $H^s$;
nie utożsamiamy automatycznie konstrukcji przez
ograniczenia funkcji z całej przestrzeni z definicją
przez słabe pochodne.
\end{uwaga}

\begin{przyklad}[Funkcja z załamaniem i funkcja ze skokiem]
Na $(-1,1)$ funkcja $u(x)=|x|$ należy do
$W^{1,p}((-1,1))$ dla każdego $1\le p\le\infty$:
jej słabą pochodną jest $\operatorname{sgn}(x)$.
Natomiast funkcja Heaviside'a ograniczona do $(-1,1)$
nie należy do $W^{1,p}((-1,1))$ dla żadnego
$1\le p\le\infty$, ponieważ jej pochodna
dystrybucyjna jest deltą i nie ma reprezentanta
lokalnie całkowalnego (dowody w sekcjach
~\ref{sec:pde7-def} i~\ref{sec:pde7-derivatives}).
Dla modułu $x$ zarówno funkcja, jak i jej pierwsza
słaba pochodna są ograniczone na przedziale
skończonej miary. Druga pochodna dystrybucyjna
jest jednak równa $2\delta_0$, więc druga słaba
pochodna nie istnieje w sensie~\eqref{eq:pde7-weak}.
Obie funkcje są niegładkie klasycznie, lecz tylko
pierwsza ma lokalnie całkowalną pierwszą słabą pochodną.
\end{przyklad}

\begin{twierdzenie}[Pochodna zerowa oznacza stałość w wymiarze jeden]
\label{thm:pde7-constant}
Niech $I\subset\R$ będzie niepustym przedziałem otwartym,
$T\in\mathcal D'(I)$ i $T'=0$. Istnieje dokładnie jedna
stała $c\in\C$ taka, że $T=T_c$ na $I$.
\end{twierdzenie}
\begin{proof}
Wybierzmy $\eta\in\mathcal D(I)$ o
$\int_I\eta(x)\,\mathrm d\lambda_1(x)=1$
(znormalizowaną funkcję testową na kuli
zawartej w $I$) i połóżmy $c=\langle T,\eta\rangle$.
Dla dowolnej $\varphi\in\mathcal D(I)$ zdefiniujmy
$\psi=\varphi-\bigl(\int_I\varphi(x)\,\mathrm d\lambda_1(x)\bigr)\eta$.
Wtedy $\int_I\psi(x)\,\mathrm d\lambda_1(x)=0$.
Przedłużamy $\psi$ zerem na $\R$ i wybieramy
domknięty przedział $[a,b]\Subset I$ zawierający
jej nośnik; jest to możliwe, ponieważ $I$ jest
przedziałem, także gdy jest nieograniczony.
Funkcja $\Psi(x)=\int_{-\infty}^x\psi(t)
\,\mathrm d\lambda_1(t)$ jest zerowa dla $x<a$,
a dla $x>b$ całka jest równa całce z $\psi$
po całej prostej, czyli również zeru.
Zatem $\supp\Psi\subset[a,b]$.
Z podstawowego twierdzenia rachunku całkowego
$\Psi'=\psi$ i gładkość $\psi$ daje gładkość
$\Psi$, więc $\Psi|_I\in\mathcal D(I)$.
Mamy $\Psi'=\psi$, a więc
$\langle T,\psi\rangle=\langle T,\Psi'\rangle
=-\langle T',\Psi\rangle=0$. Otrzymujemy
$\langle T,\varphi\rangle=c\int_I\varphi(x)\,\mathrm d\lambda_1(x)$.
Jednoznaczność $c$ wynika z wyboru testu o niezerowej całce.
\end{proof}

\begin{wniosek}[Zerowy gradient dystrybucyjny]
Niech $\Omega\subset\R^d$ będzie niepustym spójnym
zbiorem otwartym, a $T\in\mathcal D'(\Omega)$
spełnia $\partial_jT=0$ dla $j=1,\ldots,d$.
Wówczas istnieje dokładnie jedna stała $c\in\C$
taka, że $T=T_c$.
W szczególności jeśli $u\in L^1_{\mathrm{loc}}(\Omega)$
ma wszystkie pierwsze słabe pochodne równe zeru,
to $u=c$ prawie wszędzie.
Na niespójnym zbiorze otwartym wniosek zachodzi
osobno na każdej składowej spójnej.
\end{wniosek}
\begin{proof}
Wybierzmy dowolną kulę $B=B_r(a)$ z
$\overline{B_{2r}(a)}\Subset\Omega$.
Dla $0<\varepsilon<r$ wygładzenie $T_\varepsilon$
jest gładkie na $B$ i na mocy
~\eqref{eq:pde7-deriv-mollify} mamy
$\partial_jT_\varepsilon=(\partial_jT)*\rho_\varepsilon=0$.
Z całkowania wzdłuż odcinka łączącego dwa
punkty kuli (pozostaje on w kuli) wynika
$T_\varepsilon(x)=c_\varepsilon$ dla $x\in B$.
Wybierzmy test $\eta\in\mathcal D(B)$
o całce $1$. Z twierdzenia~\ref{thm:pde7-mollify}
otrzymujemy
\[
 c_\varepsilon=\int_B T_\varepsilon(x)\eta(x)\,
 \mathrm d\lambda_d(x)\longrightarrow
 c_B:=\langle T,\eta\rangle.
\]
Dla dowolnego $\varphi\in\mathcal D(B)$ to samo
twierdzenie daje
\[
 \langle T,\varphi\rangle
 =\lim_{\varepsilon\downarrow0}
 c_\varepsilon\int_B\varphi\,\mathrm d\lambda_d
 =c_B\int_B\varphi\,\mathrm d\lambda_d.
\]
Zatem $T|_B=T_{c_B}$.

Jeśli dwie takie kule przecinają się, to na ich
otwartym przecięciu wybieramy test o całce $1$.
Porównanie dwóch ograniczeń $T$ daje równość
odpowiadających im stałych. Kule tego rodzaju
pokrywają $\Omega$, zatem wyznaczają dobrze
określoną lokalnie stałą funkcję $c(x)$ na $\Omega$.
Spójność $\Omega$ wymusza jej stałość: zbiór punktów,
w których przyjmuje ustaloną wartość, oraz jego
dopełnienie są otwarte. Z lokalności równości
dystrybucji (sekcja~\ref{sec:pde7-support})
otrzymujemy $T=T_c$ na $\Omega$.
Jednoznaczność wynika z testu o całce $1$.
Dla $T=T_u$ stosujemy jeszcze iniektywność
odwzorowania $u\mapsto T_u$.
Składowe spójne zbioru otwartego w $\R^d$
są otwarte, ponieważ każdy punkt ma spójną kulę
zawartą w tym zbiorze. Możemy zatem stosować
powyższy dowód do każdej z nich.
\end{proof}

\section{Rozwiązania fundamentalne operatorów różniczkowych}
\label{sec:pde7-fundamental}

\begin{definicja}[Rozwiązanie fundamentalne]
Niech $P(\partial)=\sum_{|\alpha|\le m}a_\alpha\partial^\alpha$
będzie operatorem różniczkowym o stałych współczynnikach
na $\R^d$, gdzie $a_\alpha\in\C$.
Dystrybucję $E\in\mathcal D'(\R^d)$ nazywamy
\emph{rozwiązaniem fundamentalnym} $P$, gdy
$P(\partial)E=\delta_0$ w $\mathcal D'(\R^d)$.
\end{definicja}

Jeśli $f\in C_c^\infty(\R^d)$, możemy zdefiniować
$(E*f)(x)=\langle E(y),f(x-y)\rangle$.
Na mocy poniższego lematu funkcja ta jest gładka
i spełnia $P(\partial)(E*f)=f$.
Stąd rozwiązanie fundamentalne daje narzędzie rozwiązywania
niejednorodnego równania na całej przestrzeni.
Nie twierdzimy na tym etapie, że $E*f$ opisuje wszystkie
rozwiązania ani że bez dodatkowych warunków można w ten sposób
rozwiązywać dowolne zadanie brzegowe.


\begin{lemat}[Różniczkowanie splotu z dystrybucją]
\label{lem:pde7-convolution-derivatives}
Niech $d\ge1$, $T\in\mathcal D'(\R^d)$ i
$\psi\in C_c^\infty(\R^d)$. Funkcja
$u(x)=\langle T(y),\psi(x-y)\rangle$ należy do
$C^\infty(\R^d)$ i dla każdego wielowskaźnika
$\alpha\in\N_0^d$ zachodzi
\[
 \partial^\alpha(T*\psi)
 =T*(\partial^\alpha\psi)=(\partial^\alpha T)*\psi.
\]
W szczególności dla operatora o stałych współczynnikach
$P(\partial)$ i rozwiązania fundamentalnego $E$ mamy
$P(\partial)(E*\psi)=\psi$.
\end{lemat}
\begin{proof}
Niech $A\subset\R^d$ będzie niepustym zbiorem zwartym
i niech $B=\supp\psi$. Dla $x\in A$ nośniki
funkcji $y\mapsto\psi(x-y)$ zawierają się w zwartym
zbiorze $A-B=\{x-z:x\in A,\ z\in B\}$.
Dla różniczkowania w otoczeniu $A$ wystarczy
zastąpić $A$ jego zwartym poszerzeniem.
Stosujemy argument z dowodu twierdzenia
~\ref{thm:pde7-mollify}, w którym dla gładkości
splotu i rachunku pochodnych nie używaliśmy
ani dodatniości jądra, ani warunku całki równej $1$.
Dokładniej, po każdej pochodnej względem $y$
ilorazy różnicowe względem $x_j$ zbiegają
jednostajnie na wspólnym zwartym nośniku, dzięki
podstawowemu twierdzeniu rachunku całkowego
i jednostajnej ciągłości pochodnych $\psi$.
Zatem zbieżność zachodzi we wszystkich półnormach
przestrzeni testowej o tym nośniku.
Ciągłość $T$ i indukcja względem rzędu pochodnej dają
$\partial_x^\alpha u(x)=\langle T(y),
(\partial^\alpha\psi)(x-y)\rangle$,
a ciągłość tych pochodnych wynika z ciągłości
przesunięć testów w tych samych półnormach.
Przypadek $\psi=0$ jest oczywisty.
Dla pochodnej pierwszego rzędu korzystamy z definicji
pochodnej dystrybucyjnej:
\[
 \langle\partial_{y_j}T(y),\psi(x-y)\rangle
 =-\langle T(y),\partial_{y_j}\psi(x-y)\rangle
 =\langle T(y),\partial_{x_j}\psi(x-y)\rangle.
\]
Iteracja daje wszystkie wielowskaźniki. Na koniec
$P(\partial)(E*\psi)=(P(\partial)E)*\psi
=\delta_0*\psi=\psi$, gdzie ostatnia równość
wynika bezpośrednio z działania delty na test.
\end{proof}

\begin{twierdzenie}[Rozwiązanie fundamentalne operatora $-\Delta$]
\label{thm:pde7-laplace}
Niech $d\ge1$ i niech $\omega_d$ oznacza pole powierzchni
sfery jednostkowej $S^{d-1}\subset\R^d$ dla $d\ge2$.
Funkcja lokalnie całkowalna
\begin{equation}\label{eq:pde7-laplace-fundamental}
 \Phi_d(x)=\begin{cases}
 -|x|/2,&d=1,\\
 (2\pi)^{-1}\log(1/|x|),&d=2,\\
 ((d-2)\omega_d)^{-1}|x|^{2-d},&d\ge3,
 \end{cases}
\end{equation}
spełnia $-\Delta T_{\Phi_d}=\delta_0$ na $\R^d$.
\end{twierdzenie}
\begin{proof}
Wartość $\Phi_d(0)$ w wymiarach $d\ge2$ ustalamy
dowolnie, na przykład równą $0$.
W wymiarze jeden $\Phi_1$ jest lokalnie całkowalna
i stwierdzenie wynika z $\partial_x^2T_{|x|}=2\delta_0$
(przykład~\ref{ex:pde7-heaviside}).

Dla $d\ge3$ całkowanie w zmiennych biegunowych daje
dla $r>0$
\[
 \int_{B_r(0)}|\Phi_d(x)|\,\mathrm d\lambda_d(x)
 =\frac1{d-2}\int_0^r s\,\mathrm d\lambda_1(s)
 =\frac{r^2}{2(d-2)}.
\]
Dla $d=2$ mamy $\omega_2=2\pi$ i dla $0<r<1$
\[
 \begin{aligned}
 \int_{B_r(0)}|\Phi_2(x)|\,\mathrm d\lambda_2(x)
 &=\int_0^r s\log(1/s)\,\mathrm d\lambda_1(s)\\
 &=\frac{r^2}{2}\log(1/r)+\frac{r^2}{4}.
 \end{aligned}
\]
Granica dolnego wyrazu pierwotnej jest zerowa,
gdyż $s^2\log(1/s)\to0$ dla $s\downarrow0$.
W obu przypadkach całka dąży do zera przy $r\downarrow0$.
Z gładkości poza zerem wynika więc lokalna
całkowalność na całej $\R^d$.

Niech $F(r)$ oznacza profil radialny $\Phi_d$ dla $r>0$.
Bezpośrednie różniczkowanie obu wzorów daje
\[
 F'(r)=-\frac1{\omega_dr^{d-1}},\qquad
 F''(r)=\frac{d-1}{\omega_dr^d}.
\]
Dla $x\ne0$, $r=|x|$,
$\partial_jr=x_j/r$ oraz
$\partial_j^2r=1/r-x_j^2/r^3$.
Reguła łańcuchowa daje więc
\[
 \Delta\Phi_d(x)
 =F''(r)+\frac{d-1}{r}F'(r)=0.
\]
Dla $\varphi\in\mathcal D(\R^d)$ wybieramy $R>0$
tak, aby $\supp\varphi\Subset B_R(0)$.
Dla $0<\varepsilon<R$ stosujemy drugą tożsamość
Greena~\eqref{eq:pde3-green2} na pierścieniu
$A_{\varepsilon,R}=B_R(0)\setminus\overline{B_\varepsilon(0)}$
(do części rzeczywistej i urojonej testu osobno).
Ponieważ $\Delta\Phi_d=0$ na pierścieniu,
otrzymujemy
\[
 -\int_{A_{\varepsilon,R}}\Phi_d\Delta\varphi
 \,\mathrm d\lambda_d
 =\int_{\partial A_{\varepsilon,R}}
 (\varphi\partial_\nu\Phi_d-\Phi_d\partial_\nu\varphi)
 \,\mathrm d\sigma.
\]
Na zewnętrznej sferze znikają test i jego pochodne.
Na wewnętrznej sferze $\nu=-x/|x|$, zatem
$\partial_\nu\Phi_d=-F'(\varepsilon)
=1/(\omega_d\varepsilon^{d-1})$.
W konsekwencji
\[
 -\int_{B_R\setminus B_\varepsilon}\Phi_d(x)\Delta\varphi(x)
 \,\mathrm d\lambda_d(x)
 =\int_{\partial B_\varepsilon}\varphi(x)
 \frac{1}{\omega_d\varepsilon^{d-1}}\,\mathrm d\sigma(x)
 -\int_{\partial B_\varepsilon}\Phi_d(x)
 \partial_\nu\varphi(x)\,\mathrm d\sigma(x).
\]
Miara powierzchniowa sfery o promieniu $\varepsilon$
wynosi $\omega_d\varepsilon^{d-1}$. Po odjęciu
$\varphi(0)$ od pierwszej całki moduł różnicy
jest więc co najwyżej
$\sup_{|x|=\varepsilon}|\varphi(x)-\varphi(0)|$,
co dąży do zera z ciągłości testu.
Moduł drugiej całki nie przekracza
\[
 \|\nabla\varphi\|_\infty\,
 |F(\varepsilon)|\,\omega_d\varepsilon^{d-1}.
\]
Dla $d\ge3$ iloczyn bez normy gradientu wynosi
$\varepsilon/(d-2)$, dla $d=2$ i $\varepsilon<1$
wynosi $\varepsilon\log(1/\varepsilon)$.
Obie wielkości dążą do zera.
Ponadto moduł pominiętej całki po $B_\varepsilon(0)$
jest ograniczony przez
$\|\Delta\varphi\|_\infty
\int_{B_\varepsilon(0)}|\Phi_d|\,\mathrm d\lambda_d$,
co dąży do zera z wcześniejszych obliczeń.
Ze znaku $(-1)^2=1$ przy drugiej pochodnej
dystrybucyjnej otrzymujemy
\[
 \langle-\Delta T_{\Phi_d},\varphi\rangle
 =-\int_{\R^d}\Phi_d(x)\Delta\varphi(x)
 \,\mathrm d\lambda_d(x)=\varphi(0).
\]
Jest to żądana równość.
\end{proof}

\begin{uwaga}[Zgodność z wcześniejszą teorią potencjału]
W rozdziale~\ref{chap:pde-laplace},
w twierdzeniu~\ref{thm:pde3-phi}, otrzymaliśmy te same
funkcje dla $d\ge2$ oraz tę samą tożsamość całkową.
Tutaj interpretujemy ją w przestrzeni dystrybucji
i uwzględniamy także przypadek $d=1$.
Z lematu~\ref{lem:pde7-convolution-derivatives}
wynika w szczególności, że dla $f\in\mathcal D(\R^d)$
funkcja
\[
 u(x)=\int_{\R^d}\Phi_d(x-y)f(y)\,\mathrm d\lambda_d(y)
\]
jest gładka i spełnia $-\Delta u=f$.
Całka jest bezwzględnie zbieżna dzięki zwartemu
nośnikowi $f$ i lokalnej całkowalności $\Phi_d$;
podstawienie $z=x-y$ utożsamia ją ze splotem
$T_{\Phi_d}*f$. Nie wyciągamy stąd wniosku
o jednoznaczności rozwiązania bez dodatkowych
warunków: dodanie dowolnej gładkiej funkcji
harmonicznej daje kolejne rozwiązanie.
\end{uwaga}

\section{Równanie ciepła i delta w chwili początkowej}
\label{sec:pde7-heat}

\begin{twierdzenie}[Jądro cieplne jako rozwiązanie fundamentalne]
\label{thm:pde7-heat}
Niech $d\ge1$ i $\kappa>0$. Na $\R^{1+d}$ definiujemy
funkcję
\[
 E(t,x)=\begin{cases}
 (4\pi\kappa t)^{-d/2}
 \exp\!\bigl(-|x|^2/(4\kappa t)\bigr),&t>0,\\
 0,&t\le0.
 \end{cases}
\]
Wtedy $E\in L^1_{\mathrm{loc}}(\R^{1+d})$ i dla
operatora $L=\partial_t-\kappa\Delta_x$ zachodzi
$L T_E=\delta_{(0,0)}$ w $\mathcal D'(\R^{1+d})$.
\end{twierdzenie}
\begin{proof}
Oznaczmy przez $\Gamma_\kappa(t,x)$ jądro cieplne
dla $t>0$. Z twierdzenia~\ref{thm:pde4-gauss-basic}
wiemy, że jest ono dodatnie, gładkie, spełnia
$(\partial_t-\kappa\Delta_x)\Gamma_\kappa=0$ oraz
ma całkę przestrzenną równą $1$ dla każdego $t>0$.
Dla dowolnego zwartego zbioru
$K\subset\R^{1+d}$ wybieramy $A>0$ tak, aby
$K\subset[-A,A]\times\R^d$.
Z twierdzenia Tonellego i normalizacji jądra
otrzymujemy
\[
 \int_K|E(t,x)|\,\mathrm d\lambda_{1+d}(t,x)
 \le\int_0^A\int_{\R^d}\Gamma_\kappa(t,x)
 \,\mathrm d\lambda_d(x)\,\mathrm d\lambda_1(t)=A.
\]
Zatem $E$ jest lokalnie całkowalna mimo osobliwości
w punkcie $(0,0)$.

Ustalmy $\varphi\in\mathcal D(\R^{1+d})$.
Z definicji pochodnej dystrybucyjnej (znak minus
dla pierwszej pochodnej i znak plus dla drugiej)
wynika
\[
 \langle LT_E,\varphi\rangle
 =\int_0^\infty\int_{\R^d}\Gamma_\kappa(t,x)
 [-\partial_t\varphi(t,x)-\kappa\Delta_x\varphi(t,x)]
 \,\mathrm d\lambda_d(x)\,\mathrm d\lambda_1(t).
\]
Całka jest bezwzględnie zbieżna: pochodne testu
są ograniczone i mają wspólny zwarty nośnik,
a $E$ jest lokalnie całkowalna. Co więcej,
moduł części całki dla $0<t<\varepsilon$ nie
przekracza
\[
 \varepsilon
 \bigl(\|\partial_t\varphi\|_\infty+
             \kappa\|\Delta_x\varphi\|_\infty\bigr),
\]
ponieważ jądro ma masę przestrzenną $1$.
Wystarczy zatem obliczyć całkę $I_\varepsilon$
po $t\ge\varepsilon$ i przejść do granicy.

Wybieramy $A>\varepsilon$ oraz $R>0$ tak, aby
nośnik testu był zawarty w
$(-A,A)\times B_R(0)$. Wszystkie jego pochodne
znikają w otoczeniu górnej granicy czasowej $A$
i przestrzennego brzegu kuli.
Na $[\varepsilon,A]\times\overline{B_R(0)}$
jądro oraz jego pochodne są gładkie i ograniczone.
Możemy więc stosować twierdzenie Fubiniego,
całkowanie przez części w czasie oraz dwukrotne
całkowanie przez części względem przestrzeni.
Otrzymujemy
\[
\begin{aligned}
 I_\varepsilon
 &=\int_{\R^d}\Gamma_\kappa(\varepsilon,x)
           \varphi(\varepsilon,x)\,\mathrm d\lambda_d(x)\\
 &\quad+\int_\varepsilon^A\int_{\R^d}
  (\partial_t\Gamma_\kappa-\kappa\Delta_x\Gamma_\kappa)
       \varphi\,\mathrm d\lambda_d\,\mathrm d\lambda_1.
\end{aligned}
\]
W pierwszym składniku znak jest dodatni,
ponieważ dolna granica całkowania w czasie
występuje z przeciwnym znakiem w całce
z $-\Gamma_\kappa\partial_t\varphi$.
Nie ma składników brzegowych w przestrzeni,
gdyż test wraz z pochodnymi znika przy
brzegu wybranej kuli. Drugi składnik jest
zerowy na mocy równania cieplnego dla $t>0$.

Podstawienie $x=\sqrt{4\kappa\varepsilon}\,z$
o jakobianie $(4\kappa\varepsilon)^{d/2}$ daje
\[
 I_\varepsilon
 =\pi^{-d/2}\int_{\R^d}e^{-|z|^2}
 \varphi(\varepsilon,\sqrt{4\kappa\varepsilon}\,z)
 \,\mathrm d\lambda_d(z).
\]
Dla każdego ustalonego $z$ argument testu zmierza
do $(0,0)$, a moduł funkcji podcałkowej wraz
z czynnikiem normalizacyjnym nie przekracza
$\pi^{-d/2}e^{-|z|^2}\|\varphi\|_\infty$.
Jest to całkowalna majoranta niezależna od
$\varepsilon$. Z twierdzenia o zbieżności
zmajoryzowanej oraz normalizacji całki Gaussa
wynika $I_\varepsilon\to\varphi(0,0)$.
Ostatecznie
\[
 \langle LT_E,\varphi\rangle
 =\lim_{\varepsilon\downarrow0} I_\varepsilon
 =\varphi(0,0)
 =\langle\delta_{(0,0)},\varphi\rangle.
\]
\end{proof}

\begin{uwaga}[Nie jest to rozwiązanie klasyczne w $t=0$]
Funkcja $E$ jest klasycznie gładka dla $t>0$ oraz
dla $t<0$. Jej przedłużenie przez zero na czasy
niedodatnie nie spełnia jednak równania jednorodnego
na całej czasoprzestrzeni: w sensie dystrybucji
punkt $(0,0)$ jest źródłem o masie $1$.
Nie jest nim cała hiperpowierzchnia $t=0$.
W szczególności $LT_E=0$ na każdym otwartym
zbiorze niezawierającym $(0,0)$.

Dla ustalonego $t>0$ jądro przestrzenne
$\Gamma_\kappa(t,\cdot)$ wyznacza dystrybucję na
$\R^d$, która przy $t\downarrow0$ zbiega do
$\delta_0$: dla testu $\psi\in\mathcal D(\R^d)$
stosujemy to samo podstawienie i twierdzenie
o zbieżności zmajoryzowanej do całki
$\int_{\R^d}\Gamma_\kappa(t,x)\psi(x)\,\mathrm d\lambda_d(x)$.
Jest to granica w przestrzeni dystrybucji
przestrzennych, a nie punktowa wartość funkcji
$E(0,x)$.
\end{uwaga}

\begin{wniosek}[Czasoprzestrzenny wzór Duhamela]
Jeśli $f\in\mathcal D(\R^{1+d})$, to funkcja
$u=T_E*f$ jest gładka na całej czasoprzestrzeni,
spełnia $Lu=f$ i ma postać
\[
 u(t,x)=\int_{-\infty}^t\int_{\R^d}
 \Gamma_\kappa(t-s,x-y)f(s,y)
 \,\mathrm d\lambda_d(y)\,\mathrm d\lambda_1(s).
\]
Jeżeli $f(s,y)=0$ dla wszystkich $s<t_0$ i $y\in\R^d$,
to $u(t,x)=0$ dla $t<t_0$.
\end{wniosek}
\begin{proof}
Z lematu~\ref{lem:pde7-convolution-derivatives},
zastosowanego w wymiarze $1+d$, i twierdzenia
~\ref{thm:pde7-heat} wynika
$u\in C^\infty(\R^{1+d})$ oraz
$Lu=(LT_E)*f=\delta_{(0,0)}*f=f$.
Ponieważ $E$ jest lokalnie całkowalna, splot
z testem jest całką bezwzględnie zbieżną:
dla ustalonego $(t,x)$ argument $E$ przebiega
zwarty zbiór będący przesunięciem odbitego
nośnika $f$. Zmiana zmiennych w splocie daje
całkę z $E(t-s,x-y)f(s,y)$.
Dzięki bezwzględnej całkowalności stosujemy
twierdzenie Fubiniego, a ponieważ $E(t-s,x-y)=0$
dla $s\ge t$, pozostaje podany wzór.
Dla $t<t_0$ każda wartość $s<t$ leży w zakresie
zanikania $f$, więc $u(t,x)=0$.
Tak otrzymujemy przyczynowe rozwiązanie równania
z wymuszeniem o zwartym nośniku; nie orzekamy
jednoznaczności bez dalszych założeń.
\end{proof}

\section{Lemat Weyla: rozwiązanie dystrybucyjne jest gładkie}
\label{sec:pde7-weyl}

\begin{twierdzenie}[Lemat Weyla]
\label{thm:pde7-weyl}
Niech $d\ge2$, $\Omega\subset\R^d$ będzie otwarta,
a $T\in\mathcal D'(\Omega)$ spełnia $\Delta T=0$.
Wówczas istnieje dokładnie jedna funkcja
$h\in C^\infty(\Omega;\C)$ taka, że $T=T_h$
oraz $\Delta h=0$ punktowo.
W szczególności $h$ jest rzeczywiście analityczna
(tzn. jej część rzeczywista i urojona są
rzeczywiście analityczne).
\end{twierdzenie}
\begin{proof}
Ustalmy $a\in\Omega$ i wybierzmy $r>0$ takie, że
$\overline{B_{3r}(a)}\Subset\Omega$.
Dla $0<\varepsilon<r$ definiujemy
$u_\varepsilon=T*\rho_\varepsilon$ na $B_{2r}(a)$.
Rzeczywiście, jeśli $x\in B_{2r}(a)$, to
$\overline{B_\varepsilon(x)}\subset B_{3r}(a)$;
w szczególności $x\in\Omega_\varepsilon$.
Z twierdzenia~\ref{thm:pde7-mollify} wynika
$u_\varepsilon\in C^\infty(B_{2r}(a))$, a wzór
~\eqref{eq:pde7-deriv-mollify} daje
\[
 \Delta u_\varepsilon=(\Delta T)*\rho_\varepsilon=0
 \qquad\text{na }B_{2r}(a).
\]

Wybierzmy radialną funkcję
$\eta\in C_c^\infty(B_r(0))$, nieujemną i o całce $1$.
Możemy ją otrzymać przez skalowanie i normalizację
radialnej funkcji testowej z sekcji~\ref{sec:pde7-tests}.
Zapiszmy $\eta(z)=q(|z|)$, gdzie $q$ jest jej
profilem radialnym. Dla $x\in B_r(a)$ i $0<s<r$
mamy $\overline{B_s(x)}\subset B_{2r}(a)$.
Stosując twierdzenie~\ref{thm:pde3-mean}
do części rzeczywistej i urojonej $u_\varepsilon$,
otrzymujemy
\[
 \int_{S^{d-1}}u_\varepsilon(x-s\theta)\,
 \mathrm d\sigma(\theta)
 =\omega_d u_\varepsilon(x).
\]
Wykorzystaliśmy również niezmienniczość miary
powierzchniowej przy przekształceniu
$\theta\mapsto-\theta$.
Po całkowaniu w zmiennych biegunowych dostajemy
\[
\begin{aligned}
 \int_{B_r(0)}u_\varepsilon(x-z)\eta(z)\,
 \mathrm d\lambda_d(z)
 &=\int_0^r q(s)s^{d-1}
 \int_{S^{d-1}}u_\varepsilon(x-s\theta)\,
 \mathrm d\sigma(\theta)\,\mathrm d\lambda_1(s)\\
 &=u_\varepsilon(x)\omega_d
 \int_0^r q(s)s^{d-1}\,\mathrm d\lambda_1(s)
 =u_\varepsilon(x).
\end{aligned}
\]
Ostatnia równość wynika z normalizacji $\eta$.
Wszystkie argumenty $x-z$ występujące w całce
należą do $B_{2r}(a)$, więc nie przedłużamy
$u_\varepsilon$ na całą przestrzeń.
Po podstawieniu $y=x-z$ otrzymujemy na $B_r(a)$
\[
 u_\varepsilon(x)=
 \int_{B_{2r}(a)}u_\varepsilon(y)\eta(x-y)\,
 \mathrm d\lambda_d(y).
\]

Niech teraz $\varphi\in\mathcal D(B_r(a))$;
przedłużamy ją przez zero. Oznaczmy
$\check\eta(z)=\eta(-z)$ (tu $\check\eta=\eta$).
Splot $\check\eta*\varphi$ jest gładki
i ma nośnik zawarty w zwartym zbiorze
$\supp\varphi+\supp\check\eta\subset B_{2r}(a)$.
Gładkość wynika z różniczkowania całki o zwartym
nośniku, a inkluzja nośników z faktu, że poza
wskazaną sumą zbiory występujące w splocie
nie mają punktu wspólnego.
Jest to zatem test w $B_{2r}(a)$.
W całce podwójnej po $x,y$ niezerowe składniki
występują tylko wtedy, gdy $x\in\supp\varphi$
oraz $y\in x-\supp\eta$. Obie zmienne przebiegają
wówczas zbiory zwarte wewnątrz dziedzin
odpowiednich funkcji. Ponieważ $u_\varepsilon$
jest tam ograniczona, całka jest bezwzględnie
zbieżna i można zastosować twierdzenie Fubiniego:
\[
 \int_{B_r(a)}u_\varepsilon(x)\varphi(x)\,
 \mathrm d\lambda_d(x)
 =\int_{B_{2r}(a)}u_\varepsilon(y)
 (\check\eta*\varphi)(y)\,\mathrm d\lambda_d(y).
\]
W twierdzeniu~\ref{thm:pde7-mollify} wykazaliśmy
lokalną zbieżność dystrybucyjną $u_\varepsilon$ do $T$.
Oba testy są ustalone i mają zwarte nośniki
w $\Omega$, więc po przejściu do granicy otrzymujemy
\[
 \langle T,\varphi\rangle
 =\langle T,\check\eta*\varphi\rangle.
\]

Na $B_r(a)$ określamy funkcję
$h_a(x)=\langle T(y),\eta(x-y)\rangle$.
Z argumentu z dowodu twierdzenia
~\ref{thm:pde7-mollify}, z jądrem $\eta$ zamiast
$\rho_\varepsilon$, wynika, że $h_a$ jest gładka.
Argument ten dotyczy także ustalonego jądra:
w otoczeniu każdego $x\in B_r(a)$ wszystkie
przesunięte testy mają wspólny zwarty nośnik
w $B_{2r}(a)$.
Ponadto uzasadniona tam zamiana działania
funkcjonału z całką (przez zbieżność sum Riemanna
we wszystkich półnormach) daje
\[
 \int_{B_r(a)}h_a(x)\varphi(x)\,\mathrm d\lambda_d(x)
 =\langle T,\check\eta*\varphi\rangle
 =\langle T,\varphi\rangle.
\]
Mamy więc $T|_{B_r(a)}=T_{h_a}$.
Różniczkowanie splotu daje również
$\Delta h_a=(\Delta T)*\eta=0$ na tej kuli.

Kule tego rodzaju pokrywają $\Omega$. Jeśli dwie
kule przecinają się, to ich lokalne reprezentacje
wyznaczają tę samą dystrybucję na otwartym
przecięciu. Ze stwierdzenia~\ref{prop:pde7-injective}
są tam równe prawie wszędzie, a dzięki ciągłości
są równe w każdym punkcie (niezerowa różnica
miałaby niezerowy moduł na pewnej kuli
dodatniej miary).
Otrzymujemy zatem dobrze określoną gładką
funkcję $h$ na $\Omega$, która lokalnie jest
równa odpowiednim $h_a$.
Lokalność równości dystrybucji z sekcji
~\ref{sec:pde7-support} daje $T=T_h$,
a lokalna harmoniczność daje $\Delta h=0$.
Jednoznaczność wśród gładkich reprezentacji
wynika ponownie z iniektywności i ciągłości.
Na koniec twierdzenie~\ref{thm:pde3-analytic},
zastosowane osobno do części rzeczywistej
i urojonej, daje analityczność.
\end{proof}

\begin{uwaga}[Znaczenie założenia harmoniczności]
Lemat Weyla nie mówi, że każda dystrybucja jest
funkcją gładką: kluczowym założeniem jest
równanie $\Delta T=0$.
Na przykład rozwiązanie fundamentalne $\Phi_d$
z twierdzenia~\ref{thm:pde7-laplace} jest gładkie
i harmoniczne poza zerem, ale spełnia
$-\Delta T_{\Phi_d}=\delta_0$ na całej przestrzeni.
Nie jest więc rozwiązaniem jednorodnego równania
dystrybucyjnego na żadnym otoczeniu zera.
\end{uwaga}

\section{Ograniczenia rachunku dystrybucji}
\label{sec:pde7-limits}

Pochodne i mnożenie przez funkcje gładkie są dobrze określone
na $\mathcal D'(\Omega)$, jak wykazano w sekcji
~\ref{sec:pde7-derivatives}. Nie można jednak wprowadzić
na przestrzeni wszystkich dystrybucji iloczynu, który
jednocześnie miałby wszystkie zwykłe własności
algebraiczne, zachowywał mnożenie funkcji lokalnie
ograniczonych i spełniał regułę Leibniza.

\begin{stwierdzenie}[Przeszkoda w rozszerzeniu iloczynu]
Na $\mathcal D'(\R)$ nie istnieje przemienny, łączny
i dwuliniowy iloczyn $\star$ spełniający jednocześnie
następujące warunki:
\begin{enumerate}
\item $T_f\star T_g=T_{fg}$ dla wszystkich
$f,g\in L^\infty_{\mathrm{loc}}(\R)$;
\item $(S\star T)'=S'\star T+S\star T'$ dla wszystkich
$S,T\in\mathcal D'(\R)$.
\end{enumerate}
\end{stwierdzenie}
\begin{proof}
Niech $A=T_H$, gdzie $H$ jest funkcją Heaviside'a.
Ponieważ $H^2=H$ prawie wszędzie, pierwszy warunek
daje $A\star A=A$. Z przykładu~\ref{ex:pde7-heaviside}
mamy $A'=\delta_0$. Różniczkując równość
$A\star A=A$ za pomocą drugiego warunku i korzystając
z przemienności, otrzymujemy
\[
 2(A\star\delta_0)=\delta_0.
\]
Mnożymy tę równość przez $A$. Dwuliniowość i łączność
dają
\[
 A\star\delta_0
 =2A\star(A\star\delta_0)
 =2(A\star A)\star\delta_0
 =2A\star\delta_0.
\]
Stąd $A\star\delta_0=0$, a pierwsza równość
wymusza $\delta_0=0$. Jest to sprzeczne
z działaniem delty na test o wartości $1$ w zerze.
\end{proof}

Wynik ten jest konkretną wersją przeszkody Schwartza:
dotyczy dokładnie wymienionych założeń i nie wyklucza
wszystkich częściowo określonych iloczynów ani
innych teorii funkcji uogólnionych.
W szczególności zapis $\delta_0^2$ nie ma
standardowego znaczenia w rachunku dystrybucji.
Również wyrażenie $u\,\partial_xu$ dla nieciągłego
$u$ nie musi mieć sensu: dla $u=H$ wymagałoby
mnożenia delty przez funkcję nieciągłą w zerze.

Dla równań zachowania można ominąć tę trudność,
stosując \emph{postać dywergencyjną}. Dla funkcji
gładkiej reguła łańcuchowa daje
$u\partial_xu=\partial_x(u^2/2)$.
Jeśli rzeczywista funkcja $u$ należy do
$L^2_{\mathrm{loc}}(Q)$, gdzie $Q$ jest otwartym
podzbiorem czasoprzestrzeni $\R^2$, to
$u\in L^1_{\mathrm{loc}}(Q)$ z nierówności
Cauchy'ego--Schwarza na zbiorach zwartych,
a $u^2/2\in L^1_{\mathrm{loc}}(Q)$ z definicji.
Możemy więc nadać precyzyjny sens równaniu
\[
 \partial_tT_u+\partial_xT_{u^2/2}=0
 \quad\text{w }\mathcal D'(Q).
\]
Z definicji pochodnych dystrybucyjnych jest ono
równoważne warunkowi
\[
 \int_Q
 \left(u(t,x)\partial_t\varphi(t,x)
       +\frac{u(t,x)^2}{2}\partial_x\varphi(t,x)\right)
 \,\mathrm d\lambda_2(t,x)=0
 \qquad(\varphi\in\mathcal D(Q)).
\]
Nie wymaga on definiowania iloczynu
$u\,\partial_xT_u$. Sam warunek dystrybucyjny
nie zapewnia jednak jednoznaczności: warunki
entropijne omawiamy w teorii praw zachowania
(por. rozdział~\ref{chap:pde-characteristics}).

\begin{przyklad}[Skąd bierze się problem z $\delta^2$?]
Niech $\rho\in\mathcal D(\R)$ będzie nieujemna
i niech $\int_\R\rho\,\mathrm d\lambda_1=1$.
Dla $\varepsilon>0$ określamy
$\rho_\varepsilon(x)=\varepsilon^{-1}\rho(x/\varepsilon)$.
Wiemy z sekcji~\ref{sec:pde7-support}, że
$T_{\rho_\varepsilon}\to\delta_0$.
Wybieramy rzeczywisty, nieujemny test $\varphi$
taki, że $\varphi(0)>0$.
Podstawienie $x=\varepsilon y$ daje
\[
 \varepsilon\langle T_{\rho_\varepsilon^2},\varphi\rangle
 =\int_\R \rho(y)^2\varphi(\varepsilon y)\,
 \mathrm d\lambda_1(y).
\]
Ponieważ nośnik $\rho$ jest zwarty,
$\varphi(\varepsilon y)\to\varphi(0)$ jednostajnie
na nim. Stąd prawa strona dąży do
\[
 c=\varphi(0)\int_\R\rho(y)^2\,\mathrm d\lambda_1(y)>0.
\]
Dodatniość wynika z $\varphi(0)>0$ oraz z faktu,
że ciągła funkcja $\rho$ o całce $1$ nie jest zerowa.
Dla dostatecznie małego $\varepsilon$ otrzymujemy
zatem
$\langle T_{\rho_\varepsilon^2},\varphi\rangle
\ge c/(2\varepsilon)\to+\infty$.
Kwadraty tej aproksymacji delty nie mają więc
granicy w $\mathcal D'(\R)$, ponieważ zbieżność
dystrybucyjna wymaga skończonej granicy działania
na każdym ustalonym teście.
Jest to przeszkoda dla definiowania iloczynu
przez zwykłe przejście do granicy gładkich
aproksymacji, odrębna od powyższego argumentu
algebraicznego.
\end{przyklad}

\section[Podsumowanie: przygotowanie do analizy Fouriera]{Podsumowanie i przygotowanie do analizy Fouriera na $\R^d$}

Przestrzeń funkcji testowych ma topologię granicy
induktywnej, której zbieżność ciągów wymaga wspólnego
zwartego nośnika i jednostajnej zbieżności wszystkich
pochodnych (sekcja~\ref{sec:pde7-tests}).
Dystrybucje są ciągłymi funkcjonałami liniowymi,
opisanymi równoważnie przez lokalne oszacowania
~\eqref{eq:pde7-estimate}. Mają pochodne dowolnego
rzędu, ale pochodna słaba funkcji istnieje tylko
wtedy, gdy odpowiednia pochodna dystrybucyjna
jest regularna (sekcja~\ref{sec:pde7-weak}).

Lokalne wygładzanie dystrybucji daje funkcje gładkie
i komutuje z pochodnymi. Dla funkcji z $L^p_{\mathrm{loc}}$,
$1\le p<\infty$, przybliżone jedynki dają również
zbieżność w normie $L^p$ na zwartych podzbiorach,
natomiast w $L^\infty$ taka zbieżność nie zachodzi
na ogół (sekcja~\ref{sec:pde7-mollifiers}).

Rozwiązania fundamentalne operatorów o stałych
współczynnikach pozwalają przez splot z gładkim
wymuszeniem o zwartym nośniku konstruować rozwiązania
na całej przestrzeni (sekcja~\ref{sec:pde7-fundamental}).
Nie zapewnia to samo przez się jednoznaczności
ani spełnienia zadanych warunków brzegowych.
Lemat Weyla pokazuje, że równanie $\Delta T=0$
wymusza gładką i rzeczywiście analityczną
reprezentację dystrybucji
(twierdzenie~\ref{thm:pde7-weyl}).
Nie każdy iloczyn dystrybucji jest jednak określony;
w równaniach nieliniowych należy precyzyjnie
ustalić sens użytych wyrażeń.

W następnym rozdziale~\ref{chap:pde-fourier-Rd}
zbadamy transformatę Fouriera na przestrzeni Schwartza
i dystrybucjach temperowanych. Pozwoli to opisywać
symbole operatorów oraz wybrane jądra potencjałowe
w języku mnożników.

\chapter[Transformata Fouriera i dystrybucje temperowane]{Transformata Fouriera, przestrzeń Schwartza i dystrybucje temperowane}
\label{chap:pde-fourier-Rd}

Rozdział~\ref{chap:pde-distributions} pozwolił różniczkować funkcje
niegładkie i rozwiązywać równania o stałych współczynnikach za pomocą
rozwiązań fundamentalnych. Transformata Fouriera realizuje komplementarny
pomysł: zamienia pochodne na mnożenie przez wielomiany. Na całej
$\R^d$ nie możemy jednak bez zastrzeżeń używać klasycznej transformaty
całkowej dla wszystkich dystrybucji. Potrzebujemy więc przestrzeni
Schwartza $\mathcal S(\R^d)$ i jej przestrzeni sprzężonej $\mathcal S'(\R^d)$.

W całym rozdziale $d\ge1$, całki na $\R^d$ są liczone względem
miary Lebesgue'a $\lambda_d$, a wielowskaźniki należą do $\N_0^d$.
Dla funkcji całkowalnych przyjmujemy \emph{nieunitarną}
konwencję
\begin{equation}\label{eq:pde8-convention}
 \widehat f(\xi)=\mathcal F f(\xi)
 =\int_{\R^d}e^{-ix\cdot\xi}f(x)\,\mathrm d\lambda_d(x),\qquad
 \mathcal F^{-1}g(x)=(2\pi)^{-d}
 \int_{\R^d}e^{ix\cdot\xi}g(\xi)\,\mathrm d\lambda_d(\xi).
\end{equation}
Na razie symbol $\mathcal F^{-1}$ oznacza podaną całkę
dla $g\in L^1(\R^d)$; fakt, że jest to odwzorowanie
odwrotne do $\mathcal F$ na przestrzeni Schwartza,
wykażemy później.
Współczynniki $(2\pi)^{-d}$ zależą od konwencji i nie będą
pomijane. W szczególności nie będziemy bez przeliczenia używać
wzorów wyprowadzonych dla transformaty unitarnej.
Nie należy też utożsamiać tej transformaty całkowej na $\R^d$
z dyskretnymi współczynnikami szeregów Fouriera z
rozdziału~\ref{chap:pde-fourier}: te ostatnie są uśredniane
przez czynnik $(2\pi)^{-d}$.

\section{Przestrzeń Schwartza: po co wymagać szybkiego zaniku?}
\label{sec:pde8-schwartz}

\begin{definicja}[Przestrzeń Schwartza]
\label{def:pde8-schwartz}
Dla wielowskaźników $\alpha,\beta\in\N_0^d$ piszemy
$x^\alpha=\prod_{j=1}^d x_j^{\alpha_j}$ oraz
$\partial^\beta=\partial_1^{\beta_1}\cdots\partial_d^{\beta_d}$,
przyjmując $x_j^0=1$ także dla $x_j=0$.
Przestrzenią Schwartza na $\R^d$ nazywamy
\[
 \mathcal S(\R^d)=\bigl\{f\in C^\infty(\R^d;\C):
 p_{\alpha,\beta}(f)<\infty
 \text{ dla wszystkich }\alpha,\beta\in\N_0^d\bigr\},
\]
gdzie
\[
 p_{\alpha,\beta}(f)=\sup_{x\in\R^d}
 |x^\alpha\partial^\beta f(x)|.
\]
Wyposażamy ją w topologię wyznaczoną przez przeliczalną
rodzinę półnorm $p_{\alpha,\beta}$. Zbieżność $f_n\to f$
w $\mathcal S$ oznacza zbieżność do zera wszystkich
półnorm $p_{\alpha,\beta}(f_n-f)$.
\end{definicja}

Jest to zespolona przestrzeń wektorowa, ponieważ
liniowość pochodnych i nierówność trójkąta dają
$p_{\alpha,\beta}(af+bg)\le
|a|p_{\alpha,\beta}(f)+|b|p_{\alpha,\beta}(g)$.
Ta sama nierówność i jednorodność uzasadniają
własności półnorm. Rodzina rozdziela punkty,
gdyż $p_{0,0}(f)=0$ oznacza $f=0$.

\begin{lemat}[Równoważne półnormy i operacje]
\label{lem:pde8-seminorms}
Dla $N,m\in\N_0$ połóżmy
\[
 q_{N,m}(f)=\max_{|\beta|\le m}
 \sup_{x\in\R^d}(1+|x|)^N|\partial^\beta f(x)|.
\]
Rodzina $q_{N,m}$ określa tę samą topologię co
$p_{\alpha,\beta}$. Mnożenie przez $x_j$ oraz
różniczkowanie $\partial_j$ są ciągłymi operatorami
liniowymi $\mathcal S\to\mathcal S$.
Ponadto $\mathcal S(\R^d)\subset L^1(\R^d)\cap L^2(\R^d)$,
a oba włożenia są ciągłe.
\end{lemat}
\begin{proof}
Ponieważ $|x^\alpha|\le(1+|x|)^{|\alpha|}$, mamy
$p_{\alpha,\beta}(f)\le q_{|\alpha|,|\beta|}(f)$.
Odwrotnie, z $|x|\le\sum_{j=1}^d|x_j|$ i wzoru
wielomianowego dostajemy
\[
 (1+|x|)^N\le
 \sum_{|\alpha|\le N}
 \frac{N!}{(N-|\alpha|)!\alpha!}|x^\alpha|,
 \qquad \alpha!=\prod_{j=1}^d\alpha_j!.
\]
Mnożąc przez $|\partial^\beta f(x)|$, biorąc
supremum względem $x$ i maksimum po $|\beta|\le m$,
ograniczamy $q_{N,m}$ przez skończoną sumę półnorm
$p_{\alpha,\beta}$ z podanymi współczynnikami.
Obie rodziny wyznaczają więc tę samą topologię.

Z reguły Leibniza, dla wektora jednostkowego
$e_j$, mamy
\[
 \partial^\beta(x_j f)
 =x_j\partial^\beta f+\beta_j\partial^{\beta-e_j}f,
\]
gdzie drugi składnik pomijamy dla $\beta_j=0$.
Stąd
\[
 \begin{aligned}
 q_{N,m}(x_j f)&\le q_{N+1,m}(f)+m q_{N,m}(f),\\
 q_{N,m}(\partial_j f)&\le q_{N,m+1}(f).
 \end{aligned}
\]
Skończoność prawych stron zapewnia, że obrazy
należą do $\mathcal S$, a oszacowania dla każdej
półnormy dają ciągłość.

Dla $p=1,2$ wybieramy $N$ tak, aby $Np>d$.
Wtedy
\[
 \|f\|_{L^p}\le C_{d,N,p}q_{N,0}(f),
 \qquad
 C_{d,N,p}^p=\int_{\R^d}(1+|x|)^{-Np}
 \,\mathrm d\lambda_d(x)<\infty.
\]
Istotnie, na kuli jednostkowej całka jest skończona,
a poza nią całkowanie radialne sprowadza oszacowanie
do stałej razy $\int_1^\infty r^{d-1-Np}\,
\mathrm d\lambda_1(r)$, która jest skończona dla $Np>d$.
W wymiarze jeden całkujemy osobno na obu półprostych.
Daje to oba włożenia i ich ciągłość.
\end{proof}

Równoważne półnormy wyjaśniają szybki zanik:
dla każdego $N,\beta$ mamy
\[
 (1+|x|)^N|\partial^\beta f(x)|
 \le\frac{q_{N+1,|\beta|}(f)}{1+|x|}
 \longrightarrow 0\qquad (|x|\to\infty).
\]
Funkcja $e^{-|x|^2}$ należy do $\mathcal S$:
każda jej pochodna ma postać
$P_\beta(x)e^{-|x|^2}$ z wielomianem $P_\beta$.
Indukcja zaczyna się od $P_0=1$, a różniczkowanie
daje $P_{\beta+e_j}=\partial_j P_\beta-2x_jP_\beta$.
Dla każdej potęgi $k\ge0$ funkcja
$r^k e^{-r^2}$ jest ograniczona na $[0,\infty)$
i dąży do zera w nieskończoności (wykładnik
dominuje potęgi, jak w konstrukcji funkcji
testowej z sekcji~\ref{sec:pde7-tests}).
Każda funkcja z $C_c^\infty(\R^d)$ również należy
do $\mathcal S$, bo jej pochodne znikają poza
wspólnym zwartym nośnikiem.
Natomiast $(1+|x|^2)^{-1}$ nie należy do $\mathcal S$
w żadnym wymiarze $d\ge1$: dla
$x=te_1$ mamy $t^3/(1+t^2)\to+\infty$
przy $t\to+\infty$, więc $p_{3e_1,0}$ jest nieskończona.

\begin{stwierdzenie}[Zupełność przestrzeni Schwartza]
Przestrzeń $\mathcal S(\R^d)$ jest przestrzenią Frécheta,
czyli zupełną metryzowalną przestrzenią lokalnie wypukłą.
\end{stwierdzenie}
\begin{proof}
Rosnąca rodzina $Q_n=q_{n,n}$, $n\in\N_0$, wyznacza
tę samą topologię, ponieważ każda $q_{N,m}$ jest
ograniczona przez $Q_{\max\{N,m\}}$.
Metryka niezmiennicza względem przesunięć
\[
 D(f,g)=\sum_{n=0}^{\infty}2^{-n-1}
 \min\{1,Q_n(f-g)\}
\]
wyznacza tę topologię. Istotnie, nierówność
trójkąta wynika z półnorm i podaddytywności
$\min\{1,t\}$, a $Q_0$ rozdziela punkty.
Dla $0<\delta<1$ warunek
$D(f,g)<2^{-n-1}\delta$ daje $Q_n(f-g)<\delta$.
Odwrotnie, wybieramy $M$ z ogonem szeregu wag
mniejszym niż $\delta/2$; gdy $Q_M(f-g)<\delta/2$,
pierwsze $M+1$ składników mają sumę mniejszą
niż $\delta/2$, więc $D(f,g)<\delta$.
Te oszacowania charakteryzują również ciągi Cauchy'ego.

Niech $(f_\ell)$ będzie ciągiem Cauchy'ego w tej
metryce. Dla każdego $\beta$ ciąg
$\partial^\beta f_\ell$ spełnia jednostajny
warunek Cauchy'ego na $\R^d$ dzięki $Q_{|\beta|}$.
Ma zatem jednostajną granicę $g_\beta$, ciągłą
na $\R^d$. Dla każdego $x\in\R^d$, $t\in\R$
i $j\in\{1,\ldots,d\}$ mamy
\[
 \partial^\beta f_\ell(x+te_j)-\partial^\beta f_\ell(x)
 =\int_0^t\partial^{\beta+e_j}f_\ell(x+se_j)
 \,\mathrm d\lambda_1(s).
\]
Przechodzimy do granicy dzięki jednostajnej zbieżności
na odcinku (dla $t<0$ całka jest zorientowana).
Ciągłość $g_{\beta+e_j}$ daje
$\partial_j g_\beta=g_{\beta+e_j}$.
Indukcyjnie $f:=g_0$ jest gładka i
$\partial^\beta f=g_\beta$.

Ustalmy $n$ i $\varepsilon>0$. Dla
$\ell,k$ od pewnego indeksu mamy
$Q_n(f_\ell-f_k)<\varepsilon$. Dla każdego
$x$ i $|\beta|\le n$ przechodzimy z $k$
do nieskończoności w nierówności
\[
 (1+|x|)^n|\partial^\beta f_\ell(x)
              -\partial^\beta f_k(x)|<\varepsilon.
\]
Biorąc supremum i maksimum po przejściu do
granicy, otrzymujemy $Q_n(f_\ell-f)\le\varepsilon$.
W szczególności $Q_n(f)\le Q_n(f_\ell)+\varepsilon$
jest skończona. Ponieważ $n$ było dowolne,
$f\in\mathcal S$ i $f_\ell\to f$ we wszystkich
półnormach, a więc i w metryce $D$.
\end{proof}

\begin{twierdzenie}[Gęstość funkcji o zwartym nośniku]
\label{thm:pde8-density}
Przestrzeń $C_c^\infty(\R^d)$ jest gęsta w $\mathcal S(\R^d)$.
Jest też gęsta w $L^p(\R^d)$ dla $1\le p<\infty$.
\end{twierdzenie}
\begin{proof}
Wybierzmy funkcję odcinającą
$\psi\in C_c^\infty(\R^d)$ z $0\le\psi\le1$,
$\psi=1$ na $B_1(0)$ i $\psi=0$ poza $B_2(0)$.
Można przyjąć $\psi=\chi_{B_{3/2}(0)}*\rho_{1/4}$,
gdzie $\rho_\varepsilon$ jest przybliżoną jedynką
z sekcji~\ref{sec:pde7-mollifiers}.
Gładkość wynika z twierdzenia~\ref{thm:pde7-mollify}.
Nieujemność jądra i jego całka $1$ dają
$0\le\psi\le1$; dla $|x|\le1$ i $|y|\le1/4$
mamy $|x-y|<3/2$, a dla $|x|\ge2$ mamy
$|x-y|>3/2$. Stąd żądane wartości i zwarty nośnik.

Dla $f\in\mathcal S$ i $R\ge1$ określamy
$f_R(x)=\psi(x/R)f(x)\in C_c^\infty(\R^d)$.
Niech $\psi_R(x)=\psi(x/R)$.
Wszystkie pochodne $1-\psi_R$, włącznie z rzędem
zero, znikają dla $|x|<R$. Dla $|\gamma|>0$
ich moduły są ograniczone przez
$R^{-|\gamma|}\|\partial^\gamma\psi\|_\infty$;
dla $\gamma=0$ przez $1$.
Stosujemy wielowskaźnikową regułę Leibniza:
\[
 \partial^\beta(f-f_R)
 =\sum_{\gamma\le\beta}\binom{\beta}{\gamma}
 \partial^\gamma(1-\psi_R)\partial^{\beta-\gamma}f,
\]
gdzie $\gamma\le\beta$ oznacza nierówności
we wszystkich współrzędnych, a
$\binom{\beta}{\gamma}=\prod_j\binom{\beta_j}{\gamma_j}$.
Dla $|x|\ge R$ zachodzi
$(1+|x|)^N\le (1+R)^{-1}(1+|x|)^{N+1}$.
Zatem dla $|\beta|\le m$ każdy składnik po
pomnożeniu przez $(1+|x|)^N$ jest ograniczony
przez odpowiednią stałą zależną od $\psi,\beta,\gamma$
razy $(1+R)^{-1}q_{N+1,m}(f)$.
Sumując skończenie wiele składników, dostajemy
\[
 q_{N,m}(f-f_R)\le
 \frac{C_{\psi,m}}{1+R}q_{N+1,m}(f)
 \longrightarrow0.
\]
To daje gęstość w $\mathcal S$.

Dla $u\in L^p(\R^d)$, $1\le p<\infty$,
niech $v_R=u\chi_{B_R(0)}$. Z całkowalności
$|u|^p$ wynika
\[
 \|u-v_R\|_{L^p}^p
 =\int_{\R^d\setminus B_R(0)}|u(x)|^p
 \,\mathrm d\lambda_d(x)\longrightarrow0.
\]
Dla ustalonego $R$ funkcja
$v_R*\rho_\varepsilon$ jest gładka na $\R^d$
i dla $0<\varepsilon<1$ znika poza
$\overline{B_{R+1}(0)}$: jeśli $|x|>R+1$,
to dla $|y|\le R$ mamy $|x-y|>1>\varepsilon$.
Z twierdzenia~\ref{thm:pde7-approx-id}
na zwartym zbiorze $\overline{B_{R+1}(0)}$
wynika
\[
 \|v_R*\rho_\varepsilon-v_R\|_{L^p(\R^d)}
 \longrightarrow0\qquad(\varepsilon\downarrow0),
\]
ponieważ obie funkcje znikają poza tym zbiorem.
Dla dowolnego $\delta>0$ wybieramy najpierw
$R$ z $\|u-v_R\|_p<\delta/2$, następnie
$\varepsilon$ z błędem wygładzania mniejszym
niż $\delta/2$. Nierówność trójkąta daje
aproksymację $u$ z błędem mniejszym niż $\delta$
przez funkcję gładką o zwartym nośniku.
\end{proof}

\section{Transformata Fouriera na przestrzeni Schwartza}
\label{sec:pde8-transform}

\begin{twierdzenie}[Różniczkowanie i mnożenie]
\label{thm:pde8-rules}
Jeżeli $f\in\mathcal S(\R^d)$, to $\widehat f\in\mathcal S(\R^d)$
oraz dla każdego $\alpha\in\N_0^d$
\begin{equation}\label{eq:pde8-rules}
 \mathcal F(\partial^\alpha f)(\xi)=(i\xi)^\alpha\widehat f(\xi),
 \qquad \mathcal F(x^\alpha f)(\xi)=i^{|\alpha|}
 \partial_\xi^\alpha\widehat f(\xi).
\end{equation}
Ponadto $\mathcal F:\mathcal S(\R^d)\to\mathcal S(\R^d)$
jest ciągłym operatorem liniowym.
\end{twierdzenie}
\begin{proof}
Z lematu~\ref{lem:pde8-seminorms} wynika, że
$x^\beta f$ i $\partial^\beta f$ należą do
$\mathcal S(\R^d)$, a więc do $L^1(\R^d)$,
dla każdego wielowskaźnika $\beta$.
W szczególności całka definiująca $\widehat f$
jest bezwzględnie zbieżna.

Uzasadnijmy najpierw różniczkowanie względem $\xi$.
Dla $t\in\R$ mamy
\[
 |e^{-it}-1|
 =\left|\int_0^t(-i)e^{-is}\,\mathrm d\lambda_1(s)\right|
 \le |t|.
\]
Dla $h\ne0$ iloraz różnicowy
\[
 \frac{\widehat f(\xi+he_j)-\widehat f(\xi)}h
 =\int_{\R^d}e^{-ix\cdot\xi}
 \frac{e^{-ihx_j}-1}{h}f(x)\,\mathrm d\lambda_d(x)
\]
ma funkcję podcałkową ograniczoną co do modułu
przez całkowalną funkcję $|x_j f(x)|$.
Z twierdzenia o zbieżności zmajoryzowanej
otrzymujemy
\[
 \partial_{\xi_j}\widehat f(\xi)
 =\int_{\R^d}(-ix_j)e^{-ix\cdot\xi}f(x)
 \,\mathrm d\lambda_d(x).
\]
Ponieważ $x^\beta f\in L^1$ dla każdego $\beta$,
ten sam argument można zastosować do kolejnych
pochodnych. Indukcyjnie
\[
 \partial_\xi^\beta\widehat f(\xi)
 =(-i)^{|\beta|}\mathcal F(x^\beta f)(\xi).
\]
Każda z tych pochodnych jest ciągła: dla
$\xi_n\to\xi$ stosujemy twierdzenie o zbieżności
zmajoryzowanej z majorantą $|x^\beta f(x)|$.
Zatem $\widehat f\in C^\infty(\R^d;\C)$.
Mnożąc powyższą tożsamość przez $i^{|\beta|}$,
uzyskujemy drugą równość w~\eqref{eq:pde8-rules}.

Dla pierwszej równości ustalamy $j$ i pozostałe
współrzędne $x$. Całkowanie przez części względem
$x_j$ na przedziale $[-R,R]$ daje
\[
 \begin{aligned}
 \int_{-R}^R e^{-ix\cdot\xi}\partial_j f(x)
 \,\mathrm d\lambda_1(x_j)
 &=\left[e^{-ix\cdot\xi}f(x)\right]_{x_j=-R}^{x_j=R}\\
 &\quad+i\xi_j\int_{-R}^R e^{-ix\cdot\xi}f(x)
 \,\mathrm d\lambda_1(x_j).
 \end{aligned}
\]
Dla ustalonych pozostałych współrzędnych wyraz
brzegowy dąży do zera, ponieważ
$|f(x)|\le q_{1,0}(f)/(1+|x|)$.
Całki po prostej są bezwzględnie zbieżne dzięki
oszacowaniom $|f(x)|,|\partial_j f(x)|
\le q_{2,1}(f)(1+|x_j|)^{-2}$.
Po przejściu do granicy $R\to\infty$ całkujemy
względem pozostałych współrzędnych, stosując
twierdzenie Fubiniego: $f,\partial_j f\in L^1(\R^d)$.
Dla $d=1$ nie ma dodatkowego całkowania.
Otrzymujemy $\mathcal F(\partial_j f)=i\xi_j\widehat f$.
Ponieważ każda pochodna $f$ należy do $\mathcal S$,
iteracja daje pierwszą równość w~\eqref{eq:pde8-rules}.

Pozostaje wykazać skończoność półnorm transformaty
i jej ciągłość. Łącząc obie uzyskane tożsamości,
otrzymujemy dla dowolnych $\alpha,\beta\in\N_0^d$
\[
 \xi^\alpha\partial_\xi^\beta\widehat f(\xi)
 =i^{-|\alpha|}(-i)^{|\beta|}
 \mathcal F\bigl(\partial^\alpha(x^\beta f)\bigr)(\xi).
\]
Dla $v\in L^1(\R^d)$ zachodzi
$|\widehat v(\xi)|\le\|v\|_{L^1}$, ponieważ moduł
czynnika wykładniczego wynosi $1$.
Ustalmy całkowite $N>d$ i oznaczmy przez
$C_{d,N,1}$ stałą z lematu~\ref{lem:pde8-seminorms}.
Z wielowskaźnikowej reguły Leibniza wynika
\[
 \begin{aligned}
 p_{\alpha,\beta}(\widehat f)
 &\le \|\partial^\alpha(x^\beta f)\|_{L^1}\\
 &\le C_{d,N,1}
 \sum_{\substack{\gamma\le\alpha\\\gamma\le\beta}}
 \binom{\alpha}{\gamma}
 \frac{\beta!}{(\beta-\gamma)!}
 q_{N+|\beta|-|\gamma|,\,|\alpha-\gamma|}(f).
 \end{aligned}
\]
Tutaj pochodna $\partial^\gamma x^\beta$
jest zerowa, jeśli $\gamma\not\le\beta$;
w przeciwnym razie jest równa
$\beta!x^{\beta-\gamma}/(\beta-\gamma)!$.
Każda półnorma po prawej stronie jest
ograniczona przez $q_{N+|\beta|,|\alpha|}(f)$.
W konsekwencji istnieje skończona stała
$C_{\alpha,\beta,d,N}$, niezależna od $f$, taka że
\[
 p_{\alpha,\beta}(\widehat f)
 \le C_{\alpha,\beta,d,N}q_{N+|\beta|,|\alpha|}(f).
\]
Wszystkie półnormy $\widehat f$ są zatem skończone,
czyli $\widehat f\in\mathcal S$. Liniowość wynika
z liniowości całki, a uzyskane oszacowania
każdej półnormy dają ciągłość operatora.
\end{proof}

\begin{stwierdzenie}[Lemat Riemanna--Lebesgue'a]
\label{prop:pde8-RL}
Jeśli $f\in L^1(\R^d)$, to $\widehat f$ jest ograniczona,
jednostajnie ciągła i $\widehat f(\xi)\to0$
przy $|\xi|\to\infty$.
\end{stwierdzenie}
\begin{proof}
Całka definiująca $\widehat f$ nie zależy od wyboru
reprezentanta klasy $f$ i daje
$|\widehat f(\xi)|\le\|f\|_{L^1}$.
Dla $\xi,h\in\R^d$ mamy
\[
 |\widehat f(\xi+h)-\widehat f(\xi)|
 \le\int_{\R^d}|e^{-ix\cdot h}-1|\,|f(x)|
 \,\mathrm d\lambda_d(x).
\]
Prawa strona nie zależy od $\xi$ i dąży do zera
dla $h\to0$: funkcja podcałkowa dąży punktowo
do zera i jest ograniczona przez $2|f(x)|$.
Twierdzenie o zbieżności zmajoryzowanej daje
więc jednostajną ciągłość.

Ustalmy $\varepsilon>0$. Z twierdzenia
~\ref{thm:pde8-density} wybieramy
$g\in C_c^\infty(\R^d)$ z
$\|f-g\|_{L^1}<\varepsilon/2$.
Ponieważ $g\in\mathcal S(\R^d)$, poprzednie
twierdzenie daje $\widehat g\in\mathcal S(\R^d)$.
Istnieje zatem $R>0$ takie, że
$|\widehat g(\xi)|<\varepsilon/2$ dla $|\xi|>R$.
Dla tych $\xi$ z liniowości i oszacowania
transformaty funkcji całkowalnej wynika
\[
 |\widehat f(\xi)|
 \le\|f-g\|_{L^1}+|\widehat g(\xi)|<\varepsilon.
\]
Dowodzi to zaniku w nieskończoności.
\end{proof}

\section{Transformata funkcji Gaussa i wzór odwrócenia}
\label{sec:pde8-inversion}

\begin{lemat}[Transformata funkcji Gaussa]
\label{lem:pde8-gauss}
Dla $a>0$ i $\xi\in\R^d$ zachodzi
\begin{equation}\label{eq:pde8-gauss}
 \mathcal F(e^{-a|x|^2})(\xi)
 =\left(\frac\pi a\right)^{d/2}e^{-|\xi|^2/(4a)}.
\end{equation}
\end{lemat}
\begin{proof}
W wymiarze jeden niech $g_a(x)=e^{-ax^2}$.
Funkcja ta należy do $\mathcal S(\R)$:
argument z sekcji~\ref{sec:pde8-schwartz} stosuje się
z rekurencją wielomianów
$P_{k+1}(x)=P_k'(x)-2axP_k(x)$.
Niech $I=\widehat{g_a}$. Ponieważ
$g_a'=-2axg_a$, tożsamości~\eqref{eq:pde8-rules}
dają
\[
 i\xi I(\xi)=\widehat{g_a'}(\xi)
 =-2a\mathcal F(xg_a)(\xi)=-2aiI'(\xi).
\]
Stąd $I'(\xi)=-\xi I(\xi)/(2a)$ i
\[
 \frac{\mathrm d}{\mathrm d\xi}
 \left(e^{\xi^2/(4a)}I(\xi)\right)=0.
\]
Zatem $I(\xi)=I(0)e^{-\xi^2/(4a)}$.

Dla wyznaczenia $I(0)$ oznaczmy
$J=\int_\R e^{-t^2}\,\mathrm d\lambda_1(t)$.
Jest to liczba dodatnia i skończona, ponieważ
dla $|t|\ge1$ mamy $e^{-t^2}\le e^{-|t|}$.
Z twierdzenia Tonellego oraz całkowania w
zmiennych biegunowych na $\R^2$ otrzymujemy
\[
 J^2=\int_{\R^2}e^{-(x^2+y^2)}\,\mathrm d\lambda_2(x,y)
 =2\pi\int_0^\infty e^{-r^2}r\,\mathrm d\lambda_1(r)
 =\pi.
\]
Ponieważ $J>0$, mamy $J=\sqrt\pi$.
Podstawienie $t=\sqrt a\,x$ daje
$I(0)=a^{-1/2}J=\sqrt{\pi/a}$.
Nie korzystaliśmy tu ze wzoru odwrócenia.

W wymiarze $d$ funkcja podcałkowa jest
iloczynem czynników jednowymiarowych, a całka
z jej modułu jest skończona i wynosi
$(\pi/a)^{d/2}$ (Tonelli). Z twierdzenia
Fubiniego wynika
\[
 \begin{aligned}
 \mathcal F(e^{-a|x|^2})(\xi)
 &=\prod_{j=1}^d
 \int_\R e^{-ax_j^2}e^{-ix_j\xi_j}\,
 \mathrm d\lambda_1(x_j)\\
 &=\left(\frac{\pi}{a}\right)^{d/2}
 e^{-|\xi|^2/(4a)}.
 \end{aligned}
\]
\end{proof}

\begin{twierdzenie}[Odwrócenie transformaty]
\label{thm:pde8-inversion}
Dla $f\in\mathcal S(\R^d)$ zachodzi
\[
 f(x)=(2\pi)^{-d}\int_{\R^d}e^{ix\cdot\xi}
 \widehat f(\xi)\,\mathrm d\lambda_d(\xi),
 \qquad x\in\R^d.
\]
Zatem $\mathcal F$ jest topologicznym automorfizmem
$\mathcal S(\R^d)$, którego odwrotność podano w
\eqref{eq:pde8-convention}.
\end{twierdzenie}
\begin{proof}
Dla $\varepsilon>0$ określmy
\[
 I_\varepsilon(x)=(2\pi)^{-d}\int_{\R^d}
 e^{ix\cdot\xi}e^{-\varepsilon|\xi|^2/2}
 \widehat f(\xi)\,\mathrm d\lambda_d(\xi).
\]
Po wstawieniu definicji $\widehat f$ całka
podwójna jest bezwzględnie zbieżna, ponieważ
całka z modułu (bez stałej $(2\pi)^{-d}$) wynosi
\[
 \|f\|_{L^1}
 \int_{\R^d}e^{-\varepsilon|\xi|^2/2}
 \,\mathrm d\lambda_d(\xi)
 =\|f\|_{L^1}(2\pi/\varepsilon)^{d/2}<\infty.
\]
Stosując Fubiniego i lemat~\ref{lem:pde8-gauss}
z $a=\varepsilon/2$, otrzymujemy
\[
 \begin{aligned}
 I_\varepsilon(x)
 &=\int_{\R^d}f(y)G_\varepsilon(x-y)
 \,\mathrm d\lambda_d(y),\\
 G_\varepsilon(z)
 &:=(2\pi)^{-d}\int_{\R^d}e^{iz\cdot\xi}
 e^{-\varepsilon|\xi|^2/2}\,\mathrm d\lambda_d(\xi)\\
 &=(2\pi\varepsilon)^{-d/2}
 e^{-|z|^2/(2\varepsilon)}.
 \end{aligned}
\]
Jądro jest nieujemne i ma całkę $1$, co wynika
z obliczenia całki Gaussa. Nie ma jednak zwartego
nośnika, dlatego uzasadnimy jego własność
przybliżonej jedynki bezpośrednio.

Po podstawieniu $y=x-\sqrt\varepsilon\,z$ mamy
\[
 I_\varepsilon(x)-f(x)
 =\int_{\R^d}G_1(z)
 \bigl(f(x-\sqrt\varepsilon\,z)-f(x)\bigr)
 \,\mathrm d\lambda_d(z).
\]
Niech $M=\sup_{x\in\R^d}|\nabla f(x)|<\infty$,
gdzie dla gradientu zespolonego stosujemy
normę euklidesową. Całkowanie wzdłuż odcinka
daje $|f(x-h)-f(x)|\le M|h|$ dla $h\in\R^d$.
W konsekwencji
\[
 \sup_x|I_\varepsilon(x)-f(x)|
 \le M\sqrt\varepsilon
 \int_{\R^d}|z|G_1(z)\,\mathrm d\lambda_d(z)
 \longrightarrow0.
\]
Całka z pierwszego momentu jest skończona
dzięki wykładniczemu zanikowi Gaussa
(w zmiennych biegunowych wystarcza całkowalność
$r^d e^{-r^2/2}$ na $(0,\infty)$; dla $d=1$
całkujemy na obu półprostych).

Z drugiej strony $\widehat f\in L^1$ na mocy
twierdzenia~\ref{thm:pde8-rules} i
lematu~\ref{lem:pde8-seminorms}. Mamy
\[
 \begin{aligned}
 &\sup_x\left| I_\varepsilon(x)
 -(2\pi)^{-d}\int_{\R^d}e^{ix\cdot\xi}
 \widehat f(\xi)\,\mathrm d\lambda_d(\xi)\right|\\
 &\quad\le (2\pi)^{-d}\int_{\R^d}
 (1-e^{-\varepsilon|\xi|^2/2})|\widehat f(\xi)|
 \,\mathrm d\lambda_d(\xi)\longrightarrow0
 \end{aligned}
\]
z twierdzenia o zbieżności zmajoryzowanej
z majorantą $(2\pi)^{-d}|\widehat f(\xi)|$.
Obie granice są więc równe, co dowodzi wzoru.

Sprawdźmy jeszcze odwracalność w obu kierunkach.
Niech $\mathcal R f(x)=f(-x)$. Operator
$\mathcal R$ zachowuje każdą półnormę
$p_{\alpha,\beta}$, spełnia $\mathcal R^2=I$
i komutuje z $\mathcal F$:
podstawienie $y=-x$ daje
$\mathcal F(\mathcal R f)(\xi)=\widehat f(-\xi)$.
Operator określony całką odwrotną to
$\mathcal G=(2\pi)^{-d}\mathcal R\mathcal F$,
więc odwzorowuje $\mathcal S$ ciągle w siebie.
Udowodniony wzór mówi $\mathcal G\mathcal F=I$,
a komutowanie $\mathcal R$ i $\mathcal F$ daje
$\mathcal F\mathcal G=\mathcal G\mathcal F=I$.
Zatem $\mathcal G=\mathcal F^{-1}$ i oba
operatory są ciągłe. Otrzymujemy ponadto
$\mathcal F^2=(2\pi)^d\mathcal R$.
\end{proof}

\begin{wniosek}[Wzory na splot i iloczyn]
\label{cor:pde8-convolution}
Dla $f,g\in\mathcal S(\R^d)$ funkcje $f*g$ oraz
$fg$ należą do $\mathcal S(\R^d)$ i
\[
 \widehat{f*g}=\widehat f\,\widehat g,\qquad
 \widehat{fg}=(2\pi)^{-d}(\widehat f*\widehat g).
\]
\end{wniosek}
\begin{proof}
Przyjmujemy $(f*g)(x)=\int_{\R^d}f(y)g(x-y)
\,\mathrm d\lambda_d(y)$. Całka jest bezwzględnie
zbieżna dla każdego $x$, ponieważ $f\in L^1$,
a $g$ jest ograniczona.
Różniczkowanie pod całką daje
$\partial^\beta(f*g)=f*(\partial^\beta g)$:
ilorazy różnicowe są zmajoryzowane przez
$|f(y)|$ razy normę supremum odpowiedniej
pochodnej $g$ wyższego rzędu, a ciągłość każdej
pochodnej wynika z tej samej metody z majorantą
$|f(y)|\|\partial^\beta g\|_\infty$.
Z nierówności
$1+|x|\le(1+|y|)(1+|x-y|)$ wynika
\[
 q_{N,m}(f*g)
 \le q_{N,m}(g)
 \int_{\R^d}(1+|y|)^N|f(y)|\,\mathrm d\lambda_d(y).
\]
Ostatnia całka jest skończona, gdyż dla
całkowitego $L>d$ funkcja podcałkowa jest
ograniczona przez $q_{N+L,0}(f)(1+|y|)^{-L}$.
Zatem $f*g\in\mathcal S$.
Z reguły Leibniza, rozdzielając wagę
$(1+|x|)^N$ na pierwszy czynnik,
otrzymujemy także
$q_{N,m}(fg)\le 2^m q_{N,m}(f)q_{0,m}(g)$,
więc $fg\in\mathcal S$.

Całka podwójna z modułu $f(y)g(x-y)$
wynosi $\|f\|_{L^1}\|g\|_{L^1}$ dzięki
niezmienniczości miary względem przesunięć
i twierdzeniu Tonellego. Zatem Fubini
i podstawienie $z=x-y$ dają
\[
 \begin{aligned}
 \widehat{f*g}(\xi)
 &=\int_{\R^d}f(y)
 \int_{\R^d}e^{-ix\cdot\xi}g(x-y)
 \,\mathrm d\lambda_d(x)\,\mathrm d\lambda_d(y)\\
 &=\int_{\R^d}e^{-iy\cdot\xi}f(y)
 \widehat g(\xi)\,\mathrm d\lambda_d(y)
 =\widehat f(\xi)\widehat g(\xi).
 \end{aligned}
\]
Dla drugiej tożsamości podstawiamy wzór
odwrócenia dla $g$. Bezwzględną całkowalność
całki podwójnej zapewnia
$\|f\|_{L^1}\|\widehat g\|_{L^1}<\infty$.
Stąd
\[
 \begin{aligned}
 \widehat{fg}(\xi)
 &=(2\pi)^{-d}\int_{\R^d}\widehat g(\eta)
 \int_{\R^d}e^{-ix\cdot(\xi-\eta)}f(x)
 \,\mathrm d\lambda_d(x)\,\mathrm d\lambda_d(\eta)\\
 &=(2\pi)^{-d}\int_{\R^d}
 \widehat f(\xi-\eta)\widehat g(\eta)
 \,\mathrm d\lambda_d(\eta)
 =(2\pi)^{-d}(\widehat f*\widehat g)(\xi).
 \end{aligned}
\]
W przyjętej konwencji czynnik $(2\pi)^{-d}$
występuje we wzorze na transformatę iloczynu,
a nie splotu.
\end{proof}

\section[Twierdzenie Plancherela i przestrzeń L2]{Twierdzenie Plancherela i przestrzeń $L^2$}
\label{sec:pde8-plancherel}

\begin{twierdzenie}[Tożsamość Parsevala--Plancherela]
\label{thm:pde8-plancherel}
Dla $f,g\in\mathcal S(\R^d)$ zachodzi
\begin{equation}\label{eq:pde8-parseval}
 \int_{\R^d}f(x)\overline{g(x)}\,\mathrm d\lambda_d(x)
 =(2\pi)^{-d}\int_{\R^d}\widehat f(\xi)
 \overline{\widehat g(\xi)}\,\mathrm d\lambda_d(\xi).
\end{equation}
Transformata Fouriera przedłuża się jednoznacznie do
ciągłego operatora liniowego
$\mathcal F:L^2(\R^d)\to L^2(\R^d)$, dla którego
$\|\widehat f\|_2=(2\pi)^{d/2}\|f\|_2$.
Operator $(2\pi)^{-d/2}\mathcal F$ jest unitarny,
a tożsamość~\eqref{eq:pde8-parseval} zachodzi także
dla dowolnych $f,g\in L^2(\R^d)$.
\end{twierdzenie}
\begin{proof}
Ze wzoru odwrócenia z twierdzenia
~\ref{thm:pde8-inversion} mamy
\[
 \overline{g(x)}=(2\pi)^{-d}\int_{\R^d}
 e^{-ix\cdot\xi}\overline{\widehat g(\xi)}
 \,\mathrm d\lambda_d(\xi).
\]
Po pomnożeniu przez $f(x)$ całka podwójna
jest bezwzględnie zbieżna, ponieważ całka
z modułu, bez czynnika $(2\pi)^{-d}$, wynosi
$\|f\|_1\|\widehat g\|_1<\infty$.
Stąd z twierdzenia Fubiniego
\[
 \begin{aligned}
 \int_{\R^d}f(x)\overline{g(x)}\,\mathrm d\lambda_d(x)
 &=(2\pi)^{-d}\int_{\R^d}
 \overline{\widehat g(\xi)}
 \int_{\R^d}e^{-ix\cdot\xi}f(x)
 \,\mathrm d\lambda_d(x)\,\mathrm d\lambda_d(\xi)\\
 &=(2\pi)^{-d}\int_{\R^d}
 \widehat f(\xi)\overline{\widehat g(\xi)}
 \,\mathrm d\lambda_d(\xi).
 \end{aligned}
\]
W szczególności dla $f=g$ otrzymujemy
$\|\widehat f\|_2^2=(2\pi)^d\|f\|_2^2$.
W zespolonej przestrzeni Hilberta $L^2(\R^d)$
przyjmujemy iloczyn skalarny liniowy względem
pierwszego argumentu,
$\langle f,g\rangle_2=\int_{\R^d}f(x)\overline{g(x)}
\,\mathrm d\lambda_d(x)$.

Oznaczmy przez $U_0=(2\pi)^{-d/2}\mathcal F$
operator na $\mathcal S$. Jest on liniową izometrią
w normie $L^2$. Dla dowolnego $f\in L^2$
wybieramy ciąg $f_n\in\mathcal S$ z
$\|f_n-f\|_2\to0$, korzystając z twierdzenia
~\ref{thm:pde8-density}. Wówczas
\[
 \|U_0f_n-U_0f_k\|_2=\|f_n-f_k\|_2.
\]
Ciąg $(U_0f_n)$ jest ciągiem Cauchy'ego,
więc ma granicę w zupełnej przestrzeni $L^2$.
Definiujemy $Uf=\lim_n U_0f_n$ w $L^2$.
Jeśli $g_n\in\mathcal S$ jest innym ciągiem
aproksymującym $f$, to
$\|U_0f_n-U_0g_n\|_2=\|f_n-g_n\|_2
\le\|f_n-f\|_2+\|f-g_n\|_2\to0$.
Definicja nie zależy więc od wyboru ciągu.

Liniowość $U$ wynika z liniowości $U_0$:
dla kombinacji liniowej można użyć odpowiedniej
kombinacji ciągów aproksymujących. Ciągłość normy
daje $\|Uf\|_2=\lim_n\|f_n\|_2=\|f\|_2$.
Operator $U$ przedłuża $U_0$ (dla $f\in\mathcal S$
można wziąć ciąg stały), a jednoznaczność
ciągłego przedłużenia wynika z gęstości $\mathcal S$.

Obraz $U(L^2)$ jest domknięty: jeśli $Uf_n\to h$
w $L^2$, to $(f_n)$ jest ciągiem Cauchy'ego
dzięki izometryczności $U$. Istnieje zatem
$f\in L^2$ z $f_n\to f$, a ciągłość $U$ daje
$h=Uf$. Zarazem $U(\mathcal S)=\mathcal S$,
gdyż $\mathcal F$ jest automorfizmem $\mathcal S$
i mnożenie przez niezerową stałą nie zmienia
tej przestrzeni. Obraz $U(L^2)$ zawiera więc
gęsty podzbiór $\mathcal S$, a będąc domknięty,
jest całym $L^2$.

Dla $f,g\in L^2$ wybierzmy $f_n,g_n\in\mathcal S$
zbiegające do nich w normie $L^2$.
Z nierówności Cauchy'ego--Schwarza wynika
\[
 |\langle f_n,g_n\rangle_2-\langle f,g\rangle_2|
 \le\|f_n-f\|_2\|g_n\|_2+\|f\|_2\|g_n-g\|_2
 \longrightarrow0,
\]
ponieważ ciąg $(\|g_n\|_2)$ jest ograniczony.
Analogiczne oszacowanie dla $Uf_n,Ug_n$
uzasadnia przejście do granicy w tożsamości
$\langle U_0f_n,U_0g_n\rangle_2
=\langle f_n,g_n\rangle_2$.
Zatem $U$ zachowuje iloczyn skalarny i jest
surjektywny, czyli jest operatorem unitarnym.
Przyjmując $\mathcal F=(2\pi)^{d/2}U$ na $L^2$,
otrzymujemy wszystkie twierdzone własności.
\end{proof}

Operator odwrotny na $L^2$ ma postać
\[
 \mathcal F^{-1}=(2\pi)^{-d}\mathcal R\mathcal F,
 \qquad \mathcal R f(x)=f(-x).
\]
Rzeczywiście, $\mathcal R$ jest izometrią $L^2$
na mocy podstawienia $y=-x$. Obie strony
tożsamości są ciągłe na $L^2$ i pokrywają się
na gęstej przestrzeni $\mathcal S$ dzięki
twierdzeniu~\ref{thm:pde8-inversion}.
W szczególności także na $L^2$ mamy
$\mathcal F^2=(2\pi)^d\mathcal R$.

\begin{stwierdzenie}[Zgodność transformaty całkowej i przedłużenia]
Jeśli $f\in L^1(\R^d)\cap L^2(\R^d)$, to
transformata określona całką
\eqref{eq:pde8-convention} jest reprezentantem
klasy $\mathcal Ff\in L^2$ z twierdzenia
Plancherela. Analogiczna zgodność zachodzi
dla transformaty odwrotnej.
\end{stwierdzenie}
\begin{proof}
Potrzebujemy aproksymacji jednocześnie w obu normach.
Dla $R>0$ połóżmy $v_R=f\chi_{B_R(0)}$.
Z całkowalności $|f|$ i $|f|^2$ wynika
$\|f-v_R\|_1+\|f-v_R\|_2\to0$.
Dla ustalonego $R$ wygładzamy $v_R$ przez
jądro $\rho_\varepsilon$ z sekcji
~\ref{sec:pde7-mollifiers}. Tak jak w dowodzie
twierdzenia~\ref{thm:pde8-density},
$v_R*\rho_\varepsilon\in C_c^\infty(\R^d)$,
a dla $0<\varepsilon<1$ obie funkcje
$v_R$ i $v_R*\rho_\varepsilon$ znikają
prawie wszędzie poza $\overline{B_{R+1}(0)}$.
Twierdzenie~\ref{thm:pde7-approx-id} dla
$p=1$ i $p=2$ daje zatem zbieżność w obu
normach globalnych. Wybierając najpierw $R_n$,
a następnie $\varepsilon_n$, uzyskujemy
$f_n\in C_c^\infty$ takie, że
$\|f-f_n\|_1+\|f-f_n\|_2<1/n$.

Oznaczmy przez $h$ transformatę całkową $f$,
a przez $H$ jej transformatę w sensie $L^2$.
Z oszacowania dla transformaty całkowej
mamy $\|\widehat f_n-h\|_\infty\le\|f_n-f\|_1\to0$.
Z twierdzenia Plancherela
$\|\widehat f_n-H\|_2=(2\pi)^{d/2}\|f_n-f\|_2\to0$.
Dla każdej kuli $B_R(0)$ otrzymujemy więc
\[
 \|H-h\|_{L^2(B_R(0))}
 \le\|H-\widehat f_n\|_2
 +\lambda_d(B_R(0))^{1/2}\|\widehat f_n-h\|_\infty
 \longrightarrow0.
\]
Funkcje $H,h$ są zatem równe prawie wszędzie
na każdej takiej kuli, a przeliczalne pokrycie
$\R^d$ kulami o całkowitych promieniach daje
$H=h$ prawie wszędzie na całej przestrzeni.
Dla odwrotności stosujemy wykazaną tożsamość
$\mathcal F^{-1}=(2\pi)^{-d}\mathcal R\mathcal F$
oraz tę samą tożsamość dla całek na $L^1$.
\end{proof}

\begin{uwaga}[Właściwe znaczenie całki dla $L^2$]
Jeśli $f\in L^2\setminus L^1$, zapis $\widehat f$
oznacza przedłużenie operatora z twierdzenia
Plancherela, a nie automatycznie zbieżną całkę
z~\eqref{eq:pde8-convention}.
Dla każdego $R>0$ funkcja
$f_R=f\chi_{B_R(0)}$ należy do $L^1\cap L^2$,
gdyż $\|f_R\|_1\le
\lambda_d(B_R(0))^{1/2}\|f\|_2$.
Ponadto $f_R\to f$ w $L^2$, więc
\[
 \widehat f
 =\lim_{R\to\infty}^{L^2}
 \int_{B_R(0)} e^{-ix\cdot\xi}f(x)
 \,\mathrm d\lambda_d(x).
\]
Nie jest to twierdzenie o zbieżności punktowej
tej rodziny całek.

Na przykład $f(x)=(1+|x|)^{-3/4}$ na $\R$ spełnia
\[
 \int_\R |f(x)|^2\,\mathrm d\lambda_1(x)
 =2\int_0^\infty(1+x)^{-3/2}\,\mathrm d\lambda_1(x)=4,
\]
lecz $\int_\R |f(x)|\,\mathrm d\lambda_1(x)=+\infty$,
bo całka $\int_0^\infty(1+x)^{-3/4}\,\mathrm d\lambda_1(x)$
jest rozbieżna. W szczególności całka Fouriera
dla $\xi=0$ jest rozbieżna, mimo że transformata
w sensie $L^2$ istnieje jako klasa funkcji.
\end{uwaga}

\section[Dystrybucje temperowane]{Dystrybucje temperowane: przestrzeń sprzężona do $\mathcal S$}
\label{sec:pde8-tempered}

\begin{definicja}[Dystrybucja temperowana]
\label{def:pde8-tempered}
Przestrzenią dystrybucji temperowanych $\mathcal S'(\R^d)$
nazywamy ciągłą przestrzeń sprzężoną do $\mathcal S(\R^d)$,
złożoną ze wszystkich ciągłych zespolonych funkcjonałów
liniowych $T:\mathcal S(\R^d)\to\C$.
Działanie funkcjonału zapisujemy jako
$\langle T,\varphi\rangle=T(\varphi)$.
Jest ono liniowe względem $\varphi$, bez sprzężenia
zespolonego, zgodnie z konwencją rozdziału
~\ref{chap:pde-distributions}.
Zbieżność $T_n\to T$ w sensie dystrybucji temperowanych
oznacza zbieżność
\[
 \langle T_n,\varphi\rangle\to\langle T,\varphi\rangle
 \qquad\text{dla każdego }\varphi\in\mathcal S(\R^d).
\]
Jest to zbieżność w topologii słabej*,
czyli $\sigma(\mathcal S',\mathcal S)$.
Nie oznacza ona zbieżności jednostajnej na wszystkich
funkcjach testowych.
\end{definicja}

\begin{twierdzenie}[Globalne oszacowanie dystrybucji temperowanej]
\label{thm:pde8-growth}
Funkcjonał liniowy $T$ na $\mathcal S(\R^d)$ jest ciągły
dokładnie wtedy, gdy istnieją $C>0$ i $N,m\in\N_0$ takie, że
\begin{equation}\label{eq:pde8-growth}
 |\langle T,\varphi\rangle|\le C\,q_{N,m}(\varphi)
 \qquad(\varphi\in\mathcal S(\R^d)).
\end{equation}
\end{twierdzenie}
\begin{proof}
Jeśli $T$ jest ciągły, to przeciwobraz zbioru
$\{z\in\C:|z|<1\}$ jest otoczeniem zera.
Z definicji topologii wyznaczonej przez półnormy
istnieje skończona rodzina warunków
$q_{N_j,m_j}(\varphi)<r_j$, gdzie $r_j>0$,
których jednoczesne spełnienie zapewnia
$|\langle T,\varphi\rangle|<1$.
Wybieramy $N=\max_j N_j$, $m=\max_j m_j$
i $r=\min_j r_j>0$.
Ponieważ $q_{N_j,m_j}\le q_{N,m}$, wystarcza
$q_{N,m}(\varphi)<r$.
Dla $\varphi\ne0$ mamy $q_{N,m}(\varphi)>0$,
gdyż $q_{N,m}(\varphi)\ge\|\varphi\|_\infty$.
Stosując powyższy warunek do
$r\varphi/(2q_{N,m}(\varphi))$, otrzymujemy
\[
 |\langle T,\varphi\rangle|
 <\frac{2}{r}q_{N,m}(\varphi).
\]
Dla $\varphi=0$ nierówność nieostra zachodzi
z liniowości. To daje~\eqref{eq:pde8-growth}.
Odwrotnie, jeśli zachodzi~\eqref{eq:pde8-growth},
to dla dowolnego $\eta>0$ warunek
$q_{N,m}(\varphi)<\eta/C$ pociąga
$|T(\varphi)|<\eta$. Funkcjonał jest więc ciągły
w zerze, a z liniowości ciągły w każdym punkcie.
\end{proof}

\begin{stwierdzenie}[Granice dystrybucji temperowanych]
Jeśli $T_n\in\mathcal S'(\R^d)$ i dla każdego
$\varphi\in\mathcal S(\R^d)$ istnieje skończona
granica $T(\varphi)=\lim_n T_n(\varphi)$,
to $T\in\mathcal S'(\R^d)$. Co więcej, istnieją
$C>0$ oraz $N,m\in\N_0$ wspólne dla całego ciągu,
takie że
\[
 |T_n(\varphi)|\le C q_{N,m}(\varphi)
 \qquad(n\ge1,\ \varphi\in\mathcal S(\R^d)).
\]
\end{stwierdzenie}
\begin{proof}
Dla $k\ge1$ określmy zbiory
\[
 A_k=\{\varphi\in\mathcal S:
 |T_n(\varphi)|\le k\text{ dla każdego }n\ge1\}.
\]
Są one domknięte jako przecięcia domkniętych
przeciwobrazów zbioru $\{z:|z|\le k\}$ przez
ciągłe funkcjonały $T_n$. Ich suma jest równa
$\mathcal S$, ponieważ każdy zbieżny ciąg
liczb zespolonych jest ograniczony.
Przestrzeń $\mathcal S$ jest zupełna
metryzowalna (sekcja~\ref{sec:pde8-schwartz}),
więc z twierdzenia Baire'a pewien $A_k$
ma niepuste wnętrze. Istnieją zatem
$\varphi_0\in A_k$, $r>0$ i $N,m\in\N_0$
takie, że $q_{N,m}(\psi)<r$ pociąga
$\varphi_0+\psi\in A_k$. Możemy użyć jednej
półnormy dzięki monotoniczności rodziny
$q_{N,m}$, tak jak w poprzednim dowodzie.
Wtedy dla każdego $n$
\[
 |T_n(\psi)|\le
 |T_n(\varphi_0+\psi)|+|T_n(\varphi_0)|\le2k.
\]
Dla $\varphi\ne0$ podstawiamy
$\psi=r\varphi/(2q_{N,m}(\varphi))$,
uzyskując $|T_n(\varphi)|\le(4k/r)q_{N,m}(\varphi)$.
Dla $\varphi=0$ oszacowanie jest oczywiste.
Przechodząc do granicy, dostajemy tę samą
nierówność dla $T$. Granica punktowa
funkcjonałów liniowych jest liniowa, więc
twierdzenie~\ref{thm:pde8-growth} daje ciągłość $T$.
\end{proof}

\begin{uwaga}[Relacja między $\mathcal D'$ a $\mathcal S'$]
Ograniczenie $T\in\mathcal S'(\R^d)$ do
$C_c^\infty(\R^d)$ jest dystrybucją.
Istotnie, jeśli $\supp\varphi\subset K$ dla
niepustego zbioru zwartego $K$, to
\[
 q_{N,m}(\varphi)
 \le (1+\max_{x\in K}|x|)^N
 \max_{|\beta|\le m}\sup_{x\in K}
 |\partial^\beta\varphi(x)|.
\]
Dla $K=\varnothing$ jedyną taką funkcją jest
funkcja zerowa. Globalne oszacowanie
~\eqref{eq:pde8-growth} daje zatem lokalne
oszacowanie z twierdzenia
~\ref{thm:pde7-local-order}.
Odwzorowanie ograniczenia
$\mathcal S'\to\mathcal D'$ jest liniowe
i iniektywne: jeśli dwa ciągłe funkcjonały
na $\mathcal S$ zgadzają się na $C_c^\infty$,
to zgadzają się na całej $\mathcal S$
dzięki gęstości z twierdzenia~\ref{thm:pde8-density}.
Możemy więc utożsamiać $\mathcal S'$
z podprzestrzenią $\mathcal D'$.
Dystrybucja jest temperowana dokładnie wtedy,
gdy posiada ciągłe przedłużenie z $C_c^\infty$
na $\mathcal S$; takie przedłużenie jest jednoznaczne.
Nie wszystkie dystrybucje są temperowane,
jak pokaże ostatni z poniższych przykładów.
W szczególności transformata Fouriera dowolnej
dystrybucji z $\mathcal D'$ nie ma automatycznie
sensu w rachunku, który zbudujemy dla $\mathcal S'$.
\end{uwaga}

\begin{przyklad}[Obiekty temperowane i nietemperowane]
\begin{enumerate}
\item Niech $f$ będzie funkcją mierzalną spełniającą
$|f(x)|\le C(1+|x|)^M$ prawie wszędzie,
gdzie $C>0$, $M\in\N_0$.
Definiujemy
\[
 \langle T_f,\varphi\rangle
 =\int_{\R^d}f(x)\varphi(x)\,\mathrm d\lambda_d(x).
\]
Dla całkowitego $L>d$ mamy
\[
 \int_{\R^d}|f(x)\varphi(x)|\,\mathrm d\lambda_d(x)
 \le C C_{d,L,1}\,q_{M+L,0}(\varphi),
\]
gdzie $C_{d,L,1}$ jest stałą z lematu
~\ref{lem:pde8-seminorms}.
Całka jest zatem bezwzględnie zbieżna,
a liniowy funkcjonał $T_f$ jest ciągły
z twierdzenia~\ref{thm:pde8-growth}.
Należą tu wielomiany, stałe i funkcje istotnie
ograniczone. Punktowe oszacowanie wielomianowe
jest warunkiem wystarczającym, nie definicją
temperowaności funkcji.

\item Każda funkcja $f\in L^p(\R^d)$,
$1\le p\le\infty$, również określa
dystrybucję temperowaną. Niech $p'$ będzie
wykładnikiem sprzężonym, czyli
$1/p+1/p'=1$ z konwencją $1/\infty=0$.
Z nierówności Höldera mamy
$|\langle T_f,\varphi\rangle|\le
\|f\|_p\|\varphi\|_{p'}$.
Jeśli $p'<\infty$, wybieramy całkowite $N$
takie, że $Np'>d$. Z tego samego całkowania
radialnego co w lemacie~\ref{lem:pde8-seminorms}
wynika $\|\varphi\|_{p'}\le
C_{d,N,p'}q_{N,0}(\varphi)$, gdzie
$C_{d,N,p'}^{p'}=\int_{\R^d}(1+|x|)^{-Np'}
\,\mathrm d\lambda_d(x)<\infty$.
Jeśli $p'=\infty$, mamy
$\|\varphi\|_\infty=q_{0,0}(\varphi)$.
W obu przypadkach otrzymujemy ciągłość.
W szczególności każda funkcja z $L^2$
może być traktowana jako dystrybucja temperowana,
co wykorzystamy przy porównywaniu transformat Fouriera.

\item Dla $a\in\R^d$ oraz $\alpha\in\N_0^d$ wzór
\[
 \langle\partial^\alpha\delta_a,\varphi\rangle
 =(-1)^{|\alpha|}\partial^\alpha\varphi(a)
\]
określa ciągłe przedłużenie pochodnej delty
na $\mathcal S$. Istotnie, moduł prawej strony
jest ograniczony przez $q_{0,|\alpha|}(\varphi)$.
Przypadek $\alpha=0$ daje samą deltę.

\item Każda dystrybucja $T\in\mathcal D'(\R^d)$
o zwartym nośniku $K$ jest temperowana.
Wybieramy $R\ge1$ takie, że $K\subset B_R(0)$
i funkcję odcinającą $\psi_R(x)=\psi(x/R)$
z dowodu twierdzenia~\ref{thm:pde8-density}.
Jest ona równa $1$ na otoczeniu $K$.
Dla $\varphi\in\mathcal S$ określamy
\[
 \langle\widetilde T,\varphi\rangle
 =\langle T,\psi_R\varphi\rangle.
\]
Iloczyn $\psi_R\varphi$ jest gładki
i ma nośnik w ustalonym zbiorze zwartym
$K_1=\overline{B_{2R}(0)}$.
Lokalne oszacowanie z rozdziału
~\ref{chap:pde-distributions} daje pewne
$C>0$ i $m\in\N_0$ takie, że
\[
 \begin{aligned}
 |\langle\widetilde T,\varphi\rangle|
 &\le C\max_{|\beta|\le m}\sup_{x\in K_1}
 |\partial^\beta(\psi_R\varphi)(x)|\\
 &\le C\,2^m q_{0,m}(\psi_R)q_{0,m}(\varphi).
 \end{aligned}
\]
Druga nierówność wynika z wielowskaźnikowej
reguły Leibniza. Wobec tego $\widetilde T$
jest ciągły na $\mathcal S$.
Jeśli $\varphi\in C_c^\infty$, to
$(1-\psi_R)\varphi$ jest testem o nośniku
rozłącznym z $K$. Z definicji nośnika
dystrybucji (sekcja~\ref{sec:pde7-support})
wynika $\langle T,(1-\psi_R)\varphi\rangle=0$.
Zatem $\widetilde T$ rzeczywiście przedłuża $T$.
Nie zależy od wyboru funkcji odcinającej,
ponieważ ciągłe przedłużenie jest jednoznaczne.

\item Dystrybucja regularna wyznaczona przez
$f(x)=e^{|x|^2}$ na $\R^d$ nie jest temperowana.
Funkcja $f$ jest lokalnie całkowalna, więc
określa dystrybucję na $C_c^\infty$.
Załóżmy, że ma ona ciągłe przedłużenie
$\widetilde T$ na $\mathcal S$.
Dla $R\ge1$ bierzemy
$\varphi_R(x)=\psi_R(x)e^{-|x|^2/2}$,
gdzie $\psi_R$ jest funkcją z dowodu
twierdzenia~\ref{thm:pde8-density}.
Funkcja $e^{-|x|^2/2}$ należy do $\mathcal S$,
a konstrukcja aproksymacji w tym dowodzie
daje $\varphi_R\to e^{-|x|^2/2}$ w $\mathcal S$.
Z ciągłości $\widetilde T$ wynikałaby
zbieżność $\widetilde T(\varphi_R)$
do skończonej liczby. Tymczasem
$0\le\psi_R\le1$, $\psi_R=1$ na $B_R(0)$
i $\varphi_R\in C_c^\infty$, więc
\[
 \begin{aligned}
 \langle\widetilde T,\varphi_R\rangle
 &=\int_{\R^d}\psi_R(x)e^{|x|^2/2}
 \,\mathrm d\lambda_d(x)\\
 &\ge\lambda_d(B_R(0))
 =R^d\lambda_d(B_1(0))\longrightarrow+\infty.
 \end{aligned}
\]
Jest to sprzeczność.
\end{enumerate}
\end{przyklad}

\section{Transformata Fouriera dystrybucji temperowanych}
\label{sec:pde8-fourier-distributions}

\begin{definicja}[Transformata na $\mathcal S'$]
\label{def:pde8-distFT}
Dla $T\in\mathcal S'(\R^d)$ definiujemy $\widehat T\in\mathcal S'(\R^d)$
przez \emph{transpozycję} operatora $\mathcal F$:
\[
 \langle\widehat T,\varphi\rangle
 =\langle T,\widehat\varphi\rangle,
 \qquad\varphi\in\mathcal S(\R^d).
\]
Analogicznie definiujemy $\mathcal F^{-1}T$ przez
$\langle\mathcal F^{-1}T,\varphi\rangle
=\langle T,\mathcal F^{-1}\varphi\rangle$.
\end{definicja}

Definicja jest poprawna: z twierdzeń
~\ref{thm:pde8-rules} i~\ref{thm:pde8-inversion}
wynika ciągłość obu operatorów na $\mathcal S$,
więc ich złożenia z $T$ są ciągłymi funkcjonałami
liniowymi. Nie stosujemy sprzężenia zespolonego:
jest to transpozycja względem liniowego działania
dystrybucji, nie operator sprzężony w przestrzeni Hilberta.

\begin{stwierdzenie}[Zgodność trzech definicji transformaty]
Jeśli $f\in L^1(\R^d)$, to
$\mathcal F T_f=T_{\widehat f}$, gdzie po prawej
stronie używamy transformaty całkowej.
Jeśli $f\in L^2(\R^d)$, to ta sama równość
zachodzi dla transformaty z twierdzenia Plancherela.
\end{stwierdzenie}
\begin{proof}
Dla $f\in L^1$ i $\varphi\in\mathcal S$
całka podwójna z modułu
$f(x)e^{-ix\cdot\xi}\varphi(\xi)$ wynosi
$\|f\|_1\|\varphi\|_1<\infty$. Stąd
\[
 \begin{aligned}
 \langle\widehat{T_f},\varphi\rangle
 &=\int_{\R^d}f(x)\widehat\varphi(x)
 \,\mathrm d\lambda_d(x)\\
 &=\int_{\R^d}\widehat f(\xi)\varphi(\xi)
 \,\mathrm d\lambda_d(\xi).
 \end{aligned}
\]
Funkcja $\widehat f$ jest ograniczona, więc
wyznacza dystrybucję temperowaną.
Dla $f\in L^2$ wybieramy $f_n\in\mathcal S$
zbiegające do $f$ w $L^2$.
Z twierdzenia~\ref{thm:pde8-plancherel}
wynika $\widehat f_n\to\widehat f$ w $L^2$.
Dla każdego testu $\varphi$ obie różnice działań
szacujemy przez nierówność Cauchy'ego--Schwarza:
\[
 \begin{aligned}
 |\langle T_{f_n}-T_f,\widehat\varphi\rangle|
 &\le\|f_n-f\|_2\|\widehat\varphi\|_2,\\
 |\langle T_{\widehat f_n}-T_{\widehat f},\varphi\rangle|
 &\le\|\widehat f_n-\widehat f\|_2\|\varphi\|_2.
 \end{aligned}
\]
Prawie wszędzie określone funkcje $L^2$
wyznaczają funkcjonały temperowane na mocy
sekcji~\ref{sec:pde8-tempered}. Przechodząc
do granicy w udowodnionej równości dla $f_n$,
otrzymujemy tezę.
\end{proof}

Pochodne dystrybucji temperowanych definiujemy
przez $\langle\partial_j T,\varphi\rangle
=-\langle T,\partial_j\varphi\rangle$.
Są temperowane dzięki ciągłości różniczkowania
na $\mathcal S$. Ogólniej, jeśli $a\in C^\infty(\R^d)$
i dla każdej pochodnej istnieją
$C_\gamma>0$, $M_\gamma\in\N_0$ z
$|\partial^\gamma a(x)|\le C_\gamma(1+|x|)^{M_\gamma}$,
to mnożenie przez $a$ jest ciągłe na $\mathcal S$.
Dla $m\in\N_0$ bierzemy
$M=\max_{|\gamma|\le m} M_\gamma$ i
$C=\max_{|\gamma|\le m} C_\gamma$.
Z reguły Leibniza dostajemy
\[
 q_{N,m}(a\varphi)\le 2^m C q_{N+M,m}(\varphi).
\]
Definiujemy więc $aT$ przez
$\langle aT,\varphi\rangle=\langle T,a\varphi\rangle$.
Te definicje zgadzają się z rachunkiem
na $\mathcal D'$ po ograniczeniu do testów
o zwartym nośniku. W szczególności można
mnożyć przez wielomiany oraz funkcje Schwartza;
dowolna funkcja gładka nie musi zachowywać
przestrzeni $\mathcal S$.

\begin{twierdzenie}[Rachunek na $\mathcal S'$]
\label{thm:pde8-dist-rules}
Transformata $\mathcal F:\mathcal S'\to\mathcal S'$ jest
liniowym automorfizmem, którego odwrotnością jest
$\mathcal F^{-1}$. Dla $T\in\mathcal S'(\R^d)$ i
$\alpha\in\N_0^d$ zachodzi
\[
 \mathcal F(\partial^\alpha T)=(i\xi)^\alpha\widehat T,
 \qquad \mathcal F(x^\alpha T)=i^{|\alpha|}
 \partial_\xi^\alpha\widehat T.
\]
\end{twierdzenie}
\begin{proof}
Dla każdego testu $\varphi\in\mathcal S$ mamy
\[
 \begin{aligned}
 \langle\mathcal F^{-1}\mathcal FT,\varphi\rangle
 &=\langle T,\mathcal F\mathcal F^{-1}\varphi\rangle
 =\langle T,\varphi\rangle,\\
 \langle\mathcal F\mathcal F^{-1}T,\varphi\rangle
 &=\langle T,\mathcal F^{-1}\mathcal F\varphi\rangle
 =\langle T,\varphi\rangle.
 \end{aligned}
\]
Liniowość wynika z definicji. Oba operatory
zachowują zbieżność słabą*: dla $T_n\to T$
testujemy $T_n$ odpowiednio przez ustaloną
funkcję $\mathcal F\varphi$ lub $\mathcal F^{-1}\varphi$.
Zatem są również ciągłe w topologii słabej*,
ponieważ każdy funkcjonał współrzędny po złożeniu
z nimi jest funkcjonałem współrzędnym tej topologii.
Dla pierwszej równości i testu $\varphi$ liczymy
\[
 \langle\widehat{\partial_jT},\varphi\rangle
 =-\langle T,\partial_j\widehat\varphi\rangle
 =\langle T,\widehat{i\xi_j\varphi}\rangle
 =\langle i\xi_j\widehat T,\varphi\rangle.
\]
Dla drugiej tożsamości korzystamy z równości
$\mathcal F(\partial_{\xi_j}\varphi)(x)
=ix_j\widehat\varphi(x)$:
\[
 \begin{aligned}
 \langle\mathcal F(x_jT),\varphi\rangle
 &=\langle T,x_j\widehat\varphi\rangle\\
 &=-i\langle T,\mathcal F(\partial_{\xi_j}\varphi)\rangle\\
 &=-i\langle\widehat T,\partial_{\xi_j}\varphi\rangle
 =i\langle\partial_{\xi_j}\widehat T,\varphi\rangle.
 \end{aligned}
\]
Iteracja daje oba wzory dla wielowskaźników,
ponieważ pochodne i mnożenia współrzędniowe
odpowiedniego rodzaju komutują.
\end{proof}

\begin{przyklad}[Cztery transformaty podstawowe]
\label{ex:pde8-four}
W $\mathcal S'(\R^d)$ zachodzą wzory
\[
 \widehat{\delta_0}=1,\quad
 \widehat 1=(2\pi)^d\delta_0,\quad
 \widehat{\delta_a}(\xi)=e^{-ia\cdot\xi},\quad
 \mathcal F(e^{ia\cdot x})=(2\pi)^d\delta_a.
\]
Dla dowolnego $a\in\R^d$ i testu $\varphi$ mamy
\[
 \langle\widehat{\delta_a},\varphi\rangle
 =\widehat\varphi(a)
 =\int_{\R^d}e^{-ia\cdot\xi}\varphi(\xi)
 \,\mathrm d\lambda_d(\xi).
\]
Przy $a=0$ otrzymujemy $\widehat{\delta_0}=1$.
Natomiast wzór odwrócenia dla $\varphi$ daje
\[
 \langle\mathcal F(e^{ia\cdot x}),\varphi\rangle
 =\int_{\R^d}e^{ia\cdot x}\widehat\varphi(x)
 \,\mathrm d\lambda_d(x)
 =(2\pi)^d\varphi(a).
\]
W szczególności przy $a=0$ mamy
$\widehat1=(2\pi)^d\delta_0$. Obie całki
są bezwzględnie zbieżne, ponieważ
$\varphi,\widehat\varphi\in L^1$.
Równości po prawej należy rozumieć \emph{dystrybucyjnie},
a nie jako zwykłe całki bezwzględnie zbieżne.
\end{przyklad}

\begin{twierdzenie}[Splot dystrybucji temperowanej z funkcją Schwartza]
\label{thm:pde8-convol-temp}
Niech $T\in\mathcal S'(\R^d)$ i $\varphi\in\mathcal S(\R^d)$.
Funkcja $(T*\varphi)(x):=\langle T(y),\varphi(x-y)\rangle$
jest gładka, ma wzrost co najwyżej wielomianowy
(wraz z każdą ustaloną pochodną) i wyznacza dystrybucję
temperowaną. Ponadto
$\mathcal F(T*\varphi)=\widehat T\,\widehat\varphi$ w $\mathcal S'$.
\end{twierdzenie}
\begin{proof}
Niech $F_x(y)=\varphi(x-y)$.
Z nierówności $(1+|y|)\le(1+|x|)(1+|x-y|)$
wynika dla każdego wielowskaźnika $\alpha$
\[
 q_{N,m}(\partial_x^\alpha F_x)
 \le(1+|x|)^N q_{N,m+|\alpha|}(\varphi).
\]
W szczególności wszystkie te funkcje należą do
$\mathcal S$. Dla $|h|\le1$ i $h\ne0$ podstawowe
twierdzenie rachunku całkowego, zastosowane
dwukrotnie wzdłuż kierunku $e_j$, daje
\[
 \frac{\partial_x^\alpha F_{x+he_j}
       -\partial_x^\alpha F_x}{h}
 -\partial_x^{\alpha+e_j}F_x
 =h\int_0^1(1-t)\partial_x^{\alpha+2e_j}F_{x+the_j}
 \,\mathrm d\lambda_1(t).
\]
Po różniczkowaniu względem $y$ i oszacowaniu
wagą wynika stąd
\[
 q_{N,m}(\text{różnica po lewej})
 \le \tfrac{|h|}{2}(2+|x|)^N
 q_{N,m+|\alpha|+2}(\varphi).
\]
Analogiczne oszacowanie pierwszej różnicy
przez całkę z pierwszej pochodnej daje
ciągłość $x\mapsto\partial_x^\alpha F_x$
w każdej półnormie (dla ogólnego przesunięcia
sumujemy różnice w kolejnych współrzędnych).
Z ciągłości $T$ możemy więc różniczkować
indukcyjnie:
$\partial_x^\alpha(T*\varphi)(x)
=\langle T,\partial_x^\alpha F_x\rangle$.
Jeśli $C,N,m$ są stałymi z globalnego oszacowania
~\eqref{eq:pde8-growth} dla $T$, to
\[
 |\partial_x^\alpha(T*\varphi)(x)|
 \le C(1+|x|)^N q_{N,m+|\alpha|}(\varphi).
\]
Jest to gładka funkcja o wzroście wielomianowym
wraz z pochodnymi, a dla $\alpha=0$ oszacowanie
zapewnia temperowaność jej dystrybucji regularnej.
Dla dowodu tożsamości Fouriera ustalmy
$\psi\in\mathcal S(\R^d)$ i oznaczmy
$A_x(y)=F_x(y)\widehat\psi(x)$.
W każdej półnormie mamy oszacowanie
\[
 q_{N,m}(A_x)\le(1+|x|)^N
 q_{N,m}(\varphi)|\widehat\psi(x)|.
\]
Prawa strona jest całkowalna względem $x$,
ponieważ $\widehat\psi\in\mathcal S$.
Niech $Q_R=[-R,R]^d$ i
$B_R(y)=\int_{Q_R}A_x(y)\,\mathrm d\lambda_d(x)$.
Możemy różniczkować całkę względem $y$,
gdyż pochodne $\varphi$ są ograniczone, a
$\widehat\psi$ jest całkowalna. Otrzymane
półnormy $B_R$ są skończone przez powyższe
oszacowanie, więc $B_R\in\mathcal S$.
Odwzorowanie $x\mapsto A_x$ jest ciągłe w każdej
półnormie, jak wykazano dla $F_x$, i dlatego
jednostajnie ciągłe w tej półnormie na $Q_R$.
Dla podziału $Q_R$ na prostokąty o średnicy
co najwyżej $\delta$ błąd sumy Riemanna
w półnormie $q_{N,m}$ nie przekracza
$\lambda_d(Q_R)$ razy
\[
 \sup_{\substack{x,z\in Q_R\\|x-z|\le\delta}}
 q_{N,m}(A_x-A_z),
\]
co dąży do zera przy $\delta\downarrow0$.
Ciągłość $T$ pozwala przejść do granicy sum,
a po stronie skalarnej całkujemy funkcję ciągłą
$x\mapsto\langle T,A_x\rangle$. Otrzymujemy
$\langle T,B_R\rangle=\int_{Q_R}\langle T,A_x\rangle
\,\mathrm d\lambda_d(x)$.

Niech $B(y)=\int_{\R^d}A_x(y)\,\mathrm d\lambda_d(x)$.
Całka i każda jej pochodna względem $y$ istnieją
dzięki ograniczoności pochodnych $\varphi$
i całkowalności $\widehat\psi$. Oszacowanie wag
daje $B\in\mathcal S$ oraz
\[
 q_{N,m}(B-B_R)\le q_{N,m}(\varphi)
 \int_{\R^d\setminus Q_R}(1+|x|)^N
 |\widehat\psi(x)|\,\mathrm d\lambda_d(x)\to0.
\]
Zatem $B_R\to B$ w $\mathcal S$.
Całka skalarna po całej przestrzeni jest
bezwzględnie zbieżna przez oszacowanie
wzrostu $T*\varphi$ i szybki zanik $\widehat\psi$.
Przechodząc do granicy w równości na $Q_R$,
dostajemy
\[
 \langle\mathcal F(T*\varphi),\psi\rangle
 =\int_{\R^d}(T*\varphi)(x)\widehat\psi(x)
 \,\mathrm d\lambda_d(x)
 =\langle T,B\rangle.
\]
Po podstawieniu $z=x-y$ i użyciu wzoru
$\widehat\psi(x)=\int_{\R^d}e^{-ix\cdot\xi}\psi(\xi)
\,\mathrm d\lambda_d(\xi)$ liczymy, stosując Fubiniego
(całka podwójna z modułu dla ustalonego $y$
wynosi $\|\varphi\|_1\|\psi\|_1$),
\[
 \int_{\R^d}\varphi(x-y)\widehat\psi(x)
 \,\mathrm d\lambda_d(x)
 =\int_{\R^d}e^{-iy\cdot\xi}\widehat\varphi(\xi)
 \psi(\xi)\,\mathrm d\lambda_d(\xi)
 =\mathcal F(\widehat\varphi\,\psi)(y).
\]
Stąd $\langle\mathcal F(T*\varphi),\psi\rangle
=\langle T,\mathcal F(\widehat\varphi\psi)\rangle
=\langle\widehat T,\widehat\varphi\psi\rangle$,
czyli $\mathcal F(T*\varphi)=\widehat T\,\widehat\varphi$.
W szczególności nie pojawia się dodatkowy czynnik
$(2\pi)^{-d}$: w naszej konwencji dotyczy on
transformaty \emph{iloczynu} funkcji, a nie splotu.
\end{proof}

\section{Mnożniki Fouriera: kiedy operator jest ograniczony?}
\label{sec:pde8-multipliers}

\begin{twierdzenie}[Mnożnik w $L^2$]
\label{thm:pde8-multiplier}
Niech $m:\R^d\to\C$ będzie funkcją mierzalną i
istotnie ograniczoną. Dla $f\in L^2(\R^d)$ istnieje dokładnie
jedna funkcja $T_mf\in L^2(\R^d)$ spełniająca
$\widehat{T_mf}=m\widehat f$ w sensie $L^2$.
Operator $T_m$ jest liniowy i ograniczony oraz
$\|T_m\|_{L^2\to L^2}=\|m\|_{L^\infty}$.
\end{twierdzenie}
\begin{proof}
Mnożenie przez $m$ jest ograniczone na $L^2$, a transformata
Fouriera jest unitarnym operatorem z dokładnością do
stałego czynnika, co daje istnienie, jednoznaczność i
oszacowanie z góry. Dla dolnego oszacowania bierzemy
$\widehat f$ proporcjonalną do funkcji charakterystycznej
zbioru dodatniej i skończonej miary, na którym
$|m|>\|m\|_\infty-\varepsilon$. Taki zbiór istnieje
(po przecięciu z dostatecznie dużą kulą), a po
$\varepsilon\downarrow0$ otrzymujemy równość norm.
\end{proof}

\begin{przyklad}[Transformata Riesza w $L^2$]
Dla $1\le j\le d$ połóżmy
$m_j(\xi)=-i\xi_j/|\xi|$ dla $\xi\ne0$ i $m_j(0)=0$.
Poprzednie twierdzenie definiuje ograniczony operator
$R_j=T_{m_j}:L^2\to L^2$, tzw. transformatę Riesza,
z $\|R_j\|\le1$. Wartość mnożnika w pojedynczym punkcie
nie zmienia operatora $L^2$. Ograniczoność w $L^p$ dla
$p\ne2$ jest istotnie trudniejszym rezultatem teorii
Calderóna--Zygmunda i nie wynika z samego warunku
$m_j\in L^\infty$.
\end{przyklad}

\section{Zastosowanie I: równanie ciepła i równanie falowe}
\label{sec:pde8-evolution}

\begin{twierdzenie}[Rozwiązania ewolucyjne poprzez transformatę]
\label{thm:pde8-evolution}
Niech $\kappa,c>0$ i $f,g\in\mathcal S(\R^d)$.
Wtedy funkcje
\begin{align*}
 u(t,x)&=\mathcal F^{-1}
 \bigl(e^{-\kappa t|\xi|^2}\widehat f(\xi)\bigr)(x),&&t\ge0,\\
 v(t,x)&=\mathcal F^{-1}\biggl(
 \cos(ct|\xi|)\widehat f(\xi)
 +\frac{\sin(ct|\xi|)}{c|\xi|}\widehat g(\xi)
 \biggr)(x),&&t\in\R,
\end{align*}
gdzie iloraz w $\xi=0$ definiujemy przez wartość graniczną $t$,
są gładkimi rozwiązaniami odpowiednio zagadnień
$u_t-\kappa\Delta u=0$, $u(0,\cdot)=f$ oraz
$v_{tt}-c^2\Delta v=0$, $v(0,\cdot)=f$, $v_t(0,\cdot)=g$.
W każdym ustalonym czasie funkcje przestrzenne należą do $\mathcal S$.
\end{twierdzenie}
\begin{proof}
Z reguły różniczkowania otrzymujemy po transformacji
$\partial_t\widehat u=-\kappa|\xi|^2\widehat u$ oraz
$\partial_{tt}\widehat v=-c^2|\xi|^2\widehat v$.
Są to dla każdego $\xi$ zwyczajne równania różniczkowe
z wskazanymi danymi początkowymi. Ich rozwiązania mają
właśnie zapis z tezy. Mnożniki są gładkie w $\xi$:
zwłaszcza $\sin(ct|\xi|)/(c|\xi|)$ ma rozwinięcie w szereg
potęg $|\xi|^2$ w otoczeniu zera.
Dla $t$ w dowolnym ograniczonym przedziale ich
pochodne po $t$ i $\xi$ mają co najwyżej wielomianowy wzrost,
stąd mnożą $\mathcal S$ w siebie. Uzasadnia to wszystkie
różniczkowania i stosowanie odwrotnej transformaty.
Wartości początkowe otrzymujemy przez podstawienie $t=0$.
\end{proof}

\begin{uwaga}[Dwie różne skale propagacji]
Równanie ciepła działa w częstotliwościach przez
$m_t(\xi)=e^{-\kappa t|\xi|^2}$: dla $t>0$ tłumi wysokie
częstotliwości bardzo szybko. Równanie falowe używa
$\cos(ct|\xi|)$ i $\sin(ct|\xi|)/(c|\xi|)$; nie ma
analogicznego wygładzania wszystkich wysokich częstotliwości.
To jest spektralna wersja zjawisk z rozdziałów
\ref{chap:pde-heat} i~\ref{chap:pde-wave}.
\end{uwaga}

\section[Potencjały Bessela i przestrzenie Sobolewa]{Potencjały Bessela i pierwsze przestrzenie $H^s$}
\label{sec:pde8-bessel}

\begin{definicja}[Przestrzeń Sobolewa typu Bessela]
\label{def:pde8-Hs}
Dla dowolnego $s\in\R$ definiujemy
\[
 H^s(\R^d)=\bigl\{u\in\mathcal S'(\R^d):
 (1+|\xi|^2)^{s/2}\widehat u\in L^2(\R^d)\bigr\}.
\]
Zapis przynależności do $L^2$ oznacza tutaj, że odpowiednia
dystrybucja jest reprezentowana przez funkcję kwadratowalną.
Normę tej przestrzeni definiujemy wzorem
$\|u\|_{H^s}^2=(2\pi)^{-d}\int_{\R^d}
(1+|\xi|^2)^s|\widehat u(\xi)|^2\,\mathrm d\lambda_d(\xi)$.
W tym wzorze $\widehat u$ jest reprezentowana przez funkcję
lokalnie $L^2$: otrzymujemy ją mnożąc funkcję $L^2$
przez $(1+|\xi|^2)^{-s/2}$.
Operator $J^s=(I-\Delta)^{s/2}$ określamy na $\mathcal S'$
przez mnożnik $(1+|\xi|^2)^{s/2}$.
Dla ujemnego $s$ element $H^s$ może być autentyczną
dystrybucją, a nie funkcją lokalnie całkowalną;
wymóg $L^2$ odnosi się do ważonej transformaty,
nie do samej dystrybucji $u$.
\end{definicja}

Dla $m\in\N_0$ zapis $W^{m,2}(\R^d)$ oznacza przestrzeń
funkcji $u\in L^2(\R^d)$, których wszystkie pochodne słabe
$\partial^\alpha u$, $|\alpha|\le m$, należą do $L^2$.

\begin{uwaga}[Szeregi Fouriera na torusie a transformata Fouriera na $\R^d$]
Rozdział~\ref{chap:pde-fourier} pracował na torusie lub
na odcinku po okresowym przedłużeniu i korzystał z rozwinięć
w szeregi Fouriera, czyli dyskretnego spektrum częstości.
Tutaj pracujemy na całej przestrzeni $\R^d$, więc częstość
$\xi$ zmienia się ciągle i pojawia się całka po $\R^d$.
Obie teorie są zgodne pojęciowo, ale różnią się nośnikiem,
normalizacją oraz tym, że w części okresowej integrujemy na
$\T^d$ względem ustalonej miary, a tutaj na $\R^d$ względem
miary Lebesgue'a.
\end{uwaga}

\begin{twierdzenie}[Podstawowe własności $H^s$]
\label{thm:pde8-Hs}
Dla wszystkich $s,t\in\R$, $s\ge t$, zachodzą:
\begin{enumerate}
\item $H^s(\R^d)$ jest przestrzenią Hilberta, a
 $J^s:H^s\to L^2$ jest izometrią na całą przestrzeń $L^2$.
\item $H^s\hookrightarrow H^t$ ciągle i
 $\|u\|_{H^t}\le\|u\|_{H^s}$.
\item Dla $\alpha\in\N_0^d$, $\partial^\alpha:H^s\to
 H^{s-|\alpha|}$ jest ograniczony.
\item Dla $m\in\N_0$ zachodzi $H^m(\R^d)=W^{m,2}(\R^d)$
 ze wzajemnie równoważnymi normami.
\end{enumerate}
\end{twierdzenie}
\begin{proof}
Pierwsza teza wynika z utożsamienia $H^s$ poprzez
$J^s$ z $L^2$ i odwracalności mnożnika
$(1+|\xi|^2)^{s/2}$ na $\mathcal S'$. Pozostałe
oszacowania opierają się na nierównościach
$(1+|\xi|^2)^t\le(1+|\xi|^2)^s$ i
$|\xi^\alpha|^2\le(1+|\xi|^2)^{|\alpha|}$.
Gdy $m$ jest całkowite nieujemne, wielomian
$(1+|\xi|^2)^m$ jest równoważny
$\sum_{|\alpha|\le m}|\xi^\alpha|^2$,
ze stałymi zależnymi tylko od $m,d$. Dzięki
Plancherelowi warunek $u\in H^m$ jest równoważny
istnieniu wszystkich pochodnych dystrybucyjnych
$\partial^\alpha u\in L^2$ dla $|\alpha|\le m$.
W przypadku $m=0$ warunek ten mówi po prostu $u\in L^2$.
\end{proof}

\begin{wniosek}[Kiedy delta należy do ujemnej przestrzeni Sobolewa?]
\label{cor:pde8-delta-Hs}
Dla $s\in\R$ zachodzi
\[
 \delta_0\in H^s(\R^d)\quad\Longleftrightarrow\quad
 s<-\frac d2.
\]
\end{wniosek}
\begin{proof}
Ponieważ $\widehat{\delta_0}=1$, warunkiem jest
$\int_{\R^d}(1+|\xi|^2)^s\,\mathrm d\lambda_d(\xi)<\infty$.
Całka po kuli jednostkowej jest skończona. Poza nią,
przechodząc do współrzędnych biegunowych, porównujemy
całkę z $\int_1^\infty r^{2s+d-1}\,\mathrm d\lambda_1(r)$,
która jest skończona dokładnie wtedy, gdy $2s+d<0$.
\end{proof}

\section{Zastosowanie II: eliptyczne równanie z parametrem}
\label{sec:pde8-elliptic}

\begin{twierdzenie}[Równanie Yukawy--Helmholtza]
\label{thm:pde8-helmholtz}
Niech $a>0$ i $f\in L^2(\R^d)$. Istnieje dokładnie jedno
$u\in H^2(\R^d)$ takie, że $(-\Delta+a^2)u=f$
w sensie dystrybucji. Zachodzi
\[
 \widehat u(\xi)=\frac{\widehat f(\xi)}{|\xi|^2+a^2},
 \qquad \|u\|_{H^2}\le C_a\|f\|_{L^2},
\]
gdzie $H^2$ zdefiniowano w~\ref{def:pde8-Hs},
a $C_a=\max\{1,a^{-2}\}$ dla przyjętej tam normy.
\end{twierdzenie}
\begin{proof}
Funkcja $(|\xi|^2+a^2)^{-1}$ jest ograniczona.
Ponadto $(1+|\xi|^2)/(|\xi|^2+a^2)\le\max\{1,a^{-2}\}$.
Zatem zadany wzór wyznacza $u\in H^2$ i podane
oszacowanie przez twierdzenie Plancherela. Po
pomnożeniu przez $|\xi|^2+a^2$ otrzymujemy równanie.
Jeżeli $(-\Delta+a^2)u=0$ dla $u\in H^2$, to
$(|\xi|^2+a^2)\widehat u=0$ prawie wszędzie,
a więc $u=0$. To daje jednoznaczność.
\end{proof}

\begin{uwaga}[Dlaczego przypadek $a=0$ jest inny?]
Dla czystego równania Poissona mnożnik formalny
$|\xi|^{-2}$ ma osobliwość w zerze. Nie można więc
bez dodatkowej analizy definiować dla dowolnego
$f\in L^2$ rozwiązania $u\in H^2$ przez ten wzór.
Dokładniej, dla $1\le d\le4$ i $f\in\mathcal S(\R^d)$
z $\widehat f(0)\ne0$ iloraz
$|\xi|^{-2}\widehat f(\xi)$ nie jest funkcją $L^2$
w pobliżu zera, gdyż jego kwadrat daje po przejściu
do współrzędnych biegunowych całkę porównywalną z
$\int_0^1r^{d-5}\,\mathrm d\lambda_1(r)=\infty$.
W wymiarach $d\ge5$ ta konkretna przeszkoda
niskoczęstotliwościowa znika dla takiego gładkiego $f$;
nie zmienia to faktu, że dzielenie przez $|\xi|^2$
nie jest ograniczonym mnożnikiem $L^2\to H^2$
bez dodatkowych warunków na dane. Dodatkowo funkcje harmoniczne tworzą jądro
operatora $-\Delta$ w szerszych klasach rozwiązań.
Nie oznacza to braku rozwiązań dystrybucyjnych;
rozdział~\ref{chap:pde-distributions} konstruuje je
z użyciem funkcji fundamentalnej.
\end{uwaga}

\section{Zasada nieoznaczoności jako konsekwencja Plancherela}
\label{sec:pde8-uncertainty}

\begin{twierdzenie}[Nierówność Heisenberga w jednej współrzędnej]
\label{thm:pde8-uncertainty}
Jeśli $0\ne f\in\mathcal S(\R^d)$ i $1\le j\le d$, to
\[
 \|x_jf\|_{L^2}\,\|\xi_j\widehat f\|_{L^2}
 \ge\frac12\,\|f\|_{L^2}\,\|\widehat f\|_{L^2}.
\]
\end{twierdzenie}
\begin{proof}
Całkując pochodną $\partial_j(x_j|f|^2)$ i używając
zaniku funkcji Schwartza, dostajemy
$\|f\|_2^2=-2\operatorname{Re}\int_{\R^d}x_j\overline f\,
\partial_jf\,\mathrm d\lambda_d(x)$.
Nierówność Cauchy'ego--Schwarza daje
$\|x_jf\|_2\|\partial_jf\|_2\ge\tfrac12\|f\|_2^2$.
Z twierdzenia Plancherela i
$\widehat{\partial_jf}=i\xi_j\widehat f$ wynika
$\|\partial_jf\|_2=(2\pi)^{-d/2}\|\xi_j\widehat f\|_2$
oraz $\|f\|_2=(2\pi)^{-d/2}\|\widehat f\|_2$.
Po podstawieniu otrzymujemy tezę.
\end{proof}

\section{Podsumowanie i zakres dalszych zagadnień}
Transformata Fouriera jest automorfizmem przestrzeni Schwartza i
jej przestrzeni dualnej. Jest podstawowym narzędziem badania wielu
liniowych PDE o stałych współczynnikach, lecz sama nie zapewnia
rozwiązywalności ani jednoznaczności w dowolnej klasie funkcji;
nie rozwiązuje też automatycznie problemów brzegowych ani równań
ze zmiennymi współczynnikami. Mnożniki ograniczone działają na $L^2$,
lecz dla $L^p$ z $p\ne2$ potrzebna jest głębsza teoria.
W następnych rozdziałach zbudujemy systematycznie
przestrzenie Sobolewa na obszarach i ich twierdzenia o zanurzeniach.
W ten sposób domkniemy łańcuch pojęciowy rozpoczęty w
rozdziale~\ref{chap:pde-distributions}: od słabych pochodnych,
przez narzędzia Fouriera na całej przestrzeni, do funkcjonałowej
teorii $W^{k,p}(\Omega)$ potrzebnej dla metod wariacyjnych i
zagadnień brzegowych.




\chapter[Przestrzenie Sobolewa]{Przestrzenie Sobolewa: słabe pochodne, aproksymacja i przedłużanie}
\label{chap:pde-sobolev}

Na prostej funkcja $u(x)=|x|$ nie jest różniczkowalna w zerze,
lecz jej pochodna w sensie dystrybucji jest funkcją
$\operatorname{sgn}x$. Z drugiej strony pochodna funkcji
Heaviside'a $H=\chi_{(0,\infty)}$ jest deltą Diraca,
a nie funkcją lokalnie całkowalną (sekcja
~\ref{sec:pde7-derivatives}). Ta różnica motywuje przestrzenie Sobolewa: dopuszczamy
niegładkie funkcje, ale wymagamy, aby ich pochodne dystrybucyjne
były reprezentowane przez funkcje z określonych klas całkowalności.

Niech w całym rozdziale $d\ge1$ będzie liczbą
całkowitą, $k\in\N_0$, $\Omega\subset\R^d$ niepustym zbiorem
otwartym, a $1\le p\le\infty$. Nie zakładamy spójności $\Omega$,
o ile nie wskażemy jej jawnie. Chyba że wyraźnie zaznaczono
inaczej, funkcje są rzeczywiste. Elementy przestrzeni $L^p$
są klasami funkcji równych prawie wszędzie. Pochodne rozumiemy
w sensie dystrybucji, a ich reprezentanty funkcyjne i równości
w $L^p$ — z dokładnością do zbiorów miary zero. Wszystkie całki na $\Omega$ są liczone
względem miary Lebesgue'a $\lambda_d$. Definicje dystrybucji i
pochodnych słabych znajdują się w rozdziale~\ref{chap:pde-distributions},
a przestrzenie $H^s(\R^d)$ w sekcji~\ref{sec:pde8-bessel}.

\section{Definicje i pierwszy słownik}
\label{sec:pde9-def}

\begin{definicja}[Przestrzenie $W^{k,p}$]
\label{def:pde9-Wkp}
Dla $k\in\N_0$ i $1\le p\le\infty$ definiujemy
\[
 W^{k,p}(\Omega)=\{u\in L^p(\Omega):
  \partial^\alpha u\in L^p(\Omega)
  \text{ dla każdego }\alpha\in\N_0^d,\ |\alpha|\le k\},
\]
gdzie $\partial^\alpha u$ oznacza pochodną dystrybucyjną.
Każda funkcja $u\in L^p(\Omega)$ jest lokalnie całkowalna:
dla zwartego $K\subset\Omega$ nierówność Höldera daje
$\int_K |u|\,\mathrm d\lambda_d
\le\lambda_d(K)^{1-1/p}\|u\|_p$ dla $1<p<\infty$;
dla $p=1$ wystarcza całkowalność globalna,
a dla $p=\infty$ oszacowanie
$\lambda_d(K)\|u\|_\infty$. Definiuje więc dystrybucję regularną.
Wyrażenie $\partial^\alpha u\in L^p$ oznacza istnienie
funkcji $v_\alpha\in L^p(\Omega)$ spełniającej
\[
 \int_\Omega u(x)\partial^\alpha\varphi(x)\,\mathrm d\lambda_d(x)
 =(-1)^{|\alpha|}\int_\Omega v_\alpha(x)\varphi(x)\,
 \mathrm d\lambda_d(x)
 \qquad(\varphi\in C_c^\infty(\Omega)).
\]
Zgodnie z definicją~\ref{def:pde7-weak} i jednoznacznością
reprezentantów dystrybucji regularnych (sekcja~\ref{sec:pde7-def})
$v_\alpha$ jest wyznaczona prawie wszędzie.
Dla $\alpha=0$ mamy $v_0=u$. Normę określamy wzorami
\begin{align*}
 \|u\|_{W^{k,p}(\Omega)}
  &=\left(\sum_{|\alpha|\le k}\|\partial^\alpha u\|_{L^p(\Omega)}^p
    \right)^{1/p},&&1\le p<\infty,\\
 \|u\|_{W^{k,\infty}(\Omega)}
  &=\max_{|\alpha|\le k}\|\partial^\alpha u\|_{L^\infty(\Omega)}.
\end{align*}
Dla $k=0$ otrzymujemy $W^{0,p}=L^p$, z równością norm.
Piszemy $H^k(\Omega)=W^{k,2}(\Omega)$,
a w szczególności $H^1(\Omega)=W^{1,2}(\Omega)$.
\end{definicja}

Liniowość pochodnych dystrybucyjnych oraz przestrzeni $L^p$
pokazuje, że $W^{k,p}(\Omega)$ jest przestrzenią wektorową.
Niech $M=\#\{\alpha\in\N_0^d:|\alpha|\le k\}$.
Odwzorowanie $u\mapsto(\partial^\alpha u)_{|\alpha|\le k}$
jest liniowe i iniektywne, gdyż zawiera współrzędną $u$.
Dla $p<\infty$ podana norma jest normą $\ell^p$
wektora norm składowych, a dla $p=\infty$ ich maksimum.
Nierówność trójkąta wynika z nierówności trójkąta w $L^p$
dla każdej składowej, a następnie z nierówności Minkowskiego
dla skończonych ciągów (dla $p=\infty$ z nierówności dla maksimum).
Jednorodność wynika z liniowości pochodnych, a zerowa norma
pociąga $\|u\|_p=0$. Zatem istotnie otrzymujemy normę.
Zupełność wykażemy w następnej sekcji.

\begin{definicja}[Przestrzeń lokalna i domknięcie funkcji testowych]
\label{def:pde9-local-zero}
Niech $k\in\N_0$, $1\le p\le\infty$.
Klasa funkcji mierzalnych $u$ należy do
$W^{k,p}_{\mathrm{loc}}(\Omega)$, jeśli
$u|_U\in W^{k,p}(U)$ dla każdego niepustego otwartego $U$
takiego, że $\overline U$ jest zwartym podzbiorem $\Omega$
(zapisujemy $U\Subset\Omega$).
Dla $1\le p<\infty$ definiujemy
\[
 W^{k,p}_0(\Omega)=\overline{C_c^\infty(\Omega)}^{\,W^{k,p}(\Omega)}.
\]
Indeks $0$ oznacza w tej chwili \emph{definicję przez domknięcie},
a nie jeszcze twierdzenie o istnieniu śladu na brzegu.
W tym rozdziale zapis $W^{k,p}_0$ stosujemy tylko dla $p<\infty$;
nie przenosimy twierdzeń o gęstości w tych normach na $p=\infty$.
\end{definicja}

Ograniczenie funkcji z $W^{k,p}(\Omega)$ do otwartego
$U\subset\Omega$ należy do $W^{k,p}(U)$ i
$\partial^\alpha(u|_U)=(\partial^\alpha u)|_U$:
test z $C_c^\infty(U)$ przedłużony przez zero jest testem
na $\Omega$, więc spełnia tę samą tożsamość całkową.
Norma po ograniczeniu nie rośnie, ponieważ całkujemy
na podzbiorze (lub bierzemy istotne supremum na podzbiorze).
W szczególności $W^{k,p}(\Omega)\subset W^{k,p}_{\mathrm{loc}}(\Omega)$.

W definicji lokalnej reprezentanty pochodnych z dwóch
takich otwartych zbiorów zgadzają się prawie wszędzie
na ich przecięciu, dzięki jednoznaczności dystrybucji regularnych.
Wybierzmy przeliczalne pokrycie $(B_n)_{n\ge1}$
zbioru $\Omega$ kulami o domknięciach zwartych w $\Omega$
(można użyć kul o wymiernych środkach i promieniach).
Niech $v_{\alpha,n}$ oznacza mierzalnego reprezentanta
pochodnej na $B_n$. Na mierzalnych, parami rozłącznych
zbiorach $A_n=B_n\setminus\bigcup_{j<n}B_j$
określamy $v_\alpha=v_{\alpha,n}$.
Zbiory $A_n$ pokrywają $\Omega$, więc $v_\alpha$ jest
mierzalna. Dla ustalonego $m$ na każdym $A_n\cap B_m$
zgadza się prawie wszędzie z $v_{\alpha,m}$, a
przeliczalna suma zbiorów wyjątkowych ma miarę zero.
Wobec tego $v_\alpha|_{B_m}=v_{\alpha,m}$ prawie wszędzie.
Każdy zwarty podzbiór $\Omega$ jest zawarty w skończonej
sumie tych kul, co zapewnia lokalną przynależność do $L^p$,
również dla $p=\infty$ przez maksimum skończenie wielu norm.
Dla dowolnego testu na $\Omega$ wybieramy otwarty zbiór
$U\Subset\Omega$ zawierający jego nośnik (na przykład
skończoną sumę tych kul). Reprezentant słabej pochodnej
na $U$ zgadza się prawie wszędzie z $v_\alpha$
na każdym $U\cap B_n$, więc również na całym $U$.
Tożsamość testowa na $U$ zachodzi z $v_\alpha$
i jest zarazem tożsamością na $\Omega$, ponieważ
test i jego pochodne znikają poza $U$. Zatem $v_\alpha$ reprezentuje
$\partial^\alpha u$ na całej $\Omega$.
Nie wymaga to globalnej całkowalności na $\Omega$.

\begin{uwaga}[Wybór normy]
Dla $M$ składowych i $1\le p<\infty$ nierówność
Höldera dla skończonych ciągów daje
\[
 \|u\|_{W^{k,p}}\le\sum_{|\alpha|\le k}\|\partial^\alpha u\|_p
 \le M^{1-1/p}\|u\|_{W^{k,p}}.
\]
Pierwsza nierówność wynika z tego, że dla nieujemnych
$a_i$, przy $S=\sum_i a_i>0$, mamy
$0\le a_i/S\le1$, więc $(a_i/S)^p\le a_i/S$.
Sumując i mnożąc przez $S^p$, otrzymujemy
$\sum_i a_i^p\le S^p$; dla $S=0$ jest to równość.
Dla $p=\infty$ analogicznie maksimum jest nie większe
od sumy, a suma nie przekracza $M$ razy maksimum.
Normy te są więc równoważne.
Przy $p=2$ norma $H^1(\Omega)$ spełnia
$\|u\|_{H^1}^2=\|u\|_2^2+\|\nabla u\|_2^2$.
Zgodnie z twierdzeniem~\ref{thm:pde8-Hs}, na całej $\R^d$
przestrzenie $H^k$ zdefiniowane przez słabe pochodne pokrywają się
z przestrzeniami Bessela z rozdziału~\ref{chap:pde-fourier-Rd};
normy są równoważne, a dla $k=1$ (oraz $k=0$)
są równe. Istotnie, dla $u\in W^{1,2}(\R^d)$ reguła
różniczkowania transformaty i twierdzenie Plancherela
z rozdziału~\ref{chap:pde-fourier-Rd} dają
\[
 \begin{aligned}
 \|u\|_{W^{1,2}}^2
 &=\|u\|_2^2+\sum_{j=1}^d\|\partial_j u\|_2^2\\
 &=(2\pi)^{-d}\int_{\R^d}(1+|\xi|^2)|\widehat u(\xi)|^2
 \,\mathrm d\lambda_d(\xi).
 \end{aligned}
\]
Prawa strona jest kwadratem normy Bessela $H^1$.
\end{uwaga}

\begin{przyklad}[Załamanie nie jest skokiem]
Niech $I=(-1,1)$. Funkcja $u(x)=|x|$ należy do $W^{1,p}(I)$
dla każdego $1\le p\le\infty$, ponieważ
$u'=\operatorname{sgn}x$ prawie wszędzie i pochodna ta jest
jednocześnie pochodną dystrybucyjną, jak wykazano w sekcji
~\ref{sec:pde7-derivatives}. Obie funkcje są ograniczone na $I$
o skończonej mierze, zatem należą do każdego z rozważanych
$L^p(I)$. Natomiast funkcja
$H=\chi_{(0,1)}$ nie należy do $W^{1,p}(I)$ dla żadnego
$1\le p\le\infty$: jej pochodna dystrybucyjna to $\delta_0$
(przykład~\ref{ex:pde7-heaviside}). Delta nie jest reprezentowana
nawet przez funkcję lokalnie całkowalną na $I$:
dowód z sekcji~\ref{sec:pde7-def}, wykorzystujący testy
o nośnikach kurczących się do zera, działa również na $I$.
W szczególności nie może być reprezentowana przez funkcję z $L^p(I)$.
Ponadto $u(x)=|x|$ nie należy do $W^{2,p}(I)$ dla żadnego
$1\le p\le\infty$, ponieważ $\partial^2 u=2\delta_0$
(co wynika z $\operatorname{sgn}x=2H(x)-1$ prawie wszędzie).
\end{przyklad}

\begin{przyklad}[Osobliwość potęgowa]
\label{ex:pde9-power}
Niech $d\ge2$, $0<a<d$, a $\psi\in C_c^\infty(\R^d)$
będzie równa $1$ w otoczeniu zera. Niech $u(x)=\psi(x)|x|^{-a}$
dla $x\ne0$ i $u(0)=0$ (wartość w zerze nie wpływa
na klasę funkcji). Dla $1\le p<\infty$ zachodzi
\[
 u\in L^p(\R^d)\iff ap<d,\qquad
 u\in W^{1,p}(\R^d)\iff (a+1)p<d.
\]
Dla $p=\infty$ funkcja $u$ nie należy nawet do $L^\infty$.

Istotnie, wybieramy $\delta>0$ tak, aby $\psi=1$ na
$B_\delta(0)$. Poza tą kulą funkcja i jej klasyczne pochodne
są gładkie, ograniczone i mają zwarty nośnik w rozważanym
obszarze, więc nie sprawiają problemu z całkowalnością.
Na kuli całkowanie radialne daje
\[
 \int_{B_\delta(0)} |u(x)|^p\,\mathrm d\lambda_d(x)
 =\omega_d\int_0^\delta r^{d-1-ap}\,\mathrm d\lambda_1(r),
\]
gdzie $\omega_d>0$ jest polem powierzchni sfery jednostkowej.
Istotnie, dla $0<\varepsilon<\delta$ całka
$\int_\varepsilon^\delta r^{d-1-ap}\,\mathrm d\lambda_1(r)$
jest równa $(\delta^{d-ap}-\varepsilon^{d-ap})/(d-ap)$
dla $ap\ne d$, a $\log(\delta/\varepsilon)$ dla $ap=d$.
Granica przy $\varepsilon\downarrow0$ jest skończona
dokładnie dla $ap<d$. Przy $ap=d$ rozbieżność jest
logarytmiczna, a przy $ap>d$ potęgowa. Ponadto na przekłutej kuli
$|u(x)|=|x|^{-a}$ przekracza dowolne ustalone dodatnie
ograniczenie na kuli o dodatniej mierze, co wyklucza $L^\infty$.

Dla $x\ne0$ klasyczny gradient jest równy
\[
 v(x)=|x|^{-a}\nabla\psi(x)-a\psi(x)|x|^{-a-2}x,
 \qquad |v(x)|=a|x|^{-a-1}\quad(0<|x|<\delta).
\]
Możemy przyjąć $v(0)=0$. Wszystkie składowe $v$ należą
do $L^p$ dokładnie wtedy, gdy $|v|\in L^p$:
korzystamy z $|v_j|\le |v|\le\sum_j|v_j|$ oraz
nierówności trójkąta w $L^p$. Na kuli
\[
 \int_{B_\delta(0)} |v(x)|^p\,\mathrm d\lambda_d(x)
 =a^p\omega_d\int_0^\delta r^{d-1-(a+1)p}\,
 \mathrm d\lambda_1(r).
\]
Stosując poprzedni rachunek z $a+1$ zamiast $a$,
otrzymujemy warunek skończoności $(a+1)p<d$;
na granicy ponownie występuje rozbieżność logarytmiczna.

Załóżmy $(a+1)p<d$. Wówczas $u,v_j\in L^p$ oraz
$a<d-1$ (ponieważ $p\ge1$). Dla dowolnego
$\varphi\in C_c^\infty(\R^d)$ wybieramy $R>\delta$
tak duże, aby jej nośnik był zawarty w $B_R(0)$.
Na pierścieniu $B_R(0)\setminus\overline{B_\varepsilon(0)}$,
$0<\varepsilon<\delta$, całkowanie przez części daje
\[
 \int_{B_R\setminus\overline{B_\varepsilon}}
 u\,\partial_j\varphi\,\mathrm d\lambda_d
 =-\int_{B_R\setminus\overline{B_\varepsilon}}
 v_j\varphi\,\mathrm d\lambda_d
 +\int_{\partial B_\varepsilon}
 u\varphi\,\nu_j\,\mathrm d\sigma,
\]
gdzie $\nu=-x/|x|$ jest normalną zewnętrzną pierścienia
na wewnętrznej sferze, a $\sigma$ jest miarą powierzchniową
tej sfery. Na zewnętrznej sferze test znika.
Moduł składnika brzegowego jest nie większy niż
$\|\varphi\|_\infty\omega_d\varepsilon^{d-1-a}$,
co dąży do zera. Obie całki objętościowe zbiegają do
całek na $\R^d$ z twierdzenia o zbieżności zmajoryzowanej:
majorantami są odpowiednio
$\|\partial_j\varphi\|_\infty |u|\chi_{B_R(0)}$
i $\|\varphi\|_\infty |v_j|\chi_{B_R(0)}$,
które należą do $L^1$ dzięki lokalnej całkowalności
(lub oszacowaniu Höldera na kuli).
Stąd $\partial_j u=v_j$ w sensie słabym i $u\in W^{1,p}$.

Odwrotnie, gdyby $u\in W^{1,p}$, to ograniczenie jej
słabych pochodnych do $\R^d\setminus\{0\}$ musiałoby
zgadzać się z klasycznymi pochodnymi $v_j$ prawie
wszędzie. Ponieważ singleton ma miarę zero, mielibyśmy
$v_j\in L^p$ dla wszystkich $j$, a więc również $|v|\in L^p$.
Obliczona całka radialna wymusza $(a+1)p<d$.
\end{przyklad}

\section{Zupełność i zamkniętość operatorów pochodnych}
\label{sec:pde9-complete}

\begin{twierdzenie}[Zupełność przestrzeni Sobolewa]
\label{thm:pde9-banach}
Dla każdego niepustego otwartego $\Omega\subset\R^d$, $k\in\N_0$ i
$1\le p\le\infty$ przestrzeń $W^{k,p}(\Omega)$ jest przestrzenią
Banacha. W szczególności $H^k(\Omega)$ jest przestrzenią Hilberta
z iloczynem skalarnym
\[
 \langle u,v\rangle_{H^k}=\sum_{|\alpha|\le k}
 \int_\Omega\partial^\alpha u(x)\,\overline{\partial^\alpha v(x)}
 \,\mathrm d\lambda_d(x).
\]
W przypadku zespolonym przyjmujemy liniowość względem
pierwszego argumentu; dla funkcji rzeczywistych sprzężenie
można pominąć.
\end{twierdzenie}
\begin{proof}
Niech $(u_n)$ będzie ciągiem Cauchy'ego w $W^{k,p}(\Omega)$.
Dla każdego $|\alpha|\le k$ zachodzi
$\|\partial^\alpha(u_n-u_m)\|_p\le\|u_n-u_m\|_{W^{k,p}}$,
także dla $p=\infty$. Zatem $(\partial^\alpha u_n)$
jest ciągiem Cauchy'ego w przestrzeni Banacha $L^p(\Omega)$.
Istnieje $f_\alpha\in L^p(\Omega)$ takie, że
$\partial^\alpha u_n\to f_\alpha$ w $L^p$. Oznaczmy $u=f_0$.
Dla dowolnego $\varphi\in C_c^\infty(\Omega)$ mamy
\[
 \int_\Omega u_n\partial^\alpha\varphi\,\mathrm d\lambda_d
 =(-1)^{|\alpha|}\int_\Omega
 (\partial^\alpha u_n)\varphi\,\mathrm d\lambda_d.
\]
Niech $p'$ będzie wykładnikiem sprzężonym do $p$,
z konwencją $1/\infty=0$. Test oraz każda jego pochodna
są ograniczone i mają zwarty nośnik, więc należą do $L^{p'}$,
również gdy $p'=1$ lub $p'=\infty$. Z nierówności Höldera
otrzymujemy
\[
 \begin{aligned}
 \left|\int_\Omega(u_n-u)\partial^\alpha\varphi
 \,\mathrm d\lambda_d\right|
 &\le\|u_n-u\|_p\|\partial^\alpha\varphi\|_{p'}\to0,\\
 \left|\int_\Omega(\partial^\alpha u_n-f_\alpha)\varphi
 \,\mathrm d\lambda_d\right|
 &\le\|\partial^\alpha u_n-f_\alpha\|_p\|\varphi\|_{p'}\to0.
 \end{aligned}
\]
Przejście do granicy daje dokładnie tożsamość testową dla
$\partial^\alpha u=f_\alpha$ w $\mathcal D'(\Omega)$.
Ponieważ wszystkie $f_\alpha$ należą do $L^p$, mamy
$u\in W^{k,p}$. Dla $p<\infty$ zachodzi
\[
 \|u_n-u\|_{W^{k,p}}^p
 =\sum_{|\alpha|\le k}\|\partial^\alpha u_n-f_\alpha\|_p^p\to0,
\]
gdyż suma jest skończona. Dla $p=\infty$ maksimum skończenie
wielu norm dążących do zera także dąży do zera.
Dowodzi to zupełności.

Dla $p=2$ każda całka w podanym iloczynie skalarnym
jest bezwzględnie zbieżna z nierówności Cauchy'ego--Schwarza.
Liniowość pochodnych i własności iloczynu skalarnego $L^2$
dają liniowość w pierwszym argumencie, antyliniowość
w drugim i symetrię hermitowską. Ponadto
$\langle u,u\rangle_{H^k}=\|u\|_{W^{k,2}}^2$,
a równość z zerem daje $u=0$ w $L^2$, ponieważ suma zawiera
składnik $\alpha=0$. Zupełność wykazana powyżej
dowodzi, że jest to przestrzeń Hilberta.
\end{proof}

\begin{wniosek}[Zamkniętość pochodnej]
\label{cor:pde9-closed}
Dla ustalonego wielowskaźnika $\alpha$ określmy operator
$D_\alpha u=\partial^\alpha u$ w przestrzeni $L^p(\Omega)$
na jego naturalnej dziedzinie
\[
 \mathcal D(D_\alpha)=
 \{u\in L^p(\Omega):\partial^\alpha u
 \text{ jest reprezentowana przez funkcję z }L^p(\Omega)\}.
\]
Operator $D_\alpha$ jest zamknięty: jeżeli
$u_n\in\mathcal D(D_\alpha)$, $u_n\to u$ w $L^p$ oraz
$D_\alpha u_n\to g$ w $L^p$, to $u\in\mathcal D(D_\alpha)$
i $D_\alpha u=g$.
Nie zakładamy istnienia ani zbieżności pozostałych
pochodnych rzędu $|\alpha|$ w $L^p$.
\end{wniosek}
\begin{proof}
Dla każdego testu $\varphi$ mamy
\[
 \int_\Omega u_n\partial^\alpha\varphi\,\mathrm d\lambda_d
 =(-1)^{|\alpha|}\int_\Omega(D_\alpha u_n)\varphi\,\mathrm d\lambda_d.
\]
Oszacowania Höldera z dowodu twierdzenia
~\ref{thm:pde9-banach}, zastosowane do $u_n-u$ oraz
$D_\alpha u_n-g$, pozwalają przejść do granicy po obu stronach.
Otrzymujemy $\partial^\alpha u=g$ w sensie dystrybucji,
co daje dokładnie żądaną przynależność do dziedziny i wartość
operatora. W normowanej przestrzeni $L^p\times L^p$
domkniętość wykresu jest równoważna tej własności ciągowej.
\end{proof}

\begin{uwaga}[Dziedzina pojedynczej pochodnej a przestrzeń Sobolewa]
Zamkniętość pojedynczej pochodnej nie oznacza, że granica
należy do całej przestrzeni $W^{|\alpha|,p}$.
Na przykład na $\Omega=(0,1)^2$ funkcja
$u(x,y)=\chi_{(1/2,1)}(y)$ należy do każdego $L^p(\Omega)$
i ma pochodną $\partial_x u=0$, ponieważ całka z
$\partial_x\varphi(x,y)$ po $x\in(0,1)$ jest zerowa dla każdego
$y$. Jednak $\partial_y u$ działa przez
\[
 \langle\partial_y u,\varphi\rangle
 =\int_0^1\varphi(x,1/2)\,\mathrm d\lambda_1(x).
\]
Nie jest to dystrybucja regularna. Istotnie, wybieramy
$\eta\in C_c^\infty(0,1)$ o całce $1$ i
$\psi\in C_c^\infty(-1,1)$ z $\psi(0)=1$.
Testy $\varphi_\varepsilon(x,y)=\eta(x)\psi((y-1/2)/\varepsilon)$
dla $0<\varepsilon<1/4$ mają jednostajnie ograniczone moduły,
wspólny zwarty nośnik i nośniki zawarte w pasach,
których miary dążą do zera. Działanie $\partial_y u$
na tych testach jest stale równe $1$, natomiast działanie
każdej funkcji lokalnie całkowalnej dążyłoby do zera
z absolutnej ciągłości całki Lebesgue'a.
Stąd $u\notin W^{1,p}(\Omega)$.
Dla $p<\infty$ można też aproksymować $u$ w $L^p$
gładkimi funkcjami zależnymi tylko od $y$ i mającymi
$\partial_x u_n=0$: wystarczy zastąpić skok gładkim
przejściem z $0$ do $1$ na przedziale
$(1/2-1/n,1/2+1/n)$, dla $n\ge3$.
Błąd w $L^p$ jest nie większy niż $(2/n)^{1/p}$.
To pokazuje, dlaczego we wniosku używamy naturalnej dziedziny,
a nie ograniczenia operatora do całej przestrzeni $W^{1,p}$.

Sama norma Sobolewa nie przypisuje klasie funkcji
kanonicznych wartości punktowych. Warunki istnienia
ciągłych reprezentantów wynikają z twierdzeń o zanurzeniach
Sobolewa, omawianych w następnym rozdziale.
\end{uwaga}


\section{Wymiar jeden: absolutna ciągłość}
\label{sec:pde9-oneD}

\begin{twierdzenie}[Charakterystyka $W^{1,p}(a,b)$]
\label{thm:pde9-ac}
Niech $-\infty<a<b<\infty$ i $1\le p\le\infty$.
Klasa $u\in L^p(a,b)$ należy do $W^{1,p}(a,b)$ wtedy i tylko
wtedy, gdy ma reprezentanta absolutnie ciągłego
$\widetilde u:[a,b]\to\R$, którego pochodna klasyczna
istnieje prawie wszędzie i należy do $L^p(a,b)$.
Reprezentant ten jest jedyny wśród funkcji ciągłych na $[a,b]$.
Jego pochodna klasyczna jest słabą pochodną $u'$ prawie wszędzie
i zachodzi
\[
 \widetilde u(y)-\widetilde u(x)
 =\int_x^y u'(t)\,\mathrm d\lambda_1(t)
 \qquad(a\le x\le y\le b).
\]
\end{twierdzenie}
\begin{proof}
Załóżmy $u\in W^{1,p}(a,b)$ i niech $g=u'$ będzie jego
słabą pochodną. Ponieważ przedział ma skończoną miarę $L=b-a$,
mamy $\|g\|_1\le L^{1-1/p}\|g\|_p$ dla $p<\infty$
oraz $\|g\|_1\le L\|g\|_\infty$ dla $p=\infty$.
Zdefiniujmy na $[a,b]$
\[
 U(x)=\int_a^x g(t)\,\mathrm d\lambda_1(t).
\]
Z podstawowego twierdzenia rachunku całkowego dla całki
Lebesgue'a wynika, że $U$ jest absolutnie ciągła i
$U'=g$ prawie wszędzie. Dla przypomnienia, absolutna
ciągłość wynika z oszacowania sumy przyrostów $U$
na parami rozłącznych przedziałach przez całkę z $|g|$
na ich sumie oraz z absolutnej ciągłości tej całki.
Dla testu $\varphi\in C_c^\infty(a,b)$ całkowanie przez
części dla funkcji absolutnie ciągłych daje
\[
 \int_a^b U(t)\varphi'(t)\,\mathrm d\lambda_1(t)
 =-\int_a^b g(t)\varphi(t)\,\mathrm d\lambda_1(t).
\]
Można zastosować je do iloczynu $U\varphi$, który jest
absolutnie ciągły, ma pochodną $g\varphi+U\varphi'$
prawie wszędzie i znika w obu końcach.
Stąd dystrybucyjna pochodna $u-U$ jest zerowa.
Twierdzenie~\ref{thm:pde7-constant} zastosowane na spójnym
przedziale $(a,b)$ daje $T_{u-U}=T_C$ dla pewnej stałej $C$.
Ponieważ $U$ jest ciągła na zwartym przedziale,
jest ograniczona, więc należy do $L^p(a,b)$.
Jednoznaczność reprezentanta dystrybucji regularnej
daje $u-U=C$ prawie wszędzie; dla funkcji rzeczywistej $C$
jest rzeczywista. Wobec tego $\widetilde u=U+C$ jest
żądanym reprezentantem i ma pochodną $g$.
Dla $a\le x\le y\le b$ addytywność całki daje podany wzór.

Odwrotnie, jeśli $u$ ma reprezentanta absolutnie ciągłego
$\widetilde u$ z pochodną $g\in L^p$, to iloczyn
$\widetilde u\varphi$ jest absolutnie ciągły, ma pochodną
$g\varphi+\widetilde u\varphi'$ prawie wszędzie
i znika w $a,b$. Całka z tej pochodnej jest zatem równa zeru.
Ponieważ $u=\widetilde u$ prawie wszędzie,
otrzymujemy $\int_a^b u\varphi'\,\mathrm d\lambda_1
=-\int_a^b g\varphi\,\mathrm d\lambda_1$.
To jest definicja słabej pochodnej z $L^p$.
Jeśli dwa ciągłe reprezentanty tej samej klasy różniłyby się
w dowolnym punkcie $[a,b]$, ich różnica byłaby niezerowa
na względnie otwartym otoczeniu tego punktu, którego przecięcie
z $(a,b)$ ma dodatnią miarę. Przeczy to równości prawie wszędzie.
\end{proof}

\begin{wniosek}[Oszacowanie punktowe i ślady w jednym wymiarze]
\label{cor:pde9-trace-oneD}
Dla $I=(a,b)$, $L=b-a$ i $u\in W^{1,p}(I)$,
$1\le p<\infty$, jedyny ciągły reprezentant spełnia
\[
 \|\widetilde u\|_{C([a,b])}\le
 L^{-1/p}\|u\|_{L^p(I)}+L^{1-1/p}\|u'\|_{L^p(I)}.
\]
Dla $p=\infty$ mamy
$\|\widetilde u\|_{C([a,b])}=\|u\|_\infty
\le\|u\|_{W^{1,\infty}}$.
W szczególności odwzorowania $u\mapsto\widetilde u(a)$
i $u\mapsto\widetilde u(b)$ są ciągłymi funkcjonałami
liniowymi na $W^{1,p}(I)$.
\end{wniosek}
\begin{proof}
Dla ciągłego reprezentanta supremum na $[a,b]$ jest równy
istotnemu supremum na $(a,b)$: dowolna wartość
ściśle mniejsza od maksimum modułu jest przekroczona
na zbiorze dodatniej miary dzięki ciągłości, także gdy
maksimum jest osiągnięte na końcu przedziału.
Dla $p<\infty$ wybierzmy $x_0\in[a,b]$,
w którym $|\widetilde u|$ osiąga minimum.
Mamy $L|\widetilde u(x_0)|^p\le\|u\|_p^p$.
Z twierdzenia~\ref{thm:pde9-ac}, dla dowolnego $x\in[a,b]$,
\[
 |\widetilde u(x)-\widetilde u(x_0)|
 \le\int_{\min(x,x_0)}^{\max(x,x_0)}|u'(t)|
 \,\mathrm d\lambda_1(t)
 \le L^{1-1/p}\|u'\|_p.
\]
Dodając oszacowanie minimum i biorąc supremum,
otrzymujemy tezę. Dla $p=\infty$ korzystamy z równości
supremum i istotnego supremum oraz definicji normy.
Jednoznaczność reprezentanta zapewnia liniowość
odwzorowań ewaluacji: reprezentant kombinacji liniowej
jest kombinacją liniową reprezentantów.
Dla $p<\infty$ moduł każdej ewaluacji jest ograniczony
przez $(L^{-1/p}+L^{1-1/p})\|u\|_{W^{1,p}}$,
a dla $p=\infty$ przez $\|u\|_{W^{1,\infty}}$.
\end{proof}

\begin{wniosek}[Moduł ciągłości i przypadek Lipschitza]
Dla $u\in W^{1,p}(a,b)$ i $1<p<\infty$ zachodzi
\[
 |\widetilde u(y)-\widetilde u(x)|
 \le\|u'\|_p|y-x|^{1-1/p}\qquad(x,y\in[a,b]).
\]
Dla $p=\infty$ reprezentant jest lipschitzowski ze stałą
co najwyżej $\|u'\|_\infty$. Odwrotnie, każda funkcja
lipschitzowska na $[a,b]$ wyznacza element $W^{1,\infty}(a,b)$.
Dla $p=1$ otrzymujemy absolutną ciągłość, ale nie wynika
stąd żaden ustalony dodatni wykładnik Höldera.
\end{wniosek}
\begin{proof}
Dla $x=y$ nierówność jest natychmiastowa; zamiana
$x,y$ nie zmienia modułu różnicy. Dla $x<y$
stosujemy do wzoru z twierdzenia
~\ref{thm:pde9-ac} nierówność Höldera na $(x,y)$:
$\int_x^y|u'|\,\mathrm d\lambda_1
\le |y-x|^{1-1/p}\|u'\|_{L^p(x,y)}$.
Dla $p=\infty$ całka jest ograniczona przez
$|y-x|\|u'\|_\infty$. Funkcja lipschitzowska jest absolutnie
ciągła, gdyż suma modułów jej przyrostów jest ograniczona
przez stałą Lipschitza razy sumę długości przedziałów.
Jej pochodna istnieje zatem prawie wszędzie
i jest ograniczona przez tę stałą z oszacowania ilorazów
różnicowych; sama funkcja jest ograniczona na $[a,b]$.
Charakterystyka z twierdzenia daje przynależność
do $W^{1,\infty}$. W szczególności najlepsza stała
Lipschitza reprezentanta jest równa $\|u'\|_\infty$.

Dla $p=1$ wzór całkowy daje kontrolę przez
$\int_x^y|u'|\,\mathrm d\lambda_1$, która dąży do zera
jednostajnie względem przedziałów o długości dążącej do zera.
Na przykład na $(0,1)$ dla dowolnego $\gamma>0$ wybierzmy
$0<\theta<\min\{1,\gamma\}$. Funkcja $u(x)=x^\theta$
jest absolutnie ciągła, gdyż jest całką z
$\theta t^{\theta-1}\in L^1(0,1)$ od $0$ do $x$.
Należy do $W^{1,1}(0,1)$, lecz
$|u(x)-u(0)|/x^\gamma=x^{\theta-\gamma}\to+\infty$
przy $x\downarrow0$. Nie ma więc wspólnego dodatniego
wykładnika Höldera dla całej przestrzeni $W^{1,1}$.
\end{proof}

\begin{przyklad}[Funkcja Cantora a pochodna słaba]
Klasyczna funkcja Cantora $F$ na $[0,1]$ jest ciągła,
niemalejąca, spełnia $F(0)=0$, $F(1)=1$ i jest stała na każdym
przedziale usuniętym w konstrukcji zbioru Cantora.
Zbiór Cantora ma miarę zero, ponieważ po $n$ krokach
jest zawarty w $2^n$ przedziałach o długości $3^{-n}$,
o łącznej długości $(2/3)^n$.
Wobec tego klasyczna pochodna $F'$ istnieje i jest równa zeru
poza zbiorem miary zero. Gdyby $F$ należała do $W^{1,1}(0,1)$,
twierdzenie~\ref{thm:pde9-ac} i jednoznaczność ciągłego
reprezentanta oznaczałyby, że jest absolutnie ciągła
i jej słaba pochodna jest zerowa prawie wszędzie.
Wzór całkowy dawałby wówczas $F(1)-F(0)=0$,
wbrew $F(1)-F(0)=1$.
Zatem pochodna istniejąca prawie wszędzie nie musi być
pochodną słabą. Ponieważ na $(0,1)$ zachodzi
$W^{1,p}\subset W^{1,1}$ dla wszystkich $1\le p\le\infty$
przez nierówność Höldera, $F$ nie należy do żadnej z tych
przestrzeni Sobolewa.
\end{przyklad}

\section{Sploty i aproksymacja wewnętrzna}
\label{sec:pde9-mollify}

Ustalmy nieujemną funkcję $\rho\in C_c^\infty(B_1(0))$
o całce $1$ i połóżmy
$\rho_\varepsilon(x)=\varepsilon^{-d}\rho(x/\varepsilon)$,
jak w sekcji~\ref{sec:pde7-mollifiers}.
Niech $\Omega_\varepsilon=\{x\in\Omega:
\operatorname{dist}(x,\R^d\setminus\Omega)>\varepsilon\}$,
przy czym odległość od zbioru pustego jest nieskończona.
Dla $u\in L^1_{\mathrm{loc}}(\Omega)$ i $x\in\Omega_\varepsilon$
definiujemy
\[
 u_\varepsilon(x)=(\rho_\varepsilon*u)(x)
 =\int_\Omega \rho_\varepsilon(x-y)u(y)\,\mathrm d\lambda_d(y).
\]
Całka ma sens, ponieważ jądro jest ograniczone
i jako funkcja zmiennej $y$ ma zwarty nośnik w $\Omega$.
Nie wymagamy przedłużenia $u$ przez zero na całej przestrzeni.

\begin{twierdzenie}[Lokalne wygładzanie Sobolewa]
\label{thm:pde9-mollify}
Niech $1\le p<\infty$, $k\in\N_0$ oraz
$u\in W^{k,p}_{\mathrm{loc}}(\Omega)$.
Funkcja $u_\varepsilon$ jest gładka na $\Omega_\varepsilon$
i dla każdego $|\alpha|\le k$ zachodzi
\[
 \partial^\alpha u_\varepsilon
 =\rho_\varepsilon*(\partial^\alpha u)
 \quad\text{na }\Omega_\varepsilon.
\]
Dla każdego zwartego $K\subset\Omega$ mamy dla
dostatecznie małych $\varepsilon>0$ zbieżność
\[
 \sum_{|\alpha|\le k}
 \|\partial^\alpha u_\varepsilon-\partial^\alpha u\|_{L^p(K)}
 \longrightarrow0\qquad(\varepsilon\downarrow0).
\]
W szczególności $u_\varepsilon\to u$ w $W^{k,p}(U)$
dla każdego otwartego $U\Subset\Omega$.
\end{twierdzenie}
\begin{proof}
Z sekcji~\ref{sec:pde9-def} mamy
$\partial^\alpha u=v_\alpha$ w sensie dystrybucji,
gdzie $v_\alpha\in L^p_{\mathrm{loc}}(\Omega)$,
włącznie z $v_0=u$. Każda z tych funkcji jest lokalnie
całkowalna. Twierdzenie~\ref{thm:pde7-mollify}
zastosowane do dystrybucji regularnej $u$ daje
gładkość i tożsamość pochodnych na $\Omega_\varepsilon$.
Obie strony są w tym miejscu gładkimi funkcjami.
Jeśli $K\ne\varnothing$ jest zwarty w $\Omega$,
to dla pewnego $r>0$ zbiór
$K+\overline{B_{2r}(0)}$ jest zwarty w $\Omega$.
W szczególności $K\subset\Omega_\varepsilon$ dla
$0<\varepsilon<r$ (dla $K=\varnothing$ teza jest pusta).
Z twierdzenia~\ref{thm:pde7-approx-id}, zastosowanego
osobno do każdej funkcji $v_\alpha$, wynika
\[
 \|\partial^\alpha u_\varepsilon-v_\alpha\|_{L^p(K)}
 =\|\rho_\varepsilon*v_\alpha-v_\alpha\|_{L^p(K)}
 \longrightarrow0.
\]
Ponieważ wielowskaźników jest skończenie wiele,
ich suma dąży do zera. Dla $K=\overline U$ ograniczamy
każdą normę do $U$ i stosujemy równoważność norm
z sekcji~\ref{sec:pde9-def}.
Nie definiujemy tutaj przestrzeni Sobolewa na dowolnym
zbiorze zwartym: zapis dotyczy ograniczeń pochodnych
z otoczenia $K$ do tego zbioru.
\end{proof}

\begin{lemat}[Odcięcie i wygładzanie o zwartym nośniku]
\label{lem:pde9-compact-product}
Niech $u\in W^{k,p}_{\mathrm{loc}}(\Omega)$,
$1\le p\le\infty$, oraz $\psi\in C_c^\infty(\Omega)$.
Przedłużenie $w$ funkcji $\psi u$ przez zero poza $\Omega$
należy do $W^{k,p}(\R^d)$ i dla $|\alpha|\le k$
jego słaba pochodna jest przedłużeniem przez zero funkcji
\[
 \partial^\alpha(\psi u)=
 \sum_{\beta\le\alpha}\binom{\alpha}{\beta}
 (\partial^\beta\psi)(\partial^{\alpha-\beta}u).
\]
Jeżeli ponadto $p<\infty$, to
$\rho_\varepsilon*w\to w$ w $W^{k,p}(\R^d)$.
Dla dowolnego otwartego $V\Subset\Omega$
zawierającego $\supp\psi$ i dostatecznie małego
$\varepsilon>0$ mamy $\rho_\varepsilon*w\in C_c^\infty(V)$.
\end{lemat}
\begin{proof}
Dla $\varphi\in C_c^\infty(\R^d)$ funkcja
$\psi\varphi|_\Omega$ jest testem na $\Omega$.
W pierwszym rzędzie, dla $k\ge1$, otrzymujemy
\[
 \begin{aligned}
 \langle\partial_j w,\varphi\rangle
 &=-\int_\Omega u\psi\,\partial_j\varphi\,\mathrm d\lambda_d\\
 &=-\int_\Omega u\,\partial_j(\psi\varphi)\,\mathrm d\lambda_d
   +\int_\Omega u(\partial_j\psi)\varphi\,\mathrm d\lambda_d\\
 &=\int_\Omega(\psi\partial_j u+u\partial_j\psi)
 \varphi\,\mathrm d\lambda_d.
 \end{aligned}
\]
Wszystkie funkcje po prawej stronie mają nośniki
zawarte w $\supp\psi$. Ta tożsamość dotyczy również
gładkich czynników $\partial^\beta\psi$ oraz słabych
pochodnych $u$ rzędu mniejszego niż $k$.
Iterując ją i korzystając ze wzoru Pascala dla
współczynników dwumianowych, otrzymujemy podany
wzór wielowskaźnikowy. W każdym składniku gładki
czynnik jest ograniczony i ma zwarty nośnik w $\Omega$,
na którego otoczeniu drugi czynnik należy do $L^p$.
Zatem wszystkie przedłużone pochodne należą do
$L^p(\R^d)$, także dla $p=\infty$.
Dla $k=0$ wystarcza ten sam argument dla $w$.
Nie korzystamy z twierdzenia o przedłużaniu przez zero
dowolnego elementu przestrzeni $W^{k,p}_0$.

Niech $A=\supp\psi$. Jeśli $A=\varnothing$, to $w=0$.
W przeciwnym przypadku $A$ jest zwarty, a nośniki
$w$ i wszystkich wskazanych pochodnych są zawarte w $A$
w sensie zanikania prawie wszędzie poza $A$.
Z definicji splotu dla $0<\varepsilon\le1$
wynika, że $\rho_\varepsilon*w$ znika poza
$A+\overline{B_1(0)}$. To samo dotyczy jego pochodnych,
gdyż funkcja jest gładka. Stosujemy twierdzenie
~\ref{thm:pde9-mollify} na całej $\R^d$ ze wspólnym
zbiorem zwartym $A+\overline{B_1(0)}$.
Dla $p<\infty$ zbieżność lokalna norm staje się
globalną, ponieważ różnice pochodnych znikają poza
tym zbiorem. Wreszcie
$A+\overline{B_\varepsilon(0)}\subset V$ dla
$\varepsilon<\operatorname{dist}(A,\R^d\setminus V)$,
co daje wymagany zwarty nośnik w $V$.
\end{proof}

\begin{uwaga}[Dlaczego nie $p=\infty$?]
Konkluzja normowa na ogół zawodzi dla $p=\infty$,
nawet dla $k=1$. Na $I=(-1,1)$ funkcja $u(x)=|x|$
ma słabą pochodną $\operatorname{sgn}x$.
Dla dowolnego $0<r<1$ i każdej funkcji
$v\in C^1((-r,r))$ zachodzi
\[
 \|v'-\operatorname{sgn}\|_{L^\infty(-r,r)}\ge1.
\]
Jeżeli norma po lewej stronie jest skończona i równa $M$,
to na obu półprzedziałach nierówności
$|v'(x)+1|\le M$ dla $x<0$ i $|v'(x)-1|\le M$ dla $x>0$
zachodzą w każdym punkcie dzięki ciągłości i temu,
że każdy niepusty przedział otwarty ma dodatnią miarę.
Przechodząc do zera z obu stron, dostajemy
$2\le|v'(0)+1|+|v'(0)-1|\le2M$.
W szczególności dotyczy to każdego wygładzenia $u_\varepsilon$,
jeśli jest określone na $(-r,r)$.
Jednak dla $u\in W^{k,\infty}_{\mathrm{loc}}(\Omega)$
i $1\le q<\infty$ mamy
$u\in W^{k,q}_{\mathrm{loc}}(\Omega)$: na każdym
$U\Subset\Omega$ stosujemy oszacowanie
$\|v_\alpha\|_{L^q(U)}\le
\lambda_d(U)^{1/q}\|v_\alpha\|_{L^\infty(U)}$.
Twierdzenie~\ref{thm:pde9-mollify} daje więc lokalną
zbieżność w każdej takiej skończonej normie.
\end{uwaga}

\begin{center}
\begin{tikzpicture}[scale=.92,>=Stealth]
 \draw[thick,rounded corners=12pt,fill=blue!5,draw=blue!55]
  (-4.1,-1.55) -- (3.9,-1.55) -- (4.2,1.55) -- (-3.7,1.65) -- cycle;
 \draw[thick,rounded corners=10pt,fill=green!9,draw=green!48!black]
  (-2.9,-1.15) rectangle (2.9,1.15);
 \draw[thick,rounded corners=9pt,fill=orange!11,draw=orange!75!black]
  (-1.55,-.65) rectangle (1.55,.65);
 \node at (-3.45,1.18) {$\Omega$};
 \node at (2.48,.86) {$U$};
 \node at (-1.05,.30) {$K$};
 \coordinate (x) at (.95,-.05);
 \draw[blue!65!black,thick,dashed] (x) circle (.8);
 \fill (x) circle (1.5pt) node[below left] {$x$};
 \draw[->,blue!65!black,thick] (x)--(1.75,-.05)
  node[midway,above] {$\varepsilon$};
 \node[blue!65!black,anchor=west] at (1.82,-.38) {$B_\varepsilon(x)$};
\end{tikzpicture}

{\small\emph{Lokalność wygładzania.} Dla $K\Subset U\Subset\Omega$
wybieramy $\varepsilon<\operatorname{dist}(K,\R^d\setminus U)$.
Wówczas dla każdego $x\in K$ nośnik jądra
$y\mapsto\rho_\varepsilon(x-y)$ jest zawarty w $U$.
Splot korzysta wyłącznie z wartości funkcji w tym otoczeniu.}
\end{center}


\section{Twierdzenie Meyersa--Serrina i gęstość funkcji gładkich}
\label{sec:pde9-density}

\begin{twierdzenie}[Meyers--Serrin]
\label{thm:pde9-MeyersSerrin}
Niech $\Omega\subset\R^d$ będzie niepustym zbiorem otwartym,
$k\in\N_0$ i $1\le p<\infty$. Wówczas
$C^\infty(\Omega)\cap W^{k,p}(\Omega)$
jest gęsta w $W^{k,p}(\Omega)$.
Nie zakłada się regularności brzegu $\partial\Omega$.
\end{twierdzenie}
\begin{proof}
Najpierw skonstruujemy potrzebny rozkład jedynki
z kontrolą nośników. Dla $j\ge1$ połóżmy
\[
 K_j=\{x\in\Omega:|x|\le j,\
 \operatorname{dist}(x,\R^d\setminus\Omega)\ge1/j\},
\]
a dla $j\le0$ przyjmijmy $K_j=\varnothing$.
Dla $\Omega=\R^d$ drugi warunek pomijamy.
Zbiory $K_j$ są zwarte, gdyż są ograniczone
i domknięte w $\R^d$ (dla niepustego dopełnienia
odległość jest funkcją ciągłą, a dodatnia odległość
wyklucza punkty z dopełnienia).
Ścisłe nierówności $j<j+1$ oraz $1/j>1/(j+1)$
dają $K_j\subset\operatorname{int}K_{j+1}$,
a każdy punkt $\Omega$ należy do wnętrza pewnego $K_j$.
W szczególności każdy zwarty podzbiór $\Omega$
zawiera się we wnętrzu pewnego $K_j$, przez skończone
podpokrycie i monotoniczność zbiorów.

Dla $j\ge1$ wybierzmy $\eta_j\in C_c^\infty(\operatorname{int}K_{j+1})$
z $0\le\eta_j\le1$ i $\eta_j=1$ w otoczeniu $K_j$.
Jeśli $K_j=\varnothing$, przyjmujemy $\eta_j=0$.
W pozostałych przypadkach konstrukcja jest następująca:
wybieramy $r>0$, dla którego
$K_j+\overline{B_{3r}(0)}\subset\operatorname{int}K_{j+1}$.
Splot
$\eta_j=\rho_r*\chi_{K_j+\overline{B_{2r}(0)}}$
jest gładki dzięki twierdzeniu~\ref{thm:pde7-mollify}
na $\R^d$, przyjmuje wartości w $[0,1]$,
jest równy $1$ na $K_j+B_r(0)$ i ma nośnik
w $K_j+\overline{B_{3r}(0)}$.
Wynika to z nieujemności, normalizacji i nośnika jądra.
W szczególności $\eta_{j+1}=1$ na $\supp\eta_j$,
więc $\eta_j\le\eta_{j+1}$.
Połóżmy $\eta_0=0$ i $\psi_j=\eta_j-\eta_{j-1}$.
Funkcje $\psi_j$ są nieujemne, gładkie i mają zwarty nośnik
w $\Omega$, a przez teleskopowanie
$\sum_{j=1}^N\psi_j=\eta_N$.
Dla każdego punktu $x\in\Omega$ mamy $\eta_N(x)=1$
od pewnego $N$, stąd $\sum_{j\ge1}\psi_j=1$.

Ustalmy otwarte zbiory
$V_j=\operatorname{int}K_{j+2}\setminus K_{j-2}$,
$ j\ge1$. Mamy $V_j\Subset\Omega$, a
$\supp\psi_j\subset V_j$. Istotnie, oba składniki
$\psi_j$ mają nośniki w $\operatorname{int}K_{j+1}$
i dla $j\ge2$ oba są równe $1$ w otoczeniu $K_{j-1}$,
jeśli zbiór ten jest niepusty. Dla małych $j$
zbiory o niedodatnich indeksach są puste.
Rodzina $V_j$ jest lokalnie skończona:
jeśli zwarty $A\subset\Omega$ zawiera się
we wnętrzu $K_N$, to $A\cap V_j=\varnothing$
dla $j\ge N+2$. Zatem także rodzina nośników $\psi_j$
jest lokalnie skończona. Nie wymagaliśmy skończonej
liczby zbiorów ani jakiejkolwiek regularności brzegu.

Niech teraz $u\in W^{k,p}(\Omega)$ i $\delta>0$.
Dla każdego $j$ oznaczmy przez $w_j$ przedłużenie
$\psi_j u$ przez zero na $\R^d$.
Lemat~\ref{lem:pde9-compact-product} daje
$w_j\in W^{k,p}(\R^d)$ oraz pozwala wybrać
$\varepsilon_j>0$ tak, aby $g_j=\rho_{\varepsilon_j}*w_j$
należało do $C_c^\infty(V_j)$ i
\[
 \|g_j-w_j\|_{W^{k,p}(\R^d)}<\delta2^{-j}.
\]
Dla $\psi_j=0$ bierzemy $g_j=0$.
Ograniczenie normy do $\Omega$ nie zwiększa błędu.
Funkcja $g=\sum_{j\ge1}g_j|_\Omega$ jest gładka,
ponieważ lokalnie suma ma tylko skończenie wiele
niezerowych składników. Nie zakładamy jeszcze,
że $g\in W^{k,p}(\Omega)$.

Niech $e_j=g_j|_\Omega-w_j|_\Omega$.
Szereg $\sum_{j\ge1}e_j$ zbiega w przestrzeni Banacha
$W^{k,p}(\Omega)$ z twierdzenia~\ref{thm:pde9-banach}:
dla $M>N$ norma sumy od $N+1$ do $M$ nie przekracza
$\delta\sum_{j>N}2^{-j}$, co dąży do zera.
Oznaczmy jego granicę przez $e$; ciągłość normy daje
$\|e\|_{W^{k,p}}\le\delta$.
Wykażemy, że $e=g-u$ prawie wszędzie.
Dla dowolnego $\varphi\in C_c^\infty(\Omega)$
zbieżność w $L^p$ oraz nierówność Höldera dają
\[
 \int_\Omega e\varphi\,\mathrm d\lambda_d
 =\lim_{N\to\infty}\sum_{j=1}^N
 \int_\Omega e_j\varphi\,\mathrm d\lambda_d.
\]
Tylko skończenie wiele składników może mieć
nośnik przecinający $\supp\varphi$.
Na otoczeniu tego nośnika mamy
$\sum_j w_j|_\Omega=u$ prawie wszędzie,
ponieważ $\sum_j\psi_j=1$.
Prawa strona jest zatem równa
$\int_\Omega(g-u)\varphi\,\mathrm d\lambda_d$.
Funkcja $g-u$ jest lokalnie całkowalna, więc
jednoznaczność dystrybucji regularnych daje
$e=g-u$ prawie wszędzie. W konsekwencji
$g=u+e\in W^{k,p}(\Omega)$ oraz
$\|g-u\|_{W^{k,p}}\le\delta$.
Dla każdego zadanego błędu $\epsilon>0$
wybieramy $\delta=\epsilon/2$, otrzymując tezę.
Nie musimy zakładać zbieżności szeregu $\sum_j w_j$
w globalnej normie Sobolewa; zbieżny w tej normie
jest szereg błędów aproksymacji.
\end{proof}

\begin{uwaga}[Gładkość nie oznacza zerowania na brzegu]
Twierdzenie~\ref{thm:pde9-MeyersSerrin} nie głosi,
że $C_c^\infty(\Omega)$ jest gęsta w $W^{k,p}(\Omega)$.
Funkcje o zwartym nośniku wewnątrz $\Omega$
znikają w pewnym otoczeniu brzegu, podczas gdy
$C^\infty(\Omega)$ nie narzuca takiego warunku.
Gdy $\Omega=\R^d$ i $p<\infty$, gęstość
$C_c^\infty(\R^d)$ rzeczywiście zachodzi.
Aby to sprawdzić, wybierzmy $\psi\in C_c^\infty(\R^d)$,
$0\le\psi\le1$, równą $1$ na $B_1(0)$ i zerową
poza $B_2(0)$, jak w dowodzie twierdzenia
~\ref{thm:pde8-density}, oraz $\psi_R(x)=\psi(x/R)$.
Dla $u\in W^{k,p}(\R^d)$, $R\ge1$, z lematu
~\ref{lem:pde9-compact-product} wynika
$\psi_Ru\in W^{k,p}(\R^d)$.
W rozwinięciu
\[
 \partial^\alpha((1-\psi_R)u)=
 (1-\psi_R)\partial^\alpha u+
 \sum_{0<\beta\le\alpha}\binom{\alpha}{\beta}
 \partial^\beta(1-\psi_R)\partial^{\alpha-\beta}u
\]
pierwszy składnik ma normę $L^p$ ograniczoną przez
$\|\partial^\alpha u\|_{L^p(\R^d\setminus B_R(0))}\to0$.
W pozostałych składnikach gładki czynnik ma normę
supremum równą
$R^{-|\beta|}\|\partial^\beta\psi\|_\infty$,
a drugi ma skończoną normę $L^p$.
Każdy taki iloczyn dąży więc do zera w $L^p$.
Wzór wynika z liniowości pochodnych i lematu
~\ref{lem:pde9-compact-product}.
Otrzymujemy $\psi_Ru\to u$ w $W^{k,p}(\R^d)$.
Dla ustalonego $R$ stosujemy ponownie ten lemat,
aby aproksymować $\psi_Ru$ przez gładkie funkcje
o zwartym nośniku. Wybierając najpierw $R$,
a potem promień wygładzania z błędami mniejszymi
niż połowa zadanego błędu, dostajemy żądaną gęstość.

Dla $p=\infty$ nawet gęstość
$C^\infty(\Omega)\cap W^{k,\infty}(\Omega)$
nie zachodzi ogólnie: dla $k=1$ i $\Omega=(-1,1)$
dowiedziony w sekcji~\ref{sec:pde9-mollify}
dolny próg $1$ dla błędu pochodnej funkcji $|x|$
wyklucza taką aproksymację.
\end{uwaga}

\begin{przyklad}[Stała na przedziale nie należy do $W^{1,p}_0$]
\label{ex:pde9-constant}
Niech $\Omega=(0,1)$, $1\le p<\infty$, oraz $u\equiv1$.
Ponieważ $u\in L^p(0,1)$ i jej słaba pochodna
jest zerowa, $u\in W^{1,p}(0,1)$, ale
$u\notin W^{1,p}_0(0,1)$. Gdyby istniał ciąg
$\varphi_n\in C_c^\infty(0,1)$ zbieżny do $u$
w $W^{1,p}$, to z ciągłości śladu
(wniosek~\ref{cor:pde9-trace-oneD}) wynikałoby
\[
 |\varphi_n(0)-1|\le
 \|\varphi_n-u\|_{L^p(0,1)}+
 \|\varphi_n'-u'\|_{L^p(0,1)}\longrightarrow0.
\]
Jest to niemożliwe, ponieważ ciągły reprezentant
każdej $\varphi_n$ jest równy zero w $0$.
\end{przyklad}

\section[Domknięcie funkcji testowych i przedłużenie przez zero]{Przestrzeń $W^{k,p}_0$ i przedłużenie przez zero}
\label{sec:pde9-zero}

\begin{twierdzenie}[Przedłużenie przez zero]
\label{thm:pde9-zero-extension}
Niech $\Omega\subset\R^d$ będzie niepustym zbiorem otwartym,
$k\in\N_0$, $1\le p<\infty$ i $u\in W^{k,p}_0(\Omega)$.
Przedłużenie przez zero
\[
 E_0u(x)=\begin{cases}u(x),&x\in\Omega,\\0,&x\notin\Omega\end{cases}
\]
należy do $W^{k,p}(\R^d)$ i spełnia
\[
 \partial^\alpha(E_0u)=E_0(\partial^\alpha u)
 \quad (|\alpha|\le k),\qquad
 \|E_0u\|_{W^{k,p}(\R^d)}=\|u\|_{W^{k,p}(\Omega)}.
\]
Operator $E_0$ jest liniowym izometrycznym zanurzeniem
o domkniętym obrazie.
\end{twierdzenie}
\begin{proof}
Przedłużenie przez zero klas funkcji z $L^p(\Omega)$
jest dobrze określoną operacją liniową i zachowuje normę,
gdyż $\int_{\R^d}|E_0f|^p\,\mathrm d\lambda_d
=\int_\Omega|f|^p\,\mathrm d\lambda_d$.
Nie wynika z tego jeszcze zachowanie słabych pochodnych.

Z definicji $W^{k,p}_0$ istnieje ciąg
$\varphi_n\in C_c^\infty(\Omega)$ taki, że
$\varphi_n\to u$ w $W^{k,p}(\Omega)$.
Dla każdej $\varphi_n$ jej nośnik jest zwartym podzbiorem
$\Omega$, a zatem w otoczeniu każdego punktu brzegu
funkcja i jej przedłużenie przez zero są zerowe.
Stąd $E_0\varphi_n\in C_c^\infty(\R^d)$ i
$\partial^\alpha E_0\varphi_n=E_0(\partial^\alpha\varphi_n)$.
W szczególności dla $n,m$ mamy
\[
 \|E_0\varphi_n-E_0\varphi_m\|_{W^{k,p}(\R^d)}
 =\|\varphi_n-\varphi_m\|_{W^{k,p}(\Omega)}.
\]
Ciąg po lewej stronie jest więc ciągiem Cauchy'ego.
Z twierdzenia~\ref{thm:pde9-banach} istnieje jego
granica $v\in W^{k,p}(\R^d)$.
Jednocześnie
\[
 \|E_0\varphi_n-E_0u\|_{L^p(\R^d)}
 =\|\varphi_n-u\|_{L^p(\Omega)}\to0.
\]
Z jednoznaczności granicy w $L^p$ dostajemy $v=E_0u$.
Dla dowolnego $|\alpha|\le k$ analogicznie
\[
 \|\partial^\alpha E_0\varphi_n-E_0(\partial^\alpha u)\|_{L^p(\R^d)}
 =\|\partial^\alpha\varphi_n-\partial^\alpha u\|_{L^p(\Omega)}
 \to0.
\]
Ponieważ zbieżność w normie Sobolewa daje zarazem
$\partial^\alpha E_0\varphi_n\to\partial^\alpha v$ w $L^p$,
otrzymujemy wzór na pochodną. Sumując $p$-te potęgi norm
po wszystkich wielowskaźnikach, dostajemy równość norm.

Przestrzeń $W^{k,p}_0(\Omega)$ jest domkniętą podprzestrzenią
przestrzeni Banacha $W^{k,p}(\Omega)$, więc jest zupełna.
Jeżeli $E_0u_n$ zbiega w $W^{k,p}(\R^d)$, to z izometrii
$(u_n)$ jest ciągiem Cauchy'ego w $W^{k,p}_0(\Omega)$.
Dla jego granicy $u$ izometria daje $E_0u_n\to E_0u$.
Obraz operatora jest zatem domknięty.
\end{proof}

\begin{uwaga}[Znaczenie warunku brzegowego]
Nie każdą funkcję z przestrzeni Sobolewa można tak przedłużyć.
Dla $u\equiv1$ na $(0,1)$ mamy $u\in W^{1,p}(0,1)$,
lecz dla $\varphi\in C_c^\infty(\R)$ zachodzi
\[
 \langle (E_0u)',\varphi\rangle
 =-\int_0^1\varphi'(x)\,\mathrm d\lambda_1(x)
 =\varphi(0)-\varphi(1).
\]
Pochodną dystrybucyjną jest więc $\delta_0-\delta_1$.
Nie może jej reprezentować funkcja lokalnie całkowalna:
ograniczając testy do otoczenia zera niezawierającego $1$,
otrzymalibyśmy regularną reprezentację $\delta_0$,
co wyklucza argument z sekcji~\ref{sec:pde7-def}.
W konsekwencji $E_0u\notin W^{1,p}(\R)$.

Na dowolnym zbiorze otwartym nie zachodzi również
implikacja odwrotna do twierdzenia~\ref{thm:pde9-zero-extension}.
Niech $\Omega=(-1,0)\cup(0,1)$ oraz
$\zeta\in C_c^\infty(-1,1)$, $\zeta(0)=1$.
Dla $u=\zeta|_\Omega$ i $1\le p<\infty$ mamy
$u\in W^{1,p}(\Omega)$, a $E_0u=\zeta$ prawie wszędzie
na $\R$, ponieważ różnica może występować tylko w zerze.
Zatem $E_0u\in W^{1,p}(\R)$.
Gdyby $u\in W^{1,p}_0(\Omega)$, istniałyby
$\varphi_n\in C_c^\infty(\Omega)$ zbieżne do $u$
w $W^{1,p}(\Omega)$. Po ograniczeniu do $(0,1)$
zbieżność zachodziłaby nadal w $W^{1,p}(0,1)$.
Każda $\varphi_n$ ma jednak zerową wartość brzegową
w $0$ od prawej strony, podczas gdy $u$ ma tam wartość $1$.
Przeczy to ciągłości śladu z wniosku
~\ref{cor:pde9-trace-oneD}.
W szczególności samo przedłużenie przez zero należące
do $W^{1,p}(\R^d)$ nie charakteryzuje $W^{1,p}_0$
bez dodatkowych założeń o zbiorze otwartym.
\end{uwaga}


\section{Nierówność Poincarégo na kostce}
\label{sec:pde9-poincare}

\begin{twierdzenie}[Poincaré--Wirtinger na kostce]
\label{thm:pde9-cube-Poincare}
Niech $Q=\prod_{j=1}^d(a_j,a_j+L)$, $L>0$,
$1\le p\le\infty$. Dla $u\in W^{1,p}(Q)$ połóżmy
\[
 u_Q=\frac1{\lambda_d(Q)}\int_Q u\,\mathrm d\lambda_d.
\]
Średnia jest dobrze określona, ponieważ $Q$ ma skończoną
miarę i $L^p(Q)\subset L^1(Q)$.
Wówczas
\[
 \|u-u_Q\|_{L^p(Q)}
 \le L\sum_{j=1}^d\|\partial_j u\|_{L^p(Q)}.
\]
Prawa strona jest równa zeru dokładnie wtedy,
gdy $u$ jest prawie wszędzie stała na $Q$.
\end{twierdzenie}
\begin{proof}
Najpierw ustalmy $1\le p<\infty$ i niech $u$
będzie gładka w otoczeniu $\overline Q$.
Dla $h\in C^1([a,a+L])$ podstawowe twierdzenie
rachunku całkowego daje
\[
 \begin{aligned}
 \left|h(x)-\frac1L\int_a^{a+L}h(t)\,\mathrm d\lambda_1(t)\right|
 &\le\frac1L\int_a^{a+L}|h(x)-h(t)|\,\mathrm d\lambda_1(t)\\
 &\le\int_a^{a+L}|h'(s)|\,\mathrm d\lambda_1(s).
 \end{aligned}
\]
Ostatnia nierówność wynika stąd, że dla każdego $t$
moduł różnicy jest nie większy od całki z $|h'|$
na całym przedziale. Biorąc normę $L^p$ względem $x$
i stosując nierówność Höldera, otrzymujemy
\[
 \left\|h-\frac1L\int_a^{a+L}h\,\mathrm d\lambda_1\right\|_p
 \le L^{1/p}\|h'\|_1
 \le L\|h'\|_p.
\]

Dla $j=1,\ldots,d$ niech $I_j=(a_j,a_j+L)$ i
niech $A_j$ oznacza uśrednienie względem pierwszych
$j$ współrzędnych:
\[
 (A_j u)(x)=\frac1{L^j}
 \int_{\prod_{i=1}^j I_i}u(t_1,\ldots,t_j,x_{j+1},\ldots,x_d)
 \,\mathrm d\lambda_j(t).
\]
Traktujemy $A_j u$ jako funkcję na całej kostce,
stałą w pierwszych $j$ zmiennych. Przyjmujemy $A_0u=u$;
wówczas $A_du=u_Q$. Dla $v\in L^p(Q)$
nierówność Jensena względem znormalizowanej miary
na $\prod_{i=1}^j I_i$ daje $|A_jv|^p\le A_j(|v|^p)$
prawie wszędzie. Całkując po $Q$ i stosując Fubiniego,
otrzymujemy
\[
 \|A_jv\|_p^p\le\int_Q A_j(|v|^p)\,\mathrm d\lambda_d
 =\int_Q|v|^p\,\mathrm d\lambda_d=\|v\|_p^p.
\]
Zatem $A_j$ jest kontrakcją.

Dla ustalonych pozostałych współrzędnych stosujemy
nierówność jednowymiarową do funkcji $A_{j-1}u$
zmiennej $x_j$, której średnia w tej zmiennej jest
$A_ju$. Po podniesieniu do potęgi $p$ i scałkowaniu
względem pozostałych współrzędnych (dla $d=1$
nie ma drugiego całkowania) dostajemy
\[
 \|A_{j-1}u-A_ju\|_p^p
 \le L^p\|\partial_j A_{j-1}u\|_p^p.
\]
Różniczkowanie w $x_j$ komutuje z $A_{j-1}$,
ponieważ całkujemy względem innych współrzędnych,
a odpowiednie pochodne $u$ są ciągłe i ograniczone
na $\overline Q$. Z kontrakcyjności otrzymujemy
\[
 \|\partial_j A_{j-1}u\|_p
 =\|A_{j-1}(\partial_j u)\|_p\le\|\partial_j u\|_p.
\]
Teleskopowanie $u-u_Q=\sum_{j=1}^d(A_{j-1}u-A_ju)$
oraz nierówność trójkąta dają tezę dla gładkich $u$.

Niech teraz $u\in W^{1,p}(Q)$, nadal dla $p<\infty$.
Dla $0<\delta<L/2$ definiujemy kostkę
$Q_\delta=\prod_j(a_j+\delta,a_j+L-\delta)$.
Z twierdzenia~\ref{thm:pde9-mollify} wynika, że
dla ustalonego $\delta$ funkcje $u_\varepsilon$ są
gładkie w otoczeniu $\overline{Q_\delta}$ dla
dostatecznie małych $\varepsilon$ i zbiegają do $u$
w $W^{1,p}(Q_\delta)$. Ponadto
\[
 |(u_\varepsilon)_{Q_\delta}-u_{Q_\delta}|
 \le\lambda_d(Q_\delta)^{-1/p}
      \|u_\varepsilon-u\|_{L^p(Q_\delta)}\to0
\]
z nierówności Höldera. W szczególności
$u_\varepsilon-(u_\varepsilon)_{Q_\delta}\to
u-u_{Q_\delta}$ w $L^p(Q_\delta)$.
Stosujemy już udowodnione oszacowanie na $Q_\delta$
i przechodzimy z $\varepsilon$ do zera:
\[
 \|u-u_{Q_\delta}\|_{L^p(Q_\delta)}
 \le (L-2\delta)\sum_{j=1}^d\|\partial_j u\|_{L^p(Q_\delta)}.
\]
Teraz $u_{Q_\delta}\to u_Q$ przy $\delta\downarrow0$,
ponieważ $\lambda_d(Q_\delta)\to\lambda_d(Q)>0$
i całki z $u\in L^1(Q)$ zbiegają przez twierdzenie
o zbieżności zmajoryzowanej. Mamy też
\[
 \left|\|u-u_{Q_\delta}\|_{L^p(Q_\delta)}
       -\|u-u_Q\|_{L^p(Q_\delta)}\right|
 \le\lambda_d(Q_\delta)^{1/p}|u_{Q_\delta}-u_Q|\to0.
\]
Normy funkcji $u-u_Q$ oraz $\partial_j u$ na
$Q_\delta$ zbiegają do norm na $Q$ dzięki całkowalności
ich $p$-tych potęg i zbieżności zmajoryzowanej.
Przejście do granicy daje żądaną nierówność.

W przypadku $p=\infty$ każda funkcja
$u\in W^{1,\infty}(Q)$ należy również do $W^{1,q}(Q)$
dla każdego $1\le q<\infty$, z tymi samymi słabymi pochodnymi.
Dla dowolnego $f\in L^\infty(Q)$ mamy
$\|f\|_q\to\|f\|_\infty$ przy $q\to\infty$.
Istotnie, jeśli $M=\|f\|_\infty$, to
$\|f\|_q\le M\lambda_d(Q)^{1/q}$.
Dla $M>0$ i $0<\gamma<M$ zbiór
$A=\{x\in Q:|f(x)|>M-\gamma\}$ ma dodatnią miarę,
więc $\|f\|_q\ge(M-\gamma)\lambda_d(A)^{1/q}$.
Przejście z $q$ do nieskończoności, a następnie
z $\gamma$ do zera, daje równość granicy; dla $M=0$
jest ona oczywista. Stosujemy zatem oszacowanie dla
skończonego $q$ do $u$, a powyższą granicę do
$u-u_Q$ i każdej $\partial_j u$.

Jeśli wszystkie pierwsze pochodne są zerowe,
nierówność daje $\|u-u_Q\|_p=0$, czyli $u=u_Q$
prawie wszędzie. Odwrotnie, funkcja prawie wszędzie
stała ma zerowe pochodne dystrybucyjne, ponieważ
całka z pochodnej testu o zwartym nośniku jest zerowa.
\end{proof}

\begin{wniosek}[Wersja z zerowym warunkiem brzegowym]
\label{cor:pde9-Poincare-zero}
Dla $1\le p<\infty$ i $u\in W^{1,p}_0(Q)$
na kostce z poprzedniego twierdzenia zachodzi
\[
 \|u\|_{L^p(Q)}\le L\|\partial_j u\|_{L^p(Q)}
 \qquad (j=1,\ldots,d).
\]
W szczególności
\[
 |u|_{1,p}=\left(\sum_{j=1}^d\|\partial_j u\|_p^p\right)^{1/p}
\]
jest normą równoważną normie $W^{1,p}$ na $W^{1,p}_0(Q)$,
z oszacowaniami
$|u|_{1,p}\le\|u\|_{W^{1,p}}\le(1+L^p)^{1/p}|u|_{1,p}$.
\end{wniosek}
\begin{proof}
Dla $u\in C_c^\infty(Q)$ i ustalonych współrzędnych
poza $x_j$ funkcja znika w otoczeniu dolnego końca
odcinka $I_j=(a_j,a_j+L)$. Mamy więc
\[
 u(x)=\int_{a_j}^{x_j}\partial_j u(x_1,\ldots,t,\ldots,x_d)
 \,\mathrm d\lambda_1(t).
\]
Z nierówności Höldera wynika
\[
 |u(x)|^p\le L^{p-1}\int_{I_j}
 |\partial_j u(x_1,\ldots,t,\ldots,x_d)|^p
 \,\mathrm d\lambda_1(t).
\]
Dla $p=1$ jest to nierówność trójkąta dla całki.
Całkowanie względem $x_j$ daje dodatkowy czynnik $L$,
a całkowanie względem pozostałych współrzędnych daje
z Fubiniego $\|u\|_p^p\le L^p\|\partial_j u\|_p^p$
(dla $d=1$ tylko pierwsze całkowanie).
Bierzemy $p$-ty pierwiastek.

Jeżeli $u\in W^{1,p}_0(Q)$, to istnieje ciąg
$\varphi_n\in C_c^\infty(Q)$ zbieżny do $u$ w $W^{1,p}$.
W szczególności $\varphi_n\to u$ i
$\partial_j\varphi_n\to\partial_j u$ w $L^p$.
Ciągłość normy pozwala przejść do granicy
w oszacowaniu dla $\varphi_n$.

Z definicji normy mamy
$\|u\|_{W^{1,p}}^p=\|u\|_p^p+|u|_{1,p}^p$.
Poprzednia nierówność daje $\|u\|_p^p\le L^p|u|_{1,p}^p$.
Stąd otrzymujemy podane dwustronne oszacowanie.
Półnorma $|u|_{1,p}$ jest dodatnio określona na
$W^{1,p}_0(Q)$ dzięki temu oszacowaniu, a jednorodność
i nierówność trójkąta wynikają z normy skończonego
produktu przestrzeni $L^p$. Jest więc normą.
\end{proof}

\section{Reguły rachunkowe w przestrzeniach Sobolewa}
\label{sec:pde9-calculus}

\begin{twierdzenie}[Mnożenie przez funkcję gładką i iloczyn]
\label{thm:pde9-products}
Niech $\Omega\subset\R^d$ będzie otwarta, $1\le p\le\infty$
i $u\in W^{1,p}(\Omega)$.
Jeśli $\phi\in C^1(\Omega)$ oraz $\phi,\nabla\phi\in L^\infty(\Omega)$,
to $\phi u\in W^{1,p}(\Omega)$ i
$\partial_j(\phi u)=\phi\partial_j u+u\partial_j\phi$.
Ponadto jeśli $v\in W^{1,q}(\Omega)$, gdzie
$1/p+1/q=1$ z konwencją $1/\infty=0$,
to $uv\in W^{1,1}(\Omega)$ i
$\partial_j(uv)=v\partial_j u+u\partial_j v$.
W szczególności $W^{1,\infty}(\Omega)$ jest algebrą
względem mnożenia punktowego klas funkcji, z tą samą
regułą Leibniza.
\end{twierdzenie}
\begin{proof}
Ustalmy $\varphi\in C_c^\infty(\Omega)$ i wybierzmy
otwarty $U\Subset\Omega$ zawierający jej nośnik.
Dla $p=\infty$ funkcja $u$ należy lokalnie do $W^{1,1}$
przez oszacowanie na zbiorach skończonej miary.
Dla każdego rozważanego $p$ możemy zatem użyć lokalnych
wygładzeń $u_\varepsilon\to u$ w $W^{1,1}(U)$ na mocy
twierdzenia~\ref{thm:pde9-mollify} (dla $p<\infty$
korzystamy też z nierówności Höldera).
Iloczyn $\phi u_\varepsilon$ jest klasy $C^1$,
więc klasyczne całkowanie przez części daje
\[
 \int_U\phi u_\varepsilon\partial_j\varphi\,\mathrm d\lambda_d
 =-\int_U(\phi\partial_j u_\varepsilon
          +u_\varepsilon\partial_j\phi)\varphi\,\mathrm d\lambda_d.
\]
Nie podstawiamy $\phi\varphi$ bezpośrednio jako testu
dystrybucyjnego, ponieważ $\phi$ jest tylko klasy $C^1$.
Ograniczoność $\phi,\partial_j\phi$ i zbieżność
w $W^{1,1}(U)$ uzasadniają przejście do granicy
w każdej całce. Dostajemy żądaną tożsamość testową,
a globalnie
\[
 \begin{aligned}
 \|\phi u\|_p&\le\|\phi\|_\infty\|u\|_p,\\
 \|\partial_j(\phi u)\|_p
 &\le\|\phi\|_\infty\|\partial_j u\|_p+
      \|\partial_j\phi\|_\infty\|u\|_p .
 \end{aligned}
\]
Obie nierówności obowiązują również dla $p=\infty$.

Dla iloczynu $uv$ wybieramy wspólne lokalne wygładzenia
$u_\varepsilon,v_\varepsilon$ na otoczeniu $\overline U$.
Gdy $p,q<\infty$, dla $a=u$ lub $a=\partial_j u$
oraz $b=v$ lub $b=\partial_j v$ odpowiednie wygładzenia
$a_\varepsilon,b_\varepsilon$ zbiegają do $a,b$
w $L^p(U)$ i $L^q(U)$. Z nierówności Höldera mamy
\[
 \|a_\varepsilon b_\varepsilon-ab\|_{L^1(U)}
 \le\|a_\varepsilon-a\|_{L^p(U)}\|b_\varepsilon\|_{L^q(U)}
    +\|a\|_{L^p(U)}\|b_\varepsilon-b\|_{L^q(U)}\to0.
\]
Pierwszy składnik dąży do zera, ponieważ zbieżny
ciąg norm $\|b_\varepsilon\|_q$ jest ograniczony
dla małych $\varepsilon$.

Gdy $p=1,q=\infty$, to $\|b_\varepsilon\|_{L^\infty(U)}
\le\|b\|_{L^\infty(\Omega)}$ z nieujemności i normalizacji
jądra oraz wzoru na wygładzone pochodne.
Mamy $b_\varepsilon\to b$ w $L^1(U)$; dla pewnego
ciągu $\varepsilon_n\downarrow0$ wszystkie skończenie
liczne funkcje $v_{\varepsilon_n},\partial_jv_{\varepsilon_n}$
zbiegają do odpowiednich granic prawie wszędzie.
Korzystamy z faktu, że ze zbieżności w $L^1$ można
wybrać podciąg zbieżny prawie wszędzie (wybieramy
błędy norm sumowalne i stosujemy twierdzenie Tonellego
do sumy ich modułów).
Dla takiego ciągu mamy
\[
 \|a_{\varepsilon_n}b_{\varepsilon_n}-ab\|_{L^1(U)}
 \le \|b\|_\infty\|a_{\varepsilon_n}-a\|_{L^1(U)}
      +\int_U|a||b_{\varepsilon_n}-b|\,\mathrm d\lambda_d .
\]
Druga całka dąży do zera z twierdzenia o zbieżności
zmajoryzowanej, z majorantą $2\|b\|_\infty|a|$.
Przypadek $p=\infty,q=1$ otrzymujemy zamieniając
role $u$ i $v$.

Stosując te oszacowania do $uv$, $v\partial_j u$
i $u\partial_j v$, przechodzimy do granicy w tożsamości
testowej dla gładkiego iloczynu $u_{\varepsilon_n}v_{\varepsilon_n}$.
Globalnie $uv,v\partial_j u,u\partial_j v\in L^1(\Omega)$
dzięki nierówności Höldera. To dowodzi tezy.
Jeśli zaś $u,v\in W^{1,\infty}(\Omega)$, to lokalnie
możemy zastosować udowodniony przypadek $p=1,q=\infty$,
ponieważ na każdym $U\Subset\Omega$ mamy $u|_U\in W^{1,1}(U)$.
Otrzymana tożsamość dystrybucyjna obowiązuje na całej
$\Omega$, gdyż każdy test ma zwarty nośnik w pewnym takim $U$.
Wszystkie trzy iloczyny są globalnie ograniczone istotnie,
co daje przynależność do $W^{1,\infty}(\Omega)$.
\end{proof}

\begin{twierdzenie}[Reguła łańcuchowa dla funkcji klasy $C^1$ o ograniczonej pochodnej]
\label{thm:pde9-chain}
Niech $1\le p\le\infty$, $u\in W^{1,p}(\Omega)$
będzie rzeczywista, a $F\in C^1(\R)$ spełnia
$F(0)=0$ oraz $\|F'\|_\infty<\infty$.
Wtedy $F\circ u\in W^{1,p}(\Omega)$ i
\[
 \partial_j(F(u))=F'(u)\partial_j u
 \quad\text{prawie wszędzie w }\Omega.
\]
Jeżeli $\lambda_d(\Omega)<\infty$ lub $p=\infty$,
warunek $F(0)=0$ można pominąć.
\end{twierdzenie}
\begin{proof}
Niech $L=\|F'\|_\infty$. Z podstawowego twierdzenia
rachunku całkowego mamy $|F(t)-F(s)|\le L|t-s|$.
Przy $F(0)=0$ dostajemy $\|F(u)\|_p\le L\|u\|_p$.
Gdy $p<\infty$ i $\lambda_d(\Omega)<\infty$,
dodatkowa stała $F(0)$ ma normę
$|F(0)|\lambda_d(\Omega)^{1/p}$, a dla $p=\infty$
niezależnie od miary mamy $\|F(u)\|_\infty
\le |F(0)|+L\|u\|_\infty$.
W każdym z dopuszczalnych przypadków $F(u)\in L^p$.
Ponadto $\|F'(u)\partial_j u\|_p\le L\|\partial_j u\|_p$.

Dla $p<\infty$ ustalmy $U\Subset\Omega$ zawierający
nośnik rozważanego testu. Z twierdzenia
~\ref{thm:pde9-mollify} wybieramy ciąg wygładzeń
$u_n\to u$ w $W^{1,p}(U)$, a po przejściu
do podciągu także $u_n\to u$ prawie wszędzie na $U$.
Wybór podciągu uzasadnia sumowalność odpowiednio
wybranych $p$-tych potęg błędów w $L^p$ i Tonelli.
Z oszacowania Lipschitzowskiego wynika
$\|F(u_n)-F(u)\|_{L^p(U)}\le L\|u_n-u\|_{L^p(U)}\to0$.
Klasyczna reguła łańcuchowa dla $F(u_n)\in C^1(U)$
daje rozkład
\[
 \begin{aligned}
 \partial_j(F(u_n))-F'(u)\partial_j u
 &=F'(u_n)(\partial_j u_n-\partial_j u)\\
 &\quad +(F'(u_n)-F'(u))\partial_j u .
 \end{aligned}
\]
Norma pierwszego składnika jest nie większa niż
$L\|\partial_j u_n-\partial_j u\|_{L^p(U)}\to0$.
Drugi składnik dąży do zera prawie wszędzie
dzięki ciągłości $F'$; jego $p$-ta potęga modułu
jest ograniczona przez $(2L)^p|\partial_j u|^p$.
Twierdzenie o zbieżności zmajoryzowanej daje zatem
zbieżność w $L^p(U)$. Przechodząc do granicy
w całkowaniu przez części przeciw testowi,
uzyskujemy tożsamość słabej pochodnej na $\Omega$,
a wcześniejsze oszacowania zapewniają globalną
przynależność do $W^{1,p}$.

Dla $p=\infty$ stosujemy już udowodniony przypadek
$p=1$ na każdym $U\Subset\Omega$, gdzie $u$ i jej
pochodne należą do $L^1(U)$ i gdzie można pominąć
warunek $F(0)=0$. Tożsamości testowe obowiązują na całej
$\Omega$, a globalne oszacowania w $L^\infty$
dają tezę. Nie korzystamy ze zbieżności wygładzeń
w normie $W^{1,\infty}$.
\end{proof}

\begin{lemat}[Gradient znika na zbiorach poziomicowych]
\label{lem:pde9-levelsets}
Jeśli $u\in W^{1,p}(\Omega)$ jest rzeczywista,
$1\le p\le\infty$, i $c\in\R$, to
$\nabla u=0$ prawie wszędzie na $\{x\in\Omega:u(x)=c\}$.
\end{lemat}
\begin{proof}
Niech $\theta\in C^\infty(\R)$, $0\le\theta\le1$,
będzie równa $0$ na $[-1,1]$ i $1$ poza $(-2,2)$.
Taką funkcję otrzymujemy odejmując od $1$ gładką
funkcję odcinającą, skonstruowaną jak w dowodzie
twierdzenia~\ref{thm:pde8-density}.
Połóżmy $h_n(t)=\theta(n(t-c))$ i
$F_n(t)=\int_0^t h_n(s)\,\mathrm d\lambda_1(s)$,
gdzie dla $t<0$ całkę rozumiemy jako
$-\int_t^0 h_n(s)\,\mathrm d\lambda_1(s)$.
Wtedy $F_n(0)=0$, $F_n'=h_n$, $|F_n(t)|\le|t|$
oraz $|F_n(t)-t|\le\min\{|t|,4/n\}$,
ponieważ $1-h_n$ znika poza przedziałem
$(c-2/n,c+2/n)$ i ma wartości w $[0,1]$.

Dla $p<\infty$ z twierdzenia o zbieżności zmajoryzowanej
mamy $F_n(u)\to u$ w $L^p(\Omega)$,
z majorantą $|u|^p$ dla $p$-tych potęg błędów,
także na zbiorach nieskończonej miary.
Twierdzenie~\ref{thm:pde9-chain} daje
$\partial_j F_n(u)=h_n(u)\partial_j u$.
Ponieważ $h_n(t)\to\chi_{\R\setminus\{c\}}(t)$ dla
każdego $t$, otrzymujemy przez zbieżność zmajoryzowaną
\[
 h_n(u)\partial_j u\longrightarrow
 \chi_{\{u\ne c\}}\partial_j u \quad\text{w }L^p(\Omega).
\]
Rzeczywiście $p$-ta potęga modułu różnicy jest
ograniczona przez $|\partial_j u|^p$.
Z zamkniętości pojedynczej pochodnej
(wniosek~\ref{cor:pde9-closed}) dostajemy
$\partial_j u=\chi_{\{u\ne c\}}\partial_j u$.
Dla $p=\infty$ stosujemy przypadek $p=1$
na przeliczalnym pokryciu $\Omega$ relatywnie zwartymi
otwartymi podzbiorami (sekcja~\ref{sec:pde9-def}).
Przeliczalna suma zbiorów wyjątkowych ma miarę zero,
więc teza zachodzi na całej $\Omega$.
\end{proof}

\begin{wniosek}[Obcięcia, części dodatnie i moduł]
\label{cor:pde9-truncation}
Niech $u\in W^{1,p}(\Omega)$ będzie rzeczywista,
$1\le p\le\infty$.
Wtedy $u^+=\max(u,0)$, $u^-=\max(-u,0)$ oraz $|u|$
należą do $W^{1,p}(\Omega)$ i prawie wszędzie
\[
 \nabla u^+=\chi_{\{u>0\}}\nabla u,\qquad
 \nabla u^-=-\chi_{\{u<0\}}\nabla u,\qquad
 \nabla|u|=\operatorname{sgn}(u)\nabla u .
\]
Dla $M>0$ obcięcie $T_M(u)=\max(-M,\min(u,M))$
również należy do $W^{1,p}(\Omega)$ i
$\nabla T_M(u)=\chi_{\{|u|<M\}}\nabla u$ prawie wszędzie.
Jeżeli $p<\infty$, to $T_M(u)\to u$ w $W^{1,p}(\Omega)$
przy $M\to\infty$.
\end{wniosek}
\begin{proof}
Ustalmy jednowymiarowe nieujemne jądro wygładzające
$\eta\in C_c^\infty((-1,1))$ o całce $1$
i $\eta_n(t)=n\eta(nt)$.
Niech $G$ oznacza funkcję $t\mapsto t^+$
albo $t\mapsto T_M(t)$.
Obie funkcje są lipschitzowskie ze stałą $1$,
mają wartość $0$ w zerze i są absolutnie ciągłe
na każdym przedziale ograniczonym.
Ich słabe pochodne na $\R$ to odpowiednio
$g=\chi_{(0,\infty)}$ i $g=\chi_{(-M,M)}$,
co wynika z jednowymiarowej charakterystyki
~\ref{thm:pde9-ac}, stosowanej lokalnie.
Połóżmy
\[
 G_n(t)=(\eta_n*G)(t)-(\eta_n*G)(0).
\]
Z twierdzenia~\ref{thm:pde7-mollify} mamy
$G_n\in C^\infty(\R)$ i $G_n'=\eta_n*g$,
zatem $0\le G_n'\le1$ i $G_n(0)=0$.
Ponadto $|G_n(t)|\le|t|$ oraz
\[
 |(\eta_n*G)(t)-G(t)|
 \le\int_{\R}\eta_n(s)|s|\,\mathrm d\lambda_1(s)\le1/n.
\]
Stąd $|G_n(t)-G(t)|\le2/n$.
Poza skończonym zbiorem punktów załamania $G$
funkcja $g$ jest lokalnie stała, więc
$G_n'(t)=g(t)$ dla dostatecznie dużych $n$.

Dla $p<\infty$ twierdzenie~\ref{thm:pde9-chain} daje
$\partial_j G_n(u)=G_n'(u)\partial_j u$.
Zbieżność $G_n(u)\to G(u)$ w $L^p$ wynika
z majoranty $(2|u|)^p$ dla $p$-tych potęg błędów.
Poza przeciwobrazami punktów załamania mamy
$G_n'(u)\to g(u)$, zaś na tych przeciwobrazach
$\partial_j u=0$ prawie wszędzie z lematu
~\ref{lem:pde9-levelsets}.
Zatem $G_n'(u)\partial_j u\to g(u)\partial_j u$
prawie wszędzie na całej $\Omega$.
Ponieważ $0\le G_n',g\le1$, $p$-ta potęga modułu
różnicy jest ograniczona przez $|\partial_j u|^p$.
Zbieżność zmajoryzowana i wniosek~\ref{cor:pde9-closed}
dają $\partial_j G(u)=g(u)\partial_j u$.
Wszystkie pochodne są w $L^p$ i $|G(u)|\le|u|$.
Dla $p=\infty$ stosujemy te same tożsamości lokalnie
z $p=1$, a globalnie korzystamy z oszacowań
$|G(u)|\le|u|$ i $|\partial_j G(u)|\le|\partial_j u|$.
Wzór dla $u^-$ otrzymujemy stosując wynik dla części
dodatniej do $-u$, a dla $|u|=u^++u^-$ przez liniowość.
Przyjmujemy $\operatorname{sgn}(0)=0$.

Wreszcie dla $p<\infty$ i $M\to\infty$
mamy $|T_M(u)-u|\le|u|\chi_{\{|u|>M\}}$
oraz
$\partial_j T_M(u)-\partial_j u
=-\chi_{\{|u|\ge M\}}\partial_j u$ prawie wszędzie.
Oba błędy dążą do zera w $L^p$ dzięki zbieżności
zmajoryzowanej, odpowiednio przez $|u|^p$
i $|\partial_j u|^p$.
Skończona suma norm pochodnych i funkcji daje tezę.
\end{proof}

\begin{uwaga}[Zakres reguły łańcuchowej]
Twierdzenie~\ref{thm:pde9-chain} dotyczy funkcji
klasy $C^1$ o ograniczonej pochodnej, a nie bezpośrednio
wszystkich funkcji lipschitzowskich. Powyższy argument
dla obcięć wykorzystuje w istotny sposób skończoną
liczbę punktów załamania i lemat o zbiorach poziomicowych.
Nie należy bez dodatkowego uzasadnienia zastępować
$F'(u)$ dowolnym reprezentantem pochodnej ogólnej
funkcji lipschitzowskiej.
\end{uwaga}

\section{Przedłużanie z półprzestrzeni i rola brzegu}
\label{sec:pde9-extension}

Najpierw uzasadnimy aproksymację niezależnie od przedłużania.
Niech $P=\R^{d-1}\times(0,\infty)$, $1\le p<\infty$,
$u\in W^{k,p}(P)$, $k\in\N_0$. Dla $d=1$ pomijamy $x'$.
Dla $h>0$ określamy $u_h(x',t)=u(x',t+h)$ na
$P_h=\R^{d-1}\times(-h,\infty)$. Zmiana zmiennych w
tożsamości testowej daje $\partial^\alpha u_h(x',t)
=(\partial^\alpha u)(x',t+h)$ dla $|\alpha|\le k$.
Niech $f_\alpha$ będzie przedłużeniem przez zero
$\partial^\alpha u$ jako funkcji z $L^p(P)$.
Ciągłość przesunięć w $L^p$ daje
\[
 \|\partial^\alpha(u_h-u)\|_{L^p(P)}
 \le \|f_\alpha(\cdot+he_d)-f_\alpha\|_{L^p(\R^d)}
 \longrightarrow0.
\]
Nie zakładamy, że przedłużenia $f_\alpha$ są słabymi
pochodnymi $f_0$. Analogicznie niech $g_{h,\alpha}$
będą oddzielnymi przedłużeniami przez zero pochodnych
$u_h$ z $P_h$ jako funkcje $L^p$.
Dla $0<\varepsilon<h/2$ funkcja
$w_{h,\varepsilon}=\rho_\varepsilon*g_{h,0}$
jest gładka na $\R^d$, a lokalna reguła splotu
(twierdzenie~\ref{thm:pde9-mollify}) daje na $P$
$\partial^\alpha w_{h,\varepsilon}
=\rho_\varepsilon*g_{h,\alpha}$, ponieważ nośnik jądra
przesuniętego do punktu $P$ mieści się w $P_h$.
Z twierdzenia o przybliżonej jedynce
(rozdział~\ref{chap:pde-distributions}) mamy
\[
 \|\partial^\alpha(w_{h,\varepsilon}-u_h)\|_{L^p(P)}
 \le \|\rho_\varepsilon*g_{h,\alpha}-g_{h,\alpha}\|_{L^p(\R^d)}
 \longrightarrow0
\]
dla ustalonego $h$. Ponieważ wielowskaźników jest
skończenie wiele, wybieramy $h_n\downarrow0$ oraz
$0<\varepsilon_n<h_n/2$ z błędem w $W^{k,p}(P)$
między $w_{h_n,\varepsilon_n}$ a $u_{h_n}$ mniejszym
od $1/n$. Otrzymujemy globalną zbieżność gładkich
restrykcji do $u$ w $W^{k,p}(P)$.

\begin{twierdzenie}[Odbicie Sobolewa, rząd pierwszy]
\label{thm:pde9-reflection}
Niech $\R^d_+=\{(x',t)\in\R^{d-1}\times\R:t>0\}$,
$1\le p<\infty$. Istnieje liniowy, ograniczony operator
$E:W^{1,p}(\R^d_+)\to W^{1,p}(\R^d)$ taki, że
$(Eu)|_{\R^d_+}=u$. Można przyjąć
\[
 (Eu)(x',t)=\begin{cases}u(x',t),&t>0,\\
 u(x',-t),&t<0.\end{cases}
\]
Wówczas $\|Eu\|_{W^{1,p}(\R^d)}=2^{1/p}
\|u\|_{W^{1,p}(\R^d_+)}$.
\end{twierdzenie}
\begin{proof}
Wartości na $t=0$ są nieistotne dla klas $L^p$,
a odbicie zachowuje zbiory miary zero, więc wzór jest
dobrze określony i liniowy.
Dla restrykcji funkcji globalnie gładkiej zwykłe
różniczkowanie na obu półprzestrzeniach daje
$\partial_t Eu(x',t)=\operatorname{sgn}(t)
(\partial_tu)(x',|t|)$; pochodne styczne odbijają się
bez zmiany znaku. Całkujemy przez części przeciw
testowi z $C_c^\infty(\R^d)$ osobno na obu
półprzestrzeniach, ograniczając się do prostopadłościanu,
w którego wnętrzu mieści się nośnik testu.
Dla pochodnej normalnej składniki brzegowe na $t=0$
to całki z $-u(x',0)\varphi(x',0)$ i
$+u(x',0)\varphi(x',0)$ (dla $d=1$ same wartości).
Znaki pochodzą od przeciwnych normalnych zewnętrznych.
Składniki znoszą się; dla pochodnych stycznych nie
występują. Otrzymujemy tożsamości słabych pochodnych.
Podstawienie $s=-t$ daje dla funkcji i każdej pochodnej
podwojenie $p$-tej potęgi normy. Po zsumowaniu
$\|Eu\|_{W^{1,p}}=2^{1/p}\|u\|_{W^{1,p}}$.

Dla ogólnego $u$ wybieramy gładkie restrykcje $w_n$
z poprzedniej aproksymacji. Równość norm dla różnic
daje ciąg Cauchy'ego $Ew_n$ w $W^{1,p}(\R^d)$,
więc z zupełności istnieje jego granica $v$.
Jednocześnie sam wzór odbicia określa operator na $L^p$
o normie $2^{1/p}$, więc $Ew_n\to Eu$ w $L^p$.
Jednoznaczność granicy daje $v=Eu$ prawie wszędzie;
równość norm przechodzi do granicy.
\end{proof}

\begin{stwierdzenie}[Odbicia wyższego rzędu na półprzestrzeni]
\label{prop:pde9-highreflection}
Niech $k\in\N$, $1\le p<\infty$. Istnieje ograniczony,
liniowy operator $E_k:W^{k,p}(\R^d_+)\to W^{k,p}(\R^d)$,
który jest prawym odwrotnym do ograniczenia na $\R^d_+$.
\end{stwierdzenie}
\begin{proof}
Wybieramy liczby $a_1,\ldots,a_k$ spełniające
$\sum_{j=1}^k a_j(-j)^m=1$ dla $m=0,\ldots,k-1$.
Macierz układu jest macierzą Vandermonde'a o różnych
węzłach $-1,\ldots,-k$. Zależność liniowa jej wierszy
dawałaby wielomian stopnia co najwyżej $k-1$
z $k$ różnymi pierwiastkami. Układ ma więc dokładnie
jedno rozwiązanie. Definiujemy $E_ku=u$ dla $t>0$,
a dla $t<0$
\[
 (E_ku)(x',t)=\sum_{j=1}^k a_j u(x',-jt).
\]
Zmiany zmiennych są odwracalne i zachowują zbiory
miary zero, więc wzór jest liniowy na klasach $L^p$.
Dla gładkiej restrykcji $u$ pochodne klasyczne po
dolnej stronie mają postać
\[
 \partial_t^m\partial_{x'}^\beta E_ku(x',t)
 =\sum_{j=1}^k a_j(-j)^m
 (\partial_t^m\partial_{x'}^\beta u)(x',-jt).
\]
Jeżeli $m+|\beta|\le k-1$, to wartość graniczna
przy $t\uparrow0$ jest równa odpowiedniej pochodnej
od góry dzięki układowi na $a_j$.
Indukcyjnie dla $r=0,\ldots,k-1$ całkujemy przez części
każdą zidentyfikowaną pochodną rzędu $r$ jak w
dowodzie odbicia pierwszego rzędu. Dopasowanie znosi
składniki brzegowe i identyfikuje słabą pochodną
rzędu $r+1$ z klasyczną pochodną na obu częściach.
Nie wymagamy dopasowania pochodnych rzędu $k$.

Niech $C_m=\sum_{j=1}^k |a_j|j^{m-1/p}$.
Nierówność trójkąta i podstawienie $s=-jt$ dają
\[
 \|\partial_t^m\partial_{x'}^\beta E_ku\|_{L^p(t<0)}
 \le C_m\|\partial_t^m\partial_{x'}^\beta u\|_{L^p(P)}.
\]
Istotnie, w $p$-tej potędze normy pojedynczego składnika
różniczkowanie daje czynnik $j^{mp}$, a jakobian $j^{-1}$.
Dodając normy po górnej stronie i sumując, otrzymujemy
\[
 \|E_ku\|_{W^{k,p}(\R^d)}
 \le \max_{0\le m\le k}(1+C_m^p)^{1/p}
 \|u\|_{W^{k,p}(P)}.
\]
Dla ogólnego $u$ stosujemy poprzednią niezależną
aproksymację rzędu $k$. Oszacowanie różnic i zupełność
dają granicę przedłużeń w $W^{k,p}(\R^d)$.
Oszacowanie dla $m=0,\beta=0$ działa na samych
funkcjach $L^p$, więc granica jest równa funkcji
określonej przez jawny wzór na $E_k$. Ograniczenie do $P$
jest równe $u$, a oszacowanie przechodzi do granicy.
\end{proof}

\begin{uwaga}[Obszary Lipschitzowskie]
Jeśli ograniczony obszar $\Omega$ ma brzeg lokalnie
będący wykresem funkcji Lipschitzowskiej, można prostować
brzeg dwulipschitzowskimi zmianami współrzędnych,
zastosować odbicie pierwszego rzędu i skleić lokalne
przedłużenia podziałem jedności. Otrzymujemy ograniczony
operator $W^{1,p}(\Omega)\to W^{1,p}(\R^d)$.
Przypadek $W^{k,p}$ dla $k>1$ na dowolnym obszarze Lipschitzowskim wymaga bardziej subtelnej konstrukcji:
proste prostowanie Lipschitzowskie nie zachowuje ogólnie
drugich pochodnych w $L^p$. Pełną teorię przedłużania
i śladów rozwiniemy w następnym rozdziale.
\end{uwaga}

\begin{lemat}[Gęstość funkcji gładkich aż do płaskiego brzegu]
\label{lem:pde9-density-halfspace-slab}
Niech $1\le p<\infty$ oraz $d\ge1$. Dla
$G=\R^{d-1}\times(0,\infty)$ albo
$G=\R^{d-1}\times(0,1)$ każdą funkcję $u\in W^{1,p}(G)$
można przybliżyć w normie $W^{1,p}(G)$ restrykcjami funkcji
z $C^\infty(\R^d)$. Można przy tym wybrać przybliżenia
z $C_c^\infty(\R^d)$.
\end{lemat}
\begin{proof}
Dla półprzestrzeni korzystamy z twierdzenia
~\ref{thm:pde9-reflection}. Dla warstwy wybieramy
gładką funkcję $\theta(t)$ równą $1$ dla $t\le1/3$
i $0$ dla $t\ge2/3$, o wartościach w $[0,1]$.
Rozkładamy $u=u_0+u_1$, gdzie $u_0=\theta u$,
$u_1=(1-\theta)u$. Reguła iloczynu
(twierdzenie~\ref{thm:pde9-products}) daje ograniczone
liniowe odwzorowania $u\mapsto u_i$ w $W^{1,p}(G)$.

Przedłużamy $u_0$ przez zero dla $t\ge1$ do $P$.
Dla testu na $P$ mnożymy go przez gładką funkcję
zmiennej $t$, równą $1$ dla $t\le2/3$ i $0$ dla
$t\ge5/6$. Otrzymujemy test na $G$. Tam, gdzie
zmienia się odcięcie, $u_0$ i jego słabe pochodne są
zerowe dzięki lokalności pochodnych. Tożsamość testowa
na $G$ staje się zatem tożsamością na $P$ dla
przedłużeń $u_0$ i jego pochodnych przez zero.
Normy pozostają równe; nie zakładamy zerowego śladu
przy $t=0$. Dla $u_1$ stosujemy zmianę $s=1-t$
i tę samą konstrukcję względem $s$. Zmiana zachowuje
normę, a pochodna normalna zmienia znak, co wynika
z podstawienia w tożsamości testowej.
Obie uzyskane funkcje z $W^{1,p}(P)$ przedłużamy
przez odbicie. Drugie przedłużenie przenosimy
z powrotem do zmiennej $t=1-s$ i dodajemy do pierwszego.
Na $G$ suma jest równa $u$. Nierówność trójkąta,
oszacowania iloczynów i równość norm odbicia dają
ograniczony liniowy operator przedłużania.
Nie korzystamy z nieudowodnionej jeszcze własności
absolutnej ciągłości na prawie wszystkich prostych.

W obu przypadkach niech $v=Eu$. Gęstość
$C_c^\infty(\R^d)$ w $W^{1,p}(\R^d)$ udowodniona
w sekcji~\ref{sec:pde9-density} daje $v_n\to v$
w tej normie. Ograniczenie nie zwiększa normy,
więc $\|v_n|_G-u\|_{W^{1,p}(G)}
\le\|v_n-v\|_{W^{1,p}(\R^d)}\to0$.
\end{proof}


\begin{uwaga}[Aproksymacja gładka nie oznacza zerowego śladu]
\label{rem:pde-v42-density}
Twierdzenie Meyersa--Serrina~\ref{thm:pde9-MeyersSerrin}
nie orzeka, że $C_c^\infty(\Omega)$ jest gęste w całej
przestrzeni $W^{1,p}(\Omega)$. Na przykład funkcja
$u\equiv1$ na ograniczonym obszarze Lipschitzowskim
należy do $W^{1,p}(\Omega)$, ale jej ślad jest równy $1$,
a więc nie należy do $W_0^{1,p}(\Omega)$ na mocy
twierdzenia~\ref{thm:pde10-trace-zero}.
Są to trzy różne twierdzenia: aproksymacja wewnętrzna,
gęstość gładkich restrykcji aż do brzegu oraz
charakterystyka jądra śladu. Ten ostatni rezultat
zostanie udowodniony w rozdziale~\ref{chap:pde-embeddings}.
\end{uwaga}

\section{Ślady na płaskim brzegu: pierwszy wynik}
\label{sec:pde9-traces}

\begin{twierdzenie}[Oszacowanie śladu na warstwie]
\label{thm:pde9-trace-slab}
Niech $S=\R^{d-1}\times(0,1)$ i $1\le p<\infty$.
Istnieje dokładnie jeden ograniczony operator liniowy
$\operatorname{Tr}_0:W^{1,p}(S)\to L^p(\R^{d-1})$
zgodny z ograniczeniem do $t=0$ dla restrykcji funkcji
gładkich na otoczeniu $\overline S$, należących do $W^{1,p}(S)$.
Dla $d=1$ przestrzeń wartości jest równa $\R$ z normą
modułu (ślad w punkcie). Zachodzi
\[
 \|\operatorname{Tr}_0u\|_{L^p(\R^{d-1})}
 \le \|u\|_{L^p(S)}+\|\partial_tu\|_{L^p(S)}.
\]
Analogiczne twierdzenie zachodzi dla śladu $\operatorname{Tr}_1$
na $t=1$.
\end{twierdzenie}
\begin{proof}
Dla funkcji $h$ klasy $C^1$ na $[0,1]$ i $s\in(0,1)$
z podstawowego twierdzenia rachunku całkowego otrzymujemy
\[
 |h(0)|\le |h(s)|+\int_0^1|h'(t)|\,\mathrm d\lambda_1(t).
\]
Całkujemy po $s$ i stosujemy nierówność Höldera
(na przedziale miary $1$; dla $p=1$ nie jest potrzebna):
\[
 |h(0)|\le\|h\|_{L^p(0,1)}+\|h'\|_{L^p(0,1)}.
\]
Dla gładkiej restrykcji $u$ stosujemy to do
$h(t)=u(x',t)$. Z twierdzenia Tonellego funkcje
\[
 A(x')=\left(\int_0^1 |u(x',t)|^p\,\mathrm d\lambda_1(t)\right)^{1/p},
 \qquad
 B(x')=\left(\int_0^1 |\partial_tu(x',t)|^p\,\mathrm d\lambda_1(t)\right)^{1/p}
\]
są mierzalne, skończone prawie wszędzie, a ich normy
$L^p(\R^{d-1})$ są równe odpowiednim normom na $S$.
Nierówność punktowa $|u(x',0)|\le A(x')+B(x')$
oraz nierówność trójkąta w $L^p$ dają żądane
oszacowanie. Dla $d=1$ wystarcza bezpośrednio
oszacowanie jednowymiarowe.

Z lematu~\ref{lem:pde9-density-halfspace-slab} dla dowolnego
$u\in W^{1,p}(S)$ istnieje ciąg gładkich restrykcji $u_n$
zbieżny do $u$ w tej przestrzeni. Oszacowanie zastosowane
do $u_n-u_m$ pokazuje, że $u_n(\cdot,0)$ jest ciągiem
Cauchy'ego w przestrzeni Banacha $L^p(\R^{d-1})$
(dla $d=1$ w $\R$). Definiujemy $\operatorname{Tr}_0u$
jako jego granicę. Jeżeli $v_n$ jest drugim ciągiem
aproksymującym, to $u_n-v_n\to0$ w $W^{1,p}$, a więc
różnica ich śladów dąży do zera w $L^p$. Definicja
nie zależy od wyboru ciągu. Aproksymując kombinacje
liniowe przez te same kombinacje aproksymantów,
otrzymujemy liniowość. Przejście do granicy w
oszacowaniu dla $u_n$ daje oszacowanie dla $u$.
Dla gładkiej restrykcji można użyć ciągu stałego,
więc operator zgadza się z ograniczeniem klasycznym.
Każdy inny ograniczony operator o tej własności
musi dawać na $u$ tę samą granicę, co dowodzi
jednoznaczności. Zmiana zmiennej $t\mapsto1-t$
zachowuje normy funkcji i modułów pochodnych normalnych,
co daje twierdzenie dla $\operatorname{Tr}_1$.
\end{proof}

Jeśli $u\in W^{1,p}_0(S)$, oba ślady są zerowe:
aproksymanty z $C_c^\infty(S)$ znikają w otoczeniu
obu płaszczyzn brzegowych, a operatory śladu są
ciągłe. Nie dowodzimy tu jeszcze odwrotnej implikacji
dla ogólnych obszarów; por. rozdział~\ref{chap:pde-embeddings}.

\section{Perspektywa wariacyjna: pierwsze słabe rozwiązanie}
\label{sec:pde9-weakPDE}

\begin{przyklad}[Równanie Poissona na przedziale]
\label{ex:pde9-dirichlet}
Niech $I=(0,1)$, $f\in L^2(I)$ i rozważmy
$-u''=f$ na $I$ z warunkami $u(0)=u(1)=0$.
Jeśli $u$ jest gładkie i spełnia równanie, mnożymy je
przez $\varphi\in C_c^\infty(I)$ i całkujemy przez części:
\[
 \int_0^1u'(x)\varphi'(x)\,\mathrm d\lambda_1(x)
 =\int_0^1 f(x)\varphi(x)\,\mathrm d\lambda_1(x).
\]
Obie strony mają sens dla $u,\varphi\in H^1(I)$
dzięki nierówności Cauchy'ego--Schwarza.
Szukamy $u\in H_0^1(I)$ spełniającego tę tożsamość
dla każdego $\varphi\in H_0^1(I)$.
Jest to \emph{słabe zagadnienie Dirichleta}; zerowe
wartości brzegowe reprezentanta $u$ wynikają z
ciągłości śladów w wymiarze jeden
(wniosek~\ref{cor:pde9-trace-oneD}) i definicji $H_0^1$.

Wykażemy istnienie jawnie. Ponieważ $f\in L^1(I)$,
funkcja $F(x)=\int_0^x f(y)\,\mathrm d\lambda_1(y)$
jest absolutnie ciągła na $[0,1]$ i $F'=f$ prawie wszędzie.
Połóżmy
\[
 c=\int_0^1 F(s)\,\mathrm d\lambda_1(s),\qquad
 u(x)=cx-\int_0^x F(s)\,\mathrm d\lambda_1(s).
\]
Funkcja $F$ jest ciągła i ograniczona na $[0,1]$,
zatem $u$ jest klasy $C^1$ i $u'=c-F$ jest
absolutnie ciągła. Mamy $u''=-f$ prawie wszędzie,
a z jednowymiarowej charakterystyki słabych pochodnych
(twierdzenie~\ref{thm:pde9-ac}) wynika $u\in H^2(I)$.
Ponadto $u(0)=0$ i $u(1)=c-c=0$.

Uzasadnimy jeszcze $u\in H_0^1(I)$ bez zakładania
ogólnego twierdzenia o jądrze śladu. Dla
$0<\varepsilon<1/4$ wybierzmy gładkie odcięcie
$\eta_\varepsilon$ o wartościach w $[0,1]$, równe
$1$ na $[2\varepsilon,1-2\varepsilon]$, równe $0$
na $(0,\varepsilon)\cup(1-\varepsilon,1)$,
z $|\eta_\varepsilon'|\le C/\varepsilon$,
gdzie $C$ nie zależy od $\varepsilon$.
Oznaczmy $M=\sup_{[0,1]}|u'|$. Z zerowych wartości
brzegowych i całkowania pochodnej wynika
$|u(x)|\le M\min(x,1-x)$. Na zbiorze, gdzie
$\eta_\varepsilon'\ne0$, mamy więc
$|\eta_\varepsilon' u|\le 2CM$, a miara tego zbioru
jest nie większa niż $4\varepsilon$. Stąd
$\|\eta_\varepsilon' u\|_2^2\le16 C^2M^2\varepsilon$.
Pozostałe składniki błędu,
$(1-\eta_\varepsilon)u$ i $(1-\eta_\varepsilon)u'$,
dążą do zera w $L^2(I)$ na mocy twierdzenia o
zbieżności zmajoryzowanej. Reguła Leibniza daje więc
$\eta_\varepsilon u\to u$ w $H^1(I)$.
Każda funkcja $\eta_\varepsilon u$ ma zwarty nośnik
w $I$ i należy do $H_0^1(I)$ dzięki wygładzaniu
z zachowaniem zwartego nośnika
(lemat~\ref{lem:pde9-compact-product}, zastosowany do
$u$ i $\eta_\varepsilon$). Domkniętość $H_0^1(I)$
daje $u\in H_0^1(I)$.

Dla testów z $C_c^\infty(I)$ tożsamość słaba wynika
z równania $-u''=f$ w sensie dystrybucji. Dla
$\varphi\in H_0^1(I)$ wybieramy testy $\varphi_n\to\varphi$
w $H^1(I)$. Obie strony tożsamości zbiegają do
odpowiednich granic, ponieważ
\[
 \left|\int_0^1u'(\varphi_n'-\varphi')\,\mathrm d\lambda_1\right|
 \le \|u'\|_2\|\varphi_n'-\varphi'\|_2,\qquad
 \left|\int_0^1f(\varphi_n-\varphi)\,\mathrm d\lambda_1\right|
 \le \|f\|_2\|\varphi_n-\varphi\|_2.
\]
Otrzymaliśmy zatem rozwiązanie słabe.

Twierdzenie Fubiniego (stosowalne, gdyż $f\in L^1(I)$
i wszystkie użyte wagi są ograniczone) daje
$c=\int_0^1(1-y)f(y)\,\mathrm d\lambda_1(y)$ oraz
$\int_0^xF(s)\,\mathrm d\lambda_1(s)
=\int_0^x(x-y)f(y)\,\mathrm d\lambda_1(y)$.
Po rozdzieleniu pierwszej całki na $(0,x)$ i $(x,1)$
otrzymujemy równoważny wzór
\[
 u(x)=\int_0^1G(x,y)f(y)\,\mathrm d\lambda_1(y),\qquad
 G(x,y)=\begin{cases}y(1-x),&y\le x,\\x(1-y),&x<y.\end{cases}
\]
Dla $y\in I$ funkcja $G(\cdot,y)$ jest ciągła,
zeruje się na końcach i ma pochodną $1-y$ dla $x<y$
oraz $-y$ dla $x>y$. Całkując przeciw pochodnej
testu, otrzymujemy
\[
 \int_0^1\partial_xG(x,y)\varphi'(x)\,\mathrm d\lambda_1(x)
 =(1-y)\varphi(y)+y\varphi(y)=\varphi(y).
\]
Zatem $-\partial_x^2G(\cdot,y)=\delta_y$.
Dla $f\equiv1$ powyższa konstrukcja daje $F(x)=x$,
$c=1/2$ i $u(x)=x(1-x)/2$.

Jeśli $u_1,u_2$ są dwoma słabymi rozwiązaniami,
to $w=u_1-u_2\in H_0^1(I)$ spełnia jednorodną
tożsamość. Podstawiając $\varphi=w$ otrzymujemy
$\|w'\|_2^2=0$. Nierówność Poincarégo
(wniosek~\ref{cor:pde9-Poincare-zero}, $L=1$)
daje $\|w\|_2=0$, a więc jednoznaczność.
Podstawiając $\varphi=u$ dla dowolnego $f$, mamy
\[
 \|u'\|_2^2=\int_0^1fu\,\mathrm d\lambda_1
 \le\|f\|_2\|u\|_2\le\|f\|_2\|u'\|_2 .
\]
Jeśli $\|u'\|_2=0$, oszacowanie jest natychmiastowe;
w przeciwnym razie dzielimy przez tę normę.
Dostajemy $\|u'\|_2\le\|f\|_2$ i
$\|u\|_{H^1(I)}\le\sqrt{2}\|f\|_2$.
\end{przyklad}

\section{Skalowanie, krytyczne wykładniki i zadania na przyszłość}
\label{sec:pde9-scaling}

Niech $\phi\in C_c^\infty(\R^d)$, $\phi\not\equiv0$,
$a\in\R$, a dla $\varepsilon>0$ połóżmy
$u_\varepsilon(x)=\varepsilon^{-a}\phi(x/\varepsilon)$.
Dla $1\le p,q<\infty$ podstawienie $x=\varepsilon y$
o jakobianie $\varepsilon^d$ daje
\[
 \|u_\varepsilon\|_{L^q}^q
 =\varepsilon^{d-aq}\int_{\R^d}|\phi(y)|^q\,\mathrm d\lambda_d(y),
 \qquad
 \nabla u_\varepsilon(x)=\varepsilon^{-a-1}\nabla\phi(x/\varepsilon).
\]
Stąd, przy normie euklidesowej gradientu,
\[
 \|u_\varepsilon\|_{L^q}
 =\varepsilon^{d/q-a}\|\phi\|_{L^q},\qquad
 \|\nabla u_\varepsilon\|_{L^p}
 =\varepsilon^{d/p-a-1}\|\nabla\phi\|_{L^p}.
\]
Te same wykładniki obowiązują dla $q=\infty$ lub
$p=\infty$ z konwencją $1/\infty=0$, ponieważ
skalowanie argumentu zachowuje zbiory miary zero,
a mnożenie przez stałą skaluje istotne supremum.

Jeśli nierówność jednorodna
$\|u\|_{L^q}\le C\|\nabla u\|_{L^p}$ na $\R^d$
ma obowiązywać z jedną stałą $C$ dla wszystkich
$u\in C_c^\infty(\R^d)$, to po podstawieniu
$u_\varepsilon$ i podzieleniu przez odpowiednią
potęgę $\varepsilon$ dostajemy
\[
 \varepsilon^{d/q-d/p+1}\|\phi\|_{L^q}
 \le C\|\nabla\phi\|_{L^p}\qquad(\varepsilon>0).
\]
Norma $\|\phi\|_{L^q}$ jest dodatnia. Jeżeli wykładnik
po lewej jest ujemny, otrzymujemy sprzeczność przy
$\varepsilon\downarrow0$, a jeżeli dodatni,
przy $\varepsilon\to\infty$. Koniecznie więc
$d/q=d/p-1$. Dla $1\le p<d$ otrzymujemy wykładnik
\[
 p^*=\frac{dp}{d-p}.
\]
Warunek skalowania jest konieczny, lecz sam nie dowodzi
nierówności Sobolewa. W szczególności dla $p=d$ wskazuje
tylko potencjalny wykładnik $q=\infty$, nie gwarantuje
prawdziwości oszacowania w tym przypadku.

Skalowanie daje również kontrprzykład do ciągłego
zanurzenia w zbyt wysokim wykładniku nawet na dowolnym
niepustym zbiorze otwartym $\Omega\subset\R^d$.
Niech $1\le p<d$, $p^*<q\le\infty$,
$B_r(x_0)\subset\Omega$ i wybierzmy niezerową
$\phi\in C_c^\infty(B_1(0))$. Dla $0<\varepsilon<r$
połóżmy
$w_\varepsilon(x)=\varepsilon^{1-d/p}
\phi((x-x_0)/\varepsilon)$.
Wtedy $w_\varepsilon\in C_c^\infty(\Omega)$, a
\[
 \|w_\varepsilon\|_{L^p(\Omega)}=\varepsilon\|\phi\|_p,
 \qquad
 \|\partial_j w_\varepsilon\|_{L^p(\Omega)}
 =\|\partial_j\phi\|_p,\qquad
 \|w_\varepsilon\|_{L^q(\Omega)}
 =\varepsilon^{d/q-d/p+1}\|\phi\|_q.
\]
Pierwsze dwa wzory pokazują ograniczoność w
$W^{1,p}(\Omega)$ przy $\varepsilon\downarrow0$,
natomiast ostatnia norma dąży do nieskończoności,
ponieważ $d/q<d/p-1$. Wyklucza to ciągłe
zanurzenie $W^{1,p}(\Omega)\hookrightarrow L^q(\Omega)$
dla $q>p^*$, również w przypadku $q=\infty$.
Nie twierdzimy na podstawie samego tego argumentu,
że istnieje pojedynczy element $W^{1,p}(\Omega)$
nienależący do $L^q(\Omega)$. W następnym rozdziale
udowodnimy dodatnie twierdzenia o zanurzeniach
i zbadamy ich zwartość.

\section{Podsumowanie i zależności logiczne}
Wprowadziliśmy przestrzenie $W^{k,p}$ na dowolnych
niepustych zbiorach otwartych, wykazaliśmy zupełność
i zamkniętość pochodnych. Jednowymiarowa charakterystyka
przez absolutną ciągłość dała ślady w punktach.
Dla $p<\infty$ udowodniliśmy lokalne wygładzanie
i twierdzenie Meyersa--Serrina oraz ograniczone
przedłużenia z półprzestrzeni; przypadki $p=\infty$
rozpatrywaliśmy osobno tam, gdzie były dopuszczalne.
Oddzieliliśmy $W^{k,p}(\Omega)$ od domknięcia
$W^{k,p}_0(\Omega)$, którego nie utożsamiamy
na dowolnym zbiorze otwartym z jądrem operatora śladu.
Udowodniliśmy nierówności Poincarégo na kostkach,
reguły rachunkowe i oszacowanie śladu na warstwie,
a także skonstruowaliśmy pierwsze słabe rozwiązanie
równania Poissona. W rozdziale~\ref{chap:pde-embeddings}
wykorzystamy te narzędzia do ogólnych twierdzeń o
zanurzeniach, zwartości Rellicha--Kondrachova
i teorii śladów na brzegach Lipschitzowskich.


% Rozdział 10. Samodzielna część wykładu; włączana przed dodatkiem z kodem.

\chapter[Zanurzenia i zwartość Sobolewa]{Zanurzenia Sobolewa, nierówność Morreya, zwartość i ślady}
\label{chap:pde-embeddings}

Przestrzeń $W^{1,p}(\Omega)$ została zdefiniowana przez całkowalność
funkcji i jej słabych pochodnych. Nie wynika z tej definicji wprost,
czy funkcja jest ciągła, czy można ją przybliżać jednostajnie
ani czy ciąg ograniczony w $W^{1,p}$ ma podciąg zbieżny w $L^p$.
Odpowiedzi zależą przede wszystkim od relacji $p$ do wymiaru $d$.
Będziemy odróżniać \emph{zanurzenie ciągłe}, oznaczane
$X\hookrightarrow Y$, od \emph{zanurzenia zwartego}, oznaczanego
$X\hookrightarrow\hookrightarrow Y$. Pierwsze znaczy, że
$\|u\|_Y\le C\|u\|_X$; drugie, że każdy ciąg ograniczony w $X$
ma podciąg zbieżny w normie $Y$.

W rozdziale $d\in\N$, $d\ge1$, i $1\le p\le\infty$.
Przyjmujemy, że $\Omega$ jest niepustym otwartym podzbiorem
$\R^d$, a funkcje są rzeczywiste, o ile nie podano inaczej.
Dla gradientu używamy punktowej normy euklidesowej.
Norma ta jest równoważna normie składowych używanej
w rozdziale~\ref{chap:pde-sobolev}: dla $p<\infty$ wynika
to z równoważności norm $\ell^2$ i $\ell^p$ na $\R^d$,
po podniesieniu do potęgi $p$ i scałkowaniu, a dla
$p=\infty$ z nierówności
$\max_j|v_j|\le |v|\le\sqrt d\max_j|v_j|$.
Nie zakładamy automatycznie,
że $\Omega$ ma regularny brzeg. Kiedy korzystamy z obszaru Lipschitzowskiego, podajemy to wyraźnie. Całki na $\R^d$ i na
$\Omega$ są całkami względem $\lambda_d$; na brzegu pojawia się
miara powierzchniowa $\mathcal H^{d-1}$.

\section{Nierówność geometryczna potrzebna do dowodu Sobolewa}
\label{sec:pde10-LW}

Zaczniemy od rachunkowego lematu, który pozwala udowodnić
nierówność Sobolewa bez odwoływania się do twierdzenia
izoperymetrycznego ani do teorii potencjałów Riesza.
Dla $x=(x_1,\ldots,x_d)$ przez $\widehat x_j$ oznaczamy
wektor powstały przez opuszczenie współrzędnej $x_j$.

\begin{lemat}[Nierówność Loomisa--Whitneya w postaci całkowej]
\label{lem:pde10-LW}
Niech $d\ge2$ i $F_j:\R^{d-1}\to[0,\infty]$ będą mierzalne,
$j=1,\ldots,d$. Wówczas
\[
 \int_{\R^d}\prod_{j=1}^d F_j(\widehat x_j)^{1/(d-1)}
 \,\mathrm d\lambda_d(x)
 \le \prod_{j=1}^d
 \left(\int_{\R^{d-1}}F_j\,\mathrm d\lambda_{d-1}\right)^{1/(d-1)}.
\]
Jeżeli pewna całka po prawej stronie jest zerowa, to
lewa strona również jest zerowa, nawet gdy inne całki
są nieskończone. W tym przypadku przyjmujemy prawą
stronę równą zeru. Jeżeli wszystkie całki są dodatnie,
a przynajmniej jedna nieskończona, prawa strona jest
nieskończona. W iloczynie podcałkowym przyjmujemy wartość
zero, gdy którykolwiek czynnik jest zerowy, również
w obecności czynników nieskończonych; odpowiada to
granicy iloczynów nieujemnych obcięć.
\end{lemat}
\begin{proof}
Wystarczy najpierw rozważyć nieujemne funkcje ograniczone
o zwartych nośnikach. Dla $d=2$ nierówność jest równością:
$F_1$ zależy tylko od $x_2$, a $F_2$ tylko od $x_1$.
Z twierdzenia Tonellego otrzymujemy
\[
 \int_{\R^2}F_1(x_2)F_2(x_1)\,\mathrm d\lambda_2(x)
 =\left(\int_\R F_1\,\mathrm d\lambda_1\right)
  \left(\int_\R F_2\,\mathrm d\lambda_1\right).
\]
W rozważanym przypadku ograniczonych funkcji o zwartych
nośnikach wszystkie całki są skończone.

Załóżmy indukcyjnie tezę w wymiarze $d-1$, gdzie $d\ge3$.
Całkujemy w lewej stronie po $x_d$. Czynnik
$F_d(x_1,\ldots,x_{d-1})$ nie zależy od $x_d$.
Nierówność H\"oldera zastosowana do pozostałych $d-1$
czynników, każdy z wykładnikiem $d-1$, daje
\[
 \int_{\R}\prod_{j<d}F_j(\widehat x_j)^{1/(d-1)}
 \,\mathrm d\lambda_1(x_d)
 \le\prod_{j<d}G_j(\widehat x^{\prime}_j)^{1/(d-1)},
\]
gdzie $x'=(x_1,\ldots,x_{d-1})$, a $G_j$ powstaje przez
scałkowanie $F_j$ względem $x_d$; nie zależy więc od $x_j$.
Oznaczmy $H(x')=\prod_{j<d}G_j(\widehat x'_j)^{1/(d-2)}$.
Zapiszmy dla skrócenia wzorów
$I_d=\int_{\R^{d-1}}F_d(x')\,\mathrm d\lambda_{d-1}(x')$,
$J=\int_{\R^{d-1}}H(x')\,\mathrm d\lambda_{d-1}(x')$
oraz, dla $j<d$,
$I_j=\int_{\R^{d-2}}G_j(z)\,\mathrm d\lambda_{d-2}(z)$.
Po użyciu H\"oldera na $\R^{d-1}$ z wykładnikami
$d-1$ oraz $(d-1)/(d-2)$ otrzymujemy
\begin{align*}
 &\int_{\R^{d-1}}F_d(x')^{1/(d-1)}
 \prod_{j<d}G_j(\widehat x'_j)^{1/(d-1)}
 \,\mathrm d\lambda_{d-1}(x')\\
 &\qquad\le I_d^{1/(d-1)}J^{(d-2)/(d-1)}
 \le I_d^{1/(d-1)}\prod_{j<d} I_j^{1/(d-1)}.
\end{align*}
Druga nierówność jest dokładnie założeniem indukcyjnym.
Na mocy twierdzenia Fubiniego
$I_j=\int_{\R^{d-1}}F_j(z)\,\mathrm d\lambda_{d-1}(z)$,
co kończy krok indukcyjny. Dla dowolnych nieujemnych $F_j$
kładziemy
$F_{j,n}=\min(F_j,n)\chi_{[-n,n]^{d-1}}$.
Funkcje te są ograniczone, mają nośniki zawarte w zbiorach
zwartych i rosną punktowo do $F_j$. Iloczyny
$\prod_j F_{j,n}(\widehat x_j)^{1/(d-1)}$ również rosną
do iloczynu z tezy zgodnie z podaną konwencją.
Twierdzenie o zbieżności monotonicznej stosujemy do lewej
strony i do każdej całki po prawej stronie.
Jeżeli $\int F_j=0$ dla pewnego $j$, to wszystkie całki
$\int F_{j,n}$ są zerowe, więc lewa strona jest zerowa
w granicy. W przeciwnym razie wszystkie granice całek
są dodatnie; iloczyn prawych stron zbiega do ich iloczynu,
również gdy jedna z granic jest nieskończona.
To dowodzi nierówności w pełnej ogólności.
\end{proof}

\begin{twierdzenie}[Nierówność Sobolewa dla $p=1$]
\label{thm:pde10-Sob1}
Dla $d\ge2$ i każdej $u\in C_c^1(\R^d)$ zachodzi
\[
 \|u\|_{L^{d/(d-1)}(\R^d)}
 \le\prod_{j=1}^d\|\partial_j u\|_{L^1(\R^d)}^{1/d}
 \le \frac1d\sum_{j=1}^d\|\partial_j u\|_{L^1(\R^d)}.
\]
W szczególności $\|u\|_{L^{d/(d-1)}}\le C_d\|\nabla u\|_{L^1}$.
\end{twierdzenie}
\begin{proof}
Dla ustalonego $j$ i każdego $x\in\R^d$ funkcja jednej
zmiennej $t\mapsto u(x_1,\ldots,t,\ldots,x_d)$ jest
klasy $C^1$ i znika dla dostatecznie ujemnych $t$ dzięki
zwartości nośnika $u$. Z podstawowego twierdzenia
rachunku całkowego
\[
 u(x)=\int_{-\infty}^{x_j}
 \partial_j u(x_1,\ldots,t,\ldots,x_d)\,\mathrm d\lambda_1(t).
\]
Biorąc moduł i rozszerzając przedział całkowania
nieujemnej funkcji, uzyskujemy
\[
 |u(x)|\le A_j(\widehat x_j),\qquad
 A_j(\widehat x_j)=\int_{\R}
 |\partial_j u(x_1,\ldots,t,\ldots,x_d)|\,\mathrm d\lambda_1(t).
\]
Mnożąc $d$ nierówności i podnosząc do potęgi $1/(d-1)$,
dostajemy $|u(x)|^{d/(d-1)}\le\prod_j
A_j(\widehat x_j)^{1/(d-1)}$. Lemat~\ref{lem:pde10-LW}
daje zatem
\[
 \|u\|_{d/(d-1)}^{d/(d-1)}
 \le \prod_{j=1}^d\|\partial_j u\|_1^{1/(d-1)}.
\]
Po podniesieniu do potęgi $(d-1)/d$ otrzymujemy pierwszą
nierówność. Druga jest nierównością pomiędzy średnią
geometryczną i arytmetyczną. Ponadto z nierówności
Cauchy'ego--Schwarza punktowo
$\sum_j|\partial_j u|\le\sqrt d\,|\nabla u|$.
Całkując, możemy więc przyjąć $C_d=1/\sqrt d$.
Wszystkie całki są skończone
z uwagi na zwarty nośnik $u$.
\end{proof}

\section{Ostre wykładniki zanurzenia Sobolewa}
\label{sec:pde10-Sobolev}

\begin{twierdzenie}[Nierówność Sobolewa na całej przestrzeni]
\label{thm:pde10-Sobp}
Niech $d\ge2$ i $1\le p<d$. Przyjmijmy
$p^*=dp/(d-p)$. Istnieje stała $C(d,p)>0$ taka, że
\[
 \|u\|_{L^{p^*}(\R^d)}
 \le C(d,p)\|\nabla u\|_{L^p(\R^d)}
 \quad\text{dla każdego }u\in C_c^\infty(\R^d).
\]
W konsekwencji identyczne oszacowanie obowiązuje
(z ewentualnie inną stałą) dla $u\in W^{1,p}(\R^d)$.
\end{twierdzenie}
\begin{proof}
Dla $p=1$ jest to twierdzenie~\ref{thm:pde10-Sob1}.
Niech $1<p<d$ i $u\in C_c^\infty(\R^d)$. Wybierzmy
$\gamma=p(d-1)/(d-p)>1$. Stosujemy twierdzenie
\ref{thm:pde10-Sob1} do funkcji
$v=|u|^\gamma$. Ponieważ $\gamma>1$, funkcja
$t\mapsto |t|^\gamma$ jest klasy $C^1$ na $\R$:
dla $t\ne0$ jej pochodna wynosi
$\gamma|t|^{\gamma-1}\operatorname{sgn}(t)$,
a w zerze iloraz różnicowy dąży do zera. Pochodna
jest ciągła także w zerze. Zatem $v\in C_c^1(\R^d)$,
gdyż $u$ jest gładka i $v=0$ poza nośnikiem $u$.
Jest to dokładnie klasa funkcji dopuszczalna w twierdzeniu
dla $p=1$.
Reguła łańcuchowa oraz nierówność H\"oldera dają
\begin{align*}
 \|u\|_{L^{\gamma d/(d-1)}}^\gamma
 &\le C_d\gamma\int_{\R^d}|u|^{\gamma-1}|\nabla u|
 \,\mathrm d\lambda_d\\
 &\le C_d\gamma
 \|u\|_{L^{(\gamma-1)p'}}^{\gamma-1}
 \|\nabla u\|_{L^p},\qquad p'=\frac p{p-1}.
\end{align*}
Z wyboru $\gamma$ wynika
\[
 \gamma\frac d{d-1}=\frac{dp}{d-p}=p^*,
 \qquad
 \gamma-1=\frac{d(p-1)}{d-p},
 \qquad
 (\gamma-1)p'=\frac{dp}{d-p}=p^*.
\]
Jeśli $u\not\equiv0$, dzielimy obie strony przez
$\|u\|_{p^*}^{\gamma-1}$; przypadek $u=0$ jest oczywisty.
To daje nierówność dla funkcji gładkich.

Jeśli $u\in W^{1,p}(\R^d)$, wybieramy
$u_n\in C_c^\infty(\R^d)$ zbieżne do $u$ w $W^{1,p}$.
Istnienie takiego ciągu wynika z gęstości
$C_c^\infty(\R^d)$ w $W^{1,p}(\R^d)$ udowodnionej
w sekcji~\ref{sec:pde9-density}; tu $p<\infty$.
Stosujemy oszacowanie do różnicy $u_n-u_m$:
\[
 \|u_n-u_m\|_{p^*}
 \le C(d,p)\|\nabla u_n-\nabla u_m\|_p\longrightarrow0.
\]
Zupełność $L^{p^*}$ daje granicę $v$.
Dla dowolnego $t>0$ nierówność trójkąta i nierówność
Czebyszewa dają
\[
 \lambda_d(\{|u-v|>t\})
 \le (2/t)^p\|u-u_n\|_p^p
     +(2/t)^{p^*}\|u_n-v\|_{p^*}^{p^*}.
\]
Prawa strona dąży do zera, zatem, biorąc $t=1/m$,
$m\in\N$, otrzymujemy $u=v$ prawie wszędzie.
Ciągłość normy po obu stronach nierówności dla $u_n$
daje tezę dla $u$ z tą samą stałą $C(d,p)$.
\end{proof}

\begin{uwaga}[Dlaczego wykładnik jest krytyczny?]
\label{rem:pde10-critical}
Niech $d\ge2$, $1\le p<d$ oraz $1\le q\le\infty$.
Dla $\varphi\in C_c^\infty(B_1(0))$, $\varphi\ne0$, oraz
$x_0\in\R^d$ niech
$u_\varepsilon(x)=\varepsilon^{1-d/p}
\varphi((x-x_0)/\varepsilon)$. Wówczas
\[
 \|\nabla u_\varepsilon\|_p=\|\nabla\varphi\|_p,
 \qquad \|u_\varepsilon\|_q
 =\varepsilon^{1-d/p+d/q}\|\varphi\|_q.
\]
Wzory wynikają z podstawienia $x=x_0+\varepsilon y$,
reguły łańcuchowej oraz (dla $q=\infty$) zachowania
zbiorów miary zero, jak w sekcji~\ref{sec:pde9-scaling}.
W szczególności $\|u_\varepsilon\|_p=\varepsilon\|\varphi\|_p$,
więc normy $W^{1,p}$ są jednostajnie ograniczone dla
$0<\varepsilon\le1$. Dla $q=p^*$ drugi wykładnik
jest równy zeru; dla $q>p^*$ jest ujemny, a wówczas
$\|u_\varepsilon\|_q\to\infty$ przy $\varepsilon\downarrow0$.
Nie może więc istnieć
oszacowanie $\|u\|_q\le C\|u\|_{W^{1,p}}$ dla
wszystkich takich funkcji, nawet na dowolnym obszarze zawierającym kulę wokół $x_0$. Ten sam przykład pokaże
brak zwartości w wykładniku krytycznym.
\end{uwaga}

\begin{definicja}[Obszar Lipschitzowski]
\label{def:pde10-Lipschitz}
Ograniczony otwarty zbiór $\Omega\subset\R^d$ nazywamy
\emph{obszarem Lipschitzowskim}, jeżeli jest spójny, a jego
brzeg można pokryć skończoną liczbą otoczeń, w których
po odpowiednim sztywnym ruchu współrzędnych ma on postać
$\{(x',x_d):x_d=\phi(x')\}$ z $\phi$ Lipschitzowską,
zaś $\Omega$ leży lokalnie po jednej stronie tego wykresu.
Przyjmujemy, że otoczenia są mniejsze niż wykresowe
cylindry, aby dostępne były dodatnie odległości między
nośnikami funkcji odcinających a ich bocznymi ścianami.
\end{definicja}

\begin{lemat}[Przedłużanie pierwszego rzędu]
\label{lem:pde10-extension}
Niech $\Omega\subset\R^d$ będzie obszarem ograniczonym Lipschitzowskim i $1\le p\le\infty$. Istnieje liniowy operator
$E:W^{1,p}(\Omega)\to W^{1,p}(\R^d)$ oraz stała
$C=C(\Omega,p)>0$ takie, że $(Eu)|_\Omega=u$ i
$\|Eu\|_{W^{1,p}(\R^d)}\le C\|u\|_{W^{1,p}(\Omega)}$.
Można dodatkowo wybrać operator o obrazie w funkcjach
wspartych w jednej ustalonej kuli, niezależnej od $u$.
\end{lemat}
\begin{proof}
\emph{Krok 1. Prostowanie brzegu.}
Zmniejszając otoczenie wykresu, możemy przyjąć
$U=\{(x',x_d):x'\in D,\ |x_d-\phi(x')|<h\}$,
gdzie $D$ jest kulą, $h>0$, a $\Omega\cap U$
jest dokładnie częścią $x_d>\phi(x')$ tego otoczenia
(w razie potrzeby odwracamy ostatnią oś).
Niech $Q=D\times(-h,h)$, $Q_+=D\times(0,h)$ oraz
\[
 \Phi(x',t)=(x',t+\phi(x')), \qquad
 \Psi(x',x_d)=(x',x_d-\phi(x')).
\]
Są to wzajemnie odwrotne odwzorowania dwulipschitzowskie,
co wynika z nierówności trójkąta i własności Lipschitza
funkcji $\phi$. Mamy $\Phi(Q_+)=U\cap\Omega$.
Dla nieujemnej funkcji borelowskiej $a$ z Tonellego
i podstawienia w ostatniej współrzędnej otrzymujemy
\[
 \int_Q a\circ\Phi\,\mathrm d\lambda_d
 =\int_D\int_{\phi(x')-h}^{\phi(x')+h}a(x',s)
 \,\mathrm d\lambda_1(s)\,\mathrm d\lambda_{d-1}(x')
 =\int_U a\,\mathrm d\lambda_d.
\]
Analogicznie dla $Q_+$ i $U\cap\Omega$.
Wzór dla funkcji charakterystycznych borelowskich
zbiorów zerowych daje zachowanie zbiorów miary zero
w obu kierunkach; przez wybór reprezentantów borelowskich
rozszerzamy wzór na funkcje mierzalne w sensie Lebesgue'a.
Składanie zachowuje normy $L^p$, także dla $p=\infty$.

Dla gładkiej $v$ na $U\cap\Omega$ ograniczenia $\phi$
do prostych równoległych do osi są lipschitzowskie,
a więc absolutnie ciągłe na zwartych odcinkach.
Z twierdzenia~\ref{thm:pde9-ac} i Fubiniego pochodne
$\partial_i\phi$ istnieją prawie wszędzie na $D$
i są ograniczone co do modułu przez $L=\operatorname{Lip}\phi$.
Są mierzalne jako granice ilorazów różnicowych
(poza zbiorem istnienia przyjmujemy zero).
Na zwartych pododcinkach prostych w $Q_+$ złożenie
$v\circ\Phi$ jest absolutnie ciągłe, ponieważ krzywa
wewnętrzna jest lipschitzowska, a $v$ lokalnie
lipschitzowska na jej zwartym obrazie.
Jednowymiarowa reguła łańcuchowa daje prawie wszędzie
\[
 \partial_t(v\circ\Phi)=(\partial_d v)\circ\Phi,\qquad
 \partial_i(v\circ\Phi)=(\partial_i v)\circ\Phi+
    ((\partial_d v)\circ\Phi)\partial_i\phi\quad(i<d).
\]
Są to również słabe pochodne: całkujemy przez części
na odcinkach przeciw testowi o zwartym nośniku,
a następnie stosujemy Fubiniego. Test znika przy końcach
przekrojów; całki są bezwzględnie zbieżne dzięki
lokalnej ograniczoności czynników na nośniku testu.
Nie korzystamy z twierdzenia o przekrojach ogólnych
funkcji Sobolewa. Z zachowania miary wynika
\[
 \|\partial_t(v\circ\Phi)\|_p=\|\partial_d v\|_p,\qquad
 \|\partial_i(v\circ\Phi)\|_p
 \le\|\partial_i v\|_p+L\|\partial_d v\|_p.
\]
Normy po lewej liczymy na $Q_+$, po prawej na $U\cap\Omega$.
Dla $\Psi$ rachunek jest analogiczny ze znakiem minus
przy $\partial_i\phi$; działa też na pełnych otoczeniach.
Jeśli przekrój ma wiele składowych, całkujemy przez części
na każdej z nich, a całki sumujemy przez bezwzględną
całkowalność.

Dla $p<\infty$ i $v\in W^{1,p}$ wybieramy gładkie
$v_n$ zbieżne do $v$ w tej normie na mocy twierdzenia
~\ref{thm:pde9-MeyersSerrin}. Oszacowania dla różnic
$v_n-v_m$ dają ciąg Cauchy'ego złożeń w $W^{1,p}$.
Zupełność i zachowanie normy $L^p$ identyfikują granicę
jako złożenie z $v$; wzory pochodnych przechodzą
do granicy w $L^p$, gdyż $\partial_i\phi$ jest ograniczona.
Dotyczy to obu odwzorowań oraz pełnych i jednostronnych
otoczeń. Dla $p=\infty$ stosujemy przypadek $p=1$,
ponieważ otoczenia mają skończoną miarę. Prawe strony
wzorów należą do $L^\infty$, co daje te same oszacowania
bez aproksymacji w normie $W^{1,\infty}$ i bez korzystania
z nierówności Morreya.
Dla $d=1$ pomijamy zmienne styczne; zmiana współrzędnych
jest przesunięciem lub odbiciem.

\emph{Krok 2. Przedłużenie lokalne.}
Dla $\chi\in C_c^\infty(U)$ połóżmy
$v=(\chi u)|_{U\cap\Omega}$, $w=v\circ\Phi$ na $Q_+$.
Reguła iloczynu (twierdzenie~\ref{thm:pde9-products})
i krok 1 dają oszacowanie normy $w$ przez stałą razy
$\|u\|_{W^{1,p}(\Omega)}$.
Funkcja $w$ i jej słabe pochodne znikają prawie wszędzie
poza $A=\Psi(\operatorname{supp}\chi)\Subset Q$
(dla pochodnych wynika to z lokalności).
Przedłużamy $w$ przez zero na półprzestrzeń
$P=\R^{d-1}\times(0,\infty)$, oznaczając wynik przez $W$.
Wybierzmy $\zeta\in C_c^\infty(Q)$ równą $1$ w otoczeniu
$A$. Dla $\varphi\in C_c^\infty(P)$ funkcja
$(\zeta\varphi)|_{Q_+}$ jest testem na $Q_+$,
gdyż $\varphi$ znika blisko $t=0$, a $\zeta$ blisko
sztucznych ścian. Z zanikania $w$ i jej pochodnych
poza $A$ wynika
\[
 \begin{aligned}
 -\int_P W\partial_j\varphi\,\mathrm d\lambda_d
 &=-\int_{Q_+}w\partial_j(\zeta\varphi)\,\mathrm d\lambda_d\\
 &=\int_{Q_+}(\partial_jw)\zeta\varphi\,\mathrm d\lambda_d
 =\int_{Q_+}(\partial_jw)\varphi\,\mathrm d\lambda_d.
 \end{aligned}
\]
Pochodne $W$ są zatem przedłużeniami przez zero
pochodnych $w$, z zachowaniem normy.

Odbijamy $W$ parzyście względem $t=0$ na mocy
twierdzenia~\ref{thm:pde9-reflection}. Dla $p<\infty$
norma operatora wynosi $2^{1/p}$. Dla $p=\infty$
funkcja $W$ i jej pochodne są ograniczone i znikają
poza zbiorem ograniczonym, więc $W\in W^{1,1}(P)$.
Stosując twierdzenie dla $p=1$, dostajemy wzory
na odbite pochodne (parzyste styczne i nieparzystą
normalną), które należą do $L^\infty$; odbicie
zachowuje normę maksimum składowych $W^{1,\infty}$.
Odbita funkcja ma nośnik zawarty w
$A\cup\{(x',-t):(x',t)\in A\}\Subset Q$.
Ograniczamy ją do $Q$ i składamy z $\Psi$.
Otrzymujemy funkcję z $W^{1,p}(U)$ o zwartym nośniku
w $U$, którą przedłużamy przez zero za pomocą
lematu~\ref{lem:pde9-compact-product} (z odcięciem
równym $1$ w otoczeniu tego nośnika).
Otrzymany operator $E_\chi$ jest liniowy i ograniczony,
a jego restrykcja do $\Omega$ jest równa $\chi u$.

\emph{Krok 3. Sklejenie.}
Ze zwartości brzegu wybieramy skończenie wiele mniejszych
otoczeń brzegowych o domknięciach zawartych w odpowiednich
$U_j$, $j=1,\ldots,N$. Pozostała część $\overline\Omega$
jest zwartym podzbiorem $\Omega$; obejmujemy ją otwartym
$U_0\Subset\Omega$. Wybieramy nieujemne gładkie odcięcia
$b_j$ o zwartych nośnikach w $U_j$, równe $1$ na
odpowiednich mniejszych zbiorach (konstrukcja ze splotem
z dowodu twierdzenia~\ref{thm:pde9-MeyersSerrin}).
Suma $S=\sum_jb_j$ jest dodatnia na otoczeniu $V$
zbioru $\overline\Omega$. Wybieramy $\eta\in C_c^\infty(V)$
równą $1$ w otoczeniu $\overline\Omega$ i definiujemy
$\chi_j=\eta b_j/S$ na $V$, zero poza $V$.
Są globalnie gładkie, mają zwarte nośniki w $U_j$
i ich suma jest równa $1$ na $\Omega$.
Część $\chi_0u$ przedłużamy przez zero z lematu
~\ref{lem:pde9-compact-product}, pozostałe przez $E_{\chi_j}$,
uwzględniając sztywne ruchy zachowujące miarę
i mające stałe macierze pochodnych.
Suma $Eu$ jest liniowa, a na $\Omega$ równa $u$.
Nierówność trójkąta daje
\[
 \|Eu\|_{W^{1,p}(\R^d)}
 \le\left(\sum_{j=0}^NC_j\right)\|u\|_{W^{1,p}(\Omega)}.
\]
Nośniki składników leżą w ustalonych zbiorach zwartych,
których skończona suma mieści się w jednej kuli.
\end{proof}

\begin{wniosek}[Gęstość funkcji gładkich aż do brzegu]
\label{cor:pde10-density-boundary}
Niech $\Omega\subset\R^d$ będzie ograniczonym obszarem
Lipschitzowskim, $1\le p<\infty$. Wówczas zbiór
$\{\varphi|_\Omega:\varphi\in C_c^\infty(\R^d)\}$
jest gęsty w $W^{1,p}(\Omega)$. W szczególności jest
gęsty zbiór funkcji gładkich w otoczeniu $\overline\Omega$.
\end{wniosek}
\begin{proof}
Dla $u\in W^{1,p}(\Omega)$ przedłużenie
$v=Eu\in W^{1,p}(\R^d)$ istnieje na mocy
lematu~\ref{lem:pde10-extension}. Wybierzmy funkcje
$\chi_R$ równe $1$ w $B_R$, o nośniku w $B_{2R}$
i $\|\nabla\chi_R\|_\infty\le C/R$.
Wtedy $\chi_Rv\to v$ w $W^{1,p}(\R^d)$:
dla składników bez pochodnych korzystamy ze zbieżności
ogonów całek, a dla gradientów dodatkowy składnik
$v\nabla\chi_R$ ma normę najwyżej $C\|v\|_p/R$.
Dla ustalonego $R$ regularizujemy $\chi_Rv$ splotem
z $\rho_\varepsilon$. Ponieważ $p<\infty$,
$\rho_\varepsilon*(\chi_Rv)\to\chi_Rv$ w $W^{1,p}(\R^d)$.
Funkcje te należą do $C_c^\infty(\R^d)$.
Ograniczenie do $\Omega$ nie zwiększa normy, więc
odpowiedni ciąg restrykcji zbiega do $u$ w $W^{1,p}(\Omega)$.
\end{proof}

\begin{twierdzenie}[Zanurzenia na obszarze ograniczonym Lipschitzowskim]
\label{thm:pde10-domain-embed}
Niech $d\ge2$ i $\Omega\subset\R^d$ będzie obszarem ograniczonym Lipschitzowskim. Wówczas zachodzą zanurzenia ciągłe:
\begin{enumerate}[label=\textnormal{(\roman*)}]
\item gdy $1\le p<d$, $W^{1,p}(\Omega)\hookrightarrow
 L^q(\Omega)$ dla $1\le q\le p^*=dp/(d-p)$;
\item gdy $p=d$, $W^{1,d}(\Omega)\hookrightarrow L^q(\Omega)$
dla każdego skończonego $q\ge1$;
\item gdy $p>d$, $W^{1,p}(\Omega)\hookrightarrow
 C^{0,\alpha}(\overline\Omega)$ dla $\alpha=1-d/p$.
Dla $p=\infty$ przyjmujemy $\alpha=1$.
\end{enumerate}
W części (iii) przestrzeń $C^{0,\alpha}(\overline\Omega)$
oznacza klasy z reprezentantem ciągłym o skończonej normie
$\|u\|_\infty+[u]_{C^{0,\alpha}}$.
\end{twierdzenie}
\begin{proof}
Część (i) wynika z lematu~\ref{lem:pde10-extension},
twierdzenia~\ref{thm:pde10-Sobp} i ograniczoności miary
$\Omega$: dla $q\le p^*$ nierówność H\"oldera daje
$\|u\|_q\le|\Omega|^{1/q-1/p^*}\|u\|_{p^*}$.

Gdy $p=d$, ustalmy dowolne skończone $q>d$.
Możemy dobrać $p_0\in(1,d)$ tak, że
$p_0^*=dp_0/(d-p_0)\ge q$. Przedłużenie $Eu$ ma
zwarty nośnik zawarty w ustalonej kuli $B_R$.
Ponieważ $p_0<d$, H\"older na $B_R$ daje
$\|Eu\|_{W^{1,p_0}}\le C\|Eu\|_{W^{1,d}}$.
Twierdzenie~\ref{thm:pde10-Sobp} daje
$\|Eu\|_{p_0^*}\le C\|\nabla Eu\|_{p_0}$.
Po ograniczeniu do $\Omega$ i zastosowaniu Höldera
na tym zbiorze skończonej miary otrzymujemy
$\|u\|_q\le C_q\|u\|_{W^{1,d}(\Omega)}$.
Dla $q\le d$ wystarcza H\"older na $\Omega$.
Część (iii) wynika z nierówności Morreya,
której pełny dowód podajemy w następnej sekcji.
\end{proof}

\section{Nierówność Morreya: dowód z całkowania po odcinkach}
\label{sec:pde10-Morrey}

\begin{lemat}[Oszacowanie różnicy punktu i średniej]
\label{lem:pde10-point-average}
Niech $d\ge2$, $B=B_R(a)\subset\R^d$, a $f\in C^1(\overline B)$.
Dla każdego $x\in B$ zachodzi
\[
 |f(x)-f_B|\le C_d\int_B
 \frac{|\nabla f(z)|}{|x-z|^{d-1}}\,\mathrm d\lambda_d(z),
 \quad f_B=\frac{1}{|B|}\int_Bf\,\mathrm d\lambda_d.
\]
Stała $C_d$ zależy tylko od wymiaru.
\end{lemat}
\begin{proof}
Z wypukłości kuli wynika, że odcinek od $x$ do każdego
$y\in B$ leży w $B$. Stosujemy podstawowe twierdzenie
rachunku całkowego do funkcji $t\mapsto f(x+t(y-x))$.
Z reguły łańcuchowej i nierówności Cauchy'ego--Schwarza
otrzymujemy
$|f(x)-f(y)|\le|x-y|\int_0^1
|\nabla f(x+t(y-x))|\,\mathrm d\lambda_1(t)$.
Ponieważ
$|f(x)-f_B|\le |B|^{-1}\int_B|f(x)-f(y)|\,\mathrm d\lambda_d(y)$,
uśredniamy prawą stronę i zamieniamy całki nieujemne
na mocy Tonellego. Przy ustalonym $t\in(0,1]$ podstawiamy
$z=x+t(y-x)$, zatem $\mathrm d\lambda_d(y)=
t^{-d}\,\mathrm d\lambda_d(z)$ i
$|x-y|=|x-z|/t$. Obrazem $B$ jest kula
$x+t(B-x)\subset B$; każdy jej punkt spełnia
$|x-z|\le 2Rt$. Dla ustalonego $z\ne x$ możemy więc
powiększyć zbiór dopuszczalnych $t$ do przedziału
$[|x-z|/(2R),1]$. Ponownie stosując Tonellego,
otrzymujemy
\begin{align*}
 |f(x)-f_B|
 &\le\frac1{|B|}\int_B|\nabla f(z)|\,|x-z|
 \int_{|x-z|/(2R)}^1 t^{-d-1}\,\mathrm d\lambda_1(t)
 \,\mathrm d\lambda_d(z)\\
 &\le \frac{(2R)^d}{d|B|}
 \int_B\frac{|\nabla f(z)|}{|x-z|^{d-1}}
 \,\mathrm d\lambda_d(z).
\end{align*}
Druga nierówność wynika z obliczenia, dla $s=|x-z|>0$,
\[
 \int_{s/(2R)}^1 t^{-d-1}\,\mathrm d\lambda_1(t)
 =\frac{(2R/s)^d-1}{d}\le\frac{(2R)^d}{ds^d}.
\]
Ponieważ $|B|=\omega_dR^d$, gdzie
$\omega_d=\lambda_d(B_1(0))$, można przyjąć
$C_d=2^d/(d\omega_d)$. Punkt $z=x$ pomijamy jako zbiór
miary zero. Jądro $|x-z|^{1-d}$ jest lokalnie całkowalne:
całkowanie w biegunowych na kuli o promieniu $2R$
daje $d\omega_d\int_0^{2R}1\,\mathrm d\lambda_1(s)$.
W szczególności prawa strona jest skończona
dla rozważanej funkcji klasy $C^1$.
\end{proof}

\begin{twierdzenie}[Morrey na $\R^d$]
\label{thm:pde10-Morrey}
Niech $d\ge2$ i $d<p<\infty$. Każda klasa
$u\in W^{1,p}(\R^d)$ ma jednoznacznego ciągłego
reprezentanta $\widetilde u$ i istnieje stała $C(d,p)$ taka, że
\[
 |\widetilde u(x)-\widetilde u(y)|
 \le C(d,p)|x-y|^{1-d/p}\|\nabla u\|_{L^p(\R^d)}
 \quad(x,y\in\R^d).
\]
Dodatkowo $\|\widetilde u\|_\infty\le
 C(d,p)\|u\|_{W^{1,p}(\R^d)}$.
Dla $p=\infty$ każda klasa $u\in W^{1,\infty}(\R^d)$
ma jedyny ciągły reprezentant globalnie lipschitzowski
z $\operatorname{Lip}\widetilde u\le\|\nabla u\|_\infty$
oraz $\|\widetilde u\|_\infty=\|u\|_{L^\infty}$.
\end{twierdzenie}
\begin{proof}
Najpierw rozważmy $u\in C_c^\infty(\R^d)$. Ustalmy
$x\ne y$ i oznaczmy $r=|x-y|$.
Wybieramy kulę $B=B_{2r}(x)$, która zawiera oba punkty.
Z lematu~\ref{lem:pde10-point-average} otrzymujemy
$|u(x)-u(y)|\le|u(x)-u_B|+|u(y)-u_B|$.
Każdy składnik szacujemy nierównością Höldera
z wykładnikami $p$ oraz $p'=p/(p-1)$.
Dla $z_0=x$ lub $z_0=y$ liczymy
\[
 \left(\int_B |z-z_0|^{-(d-1)p'}
 \,\mathrm d\lambda_d(z)\right)^{1/p'}
 \le C(d,p)r^{1-d/p}.
\]
Istotnie, $B\subset B_{4r}(z_0)$, a w biegunowych
całka jest ograniczona przez
$C_d\int_0^{4r}s^{d-1-(d-1)p'}\,\mathrm d\lambda_1(s)$.
Mamy
\[
 d-(d-1)p'=\frac{p-d}{p-1}>0,\qquad
 \frac{d-(d-1)p'}{p'}=1-\frac dp.
\]
Całka potęgowa jest więc skończona i po wyciągnięciu
$p'$-tego pierwiastka daje dokładnie czynnik
$r^{1-d/p}$. Z lematu i Höldera otrzymujemy
\[
 |u(z_0)-u_B|\le C(d,p)r^{1-d/p}\|\nabla u\|_{L^p(B)}.
\]
Sumując dla $z_0=x,y$, dostajemy oszacowanie z tezy;
dla $x=y$ obie strony są zerowe.

Aby kontrolować wartość w jednym punkcie,
uśredniamy $u$ na kuli jednostkowej $B_1(x)$.
Stosujemy bezpośrednio lemat~\ref{lem:pde10-point-average}
na $B_1(x)$ i Höldera. Rachunek całki jądra dla promienia
$1$ daje $|u(x)-u_{B_1(x)}|\le C(d,p)\|\nabla u\|_p$,
a z Höldera $|u_{B_1(x)}|\le\omega_d^{-1/p}\|u\|_p$.
Dodając oszacowania i biorąc supremum po $x$, uzyskujemy
$\|u\|_\infty\le C(d,p)\|u\|_{W^{1,p}}$,
z uwzględnieniem równoważności norm gradientu.

Dla ogólnego $u\in W^{1,p}(\R^d)$ wybieramy
gładkie funkcje o zwartym nośniku $u_n\to u$ w $W^{1,p}$
na mocy gęstości z sekcji~\ref{sec:pde9-density}.
Oszacowanie supremum zastosowane do $u_n-u_m$ daje
ciąg Cauchy'ego w normie jednostajnej. Jego granica
$\widetilde u$ jest ograniczona i ciągła. Wybieramy
podciąg z $\|u_{n_j}-u\|_p^p<2^{-j}$. Z Tonellego
$\sum_j|u_{n_j}-u|^p<\infty$ prawie wszędzie, więc
$u_{n_j}\to u$ prawie wszędzie. Jednostajna granica
$\widetilde u$ reprezentuje zatem klasę $u$.
Dwa ciągłe reprezentanty równe prawie wszędzie są równe
wszędzie: niezerowa różnica w punkcie byłaby niezerowa
na otwartej kuli dodatniej miary.
W oszacowaniu przyrostów dla $u_n$ przechodzimy
do granicy dla każdego $x,y$, ponieważ zbieżność w
$W^{1,p}$ daje zbieżność norm gradientów. Analogicznie
przechodzimy do granicy w oszacowaniu supremum.
Jednoznaczność reprezentanta zapewnia liniowość
przyporządkowania $u\mapsto\widetilde u$.

Dla $p=\infty$ niech $u_\varepsilon=\rho_\varepsilon*u$
na $\R^d$, z nieujemnym jądrem o całce $1$ z rozdziału
~\ref{chap:pde-sobolev}. Z twierdzenia
~\ref{thm:pde9-mollify}, stosowanego lokalnie przy $p=1$,
$\nabla u_\varepsilon=\rho_\varepsilon*(\nabla u)$.
Nieujemność jądra i nierówność trójkąta dla normy
euklidesowej dają w każdym punkcie
\[
 |u_\varepsilon(x)|\le\|u\|_\infty,\qquad
 |\nabla u_\varepsilon(x)|\le\|\nabla u\|_\infty.
\]
Całkowanie po odcinku daje
$|u_\varepsilon(x)-u_\varepsilon(y)|
 \le\|\nabla u\|_\infty|x-y|$.
Z lokalnej zbieżności w $L^1$ wybieramy
$\varepsilon_n\downarrow0$ tak, aby
$\|u_{\varepsilon_n}-u\|_{L^1(B_n(0))}<2^{-n}$.
Dla ustalonego $m$ suma tych błędów na $B_m(0)$
od $n=m$ jest skończona; Tonelli daje zbieżność
punktową prawie wszędzie na $B_m(0)$.
Po przeliczalnej sumie zbiorów wyjątkowych otrzymujemy
zbiór $S\subset\R^d$ pełnej miary, na którym
$u_{\varepsilon_n}\to u$. Zbiór $S$ jest gęsty,
gdyż każda niepusta kula ma dodatnią miarę.
Przejście do granicy dla $x,y\in S$ daje na $S$
to samo oszacowanie lipschitzowskie i ograniczenie
modułu przez $\|u\|_\infty$. Dla każdego $x\in\R^d$
wybieramy $x_j\in S$, $x_j\to x$, i definiujemy
$\widetilde u(x)=\lim_j u(x_j)$.
Oszacowanie przyrostów daje istnienie granicy
i jej niezależność od wyboru ciągu, a po przejściu
do granicy również globalną stałą Lipschitza.
Na $S$ reprezentant zgadza się z $u$.
Jego supremum jest równe istotnemu supremum dzięki
ciągłości, a jednoznaczność wynika jak wyżej.
Nie korzystaliśmy ze zbieżności w $W^{1,\infty}$.
\end{proof}

\begin{wniosek}[Morrey na obszarze i wyższe rzędy]
\label{cor:pde10-Morrey-domain}
Jeśli $\Omega\subset\R^d$ jest obszarem ograniczonym Lipschitzowskim i $d<p\le\infty$, to
$W^{1,p}(\Omega)\hookrightarrow C^{0,1-d/p}(\overline\Omega)$,
z naturalną interpretacją wykładnika $1$ dla $p=\infty$.
Jeśli $k\ge2$ i $u\in W^{k,p}(\Omega)$, $d<p\le\infty$,
to pochodne słabe $\partial^\alpha u$ dla $|\alpha|\le k-1$
mają ciągłe reprezentanty H\"olderowskie wykładnika
$1-d/p$; stąd $u$ ma reprezentanta w $C^{k-1,1-d/p}$.
\end{wniosek}
\begin{proof}
Dla $k=1$, $d\ge2$ bierzemy reprezentanta ciągłego
$Eu$ z twierdzenia~\ref{thm:pde10-Morrey}, gdzie $E$
jest operatorem z lematu~\ref{lem:pde10-extension}.
Jego ograniczenie do $\overline\Omega$ reprezentuje $u$
na $\Omega$. Norma supremum i półnorma Höldera
na $\overline\Omega$ nie przekraczają odpowiednich norm
na $\R^d$. Otrzymujemy, przy $\alpha=1-d/p$,
\[
 \|\widetilde u\|_{C^{0,\alpha}(\overline\Omega)}
 \le C(d,p)\|Eu\|_{W^{1,p}(\R^d)}
 \le C(\Omega,p)\|u\|_{W^{1,p}(\Omega)}.
\]
Dla $p=\infty$ używamy $\alpha=1$ i oszacowania
Lipschitza wraz z normą supremum. Dla $d=1$ obszar
jest ograniczonym przedziałem; stosujemy twierdzenie
~\ref{thm:pde9-ac} i wniosek~\ref{cor:pde9-trace-oneD}
oraz oszacowanie przyrostów z tego twierdzenia i Höldera
(dla $p=\infty$ z normą supremum pochodnej).
Jednoznaczność wynika z równości prawie wszędzie
funkcji ciągłych na $\Omega$ i z gęstości $\Omega$
w $\overline\Omega$. Liniowość wynika z jednoznaczności.

Dla $k\ge2$ każdy $\partial^\beta u$, $|\beta|\le k-1$,
należy do $W^{1,p}(\Omega)$, ponieważ jego słabe
pochodne pierwszego rzędu są pochodnymi $u$ rzędu
$|\beta|+1\le k$. Niech $v_\beta$ będzie jego ciągłym
reprezentantem na $\overline\Omega$. Wszystkie te
funkcje spełniają oszacowania Höldera i supremum
przez stałą razy $\|u\|_{W^{k,p}(\Omega)}$.
Dla $|\beta|\le k-2$ i $x\in\Omega$ wybieramy kulę
$B_{2r}(x)\Subset\Omega$. Na $B_r(x)$ lokalne sploty
$u_\varepsilon$ oraz ich pochodne do rzędu $k-1$
są splotami odpowiednich $v_\gamma$ z jądrem, na mocy
twierdzenia~\ref{thm:pde9-mollify} (dla $p=\infty$
stosujemy lokalnie przypadek $p=1$).
Z jednostajnej ciągłości $v_\gamma$ na
$\overline{B_{2r}(x)}$ otrzymujemy dla $\varepsilon<r$
\[
 \sup_{z\in B_r(x)}
 |\partial^\gamma u_\varepsilon(z)-v_\gamma(z)|
 \le\sup_{\substack{z\in B_r(x)\\|h|\le\varepsilon}}
 |v_\gamma(z-h)-v_\gamma(z)|\longrightarrow0.
\]
Dla $|t|<r$ podstawowe twierdzenie rachunku całkowego
daje
\[
 \partial^\beta u_\varepsilon(x+te_i)
 -\partial^\beta u_\varepsilon(x)
 =\int_0^t\partial^{\beta+e_i}u_\varepsilon(x+se_i)
             \,\mathrm d\lambda_1(s).
\]
Przechodząc do granicy jednostajnej po obu stronach,
otrzymujemy tę samą tożsamość dla $v_\beta$ oraz
$v_{\beta+e_i}$. Dzieląc przez $t\ne0$ i korzystając
z ciągłości $v_{\beta+e_i}$ przy $t\to0$, mamy
klasyczną pochodną $\partial_i v_\beta=v_{\beta+e_i}$.
Wszystkie te pochodne są ciągłe, więc indukcyjnie
$v_0\in C^{k-1}(\Omega)$, a jego klasyczne pochodne
przedłużają się ciągle i hölderowsko na
$\overline\Omega$. To właśnie oznacza tu
$C^{k-1,\alpha}(\overline\Omega)$; nie wymaga
istnienia klasycznego przedłużenia wyższego rzędu
poza obszar.
\end{proof}

\begin{uwaga}[Wymiar $d=1$]
Na ograniczonym przedziale wyniki są silniejsze:
$W^{1,1}(a,b)$ ma reprezentanta absolutnie ciągłego,
a dla $1<p<\infty$ nierówność Höldera i twierdzenie
~\ref{thm:pde9-ac} dają
$|u(x)-u(y)|\le |x-y|^{1-1/p}\|u'\|_{L^p(a,b)}$.
Dla $p=\infty$ otrzymujemy oszacowanie lipschitzowskie
z wykładnikiem $1$. Dla $p=1$ nie ma wspólnego
dodatniego wykładnika Höldera (przykłady $u(x)=x^\theta$
na $(0,1)$, $0<\theta<1$, omówione w rozdziale
~\ref{chap:pde-sobolev}).
Nie wolno utożsamiać tego przypadku z uogólnieniem wzoru
$p^*=dp/(d-p)$ na $d=1$.
\end{uwaga}

\section{Zwartość: od średnich po kostkach do twierdzenia Rellicha}
\label{sec:pde10-Rellich}

\begin{lemat}[Aproksymacja funkcjami stałymi na kostkach]
\label{lem:pde10-cube-approx}
Niech $Q\subset\R^d$ będzie kostką o krawędzi $L$, a
$1\le p<\infty$. Dla każdego $N\in\N$ podzielmy $Q$
na $N^d$ równych kostek $Q_j$ o krawędzi $h=L/N$,
a dla $u\in L^1(Q)$ połóżmy
$P_Nu=\sum_j u_{Q_j}\chi_{Q_j}$, gdzie
$u_{Q_j}=|Q_j|^{-1}\int_{Q_j}u\,\mathrm d\lambda_d$.
Dla wszystkich $u\in W^{1,p}(Q)$ zachodzi
\[
 \|u-P_Nu\|_{L^p(Q)}\le dh\|\nabla u\|_{L^p(Q)}.
\]
\end{lemat}
\begin{proof}
Stosujemy nierówność Poincar\'ego--Wirtingera
z twierdzenia~\ref{thm:pde9-cube-Poincare} osobno na
każdej $Q_j$:
$\|u-u_{Q_j}\|_{L^p(Q_j)}\le
 h\sum_{i=1}^d\|\partial_i u\|_{L^p(Q_j)}$.
Ponieważ $(\sum_i a_i)^p\le d^{p-1}\sum_i a_i^p$
dla $a_i\ge0$, po podniesieniu do potęgi $p$
i zsumowaniu po rozłącznych kostkach otrzymujemy
$\|u-P_Nu\|_p^p\le d^{p-1}h^p\sum_i\|\partial_i u\|_p^p$.
Po wyciągnięciu pierwiastka mamy w istocie stałą
$d^{1-1/p}$; oszacowanie przez $d$ jest wystarczające.
\end{proof}

\begin{twierdzenie}[Rellich--Kondrachov: podstawowa zwartość]
\label{thm:pde10-Rellich-p}
Niech $\Omega\subset\R^d$ będzie obszarem ograniczonym Lipschitzowskim i $1\le p<\infty$. Wtedy zanurzenie
$W^{1,p}(\Omega)\hookrightarrow\hookrightarrow L^p(\Omega)$
jest zwarte.
\end{twierdzenie}
\begin{proof}
Niech $(u_n)$ będzie ograniczony w $W^{1,p}(\Omega)$.
Przedłużamy każdą funkcję operatorem
$E$ z lematu~\ref{lem:pde10-extension}. Wszystkie
$Eu_n$ mają nośniki zawarte w jednej ustalonej kostce $Q$
i są ograniczone w $W^{1,p}(Q)$ przez stałą $M$.

Ustalmy $\varepsilon>0$. Z lematu~\ref{lem:pde10-cube-approx}
wybieramy $N$ tak duże, aby
$\|Eu_n-P_NEu_n\|_{L^p(Q)}<\varepsilon/3$
jednocześnie dla wszystkich $n$.
Funkcje $P_NEu_n$ należą do skończeniewymiarowej przestrzeni
funkcji stałych na kostkach $Q_j$, a ich normy są
jednostajnie ograniczone przez $M$ (Jensen daje
$\|P_Nf\|_p\le\|f\|_p$). Zbiór ograniczony w przestrzeni
skończeniewymiarowej jest całkowicie ograniczony,
więc ma skończoną $\varepsilon/3$-sieć.
Z nierówności trójkąta wynika, że zbiór $\{Eu_n\}$
ma skończoną $\varepsilon$-sieć w $L^p(Q)$.

Dla pełnej ścisłości, wybierając kolejno sieci
dla $\varepsilon=2^{-m}$ i zagnieżdżone podciągi,
konstrukcja przekątniowa daje podciąg $(Eu_{n_k})$
Cauchy'ego w $L^p(Q)$; zupełność $L^p$
zapewnia jego zbieżność. Ograniczenie do $\Omega$
nie zwiększa normy, więc $u_{n_k}$ zbiega w
$L^p(\Omega)$. Nie wykorzystaliśmy
refleksywności $W^{1,p}$, dlatego dowód działa również
przy $p=1$.
\end{proof}

\begin{lemat}[Interpolacja elementarna norm $L^p$]
\label{lem:pde10-interpolation}
Niech $1\le a<q<b<\infty$, a $f\in L^a(\Omega)\cap L^b(\Omega)$.
Jeśli $0<\theta<1$ spełnia $1/q=\theta/a+(1-\theta)/b$,
to
$\|f\|_{L^q}\le\|f\|_{L^a}^\theta
\|f\|_{L^b}^{1-\theta}$.
\end{lemat}
\begin{proof}
Piszemy $|f|^q=(|f|^a)^{q\theta/a}
(|f|^b)^{q(1-\theta)/b}$.
Dwa wykładniki sumują się do $1$.
Oznaczmy $J_r=\int_\Omega|f(x)|^r\,\mathrm d\lambda_d(x)$
dla $r\in\{a,q,b\}$. Nierówność H\"oldera z wykładnikami
$a/(q\theta)$ i $b/(q(1-\theta))$ daje
\[
 J_q\le J_a^{q\theta/a}J_b^{q(1-\theta)/b},
\]
co kończy dowód.
\end{proof}

\begin{twierdzenie}[Rellich--Kondrachov: wykładniki podkrytyczne]
\label{thm:pde10-Rellich-q}
Niech $\Omega\subset\R^d$ będzie obszarem ograniczonym Lipschitzowskim, $d\ge2$.
\begin{enumerate}[label=\textnormal{(\roman*)}]
\item Jeśli $1\le p<d$, to
$W^{1,p}(\Omega)\hookrightarrow\hookrightarrow L^q(\Omega)$
dla $1\le q<p^*=dp/(d-p)$.
\item Jeśli $p=d$, to to samo zanurzenie jest zwarte dla
każdego skończonego $q\ge1$.
\item Jeśli $d<p\le\infty$, to
$W^{1,p}(\Omega)\hookrightarrow\hookrightarrow
C^{0,\beta}(\overline\Omega)$
dla każdego $0\le\beta<1-d/p$; przy $p=\infty$ wolno
przyjąć dowolne $\beta<1$.
\end{enumerate}
\end{twierdzenie}
\begin{proof}
W (i) niech $(u_n)$ będzie ograniczony w $W^{1,p}$.
Z twierdzenia~\ref{thm:pde10-Rellich-p} wybieramy
podciąg Cauchy'ego w $L^p$. Twierdzenie
\ref{thm:pde10-domain-embed} daje jednostajne
ograniczenie norm $L^{p^*}$, zatem różnice
$u_n-u_m$ mają jednostajnie ograniczone normy $L^{p^*}$.
Dla $p<q<p^*$ lemat~\ref{lem:pde10-interpolation}
implikuje, że podciąg jest Cauchy'ego w $L^q$.
Dla $q\le p$ wystarcza nierówność H\"oldera na
zbiorze o skończonej mierze. Tak uzyskujemy (i).

W (ii) jeśli $q>p=d$, wybieramy ustalone skończone
$r>q$. Zanurzenie $W^{1,d}\hookrightarrow L^r$ daje
ograniczenie norm $L^r$, a interpolacja między $L^d$
i $L^r$ daje zbieżność w $L^q$. Dla $q\le d$
postępujemy jak wcześniej.

W (iii) nierówność Morreya daje jednostajne ograniczenie
norm $C^{0,\alpha}(\overline\Omega)$, gdzie
$\alpha=1-d/p>0$ (albo $\alpha=1$ dla $p=\infty$).
Rodzina jest więc jednostajnie ograniczona i jednakowo
ciągła na zwartym $\overline\Omega$.
Twierdzenie Arzel\`a--Ascoliego zapewnia podciąg zbieżny
jednostajnie do ciągłej funkcji $u$.
Jeśli $f_n=u_n-u$, to $\|f_n\|_\infty\to0$
i $[f_n]_{C^{0,\alpha}}\le C$.
Dla $0<\beta<\alpha$ i każdego $\delta>0$ mamy
\[
 [f_n]_{C^{0,\beta}}
 \le\max\{2\|f_n\|_\infty\delta^{-\beta},
 C\delta^{\alpha-\beta}\}.
\]
Najpierw wybieramy $\delta$ małe, by drugi wyraz
był mały, potem $n$ duże, by pierwszy też był mały.
To dowodzi zbieżności w $C^{0,\beta}$.
Dla $\beta=0$ oznacza ona zbieżność jednostajną.
\end{proof}

\begin{kontrprzyklad}[Brak zwartości dla wykładnika krytycznego]
\label{ex:pde10-noncompact-critical}
Niech $1\le p<d$, $x_0\in\Omega$ oraz
$\varphi\in C_c^\infty(B_1(0))$ będzie niezerowa.
Funkcje $u_\varepsilon(x)=\varepsilon^{1-d/p}
\varphi((x-x_0)/\varepsilon)$, dla dostatecznie małych
$\varepsilon>0$, mają jednostajnie ograniczone normy
$W^{1,p}(\Omega)$ i stałą dodatnią normę $L^{p^*}(\Omega)$.
Jednocześnie $u_\varepsilon\to0$ prawie wszędzie, gdy
$\varepsilon\downarrow0$. Gdyby istniał podciąg zbieżny
silnie w $L^{p^*}$, jego granica musiałaby być równa zeru,
co przeczyłoby stałości norm. Dlatego zanurzenie przy
$q=p^*$ nie jest zwarte.
\end{kontrprzyklad}

\begin{kontrprzyklad}[Brak zwartości na całej $\R^d$]
\label{ex:pde10-translation}
Niech $\varphi\in C_c^\infty(B_1(0))$ będzie niezerowa
i $u_n(x)=\varphi(x-3ne_1)$.
Wówczas $\|u_n\|_{W^{1,p}(\R^d)}$ jest stała dla
każdego $1\le p<\infty$, lecz dla $m\ne n$ nośniki są
rozłączne, więc
$\|u_n-u_m\|_p^p=2\|\varphi\|_p^p$.
Żaden podciąg nie może być Cauchy'ego w $L^p$.
Ograniczoność obszaru w twierdzeniu Rellicha nie jest
zatem przypadkowym warunkiem technicznym.
\end{kontrprzyklad}

\section[Przypadek graniczny w zanurzeniach Sobolewa]{Przypadek graniczny $p=d$: dlaczego nie dostajemy $L^\infty$?}
\label{sec:pde10-Moser}

\begin{kontrprzyklad}[Ciąg koncentracyjny typu Mosera]
\label{ex:pde10-Moser}
Niech $d\ge2$, $B=B_1(0)$ oraz $0<\varepsilon<e^{-1}$.
Połóżmy $L=\log(1/\varepsilon)$ i
\[
 u_\varepsilon(x)=
 \begin{cases}
 L^{1-1/d}, & |x|\le\varepsilon,\\
 L^{-1/d}\log(1/|x|), & \varepsilon<|x|<1,\\
 0,&|x|\ge1.
 \end{cases}
\]
Funkcja jest ciągła i należy do $W^{1,d}(\R^d)$.
W istocie dla $\varepsilon<|x|<1$ mamy
$|\nabla u_\varepsilon(x)|=L^{-1/d}/|x|$, skąd
\[
 \|\nabla u_\varepsilon\|_{L^d(\R^d)}^d
 =|S^{d-1}|L^{-1}\int_\varepsilon^1\frac{\mathrm dr}{r}
 =|S^{d-1}|.
\]
Natomiast $\|u_\varepsilon\|_{L^d}$ jest jednostajnie
ograniczona: na $\varepsilon<|x|<1$ wystarcza
$L^{-1}\int_0^\infty t^de^{-dt}\,\mathrm dt<\infty$,
a na kuli $B_\varepsilon$ wkład wynosi
$\omega_d\varepsilon^dL^{d-1}$ i pozostaje ograniczony.
Pomimo tego
$\|u_\varepsilon\|_\infty=L^{1-1/d}\to\infty$.
Nie istnieje więc ciągłe zanurzenie
$W^{1,d}(B)\hookrightarrow L^\infty(B)$.
\end{kontrprzyklad}

\begin{uwaga}[Ostrość wykładnika Morreya]
Niech $p>d$ i $\alpha=1-d/p$. Dla tej samej skalowanej
rodziny $u_\varepsilon(x)=\varepsilon^\alpha
\varphi((x-x_0)/\varepsilon)$ normy $W^{1,p}$ są
ograniczone, lecz dla $\beta>\alpha$ ilorazy H\"oldera
rosną jak $\varepsilon^{\alpha-\beta}$.
Nie można więc ogólnie zastąpić wykładnika Morreya
większym wykładnikiem, nawet na kuli.
\end{uwaga}

\section{Ślady na brzegu Lipschitzowskim}
\label{sec:pde10-traces}

\begin{definicja}[Miara i przestrzeń $L^p$ na brzegu]
Jeśli $\Omega\subset\R^d$ jest obszarem ograniczonym Lipschitzowskim, przyjmujemy na $\partial\Omega$ miarę
$\mathcal H^{d-1}|_{\partial\Omega}$, czyli ograniczenie
$(d-1)$-wymiarowej miary Hausdorffa. Na każdym lokalnym
wykresie $x_d=\phi(x')$ miara ta ma element
$\sqrt{1+|\nabla\phi(x')|^2}\,\mathrm d\lambda_{d-1}(x')$
prawie wszędzie. Przestrzeń $L^p(\partial\Omega)$
rozumiemy względem tej miary.
\end{definicja}

\begin{twierdzenie}[Twierdzenie o śladzie]
\label{thm:pde10-trace}
Niech $\Omega\subset\R^d$ będzie obszarem ograniczonym Lipschitzowskim, $1\le p<\infty$. Istnieje dokładnie jeden
ograniczony operator liniowy
$\operatorname{Tr}:W^{1,p}(\Omega)\to L^p(\partial\Omega)$,
który dla każdej $u\in C^\infty(\overline\Omega)$
(gładkiej w otoczeniu $\overline\Omega$) daje zwykłą
restrykcję $u|_{\partial\Omega}$. Ponadto
\[
 \|\operatorname{Tr}u\|_{L^p(\partial\Omega)}
 \le C(\Omega,p)\|u\|_{W^{1,p}(\Omega)}.
\]
\end{twierdzenie}
\begin{proof}
Najpierw udowodnimy oszacowanie na płaskim brzegu.
Dla $w\in C^1(\R^{d-1}\times[0,a])$ o odpowiedniej
całkowalności, $t\in(0,a)$, podstawowe twierdzenie
rachunku daje
$|w(x',0)|\le|w(x',t)|+
\int_0^a|\partial_tw(x',s)|\,\mathrm d\lambda_1(s)$.
Całkujemy nierówność po $t\in(0,a)$, dzielimy przez $a$
i stosujemy H\"oldera. Otrzymujemy punktowo
\[
 |w(x',0)|\le a^{-1/p}\|w(x',\cdot)\|_{L^p(0,a)}
 +a^{1-1/p}\|\partial_tw(x',\cdot)\|_{L^p(0,a)}.
\]
Po podniesieniu do potęgi $p$ i scałkowaniu po $x'$
dostajemy
\[
 \|w(\cdot,0)\|_{L^p(\R^{d-1})}^p
 \le 2^{p-1}\bigl(a^{-1}\|w\|_{L^p}^p
 +a^{p-1}\|\partial_tw\|_{L^p}^p\bigr).
\]
W lokalnym wykresie brzegu prostujemy go odwzorowaniem
$\Phi(x',t)=(x',t+\phi(x'))$. Funkcję $\chi u$,
gdzie $\chi$ jest odcięciem wewnątrz mniejszego cylindra,
przenosimy do płaskiej warstwy. Zmiana zmiennych
ma jakobian $1$, a $\partial_t[(\chi u)\circ\Phi]$
zgadza się z $\partial_{x_d}(\chi u)\circ\Phi$.
Na wykresie element powierzchniowy jest ograniczony
z góry przez $\sqrt{1+\operatorname{Lip}(\phi)^2}$
razy $\mathrm d\lambda_{d-1}$.
Otrzymujemy więc lokalne oszacowanie śladu przez
$\|\chi u\|_{W^{1,p}(\Omega)}$.
Skończone pokrycie brzegu i podział jedności pozwalają
zsumować te nierówności; stąd globalne oszacowanie
na gładkich $u$.

Aby zdefiniować ślad ogólnej funkcji, wybieramy
ciąg restrykcji funkcji $C_c^\infty(\R^d)$ zbieżny
w $W^{1,p}(\Omega)$, którego istnienie wykazano
we wniosku~\ref{cor:pde10-density-boundary}.
Otrzymane oszacowanie pokazuje, że restrykcje przybliżeń
tworzą ciąg Cauchy'ego w $L^p(\partial\Omega)$;
jego granica nie zależy od wyboru przybliżeń.
Przedłużenie operatora jest liniowe i ograniczone,
a gęstość funkcji gładkich aż do brzegu zapewnia
jednoznaczność.
\end{proof}

\begin{definicja}[Przestrzeń danych brzegowych $H^{1/2}(\partial\Omega)$]
\label{def:pde10-halftrace}
Niech $\Omega\subset\R^d$ będzie ograniczonym obszarem Lipschitzowskim,
a $\operatorname{Tr}:H^1(\Omega)\to L^2(\partial\Omega)$
operatorem z twierdzenia~\ref{thm:pde10-trace}. Definiujemy
\[
 H^{1/2}(\partial\Omega):=\operatorname{Tr}(H^1(\Omega)),\qquad
 \|g\|_{H^{1/2}(\partial\Omega)}
 :=\inf_{\substack{G\in H^1(\Omega)\\\operatorname{Tr}G=g}}
 \|G\|_{H^1(\Omega)}.
\]
Jest to definicja przez przestrzeń śladów, z normą ilorazową.
Na ograniczonym brzegu Lipschitzowskim jest ona równoważna
zwyczajowej definicji przez ułamkowe seminormy Gagliarda;
dowód tej identyfikacji wymaga oddzielnego twierdzenia o śladach.
\end{definicja}

\begin{lemat}[Liniowy operator podnoszenia śladu]
\label{lem:pde10-halftrace-rightinverse}
Dla każdego ograniczonego obszaru Lipschitzowskiego $\Omega$
istnieje liniowy, ograniczony operator
$E_{\operatorname{Tr}}:H^{1/2}(\partial\Omega)\to H^1(\Omega)$
taki, że $\operatorname{Tr}E_{\operatorname{Tr}}g=g$ oraz
$\|E_{\operatorname{Tr}}g\|_{H^1(\Omega)}
=\|g\|_{H^{1/2}(\partial\Omega)}$.
\end{lemat}
\begin{proof}
Niech $H=H^1(\Omega)$ i $N=\ker\operatorname{Tr}$.
Ponieważ $\operatorname{Tr}:H\to L^2(\partial\Omega)$
jest ciągły, $N$ jest domkniętą podprzestrzenią Hilberta.
Mamy więc rozkład ortogonalny $H=N\oplus N^\perp$.
Dla $g\in\operatorname{Tr}(H)$ wybieramy dowolny $G\in H$
o śladzie $g$ i definiujemy
$E_{\operatorname{Tr}}g=P_{N^\perp}G$.
Jeżeli $G_1$ i $G_2$ mają ten sam ślad, to
$G_1-G_2\in N$, a zatem ich rzuty na $N^\perp$ są równe.
Konstrukcja jest więc niezależna od wyboru $G$ i liniowa.
Ponadto $G-P_{N^\perp}G\in N$, skąd
$\operatorname{Tr}E_{\operatorname{Tr}}g=g$.
Każde inne podniesienie jest postaci
$E_{\operatorname{Tr}}g+n$ dla $n\in N$, więc
z twierdzenia Pitagorasa
$\|E_{\operatorname{Tr}}g+n\|_H^2
=\|E_{\operatorname{Tr}}g\|_H^2+\|n\|_H^2$.
To dowodzi równości norm i ograniczoności operatora.
W szczególności norma ilorazowa czyni
$H^{1/2}(\partial\Omega)$ przestrzenią Hilberta.
\end{proof}

\begin{twierdzenie}[Zerowy ślad i domknięcie funkcji testowych]
\label{thm:pde10-trace-zero}
Niech $\Omega\subset\R^d$ będzie obszarem ograniczonym Lipschitzowskim, $1\le p<\infty$. Wówczas
\[
 W_0^{1,p}(\Omega)=\ker\operatorname{Tr}.
\]
\end{twierdzenie}
\begin{proof}
Jeśli $u\in W_0^{1,p}(\Omega)$, to z definicji jest
granicą w $W^{1,p}$ funkcji $C_c^\infty(\Omega)$.
Te funkcje mają zerowy ślad, więc przez ciągłość
$\operatorname{Tr}u=0$.

Dla implikacji odwrotnej niech $\operatorname{Tr}u=0$.
Rozważmy lokalny fragment brzegu wyprostowany do
półwarstwy $t>0$. Z twierdzenia o absolutnej ciągłości
na prawie wszystkich pionowych odcinkach dla funkcji
$W^{1,p}$ wynika, że dla prawie wszystkich $x'$
wartość brzegowa $w(x',0)$ istnieje i pokrywa się
ze śladem (dowodzi się tego przez gęstość w normie
$W^{1,p}$ i jednowymiarowe oszacowanie z poprzedniego
dowodu). Jest zatem równa zero. Przedłużenie $w$
przez zero na $t<0$ ma słabą pochodną normalną
równą przedłużeniu $\partial_tw$ przez zero:
podczas całkowania przez części w $t$ składnik brzegowy
$w(x',0)$ znika. Pochodne styczne również
przedłużają się przez zero, gdyż granica całkowania
w tych współrzędnych jest niezależna od $x'$.

Wybierzmy funkcję odcinającą $\chi$ o nośniku w
mniejszym cylindrze oraz dodatni kierunek pionowy
skierowany do wnętrza obszaru. Przesuwamy
przedłużenie $\chi w$ o odległość $a\varepsilon$
do wnętrza, a następnie wygładzamy jądrem o promieniu
$\varepsilon$; wybieramy stałą $a>1$ tak, aby wygładzenie
miało nośnik po stronie $t>0$. Translacja jest ciągła
w $W^{1,p}(\R^d)$ dla $p<\infty$, a splot
z przybliżoną jedynką zbiega w tej samej normie.
Wszystkie uzyskane funkcje są zatem gładkie, wspierają
się we wnętrzu odpowiedniego fragmentu obszaru
i zbiegają do $\chi w$ w $W^{1,p}$.

Po powrocie do pierwotnych współrzędnych na wykresie
Lipschitzowskim gładkość nie musi być zachowana.
Dlatego konstrukcję wykonujemy geometrycznie przed
prostowaniem: przesuwamy fragment $\chi u$ w stałym
kierunku pionowym oryginalnego cylindra na odległość
$a\varepsilon$, gdzie $a>1+\operatorname{Lip}\phi$,
a następnie wygładzamy w oryginalnych współrzędnych
promieniem $\varepsilon$. Warunek na $a$ gwarantuje,
że całe kule używane w splocie pozostają po stronie
wewnętrznej grafu. Zbieżność w $W^{1,p}$ wynika
z własności translacji w $\R^d$ i ograniczoności
operatorów mnożenia przez odcięcia.
Części podziału jedności wspierające się ściśle
wewnątrz $\Omega$ wygładzamy bez przesunięć.
Sumując skończenie wiele aproksymacji, otrzymujemy
ciąg w $C_c^\infty(\Omega)$ zbieżny do $u$
w $W^{1,p}(\Omega)$. To kończy dowód.
\end{proof}

\begin{wniosek}[Nierówność Poincar\'ego na obszarze]
\label{cor:pde10-Poincare}
Niech $\Omega\subset\R^d$ będzie obszarem ograniczonym Lipschitzowskim, $1\le p<\infty$.
Istnieje stała $C=C(\Omega,p)$ taka, że
\[
 \|u-u_\Omega\|_{L^p(\Omega)}
 \le C\|\nabla u\|_{L^p(\Omega)},\qquad
 u_\Omega=|\Omega|^{-1}\int_\Omega u\,\mathrm d\lambda_d.
\]
Dla $u\in W_0^{1,p}(\Omega)$ istnieje również stała
$C'$ z $\|u\|_p\le C'\|\nabla u\|_p$.
\end{wniosek}
\begin{proof}
Gdyby pierwsza nierówność była fałszywa, istniałyby
$u_n\in W^{1,p}(\Omega)$ takie, że
$(u_n)_\Omega=0$, $\|u_n\|_p=1$ i
$\|\nabla u_n\|_p\to0$ (odejmujemy średnią i
normalizujemy przez normę $L^p$).
Ciąg jest ograniczony w $W^{1,p}$, więc twierdzenie
\ref{thm:pde10-Rellich-p} daje podciąg zbieżny w
$L^p$ do pewnego $u$ o normie $1$ i średniej $0$.
Z zamkniętości pochodnej dystrybucyjnej
(wniosek~\ref{cor:pde9-closed}) otrzymujemy
$\nabla u=0$. Na obszarze funkcja lokalnie
całkowalna o zerowych pochodnych dystrybucyjnych
jest stała prawie wszędzie; jej średnia jest zerowa,
więc $u=0$, sprzeczność.
Dla drugiej nierówności przedłużamy funkcję z
$W_0^{1,p}(\Omega)$ przez zero na kostkę $Q\supset\Omega$
i stosujemy wniosek~\ref{cor:pde9-Poincare-zero}.
\end{proof}

% Audyt rozdziału 10: jawny krok techniczny używany w dalszej części tomu.
\begin{wniosek}[Przedłużenie przez zero funkcji o zerowym śladzie]
\label{cor:pde10-zeroextension}
Niech $\Omega\subset\R^d$ będzie obszarem ograniczonym Lipschitzowskim, $1\le p<\infty$ i
$u\in W^{1,p}(\Omega)$ z $\operatorname{Tr}u=0$.
Niech $E_0u$ oznacza funkcję równą $u$ prawie wszędzie
na $\Omega$ i $0$ poza $\Omega$. Wówczas
$E_0u\in W^{1,p}(\R^d)$ oraz
\[ \partial_j(E_0u)=E_0(\partial_j u)
\quad\text{w }\mathcal D'(\R^d),\qquad j=1,\ldots,d. \]
W konsekwencji normy $L^p$ funkcji i jej gradientu
są zachowane przez przedłużenie.
\end{wniosek}
\begin{proof}
Twierdzenie~\ref{thm:pde10-trace-zero} pozwala wybrać
$u_n\in C_c^\infty(\Omega)$ zbieżne do $u$ w $W^{1,p}(\Omega)$.
Przedłużenia przez zero $E_0u_n$ należą do
$C_c^\infty(\R^d)$ i spełniają
$\partial_j(E_0u_n)=E_0(\partial_j u_n)$ punktowo.
Z definicji przedłużenia
$\|E_0u_n-E_0u_m\|_{L^p(\R^d)}=
\|u_n-u_m\|_{L^p(\Omega)}$;
taka sama równość zachodzi dla gradientów.
Ciąg $(E_0u_n)$ jest zatem Cauchy'ego w $W^{1,p}(\R^d)$
i zbiega do pewnej funkcji $w$. Zbieżność w $L^p$ pokazuje,
że $w=E_0u$ prawie wszędzie. Zamkniętość pochodnej słabej
daje $\partial_jw=E_0(\partial_j u)$.
Równości norm wynikają z równości funkcji prawie wszędzie.
\end{proof}

\begin{uwaga}[Zależności twierdzeń i brzegi nieregularne]
Powyższy dowód świadomie korzysta z ustalonej wcześniej
równości $W_0^{1,p}(\Omega)=\ker\operatorname{Tr}$;
nie może więc zastąpić jej dowodu bez błędnego koła.
Na obszarze o nieregularnym brzegu nie należy
identyfikować obu przestrzeni bez dodatkowych hipotez.
\end{uwaga}

\section{Schemat związków i ograniczenia twierdzeń}
\label{sec:pde10-summary}
\begin{center}
\begin{tikzpicture}[node distance=8mm and 10mm,>=Stealth,
 every node/.style={font=\small,align=center}]
 \node[draw,rounded corners,fill=blue!5,text width=3.3cm] (a)
  {$W^{1,p}(\Omega)$\\$\Omega$ Lipschitzowska};
 \node[draw,rounded corners,fill=green!8,text width=3cm,
 below left=of a] (b) {$p<d$\\$L^{p^*}$ ciągle\\$L^q$, $q<p^*$ zwarto};
 \node[draw,rounded corners,fill=orange!10,text width=3cm,
 below=of a] (c) {$p=d$\\$L^q$ dla $q<\infty$\\zwarto dla $q<\infty$};
 \node[draw,rounded corners,fill=cyan!8,text width=3cm,
 below right=of a] (d) {$p>d$\\$C^{0,1-d/p}$ ciągle\\$C^{0,\beta}$ zwarto};
 \draw[->] (a)--(b);
 \draw[->] (a)--(c);
 \draw[->] (a)--(d);
\end{tikzpicture}

{\small\emph{Rysunek:} Zachowanie funkcji Sobolewa zależy od relacji
wykładnika całkowalności $p$ do wymiaru $d$. Zwarty charakter
zanurzeń zanika przy niektórych wykładnikach granicznych.}
\end{center}

Kluczowa różnica między zanurzeniem ciągłym a zwartym
polega na tym, że oszacowanie normy nie zapewnia istnienia
silnie zbieżnego podciągu. Przykład skupiania masy
pokazuje utratę zwartości przy $q=p^*$, a translacje
pokazują jej utratę na nieograniczonej przestrzeni.
Te dwa mechanizmy będą ważne w metodach wariacyjnych.


% --- Rozdział 11. Metody wariacyjne ---

\chapter[Metody wariacyjne i twierdzenie Laxa--Milgrama]{Metody wariacyjne: słabe rozwiązania, twierdzenie Laxa--Milgrama i minimalizacja energii}
\label{chap:pde-variational}

Równania eliptyczne rzadko dają się rozwiązać za pomocą jawnych wzorów.
Wcześniej nauczyliśmy się jednak kontrolować słabe pochodne,
operator śladu i normy przestrzeni Sobolewa. Teraz pokażemy, że
te narzędzia pozwalają udowadniać istnienie oraz jednoznaczność
rozwiązań bez konstruowania ich punkt po punkcie.

Kluczowe będą dwa podejścia. Pierwsze polega na rozwiązywaniu
równania liniowego w przestrzeni Hilberta; jego podstawą jest
twierdzenie Laxa--Milgrama. Drugie polega na minimalizacji energii,
czasem również nieliniowej. Będziemy wyraźnie odróżniać
\emph{formę dwuliniową} od \emph{funkcjonału energii}: jeśli forma nie
jest symetryczna, to zwykły funkcjonał kwadratowy nie daje
na ogół wyjściowego równania.

W całym rozdziale przestrzenie Hilberta oraz formy dwuliniowe są
rzeczywiste. Całkujemy względem miary Lebesgue'a
$\lambda_d$ na $\Omega\subset\R^d$, zaś na brzegu Lipschitzowskim
względem $\mathcal H^{d-1}$. Zapisy $\int_\Omega f\,\mathrm d\lambda_d(x)$
i $\int_{\partial\Omega}g\,\mathrm d\mathcal H^{d-1}(x)$
określają za każdym razem także miarę całkowania.

\section{Przestrzeń energetyczna i sformułowanie słabe}
\label{sec:pde11-weak}

Niech $\Omega\subset\R^d$ będzie obszarem ograniczonym Lipschitzowskim. Kładziemy $V=H_0^1(\Omega)=W_0^{1,2}(\Omega)$.
Ustalmy stałą $C_P>0$ taką, że
$\|v\|_2\le C_P\|\nabla v\|_2$ dla każdego $v\in V$
(wniosek~\ref{cor:pde10-Poincare}). Z tej nierówności wynika, że
\[ (u,v)_V=\int_\Omega\nabla u\cdot\nabla v\,\mathrm d\lambda_d(x) \]
wyznacza iloczyn skalarny: jest symetryczna i dwuliniowa,
$(u,u)_V=\|\nabla u\|_2^2\ge0$, a gdy $(u,u)_V=0$,
nierówność Poincarégo daje $\|u\|_2=0$, czyli $u=0$ w $V$.
Norma $\|u\|_V=\sqrt{(u,u)_V}$ jest równoważna normie
$H^1(\Omega)$. W istocie $\|u\|_{H^1}^2=\|u\|_2^2+
\|\nabla u\|_2^2\le(1+C_P^2)\|\nabla u\|_2^2$,
a $\|u\|_V\le\|u\|_{H^1}$ wynika z definicji normy.
Przestrzeń $H_0^1$ jest domknięciem liniowej przestrzeni
$C_c^\infty(\Omega)$ w $H^1$, a więc domkniętą podprzestrzenią
przestrzeni Hilberta $H^1$ (twierdzenie~\ref{thm:pde9-banach}).
Ciąg Cauchy'ego w normie $\|\cdot\|_V$ jest dzięki pierwszemu
oszacowaniu ciągiem Cauchy'ego w $H^1$, ma tam granicę w $V$
i zbiega do niej także w normie $\|\cdot\|_V$.
Zatem $(V,(\cdot,\cdot)_V)$ jest przestrzenią Hilberta.

\begin{definicja}[Przestrzeń dualna i słabe rozwiązanie]
\label{def:pde11-dual}
Niech $V$ będzie rzeczywistą przestrzenią Hilberta.
Jej \emph{przestrzenią dualną} $V^*$ nazywamy przestrzeń wszystkich
ciągłych funkcjonałów liniowych $F:V\to\R$, wyposażoną w normę
$\|F\|_{V^*}=\sup_{\|v\|_V\le1}|F(v)|$.
Ciągłość liniowego funkcjonału jest równoważna skończoności
tego supremum: z ciągłości w zerze istnieje $\delta>0$,
dla którego $\|v\|_V\le\delta$ implikuje $|F(v)|\le1$,
a po przeskalowaniu otrzymujemy $\|F\|_{V^*}\le1/\delta$.
W drugą stronę dla $v\ne0$ stosujemy definicję do
$v/\|v\|_V$, uzyskując $|F(v)|\le\|F\|_{V^*}\|v\|_V$;
dla $v=0$ nierówność jest równością.
Jednorodność i nierówność trójkąta wynikają z własności
modułu i supremum, a norma zerowa oznacza $F=0$.
Jeżeli $a:V\times V\to\R$ jest formą dwuliniową i $F\in V^*$,
to \emph{słabym rozwiązaniem zagadnienia} jest element $u\in V$
spełniający $a(u,v)=F(v)$ dla każdego $v\in V$.
Dla $V=H_0^1(\Omega)$ oznaczamy też $V^*=H^{-1}(\Omega)$;
w tej notacji $H^{-1}(\Omega)$ oznacza przestrzeń dualną do
$H_0^1(\Omega)$, a nie przestrzeń wszystkich dystrybucji
z ujemną regularnością na całej $\R^d$.
W tym rozdziale normę $H^{-1}$ liczymy względem normy
energetycznej $\|\cdot\|_V$. Jeśli $\|F\|_{(H^1)^*_{0}}$
oznacza normę funkcjonału na $H_0^1$ liczoną względem
pełnej normy $H^1$, to z zawierania odpowiednich kul
jednostkowych i ich skalowania mamy
\[
 \|F\|_{(H^1)^*_{0}}\le\|F\|_{V^*}
 \le\sqrt{1+C_P^2}\,\|F\|_{(H^1)^*_{0}}.
\]
Obie normy określają te same ciągłe funkcjonały.
Każdy $F\in H^{-1}(\Omega)$ określa dystrybucję przez
ograniczenie do $C_c^\infty(\Omega)$: dla testu o nośniku
w ustalonym zwartym $K\subset\Omega$ zachodzi
$|F(\varphi)|\le\|F\|_{V^*}\sqrt{d\lambda_d(K)}
\max_i\|\partial_i\varphi\|_\infty$.
Jest to oszacowanie ciągłości w topologii testów.
Ograniczenie jest iniektywne z gęstości testów w $V$.
\end{definicja}

\begin{przyklad}[Od równania Poissona do formy słabej]
\label{ex:pde11-poisson-weak}
Załóżmy chwilowo, że $u\in C^2(\overline\Omega)$,
$u|_{\partial\Omega}=0$, $f\in L^2(\Omega)$ oraz
$-\Delta u=f$ prawie wszędzie. Przez $C^2(\overline\Omega)$
rozumiemy tu funkcje klasy $C^2$ w otoczeniu domknięcia.
Funkcja $u$ i jej pochodne są ograniczone na zwartym
$\overline\Omega$, stąd $u\in H^1(\Omega)$.
Jej ślad jest równy restrykcji klasycznej: wybieramy gładkie
odcięcie równe jeden na otoczeniu $\overline\Omega$,
o zwartym nośniku w otoczeniu, gdzie $u$ jest klasy $C^2$.
Iloczyn przedłużamy przez zero na $\R^d$ i wygładzamy
splotem; przybliżenia zbiegają wraz z pierwszymi pochodnymi
jednostajnie na $\overline\Omega$ dzięki jednostajnej ciągłości.
Z ciągłości śladu w $H^1$ i skończonej miary brzegu
(twierdzenie~\ref{thm:pde10-trace}) ślad $u$ jest zatem zerowy.
Twierdzenie~\ref{thm:pde10-trace-zero} daje $u\in V$.
Dla $v\in C_c^\infty(\Omega)$ całkujemy przez części
w każdej współrzędnej (bez składników brzegowych z uwagi
na zwarty nośnik testu) i sumujemy. Otrzymujemy
$\int_\Omega\nabla u\cdot\nabla v\,\mathrm d\lambda_d(x)
=\int_\Omega fv\,\mathrm d\lambda_d(x)$.
Lewa strona jest ciągła względem $v$ w normie $H_0^1$,
ponieważ ogranicza ją $\|\nabla u\|_2\|\nabla v\|_2$.
Prawa strona również jest ciągła:
$|\int_\Omega fv\,\mathrm d\lambda_d|
\le C_P\|f\|_2\|\nabla v\|_2$.
Dla dowolnego $v\in V$ wybieramy $v_n\in C_c^\infty(\Omega)$
zbieżne w $H^1$, a więc także w normie $V$. Różnice
całek dla $v_n$ i $v$ są ograniczone odpowiednio przez
$\|\nabla u\|_2\|v_n-v\|_V$ i $C_P\|f\|_2\|v_n-v\|_V$.
Obie dążą do zera, co daje tożsamość dla wszystkich $v\in V$.
\end{przyklad}

W odwrotną stronę słaba tożsamość daje równość
$-\Delta u=f$ w sensie dystrybucji. Istotnie, dla
$v\in C_c^\infty(\Omega)$ z definicji pochodnej słabej
i pochodnej dystrybucji mamy
\[
 \langle-\Delta u,v\rangle
 =-\sum_{i=1}^d\langle\partial_i^2u,v\rangle
 =\sum_{i=1}^d\int_\Omega\partial_i u\,\partial_i v
       \,\mathrm d\lambda_d(x)
 =\int_\Omega fv\,\mathrm d\lambda_d(x).
\]
Dla ogólnego $F\in V^*$ ostatnią całkę zastępujemy
przez $F(v)$, czyli działanie dystrybucji zdefiniowanej
przez ograniczenie $F$ do testów.
Warunek $u|_{\partial\Omega}=0$ jest zakodowany w wyborze $V$;
na obszarze Lipschitzowskim jest równoważny
$\operatorname{Tr}u=0$ (twierdzenie~\ref{thm:pde10-trace-zero}).
Nie wynika stąd jeszcze, że $u\in H^2$ lub $u\in C^2$.
\emph{Istnienie słabe i regularność klasyczna to dwa osobne twierdzenia.}

\section{Dwa lematy o przestrzeniach Hilberta}
\label{sec:pde11-hilbert}

Najpierw udowodnimy fakty, od których zależy twierdzenie Laxa--Milgrama.
Nie zakładamy z góry, że każdy funkcjonał można zapisać
w postaci iloczynu skalarnego.

\begin{lemat}[Rzut na domknięty zbiór wypukły]
\label{lem:pde11-projection}
Niech $H$ będzie rzeczywistą przestrzenią Hilberta,
$K\subset H$ zbiorem niepustym, domkniętym i wypukłym,
a $x\in H$. Istnieje dokładnie jeden $y\in K$ taki, że
$\|x-y\|=\inf_{z\in K}\|x-z\|$. Ponadto
\[ (x-y,z-y)_H\le0\qquad\text{dla każdego }z\in K. \]
\end{lemat}
\begin{proof}
Oznaczmy $r=\inf_{z\in K}\|x-z\|$. Ponieważ $K$ jest
niepusty, $0\le r<\infty$. Z definicji infimum wybieramy
$y_n\in K$ z $r\le\|x-y_n\|<r+1/n$. Ponieważ $K$ jest wypukły,
$(y_n+y_m)/2\in K$. Tożsamość równoległoboku daje
\begin{align*}
\|y_n-y_m\|^2
&=2\|x-y_n\|^2+2\|x-y_m\|^2
 -4\left\|x-\frac{y_n+y_m}{2}\right\|^2\\
&\le 2\|x-y_n\|^2+2\|x-y_m\|^2-4r^2\longrightarrow0.
\end{align*}
Zatem $(y_n)$ jest ciągiem Cauchy'ego. Zupełność i domkniętość
zapewniają istnienie granicy $y\in K$; ciągłość normy
daje $\|x-y\|=r$.
Gdyby $y'\in K$ było drugim minimizatorem, ten sam rachunek
z $y_n=y$, $y_m=y'$ dałby $\|y-y'\|=0$.

Dla dowolnego $z\in K$ i $0\le t\le1$ mamy
$y+t(z-y)\in K$. Z minimalności $y$ wynika
$\|x-y-t(z-y)\|^2\ge\|x-y\|^2$.
Rozwijając kwadrat, dzieląc przez $t>0$ i przechodząc
do granicy $t\downarrow0$, otrzymujemy
$-2(x-y,z-y)_H\ge0$, czyli tezę.
Odwrotnie, jeżeli $y\in K$ spełnia tę nierówność
rzutową dla wszystkich $z\in K$, to
\[
 \|x-z\|^2=\|x-y\|^2-2(x-y,z-y)_H+\|z-y\|^2
 \ge\|x-y\|^2.
\]
Zatem nierówność charakteryzuje jedyny punkt minimalizujący.
Gdy $K$ jest podprzestrzenią liniową, podstawienie
$z=y+k$ i $z=y-k$ dla $k\in K$ daje $(x-y,k)_H=0$.
\end{proof}

\begin{twierdzenie}[Twierdzenie Riesza o reprezentacji]
\label{thm:pde11-riesz}
Niech $H$ będzie rzeczywistą przestrzenią Hilberta.
Dla każdego $F\in H^*$ istnieje dokładnie jeden $w_F\in H$ taki, że
$F(v)=(w_F,v)_H$ dla wszystkich $v\in H$.
Ponadto $\|w_F\|_H=\|F\|_{H^*}$.
\end{twierdzenie}
\begin{proof}
Dla $F=0$ wektor $w_F=0$ spełnia żądaną tożsamość
i równość norm; jeśli inny wektor ją spełnia, podstawienie
tego wektora za $v$ daje jego zerową normę.
Jeżeli $F\ne0$, to $K=\ker F$ jest podprzestrzenią
liniową, domkniętą jako przeciwobraz $\{0\}$ przez
ciągłe $F$ i właściwą, ponieważ $F$ nie jest zerowy. Wybierzmy $x\notin K$.
Z lematu~\ref{lem:pde11-projection} otrzymujemy rzut
$y\in K$ punktu $x$ i kładziemy $e=x-y\ne0$.
Z charakterystyki rzutu, podstawiając $z=y\pm k$,
$k\in K$, dostajemy $(e,k)_H=0$.
Dla dowolnego $v\in H$ element
$k=v-[F(v)/F(e)]e$ należy do $K$, ponieważ
$F(e)=F(x)\ne0$. Stąd
$0=(e,k)_H=(e,v)_H-[F(v)/F(e)]\|e\|_H^2$.
Otrzymujemy $F(v)=(w_F,v)_H$ dla
$w_F=F(e)e/\|e\|_H^2$.
Jeżeli $w_1,w_2$ reprezentują $F$, to
$(w_1-w_2,v)_H=0$ dla każdego $v\in H$; podstawiając
$v=w_1-w_2$, dostajemy $\|w_1-w_2\|^2=0$.
Nierówność Cauchy'ego--Schwarza daje
$|F(v)|\le\|w_F\|_H$ dla $\|v\|_H\le1$, więc
$\|F\|_{H^*}\le\|w_F\|_H$; biorąc
$v=w_F/\|w_F\|$ w przypadku $w_F\ne0$,
uzyskujemy $|F(v)|=\|w_F\|_H$, czyli nierówność odwrotną.
Jednoznaczność zapewnia, że przyporządkowanie
$F\mapsto w_F$ jest liniowe: kombinacja liniowa
reprezentantów reprezentuje tę samą kombinację funkcjonałów.
Jest ono izometrią na $H$, ponieważ każdy $w\in H$
reprezentuje ciągły funkcjonał $v\mapsto(w,v)_H$.
\end{proof}

\section{Twierdzenie Laxa--Milgrama z pełnym dowodem}
\label{sec:pde11-lax}

\begin{definicja}[Ograniczoność i koercywność]
\label{def:pde11-coercivity}
Niech $H$ będzie rzeczywistą przestrzenią Hilberta.
Formę dwuliniową $a:H\times H\to\R$ nazywamy
\emph{ograniczoną}, jeżeli istnieje $M\ge0$ takie, że
$|a(u,v)|\le M\|u\|_H\|v\|_H$ dla wszystkich $u,v\in H$.
Nazywamy ją \emph{koercywną}, jeżeli istnieje $\alpha>0$,
dla którego $a(v,v)\ge\alpha\|v\|_H^2$ dla każdego $v\in H$.
Koercywność nie jest tym samym co symetria: forma
może być koercywna i jednocześnie niesymetryczna.
Dla przykładu na $\R^2$ forma
$a(u,v)=u\cdot v+u_1v_2-u_2v_1$ spełnia
$a(v,v)=|v|^2$, ale $a(e_1,e_2)=1$ i $a(e_2,e_1)=-1$.
Jest ograniczona, gdyż
$|u_1v_2-u_2v_1|\le |u|\,|v|$ z nierówności
Cauchy'ego--Schwarza dla $(u_1,u_2)$ i $(v_2,-v_1)$.
Można więc przyjąć $M=2$ i $\alpha=1$.
\end{definicja}

\begin{twierdzenie}[Twierdzenie Laxa--Milgrama]
\label{thm:pde11-LM}
Niech $H$ będzie rzeczywistą przestrzenią Hilberta,
$a:H\times H\to\R$ ograniczoną formą dwuliniową
ze stałą $M>0$, koercywną ze stałą $\alpha>0$,
a $F\in H^*$. Wtedy istnieje dokładnie jeden element $u\in H$
spełniający
\[ a(u,v)=F(v)\qquad\text{dla każdego }v\in H. \]
Ponadto $\|u\|_H\le\alpha^{-1}\|F\|_{H^*}$.
Nie wymaga się symetrii formy $a$.
\end{twierdzenie}
\begin{proof}
Jeśli $H=\{0\}$, jedynym rozwiązaniem jest $u=0$,
a jedyny funkcjonał to $F=0$; oszacowanie zachodzi.
Dalej zakładamy $H\ne\{0\}$. Dla dowolnego $v\ne0$
koercywność i ograniczoność dają
$\alpha\|v\|_H^2\le a(v,v)\le M\|v\|_H^2$,
a zatem $0<\alpha\le M$.

\textbf{Krok 1: operator reprezentujący formę.}
Dla ustalonego $u\in H$ odwzorowanie $v\mapsto a(u,v)$
jest funkcjonałem liniowym i ograniczonym o normie
co najwyżej $M\|u\|_H$. Z twierdzenia Riesza
(\ref{thm:pde11-riesz}) istnieje dokładnie jeden
$Tu\in H$ spełniający
$(Tu,v)_H=a(u,v)$ dla każdego $v\in H$.
Dla $u,w\in H$ i $s,r\in\R$ mamy
$(T(su+rw),v)_H=(sTu+rTw,v)_H$ dla każdego $v$,
z liniowości $a$ w pierwszej zmiennej.
Podstawiając za $v$ różnicę tych wektorów, otrzymujemy
$T(su+rw)=sTu+rTw$, czyli liniowość $T$.
Równość norm w twierdzeniu Riesza daje
$\|Tu\|_H\le M\|u\|_H$, a zatem $T$ jest ograniczony.
Koercywność oznacza $(Tu,u)_H\ge\alpha\|u\|_H^2$.
Z kolei $F(v)=(f,v)_H$ dla jednoznacznie wyznaczonego
$f\in H$ z $\|f\|_H=\|F\|_{H^*}$.
Równanie słabe jest równoważne $(Tu-f,v)_H=0$
dla każdego $v\in H$, czyli $Tu=f$ (bierzemy $v=Tu-f$).

\textbf{Krok 2: konstrukcja kontrakcji.}
Dla parametru $t>0$ zdefiniujmy
$S_t(u)=u-t(Tu-f)$. Dla $u,w\in H$, oznaczając $z=u-w$,
obliczamy dokładnie
\begin{align*}
\|S_t(u)-S_t(w)\|_H^2
&=\|z-tTz\|_H^2\\
&=\|z\|_H^2-2t(Tz,z)_H+t^2\|Tz\|_H^2\\
&\le (1-2t\alpha+t^2M^2)\|z\|_H^2.
\end{align*}
Wybierzmy $0<t<2\alpha/M^2$.
Z nierówności $\alpha\le M$ wynika
\[
 1-2t\alpha+t^2M^2
 =M^2(t-\alpha/M^2)^2+1-\alpha^2/M^2\ge0.
\]
Warunek $0<t<2\alpha/M^2$ daje zarazem
$t(tM^2-2\alpha)<0$, więc
$q=\sqrt{1-2t\alpha+t^2M^2}\in[0,1)$.
Po wyciągnięciu pierwiastka z oszacowania dostajemy
$\|S_t(u)-S_t(w)\|_H\le q\|u-w\|_H$,
czyli $S_t$ jest kontrakcją.
Najmniejszą wartość tego oszacowania stałej kontrakcji
otrzymujemy dla $t=\alpha/M^2$, mianowicie
$q=\sqrt{1-\alpha^2/M^2}$; nie musi to być rzeczywista
najmniejsza stała Lipschitza odwzorowania $S_t$.

\textbf{Krok 3: punkt stały bez pomijania granicy.}
Ustalmy dowolne $u_0\in H$ i zdefiniujmy
$u_{n+1}=S_t(u_n)$.
Jeśli $q=0$, odwzorowanie $S_t$ jest stałe,
zatem $u_2=S_t(u_1)=S_t(u_0)=u_1$ i $u_1$ jest
punktem stałym. Niech teraz $0<q<1$.
Z kontrakcji, przez indukcję względem $n$, mamy
$\|u_{n+1}-u_n\|_H\le q^n\|u_1-u_0\|_H$:
dla $n=0$ jest równość, a krok indukcyjny wynika
z zastosowania $S_t$ do $u_{n+1}$ i $u_n$.
Dla $m>n$ sumujemy te oszacowania:
\[
 \|u_m-u_n\|_H\le
 \sum_{j=n}^{m-1}\|u_{j+1}-u_j\|_H
 \le\frac{q^n}{1-q}\|u_1-u_0\|_H.
\]
Prawa strona dąży do zera niezależnie od $m>n$.
Zupełność $H$ zapewnia istnienie $u=\lim_nu_n$.
Ponieważ $S_t$ jest ciągłe, $S_t(u_n)\to S_t(u)$,
a jednocześnie $S_t(u_n)=u_{n+1}\to u$.
Stąd $u=S_t(u)$, a ponieważ $t>0$, mamy $Tu=f$.
Przejście z $m$ do nieskończoności daje dodatkowo
oszacowanie błędu
$\|u-u_n\|_H\le q^n(1-q)^{-1}\|u_1-u_0\|_H$.
Jeżeli $Tu=f=Tw$, to z koercywności
$\alpha\|u-w\|_H^2\le(T(u-w),u-w)_H=0$;
stąd $u=w$.

\textbf{Krok 4: oszacowanie rozwiązania.}
Podstawiając $v=u$ w tożsamości słabej, otrzymujemy
$\alpha\|u\|_H^2\le a(u,u)=F(u)
\le\|F\|_{H^*}\|u\|_H$.
Dla $u\ne0$ dzielimy przez dodatnią liczbę $\alpha\|u\|_H$;
dla $u=0$ lewa strona żądanego oszacowania jest zerowa.
Ponadto $T$ jest bijekcją: dla dowolnego $f\in H$
stosujemy udowodnione twierdzenie do funkcjonału
$v\mapsto(f,v)_H$, a jednoznaczność daje iniektywność.
Z równości norm w twierdzeniu Riesza
$\|T^{-1}f\|_H\le\alpha^{-1}\|f\|_H$.
Odwzorowanie odwrotne jest liniowe jako odwrotność
bijekcji liniowej i jest ograniczone.
\end{proof}

\begin{wniosek}[Stabilność względem prawej strony]
\label{cor:pde11-stability}
Przy założeniach twierdzenia~\ref{thm:pde11-LM},
niech $u_i\in H$ będą rozwiązaniami odpowiadającymi
$F_i\in H^*$, $i=1,2$. Wówczas
$\|u_1-u_2\|_H\le\alpha^{-1}\|F_1-F_2\|_{H^*}$.
\end{wniosek}
\begin{proof}
Odejmujemy słabe równania: $a(u_1-u_2,v)=(F_1-F_2)(v)$.
Dla $w=u_1-u_2$ bierzemy $v=w$, otrzymując
$\alpha\|w\|_H^2\le a(w,w)=(F_1-F_2)(w)
\le\|F_1-F_2\|_{H^*}\|w\|_H$.
Jeśli $w\ne0$, dzielimy przez $\alpha\|w\|_H$;
w przeciwnym przypadku oszacowanie zachodzi.
Jest to ilościowa stabilność względem danych:
wraz z istnieniem i jednoznacznością oznacza poprawność
zagadnienia w sensie Hadamarda dla ustalonej formy $a$.
\end{proof}

\begin{uwaga}[Metoda iteracyjna a dowód istnienia]
Punkt stały w dowodzie dostarcza jednocześnie algorytmu:
$u_{n+1}=u_n-t(Tu_n-f)$, o zbieżności geometrycznej
przy podanym wyborze $t$. Sam operator $T$ zależy jednak
od wybranego iloczynu skalarnego w $H$, toteż bez
skończeniewymiarowej reprezentacji nie jest to automatycznie
praktyczna metoda numeryczna.
\end{uwaga}

\section{Równanie eliptyczne z mierzalnymi współczynnikami}
\label{sec:pde11-elliptic}

\begin{twierdzenie}[Zagadnienie Dirichleta w postaci dywergencyjnej]
\label{thm:pde11-elliptic}
Niech $\Omega\subset\R^d$ będzie obszarem ograniczonym Lipschitzowskim. Załóżmy, że $A:\Omega\to\R^{d\times d}$
ma mierzalne, istotnie ograniczone współczynniki,
a $c\in L^\infty(\Omega)$ jest nieujemna prawie wszędzie.
Dla macierzy używamy normy operatorowej względem
normy euklidesowej, a dla wektorów normy euklidesowej.
Ich normy $L^\infty$ oznaczają istotne suprema
tych norm punktowych.
Załóżmy ponadto, że dla pewnego $\nu>0$ oraz
prawie wszystkich $x\in\Omega$ zachodzi
$\xi^TA(x)\xi\ge\nu|\xi|^2$ dla każdego $\xi\in\R^d$.
Wówczas dla każdego $F\in H^{-1}(\Omega)$ istnieje
jednoznaczne $u\in H_0^1(\Omega)$ spełniające
\[
\int_\Omega\big(A\nabla u\cdot\nabla v+cuv\big)
\,\mathrm d\lambda_d(x)=F(v)
\quad\text{dla wszystkich }v\in H_0^1(\Omega).
\]
Przy normie $\|v\|_V=\|\nabla v\|_2$ zachodzi
$\|u\|_V\le\nu^{-1}\|F\|_{V^*}$.
Nie zakłada się symetrii macierzy $A(x)$.
\end{twierdzenie}
\begin{proof}
Na $V=H_0^1(\Omega)$ zdefiniujmy
$a(u,v)=\int_\Omega(A\nabla u\cdot\nabla v+cuv)
\,\mathrm d\lambda_d(x)$.
Oznaczmy $K=\operatorname*{ess\,sup}_{x\in\Omega}
\|A(x)\|_{\mathrm{op}}$ oraz $C=\|c\|_\infty$.
Norma macierzy jest funkcją ciągłą jej współczynników,
więc $x\mapsto\|A(x)\|_{\mathrm{op}}$ jest mierzalna.
Ponadto $\|A(x)\|_{\mathrm{op}}\le
(\sum_{i,j}|A_{ij}(x)|^2)^{1/2}$, z nierówności
Cauchy'ego--Schwarza dla każdego wiersza i sumowania
kwadratów. Zatem $K<\infty$.
Poza jednym zbiorem miary zero mamy
$|A(x)\xi|\le K|\xi|$ dla wszystkich $\xi\in\R^d$.
Składniki podcałkowe są mierzalne, niezależne od
reprezentantów i całkowalne, ponieważ ich moduły
ograniczają odpowiednio $K|\nabla u||\nabla v|$
oraz $C|u||v|$, należące do $L^1$ z Höldera.
Liniowość całki zapewnia dwuliniowość formy.
Z Cauchy'ego--Schwarza i nierówności Poincar\'ego
otrzymujemy
\begin{align*}
|a(u,v)|
&\le K\|\nabla u\|_2\|\nabla v\|_2
+C\|u\|_2\|v\|_2\\
&\le(K+CC_P^2)\|u\|_V\|v\|_V.
\end{align*}
Tak więc $a$ jest formą ograniczoną.
Ponieważ $c\ge0$, eliptyczność zapewnia
$a(v,v)\ge\nu\int_\Omega|\nabla v|^2
\,\mathrm d\lambda_d=\nu\|v\|_V^2$.
Wszystkie założenia twierdzenia~\ref{thm:pde11-LM} są spełnione,
więc otrzymujemy istnienie, jednoznaczność i oszacowanie.
\end{proof}

\begin{uwaga}[Znaczenie sformułowania dystrybucyjnego]
Dla $F_f(v)=\int_\Omega fv\,\mathrm d\lambda_d(x)$,
$f\in L^2(\Omega)$, zachodzi
$\|F_f\|_{V^*}\le C_P\|f\|_2$.
W tym przypadku rozwiązanie spełnia w $\mathcal D'(\Omega)$
$-\operatorname{div}(A\nabla u)+cu=f$.
Dokładniej, dla $\varphi\in C_c^\infty(\Omega)$
definicja pochodnej dystrybucyjnej daje
\[
 \langle-\operatorname{div}(A\nabla u)+cu,\varphi\rangle
 =\int_\Omega(A\nabla u\cdot\nabla\varphi+cu\varphi)
       \,\mathrm d\lambda_d(x)
 =F_f(\varphi).
\]
Dla ogólnego $F\in V^*$ ostatnia wartość to $F(\varphi)$,
działanie dystrybucji z definicji~\ref{def:pde11-dual}.
Odwrotnie, jeśli $u\in V$ spełnia tę równość
dystrybucyjnie, to gęstość testów w $V$ oraz
ograniczoność formy i funkcjonału rozszerzają ją na $V$.
Nie ma tu obowiązku definiowania klasycznej dywergencji
$A\nabla u$: przy mierzalnych współczynnikach produkt
$A\nabla u$ może nie mieć słabych pochodnych w $L^2$,
ale ma sens jako wektor z $L^2$, którego dywergencja
jest dystrybucją.
\end{uwaga}

\begin{przyklad}[Jawne rozwiązanie w jednym wymiarze]
\label{ex:pde11-1d}
Dla $\Omega=(0,1)$, $A=1$, $c=0$ i $f=1$
szukamy $u\in H_0^1(0,1)$ spełniającego
$\int_0^1u'(x)v'(x)\,\mathrm d\lambda_1(x)
=\int_0^1v(x)\,\mathrm d\lambda_1(x)$.
Dla $u(x)=x(1-x)/2$ mamy $u'(x)=1/2-x$,
$u''(x)=-1$ i $u(0)=u(1)=0$.
Funkcja jest w $H^1(0,1)$, a z jednowymiarowego
oszacowania śladu i twierdzenia~\ref{thm:pde10-trace-zero}
należy do $H_0^1(0,1)$.
Dla testów $v$ całkowanie przez części daje
$\int_0^1 u'v'\,\mathrm d\lambda_1
=-\int_0^1 u''v\,\mathrm d\lambda_1
=\int_0^1 v\,\mathrm d\lambda_1$.
Przejście przez gęstość jak w przykładzie
~\ref{ex:pde11-poisson-weak} daje tożsamość dla wszystkich $v\in V$. Twierdzenie~\ref{thm:pde11-elliptic}
wykazuje, że nie istnieje żadne inne rozwiązanie w $H_0^1$.
\end{przyklad}

\begin{twierdzenie}[Koercywność przy wyrazach niższego rzędu]
\label{thm:pde11-lower-order}
Niech $\Omega\subset\R^d$ będzie ograniczonym obszarem
Lipschitzowskim, a $V=H_0^1(\Omega)$ z normą
$\|u\|_V=\|\nabla u\|_{L^2(\Omega)}$. Niech
$A\in L^\infty(\Omega;\R^{d\times d})$ spełnia
$\xi^TA(x)\xi\ge\nu|\xi|^2$ dla pewnego $\nu>0$,
wszystkich $\xi\in\R^d$ i prawie wszystkich $x\in\Omega$.
Niech $b\in L^\infty(\Omega;\R^d)$ i
$c\in L^\infty(\Omega;\R)$; oznaczmy
$c_-(x)=\max\{-c(x),0\}$ i przez $C_P$ dowolną
stałą w nierówności $\|u\|_2\le C_P\|\nabla u\|_2$
dla $u\in V$. Jeśli
\[
\alpha:=\nu-C_P\|b\|_\infty-C_P^2\|c_-\|_\infty>0,
\]
to forma
\[
 a(u,v)=\int_\Omega\bigl(A\nabla u\cdot\nabla v
 +(b\cdot\nabla u)v+cuv\bigr)\,\mathrm d\lambda_d(x)
\]
jest ograniczona i koercywna na $V$ ze stałą co
najmniej $\alpha$. Dla każdego $F\in V^*$ równanie
$a(u,v)=F(v)$, $v\in V$, ma dokładnie jedno rozwiązanie
i zachodzi $\|u\|_V\le\alpha^{-1}\|F\|_{V^*}$.
\end{twierdzenie}
\begin{proof}
Przyjmujemy tę samą konwencję norm macierzy i wektorów
co w twierdzeniu~\ref{thm:pde11-elliptic}.
Eliptyczność rozumiemy poza wspólnym zbiorem miary zero;
jeśli nierówność przyjęto osobno dla każdego $\xi$,
wystarczy najpierw przeciąć zbiory pełnej miary
dla $\xi\in\mathbb Q^d$, a potem użyć ciągłości
funkcji kwadratowej względem $\xi$.
Z nierówności Cauchy'ego--Schwarza i Höldera mamy
\[
 \begin{aligned}
 |a(u,v)|&\le\|A\|_\infty\|\nabla u\|_2\|\nabla v\|_2\\
 &\quad+\|b\|_\infty\|\nabla u\|_2\|v\|_2
       +\|c\|_\infty\|u\|_2\|v\|_2\\
 &\le(\|A\|_\infty+C_P\|b\|_\infty
       +C_P^2\|c\|_\infty)\|u\|_V\|v\|_V,
 \end{aligned}
\]
gdzie ostatni krok to nierówność Poincarégo.
Oszacowanie uzasadnia również całkowalność każdego
składnika i ograniczoność formy.
Z kolei dla $u\in V$ dodatnia część $c$ nie może
obniżać wartości $a(u,u)$, a pozostałe wyrazy dają
\begin{align*}
 a(u,u)&\ge\nu\|\nabla u\|_2^2
 -\|b\|_\infty\|\nabla u\|_2\|u\|_2
 -\|c_-\|_\infty\|u\|_2^2\\
 &\ge\alpha\|u\|_V^2.
\end{align*}
Twierdzenie Laxa--Milgrama~\ref{thm:pde11-LM}
kończy dowód. Warunek ma charakter wystarczający;
ujemna część potencjału lub dryft mogą być większe,
jeżeli znamy dodatkową strukturę formy.
\end{proof}

\begin{przyklad}[Dryft bez ograniczenia małości]
\label{ex:pde11-divfree-drift}
Niech $b\in W^{1,\infty}(\Omega;\R^d)$ spełnia
$\operatorname{div}b=0$ prawie wszędzie oraz
$A=I$, $c=0$. Dla każdego $u\in C_c^\infty(\Omega)$
całkowanie przez części daje
$\int_\Omega(b\cdot\nabla u)u\,\mathrm d\lambda_d
=\frac12\int_\Omega b\cdot\nabla(u^2)\,\mathrm d\lambda_d=0$.
Istotnie, $u^2$ jest testem i z definicji
dywergencji dystrybucyjnej
$\int_\Omega b\cdot\nabla(u^2)\,\mathrm d\lambda_d
=-\langle\operatorname{div}b,u^2\rangle=0$.
Dla dowolnego $u\in H_0^1(\Omega)$ wybieramy
$u_n\in C_c^\infty(\Omega)$ zbieżne do $u$ w $H^1$.
Z nierówności Höldera
\[
 \begin{aligned}
 &\left|\int_\Omega
 [(b\cdot\nabla u_n)u_n-(b\cdot\nabla u)u]
               \,\mathrm d\lambda_d\right|\\
 &\qquad\le\|b\|_\infty
 (\|\nabla(u_n-u)\|_2\|u_n\|_2
   +\|\nabla u\|_2\|u_n-u\|_2)\longrightarrow0.
 \end{aligned}
\]
Ciąg norm $\|u_n\|_2$ jest ograniczony dzięki zbieżności.
Otrzymujemy tożsamość dla wszystkich $u\in V$. Dlatego
$a(u,u)=\|\nabla u\|_2^2$ i forma jest koercywna
ze stałą $1$ bez założenia, że $\|b\|_\infty$ jest mała.
Sam dryft nie jest więc przeszkodą: znaczenie ma również
jego struktura i warunki brzegowe.
\end{przyklad}

\section{Niezerowe dane brzegowe: metoda podnoszenia}
\label{sec:pde11-lifting}

\begin{twierdzenie}[Warunek Dirichleta z danymi brzegowymi]
\label{thm:pde11-lifting}
Niech $\Omega,A,c,\nu$ spełniają założenia
twierdzenia~\ref{thm:pde11-elliptic}. Niech
$f\in H^{-1}(\Omega)$ oraz $g\in H^1(\Omega)$.
Istnieje dokładnie jedna funkcja $u\in H^1(\Omega)$
taka, że $u-g\in H_0^1(\Omega)$ oraz
\[
\int_\Omega(A\nabla u\cdot\nabla v+cuv)
\,\mathrm d\lambda_d(x)=f(v)
\quad(v\in H_0^1(\Omega)).
\]
Jej ślad jest równy $\operatorname{Tr}g$.
Istnieje stała $C$, zależna wyłącznie od
$\Omega,\nu,\|A\|_\infty,\|c\|_\infty$, taka, że
$\|u\|_{H^1}\le C(\|f\|_{H^{-1}}+\|g\|_{H^1})$.
\end{twierdzenie}
\begin{proof}
Podstawiamy $u=w+g$. Żądany warunek brzegowy jest równoważny
$w\in V=H_0^1(\Omega)$, a równanie przyjmuje postać
$a(w,v)=G(v)$, gdzie $G(v)=f(v)-a(g,v)$.
Ograniczoność współczynników, Cauchy--Schwarz i
Poincar\'e dla $v\in V$ zapewniają
$|a(g,v)|\le C_0\|g\|_{H^1}\|v\|_V$.
Wobec tego $G\in V^*$ i
$\|G\|_{V^*}\le\|f\|_{V^*}+C_0\|g\|_{H^1}$.
Twierdzenie Laxa--Milgrama daje jednoznaczne $w\in V$ oraz
$\|w\|_V\le\nu^{-1}\|G\|_{V^*}$.
Ponieważ $\|w\|_{H^1}\le(1+C_P^2)^{1/2}\|w\|_V$,
uzyskujemy żądane oszacowanie dla $u=w+g$.
Równość śladów wynika z
$\operatorname{Tr}(u-g)=0$ i twierdzenia
\ref{thm:pde10-trace-zero}.
\end{proof}

\begin{uwaga}[Dane na samym brzegu]
Dane $\varphi$ na $\partial\Omega$ można podać niezależnie,
ale nie każde $\varphi\in L^2(\partial\Omega)$ jest śladem
funkcji z $H^1(\Omega)$. W twierdzeniu zakładamy explicite
istnienie podniesienia $g\in H^1(\Omega)$.
Przestrzeń danych brzegowych $H^{1/2}(\partial\Omega)$
z normą ilorazową i jej liniowe podnoszenie zdefiniowano
już w~\ref{def:pde10-halftrace} i
\ref{lem:pde10-halftrace-rightinverse}. Równoważność tej
normy z normą Gagliarda pozostaje oddzielnym twierdzeniem.
\end{uwaga}


\begin{uwaga}[Wersja z danymi zadanymi na brzegu]
Twierdzenie~\ref{thm:pde11-lifting} sformułowano
z wybranym podniesieniem $g\in H^1(\Omega)$.
Z lematu~\ref{lem:pde10-halftrace-rightinverse} można
bez nowych argumentów otrzymać dane
$\varphi\in H^{1/2}(\partial\Omega)$ i oszacowanie
przez ich normę śladową. Uogólnienie to, także
z wyrazami niższego rzędu, zostaje jawnie sformułowane
w twierdzeniu~\ref{thm:pde13-nonhom}; nie powielamy
tutaj jego dowodu.
\end{uwaga}

\section{Zagadnienie Neumanna i warunek zgodności}
\label{sec:pde11-neumann}

Dla laplasjanu z warunkiem Neumanna nie można oczekiwać
jednoznaczności: dodanie stałej do rozwiązania nie zmienia
gradientu ani pochodnej normalnej. Warunek całkowy
pojawia się natomiast już w sformułowaniu słabym.

\begin{twierdzenie}[Słabe zagadnienie Neumanna]
\label{thm:pde11-neumann}
Niech $\Omega\subset\R^d$ będzie obszarem ograniczonym Lipschitzowskim, $f\in L^2(\Omega)$ oraz
$h\in L^2(\partial\Omega)$.
Rozważmy równanie słabe
\[
\int_\Omega\nabla u\cdot\nabla v\,\mathrm d\lambda_d(x)
=\int_\Omega fv\,\mathrm d\lambda_d(x)
+\int_{\partial\Omega}h\operatorname{Tr}v
\,\mathrm d\mathcal H^{d-1}(x)
\quad(v\in H^1(\Omega)).
\]
Istnieje rozwiązanie $u\in H^1(\Omega)$ wtedy i tylko wtedy, gdy
\[
\int_\Omega f\,\mathrm d\lambda_d(x)
+\int_{\partial\Omega}h\,\mathrm d\mathcal H^{d-1}(x)=0.
\]
Gdy ono istnieje, wszystkie rozwiązania różnią się o stałą.
Po narzuceniu normalizacji $\int_\Omega u\,\mathrm d\lambda_d=0$
rozwiązanie jest jednoznaczne.
\end{twierdzenie}
\begin{proof}
\textbf{Konieczność.} Funkcja $v\equiv1$ należy do $H^1(\Omega)$
i ma zerowy gradient, więc jej podstawienie daje warunek całkowy.

\textbf{Wystarczalność.} Niech
$V_\circ=\{v\in H^1(\Omega):\int_\Omega v\,\mathrm d\lambda_d=0\}$.
Jest to domknięta podprzestrzeń $H^1$, więc przestrzeń Hilberta.
Nierówność Poincar\'ego z odjęciem średniej
(wniosek~\ref{cor:pde10-Poincare}) pokazuje, że
$\|v\|_{V_\circ}=\|\nabla v\|_2$ jest na niej normą
równoważną normie $H^1$.
Forma $a(u,v)=\int_\Omega\nabla u\cdot\nabla v\,
\mathrm d\lambda_d$ jest ograniczona i koercywna w tej normie.
Prawa strona $F(v)$ jest liniowa i ciągła na $H^1(\Omega)$:
$|\int_\Omega fv|\le\|f\|_2\|v\|_2$,
a twierdzenie o śladzie~\ref{thm:pde10-trace} daje
$|\int_{\partial\Omega}h\operatorname{Tr}v|
\le C\|h\|_{L^2(\partial\Omega)}\|v\|_{H^1}$.
Z twierdzenia Laxa--Milgrama istnieje jednoznaczne $u\in V_\circ$
z $a(u,v)=F(v)$ dla wszystkich $v\in V_\circ$.
Dla dowolnego $v\in H^1(\Omega)$ rozkładamy
$v=(v-v_\Omega)+v_\Omega$.
Pierwszy składnik należy do $V_\circ$; dla stałej
$v_\Omega$ obie strony równania są równe zeru dzięki
warunkowi zgodności. Dostajemy równość dla każdego $v\in H^1$.

\textbf{Jednoznaczność modulo stałe.} Jeśli $u_1,u_2$
są rozwiązaniami, różnica $w=u_1-u_2$ spełnia
$\int_\Omega|\nabla w|^2\,\mathrm d\lambda_d=0$
po podstawieniu $v=w$. Wobec tego $\nabla w=0$ prawie wszędzie,
a spójność $\Omega$ implikuje, że $w$ jest stała prawie wszędzie.
\end{proof}

\section{Metoda bezpośrednia rachunku wariacyjnego}
\label{sec:pde11-direct}

Założenie, że rozwiązanie jest punktem minimum odpowiedniej energii,
pozwala badać również równania nieliniowe.
Do tego potrzebujemy precyzyjnego rozróżnienia zbieżności
silnej i słabej.

\begin{definicja}[Zbieżność słaba]
\label{def:pde11-weakconv}
W przestrzeni Hilberta $H$ piszemy $u_n\rightharpoonup u$,
jeżeli $(u_n,v)_H\to(u,v)_H$ dla każdego $v\in H$.
W przestrzeni Banacha $X$ oznacza to warunek
$F(u_n)\to F(u)$ dla każdego $F\in X^*$.
Funkcjonał $J:X\to(-\infty,+\infty]$ jest
\emph{słabo półciągły z dołu w sensie ciągów}, jeżeli
$u_n\rightharpoonup u$ implikuje
$J(u)\le\liminf_n J(u_n)$.
Jest \emph{koercywny}, jeżeli $J(u)\to+\infty$,
gdy $\|u\|_X\to\infty$.
\end{definicja}

\begin{lemat}[Słaba zwartość ograniczonych ciągów w Hilbercie]
\label{lem:pde11-weakcompact}
Każdy ciąg ograniczony w rzeczywistej przestrzeni Hilberta
$H$ zawiera podciąg słabo zbieżny.
\end{lemat}
\begin{proof}
Niech $(u_n)$ będzie ciągiem ograniczonym przez $C$.
Jego domknięta otoczka liniowa
$H_0=\overline{\operatorname{span}\{u_n:n\ge1\}}$
jest separowalną przestrzenią Hilberta.
Jeżeli $H_0$ jest nieskończeniewymiarowa,
wybieramy przeliczalną bazę ortonormalną $(e_j)$.
Każdy ciąg liczb $((u_n,e_j))_n$ jest ograniczony;
argument przekątniowy daje podciąg $(u_{n_k})$ taki, że
$(u_{n_k},e_j)\to a_j$ dla każdego $j$.
Z nierówności Bessela
$\sum_{j=1}^N|a_j|^2
=\lim_{k\to\infty}\sum_{j=1}^N|(u_{n_k},e_j)|^2\le C^2$
dla każdego $N$. Zatem $\sum_j|a_j|^2\le C^2$,
więc istnieje $u\in H_0$ o współrzędnych $(a_j)$.
Dla skończonych kombinacji wektorów $e_j$ zachodzi
$(u_{n_k},v)\to(u,v)$.
Dla ogólnego $v\in H_0$ przybliżamy go takimi kombinacjami:
$|(u_{n_k}-u,v-v_N)|\le 2C\|v-v_N\|$,
co pozwala najpierw przejść $k\to\infty$, a potem $N\to\infty$.
Dla $v\in H$ rozkład $v=P_{H_0}v+(v-P_{H_0}v)$
redukuje twierdzenie do $H_0$, bo $u_n,u\in H_0$.
W przypadku skończeniewymiarowym zwykła zwartość
ograniczonego zbioru dostarcza nawet zbieżności silnej.
\end{proof}

\begin{lemat}[Słaba półciągłość normy]
\label{lem:pde11-weaklsc}
Niech $H$ będzie rzeczywistą przestrzenią Hilberta.
Jeżeli $u_n\rightharpoonup u$ w $H$, to
$\|u\|_H\le\liminf_n\|u_n\|_H$.
Ponadto dla każdego $F\in H^*$ mamy $F(u_n)\to F(u)$.
\end{lemat}
\begin{proof}
Jeżeli $u=0$, nierówność jest oczywista.
W przeciwnym razie słaba zbieżność daje
$\|u\|_H^2=\lim_n(u_n,u)_H
\le(\liminf_n\|u_n\|_H)\|u\|_H$,
po czym dzielimy przez $\|u\|_H$.
Druga własność wynika z definicji i twierdzenia Riesza.
\end{proof}

\begin{twierdzenie}[Metoda bezpośrednia w przestrzeni Hilberta]
\label{thm:pde11-direct}
Niech $H$ będzie rzeczywistą przestrzenią Hilberta,
a $J:H\to\R$ funkcjonałem ograniczonym z dołu,
koercywnym i słabo półciągłym z dołu w sensie ciągów.
Wówczas istnieje $u\in H$ taki, że
$J(u)=\inf_{v\in H}J(v)$.
Jeżeli $J$ jest ponadto ściśle wypukły,
to minimizator jest jedyny.
\end{twierdzenie}
\begin{proof}
Niech $m=\inf_HJ\in\R$ i wybierzmy $u_n\in H$ z
$m\le J(u_n)<m+1/n$. Koercywność zapewnia
ograniczoność $(u_n)$: gdyby istniał podciąg
o normach dążących do nieskończoności,
to jego energie dążyłyby do $+\infty$, w sprzeczności z
$J(u_n)\to m$.
Na mocy lematu~\ref{lem:pde11-weakcompact}
istnieje podciąg $u_{n_k}\rightharpoonup u$.
Słaba półciągłość z dołu daje
$J(u)\le\liminf_kJ(u_{n_k})=m$.
Z definicji infimum również $m\le J(u)$;
zatem $J(u)=m$.
Jeśli istniałyby dwa różne minimizatory $u,w$,
ścisła wypukłość zapewniłaby
$J((u+w)/2)<(J(u)+J(w))/2=m$, sprzeczność.
\end{proof}

\begin{uwaga}[Dlaczego słaba półciągłość jest niezbędna?]
Sama koercywność nie wystarcza. Przykładowo na $H=\R$
zdefiniujmy $J(0)=1$, $J(x)=x^2$ dla $x\ne0$.
Funkcjonał jest ograniczony z dołu i koercywny,
lecz jego infimum wynosi $0$ i nie jest osiągane.
Nie jest też półciągły z dołu w punkcie $0$.
\end{uwaga}

\section{Energia kwadratowa i konieczność symetrii}
\label{sec:pde11-energy}

\begin{twierdzenie}[Zasada Dirichleta dla formy symetrycznej]
\label{thm:pde11-dirichlet}
Niech $H$ będzie rzeczywistą przestrzenią Hilberta,
$a:H\times H\to\R$ ograniczoną, \emph{symetryczną}
i koercywną formą dwuliniową, $F\in H^*$.
Zdefiniujmy $J:H\to\R$ wzorem
$J(v)=\frac12 a(v,v)-F(v)$.
Wtedy $J$ ma dokładnie jeden minimizator $u$, który jest
jednocześnie jedynym rozwiązaniem słabym
$a(u,v)=F(v)$ dla wszystkich $v\in H$.
Dla każdego $w\in H$ zachodzi tożsamość
\[ J(w)-J(u)=\frac12a(w-u,w-u)
\ge\frac\alpha2\|w-u\|_H^2. \]
\end{twierdzenie}
\begin{proof}
\textbf{Koercywność funkcjonału.} Z założeń otrzymujemy
\[ J(v)\ge\frac\alpha2\|v\|_H^2
 -\|F\|_{H^*}\|v\|_H\longrightarrow+\infty
 \quad\text{gdy }\|v\|_H\to\infty. \]

\textbf{Półciągłość z dołu.} Koercywna forma symetryczna
$a$ definiuje iloczyn skalarny $(v,w)_a=a(v,w)$.
Ograniczoność i koercywność dają równoważność norm:
$\sqrt\alpha\|v\|_H\le\|v\|_a\le\sqrt M\|v\|_H$.
Z twierdzenia Riesza w $(H,(\cdot,\cdot)_a)$ i
równoważności norm wynika, że zbieżność słaba w $H$
jest taka sama jak zbieżność słaba w nowym iloczynie.
Lemat~\ref{lem:pde11-weaklsc} daje więc
$\|v\|_a^2\le\liminf_n\|v_n\|_a^2$
dla $v_n\rightharpoonup v$; składnik $F(v_n)$ zbiega do $F(v)$.
Stąd $J$ jest słabo półciągły z dołu.
Twierdzenie~\ref{thm:pde11-direct} daje minimizator.

\textbf{Równanie Eulera.} Dla dowolnych $u,v\in H$
i $t\in\R$ wykorzystujemy symetrię:
$J(u+tv)=J(u)+t[a(u,v)-F(v)]+\frac{t^2}{2}a(v,v)$.
Jeżeli $u$ minimalizuje $J$, wielomian po prawej ma
minimum w $t=0$, więc jego współczynnik liniowy znika:
$a(u,v)=F(v)$. Odwrotnie, jeśli $u$ spełnia tę równość,
to po rozwinięciu $J(u+(w-u))$ otrzymujemy
$J(w)-J(u)=\frac12a(w-u,w-u)\ge0$.
Koercywność czyni nierówność ostrą dla $w\ne u$,
dając jednoznaczność. Ostatnia nierówność daje również
ilościową kontrolę błędu energetycznego.
\end{proof}

\begin{kontrprzyklad}[Dlaczego symetrii nie wolno pominąć]
\label{ex:pde11-nonsymmetric}
Na $H=\R^2$ przyjmijmy
$B=\begin{pmatrix}1&-1\\1&1\end{pmatrix}$,
$a(u,v)=(Bu)\cdot v$ oraz $F(v)=(1,0)\cdot v$.
Wtedy $a(u,u)=|u|^2$, więc $a$ jest koercywna i ograniczona.
Minimizer funkcjonału $J(u)=|u|^2/2-(1,0)\cdot u$
wynosi $u_J=(1,0)$. Tymczasem równanie
$a(u,v)=F(v)$ dla wszystkich $v$ jest równoważne
$Bu=(1,0)$ i ma rozwiązanie $u_*=\frac12(1,-1)$.
Zatem $u_J\ne u_*$. Część antysymetryczna formy znika
w $a(v,v)$, ale nie znika z samego równania słabego.
\end{kontrprzyklad}

\section[Operator p-Laplace'a]{Pierwszy problem nieliniowy: operator $p$-Laplace'a}
\label{sec:pde11-plaplace}

W części nieliniowej musimy użyć faktu, że dla
$1<p<\infty$ przestrzeń $W_0^{1,p}(\Omega)$ jest refleksywna.
Jest ona domkniętą podprzestrzenią produktu skończonej liczby
przestrzeni $L^p$, a przestrzenie $L^p$ są refleksywne w tym zakresie.
Będziemy korzystać ze standardowego twierdzenia, że każdy ciąg
ograniczony w przestrzeni refleksywnej ma słabo zbieżny podciąg
(twierdzenie Eberleina--\v Smuliana); dowód tego ogólnego wyniku
należy do analizy funkcjonalnej, nie do konstrukcji samego PDE.

\begin{twierdzenie}[Istnienie i jednoznaczność dla $p$-laplasjanu]
\label{thm:pde11-plaplace}
Niech $\Omega\subset\R^d$ będzie obszarem ograniczonym Lipschitzowskim, $1<p<\infty$, $p'=p/(p-1)$,
a $f\in L^{p'}(\Omega)$. Istnieje dokładnie jeden
$u\in W_0^{1,p}(\Omega)$ taki, że
\[
\int_\Omega |\nabla u|^{p-2}\nabla u\cdot\nabla v
\,\mathrm d\lambda_d(x)
=\int_\Omega fv\,\mathrm d\lambda_d(x)
\quad\text{dla każdego }v\in W_0^{1,p}(\Omega).
\]
W szczególności $u$ jest rozwiązaniem dystrybucyjnym
$-\operatorname{div}(|\nabla u|^{p-2}\nabla u)=f$.
Ponadto
$\|\nabla u\|_p^{p-1}\le C_P\|f\|_{p'}$.
\end{twierdzenie}
\begin{proof}
\textbf{Krok 1: energia i ograniczenie z dołu.}
Na $X=W_0^{1,p}(\Omega)$ przyjmijmy
$\|v\|_X=\|\nabla v\|_p$, co jest normą równoważną
$W^{1,p}$ dzięki nierówności Poincar\'ego.
Definiujemy
\[ J_p(v)=\frac1p\int_\Omega|\nabla v|^p
\,\mathrm d\lambda_d(x)-\int_\Omega fv
\,\mathrm d\lambda_d(x). \]
H\"older i Poincar\'e dają
$|\int_\Omega fv|\le\|f\|_{p'}\|v\|_p
\le C_P\|f\|_{p'}\|\nabla v\|_p$.
Zatem
$J_p(v)\ge\frac1p\|v\|_X^p-C_P\|f\|_{p'}\|v\|_X$,
co implikuje koercywność i ograniczenie z dołu,
ponieważ $p>1$.

\textbf{Krok 2: minimalizator.}
Niech $(v_n)$ będzie ciągiem minimalizującym.
Koercywność zapewnia ograniczoność w $X$,
a refleksywność pozwala wybrać podciąg
$v_{n_k}\rightharpoonup u$ w $X$.
Operator gradientu jest ograniczony z $X$ do $L^p(\Omega)^d$,
więc $\nabla v_{n_k}\rightharpoonup\nabla u$ w $L^p$.
Norma w przestrzeni Banacha jest słabo półciągła z dołu:
$\|z\|=\sup_{\|\ell\|_{X^*}\le1}\ell(z)$,
a przy zbieżności słabej każdy składnik supremum
przechodzi do granicy; stąd
$\|\nabla u\|_p^p\le\liminf_k\|\nabla v_{n_k}\|_p^p$.
Ponieważ funkcjonał $v\mapsto\int_\Omega fv$ jest
ciągły i liniowy na $X$, jego wartości na $v_{n_k}$
zbiegają do wartości na $u$. Zatem
$J_p(u)\le\liminf_kJ_p(v_{n_k})=\inf_XJ_p$.
Odwrotna nierówność wynika z definicji infimum,
więc $u$ minimalizuje $J_p$.

\textbf{Krok 3: wyprowadzenie słabego równania.}
Dla $u,v\in X$ rozważamy $\psi(t)=J_p(u+tv)$.
Funkcja $\xi\mapsto|\xi|^p/p$ jest klasy $C^1$ na $\R^d$
dla każdego $p>1$ i ma gradient
$|\xi|^{p-2}\xi$ (w zerze przyjmujemy wartość $0$).
W celu uzasadnienia zamiany granicy z całkowaniem
ustalmy $|t|\le1$. Twierdzenie o wartości średniej daje
\[
\left|\frac{|\nabla u+t\nabla v|^p-|\nabla u|^p}{pt}\right|
\le C_p\big(|\nabla u|^{p-1}|\nabla v|+|\nabla v|^p\big).
\]
Funkcja po prawej jest całkowalna: pierwszy składnik
należy do $L^1$ dzięki H\"olderowi z wykładnikami
$p'$ i $p$, a drugi dzięki $v\in W^{1,p}$.
Twierdzenie o zbieżności zdominowanej daje
\[ \psi'(0)=\int_\Omega |\nabla u|^{p-2}\nabla u
\cdot\nabla v\,\mathrm d\lambda_d(x)
-\int_\Omega fv\,\mathrm d\lambda_d(x). \]
Ponieważ $u$ jest minimizatorem, $\psi'(0)=0$
dla każdego $v\in X$, co daje słabe równanie.

\textbf{Krok 4: jednoznaczność.}
Funkcja $\xi\mapsto|\xi|^p/p$ jest ściśle wypukła dla
$p>1$. Jej gradient jest zatem ściśle monotoniczny:
$(|\xi|^{p-2}\xi-|\eta|^{p-2}\eta)\cdot(\xi-\eta)>0$
dla $\xi\ne\eta$. Jeśli $u_1,u_2$ są rozwiązaniami,
testujemy różnicę ich równań funkcją $v=u_1-u_2$.
Wtedy całka z nieujemnej funkcji jest równa $0$,
więc $\nabla u_1=\nabla u_2$ prawie wszędzie.
Nierówność Poincar\'ego na $X$ daje $u_1=u_2$.

\textbf{Krok 5: oszacowanie.}
W słabym równaniu podstawiamy $v=u$.
Dostajemy $\|\nabla u\|_p^p=\int_\Omega fu
\,\mathrm d\lambda_d\le C_P\|f\|_{p'}\|\nabla u\|_p$.
Dla $\|\nabla u\|_p>0$ dzielimy przez tę normę;
przypadek zerowy jest oczywisty.
\end{proof}

\begin{uwaga}[Dlaczego pominięto przypadek $p=1$?]
Przy $p=1$ energia $\int_\Omega|\nabla u|\,\mathrm d\lambda_d(x)$ nie jest ściśle wypukła,
a $W_0^{1,1}$ nie jest refleksywna. Nie można zatem
bezkrytycznie powtórzyć dowodu; pojawiają się przestrzenie
funkcji o ograniczonej wariacji ($BV$) oraz inna postać
równań Eulera--Lagrange'a.
\end{uwaga}

\section{Metoda Galerkina i pierwsze oszacowanie błędu}
\label{sec:pde11-galerkin}

\begin{twierdzenie}[Lemat C\'ei]
\label{thm:pde11-cea}
Niech $H$ będzie rzeczywistą przestrzenią Hilberta,
$a:H\times H\to\R$ formą ograniczoną ze stałą $M$
i koercywną ze stałą $\alpha>0$, a $F\in H^*$.
Niech $H_h\subset H$ będzie skończeniewymiarową
podprzestrzenią. Jeśli $u\in H$ spełnia
$a(u,v)=F(v)$ dla wszystkich $v\in H$, a $u_h\in H_h$
spełnia $a(u_h,v_h)=F(v_h)$ dla wszystkich $v_h\in H_h$,
to $u_h$ istnieje i jest jednoznaczne oraz
\[
\|u-u_h\|_H\le\frac M\alpha
\inf_{w_h\in H_h}\|u-w_h\|_H.
\]
Gdy $a$ jest symetryczna, $u_h$ jest rzutem ortogonalnym
$u$ na $H_h$ w iloczynie skalarnym $a$, a więc
$\|u-u_h\|_a=\inf_{w_h\in H_h}\|u-w_h\|_a$.
\end{twierdzenie}
\begin{proof}
Ograniczenie formy $a$ do $H_h\times H_h$ zachowuje
ograniczoność i koercywność. Z twierdzenia Laxa--Milgrama
na $H_h$ otrzymujemy dokładnie jedno $u_h$.
Odejmując równania, dostajemy \emph{ortogonalność Galerkina}:
$a(u-u_h,v_h)=0$ dla wszystkich $v_h\in H_h$.
Niech $e=u-u_h$ i $w_h\in H_h$ będzie dowolny.
Ponieważ $w_h-u_h\in H_h$, mamy
\begin{align*}
\alpha\|e\|_H^2
&\le a(e,e)=a(e,u-w_h)+a(e,w_h-u_h)\\
&=a(e,u-w_h)\le M\|e\|_H\|u-w_h\|_H.
\end{align*}
Po podzieleniu przez $\alpha\|e\|_H$ i wzięciu infimum
po $w_h$ otrzymujemy tezę; dla $e=0$ jest oczywista.
Przy symetrii $a$ jest iloczynem skalarnym.
Ortogonalność Galerkina staje się
$(u-u_h,v_h)_a=0$, więc twierdzenie Pitagorasa
daje $\|u-w_h\|_a^2=\|u-u_h\|_a^2+
\|u_h-w_h\|_a^2$ dla wszystkich $w_h\in H_h$,
z czego wynika minimalność błędu.
\end{proof}

\begin{wniosek}[Zbieżność przy zagęszczaniu przestrzeni]
\label{cor:pde11-galerkin-density}
Przy założeniach twierdzenia~\ref{thm:pde11-cea},
niech $H_n\subset H$ będą skończeniewymiarowymi
podprzestrzeniami takimi, że
$\inf_{w\in H_n}\|u-w\|_H\to0$ dla rozważanego rozwiązania $u$.
Wówczas $\|u-u_n\|_H\to0$, gdzie $u_n$ jest rozwiązaniem otrzymanym metodą Galerkina
w $H_n$.
\end{wniosek}
\begin{proof}
Prawa strona oszacowania C\'ei dąży do zera.
Nie jest potrzebne założenie, że $H_n\subset H_{n+1}$,
lecz sama gęstość \emph{sumy} podprzestrzeni bez
kontroli dla kolejnych $n$ nie wystarczyłaby.
\end{proof}

\section{Metoda elementów skończonych na odcinku}
\label{sec:pde11-fem}

Pokażemy, że teoria wariacyjna prowadzi do konkretnego układu
macierzowego. Wróćmy do $-u''=1$ na $(0,1)$ z
warunkami $u(0)=u(1)=0$. Dzielimy $[0,1]$ na
$N\ge2$ równych pododcinków długości $h=1/N$,
a węzły oznaczamy $x_i=ih$, $i=0,\ldots,N$.

Niech $V_h\subset H_0^1(0,1)$ będzie przestrzenią funkcji
ciągłych, odcinkami afinicznych, zerujących się na końcach.
Funkcje bazowe $\phi_i$, $i=1,\ldots,N-1$, spełniają
$\phi_i(x_j)=\delta_{ij}$. Przyjmując
$u_h=\sum_{i=1}^{N-1}U_i\phi_i$,
otrzymujemy układ $KU=b$, gdzie
\[
K_{ij}=\int_0^1\phi_i'\phi_j'\,\mathrm d\lambda_1,
\qquad b_i=\int_0^1\phi_i\,\mathrm d\lambda_1.
\]
Rachunek na dwóch przedziałach, na których każda
$\phi_i$ jest niezerowa, daje
$K_{ii}=2/h$, $K_{i,i+1}=K_{i+1,i}=-1/h$,
$K_{ij}=0$ dla $|i-j|>1$ oraz $b_i=h$.

\begin{stwierdzenie}[Rozwiązanie i dokładne błędy na siatce]
\label{prop:pde11-fem-errors}
Przy powyższych definicjach rozwiązaniem macierzowym jest
$U_i=u(x_i)=x_i(1-x_i)/2$, $i=1,\ldots,N-1$.
Jeśli $u_h$ jest funkcją odcinkami afiniczną o tych wartościach,
to
\[
\|u'-u_h'\|_{L^2(0,1)}=\frac h{\sqrt{12}},
\qquad
\|u-u_h\|_{L^2(0,1)}=\frac{h^2}{\sqrt{120}}.
\]
\end{stwierdzenie}
\begin{proof}
Niech $U_0=U_N=0$ oraz
$U_i=(ih)(1-ih)/2$. Ponieważ $U_i$ jest
kwadratowym wielomianem w $i$, mamy
$-U_{i-1}+2U_i-U_{i+1}=h^2$.
Po pomnożeniu przez $1/h$ dostajemy
$-U_{i-1}/h+2U_i/h-U_{i+1}/h=h=b_i$,
więc wektor $U$ rozwiązuje układ. Dodatnia określoność
macierzy $K$ wynika z
$U^TKU=\int_0^1|\sum_iU_i\phi_i'|^2\,
\mathrm d\lambda_1>0$ dla $U\ne0$;
zatem rozwiązanie jest jednoznaczne.

Na komórce $I_i=[x_i,x_{i+1}]$ wielomian $u$
ma pochodną $u'(x)=1/2-x$, a pochodna interpolantu
$u_h'$ jest stała i równa $u'(m_i)$,
gdzie $m_i=(x_i+x_{i+1})/2$.
W konsekwencji
\[
\int_{I_i}|u'-u_h'|^2\,\mathrm d\lambda_1(x)
=\int_{-h/2}^{h/2}s^2\,\mathrm d\lambda_1(s)=h^3/12.
\]
Ponieważ jest $N=1/h$ komórek, kwadrat całkowitego błędu
jest równy $h^2/12$.
Dalej na $I_i$ zachodzi dokładnie
$u(x)-u_h(x)=(x-x_i)(x_{i+1}-x)/2$.
Dlatego
\[
\int_{I_i}|u-u_h|^2\,\mathrm d\lambda_1(x)
=\frac14\int_0^h s^2(h-s)^2\,\mathrm d\lambda_1(s)
=\frac{h^5}{120}.
\]
Po zsumowaniu $N$ identycznych wkładów otrzymujemy
$\|u-u_h\|_2^2=h^4/120$, jak w tezie.
\end{proof}

\begin{center}
\begin{tikzpicture}[x=9cm,y=7cm,>=Stealth]
 \draw[->] (-.02,0)--(1.04,0) node[right] {$x$};
 \draw[->] (0,-.01)--(0,.17) node[above] {$u$};
 \draw[thick,blue,domain=0:1,samples=80] plot (\x,{\x*(1-\x)/2});
 \draw[orange!85!black,thick] (0,0)--(.25,.09375)--(.5,.125)--(.75,.09375)--(1,0);
 \foreach \x in {0,.25,.5,.75,1}{\fill[orange!85!black] (\x,{\x*(1-\x)/2}) circle[radius=.5pt];}
 \draw[blue,thick] (.13,-.033)--(.22,-.033);
 \node[black,anchor=west,font=\footnotesize] at (.23,-.033) {dokładne};
 \draw[orange!85!black,thick] (.56,-.033)--(.65,-.033);
 \node[black,anchor=west,font=\footnotesize] at (.66,-.033) {Galerkin};
\end{tikzpicture}

{\small\emph{Rysunek:} Dokładne rozwiązanie $-u''=1$ oraz
rozwiązanie Galerkina dla czterech elementów liniowych.}
\end{center}

\section{Zakres zastosowań i granice metody}
\label{sec:pde11-limits}

\begin{uwaga}[Koercywność jest warunkiem wystarczającym]
Twierdzenie Laxa--Milgrama nie obejmuje każdej formy eliptycznej.
Na przykład forma $a(u,v)=\int_\Omega\nabla u\cdot\nabla v
-\lambda\int_\Omega uv$ może przestać być koercywna
przy dużym $\lambda$. Nie oznacza to automatycznie braku
rozwiązania: zagadnienie wymaga analizy widma lub
alternatywy Fredholma. Zbadamy to w rozdziale 12.
\end{uwaga}

\begin{uwaga}[Od istnienia do regularności]
W dowodzie słabego zagadnienia Dirichleta nie założyliśmy
gładkości współczynników $A$ ani brzegu wyższej niż
Lipschitzowska. To wystarcza do istnienia rozwiązania
w $H_0^1$, ale nie gwarantuje jego regularności $H^2$.
W rozdziałach o regularności eliptycznej będziemy analizować,
które dodatkowe założenia pozwalają odzyskać drugie
pochodne i oszacowania wyższych rzędów.
\end{uwaga}

\begin{uwaga}[Co zostanie rozwinięte dalej]
Założenie $c\ge0$ i brak wyrazów pierwszego rzędu
upraszczają dowód. Dla członów dryftu potrzebne bywają
oszacowania perturbacyjne lub warunki dodatkowe.
W następnym rozdziale przejdziemy do operatorów nieograniczonych,
teorii spektralnej i półgrup; wcześniej otrzymane twierdzenia
pozwolą m.in. skonstruować odwrotność laplasjanu Dirichleta.
\end{uwaga}

% ==============================================================
% ROZDZIAŁ 12. OPERATORY NIEOGRANICZONE I METODY SPEKTRALNE
% ==============================================================

\chapter[Operatory nieograniczone i teoria spektralna]{Operatory nieograniczone, samosprzężoność, spektrum i półgrupy}
\label{chap:pde12}

\providecommand{\atlasref}[2]{#2}

W rozdziale~\ref{chap:pde-variational} zagadnienie Dirichleta rozwiązaliśmy
w przestrzeni energetycznej $H_0^1(\Omega)$. Teraz zapytamy, kiedy
formalny operator $-\Delta$ jest operatorem samosprzężonym w
$L^2(\Omega)$, jak wyglądają jego wartości własne i dlaczego
jego widmo steruje ewolucją cieplną i falową. Dalsze rachunki
zostaną zapisane najpierw dla rzeczywistych przestrzeni Hilberta.
W argumentach z operatorami unitarnymi i widmem zespolonym
korzystamy z ich zespolenia. W zespolonej przestrzeni Hilberta
iloczyn skalarny przyjmujemy liniowy w pierwszym argumencie.

\section{Dziedzina jest częścią operatora}
\label{sec:pde12-domain}

\begin{definicja}[Operator liniowy, domknięty i domykalny]
\label{def:pde12-operator}
Niech $H$ będzie przestrzenią Hilberta. Operatorem liniowym
$T$ w $H$ nazywamy parę składającą się z podprzestrzeni
$D(T)\subset H$, zwanej \emph{dziedziną}, oraz odwzorowania
liniowego $T:D(T)\to H$. Wykresem $T$ nazywamy
$\mathcal G(T)=\{(u,Tu):u\in D(T)\}\subset H\times H$.
Wykres jest podprzestrzenią liniową $H\times H$:
dla $u,v\in D(T)$ i skalarów $a,b$ mamy
$a(u,Tu)+b(v,Tv)=(au+bv,T(au+bv))$.
Operator $T$ jest \emph{domknięty}, jeżeli $\mathcal G(T)$ jest
domknięty w $H\times H$. Jest \emph{domykalny}, jeżeli
domknięcie $\mathcal G(T)$ jest wykresem pewnego operatora;
operator ten oznaczamy przez $\overline T$.
Równoważnie $T$ jest domykalny wtedy i tylko wtedy, gdy
dla każdego ciągu $(u_n)$ z $D(T)$ warunki
$u_n\to0$ i $Tu_n\to v$ implikują $v=0$.
Istotnie, domknięcie wykresu jest podprzestrzenią liniową
i jest wykresem dokładnie wtedy, gdy nad pierwszą współrzędną
$0$ zawiera jedynie parę $(0,0)$: dwie pary $(u,v_1)$
i $(u,v_2)$ dałyby przez odjęcie $(0,v_1-v_2)$.
Odwrotnie, jeśli warunek ten zachodzi, na rzucie domknięcia
wykresu na pierwszą współrzędną (podprzestrzeni liniowej)
możemy jednoznacznie przypisać każdemu $u$ drugą współrzędną.
Liniowość domknięcia wykresu daje liniowość tego przypisania.
Ponieważ $H\times H$ jest przestrzenią metryczną,
$(0,v)$ należy do domknięcia wykresu dokładnie wtedy,
gdy jest granicą pewnego ciągu $(u_n,Tu_n)$.
W przypadku domykalności operator $\overline T$ jest
liniowy i domknięty, rozszerza $T$ i zawiera się w każdym
domkniętym rozszerzeniu $T$, ponieważ wykres takiego
rozszerzenia zawiera domknięcie wykresu $T$.
\end{definicja}

\begin{lemat}[Kryterium domkniętości i norma grafowa]
\label{lem:pde12-graph}
Dla operatora liniowego $T:D(T)\subset H\to H$
następujące warunki są równoważne:
(i) $T$ jest domknięty; (ii) ilekroć $u_n\in D(T)$,
$u_n\to u$ i $Tu_n\to v$ w $H$, to $u\in D(T)$ oraz $Tu=v$;
(iii) przestrzeń $D(T)$ z normą
$\|u\|_{D(T)}=(\|u\|_H^2+\|Tu\|_H^2)^{1/2}$ jest zupełna.
\end{lemat}
\begin{proof}
Zbieżność $(u_n,Tu_n)\to(u,v)$ w iloczynie $H\times H$
jest równoważna dwóm zbieżnościom z (ii), więc (i) i (ii)
są równoważne. Odwzorowanie $u\mapsto(u,Tu)$ jest izometrią
$D(T)$ z normą grafową na $\mathcal G(T)$ z normą iloczynową:
jest liniowe, różnowartościowe dzięki pierwszej współrzędnej
i na mocy definicji wykresu surjektywne. Norma grafowa
pochodzi więc od iloczynu skalarnego
$(u,v)_{D(T)}=(u,v)_H+(Tu,Tv)_H$.
Domknięta podprzestrzeń przestrzeni zupełnej jest zupełna.
Odwrotnie, jeśli $(u_n,Tu_n)$ zbiega w $H\times H$,
to $(u_n)$ jest ciągiem Cauchy'ego w normie grafowej.
Z zupełności $D(T)$ ma granicę $w\in D(T)$ w tej normie,
więc $(u_n,Tu_n)\to(w,Tw)$ w $H\times H$.
Jednoznaczność granicy daje przynależność pierwotnej
granicy do wykresu, który jest zatem domknięty. To daje (i)$\Leftrightarrow$(iii).
\end{proof}

\begin{przyklad}[Pochodna i dwie różne dziedziny]
\label{ex:pde12-derivative}
W $H=L^2(0,1;\C)$ rozważmy operator $D_0u=-iu'$ z dziedziną
$D(D_0)=H_0^1(0,1;\C)$. Dziedzina jest gęsta, ponieważ
zawiera $C_c^\infty(0,1;\C)$. Norma grafowa jest dokładnie
normą $H^1$, gdyż $\|-iu'\|_2=\|u'\|_2$.
Przestrzeń $H_0^1$ jest domknięta w zupełnej przestrzeni
$H^1$ (twierdzenie~\ref{thm:pde9-banach}), więc
lemat~\ref{lem:pde12-graph} daje domkniętość $D_0$.
Operator jest też symetryczny:
dla testów $u,v$ całkowanie przez części daje
\[
 (D_0u,v)=-i\int_0^1u'\overline v\,\mathrm d\lambda_1
 =i\int_0^1u\overline{v'}\,\mathrm d\lambda_1
 =(u,D_0v).
\]
Obie skrajne formy są ciągłe na $H^1\times H^1$
z nierówności Cauchy'ego--Schwarza, więc przez gęstość
testów w $H_0^1$ równość zachodzi dla całej dziedziny.
Nie jest jednak samosprzężony. W istocie
$D(D_0^*)=H^1(0,1;\C)$ i $D_0^*v=-iv'$.
Aby wykazać to ostatnie, dla $v\in D(D_0^*)$ podstawiamy
najpierw $u\in C_c^\infty(0,1)$: istnienie $g\in L^2$
takiego, że $\int_0^1(-iu')\overline v\,\mathrm d\lambda_1
=\int_0^1u\overline g\,\mathrm d\lambda_1$ oznacza
$i(\overline v)'=\overline g$ dystrybucyjnie: z definicji
pochodnej dystrybucji lewa strona działa na test $u$ przez
$-i\int_0^1u'\overline v\,\mathrm d\lambda_1$.
Po sprzężeniu mamy $-iv'=g\in L^2$, a więc $v\in H^1$.
Odwrotnie, dla $v\in H^1$ definicja słabej pochodnej daje
$(D_0u,v)=(u,-iv')$ najpierw dla testów $u$.
Dla $u\in H_0^1$ wybieramy testy $u_n\to u$ w $H^1$.
Różnice obu stron ograniczają odpowiednio
$\|u_n'-u'\|_2\|v\|_2$ i $\|u_n-u\|_2\|v'\|_2$,
więc tożsamość przechodzi na całą dziedzinę $D_0$.
To dowodzi obu inkluzji dziedzin i wzoru na sprzężenie.
W szczególności stała funkcja $1$ należy do $D(D_0^*)$,
ale nie do $D(D_0)$, gdyż jej ślad na końcach odcinka
nie jest zerowy (twierdzenie~\ref{thm:pde10-trace-zero}).
\end{przyklad}

Ogólne kryterium samosprzężoności gęsto określonych operatorów symetrycznych przez surjektywność operatorów przesuniętych podano w \atlasref{FA-33::thm:selfadjoint-range33}{FA-33: kryterium samosprzężoności}.

\begin{definicja}[Operator sprzężony i samosprzężoność]
\label{def:pde12-adjoint}
Niech $T:D(T)\subset H\to H$ ma gęstą dziedzinę.
Element $v\in H$ należy do $D(T^*)$ wtedy i tylko wtedy,
gdy istnieje $g\in H$ takie, że
$(Tu,v)_H=(u,g)_H$ dla wszystkich $u\in D(T)$.
Wówczas $T^*v=g$; gęstość $D(T)$ zapewnia
jednoznaczność $g$: różnica dwóch takich wektorów
jest ortogonalna do gęstego $D(T)$, a więc do całego $H$.
Jeśli $v,w\in D(T^*)$ i $a,b$ są skalarami, to
$(Tu,av+bw)_H=\overline a(u,T^*v)_H+
\overline b(u,T^*w)_H=(u,aT^*v+bT^*w)_H$
dla każdego $u\in D(T)$ (w przypadku rzeczywistym
sprzężenia są zbędne). Zatem $av+bw\in D(T^*)$ oraz
$T^*(av+bw)=aT^*v+bT^*w$.
Równoważnie $v\in D(T^*)$ wtedy i tylko wtedy, gdy
istnieje $C_v\ge0$ takie, że
$|(Tu,v)_H|\le C_v\|u\|_H$ dla każdego $u\in D(T)$.
Konieczność wynika z nierówności Cauchy'ego--Schwarza.
Dla dostateczności ograniczony funkcjonał $u\mapsto(Tu,v)_H$
przedłuża się jednoznacznie przez gęstość na $H$, a twierdzenie
Riesza o reprezentacji funkcjonału daje wektor $g\in H$
spełniający definicję sprzężenia.
Operator $T^*$ jest domknięty: jeśli $v_n\to v$
i $T^*v_n\to g$, to dla ustalonego $u\in D(T)$ przechodzimy
do granicy w $(Tu,v_n)_H=(u,T^*v_n)_H$, dostając
$(Tu,v)_H=(u,g)_H$. Z definicji $v\in D(T^*)$ i $T^*v=g$.
Operator jest \emph{symetryczny}, gdy
$(Tu,v)=(u,Tv)$ dla wszystkich $u,v\in D(T)$,
czyli $T\subset T^*$. Jest \emph{samosprzężony}, gdy
$T=T^*$, łącznie z równością dziedzin.
\end{definicja}

\begin{uwaga}
Symetria nie gwarantuje samosprzężoności; przykład
\ref{ex:pde12-derivative} pokazuje to nawet dla domkniętego
operatora różniczkowego pierwszego rzędu. Warunki brzegowe
mogą radykalnie zmienić dziedzinę operatora sprzężonego,
a więc także jego widmo i ewolucję.
\end{uwaga}

\section{Konstrukcja laplasjanu Dirichleta przez twierdzenie Laxa--Milgrama}
\label{sec:pde12-dirichlet}

W konstrukcji laplasjanu Dirichleta i jej zastosowaniach
przyjmujemy, o ile nie zaznaczono inaczej, że
$\Omega\subset\R^d$ jest niepustym obszarem ograniczonym
Lipschitzowskim, $H=L^2(\Omega;\R)$ i $V=H_0^1(\Omega;\R)$.
Przechodząc do widma zespolonego i grup unitarnych,
zespolimy te przestrzenie i operator. Przestrzeń $V$
wyposażamy w normę
$\|u\|_V^2=\|\nabla u\|_2^2+\|u\|_2^2$.
Rozważmy formę $b(u,v)=\int_\Omega(\nabla u\cdot\nabla v+uv)
\,\mathrm d\lambda_d(x)$. Jest ona iloczynem skalarnym
pełnej normy $H^1$ na $V$. Przestrzeń $V$ jest zupełna,
gdyż jest domkniętą podprzestrzenią $H^1$ na mocy
twierdzenia~\ref{thm:pde9-banach} i definicji $H_0^1$.
Jest też gęsta w $H$, ponieważ zawiera funkcje testowe.

\begin{twierdzenie}[Rezolwenta z twierdzenia Laxa--Milgrama]
\label{thm:pde12-resolvent}
Dla każdego $f\in H$ istnieje dokładnie jedno $u=Bf\in V$
takie, że $b(u,v)=(f,v)_H$ dla wszystkich $v\in V$.
Operator $B:H\to H$ jest liniowy, ograniczony, zwarty,
samosprzężony, dodatni i różnowartościowy. Ponadto
$\|B\|_{H\to H}\le1$ oraz $\overline{\operatorname{Ran}B}=H$.
\end{twierdzenie}
\begin{proof}
\textbf{Istnienie i oszacowanie.} Forma $b$ jest
iloczynem skalarnym na $V$, a więc jest ciągła i koercywna
ze stałą $1$. Funkcjonał $v\mapsto(f,v)_H$ jest ciągły na $V$:
$|(f,v)_H|\le\|f\|_2\|v\|_V$. Z twierdzenia
Laxa--Milgrama~\ref{thm:pde11-LM} otrzymujemy jedyne $u$.
Przy $v=u$ mamy
$\|u\|_V^2=(f,u)_H\le\|f\|_2\|u\|_2\le\|f\|_2\|u\|_V$.
Jeżeli $u\ne0$, dzielimy przez $\|u\|_V>0$;
gdy $u=0$, nierówność zachodzi bez dzielenia.
Stąd $\|u\|_V\le\|f\|_2$, a tym bardziej
$\|Bf\|_2\le\|f\|_2$. Dla $f,g\in H$ i $s,t\in\R$
dwuliniowość formy daje
$b(sBf+tBg,v)=(sf+tg,v)_H$ dla każdego $v\in V$.
Jednoznaczność implikuje $B(sf+tg)=sBf+tBg$.
Otrzymujemy ciągłość zarówno $B:H\to V$, jak i $B:H\to H$.

\textbf{Zwarty charakter.} Jeśli $(f_n)$ jest ograniczony w $H$,
to $(Bf_n)$ jest ograniczony w $V$. Twierdzenie
Rellicha--Kondrachova~\ref{thm:pde10-Rellich-p} (w postaci
zanurzenia $H_0^1(\Omega)\hookrightarrow L^2(\Omega)$)
dostarcza podciągu zbieżnego w $H$. To znaczy, że $B$
jest zwarty.

\textbf{Samosprzężoność i dodatniość.} Dla $u=Bf$, $w=Bg$
możemy wybrać odpowiednio $v=w$ i $v=u$ w równaniach
wariacyjnych. Symetria $b$ daje
$(Bf,g)_H=(g,u)_H=b(w,u)=b(u,w)=(f,Bg)_H$.
Ponieważ $B$ jest ograniczony i określony na całym $H$,
tożsamość ta daje $D(B^*)=H$ oraz $B^*=B$.
Dalej $(Bf,f)_H=b(Bf,Bf)=\|Bf\|_V^2\ge0$.

\textbf{Jądro i gęstość obrazu.} Gdy $Bf=0$, jego równanie
wariacyjne daje $(f,v)_H=0$ dla wszystkich $v\in V$.
Ponieważ $C_c^\infty(\Omega)\subset V$ jest gęste w $H$,
to $f=0$. Wreszcie
$y\perp\operatorname{Ran}B$ wtedy i tylko wtedy, gdy
$(Bf,y)_H=(f,By)_H=0$ dla każdego $f\in H$, czyli $By=0$.
Zatem $(\operatorname{Ran}B)^\perp=\ker B=\{0\}$.
Twierdzenie o rzucie ortogonalnym na domkniętą podprzestrzeń
(lemat~\ref{lem:pde11-projection}) daje
$\overline{\operatorname{Ran}B}=H$.
\end{proof}


\begin{uwaga}[Z czego wynika zwartość rezolwenty?]
\label{rem:pde-v42-compact-chain}
Zależność logiczna w poprzednim dowodzie jest istotna:
koercywność formy i twierdzenie Laxa--Milgrama
(\ref{thm:pde11-LM}) zapewniają oszacowanie
$B:L^2(\Omega)\to H_0^1(\Omega)$, a dopiero
zwartość zanurzenia Rellicha--Kondrachova
(\ref{thm:pde10-Rellich-p}) zapewnia zwartość
$B:L^2(\Omega)\to L^2(\Omega)$. Samo istnienie
ograniczonej rezolwenty nie daje zwartości.
W szczególności ograniczoność obszaru i warunki brzegowe
odgrywają rolę przy późniejszym rozkładzie spektralnym.
\end{uwaga}

\begin{definicja}[Laplasjan Dirichleta w $L^2$]
\label{def:pde12-laplacian}
Operatorem Dirichleta $A=-\Delta_D$ nazywamy operator w $H$
o dziedzinie
\[
D(A)=\big\{u\in V:\ \exists g\in H\quad
\int_\Omega\nabla u\cdot\nabla v\,\mathrm d\lambda_d(x)
=(g,v)_H\ \text{dla każdego }v\in V\big\},
\]
przy czym $Au=g$. Element $g$ jest wyznaczony
jednoznacznie dzięki gęstości $V$ w $H$: różnica dwóch
reprezentantów jest ortogonalna do całego $V$.
Jeśli $u_1,u_2\in D(A)$ mają reprezentantów $g_1,g_2$,
to dla skalarów $s,t\in\R$ dwuliniowość całki pokazuje,
że $su_1+tu_2\in D(A)$ z reprezentantem $sg_1+tg_2$.
Zatem $D(A)$ jest podprzestrzenią liniową i $A$ jest liniowy.
Dla każdego $u\in D(A)$, podstawiając test
$v\in C_c^\infty(\Omega)$, dostajemy
$\langle-\Delta u,v\rangle=
\int_\Omega\nabla u\cdot\nabla v\,\mathrm d\lambda_d(x)
=(Au,v)_H$. Odwrotnie, jeśli $u\in V$ i dystrybucja
$-\Delta u$ jest reprezentowana przez $g\in L^2$, to ta
sama tożsamość dla testów rozszerza się na $V$ przez
gęstość w $H^1$ i nierówność Cauchy'ego--Schwarza.
Wobec tego
$D(A)=\{u\in H_0^1(\Omega):\Delta u\in L^2(\Omega)
\text{ w sensie dystrybucji}\}$.
Warunek $u\in H_0^1(\Omega)$ narzuca zerowy ślad
(twierdzenie~\ref{thm:pde10-trace-zero}).
\end{definicja}

\begin{twierdzenie}[Samosprzężoność laplasjanu Dirichleta]
\label{thm:pde12-A-selfadjoint}
Operator $A$ z definicji~\ref{def:pde12-laplacian} jest
gęsto określony, nieujemny i samosprzężony.
Ponadto $D(A)=\operatorname{Ran}B$, a na tej dziedzinie
$A=B^{-1}-I$, więc $(A+I)^{-1}=B$.
\end{twierdzenie}
\begin{proof}
Jeśli $u=Bf$, to
$\int_\Omega\nabla u\cdot\nabla v\,\mathrm d\lambda_d(x)
=(f-u,v)_H$, czyli $u\in D(A)$ i $Au=f-u$.
Odwrotnie, jeśli $u\in D(A)$ i $Au=g$, to
$b(u,v)=(g+u,v)_H$, a z jednoznaczności równania
dla $B$ wynika $u=B(g+u)$. Tak otrzymujemy
$D(A)=\operatorname{Ran}B$, a więc gęstość dziedziny
z twierdzenia~\ref{thm:pde12-resolvent}.
W szczególności $B:H\to D(A)$ jest bijekcją,
$(A+I)Bf=f$ dla $f\in H$ i $B(A+I)u=u$ dla $u\in D(A)$.
Oznacza to $(A+I)^{-1}=B$ i $A=B^{-1}-I$ na $D(A)$.
Wybór $v=u$ w definicji daje
$(Au,u)_H=\|\nabla u\|_2^2\ge0$.

Aby wykazać samosprzężoność bez ukrywania dziedziny,
zauważmy, że $A$ jest symetryczny:
$(Au,v)_H=(u,Av)_H$ dla $u,v\in D(A)$
przez symetrię całki z gradientów. Jeśli $w\in D(A^*)$,
to dla każdego $u=Bf\in D(A)$ mamy
$(Au,w)_H=(u,A^*w)_H$, a więc
$(f-Bf,w)_H=(Bf,A^*w)_H=(f,BA^*w)_H$.
Samosprzężoność $B$ daje
$(f,w)_H=(f,Bw+BA^*w)_H$ dla wszystkich $f\in H$.
Zatem $w=B(w+A^*w)\in\operatorname{Ran}B=D(A)$.
Otrzymaliśmy $D(A^*)\subset D(A)$; inkluzja przeciwna
wynika z symetrii. Tym samym $A=A^*$. Operator jest również domknięty
jako operator sprzężony, zgodnie z definicją
~\ref{def:pde12-adjoint}. Wprost: gdy $u_n\in D(A)$,
$u_n\to u$ i $Au_n\to g$ w $H$, ciągłość $B$ daje
$u_n=B(Au_n+u_n)\to B(g+u)$, więc
$u=B(g+u)\in D(A)$ oraz $Au=g$.
\end{proof}

\begin{uwaga}[Dziedzina nie zawsze jest $H^2\cap H_0^1$]
\label{rem:pde12-corners}
Jeśli brzeg $\partial\Omega$ jest klasy $C^{1,1}$, to
$D(A)=H^2(\Omega)\cap H_0^1(\Omega)$: inkluzja z prawej
strony wynika z dystrybucyjnej charakterystyki $D(A)$,
a odwrotna z globalnej regularności eliptycznej
(twierdzenie~\ref{thm:pde15-boundary} dla macierzy jednostkowej).
Nie korzystamy z regularności $H^2$ w dowodzie samosprzężoności.

Dla ogólnego obszaru Lipschitzowskiego równość może zawieść.
Niech $\Omega=\{(r\cos\theta,r\sin\theta):0<r<1,\
0<\theta<\omega\}$, gdzie $\pi<\omega<2\pi$, i
$\beta=\pi/\omega\in(1/2,1)$. Jest to obszar Lipschitzowski:
przy wierzchołku wybierzmy oś $y$ wzdłuż dwusiecznej
dopełniającego sektora o kącie $\gamma=2\pi-\omega\in(0,\pi)$.
W tych współrzędnych, w dostatecznie małym otoczeniu zera,
obszar leży poniżej wykresu $y=|x|\cot(\gamma/2)$,
który jest Lipschitzowski. Poza wierzchołkiem brzeg
składa się z gładkich łuków i odcinków, a w punktach
ich przecięcia kąt wewnętrzny wynosi $\pi/2$, więc
po obrocie także tam otrzymujemy lokalny wykres
Lipschitzowski.
Wybierzmy gładkie odcięcie radialne $\eta$, równe $1$
dla $0\le r\le1/3$ i $0$ dla $r\ge2/3$, oraz połóżmy
$w(r,\theta)=r^\beta\sin(\beta\theta)$ i $u=\eta w$.
Dla $r>0$ mamy
\[
 \begin{aligned}
 \Delta w&=\partial_r^2 w+r^{-1}\partial_r w+
 r^{-2}\partial_\theta^2 w\\
 &= [\beta(\beta-1)+\beta-\beta^2]r^{\beta-2}
 \sin(\beta\theta)=0,\\
 |\nabla w|^2&=|\partial_r w|^2+r^{-2}|\partial_\theta w|^2
 =\beta^2r^{2\beta-2}.
 \end{aligned}
\]
W pobliżu wierzchołka zarówno $u$, jak i $\nabla u$
należą do $L^2$, ponieważ w biegunowych współrzędnych
ich kwadraty całkują się z czynnikiem $r$ i wykładniki
$2\beta+1$, $2\beta-1$ są większe od $-1$.
Z dala od wierzchołka funkcja jest gładka z ograniczonymi
pochodnymi. Pochodne klasyczne we wnętrzu są słabymi
pochodnymi, więc $u\in H^1(\Omega)$.

Aby uzasadnić zerowy ślad, wybierzmy radialne gładkie
odcięcie $\eta_\varepsilon$ znikające dla $r\le\varepsilon$,
równe $1$ dla $r\ge2\varepsilon$ i spełniające
$|\eta_\varepsilon'|\le C/\varepsilon$ (przez przeskalowanie
ustalonej funkcji). Funkcje $u_\varepsilon=\eta_\varepsilon u$
są gładkie aż do brzegu i znikają na nim, więc mają
zerowy ślad. Zbieżność $u_\varepsilon\to u$ w $H^1$
wynika z całkowalności $u,\nabla u$ oraz oszacowania
\[
 \|u\nabla\eta_\varepsilon\|_2^2
 \le C_1\varepsilon^{-2}\int_\varepsilon^{2\varepsilon}
 r^{2\beta+1}\,\mathrm d\lambda_1(r)
 \le C_2\varepsilon^{2\beta}\longrightarrow0.
\]
Z ciągłości śladu i twierdzenia~\ref{thm:pde10-trace-zero}
otrzymujemy $u\in H_0^1(\Omega)$.
Ponieważ $\Delta u=2\nabla\eta\cdot\nabla w+w\Delta\eta$
ma nośnik z dala od wierzchołka i jest ograniczona,
$\Delta u\in L^2(\Omega)$. Nie występuje dodatkowy składnik
dystrybucyjny w wierzchołku, który nie należy do $\Omega$.
Z definicji~\ref{def:pde12-laplacian} mamy $u\in D(A)$.

Tymczasem dla $r<1/3$ zachodzi
$\partial_r^2u=\beta(\beta-1)r^{\beta-2}\sin(\beta\theta)$.
Na podsektorze $\omega/4<\theta<3\omega/4$ moduł sinusa
jest nie mniejszy niż $1/\sqrt2$. Ponieważ dla jednostkowego
wektora radialnego $e_r$ mamy
$\partial_r^2u=e_r^T(\nabla^2u)e_r$,
moduł tej pochodnej nie przekracza normy Frobeniusa
macierzy drugich pochodnych. Wobec tego
\[
 \int_{\Omega\cap B(0,1/3)}|\nabla^2u|^2\,\mathrm d\lambda_2
 \ge c\int_0^{1/3}r^{2\beta-3}\,\mathrm d\lambda_1(r)
 =+\infty,
\]
gdyż $\beta<1$ i $c>0$. Gdyby $u\in H^2$, jej słabe
drugie pochodne zgadzałyby się we wnętrzu z klasycznymi,
wbrew tej rozbieżności. Zatem $u\notin H^2(\Omega)$.
Szczególny przypadek $\omega=3\pi/2$ powróci w
kontrprzykładzie~\ref{ex:pde15-corner}.
\end{uwaga}

\section{Twierdzenie spektralne dla zwartego operatora dodatniego}
\label{sec:pde12-compact-spectrum}

\begin{twierdzenie}[Rozkład spektralny zwartego operatora dodatniego]
\label{thm:pde12-compact-spectral}
Niech $H$ będzie rzeczywistą lub zespoloną,
separowalną przestrzenią Hilberta, a $K:H\to H$
operatorem zwartym, samosprzężonym i nieujemnym.
Wówczas istnieje ortonormalny układ wektorów własnych
$(e_j)$ odpowiadających dodatnim wartościom własnym
$\mu_j$, uporządkowanym nierosnąco i z uwzględnieniem
krotności. Układ może być pusty lub skończony.
Każda dodatnia wartość własna ma skończoną krotność.
Jeśli układ jest nieskończony, to $\mu_j\to0$. Ponadto
$H=\ker K\oplus\overline{\operatorname{span}}\{e_j\}$,
a dla $x\in H$ zachodzi
$Kx=\sum_j\mu_j(x,e_j)e_j$ zbieżnie w $H$.
Jeżeli $K$ jest różnowartościowy, $(e_j)$ jest bazą
ortonormalną całego $H$.
\end{twierdzenie}
\begin{proof}
Jeżeli $K=0$ (w szczególności gdy $H=\{0\}$), wybieramy
pusty układ, a wszystkie twierdzenia wynikają z $\ker K=H$.
Zakładamy dalej $K\ne0$.

\textbf{Krok 1: maksimum ilorazu Rayleigha.}
Połóżmy $m=\sup_{\|x\|=1}(Kx,x)$.
Pokażemy, że $m=\|K\|>0$. Z nierówności
Cauchy'ego--Schwarza w $H$ mamy
$0\le (Kx,x)\le\|Kx\|\|x\|\le\|K\|$ dla $\|x\|=1$,
więc $0\le m\le\|K\|$.
Dodatniość daje nierówność Cauchy'ego--Schwarza
dla formy $(Kx,y)$:
$|(Kx,y)|^2\le(Kx,x)(Ky,y)$.
Istotnie, dla $q(x,y)=(Kx,y)$ forma $q$ jest hermitowska
(w przypadku rzeczywistym symetryczna).
Niech $a=q(x,y)$ i $b=q(y,y)\ge0$.
Dla dowolnego skalara $t$ mamy
$q(x-ty,x-ty)=q(x,x)-2\operatorname{Re}(\overline t a)
+|t|^2b\ge0$. Gdy $b>0$, wybór $t=a/b$ daje żądaną
nierówność; gdy $b=0$, wybór $t=sa$ dla dowolnego
$s>0$ wymusza $a=0$.
Dla wektorów jednostkowych $x,y$ otrzymujemy
$|(Kx,y)|\le m$; dla $Kx\ne0$ wybieramy
$y=Kx/\|Kx\|$, otrzymując $\|Kx\|\le m$.
Dla $Kx=0$ ta nierówność również zachodzi.
Biorąc supremum po jednostkowych $x$, dostajemy
$\|K\|\le m$. Zatem $m=\|K\|>0$, ponieważ $K\ne0$.
Wybierzmy $\|x_n\|=1$ z $(Kx_n,x_n)\to m$.
Z ograniczoności i słabej zwartości ciągowej kul
(lemat~\ref{lem:pde11-weakcompact} w przypadku rzeczywistym;
dla zespolonego stosujemy go do przestrzeni rzeczywistej
z iloczynem $\operatorname{Re}(\cdot,\cdot)$: testowanie
wektorami $y$ i $iy$ daje zbieżność odpowiednio części
rzeczywistych i urojonych iloczynów skalarnych)
wybieramy podciąg $x_n\rightharpoonup x$,
zachowując oznaczenie indeksów.
Mamy $\|x\|\le1$, ponieważ dla $x\ne0$
$\|x\|^2=\lim_n |(x_n,x)|\le\|x\|$,
a dla $x=0$ nierówność jest spełniona.
Zwarty operator przenosi ten ciąg w ciąg silnie zbieżny:
gdyby $Kx_n\not\to Kx$, pewien podciąg miałby odległość
od $Kx$ co najmniej $\varepsilon>0$. Z jego zwartości
wybieramy dalszy podciąg zbieżny w normie do $z$.
Dla każdego $y\in H$ zachodzi
$(z,y)=\lim_k(x_{n_k},Ky)=(x,Ky)=(Kx,y)$,
więc $z=Kx$, sprzeczność. Ponadto
\[
 |(Kx_n,x_n)-(Kx,x)|\le
 \|Kx_n-Kx\|\|x_n\|+|(Kx,x_n-x)|\longrightarrow0.
\]
Stąd $(Kx,x)=m$.
Ponieważ $(Kx,x)\le m\|x\|^2$ i $\|x\|\le1$,
otrzymujemy $\|x\|=1$.

\textbf{Krok 2: maksimum daje wektor własny.}
Dla $y\perp x$ rozpatrujemy
$\varphi(t)=\frac{(K(x+ty),x+ty)}{\|x+ty\|^2}$,
gdzie $t\in\R$. Ortogonalność daje mianownik
$1+t^2\|y\|^2$, zaś samosprzężoność daje licznik
$m+2t\operatorname{Re}(Kx,y)+t^2(Ky,y)$.
Funkcja osiąga maksimum przy $t=0$, więc
$0=\varphi'(0)=2\operatorname{Re}(Kx,y)$.
W przypadku zespolonym powtarzamy argument dla $iy$:
$\operatorname{Re}(Kx,iy)=\operatorname{Re}(-i(Kx,y))
=\operatorname{Im}(Kx,y)$, więc także część urojona
jest zerowa. Zatem $Kx-mx$ jest ortogonalny do
$x^\perp$ oraz do $x$, ponieważ $(Kx,x)=m$ i $\|x\|=1$.
W konsekwencji $Kx-mx=0$.

\textbf{Krok 3: iteracja na ortogonalnym dopełnieniu.}
Przestrzeń $x^\perp$ jest niezmiennicza względem $K$,
bo $(Ky,x)=(y,Kx)=m(y,x)=0$. Ograniczenie $K$ do niej
pozostaje zwarte, samosprzężone i dodatnie. Powtarzamy
konstrukcję, dopóki ograniczenie nie stanie się zerowe.
Ściślej, po wyborze $e_1,\ldots,e_j$ rozpatrujemy
$M_j=\operatorname{span}\{e_1,\ldots,e_j\}^\perp$.
Jest to domknięta podprzestrzeń niezmiennicza, gdyż
dla $y\in M_j$ mamy $(Ky,e_\ell)=\mu_\ell(y,e_\ell)=0$
dla $\ell=1,\ldots,j$. Ograniczenie $K|_{M_j}$
jest zwarte (granice podciągów leżą w domkniętym $M_j$),
nieujemne i samosprzężone jako operator na całym $M_j$.
Jeśli $K|_{M_j}\ne0$, stosujemy kroki 1--2 w $M_j$
i otrzymujemy $\mu_{j+1}=\|K|_{M_j}\|>0$ i jednostkowy
$e_{j+1}\in M_j$. Ponieważ $M_j\subset M_{j-1}$,
wartości te są nierosnące, a wektory ortonormalne.
W wymiarze skończonym konstrukcja kończy się najpóźniej
po $\dim H$ krokach.
Jeśli układ jest nieskończony i $\mu_j$ nie dążą do zera,
istnieje $c>0$ i nieskończenie wiele indeksów, dla których
$\|Ke_j\|=\mu_j\ge c$. Dla dwóch różnych takich indeksów
$\|Ke_j-Ke_k\|^2=\mu_j^2+\mu_k^2\ge2c^2$.
Zatem obrazy odpowiedniego podciągu jednostkowych $e_j$
nie mają podciągu Cauchy'ego, co przeczy zwartości $K$.

\textbf{Krok 4: zupełność na $(\ker K)^\perp$.}
Jeżeli $z\perp e_j$ dla wszystkich $j$, to $z$
należy do kolejnych ortogonalnych dopełnień, na których
norma ograniczenia $K$ jest co najwyżej $\mu_{j+1}$.
Zatem $\|Kz\|\le\mu_{j+1}\|z\|\to0$;
gdy konstrukcja skończyła się po skończonej liczbie kroków,
operator na ostatnim dopełnieniu jest z definicji zerowy.
W obu przypadkach $Kz=0$. Odwrotnie, jeśli $Kz=0$,
to $\mu_j(z,e_j)=(z,Ke_j)=(Kz,e_j)=0$, więc $z$
jest ortogonalny do wszystkich $e_j$. Otrzymujemy
$\ker K=(\overline{\operatorname{span}}\{e_j\})^\perp$,
co daje rozkład ortogonalny. Dla $x\in H$ rozwinięcie
$x=x_0+\sum_j(x,e_j)e_j$ z $x_0\in\ker K$ jest
zbieżne w $H$; stosując ograniczony operator $K$
do sum częściowych, uzyskujemy żądany wzór.
Jeśli $Ky=\mu y$ dla $\mu>0$, zapiszmy
$y=y_0+\sum_j(y,e_j)e_j$, gdzie $y_0\in\ker K$.
Rzut równania na $\ker K$ daje $0=\mu y_0$, więc $y_0=0$.
Sparowanie równania z $e_j$ daje
$(\mu_j-\mu)(y,e_j)=0$ dla wszystkich $j$.
Zatem przestrzeń własna jest rozpięta przez dokładnie
te $e_j$, dla których $\mu_j=\mu$; jest skończeniewymiarowa,
gdyż $\mu_j\to0$ w przypadku nieskończonego układu.
Przy $\ker K=\{0\}$ układ jest bazą całego $H$.
\end{proof}

\begin{twierdzenie}[Twierdzenie spektralne dla zwartego operatora samosprzężonego]
\label{thm:pde12-compact-general-spectral}
Niech $H$ będzie rzeczywistą lub zespoloną separowalną
przestrzenią Hilberta, a $K\in\mathcal L(H)$ operatorem
zwartym i samosprzężonym, niekoniecznie dodatnim.
Wówczas $H$ jest domkniętą sumą ortogonalną $\ker K$ oraz
podprzestrzeni własnych odpowiadających niezerowym
rzeczywistym wartościom własnym $\kappa_j$.
Każda niezerowa wartość własna ma skończoną krotność,
a jeśli jest ich nieskończenie wiele (licząc z krotnościami),
to $|\kappa_j|\to0$. Istnieje ortonormalny układ
$(e_j)$ tych wektorów własnych taki, że
\[
 Kx=\sum_j\kappa_j(x,e_j)_He_j\quad(x\in H),
\]
gdzie szereg jest zbieżny w $H$; jeżeli $\ker K=\{0\}$,
to $(e_j)$ jest bazą ortonormalną całego $H$.
\end{twierdzenie}
\begin{proof}
Operator $K^2$ jest zwarty, samosprzężony i nieujemny:
$K^2$ jest zwarty jako złożenie operatora zwartego
i ograniczonego, a $(K^2x,x)=\|Kx\|^2\ge0$.
Ponadto $\ker K^2=\ker K$: jeśli $K^2x=0$, to
$\|Kx\|^2=(K^2x,x)=0$; druga inkluzja jest bezpośrednia.
Twierdzenie~\ref{thm:pde12-compact-spectral} zastosowane
do $K^2$ daje rozkład ortogonalny
$H=\ker K\oplus\overline{\bigoplus_{\mu>0}E_\mu}$,
gdzie $E_\mu=\ker(K^2-\mu I)$, $\mu>0$.
Każde $E_\mu$ jest skończeniewymiarowe, gdyż nieskończony
układ ortonormalny w $E_\mu$ miałby obrazy przez $K^2$
bez podciągu zbieżnego, co przeczyłoby zwartości.

Ponieważ $K^2(Kx)=K(K^2x)=\mu Kx$ dla $x\in E_\mu$,
każda przestrzeń $E_\mu$ jest niezmiennicza względem $K$.
Na $E_\mu$ operator $J=\mu^{-1/2}K|_{E_\mu}$
jest samosprzężony i spełnia $J^2=I_{E_\mu}$.
Każdy $x\in E_\mu$ ma rozkład
$x=x_++x_-$, gdzie $x_\pm=(x\pm Jx)/2$.
Bezpośrednio $Jx_\pm=\pm x_\pm$, więc
$Kx_\pm=\pm\sqrt\mu\,x_\pm$.
Podprzestrzenie odpowiadające obu znakom są ortogonalne:
dla $Jv=v$ i $Jw=-w$ mamy
$(v,w)=(Jv,w)=(v,Jw)=-(v,w)$.
W każdej z nich wybieramy bazę ortonormalną
(może być pusta). W ten sposób otrzymujemy bazę własną
całego $E_\mu$ bez odwołania do odrębnego twierdzenia
o diagonalizacji macierzy.
Bazy wszystkich $E_\mu$ tworzą wraz z bazą jądra
ortogonalny rozkład $H$ na wektory własne $K$.
Każda niezerowa wartość własna $\kappa$ operatora $K$
ma przestrzeń własną zawartą w skończeniewymiarowym
$E_{\kappa^2}$. Dla każdego $\varepsilon>0$ istnieje
tylko skończenie wiele wektorów w uzyskanej bazie własnej
z $|\kappa_j|\ge\varepsilon$, ponieważ odpowiadają
wartościom własnym $K^2$ nie mniejszym niż
$\varepsilon^2$, liczonym z krotnościami.
Jeśli układ niezerowych modów jest nieskończony,
to $|\kappa_j|\to0$.
Dla sum skończonych podany wzór na $K$ wynika
z diagonalizacji. Dla dowolnego $x\in H$ przechodzimy
do granicy, ponieważ rzuty na skończone sumy podprzestrzeni
własnych zbiegają silnie do rzutu na $(\ker K)^\perp$,
a $K$ jest ograniczony i znika na swoim jądrze.
\end{proof}

\begin{wniosek}[Dyskretne widmo laplasjanu Dirichleta]
\label{cor:pde12-eigen}
Dla obszaru ograniczonego Lipschitzowskiego $\Omega$ operator
$A=-\Delta_D$ ma bazę ortonormalną $(\phi_j)_{j\ge1}$ w
$L^2(\Omega)$ złożoną z wektorów własnych:
$A\phi_j=\lambda_j\phi_j$, gdzie
$0<\lambda_1\le\lambda_2\le\cdots$, $\lambda_j\to\infty$.
Każda wartość własna ma skończoną krotność,
a w sensie słabym
$\int_\Omega\nabla\phi_j\cdot\nabla v\,\mathrm d\lambda_d(x)
=\lambda_j\int_\Omega\phi_jv\,\mathrm d\lambda_d(x)$
dla wszystkich $v\in H_0^1(\Omega)$.
\end{wniosek}
\begin{proof}
Dla operatora $B=(A+I)^{-1}$ z
\ref{thm:pde12-resolvent} jądro jest zerowe,
więc twierdzenie~\ref{thm:pde12-compact-spectral}
daje bazę $(\phi_j)$ i liczby $\mu_j>0$ z
$B\phi_j=\mu_j\phi_j$.
Przestrzeń $L^2(\Omega)$ jest separowalna i
nieskończeniewymiarowa: w kuli zawartej w niepustym
$\Omega$ można wybrać nieskończenie wiele parami
rozłącznych mniejszych kul i znormalizować ich funkcje
charakterystyczne, uzyskując układ ortonormalny.
Ponieważ $B$ jest różnowartościowy, baza nie może
być skończona, a $\mu_j\downarrow0$ (licząc krotności).
Ponadto $\phi_j=\mu_j^{-1}B\phi_j\in\operatorname{Ran}B=D(A)$.
Z równości $A=B^{-1}-I$ otrzymujemy
$A\phi_j=(\mu_j^{-1}-1)\phi_j$, czyli
$\lambda_j=\mu_j^{-1}-1$. Nieujemność $A$ daje
$\lambda_j\ge0$. Gdyby $\lambda_j=0$, to
$\|\nabla\phi_j\|_2^2=(A\phi_j,\phi_j)=0$.
Nierówność Poincarégo w $H_0^1(\Omega)$
(wniosek~\ref{cor:pde10-Poincare} dla $p=2$) daje
$\phi_j=0$, sprzeczność z normalizacją.
Zatem $\lambda_j>0$, a $\lambda_j\to+\infty$,
ponieważ $\mu_j\to0$. Dla $\lambda\ge0$ równanie
$Au=\lambda u$ z $u\ne0$ jest równoważne
$Bu=(1+\lambda)^{-1}u$: w drugim kierunku
$u=(1+\lambda)Bu\in D(A)$. Zatem krotności są równe
krotnościom odpowiednich dodatnich wartości własnych
$B$, skończonym na mocy twierdzenia spektralnego.
Tożsamość słaba wynika z definicji $A$.
Brak innych punktów widma $A$ niż wartości własne
wykażemy w twierdzeniu~\ref{thm:pde12-fredholm}.
\end{proof}

\section{Dziedzina formy, zasada min--max i oszacowania widma}
\label{sec:pde12-minmax}

\begin{twierdzenie}[Rozwinięcie spektralne formy energii]
\label{thm:pde12-formdomain}
Niech $(\lambda_j,\phi_j)$ będą wartościami i funkcjami
własnymi z wniosku~\ref{cor:pde12-eigen}. Dla $u\in L^2(\Omega)$
niech $c_j=(u,\phi_j)_{L^2}$. Wówczas
\begin{align*}
u\in H_0^1(\Omega)
&\quad\Longleftrightarrow\quad
\sum_{j=1}^\infty\lambda_j|c_j|^2<\infty,\\
u\in D(A)
&\quad\Longleftrightarrow\quad
\sum_{j=1}^\infty\lambda_j^2|c_j|^2<\infty.
\end{align*}
Dla $u\in H_0^1(\Omega)$ zachodzi
$\|\nabla u\|_2^2=\sum_j\lambda_j|c_j|^2$,
a dla $u\in D(A)$ zachodzi
$Au=\sum_j\lambda_jc_j\phi_j$ zbieżnie w $L^2$.
\end{twierdzenie}
\begin{proof}
Piszemy $\mu_j=(1+\lambda_j)^{-1}$.
Ponieważ $\phi_j=\mu_j^{-1}B\phi_j$, mamy
$\phi_j\in H_0^1(\Omega)$. Równanie wariacyjne dla $B$
daje $b(\phi_i,v)=\mu_i^{-1}(\phi_i,v)_{L^2}$
dla każdego $v\in V$. W szczególności
$b(\sqrt{\mu_i}\phi_i,\sqrt{\mu_j}\phi_j)
=\sqrt{\mu_j/\mu_i}\,\delta_{ij}=\delta_{ij}$.
Wektory $\sqrt{\mu_j}\phi_j$ tworzą więc układ
ortonormalny w $V=H_0^1(\Omega)$ z iloczynem skalarnym $b$. Układ ten jest zupełny w $V$:
jeśli $v\in V$ jest do niego ortogonalny, to
$0=b(v,\phi_j)=\mu_j^{-1}(v,\phi_j)_{L^2}$
dla wszystkich $j$; zupełność $(\phi_j)$ w $L^2$ daje $v=0$.
Z twierdzenia Parsevala w $V$ dostajemy
\[
\|v\|_V^2=\sum_j|b(v,\sqrt{\mu_j}\phi_j)|^2
=\sum_j\mu_j^{-1}|(v,\phi_j)_{L^2}|^2
=\sum_j(1+\lambda_j)|(v,\phi_j)_{L^2}|^2.
\]
Dla ogólnego $u\in L^2$ z Parsevala w $L^2$ mamy
$\sum_j|c_j|^2=\|u\|_2^2<\infty$, więc warunek
$\sum_j\lambda_j|c_j|^2<\infty$ jest równoważny
sumowalności z wagami $1+\lambda_j$.
Dla $u_N=\sum_{j\le N}c_j\phi_j$ i $M>N$ ortogonalność
w $V$ daje
$\|u_M-u_N\|_V^2=\sum_{j=N+1}^M(1+\lambda_j)|c_j|^2$.
Stąd sumy częściowe są ciągiem Cauchy'ego w $V$. Ich granica $v\in V$
równa się $u$, gdyż zbieżność w $V$ implikuje zbieżność
w $L^2$, a w $L^2$ mamy $u_N\to u$ z Parsevala.
Tak wykazaliśmy równoważność dla $H_0^1$.
Dla $u\in V$ stosujemy powyższą tożsamość Parsevala
z $v=u$ i odejmujemy $\|u\|_2^2=\sum_j|c_j|^2$,
otrzymując równość energii.
Ponadto dla $u,v\in V$, z tożsamości polaryzacyjnej
lub rozwinięcia w bazie $V$, wynika
\[
 \int_\Omega\nabla u\cdot\nabla v\,\mathrm d\lambda_d
 =\sum_j\lambda_j(u,\phi_j)_H(v,\phi_j)_H .
\]
Szereg jest bezwzględnie zbieżny z nierówności
Cauchy'ego--Schwarza dla ważonych współczynników.
W zespoleniu drugie czynniki po prawej stronie oraz
gradient $v$ po lewej zastępujemy ich sprzężeniami.

Jeśli $u\in D(A)$, to $Au\in L^2$, a dla każdego $j$
$(Au,\phi_j)=(u,A\phi_j)=\lambda_jc_j$.
Parseval daje sumowalność $\lambda_j^2|c_j|^2$ i wzór
na $Au$. Odwrotnie, jeśli ta suma jest skończona,
to $\sum_j\lambda_j|c_j|^2
\le\lambda_1^{-1}\sum_j\lambda_j^2|c_j|^2<\infty$
dzięki $\lambda_1>0$, więc $u\in H_0^1$ z pierwszej
części twierdzenia. Niech $g=\sum_j\lambda_jc_j\phi_j\in L^2$.
Dla $v=\phi_j$ symetria formy i równanie własne
z testem $u$ dają
$\int_\Omega\nabla u\cdot\nabla\phi_j\,\mathrm d\lambda_d
=\lambda_j(u,\phi_j)_{L^2}=(g,\phi_j)_{L^2}$.
Z liniowości równość zachodzi dla skończonych sum.
Takie sumy są gęste w $V$, ponieważ
$\sqrt{\mu_j}\phi_j$ tworzą bazę ortonormalną $V$.
Jeżeli $v_n\to v$ w $V$, to błędy obu stron nie przekraczają
$\|\nabla u\|_2\|\nabla(v_n-v)\|_2$ i
$\|g\|_2\|v_n-v\|_2$. Obie liczby dążą do zera.
Dostajemy równość dla każdego $v\in V$, więc
$u\in D(A)$ i $Au=g$.
\end{proof}

\begin{twierdzenie}[Zasada Couranta--Fischera]
\label{thm:pde12-minmax}
Przy założeniach wniosku~\ref{cor:pde12-eigen},
dla każdej liczby całkowitej $k\ge1$
\[
\lambda_k=
\min_{\substack{E\subset H_0^1(\Omega)\\\dim E=k}}
\ \max_{\substack{u\in E\\u\ne0}}
\frac{\displaystyle\int_\Omega|\nabla u|^2\,\mathrm d\lambda_d(x)}
{\displaystyle\int_\Omega|u|^2\,\mathrm d\lambda_d(x)}.
\]
Minimum jest osiągane dla $E=\operatorname{span}
\{\phi_1,\ldots,\phi_k\}$.
\end{twierdzenie}
\begin{proof}
Dla dowolnej $k$-wymiarowej podprzestrzeni $E\subset V$
iloraz Rayleigha można ograniczyć do sfery
$\{u\in E:\|u\|_2=1\}$. Jest ona zwarta w $E$,
ponieważ wszystkie normy w skończonym wymiarze są
równoważne. Funkcja $u\mapsto\|\nabla u\|_2^2$ jest
ciągła w normie $V$, więc maksimum na tej sferze
jest osiągane.
Dla wskazanego $E$ i $u=\sum_{j\le k}c_j\phi_j$
twierdzenie~\ref{thm:pde12-formdomain} daje
$\|\nabla u\|_2^2=\sum_{j\le k}\lambda_j|c_j|^2
\le\lambda_k\sum_{j\le k}|c_j|^2=\lambda_k\|u\|_2^2$.
Wektor $u=\phi_k$ daje równość, więc maksimum na tym $E$
wynosi $\lambda_k$. Z drugiej strony weźmy dowolne
$k$-wymiarowe $E\subset H_0^1$. Odwzorowanie liniowe
$E\to\R^{k-1}$, $u\mapsto((u,\phi_j)_H)_{j<k}$,
ma jądro wymiaru co najmniej $1$ na mocy twierdzenia
o wymiarze jądra i obrazu (dla $k=1$ przestrzeń docelowa
jest zerowa). Wybieramy $u\ne0$ z tego jądra.
Dla takiego $u$ w rozwinięciu $L^2$ znikają wszystkie
współczynniki o indeksach mniejszych od $k$, wobec czego
$\|\nabla u\|_2^2=\sum_{j\ge k}\lambda_j|c_j|^2
\ge\lambda_k\sum_{j\ge k}|c_j|^2=\lambda_k\|u\|_2^2$.
Maksimum na dowolnym $E$ jest zatem co najmniej $\lambda_k$.
Łącząc oba oszacowania, otrzymujemy wzór.
\end{proof}

\begin{wniosek}[Pierwsza wartość własna, Poincaré i monotoniczność]
\label{cor:pde12-first}
Jeśli $\Omega\subset\R^d$ jest obszarem ograniczonym Lipschitzowskim, to
$\lambda_1=\inf_{u\in H_0^1(\Omega)\setminus\{0\}}
\|\nabla u\|_2^2/\|u\|_2^2$.
Najlepsza stała $C_P$ w nierówności
$\|u\|_2\le C_P\|\nabla u\|_2$ wynosi
$C_P=\lambda_1^{-1/2}$. Ponadto jeżeli
$\Omega_1\subset\Omega_2$ są ograniczonymi obszarami
Lipschitzowskimi, to
$\lambda_k(\Omega_2)\le\lambda_k(\Omega_1)$ dla każdego $k$.
\end{wniosek}
\begin{proof}
Pierwsza równość wynika z twierdzenia
~\ref{thm:pde12-minmax} dla $k=1$, gdyż iloraz jest
stały na każdej jednowymiarowej podprzestrzeni
poza zerem. Daje ona $\|u\|_2\le
\lambda_1^{-1/2}\|\nabla u\|_2$ dla $u\ne0$,
a dla $u=0$ nierówność też zachodzi.
Wektor $\phi_1$ osiąga równość, więc żadna mniejsza
stała nie może wystarczyć.
Dla monotoniczności rozważmy przedłużenie przez zero:
funkcja z $C_c^\infty(\Omega_1)$ po przedłużeniu jest
w $C_c^\infty(\Omega_2)$ i zachowuje normy $L^2$
oraz gradientu. Dla $u\in H_0^1(\Omega_1)$ wybieramy
testy $u_n\to u$ w $H^1(\Omega_1)$. Ich przedłużenia
tworzą ciąg Cauchy'ego w $H_0^1(\Omega_2)$, a granica
w $L^2$ jest przedłużeniem $u$ przez zero. Przechodząc
do granicy w obu tożsamościach norm, otrzymujemy
liniowe izometryczne zanurzenie
$H_0^1(\Omega_1)\hookrightarrow H_0^1(\Omega_2)$,
zachowujące także osobno normę gradientu.
Każda $k$-wymiarowa podprzestrzeń pierwszej przestrzeni
staje się więc dopuszczalną podprzestrzenią drugiej,
z tym samym maksimum ilorazu Rayleigha. Przyjmując minima
otrzymujemy podaną nierówność.
\end{proof}

\begin{przyklad}[Odcinek i prostokąt]
Na odcinku $(0,L)$ rozwiązania problemu
$-u''=\lambda u$, $u(0)=u(L)=0$ mają postać
$\phi_n(x)=\sqrt{2/L}\sin(n\pi x/L)$ oraz
$\lambda_n=(n\pi/L)^2$, $n\ge1$.
Wynik dla klasycznego zagadnienia brzegowego wraz
z zupełnością bazy sinusowej udowodniono
w twierdzeniu~\ref{thm:pde6-spectrum}.
Funkcje $\phi_n$ należą do $D(A)$ na mocy charakterystyki
dystrybucyjnej i zerowego śladu. Ponieważ tworzą pełną
bazę $L^2(0,L)$, nie ma dodatkowych wartości własnych:
sparowanie równania $Au=\lambda u$ z $\phi_n$ daje
$(\lambda_n-\lambda)(u,\phi_n)_H=0$ dla każdego $n$.
Dla $u\ne0$ co najmniej jeden współczynnik jest niezerowy.
W prostokącie $(0,L_1)\times(0,L_2)$ funkcje
\[
 \phi_{m,n}(x,y)=\frac{2}{\sqrt{L_1L_2}}
 \sin(m\pi x/L_1)\sin(n\pi y/L_2),\qquad m,n\ge1,
\]
mają zerowy ślad i należą do $D(A)$, ponieważ są
gładkie na domknięciu prostokąta i ich laplasjany są
w $L^2$. Dwukrotne różniczkowanie w obu zmiennych daje
$\lambda_{m,n}=\pi^2(m^2/L_1^2+n^2/L_2^2)$.
Ortogonalność i normalizacja wynikają z twierdzenia
Fubiniego i jednowymiarowych tożsamości
$\int_0^{L_i}\sin(m\pi x/L_i)\sin(n\pi x/L_i)
\,\mathrm d\lambda_1(x)=(L_i/2)\delta_{mn}$.
Stąd iloczyn skalarny dwóch znormalizowanych funkcji
prostokąta wynosi $\delta_{mm'}\delta_{nn'}$. Aby sprawdzić zupełność,
niech $f\in L^2$ będzie ortogonalna do wszystkich
$\phi_{m,n}$. Dla każdego $m$ funkcja
$g_m(y)=\int_0^{L_1} f(x,y)\sin(m\pi x/L_1)\,
\mathrm d\lambda_1(x)$ należy do $L^2(0,L_2)$
z nierówności Cauchy'ego--Schwarza i twierdzenia Fubiniego.
Z ortogonalności $f$ wynika, że wszystkie współczynniki
sinusowe $g_m$ są zerowe, zatem $g_m=0$ prawie wszędzie.
Po odrzuceniu sumy przeliczalnie wielu zbiorów miary zero
dla prawie każdego $y$ wszystkie współczynniki funkcji
$f(\cdot,y)\in L^2(0,L_1)$ znikają. Zupełność bazy
jednowymiarowej daje $f(\cdot,y)=0$ prawie wszędzie,
a Fubini daje $f=0$ w $L^2$ prostokąta.
W konsekwencji wymienione funkcje tworzą pełną bazę własną.
Wartości te należy uporządkować rosnąco z krotnościami.
Na kwadracie o boku $L$ pierwsza wartość wynosi
$2\pi^2/L^2$ i jest prosta, a druga i trzecia są równe
$5\pi^2/L^2$ (pary $(1,2)$ i $(2,1)$).
Dalsze krotności mogą wynikać nie tylko z symetrii
$\lambda_{m,n}=\lambda_{n,m}$, lecz także z różnych
przedstawień liczby jako sumy dwóch dodatnich kwadratów.
\end{przyklad}

\section{Widmo, rezolwenta i alternatywa Fredholma}
\label{sec:pde12-fredholm}

\begin{definicja}[Zbiór rezolwenty i widmo]
\label{def:pde12-spectrum}
Dla domkniętego gęsto określonego operatora $T$ w
zespolonej przestrzeni Hilberta jego \emph{zbiorem rezolwenty}
$\rho(T)$ nazywamy zbiór liczb $z\in\C$, dla których
$T-zI:D(T)\to H$ jest bijekcją, a odwrotność
$(T-zI)^{-1}:H\to H$ jest ograniczona.
\emph{Widmo} to $\sigma(T)=\C\setminus\rho(T)$.
Liczba $z$ jest wartością własną, jeżeli
$0\ne u\in D(T)$ spełnia $Tu=zu$.
Jeżeli $T-zI$ jest bijekcją, ograniczoność odwrotności
wynika już z twierdzenia o wykresie domkniętym:
z $f_n\to f$ i $u_n=(T-zI)^{-1}f_n\to u$ wynika
$Tu_n=f_n+zu_n\to f+zu$. Domkniętość $T$ daje
$u\in D(T)$ i $(T-zI)u=f$, czyli $u=(T-zI)^{-1}f$.
Wykres odwrotności jest zatem domknięty w $H\times H$.
Jest ona liniowa i określona na całej przestrzeni Banacha
$H$, więc twierdzenie o wykresie domkniętym daje
jej ograniczoność. Ponadto dla $R_z=(T-zI)^{-1}$ mamy
$TR_zf=f+zR_zf$, skąd
\[
 \|R_zf\|_{D(T)}^2
 \le\bigl(\|R_z\|^2+(1+|z|\|R_z\|)^2\bigr)\|f\|^2.
\]
Zatem $R_z:H\to D(T)$ jest ograniczony również
w normie grafowej. Wyróżniamy ponadto trzy rozłączne
części widma: punktową (brak iniektywności $T-zI$),
ciągłą (iniektywność i gęsty, lecz niepełny obraz)
oraz residualną (iniektywność i obraz, który nie jest
gęsty). Wyczerpują one widmo, ponieważ bijekcja
$T-zI$ ma na mocy powyższego argumentu ograniczoną
odwrotność.
\end{definicja}

\begin{uwaga}[Tożsamość rezolwenty i otwartość zbioru rezolwenty]
Jeżeli $z\in\rho(T)$ i $|w-z|\|R_z\|<1$, to
na $D(T)$ mamy faktoryzację
$T-wI=(I-(w-z)R_z)(T-zI)$.
Szereg Neumanna dla operatora $(w-z)R_z$ zapewnia
odwracalność pierwszego czynnika w $\mathcal L(H)$.
Drugi czynnik jest bijekcją $D(T)\to H$.
Stąd $T-wI$ jest bijekcją i jego ograniczona odwrotność
ma postać $R_w=R_z(I-(w-z)R_z)^{-1}$, z obrazem
w $D(T)$. Zatem $w\in\rho(T)$. W szczególności
$\rho(T)$ jest otwarty, a $\sigma(T)$ domknięte.
Dla dowolnych $z,w\in\rho(T)$ mamy
\[
 R_z-R_w=R_z[(T-wI)-(T-zI)]R_w
 =(z-w)R_zR_w.
\]
Rachunek jest poprawny dziedzinowo, ponieważ
$R_w(H)\subset D(T)$, a $R_z(T-zI)$ jest identycznością
na $D(T)$. Zamieniając $z$ z $w$, otrzymujemy
także $(z-w)R_wR_z=R_z-R_w$, więc rezolwenty komutują
(dla $z=w$ jest to bezpośrednie).
Jeżeli $R_z$ jest zwarty dla jednego $z$,
to $R_w=R_z+(w-z)R_zR_w$ jest zwarty dla każdego
$w\in\rho(T)$: złożenie zwartego $R_z$ z ograniczonym
$R_w$ jest zwarte, a suma operatorów zwartych
jest zwarta.
\end{uwaga}

\begin{twierdzenie}[Alternatywa Fredholma dla laplasjanu]
\label{thm:pde12-fredholm}
Niech $A=-\Delta_D$ będzie operatorem na obszarze ograniczonym Lipschitzowskim $\Omega$, a $\zeta\in\C$.
Jeżeli $\zeta\notin\{\lambda_j:j\ge1\}$, to dla
każdego $f\in L^2(\Omega;\C)$ istnieje dokładnie jedno
$u\in D(A)$ spełniające $(A-\zeta I)u=f$, i
\[
\|u\|_2\le\frac{\|f\|_2}
{\operatorname{dist}(\zeta,\{\lambda_j\})}.
\]
Jeżeli $\zeta$ jest wartością własną, to równanie
$(A-\zeta I)u=f$ jest rozwiązywalne wtedy i tylko wtedy,
gdy $f\perp\ker(A-\zeta I)$; rozwiązanie jest wówczas
jednoznaczne z dokładnością do dodania elementu jądra.
Rezolwenta $(A-\zeta I)^{-1}$ w każdym punkcie
$\zeta\in\rho(A)$ jest zwarta, a
$\sigma(A)=\{\lambda_j:j\ge1\}$ jako zbiór (bez krotności).
\end{twierdzenie}
\begin{proof}
Przechodzimy do zespolenia przestrzeni i operatora:
$D(A_\C)=D(A)+iD(A)$ oraz
$A_\C(v+iw)=Av+iAw$. Zachowujemy oznaczenie $A$.
Rzeczywista baza własna pozostaje bazą ortonormalną
po zespoleniu, gdyż rozwijamy oddzielnie część rzeczywistą
i urojoną. Charakterystyki z twierdzenia
~\ref{thm:pde12-formdomain} obowiązują z modułami
współczynników zespolonych przez dodanie odpowiednich
tożsamości dla obu części. W szczególności operator
jest domknięty także po zespoleniu (stosujemy kryterium
ciągowe osobno do obu części).
Operator zespolony jest także samosprzężony.
Istotnie, jest symetryczny przez rozwinięcie zespolonego
iloczynu skalarnego i rzeczywistą symetrię $A$.
Jeśli $v+iw\in D(A_\C^*)$, $v,w\in H$ rzeczywistym,
i $A_\C^*(v+iw)=p+iq$, to testowanie definicji
sprzężenia rzeczywistymi $h\in D(A)$ daje
$(Ah,v)=(h,p)$ oraz $(Ah,w)=(h,q)$.
Rzeczywista samosprzężoność implikuje $v,w\in D(A)$
oraz $Av=p$, $Aw=q$. Wraz z symetrią daje to równość
dziedzin i operatorów.
Piszemy $f=\sum_jf_j\phi_j$, a dla potencjalnego
rozwiązania $u=\sum_ju_j\phi_j$, gdzie
$f_j=(f,\phi_j)_H$, $u_j=(u,\phi_j)_H$.
Z charakterystyki $D(A)$ wynika, że równanie jest
równoważne warunkom $(\lambda_j-\zeta)u_j=f_j$
dla wszystkich $j$, o ile $u\in D(A)$.
Jeżeli $\zeta$ nie jest wartością własną, to odległość
$\delta=\inf_j|\lambda_j-\zeta|$ jest dodatnia:
$\lambda_j\to\infty$, więc istnieje $N$ takie, że
$|\lambda_j-\zeta|>|\lambda_1-\zeta|+1$ dla $j>N$.
Minimum odległości jest zatem osiągane wśród pierwszych
$N$ wyrazów i jest dodatnie. Definiujemy $u_j=f_j/(\lambda_j-\zeta)$.
Z Parsevala
$\|u\|_2^2=\sum_j|f_j|^2/|\lambda_j-\zeta|^2
\le\delta^{-2}\|f\|_2^2$.
Dodatkowo z nierówności trójkąta
$\lambda_j\le|\lambda_j-\zeta|+|\zeta|$ otrzymujemy
jednostajne oszacowanie
$\lambda_j/|\lambda_j-\zeta|\le1+|\zeta|/\delta$.
W szczególności $\sum_j|u_j|^2<\infty$ pozwala
zdefiniować $u$ jako granicę szeregu ortogonalnego. Stąd $\sum\lambda_j^2|u_j|^2<\infty$,
więc $u\in D(A)$ na mocy~\ref{thm:pde12-formdomain},
i rzeczywiście rozwiązuje równanie. Każde inne rozwiązanie
ma te same współczynniki, co dowodzi jednoznaczności.
Oszacowanie normy rezolwenty jest ostre: wybierając
$f=\phi_{j_0}$ dla indeksu realizującego $\delta$,
otrzymujemy $\|f\|_2=1$ oraz $\|R_\zeta f\|_2=\delta^{-1}$.

Jeśli $\zeta$ jest wartością własną, to dla indeksów
$J=\{j:\lambda_j=\zeta\}$ równanie daje warunek
$f_j=0$ dla każdego $j\in J$.
Z charakterystyki spektralnej dziedziny $A$ wynika
$\ker(A-\zeta I)=\operatorname{span}\{\phi_j:j\in J\}$:
wektory z jądra mają zerowe współczynniki poza $J$,
a każda taka suma skończona należy do $D(A)$ i spełnia
równanie własne. Warunek na $f_j$ jest zatem dokładnie
ortogonalnością $f$ do tego jądra.
Dla $j\notin J$ możemy określić
$u_j=f_j/(\lambda_j-\zeta)$, a współczynniki $u_j$
z $j\in J$ pozostają dowolne. Zbiór $J$ jest skończony
na mocy skończonej krotności, a jego dopełnienie
jest nieskończone. Na tym dopełnieniu
$\delta'=\min_{j\notin J}|\lambda_j-\zeta|>0$
z tego samego argumentu co poprzednio, rozpoczętego
od dowolnego indeksu spoza $J$.
Mamy
\[
 \sum_{j\notin J}\lambda_j^2|u_j|^2
 \le(1+|\zeta|/\delta')^2\sum_{j\notin J}|f_j|^2.
\]
Pozostałe składniki są skończoną sumą, więc
otrzymujemy $u\in D(A)$ i równanie dla wszystkich
współczynników. Wybór $u_j=0$ dla $j\in J$ daje jedyne
rozwiązanie $u_0$ ortogonalne do jądra, z oszacowaniem
$\|u_0\|_2\le(\delta')^{-1}\|f\|_2$.
Wszystkie rozwiązania mają postać $u_0+w$,
$w\in\ker(A-\zeta I)$; z ortogonalności
$\|u_0+w\|_2^2=\|u_0\|_2^2+\|w\|_2^2$.
Zatem $u_0$ jest jedynym rozwiązaniem o najmniejszej
normie $L^2$.

Dla $\zeta$ poza zbiorem wartości własnych definiujemy
na całym $H$ operator skończonego rzędu
$R_{\zeta,N}f=\sum_{j\le N} f_j(\lambda_j-\zeta)^{-1}\phi_j$.
Z Parsevala
\[
 \|(R_\zeta-R_{\zeta,N})f\|_2^2
 \le \sup_{j>N}|\lambda_j-\zeta|^{-2}\|f\|_2^2.
\]
Supremum dąży do zera, bo $\lambda_j\to\infty$.
Każdy $R_{\zeta,N}$ jest zwarty jako ograniczony operator
skończonego rzędu. Dla dowolnego $\varepsilon>0$ wybieramy
$N$ tak, aby norma różnicy była mniejsza niż $\varepsilon/2$.
Obraz kuli jednostkowej przez $R_{\zeta,N}$ ma skończoną
$(\varepsilon/2)$-sieć, która jest $\varepsilon$-siecią
obrazu przez $R_\zeta$. Z całkowitej ograniczoności
i zupełności $H$ wynika zwartość $R_\zeta$.
Każda $\zeta$ spoza zbioru wartości własnych należy do
$\rho(A)$, a każda wartość własna należy do $\sigma(A)$
z powodu braku iniektywności. To daje równość zbiorów.
\end{proof}

Ogólny mechanizm alternatywy Fredholma dla operatorów
zwartych omówiono również w \atlasref{FA-30::thm:fredholm-hilbert30}{FA-30: alternatywa Fredholma dla operatora zwartego}.

\begin{twierdzenie}[Alternatywa Fredholma dla operatorów zwartych]
\label{thm:pde12-general-fredholm}
Niech $H$ będzie zespoloną przestrzenią Hilberta,
a $K:H\to H$ operatorem zwartym. Wówczas
$T=I-K$ ma domknięty obraz oraz skończeniewymiarowe
jądra $\ker T$ i $\ker T^*$. Zachodzi
$\dim\ker T=\dim\ker T^*$, a dla $f\in H$
równanie $(I-K)u=f$ ma rozwiązanie wtedy i tylko wtedy,
gdy $f\perp\ker(I-K^*)$. Rozwiązanie jest jednoznaczne
z dokładnością do dodania elementu $\ker(I-K)$.
\end{twierdzenie}
\begin{proof}
\textbf{Krok 1: skończony wymiar jądra.}
Jeśli $\ker T$ byłoby nieskończeniewymiarowe,
wybralibyśmy w nim nieskończony układ ortonormalny
$(e_n)$. Równość $Te_n=0$ oznacza $Ke_n=e_n$.
Ponieważ $(e_n)$ nie ma podciągu zbieżnego, byłoby
to sprzeczne ze zwartością $K$. Skończoną wymiarowość
$\ker T^*$ otrzymamy w kroku 4.

\textbf{Krok 2: domkniętość obrazu.} Niech
$N=\ker T$ i $M=N^\perp$. Jądro jest domknięte
jako jądro ciągłego operatora, więc $H=N\oplus M$.
Pokażemy, że istnieje $c>0$
takie, że $\|Tx\|\ge c\|x\|$ dla wszystkich $x\in M$.
Dla $M=\{0\}$ można wybrać dowolne $c>0$.
W przeciwnym razie, gdyby oszacowanie było fałszywe, można wybrać $x_n\in M$,
$\|x_n\|=1$, $Tx_n\to0$. Zwarty operator $K$
daje podciąg $(Kx_{n_k})$ zbieżny w $H$.
Ponieważ $x_{n_k}=Tx_{n_k}+Kx_{n_k}$,
ten sam podciąg $x_{n_k}$ zbiega do pewnego $x\in M$
o normie $1$. Ciągłość $T$ daje $Tx=0$,
czyli $x\in M\cap N=\{0\}$; sprzeczność.
Teraz niech $Tx_n\to y$ w $H$.
Zastępując $x_n$ przez jego rzut na $M$, nie zmieniamy
$Tx_n$. Uzyskane $x_n\in M$ są Cauchy'ego,
bo $\|x_n-x_m\|\le c^{-1}\|Tx_n-Tx_m\|$.
Zatem $x_n\to x\in M$ i $Tx=y$, dowodząc domkniętości
$\operatorname{Ran}T$.

\textbf{Krok 3: równość wymiarów jąder.}
Pokażemy najpierw, że $K$ można aproksymować
w normie operatorowej operatorami skończonego rzędu.
Dla $0<\varepsilon<1$ zbiór $K(B_H)$ jest względnie
zwarty, więc można znaleźć skończony zbiór
$y_1,\dots,y_m$ w $H$ taki, że każdy $Kx$, $\|x\|\le1$,
leży w odległości mniejszej niż $\varepsilon$ od
jednego z $y_j$. Jeśli $W=\operatorname{span}\{y_j\}$,
a $P_W$ oznacza rzut ortogonalny na $W$, to
$F=P_WK$ jest skończonego rzędu i
$\|K-F\|\le\varepsilon$, gdyż
$\|Kx-P_WKx\|=\operatorname{dist}(Kx,W)
\le\min_j\|Kx-y_j\|<\varepsilon$ dla $\|x\|\le1$.
Zapiszmy $K=F+E$, gdzie $\|E\|\le\varepsilon<1$.
Szereg $\sum_{n=0}^\infty E^n$ zbiega w normie operatorowej
z oszacowania $\|E^n\|\le\|E\|^n$.
Mnożąc sumy częściowe z obu stron przez $I-E$
i przechodząc do granicy, otrzymujemy odwrotność $I-E$,
a faktoryzacja
\[
I-K=(I-E)(I-G),\qquad G=(I-E)^{-1}F
\]
przedstawia $T$ jako iloczyn operatora odwracalnego
i operatora $I-G$ z $G$ skończonego rzędu.
Niech $W'=\operatorname{Ran}G$ (przestrzeń skończeniewymiarowa).
W rozkładzie $H=W'\oplus(W')^\perp$ operator $I-G$
ma postać blokową
\[
(I-G)(w,z)=((I-G|_{W'})w-Gz,\ z).
\]
Przez $I_{W'}$ oznaczamy identyczność na $W'$.
Druga współrzędna wektora z jądra musi być zerowa,
więc jądro jest równe
$\ker(I_{W'}-G|_{W'})\oplus\{0\}$.
Oznaczmy $S=I_{W'}-G|_{W'}$ i $L=\operatorname{Ran}S$.
Wprost z postaci blokowej $(a,z)$ należy do obrazu
$I-G$ wtedy i tylko wtedy, gdy $a+Gz\in L$.
Odwzorowanie $(a,z)\mapsto a+Gz+L$ do przestrzeni
ilorazowej $W'/L$ jest liniowe i surjektywne
(bierzemy $z=0$); jego jądrem jest właśnie
$\operatorname{Ran}(I-G)$. Z twierdzenia o izomorfizmie
przestrzeń ilorazowa $H/\operatorname{Ran}(I-G)$
jest izomorficzna z $W'/L$, więc kowymiar obrazu
wynosi $\dim W'-\dim L$.
Z twierdzenia o wymiarze jądra i obrazu w $W'$
ten kowymiar równa się $\dim\ker S$.
Mnożenie z lewej przez odwracalny $I-E$ nie zmienia
jądra, a na przestrzeniach ilorazowych indukuje izomorfizm
$H/\operatorname{Ran}(I-G)\to H/\operatorname{Ran}T$,
$[x]\mapsto[(I-E)x]$, z odwrotnością indukowaną przez
$(I-E)^{-1}$. Nie zmienia zatem kowymiaru obrazu. Otrzymujemy
$\dim\ker T=\operatorname{codim}\operatorname{Ran}T$.

\textbf{Krok 4: warunek rozwiązywalności.}
Dla dowolnego ograniczonego operatora w Hilbercie
$(\operatorname{Ran}T)^\perp=\ker T^*$,
gdyż $(Tu,v)=0$ dla wszystkich $u$ jest dokładnie
warunkiem $T^*v=0$. Obraz $T$ jest domknięty
z kroku 2, stąd
$\operatorname{Ran}T=(\ker T^*)^\perp$.
Ponieważ $\operatorname{Ran}T$ jest domknięty,
rozkład ortogonalny
$H=\operatorname{Ran}T\oplus(\operatorname{Ran}T)^\perp$
identyfikuje przestrzeń ilorazową
$H/\operatorname{Ran}T$ z $(\operatorname{Ran}T)^\perp$
przez rzut ortogonalny.
Stąd $\dim\ker T^*=\operatorname{codim}\operatorname{Ran}T
=\dim\ker T$ na mocy kroku 3. Dowodzi to także
warunku rozwiązywalności dla $f$.
Każde dwa rozwiązania różnią się o element jądra,
a dodanie dowolnego jego elementu do rozwiązania zachowuje
równanie. Jeśli warunek zgodności jest spełniony, rozkładamy
dowolne rozwiązanie $u=u_N+u_M$ w $N\oplus M$.
Ponieważ $Tu_N=0$, wektor $u_M$ również jest rozwiązaniem.
Dwa rozwiązania w $M$ różniłyby się o wektor
z $M\cap N=\{0\}$, więc rozwiązanie w $M$ jest jedyne.
Z kroku 2 ma ono normę co najwyżej $c^{-1}\|f\|$;
jest też rozwiązaniem o najmniejszej normie,
gdyż $\|u_M+n\|^2=\|u_M\|^2+\|n\|^2$ dla $n\in N$.
W szczególności jeśli $\ker T=\{0\}$, równość wymiarów
daje $\ker T^*=\{0\}$, a więc $T$ jest bijekcją
z ograniczoną odwrotnością. Jeśli $T$ jest surjektywny,
to $\ker T^*=\{0\}$, a równość wymiarów daje iniektywność.
\end{proof}

\begin{uwaga}
W odróżnieniu od twierdzenia~\ref{thm:pde12-fredholm},
nie założyliśmy tutaj samosprzężoności $K$.
Dlatego w warunku rozwiązywalności pojawia się
$\ker(I-K^*)$, a niekoniecznie $\ker(I-K)$.
W terminologii teorii operatorów wynik stwierdza,
że $I-K$ jest operatorem Fredholma o indeksie $0$,
czyli $\dim\ker T-\operatorname{codim}\operatorname{Ran}T=0$.
Na przykład na $\C^2$ dla
$K=\begin{pmatrix}1&-1\\0&1\end{pmatrix}$
mamy $T=I-K=\begin{pmatrix}0&1\\0&0\end{pmatrix}$.
W standardowej bazie ortonormalnej $e_1,e_2$ mamy
$T(u_1,u_2)=(u_2,0)$ oraz $T^*(v_1,v_2)=(0,v_1)$.
Zatem $\ker T=\C e_1$, $\ker T^*=\C e_2$,
a równanie $Tu=f$ jest rozwiązywalne dokładnie dla
$f=(f_1,0)$ (wówczas $u_2=f_1$, $u_1$ dowolne).
Na przykład $f=e_2$ jest ortogonalny do $\ker T$,
ale równanie jest nierozwiązywalne, natomiast $f=e_1$
jest rozwiązywalny mimo braku ortogonalności do $\ker T$.
To pokazuje, dlaczego sprzężenie w warunku zgodności
jest niezbędne.
\end{uwaga}


\begin{kontrprzyklad}[Brak zwartości na całej przestrzeni]
\label{ex:pde12-unbounded-domain}
W $H=L^2(\R^d;\C)$ rozważmy $A=-\Delta$ z dziedziną
$H^2(\R^d;\C)$. Na mocy twierdzenia
~\ref{thm:pde8-plancherel} transformata
$U=(2\pi)^{-d/2}\mathcal F$ jest unitarna w konwencji
rozdziału 8. Z reguł różniczkowania dystrybucji
(twierdzenie~\ref{thm:pde8-dist-rules}) wynika
$UAU^{-1}=M_m$, gdzie $m(\xi)=|\xi|^2$ i
$D(M_m)=\{g\in L^2:mg\in L^2\}$.
Istotnie, jeśli $u\in H^2$, to $\partial_j^2u\in L^2$
dla każdego $j$, a w sensie dystrybucji temperowanych
$U(\partial_j^2u)=-\xi_j^2Uu$.
Zatem $|\xi|^2Uu=-\sum_{j=1}^d U(\partial_j^2u)\in L^2$.
Odwrotnie, jeśli $g$ i $|\xi|^2 g$ należą do $L^2$,
to wszystkie iloczyny $(i\xi)^\alpha g$ dla
$|\alpha|\le2$ należą do $L^2$, ponieważ
$|\xi^\alpha|\le 1+|\xi|^2$. Odwrotna transformata
tych iloczynów daje słabe pochodne $U^{-1}g$.

Operator $M_m$ jest gęsto określony, gdyż zawiera
wszystkie funkcje $L^2$ o ograniczonym nośniku, a ich
gęstość wynika z obcięcia do kul i zbieżności w $L^2$.
Jest symetryczny, ponieważ $m$ jest rzeczywista.
Jeśli $h\in D(M_m^*)$ oraz $k=M_m^*h$, to dla każdego
$g\in L^2$ o ograniczonym nośniku
$\int mg\overline h\,\mathrm d\lambda_d
=\int g\overline k\,\mathrm d\lambda_d$.
Ponieważ $m$ jest ograniczona na każdej kuli,
$g=\chi_{B(0,R)}(mh-k)$ należy do $L^2$
i do $D(M_m)$. Wstawiając ten test, otrzymujemy
$\|\chi_{B(0,R)}(mh-k)\|_2^2=0$ dla każdego $R>0$.
Stąd $mh=k\in L^2$, czyli $h\in D(M_m)$.
Wraz z symetrią daje to samosprzężoność $M_m$.
Ponieważ $U$ jest unitarny, definicja sprzężenia daje
$(U^{-1}M_mU)^*=U^{-1}M_m^*U$, z dziedziną
$U^{-1}D(M_m^*)$. Zatem $A$ również jest samosprzężony,
a oba operatory są domknięte.

Dla $z\notin[0,\infty)$ mnożnik $(m-z)^{-1}$ jest
ograniczony przez $\operatorname{dist}(z,[0,\infty))^{-1}$,
a $m/(m-z)=1+z/(m-z)$ także jest ograniczony.
Mnożenie przez $(m-z)^{-1}$ jest więc ograniczoną
odwrotnością $M_m-zI$ o obrazie w $D(M_m)$.
Dla $s\ge0$ zbiór
$E_n=\{\xi:\big||\xi|^2-s\big|<1/n\}$ ma dodatnią
skończoną miarę: jest otwarty i niepusty (zawiera
punkt o normie $\sqrt{s}$) oraz zawiera się w kuli
o promieniu $\sqrt{s+1/n}$. Funkcje $g_n=\chi_{E_n}/\sqrt{\lambda_d(E_n)}$
należą do $D(M_m)$, mają normę $1$ i spełniają
$\|(M_m-sI)g_n\|_2\le1/n$. Wyklucza to istnienie
ograniczonej odwrotności. Jednocześnie zbiór
$\{\xi:|\xi|^2=s\}$ ma miarę zero (dla $s=0$ jest
punktem, dla $s>0$ sferą; miara sfery jest zerowa
przez oszacowanie miarą pierścieni o promieniach
dążących z obu stron do $\sqrt{s}$).
Jeśli $(M_m-sI)g=0$, to $g$ znika prawie wszędzie
poza tym zbiorem, a więc $g=0$ w $L^2$.
Dopełnienie ortogonalne obrazu to jądro sprzężenia
$(M_m-sI)^*=M_m-sI$, z definicji sprzężenia
i rzeczywistości $s$. Obraz jest zatem gęsty,
ale nie jest całym $H$: gdyby był całym $H$,
operator byłby bijekcją i miałby ograniczoną odwrotność,
co przeczy ciągowi $g_n$.
Z definicji każdy $s\ge0$ należy do widma ciągłego.
Unitarna równoważność przenosi iniektywność, gęstość
i surjektywność obrazu oraz ograniczoność odwrotności,
więc $\sigma(A)=[0,\infty)$ jest w całości widmem ciągłym.

Pokażmy niezależnie brak zwartości rezolwenty $B=(A+I)^{-1}$.
Wybierzmy $0\ne f\in C_c^\infty(\R^d)$ oraz $L>0$
tak duże, aby przesunięcia
$f_n(x)=f(x-nLe_1)$ miały parami rozłączne nośniki.
Mamy $\|f_n\|_2=\|f\|_2$, zatem ciąg jest ograniczony.
Operator $B=U^{-1}M_{(1+m)^{-1}}U$ komutuje
z translacjami: transformata przesunięcia o $a$
jest mnożeniem przez $e^{-ia\cdot\xi}$, co wynika
ze zmiany zmiennych dla funkcji całkowalnych
i przez gęstość oraz unitarność dla $L^2$.
Stąd $u_n=Bf_n=u(\cdot-nLe_1)$, gdzie $u=Bf\ne0$
dzięki iniektywności $B$.

Dla dowolnego $v\in C_c^\infty(\R^d)$ o nośniku $Q$
nierówność Cauchy'ego--Schwarza daje
\[
 |(u_n,v)|\le\|v\|_2
 \left(\int_{Q-nLe_1}|u(y)|^2\,\mathrm d\lambda_d(y)
 \right)^{1/2}\longrightarrow0,
\]
gdyż zbiory $Q-nLe_1$ opuszczają każdą ustaloną kulę,
a $|u|^2$ jest całkowalna. Dla ogólnego $v\in L^2$
przybliżamy go testem $v_\varepsilon$ z błędem
$\|v-v_\varepsilon\|_2<\varepsilon$;
$|(u_n,v-v_\varepsilon)|\le\|u\|_2\varepsilon$
jednostajnie względem $n$. Otrzymujemy $u_n\rightharpoonup0$.
Ponieważ $\|u_n\|_2=\|u\|_2>0$, żaden podciąg nie może
zbiegać silnie: jego granica musiałaby być jednocześnie
granicą słabą $0$. Zwartość $B$ przeczyłaby temu faktowi.
Nie należy więc oczekiwać dyskretnego rozwinięcia
po funkcjach własnych jak na obszarze ograniczonym.
\end{kontrprzyklad}

\section{Półgrupa cieplna zbudowana z widma}
\label{sec:pde12-heat}

\begin{definicja}[Silnie ciągła półgrupa]
Niech $H$ będzie przestrzenią Banacha. Rodzinę
$S(t)\in\mathcal L(H)$, $t\ge0$, nazywamy
\emph{silnie ciągłą półgrupą}, gdy $S(0)=I$,
$S(t+s)=S(t)S(s)$ dla wszystkich $s,t\ge0$, a dla
każdego $u\in H$ zachodzi $\|S(t)u-u\|_H\to0$
przy $t\downarrow0$. Jej \emph{generator} $G$ jest
dany przez $Gu=\lim_{t\downarrow0}(S(t)u-u)/t$
na dziedzinie tych $u$, dla których granica istnieje w $H$.
Dziedzina $D(G)$ jest podprzestrzenią liniową, a $G$ jest liniowy,
gdyż ilorazy różnicowe są liniowe i granica w normie
zachowuje kombinacje liniowe.
\end{definicja}

\begin{twierdzenie}[Półgrupa generowana przez $-A$]
\label{thm:pde12-heat}
Niech $A=-\Delta_D$ dla obszaru ograniczonego Lipschitzowskiego $\Omega$, a $H=L^2(\Omega)$.
Dla $t\ge0$ określmy
\[
S(t)u=\sum_{j=1}^\infty e^{-\lambda_jt}(u,\phi_j)_H\phi_j.
\]
Wtedy $(S(t))_{t\ge0}$ jest silnie ciągłą półgrupą na $H$
o generatorze dokładnie $-A$ i
$\|S(t)\|_{H\to H}=e^{-\lambda_1t}$.
Dla $t>0$, całkowitego $m\ge1$ i $u\in H$ zachodzi
$S(t)u\in D(A^m)$ oraz
\[
\|A^mS(t)u\|_H
\le\bigg(\frac{m}{e t}\bigg)^m\|u\|_H.
\]
Dla $u_0\in D(A)$ funkcja $u(t)=S(t)u_0$
należy do $C^1([0,\infty);H)\cap C([0,\infty);D(A))$
(z normą grafową na $D(A)$) i spełnia $u'(t)+Au(t)=0$, $u(0)=u_0$.
Dla dowolnego $u_0\in H$ to samo równanie zachodzi dla
$t>0$, a $u(t)\to u_0$ w $H$ przy $t\downarrow0$.
\end{twierdzenie}
\begin{proof}
\textbf{Definicja i własność półgrupy.}
Niech $u_j=(u,\phi_j)_H$. Ponieważ
$\sum_j e^{-2\lambda_jt}|u_j|^2\le\sum_j|u_j|^2<\infty$,
sumy częściowe są ciągiem Cauchy'ego w $H$ dzięki
ortonormalności. Ich granica definiuje liniowy operator.
Z Parsevala
$\|S(t)u\|_H^2=\sum_je^{-2\lambda_jt}|u_j|^2
\le e^{-2\lambda_1t}\|u\|_H^2$, więc jest on ograniczony. Dla $u=\phi_1$ mamy
równość, co daje dokładną normę. Mnożenie współczynników
$e^{-\lambda_js}e^{-\lambda_jt}=e^{-\lambda_j(s+t)}$
pokazuje $S(t)S(s)=S(t+s)$, a $S(0)=I$.

\textbf{Ciągłość silna.} Dla $u\in H$ mamy
$\|S(t)u-u\|_H^2=\sum_j|e^{-\lambda_jt}-1|^2|u_j|^2$.
Dla każdego $j$ czynnik w module dąży do zera,
a jest ograniczony przez $1$. Ponieważ
$\sum_j|u_j|^2<\infty$, z twierdzenia o zbieżności
zmajoryzowanej dla szeregu otrzymujemy zbieżność do zera.
Dla $h>0$ mamy
$\|S(t+h)u-S(t)u\|\le\|S(h)u-u\|$,
a dla $0<h\le t$ również
$\|S(t)u-S(t-h)u\|\le\|S(h)u-u\|$.
W obu nierównościach zastosowaliśmy własność półgrupy
i $\|S(r)\|\le1$ dla $r\ge0$. Zatem $S(\cdot)u$
jest ciągła w każdym $t\ge0$.

\textbf{Wygładzanie.} Najpierw uzasadnijmy charakterystykę
potęg operatora:
\[
 D(A^m)=\{u\in H:\sum_j\lambda_j^{2m}|u_j|^2<\infty\},
 \qquad A^mu=\sum_j\lambda_j^m u_j\phi_j .
\]
Dla $m=1$ jest to twierdzenie~\ref{thm:pde12-formdomain}.
Jeżeli wzór zachodzi dla $m$, to z definicji
$u\in D(A^{m+1})$ dokładnie wtedy, gdy $u\in D(A^m)$
i $A^mu\in D(A)$. Jeśli $u\in D(A^{m+1})$, z założenia indukcyjnego
współczynniki $A^mu$ wynoszą $\lambda_j^mu_j$,
a warunek $A^mu\in D(A)$ daje
$\sum_j\lambda_j^{2m+2}|u_j|^2<\infty$.
Odwrotnie, jeśli ta suma jest skończona, to również
$\sum_j\lambda_j^{2m}|u_j|^2<\infty$, gdyż
$\lambda_j^{-2}\le\lambda_1^{-2}$. Najpierw otrzymujemy
zatem $u\in D(A^m)$, a następnie $A^mu\in D(A)$
z charakterystyki $D(A)$, zastosowanej do współczynników
$\lambda_j^mu_j$. Zastosowanie $A$ do $A^mu$
daje wzór dla $A^{m+1}$.
Dla $m\ge1$ i $t>0$
$\sup_{x\ge0}x^me^{-tx}=(m/(et))^m$:
pochodna dla $x>0$ wynosi
$x^{m-1}e^{-tx}(m-tx)$, więc funkcja rośnie do $m/t$
i maleje dalej, a w zerze ma wartość $0$.
Wobec tego
$\sum_j\lambda_j^{2m}e^{-2\lambda_jt}|u_j|^2
\le(m/(et))^{2m}\|u\|_H^2$.
Z wykazanej charakterystyki dziedzin potęg wynika
$S(t)u\in D(A^m)$, a Parseval daje podane oszacowanie.
Dokładna norma wynosi
$\|A^mS(t)\|=\sup_j\lambda_j^me^{-t\lambda_j}$:
nierówność w jedną stronę wynika z Parsevala, a w drugą
z podstawienia kolejnych jednostkowych $\phi_j$.
Oszacowanie przez maksimum po wszystkich $x\ge0$
nie musi być równością dla dyskretnego widma.

\textbf{Generator.} Dla $u\in D(A)$ iloraz
$(e^{-\lambda_jt}-1)/t$ dąży do $-\lambda_j$, a
$|(e^{-\lambda_jt}-1)/t|\le\lambda_j$,
ponieważ $1-e^{-a}=\int_0^a e^{-r}\,\mathrm d\lambda_1(r)
\le a$ dla $a\ge0$. Kwadrat modułu różnicy
między współczynnikiem ilorazu a $-\lambda_j u_j$
jest ograniczony przez $4\lambda_j^2|u_j|^2$.
Sumowalność tej majoranty pozwala zastosować zbieżność
zmajoryzowaną i otrzymać
$\lim_{t\downarrow0}(S(t)u-u)/t=-Au$ w $H$.
Odwrotnie, gdy granica ilorazu istnieje w $H$ i równa się
$g$, jej $j$-ty współczynnik musi równać się
$-\lambda_ju_j$. Stąd z Parsevala
$\sum_j\lambda_j^2|u_j|^2=\|g\|_H^2<\infty$,
a więc $u\in D(A)$. To dowodzi równości dziedzin generatora.
Dla $u_0\in D(A)$ charakterystyka spektralna daje
$S(t)u_0\in D(A)$ oraz $AS(t)u_0=S(t)Au_0$.
Ciągłość w normie grafowej wynika więc z silnej ciągłości
$S(\cdot)$ zastosowanej do $u_0$ i $Au_0$.
Prawostronny iloraz różnicowy w chwili $t$ jest równy
$S(t)(S(h)-I)u_0/h$ i dąży do $-S(t)Au_0$.
Dla $t>0$ lewostronny iloraz jest równy
$S(t-h)(S(h)-I)u_0/h$ (jest to iloraz
$(u(t)-u(t-h))/h$ dla $0<h<t$), a jego różnica względem
$-S(t)Au_0$ jest ograniczona przez
\[
 \|(S(h)-I)u_0/h+Au_0\|
 +\|(S(t)-S(t-h))Au_0\|\longrightarrow0.
\]
Pochodna $-S(t)Au_0$ jest ciągła również w zerze
(w zerze rozumiemy pochodną prawostronną).
Dla dowolnego $u_0\in H$ i $t_0>0$ możemy zapisać
$u(t)=S(t-t_0/2)S(t_0/2)u_0$ dla $t\ge t_0/2$.
Wygładzanie daje $S(t_0/2)u_0\in D(A)$,
więc udowodniony wynik stosuje się w otoczeniu $t_0$.
W ten sposób otrzymujemy równanie i ciągłą pochodną
dla wszystkich $t>0$, a w zerze silną ciągłość.
Dla $t>0$ można iterować ten argument: dla każdego
całkowitego $k\ge1$ mamy $u^{(k)}(t)=(-A)^kS(t)u_0$,
ponieważ na otoczeniu dowolnego $t_0>0$ początkowy
wektor $S(t_0/2)u_0$ należy do wszystkich $D(A^k)$.
Indukcyjnie różniczkujemy $S(t-t_0/2)A^kS(t_0/2)u_0$,
korzystając z wykazanej komutacji i przynależności
wektora po prawej stronie do $D(A)$.
Jednoznaczność w klasie rozwiązań ciągłych w $H$
od $t=0$ i klasycznych dla $t>0$ wynika z równań
skalarnych: dla $v'(t)+Av(t)=0$ współczynniki
$v_j(t)=(v(t),\phi_j)_H$ spełniają
$v_j'+\lambda_jv_j=0$. Zatem
$e^{\lambda_jt}v_j(t)=v_j(0)$ przez ciągłość w zerze,
a zupełność bazy daje $v(t)=S(t)v(0)$.
\end{proof}

\begin{twierdzenie}[Wzór Duhamela w $L^2$]
\label{thm:pde12-duhamel}
Niech $T>0$, $u_0\in H=L^2(\Omega)$ i
$f\in C([0,T];H)$, gdzie $S(t)$ jest półgrupą
z twierdzenia~\ref{thm:pde12-heat}. Funkcja
\[
u(t)=S(t)u_0+\int_0^tS(t-s)f(s)\,\mathrm d\lambda_1(s),
\qquad 0\le t\le T,
\]
gdzie całkę względem miary Lebesgue'a na osi czasu
rozumiemy w sensie Bochnera, jest ciągła
z wartościami w $H$ i spełnia oszacowanie
$\|u(t)\|_H\le e^{-\lambda_1t}\|u_0\|_H
+\int_0^te^{-\lambda_1(t-s)}\|f(s)\|_H\,\mathrm d\lambda_1(s)$.
Jest jedynym rozwiązaniem \emph{łagodnym} zagadnienia
$u'(t)+Au(t)=f(t)$, $u(0)=u_0$, zdefiniowanym
przez powyższy wzór całkowy. Jeśli dodatkowo
$u_0\in D(A)$ i $f\in C^1([0,T];H)$, to
$u\in C^1([0,T];H)\cap C([0,T];D(A))$
(z normą grafową na $D(A)$) i równanie zachodzi w $H$
dla każdego $t\in[0,T]$, z pochodnymi jednostronnymi
w punktach końcowych.
\end{twierdzenie}
\begin{proof}
Mapowanie $(t,s)\mapsto S(t-s)f(s)$ jest silnie ciągłe
na trójkącie $0\le s\le t\le T$: dla $(t_n,s_n)\to(t,s)$
różnica wartości jest ograniczona przez
$\|f(s_n)-f(s)\|+\|(S(t_n-s_n)-S(t-s))f(s)\|$,
co dąży do zera. Ciągła funkcja o wartościach w przestrzeni
Banacha na przedziale zwartym jest całkowalna w sensie
Bochnera. Korzystamy tu z kryterium całkowalności
i oszacowania normy całki z analizy funkcjonalnej
(\atlasref{FA-15B::thm:bochner-criterion15b}{kryterium
całkowalności Bochnera} oraz
\atlasref{FA-15B::prop:riemann-bochner15b}{zgodność
całek Riemanna i Bochnera dla funkcji ciągłych}). Stosując dokładną normę $S(t-s)$, dostajemy
podane oszacowanie.
Niech $M=\max_{s\in[0,T]}\|f(s)\|$. Dla $t_n\to t$
na wspólnym przedziale $[0,\min(t_n,t)]$ różnica
całkowanych funkcji dąży jednostajnie do zera
z jednostajnej ciągłości na zwartym trójkącie.
Norma całki na pozostałym odcinku jest nie większa
niż $M|t_n-t|$. To dowodzi ciągłości składnika całkowego;
w szczególności jego wartość w zerze wynosi $0$.

Dla silniejszych danych, po zmianie zmiennej $r=t-s$,
możemy pisać $F(t)=\int_0^tS(r)f(t-r)\,\mathrm d\lambda_1(r)$.
Dla $h>0$ i $t+h\le T$ rozkładamy iloraz różnicowy:
\[
 \begin{aligned}
 \frac{F(t+h)-F(t)}h
 &=\int_0^t S(r)\frac{f(t+h-r)-f(t-r)}h
 \,\mathrm d\lambda_1(r)\\
 &\quad+\frac1h\int_t^{t+h}S(r)f(t+h-r)
 \,\mathrm d\lambda_1(r).
 \end{aligned}
\]
Na pierwszym przedziale ilorazy
$(f(t+h-r)-f(t-r))/h$ zbiegają jednostajnie
do $f'(t-r)$ z jednostajnej ciągłości $f'$ i wzoru
Newtona--Leibniza w przestrzeni Banacha.
Mnożenie przez $S(r)$ nie powiększa błędów.
Średnia całkowa na krótkim odcinku zbiega do $S(t)f(0)$
przez wspólną ciągłość $S(r)f(t+h-r)$.
Dla ilorazu lewostronnego, przy $0<h\le t$, piszemy
\[
 \begin{aligned}
 \frac{F(t)-F(t-h)}h
 &=\int_0^{t-h} S(r)\frac{f(t-r)-f(t-h-r)}h
 \,\mathrm d\lambda_1(r)\\
 &\quad+\frac1h\int_{t-h}^tS(r)f(t-r)
 \,\mathrm d\lambda_1(r).
 \end{aligned}
\]
Na pierwszym przedziale błąd względem całki z $S(r)f'(t-r)$
dąży do zera jednostajnie, a norma brakującej całki po
$[t-h,t]$ nie przekracza $h\max\|f'\|$. Drugi składnik
dąży do $S(t)f(0)$ przez wspólną ciągłość. Otrzymujemy
więc tę samą granicę.
Stąd
$F'(t)=S(t)f(0)+\int_0^tS(r)f'(t-r)\,\mathrm d\lambda_1(r)$,
z pochodnymi jednostronnymi na końcach przedziału.
Pochodna jest ciągła przez ten sam argument co dla $F$,
z $f'$ zamiast $f$.
Ustalmy $t>0$ i $0<\varepsilon<t$.
Na $[\varepsilon,t]$ funkcja $r\mapsto S(r)f(t-r)$
jest ciągła w normie grafowej: zapisujemy
$S(r)=S(\varepsilon/2)S(r-\varepsilon/2)$, gdzie
$S(\varepsilon/2):H\to D(A)$ jest ograniczony dzięki
oszacowaniu wygładzania dla $m=1$. Dziedzina $D(A)$
z normą grafową jest przestrzenią Banacha, więc całka
$F_\varepsilon(t)=\int_\varepsilon^t S(r)f(t-r)\,
\mathrm d\lambda_1(r)$ istnieje w tej przestrzeni.
Odwzorowania $D(A)\to H$, $v\mapsto v$ i $v\mapsto Av$
są ograniczone w normie grafowej, zatem można je
przenieść pod całkę Bochnera (np. przez granicę sum
Riemanna dla funkcji ciągłej).
Z poprzedniego twierdzenia i reguły iloczynu dla
silnej pochodnej mamy dla $r\in[\varepsilon,t]$
\[
 \frac{\mathrm d}{\mathrm dr}[S(r)f(t-r)]
 =-AS(r)f(t-r)-S(r)f'(t-r).
\]
Reguła iloczynu wynika z rozdzielenia ilorazu różnicowego
na zmianę $S(r)$ przy ustalonym wektorze i zmianę
wektora $f(t-r)$; stosujemy $\|S(r)\|\le1$ oraz silną
ciągłość. Pochodna jest ciągła na $[\varepsilon,t]$:
pierwszy składnik jest ciągły dzięki wykazanej ciągłości
grafowej, drugi dzięki wspólnej ciągłości $S(r)f'(t-r)$.
Można więc zastosować wzór Newtona--Leibniza w $H$.
Całkując tę tożsamość, dostajemy
\[
 AF_\varepsilon(t)=S(\varepsilon)f(t-\varepsilon)
 -S(t)f(0)-\int_\varepsilon^t S(r)f'(t-r)\,
 \mathrm d\lambda_1(r).
\]
Gdy $\varepsilon\downarrow0$, mamy
$F_\varepsilon(t)\to F(t)$ w $H$, a prawa strona
zbiega w $H$ do $f(t)-F'(t)$: pierwszy składnik
zbiega przez wspólną ciągłość, a pominięta całka
ma normę co najwyżej $\varepsilon\max\|f'\|$.
Domkniętość $A$ daje $F(t)\in D(A)$ i
$AF(t)=f(t)-F'(t)$. W zerze $F(0)=0$ oraz
$F'(0)=f(0)$, więc tożsamość zachodzi także tam.
Ponieważ $F$ i $AF=f-F'$ są ciągłe w $H$,
$F$ jest ciągła w normie grafowej.
Dla $S(t)u_0$ mamy $\frac{\mathrm d}{\mathrm dt}S(t)u_0
=-AS(t)u_0$ z poprzedniego twierdzenia.
Zsumowanie obu rachunków daje $u'+Au=f$.

Dla dowolnego $u_0\in H$ i ciągłego $f$ wynik definiuje
rozwiązanie łagodne. Jednoznaczność w tej klasie wynika
bezpośrednio z wymogu spełniania wzoru całkowego.
Jeśli $v$ jest rozwiązaniem klasycznym, to
dla każdego $j$ funkcja $v_j(t)=(v(t),\phi_j)_H$
spełnia $v_j'(t)+\lambda_j v_j(t)=f_j(t)$, gdzie
$f_j(t)=(f(t),\phi_j)_H$. Pomnożenie przez
$e^{\lambda_jt}$ i całkowanie daje
\[
 v_j(t)=e^{-\lambda_jt}(u_0,\phi_j)_H+
 \int_0^t e^{-\lambda_j(t-s)} f_j(s)\,
 \mathrm d\lambda_1(s).
\]
Prawa strona jest $j$-tym współczynnikiem wzoru Duhamela,
ponieważ ograniczony funkcjonał $(\cdot,\phi_j)_H$
przenosi się pod całkę Bochnera. Z zupełności bazy
otrzymujemy $v=u$, czyli rozwiązanie łagodne zgadza się
z klasycznym, ilekroć drugie istnieje. W szczególności
sama ciągłość $f$ nie wystarcza na ogół do
przynależności rozwiązania do $D(A)$ w każdej chwili,
jak pokazuje poniższy kontrprzykład.
\end{proof}

\begin{kontrprzyklad}[Ciągła prawa strona bez regularności w dziedzinie operatora]
Ustalmy $T=1$ i $u_0=0$. Wybierzmy rosnący ciąg indeksów
$j_n$ taki, że $L_n=\lambda_{j_n}>4$ oraz
$L_{n+1}\ge4L_n$; jest to możliwe, gdyż $\lambda_j\to\infty$.
Niech $\eta:\R\to[0,1]$ będzie ciągła, niezerowa i z nośnikiem
w $[1,2]$, np. $\eta(r)=\max(0,1-2|r-3/2|)$.
Dla $0\le s\le1$ określmy
\[
 f(s)=\sum_{n\ge1}n^{-1/2}\eta(L_n(1-s))\phi_{j_n}.
\]
Niezerowe składniki mogą występować tylko na przedziałach
$[1-2/L_n,1-1/L_n]$, które są parami rozłączne.
W każdej chwili co najwyżej jeden składnik jest niezerowy,
a przedziały gromadzą się tylko przy $s=1$.
Funkcja jest ciągła na $[0,1)$ dzięki ciągłości $\eta$,
a w $s=1$ również, ponieważ $f(1)=0$ i amplitudy
$n^{-1/2}$ dążą do zera. Zatem $f\in C([0,1];H)$.
W chwili $t=1$ współczynniki rozwiązania łagodnego wynoszą
\[
 \begin{aligned}
 (u(1),\phi_{j_n})_H
 &=n^{-1/2}\int_0^1 e^{-L_n(1-s)}
 \eta(L_n(1-s))\,\mathrm d\lambda_1(s)\\
 &=\frac{C_\eta}{\sqrt n\,L_n},
 \qquad
 C_\eta=\int_1^2 e^{-r}\eta(r)\,\mathrm d\lambda_1(r)>0 .
 \end{aligned}
\]
Pozostałe współczynniki znikają. Druga równość wynika
z podstawienia $r=L_n(1-s)$ i nośnika $\eta$.
Zatem $\sum_j\lambda_j^2|(u(1),\phi_j)_H|^2
=C_\eta^2\sum_{n\ge1}1/n=\infty$.
Twierdzenie~\ref{thm:pde12-formdomain} daje
$u(1)\notin D(A)$ mimo ciągłości $f$ i zerowego $u_0$.
\end{kontrprzyklad}

\begin{uwaga}[Regularność w czasie i na brzegu]
Oszacowanie $\|A^mS(t)\|\lesssim t^{-m}$ jest
oszacowaniem operatorowym. Nie daje automatycznie
$C^\infty(\overline\Omega)$ na dowolnym obszarze Lipschitzowskim: w narożach może zabraknąć regularności
brzegowej. Dla jednorodnego równania ciepła i $t>0$
$u(t)=S(t)u_0$ jest jednak gładka przestrzennie we wnętrzu.
Istotnie, dla każdego $m\ge1$ wektory $u(t),Au(t),\ldots,
A^{m-1}u(t)$ należą do $D(A)\subset H_0^1$,
a $A^mu(t)\in L^2$. Stosując wewnętrzną regularność
eliptyczną (twierdzenia~\ref{thm:pde15-h2}
i~\ref{thm:pde15-bootstrap}) kolejno do
$A^{m-1}u(t),\ldots,u(t)$, otrzymujemy
$A^{m-1}u(t)\in H^2_{\mathrm{loc}}$,
$A^{m-2}u(t)\in H^4_{\mathrm{loc}}$ i w końcu
$u(t)\in H^{2m}_{\mathrm{loc}}$.
Ponieważ $m$ jest dowolne, dla każdego zwartego
podzbioru wnętrza i dowolnego $k\ge0$ wybieramy
$2m>k+d/2$. Lokalne zanurzenia Sobolewa dają
reprezentanta klasy $C^k$ w otoczeniu tego zbioru.
Reprezentanci dla różnych $k$ i otoczeń zgadzają się
na częściach wspólnych, gdyż są funkcjami ciągłymi
zgodnymi prawie wszędzie. Otrzymujemy reprezentanta klasy $C^\infty(\Omega)$.
Ten argument nie dotyczy dowolnej niegładkiej prawej strony.
\end{uwaga}

\section{Równanie falowe: propagacja i zachowanie energii}
\label{sec:pde12-wave}

Piszemy $V=H_0^1(\Omega)$, $H=L^2(\Omega)$ i utożsamiamy
$H$ z podprzestrzenią $V^*=H^{-1}(\Omega)$ za pomocą
odwzorowania $h\mapsto(v\mapsto(h,v)_H)$.
Jest ono ciągłe z nierówności Poincarégo i różnowartościowe
z gęstości $V$ w $H$. W tej sekcji pracujemy w przestrzeniach
rzeczywistych; wyniki dla danych zespolonych otrzymujemy
przez rozdzielenie części rzeczywistej i urojonej.
Rozszerzenie energetyczne operatora $A$ oznaczamy przez
$\mathcal A:V\to V^*$, gdzie
\[
 \langle\mathcal Au,v\rangle
 =\int_\Omega\nabla u(x)\cdot\nabla v(x)\,\mathrm d\lambda_d(x).
\]
Jest ono ograniczone z nierówności Cauchy'ego--Schwarza
i zgodne z $A$ na $D(A)$.
Przez rozwiązanie energetyczne rozumiemy funkcję
$u\in C(\R;V)\cap C^1(\R;H)$ spełniającą dane początkowe
oraz dla każdego $v\in V$ i $\psi\in C_c^\infty(\R)$
tożsamość
\[
 \int_\R (u(t),v)_H\psi''(t)\,\mathrm d\lambda_1(t)
 +c^2\int_\R\langle\mathcal Au(t),v\rangle\psi(t)\,
 \mathrm d\lambda_1(t)=0.
\]
W tym sensie $u_{tt}+c^2\mathcal Au=0$ jest równaniem
o wartościach w $V^*$; nie wymagamy $u(t)\in D(A)$.

\begin{twierdzenie}[Spektralne rozwiązanie równania falowego]
\label{thm:pde12-wave}
Niech $A=-\Delta_D$ jak wyżej, $c>0$,
$u_0\in H_0^1(\Omega)$ i $u_1\in L^2(\Omega)$.
Oznaczmy $a_j=(u_0,\phi_j)$, $b_j=(u_1,\phi_j)$.
Wtedy szereg
\[
u(t)=\sum_{j=1}^\infty
\left(a_j\cos(c\sqrt{\lambda_j}t)
+\frac{b_j}{c\sqrt{\lambda_j}}\sin(c\sqrt{\lambda_j}t)\right)\phi_j
\]
wyznacza jedyne rozwiązanie energetyczne równania
$u_{tt}+c^2Au=0$ z danymi $u(0)=u_0$, $u_t(0)=u_1$,
przy czym
$u\in C(\R;H_0^1(\Omega))\cap C^1(\R;L^2(\Omega))$.
Energia
\[
E(t)=\frac12\|u_t(t)\|_2^2+
\frac{c^2}{2}\|\nabla u(t)\|_2^2
\]
jest stała i równa $E(0)$.
\end{twierdzenie}
\begin{proof}
Z twierdzenia~\ref{thm:pde12-formdomain} mamy
$\sum_j\lambda_j|a_j|^2<\infty$ i
$\sum_j|b_j|^2<\infty$. Dla współczynnika
$u_j(t)=a_j\cos(c\sqrt{\lambda_j}t)
+b_j\sin(c\sqrt{\lambda_j}t)/(c\sqrt{\lambda_j})$
zachodzi
\[
|u_j'(t)|^2+c^2\lambda_j|u_j(t)|^2
=|b_j|^2+c^2\lambda_j|a_j|^2.
\]
Istotnie, dla $\omega_j=c\sqrt{\lambda_j}>0$
para $(\omega_j u_j(t),u_j'(t))$ jest równa
\[
 \begin{pmatrix}\cos(\omega_jt)&\sin(\omega_jt)\\
 -\sin(\omega_jt)&\cos(\omega_jt)\end{pmatrix}
 \begin{pmatrix}\omega_j a_j\\b_j\end{pmatrix}.
\]
Macierz jest ortogonalna, więc zachowuje sumę kwadratów
modułów obu współrzędnych (również zespolonych).
Oznaczmy przez $u_N$ sumę pierwszych $N$ modów.
Dla $M>N$, z Parsevala i charakterystyki dziedziny formy,
\begin{align*}
 \|u_M'(t)-u_N'(t)\|_2^2
 &+c^2\|\nabla(u_M(t)-u_N(t))\|_2^2\\
 &=\sum_{j=N+1}^{M}(|b_j|^2+c^2\lambda_j|a_j|^2).
\end{align*}
Prawa strona nie zależy od $t$ i dąży do zera
jednostajnie względem $M>N$. Ponieważ norma gradientu
jest równoważna pełnej normie $V$, z zupełności
$C([-T,T];V)$ i $C([-T,T];H)$ uzyskujemy jednostajne
granice $u_N\to u$ i $u_N'\to p$ na każdym takim przedziale.
Granice są zgodne na częściach wspólnych przedziałów.
W tożsamości
$u_N(t)-u_N(0)=\int_0^t u_N'(s)\,\mathrm d\lambda_1(s)$
przechodzimy do granicy w $H$ (dla $t<0$ używamy
całki zorientowanej). Z ciągłości $p$ i twierdzenia
Newtona--Leibniza wynika $u'=p$ w $H$.
Dane początkowe otrzymujemy z rozwinięć $u_N(0)$ w $V$
oraz $u_N'(0)$ w $H$.

Każde $u_N$ jest gładkie w czasie i spełnia
$u_N''+c^2Au_N=0$ w $H$, ponieważ każdy współczynnik
spełnia skalarne równanie oscylatora.
Dla $v\in V$ mnożymy równanie przez $\psi$ i całkujemy
dwukrotnie przez części w czasie. W otrzymanej tożsamości
błędy przy zastąpieniu $u_N$ przez $u$ są ograniczone
przez
$\sup_{|t|\le T}\|u_N(t)-u(t)\|_H\|v\|_H\|\psi''\|_{L^1}$
oraz
$c^2\sup_{|t|\le T}\|u_N(t)-u(t)\|_V\|v\|_V\|\psi\|_{L^1}$,
gdy nośnik $\psi$ leży w $[-T,T]$.
Oba błędy znikają, co daje tożsamość definiującą
rozwiązanie energetyczne.

Dla jednoznaczności niech $w$ będzie różnicą dwóch
rozwiązań energetycznych z zerowymi danymi. Wybierając
$v=\phi_j$ w tożsamości słabej, z równania własnego
otrzymujemy $w_j''=-c^2\lambda_j w_j$ dystrybucyjnie,
gdzie $w_j=(w,\phi_j)_H\in C^1(\R)$.
Prawa strona jest ciągła, więc $w_j'$ jest klasycznie
różniczkowalna: różnica między $w_j'$ a funkcją
$t\mapsto\int_0^t(-c^2\lambda_jw_j(s))\,\mathrm d\lambda_1(s)$
ma zerową pochodną dystrybucyjną, a zatem jest stała.
W konsekwencji $w_j\in C^2(\R)$ i spełnia klasyczne
równanie oscylatora z $w_j(0)=w_j'(0)=0$.
Jednoznaczność tego równania daje $w_j=0$ dla każdego $j$.
Zupełność bazy daje $w=0$.
Dla sum częściowych energia jest stała dzięki tożsamości
modalnej. Silna zbieżność obu składników energii oraz
danych początkowych daje po przejściu do granicy
$E(t)=E(0)$ dla każdego $t\in\R$.
\end{proof}

\begin{uwaga}[Silniejsze dane początkowe]
Jeśli ponadto $u_0\in D(A)$ i $u_1\in V$, to
$u\in C(\R;D(A))\cap C^1(\R;V)\cap C^2(\R;H)$
i równanie zachodzi w $H$. Istotnie, mnożymy tożsamość
energii każdego modu przez $\lambda_j$.
Sumowalność $\sum_j(\lambda_j|b_j|^2+
c^2\lambda_j^2|a_j|^2)$ daje jednostajną zbieżność
ogonów $u_N$ w normie grafowej $A$ oraz $u_N'$ w $V$;
normę $L^2$ kontroluje $\lambda_1>0$.
Ponadto $u_N''=-c^2Au_N$ zbiega jednostajnie w $H$
na każdym przedziale zwartym. Dwukrotne zastosowanie
argumentu z całkowaniem po czasie dowodzi podanej
różniczkowalności i równania.
\end{uwaga}

\begin{uwaga}[Czego analiza spektralna nie zastępuje?]
Rozwinięcie w funkcje własne pokazuje istnienie i
zachowanie globalnej energii, ale nie ujawnia bezpośrednio
geometrii stożków charakterystycznych. Skończona prędkość
propagacji, uzyskana w twierdzeniu~\ref{thm:pde5-cone}
metodą energii lokalnej (dla rozważanych tam założeń),
pozostaje niezależną, silniejszą lokalnie informacją.
\end{uwaga}

\section{Przepływ Schrödingera i związek z operatorem unitarnym}
\label{sec:pde12-schroedinger}

\begin{twierdzenie}[Ewolucja unitarna laplasjanu Dirichleta]
\label{thm:pde12-unitary}
Dla laplasjanu Dirichleta $A$ na ograniczonym obszarze
Lipschitzowskim $\Omega$ przechodzimy do zespolonej przestrzeni
$H=L^2(\Omega;\C)$ i zespolonego operatora $A$.
Zdefiniujmy $U(t)u=\sum_je^{-it\lambda_j}(u,\phi_j)_H\phi_j$,
$t\in\R$. Wtedy $(U(t))_{t\in\R}$ jest silnie
ciągłą grupą unitarną z generatorem $-iA$.
Dla każdego $u_0\in H$ zachodzi $\|U(t)u_0\|_2=\|u_0\|_2$.
Jeśli $u_0\in D(A)$, to $u(t)=U(t)u_0$ jest jedynym
rozwiązaniem $iu_t=Au$, $u(0)=u_0$ w klasie
$C^1(\R;H)\cap C(\R;D(A))$, gdzie $D(A)$ ma normę grafową.
\end{twierdzenie}
\begin{proof}
Piszemy $u_j=(u,\phi_j)_H$. Ponieważ $|e^{-it\lambda_j}|=1$
i $\sum_j|u_j|^2<\infty$, szereg jest zbieżny w $H$
z kryterium Cauchy'ego dla szeregu ortogonalnego.
Z liniowości granicy otrzymujemy operator liniowy,
a Parseval daje $\|U(t)u\|_2=\|u\|_2$.
Na współczynnikach odczytujemy
$U(t)U(s)=U(t+s)$ oraz $U(0)=I$.
Zatem $U(t)$ jest surjektywną izometrią liniową
z odwrotnością $U(-t)$, a więc operatorem unitarnym
(z tożsamości polaryzacyjnej).
Silna ciągłość w zerze wynika z
$\sum_j|e^{-it\lambda_j}-1|^2|u_j|^2\to0$
przez zbieżność zmajoryzowaną, z majorantą
$4|u_j|^2$. Ciągłość w pozostałych punktach wynika
z własności grupy i izometrii.

Dla $u\in D(A)$ i $h\ne0$ mamy
\[
 \left|\frac{e^{-ih\lambda_j}-1}{h}\right|
 =\left|\frac{1}{h}\int_0^h(-i\lambda_j)e^{-ir\lambda_j}
 \,\mathrm d\lambda_1(r)\right|\le\lambda_j .
\]
Iloraz dąży do $-i\lambda_j$ dla $h\to0$.
Kwadrat modułu różnicy odpowiednich współczynników
jest ograniczony przez $4\lambda_j^2|u_j|^2$,
co jest sumowalne. Zatem
$\lim_{h\to0}(U(h)u-u)/h=-iAu$ w $H$.
Odwrotnie, jeśli iloraz ma granicę $g\in H$,
ciągłość funkcjonałów $(\cdot,\phi_j)_H$ daje
$(g,\phi_j)_H=-i\lambda_j u_j$ dla każdego $j$.
Parseval daje $\sum_j\lambda_j^2|u_j|^2=\|g\|_H^2<\infty$,
czyli $u\in D(A)$.
Generator jest więc dokładnie $-iA$.

Z charakterystyki spektralnej wynika, że $U(t)$ zachowuje
$D(A)$ i $AU(t)u=U(t)Au$ na $D(A)$.
Ciągłość grafowa wynika z silnej ciągłości zastosowanej
do $u$ i $Au$. Iloraz różnicowy w dowolnej chwili $t$
jest równy $U(t)(U(h)u-u)/h$; stąd pochodna
$-iU(t)Au$ istnieje i jest ciągła w $H$.
To dowodzi równania Schrödingera. Dla innego rozwiązania
$v$ w podanej klasie współczynniki spełniają
$v_j'=-i\lambda_jv_j$, zatem
$v_j(t)=e^{-it\lambda_j}(u_0,\phi_j)_H$.
Zupełność bazy dowodzi jednoznaczności.
Ponadto dla każdego $u\in V$ tożsamość spektralna energii
daje $\|\nabla U(t)u\|_2=\|\nabla u\|_2$; dla $u\in D(A)$
zachodzi $\|AU(t)u\|_2=\|Au\|_2$.
Przepływ zachowuje więc również normę formy i normę grafową.
\end{proof}

\begin{uwaga}
Dla ogólnego operatora samosprzężonego, niekoniecznie
mającego zwartą rezolwentę, rachunek spektralny daje
silnie ciągłą grupę unitarną $e^{-itA}$.
Półgrupa $e^{-tA}$ składa się z operatorów ograniczonych
dla wszystkich $t\ge0$, jeśli $A$ jest ograniczony z dołu;
jeżeli $\sigma(A)\subset[a,\infty)$, to
$\|e^{-tA}\|\le e^{-ta}$. Bez dolnego ograniczenia widma
nie można w ogólności definiować w ten sposób półgrupy
ograniczonych operatorów na całym $H$.
Twierdzenie Stone'a identyfikuje generator każdej silnie
ciągłej grupy unitarnej jako $-iA$ dla dokładnie jednego
operatora samosprzężonego $A$ (konwencja $e^{-itA}$;
dla $e^{itA}$ generator ma znak przeciwny).
Ogólny dowód i rachunek spektralny podano w analizie
funkcjonalnej, rozdział 33
(\atlasref{FA-33::thm:stone33}{twierdzenie Stone'a}).
Tutaj dowiedliśmy bezpośrednio konstrukcji dla laplasjanu
Dirichleta z dyskretnym widmem.
\end{uwaga}

\section{Podsumowanie: co wynika z założeń, a co nie?}
\label{sec:pde12-summary}

Na obszarze ograniczonym Lipschitzowskim
słaba forma Dirichleta definiuje samosprzężony
operator $A=-\Delta_D$. Koercywność formy przesuniętej zapewnia ciągłość
$(A+I)^{-1}:L^2\to H_0^1$, a zwartość zanurzenia
$H_0^1\hookrightarrow L^2$ zapewnia zwartość rezolwenty.
W połączeniu z samosprzężonością i różnowartościowością
rezolwenty otrzymujemy dyskretne widmo i ortonormalną
bazę funkcji własnych. Stąd otrzymujemy alternatywę
Fredholma, rozwinięcia dla równania ciepła i falowego
oraz przepływ unitarny równania Schrödingera.

Nie dowiedliśmy tu ogólnego twierdzenia spektralnego dla
nieograniczonych operatorów samosprzężonych ani
pełnego twierdzenia Hille'a--Yosidy. W kolejnych rozdziałach
wykorzystamy jednak już udowodnione przypadki oraz
wyraźnie oddzielimy twierdzenia w pełnej ogólności od
wyników obowiązujących dla zwartej rezolwenty.

% Rozdział 13. Tekst przeznaczony do włączenia przed \appendix.

\chapter[Eliptyczne zagadnienia brzegowe]{Eliptyczne zagadnienia brzegowe: formy dywergencyjne, nierówność Gårdinga i alternatywa Fredholma}
\label{chap:pde13}

W rozdziałach~\ref{chap:pde-variational} i~\ref{chap:pde12} zbudowaliśmy
słabą teorię laplasjanu. Teraz współczynniki będą zależały od punktu,
możliwe będą człony pierwszego i zerowego rzędu, a forma nie musi być
symetryczna. Rozdzielimy trzy pytania: czy operator jest eliptyczny,
czy związana z nim forma jest koercywna, i czy odpowiednie zagadnienie
brzegowe ma dokładnie jedno rozwiązanie. Wykazanie samej eliptyczności
nie rozstrzyga ostatniego pytania.

\section{Postać dywergencyjna a niedywergencyjna}
\label{sec:pde13-operators}
Niech $d\geq1$, a $\Omega\subset\R^d$ będzie obszarem ograniczonym Lipschitzowskim. Wszystkie współczynniki w tym rozdziale są rzeczywiste;
analogiczne sformułowania zespolone wymagają form półtoraliniowych.
Dla funkcji $u$ rozważmy formalne wyrażenie różniczkowe
\begin{equation}\label{eq:pde13-div}
 Lu=-\sum_{i,j=1}^d\partial_i(a_{ij}\partial_j u)
       +\sum_{i=1}^d b_i\partial_i u+cu.
\end{equation}
Zakładamy $a_{ij},b_i,c\in L^\infty(\Omega)$; piszemy
$A(x)=(a_{ij}(x))$, $b(x)=(b_i(x))$. Ponadto istnieje $\lambda>0$
takie, że dla prawie każdego $x\in\Omega$ i każdego $\xi\in\R^d$
\begin{equation}\label{eq:pde13-elliptic}
 \xi^{\mathsf T}A(x)\xi\geq\lambda|\xi|^2.
\end{equation}
Warunek ten nazywamy \emph{jednostajną eliptycznością}.
Wybieramy wspólny zbiór pełnej miary, na którym nierówność
zachodzi dla wszystkich $\xi$; wystarczy jej spełnienie dla
wektorów wymiernych poza przeliczalną sumą zbiorów miary zero
i ciągłość obu stron względem $\xi$.
Niech $A_s=(A+A^{\mathsf T})/2$ oraz
$A_a=(A-A^{\mathsf T})/2$. Ponieważ $A_a^{\mathsf T}=-A_a$,
skalar $\xi^{\mathsf T}A_a\xi$ jest równy swojej transpozycji
$-\xi^{\mathsf T}A_a\xi$, a więc wynosi zero.
Zatem warunek zależy tylko od $A_s$, która jest jednostajnie
dodatnio określona; nie wymaga on symetrii $A$.
Oznaczamy przez $\|A\|_\infty$ kres istotny normy operatorowej
macierzy, a przez $\|b\|_\infty$ kres istotny normy euklidesowej
wektora. Obie liczby są skończone dzięki założeniom.

\begin{definicja}[Symbol główny i postać niedywergencyjna]
\label{def:pde13-symbol}
Symbolem głównym wyrażenia $L$ z~\eqref{eq:pde13-div} jest
$p_2(x,\xi)=\sum_{i,j}a_{ij}(x)\xi_i\xi_j$.
Formalny operator w \emph{postaci niedywergencyjnej} ma postać
$\widetilde Lu=-\sum_{i,j}a_{ij}\partial_{ij}u+
\sum_i\widetilde b_i\partial_i u+\widetilde c u$.
Jeśli $a_{ij}\in W^{1,\infty}(\Omega)$ i $u\in H^2_{\mathrm{loc}}(\Omega)$,
to lokalna reguła Leibniza dla słabych pochodnych
(twierdzenie~\ref{thm:pde9-products}, zastosowane z $p=q=2$
na kulach o domknięciu zawartym w $\Omega$, na których
$a_{ij}\in H^1$ i $\partial_j u\in H^1$) daje dla każdego $i,j$
\[
 -\partial_i(a_{ij}\partial_j u)
 =-a_{ij}\partial_{ij}u-(\partial_i a_{ij})\partial_j u
\]
w sensie dystrybucji; oba składniki po prawej są lokalnie w $L^2$.
Po zsumowaniu otrzymujemy $\widetilde L=L$ na tej klasie funkcji,
gdy $\widetilde b_j=b_j-\sum_i\partial_i a_{ij}$ oraz
$\widetilde c=c$. Przy samych współczynnikach mierzalnych
nie wolno używać tego przekształcenia jako tożsamości operatorów
na $H^1$: drugie pochodne $u$ oraz pochodne współczynników
mogą być dystrybucjami, których wskazane iloczyny nie są określone.
\end{definicja}

\begin{przyklad}[Ta sama eliptyczność, inne pojęcie rozwiązania]
\label{ex:pde13-interface}
Na $\Omega=(-1,1)$ niech $a(x)=1$ dla $x<0$ i $a(x)=4$ dla $x>0$.
Funkcja $u(x)=x$ dla $x\le0$ oraz $u(x)=x/4$ dla $x\ge0$
jest ciągła i kawałkami liniowa. Całkując przez części na obu
podprzedziałach dla dowolnej $\varphi\in C_c^\infty(-1,1)$,
otrzymujemy
\[
 -\int_{-1}^1u(x)\varphi'(x)\,\mathrm d\lambda_1(x)
 =\int_{-1}^0\varphi(x)\,\mathrm d\lambda_1(x)
 +\frac14\int_0^1\varphi(x)\,\mathrm d\lambda_1(x),
\]
gdyż człony brzegowe w zerze znikają. Zatem
$u'=\chi_{(-1,0)}+\tfrac14\chi_{(0,1)}\in L^2(-1,1)$
i $u\in H^1(-1,1)$. Ponieważ $au'=1$ prawie wszędzie,
mamy $-(au')'=0$ dystrybucyjnie.
Ponowne różniczkowanie, jak w przykładzie~\ref{ex:pde7-heaviside},
daje
\[
 \langle u'',\varphi\rangle
 =-\int_{-1}^0\varphi'(x)\,\mathrm d\lambda_1(x)
 -\frac14\int_0^1\varphi'(x)\,\mathrm d\lambda_1(x)
 =-\frac34\varphi(0).
\]
Stąd $u''=-\tfrac34\delta_0$, więc $u\notin H^2(-1,1)$.
Iloczyn $a u''$ nie jest kanonicznie określony przez klasę
$a\in L^\infty$: wybór wartości reprezentanta $a(0)$
zmieniałby formalny iloczyn z deltą, choć nie zmienia
współczynnika jako klasy funkcji. Postać niedywergencyjna
nie jest więc jedynie innym zapisem tego samego zagadnienia.
\end{przyklad}

\section{Rozwiązanie słabe i funkcjonał prawej strony}
\label{sec:pde13-weak}
Przez $H^{-1}(\Omega)$ rozumiemy przestrzeń dualną
$(H_0^1(\Omega))^*$ (definicja~\ref{def:pde11-dual}).
Niech $V=H_0^1(\Omega)$, z pełną normą Sobolewa
$\|u\|_V^2=\|u\|_{L^2}^2+\|\nabla u\|_{L^2}^2$.
W rozdziale 11 używano normy gradientu; obie normy są
równoważne na mocy wniosku~\ref{cor:pde10-Poincare},
ale tutaj normę $\|F\|_{V^*}$ liczymy względem pełnej normy.
Przestrzeń $V$ jest przestrzenią Hilberta jako domknięta
podprzestrzeń $H^1(\Omega)$
(twierdzenie~\ref{thm:pde9-banach}).
Definiujemy formę dwuliniową na $H^1(\Omega)\times H^1(\Omega)$,
w szczególności na $V\times V$,
\begin{equation}\label{eq:pde13-form}
 a(u,v)=\int_\Omega\big(A\nabla u\cdot\nabla v
 +(b\cdot\nabla u)v+cuv\big)\,\mathrm d\lambda_d(x).
\end{equation}
Z nierówności Cauchy'ego--Schwarza i ograniczoności
współczynników dla $u,v\in H^1(\Omega)$ otrzymujemy
\begin{align*}
 |a(u,v)|&\le \|A\|_\infty\|\nabla u\|_2\|\nabla v\|_2\\
 &\quad+\|b\|_\infty\|\nabla u\|_2\|v\|_2
 +\|c\|_\infty\|u\|_2\|v\|_2\\
 &\le M\|u\|_{H^1}\|v\|_{H^1},
\end{align*}
gdzie $M=\|A\|_\infty+\|b\|_\infty+\|c\|_\infty$.
Na $V$ jest to oszacowanie względem normy $\|\cdot\|_V$.
Nie zakładamy symetrii formy.

\begin{definicja}[Zagadnienie Dirichleta w sensie słabym]
\label{def:pde13-weak}
Dla $F\in V^*=H^{-1}(\Omega)$ funkcja $u\in V$ jest
\emph{słabym rozwiązaniem} zagadnienia $Lu=F$ z jednorodnym
warunkiem Dirichleta, jeżeli
\begin{equation}\label{eq:pde13-weak}
 a(u,v)=\langle F,v\rangle_{V^*,V}\qquad(v\in V).
\end{equation}
Jeśli $f\in L^2(\Omega)$, to utożsamiamy $f$ z elementem $F\in V^*$
przez $\langle F,v\rangle=\int_\Omega f(x)v(x)\,\mathrm d\lambda_d(x)$.
Z nierówności Cauchy'ego--Schwarza mamy
$|\langle F,v\rangle|\le\|f\|_2\|v\|_V$, zatem
$\|F\|_{V^*}\le\|f\|_2$. To utożsamienie jest różnowartościowe,
gdyż funkcje testowe są gęste w $L^2(\Omega)$.
Warunek brzegowy jest zakodowany w przynależności $u\in H_0^1$,
równoważnej zerowemu śladowi na mocy
twierdzenia~\ref{thm:pde10-trace-zero}.
\end{definicja}

\begin{stwierdzenie}[Zgodność ze znaczeniem dystrybucyjnym]
\label{prop:pde13-distributional}
Jeżeli $u\in H_0^1(\Omega)$ spełnia~\eqref{eq:pde13-weak}, to
$Lu=F$ jako dystrybucje na $\Omega$. Odwrotnie, gdy
$u\in H_0^1(\Omega)$ spełnia $Lu=F$ dystrybucyjnie dla
$F\in H^{-1}(\Omega)$, to spełnia~\eqref{eq:pde13-weak}.
\end{stwierdzenie}
\begin{proof}
Element $F\in V^*$ definiuje dystrybucję przez ograniczenie
do $C_c^\infty(\Omega)$: dla testów o nośniku w ustalonym
zbiorze zwartym norma $H^1$ jest ograniczona przez stałą
razy supremum modułów funkcji i jej pierwszych pochodnych.
Ograniczenie to wyznacza $F$ jednoznacznie przez gęstość testów w $V$.
Dla $\varphi\in C_c^\infty(\Omega)$ definicja pochodnej
dystrybucyjnej daje
\[
 \langle-\partial_i(a_{ij}\partial_j u),\varphi\rangle
 =\int_\Omega a_{ij}(x)(\partial_j u)(x)(\partial_i\varphi)(x)
 \,\mathrm d\lambda_d(x).
\]
Iloczyn $a_{ij}\partial_j u$ należy do $L^2$, podobnie
$b_i\partial_i u$ i $cu$, więc wszystkie wyrażenia definiują
odpowiednie dystrybucje. Po zsumowaniu po indeksach otrzymujemy
$\langle Lu,\varphi\rangle=a(u,\varphi)$.
Jeśli zachodzi~\eqref{eq:pde13-weak}, równość z $F$ zachodzi
w szczególności dla testów, co oznacza $Lu=F$ dystrybucyjnie.
Odwrotnie, wybieramy dla $v\in V$ ciąg testów $\varphi_n\to v$
w normie $V$. Z ciągłości formy i $F$ mamy
\[
 |a(u,v)-\langle F,v\rangle|
 \le (M\|u\|_V+\|F\|_{V^*})\|v-\varphi_n\|_V\longrightarrow0,
\]
gdyż równość zachodzi dla $\varphi_n$. To daje~\eqref{eq:pde13-weak}.
\end{proof}

\section{Nierówność Gårdinga i wymuszenie koercywności}
\label{sec:pde13-garding}
\begin{twierdzenie}[Elementarna nierówność Gårdinga]
\label{thm:pde13-garding}
Przy założeniach~\eqref{eq:pde13-elliptic} oraz
$a_{ij},b_i,c\in L^\infty(\Omega)$ forma~\eqref{eq:pde13-form}
spełnia
\begin{equation}\label{eq:pde13-garding}
 a(u,u)\geq\frac\lambda2\|\nabla u\|_2^2
 -\Gamma\|u\|_2^2,\qquad u\in V,
\end{equation}
gdzie można przyjąć
$\Gamma=\|b\|_\infty^2/(2\lambda)+\|c^-\|_\infty$,
$c^-=\max\{-c,0\}$. W szczególności dla
$\mu\geq\Gamma+\lambda/2$ forma
$a_\mu(u,v)=a(u,v)+\mu(u,v)_{L^2}$ jest koercywna na $V$:
$a_\mu(u,u)\geq(\lambda/2)\|u\|_V^2$.
\end{twierdzenie}
\begin{proof}
Z eliptyczności, stosowanej do $\xi=\nabla u(x)$ prawie wszędzie,
otrzymujemy
$\int_\Omega A(x)\nabla u(x)\cdot\nabla u(x)\,\mathrm d\lambda_d(x)
\geq\lambda\|\nabla u\|_2^2$.
Ponadto $c(x)u(x)^2\geq-\|c^-\|_\infty |u(x)|^2$ prawie wszędzie,
więc po scałkowaniu człon zerowego rzędu jest nie mniejszy
niż $-\|c^-\|_\infty\|u\|_2^2$.
Człon pierwszego rzędu szacujemy przez nierówność Cauchy'ego--Schwarza
i nierówność Younga $st\leq\eps s^2/2+t^2/(2\eps)$
dla $s,t\geq0$, $\eps>0$ (wynikającą z
$(\sqrt{\eps}s-t/\sqrt{\eps})^2\geq0$), wybierając $\eps=\lambda$:
\[
 \left|\int_\Omega(b(x)\cdot\nabla u(x))u(x)\,\mathrm d\lambda_d(x)\right|
 \leq \|b\|_\infty\|\nabla u\|_2\|u\|_2
 \leq\frac\lambda2\|\nabla u\|_2^2+
 \frac{\|b\|_\infty^2}{2\lambda}\|u\|_2^2.
\]
Dodanie trzech oszacowań daje~\eqref{eq:pde13-garding}.
Po dodaniu $\mu\|u\|_2^2$ otrzymujemy
$a_\mu(u,u)\geq(\lambda/2)\|\nabla u\|_2^2+
(\mu-\Gamma)\|u\|_2^2$, co jest co najmniej
$(\lambda/2)\|u\|_V^2$ przy podanym wyborze $\mu$.
\end{proof}

\begin{twierdzenie}[Istnienie dla przesuniętego operatora]
\label{thm:pde13-shift}
Przy założeniach twierdzenia~\ref{thm:pde13-garding} i dla
$\mu\geq\Gamma+\lambda/2$ dla każdego $F\in V^*$ istnieje
dokładnie jedno $u\in V$ spełniające
$a(u,v)+\mu(u,v)_{L^2}=\langle F,v\rangle$ dla wszystkich $v\in V$.
Ponadto $\|u\|_V\leq(2/\lambda)\|F\|_{V^*}$.
\end{twierdzenie}
\begin{proof}
Forma $a_\mu$ jest ograniczona ze stałą $M+\mu$, gdyż
$|(u,v)_{L^2}|\le\|u\|_V\|v\|_V$ oraz $\mu>0$.
Z poprzedniego twierdzenia jest koercywna ze stałą $\lambda/2$.
Przestrzeń $V$ jest rzeczywistą przestrzenią Hilberta, więc
twierdzenie Laxa--Milgrama~\ref{thm:pde11-LM}, które nie wymaga
symetrii formy, daje istnienie i jednoznaczność. Podstawiając $v=u$ do równania słabego,
uzyskujemy $(\lambda/2)\|u\|_V^2\leq
\langle F,u\rangle\leq\|F\|_{V^*}\|u\|_V$.
Jeżeli $u=0$, oszacowanie jest trywialne; w przeciwnym razie
dzielimy przez $\|u\|_V$.
\end{proof}

\begin{uwaga}[Dlaczego samo przesunięcie nie rozwiązuje zagadnienia?]
Wynik dotyczy operatora przesuniętego $L+\mu J:V\to V^*$,
gdzie $\langle Ju,v\rangle=(u,v)_{L^2}$, a niekoniecznie $L$.
W realizacji w $L^2$ odpowiada to przesunięciu $L+\mu I$.
Przykładowo na $(0,1)$ dla $L=-\mathrm d^2/\mathrm dx^2-\pi^2$
funkcja $\phi(x)=\sin(\pi x)$ należy do $H_0^1(0,1)$
i spełnia $L\phi=0$. Dla każdego $u\in H_0^1(0,1)$
całkowanie przez części (najpierw dla testów $u$, potem przez
gęstość) daje $a(u,\phi)=0$.
Gdyby istniało słabe rozwiązanie z prawą stroną $f=\phi$,
test $v=\phi$ dawałby $0=\|\phi\|_2^2>0$, sprzeczność.
Samo przesunięcie nie zapewnia więc rozwiązywalności
pierwotnego równania.
\end{uwaga}


\begin{uwaga}[Trzy różne wnioski z oszacowań formy]
\label{rem:pde-v42-garding}
Koercywność z twierdzenia~\ref{thm:pde11-lower-order}
pozwala bezpośrednio zastosować twierdzenie Laxa--Milgrama.
Sama nierówność Gårdinga~\eqref{eq:pde13-garding}
nie daje tego wniosku dla operatora $L$, lecz zapewnia
koercywność odpowiednio przesuniętego operatora $L+\mu I$.
Dopiero argument Fredholma w sekcji~\ref{sec:pde13-fredholm}
pozwoli wrócić do $L$, z możliwym warunkiem ortogonalności
do jądra operatora sprzężonego. Tych trzech kroków
nie należy utożsamiać.
\end{uwaga}

\section{Dirichlet: jednoznaczność, dane brzegowe i stabilność}
\label{sec:pde13-dirichlet}
\begin{twierdzenie}[Koercywność bez przesunięcia]
\label{thm:pde13-coercive}
Załóżmy, że $b=0$ i $c\geq0$ prawie wszędzie, a macierz $A$
spełnia~\eqref{eq:pde13-elliptic}. Wówczas dla każdego $F\in V^*$
zagadnienie~\eqref{eq:pde13-weak} ma dokładnie jedno rozwiązanie,
a $\|u\|_V\leq(1+C_P^2)\|F\|_{V^*}/\lambda$, gdzie
$C_P$ jest stałą nierówności Poincarégo
$\|v\|_2\leq C_P\|\nabla v\|_2$ na $V$.
Przy ogólnych $b,c$, jeśli $a(v,v)\geq\alpha\|v\|_V^2$
dla wszystkich $v\in V$ i pewnego $\alpha>0$, również istnieje
dokładnie jedno rozwiązanie, lecz oszacowanie ma wtedy postać
$\|u\|_V\leq\alpha^{-1}\|F\|_{V^*}$.
W szczególności, gdy $\Gamma C_P^2<\lambda/2$,
można przyjąć
\[
 \alpha=\frac{\lambda/2-\Gamma C_P^2}{1+C_P^2}.
\]
\end{twierdzenie}
\begin{proof}
Z wniosku~\ref{cor:pde10-Poincare} mamy
$\|v\|_V^2\leq(1+C_P^2)\|\nabla v\|_2^2$.
Dla $b=0$, $c\geq0$ eliptyczność daje
$a(v,v)\geq\lambda\|\nabla v\|_2^2
\geq\lambda(1+C_P^2)^{-1}\|v\|_V^2$.
Forma jest ograniczona na rzeczywistej przestrzeni Hilberta $V$,
więc twierdzenie Laxa--Milgrama~\ref{thm:pde11-LM} daje
istnienie i jednoznaczność.
Test $v=u$ prowadzi do
\[
 \frac{\lambda}{1+C_P^2}\|u\|_V^2
 \leq a(u,u)=\langle F,u\rangle
 \leq\|F\|_{V^*}\|u\|_V.
\]
Jeśli $u\ne0$, dzielimy przez $\|u\|_V$; dla $u=0$
oszacowanie zachodzi bez dzielenia.
Przy ogólnej koercywności stosujemy to samo twierdzenie
ze stałą $\alpha$, co daje oszacowanie $\alpha^{-1}\|F\|_{V^*}$.
W szczególnym przypadku z~\eqref{eq:pde13-garding}
i Poincarégo mamy
$a(v,v)\geq(\lambda/2-\Gamma C_P^2)\|\nabla v\|_2^2$.
Gdy współczynnik w nawiasie jest dodatni, ponowne użycie
Poincarégo daje wskazaną w tezie stałą $\alpha$.
\end{proof}

\begin{twierdzenie}[Niejednorodny warunek Dirichleta]
\label{thm:pde13-nonhom}
Niech $A,b,c$ spełniają założenia sekcji~\ref{sec:pde13-operators},
a forma $a$ będzie koercywna: $a(v,v)\geq\alpha\|v\|_V^2$
dla $v\in V$, gdzie $\alpha>0$.
Dla każdego $F\in V^*$ i $g\in H^{1/2}(\partial\Omega)$
istnieje dokładnie jedna funkcja $u\in H^1(\Omega)$ o śladzie
$\operatorname{Tr}u=g$ taka, że
$a(u,v)=\langle F,v\rangle$ dla wszystkich $v\in V$.
Ponadto
\[
 \|u\|_{H^1(\Omega)}\leq C\big(\|F\|_{V^*}
 +\|g\|_{H^{1/2}(\partial\Omega)}\big),
\]
gdzie $C$ zależy od $\alpha$, ciągłości formy i operatora
przedłużania śladu, ale nie od $F,g$.
\end{twierdzenie}
\begin{proof}
Definicja~\ref{def:pde10-halftrace} i lemat~\ref{lem:pde10-halftrace-rightinverse}
zapewniają ciągły liniowy operator
$E_{\operatorname{Tr}}:H^{1/2}(\partial\Omega)\to H^1(\Omega)$,
prawostronnie odwrotny do operatora śladu.
Kładziemy $G=E_{\operatorname{Tr}}g$; wtedy $\operatorname{Tr}G=g$
oraz $\|G\|_{H^1}\leq C_{\operatorname{Tr}}\|g\|_{H^{1/2}}$
(dla normy ilorazowej przyjętej w rozdziale 10 można wziąć
$C_{\operatorname{Tr}}=1$).
Z twierdzenia~\ref{thm:pde10-trace-zero} zbiór funkcji
o śladzie $g$ jest dokładnie $G+V$: jeśli $\operatorname{Tr}u=g$,
to $\operatorname{Tr}(u-G)=0$, a odwrotnie dodanie elementu $V$
nie zmienia śladu.
Szukamy zatem $u=G+w$, gdzie $w\in V$.
Równanie dla $w$ przyjmuje postać
$a(w,v)=\langle F,v\rangle-a(G,v)$.
Prawa strona jest ciągłym funkcjonałem na $V$ o normie
co najwyżej $\|F\|_{V^*}+M\|G\|_{H^1}$.
Z koercywności i twierdzenia Laxa--Milgrama istnieje
jedno $w$ i
$\|w\|_V\leq\alpha^{-1}
(\|F\|_{V^*}+M\|G\|_{H^1})$.
Funkcja $u=G+w$ ma ślad $g$ i przez dwuliniowość formy
spełnia równanie z tezy. Z nierówności trójkąta
\[
 \|u\|_{H^1}
 \leq \alpha^{-1}\|F\|_{V^*}
 +(1+M/\alpha)C_{\operatorname{Tr}}\|g\|_{H^{1/2}}.
\]
Można więc przyjąć
$C=\max\{\alpha^{-1},(1+M/\alpha)C_{\operatorname{Tr}}\}$.
Jeśli $u_1,u_2$ są rozwiązaniami o tych samych danych, to
$z=u_1-u_2\in V$ z twierdzenia o jądrze śladu.
Odejmując równania i testując $v=z$, dostajemy
$0=a(z,z)\geq\alpha\|z\|_V^2$, więc $z=0$.
W szczególności rozwiązanie nie zależy od wyboru funkcji $G$.
Dla dwóch różnych par danych $(F_1,g_1)$ i $(F_2,g_2)$
różnica rozwiązań spełnia to samo zagadnienie z danymi
$(F_1-F_2,g_1-g_2)$; powyższe oszacowanie daje zatem
ciągłą zależność od obu danych.
\end{proof}

\begin{lemat}[Ślad dodatniej i ujemnej części funkcji]
\label{lem:pde13-trace-parts}
Dla rzeczywistej funkcji $w\in H^1(\Omega)$ zachodzi
\[
 \operatorname{Tr}(w^+)=(\operatorname{Tr}w)^+,\qquad
 \operatorname{Tr}(w^-)=(\operatorname{Tr}w)^-
\]
w $L^2(\partial\Omega)$.
\end{lemat}
\begin{proof}
Z wniosku~\ref{cor:pde9-truncation} obie części należą do $H^1$.
Wybierzmy $w_n$ gładkie w otoczeniu $\overline\Omega$,
$w_n\to w$ w $H^1$ (wniosek~\ref{cor:pde10-density-boundary}).
Niech $\varepsilon_n=1/n$ i
$P_\varepsilon(t)=(t+\sqrt{t^2+\varepsilon^2}-\varepsilon)/2$.
Wtedy $P_\varepsilon(0)=0$, $0\le P_\varepsilon'\le1$
i $|P_\varepsilon(t)-t^+|\le\varepsilon/2$.
Funkcje $z_n=P_{\varepsilon_n}(w_n)$ są gładkie w otoczeniu
domknięcia obszaru i ograniczone w $H^1$, gdyż
$|P_\varepsilon(t)|\le|t|$ oraz
$|\nabla z_n|\le|\nabla w_n|$.
Z nierówności Lipschitzowskiej dla $t^+$ mamy
\[
 \|z_n-w^+\|_2
 \le \tfrac{\varepsilon_n}{2}\lambda_d(\Omega)^{1/2}
 +\|w_n-w\|_2\longrightarrow0.
\]
Ze słabej zwartości kul (lemat~\ref{lem:pde11-weakcompact})
wybieramy podciąg słabo zbieżny w $H^1$; silna zbieżność
w $L^2$ identyfikuje jego granicę jako $w^+$.
Ciągłość operatora śladu (twierdzenie~\ref{thm:pde10-trace})
daje słabą zbieżność śladów tego podciągu do
$\operatorname{Tr}(w^+)$.
Jednocześnie z gładkości mamy
$\operatorname{Tr}z_n=P_{\varepsilon_n}(\operatorname{Tr}w_n)$.
Ponieważ miara powierzchniowa brzegu ograniczonego obszaru
Lipschitzowskiego jest skończona, oszacowanie błędu
$P_{\varepsilon_n}$ i silna zbieżność
$\operatorname{Tr}w_n\to\operatorname{Tr}w$ w $L^2$
dają silną zbieżność $\operatorname{Tr}z_n$ do
$(\operatorname{Tr}w)^+$. Jednoznaczność granicy słabej
dowodzi pierwszej tożsamości. Drugą uzyskujemy dla $-w$.
\end{proof}

\begin{twierdzenie}[Słaba zasada minimum bez członu transportowego]
\label{thm:pde13-positivity}
Niech $b=0$, $c\geq0$ i $A$ spełnia~\eqref{eq:pde13-elliptic};
nie wymagamy symetrii macierzy $A$. Niech $u\in H^1(\Omega)$
spełnia $a(u,v)=\langle F,v\rangle$ dla wszystkich $v\in V$,
przy czym $\operatorname{Tr}u\geq0$ na $\partial\Omega$
(w sensie prawie wszędzie) i
$\langle F,v\rangle\geq0$ dla wszystkich nieujemnych $v\in V$.
Wtedy $u\geq0$ prawie wszędzie w $\Omega$.
\end{twierdzenie}
\begin{proof}
Ujemna część $u^-=\max\{-u,0\}$ należy do $H^1(\Omega)$.
Z wniosku~\ref{cor:pde9-truncation} mamy
$\nabla u^-=-\chi_{\{u<0\}}\nabla u$ prawie wszędzie.
Ponieważ ślad $u$ jest nieujemny, lemat~\ref{lem:pde13-trace-parts}
daje $\operatorname{Tr}(u^-)=0$; z twierdzenia
~\ref{thm:pde10-trace-zero} wynika $u^-\in V$.
Testujemy równanie $v=u^-$. Na zbiorze $\{u<0\}$ mamy
$u=-u^-$, $\nabla u=-\nabla u^-$, a poza nim
$u^-=\nabla u^-=0$ prawie wszędzie. Zatem
\[
 0\leq\langle F,u^-\rangle=a(u,u^-)
 =-\int_\Omega \big(A\nabla u^-\cdot\nabla u^-+c(u^-)^2\big)
 \,\mathrm d\lambda_d(x)
 \leq-\lambda\|\nabla u^-\|_2^2.
\]
Ponieważ $\lambda>0$, otrzymujemy $\|\nabla u^-\|_2=0$,
a nierówność Poincarégo na $V$ (wniosek~\ref{cor:pde10-Poincare})
daje $\|u^-\|_2=0$. Z definicji ujemnej części jest to
równoważne $u\geq0$ prawie wszędzie.
\end{proof}

\begin{twierdzenie}[Zasada porównania z dryfem o kontrolowanej dywergencji]
\label{thm:pde13-drift-comparison}
Niech $\Omega\subset\R^d$ będzie ograniczonym obszarem
Lipschitzowskim, $A\in L^\infty(\Omega;\R^{d\times d})$
spełnia~\eqref{eq:pde13-elliptic}, $b\in W^{1,\infty}(\Omega;\R^d)$,
$c\in L^\infty(\Omega)$ oraz
$c-\frac12\diver b\ge0$ prawie wszędzie. Jeżeli
$u_1,u_2\in H^1(\Omega)$ mają ślady spełniające
$\operatorname{Tr}u_1\le\operatorname{Tr}u_2$ prawie wszędzie
na $\partial\Omega$ oraz
\[
 a(u_1,\varphi)\le a(u_2,\varphi)
 \quad\text{dla wszystkich }\varphi\in H_0^1(\Omega),\ \varphi\ge0,
\]
to $u_1\le u_2$ prawie wszędzie w $\Omega$.
W szczególności założenie $\diver b=0$, $c\ge0$ zapewnia
porównanie niezależnie od wielkości normy $b$.
\end{twierdzenie}
\begin{proof}
Kładziemy $w=u_1-u_2$ i $v=w^+$.
Wniosek~\ref{cor:pde9-truncation} daje $v\in H^1$
oraz $\nabla v=\chi_{\{w>0\}}\nabla w$.
Z liniowości śladu i założenia $\operatorname{Tr}w\le0$,
więc lemat~\ref{lem:pde13-trace-parts} daje $\operatorname{Tr}v=0$.
Stąd $v\in H_0^1(\Omega)$ na mocy twierdzenia o jądrze śladu.
Ponieważ $v\ge0$, wolno wybrać $\varphi=v$, co daje
$a(w,v)\le0$.
Człon pierwszego rzędu obliczamy bez założenia małości $b$:
\[
 \int_\Omega (b\cdot\nabla w)v\,\mathrm d\lambda_d(x)
 =\frac12\int_\Omega b\cdot\nabla(v^2)\,\mathrm d\lambda_d(x)
 =-\frac12\int_\Omega(\diver b)v^2\,\mathrm d\lambda_d(x).
\]
Pierwsza równość wynika z $\nabla w=\nabla v$ na $\{w>0\}$
i $v=0$ poza tym zbiorem oraz z reguły Leibniza
(twierdzenie~\ref{thm:pde9-products}): $v^2\in W^{1,1}$
i $\nabla(v^2)=2v\nabla v$.
Dla ostatniej równości wybieramy $v_n\in C_c^\infty(\Omega)$
zbieżne do $v$ w $H^1$. Mamy
\begin{align*}
 \|v_n^2-v^2\|_1
 &\le\|v_n-v\|_2\|v_n+v\|_2,\\
 \|\nabla(v_n^2)-\nabla(v^2)\|_1
 &\le 2\|v_n-v\|_2\|\nabla v_n\|_2
 +2\|v\|_2\|\nabla v_n-\nabla v\|_2.
\end{align*}
Obie prawe strony dążą do zera, zatem $v_n^2\to v^2$
w $W^{1,1}$. Dla testu $v_n^2$ definicja słabej dywergencji
daje
$\int_\Omega b\cdot\nabla(v_n^2)\,\mathrm d\lambda_d(x)
=-\int_\Omega(\diver b)v_n^2\,\mathrm d\lambda_d(x)$.
Ponieważ $b,\diver b$ są ograniczone, możemy przejść
do granicy w obu całkach. W ten sposób uzyskujemy równość
bez członu brzegowego, nie zakładając $v^2\in H^1$.
Wobec tego
\[
 0\ge a(w,v)=\int_\Omega\left(A\nabla v\cdot\nabla v+
 (c-\tfrac12\diver b)v^2\right)\,\mathrm d\lambda_d(x)
 \ge\lambda\|\nabla v\|_2^2.
\]
Nierówność Poincarégo daje $\|v\|_2=0$, czyli
$(u_1-u_2)^+=0$ prawie wszędzie. Otrzymujemy żądane porównanie.
Jest to argument energetyczny, a nie twierdzenie o punktowej
regularności rozwiązań. W szczególności dla dwóch rozwiązań
o tych samych danych zastosowanie porównania w obu kierunkach
daje jednoznaczność.

Przy tych samych założeniach znakowych forma jest również
koercywna na $V$. Istotnie, dla dowolnego $z\in V$ argument
aproksymacji kwadratów zastosowany do $z$ daje
\[
 a(z,z)=\int_\Omega
 \big(A\nabla z\cdot\nabla z+(c-\tfrac12\diver b)z^2\big)
 \,\mathrm d\lambda_d(x)
 \ge\lambda\|\nabla z\|_2^2.
\]
Z Poincarégo otrzymujemy koercywność ze stałą
$\lambda/(1+C_P^2)$, bez warunku małości $b$.
Twierdzenia~\ref{thm:pde13-coercive} i~\ref{thm:pde13-nonhom}
zapewniają więc także istnienie przy jednorodnych i niejednorodnych
danych Dirichleta.
\end{proof}



\section{Alternatywa Fredholma dla eliptycznych zagadnień Dirichleta}
\label{sec:pde13-fredholm}
Dla współczynników dolnego rzędu bez warunku znakowego forma $a$
może nie być koercywna. Mimo to operator $L:V\to V^*$,
$\langle Lu,v\rangle=a(u,v)$, ma precyzyjną własność Fredholma.
Nie należy go mylić z realizacją $L$ w $L^2$, której dziedzina
wymaga dodatkowej definicji.

\begin{twierdzenie}[Alternatywa Fredholma dla $L:V\to V^*$]
\label{thm:pde13-fredholm}
Niech $\Omega\subset\R^d$ będzie obszarem ograniczonym Lipschitzowskim, a $A,b,c\in L^\infty$ niech spełniają
jednostajną eliptyczność~\eqref{eq:pde13-elliptic}.
Wówczas $L:V\to V^*$ jest operatorem Fredholma indeksu $0$:
ma domknięty obraz, skończeniewymiarowe jądro i skończony
kowymiar obrazu, równy wymiarowi jądra.
Definiując operator transponowany $L^{\mathrm t}:V\to V^*$ przez
$\langle L^{\mathrm t}z,v\rangle=a(v,z)$, mamy
$\dim\ker L=\dim\ker L^{\mathrm t}<\infty$.
Równanie $Lu=F$, $F\in V^*$, ma rozwiązanie wtedy i tylko wtedy, gdy
\begin{equation}\label{eq:pde13-compat}
 \langle F,z\rangle=0\qquad(z\in\ker L^{\mathrm t}).
\end{equation}
Rozwiązania różnią się o elementy $\ker L$. W szczególności
gdy $\ker L=\{0\}$, dla każdego $F\in V^*$ istnieje jedno rozwiązanie.
\end{twierdzenie}
\begin{proof}
\textbf{Krok 1: odwracalny operator przesunięty.}
Wybieramy $\mu\geq\Gamma+\lambda/2$.
Z twierdzenia~\ref{thm:pde13-shift} operator
$L_\mu=L+\mu J:V\to V^*$ jest izomorfizmem ograniczonym,
przy czym $J:V\to V^*$ definiujemy wzorem
$\langle Ju,v\rangle=\int_\Omega uv\,\mathrm d\lambda_d(x)$.

\textbf{Krok 2: zwartość zaburzenia.}
Zanurzenie $i:V\hookrightarrow L^2(\Omega)$ jest zwarte
na mocy twierdzenia Rellicha--Kondrachova. Operator
$j:L^2(\Omega)\to V^*$, $\langle jf,v\rangle=(f,v)_{L^2}$,
jest ograniczony, ponieważ $\|v\|_2\leq\|v\|_V$.
Zatem $J=j\circ i$ jest zwarty jako złożenie operatora
zwartego z ograniczonym. W konsekwencji operator
$K=\mu L_\mu^{-1}J:V\to V$ jest zwarty.

\textbf{Krok 3: sprowadzenie do alternatywy dla operatora zwartego.}
Mamy równość operatorów
$L=L_\mu(I-K)$; istotnie
$L_\mu K=\mu J$.
Twierdzenie~\ref{thm:pde12-general-fredholm} zostało
udowodnione dla zespolonych przestrzeni Hilberta;
jego dowód oparty na aproksymacji operatorami skończonego rzędu
pozostaje bez zmian prawdziwy w rzeczywistych przestrzeniach
Hilberta. Daje to, że $I-K:V\to V$ ma domknięty obraz,
skończeniewymiarowe jądro i kowymiar równy wymiarowi jądra.
Ponieważ $L_\mu$ jest izomorfizmem, własności te przechodzą
na $L:V\to V^*$, a zatem $\operatorname{ind}L=0$.

\textbf{Krok 4: dokładny warunek zgodności.}
Dla $z\in V$ warunek $z\in\ker L^{\mathrm t}$ oznacza
$a(v,z)=0$ dla wszystkich $v\in V$. To jest równoważne temu,
że ciągły funkcjonał na $V^*$ zadany przez
$F\mapsto\langle F,z\rangle$ znika na $\operatorname{Ran}L$.
Każdy ciągły funkcjonał na $V^*$ ma tę postać, ponieważ
$V$ jest refleksywna (jako przestrzeń Hilberta).
Przestrzeń funkcjonałów zerujących obraz $L$ ma zatem
wymiar $\dim\ker L^{\mathrm t}$.
Ponieważ obraz jest domknięty, twierdzenie Hahna--Banacha
lub identyfikacja przestrzeni ilorazowej daje
$\operatorname{Ran}L=\{F\in V^*: \langle F,z\rangle=0
\text{ dla każdego } z\in\ker L^{\mathrm t}\}$.
Kowymiar obrazu jest równy $\dim\ker L^{\mathrm t}$,
a z kroku 3 także $\dim\ker L$. Wreszcie różnica rozwiązań
jednego równania należy do jądra $L$.
\end{proof}

\begin{przyklad}[Zgodność przy rezonansie]
\label{ex:pde13-resonance}
Dla $\Omega=(0,1)$ i $L=-\mathrm d^2/\mathrm dx^2-\pi^2$
z warunkiem Dirichleta mamy
$\ker L=\ker L^{\mathrm t}=\operatorname{span}\{\sin(\pi x)\}$.
Równanie $-u''-\pi^2u=f$, $u(0)=u(1)=0$,
ma słabe rozwiązanie dokładnie wtedy, gdy
$\int_0^1 f(x)\sin(\pi x)\,\mathrm d\lambda_1(x)=0$.
Gdy ten warunek zachodzi, rozwiązanie jest wyznaczone
z dokładnością do dodania wielokrotności $\sin(\pi x)$.
W szczególności dla $f=\sin(\pi x)$ rozwiązanie nie istnieje.
\end{przyklad}

\section{Zagadnienie Neumanna: warunek zgodności i stałe}
\label{sec:pde13-neumann}
Dla $b=c=0$ i jednorodnej eliptyczności rozważmy
$-\diver(A\nabla u)=f$ wraz z zadanym strumieniem na brzegu.
Przy słabych współczynnikach nie musimy definiować strumienia
punktowo; dane brzegowe można wprowadzić przez funkcjonał śladu.

\begin{lemat}[Nierówność Poincarégo--Wirtingera]
\label{lem:pde13-wirtinger}
Jeżeli $\Omega$ jest obszarem ograniczonym Lipschitzowskim,
to istnieje $C_W>0$ takie, że dla każdego $u\in H^1(\Omega)$
\[
\|u-u_\Omega\|_{L^2(\Omega)}\le C_W\|\nabla u\|_{L^2(\Omega)},
\qquad u_\Omega=\frac1{\lambda_d(\Omega)}
\int_\Omega u\,\mathrm d\lambda_d(x).
\]
\end{lemat}
\begin{proof}
Gdyby twierdzenie było fałszywe, znaleźlibyśmy $v_n\in H^1$
o średniej zero takie, że $\|v_n\|_2=1$ i
$\|\nabla v_n\|_2\to0$ (odejmujemy średnią i normalizujemy).
Ciąg $v_n$ jest ograniczony w $H^1$, więc dzięki
zwartości $H^1(\Omega)\hookrightarrow L^2(\Omega)$
ma podciąg zbieżny silnie w $L^2$ do $v$. Po wybraniu
podciągu słabo zbieżnego w $H^1$ otrzymujemy
$\nabla v=0$ w sensie dystrybucyjnym.
Ponieważ $\Omega$ jest spójna, $v$ jest prawie wszędzie stała.
Ciągłość funkcjonału $w\mapsto\int_\Omega w$ oraz średnie
zerowe pokazują, że ta stała wynosi zero.
Ale silna zbieżność w $L^2$ wymusza $\|v\|_2=1$,
co jest sprzecznością.
\end{proof}

\begin{twierdzenie}[Neumann: istnienie modulo stałe]
\label{thm:pde13-neumann}
Niech $\Omega$ będzie obszarem ograniczonym Lipschitzowskim, a $A\in L^\infty(\Omega;\R^{d\times d})$
spełnia~\eqref{eq:pde13-elliptic}. Niech
$f\in L^2(\Omega)$ i $h\in L^2(\partial\Omega)$,
przy czym na brzegu całkujemy względem miary powierzchniowej
$\sigma=\mathcal H^{d-1}|_{\partial\Omega}$.
Istnieje $u\in H^1(\Omega)$ spełniające
\begin{equation}\label{eq:pde13-neumann}
\int_\Omega A\nabla u\cdot\nabla v\,\mathrm d\lambda_d(x)
=\int_\Omega fv\,\mathrm d\lambda_d(x)
+\int_{\partial\Omega} h\,\operatorname{Tr}v\,\mathrm d\sigma
\quad(v\in H^1(\Omega))
\end{equation}
wtedy i tylko wtedy, gdy
\begin{equation}\label{eq:pde13-neumann-comp}
\int_\Omega f\,\mathrm d\lambda_d(x)
+\int_{\partial\Omega}h\,\mathrm d\sigma=0.
\end{equation}
Rozwiązanie jest jednoznaczne modulo stałe. Istnieje dokładnie
jedno rozwiązanie o średniej zero; spełnia ono
$\|u\|_{H^1}\leq C(\|f\|_2+\|h\|_{L^2(\partial\Omega)})$.
\end{twierdzenie}
\begin{proof}
\textbf{Konieczność.} Podstawiamy $v=1$ do
\eqref{eq:pde13-neumann}; gradient stałej znika i
otrzymujemy dokładnie~\eqref{eq:pde13-neumann-comp}.

\textbf{Wystarczalność.} Rozważmy domkniętą podprzestrzeń
$V_\diamond=\{v\in H^1(\Omega):v_\Omega=0\}$.
Na niej norma $\|\nabla v\|_2$ jest równoważna normie
$H^1$ na mocy lematu~\ref{lem:pde13-wirtinger}.
Forma $B(u,v)=\int_\Omega A\nabla u\cdot\nabla v
\,\mathrm d\lambda_d(x)$ jest ciągła oraz koercywna
względem tej normy, ponieważ
$B(v,v)\geq\lambda\|\nabla v\|_2^2$.
Funkcjonał po prawej stronie~\eqref{eq:pde13-neumann}
jest ciągły na $V_\diamond$: przez nierówność Höldera
oraz ciągłość śladu jego moduł jest nie większy niż
$C(\|f\|_2+\|h\|_{L^2(\partial\Omega)})\|v\|_{H^1}$.
Z twierdzenia Laxa--Milgrama istnieje jedno $u\in V_\diamond$
spełniające równanie dla wszystkich $v\in V_\diamond$.
Dla dowolnego $v\in H^1$ piszemy $v=(v-v_\Omega)+v_\Omega$.
Pierwsza część jest w $V_\diamond$, a prawa strona na stałej
znika dzięki zgodności~\eqref{eq:pde13-neumann-comp}.
Równanie jest zatem prawdziwe dla całego $H^1$.
Oszacowanie wynika z Laxa--Milgrama i lematu Wirtingera.
Jeśli $u_1,u_2$ są rozwiązaniami, ich różnica $w$ spełnia
$B(w,v)=0$ dla wszystkich $v\in H^1$. Biorąc $v=w$,
otrzymujemy $\nabla w=0$, a spójność $\Omega$ implikuje,
że $w$ jest prawie wszędzie stała.
\end{proof}

\begin{uwaga}[Znak strumienia]
Dla gładkiego $u$ i $A$ całkowanie przez części daje
$\int_\Omega A\nabla u\cdot\nabla v
=\int_\Omega[-\diver(A\nabla u)]v
+\int_{\partial\Omega}(A\nabla u\cdot\nu)v$.
Dlatego $h$ w~\eqref{eq:pde13-neumann} oznacza
\emph{zewnętrzny} strumień $A\nabla u\cdot\nu$.
Przy danych szorstkich jest to jedynie interpretacja
funkcjonałowa; nie zakładamy istnienia punktowego śladu
wektora $A\nabla u$.
\end{uwaga}

\section{Robin: człon brzegowy przywracający jednoznaczność}
\label{sec:pde13-robin}
\begin{lemat}[Nierówność Poincarégo z wyrazem brzegowym]
\label{lem:pde13-boundary-poincare}
Na obszarze ograniczonym Lipschitzowskim istnieje
$C>0$ takie, że dla wszystkich $u\in H^1(\Omega)$
\[
\|u\|_{H^1}^2\le C\big(\|\nabla u\|_2^2+
 \|\operatorname{Tr}u\|_{L^2(\partial\Omega)}^2\big).
\]
\end{lemat}
\begin{proof}
Gdyby oszacowanie nie zachodziło, mielibyśmy ciąg
$\|u_n\|_{H^1}=1$ i $\|\nabla u_n\|_2^2+
\|\operatorname{Tr}u_n\|_{L^2(\partial\Omega)}^2\to0$.
Ze zwartości zanurzenia $H^1\hookrightarrow L^2$ i słabej
zwartości kuli w $H^1$ otrzymujemy podciąg słabo zbieżny
w $H^1$, silnie w $L^2$ do pewnego $u$.
Ponieważ $\nabla u_n\to0$ w $L^2$, mamy $\nabla u=0$,
a więc $u$ jest stała. Z ciągłości śladu $H^1\to
L^2(\partial\Omega)$ otrzymujemy słabą zbieżność śladów
do $\operatorname{Tr}u$, lecz ślady dążą silnie do zera.
Zatem stała $u$ ma zerowy ślad, czyli $u=0$.
Wówczas $\|u_n\|_2\to0$ oraz $\|\nabla u_n\|_2\to0$,
wbrew $\|u_n\|_{H^1}=1$.
\end{proof}

\begin{twierdzenie}[Zagadnienie Robina]
\label{thm:pde13-robin}
Niech $\Omega$ będzie obszarem ograniczonym Lipschitzowskim, $A\in L^\infty$ spełnia
\eqref{eq:pde13-elliptic}, a $\beta\in L^\infty(\partial\Omega)$
spełnia $\beta\geq\beta_0>0$ prawie wszędzie.
Dla każdego $f\in L^2(\Omega)$ i $h\in L^2(\partial\Omega)$
istnieje dokładnie jedno $u\in H^1(\Omega)$ spełniające
\begin{equation}\label{eq:pde13-robin}
\int_\Omega A\nabla u\cdot\nabla v\,\mathrm d\lambda_d(x)
+\int_{\partial\Omega}\beta\,\operatorname{Tr}u\,
\operatorname{Tr}v\,\mathrm d\sigma
=\int_\Omega fv\,\mathrm d\lambda_d(x)
+\int_{\partial\Omega}h\,\operatorname{Tr}v\,\mathrm d\sigma
\end{equation}
dla wszystkich $v\in H^1(\Omega)$, a ponadto
$\|u\|_{H^1}\leq C(\|f\|_2+\|h\|_{L^2(\partial\Omega)})$.
Formalnie warunek brzegowy ma postać
$A\nabla u\cdot\nu+\beta u=h$.
\end{twierdzenie}
\begin{proof}
Ciągłość formy w lewej stronie wynika z ograniczoności $A,\beta$
oraz ciągłości operatora śladu $H^1\to L^2(\partial\Omega)$.
Jej wartość dla $v=u$ jest nie mniejsza niż
$\lambda\|\nabla u\|_2^2+
\beta_0\|\operatorname{Tr}u\|_{L^2(\partial\Omega)}^2$,
a z lematu~\ref{lem:pde13-boundary-poincare} wynika koercywność
na całym $H^1$. Prawa strona jest ciągłym funkcjonałem
na $H^1$ przez Höldera i twierdzenie o śladzie.
Wszystkie założenia Laxa--Milgrama są więc spełnione;
otrzymujemy istnienie, jednoznaczność i oszacowanie.
\end{proof}

\section{Przykład obliczeniowy: strumień przez granicę ośrodków}
\label{sec:pde13-transmission}
Niech $\Omega=(-1,1)$, $f=1$ i
$a(x)=1$ dla $x<0$, $a(x)=4$ dla $x>0$.
Rozważamy $-(a u')'=1$ wraz z $u(-1)=u(1)=0$.
Na obu podprzedziałach $u''=-1/a$, a zatem
\begin{equation}\label{eq:pde13-exact}
 u(x)=\begin{cases}
 -x^2/2-3x/10+1/5,& -1\le x\le0,\\
 -x^2/8-3x/40+1/5,& 0\le x\le1.
 \end{cases}
\end{equation}
\begin{center}
\begin{tikzpicture}[x=3.1cm,y=0.7cm,>=Stealth]
  \fill[blue!8] (-1,0) rectangle (0,1.35);
  \fill[orange!13] (0,0) rectangle (1,1.35);
  \draw[thick] (-1,0) rectangle (1,1.35);
  \draw[thick,dashed] (0,0)--(0,1.35);
  \node at (-.52,.87) {$a_-=1$};
  \node at (.52,.87) {$a_+=4$};
  \node at (-.52,.38) {ośrodek I};
  \node at (.52,.38) {ośrodek II};
  \draw[->] (-1.15,-.24)--(1.18,-.24) node[right] {$x$};
  \foreach \x/\lab in {-1/-1,0/0,1/1}{
    \draw (\x,-.31)--(\x,-.17);\node[below] at (\x,-.31) {$\lab$};}
  \node[above] at (0,1.35) {interfejs};
\end{tikzpicture}

\small\emph{Dwa ośrodki o różnych współczynnikach: ciągłe $u$ i $a u'$, lecz skok $u'$.}
\end{center}

Wartość $u(0^-)=u(0^+)=1/5$ jest ciągła, natomiast
$u'(0^-)=-3/10$ i $u'(0^+)=-3/40$. Mimo skoku pochodnej
strumień $a(x)u'(x)=-x-3/10$ jest ciągły w $0$.

\begin{stwierdzenie}[Warunki transmisyjne z równania słabego]
\label{prop:pde13-transmission}
Niech $a$ jest dodatnią stałą $a_-$ na $(-1,0)$ i
$a_+$ na $(0,1)$, $f\in L^2(-1,1)$, a $u\in H^1(-1,1)$
ma na obu podprzedziałach regularność $H^2$.
Wówczas $-(au')'=f$ w $\mathcal D'(-1,1)$ wtedy i tylko wtedy,
gdy $-a_\pm u''=f$ prawie wszędzie po odpowiednich stronach,
a ślady strumienia spełniają
$a_-u'(0^-)=a_+u'(0^+)$.
\end{stwierdzenie}
\begin{proof}
Z ciągłości funkcji z $H^1(-1,1)$ w wymiarze jednym
wynika $u(0^-)=u(0^+)$. Testując równanie funkcjami o nośniku
w każdym z podprzedziałów, dostajemy tam równania klasyczne
prawie wszędzie. Dla dowolnej $\varphi\in C_c^\infty(-1,1)$
całkujemy osobno po obu przedziałach:
\[
\int_{-1}^1 au'\varphi'\,\mathrm d\lambda_1
=\int_{-1}^1 f\varphi\,\mathrm d\lambda_1
 +\big(a_-u'(0^-)-a_+u'(0^+)\big)\varphi(0).
\]
Zatem równość dystrybucyjna na całym przedziale jest
równoważna znikaniu współczynnika przy $\varphi(0)$.
W drugą stronę te same rachunki, przy zerowym skoku strumienia,
dają równanie globalne. Dla~\eqref{eq:pde13-exact} skok strumienia
wynosi zero, co potwierdza poprawność rozwiązania.
\end{proof}

\section{Zakres otrzymanych rezultatów}
\label{sec:pde13-summary}
Jednostajna eliptyczność zapewnia kontrolę części gradientowej,
lecz człony dolnego rzędu mogą zniszczyć koercywność.
Nierówność Gårdinga pozwala odzyskać ją po przesunięciu,
a zwartość zanurzenia Sobolewa prowadzi do alternatywy Fredholma
dla problemu pierwotnego. W zagadnieniu Neumanna przeszkodą
jest przestrzeń stałych, zaś dodatni człon brzegowy w problemie
Robina usuwa tę niejednoznaczność.

Wszystkie twierdzenia o istnieniu dotyczą rozwiązań słabych.
Nie twierdzimy, że sam warunek eliptyczności i ograniczoność
mierzalnych współczynników dają $H^2$ lub gładkość klasyczną.
Przykład transmisyjny pokazuje dokładnie, gdzie taka
regularność może zawieść. Regularność wewnętrzną i brzegową
rozwiniemy w kolejnych rozdziałach.

% Iteracja 14: wykład, do włączenia przed dodatkiem z kodem.

\chapter[Zasady maksimum i teoria potencjału]{Zasady maksimum, lemat Hopfa, nierówności Harnacka i podstawy teorii potencjału}
\label{chap:pde14}

W rozdziale~\ref{chap:pde13} uzyskaliśmy rozwiązania słabe
zagadnień eliptycznych metodami wariacyjnymi. Samo istnienie w
$H^1$ nie mówi jeszcze, jaki znak ma rozwiązanie, czy osiąga ono
maksimum wewnątrz obszaru ani jak wartości w różnych punktach
zależą od siebie. Odpowiedzi na te pytania dostarczają zasady
maksimum i nierówności Harnacka. Zaczniemy od współczynników
mierzalnych, dla których wystarczy argument energetyczny, a potem
przejdziemy do silniejszych twierdzeń dla funkcji harmonicznych.

\section{Słaba zasada maksimum dla rozwiązań słabych}
\label{sec:pde14-weak}
Niech $\Omega\subset\R^d$ będzie obszarem ograniczonym Lipschitzowskim. Załóżmy, że $A\in L^\infty(\Omega;\R^{d\times d})$
spełnia $\xi^{\mathsf T}A(x)\xi\ge\lambda|\xi|^2$
dla pewnego $\lambda>0$, prawie każdego $x\in\Omega$ i każdego
$\xi\in\R^d$. Nie zakładamy symetrii $A$.
Rozważamy operator $Lu=-\diver(A\nabla u)+cu$, gdzie
$c\in L^\infty(\Omega)$ i $c\ge0$ prawie wszędzie.
Wszystkie funkcje i współczynniki są rzeczywiste.
Stosujemy konwencje rozdziału~\ref{chap:pde13}: nierówność
eliptyczności zachodzi na wspólnym zbiorze pełnej miary,
a formę $a$ definiujemy na $H^1(\Omega)\times H^1(\Omega)$.

\begin{definicja}[Słabe rozwiązanie nierówności]
\label{def:pde14-subsolution}
Funkcja $u\in H^1(\Omega)$ jest \emph{słabym podrozwiązaniem}
$Lu\le0$ w $\Omega$, jeżeli
\[
 \int_\Omega\big(A\nabla u\cdot\nabla\varphi
 +cu\varphi\big)\,\mathrm d\lambda_d(x)\le0
 \quad\text{dla każdego }\varphi\in H_0^1(\Omega),\ \varphi\ge0.
\]
Analogicznie definiujemy nadrozwiązanie, odwracając nierówność.
Całka istnieje i jest ciągła względem $\varphi$ w normie
$H^1$, na mocy oszacowania formy~\eqref{eq:pde13-form}.
Jeżeli $u$ jest jednocześnie podrozwiązaniem i nadrozwiązaniem,
to $a(u,\varphi)=0$ dla nieujemnych $\varphi\in H_0^1$.
Dla dowolnego rzeczywistego $\varphi\in H_0^1$ piszemy
$\varphi=\varphi^+-\varphi^-$; obie części należą do $H_0^1$
na mocy wniosku~\ref{cor:pde9-truncation},
lematu~\ref{lem:pde13-trace-parts} i twierdzenia
~\ref{thm:pde10-trace-zero}. Z liniowości otrzymujemy
równość dla wszystkich testów, czyli rozwiązanie słabe
$Lu=0$ (bez narzucania w tej definicji danych brzegowych).
Odwrotna implikacja wynika przez ograniczenie do testów nieujemnych.
W definicji nie zakładamy istnienia klasycznych pochodnych drugiego rzędu.
\end{definicja}

\begin{lemat}[Obcięcie poziomowe]
\label{lem:pde14-trunc}
Jeżeli $u\in H^1(\Omega)$ i $k\in\R$, to
$w=(u-k)^+:=\max\{u-k,0\}\in H^1(\Omega)$ oraz
$\nabla w=\chi_{\{u>k\}}\nabla u$ prawie wszędzie.
Jeżeli ponadto $\operatorname{Tr}u\le k$ prawie wszędzie na
$\partial\Omega$, to $w\in H_0^1(\Omega)$.
\end{lemat}
\begin{proof}
Ponieważ $\Omega$ jest ograniczona, stała $k$ należy do
$H^1(\Omega)$, ma zerowy słaby gradient i ślad równy $k$.
Stosujemy wniosek~\ref{cor:pde9-truncation} z $p=2$
do funkcji $u-k$. Otrzymujemy $w\in H^1$ i
$\nabla w=\chi_{\{u-k>0\}}\nabla(u-k)=\chi_{\{u>k\}}\nabla u$
prawie wszędzie.
Z lematu~\ref{lem:pde13-trace-parts} oraz liniowości śladu
wynika $\operatorname{Tr}w=(\operatorname{Tr}u-k)^+$.
Przy założeniu na ślad prawa strona znika.
Twierdzenie~\ref{thm:pde10-trace-zero} daje $w\in H_0^1$.
Nie powtarzamy dowodów obcięcia i zgodności śladu z częściami
dodatnimi z wcześniejszych rozdziałów.
\end{proof}

Poniższy wynik uogólnia elementarną zasadę maksimum
z twierdzenia~\ref{thm:laplace-max}: tutaj współczynniki
są jedynie mierzalne, a równanie i warunek brzegowy
interpretujemy w sensie słabym. Hipotezy obu twierdzeń
są różne, więc nie utożsamiamy ich dowodów.

\begin{twierdzenie}[Słaba zasada maksimum]
\label{thm:pde14-weak-max}
Niech $u\in H^1(\Omega)$ będzie podrozwiązaniem $Lu\le0$,
przy powyższych założeniach o $A$ i $c$. Niech $M\ge0$ będzie
liczbą taką, że $\operatorname{Tr}u\le M$ prawie wszędzie na
$\partial\Omega$. Wówczas $u\le M$ prawie wszędzie w $\Omega$.
Jeżeli $c=0$, ten sam wniosek jest prawdziwy dla dowolnego
$M\in\R$ spełniającego warunek na ślad.
\end{twierdzenie}
\begin{proof}
Kładziemy $w=(u-M)^+$. Z lematu~\ref{lem:pde14-trunc}
$w\in H_0^1(\Omega)$ i możemy użyć go jako nieujemnej
funkcji testowej. Na zbiorze $\{u>M\}$ mamy
$\nabla w=\nabla u$ i $u=M+w$, natomiast poza nim
$w=\nabla w=0$ prawie wszędzie. Zatem
\[
 0\ge \int_\Omega(A\nabla u\cdot\nabla w+cuw)\,
 \mathrm d\lambda_d(x)
 \ge \lambda\|\nabla w\|_2^2
      +\int_{\{u>M\}}c(M+w)w\,\mathrm d\lambda_d(x).
\]
Ponieważ $c\ge0$, $M\ge0$ i $w\ge0$, ostatnia całka jest
nieujemna. Ponieważ $\lambda>0$, mamy $\|\nabla w\|_2=0$,
a nierówność Poincarégo dla $H_0^1(\Omega)$
(wniosek~\ref{cor:pde10-Poincare}) daje $\|w\|_2=0$,
czyli $u\le M$ prawie wszędzie. Jeśli $c=0$, składnik
zerowego rzędu znika, więc znak $M$ nie odgrywa roli.
\end{proof}

\begin{wniosek}[Porównanie]
\label{cor:pde14-compare}
Niech $u,v\in H^1(\Omega)$, $c\ge0$, a
$L=-\diver(A\nabla\cdot)+c$. Jeżeli $Lu\le Lv$ słabo w
$\Omega$ i $\operatorname{Tr}u\le\operatorname{Tr}v$ prawie
wszędzie na $\partial\Omega$, to $u\le v$ prawie wszędzie
w $\Omega$.
\end{wniosek}
\begin{proof}
Różnica $w=u-v$ jest podrozwiązaniem $Lw\le0$ i ma
niedodatni ślad. Stosujemy twierdzenie~\ref{thm:pde14-weak-max}
z $M=0$. Nie wykorzystujemy symetrii macierzy $A$.
\end{proof}

\begin{uwaga}[Granice wniosku]
Przy członie $b\cdot\nabla u$ poprzedni test zawiera
$\int_\Omega(b\cdot\nabla w)w\,\mathrm d\lambda_d(x)$,
którego znak nie jest określony. Dodatkowy argument
energetyczny podano już w twierdzeniu
~\ref{thm:pde13-drift-comparison}: dla $b\in L^\infty$
o dystrybucyjnej dywergencji w $L^\infty$ wystarcza
$c-\tfrac12\diver b\ge0$ do porównania.
Jeśli dodatkowo $c\ge0$ i $M\ge0$, stała $M$ jest
nadrozwiązaniem, gdyż $LM=cM\ge0$; porównanie daje
wtedy również $u\le M$ dla podrozwiązania $Lu\le0$.
Bez dalszych założeń nie wolno pomijać warunku znakowego
na $c$, co pokazuje poniższy przykład.
\end{uwaga}

\begin{kontrprzyklad}[Ujemny człon zerowego rzędu]
\label{ex:pde14-negative-c}
Na $\Omega=(0,1)$ funkcja $u(x)=\sin(\pi x)$ jest dodatnia
we wnętrzu i ma zerowy ślad, choć spełnia
$-u''-\pi^2u=0$. Gdyby dla $c=-\pi^2$ zachodziła zasada
maksimum z twierdzenia~\ref{thm:pde14-weak-max},
wynikałoby $u\le0$, co jest fałszem. Jest to jednocześnie
przykład rezonansu z rozdziału~\ref{chap:pde13}.
\end{kontrprzyklad}

\section{Funkcje harmoniczne i mocna zasada maksimum}
\label{sec:pde14-strong}
Od tej sekcji przyjmujemy $d\ge2$, chyba że zaznaczono inaczej.
Funkcją harmoniczną w otwartym zbiorze $U\subset\R^d$
nazywamy rzeczywistą funkcję $u\in C^2(U)$ spełniającą
$\Delta u=0$. Z lematu Weyla~\ref{thm:pde7-weyl} wiemy,
że każda rzeczywista dystrybucja $T$ spełniająca $\Delta T=0$
jest dystrybucją generowaną przez jedyną gładką funkcję
harmoniczną. Jest ona rzeczywista, ponieważ część urojona
generuje dystrybucję zerową, a stwierdzenie
~\ref{prop:pde7-injective} zapewnia jej znikanie prawie
wszędzie; z ciągłości wynika znikanie wszędzie.
Twierdzenia punktowe stosujemy do tego reprezentanta.
W wymiarze jeden funkcje harmoniczne na przedziałach są
afiniczne ($u''=0$), więc podstawowe zasady maksimum
wynikają bezpośrednio z tej postaci.

\begin{twierdzenie}[Mocna zasada maksimum dla laplasjanu]
\label{thm:pde14-strong}
Niech $U\subset\R^d$ będzie obszarem, a $u\in C^2(U)$ niech spełnia $\Delta u=0$ w $U$.
Jeżeli $u$ osiąga globalne maksimum lub globalne minimum
w pewnym punkcie $x_0\in U$, to $u$ jest stała na $U$.
Silniejszą wersję dla ekstremum lokalnego udowodniono
w twierdzeniu~\ref{thm:pde3-strong-max}.
\end{twierdzenie}
\begin{proof}
Załóżmy, że $u(x_0)=M=\sup_U u$. Dla dowolnego
$r>0$ takiego, że $\overline{B_r(x_0)}\subset U$, własność
wartości średniej (twierdzenie~\ref{thm:pde3-mean}) daje
\[
 M=u(x_0)=\frac{1}{\sigma(\partial B_r(x_0))}
 \int_{\partial B_r(x_0)}u(y)\,\mathrm d\sigma(y),
\]
gdzie $\sigma$ jest miarą powierzchniową sfery.
Funkcja $M-u$ jest nieujemna i ciągła na sferze.
Jeżeli byłaby dodatnia w choć jednym jej punkcie, to
byłaby dodatnia na pewnym otwartym fragmencie sfery,
a średnia $u$ byłaby ściśle mniejsza od $M$.
Tak więc $u\equiv M$ na każdej takiej sferze, a zatem
również w całej małej kuli.
Zbiór $E=\{x\in U:u(x)=M\}$ jest niepusty i domknięty
w $U$ z ciągłości. Ten sam argument zastosowany do każdego
$x\in E$ pokazuje, że $E$ jest też otwarty w $U$.
Spójność $U$ wymusza $E=U$. Dla minimum stosujemy wynik do $-u$.
\end{proof}

\begin{wniosek}[Dodatniość i jednoznaczność]
\label{cor:pde14-positive}
Jeżeli $U$ jest spójna, a $u\ge0$ jest harmoniczna w $U$,
to $u\equiv0$ albo $u(x)>0$ dla wszystkich $x\in U$.
Jeżeli $U$ jest ponadto ograniczona oraz
$u,v\in C^2(U)\cap C(\overline U)$ są harmoniczne i mają
te same wartości na $\partial U$, to $u=v$ na $U$.
\end{wniosek}
\begin{proof}
Jeżeli $u(x_0)=0$ we wnętrzu, funkcja osiąga tam minimum,
więc jest stała, a jej wartość w $x_0$ wynosi zero,
zatem $u\equiv0$.
Dla jednoznaczności $w=u-v$ jest harmoniczna i znika na
brzegu. Ciągłość na zwartym $\overline U$ daje istnienie
maksimum i minimum; gdyby którakolwiek ekstremalna
wartość była różna od zera, ekstremum osiągnięte byłoby
we wnętrzu, co w myśl twierdzenia~\ref{thm:pde14-strong}
wymusiłoby stałość $w$ na $U$. Ponieważ niepusty
ograniczony obszar ma niepusty brzeg, ciągłość $w$
na $\overline U$ i zerowe wartości brzegowe wymuszają
wtedy stałą równą zero, sprzecznie z niezerowym ekstremum.
Stąd oba ekstrema są zerowe, czyli $u=v$.
\end{proof}

\section{Lemat Hopfa: znak pochodnej normalnej}
\label{sec:pde14-hopf}
Słaba zasada maksimum mówi, gdzie może znajdować się
minimum. Lemat Hopfa wyjaśnia, jak funkcja zachowuje się
przy zbliżaniu do punktu minimum na brzegu.

\begin{definicja}[Warunek kuli wewnętrznej w punkcie]
Niech $U\subset\R^d$ będzie otwarta, a $x_0\in\partial U$.
Mówimy, że w $x_0$ zachodzi \emph{warunek kuli wewnętrznej},
jeżeli istnieją $z\in U$ i $R>0$ takie, że
$B_R(z)\subset U$ oraz $x_0\in\partial B_R(z)$.
Wektor $\nu=(x_0-z)/R$ jest zewnętrzną normalną do tej kuli.
\end{definicja}

\begin{twierdzenie}[Lemat Hopfa dla funkcji harmonicznych]
\label{thm:pde14-hopf}
Niech $U\subset\R^d$ będzie obszarem ograniczonym,
$x_0\in\partial U$ i niech w $x_0$ zachodzi warunek kuli
wewnętrznej $B_R(z)\subset U$. Załóżmy, że
$u\in C^2(U)\cap C(\overline U)$ jest harmoniczna, nie jest
stała, a $u(x_0)=\min_{\overline U}u$.
Wówczas dla kierunku $-\nu=(z-x_0)/R$ zachodzi
\[
 \liminf_{t\downarrow0}\frac{u(x_0-t\nu)-u(x_0)}{t}>0.
\]
Jeżeli dodatkowo istnieje skończona granica
\[
 \partial_\nu u(x_0):=
 \lim_{t\downarrow0}\frac{u(x_0)-u(x_0-t\nu)}{t},
\]
to $\partial_\nu u(x_0)<0$. Jest to pochodna w kierunku
normalnej do kuli, obliczana od strony wnętrza; gdy $u$
ma przedłużenie klasy $C^1$ do otoczenia $x_0$,
jest równa $\nabla u(x_0)\cdot\nu$.
Nie zakładamy istnienia normalnej do samego $\partial U$.
\end{twierdzenie}
\begin{proof}
Odejmując stałą, zakładamy $u(x_0)=0$ i $u\ge0$ na
$\overline U$. Z mocnej zasady minimum
$u>0$ wewnątrz $U$. W pierścieniu
$\mathcal A=\{x:R/2<|x-z|<R\}$ rozważmy funkcję
\[
 w(x)=e^{-\alpha|x-z|^2}-e^{-\alpha R^2},\qquad \alpha>0.
\]
Jest ona dodatnia wewnątrz pierścienia, a na jego zewnętrznej
sferze wynosi zero. Różniczkowanie daje dla $j=1,\ldots,d$
\[
 \partial_j w=-2\alpha(x_j-z_j)e^{-\alpha|x-z|^2},
 \qquad
 \partial_{jj}w=(4\alpha^2(x_j-z_j)^2-2\alpha)
 e^{-\alpha|x-z|^2}.
\]
Po zsumowaniu $\Delta w=(4\alpha^2|x-z|^2-2d\alpha)
e^{-\alpha|x-z|^2}$. Jeśli $\alpha>2d/R^2$, to dla
$|x-z|\ge R/2$ mamy
$4\alpha^2|x-z|^2-2d\alpha\ge\alpha(\alpha R^2-2d)>0$.
Stąd $\Delta w>0$ na $\overline{\mathcal A}$. Z ciągłości i dodatniości $u$
na zwartej wewnętrznej sferze
$m:=\min_{|x-z|=R/2}u(x)>0$. Wybieramy
$0<\varepsilon\le m/\max_{|x-z|=R/2}w(x)$.
Wtedy $u-\varepsilon w\ge0$ na obu składowych brzegu
pierścienia: na wewnętrznej z wyboru $\varepsilon$, na
zewnętrznej z $u\ge0$ oraz $w=0$. Ponieważ
$\Delta(u-\varepsilon w)=-\varepsilon\Delta w<0$,
funkcja $u-\varepsilon w$ nie może osiągnąć ujemnego
minimum we wnętrzu $\mathcal A$: w takim punkcie hesjan
byłby dodatnio półokreślony i laplasjan nieujemny.
Ciągłość na zwartym $\overline{\mathcal A}$ oraz nieujemne
wartości brzegowe dają więc $u\ge\varepsilon w$.
Dla $0<t<R/2$ punkt $x=x_0-t\nu$ leży w $\mathcal A$,
a $|x-z|=R-t$. Otrzymujemy
\[
 \frac{u(x_0-t\nu)}t\ge
 \varepsilon\frac{e^{-\alpha(R-t)^2}-e^{-\alpha R^2}}t
 \longrightarrow 2\varepsilon\alpha R e^{-\alpha R^2}>0.
\]
W konsekwencji granica dolna lewej strony jest nie mniejsza
niż $2\varepsilon\alpha R e^{-\alpha R^2}>0$.
Jeśli istnieje pochodna zdefiniowana w tezie, to lewa strona
zbiega do $-\partial_\nu u(x_0)$ (po przywróceniu odjętej
wartości $u(x_0)$), co daje tezę.
\end{proof}

\begin{przyklad}[Konieczność geometrii brzegu]
\label{ex:pde14-cusp}
W klinie płaskim $U=\{(r\cos\theta,r\sin\theta):
0<r<1,\ 0<\theta<\omega\}$, gdzie $0<\omega<\pi$,
funkcja $u(r,\theta)=r^{\pi/\omega}\sin(\pi\theta/\omega)$
jest dodatnia we wnętrzu, ponieważ $0<\pi\theta/\omega<\pi$.
Niech $p=\pi/\omega>1$. Ze wzoru laplasjanu we współrzędnych
biegunowych otrzymujemy
\[
 \Delta u=\partial_{rr}u+r^{-1}\partial_r u
 +r^{-2}\partial_{\theta\theta}u
 =[p(p-1)+p-p^2]r^{p-2}\sin(p\theta)=0.
\]
Funkcja przedłuża się ciągle na domknięcie (w zerze
$|u|\le r^p\to0$), jest nieujemna na brzegu i znika
na obu ramionach. Przy wierzchołku ma minimum, lecz
$0\le u(r,\theta)/r\le r^{p-1}\to0$, nawet jednostajnie
w kierunku $\theta$.
W zerze nie zachodzi warunek kuli wewnętrznej.
Gdyby $B_R(z)\subset U$ i $|z|=R$, to środek $z$
o kącie $\theta_z\in(0,\omega)$ musiałby mieć odległość
od każdej z prostych ograniczających klin co najmniej $R$,
ponieważ klin jest zawarty w przecięciu odpowiednich
półpłaszczyzn. Odległości te wynoszą
$R\sin\theta_z$ i $R\sin(\omega-\theta_z)$.
Co najmniej jeden z kątów $\theta_z,\omega-\theta_z$
jest mniejszy niż $\pi/2$, więc jedna odległość jest
ściśle mniejsza niż $R$, sprzeczność.
Nie wolno zatem stosować lematu Hopfa z dodatnią liniową
dolną granicą przy wierzchołku.
\end{przyklad}

\section{Jądro Poissona i ilościowe porównywanie wartości}
\label{sec:pde14-poisson}
Oznaczmy przez $\omega_{d-1}$ pole powierzchni
jednostkowej sfery $S^{d-1}$ dla $d\ge2$ (w rozdziale
~\ref{chap:pde-laplace} tę samą wielkość oznaczono przez
$\omega_d$). Przez $\sigma$ rozumiemy miarę powierzchniową.
W kuli
$B_R=B_R(0)\subset\R^d$ jądro Poissona ma postać
\[
 P_R(x,y)=\frac{R^2-|x|^2}{\omega_{d-1}R|x-y|^d},
 \qquad x\in B_R,\quad y\in\partial B_R.
\]
Jest dodatnie i spełnia
$\int_{\partial B_R}P_R(x,y)\,\mathrm d\sigma(y)=1$.
Dla $g\in C(\partial B_R)$ funkcja
$u(x)=\int_{\partial B_R}P_R(x,y)g(y)\,\mathrm d\sigma(y)$
jest harmoniczna i przyjmuje dane brzegowe $g$ w sensie
ciągłego przedłużenia na $\overline{B_R}$; są to wyniki
twierdzenia~\ref{thm:pde3-poisson}. Jeżeli $u$ jest już
harmoniczna w $B_R$ i ciągła na $\overline{B_R}$,
stosujemy ten wynik do $g=u|_{\partial B_R}$.
Różnica między $u$ a całką Poissona jest harmoniczna
i ma zerowe wartości brzegowe, więc znika na mocy
wniosku~\ref{cor:pde14-positive}. Otrzymujemy formułę
Poissona dla dowolnej takiej funkcji $u$.

\begin{twierdzenie}[Nierówność Harnacka w kuli]
\label{thm:pde14-harnack}
Niech $d\ge2$, $0<r<R$, a $u\ge0$ będzie funkcją
harmoniczną w otoczeniu $\overline{B_R(0)}$.
Wtedy dla każdego $x\in\overline{B_r(0)}$ zachodzi
\begin{equation}\label{eq:pde14-harnack}
 \frac{R^2-r^2}{R^2}\left(\frac{R}{R+r}\right)^d u(0)
 \le u(x)\le
 \frac{R^{d-2}(R+r)}{(R-r)^{d-1}}u(0).
\end{equation}
W szczególności $\sup_{B_r}u\le C_{d,r/R}\inf_{B_r}u$
przy stałej zależnej wyłącznie od $d$ i $r/R$.
\end{twierdzenie}
\begin{proof}
Stosujemy stwierdzenie~\ref{prop:pde3-harnack} z $n=d$;
jego założenia wynikają z harmoniczności w otoczeniu
domknięcia kuli i nieujemności $u$. Przypomnijmy rachunek
stałych. Dla $s=|x|\le r$ i $|y|=R$ mamy
$R-s\le|x-y|\le R+s$ oraz
\[
 \frac{P_R(x,y)}{P_R(0,y)}
 =\frac{R^{d-2}(R-s)(R+s)}{|x-y|^d}.
\]
Dolna granica ilorazu wynosi co najmniej
$R^{d-2}(R-s)/(R+s)^{d-1}$, a górna nie przekracza
$R^{d-2}(R+s)/(R-s)^{d-1}$. Pierwsza funkcja $s$
jest malejąca, druga rosnąca na $[0,R)$, ponieważ
odpowiednio licznik maleje, a mianownik rośnie, lub
odwrotnie. Zastępujemy zatem $s$ przez $r$.
Mnożymy nierówności przez $P_R(0,y)u(y)\ge0$
i całkujemy względem $\mathrm d\sigma(y)$ na $\partial B_R$.
Formuła Poissona daje~\eqref{eq:pde14-harnack}; dolna
stała jest równa $R^{d-2}(R-r)/(R+r)^{d-1}$.
Oznaczmy dolną i górną stałą przez $a_r>0$ i $b_r>0$.
Wtedy $\inf_{B_r}u\ge a_ru(0)$ oraz
$\sup_{B_r}u\le b_ru(0)$, więc bez dzielenia przez
wartość funkcji otrzymujemy
\[
 \sup_{B_r}u\le\frac{b_r}{a_r}\inf_{B_r}u
 =\left(\frac{R+r}{R-r}\right)^d\inf_{B_r}u.
\]
Możemy wziąć $C_{d,r/R}=((1+r/R)/(1-r/R))^d$.
W szczególności dla $r=R/2$ otrzymujemy stałą $3^d$.
\end{proof}

\begin{uwaga}[Usunięcie założenia o przedłużeniu]
Jeżeli $u\ge0$ jest harmoniczna jedynie w $B_R$, bez
ciągłości na $\partial B_R$, stosujemy twierdzenie w
kulach $B_\rho$ o $r<\rho<R$, gdyż ich domknięcia
są zawarte w $B_R$. Dla ustalonego $r$ i $x\in\overline{B_r}$
wartości $u(0),u(x)$ nie zależą od $\rho$. Przechodząc
do granicy $\rho\uparrow R$ po obu stronach oszacowania,
korzystamy tylko z ciągłości jawnych stałych dla $\rho>r$.
Otrzymujemy te same nierówności bez żadnego założenia
o istnieniu wartości brzegowych na $\partial B_R$.
\end{uwaga}

\begin{wniosek}[Nierówność Harnacka na zwartych podzbiorach]
\label{cor:pde14-harnack-chain}
Niech $U\subset\R^d$ będzie obszarem,
$K\subset U$ będzie zwarty i niepusty, a $u\ge0$ harmoniczna
w $U$. Istnieje stała $C=C(U,K,d)>0$, niezależna od $u$,
taka że $\sup_Ku\le C\inf_Ku$.
\end{wniosek}
\begin{proof}
Wybieramy punkt $x_*\in K$. Zbiór punktów $U$, które
można połączyć z $x_*$ łamaną zawartą w $U$, jest niepusty,
otwarty i względnie domknięty: dla dowolnego punktu $z\in U$
istnieje kula o środku $z$ zawarta w $U$, a dowolne dwa
jej punkty można połączyć odcinkiem. Jeśli jeden punkt
kuli jest osiągalny łamaną, osiągalne są wszystkie.
Ze spójności $U$ każdy $x\in K$ można więc połączyć
z $x_*$ taką łamaną $\gamma:[0,1]\to U$.

Obraz $\gamma([0,1])$ jest zwarty. Istnieje $\rho>0$
takie, że $\overline{B_{2\rho}(z)}\subset U$ dla każdego
$z$ z obrazu łamanej: jeśli $\R^d\setminus U$ jest
niepusty, wybieramy $2\rho$ mniejsze od dodatniej odległości
tego zwartego zbioru od domkniętego dopełnienia $U$;
dla $U=\R^d$ wystarczy dowolne $\rho>0$.
Z jednostajnej ciągłości $\gamma$ wybieramy podział
$0=t_0<\cdots<t_N=1$ taki, że
$|z_{j+1}-z_j|<\rho$, gdzie $z_j=\gamma(t_j)$.
Mamy $z_0=x_*$, $z_N=x$.
Nierówność Harnacka w każdej kuli $B_{2\rho}(z_j)$
z promieniem wewnętrznym $\rho$ porównuje $u(z_j)$
i $u(z_{j+1})$ w obu kierunkach ze stałą $3^d$.
Iterując, otrzymujemy obie nierówności
$u(x)\le 3^{dN}u(x_*)$ i $u(x_*)\le 3^{dN}u(x)$.
Ostatnia kula $B_{2\rho}(x)$ pozwala porównać jeszcze
$u(x)$ z $u(y)$ dla każdego $y\in B_\rho(x)$.
Zatem dla wszystkich takich $y$ zachodzi
\[
 C_x^{-1}u(x_*)\le u(y)\le C_xu(x_*),
 \qquad C_x=3^{d(N+1)}.
\]
Kule $B_\rho(x)$, z promieniem zależnym od wybranej łamanej,
pokrywają $K$. Wybieramy skończone podpokrycie i oznaczamy
przez $C_0$ maksimum odpowiadających stałych $C_x$.
Wszystkie konstrukcje zależą od $U,K,d$, a nie od $u$.
Dla dowolnych $y,z\in K$ mamy
$u(y)\le C_0u(x_*)\le C_0^2u(z)$.
Biorąc supremum po $y$ i infimum po $z$ otrzymujemy tezę
z $C=C_0^2$, również gdy $u(x_*)=0$.
\end{proof}

\begin{wniosek}[Twierdzenie Liouville'a z ograniczeniem jednostronnym]
\label{cor:pde14-liouville}
Jeżeli $u$ jest harmoniczna na $\R^d$, $d\ge2$, i jest
ograniczona z dołu albo z góry, to $u$ jest stała.
\end{wniosek}
\begin{proof}
Jeżeli $u\ge m$ dla pewnego $m\in\R$, zastępujemy
$u$ przez $u-m$, która jest nadal harmoniczna i nieujemna.
Wystarczy więc udowodnić wynik dla $u\ge0$.
Dla ustalonego $x\in\R^d$ wybieramy $R>|x|$ i stosujemy
\eqref{eq:pde14-harnack} z $r=|x|$ (dla $x=0$ teza jest
oczywista). Gdy $R\to\infty$, oba współczynniki przy $u(0)$
dążą do $1$. Z nierówności i przejścia do granicy otrzymujemy
$u(x)=u(0)$. Jeśli pierwotna funkcja jest ograniczona
z góry przez $M$, stosujemy dowód do $M-u\ge0$.
W obu przypadkach pierwotna funkcja jest stała.
Jest to wzmocnienie wniosku~\ref{cor:pde3-liouville},
w którym zakładano ograniczenie obustronne.
\end{proof}

\section{Oszacowania pochodnych i lokalne wygładzanie}
\label{sec:pde14-deriv}
\begin{twierdzenie}[Oszacowanie gradientu funkcji harmonicznej]
\label{thm:pde14-gradient}
Niech $d\ge2$, $R>0$, $x_0\in\R^d$ i niech $u$ będzie
harmoniczna w otoczeniu $\overline{B_R(x_0)}$.
Wtedy
\[
 |\nabla u(x_0)|\le \frac dR
 \sup_{B_R(x_0)}|u|.
\]
W szczególności, jeżeli $0<r<R$ i $u$ jest harmoniczna
w $B_R(x_0)$, to
$\sup_{B_r(x_0)}|\nabla u|\le
\frac{2d}{R-r}\sup_{B_R(x_0)}|u|$.
\end{twierdzenie}
\begin{proof}
Przez translację możemy przyjąć $x_0=0$.
Uzupełniamy składowe oszacowanie ze stwierdzenia
~\ref{prop:pde3-gradient} do oszacowania normy euklidesowej
całego gradientu. Różniczkujemy formułę Poissona:
na $\overline{B_{R/2}}\times\partial B_R$ mamy $|x-y|\ge R/2$,
więc wszystkie pierwsze pochodne jądra są ciągłe
i ograniczone. Ilorazy różnicowe względem każdej współrzędnej
w punkcie $0$ są jednostajnie ograniczone przez supremum
odpowiedniej pochodnej (twierdzenie o wartości średniej).
Po pomnożeniu przez ograniczoną funkcję $u(y)$ otrzymujemy
majorantę całkowalną na sferze. Twierdzenie o zbieżności
zmajoryzowanej uzasadnia różniczkowanie pod całką.
Rachunek daje
\[
 \nabla_x P_R(x,y)=\frac{1}{\omega_{d-1}R}
 \left[-2x|x-y|^{-d}
 -d(R^2-|x|^2)(x-y)|x-y|^{-d-2}\right].
\]
Dla $|y|=R$ mamy zatem
$P_R(0,y)=1/(\omega_{d-1}R^{d-1})$ oraz
$\nabla_x P_R(0,y)=d y/(\omega_{d-1}R^{d+1})$. Stąd
\[
 \nabla u(0)=\frac{d}{\omega_{d-1}R^{d+1}}
 \int_{\partial B_R}y\,u(y)\,\mathrm d\sigma(y).
\]
Wykorzystując $|y|=R$ oraz
$\sigma(\partial B_R)=\omega_{d-1}R^{d-1}$,
z nierówności trójkąta dla całki wektorowej otrzymujemy
\[
 |\nabla u(0)|\le
 \frac{d}{\omega_{d-1}R^{d+1}}\,
 R\,\omega_{d-1}R^{d-1}\sup_{\partial B_R}|u|
 =\frac dR\sup_{\partial B_R}|u|.
\]
Z ciągłości na domknięciu $\sup_{\partial B_R}|u|
\le\sup_{B_R}|u|$. Dla $x\in B_r(x_0)$ kula
$\overline{B_{(R-r)/2}(x)}$ jest zawarta w $B_R(x_0)$
z nierówności trójkąta. Stosując pierwsze oszacowanie
w tej kuli i biorąc supremum po $x$, dostajemy drugie
oszacowanie z tezy.
W istocie można wziąć nawet promień $R-r$, ponieważ
$|x-x_0|<r$; daje to stałą $d/(R-r)$ zamiast $2d/(R-r)$.

Na mocy twierdzenia~\ref{thm:pde7-weyl} funkcja harmoniczna
jest gładka, a każda jej pochodna jest harmoniczna, gdyż
pochodne komutują z laplasjanem. Dla wielowskaźnika
$\alpha$ o długości $m\ge1$ dzielimy odstęp $R-r$
na $m$ równych części i kładziemy
$r_j=R-j(R-r)/m$ dla $j=0,\ldots,m$.
Stosujemy drugie oszacowanie kolejno do pochodnych
rzędów $0,1,\ldots,m-1$ na kulach o promieniach $r_j$.
Każde przejście daje czynnik co najwyżej $dm/(R-r)$,
gdyż moduł pojedynczej składowej gradientu nie przekracza
jego normy euklidesowej. W rezultacie
\[
 \sup_{B_r(x_0)}|\partial^\alpha u|
 \le \left(\frac{dm}{R-r}\right)^m\sup_{B_R(x_0)}|u|.
\]
Jeśli supremum po prawej jest nieskończone, nierówność
rozumiemy w rozszerzonym sensie; na każdej kuli o domknięciu
zawartym w dziedzinie supremum jest skończone.
\end{proof}

\section{Subharmoniczność i nierówność wartości średniej}
\label{sec:pde14-subharmonic}
\begin{definicja}[Funkcja subharmoniczna w sensie klasycznym]
Niech $U\subset\R^d$ będzie otwarta, a $u\in C^2(U)$.
Nazywamy $u$ \emph{subharmoniczną}, jeśli $\Delta u\ge0$
w $U$, zaś \emph{superharmoniczną}, jeśli $\Delta u\le0$.
Istnieją również ogólniejsze pojęcia subharmoniczności
dla funkcji półciągłych z góry oraz superharmoniczności
dla funkcji półciągłych z dołu, związane z odpowiednimi
nierównościami dystrybucyjnymi. W tej sekcji twierdzenia
punktowe dowodzimy w klasie rzeczywistych funkcji $C^2$.
\end{definicja}

\begin{twierdzenie}[Nierówność wartości średniej]
\label{thm:pde14-sub-mean}
Jeżeli $u\in C^2(U)$ jest subharmoniczna i
$\overline{B_R(x_0)}\subset U$, to dla każdego $0<r\le R$
\[
 u(x_0)\le \frac{1}{\sigma(\partial B_r(x_0))}
 \int_{\partial B_r(x_0)}u(y)\,\mathrm d\sigma(y).
\]
Zachodzi również nierówność dla średniej po kuli:
\[
 u(x_0)\le \frac1{\lambda_d(B_r(x_0))}
 \int_{B_r(x_0)}u(y)\,\mathrm d\lambda_d(y).
\]
Dla funkcji superharmonicznej obie nierówności mają przeciwny zwrot.
\end{twierdzenie}
\begin{proof}
Niech $m(r)$ oznacza średnią sferyczną z tezy.
Korzystamy z rachunku przeprowadzonego w dowodzie
twierdzenia~\ref{thm:pde3-mean}, bez zakładania $\Delta u=0$.
Po podstawieniu $y=x_0+r\theta$, $\theta\in S^{d-1}$,
mamy $m(r)=\omega_{d-1}^{-1}
\int_{S^{d-1}}u(x_0+r\theta)\,\mathrm d\sigma(\theta)$.
Na zwartym otoczeniu przedziału promieni $[0,R]$
dopuszczalnym dzięki $\overline{B_R(x_0)}\subset U$
funkcja podcałkowa i jej pochodna względem $r$
są jednostajnie ograniczone na sferze. Twierdzenie
o wartości średniej i zbieżność zmajoryzowana
uzasadniają różniczkowanie pod całką. Otrzymujemy
\[
 m'(r)=\frac1{\omega_{d-1}}
 \int_{S^{d-1}}\nabla u(x_0+r\theta)\cdot\theta\,
 \mathrm d\sigma(\theta)
 =\frac1{\omega_{d-1}r^{d-1}}
 \int_{B_r(x_0)}\Delta u(y)\,\mathrm d\lambda_d(y).
\]
Druga równość wynika ze skalowania miary powierzchniowej
i twierdzenia Gaussa zastosowanego do pola $\nabla u$,
klasy $C^1$ w otoczeniu domkniętej kuli; normalna
zewnętrzna w punkcie $x_0+r\theta$ jest równa $\theta$.
Dla $\Delta u\ge0$ mamy $m'(r)\ge0$.
Z ciągłości $u$ w $x_0$ otrzymujemy
$|m(r)-u(x_0)|\le
\sup_{\theta\in S^{d-1}}|u(x_0+r\theta)-u(x_0)|\to0$,
gdy $r\downarrow0$. Monotoniczność $m$ daje pierwszą
nierówność. Całkując ją dla $0<s<r$ z wagą
$\omega_{d-1}s^{d-1}$, otrzymujemy
\[
 \int_{B_r(x_0)}u(y)\,\mathrm d\lambda_d(y)
 =\omega_{d-1}\int_0^r m(s)s^{d-1}\,\mathrm d\lambda_1(s)
 \ge u(x_0)\frac{\omega_{d-1}r^d}{d}.
\]
Ostatni współczynnik jest miarą kuli. Dzielimy przez niego.
Dla $\Delta u\le0$ stosujemy obie nierówności do $-u$.
\end{proof}

\begin{uwaga}[Wersja dystrybucyjna]
Dla $u\in L^1_{\mathrm{loc}}(U)$ warunek
$\Delta u\ge0$ w $\mathcal D'(U)$ oznacza
$\langle\Delta u,\varphi\rangle\ge0$ dla nieujemnych
$\varphi\in C_c^\infty(U)$. Dla nieujemnego standardowego
jądra wygładzającego $\rho_\varepsilon$ o całce $1$
i nośniku w $\overline{B_\varepsilon(0)}$ określamy
$u_\varepsilon=u*\rho_\varepsilon$ na
$U_\varepsilon=\{x\in U:\operatorname{dist}(x,\R^d\setminus U)>
\varepsilon\}$ (dla $U=\R^d$ przyjmujemy $U_\varepsilon=U$).
Z twierdzenia~\ref{thm:pde7-mollify} mamy
\[
 \Delta u_\varepsilon(x)
 =\langle\Delta u,\rho_\varepsilon(x-\cdot)\rangle\ge0,
 \qquad x\in U_\varepsilon,
\]
gdyż funkcja po prawej jest nieujemnym testem o zwartym
nośniku w $U$. Zatem $u_\varepsilon$ jest gładka
i subharmoniczna. Twierdzenie~\ref{thm:pde7-approx-id}
daje zbieżność do $u$ w $L^1_{\mathrm{loc}}(U)$.
Taka zbieżność nie wystarcza sama w sobie do przejścia
do nierówności punktowej dla dowolnego reprezentanta.
Pełna teoria wyróżnionych reprezentantów półciągłych
z góry (mogących przyjmować wartość $-\infty$)
wymaga dodatkowego twierdzenia, którego tutaj nie dowodzimy.
\end{uwaga}

\begin{twierdzenie}[Silna zasada minimum dla funkcji superharmonicznych]
\label{thm:pde14-strong-super}
Niech $U\subset\R^d$ będzie obszarem i niech
$u\in C^2(U)$ spełnia $\Delta u\le0$ w $U$.
Jeżeli $u$ osiąga minimum globalne względem $U$ w punkcie
$x_0\in U$, to jest stała na $U$.
\end{twierdzenie}
\begin{proof}
Niech $m=u(x_0)=\inf_Uu$. Z twierdzenia
~\ref{thm:pde14-sub-mean} w wersji superharmonicznej wynika
\[
 u(x_0)\ge\frac{1}{\sigma(\partial B_r(x_0))}
 \int_{\partial B_r(x_0)}u(y)\,\mathrm d\sigma(y)
\]
dla każdej kuli o domknięciu zawartym w $U$.
Ponieważ $u\ge m$, średnia wynosi jednocześnie co
najmniej $m$. Średnia funkcji nieujemnej $u-m$
na sferze jest więc zerowa. Gdyby w pewnym punkcie
sfery $u-m>0$, z ciągłości byłaby dodatnia, z dodatnią
dolną granicą, na fragmencie sfery o dodatniej mierze,
co przeczyłoby zerowej średniej. Stąd $u=m$ na każdej
takiej sferze, a więc w całej małej kuli o środku $x_0$.
Dla każdego $z$ ze zbioru $E=\{x\in U:u(x)=m\}$
powtarzamy argument z $z$ zamiast $x_0$. Zbiór $E$ jest
zatem otwarty w $U$, a dzięki ciągłości również względnie
domknięty. Jest niepusty, więc ze spójności $U$ mamy $E=U$.
Stosując wynik do $-u$, otrzymujemy analogiczną mocną
zasadę maksimum dla funkcji subharmonicznych.
\end{proof}

\begin{wniosek}[Lemat Hopfa dla nierówności $-\Delta u\ge0$]
\label{cor:pde14-hopf-super}
Niech $U\subset\R^d$ będzie obszarem ograniczonym,
$u\in C^2(U)\cap C(\overline U)$, $\Delta u\le0$ w $U$.
Załóżmy, że $u$ nie jest stała,
$u(x_0)=\min_{\overline U}u$ dla $x_0\in\partial U$
i istnieje kula $B_R(z)\subset U$ styczna do brzegu w $x_0$.
Wówczas, przy $\nu=(x_0-z)/R$,
\[
 \liminf_{t\downarrow0}\frac{u(x_0-t\nu)-u(x_0)}t>0.
\]
Jeśli istnieje skończona pochodna w kierunku normalnej
do kuli, obliczana od strony wnętrza zgodnie z twierdzeniem
~\ref{thm:pde14-hopf}, jest ona ujemna.
\end{wniosek}
\begin{proof}
Z twierdzenia~\ref{thm:pde14-strong-super} wynika,
że $u(x)>u(x_0)$ dla każdego $x\in U$. Odejmujemy
$u(x_0)$ i używamy tej samej bariery
$w(x)=e^{-\alpha|x-z|^2}-e^{-\alpha R^2}$,
co w dowodzie twierdzenia~\ref{thm:pde14-hopf}.
Na pierścieniu $R/2<|x-z|<R$ mamy $\Delta w>0$
dla $\alpha>2d/R^2$, a więc
$\Delta(u-\varepsilon w)\le-\varepsilon\Delta w<0$
dla $\varepsilon>0$. Po odjęciu minimum funkcja $u$
jest dodatnia na zwartej wewnętrznej sferze; oznaczamy
jej minimum przez $m>0$ i wybieramy
$0<\varepsilon\le m/(e^{-\alpha R^2/4}-e^{-\alpha R^2})$.
Wtedy $u-\varepsilon w\ge0$ na wewnętrznej sferze,
a na zewnętrznej również, gdyż $w=0$ i $u\ge0$.
Ujemne minimum na domkniętym pierścieniu musiałoby
być osiągnięte we wnętrzu, gdzie laplasjan byłby
nieujemny, sprzecznie z powyższą ścisłą nierównością.
Stąd $u\ge\varepsilon w$.
Dla $0<t<R/2$ podstawiamy $x=x_0-t\nu$ i otrzymujemy
\[
 \liminf_{t\downarrow0}
 \frac{u(x_0-t\nu)-u(x_0)}{t}
 \ge 2\varepsilon\alpha R e^{-\alpha R^2}>0,
\]
gdzie po lewej przywrócono odjętą wartość minimum.
Istnienie pochodnej z tezy daje jej ścisłą ujemność.
\end{proof}



\section{Nierówność Caccioppolego dla współczynników mierzalnych}
\label{sec:pde14-caccioppoli}
Przechodzimy ponownie do form dywergencyjnych, ale
bez założenia gładkości współczynników.

\begin{twierdzenie}[Oszacowanie wewnętrzne Caccioppolego]
\label{thm:pde14-caccioppoli}
Niech $\Omega\subset\R^d$ będzie otwarta,
$A\in L^\infty(\Omega;\R^{d\times d})$,
$\xi^{\mathsf T}A(x)\xi\ge\lambda|\xi|^2$ i
$|A(x)\xi|\le\Lambda|\xi|$ dla wszystkich $\xi\in\R^d$
na wspólnym zbiorze pełnej miary, gdzie $\lambda,\Lambda>0$.
Jeżeli rzeczywista funkcja $u\in H^1_{\mathrm{loc}}(\Omega)$
spełnia $-\diver(A\nabla u)=0$ dystrybucyjnie, to dla każdej
rzeczywistej $\eta\in C_c^\infty(\Omega)$ oraz $k\in\R$ zachodzi
\begin{equation}\label{eq:pde14-caccioppoli}
 \int_\Omega\eta^2|\nabla u|^2\,\mathrm d\lambda_d(x)
 \le \frac{4\Lambda^2}{\lambda^2}
 \int_\Omega (u-k)^2|\nabla\eta|^2\,\mathrm d\lambda_d(x).
\end{equation}
W szczególności, gdy $B_{2r}(x_0)\Subset\Omega$,
\[
 \int_{B_r(x_0)}|\nabla u|^2\,\mathrm d\lambda_d(x)
 \le \frac{C\Lambda^2}{\lambda^2r^2}
 \int_{B_{2r}(x_0)}|u-k|^2\,\mathrm d\lambda_d(x),
\]
gdzie $C$ zależy wyłącznie od wyboru standardowego obcięcia.
\end{twierdzenie}
\begin{proof}
Ponieważ operator nie ma członu zerowego rzędu,
$v=u-k$ spełnia to samo równanie i $\nabla v=\nabla u$.
Wybieramy otwarte $W\Subset\Omega$ zawierające
$\operatorname{supp}\eta$. Z lokalnej przynależności
$v\in H^1(W)$ i reguły iloczynu
~\ref{thm:pde9-products} mamy
$\varphi=\eta^2v\in H^1(W)$ oraz
$\nabla\varphi=\eta^2\nabla v+2\eta v\nabla\eta$.
Funkcja $\varphi$ znika poza zwartym podzbiorem $W$.
Jej przedłużenie przez zero na $\R^d$ należy do $H^1(\R^d)$:
dla testu na $\R^d$ mnożymy go przez gładkie obcięcie
równe $1$ w otoczeniu $\operatorname{supp}\eta$
i o nośniku w $W$, a następnie stosujemy tożsamość
słabej pochodnej na $W$. Otrzymujemy pochodne
przedłużone przez zero. Wygładzenie tego przedłużenia
przez splot daje testy o nośnikach w $W$ dla dostatecznie
małego parametru, zbieżne do $\varphi$ w $H^1$
(z twierdzeń~\ref{thm:pde7-mollify}
i~\ref{thm:pde7-approx-id} zastosowanych do funkcji
i jej pierwszych pochodnych na ustalonym zwartym
otoczeniu ich nośników).
Wobec tego $\varphi\in H_0^1(W)$.
Ponieważ $A\nabla u\in L^2(W)$, całka
$\int_W A\nabla u\cdot\nabla\psi\,\mathrm d\lambda_d(x)$
jest ciągła w normie $H^1(W)$. Przechodząc do granicy
z testami gładkimi, wolno więc użyć $\varphi$ w równaniu
dystrybucyjnym. Otrzymujemy
\[
 \int_\Omega\eta^2 A\nabla u\cdot\nabla u\,\mathrm d\lambda_d(x)
 =-2\int_\Omega\eta (u-k) A\nabla u\cdot\nabla\eta\,\mathrm d\lambda_d(x).
\]
Z eliptyczności, oszacowania normy operatorowej $A$
i nierówności Cauchy'ego--Schwarza otrzymujemy
$\lambda X^2\le2\Lambda XY$, gdzie
\[
 X^2=\int_\Omega\eta^2|\nabla u|^2\,\mathrm d\lambda_d(x),
 \qquad
 Y^2=\int_\Omega (u-k)^2|\nabla\eta|^2\,\mathrm d\lambda_d(x).
\]
Gdy $X=0$, teza jest oczywista. W przeciwnym razie
$X\le 2\Lambda Y/\lambda$, co po podniesieniu do kwadratu
daje~\eqref{eq:pde14-caccioppoli}.
Niech $\chi\in C^\infty(\R)$, $0\le\chi\le1$,
będzie równa $1$ dla $t\le1$ i równa $0$ dla
$t\ge3/2$. Kładziemy $\eta(x)=\chi(|x-x_0|/r)$.
Funkcja jest gładka również w środku (jest tam lokalnie
stała), ma nośnik w $\overline{B_{3r/2}(x_0)}$,
jest równa $1$ na $B_r(x_0)$ oraz
$|\nabla\eta|\le C_0/r$, gdzie $C_0=\|\chi'\|_\infty$.
Podstawienie do~\eqref{eq:pde14-caccioppoli} daje
drugie oszacowanie z $C=4C_0^2$.
W szczególności możemy wybrać $k=0$ lub średnią
$u_{B_{2r}(x_0)}$, uzyskując oszacowanie przez oscylację
w normie $L^2$ bez dodatkowego założenia na znak $u$.
\end{proof}

\begin{uwaga}[Dlaczego nie dowodzimy tu ogólnej nierówności Harnacka?]
Oszacowanie Caccioppolego jest punktem wyjścia do iteracji
De Giorgiego i Mosera. Pełna nierówność Harnacka dla
nieujemnych słabych rozwiązań przy mierzalnych współczynnikach
jest wynikiem znacznie głębszym niż nierówność dla funkcji
harmonicznych wyprowadzona z jądra Poissona. Przedstawimy
jej dowód w twierdzeniu~\ref{thm:pde17-harnack}
w rozdziale o regularności De Giorgiego--Nasha--Mosera;
nie używamy jej tutaj jako przesłanki.
\end{uwaga}

\section{Funkcja Greena i porównanie dla problemu Dirichleta}
\label{sec:pde14-green}
Niech $d\ge2$ oraz $\Phi$ będzie rozwiązaniem fundamentalnym
operatora $-\Delta$ w $\R^d$:
\[
 \Phi(z)=\begin{cases}
 \dfrac{|z|^{2-d}}{(d-2)\omega_{d-1}},&d\ge3,\\[4pt]
 -\dfrac1{2\pi}\log|z|,&d=2.
 \end{cases}
\]
W obu przypadkach $-\Delta\Phi=\delta_0$ w sensie
dystrybucyjnym. Nie utożsamiamy jądra Greena z samym
rozwiązaniem fundamentalnym: warunek brzegowy wymaga
odjęcia części harmonicznej.

\begin{definicja}[Funkcja Greena Dirichleta]
Niech $\Omega\subset\R^d$ będzie obszarem ograniczonym,
w którym dla każdego $y\in\Omega$ istnieje funkcja
$h_y\in C(\overline\Omega)\cap C^2(\Omega)$ harmoniczna
w $\Omega$ i taka, że $h_y(x)=\Phi(x-y)$ dla
$x\in\partial\Omega$. Definiujemy
$G_\Omega(x,y)=\Phi(x-y)-h_y(x)$ dla $x\ne y$.
Wówczas $-\Delta_xG_\Omega(x,y)=\delta_y$ dystrybucyjnie,
a $G_\Omega(\cdot,y)=0$ na brzegu.
\end{definicja}

\begin{twierdzenie}[Dodatniość funkcji Greena]
\label{thm:pde14-green-pos}
Przy powyższych założeniach dla każdego $y\in\Omega$
$G_\Omega(x,y)>0$ dla $x\in\Omega\setminus\{y\}$.
\end{twierdzenie}
\begin{proof}
Dla małego $\rho>0$ kula $\overline{B_\rho(y)}$ mieści
się w $\Omega$. Ponieważ $h_y$ jest ograniczona na
$\overline\Omega$, podczas gdy $\Phi(x-y)\to+\infty$
gdy $x\to y$ (również przy $d=2$), istnieje $\rho_0>0$
takie, że $G_\Omega(x,y)>0$ dla $|x-y|=\rho\le\rho_0$.
Na zewnętrznej granicy $\partial\Omega$ funkcja $G$ znika.
W obszarze pierścieniowym
$\Omega\setminus\overline{B_\rho(y)}$ jest harmoniczna;
słaba zasada minimum daje $G\ge0$ wszędzie w tym obszarze. Mocna zasada minimum daje ścisłą dodatniość.
Istotnie, $\Omega\setminus\overline{B_\rho(y)}$ jest spójna:
dowolną łamaną w $\Omega$ przecinającą kulę można zmodyfikować
wzdłuż sfery o promieniu nieco większym niż $\rho$,
która nadal leży w $\Omega$. Dowolny $x\ne y$
możemy objąć takim pierścieniem, wybierając $\rho<|x-y|$.
\end{proof}

\begin{lemat}[Jawna funkcja Greena kuli i jej pochodna normalna]
\label{lem:pde14-green-ball}
Niech $d\ge2$, $R>0$ i $x\in B_R\setminus\{0\}$.
Zapiszmy $x^*=R^2x/|x|^2$, a więc $|x^*|>R$. Dla
$y\in B_R$, $y\ne x$, funkcja Greena ma postać
\[
 G_{B_R}(x,y)=\begin{cases}
 \Phi(x-y)-(R/|x|)^{d-2}\Phi(x^*-y),&d\ge3,\\[3pt]
 \dfrac1{2\pi}\log\left(\dfrac{|x|\,|x^*-y|}
 {R|x-y|}\right),&d=2.
 \end{cases}
\]
W punkcie $x=0$ rozumiemy ją przez granicę. Jądro jest
symetryczne: $G_{B_R}(x,y)=G_{B_R}(y,x)$. Ponadto
$G_{B_R}(x,y)=0$ dla $|y|=R$, a na brzegu zachodzi
$-\partial_{\nu_y}G_{B_R}(x,y)=P_R(x,y)$, gdzie
$\nu_y=y/R$.
\end{lemat}
\begin{proof}
Dla $|y|=R$ otrzymujemy tożsamość geometryczną
$|x^*-y|=(R/|x|)|x-y|$. Po jej podstawieniu oba wzory
na $G$ przyjmują wartość zero. W zmiennej $y$ pierwszy
składnik jest harmoniczny poza biegunem $x$, drugi zaś
jest harmoniczny w całej kuli, ponieważ jego biegun
$x^*$ leży na zewnątrz. W pobliżu $y=x$ pierwszy
składnik jest rozwiązaniem fundamentalnym $-\Delta$,
więc $-\Delta_yG=\delta_x$.
Symetria wynika z równości
\[
 |x|^2|x^*-y|^2=R^4-2R^2x\cdot y+|x|^2|y|^2,
\]
której prawa strona jest symetryczna w $x,y$.

Dla $d\ge3$ wykorzystujemy
$\nabla_y\Phi(x-y)=(x-y)/(\omega_{d-1}|x-y|^d)$.
Różniczkując na zewnętrznej sferze i używając ponownie
$|x^*-y|=(R/|x|)|x-y|$, otrzymujemy
\begin{align*}
 \partial_{\nu_y}G(x,y)
 &=\frac{x\cdot y/R-R}{\omega_{d-1}|x-y|^d}
 -\frac{(|x|/R)^2(x^*\cdot y/R-R)}
 {\omega_{d-1}|x-y|^d}\\
 &=\frac{|x|^2-R^2}{\omega_{d-1}R|x-y|^d}
 =-P_R(x,y).
\end{align*}
Dla $d=2$ pochodna logarytmu spełnia tę samą
formułę, gdyż $\omega_1=2\pi$ i wykładnik $d-2$ wynosi
zero; składnik stały względem $y$ nie wnosi pochodnej.
Przejście $x\to0$ daje pozostały przypadek.
\end{proof}

\begin{twierdzenie}[Reprezentacja Greena w kuli]
\label{thm:pde14-green-repr}
Niech $d\ge2$, $R>0$, $f\in C(\overline{B_R})$,
$g\in C^2(\partial B_R)$, a
$u\in C^2(B_R)\cap C^1(\overline{B_R})$ spełnia
$-\Delta u=f$ w $B_R$ i $u=g$ na $\partial B_R$.
Wówczas dla $x\in B_R$ zachodzi
\begin{equation}\label{eq:pde14-green-repr}
 u(x)=\int_{B_R}G_{B_R}(x,y)f(y)\,\mathrm d\lambda_d(y)
 +\int_{\partial B_R}P_R(x,y)g(y)\,\mathrm d\sigma(y).
\end{equation}
W szczególności dla $f\ge0$ i $g\ge0$ mamy $u\ge0$.
\end{twierdzenie}
\begin{proof}
Dla ustalonego $x$ stosujemy drugą tożsamość Greena
w $B_R\setminus\overline{B_\rho(x)}$ do pary
$u(y)$ i $G_{B_R}(x,y)$. W tym obszarze $G$ jest
harmoniczna i $\Delta u=-f$. Otrzymujemy
\[
 \int_{B_R\setminus B_\rho(x)}G_{B_R}(x,y)f(y)
 \,\mathrm d\lambda_d(y)
 =\int_{\partial B_R}g\,\partial_{\nu_y}G\,
 \mathrm d\sigma+I_\rho,
\]
gdzie $I_\rho$ to wyraz z małej sfery, z normalną
skierowaną do jej środka. Ponieważ
$G(x,y)=\Phi(x-y)+O(1)$ w pobliżu $y=x$ oraz
$\partial_{\nu}G=\partial_{\nu}\Phi+O(1)$,
bezpośrednie obliczenie strumienia $\Phi$ daje
$\lim_{\rho\downarrow0}I_\rho=u(x)$.
Ponadto całka objętościowa przechodzi do granicy:
osobliwość $\Phi$ jest lokalnie całkowalna dla $d\ge2$.
Tożsamość $-\partial_{\nu_y}G_{B_R}(x,y)=P_R(x,y)$
została jawnie wyprowadzona w lemacie~\ref{lem:pde14-green-ball}.
Po przeniesieniu wyrazów otrzymujemy~\eqref{eq:pde14-green-repr}.
Wniosek o znaku wynika z dodatniości obu jąder.
\end{proof}

\begin{uwaga}[Uwaga o znakach we wzorze Greena]
W powyższym dowodzie normalna na \emph{wewnętrznej}
sferze $\partial B_\rho(x)$ jest skierowana \emph{do środka}
wyciętej kuli. To właśnie dlatego wyraz z tej sfery
ma znak prowadzący do $+u(x)$ w reprezentacji.
Rozpisanie orientacji na każdej granicy chroni przed
częstą pomyłką znaku przy przejściu od $\Phi$ do $G$.
\end{uwaga}

\section{Rysunek: bariera przy brzegu i geometria nierówności}
\label{sec:pde14-figure}
\begin{center}
\begin{tikzpicture}[scale=0.88,>=Latex]
  % Dwie kule styczne wewnetrznie: x0 nalezy do obu brzegow.
  \fill[blue!5] (-0.6,0) circle (2.45);
  \draw[thick,blue!65!black] (-0.6,0) circle (2.45);
  \draw[thick,orange!90!black] (0,0) circle (1.85);
  \fill (0,0) circle(2pt) node[below left] {$z$};
  \fill (1.85,0) circle(2.5pt) node[above right] {$x_0$};
  \draw[dashed] (0,0)--(1.85,0) node[midway,above] {$R$};
  \draw[-{Latex[length=2.5mm]},very thick] (1.85,0) -- (2.65,0)
    node[right] {$\nu$};
  \draw[-{Latex[length=2.5mm]},very thick,red!75!black]
    (1.78,-0.28) -- (0.90,-0.28);
  \node[red!70!black,below] at (0.15,-0.65) {kierunek $-\nu$};
  \node[blue!70!black] at (-2.3,1.25) {$\Omega$};
  \node[orange!85!black] at (-0.7,1.20) {$B_R(z)$};
\end{tikzpicture}

{\small Geometria lematu Hopfa: kula wewnętrzna dotyka brzegu
$\partial\Omega$ w $x_0$, a bariera wykrywa dodatni przyrost
w kierunku do wnętrza $-\nu$.}
\end{center}

\section{Podsumowanie i przejście do teorii regularności}
\label{sec:pde14-summary}
Wykazaliśmy słabą zasadę maksimum przy mierzalnych
współczynnikach i $c\ge0$, nie wykorzystując gładkości
ani symetrii części głównej. Dla laplasjanu udowodniliśmy
mocną zasadę maksimum, lemat Hopfa, ilościową nierówność
Harnacka, oszacowania gradientu i własności subharmoniczne.
Nierówność Caccioppolego pozostaje prawdziwa dla słabych
rozwiązań z mierzalną macierzą współczynników, ale sama
nie zapewnia od razu ciągłości rozwiązań ani nierówności
Harnacka. Te głębsze zagadnienia prowadzą do
teorii De Giorgiego--Nasha--Mosera oraz regularności
eliptycznej w kolejnych rozdziałach.

% Iteracja 15: regularność eliptyczna, całość rozdziału.

\chapter[Regularność Sobolewa rozwiązań eliptycznych]{Regularność Sobolewa rozwiązań eliptycznych: ilorazy różnicowe, oszacowania wewnętrzne i brzegowe}
\label{chap:pde15}

W rozdziale~\ref{chap:pde13} skonstruowaliśmy rozwiązania słabe
$u\in H^1$. Teraz zbadamy, kiedy równanie pozwala odzyskać drugie
pochodne $u$ w $L^2$. Zasadnicza trudność jest logiczna: nie wolno
różniczkować równania, zanim nie udowodnimy, że odpowiednie pochodne
istnieją. Ominiemy ją dzięki ilorazom różnicowym, których zastosowanie
wymaga jedynie $H^1$. W końcowej części zobaczymy, dlaczego gładkość
współczynników i brzegu naprawdę jest potrzebna.

W całym rozdziale $d\ge1$, a obszarem w $\R^d$ nazywamy niepusty,
otwarty i spójny podzbiór $\R^d$. Dla otwartych zbiorów
$U,V$ zapis $U\Subset V$ oznacza, że $\overline U$ jest zwartym
podzbiorem $V$. Symbole $H^k(U)=W^{k,2}(U)$ oraz
$\|\cdot\|_{2,U}=\|\cdot\|_{L^2(U)}$ mają znaczenie ustalone
w rozdziale~\ref{chap:pde-sobolev}. Całki po $U$ zapisujemy względem
miary Lebesgue'a $\mathrm d\lambda_d(x)$.

\section{Eliptyczność a możliwość różniczkowania rozwiązań}
\label{sec:pde15-model}
Rozważamy równanie w postaci dywergencyjnej
\begin{equation}\label{eq:pde15-basic}
 -\diver(A\nabla u)=f\qquad\text{w }\Omega,
\end{equation}
gdzie $\Omega\subset\R^d$ jest obszarem, $A=(a_{ij})$ ma współczynniki
rzeczywiste, a $f\in L^2_{\mathrm{loc}}(\Omega)$.
W całym rozdziale funkcje $u,f$ i funkcje testowe są rzeczywiste.

\begin{definicja}[Jednostajna eliptyczność i rozwiązanie słabe]
\label{def:pde15-elliptic}
Mówimy, że $A\in L^\infty_{\mathrm{loc}}(\Omega;\R^{d\times d})$
jest jednostajnie eliptyczna i ograniczona na mierzalnym
$U\subset\Omega$, jeżeli istnieją $\lambda,\Lambda>0$
oraz zbiór $E\subset U$ pełnej miary w $U$ takie, że
dla każdego $x\in E$ i każdego $\xi\in\R^d$ zachodzi
$\xi^{\mathsf T}A(x)\xi\ge\lambda|\xi|^2$ oraz
$|A(x)\xi|\le\Lambda|\xi|$. Nie wymagamy symetrii $A$.
Funkcja $u\in H^1_{\mathrm{loc}}(\Omega)$ jest rozwiązaniem słabym
równania~\eqref{eq:pde15-basic}, jeżeli dla każdego
$\varphi\in C_c^\infty(\Omega)$ zachodzi
\[
 \int_\Omega A\nabla u\cdot\nabla\varphi\,\mathrm d\lambda_d(x)
 =\int_\Omega f\varphi\,\mathrm d\lambda_d(x).
\]
Równość rozciąga się na funkcje $\varphi\in H^1(\Omega)$
o zwartym nośniku $K\subset\Omega$. Istotnie, wybieramy
otwarty $W\Subset\Omega$ zawierający $K$ i funkcję
$\psi\in C_c^\infty(W)$ równą $1$ w otoczeniu $K$.
Z lematu~\ref{lem:pde9-compact-product} dla $\psi\varphi=\varphi$
otrzymujemy ciąg $\varphi_n\in C_c^\infty(W)$ zbieżny do
$\varphi$ w $H^1(W)$. Ponieważ $A$ jest ograniczona na $W$,
a $\nabla u,f\in L^2(W)$, mamy
\[
 \left|\int_W A\nabla u\cdot\nabla(\varphi_n-\varphi)\,
 \mathrm d\lambda_d(x)\right|
 \le \|A\|_{\infty,W}\|\nabla u\|_{2,W}
       \|\nabla(\varphi_n-\varphi)\|_{2,W},
\]
\[
 \left|\int_W f(\varphi_n-\varphi)\,\mathrm d\lambda_d(x)\right|
 \le \|f\|_{2,W}\|\varphi_n-\varphi\|_{2,W}.
\]
Obie prawe strony dążą do zera, co uzasadnia przejście
do granicy w równaniu słabym.
\end{definicja}

\begin{uwaga}[Dlaczego samo $A\in L^\infty$ nie wystarcza?]
Jeśli $u$ ma tylko jedną słabą pochodną, to iloczyn $A\nabla u$
należy jedynie do $L^2_{\mathrm{loc}}$. Wyrażenie
$\diver(A\nabla u)$ jest jednak poprawnie określoną dystrybucją.
Przejście do $\nabla^2u\in L^2$ wymaga dodatkowej kontroli zmian
współczynników $A$. W poniższym twierdzeniu wystarczy warunek
$A\in W^{1,\infty}_{\mathrm{loc}}$.
\end{uwaga}

\section{Ilorazy różnicowe i ich podstawowe własności}
\label{sec:pde15-differences}
Dla $1\le j\le d$ i $h\ne0$ wprowadzamy
$\delta_h^jv(x)=\bigl(v(x+he_j)-v(x)\bigr)/h$, o ile oba punkty
należą do zbioru, na którym określona jest $v$.

\begin{lemat}[Oszacowania i całkowanie przez części]
\label{lem:pde15-difference}
Niech $U\subset\R^d$ będzie otwarty. Dla $h\ne0$ i ustalonego
$j$ oznaczamy $U_h=U\cap(U-he_j)$.
Jeżeli $v\in H^1(U)$, $K\subset U$ jest niepusty i zwarty oraz
$|h|<\operatorname{dist}(K,\R^d\setminus U)$
(przy $U=\R^d$ przyjmujemy odległość równą $+\infty$),
to $\|\delta_h^jv\|_{L^2(K)}\le\|\partial_jv\|_{L^2(U)}$.
Jeżeli $v\in L^2_{\mathrm{loc}}(U)$, $w\in L^2(U)$ ma
zwarty nośnik $K\subset U$ i $K+he_j\subset U$, to
\[
 \int_{U_h}(\delta_h^jv)w\,\mathrm d\lambda_d(x)
 =-\int_U v(\delta_{-h}^jw)\,\mathrm d\lambda_d(x),
\]
gdzie po prawej $w$ przedłużamy przez zero poza $U$.
Nie wymagamy zwartego nośnika $v$. Dla
$v\in H^1_{\mathrm{loc}}(U)$ zachodzi ponadto
$\delta_h^jv\in H^1_{\mathrm{loc}}(U_h)$ i
$\partial_k(\delta_h^jv)=\delta_h^j(\partial_kv)$
w $U_h$ dla $1\le k\le d$.
\end{lemat}
\begin{proof}
Dla funkcji gładkiej podstawowe twierdzenie rachunku całkowego
wzdłuż odcinka (również dla $h<0$) daje
$\delta_h^jv(x)=\int_0^1\partial_jv(x+the_j)\,
\mathrm d\lambda_1(t)$. Każdy przesunięty zbiór
$K+the_j$ jest zawarty w $U$. Z nierówności Cauchy'ego--Schwarza,
twierdzenia Tonellego i niezmienniczości miary Lebesgue'a
względem translacji otrzymujemy
\[
 \int_K|\delta_h^jv(x)|^2\,\mathrm d\lambda_d(x)
 \le\int_0^1\int_{K+the_j}|\partial_jv(y)|^2\,
 \mathrm d\lambda_d(y)\,\mathrm d\lambda_1(t)
 \le\int_U|\partial_jv(y)|^2\,\mathrm d\lambda_d(y).
\]
Dla ogólnej funkcji Sobolewa zbiór
$S=\{x+the_j:x\in K,\ 0\le t\le1\}$ jest zwarty jako obraz
$K\times[0,1]$ przez odwzorowanie ciągłe i zawarty w $U$.
Wybieramy otwarty $W$ taki, że $S\subset W\Subset U$.
Z twierdzenia~\ref{thm:pde9-mollify} otrzymujemy funkcje
gładkie $v_n$ na $W$, zbieżne do $v$ w $H^1(W)$.
Dla ustalonego $h$ mamy
\[
 \|\delta_h^j(v_n-v)\|_{2,K}
 \le |h|^{-1}(\|v_n-v\|_{2,K+he_j}
                  +\|v_n-v\|_{2,K})
 \le 2|h|^{-1}\|v_n-v\|_{2,W}\longrightarrow0.
\]
Stosując poprzedni rachunek do $v_n$ na $W$, a następnie
przechodząc do granicy, dostajemy
$\|\delta_h^jv\|_{2,K}\le\|\partial_jv\|_{2,W}
\le\|\partial_jv\|_{2,U}$.

W drugiej tożsamości wszystkie całki istnieją z nierówności
Cauchy'ego--Schwarza, ponieważ $v$ jest lokalnie kwadratowo
całkowalna, a $K$ oraz $K+he_j$ są zwarte w $U$.
Mamy $K\subset U_h$. Po przedłużeniu $w$ przez zero
podstawienie $y=x+he_j$ daje
\[
 \int_{U_h}(\delta_h^jv)w\,\mathrm d\lambda_d(x)
 =\frac1h\int_U v(y)\bigl(w(y-he_j)-w(y)\bigr)
 \,\mathrm d\lambda_d(y)
 =-\int_U v\,\delta_{-h}^jw\,\mathrm d\lambda_d.
\]
Wreszcie dla $\zeta\in C_c^\infty(U_h)$, po przedłużeniu
jej przez zero, z tej samej tożsamości otrzymujemy
\begin{align*}
 \int_{U_h}\delta_h^jv\,\partial_k\zeta\,\mathrm d\lambda_d
 &=-\int_U v\,\partial_k(\delta_{-h}^j\zeta)\,\mathrm d\lambda_d\\
 &=\int_U \partial_kv\,\delta_{-h}^j\zeta\,\mathrm d\lambda_d\\
 &=-\int_{U_h}\delta_h^j(\partial_kv)\zeta\,\mathrm d\lambda_d.
\end{align*}
Na każdym zwartym podzbiorze $U_h$ obie translacje
$v$ i $\partial_kv$ są kwadratowo całkowalne.
Ostatnia tożsamość daje więc postulowaną słabą pochodną.
\end{proof}

\begin{lemat}[Kryterium istnienia słabej pochodnej]
\label{lem:pde15-criterion}
Niech $U\subset\R^d$ będzie otwarty i $v\in L^2_{\mathrm{loc}}(U)$.
Jeśli dla każdego niepustego zwartego $K\subset U$ istnieją
$C_K,h_K>0$ takie, że $K+he_j\subset U$ oraz
$\|\delta_h^jv\|_{L^2(K)}\le C_K$ dla $0<|h|<h_K$, to
$\partial_jv\in L^2_{\mathrm{loc}}(U)$ i
$\|\partial_jv\|_{L^2(K)}\le C_{K'}$ dla zwartych
$K,K'\subset U$ spełniających $K\subset\operatorname{int}K'$.
\end{lemat}
\begin{proof}
Ustalmy takie $K,K'$. Wybieramy otwarty $V$ z
$K\subset V\subset\overline V\subset\operatorname{int}K'$;
jest to możliwe dzięki zwartości $K$ i otwartości
$\operatorname{int}K'$. Dla dowolnego ciągu niezerowych
$h_n\to0$ z $|h_n|<h_{K'}$ mamy
$\|\delta_{h_n}^jv\|_{2,V}\le C_{K'}$.
Z twierdzenia o słabej ciągowej zwartości ograniczonych
ciągów w przestrzeni Hilberta $L^2(V)$ wybieramy podciąg
(oznaczany tak samo), dla którego
$\delta_{h_n}^jv\rightharpoonup w$ w $L^2(V)$.
Dla $\varphi\in C_c^\infty(V)$ stosujemy poprzedni lemat
dla dostatecznie małych $|h_n|$:
\[
 \int_V\delta_{h_n}^jv\,\varphi\,\mathrm d\lambda_d
 =-\int_U v\,\delta_{-h_n}^j\varphi\,\mathrm d\lambda_d .
\]
Nośniki $\delta_{-h_n}^j\varphi$ dla dużych $n$
leżą w jednym zwartym podzbiorze $V$.
Z podstawowego twierdzenia rachunku całkowego i jednostajnej
ciągłości $\partial_j\varphi$ na $\R^d$ wynika jednostajna
zbieżność $\delta_{-h_n}^j\varphi\to\partial_j\varphi$,
a z ograniczenia nośników także zbieżność w $L^2(V)$.
Z nierówności Cauchy'ego--Schwarza prawa strona zbiega do
$-\int_V v\partial_j\varphi\,\mathrm d\lambda_d$;
lewa zbiega do $\int_V w\varphi\,\mathrm d\lambda_d$
na mocy słabej zbieżności. Stąd $w$ reprezentuje
$\partial_jv$ na $V$. Półciągłość normy przy słabej zbieżności
daje
\[
 \|w\|_{2,K}\le\|w\|_{2,V}
 \le\liminf_n\|\delta_{h_n}^jv\|_{2,V}\le C_{K'}.
\]
Każdy punkt $U$ ma takie otoczenie $V$. Na przecięciach
reprezentanty są równe prawie wszędzie z jednoznaczności
dystrybucji regularnych (stwierdzenie~\ref{prop:pde7-injective}).
Wybierając przeliczalne podpokrycie tych otoczeń (przestrzeń
euklidesowa ma przeliczalną bazę), otrzymujemy jeden
mierzalny reprezentant $\partial_jv$, lokalnie kwadratowo
całkowalny. Istotnie, każdy zwarty podzbiór $U$ pokrywają
skończenie liczne z tych otoczeń i na każdym z nich
odpowiednia całka kwadratu jest skończona.
\end{proof}

\section[Wewnętrzna regularność Sobolewa drugiego rzędu]{Wewnętrzna regularność $H^2$}
\label{sec:pde15-h2}
\begin{twierdzenie}[Wewnętrzna regularność drugiego rzędu]
\label{thm:pde15-h2}
Niech $\Omega\subset\R^d$ będzie obszarem,
$A\in W^{1,\infty}_{\mathrm{loc}}(\Omega;\R^{d\times d})$ będzie
jednostajnie eliptyczna i ograniczona na zwartych podzbiorach,
$f\in L^2_{\mathrm{loc}}(\Omega)$, a $u\in H^1_{\mathrm{loc}}(\Omega)$
rozwiązaniem słabym równania~\eqref{eq:pde15-basic}.
Wówczas $u\in H^2_{\mathrm{loc}}(\Omega)$. Dla każdych
niepustych otwartych zbiorów $U,V$ takich, że
$U\Subset V\Subset\Omega$, zachodzi
\begin{equation}\label{eq:pde15-h2-est}
 \|u\|_{H^2(U)}\le C\bigl(\|f\|_{L^2(V)}
 +\|u\|_{H^1(V)}\bigr),
\end{equation}
gdzie $C$ zależy od $U,V,d$, stałych eliptyczności na $V$
oraz $\|A\|_{W^{1,\infty}(V)}$.
\end{twierdzenie}
\begin{proof}
\emph{Krok 1: dziedziny przesunięć i równanie różnicowe.}
Wybieramy $\eta\in C_c^\infty(V)$ taką, że
$0\le\eta\le1$ i $\eta=1$ w otoczeniu $\overline U$.
Niech $K=\operatorname{supp}\eta$ i wybierzmy $\rho>0$ tak,
aby $K+\overline{B_{2\rho}(0)}\subset V$.
Ustalamy $j\in\{1,\ldots,d\}$ i $0<|h|<\rho$.
Oznaczamy $V_h=V\cap(V-he_j)$ i $w_h=\delta_h^ju$ na $V_h$.
Z lematu~\ref{lem:pde15-difference} mamy
$w_h\in H^1_{\mathrm{loc}}(V_h)$ oraz
$\nabla w_h=\delta_h^j\nabla u$.
W szczególności $K\Subset V_h$.

Dla $\varphi\in C_c^\infty(V_h)$ funkcja
$-\delta_{-h}^j\varphi$ (po przedłużeniu $\varphi$
przez zero) jest gładkim testem o zwartym nośniku w $V$.
Podstawiamy ją do równania słabego
i stosujemy dyskretne całkowanie przez części
z lematu~\ref{lem:pde15-difference}, osobno do każdej
składowej pola $A\nabla u$ oraz do $f$.
Otrzymujemy
\[
 \int_{V_h}\delta_h^j(A\nabla u)\cdot\nabla\varphi\,
 \mathrm d\lambda_d(x)
 =\int_{V_h}(\delta_h^jf)\varphi\,\mathrm d\lambda_d(x).
\]
Rozwinięcie różnicy iloczynu daje prawie wszędzie w $V_h$
\[
 \delta_h^j(A\nabla u)(x)
 =A(x)\nabla w_h(x)
   +(\delta_h^jA)(x)\nabla u(x+he_j).
\]
Po przeniesieniu drugiego składnika uzyskujemy
\begin{equation}\label{eq:pde15-diff-eq}
 \begin{aligned}
 \int_{V_h} A\nabla w_h\cdot\nabla\varphi\,\mathrm d\lambda_d(x)
 &=\int_{V_h}(\delta_h^jf)\varphi\,\mathrm d\lambda_d(x)\\
 &\quad-\int_{V_h}(\delta_h^jA)(x)\nabla u(x+he_j)
       \cdot\nabla\varphi\,\mathrm d\lambda_d(x).
 \end{aligned}
\end{equation}
Dla ustalonego $h$ współczynniki $\delta_h^jA$ są ograniczone,
a $\delta_h^jf$ i $\nabla w_h$ należą do $L^2(V_h)$
(z oszacowania norm obu translacji przez normy na $V$).
Zatem obie strony równania są ciągłe względem normy
$H^1(V_h)$. Na mocy lematu~\ref{lem:pde9-compact-product}
funkcję $\varphi=\eta^2w_h$ można aproksymować testami
gładkimi o nośnikach w ustalonym zwartym podzbiorze $V_h$.
Jej przedłużenie przez zero na $\R^d$ należy do
$H^1(\R^d)$, a słaby gradient jest przedłużeniem przez zero
pola
\[
 \nabla\varphi=\eta^2\nabla w_h+2\eta w_h\nabla\eta.
\]
Wolno więc użyć tej funkcji w~\eqref{eq:pde15-diff-eq}.

\emph{Krok 2: jednostajna kontrola różnic współczynników.}
Niech $D=\|\partial_jA\|_{L^\infty(V;\mathrm{op})}$,
gdzie używamy normy operatorowej macierzy. Dla ustalonego
$h$ mamy $|\delta_h^jA(x)|_{\mathrm{op}}\le D$ prawie
wszędzie na $K$. Istotnie, wygładzamy $A$ przez splot
na otoczeniu zwartego zbioru $K+\overline{B_\rho(0)}$
zawartego w $V$. Dla dostatecznie małego parametru
jądro nieujemne o całce $1$ i komutacja pochodnych
ze splotem (twierdzenie~\ref{thm:pde7-mollify})
dają $\|\partial_j A_\varepsilon(x)\|_{\mathrm{op}}\le D$
na tym otoczeniu. Całkując po odcinku od $x$ do $x+he_j$,
otrzymujemy analogiczne ograniczenie dla
$\delta_h^j A_\varepsilon(x)$. Z twierdzenia
~\ref{thm:pde9-mollify} dla $p=2$ (na relatywnie zwartym
otoczeniu) mamy zbieżność $A_\varepsilon$ do $A$ w $L^2$
na $K$ i $K+he_j$. W konsekwencji
$\delta_h^j A_\varepsilon\to\delta_h^j A$ w $L^2(K)$.
Po przejściu do podciągu zbieżnego prawie wszędzie
otrzymujemy żądane ograniczenie. Stała $D$ nie zależy od $h$.

\emph{Krok 3: oszacowanie energetyczne.}
Oznaczmy
\[
 X=\|\eta\nabla w_h\|_{2,V_h},\quad
 Y=\|\nabla u\|_{2,V},\quad F=\|f\|_{2,V},
 \quad M=\|\nabla\eta\|_\infty .
\]
Lemat~\ref{lem:pde15-difference} i niezmienniczość miary
względem translacji dają
$\|w_h\|_{2,K}\le Y$ oraz
$\|\nabla u(\cdot+he_j)\|_{2,K}
=\|\nabla u\|_{2,K+he_j}\le Y$.
Ponieważ $0\le\eta\le1$,
\[
 \|\nabla\varphi\|_{2,V_h}\le X+2MY.
\]
Nie szacujemy normy $\delta_h^jf$ przez stałą niezależną
od $h$, gdyż $f$ może nie mieć żadnej słabej pochodnej.
Zamiast tego, na mocy dyskretnego całkowania przez części,
\[
 \left|\int_{V_h}(\delta_h^jf)\varphi\,\mathrm d\lambda_d(x)\right|
 =\left|\int_V f\,\delta_{-h}^j\varphi\,\mathrm d\lambda_d(x)\right|
 \le F\|\partial_j\varphi\|_{2,V_h}
 \le F(X+2MY).
\]
W przedostatniej nierówności zastosowaliśmy oszacowanie
ilorazu różnicowego z lematu~\ref{lem:pde15-difference}
do przedłużenia $\varphi$ przez zero na $\R^d$:
biorąc rosnące zwarte kule i przechodząc do granicy,
dostajemy oszacowanie na całej $\R^d$ przez normę
$\partial_j\varphi$, która znika poza $K$.

Rozwijając lewą stronę~\eqref{eq:pde15-diff-eq}
i korzystając z eliptyczności na $V$, otrzymujemy
\[
 \lambda X^2
 \le 2\Lambda MXY+F(X+2MY)+DY(X+2MY).
\]
Pierwszy składnik oszacowano przez Cauchy'ego--Schwarza
i $|A\xi|\le\Lambda|\xi|$, ostatni przez ograniczenie
$\delta_h^j A$ i normę przesuniętego gradientu na $K$.
Nie wykorzystujemy symetrii macierzy $A$.
Kładąc $T=F+(2\Lambda M+D)Y$, otrzymujemy
$\lambda X^2\le TX+2MY(F+DY)$.
Nierówność Younga daje
$TX\le\lambda X^2/2+T^2/(2\lambda)$, zatem
\[
 X^2\le \frac{T^2}{\lambda^2}
       +\frac{4M}{\lambda}Y(F+DY)
 \le C_1(F^2+Y^2).
\]
Ostatnia nierówność wynika z
$(a+b)^2\le2a^2+2b^2$ oraz $2FY\le F^2+Y^2$;
$C_1$ zależy wyłącznie od $\lambda,\Lambda,M,D$.
Wybierając maksimum po skończenie wielu kierunkach $j$,
możemy przyjąć jedną stałą $C_1$ niezależną od $j,h$.

\emph{Krok 4: odzyskanie drugich słabych pochodnych.}
Ponieważ $\eta=1$ na $U$, dla każdego $j,k$ mamy
\[
 \|\delta_h^j\partial_ku\|_{2,U}^2
 =\|\partial_kw_h\|_{2,U}^2\le C_1(F^2+Y^2).
\]
Ustalmy $j,k$ i ciąg niezerowych $h_n\to0$ z $|h_n|<\rho$.
Z ograniczoności w przestrzeni Hilberta $L^2(U)$
wybieramy słabo zbieżny podciąg. Dla każdego
$\zeta\in C_c^\infty(U)$ dyskretne całkowanie przez części
oraz zbieżność $\delta_{-h_n}^j\zeta\to\partial_j\zeta$
w $L^2(U)$ (z nośnikami w $U$ dla dużych $n$)
identyfikują granicę słabą z $\partial_j\partial_ku$
dokładnie jak w dowodzie lematu~\ref{lem:pde15-criterion}.
Półciągłość normy daje
$\|\partial_j\partial_ku\|_{2,U}^2\le C_1(F^2+Y^2)$.
Dotyczy to wszystkich $d^2$ par indeksów, więc $u\in H^2(U)$.
Dodając normy pochodnych rzędów zero i jeden, otrzymujemy
\[
 \|u\|_{H^2(U)}^2
 \le\|u\|_{H^1(V)}^2+d^2C_1
       \bigl(\|f\|_{2,V}^2+\|\nabla u\|_{2,V}^2\bigr).
\]
Po wyciągnięciu pierwiastka dostajemy
~\eqref{eq:pde15-h2-est}. Wybór $U,V$ był dowolny,
więc $u\in H^2_{\mathrm{loc}}(\Omega)$.
\end{proof}

\begin{uwaga}[Zależność od $V$ w oszacowaniu]
Wszystkie przesunięcia użyte do oszacowania normy
na $U$ mieszczą się w $V$. Nie jest potrzebna wypukłość
ani regularność brzegu $V$. Stała zależy między innymi
od normy gradientu wybranej funkcji odcinającej,
która może wzrastać przy zmniejszaniu odstępu
między $\overline U$ i $\R^d\setminus V$.
\end{uwaga}

\section{Oszacowanie Caccioppolego z prawą stroną}
\label{sec:pde15-caccioppoli}
\begin{stwierdzenie}[Wersja niejednorodna]
Niech $r>0$, $B_{2r}(x_0)\Subset\Omega$,
$A\in L^\infty(B_{2r};\R^{d\times d})$ spełnia warunki
eliptyczności i ograniczoności z definicji~\ref{def:pde15-elliptic},
$f\in L^2(B_{2r})$, a $u\in H^1(B_{2r})$ jest słabym
rozwiązaniem $-\diver(A\nabla u)=f$. Istnieje $C=C(d,\lambda,\Lambda)>0$
taka, że dla każdego $k\in\R$
\[
 \int_{B_r}|\nabla u|^2\,\mathrm d\lambda_d(x)
 \le C\left(r^{-2}\int_{B_{2r}}|u-k|^2\,\mathrm d\lambda_d(x)
 +r^2\int_{B_{2r}}|f|^2\,\mathrm d\lambda_d(x)\right).
\]
W szczególności można wziąć $k=0$ lub średnią $u$ po $B_{2r}$.
\end{stwierdzenie}
\begin{proof}
Wszystkie kule mają środek $x_0$. Kładziemy $v=u-k$;
ponieważ $\nabla v=\nabla u$, funkcja $v$ spełnia to samo
równanie słabe. Wybieramy $\eta\in C_c^\infty(B_{2r})$,
$0\le\eta\le1$, równą $1$ na $B_r$ i spełniającą
$|\nabla\eta|\le C_0/r$, jak w dowodzie twierdzenia
~\ref{thm:pde14-caccioppoli}. Stała $C_0$ jest niezależna
od $r$ i $x_0$. Z lematu~\ref{lem:pde9-compact-product}
funkcja $\varphi=\eta^2v$ należy do $H^1_0(B_{2r})$
i może być aproksymowana gładkimi testami w $H^1$.
Ciągłość obu stron równania słabego, jak w definicji
~\ref{def:pde15-elliptic}, pozwala użyć $\varphi$ jako testu.
Rozwijając jej gradient, otrzymujemy
\[
 \int_{B_{2r}}\eta^2 A\nabla u\cdot\nabla u\,\mathrm d\lambda_d(x)
 =\int_{B_{2r}}f\eta^2v\,\mathrm d\lambda_d(x)
 -2\int_{B_{2r}}\eta v A\nabla u\cdot\nabla\eta\,\mathrm d\lambda_d(x).
\]
Oznaczmy $X=\|\eta\nabla u\|_{2,B_{2r}}$,
$Y=\|v\|_{2,B_{2r}}$ i $F=\|f\|_{2,B_{2r}}$.
Eliptyczność, ograniczoność $A$ i nierówność
Cauchy'ego--Schwarza dają
\[
 \lambda X^2\le FY+2\Lambda C_0r^{-1}XY.
\]
W szczególności nie musimy stosować nierówności
Poincarégo do $\eta v$. Z nierówności Younga mamy
\[
 FY\le\tfrac12 r^2 F^2+\tfrac12 r^{-2}Y^2,\qquad
 2\Lambda C_0r^{-1}XY
 \le\tfrac{\lambda}{2}X^2+
       \tfrac{2\Lambda^2 C_0^2}{\lambda r^2}Y^2.
\]
Przenosząc $\lambda X^2/2$ na lewą stronę, otrzymujemy
\[
 X^2\le \lambda^{-1}r^2 F^2+
 \left(\lambda^{-1}+\frac{4\Lambda^2 C_0^2}{\lambda^2}\right)
 r^{-2}Y^2.
\]
Ponieważ $\eta=1$ na $B_r$, daje to tezę z
$C=\lambda^{-1}+4\Lambda^2C_0^2/\lambda^2$.
Średnia $u$ na $B_{2r}$ minimalizuje $Y^2$ względem $k$,
jak wykazano w dowodzie twierdzenia~\ref{thm:pde14-caccioppoli}.
\end{proof}

\section{Regularność wyższych rzędów i bootstrap}
\label{sec:pde15-bootstrap}
\begin{twierdzenie}[Wewnętrzny bootstrap Sobolewa]
\label{thm:pde15-bootstrap}
Niech $m\in\N_0$, $\Omega\subset\R^d$ będzie obszarem,
$A\in W^{m+1,\infty}_{\mathrm{loc}}(\Omega;\R^{d\times d})$
będzie jednostajnie eliptyczna i ograniczona na zwartych
podzbiorach, $f\in H^m_{\mathrm{loc}}(\Omega)$,
a $u\in H^1_{\mathrm{loc}}(\Omega)$ niech spełnia
$-\diver(A\nabla u)=f$ słabo. Wtedy
$u\in H^{m+2}_{\mathrm{loc}}(\Omega)$.
Dla niepustych otwartych zbiorów $U\Subset V\Subset\Omega$
istnieje $C>0$ takie, że
\[
 \|u\|_{H^{m+2}(U)}
 \le C(\|f\|_{H^m(V)}+\|u\|_{H^1(V)}).
\]
Stała zależy od $m,d$, $U,V$, stałych eliptyczności na $V$
i normy $A$ w $W^{m+1,\infty}(V)$.
\end{twierdzenie}
\begin{proof}
Dowodzimy indukcyjnie jednocześnie regularności i oszacowania.
Dla $m=0$ jest to twierdzenie~\ref{thm:pde15-h2}.
Niech $m\ge1$ i załóżmy oba wnioski dla $m-1$.
Ponieważ obecne założenia implikują założenia dla $m-1$,
mamy $u\in H^{m+1}_{\mathrm{loc}}(\Omega)$.

\emph{Krok 1: różniczkowanie dystrybucyjne iloczynu.}
Ustalmy wielowskaźnik $\alpha$, $|\alpha|=m$, i połóżmy
$v_\alpha=\partial^\alpha u\in H^1_{\mathrm{loc}}(\Omega)$.
Uzasadnimy wzór Leibniza dla $A\nabla u$.
Jeżeli $a\in W^{1,\infty}(W)$ oraz $z\in H^1(W)$,
gdzie $W\Subset\Omega$ jest otwarty, to obie funkcje
należą do $H^1(W)$, gdyż $W$ ma skończoną miarę.
Z twierdzenia~\ref{thm:pde9-products} dla $p=q=2$
otrzymujemy słabą regułę iloczynu
$\partial_j(az)=(\partial_j a)z+a\partial_j z$.
Iloczyn i oba składniki pochodnej należą do $L^2(W)$,
ponieważ $a,\partial_j a\in L^\infty(W)$.
Stąd $az\in H^1(W)$ i
\[
 \|az\|_{H^1(W)}
 \le C_d\|a\|_{W^{1,\infty}(W)}\|z\|_{H^1(W)}.
\]
Stosujemy ten argument składowo również do macierzy.
Powtarzając regułę iloczynu, dla $|\gamma|\le m$ dostajemy
dystrybucyjnie
\[
 \partial^\gamma(A\nabla u)=
 \sum_{\beta\le\gamma}\binom{\gamma}{\beta}
 (\partial^\beta A)\nabla\partial^{\gamma-\beta}u .
\]
Istotnie, dla $|\gamma|<m$ każdy składnik ma oba czynniki
o regularności wystarczającej do kolejnego zastosowania
reguły: $\partial^\beta A\in W^{1,\infty}(W)$ i
$\nabla\partial^{\gamma-\beta}u\in H^1(W)$.
Po zróżniczkowaniu względem jednej współrzędnej
współczynniki sumują się według tożsamości Pascala
dla odpowiedniej składowej wielowskaźnika.
Wszystkie potrzebne pochodne istnieją już na mocy
założenia indukcyjnego, nie zakładamy więc z góry
regularności $H^{m+2}$.

\emph{Krok 2: równanie dla pochodnej i regularność.}
Ponieważ pochodne dystrybucyjne są przemienne,
z równania wyjściowego i powyższego wzoru wynika
\[
 -\diver(A\nabla v_\alpha)=\partial^\alpha f+\diver Q_\alpha,
 \qquad
 Q_\alpha=\sum_{0<\beta\le\alpha}\binom{\alpha}{\beta}
 (\partial^\beta A)\nabla\partial^{\alpha-\beta}u .
\]
Dla $0<\beta\le\alpha$ mamy
$\partial^\beta A\in W^{1,\infty}_{\mathrm{loc}}$
oraz $\nabla\partial^{\alpha-\beta}u\in
H^{|\beta|}_{\mathrm{loc}}\subset H^1_{\mathrm{loc}}$.
Z oszacowania iloczynu z kroku 1 wynika dla każdego
$W\Subset\Omega$
\[
 \|Q_\alpha\|_{H^1(W)}
 \le C_{d,m}\|A\|_{W^{m+1,\infty}(W)}\|u\|_{H^{m+1}(W)}.
\]
Dywergencja pola $H^1(W;\R^d)$ należy do $L^2(W)$
i jej norma nie przekracza $\sqrt d$ razy normy
macierzy pierwszych pochodnych w $L^2(W)$
(z nierówności Cauchy'ego--Schwarza dla sumy $d$ składników).
Zatem prawa strona $F_\alpha=\partial^\alpha f+\diver Q_\alpha$
jest lokalnie kwadratowo całkowalna.
Twierdzenie~\ref{thm:pde15-h2} zastosowane do $v_\alpha$
daje $v_\alpha\in H^2_{\mathrm{loc}}(\Omega)$.
Każdy wielowskaźnik długości $m+2$ jest sumą
wielowskaźnika długości $m$ i wielowskaźnika długości $2$,
więc wraz z założeniem indukcyjnym daje to
$u\in H^{m+2}_{\mathrm{loc}}(\Omega)$.

\emph{Krok 3: oszacowanie normy.}
Dla ustalonych $U,V$ wybieramy otwarty zbiór $W$ taki,
że $U\Subset W\Subset V$. Wtedy twierdzenie
~\ref{thm:pde15-h2} oraz oszacowanie $Q_\alpha$ dają
\[
 \|v_\alpha\|_{H^2(U)}
 \le C_1(\|F_\alpha\|_{2,W}+\|v_\alpha\|_{H^1(W)})
 \le C_2(\|f\|_{H^m(W)}+\|u\|_{H^{m+1}(W)}).
\]
Stałe są niezależne od $u,f$; maksimum po skończenie wielu
$\alpha$ można przyjąć jako wspólną stałą.
Z oszacowania indukcyjnego dla pary $W,V$ mamy
\[
 \|u\|_{H^{m+1}(W)}
 \le C_3(\|f\|_{H^{m-1}(V)}+\|u\|_{H^1(V)}).
\]
Normy $f$ na $W$ i rzędu $m-1$ na $V$ nie przekraczają
$\|f\|_{H^m(V)}$. Sumując kwadraty oszacowań
$\|v_\alpha\|_{H^2(U)}$ oraz dodając oszacowanie indukcyjne
dla pochodnych niższych rzędów, otrzymujemy tezę
po wyciągnięciu pierwiastka. Liczba wielowskaźników
zależy tylko od $d,m$; wybór $W$ zależy od $U,V$.
To zamyka indukcję również dla zależności stałych.
\end{proof}

\begin{wniosek}[Gładkość wewnętrzna]
Jeżeli $A\in C^\infty(\Omega;\R^{d\times d})$ jest lokalnie
jednostajnie eliptyczna, $f\in C^\infty(\Omega)$ i
$u\in H^1_{\mathrm{loc}}(\Omega)$ jest słabym rozwiązaniem,
to klasa $u$ ma jedyny gładki reprezentant na $\Omega$.
Reprezentant ten spełnia równanie również klasycznie.
\end{wniosek}
\begin{proof}
Każda pochodna współczynników i prawej strony jest
ograniczona na domknięciu zbioru $V\Subset\Omega$.
Poprzednie twierdzenie stosujemy dla każdego $m$
i otrzymujemy $u\in H^k_{\mathrm{loc}}(\Omega)$
dla wszystkich $k\in\N$.
Ustalmy kulę $B\Subset\Omega$. Dla $d\ge2$ zastosujemy
twierdzenie~\ref{thm:pde10-domain-embed} do wszystkich
słabych pochodnych na $B$. Jeżeli wszystkie należą do
$L^p(B)$ dla pewnego skończonego $p\ge2$, to wszystkie
należą do $W^{1,p}(B)$, bo ich kolejne słabe pochodne
również należą do $L^p(B)$. Gdy $p<d$, możemy zastąpić
$p$ przez $dp/(d-p)$, czyli zmniejszyć $1/p$ o $1/d$.
Po skończenie wielu krokach osiągamy $p\ge d$;
jeśli $p=d$, stosujemy zanurzenie do $L^{d+1}(B)$.
W ten sposób wszystkie pochodne należą do $L^p(B)$
dla pewnego $p>d$, więc $u\in W^{k,p}(B)$ dla każdego $k$.
Wniosek~\ref{cor:pde10-Morrey-domain} daje dla każdego
$\ell\ge0$ reprezentanta klasy $C^\ell$ na $B$.
Dla $d=1$ stosujemy do każdej pochodnej na przedziale $B$
jednowymiarowe twierdzenie o reprezentancie Sobolewa
(w szczególności wniosek~\ref{cor:pde9-trace-oneD}):
każda pochodna jest ciągła, a reprezentant poprzedniej
jest absolutnie ciągły i ma tę pochodną prawie wszędzie.
Podstawowe twierdzenie rachunku całkowego daje wtedy
klasyczną pochodną równą ciągłemu reprezentantowi.
Iteracja daje gładkość. W obu przypadkach ciągłe
reprezentanty tej samej klasy są równe punktowo na
przecięciach otwartych zbiorów, gdyż są równe prawie
wszędzie. Sklejamy je więc do jednego reprezentanta
$C^\infty(\Omega)$. Różnica $-\diver(A\nabla u)-f$
dla tego reprezentanta jest gładka i generuje dystrybucję
zerową; z ciągłości i iniektywności dystrybucji regularnych
(stwierdzenie~\ref{prop:pde7-injective}) znika wszędzie.
\end{proof}

\section{Regularność przy brzegu: metoda spłaszczania}
\label{sec:pde15-boundary}

\begin{lemat}[Forma dywergencyjna po spłaszczeniu brzegu]
\label{lem:pde15-flatten}
Niech $V,W\subset\R^d$ będą zbiorami otwartymi i
$\Phi:W\to V$ będzie dyfeomorfizmem klasy $C^{1,1}$
(tj. klasy $C^1$ z lokalnie lipschitzowską pochodną).
Piszemy $J=D\Phi$ i $q=|\det J|$. Jeżeli
$u\in H^1_{\mathrm{loc}}(V)$ jest słabym rozwiązaniem
$-\diver(A\nabla u)=f$ na $V$, gdzie
$A\in W^{1,\infty}_{\mathrm{loc}}(V;\R^{d\times d})$
i $f\in L^2_{\mathrm{loc}}(V)$, to
$v=u\circ\Phi\in H^1_{\mathrm{loc}}(W)$ i spełnia równanie
$-\diver(B\nabla v)=F$ słabo, gdzie
\[
 B(y)=q(y)J(y)^{-1}A(\Phi(y))J(y)^{-\mathsf T},
 \qquad F(y)=q(y)f(\Phi(y)).
\]
Mamy $B\in W^{1,\infty}_{\mathrm{loc}}(W)$ i
$F\in L^2_{\mathrm{loc}}(W)$.
Jeżeli na mierzalnym $E\subset W$ zachodzą ograniczenia
$\|J\|_{\mathrm{op}}\le M$, $\|J^{-1}\|_{\mathrm{op}}\le N$
i na $\Phi(E)$ macierz $A$ spełnia warunki z definicji
~\ref{def:pde15-elliptic} ze stałymi $\lambda,\Lambda$,
to na $E$ macierz $B$ spełnia te same warunki ze stałymi
$\lambda'= \lambda/(N^d M^2)$ i
$\Lambda'=\Lambda M^d N^2$.
Jeżeli powyższe ograniczenia $M,N$ zachodzą na całym $W$,
to składanie z $\Phi$ jest ograniczonym operatorem
$H^1(V)\to H^1(W)$ oraz $F\in L^2(W)$ dla $f\in L^2(V)$.
Jeżeli dodatkowo $D^2\Phi$ jest słabo ograniczone na $W$,
to składanie z $\Phi$ jest ograniczone również
$H^2(V)\to H^2(W)$. Analogiczne twierdzenia stosują się
do $\Phi^{-1}$ przy odpowiednich ograniczeniach jej pochodnych.
\end{lemat}
\begin{proof}
\emph{Krok 1: słaba reguła łańcuchowa.}
Z klasycznej zmiany zmiennych dla dyfeomorfizmu $C^1$
wynika zachowywanie zbiorów miary zero przez $\Phi$
i $\Phi^{-1}$. Złożenie jest więc dobrze określone dla
klas funkcji równych prawie wszędzie. Niech $W_0\Subset W$
będzie otwarty. Zwartość $\Phi(\overline{W_0})\subset V$
pozwala wybrać otwarty $V_0\Subset V$ zawierający ten obraz.
Na $\overline{W_0}$ macierze $J,J^{-1}$ są ograniczone
i $q$ jest oddzielone od zera. Jeżeli $M,N$ są ich
odpowiednimi ograniczeniami, to z oszacowań wartości
osobliwych mamy $N^{-d}\le q\le M^d$.
Dla każdej funkcji $z\in L^2(V_0)$ wzór zmiany zmiennych daje
\[
 \|z\circ\Phi\|_{2,W_0}^2
 =\int_{\Phi(W_0)}|z(x)|^2
        q(\Phi^{-1}(x))^{-1}\,\mathrm d\lambda_d(x)
 \le N^d\|z\|_{2,V_0}^2.
\]
Oszacowanie stosuje się składowo także do pól wektorowych.
Z twierdzenia~\ref{thm:pde9-mollify} wybieramy funkcje
$u_n$ gładkie na $V_0$, zbieżne do $u$ w $H^1(V_0)$.
Klasyczna reguła łańcuchowa daje
$\nabla(u_n\circ\Phi)=J^{\mathsf T}(\nabla u_n)\circ\Phi$.
Powyższe oszacowanie zapewnia zbieżność złożeń w $L^2(W_0)$
do $v=u\circ\Phi$ i ich gradientów do
$J^{\mathsf T}(\nabla u)\circ\Phi$ w $L^2(W_0)$.
Przechodząc do granicy w tożsamości definiującej słabe
pochodne, otrzymujemy
\[
 \nabla v=J^{\mathsf T}(\nabla u)\circ\Phi,
 \qquad (\nabla u)\circ\Phi=J^{-\mathsf T}\nabla v.
\]
Ponieważ $W_0$ był dowolny, wzory zachodzą lokalnie
na całym $W$. Przy globalnych ograniczeniach $M,N$
ten sam wzór zmiany zmiennych na $W$ daje
\[
 \|v\|_{2,W}\le N^{d/2}\|u\|_{2,V},\qquad
 \|\nabla v\|_{2,W}\le MN^{d/2}\|\nabla u\|_{2,V}.
\]
W szczególności $v\in H^1(W)$ dla $u\in H^1(V)$.

\emph{Krok 2: transformacja równania.}
Dla $\psi\in C_c^\infty(W)$ funkcja
$\varphi=\psi\circ\Phi^{-1}$ jest klasy $C^1$,
ma zwarty nośnik w $V$ i należy do $H^1(V)$.
Wolno jej użyć na mocy rozszerzenia klasy testów
z definicji~\ref{def:pde15-elliptic}. Reguła łańcuchowa
daje $(\nabla\varphi)\circ\Phi=J^{-\mathsf T}\nabla\psi$.
W równaniu słabym zmieniamy zmienne $x=\Phi(y)$:
\[
 \int_W q\bigl(A(\Phi)J^{-\mathsf T}\nabla v\bigr)
      \cdot J^{-\mathsf T}\nabla\psi\,\mathrm d\lambda_d(y)
 =\int_W q(f\circ\Phi)\psi\,\mathrm d\lambda_d(y).
\]
Z tożsamości $a\cdot J^{-\mathsf T}b=(J^{-1}a)\cdot b$
otrzymujemy dokładnie $\int_W B\nabla v\cdot\nabla\psi
\,\mathrm d\lambda_d(y)=\int_W F\psi\,\mathrm d\lambda_d(y)$.
Ponadto przy ograniczeniach $M,N$ na $W_0$ mamy
\[
 \|F\|_{2,W_0}^2
 =\int_{\Phi(W_0)}q(\Phi^{-1}(x))|f(x)|^2\,\mathrm d\lambda_d(x)
 \le M^d\|f\|_{2,\Phi(W_0)}^2.
\]
To daje lokalną kwadratową całkowalność $F$, a dla
globalnego $M$ również oszacowanie na całym $W$.

\emph{Krok 3: eliptyczność i regularność współczynników.}
Na $E$ mamy prawie wszędzie, dla każdego $\xi\in\R^d$,
\[
 \xi^{\mathsf T}B\xi
 =q(J^{-\mathsf T}\xi)^{\mathsf T}A(\Phi)(J^{-\mathsf T}\xi)
 \ge N^{-d}\lambda M^{-2}|\xi|^2,
 \qquad |B\xi|\le M^d N^2\Lambda|\xi|.
\]
W pierwszej nierówności użyto
$|\xi|=|J^{\mathsf T}J^{-\mathsf T}\xi|
\le M|J^{-\mathsf T}\xi|$. Zbiór wyjątkowy dla $A$
po przeniesieniu przez $\Phi^{-1}$ pozostaje zbiorem miary
zero, wspólnym dla wszystkich $\xi$; symetria $A$
nie jest potrzebna.

Lokalnie współczynniki $A$ mają lipschitzowskie
reprezentanty: dla $d\ge2$ stosujemy
twierdzenie~\ref{thm:pde10-domain-embed} na małych kulach
z $p=\infty$, a dla $d=1$ jednowymiarową teorię Sobolewa.
Na przecięciach reprezentanty zgadzają się punktowo,
gdyż są ciągłe i równe prawie wszędzie.
Złożenie $A\circ\Phi$ jest lokalnie lipschitzowskie,
ponieważ $\Phi$ jest lokalnie lipschitzowska jako
odwzorowanie klasy $C^1$. Macierz $J$ jest lokalnie
lipschitzowska z założenia; wyznacznik jest wielomianem
jej współczynników, a wartość bezwzględna jest
lipschitzowska. Macierz odwrotna jest lokalnie
lipschitzowska na mocy tożsamości
\[
 J(y)^{-1}-J(z)^{-1}
 =J(y)^{-1}(J(z)-J(y))J(z)^{-1}
\]
i lokalnej ograniczoności $J^{-1}$.
W konsekwencji wszystkie czynniki we wzorze na $B$
są lokalnie lipschitzowskie. Z twierdzenia Rademachera
i charakterystyki słabych pochodnych funkcji
lipschitzowskich (stosując na odcinkach jednowymiarową
regułę całkowania przez części i twierdzenie Fubiniego)
wynika $B\in W^{1,\infty}_{\mathrm{loc}}(W)$.

\emph{Krok 4: przenoszenie regularności drugiego rzędu.}
Dla $u\in H^2_{\mathrm{loc}}(V)$ stosujemy krok 1
także do każdej składowej $\partial_i u$. Z reguły
iloczynu dla mnożnika $W^{1,\infty}$ udowodnionej
w dowodzie twierdzenia~\ref{thm:pde15-bootstrap} wynika
\[
 \partial_{ab}v
 =\sum_{i,k=1}^d(\partial_{ik}u)\circ\Phi\,
       \partial_a\Phi_i\,\partial_b\Phi_k
   +\sum_{i=1}^d(\partial_i u)\circ\Phi\,\partial_{ab}\Phi_i.
\]
Wszystkie składniki są lokalnie kwadratowo całkowalne,
więc $v\in H^2_{\mathrm{loc}}(W)$.
Jeżeli $M,N$ i $L=\|D^2\Phi\|_\infty$ są skończone
na całym $W$, wzór zmiany zmiennych z kroku 1
i powyższa skończona suma dają
\[
 \|D^2 v\|_{2,W}
 \le C_d N^{d/2}
      \bigl(M^2\|D^2u\|_{2,V}+L\|\nabla u\|_{2,V}\bigr).
\]
Wraz z oszacowaniem $H^1$ dowodzi to ograniczoności
składania w $H^2$. Dla odwrotnej zmiany zmiennych
powtarzamy argument z $\Phi^{-1}$.
\end{proof}

\begin{lemat}[Oszacowanie pochodnych stycznych w półkuli]
\label{lem:pde15-tangential}
Niech $d\ge2$, $B_2^+=\{x\in\R^d:|x|<2,\ x_d>0\}$,
$A\in W^{1,\infty}(B_2^+;\R^{d\times d})$ będzie
jednostajnie eliptyczna, $f\in L^2(B_2^+)$ oraz
$v\in H^1(B_2^+)$ niech spełnia
\[
 \int_{B_2^+}A\nabla v\cdot\nabla\varphi\,\mathrm d\lambda_d(x)
 =\int_{B_2^+}f\varphi\,\mathrm d\lambda_d(x)
 \quad(\varphi\in H_0^1(B_2^+)).
\]
Jeżeli ślad $v$ na płaskiej części brzegu $\{x_d=0,|x|<2\}$
wynosi zero, to dla każdego $j<d$
$\partial_j\nabla v\in L^2(B_1^+)$ i
\[
 \|\partial_j\nabla v\|_{L^2(B_1^+)}
 \le C\bigl(\|f\|_{L^2(B_2^+)}+\|v\|_{H^1(B_2^+)}\bigr),
\]
gdzie $C$ zależy od $d$, eliptyczności oraz
$\|A\|_{W^{1,\infty}(B_2^+)}$.
\end{lemat}
\begin{proof}
Ustalmy $j<d$. Wybieramy $\eta\in C_c^\infty(B_{3/2})$
równą $1$ na $B_1$, i dla małych $|h|$ kładziemy
$w_h=\delta_h^jv$. Ponieważ przesunięcia są styczne do
płaskiego brzegu, $w_h$ ma tam zerowy ślad.
W szczególności $\eta^2w_h\in H_0^1(B_2^+)$.
Ponadto $-\delta_{-h}^j(\eta^2w_h)$ również należy
do $H_0^1(B_2^+)$: jego nośnik omija część sferyczną,
a ślad na płaskiej jest zerowy.
Testujemy słabe równanie ostatnią funkcją i wykonujemy
dyskretne całkowanie przez części. Otrzymujemy
\[
 \int_{B_2^+}A\nabla w_h\cdot\nabla(\eta^2w_h)
 \,\mathrm d\lambda_d(x)
 =\int_{B_2^+}(\delta_h^jf)\eta^2w_h\,\mathrm d\lambda_d(x)
 -\int_{B_2^+}(\delta_h^jA)\nabla v(\cdot+he_j)
 \cdot\nabla(\eta^2w_h)\,\mathrm d\lambda_d(x).
\]
W pierwszym członie po prawej nie różniczkujemy $f$:
\[
 \left|\int_{B_2^+}(\delta_h^jf)\eta^2w_h
 \,\mathrm d\lambda_d\right|
 =\left|\int_{B_2^+}f\,\delta_{-h}^j(\eta^2w_h)
 \,\mathrm d\lambda_d\right|
 \le\|f\|_2\|\partial_j(\eta^2w_h)\|_2.
\]
Własności ilorazów różnicowych dają
$\|w_h\|_{L^2(\operatorname{supp}\eta)}\le\|\nabla v\|_2$
i $\|\delta_h^jA\|_\infty\le\|\partial_jA\|_\infty$.
Rozwijamy $\nabla(\eta^2w_h)=\eta^2\nabla w_h+
2\eta w_h\nabla\eta$, stosujemy eliptyczność do
pierwszego członu i nierówność Younga do pozostałych.
Po pochłonięciu otrzymujemy jednostajnie w $h$
\[
 \|\eta\nabla w_h\|_2^2
 \le C(\|f\|_2^2+\|v\|_{H^1(B_2^+)}^2).
\]
Kryterium z lematu~\ref{lem:pde15-criterion}, stosowane
lokalnie w półkuli i z oszacowaniem jednostajnym aż do
płaskiego brzegu, daje $\partial_j\partial_kv\in L^2(B_1^+)$
dla każdego $k$. Nie rozciągamy funkcji przez zero
w kierunku normalnym, ponieważ takie przesunięcie
nie zachowuje równania słabego.
\end{proof}



\begin{twierdzenie}[Globalne oszacowanie $H^2$ dla Dirichleta]
\label{thm:pde15-boundary}
Niech $\Omega\subset\R^d$ będzie ograniczonym obszarem
klasy $C^{1,1}$, $A\in W^{1,\infty}(\Omega;\R^{d\times d})$
będzie jednostajnie eliptyczna i ograniczona, a
$f\in L^2(\Omega)$. Jeżeli $u\in H_0^1(\Omega)$ spełnia
$-\diver(A\nabla u)=f$ słabo, to $u\in H^2(\Omega)$ oraz
\begin{equation}\label{eq:pde15-boundary}
 \|u\|_{H^2(\Omega)}\le C\bigl(\|f\|_{L^2(\Omega)}
 +\|u\|_{H^1(\Omega)}\bigr),
\end{equation}
gdzie $C$ zależy od $d,\Omega$, stałych eliptyczności
i $\|A\|_{W^{1,\infty}(\Omega)}$.
\end{twierdzenie}
\begin{proof}
Podzielimy dowód na kroki, rozdzielając dwa różne zagadnienia:
kontrolę pochodnych stycznych i odzyskanie pochodnej normalnej.

\emph{Krok 1: skończone pokrycie brzegu.}
Ponieważ $\partial\Omega$ jest zwarty i klasy $C^{1,1}$,
można pokryć go skończenie wieloma otoczeniami, w których
po obrocie i translacji $\Omega$ ma postać
$x_d>\gamma(x')$, gdzie $\gamma\in C^{1,1}$.
Odwzorowanie $y' =x'$, $y_d=x_d-\gamma(x')$
i jego odwrotność są klasy $C^{1,1}$, a ich pierwsze
pochodne są Lipschitzowskie.
Dokładny wzór na nową macierz $B$ i prawą stronę $F$
podaje lemat~\ref{lem:pde15-flatten}; w tej zmianie zmiennych
$\det D\Phi=1$. Otrzymujemy jednostajnie eliptyczne
współczynniki $B\in W^{1,\infty}$ na mniejszym otoczeniu.
Na płaskim brzegu $\{y_d=0\}$ funkcja $v$ ma ślad zerowy.

\emph{Krok 2: pochodne styczne.}
Ustalmy funkcję odcinającą $\eta$ w półkuli $B_R^+$.
Dla $j<d$ przesunięcia $y\mapsto y+he_j$ nie przecinają
płaskiego brzegu. Iloraz różnicowy $w_h=\delta_h^jv$
ma zatem zerowy ślad na tym brzegu. Dopuszczalny jest test
$\varphi=-\delta_{-h}^j(\eta^2w_h)$ w słabym równaniu;
w przypadku $\eta^2w_h$ o zwartym nośniku wewnątrz
półkuli z wyjątkiem płaskiej części brzegu test należy do
$H_0^1(B_R^+)$.
Stosujemy lemat~\ref{lem:pde15-tangential} po przeskalowaniu
półkuli i odcięciu równania w mniejszym otoczeniu.
Otrzymujemy $\partial_j\nabla v\in L^2$ lokalnie do
płaskiego brzegu dla $j<d$ i kontrolę tych pochodnych
przez normy $\|f\|_2+\|v\|_{H^1}$.
Lemat uzasadnia dokładnie, dlaczego wolno przenieść
iloraz różnicowy na test, bez zakładania pochodnych $f$.

\emph{Krok 3: druga pochodna normalna.}
Dla $d=1$ nie ma kierunków stycznych: z równania
$-(a v')'=F$ mamy $(a v')'=-F\in L^2$, a ponieważ
$a\in W^{1,\infty}$ i $a\ge\lambda>0$, wynika
$v''\in L^2$. Dalej rozważamy $d\ge2$.
Rozpiszmy równanie (w zmiennych $y$) jako
\[
 -a_{dd}\partial_{dd}v=f+
 \sum_{(i,j)\ne(d,d)}\partial_i(a_{ij}\partial_jv)
 +(\partial_da_{dd})\partial_dv.
\]
Wszystkie składniki prawej strony oprócz
$\partial_{dd}v$ są już w $L^2$: gdy $i<d$,
korzystamy z pochodnych stycznych $\partial_i\nabla v$,
a gdy $i=d$ i $j<d$, korzystamy z przemienności
słabych pochodnych mieszanych.
Ponadto $a_{dd}\ge\lambda>0$ prawie wszędzie.
Stąd $\partial_{dd}v\in L^2$, a więc $v\in H^2$
lokalnie aż do spłaszczonego brzegu.

\emph{Krok 4: sklejenie oszacowań.}
Wnętrze $\Omega$ pokrywamy skończenie wieloma mniejszymi
zbiorami, a brzeg powyższymi mapami. Wybieramy partycję jedności
i sumujemy lokalne oszacowania. Ponieważ zmiany zmiennych
są klasy $C^{1,1}$, przeniesienie drugich słabych pochodnych
z powrotem do współrzędnych $x$ zachowuje $H^2$.
Daje to~\eqref{eq:pde15-boundary}.
\end{proof}

\begin{uwaga}[Niejednorodne dane Dirichleta]
Jeśli $g$ jest śladem pewnej funkcji $G\in H^2(\Omega)$,
to dla $v=u-G\in H_0^1(\Omega)$ otrzymujemy
$-\diver(A\nabla v)=f+\diver(A\nabla G)\in L^2(\Omega)$.
Stosując twierdzenie~\ref{thm:pde15-boundary} do $v$,
otrzymujemy $u\in H^2(\Omega)$. Zamiast wymagania
$G\in H^2$ można sformułować warunek na $g$ w odpowiedniej
przestrzeni śladów wyższego rzędu; wymaga to osobnego
twierdzenia o operatorze śladu.
\end{uwaga}


\begin{uwaga}[Granica obowiązywania identyfikacji dziedziny]
Dla $A=I$ twierdzenie~\ref{thm:pde15-boundary}
wyjaśnia identyfikację $D(-\Delta_D)=H^2(\Omega)\cap
H_0^1(\Omega)$ na obszarze klasy $C^{1,1}$, zapowiedzianą
w uwadze~\ref{rem:pde12-corners}. Na ogólnym obszarze
Lipschitzowskim taka identyfikacja może być fałszywa;
nie wolno jej używać w argumentach spektralnych
bez odpowiednich dodatkowych założeń.
\end{uwaga}

\section{Dwa kontrprzykłady: współczynniki i geometria brzegu}
\label{sec:pde15-counterexamples}
\begin{kontrprzyklad}[Skok współczynnika: $H^1$ bez $H^2$]
\label{ex:pde15-jump}
Na $\Omega=(-1,1)$ połóżmy
\[
 a(x)=\begin{cases}1,&x<0,\\2,&x\ge0,\end{cases}
 \qquad
 u(x)=\begin{cases}x,&x\le0,\\x/2,&x>0.\end{cases}
\]
Funkcja $u$ jest Lipschitzowska, więc należy do $H^1(\Omega)$.
Mamy $a(x)u'(x)=1$ prawie wszędzie, zatem
$-(a u')'=0$ dystrybucyjnie.
Jednak $u'$ ma skok równy $-1/2$ w zerze, a więc
$u''=-\tfrac12\delta_0$ jako dystrybucja; stąd
$u\notin H^2(\Omega)$. Samo ograniczenie i eliptyczność
$A$ nie gwarantują drugich pochodnych w $L^2$.
\end{kontrprzyklad}

\begin{kontrprzyklad}[Narożnik wklęsły]
\label{ex:pde15-corner}
Niech $\Omega$ będzie ograniczonym obszarem płaskim,
który w sąsiedztwie zera jest sektorem o kącie rozwarcia
$3\pi/2$, np. obszarem w kształcie litery L.
Współrzędne biegunowe dobieramy tak, że w tym sektorze
$0<\theta<3\pi/2$. Funkcja
$w(r,\theta)=r^{2/3}\sin(2\theta/3)$ jest harmoniczna dla $r>0$
i znika na ramionach sektora. Wybieramy gładką radialną
funkcję $\chi$, równą $1$ w otoczeniu zera i zerową poza
małą kulą zawartą w sąsiedztwie narożnika. Wtedy
$u=\chi w\in H_0^1(\Omega)$ oraz
$f=-\Delta u\in L^2(\Omega)$: blisko zera mamy $f=0$,
a tam, gdzie $\nabla\chi\ne0$, wszystkie składniki są gładkie.
Tymczasem $|\nabla^2u|\asymp r^{-4/3}$ blisko zera na
podsektorze dodatniej miary kątowej i dlatego
\[
 \int_{\Omega\cap B_\eps(0)}|\nabla^2u|^2\,\mathrm d\lambda_2
 \ge c\int_0^\eps r^{-8/3}r\,\mathrm dr=+\infty.
\]
Zatem $u\notin H^2(\Omega)$ mimo gładkiego operatora
$-\Delta$ i prawej strony $f\in L^2$. Obszar Lipschitzowski
nie musi mieć globalnej regularności $H^2$.
\end{kontrprzyklad}

\section{Wnioski i zakres dalszych twierdzeń}
\label{sec:pde15-summary}
Regularność $H^2$ we wnętrzu wynika z eliptyczności i
Lipschitzowskich współczynników; nie trzeba zakładać
$C^2$-regularności rozwiązania. Regularność przy brzegu
wymaga natomiast dodatkowej geometrii brzegu i poprawnie
postawionych warunków brzegowych.
W kolejnym rozdziale przejdziemy od oszacowań $L^2$ do
regularności $W^{2,p}$ oraz oszacowań Schaudera.
Nie należy przy tym mylić regularności Sobolewa z analitycznością:
$C^\infty$-gładkość współczynników gwarantuje gładkość rozwiązania
w powyższej sytuacji, lecz sama nie gwarantuje analityczności.


\chapter[Oszacowania Calderóna--Zygmunda i Schaudera]{Regularność $W^{2,p}$, całki osobliwe Calderóna--Zygmunda i oszacowania Schaudera}
\label{chap:pde16}

Rozdział~\ref{chap:pde15} dostarczył oszacowań drugich pochodnych w $L^2$.
Nie jest to koniec regularności. Dla potencjału Newtona z odpowiednimi danymi $f\in L^p$
uzyskamy oszacowania $D^2u\in L^p$ dla $1<p<\infty$,
a przy lokalnie Hölderowskich danych równania Poissona
wykażemy wewnętrzną regularność $C^{2,\alpha}$ rozwiązań klasy $C^2$.
Przejścia te nie wynikają z różniczkowania funkcji $f$: mechanizm
opiera się na strukturze całek osobliwych i na wzajemnym znoszeniu
się ich dodatnich i ujemnych części. Właśnie dlatego końcowe
przypadki $p=1$, $p=\infty$ oraz $\alpha=1$ są subtelne.

Przez $\Omega$ rozumiemy obszar w $\R^d$; $U\Subset V$ oznacza
$\overline U\subset V$ i zwartość $\overline U$. W tym rozdziale
$d\ge2$, chyba że jawnie powiemy inaczej. Wszystkie całki po
$\R^d$ i jego podzbiorach zapisujemy względem $\mathrm d\lambda_d$.
Przyjmujemy konwencję Fouriera z~\eqref{eq:pde8-convention}.

\section[Dlaczego pojawiają się całki osobliwe?]{Dlaczego $D^2(-\Delta)^{-1}$ nie jest operatorem splotu z funkcją całkowalną?}
Niech $\omega_{d-1}=\mathcal H^{d-1}(S^{d-1})$ oznacza pole
jednostkowej sfery, a $\Gamma$ rozwiązanie fundamentalne $-\Delta$:
$-\Delta\Gamma=\delta_0$. Z twierdzenia~\ref{thm:pde7-laplace} dla $d\ge3$ można wziąć
$\Gamma(x)=\bigl((d-2)\omega_{d-1}\bigr)^{-1}|x|^{2-d}$,
a dla $d=2$ funkcję $\Gamma(x)=-(2\pi)^{-1}\log|x|$.
Dla $x\ne0$ w każdym z obu przypadków
$\partial_i\Gamma(x)=-x_i/(\omega_{d-1}|x|^d)$.
Różniczkując względem $x_j$, otrzymujemy
\begin{equation}\label{eq:pde16-kernel}
 K_{ij}(x):=\partial_{ij}\Gamma(x)
 =\frac{d x_ix_j-\delta_{ij}|x|^2}
 {\omega_{d-1}|x|^{d+2}},\qquad 1\le i,j\le d,
\end{equation}
gdzie $\delta_{ij}$ oznacza deltę Kroneckera.
Jądro jest gładkie poza zerem i spełnia
$K_{ij}(tx)=t^{-d}K_{ij}(x)$ dla $t>0$.
Różniczkowanie tej tożsamości daje jednorodność gradientu
stopnia $-d-1$. Z ograniczoności obu funkcji na zwartej
sferze wynika $|K_{ij}(x)|\le C_d|x|^{-d}$ oraz
$|\nabla K_{ij}(x)|\le C_d|x|^{-d-1}$.
Dla $i\ne j$ odbicie względem $\theta_i=0$ zmienia znak
$\theta_i\theta_j$ i zachowuje miarę powierzchniową.
Dla $i=j$ całki z $\theta_i^2$ są równe przez symetrię
permutacji współrzędnych, a ich suma wynosi $\omega_{d-1}$.
Stąd
\[
 \int_{S^{d-1}}\theta_i\theta_j\,\mathrm d\mathcal H^{d-1}
 =\frac{\omega_{d-1}}d\delta_{ij},\qquad
 \int_{S^{d-1}}K_{ij}(\theta)\,\mathrm d\mathcal H^{d-1}(\theta)=0.
\]
Każde jądro jest niezerowe na pewnym otwartym podzbiorze
sfery (dla $i=j$ weźmy $\theta=e_i$, dla $i\ne j$
$\theta=(e_i+e_j)/\sqrt2$). Współrzędne biegunowe dają
\[
 \int_{B_1(0)}|K_{ij}(z)|\,\mathrm d\lambda_d(z)
 =\int_{S^{d-1}}|K_{ij}(\theta)|\,\mathrm d\mathcal H^{d-1}(\theta)
 \int_0^1 r^{-1}\,\mathrm d\lambda_1(r)=+\infty .
\]
Jeżeli $f$ jest ciągła i $f(x)\ne0$, to $|f(x-z)|$
jest oddzielone od zera w otoczeniu zera. Zwykła całka
Lebesgue'a z $K_{ij}(z)f(x-z)$ nie jest więc bezwzględnie
zbieżna, nawet dla takich $f\in C_c^\infty$.

\begin{definicja}[Główna wartość całki]
Niech $f\in C_c^1(\R^d)$, a $K$ będzie jądrem klasy $C^1$ poza
zerem, jednorodnym stopnia $-d$ i o zerowej średniej sferycznej.
Definiujemy
\[
 T_Kf(x)=\operatorname{p.v.}\!\int_{\R^d}K(x-y)f(y)\,\mathrm d\lambda_d(y)
 :=\lim_{\varepsilon\downarrow0}
 \int_{|x-y|>\varepsilon}K(x-y)f(y)\,\mathrm d\lambda_d(y).
\]
Dla $0<\varepsilon<1$ podstawiamy $z=x-y$.
Jednorodność i zerowa średnia dają zerową całkę z $K$
po każdym pierścieniu. Zatem po odjęciu $f(x)$
w pierścieniu $\varepsilon<|z|<1$ dostajemy w granicy
\[
 \begin{aligned}
 T_Kf(x)&=\int_{|z|<1}K(z)(f(x-z)-f(x))\,\mathrm d\lambda_d(z)\\
 &\quad+\int_{|z|\ge1}K(z)f(x-z)\,\mathrm d\lambda_d(z).
 \end{aligned}
\]
Obie całki są bezwzględnie zbieżne: w pierwszej
$|f(x-z)-f(x)|\le\|\nabla f\|_\infty |z|$,
więc całka radialna jest ograniczona przez stałą razy
$\int_0^1\mathrm d\lambda_1(r)$; druga ma moduł
nie większy niż $C\|f\|_1$. Uzasadnia to istnienie
granicy dla każdego $x$ oraz
$|T_Kf(x)|\le C(\|\nabla f\|_\infty+\|f\|_1)$.
\end{definicja}

Dystrybucję $\operatorname{p.v.}K$ określamy przez
\[
 \begin{aligned}
 \langle\operatorname{p.v.}K,\varphi\rangle
 &=\int_{|z|<1}K(z)(\varphi(z)-\varphi(0))\,\mathrm d\lambda_d(z)\\
 &\quad+\int_{|z|\ge1}K(z)\varphi(z)\,\mathrm d\lambda_d(z).
 \end{aligned}
\]
Dla testów o nośniku w ustalonym zwartym $L$ moduł
działania jest ograniczony przez
$C\|\nabla\varphi\|_\infty+C\lambda_d(L)\|\varphi\|_\infty$.
Liniowość i oszacowanie przez półnormy na każdym zwartym
nośniku dowodzą ciągłości dystrybucyjnej.

\begin{lemat}[Hessian rozwiązania fundamentalnego jako dystrybucja]
\label{lem:pde16-distr-kernel}
Dla każdego $1\le i,j\le d$ zachodzi w $\mathcal D'(\R^d)$
\begin{equation}\label{eq:pde16-distr-kernel}
 \partial_{ij}\Gamma
 =\operatorname{p.v.}K_{ij}-\frac{\delta_{ij}}d\delta_0.
\end{equation}
W szczególności dla $f\in C_c^\infty(\R^d)$ i $u=\Gamma*f$,
\begin{equation}\label{eq:pde16-repr}
 \partial_{ij}u=T_{ij}f-\frac{\delta_{ij}}d f,
 \qquad T_{ij}=T_{K_{ij}}.
\end{equation}
\end{lemat}
\begin{proof}
Niech $\varphi\in C_c^\infty(\R^d)$, $R>1$ tak duże,
aby nośnik $\varphi$ leżał w $B_R(0)$, oraz
$D_\varepsilon=\{x:\varepsilon<|x|<R\}$.
Z lokalnej całkowalności $\Gamma$ działanie
$\partial_{ij}\Gamma$ jest granicą całki
$\int_{D_\varepsilon}\Gamma\partial_{ij}\varphi\,\mathrm d\lambda_d$.
Całkując przez części najpierw względem $x_i$, potem
$x_j$, otrzymujemy
\[
 \begin{aligned}
 \int_{D_\varepsilon}\Gamma\partial_{ij}\varphi\,\mathrm d\lambda_d
 &=\int_{D_\varepsilon}K_{ij}\varphi\,\mathrm d\lambda_d\\
 &\quad+\int_{\partial D_\varepsilon}
 \Gamma n_i\partial_j\varphi\,\mathrm d\mathcal H^{d-1}\\
 &\quad-\int_{\partial D_\varepsilon}
 (\partial_i\Gamma)n_j\varphi\,\mathrm d\mathcal H^{d-1}.
 \end{aligned}
\]
Na zewnętrznej sferze oba wyrazy brzegowe znikają.
Na wewnętrznej sferze $x=\varepsilon\theta$ mamy
$n=-\theta$, $\partial_i\Gamma(x)=
-\theta_i/(\omega_{d-1}\varepsilon^{d-1})$,
a miara powierzchniowa skaluje się przez
$\varepsilon^{d-1}$. Ostatni wyraz wynosi więc
\[
 -\frac1{\omega_{d-1}}\int_{S^{d-1}}\theta_i\theta_j
 \varphi(\varepsilon\theta)\,\mathrm d\mathcal H^{d-1}(\theta)
 \longrightarrow-\frac{\delta_{ij}}d\varphi(0).
\]
Przedostatni wyraz jest ograniczony co do modułu przez
$C\|\partial_j\varphi\|_\infty\varepsilon|\log\varepsilon|$
dla $d=2$ i przez $C_d\|\partial_j\varphi\|_\infty\varepsilon$
dla $d\ge3$, więc znika w granicy.
Pierwsza całka dąży do działania $\operatorname{p.v.}K_{ij}$.
Daje to~\eqref{eq:pde16-distr-kernel}.
Z lematu~\ref{lem:pde7-convolution-derivatives} splot
$u=\Gamma*f$ jest gładki, $-\Delta u=f$ oraz
$\partial_{ij}u=(\partial_{ij}\Gamma)*f$.
Działanie $\operatorname{p.v.}K_{ij}$ na teście
$z\mapsto f(x-z)$ jest dokładnie $T_{ij}f(x)$,
a $\delta_0*f=f$. Otrzymujemy~\eqref{eq:pde16-repr}.
\end{proof}

\section[Oszacowanie Hilbertowskie i rozkład Calderóna--Zygmunda]{Oszacowanie $L^2$ i rozkład Calderóna--Zygmunda}
\begin{lemat}[Oszacowanie w $L^2$]
\label{lem:pde16-l2}
Dla $f\in C_c^\infty(\R^d)$ zachodzi
$\|T_{ij}f\|_2\le C_d\|f\|_2$. Operator $T_{ij}$
przedłuża się jednoznacznie do ograniczonego operatora na
$L^2(\R^d)$.
\end{lemat}
\begin{proof}
Najpierw wykażemy przynależność do $L^2$, bez jednostajnego
oszacowania normy przez $\|f\|_2$. Z definicji głównej
wartości $T_{ij}f$ jest ograniczona na całej przestrzeni.
Niech $\operatorname{supp}f\subset B_R(0)$, $R\ge1$.
Dla $|x|>2R$ mamy $|x-y|\ge|x|/2$ na nośniku $f$,
więc całka nie jest osobliwa i
$|T_{ij}f(x)|\le C_d|x|^{-d}\|f\|_1$.
Całka kwadratu tego ograniczenia poza $B_{2R}$ jest
skończona, gdyż $\int_{2R}^\infty r^{-d-1}
\,\mathrm d\lambda_1(r)<\infty$.
Na $B_{2R}$ wystarcza ograniczoność. Stąd $T_{ij}f\in L^2$,
a z~\eqref{eq:pde16-repr} także $H_{ij}:=\partial_{ij}u\in L^2$.

Z $-\Delta u=f$ i przemienności pochodnych mamy
$-\Delta H_{ij}=\partial_{ij}f$ dystrybucyjnie.
Obie strony są dystrybucjami temperowanymi; równość
w $\mathcal S'$ wynika z gęstości $C_c^\infty$ w $\mathcal S$
(twierdzenie~\ref{thm:pde8-density}). Twierdzenie~\ref{thm:pde8-dist-rules} daje
\[
 |\xi|^2\widehat H_{ij}=-\xi_i\xi_j\widehat f
 \qquad\text{w }\mathcal S'(\R^d).
\]
Obie strony są funkcjami lokalnie całkowalnymi, ponieważ
$\widehat H_{ij},\widehat f\in L^2$ z twierdzenia
Plancherela. Z jednoznaczności dystrybucji regularnych
(stwierdzenie~\ref{prop:pde7-injective}) mamy równość
prawie wszędzie. Dzieląc dla $\xi\ne0$, otrzymujemy
$\widehat H_{ij}(\xi)=-(\xi_i\xi_j/|\xi|^2)\widehat f(\xi)$
prawie wszędzie. Nie ma dodatkowego składnika skupionego
w zerze, bo $\widehat H_{ij}$ jest funkcją $L^2$.
Z~\eqref{eq:pde16-repr} wynika
$\widehat{T_{ij}f}=m_{ij}\widehat f$, gdzie
$m_{ij}(\xi)=\delta_{ij}/d-\xi_i\xi_j/|\xi|^2$
poza zerem, a wartość w zerze jest nieistotna.
Mamy $|m_{ij}|\le1+1/d\le2$. Z twierdzenia
~\ref{thm:pde8-plancherel} przy przyjętej konwencji
\[
 \|T_{ij}f\|_2=(2\pi)^{-d/2}\|m_{ij}\widehat f\|_2
 \le 2(2\pi)^{-d/2}\|\widehat f\|_2=2\|f\|_2.
\]
Argument działa również dla $d=2$, bez wymagania, aby
sam potencjał $u$ należał do $L^2$. Liniowość, gęstość
$C_c^\infty$ w $L^2$ i zupełność $L^2$ dają jednoznaczne
ciągłe przedłużenie (gęstość w $L^2$: twierdzenie
~\ref{thm:pde8-density}).
\end{proof}

\begin{lemat}[Rozkład Calderóna--Zygmunda]
\label{lem:pde16-cz-decomp}
Niech $f\in L^1(\R^d)\cap L^2(\R^d)$ i $s>0$. Istnieje
co najwyżej przeliczalna rodzina parami rozłącznych półotwartych
kostek dyadycznych $Q_k$ oraz rozkład $f=g+\sum_k b_k$
zbieżny w $L^1$ i $L^2$, taki że
\begin{enumerate}[label=(\roman*)]
\item $|g|\le 2^ds$ prawie wszędzie i $\|g\|_1\le\|f\|_1$;
\item $b_k=0$ prawie wszędzie poza $Q_k$,
 $\int_{Q_k} b_k\,\mathrm d\lambda_d=0$,
 $\sum_k\|b_k\|_1\le2\|f\|_1$;
\item $\sum_k|Q_k|\le s^{-1}\|f\|_1$ oraz
 $\|g\|_2^2\le 2^ds\|f\|_1$.
\end{enumerate}
\end{lemat}
\begin{proof}
Ustalamy siatkę kostek $2^m(k+[0,1)^d)$,
gdzie $m\in\mathbb Z$, $k\in\mathbb Z^d$.
Dwie kostki z tej siatki są rozłączne lub jedna
zawiera się w drugiej. Każda ma jedynego rodzica
o dwukrotnie większym boku. Dla $w\in L^1$
i $s>0$ nazywamy kostkę złą, jeśli średnia $|w|$
na niej przekracza $s$. Każda zła kostka zawiera
się w maksymalnej złej kostce: w łańcuchu rodziców
objętości rosną do nieskończoności, a średnie
nie przekraczają $\|w\|_1/|Q|$, więc tylko skończenie
wiele z nich może być złych. Maksymalne złe kostki
są parami rozłączne. Oznaczając przez $M_{\mathrm{dy}}w(x)$
supremum średnich $|w|$ po kostkach dyadycznych
zawierających $x$, dostajemy
\[
 \lambda_d(\{M_{\mathrm{dy}}w>s\})
 =\sum_{Q\ \mathrm{maks.}}|Q|
 \le s^{-1}\sum_{Q\ \mathrm{maks.}}\int_Q|w|\,\mathrm d\lambda_d
 \le s^{-1}\|w\|_1.
\]
Supremum jest mierzalne, gdyż siatka jest przeliczalna.

Uzasadnijmy także różniczkowanie całek po tych kostkach.
Niech $Q_n(x)$ oznacza jedyną kostkę o boku $2^{-n}$
zawierającą $x$ i $A_n w(x)$ średnią $w$ na $Q_n(x)$.
Dla $v\in C_c^\infty(\R^d)$ mamy $A_nv(x)\to v(x)$
w każdym punkcie przez ciągłość i średnicę kostki dążącą
do zera. Dla dowolnego $w\in L^1$ wybieramy
mierzalnego reprezentanta; zachodzi prawie wszędzie
\[
 \limsup_n |A_nw(x)-w(x)|
 \le M_{\mathrm{dy}}(w-v)(x)+|w(x)-v(x)|.
\]
Miara zbioru, gdzie lewa strona przekracza $\tau>0$,
jest więc ograniczona przez $4\|w-v\|_1/\tau$,
na mocy powyższej nierówności i nierówności Czebyszewa.
Z gęstości $C_c^\infty$ w $L^1$
(twierdzenie~\ref{thm:pde8-density}) miara ta wynosi zero.
Biorąc $\tau=1/l$ dla $l\in\N$, otrzymujemy
$A_nw\to w$ prawie wszędzie.

Wracamy do $w=f$ i oznaczamy maksymalne złe kostki
przez $Q_k$. Ich rodzice nie są źli, więc
\[
 s<\frac1{|Q_k|}\int_{Q_k}|f|\,\mathrm d\lambda_d
 \le 2^ds.
\]
Poza sumą tych kostek wszystkie średnie $|f|$ są
nie większe niż $s$; zastosowane do $|f|$ różniczkowanie
daje $|f|\le s$ prawie wszędzie. Kładziemy
$g=f$ poza kostkami oraz $g=f_{Q_k}$ na $Q_k$, gdzie
$f_{Q_k}=|Q_k|^{-1}\int_{Q_k}f\,\mathrm d\lambda_d$,
a $b_k=(f-f_{Q_k})\chi_{Q_k}$.
Z definicji średniej wynika zerowa całka $b_k$.
Ponadto $|f_{Q_k}|\le 2^ds$ i
\[
 \|g\|_1
 \le\int_{\R^d\setminus\bigcup_k Q_k}|f|\,\mathrm d\lambda_d
    +\sum_k\int_{Q_k}|f|\,\mathrm d\lambda_d
 =\|f\|_1,\qquad
 \|b_k\|_1\le2\int_{Q_k}|f|\,\mathrm d\lambda_d.
\]
Otrzymujemy $\|g\|_2^2\le2^ds\|f\|_1$ oraz
oszacowanie sumy norm $b_k$. Oszacowanie objętości
wynika z nierówności dla $M_{\mathrm{dy}}f$.
Każdy $b_k$ należy do $L^2$, a punktowo prawie wszędzie
$\sum_k b_k=f-g$, przy czym w każdym punkcie co najwyżej
jeden składnik jest niezerowy.
Dla $q=1,2$ z rozłączności i twierdzenia Tonellego
\[
 \sum_k\|b_k\|_q^q=\|f-g\|_q^q<\infty .
\]
Potęga $q$ normy ogona jest sumą odpowiednich potęg
norm, więc szeregi zbiegają w obu przestrzeniach.
Wybierając reprezentanta $b_k$ równego zero poza $Q_k$,
mamy jego nośnik topologiczny zawarty w $\overline{Q_k}$,
niekoniecznie w samej półotwartej kostce.
\end{proof}

\begin{twierdzenie}[Słaby typ $(1,1)$ operatorów $T_{ij}$]
\label{thm:pde16-weak11}
Istnieje $C_d>0$ taka, że dla każdego
$f\in L^1(\R^d)\cap L^2(\R^d)$ i $t>0$ zachodzi
\begin{equation}\label{eq:pde16-weak11}
 \lambda_d\bigl(\{x:|T_{ij}f(x)|>t\}\bigr)
 \le \frac{C_d}{t}\,\|f\|_{L^1(\R^d)}.
\end{equation}
\end{twierdzenie}
\begin{proof}
Piszemy $T=T_{ij}$, $K=K_{ij}$ i stosujemy
lemat~\ref{lem:pde16-cz-decomp} z $s=t$.
Niech $Q_k^*$ oznacza domkniętą kostkę o tym samym środku
$c_k$ co $Q_k$ i boku $3\sqrt d\,\ell_k$,
gdzie $\ell_k$ jest bokiem $Q_k$. Dla
$E=\bigcup_k Q_k^*$ mamy
$|E|\le(3\sqrt d)^d\sum_k|Q_k|\le C_d\|f\|_1/t$.
Dla części dobrej nierówność Czebyszewa i
lemat~\ref{lem:pde16-l2} dają
\[
 \lambda_d(\{|Tg|>t/2\})
 \le 4t^{-2}\|Tg\|_2^2
 \le C_dt^{-1}\|f\|_1.
\]

Uzasadnijmy reprezentację przedłużonego operatora $T$
dla $b_k$, które nie muszą być gładkie.
Wygładzamy $b_k$ standardowymi jądrami
$\rho_\varepsilon$ o nośnikach w $B_\varepsilon(0)$.
Funkcje $b_k*\rho_\varepsilon\in C_c^\infty$
zbiegają do $b_k$ w $L^1$ i $L^2$; wynika to
z twierdzenia~\ref{thm:pde7-approx-id} na ustalonym
zwartym otoczeniu $\overline{Q_k}$, gdyż wszystkie
nośniki są w nim zawarte dla małych $\varepsilon$.
Z ciągłości $T$ na $L^2$ i wyboru podciągu
$T(b_k*\rho_\varepsilon)\to Tb_k$ prawie wszędzie.
Dla ustalonego $x\notin Q_k^*$ odległość od
$\overline{Q_k}$ jest dodatnia; dla małych
$\varepsilon$ jądro $K(x-y)$ jest ograniczone
na wspólnym zwartym otoczeniu nośników.
Zbieżność w $L^1$ daje tam granicę zwykłych całek.
Stąd dla prawie każdego $x\notin Q_k^*$, dzięki zerowej
całce $b_k$,
\[
 Tb_k(x)=\int_{Q_k}[K(x-y)-K(x-c_k)]
 b_k(y)\,\mathrm d\lambda_d(y).
\]
Dla $y\in Q_k$ mamy $|y-c_k|\le\sqrt d\,\ell_k/2$,
a dla $x\notin Q_k^*$ zachodzi
$|x-c_k|\ge3\sqrt d\,\ell_k/2$.
Każdy punkt odcinka od $x-c_k$ do $x-y$ ma więc
moduł co najmniej $2|x-c_k|/3$. Stosując podstawowe
twierdzenie rachunku całkowego wzdłuż odcinka oraz
oszacowanie gradientu jądra, dostajemy
\[
 |Tb_k(x)|\le C_d\int_{Q_k}
 \frac{|y-c_k|\,|b_k(y)|}{|x-c_k|^{d+1}}
 \,\mathrm d\lambda_d(y)
 \le C_d\ell_k |x-c_k|^{-d-1}\|b_k\|_1.
\]
Całkując i powiększając zbiór całkowania do
$\{|x-c_k|\ge3\sqrt d\,\ell_k/2\}$, otrzymujemy
\[
 \int_{\R^d\setminus Q_k^*}|Tb_k|\,\mathrm d\lambda_d
 \le C_d\ell_k\|b_k\|_1
 \int_{3\sqrt d\,\ell_k/2}^\infty r^{-2}\,\mathrm d\lambda_1(r)
 \le C_d\|b_k\|_1.
\]
Dla skończonych sum nierówność trójkąta daje
$\int_{\R^d\setminus E}|T(\sum_{k\le N}b_k)|
\,\mathrm d\lambda_d\le C_d\sum_k\|b_k\|_1$.
Ponieważ $\sum_{k\le N}b_k\to f-g$ w $L^2$,
obrazy zbiegają w $L^2$ do $T(f-g)$.
Wybieramy podciąg zbieżny prawie wszędzie
i stosujemy lemat Fatou. Otrzymujemy
\[
 \int_{\R^d\setminus E}|T(f-g)|\,\mathrm d\lambda_d
 \le C_d\sum_k\|b_k\|_1\le C_d\|f\|_1.
\]
Z liniowości $Tf=Tg+T(f-g)$ prawie wszędzie.
Zbiór $\{|Tf|>t\}$ zawiera się z dokładnością
do zbioru miary zero w sumie $E$,
$\{|Tg|>t/2\}$ i
$\{x\notin E:|T(f-g)(x)|>t/2\}$.
Ostatnią miarę szacujemy przez $2/t$ razy
poprzednią całkę. To kończy dowód.
\end{proof}

\begin{lemat}[Interpolacja pomiędzy słabym $L^1$ i silnym $L^2$]
\label{lem:pde16-interp}
Niech $T$ będzie ograniczonym operatorem liniowym
na $L^2(\R^d)$, spełniającym oszacowanie słabego typu
$(1,1)$ na $L^1\cap L^2$ ze stałą $C_1$.
Dla każdego $1<p<2$ jego ograniczenie do
$L^1\cap L^2\cap L^p$ spełnia
$\|Tf\|_p\le C_p\|f\|_p$, gdzie $C_p$ zależy również
od stałych obu oszacowań. Ma ono jednoznaczne
ograniczone przedłużenie na $L^p$.
\end{lemat}
\begin{proof}
Dla $f\in L^1\cap L^2\cap L^p$ i $t>0$ kładziemy
$f_t^{>}=f\chi_{\{|f|>t\}}$ i
$f_t^{\le}=f\chi_{\{|f|\le t\}}$.
Obie funkcje należą do $L^1\cap L^2$.
Z liniowości i nierówności trójkąta wynika
$\{|Tf|>2t\}\subset\{|Tf_t^{>}|>t\}
\cup\{|Tf_t^{\le}|>t\}$ z dokładnością do zbioru
miary zero. Oznaczając $\|T\|_{L^2\to L^2}=C_2$,
otrzymujemy
\[
 \lambda_d(\{|Tf|>2t\})
 \le \frac{C_1}{t}\int_{\{|f|>t\}}|f|\,\mathrm d\lambda_d
 +\frac{C_2^2}{t^2}\int_{\{|f|\le t\}}|f|^2\,\mathrm d\lambda_d .
\]
Tożsamość
$|Tf(x)|^p=\int_0^{|Tf(x)|}p s^{p-1}\,\mathrm d\lambda_1(s)$
i twierdzenie Tonellego dają tożsamość warstwową
\[
 \|Tf\|_p^p=p\int_0^\infty s^{p-1}
 \lambda_d(\{|Tf|>s\})\,\mathrm d\lambda_1(s)
\]
również wtedy, gdy początkowo obie strony są nieskończone.
Po podstawieniu $s=2t$ i użyciu Tonellego otrzymujemy
\[
 \begin{aligned}
 \|Tf\|_p^p
 &\le p2^p C_1\int_{\{|f|>0\}} |f(x)|
       \int_0^{|f(x)|}t^{p-2}\,\mathrm d\lambda_1(t)
       \,\mathrm d\lambda_d(x)\\
 &\quad+p2^p C_2^2\int_{\{|f|>0\}} |f(x)|^2
       \int_{|f(x)|}^\infty t^{p-3}\,\mathrm d\lambda_1(t)
       \,\mathrm d\lambda_d(x)\\
 &=p2^p\left(\frac{C_1}{p-1}+\frac{C_2^2}{2-p}\right)
   \|f\|_p^p .
 \end{aligned}
\]
Punkty $f=0$ pomijamy, gdyż pierwotne całki po $t$
są w nich zerowe. Warunki $p>1$ i $p<2$ odpowiadają
zbieżności obu całek wewnętrznych.
Gęstość $C_c^\infty$ w $L^p$
(twierdzenie~\ref{thm:pde8-density}) i zupełność $L^p$
dają jednoznaczne przedłużenie.
Jest ono zgodne z wyjściowym operatorem $T$ na
$L^2\cap L^p$: dla $f$ z tej przestrzeni funkcje
$f_n=f\chi_{B_n(0)}\chi_{\{|f|\le n\}}$
należą do $L^1\cap L^2\cap L^p$ i zbiegają do $f$
w obu normach przez zbieżność zmajoryzowaną.
Ich obrazy zbiegają w $L^2$ do $Tf$ oraz w $L^p$
do wartości przedłużenia. Z obu zbieżności wybieramy
wspólny podciąg zbieżny prawie wszędzie, więc granice
są równe prawie wszędzie.
\end{proof}

\begin{twierdzenie}[Oszacowanie Calderóna--Zygmunda dla laplasjanu]
\label{thm:pde16-cz}
Niech $d\ge2$ i $1<p<\infty$. Dla
$f\in C_c^\infty(\R^d)$ oraz $u=\Gamma*f$ zachodzi
\begin{equation}\label{eq:pde16-cz}
 \|D^2u\|_{L^p(\R^d)}\le C_{d,p}\|f\|_{L^p(\R^d)}.
\end{equation}
Oszacowanie przedłuża operator $f\mapsto D^2(\Gamma*f)$
do $L^p$; nie oznacza ono, że samo $u$ musi należeć
do $L^p(\R^d)$.
\end{twierdzenie}
\begin{proof}
Z twierdzenia~\ref{thm:pde16-weak11},
lematu~\ref{lem:pde16-l2} i
lematu~\ref{lem:pde16-interp} wynika ograniczoność
$T_{ij}$ w $L^p$ dla $1<p\le2$.
Dla $p>2$ kładziemy $p'=p/(p-1)\in(1,2)$.
Mnożnik $m_{ij}$ wyprowadzony w dowodzie
lematu~\ref{lem:pde16-l2} jest rzeczywisty.
Z tożsamości Parsevala--Plancherela
~\ref{thm:pde8-plancherel} wynika więc
\[
 \int_{\R^d}(T_{ij}f)\overline h\,\mathrm d\lambda_d
 =\int_{\R^d}f\,\overline{T_{ij}h}\,\mathrm d\lambda_d
 \qquad(f,h\in L^2).
\]
Najpierw bierzemy $f,h\in C_c^\infty$.
Z nierówności Höldera i już udowodnionej ograniczoności
w $L^{p'}$ dostajemy
\[
 \left|\int_{\R^d}(T_{ij}f)\overline h\,\mathrm d\lambda_d\right|
 \le\|f\|_p\|T_{ij}h\|_{p'}
 \le C_{d,p'}\|f\|_p\|h\|_{p'}.
\]
Dla ustalonego testu $f$ wiemy już, że $T_{ij}f$ jest
ograniczona i zanika jak $O(|x|^{-d})$ w nieskończoności
(z dowodu lematu~\ref{lem:pde16-l2}).
Stąd należy do $L^p$ dla każdego $p>1$, gdyż
$\int_1^\infty r^{d-1-dp}\,\mathrm d\lambda_1(r)<\infty$.
Wolno zatem obliczyć jej normę przez dualność z $L^{p'}$.
Dokładniej, dla niezerowej funkcji $q=T_{ij}f\in L^p$
funkcja $h_0=|q|^{p-2}q/\|q\|_p^{p-1}$ ma normę
$L^{p'}$ równą jeden i
$\int q\overline{h_0}\,\mathrm d\lambda_d=\|q\|_p$.
Nierówność Höldera daje oszacowanie przeciwne dla
dowolnego $h$ o normie co najwyżej jeden.
Gęstość testów w $L^{p'}$ (twierdzenie
~\ref{thm:pde8-density}) i w razie potrzeby
normalizacja przybliżeń $h_0$ pozwalają ograniczyć
supremum do testów o normie nie większej niż jeden.
Otrzymujemy $\|T_{ij}f\|_p\le C_{d,p'}\|f\|_p$,
a następnie przedłużenie przez gęstość.
Także dla $p>2$ przedłużenie jest zgodne z operatorem
$L^2$ na $L^p\cap L^2$. Istotnie, obcinamy funkcję
jak w dowodzie lematu~\ref{lem:pde16-interp},
a każde obcięcie wygładzamy w obu normach
(twierdzenie~\ref{thm:pde7-approx-id} i wspólny
zwarty nośnik). Wybierając dla $n$ parametr wygładzania
tak, aby oba błędy były mniejsze niż $1/n$,
otrzymujemy ciąg testów zbieżny w obu normach.
Zbieżność obrazów i wspólny podciąg prawie wszędzie
dają zgodność przedłużeń.
Ze wzoru~\eqref{eq:pde16-repr} mamy
$\|\partial_{ij}u\|_p\le(C_{d,p}+1/d)\|f\|_p$.
Norma Frobeniusa Hessianu jest punktowo nie większa
niż suma modułów jego $d^2$ składowych, więc nierówność
trójkąta w $L^p$ daje~\eqref{eq:pde16-cz}.
Zupełność przestrzeni pól macierzowych $L^p$
uzasadnia ostatnie zdanie tezy.
\end{proof}

\begin{uwaga}[Dlaczego przypadki $1$ i $\infty$ są osobne?]
\label{rem:pde16-endpoints}
Silne oszacowanie $L^1\to L^1$ jest fałszywe nawet
dla gładkich danych. Istotnie, wybierzmy
$f\in C_c^\infty(B_1(0))$, $f\ge0$, $\int f\,\mathrm d\lambda_d=1$.
Dla $|x|>2$ mamy
\[
 |T_{12}f(x)-K_{12}(x)|
 \le C_d|x|^{-d-1}\int |y|f(y)\,\mathrm d\lambda_d(y)
 \le C_d|x|^{-d-1}.
\]
Wynika to z oszacowania gradientu jądra na odcinku
od $x$ do $x-y$, którego punkty mają moduł co najmniej
$|x|/2$. Na pewnym otwartym fragmencie sfery wokół
$(e_1+e_2)/\sqrt2$ mamy $K_{12}(\theta)\ge c_d>0$.
Dzięki jednorodności, dla odpowiednio dużego $r$ i takich
$\theta$, $T_{12}f(r\theta)\ge(c_d/2)r^{-d}$.
Całka po tym stożku jest rozbieżna, ponieważ
$\int_R^\infty r^{-1}\,\mathrm d\lambda_1(r)=\infty$.
Ponieważ $\partial_{12}(\Gamma*f)=T_{12}f$,
nie zachodzi też silne oszacowanie Hessianu w $L^1$.
Oszacowanie~\eqref{eq:pde16-weak11} jest słabszym
oszacowaniem końcowym.

Dla końca $p=\infty$ w wymiarze $d=2$ weźmy
$f=\chi_{(0,1)^2}$ i $u=\Gamma*f$.
Poza domkniętym kwadratem wolno różniczkować pod całką
(lokalnie odległość od nośnika jest dodatnia).
Dla $x=(-t,-t)$, $0<t<1/8$, otrzymujemy
\[
 \partial_{12}u(x)=\frac1\pi
 \int_{(0,1)^2}
 \frac{(t+y_1)(t+y_2)}
 {((t+y_1)^2+(t+y_2)^2)^2}\,\mathrm d\lambda_2(y).
\]
Całkowana funkcja jest nieujemna. Na podsektorze
$\pi/6<\theta<\pi/3$, $2t<r<1/2$, gdzie
$y=(r\cos\theta,r\sin\theta)$, obie współrzędne $y$
są co najmniej $r/2$, a $|x-y|\le r+\sqrt2t\le2r$.
Stąd całkowana funkcja jest nie mniejsza niż $c r^{-2}$,
z $c>0$ niezależnym od $t$. Zatem
$\partial_{12}u(-t,-t)\ge c'\log(1/(4t))\to+\infty$.
Pochodna jest ciągła w otwartym dopełnieniu kwadratu,
więc każda dostatecznie duża wartość daje także
przekroczenie zadanego progu na otwartym zbiorze
dodatniej miary. Jest to brak ograniczoności istotnej,
a nie jedynie rozbieżność na ciągu punktów.
Przykład można przenieść na gładkie dane: wygładzenia
$f_\varepsilon=f*\rho_\varepsilon$ mają normę
$L^\infty$ nie większą niż $1$, a w każdym ustalonym
punkcie poza domkniętym kwadratem
$\partial_{12}(\Gamma*f_\varepsilon)$ zbiega
do powyższej całki dzięki zbieżności w $L^1$
i temu, że dla dostatecznie małych $\varepsilon$
wszystkie nośniki leżą w ustalonym zwartym otoczeniu
kwadratu, oddalonym od rozważanego punktu.
Jednostajne oszacowanie $L^\infty$ dla gładkich danych
prowadziłoby zatem do sprzeczności.
\end{uwaga}

\section[Lokalna regularność Sobolewa i zamrażanie współczynników]{Lokalna regularność $W^{2,p}$ i zamrażanie współczynników}
W tej sekcji kule $B_r,B_R$ o promieniach $0<r<R$
są współśrodkowe. Normę macierzy drugich pochodnych
liczymy za pomocą normy Frobeniusa.

\begin{lemat}[Lokalizacja oszacowania dla laplasjanu]
\label{lem:pde16-localization}
Niech $1<p<\infty$, $B_r\Subset B_R\subset\R^d$ i
$u\in W^{2,p}(B_R)$. Wówczas
\[
 \|D^2u\|_{L^p(B_r)}\le C_{d,p}\bigl(
 \|\Delta u\|_{L^p(B_R)}+
 (R-r)^{-2}\|u\|_{L^p(B_R)}+
 (R-r)^{-1}\|\nabla u\|_{L^p(B_R)}\bigr).
\]
\end{lemat}
\begin{proof}
Najpierw wyprowadzimy z twierdzenia~\ref{thm:pde16-cz}
oszacowanie
\[
 \|D^2v\|_{L^p(\R^d)}\le C_{d,p}\|\Delta v\|_{L^p(\R^d)}
 \quad(v\in W^{2,p}(\R^d),\ \operatorname{supp}v
 \text{ zwarty}).
\]
Dla $v\in C_c^\infty(\R^d)$ z lematu
~\ref{lem:pde7-convolution-derivatives} i tożsamości
$-\Delta\Gamma=\delta_0$ mamy
$\Gamma*(-\Delta v)=(-\Delta\Gamma)*v=v$.
Stosując~\eqref{eq:pde16-cz} do $f=-\Delta v$,
otrzymujemy żądane oszacowanie dla testów.
Dla $v$ o zwartej podporze wygładzenia
$v*\rho_\varepsilon$ są testami o wspólnym zwartym
nośniku i zbiegają do $v$ w $W^{2,p}(\R^d)$
na mocy twierdzenia~\ref{thm:pde9-mollify}
(zbieżność lokalna wystarcza dzięki wspólnemu nośnikowi).
W obu normach przechodzimy do granicy.

Niech $\sigma=(R-r)/4$ i $\rho_\sigma$ będzie
nieujemnym standardowym jądrem wygładzającym
o całce $1$ i nośniku w $B_\sigma(0)$.
Kładziemy $\eta=\chi_{B_{(r+R)/2}}*\rho_\sigma$.
Wówczas $0\le\eta\le1$, $\eta=1$ w otoczeniu
$\overline{B_r}$, a nośnik $\eta$ mieści się
w $\overline{B_{(r+R)/2+\sigma}}\subset B_R$.
Dla $|\beta|=1,2$ z różniczkowania splotu mamy
$\|\partial^\beta\eta\|_\infty\le
\|\partial^\beta\rho_\sigma\|_1
=\sigma^{-|\beta|}\|\partial^\beta\rho_1\|_1$.
Zatem $|D^k\eta|\le C_d(R-r)^{-k}$ dla $k=1,2$.
Z lematu~\ref{lem:pde9-compact-product} przedłużenie
$v$ funkcji $\eta u$ przez zero należy do
$W^{2,p}(\R^d)$ i ma zwarty nośnik. Reguła Leibniza
daje w $B_R$
\[
 \Delta v=\eta\Delta u+2\nabla\eta\cdot\nabla u+u\Delta\eta,
\]
a poza $B_R$ wszystkie składniki są zerowe
(po przedłużeniu przez zero).
Ponieważ $D^2v=D^2u$ na $B_r$, poprzednie oszacowanie,
nierówność trójkąta i ograniczenia pochodnych $\eta$
dają tezę.
\end{proof}

\begin{lemat}[Interpolacja pierwszych pochodnych]
\label{lem:pde16-gradinterp}
Niech $1<p<\infty$, $B_r\Subset B_R$, $u\in W^{2,p}(B_R)$.
Dla każdego $0<h<R-r$ zachodzi
\[
 \|\nabla u\|_{L^p(B_r)}
 \le \frac{2d}{h}\|u\|_{L^p(B_R)}
 +\frac{dh}{2}\|D^2u\|_{L^p(B_R)}.
\]
W szczególności dla każdego $\varepsilon>0$ mamy
\[
 \|\nabla u\|_{L^p(B_r)}\le
 \varepsilon\|D^2u\|_{L^p(B_R)}
 +C_{\varepsilon,r,R}\|u\|_{L^p(B_R)}.
\]
\end{lemat}
\begin{proof}
Dla funkcji gładkiej i $x\in B_r$ cały odcinek
$x+[0,h]e_j$ mieści się w $B_R$.
Dwukrotne użycie podstawowego twierdzenia rachunku
całkowego i zamiana kolejności całkowania dają
\[
 \begin{aligned}
 u(x+he_j)-u(x)
 &=\int_0^h\partial_ju(x+se_j)\,\mathrm d\lambda_1(s)\\
 &=h\partial_ju(x)+\int_0^h(h-t)
       \partial_{jj}u(x+te_j)\,\mathrm d\lambda_1(t).
 \end{aligned}
\]
Po podzieleniu przez $h$ i przeniesieniu drugiego
składnika otrzymujemy wzór na $\partial_ju(x)$.
Nierówność Minkowskiego dla całki i translacja zmiennych
dają
\[
 \|\partial_ju\|_{L^p(B_r)}
 \le \frac2h\|u\|_{L^p(B_{r+h})}
 +\frac1h\int_0^h(h-t)\,\mathrm d\lambda_1(t)
       \|\partial_{jj}u\|_{L^p(B_{r+h})}.
\]
Ponieważ całka z wagi wynosi $h^2/2$, po zsumowaniu
po $j$ mamy pierwsze oszacowanie, nawet z normami
na $B_{r+h}$ po prawej stronie. Dla dowolnego
$u\in W^{2,p}(B_R)$ wybieramy $R'$ z $r+h<R'<R$.
Z twierdzenia~\ref{thm:pde9-mollify} wygładzenia
$u_\varepsilon$ są gładkie na $B_{R'}$ dla małych
parametrów i zbiegają tam w $W^{2,p}$.
Stosujemy oszacowanie z prawą stroną na $B_{r+h}$
do wygładzeń i przechodzimy w normach do granicy.
Powiększenie prawej strony do $B_R$ daje żądany wzór.
Wybór $h=\min\{(R-r)/2,2\varepsilon/d\}$
daje drugie oszacowanie z $C_{\varepsilon,r,R}=2d/h$.
\end{proof}

\begin{twierdzenie}[Wewnętrzne oszacowanie dla współczynników ciągłych]
\label{thm:pde16-variable}
Niech $\Omega\subset\R^d$ będzie obszarem, $1<p<\infty$,
$a_{ij}\in C(\Omega)$ i niech symetryczna macierz
$A(x)=(a_{ij}(x))$ będzie lokalnie jednostajnie eliptyczna:
na każdym zwartym podzbiorze $\Omega$ istnieją
$0<\lambda\le\Lambda$ takie, że
$\lambda|\xi|^2\le\xi^{\mathsf T}A(x)\xi\le\Lambda|\xi|^2$
dla wszystkich $\xi\in\R^d$.
Piszemy $A:D^2u=\sum_{i,j=1}^d a_{ij}\partial_{ij}u$.
Jeżeli $u\in W^{2,p}_{\mathrm{loc}}(\Omega)$ i
$A:D^2u=f\in L^p_{\mathrm{loc}}(\Omega)$ prawie wszędzie,
to dla niepustych otwartych zbiorów $U\Subset V\Subset\Omega$
\begin{equation}\label{eq:pde16-variable}
 \|u\|_{W^{2,p}(U)}\le C\bigl(
 \|f\|_{L^p(V)}+\|u\|_{L^p(V)}\bigr).
\end{equation}
Stała zależy od $d,p,U,V$, eliptyczności i modułu ciągłości $A$
na $V$. Twierdzenie jest tu \emph{oszacowaniem a priori}: samo
w sobie nie konstruuje rozwiązania $W^{2,p}$.
\end{twierdzenie}
\begin{proof}
\emph{Krok 1: operator o stałych współczynnikach.}
Niech $A_0=A(x_0)$, $S=A_0^{1/2}$ i
$w(z)=v(x_0+Sz)$ dla $v\in C_c^\infty(\R^d)$.
Symetria i dodatnia określoność dają
$\|S\|_{\mathrm{op}}\le\sqrt\Lambda$ oraz
$\|S^{-1}\|_{\mathrm{op}}\le\lambda^{-1/2}$.
Klasyczna reguła łańcuchowa daje
$D_z^2w=S^{\mathsf T}(D_x^2v)(x_0+Sz)S$,
więc $\Delta_z w=(A_0:D_x^2v)(x_0+Sz)$.
Z oszacowania globalnego udowodnionego w lemacie
~\ref{lem:pde16-localization} i zmiany zmiennych
$x=x_0+Sz$ mamy
\[
 \begin{aligned}
 \|D_x^2v\|_p
 &\le\lambda^{-1}|\det S|^{1/p}\|D_z^2w\|_p\\
 &\le\lambda^{-1}C_{d,p}|\det S|^{1/p}\|\Delta_z w\|_p\\
 &=\lambda^{-1}C_{d,p}\|A_0:D_x^2v\|_p.
 \end{aligned}
\]
Wyznacznik się skraca. Przez wygładzanie jak
w dowodzie lematu~\ref{lem:pde16-localization}
oszacowanie obowiązuje także dla $v\in W^{2,p}(\R^d)$
o zwartym nośniku. Oznaczmy jego stałą przez $C_0$.

\emph{Krok 2: zamrażanie i lokalizacja.}
Wybieramy $L>0$ takie, że dla każdego
$x_0\in\overline U$ mamy $\overline{B_L(x_0)}\subset V$,
a ponadto
$\sup_{B_L(x_0)}|A(x)-A(x_0)|_{\mathrm F}\le1/(2C_0)$.
Taki wspólny promień istnieje z dodatniej odległości
$\overline U$ od dopełnienia $V$ i jednostajnej
ciągłości $A$ na $\overline V$.
Dla $v\in W^{2,p}(\R^d)$ o zwartym nośniku
w $B_L(x_0)$, z nierówności Cauchy'ego--Schwarza
dla iloczynu Frobeniusa, dostajemy
\[
 \begin{aligned}
 \|D^2v\|_p
 &\le C_0\|A:D^2v\|_{p,B_L}
      +C_0\|(A_0-A):D^2v\|_{p,B_L}\\
 &\le C_0\|A:D^2v\|_{p,B_L}+\tfrac12\|D^2v\|_p.
 \end{aligned}
\]
Po pochłonięciu mamy $\|D^2v\|_p\le
2C_0\|A:D^2v\|_{p,B_L}$.
Ustalmy środek $x_0$ i pomijajmy go w oznaczeniach kul.
Dla $0<s<t<L$ używamy odcięcia równego $1$ na $B_s$
i o zwartym nośniku w $B_t$, skonstruowanego w lemacie
~\ref{lem:pde16-localization}. Dla $v=\eta u$
przedłużonego przez zero reguła Leibniza daje,
dzięki symetrii $A$,
\[
 A:D^2v=\eta f+2(A\nabla\eta)\cdot\nabla u+u\,A:D^2\eta.
\]
Stąd, używając $\|A\|_{\mathrm{op}}\le\Lambda$
i $|A|_{\mathrm F}\le\sqrt d\,\Lambda$, otrzymujemy
\[
 \|D^2u\|_{p,B_s}\le C_1\bigl(
 \|f\|_{p,B_t}+(t-s)^{-1}\|\nabla u\|_{p,B_t}
 +(t-s)^{-2}\|u\|_{p,B_t}\bigr),
\]
gdzie $C_1$ jest niezależna od $s,t,x_0,u,f$.

\emph{Krok 3: pochłanianie na zagnieżdżonych kulach.}
Niech $H(s)=\|D^2u\|_{p,B_s}$,
$F=\|f\|_{p,B_L}$ i $Y=\|u\|_{p,B_L}$.
Dla $0<s<t<L$ stosujemy poprzednie oszacowanie
z promieniem pośrednim $m=(s+t)/2$.
Z lematu~\ref{lem:pde16-gradinterp}, dla
$h=\kappa(t-s)$, gdzie $0<\kappa<1/2$, mamy
\[
 \|\nabla u\|_{p,B_m}
 \le \frac{2d}{\kappa(t-s)}Y
       +\frac{d\kappa(t-s)}2 H(t).
\]
W konsekwencji
\[
 H(s)\le C_1 d\kappa H(t)+C_1 F
       +C_1(4+4d/\kappa)(t-s)^{-2}Y .
\]
Ustalamy $\kappa=\min\{1/4,1/(8dC_1)\}$,
tak że współczynnik przy $H(t)$ jest nie większy
niż $\theta=1/8$. Zatem z pewną stałą $C_2$,
\[
 H(s)\le\theta H(t)+C_2\bigl(F+(t-s)^{-2}Y\bigr).
\]
Bierzemy $s_n=L(1-2^{-n-1})$, $n\ge0$,
więc $s_0=L/2$ i $s_{n+1}-s_n=L2^{-n-2}$.
Iterując nierówność od $n=0$ do $N-1$, dostajemy
\[
 H(L/2)\le\theta^N H(s_N)
 +C_2 F\sum_{n=0}^{N-1}\theta^n
 +16C_2 L^{-2}Y\sum_{n=0}^{N-1}(4\theta)^n.
\]
Ponieważ $u\in W^{2,p}(B_L)$, mamy $H(s_N)\le H(L)<\infty$,
więc pierwszy składnik dąży do zera.
Oba szeregi geometryczne są zbieżne, gdyż $4\theta=1/2$.
Otrzymujemy
$H(L/2)\le C_3(F+L^{-2}Y)$.
Stosując lemat~\ref{lem:pde16-gradinterp}
do $B_{L/4}\Subset B_{L/2}$ z $h=L/8$, dostajemy
także $\|\nabla u\|_{p,B_{L/4}}\le
C_d(L^{-1}Y+L H(L/2))$.
Łącznie
\[
 \|u\|_{W^{2,p}(B_{L/4})}
 \le C_{d,p,\lambda,\Lambda,L}(F+Y).
\]

\emph{Krok 4: skończone pokrycie.}
Kule $B_{L/4}(x_0)$, $x_0\in\overline U$,
pokrywają zwarty zbiór $\overline U$; wybieramy
skończone podpokrycie. Na każdej kuli prawe strony
powyższych oszacowań są ograniczone przez wspólną
stałą razy $\|f\|_{p,V}+\|u\|_{p,V}$.
Całka po $U$ modułu do potęgi $p$ każdej pochodnej
rzędu co najwyżej dwa jest nie większa od sumy
analogicznych całek po kulach pokrycia.
Po zsumowaniu i wyciągnięciu pierwiastka stopnia $p$
otrzymujemy~\eqref{eq:pde16-variable}.
Pokrycie oraz $L$ wybrano niezależnie od $u,f$,
co daje deklarowaną zależność stałej.
\end{proof}

\begin{uwaga}[Postać dywergencyjna a niedywergencyjna]
Niech $A\in W^{1,\infty}_{\mathrm{loc}}$ będzie lokalnie
jednostajnie eliptyczna, $f\in L^2_{\mathrm{loc}}$ i
$u\in H^1_{\mathrm{loc}}$ niech spełnia
$-\diver(A\nabla u)=f$ słabo. Z twierdzenia
~\ref{thm:pde15-h2} mamy $u\in H^2_{\mathrm{loc}}$.
Reguła iloczynu dla mnożnika $W^{1,\infty}$,
uzasadniona w dowodzie twierdzenia
~\ref{thm:pde15-bootstrap}, daje wtedy prawie wszędzie
\[
 -\sum_{i,j=1}^d a_{ij}\partial_{ij}u
 =f+\sum_{i,j=1}^d(\partial_i a_{ij})\partial_j u .
\]
Obie strony należą lokalnie do $L^2$. Dla wykładnika
$p>2$ sama informacja $f\in L^p_{\mathrm{loc}}$
nie uzasadnia przynależności całej prawej strony
do $L^p_{\mathrm{loc}}$ na podstawie znanej jedynie
lokalnej regularności $H^2$. Ponadto nawet gdy
prawa strona należy do $L^p_{\mathrm{loc}}$, twierdzenie
~\ref{thm:pde16-variable} jest oszacowaniem a priori
dla funkcji już należących do $W^{2,p}_{\mathrm{loc}}$.
Do wykazania nowej regularności potrzebny jest odrębny
argument; nie można zakładać jej w dowodzie samej siebie.
Jeżeli $A$ nie jest symetryczna, w członie drugiego
rzędu można zastąpić ją przez $(A+A^{\mathsf T})/2$,
gdyż Hessian jest symetryczny prawie wszędzie.
\end{uwaga}

\section{Przestrzenie Höldera i regularność w przestrzeniach Höldera}
\begin{definicja}[Norma Höldera]
\label{def:pde16-holder}
Dla $0<\alpha<1$ i niepustego zbioru $U\subset\R^d$
definiujemy
\[
 [f]_{C^{0,\alpha}(U)}=
 \sup_{\substack{x,y\in U\\x\ne y}}
 \frac{|f(x)-f(y)|}{|x-y|^\alpha}.
\]
Supremum pustej rodziny przyjmujemy za zero.
Jeżeli $U$ jest ograniczonym zbiorem otwartym i
$k\in\N_0$, przestrzeń $C^{k,\alpha}(\overline U)$
składa się z funkcji klasy $C^k(U)$, których pochodne
do rzędu $k$ przedłużają się ciągle na $\overline U$,
a pochodne rzędu $k$ mają skończoną półnormę Höldera.
Normą jest
\[
 \|f\|_{C^{k,\alpha}(\overline U)}
 =\sum_{|\beta|\le k}\|\partial^\beta f\|_{\infty,U}
 +\sum_{|\beta|=k}[\partial^\beta f]_{C^{0,\alpha}(U)}.
\]
Półnormy i normy supremum można równoważnie liczyć
na domknięciu, przez gęstość $U$ w $\overline U$
i ciągłość przedłużeń. Na $\R^d$ przez
$C^{0,\alpha}(\R^d)$ rozumiemy przestrzeń funkcji
ograniczonych o skończonej globalnej półnormie Höldera,
z normą $\|f\|_\infty+[f]_\alpha$; indeks $c$
oznacza zwarty nośnik. Dla pól o wartościach w skończenie
wymiarowej przestrzeni stosujemy analogiczne definicje
z normą tej przestrzeni.
\end{definicja}

Będziemy używali nierówności
\[
 [ab]_\alpha\le\|a\|_\infty[b]_\alpha+
                    \|b\|_\infty[a]_\alpha,
\]
wynikającej z rozpisania
$a(x)b(x)-a(y)b(y)=a(x)(b(x)-b(y))+
b(y)(a(x)-a(y))$. W szczególności pełna norma
iloczynu jest nie większa od iloczynu pełnych norm.
Dla funkcji wektorowych i macierzowych analogiczne
oszacowania otrzymujemy składowo.

\begin{lemat}[Potencjał Newtona: oszacowanie Höldera]
\label{lem:pde16-newton-holder}
Niech $0<\alpha<1$, $f\in C_c^{0,\alpha}(\R^d)$,
$\operatorname{supp}f\subset B_R(0)$ i $u=\Gamma*f$.
Wtedy $u\in C^{2,\alpha}_{\mathrm{loc}}(\R^d)$,
$-\Delta u=f$ klasycznie i dla każdego $M>0$
\[
 \|u\|_{C^{2,\alpha}(\overline{B_M(0)})}
 \le C(d,\alpha,R,M)\|f\|_{C^{0,\alpha}(\R^d)}.
\]
\end{lemat}
\begin{proof}
\emph{Krok 1: przypadek gładkich danych i norma supremum.}
Zaczynamy od $f\in C_c^\infty$. Piszemy $K=K_{ij}$.
Równość~\eqref{eq:pde16-repr} oraz zerowa średnia jądra
po sferach dają
\begin{equation}\label{eq:pde16-cancel}
 \begin{aligned}
 \partial_{ij}u(x)&=-\frac{\delta_{ij}}d f(x)\\
 &\quad+\int_{|z|<1}K(z)[f(x-z)-f(x)]\,\mathrm d\lambda_d(z)\\
 &\quad+\int_{|z|\ge1}K(z)f(x-z)\,\mathrm d\lambda_d(z).
 \end{aligned}
\end{equation}
Pierwsza całka jest bezwzględnie zbieżna i ograniczona
co do modułu przez
$C_d[f]_\alpha\int_0^1 r^{\alpha-1}\,\mathrm d\lambda_1(r)
=C_d[f]_\alpha/\alpha$.
Druga ma moduł nie większy niż $C_d\|f\|_1$,
gdyż $|K(z)|\le C_d$ dla $|z|\ge1$.
Zatem dla wszystkich $x$ mamy
$|\partial_{ij}u(x)|\le C_{d,\alpha,R}
\|f\|_{C^{0,\alpha}(\R^d)}$.

\emph{Krok 2: półnorma Höldera.}
Niech $x,x'\in B_M(0)$ i $h=|x-x'|\in(0,1/4)$.
Wybieramy $L=R+M+2$ oraz $D=B_L(x)$, które
zawiera nośnik $f$. Niech $T=T_{ij}$.
Ze znoszenia całek jądra po pierścieniach mamy
\[
 \begin{aligned}
 Tf(x)&=\int_D K(x-y)(f(y)-f(x))\,\mathrm d\lambda_d(y),\\
 Tf(x')&=\int_D K(x'-y)(f(y)-f(x'))\,\mathrm d\lambda_d(y)
           +f(x')J,
 \end{aligned}
\]
gdzie $J=\operatorname{p.v.}\int_D K(x'-y)\,\mathrm d\lambda_d(y)$.
Obie całki z różnicami funkcji są bezwzględnie
zbieżne. Ponieważ
$B_{L-h}(x')\subset D\subset B_{L+h}(x')$,
wartość główna całki po pierwszej z tych kul jest
zerowa, a resztę ograniczamy przez
\[
 |J|\le C_d\int_{L-h}^{L+h}r^{-1}\,\mathrm d\lambda_1(r)
 \le C_d h/L .
\]
Rozdzielamy $D$ na $N=B_{2h}(x)$ i $E=D\setminus N$.
Całka z $K(x-y)(f(y)-f(x))$ po $N$ ma moduł
nie większy niż $C_{d,\alpha}[f]_\alpha h^\alpha$.
To samo zachodzi dla całki z $K(x'-y)(f(y)-f(x'))$
po $N$, gdyż $N\subset B_{3h}(x')$.

Na $E$ rozpisujemy różnicę całkowanych funkcji
jako sumę
\[
 [K(x-y)-K(x'-y)](f(y)-f(x))
       +K(x'-y)(f(x')-f(x)).
\]
Dla $y\in E$ odcinek między $x-y$ i $x'-y$
jest oddalony od zera o co najmniej $|x-y|/2$.
Oszacowanie gradientu jądra i twierdzenie
o wartości średniej dają
$|K(x-y)-K(x'-y)|\le C_d h|x-y|^{-d-1}$.
Pierwszy składnik po scałkowaniu ma zatem moduł
nie większy niż
\[
 C_d h[f]_\alpha\int_{2h}^{L}r^{\alpha-2}\,
          \mathrm d\lambda_1(r)
 \le \frac{C_d}{1-\alpha}[f]_\alpha h^\alpha .
\]
Dla drugiego składnika potrzebujemy oszacowania
całki z jądra, a nie całki z jego modułu na całym $E$.
Zachodzą inkluzje (z dokładnością do brzegów)
\[
 \{y:3h<|y-x'|<L-h\}\subset E
 \subset\{y:h<|y-x'|<L+h\}.
\]
Całka po wewnętrznym pierścieniu jest zerowa.
Pozostałe części leżą w pierścieniach o promieniach
$h,3h$ oraz $L-h,L+h$. Wobec tego
\[
 \left|\int_E K(x'-y)\,\mathrm d\lambda_d(y)\right|
 \le C_d\left(\log3+\log\frac{L+h}{L-h}\right)\le C_d'.
\]
Mnożenie przez $|f(x')-f(x)|\le[f]_\alpha h^\alpha$
daje wymagane oszacowanie. Wreszcie wyraz
$f(x')J$ ma moduł nie większy niż
$C_d\|f\|_\infty h/L\le C_d\|f\|_\infty h^\alpha$.
Uwzględniając składnik $-\delta_{ij}f/d$
z~\eqref{eq:pde16-cancel}, otrzymujemy
\[
 |\partial_{ij}u(x)-\partial_{ij}u(x')|
 \le C_{d,\alpha}\|f\|_{C^{0,\alpha}} h^\alpha .
\]
Dla $h\ge1/4$ dzielimy dwukrotne oszacowanie
supremum przez $h^\alpha\ge4^{-\alpha}$.
To kończy kontrolę drugich pochodnych.

\emph{Krok 3: niższe pochodne i przejście do granicy.}
Funkcje $\Gamma$ oraz $\partial_i\Gamma$ są lokalnie
całkowalne. Dla pochodnej pierwszego rzędu wynika
to z oszacowania $|\partial_i\Gamma(z)|\le C_d|z|^{1-d}$.
Pochodna dystrybucyjna $\Gamma$ jest reprezentowana
przez tę pochodną klasyczną poza zerem: całkowanie
przez części poza $B_\varepsilon(0)$ daje wyraz brzegowy
rzędu $\varepsilon$ dla $d\ge3$ lub
$\varepsilon|\log\varepsilon|$ dla $d=2$, który zanika.
Zatem dla gładkiego $f$,
$\partial_i u=(\partial_i\Gamma)*f$.
Dla $x\in B_M(0)$ i $y\in B_R(0)$ mamy
$|x-y|<M+R$. Wobec tego
\[
 \|u\|_{\infty,B_M}+\sum_i\|\partial_i u\|_{\infty,B_M}
 \le\left(\int_{B_{M+R}(0)}
 (|\Gamma|+\sum_i|\partial_i\Gamma|)\,\mathrm d\lambda_d\right)
 \|f\|_\infty .
\]

Teraz niech $f\in C_c^{0,\alpha}$.
Bierzemy $f_n=f*\rho_{1/n}$ z nieujemnym jądrem
o całce jeden. Wszystkie nośniki leżą w $B_{R+2}(0)$,
$\|f_n\|_\infty\le\|f\|_\infty$,
$[f_n]_\alpha\le[f]_\alpha$ (przez całkowanie
różnicy wartości $f$), a
$\|f_n-f\|_\infty\le n^{-\alpha}[f]_\alpha$.
Dla $u_n=\Gamma*f_n$ powyższe oszacowanie lokalnych
całek daje jednostajną zbieżność $u_n$ do $u$
oraz $\partial_i u_n$ do $(\partial_i\Gamma)*f$
na każdej zwartej kuli. Drugie pochodne $u_n$
są jednostajnie ograniczone i mają jednostajnie
ograniczone półnormy Höldera na każdej takiej kuli.
Z twierdzenia Arzelà--Ascoliego i wyboru przekątniowego
na kulach o całkowitych promieniach otrzymujemy
podciąg, dla którego wszystkie drugie pochodne
zbiegają lokalnie jednostajnie do ciągłych funkcji.
Z podstawowego twierdzenia rachunku całkowego
na odcinkach równoległych do osi, przez przejście
do granicy najpierw dla $u_n$, a następnie dla
$\partial_i u_n$, wynika, że granice są odpowiednio
pierwszymi i drugimi pochodnymi klasycznymi $u$.
W szczególności $u\in C^2$. Nierówności Höldera
i oszacowania supremum przechodzą do granicy,
co daje deklarowaną normę $C^{2,\alpha}$.
Z równości $-\Delta u_n=f_n$ i lokalnie jednostajnej
zbieżności Hessianów wynika $-\Delta u=f$ klasycznie.
\end{proof}

\begin{twierdzenie}[Wewnętrzne oszacowanie Schaudera: model Poissona]
\label{thm:pde16-schauder-poisson}
Niech $0<\alpha<1$, $B_r\Subset B_R\subset\R^d$
będą współśrodkowymi kulami, a
$u\in C^2(B_R)\cap C^0(\overline{B_R})$ spełnia
$-\Delta u=f$ w $B_R$ dla $f\in C^{0,\alpha}(\overline{B_R})$.
Wówczas $u\in C^{2,\alpha}(\overline{B_r})$ i
\begin{equation}\label{eq:pde16-schauder-int}
 \|u\|_{C^{2,\alpha}(\overline{B_r})}
 \le C\bigl(\|u\|_{C^0(\overline{B_R})}
 +\|f\|_{C^{0,\alpha}(\overline{B_R})}\bigr).
\end{equation}
Stała zależy od $d,\alpha,r,R$.
\end{twierdzenie}
\begin{proof}
Przez translację zakładamy, że środek kul jest zerem.
Wybieramy $\rho=(r+R)/2$ i odcięcie
$\eta\in C_c^\infty(B_R)$ równe jeden na $B_\rho$,
z konstrukcji w lemacie~\ref{lem:pde16-localization}.
Niech $F$ będzie przedłużeniem $\eta f$ przez zero
na $\R^d$. Dla dwóch punktów w $B_R$ stosujemy
nierówność dla iloczynu z definicji
~\ref{def:pde16-holder}. Jeżeli $x\in B_R$ i
$y\notin B_R$, to $\eta(y)=0$ i
$|F(x)-F(y)|\le\|f\|_\infty|\eta(x)-\eta(y)|$.
Dla dwóch punktów poza $B_R$ różnica jest zerowa.
Gładkość i zwarty nośnik odcięcia dają skończoną
globalną półnormę $[\eta]_\alpha$ (dla odległości
co najwyżej jeden używamy gradientu, dla większych
normy supremum). Stąd
$\|F\|_{C^{0,\alpha}(\R^d)}
\le C_{d,\alpha,r,R}\|f\|_{C^{0,\alpha}(\overline{B_R})}$.

Z lematu~\ref{lem:pde16-newton-holder}
$v=\Gamma*F$ jest klasy $C^{2,\alpha}$ lokalnie,
$-\Delta v=F$ i pełna norma $C^{2,\alpha}$
na $\overline{B_\rho}$ jest kontrolowana przez
powyższą normę $F$.
W $B_\rho$ funkcja $h=u-v$ jest harmoniczna,
a z twierdzenia~\ref{thm:pde3-analytic} jest gładka.
Wszystkie jej pochodne są harmoniczne przez
przemienność pochodnych klasycznych. Ustalmy
$\sigma=(\rho-r)/4$ i zastosujmy oszacowanie
~\eqref{eq:pde3-derivative-bound}
do kul o promieniu $\sigma$ ze środkami
w $B_{\rho-\sigma}$. Otrzymujemy
$\|\partial_jh\|_{\infty,B_{\rho-\sigma}}
\le d\sigma^{-1}\|h\|_{\infty,B_\rho}$.
Następnie stosujemy to samo oszacowanie do pochodnych
pierwszego rzędu na $B_{\rho-2\sigma}$ i pochodnych
drugiego rzędu na $B_{\rho-3\sigma}$.
Ponieważ $r<\rho-3\sigma$, dla $1\le k\le3$ mamy
\[
 \|D^kh\|_{\infty,B_r}
 \le C_{d,k}(\rho-r)^{-k}\|h\|_{\infty,B_\rho}.
\]
Z wypukłości $B_r$ i całkowania wzdłuż odcinka
$[D^2h]_{\alpha,B_r}
\le (2r)^{1-\alpha}\|D^3h\|_{\infty,B_r}$
(z ewentualną stałą zależną od normy tensorowej).
Wreszcie
$\|h\|_{\infty,B_\rho}\le\|u\|_{\infty,B_R}
+\|v\|_{\infty,B_\rho}$.
Sumujemy oszacowania dla $u=h+v$,
włączając niższe pochodne i półnormę drugich pochodnych.
Ciągłość na domknięciu $B_r$ wynika z regularności
w jego otoczeniu. Otrzymujemy~\eqref{eq:pde16-schauder-int}
bez zakładania z góry regularności Hölderowskiej Hessianu.
\end{proof}

\begin{twierdzenie}[Skalarny współczynnik w postaci niedywergencyjnej]
\label{thm:pde16-scalar-schauder}
Niech $d\ge2$, $0<\alpha<1$, $\Omega\subset\R^d$ będzie obszarem,
$a,f\in C^{0,\alpha}_{\mathrm{loc}}(\Omega)$, przy czym
$a(x)\ge\lambda>0$ dla każdego $x\in\Omega$.
Jeżeli $u\in C^2(\Omega)$ spełnia $a(x)\Delta u=f$ w $\Omega$,
to $u\in C^{2,\alpha}_{\mathrm{loc}}(\Omega)$.
Dla niepustych otwartych zbiorów $U\Subset V\Subset\Omega$
zachodzi
\[
 \|u\|_{C^{2,\alpha}(\overline U)}
 \le C\bigl(\|f\|_{C^{0,\alpha}(\overline V)}
 +\|u\|_{C^0(\overline V)}\bigr),
\]
gdzie $C$ zależy tylko od $d,\alpha,U,V,\lambda$ i
$\|a\|_{C^{0,\alpha}(\overline V)}$.
\end{twierdzenie}
\begin{proof}
Na $\overline V$ funkcja $1/a$ należy do $C^{0,\alpha}$,
gdyż $\|1/a\|_\infty\le\lambda^{-1}$ oraz
$|a(x)^{-1}-a(y)^{-1}|
\le\lambda^{-2}[a]_\alpha|x-y|^\alpha$.
Z nierówności dla iloczynu wynika dla $g=f/a$
\[
 \|g\|_{C^{0,\alpha}(\overline V)}
 \le\lambda^{-1}\|f\|_{C^{0,\alpha}(\overline V)}
 +\lambda^{-2}\|f\|_\infty[a]_{C^{0,\alpha}(\overline V)}.
\]
Mamy $-\Delta u=-g$. Niech
$\ell=\operatorname{dist}(\overline U,\R^d\setminus V)/4>0$.
Dla każdego $x\in\overline U$ kula
$\overline{B_{2\ell}(x)}$ mieści się w $V$.
Twierdzenie~\ref{thm:pde16-schauder-poisson}
na parze $B_\ell(x),B_{2\ell}(x)$ daje jednolite
oszacowanie normy $C^{2,\alpha}$ na $B_\ell(x)$
przez $C(\|g\|_{C^{0,\alpha}(\overline V)}
+\|u\|_{\infty,V})$, niezależnie od $x$.
W szczególności wszystkie pochodne do rzędu dwa
są jednostajnie ograniczone na $\overline U$.
Dla $x,y\in\overline U$ z $0<|x-y|<\ell$
oba punkty leżą w $B_\ell(x)$, więc iloraz Höldera
pochodnych rzędu dwa ma powyższe ograniczenie.
Dla $|x-y|\ge\ell$ iloraz ograniczamy przez
$2\ell^{-\alpha}$ razy odpowiednią normę supremum.
To daje pełną normę na $\overline U$,
również gdy $U$ nie jest wypukły lub spójny.
Podstawiając oszacowanie $g$, otrzymujemy tezę.
Dowolność $U,V$ daje regularność lokalną.
\end{proof}

\begin{lemat}[Oszacowanie Schaudera \emph{a priori} i lokalizacja]
\label{lem:pde16-schauder-apriori}
Niech $0<\alpha<1$, $B_r\Subset B_R\Subset\Omega$
będą współśrodkowymi kulami, a
$A=(a_{ij})$, $b=(b_i)$, $c$ należą do
$C^{0,\alpha}(\overline{B_R})$; załóżmy, że $A$ jest
symetryczna i jednostajnie eliptyczna.
Jeżeli $v\in C^{2,\alpha}(\overline{B_R})$, to
\begin{equation}\label{eq:pde16-schauder-apriori}
 \|v\|_{C^{2,\alpha}(\overline{B_r})}
 \le C\bigl(\|A:D^2v+b\cdot\nabla v+cv\|_{C^{0,\alpha}(\overline{B_R})}
 +\|v\|_{C^0(\overline{B_R})}\bigr),
\end{equation}
gdzie $C$ zależy od $d,\alpha,r,R$, eliptyczności i norm
Hölderowskich współczynników, ale nie od $v$.
\end{lemat}
\begin{proof}
\emph{Krok 1: oszacowanie dla zamrożonego operatora.}
Dla kuli $B_\rho(x_0)$ i dowolnego pola $g$ oznaczamy
\[
 |g|_{\alpha,\rho;E}=\|g\|_{\infty,E}
                    +\rho^\alpha[g]_{\alpha,E},
 \qquad
 |w|_{k,\alpha,\rho;E}
 =\sum_{j=0}^k\rho^j\|D^jw\|_{\infty,E}
       +\rho^{k+\alpha}[D^kw]_{\alpha,E}.
\]
W tych oznaczeniach skala $\rho$ jest ustalona również
dla mniejszych zbiorów $E$. Dla iloczynów skalarnych
$|gh|_{\alpha,\rho;E}\le
|g|_{\alpha,\rho;E}|h|_{\alpha,\rho;E}$ na mocy
nierówności po definicji~\ref{def:pde16-holder}.
Dla iloczynów i kontrakcji pól o wartościach w skończenie
wymiarowych przestrzeniach wolno wprowadzić stałą
zależną jedynie od wymiaru.

Niech $w\in C^{2,\alpha}(\R^d)$ ma zwarty nośnik
w $B_\rho(x_0)$, a $A_0$ będzie stałą symetryczną
macierzą o wartościach własnych w $[\lambda,\Lambda]$.
Wówczas
\begin{equation}\label{eq:pde16-frozen}
 \rho^2|D^2w|_{\alpha,\rho;B_\rho(x_0)}
 \le C_0\rho^2|A_0:D^2w|_{\alpha,\rho;B_\rho(x_0)},
\end{equation}
gdzie $C_0$ zależy tylko od $d,\alpha,\lambda,\Lambda$.
Aby to wykazać, podstawiamy $x=x_0+\rho S z$,
$S=A_0^{1/2}$ i $\widetilde w(z)=w(x_0+\rho Sz)$.
Reguła łańcuchowa z kroku 1 dowodu twierdzenia
~\ref{thm:pde16-variable} daje
$\Delta_z\widetilde w=\rho^2(A_0:D^2w)(x_0+\rho Sz)$.
Nośnik $\widetilde w$ leży w kuli o promieniu
$\lambda^{-1/2}$, a
$B_\rho(x_0)$ po zmianie zmiennych mieści się
w tej samej kuli.
Dla $F=-\Delta_z\widetilde w$ mamy
$\Gamma*F=\widetilde w$: dla funkcji testowych wynika to
z lematu~\ref{lem:pde7-convolution-derivatives}
i $-\Delta\Gamma=\delta_0$, a dla $\widetilde w$
wygładzamy ją z jednostajną zbieżnością pochodnych
do rzędu dwa i wspólnym zwartym nośnikiem.
Lokalna całkowalność $\Gamma$ pozwala przejść
do granicy w całce potencjału.

Niech $G=A_0:D^2w$. Przedłużenie zerowe $G$
ma tę samą półnormę Höldera co ograniczenie do kuli:
dla pary punktów po przeciwnych stronach sfery
wybieramy punkt przecięcia łączącego je odcinka
ze sferą, gdzie $G=0$, i stosujemy oszacowanie
wewnątrz domkniętej kuli. Pary punktów poza kulą
dają zerową różnicę. Stąd
\[
 \|F\|_{C^{0,\alpha}(\R^d)}
 \le \rho^2\bigl(\|G\|_{\infty,B_\rho}
             +\Lambda^{\alpha/2}\rho^\alpha[G]_{\alpha,B_\rho}\bigr).
\]
Z lematu~\ref{lem:pde16-newton-holder} dla $F$
otrzymujemy oszacowanie normy $D_z^2\widetilde w$
w ustalonej kuli o promieniu $\lambda^{-1/2}+1$.
Z tożsamości
$D_x^2w(x_0+\rho Sz)=
\rho^{-2}S^{-\mathsf T}D_z^2\widetilde w(z)S^{-1}$
i nierówności
$|z-z'|\le\lambda^{-1/2}\rho^{-1}|x-x'|$
wynika~\eqref{eq:pde16-frozen}.

\emph{Krok 2: interpolacja na całej kuli.}
Wykażemy, że dla każdego $\varepsilon>0$ istnieje
$C_\varepsilon=C(d,\alpha,\varepsilon)$ takie, że
\[
 |v|_{\alpha,\rho;B_\rho}
 +\rho|\nabla v|_{\alpha,\rho;B_\rho}
 \le \varepsilon\rho^2\|D^2v\|_{\infty,B_\rho}
       +C_\varepsilon\|v\|_{\infty,B_\rho}.
\]
Wystarczy rozważyć kulę jednostkową, gdyż dla
$V(z)=v(x_0+\rho z)$ wszystkie składniki
przekształcają się zgodnie z powyższymi wagami.
Oznaczmy $Y=\|V\|_\infty$, $G=\|\nabla V\|_\infty$,
$H=\|D^2V\|_\infty$ na $B_1$.
Dla $x\in B_1$ i $0<t<1/2$ połóżmy $x_t=(1-t)x$.
Mamy $B_{t/2}(x_t)\subset B_1$.
Wzór Taylora z resztą całkową na odcinku
$x_t+[0,t/2]e_j$ daje
$|\partial_jV(x_t)|\le 4Y/t+(t/4)H$.
Ponadto z całkowania gradientu wzdłuż odcinka
$|\nabla V(x)-\nabla V(x_t)|\le tH$.
Po zsumowaniu składowych mamy
$G\le C_d(t^{-1}Y+tH)$.

Z wypukłości kuli
$|\nabla V(x)-\nabla V(y)|\le H|x-y|$,
a także ta różnica jest nie większa niż $2G$.
Dla liczb nieujemnych $a,b$ zachodzi
$\min(a,b)\le a^\alpha b^{1-\alpha}$.
Stąd $[\nabla V]_\alpha\le H^\alpha(2G)^{1-\alpha}$;
analogicznie $[V]_\alpha\le G^\alpha(2Y)^{1-\alpha}$.
W szczególności $[V]_\alpha\le G+2Y$.
Z ważonej nierówności Younga, dla każdego $\tau>0$,
mamy $H^\alpha(2G)^{1-\alpha}\le\tau H+C_{\alpha,\tau}G$.
Wybieramy najpierw $\tau=\varepsilon/2$, a następnie
$t$ tak małe, by po podstawieniu oszacowania $G$
całkowity współczynnik przy $H$ był nie większy
niż $\varepsilon$. Pozostałe wyrazy są ograniczone
przez $C_\varepsilon Y$. To dowodzi interpolacji
bez przedłużania funkcji poza kulę.

\emph{Krok 3: lokalne oszacowanie i pochłonięcie różnicy
współczynników.}
Piszemy $Lv=A:D^2v+b\cdot\nabla v+cv$ i wybieramy
kulę $\overline{B_\rho(x_0)}\subset B_R$, $\rho\le1$.
Niech $\eta$ będzie odcięciem równym jeden
na $B_{\rho/2}(x_0)$ i o zwartym nośniku w $B_\rho(x_0)$,
z konstrukcji w lemacie~\ref{lem:pde16-localization}.
Ta sama konstrukcja daje
$\|D^j\eta\|_\infty\le C_{d,j}\rho^{-j}$ dla $j\le3$.
Z wypukłości kuli i oszacowania pochodnej rzędu $j+1$
dla $j\le2$ wynika
$\rho^j|D^j\eta|_{\alpha,\rho;B_\rho}\le C_{d,\alpha}$.
Przedłużone przez zero $w=\eta v$ należy do
$C^{2,\alpha}(\R^d)$ i ma zwarty nośnik w $B_\rho$.
Istotnie, pochodne z reguły Leibniza znikają w otoczeniu
sfery, a półnormę dla par po przeciwnych stronach
sfery kontrolujemy przez punkt przecięcia odcinka
ze sferą jak w kroku 1.

Dla $A_0=A(x_0)$ mamy dokładną tożsamość
\[
 \begin{aligned}
 A_0:D^2w
 &=\eta Lv+(A_0-A):D^2w\\
 &\quad+2(A\nabla\eta)\cdot\nabla v+v A:D^2\eta\\
 &\quad-\eta b\cdot\nabla v-\eta cv .
 \end{aligned}
\]
Na $B_\rho$ zachodzi
$|A-A_0|_{\alpha,\rho;B_\rho}\le2\rho^\alpha[A]_{\alpha,B_R}$.
Zatem stosując~\eqref{eq:pde16-frozen} i nierówności
dla iloczynów, otrzymujemy dla $X=\rho^2|D^2w|_{\alpha,\rho;B_\rho}$
\[
 \begin{aligned}
 X &\le C_1\rho^2|Lv|_{\alpha,\rho;B_\rho}
       +C_1\rho^\alpha[A]_{\alpha,B_R} X\\
   &\quad+C_2\bigl(|v|_{\alpha,\rho;B_\rho}
                         +\rho|\nabla v|_{\alpha,\rho;B_\rho}\bigr).
 \end{aligned}
\]
Rzeczywiście, wyraz z $D\eta$ ma po przemnożeniu
przez $\rho^2$ czynnik co najwyżej $C\rho$,
wyraz z $D^2\eta$ czynnik $C$, a wyrazy z $b,c$
czynniki odpowiednio $C\rho^2$ przy normie gradientu
i $C\rho^2$ przy normie funkcji. Ponieważ $\rho\le1$,
ostatnie dwa wyrazy mieszczą się w zapisanym oszacowaniu.
Stałe $C_1,C_2$ zależą jedynie od danych twierdzenia,
nie od środka, promienia ani $v$.
Wybieramy $\rho_0\in(0,1]$ tak, by
$C_1\rho_0^\alpha[A]_{\alpha,B_R}\le1/2$.
Dla $\rho\le\rho_0$ pochłaniamy wyraz z $X$.
Na $B_{\rho/2}$ mamy $D^2w=D^2v$, a z kroku 2,
z parametrem odpowiednio zmniejszonym przez stałą $2C_2$,
otrzymujemy dla dowolnego $\varepsilon>0$
\[
 \begin{aligned}
 \rho^2|D^2v|_{\alpha,\rho;B_{\rho/2}}
 &\le C_3\rho^2|Lv|_{\alpha,\rho;B_\rho}
     +\varepsilon\rho^2\|D^2v\|_{\infty,B_\rho}\\
 &\quad+C_{4,\varepsilon}\|v\|_{\infty,B_\rho}.
 \end{aligned}
\]
Zauważmy, że mały składnik po prawej zawiera
jedynie normę supremum Hessianu, nie jego półnormę.

\emph{Krok 4: ważona norma i kontrola całego wnętrza.}
Dla $x\in B_R$ definiujemy
\[
 \delta(x)=\tfrac18\min\{\rho_0,
                 \operatorname{dist}(x,\partial B_R)\},\qquad
 M=\sup_{x\in B_R}\delta(x)^2|D^2v(x)|.
\]
Liczba $M$ jest skończona z założenia
$v\in C^{2,\alpha}(\overline{B_R})$.
Funkcja $\delta$ jest Lipschitzowska ze stałą $1/8$.
Kula $\overline{B_{2\delta(x)}(x)}$ mieści się w $B_R$,
jej promień nie przekracza $\rho_0$, a dla $y$
w tej kuli mamy $\delta(y)\ge3\delta(x)/4$.
Stąd
\[
 \delta(x)^2\|D^2v\|_{\infty,B_{2\delta(x)}(x)}
 \le (4/3)^2 M .
\]
W oszacowaniu z kroku 3 bierzemy $\rho=2\delta(x)$
i zachowujemy tylko normę supremum po lewej.
Po podzieleniu przez $4$ dostajemy
\[
 \delta(x)^2|D^2v(x)|
 \le C_\varepsilon(\|Lv\|_{C^{0,\alpha}(\overline{B_R})}
                  +\|v\|_{\infty,B_R})
       +\varepsilon(4/3)^2 M,
\]
ponieważ $\rho\le1$. Wybieramy na przykład
$\varepsilon=1/8$ i bierzemy supremum po $x$.
Współczynnik przy $M$ jest mniejszy niż $1/2$,
więc po pochłonięciu
\[
 M\le C(\|Lv\|_{C^{0,\alpha}(\overline{B_R})}
                  +\|v\|_{\infty,B_R})=:C Z.
\]
Na $\overline{B_r}$ zachodzi
$\delta(x)\ge\delta_*=\min\{\rho_0,R-r\}/8>0$.
W szczególności $\|D^2v\|_{\infty,B_r}\le C\delta_*^{-2}Z$.
Dla pary $x,y\in B_r$, $0<|x-y|<\delta_*$,
oba punkty leżą w $B_{\delta(x)}(x)$.
Ponowne zastosowanie kroku 3 z $\rho=2\delta(x)$
oraz oszacowanie ważonej normy dają
$\rho^{2+\alpha}[D^2v]_{\alpha,B_{\delta(x)}(x)}
\le C Z$. Iloraz Höldera tej pary jest więc
nie większy niż $C(2\delta_*)^{-2-\alpha}Z$.
Dla $|x-y|\ge\delta_*$ używamy dwukrotnej normy
supremum Hessianu pomnożonej przez $\delta_*^{-\alpha}$.
Otrzymujemy pełną kontrolę półnormy na $B_r$.
Z kroku 2 na $B_r$ wynika także oszacowanie gradientu
przez $\|D^2v\|_{\infty,B_r}$ i $\|v\|_{\infty,B_r}$,
ze stałą zależną od $r$. Przez ciągłość wszystkie
oszacowania obowiązują na domknięciu.
To dowodzi~\eqref{eq:pde16-schauder-apriori}.
\end{proof}

\begin{lemat}[Komutator wygładzania i mnożenia]
\label{lem:pde16-commutator}
Niech $0<\alpha<1$, $a\in C^{0,\alpha}(\R^d)$,
$g\in L^\infty(\R^d)$, a $\rho_\varepsilon(x)
=\varepsilon^{-d}\rho(x/\varepsilon)$, gdzie
$\rho\in C_c^\infty(B_1(0))$ jest nieujemna i ma całkę $1$.
Dla
\[
 \mathcal C_\varepsilon[a,g](x)
 =a(x)(g*\rho_\varepsilon)(x)
 -(ag)*\rho_\varepsilon(x)
\]
zachodzi jednostajnie względem $0<\varepsilon\le1$
\[
 \|\mathcal C_\varepsilon[a,g]\|_{C^{0,\alpha}(\R^d)}
 \le C_{d,\alpha,\rho}[a]_{C^{0,\alpha}(\R^d)}\|g\|_\infty.
\]
Ponadto $\|\mathcal C_\varepsilon[a,g]\|_\infty
\le \varepsilon^\alpha[a]_\alpha\|g\|_\infty$.
\end{lemat}
\begin{proof}
Oba sploty są dobrze określone w każdym punkcie:
całkujemy funkcje istotnie ograniczone z jądrem
z $L^1$. Zmiana reprezentanta $g$ na zbiorze miary zero
nie zmienia żadnej z tych całek.
Z definicji otrzymujemy
\[
 \mathcal C_\varepsilon[a,g](x)
 =\int_{\R^d}(a(x)-a(y))g(y)\rho_\varepsilon(x-y)
 \,\mathrm d\lambda_d(y).
\]
Na nośniku jądra $|x-y|\le\varepsilon$.
Z nieujemności jądra i jego całki równej jeden
wynika oszacowanie supremum
$\|\mathcal C_\varepsilon[a,g]\|_\infty
\le\varepsilon^\alpha[a]_\alpha\|g\|_\infty$.

Niech $h=|x-x'|>0$. Jeśli $h\ge\varepsilon$,
to różnica wartości komutatora jest co do modułu
nie większa niż
$2\varepsilon^\alpha[a]_\alpha\|g\|_\infty
\le2h^\alpha[a]_\alpha\|g\|_\infty$.
Dla $0<h<\varepsilon$ dodajemy i odejmujemy
$a(x')(g*\rho_\varepsilon)(x)$, uzyskując
\begin{align*}
 \mathcal C_\varepsilon[a,g](x)-\mathcal C_\varepsilon[a,g](x')
 &= (a(x)-a(x'))(g*\rho_\varepsilon)(x)\\
 &\quad+\int_{\R^d}(a(x')-a(y))g(y)
 [\rho_\varepsilon(x-y)-\rho_\varepsilon(x'-y)]
 \,\mathrm d\lambda_d(y).
\end{align*}
Pierwszy wyraz jest ograniczony przez
$[a]_\alpha\|g\|_\infty h^\alpha$.
W drugim $|x'-y|\le2\varepsilon$ wszędzie tam,
gdzie różnica jąder jest niezerowa.
Ponadto z podstawowego twierdzenia rachunku całkowego
wzdłuż odcinka między $x'-y$ i $x-y$ mamy
\[
 \begin{aligned}
 &\int_{\R^d}|\rho_\varepsilon(x-y)
               -\rho_\varepsilon(x'-y)|\,\mathrm d\lambda_d(y)\\
 &\quad\le h\int_0^1\int_{\R^d}
 |\nabla\rho_\varepsilon(x'+t(x-x')-y)|
 \,\mathrm d\lambda_d(y)\,\mathrm d\lambda_1(t)\\
 &\quad=h\|\nabla\rho_\varepsilon\|_1
 =h\varepsilon^{-1}\|\nabla\rho\|_1.
 \end{aligned}
\]
W drugiej linii zastosowano twierdzenie Tonellego,
a w trzeciej niezmienniczość miary Lebesgue'a
względem translacji oraz zmianę skali.
Zatem drugi wyraz jest ograniczony przez
$2^\alpha[a]_\alpha\|g\|_\infty\|\nabla\rho\|_1
\varepsilon^{\alpha-1}h$.
Ponieważ $\varepsilon^{\alpha-1}h
=h^\alpha(h/\varepsilon)^{1-\alpha}\le h^\alpha$,
otrzymujemy jednostajne oszacowanie półnormy Höldera.
Razem z normą supremum i $\varepsilon\le1$
daje ono deklarowaną pełną normę.
\end{proof}

\begin{twierdzenie}[Oszacowania Schaudera dla zmiennych współczynników]
\label{thm:pde16-schauder-variable}
Niech $0<\alpha<1$, $\Omega\subset\R^d$ będzie obszarem,
$a_{ij},b_i,c,f\in C^{0,\alpha}_{\mathrm{loc}}(\Omega)$,
a macierz $A=(a_{ij})$ będzie symetryczna i lokalnie
jednostajnie eliptyczna. Jeżeli $u\in C^2(\Omega)$ spełnia
\[
 Lu:=A:D^2u+b\cdot\nabla u+cu=f
 \qquad\text{w }\Omega,
\]
to $u\in C^{2,\alpha}_{\mathrm{loc}}(\Omega)$ oraz dla
niepustych otwartych zbiorów $U\Subset V\Subset\Omega$
\[
 \|u\|_{C^{2,\alpha}(\overline U)}
 \le C\bigl(\|f\|_{C^{0,\alpha}(\overline V)}
 +\|u\|_{C^0(\overline V)}\bigr),
\]
gdzie $C$ zależy od $U,V,d,\alpha$, norm
$C^{0,\alpha}(\overline V)$ współczynników i stałych
eliptyczności na $\overline V$, ale nie od $u,f$.
\end{twierdzenie}
\begin{proof}
\emph{Krok 1: lokalne przedłużenia i regularizacja.}
Wybierzmy współśrodkowe kule
$B_r\Subset B_R\Subset B_S\Subset\Omega$.
Niech $\chi\in C_c^\infty(B_S)$ będzie odcięciem
równym jeden na otoczeniu $\overline{B_R}$,
a $0\le\chi\le1$. Możemy użyć konstrukcji
z lematu~\ref{lem:pde16-localization}.
Niech $N$ będzie otwartym otoczeniem $\overline{B_R}$,
na którym $\chi=1$, zawartym w $B_S$.
Przedłużamy $\chi u$ przez zero na $\R^d$;
otrzymana funkcja $w$ jest klasy $C_c^2(\R^d)$,
gdyż $\chi$ znika w otoczeniu brzegu $B_S$.
Dla każdego współczynnika $a$ (dowolnej składowej $A,b,c$)
definiujemy $\widetilde a$ jako zerowe przedłużenie
$\chi a$, a $\widetilde f$ jako zerowe przedłużenie $\chi f$.
Funkcje te należą do $C_c^{0,\alpha}(\R^d)$:
wewnątrz $B_S$ stosujemy nierówność dla iloczynu
po definicji~\ref{def:pde16-holder}, natomiast
dla $x\in B_S$, $y\notin B_S$ mamy
$|\widetilde a(x)-\widetilde a(y)|
\le\|a\|_{\infty,B_S}|\chi(x)-\chi(y)|$.
Dla dwóch punktów poza $B_S$ różnica jest zerowa.
Analogicznie dla $\widetilde f$.
Lokalna regularność Höldera zapewnia skończoność
pełnych norm na $\overline{B_S}$: na zwartym
podzbiorze wybieramy skończone pokrycie kulami,
na których półnormy są skończone, i dla par punktów
bliskich stosujemy wspólne otoczenie z pokrycia,
a dla pozostałych ograniczenie supremum i dodatnią
liczbę Lebesgue'a pokrycia.

Kładziemy $u_\varepsilon=w*\rho_\varepsilon$.
Dla każdego wielowskaźnika $\beta$, $|\beta|\le2$,
$\partial^\beta u_\varepsilon
=(\partial^\beta w)*\rho_\varepsilon$:
różniczkowanie pod całką uzasadnia ograniczoność
ciągłych pochodnych $w$ i zwarty nośnik jądra.
Z ich jednostajnej ciągłości wynika
$\|\partial^\beta u_\varepsilon-\partial^\beta w\|_\infty
\to0$, ponieważ różnicę szacujemy przez supremum
$|\partial^\beta w(x-z)-\partial^\beta w(x)|$
dla $|z|\le\varepsilon$.
W szczególności $u_\varepsilon\to u$
w $C^2(\overline{B_R})$.

\emph{Krok 2: dokładne równanie po wygładzeniu.}
Bierzemy $0<\varepsilon<\min\{1,
\operatorname{dist}(\overline{B_R},\R^d\setminus N)\}$.
Dla $x\in\overline{B_R}$ nośnik jądra w zmiennej
całkowania $y$ leży w $N$, gdzie $w=u$ i
$\widetilde a=a$, $\widetilde f=f$.
W szczególności, dla $|\beta|\le2$,
\[
 a(x)\partial^\beta u_\varepsilon(x)
 =((\widetilde a\partial^\beta w)*\rho_\varepsilon)(x)
    +\mathcal C_\varepsilon[\widetilde a,\partial^\beta w](x).
\]
Sumując tożsamości dla wszystkich współczynników,
w pierwszych składnikach wykorzystujemy $Lu=f$
na $N$. Dostajemy na $\overline{B_R}$
\[
 \begin{aligned}
 Lu_\varepsilon
 &=\widetilde f*\rho_\varepsilon\\
 &\quad+\sum_{i,j=1}^d
       \mathcal C_\varepsilon[\widetilde a_{ij},\partial_{ij}w]\\
 &\quad+\sum_{i=1}^d
       \mathcal C_\varepsilon[\widetilde b_i,\partial_i w]
       +\mathcal C_\varepsilon[\widetilde c,w].
 \end{aligned}
\]
Współczynniki operatora po lewej stronie
nie zostały wygładzone.
Wszystkie drugie argumenty komutatorów są globalnie
ograniczone, gdyż $w\in C_c^2$.
Lemat~\ref{lem:pde16-commutator} daje jednostajne
oszacowania ich norm $C^{0,\alpha}(\R^d)$.
Także $\|\widetilde f*\rho_\varepsilon\|_\infty
\le\|\widetilde f\|_\infty$, a dla dowolnych $x,x'$
\[
 \begin{aligned}
 &|(\widetilde f*\rho_\varepsilon)(x)
       -(\widetilde f*\rho_\varepsilon)(x')|\\
 &\quad\le\int_{\R^d}
 |\widetilde f(x-z)-\widetilde f(x'-z)|
 \rho_\varepsilon(z)\,\mathrm d\lambda_d(z)
 \le[\widetilde f]_\alpha|x-x'|^\alpha .
 \end{aligned}
\]
Zatem prawa strona dokładnego równania jest
jednostajnie ograniczona w $C^{0,\alpha}(\overline{B_R})$.
Stała na tym etapie może zależeć od normy $C^2$
funkcji $w$, ale nie od $\varepsilon$.
Nie zakładamy Hölderowskiej ciągłości $D^2u$.

\emph{Krok 3: uzyskanie regularności.}
Każde $u_\varepsilon$ jest gładkie na $\R^d$,
a więc należy do $C^{2,\alpha}(\overline{B_R})$.
Możemy zastosować lemat
~\ref{lem:pde16-schauder-apriori} do $u_\varepsilon$,
z tymi samymi współczynnikami i stałą niezależną
od parametru wygładzania. Otrzymujemy
\[
 \|u_\varepsilon\|_{C^{2,\alpha}(\overline{B_r})}
 \le C\bigl(\|Lu_\varepsilon\|_{C^{0,\alpha}(\overline{B_R})}
                  +\|w\|_\infty\bigr)\le C_u.
\]
Dla każdych różnych $x,y\in\overline{B_r}$
przechodzimy do granicy w nierówności
$|D^2u_\varepsilon(x)-D^2u_\varepsilon(y)|
\le C_u|x-y|^\alpha$ dzięki jednostajnej zbieżności
drugich pochodnych. W ten sposób
$u\in C^{2,\alpha}(\overline{B_r})$.
Dowolność wewnętrznych kul daje
$u\in C^{2,\alpha}_{\mathrm{loc}}(\Omega)$.

\emph{Krok 4: oszacowanie niezależne od normy $C^2$ rozwiązania.}
W poprzednim kroku nie uzyskaliśmy jeszcze żądanej
zależności stałej, ponieważ oszacowanie komutatorów
zawierało normę $C^2$ funkcji $w$.
Teraz stosujemy lemat~\ref{lem:pde16-schauder-apriori}
bezpośrednio do $u$, którego regularność Höldera
już wykazano. Niech
$\ell=\operatorname{dist}(\overline U,\R^d\setminus V)/4$.
Dla każdego $x\in\overline U$ mamy
$\overline{B_{2\ell}(x)}\subset V$.
Na tej kuli $u$ należy do $C^{2,\alpha}$:
wynika to z uzyskanej regularności lokalnej
i zwartości domknięcia kuli, jak w kroku 1.
Ponieważ $Lu=f$, lemat daje
\[
 \|u\|_{C^{2,\alpha}(\overline{B_\ell(x)})}
 \le C\bigl(\|f\|_{C^{0,\alpha}(\overline V)}
                 +\|u\|_{\infty,V}\bigr).
\]
Stała jest niezależna od $x$: promienie są ustalone,
a normy współczynników i stałe eliptyczności
kontrolujemy na $\overline V$.
Otrzymujemy oszacowanie supremum wszystkich pochodnych
do rzędu dwa na $\overline U$.
Dla $x,y\in\overline U$, $0<|x-y|<\ell$, oba punkty
leżą w $B_\ell(x)$, gdzie obowiązuje powyższa
kontrola półnormy Höldera.
Dla $|x-y|\ge\ell$ iloraz Höldera ograniczamy
przez dwukrotne supremum odpowiedniej pochodnej
pomnożone przez $\ell^{-\alpha}$.
Sumujemy oszacowania składowych, otrzymując
żądaną pełną normę na $\overline U$.
\end{proof}

\begin{uwaga}[Granice formułowanych twierdzeń]
\label{rem:pde16-boundary}
W tym rozdziale udowodniliśmy oszacowania wewnętrzne,
nie zaś globalne twierdzenia regularności brzegowej
w $W^{2,p}$ i $C^{2,\alpha}$.
Typowymi założeniami w globalnej teorii Dirichleta
są brzeg klasy $C^{1,1}$ dla oszacowań $W^{2,p}$
oraz brzeg klasy $C^{2,\alpha}$ i dane brzegowe będące
śladem funkcji klasy $C^{2,\alpha}$ dla oszacowań
Schaudera; potrzebne są także odpowiednie założenia
o operatorze i prawej stronie aż do brzegu.
Nie są to twierdzenia o konieczności tych warunków
dla każdego rozwiązania.
Dla $p=2$ i postaci dywergencyjnej pełne twierdzenie
brzegowe z dokładnymi hipotezami udowodniono
w rozdziale poprzednim (twierdzenie~\ref{thm:pde15-boundary}).
Dowód ogólnych globalnych oszacowań wymaga odrębnej
kontroli przy brzegu, której nie zastępuje rezultat
wewnętrzny: stałe tego ostatniego mogą rosnąć
przy zmniejszaniu odległości od brzegu.
\end{uwaga}

\begin{uwaga}[Postać równania i przestrzeń rozwiązania]
\label{rem:pde-v42-two-regularity-classes}
Twierdzenie~\ref{thm:pde15-boundary} dotyczy słabych
rozwiązań $u\in H_0^1(\Omega)$ równania w postaci
dywergencyjnej $-\diver(A\nabla u)=f$ na ograniczonym
obszarze klasy $C^{1,1}$, przy jednostajnej eliptyczności,
$A\in W^{1,\infty}$ i $f\in L^2$.
Natomiast twierdzenie~\ref{thm:pde16-schauder-variable}
dotyczy rozwiązań klasy $C^2$ równania
w postaci niedywergencyjnej $Lu=f$ przy lokalnie
jednostajnie eliptycznej symetrycznej macierzy $A$
oraz lokalnie Hölderowskich współczynnikach i prawej
stronie, a jego wniosek jest lokalny.
Nie wolno przenosić tezy $C^{2,\alpha}$ na słabe
rozwiązania postaci dywergencyjnej z jedynie
$C^{0,\alpha}$-współczynnikami: rozwinięcie dywergencji
do równania niedywergencyjnego wymaga osobnego
uzasadnienia, gdyż pochodne współczynników nie muszą
być funkcjami ograniczonymi. Samo oszacowanie
$W^{2,p}$ z twierdzenia~\ref{thm:pde16-variable}
również nie zastępuje argumentu uzyskania regularności.
\end{uwaga}

\section{Kontrprzykłady i granice regularności}
\begin{kontrprzyklad}[Brak ciągłości $D^2u$ przy ograniczonej prawej stronie]
Na przedziale $(-1,1)$ niech $u(x)=-\tfrac12x|x|$.
Dla $x>0$ mamy $u(x)=-x^2/2$, a dla $x<0$
mamy $u(x)=x^2/2$. W zerze iloraz różnicowy
$u(h)/h=-|h|/2$ dąży do zera.
Zatem $u$ jest klasy $C^1$, a $u'(x)=-|x|$ na całym
przedziale. Funkcja $u'$ jest Lipschitzowska i jej
słaba pochodna wynosi $-\operatorname{sgn}(x)$:
dla testu $\varphi\in C_c^\infty((-1,1))$
całkujemy przez części na obu półprzedziałach,
a wyrazy w zerze znikają, ponieważ $u'(0)=0$.
W konsekwencji
\[
 -u''=\operatorname{sgn}(x)\quad\text{prawie wszędzie},
 \qquad u\in W^{2,\infty}((-1,1)).
\]
Na ograniczonym przedziale wynika stąd również
$u\in W^{2,p}((-1,1))$ dla każdego $1\le p<\infty$.
Jednostronne ilorazy różnicowe funkcji $u'$ w zerze
mają granice $-1$ i $1$, więc druga pochodna klasyczna
w zerze nie istnieje. Tym samym $u\notin C^2((-1,1))$.
W wymiarze $d\ge2$ na kostce $(-1,1)^d$ bierzemy
$U(x)=-x_1|x_1|/2$. Z twierdzenia Fubiniego
i jednowymiarowego rachunku słabych pochodnych
wynika $\partial_{11}U=-\operatorname{sgn}(x_1)$
oraz zerowe pozostałe drugie pochodne.
Otrzymujemy $-\Delta U=\operatorname{sgn}(x_1)$,
$U\in W^{2,\infty}$, ale $U\notin C^2$.
\end{kontrprzyklad}

\begin{kontrprzyklad}[Ostrość wykładnika Höldera]
Dla $0<\alpha<1$ połóżmy na $(-1,1)$
$u(x)=-|x|^{2+\alpha}/((1+\alpha)(2+\alpha))$.
Dla $x\ne0$ bezpośrednio otrzymujemy
\[
 u'(x)=-\frac{\operatorname{sgn}(x)|x|^{1+\alpha}}{1+\alpha},
 \qquad u''(x)=-|x|^\alpha.
\]
Iloraz $u(h)/h$ dąży do zera, więc $u'(0)=0$.
Następnie $u'(h)/h=-|h|^\alpha/(1+\alpha)\to0$,
co daje $u''(0)=0$ i ciągłość drugiej pochodnej.
Zatem $-\Delta u=|x|^\alpha$ (w wymiarze jeden)
także w sensie klasycznym w zerze.

Dla $s,t\ge0$ zachodzi $(s+t)^\alpha\le s^\alpha+t^\alpha$:
jeśli $s>0$, wynika to z monotoniczności funkcji
$z\mapsto (1+z)^\alpha-z^\alpha$ dla $z>0$,
której pochodna jest ujemna, a granica przy $z\downarrow0$
wynosi $1$; przypadek $s=0$ daje równość.
Stąd dla $x,y\in[-1,1]$,
\[
 \bigl||x|^\alpha-|y|^\alpha\bigr|
 \le\bigl||x|-|y|\bigr|^\alpha
 \le |x-y|^\alpha.
\]
Zatem $u\in C^{2,\alpha}([-1,1])$.
Dla $\alpha<\beta\le1$ mamy natomiast
$|u''(h)-u''(0)|/h^\beta=h^{\alpha-\beta}\to+\infty$
przy $h\downarrow0$, więc druga pochodna nie jest
Hölderowska z wykładnikiem $\beta$ nawet lokalnie
w otoczeniu zera. W wymiarze $d\ge2$ funkcja
$U(x)=u(x_1)$ na kostce $(-1,1)^d$ daje ten sam
kontrprzykład dla laplasjanu: jedyną niezerową
składową Hessianu jest $u''(x_1)$, a ilorazy
wzdłuż pierwszej osi mają powyższą postać.
\end{kontrprzyklad}

\section{Podsumowanie i miejsce w dalszej teorii}
Twierdzenie~\ref{thm:pde16-cz} kontroluje drugie
pochodne potencjału Newtona przez normę prawej
strony w $L^p$ dla $1<p<\infty$; nie zapewnia
całkowalności samego potencjału na $\R^d$.
Silne oszacowania dla $p=1$ i $p=\infty$ zawodzą
w ogólności (uwaga~\ref{rem:pde16-endpoints}),
a oszacowanie~\eqref{eq:pde16-weak11} dotyczy słabego
typu $(1,1)$ na $L^1\cap L^2$.
Decydujące są znoszenie się jądra o średniej zero,
rozkład Calderóna--Zygmunda, interpolacja i dualność.
Twierdzenie~\ref{thm:pde16-variable} dla ciągłych
współczynników jest oszacowaniem a priori
w klasie $W^{2,p}_{\mathrm{loc}}$, nie samodzielnym
twierdzeniem o istnieniu lub uzyskaniu tej regularności.

W teorii Schaudera odejmowanie $f(x)$ pod całką
osobliwą daje zbieżność dzięki Hölderowskiemu
oszacowaniu różnicy. Twierdzenia
~\ref{thm:pde16-schauder-poisson}
i~\ref{thm:pde16-schauder-variable} dają wewnętrzną
regularność $C^{2,\alpha}$ dla $0<\alpha<1$
rozwiązań klasy $C^2$ przy odpowiednich hipotezach.
Nie obejmują końca $\alpha=1$ ani dowolnych
ograniczonych prawych stron. Nie dają też
oszacowania globalnego aż do brzegu.
W kolejnym rozdziale~\ref{chap:pde17} rozpatrzymy
regularność De Giorgiego--Nasha--Mosera dla słabych
rozwiązań skalarnych jednorodnych równań eliptycznych
w postaci dywergencyjnej, przy ograniczonych mierzalnych
i jednostajnie eliptycznych współczynnikach.
Regularność Hölderowska dotyczy tam samego rozwiązania,
nie automatycznie jego drugich pochodnych.

\chapter[Regularność De Giorgiego--Nasha--Mosera]{Teoria De Giorgiego--Nasha--Mosera: regularność przy mierzalnych współczynnikach}
\label{chap:pde17}

W rozdziałach~\ref{chap:pde15}--\ref{chap:pde16} uzyskiwaliśmy regularność
pochodnych przy odpowiedniej gładkości współczynników. Teraz przyjmujemy
\emph{jedynie mierzalność} macierzy $A$. Mimo to rozwiązanie słabe
jednostajnie eliptycznego równania skalarnego ma lokalnie reprezentanta
Hölderowskiego. Jest to szczególny przypadek twierdzenia
De Giorgiego--Nasha--Mosera. Przedstawimy metodę iteracji Mosera,
od oszacowań energetycznych aż do nierówności Harnacka i regularności.
Nie będziemy utożsamiać tej teorii z wyższą regularnością $W^{2,p}$:
przy skokowych współczynnikach pochodna rozwiązania może mieć skok.

\begin{uwaga}[Mapa założeń: zasada maksimum, Schauder i Harnack]
\label{rem:pde-v43-regularity-map}
Klasyczne zasady maksimum z rozdziału~\ref{chap:pde14}
wykorzystują regularność klasyczną albo dodatkowe argumenty
aproksymacyjne. Oszacowania Schaudera z
rozdziału~\ref{chap:pde16} wymagają regularności H\"olderowskiej
współczynników i dotyczą pochodnych drugiego rzędu.
W niniejszym rozdziale pracujemy z równaniem \emph{skalarnym}
w postaci dywergencyjnej, z mierzalnym, jednostajnie
eliptycznym $A$; regularność H\"olderowską otrzymujemy
wyłącznie dla samego rozwiązania. Nie przenosimy automatycznie
nierówności Harnacka na układy ani na równania z niekontrolowanymi
wyrazami niższego rzędu.
\end{uwaga}

Przyjmujemy oznaczenia $u_+=\max\{u,0\}$,
$\operatorname*{ess\,osc}_E u=
\operatorname*{ess\,sup}_E u-\operatorname*{ess\,inf}_E u$ oraz
$\fint_E h\,\mathrm d\lambda_d=|E|^{-1}
\int_E h\,\mathrm d\lambda_d$ dla zbiorów mierzalnych o dodatniej skończonej mierze;
$|E|=\lambda_d(E)$, a $h_E$ oznacza średnią funkcji
całkowalnej na $E$. Porównania wartości funkcji Sobolewa rozumiemy prawie
wszędzie, natomiast słabe nierówności różniczkowe
definiujemy za pomocą funkcji testowych.

\section{Założenia, słabe nierówności i skalowanie}
\label{sec:pde17-assumptions}
\begin{definicja}[Jednostajna eliptyczność i rozwiązania słabe]
Niech $d\ge2$, $\Omega\subset\R^d$ będzie obszarem, a
$A\colon\Omega\to\R^{d\times d}$ mierzalną macierzą, taką że dla
prawie każdego $x\in\Omega$, $\xi,\zeta\in\R^d$ zachodzą
\begin{equation}\label{eq:pde17-ellipticity}
 \xi^{\mathsf T}A(x)\xi\ge\lambda|\xi|^2,
 \qquad |A(x)\xi\cdot\zeta|\le\Lambda|\xi||\zeta|,
 \qquad 0<\lambda\le\Lambda<\infty.
\end{equation}
Nie wymagamy symetrii $A$. Rzeczywistą funkcję
$u\in H^1_{\rm loc}(\Omega)$
nazywamy słabym rozwiązaniem równania
$-\operatorname{div}(A\nabla u)=0$, jeżeli
\begin{equation}\label{eq:pde17-weak}
 \int_\Omega A\nabla u\cdot\nabla\varphi\,\mathrm d\lambda_d=0
 \quad\text{dla każdej }\varphi\in C_c^\infty(\Omega).
\end{equation}
Słabym podrozwiązaniem (odpowiednio nadrozwiązaniem)
nazywamy rzeczywistą funkcję $u\in H^1_{\rm loc}(\Omega)$,
dla której lewa strona~\eqref{eq:pde17-weak} jest $\le0$
(odpowiednio $\ge0$) dla każdej nieujemnej funkcji testowej.
\end{definicja}

\begin{uwaga}[Dopuszczalne testy]
Niech $\varphi\in H^1(\Omega)$ ma zwarty nośnik $K\Subset\Omega$.
Wybieramy zbiór otwarty $V$ taki, że $K\subset V\Subset\Omega$.
Przedłużenie $\varphi$ przez zero poza $\Omega$ należy do
$H^1(\R^d)$: można najpierw pomnożyć $\varphi$ przez gładką
funkcję odcięcia równą $1$ w otoczeniu $K$ i o nośniku w $V$,
a następnie przedłużyć iloczyn przez zero.
Sploty tego przedłużenia z nieujemnym gładkim jądrem
aproksymującym jedynkę, o dostatecznie małym nośniku, należą
do $C_c^\infty(V)$ i zbiegają do $\varphi$ w $H^1(V)$.
Jeśli $\varphi\ge0$ prawie wszędzie, także sploty są nieujemne.
Ponieważ
\[
 \left|\int_V A\nabla u\cdot\nabla(\varphi_j-\varphi)\,
 \mathrm d\lambda_d\right|
 \le\Lambda\|\nabla u\|_{L^2(V)}
       \|\nabla(\varphi_j-\varphi)\|_{L^2(V)}\longrightarrow0,
\]
równość oraz odpowiednie słabe nierówności przechodzą na takie
testy, z zachowaniem wymaganego znaku.
Dopuszczalność obcięć i potęg sprawdzamy za każdym razem za
pomocą reguły łańcuchowej dla funkcji Sobolewa i zwartego
nośnika funkcji odcięcia. Nie zakładamy gładkości $u$
ani nie różniczkujemy współczynników $A$.
\end{uwaga}

Jeżeli $u$ rozwiązuje równanie w $B_{2R}(x_0)$, funkcja
$v(y)=u(x_0+Ry)$ rozwiązuje je w $B_2(0)$ z macierzą
$A_R(y)=A(x_0+Ry)$ i \emph{tymi samymi} stałymi $\lambda,\Lambda$.
Istotnie, dla $\psi\in C_c^\infty(B_2)$ testujemy równanie
funkcją $\varphi(x)=\psi((x-x_0)/R)$. Z reguły łańcuchowej
dla słabych pochodnych i podstawienia $x=x_0+Ry$ otrzymujemy
\[
 \int_{B_2} A_R\nabla v\cdot\nabla\psi\,\mathrm d\lambda_d
 =R^{2-d}\int_{B_{2R}(x_0)}
 A\nabla u\cdot\nabla\varphi\,\mathrm d\lambda_d=0.
\]
Dodatni czynnik $R^{2-d}$ zachowuje także znak słabej
nierówności. Dzieląc $A_R$ przez $\lambda$, normalizujemy
dolną stałą eliptyczności do $1$, a górną do $\Lambda/\lambda$.
Stałe w znormalizowanych lokalnych oszacowaniach zależą
więc od $d$ i $\Lambda/\lambda$, a odpowiednie potęgi
promienia wynikają z powrotu do pierwotnej skali.

\section{Nierówność Caccioppolego dla obcięć De Giorgiego}
\label{sec:pde17-caccioppoli}
\begin{lemat}[Energia części dodatniej]
\label{lem:pde17-caccioppoli}
Niech $u\in H^1(B_R)$ będzie podrozwiązaniem w $B_R\subset\Omega$.
Dla $k\in\R$ i $\eta\in C_c^1(B_R)$, $0\le\eta\le1$, zachodzi
\begin{equation}\label{eq:pde17-cacc}
 \int_{B_R}\eta^2|\nabla(u-k)_+|^2\,\mathrm d\lambda_d
 \le \frac{4\Lambda^2}{\lambda^2}
 \int_{B_R}(u-k)_+^2|\nabla\eta|^2\,\mathrm d\lambda_d.
\end{equation}
Analogicznie, dla nadrozwiązania oszacowanie dotyczy $(k-u)_+$.
\end{lemat}
\begin{proof}
Oznaczmy $w=(u-k)_+$. Funkcja $w$ należy do $H^1$;
$\nabla w=\chi_{\{u>k\}}\nabla u$ prawie wszędzie,
co wynika z wniosku~\ref{cor:pde9-truncation}.
Testujemy nierówność podrozwiązania funkcją $\varphi=\eta^2w\ge0$.
Ponieważ $\nabla\varphi=\eta^2\nabla w+2\eta w\nabla\eta$
prawie wszędzie, a oba składniki znikają na $\{u\le k\}$,
dostajemy
\[
 \lambda\int_{B_R}\eta^2|\nabla w|^2\,\mathrm d\lambda_d
 \le 2\Lambda\int_{B_R}\eta w|\nabla w||\nabla\eta|\,\mathrm d\lambda_d.
\]
Nierówność Younga
$2\Lambda ab\le(\lambda/2)a^2+(2\Lambda^2/\lambda)b^2$
daje tezę. Jeżeli $u$ jest nadrozwiązaniem, to $-u$
jest podrozwiązaniem, gdyż forma jest liniowa względem
pierwszego argumentu. Stosujemy już udowodnione oszacowanie
do $-u$ i poziomu $-k$, zauważając, że
$(-u-(-k))_+=(k-u)_+$.
\end{proof}

\begin{uwaga}[Kontrast z metodą ilorazów różnicowych]
W poprzednich rozdziałach do regularności $H^2$ wykorzystywaliśmy
pochodne $A$. W~\eqref{eq:pde17-cacc} wystarcza informacja o
$A\in L^\infty$ i dodatniości części symetrycznej macierzy.
To właśnie pozwala prowadzić argument dla współczynników mierzalnych.
\end{uwaga}

\section{Lemat iteracyjny De Giorgiego i lokalne ograniczenie}
\label{sec:pde17-degiorgi}
\begin{lemat}[Zbieżność iteracji superliniowej]
\label{lem:pde17-iteration}
Niech $Y_n\ge0$, $a>1$, $b>1$, $\delta>0$ oraz
$Y_{n+1}\le a b^n Y_n^{1+\delta}$ dla $n\ge0$.
Jeżeli $Y_0\le a^{-1/\delta}b^{-1/\delta^2}$,
to $Y_n\to0$.
\end{lemat}
\begin{proof}
Niech $Z_n=a^{1/\delta}b^{n/\delta+1/\delta^2}Y_n$.
Podstawienie do rekurencji daje
\[
 Z_{n+1}\le
 b^{(n+1)/\delta+1/\delta^2+n
 -(1+\delta)(n/\delta+1/\delta^2)}
 Z_n^{1+\delta}=Z_n^{1+\delta},
\]
gdyż potęgi $a$ skracają się, a wykładnik $b$ jest zerowy.
Z założenia $0\le Z_0\le1$, więc przez indukcję
$0\le Z_n\le1$ dla wszystkich $n$, a zatem
$Y_n\le a^{-1/\delta}b^{-n/\delta-1/\delta^2}\to0$.
\end{proof}

\begin{twierdzenie}[Oszacowanie $L^2\to L^\infty$ podrozwiązania]
\label{thm:pde17-l2linfty}
Niech $B_{2R}(x_0)\Subset\Omega$ i $u\in H^1(B_{2R}(x_0))$
będzie podrozwiązaniem równania z~\eqref{eq:pde17-ellipticity}.
Wtedy
\begin{equation}\label{eq:pde17-l2linfty}
 \operatorname*{ess\,sup}_{B_R(x_0)}u_+
 \le C\left(\fint_{B_{2R}(x_0)}u_+^2\,\mathrm d\lambda_d\right)^{1/2},
 \qquad C=C(d,\Lambda/\lambda).
\end{equation}
\end{twierdzenie}
\begin{proof}
Skalowanie pozwala przyjąć $R=1$, $x_0=0$. Ustalmy
$K>0$ i ciągi
$r_n=1+2^{-n}$, $k_n=K(1-2^{-n-1})$, $n\ge0$.
Niech $w_n=(u-k_n)_+$ oraz
$Y_n=\int_{B_{r_n}}w_n^2\,\mathrm d\lambda_d$.
Wybieramy $\eta_n=1$ na $B_{r_{n+1}}$, o nośniku w $B_{r_n}$,
z $|\nabla\eta_n|\le C2^n$.
Przedłużamy $\eta_n w_n$ przez zero do $B_2$.
Z nierówności Sobolewa (twierdzenie~\ref{thm:pde10-Sobp})
dla $d>2$, a dla $d=2$ z tej samej nierówności dla $p=4/3$
i nierówności Höldera na $B_2$, otrzymujemy dla odpowiedniego
$\chi>1$
\[
 \|\eta_n w_n\|_{L^{2\chi}(B_2)}^2
 \le C_d\|\nabla(\eta_n w_n)\|_{L^2(B_2)}^2.
\]
Stosujemy te oszacowania do przedłużeń przez zero, korzystając
z gęstości $C_c^\infty(B_2)$ w $H_0^1(B_2)$.
Z reguły Leibniza i lematu~\ref{lem:pde17-caccioppoli}
prawa strona jest nie większa od
$C\int_{B_{r_n}}w_n^2|\nabla\eta_n|^2\,\mathrm d\lambda_d$,
a więc
\begin{equation}\label{eq:pde17-sobstep}
 \left(\int_{B_{r_n}}|\eta_nw_n|^{2\chi}\,\mathrm d\lambda_d
 \right)^{1/\chi}\le C4^nY_n.
\end{equation}
Dla $d>2$ bierzemy $\chi=d/(d-2)$, a dla $d=2$
możemy wybrać $\chi=2$ (zanurzenie $H_0^1(B_2)\hookrightarrow L^4$).
Na zbiorze $E_{n+1}=B_{r_{n+1}}\cap\{u>k_{n+1}\}$ mamy
$w_n\ge k_{n+1}-k_n=K2^{-n-2}$, zatem
$|E_{n+1}|\le C4^nK^{-2}Y_n$. Nierówność Höldera i
\eqref{eq:pde17-sobstep} dają
\[
 Y_{n+1}\le
 \left(\int_{B_{r_n}}|\eta_nw_n|^{2\chi}\,\mathrm d\lambda_d\right)^{1/\chi}
 |E_{n+1}|^{1-1/\chi}
 \le Cb^nK^{-2\delta}Y_n^{1+\delta},
 \quad\delta=1-1/\chi>0,
\]
przy pewnym $b=b(d)>1$. Po podzieleniu przez $K^2$
otrzymujemy rekurencję dla $Z_n=Y_n/K^2$ postaci
z lematu~\ref{lem:pde17-iteration}. Jeżeli $\|u_+\|_{L^2(B_2)}=0$, teza jest natychmiastowa.
W przeciwnym razie bierzemy
$K=C_0\|u_+\|_{L^2(B_2)}>0$.
Ponieważ $Y_0\le\|u_+\|_{L^2(B_2)}^2$, wybór dostatecznie
dużego $C_0$ zapewnia próg małości dla $Z_0$ z lematu
~\ref{lem:pde17-iteration}. Otrzymujemy $Y_n\to0$.
Na $B_1\subset B_{r_n}$ mamy $(u-K)_+\le w_n$, więc
$\int_{B_1}(u-K)_+^2\,\mathrm d\lambda_d\le Y_n\to0$,
czyli $u\le K$ prawie wszędzie na $B_1$.
Dla $v(y)=u(x_0+Ry)$ zachodzi
\[
 \|v_+\|_{L^2(B_2)}
 =|B_2|^{1/2}
 \left(\fint_{B_{2R}(x_0)}u_+^2\,\mathrm d\lambda_d\right)^{1/2}.
\]
Wchłaniając $|B_2|^{1/2}$ do stałej zależnej od $d$,
otrzymujemy~\eqref{eq:pde17-l2linfty}.
\end{proof}

\begin{wniosek}[Lokalna ograniczoność rozwiązania]
Jeżeli $u$ jest rozwiązaniem słabym w $B_{2R}$, to
$u$ i $-u$ są podrozwiązaniami. Stąd
$\|u\|_{L^\infty(B_R)}\le C(\fint_{B_{2R}}|u|^2\,\mathrm d\lambda_d)^{1/2}$.
W szczególności każde rozwiązanie słabe jest lokalnie
istotnie ograniczone.
\end{wniosek}

\section{Energia potęg i mechanizm iteracji Mosera}
\label{sec:pde17-moser}
W dalszej części zakładamy $v\ge\eps>0$. Dla nieujemnego
rozwiązania $u$ otrzymujemy to, biorąc $v=u+\eps$,
a po udowodnieniu oszacowań przechodzimy do granicy
$\eps\downarrow0$. Takie przesunięcie jest rozwiązaniem,
ponieważ równanie nie zawiera wyrazu zerowego rzędu.
Ponadto z lokalnej ograniczoności rozwiązania otrzymujemy
$v\in L^\infty_{\rm loc}$. Na zwartym nośniku każdego testu
wartości $v$ należą prawie wszędzie do pewnego przedziału
$[\eps,M]$. Funkcję potęgową można przedłużyć z tego przedziału
do funkcji klasy $C^1$ o ograniczonej pochodnej na $\R$
(np.\ przez liniową kontynuację z obu końców przedziału).
Stosujemy wówczas twierdzenie~\ref{thm:pde9-chain}
na otoczeniu nośnika testu, na którym $v$ jest ograniczona. Stałe w uzyskanych oszacowaniach
nie będą zależeć od $\eps$ ani od $M$.

\begin{lemat}[Oszacowanie dla potęg dodatnich i ujemnych]
\label{lem:pde17-powers}
Niech $v\in H^1(B_R)\cap L^\infty(B_R)$, $v\ge\eps>0$,
będzie słabym rozwiązaniem, $\beta\in(-\infty,0)\cup(1,\infty)$,
a $\eta\in C_c^1(B_R)$, $0\le\eta\le1$.
Wówczas
\begin{equation}\label{eq:pde17-powerenergy}
 \int_{B_R}|\nabla(\eta v^{\beta/2})|^2\,\mathrm d\lambda_d
 \le C\left(1+\frac{\beta^2}{(\beta-1)^2}\right)
 \int_{B_R}v^\beta|\nabla\eta|^2\,\mathrm d\lambda_d,
\end{equation}
gdzie $C$ zależy tylko od $\Lambda/\lambda$.
W szczególności stała jest ograniczona jednostajnie dla
$\beta\ge2$ lub $\beta<0$.
\end{lemat}
\begin{proof}
Testujemy równanie funkcją $\eta^2v^{\beta-1}$.
Jest to dopuszczalne, ponieważ $\eps\le v\le M$ na nośniku
$\eta$; mapy potęgowe mają tam ograniczone pochodne.
Reguła łańcuchowa i reguła Leibniza dają
\[
 0=(\beta-1)\int_{B_R}\eta^2v^{\beta-2}
 A\nabla v\cdot\nabla v\,\mathrm d\lambda_d
 +2\int_{B_R}\eta v^{\beta-1}
 A\nabla v\cdot\nabla\eta\,\mathrm d\lambda_d.
\]
Pierwsza całka bez czynnika $\beta-1$ jest nieujemna.
Biorąc wartości bezwzględne po przeniesieniu drugiego
składnika na prawą stronę, otrzymujemy niezależnie od znaku
$\beta-1$
\[
 |\beta-1|\lambda\int_{B_R}\eta^2v^{\beta-2}|\nabla v|^2
 \,\mathrm d\lambda_d
 \le2\Lambda\int_{B_R}\eta v^{\beta-1}|\nabla v||\nabla\eta|
 \,\mathrm d\lambda_d.
\]
Stosujemy nierówność Younga do iloczynu
$a=\eta v^{(\beta-2)/2}|\nabla v|$ oraz
$b=v^{\beta/2}|\nabla\eta|$:
\[
 2\Lambda ab\le\frac{\lambda|\beta-1|}{2}a^2+
 \frac{2\Lambda^2}{\lambda|\beta-1|}b^2.
\]
Po absorpcji pierwszego składnika otrzymujemy
\[
 \int_{B_R}a^2\,\mathrm d\lambda_d
 \le\frac{4(\Lambda/\lambda)^2}{(\beta-1)^2}
 \int_{B_R}b^2\,\mathrm d\lambda_d.
\]
Ponieważ $|\nabla(v^{\beta/2})|^2
=(\beta^2/4)v^{\beta-2}|\nabla v|^2$,
zaś $|\nabla(\eta v^{\beta/2})|^2$ jest co najwyżej
$2\eta^2|\nabla(v^{\beta/2})|^2+2v^\beta|\nabla\eta|^2$,
otrzymujemy~\eqref{eq:pde17-powerenergy}.
\end{proof}

\begin{lemat}[Iteracja $L^q\to L^\infty$ i potęg odwrotnych]
\label{lem:pde17-moser-iteration}
Niech $B_{2R}\Subset\Omega$, $v$ będzie rozwiązaniem słabym
w $\Omega$ oraz $v\ge\eps>0$ na $B_{2R}$, a $0<r<s\le2R$.
Dla każdego $q>0$ istnieje $C=C(d,\Lambda/\lambda,q)$ takie, że
\begin{align}
 \operatorname*{ess\,sup}_{B_r}v
 &\le C(s-r)^{-d/q}\|v\|_{L^q(B_s)},\label{eq:pde17-moser-up}\\
 \operatorname*{ess\,sup}_{B_r}v^{-1}
 &\le C(s-r)^{-d/q}\|v^{-1}\|_{L^q(B_s)}.
 \label{eq:pde17-moser-inv}
\end{align}
Dla $0<q<1$ zapis $\|w\|_{L^q(E)}$ oznacza wielkość
$(\int_E|w|^q\,\mathrm d\lambda_d)^{1/q}$, nie normę.
Lokalna ograniczoność $v$ wynika z twierdzenia
~\ref{thm:pde17-l2linfty}; dla $u\ge0$ stosujemy
nierówności do $v=u+\eps$.
\end{lemat}
\begin{proof}
Najpierw dowodzimy oszacowania na parach kul $B_1\subset B_2$
dla $q\ge2$. Przy $d>2$ bierzemy $\chi=d/(d-2)$, a przy
$d=2$ wybieramy $\chi=2$. Na ustalonej kuli $B_2$ stosujemy to samo oszacowanie
Sobolewa dla $H_0^1(B_2)$ co w dowodzie twierdzenia
~\ref{thm:pde17-l2linfty}:
$\|h\|_{L^{2\chi}(B_2)}^2\le C_d\|\nabla h\|_{L^2(B_2)}^2$.
Dla $1\le\rho<\sigma\le2$ wybieramy odcięcie równe $1$
na $B_\rho$, z nośnikiem w $B_\sigma$ i pochodną ograniczoną
przez $C/(\sigma-\rho)$. Z~\eqref{eq:pde17-powerenergy},
dla $|\beta|\ge2$, wynika
\[
 \left(\int_{B_\rho}v^{\beta\chi}\,\mathrm d\lambda_d\right)^{1/\chi}
 \le\frac{C}{(\sigma-\rho)^2}
 \int_{B_\sigma}v^\beta\,\mathrm d\lambda_d.
\]
Niech $w=v$ dla $\beta>0$ i $w=v^{-1}$ dla $\beta<0$.
Po podniesieniu do dodatniej potęgi $1/|\beta|$ otrzymujemy
\[
 \|w\|_{L^{|\beta|\chi}(B_\rho)}
 \le\left(\frac{C}{(\sigma-\rho)^2}\right)^{1/|\beta|}
       \|w\|_{L^{|\beta|}(B_\sigma)}.
\]
Bierzemy $p_k=q\chi^k$, $\rho_k=1+2^{-k}$.
Kolejne zastosowania z $|\beta|=p_k$ dają
\[
 \|w\|_{L^{p_n}(B_{\rho_n})}
 \le \prod_{k=0}^{n-1}(C4^{k+1})^{1/p_k}
        \|w\|_{L^q(B_2)}.
\]
Iloczyn jest ograniczony niezależnie od $n$, ponieważ
$\sum_{k\ge0}(1+k)/\chi^k<\infty$.
Na $B_1$ norma $L^{p_n}$ jest nie większa od normy na
$B_{\rho_n}$. Dla każdego $a>0$, dla którego zbiór
$E_a=B_1\cap\{w>a\}$ ma dodatnią miarę, mamy
$a|E_a|^{1/p_n}\le\|w\|_{L^{p_n}(B_1)}
\le C\|w\|_{L^q(B_2)}$.
Ponieważ $|E_a|^{1/p_n}\to1$, dostajemy
$a\le C\|w\|_{L^q(B_2)}$. Biorąc supremum po takich
$a$ (przypadek $w=0$ prawie wszędzie jest natychmiastowy),
otrzymujemy $\|w\|_{L^\infty(B_1)}\le C\|w\|_{L^q(B_2)}$.

Przeskalowanie pary $B_t(y)\subset B_{2t}(y)$ do powyższych
kul daje
\[
 \|w\|_{L^\infty(B_t(y))}
 \le C t^{-d/q}\|w\|_{L^q(B_{2t}(y))}.
\]
W szczególności wykładnik $d/q$ jest poprawny także dla $d=2$:
skalujemy ustaloną parę kul, a nie stosujemy jednorodnej
nierówności Sobolewa z wykładnikiem $4$ na dowolnej skali.
Dla $0<r<s$ kładziemy $t=(s-r)/2$. Dla każdego $y\in B_r$
mamy $B_{2t}(y)\subset B_s$. Przeliczalne pokrycie $B_r$
kulami $B_t(y)$ (np.\ o środkach z przeliczalnego zbioru
gęstego w $B_r$) pozwala połączyć oszacowania prawie wszędzie:
\[
 \|w\|_{L^\infty(B_r)}
 \le C(s-r)^{-d/q}\|w\|_{L^q(B_s)}.
\]

Dla $0<q<2$ z punktowej nierówności
$|w|^2\le\|w\|_{L^\infty(B_\sigma)}^{2-q}|w|^q$
po scałkowaniu wynika
$\|w\|_{L^2(B_\sigma)}\le
\|w\|_{L^\infty(B_\sigma)}^{1-q/2}
\|w\|_{L^q(B_\sigma)}^{q/2}$,
gdzie $w=v$ lub $w=v^{-1}$. Aby uzasadnić absorpcję,
ustalmy $r_j=r+(s-r)(1-2^{-j})$, $j\ge0$, i
$M_j=\|w\|_{L^\infty(B_{r_j})}$ oraz
$Q=\|w\|_{L^q(B_s)}$. Stosując oszacowanie $L^2\to
L^\infty$ pomiędzy $B_{r_j}$ a $B_{r_{j+1}}$
otrzymujemy
\[
 M_j\le C2^{jd/2}(s-r)^{-d/2}M_{j+1}^{1-q/2}Q^{q/2}.
\]
Dla dowolnego $0<\delta<1$ nierówność Younga z wykładnikami
$2/(2-q)$ i $2/q$ daje
\[
 M_j\le\delta M_{j+1}
 +C_{d,\Lambda/\lambda,q,\delta}2^{jd/q}(s-r)^{-d/q}Q.
\]
Wybieramy $\delta$ tak małe, aby $\delta2^{d/q}<1$,
i iterujemy nierówność $n$ razy:
\[
 M_0\le\delta^n M_n+
 C(s-r)^{-d/q}Q\sum_{j=0}^{n-1}(\delta2^{d/q})^j.
\]
Szereg geometryczny jest zbieżny. Pozostałość
$\delta^nM_n$ dąży do zera: dzięki lokalnej ograniczoności
$w$ jest ograniczone na $B_s\Subset\Omega$. Uzyskujemy
$M_0\le C(s-r)^{-d/q}Q$, czyli
\eqref{eq:pde17-moser-up}--\eqref{eq:pde17-moser-inv}.
Dowodzi to oszacowań także dla $0<q<2$.
\end{proof}


\begin{uwaga}[Jawny zakres stosowania potęg ujemnych]
\label{rem:pde17-negpower}
W oszacowaniu~\eqref{eq:pde17-moser-inv} nie wolno od razu
podstawić $u^{-1}$, jeśli nie wykazano dodatniej dolnej granicy
rozwiązania. Dla $u\ge0$ pracujemy z $v=u+\eps$;
wówczas $v^{-1}\in L^\infty$ na każdej kuli, a reguła
łańcuchowa jest dopuszczalna. Najpierw dowodzimy obu
oszacowań dla każdego $\eps>0$, przy tych samych stałych,
a dopiero na końcu przechodzimy do granicy monotonnej.
Nie otrzymujemy w ten sposób osobnego oszacowania
$\|u^{-1}\|_{L^q}$, które mogłoby być nieskończone.
To właśnie połączenie oszacowania dla $v^{-1}$
z mostem~\eqref{eq:pde17-bridge} usuwa tę trudność.
\end{uwaga}

\section{Logarytm rozwiązania i wykładnicza całkowalność}
\label{sec:pde17-log}
\begin{lemat}[Energia logarytmu]
\label{lem:pde17-log}
Niech $v\ge\eps>0$ będzie słabym dodatnim nadrozwiązaniem
w $B_{2r}(z)$. Wówczas
\begin{equation}\label{eq:pde17-log}
 \int_{B_r(z)}|\nabla\log v|^2\,\mathrm d\lambda_d
 \le C r^{d-2},\qquad C=C(d,\Lambda/\lambda).
\end{equation}
Dla każdej kuli $B_s(y)\subset B_r(z)$ otrzymujemy w konsekwencji
$\fint_{B_s(y)}|\log v-(\log v)_{B_s(y)}|\,\mathrm d\lambda_d
\le C(d,\Lambda/\lambda)$, o ile równanie obowiązuje
na kuli o podwójnym promieniu.
\end{lemat}
\begin{proof}
Niech $\eta\in C_c^1(B_{2r}(z))$, $0\le\eta\le1$.
Funkcje $t\mapsto t^{-1}$ i $t\mapsto\log t$ na
$[\eps,\infty)$ przedłużamy liniowo poniżej $\eps$,
otrzymując funkcje klasy $C^1$ o ograniczonych pochodnych.
Z twierdzenia~\ref{thm:pde9-chain}, stosowanego lokalnie,
wynika $\log v\in H^1_{\rm loc}$, z
$\nabla\log v=v^{-1}\nabla v$ prawie wszędzie, oraz
$\varphi=\eta^2v^{-1}$ jest dopuszczalnym nieujemnym
testem o zwartym nośniku. Z nierówności nadrozwiązania
uzyskujemy
\[
 \int_{B_{2r}(z)}\eta^2v^{-2}A\nabla v\cdot\nabla v\,\mathrm d\lambda_d
 \le 2\int_{B_{2r}(z)}\eta v^{-1}A\nabla v\cdot\nabla\eta\,\mathrm d\lambda_d.
\]
Eliptyczność, ograniczenie $A$ i nierówność Younga
z $a=\eta|\nabla\log v|$, $b=|\nabla\eta|$ dają
\[
 \int_{B_{2r}(z)}\eta^2|\nabla\log v|^2\,\mathrm d\lambda_d
 \le 4(\Lambda/\lambda)^2
 \int_{B_{2r}(z)}|\nabla\eta|^2\,\mathrm d\lambda_d.
\]
Wybieramy $\eta=1$ na $B_r(z)$ z nośnikiem w $B_{2r}(z)$,
$|\nabla\eta|\le C_d/r$. Całka po prawej stronie jest
nie większa niż $C_d r^{d-2}$.
Dla $B=B_s(y)$, $g=\log v$ i
$g_B=\fint_B g\,\mathrm d\lambda_d$ korzystamy z nierówności
Poincarégo (wniosek~\ref{cor:pde10-Poincare}, zastosowany
na kuli jednostkowej i przeskalowany) i nierówności Höldera:
\[
 \fint_B |g-g_B|\,\mathrm d\lambda_d
 \le C_d s\left(\fint_B|\nabla g|^2\,\mathrm d\lambda_d\right)^{1/2}
 \le C(d,\Lambda/\lambda).
\]
Ostatnia nierówność wynika z oszacowania energii na $B_s(y)$,
otrzymanego przy użyciu odcięcia na $B_{2s}(y)$.
\end{proof}

Potrzebujemy teraz jednego lematu analizy harmonicznej.
Nie traktujemy go jako automatycznej własności funkcji $H^1$:
przydatna jest \emph{jednostajna} kontrola średnich oscylacji
na wszystkich podkulach.

\begin{lemat}[Lemat Johna--Nirenberga; lokalna postać]
\label{lem:pde17-jn}
Niech $B=B_R(z)\subset\R^d$, a
$g\in L^1(B_{3R/2}(z))$. Zakładamy, że
$\sup_{B'\subset B_{3R/2}(z)}\fint_{B'}
|g-g_{B'}|\,\mathrm d\lambda_d\le M$,
gdzie supremum biegnie po kulach zawartych w większej kuli.
Jeśli $M>0$, istnieją $c_d,C_d>0$, dla których
\begin{equation}\label{eq:pde17-jn}
 \fint_B\exp\left(\frac{c_d}{M}|g-g_B|\right)
 \,\mathrm d\lambda_d\le C_d
\end{equation}
(gdy $M=0$, funkcja jest stała prawie wszędzie).
\end{lemat}
\begin{proof}[Dowód metodą kostek zatrzymania]
Gdy $M=0$, możemy zastosować założenie do samej kuli $B$;
wtedy $\int_B|g-g_B|\,\mathrm d\lambda_d=0$, więc $g=g_B$
prawie wszędzie na $B$. Dalej zakładamy $M>0$.
Najpierw ustalmy kostkę $Q$ i załóżmy, że średnia oscylacja
jest ograniczona przez $M$ na wszystkich jej podkostkach.
Rozkładamy $Q$ na kostki dyadyczne, pomijając ich brzegi,
których przeliczalna suma ma miarę zero.
Wybieramy $L=2^{d+2}M$. Ponieważ średnia oscylacja
na całej $Q$ jest nie większa niż $M<L$, każda kostka
przekraczająca próg zawiera się w maksymalnej takiej
kostce (wystarczy przejść po skończonym ciągu jej przodków).
Maksymalne kostki dyadyczne
$Q_j\subset Q$, dla których
$\fint_{Q_j}|g-g_Q|>L$, są rozłączne. Ponieważ
$\int_Q|g-g_Q|\le M|Q|$, łączna miara tych kostek
jest nie większa od $(M/L)|Q|\le 2^{-d-2}|Q|$.
Maksymalność i porównanie z rodzicem dają
$|g_{Q_j}-g_Q|\le2^dL$. Poza ich sumą średnia $|g-g_Q|$ na każdej kostce
dyadycznej zawierającej dany punkt nie przekracza $L$.
W punkcie Lebesgue'a funkcji $|g-g_Q|$ średnie te dążą
do jej wartości: kostka boku $b$ zawierająca ten punkt
mieści się w kuli o środku w tym punkcie i promieniu
$\sqrt d\,b$, a iloraz objętości kuli i kostki zależy
tylko od $d$. Z twierdzenia o różniczkowaniu całek
Lebesgue'a wynika więc $|g-g_Q|\le L$ prawie wszędzie
poza sumą wybranych kostek. W każdej złej kostce powtarzamy procedurę
z $g_{Q_j}$ zamiast $g_Q$. Oznaczmy przez $G_k$ sumę
złych kostek $k$-tego pokolenia, przyjmując $G_0=Q$.
W każdej kostce poprzedniego pokolenia suma miar
wybranych potomków stanowi co najwyżej część $M/L$
jej miary. Rozłączność kostek w danym pokoleniu i indukcja
dają $|G_k|\le (M/L)^k|Q|=2^{-k(d+2)}|Q|$.
Dla prawie każdego $x\in Q\setminus G_k$ łańcuch
wybranych kostek zawierających $x$ kończy się
po $j<k$ krokach. W ostatniej kostce $P$ zachodzi
$|g(x)-g_P|\le L$, a różnica średnich między każdą
parą kolejnych kostek w tym łańcuchu jest ograniczona
przez $2^dL$. Nierówność trójkąta daje więc
$|g(x)-g_Q|\le(1+k2^d)L$. Dla $k=0$ oszacowanie
na dopełnieniu jest puste.
Dla $t\ge L$ wybieramy największe $k\ge0$ takie, że
$(1+k2^d)L\le t$. Z oszacowania miary $k$-tego pokolenia,
uwzględniając $L=2^{d+2}M$, wynika
$|\{x\in Q:|g(x)-g_Q|>t\}|\le C|Q|e^{-ct/M}$.
Po zwiększeniu $C$ oszacowanie obowiązuje także dla $t<L$.
Dla $X=|g-g_Q|/M$ i $0<c'<c$ z tożsamości
$e^{c'X}=1+\int_0^X c'e^{c't}\,\mathrm dt$ oraz
twierdzenia Tonellego otrzymujemy
\[
 \fint_Q e^{c'X}\,\mathrm d\lambda_d
 =1+\frac{1}{|Q|}\int_0^\infty c'e^{c't}
      |\{x\in Q:X(x)>t\}|\,\mathrm dt
 \le 1+\frac{Cc'}{c-c'}.
\]
Przejście do kuli wymaga jeszcze uzasadnienia. Ustalmy
$B=B_R(z)$ i wybierzmy siatkę kostek o boku
$\ell=R/(100\sqrt d)$. Niech $h=\ell/2$. Wybieramy wszystkie otwarte kostki
o boku $\ell$, o środkach $y\in z+h\mathbb Z^d$ spełniających
$|y-z|\le R+\sqrt d\,h/2$. Dla każdego punktu
$\overline B$ zaokrąglenie współrzędnych do najbliższego
punktu siatki daje taki środek, z odległością w każdej
współrzędnej co najwyżej $h/2<\ell/2$. Kostki pokrywają
więc $\overline B$. Z każdego środka można dojść do $z$
krokami długości $h$ wzdłuż osi, zmniejszając kolejno
wartości bezwzględne współrzędnych względem $z$.
Odległość od $z$ nie rośnie, zatem wszystkie środki
pośrednie należą do pokrycia. Kostki odpowiadające
kolejnym krokom mają przecięcie objętości $\ell^d/2$.
Każde trzykrotne powiększenie kostki mieści się
w $B_{3R/2}(z)$, gdyż
$R+\sqrt d\,h/2+3\sqrt d\,\ell/2<3R/2$.
Liczba kostek $N$ jest ograniczona przez stałą zależną
tylko od $d$, ponieważ $R/h=200\sqrt d$.

Na każdej podkostce $Q'\subset Q_j$ średnia oscylacja
$g$ nie przekracza $C_dM$. Rzeczywiście, otaczamy $Q'$
kulą $B'$ o środku w środku $Q'$ i promieniu
$\sqrt d$ razy bok $Q'$ (można też wybrać promień
$\sqrt d/2$ razy bok $Q'$ dla kostki otwartej).
Jeśli $b$ jest bokiem $Q'$, to odległość środka
$Q'$ od środka $Q_j$ jest nie większa niż
$\sqrt d(\ell-b)/2$. Zatem $B'$ mieści się w kuli
o środku w środku $Q_j$ i promieniu
$\sqrt d(\ell+b)/2\le\sqrt d\,\ell$.
Ta kula wraz z $Q_j$ znajduje się w $B_{3R/2}(z)$,
na mocy ograniczenia odległości środków od $z$. Z nierówności
$\fint_{Q'}|g-g_{Q'}|\le
2\fint_{Q'}|g-g_{B'}|\le
2(|B'|/|Q'|)\fint_{B'}|g-g_{B'}|$
dostajemy żądane ograniczenie, ponieważ iloraz objętości
zależy tylko od $d$. Dyadyczny dowód dla kostek
zastosowany do każdego $Q_j$ zapewnia zatem
\[
 \fint_{Q_j}\exp\bigl(c_d|g-g_{Q_j}|/M\bigr)
 \,\mathrm d\lambda_d\le C_d.
\]
Jeżeli sąsiednie kostki $Q_i,Q_j$ mają część wspólną
o objętości co najmniej $c_d|Q_i|$, to
$|g_{Q_i}-g_{Q_j}|\le C_dM$: wystarczy całkować
$|g_{Q_i}-g_{Q_j}|\le
|g-g_{Q_i}|+|g-g_{Q_j}|$ po ich przecięciu,
a oba składniki oszacować przez oscylacje na kostkach.
Każdą parę kostek pokrycia łączy łańcuch co najwyżej
$N-1$ takich sąsiadów, stąd
$|g_{Q_j}-g_{Q_1}|\le C_dM$ dla wszystkich $j$.
Z nierówności trójkąta mamy dla $x\in Q_j$
$e^{c_d|g(x)-g_{Q_1}|/M}
\le e^{c_d C_d}e^{c_d|g(x)-g_{Q_j}|/M}$.
Ponieważ $B\subset\bigcup_{j=1}^N Q_j$, całkę po $B$
ograniczamy sumą całek po kostkach. Po podzieleniu
przez $|B|$ otrzymujemy
\[
 \fint_B e^{c_d|g-g_{Q_1}|/M}\,\mathrm d\lambda_d
 \le C_d\frac{N\ell^d}{|B|}\le C'_d,
\]
gdzie ostatnia stała zależy wyłącznie od $d$,
na mocy konstrukcji pokrycia.
Z nierówności Jensena lub bezpośrednio z oszacowania
pierwszego momentu wynika
$|g_B-g_{Q_1}|\le C_dM$. Zastępując $g_{Q_1}$
przez $g_B$, kończymy dowód~\eqref{eq:pde17-jn}.
\end{proof}

\begin{wniosek}[Most między potęgami dodatnimi i ujemnymi]
\label{cor:pde17-bridge}
Jeżeli $v\ge\eps>0$ jest nadrozwiązaniem w $B_{4R}(x_0)$,
to istnieje $q_0=q_0(d,\Lambda/\lambda)>0$, dla którego
\begin{equation}\label{eq:pde17-bridge}
 \left(\fint_{B_{2R}}v^{q_0}\,\mathrm d\lambda_d\right)
 \left(\fint_{B_{2R}}v^{-q_0}\,\mathrm d\lambda_d\right)
 \le C(d,\Lambda/\lambda).
\end{equation}
\end{wniosek}
\begin{proof}
Bierzemy $g=\log v$. Z lematu~\ref{lem:pde17-log}
i nierówności Poincarégo wynika, że średnia oscylacja $g$
na każdej podkuli $B_{2R}$ jest jednostajnie ograniczona.
Uściślijmy użycie kul. Dla dowolnej
$B_\rho(y)\subset B_{3R}(x_0)$ i $\rho\le R$
kula $B_{2\rho}(y)$ mieści się w $B_{4R}(x_0)$,
więc testowanie logarytmu na podwójnej kuli daje
jednostajną kontrolę oscylacji. Jeśli $R<\rho\le3R$, kula $B_\rho(y)$ mieści się w
$B_{3R}(x_0)$. W dowodzie lematu~\ref{lem:pde17-log} wybieramy teraz
odcięcie równe $1$ na $B_{3R}(x_0)$, o zwartym nośniku
w $B_{4R}(x_0)$ i $|\nabla\eta|\le C_d/R$. Daje to
$\int_{B_{3R}(x_0)}|\nabla g|^2\,\mathrm d\lambda_d(x)
\le C R^{d-2}$.
Nierówność Poincarégo na $B_\rho(y)$ i $\rho>R$
prowadzą zatem do
$\frac{1}{\lambda_d(B_\rho(y))}\int_{B_\rho(y)}|g(x)-g_{B_\rho(y)}|\,\mathrm d\lambda_d(x)
\le C(d,\Lambda/\lambda)(R/\rho)^{(d-2)/2}
\le C(d,\Lambda/\lambda)$.
W obu przypadkach
otrzymujemy ograniczenie średnich oscylacji zależne
tylko od $d$ i $\Lambda/\lambda$. Stosujemy zatem
zlokalizowaną wersję lematu Johna--Nirenberga z kulą
mniejszą $B_{2R}(x_0)$ i większą $B_{3R}(x_0)$.
Jako $M$ przyjmujemy dodatnią stałą nie mniejszą niż $1$
i nie mniejszą niż uzyskane ograniczenie średnich oscylacji.
W lemacie~\ref{lem:pde17-jn} dobieramy $q_0=c_d/M>0$.
Zachodzą tożsamości
$v^{\pm q_0}=\exp(\pm q_0(g-g_{B_{2R}}))
\exp(\pm q_0g_{B_{2R}})$.
Mnożąc całki i używając~\eqref{eq:pde17-jn}, usuwamy
przeciwne czynniki zależne od średniej $g$.
\end{proof}

\section{Nierówność Harnacka i regularność Höldera}
\label{sec:pde17-harnack}
\begin{twierdzenie}[Nierówność Harnacka dla słabych rozwiązań]
\label{thm:pde17-harnack}
Niech $d\ge2$, $A$ spełnia~\eqref{eq:pde17-ellipticity},
$B_{4R}(x_0)\subset\Omega$, zaś $u\in H^1(B_{4R})$
będzie nieujemnym słabym rozwiązaniem
$-\operatorname{div}(A\nabla u)=0$.
Wówczas istnieje stała $C=C(d,\Lambda/\lambda)$ taka, że
\begin{equation}\label{eq:pde17-harnack}
 \operatorname*{ess\,sup}_{B_R(x_0)}u
 \le C\operatorname*{ess\,inf}_{B_R(x_0)}u.
\end{equation}
\end{twierdzenie}
\begin{proof}
Pracujemy na obszarze $B_{4R}(x_0)$; wszystkie funkcje
testowe mają zwarty nośnik w tej kuli. W szczególności
$\overline{B_{2R}(x_0)}\subset B_{4R}(x_0)$, więc
$v=u+\eps$ jest ograniczona na $B_{2R}$ i spełnia
hipotezy lematu~\ref{lem:pde17-moser-iteration}.
Do $v$ stosujemy kolejno iterację
\eqref{eq:pde17-moser-up} dla potęgi dodatniej i
\eqref{eq:pde17-moser-inv} dla potęgi odwrotnej,
z wykładnikiem $q_0$ z~\eqref{eq:pde17-bridge}:
\[
 \operatorname*{ess\,sup}_{B_R}v
 \le C\left(\fint_{B_{2R}}v^{q_0}\,\mathrm d\lambda_d\right)^{1/q_0},
 \qquad
 \left(\fint_{B_{2R}}v^{-q_0}\,\mathrm d\lambda_d\right)^{-1/q_0}
 \le C\operatorname*{ess\,inf}_{B_R}v.
\]
Przy przejściu od norm do średnich wykorzystujemy
$R^{-d/q_0}|B_{2R}|^{1/q_0}=|B_2|^{1/q_0}$;
w drugim oszacowaniu odwracamy dodatnie wielkości
i stosujemy tożsamość
$\operatorname*{ess\,sup}v^{-1}=1/\operatorname*{ess\,inf}v$.
Oznaczmy $P=\fint_{B_{2R}}v^{q_0}\,\mathrm d\lambda_d$,
$N=\fint_{B_{2R}}v^{-q_0}\,\mathrm d\lambda_d$.
Obie wielkości są dodatnie i skończone. Z mostu
~\eqref{eq:pde17-bridge} mamy $PN\le C_0$, skąd
$P^{1/q_0}\le C_0^{1/q_0}N^{-1/q_0}$.
Łącząc trzy oszacowania, otrzymujemy
$\operatorname*{ess\,sup}_{B_R}v\le
C\operatorname*{ess\,inf}_{B_R}v$ ze stałą niezależną od
$\eps$. Ponieważ dodanie stałej $\eps$ przesuwa zarówno
kres istotny górny, jak i dolny dokładnie o $\eps$,
przejście $\eps\downarrow0$ kończy dowód.
\end{proof}

\begin{twierdzenie}[Twierdzenie De Giorgiego--Nasha--Mosera o regularności lokalnej]
\label{thm:pde17-holder}
Niech $d\ge2$, $\Omega\subset\R^d$ będzie obszarem,
macierz $A$ spełnia~\eqref{eq:pde17-ellipticity}, a
$u\in H^1_{\rm loc}(\Omega)$ niech będzie słabym rozwiązaniem
$-\operatorname{div}(A\nabla u)=0$.
Istnieje $\alpha=\alpha(d,\Lambda/\lambda)\in(0,1)$
taka, że $u$ ma reprezentanta należącego do $C^{0,\alpha}_{\rm loc}(\Omega)$.
Dla $B_{4R}(x_0)\Subset\Omega$ zachodzi oszacowanie
\begin{equation}\label{eq:pde17-holder}
 [u]_{C^{0,\alpha}(\overline{B_R(x_0)})}
 \le C R^{-\alpha}
 \left(\fint_{B_{4R}(x_0)}|u|^2\,\mathrm d\lambda_d\right)^{1/2},
\end{equation}
z $C$ zależną tylko od $d$ i $\Lambda/\lambda$.
\end{twierdzenie}
\begin{proof}
Najpierw uzasadnimy redukcję oscylacji. Oznaczmy
$m=\operatorname*{ess\,inf}_{B_{4r}}u$,
$M=\operatorname*{ess\,sup}_{B_{4r}}u$, dla kuli
$B_{4r}\Subset\Omega$. Z lokalnego ograniczenia
oba wyrazy są skończone. Funkcje $u-m$ i $M-u$ są
nieujemnymi rozwiązaniami w $B_{4r}$, a po zastosowaniu
nierówności Harnacka z obszarem $B_{4r}$ otrzymujemy,
oznaczając
$m_r=\operatorname*{ess\,inf}_{B_r}u$,
$M_r=\operatorname*{ess\,sup}_{B_r}u$,
\[
 M_r-m\le C(m_r-m),\qquad M-m_r\le C(M-M_r).
\]
Przyjmijmy od początku $C>1$. Jeżeli $a=m_r-m$,
$b=M-M_r$, $c=M_r-m_r$, to
$c\le(C-1)a$ i $c\le(C-1)b$. Dodając, dostajemy
$2c\le(C-1)(a+b)=(C-1)(M-m-c)$, czyli
$(C+1)c\le(C-1)(M-m)$. Stąd
\begin{equation}\label{eq:pde17-osc-decay}
 \operatorname*{ess\,osc}_{B_r}u
 \le\vartheta\operatorname*{ess\,osc}_{B_{4r}}u,
 \qquad \vartheta=\frac{C-1}{C+1}<1.
\end{equation}
Z wyboru stałej Harnacka mamy $0<\vartheta<1$. Kładziemy
$\alpha=\min\{1/2,-\log_4\vartheta\}>0$.
Dla $0<\rho\le R$ wybieramy $k\ge0$ takie, że
$4^{-k-1}R<\rho\le4^{-k}R$. Iteracja
~\eqref{eq:pde17-osc-decay} oraz monotoniczność oscylacji
względem zawierania kul dają
\[
 \operatorname*{ess\,osc}_{B_\rho(z)}u
 \le\vartheta^k\operatorname*{ess\,osc}_{B_R(z)}u
 \le4^\alpha(\rho/R)^\alpha
       \operatorname*{ess\,osc}_{B_R(z)}u
\]
dla każdej kuli $B_R(z)\Subset\Omega$.

Ustalmy teraz $B_{4R}(x_0)\Subset\Omega$ i oznaczmy
$H=(\fint_{B_{4R}(x_0)}|u|^2\,\mathrm d\lambda_d)^{1/2}$.
Stosując~\eqref{eq:pde17-l2linfty} do $u$ i $-u$,
otrzymujemy $\|u\|_{L^\infty(B_{2R}(x_0))}\le C H$.
Niech $D$ będzie zbiorem punktów Lebesgue'a funkcji $u$
w $B_R(x_0)$; jego dopełnienie ma miarę zero.
Dla $x,y\in D$, $t=|x-y|>0$, rozważamy dwa przypadki.
Gdy $t<R/2$, oba punkty należą do $B_{2t}(x)$, a
$B_R(x)\subset B_{2R}(x_0)$. Wartości w punktach Lebesgue'a
leżą między kresami istotnymi na każdej kuli zawierającej
te punkty we wnętrzu (wynika to z granicy średnich po
małych kulach). Zatem
\[
 |u(x)-u(y)|
 \le\operatorname*{ess\,osc}_{B_{2t}(x)}u
 \le C(t/R)^\alpha H.
\]
Gdy $t\ge R/2$, korzystamy z
$|u(x)-u(y)|\le2CH\le 2^{1+\alpha}C(t/R)^\alpha H$.
Przypadek $x=y$ jest natychmiastowy.

Zbiór $D$ jest gęsty w $\overline{B_R(x_0)}$, ponieważ
każda niepusta otwarta podkula $B_R(x_0)$ ma dodatnią
miarę. Dla dowolnego punktu domknięcia wybieramy ciąg
punktów z $D$ do niego zbieżny. Powyższa nierówność
zapewnia, że ciąg wartości jest ciągiem Cauchy'ego w
$\R$ oraz że jego granica nie zależy od wyboru ciągu.
Definiujemy reprezentanta przez te granice; przechodząc
do granicy w nierówności dla par punktów z $D$,
otrzymujemy~\eqref{eq:pde17-holder} na domkniętej kuli.
Reprezentant zgadza się z $u$ prawie wszędzie.
Na przecięciach otwartych kul reprezentanci pokrywają się,
gdyż są ciągłe i równe prawie wszędzie, a zbiór pełnej
miary w otwartym przecięciu jest gęsty. Kule $B_R(x_0)$ takie, że $B_{4R}(x_0)\Subset\Omega$,
pokrywają $\Omega$. Wybieramy z nich przeliczalne
podpokrycie, korzystając z drugiego aksjomatu przeliczalności
$\R^d$. Sklejona funkcja jest równa $u$ prawie wszędzie
na całej $\Omega$, ponieważ przeliczalna suma zbiorów
miary zero ma miarę zero. Jest ciągła i lokalnie Hölderowska
z tym samym wykładnikiem; jej jednoznaczność wynika ponownie
z gęstości zbioru pełnej miary w obszarze.
\end{proof}

\begin{wniosek}[Mocna zasada maksimum przy mierzalnych współczynnikach]
Jeżeli $u$ jest ciągłym reprezentantem słabego rozwiązania
w spójnym obszarze $\Omega$ i osiąga w punkcie $\Omega$
maksimum lub minimum globalne na $\Omega$, to jest stałe
w $\Omega$.
\end{wniosek}
\begin{proof}
Niech $M=\max_\Omega u$ i $S=\{x\in\Omega:u(x)=M\}$.
Zbiór $S$ jest niepusty i względnie domknięty z ciągłości.
Dla każdego $z\in S$ wybieramy $B_{4r}(z)\Subset\Omega$.
Funkcja $M-u$ jest nieujemnym rozwiązaniem w tej kuli
i przyjmuje w $z$ wartość zero. Z ciągłości i dodatniej
miary małych kul wynika, że jej kres istotny dolny
na $B_r(z)$ jest zerowy. Nierówność Harnacka daje zerowy
kres istotny górny, a ciągłość wymusza $u=M$ na całej
$B_r(z)$. Zatem $S$ jest także otwarty. Spójność daje
$S=\Omega$; minimum rozpatrujemy, stosując wynik do $-u$.
\end{proof}

\section[Współczynniki skokowe i nieciągłość pochodnych]{Współczynniki skokowe: dlaczego nie otrzymujemy $C^1$?}
\label{sec:pde17-counterexample}
\begin{kontrprzyklad}[Brak różniczkowalności pomimo jednostajnej eliptyczności]
Dla $x\in(-1,1)$ przyjmijmy $a(x)=1$ gdy $x<0$ i
$a(x)=2$ gdy $x\ge0$. Funkcja
$u(x)=x$ dla $x<0$ oraz $u(x)=x/2$ dla $x\ge0$ należy do
$H^1(-1,1)\cap C^{0,1}([-1,1])$, ale nie do $C^1(-1,1)$.
Słaba pochodna $u'$ jest równa $1$ na $(-1,0)$ i $1/2$
na $(0,1)$: całkując przez części na obu przedziałach,
nie otrzymujemy skoku wartości $u$ w zerze, ponieważ $u(0)=0$.
Jednostronne pochodne w zerze są różne.
Ponieważ $a(x)u'(x)=1$ prawie wszędzie, mamy
$-(au')'=0$ w sensie dystrybucji. W wymiarach $d\ge2$
na kostce $\Omega=(-1,1)^d$ bierzemy
$A(x)=\operatorname{diag}(a(x_1),1,\ldots,1)$ i
$U(x)=u(x_1)$. Macierz jest mierzalna i spełnia
~\eqref{eq:pde17-ellipticity} z $\lambda=1$, $\Lambda=2$.
Z twierdzenia Fubiniego wynika
$\nabla U=(u'(x_1),0,\ldots,0)$ prawie wszędzie.
Wobec $A\nabla U=(1,0,\ldots,0)$ prawie wszędzie
dla $\varphi\in C_c^\infty(\Omega)$ mamy
$\int_\Omega A\nabla U\cdot\nabla\varphi\,\mathrm d\lambda_d=0$,
ponownie z twierdzenia Fubiniego i jednowymiarowego
twierdzenia podstawowego rachunku całkowego.
\end{kontrprzyklad}

\begin{kontrprzyklad}[Ostrość wykładnika Höldera dla macierzy mierzalnych]
\label{ex:pde17-radial}
Niech $0<\gamma<1$ i $\Omega=B_1(0)\subset\R^2$. Dla $x\ne0$
niech $e_r=x/|x|$, $e_\theta=(-x_2,x_1)/|x|$ i
\[
 A(x)=e_r\otimes e_r+\gamma^2 e_\theta\otimes e_\theta,
 \qquad u(r,\theta)=r^\gamma\cos\theta.
\]
Definiujemy $A(0)=I$. Wektory $e_r,e_\theta$ tworzą
bazę ortonormalną $\R^2$, w której macierz $A$ ma wartości
własne $1$ i $\gamma^2$. Stąd $\gamma^2I\le A\le I$,
a macierz jest mierzalna (ciągła poza zerem).
Wykażemy, że $u\in H^1(B_1)$ i jest słabym rozwiązaniem
$-\operatorname{div}(A\nabla u)=0$ na całej kuli.
Rzeczywiście, dla $r>0$ rachunek biegunowy daje
\[
 \operatorname{div}(A\nabla u)=
 \frac1r\partial_r(r\partial_ru)
 +\frac{\gamma^2}{r^2}\partial_{\theta\theta}u
 =\gamma^2r^{\gamma-2}\cos\theta
 -\gamma^2r^{\gamma-2}\cos\theta=0.
\]
Definiujemy $u(0)=0$. Dla $r>0$ mamy
\[
 \nabla u=\partial_r u\,e_r+r^{-1}\partial_\theta u\,e_\theta
 =r^{\gamma-1}(\gamma\cos\theta\,e_r-\sin\theta\,e_\theta).
\]
Z ortonormalności wynika $|\nabla u|\le r^{\gamma-1}$,
zatem
\[
 \int_{B_1}|\nabla u|^2\,\mathrm d\lambda_2
 \le 2\pi C^2\int_0^1 r^{2\gamma-1}\,\mathrm dr<\infty.
\]
Ponadto $|u|\le1$. Aby utożsamić gradient klasyczny poza
zerem z gradientem słabym, całkujemy przez części na
$B_1\setminus\overline{B_\eps}$ z testem
$\varphi\in C_c^\infty(B_1)$. Wyraz brzegowy
dla pierwszej pochodnej jest ograniczony przez
$2\pi\eps^{1+\gamma}\|\varphi\|_\infty$, więc znika.
Całki objętościowe przechodzą do granicy z całkowalności
$u$ i gradientu. To dowodzi $u\in H^1(B_1)$.
W analogicznym rachunku dla strumienia $A\nabla u$
wyraz brzegowy ma moduł co najwyżej
$2\pi C\eps^\gamma\|\varphi\|_\infty$, a całka
objętościowa z dywergencją poza zerem jest zerowa.
Przejście do granicy daje~\eqref{eq:pde17-weak}.

Sprawdźmy dokładnie regularność Höldera. Dla
$x,y\in B_1$ przyjmijmy $r=|x|\ge|y|$ i $t=|x-y|>0$.
Jeśli $t\ge r/2$, to
$|u(x)-u(y)|\le2r^\gamma\le 2^{1+\gamma}t^\gamma$.
Jeśli $t<r/2$, każdy punkt odcinka $[x,y]$
ma normę co najmniej $r-t>r/2$. Z oszacowania
gradientu i całkowania wzdłuż odcinka otrzymujemy
$|u(x)-u(y)|\le C (r/2)^{\gamma-1}t\le C t^\gamma$.
Zatem $u\in C^{0,\gamma}(B_1)$.
Nie należy do $C^{0,\beta}$ w żadnym otoczeniu zera
dla $\beta>\gamma$, gdyż dla $x=(r,0)$ mamy
$|u(x)-u(0)|/|x|^\beta=r^{\gamma-\beta}\to\infty$.
W szczególności nie istnieje jeden dodatni wykładnik
Höldera obowiązujący dla wszystkich ilorazów
eliptyczności: w tym przykładzie $\Lambda/\lambda=\gamma^{-2}$.
Nie twierdzimy, że przykład wyznacza optymalny wykładnik
dla każdej klasy macierzy w dowolnym wymiarze.
\end{kontrprzyklad}

\section{Granice teorii i dalsze kierunki}
\label{sec:pde17-final}
Dowiedliśmy nierówności Harnacka oraz regularności Höldera
dla jednorodnych, \emph{skalarnych} równań dywergencyjnych
przy jedynie mierzalnych i jednostajnie eliptycznych
współczynnikach. Kluczowy mechanizm to oszacowanie energii,
zanurzenie Sobolewa, iteracja i kontrola $\log u$.
Dla $-\operatorname{div}(A\nabla u)=f$ prawa strona wnosi
nowe człony: typowym warunkiem regularności jest
$f\in L^q_{\rm loc}$ z $q>d/2$ (dla odpowiednio dobranej
klasy rozwiązań). W twierdzeniach dla układów równań
sytuacja może być zasadniczo gorsza. Rozdział kolejny
zajmie się teorią spektralną, ale osobno zachowamy miejsce
na pełne warianty niejednorodne, paraboliczne i kontrprzykłady.

% -----------------------------------------------------------------------
% Rozdział 18 — do włączenia PRZED \appendix w pliku wykładu
% -----------------------------------------------------------------------

\chapter[Widmo operatorów eliptycznych]{Widmo operatorów eliptycznych: zasada min--max, asymptotyka Weyla i oszacowania wartości własnych}
\label{chap:pde18}

\section{Cel i zakres rozdziału}
\label{sec:pde18-intro}
W rozdziale~\ref{chap:pde12} skonstruowaliśmy laplasjan Dirichleta
jako operator związany z domkniętą formą i uzyskaliśmy bazę
funkcji własnych. Teraz chcemy odpowiedzieć na pytania ilościowe:
jak duże są kolejne wartości własne, jak zależą od geometrii
obszaru i współczynników oraz co dokładnie oznacza asymptotyka
liczby modów o energii poniżej ustalonego progu? Podamy dowody
zasady min--max, oszacowania Li--Yau, prawa Weyla dla laplasjanu
Dirichleta i wzoru na wiodący człon śladu półgrupy cieplnej.

W całym rozdziale $d\ge1$, a $\Omega\subset\R^d$ oznacza
\emph{ograniczony obszar}. Zwartość zanurzenia
$H_0^1(\Omega)\hookrightarrow L^2(\Omega)$ wynika już z
ograniczoności obszaru i nie wymaga regularnego brzegu.
Istotnie, wybieramy otwartą kostkę $Q$ zawierającą
$\overline\Omega$. Przedłużenie przez zero funkcji z
$C_c^\infty(\Omega)$ zachowuje normę $H^1$, a przez
gęstość wyznacza izometryczne zanurzenie
$H_0^1(\Omega)\to H_0^1(Q)$. Z twierdzenia Rellicha--Kondrachova
~\ref{thm:pde10-Rellich-p} na kostce i ciągłości ograniczenia
$L^2(Q)\to L^2(\Omega)$ wynika żądana zwartość.
Ten argument działa także dla ograniczonych, niekoniecznie
spójnych zbiorów otwartych.

Formy symetryczne rozpatrujemy w rzeczywistych przestrzeniach
Hilberta. Widmo operatora rozumiemy jako widmo jego
zespolenia; dla rozpatrywanych operatorów samosprzężonych
rzeczywista baza funkcji własnych jest również bazą
po zespoleniu, z tymi samymi wartościami własnymi
i krotnościami. W rachunkach z transformatą Fouriera
korzystamy z jej zwykłej zespolonej postaci.
Silniejsze założenia geometryczne stawiamy jawnie tam,
gdzie potrzebujemy śladów lub twierdzeń o regularności.
Dla prawa Weyla wystarcza ograniczoność i warunek
$\lambda_d(\partial\Omega)=0$; w dowodzie zostanie
to wskazane osobno.

\begin{uwaga}[Od widma do ewolucji — zakres zastosowania]
\label{rem:pde-v43-spectral-heat}
Rozkład spektralny dotyczy operatora samosprzężonego,
niekoniecznie formalnego wyrażenia $-\operatorname{div}(A\nabla)+V$
na dowolnie przyjętej dziedzinie. W rozdziale~\ref{chap:pde20}
półgrupa $e^{-tL}$ i jej rozwinięcie w funkcje własne
są używane w przypadku \emph{autonomicznego}, dodatniego
operatora Dirichleta. Dla macierzy $A(t,x)$ zależnej od
czasu, zwłaszcza niesymetrycznej, nie wolno stosować
tego samego rozwinięcia do operatora chwilowego;
konstrukcja Galerkina wykorzysta wówczas bazę
\emph{pomocniczego} laplasjanu Dirichleta.
\end{uwaga}

\section{Od formy do operatora z dyskretnym widmem}
\label{sec:pde18-form}
\begin{definicja}[Symetryczna forma eliptyczna]
Niech $\Omega\subset\R^d$ będzie niepustym ograniczonym
zbiorem otwartym. Niech $A\in L^\infty(\Omega;\R^{d\times d})$
będzie symetryczna prawie wszędzie, a $V\in L^\infty(\Omega;\R)$.
Zakładamy, że istnieją $0<a_0\le a_1<\infty$ takie, iż
$ a_0|\xi|^2\le\xi^{\mathsf T}A(x)\xi\le a_1|\xi|^2$
dla prawie każdego $x\in\Omega$ i każdego $\xi\in\R^d$.
Na $H_0^1(\Omega)$ określamy formę
\begin{equation}\label{eq:pde18-form}
 q(u,v)=\int_\Omega \bigl(A\nabla u\cdot\nabla v+Vuv\bigr)
 \,\mathrm d\lambda_d(x),\qquad q[u]=q(u,u).
\end{equation}
Przez $L_{A,V}$ oznaczamy operator w $L^2(\Omega)$
związany z tą formą (jego samosprzężoność wykażemy poniżej): $u\in D(L_{A,V})$
wtedy i tylko wtedy, gdy $u\in H_0^1(\Omega)$ oraz istnieje
$f\in L^2(\Omega)$ takie, że $q(u,v)=(f,v)_{L^2}$
dla każdego $v\in H_0^1(\Omega)$; wówczas $L_{A,V}u=f$.
Funkcja $f$ jest jednoznacznie wyznaczona, ponieważ
$C_c^\infty(\Omega)\subset H_0^1(\Omega)$ jest gęsty w
$L^2(\Omega)$.
\end{definicja}

\begin{uwaga}[Istotność dziedziny operatora]
Symbol $-\operatorname{div}(A\nabla u)+Vu$ należy na początku
rozumieć dystrybucyjnie. Przy $A\in L^\infty$
nie można bez dowodu przyjąć $D(L_{A,V})=H^2(\Omega)\cap
H_0^1(\Omega)$. Definicja przez formę jest poprawna również dla
współczynników skokowych. W tym rozdziale \emph{obszar} to
otwarty zbiór spójny, a \emph{dziedzina operatora} to $D(L_{A,V})$.
\end{uwaga}

\begin{twierdzenie}[Samosprzężoność i zwarta rezolwenta]
\label{thm:pde18-compact}
Przy założeniach poprzedniej definicji operator $L_{A,V}$
jest samosprzężony i ograniczony z dołu. Jego widmo składa się
z rzeczywistych wartości własnych, liczonych z krotnościami,
$\lambda_1\le\lambda_2\le\cdots$, przy czym
$\lambda_k\to+\infty$. Istnieje ortonormalna baza
$L^2(\Omega)$ złożona z funkcji własnych $L_{A,V}$.
\end{twierdzenie}
\begin{proof}
Oznaczmy $M=\|V_-\|_{L^\infty}$, gdzie
$V_-=\max\{-V,0\}$. Przy $c=M+1$ forma
$q_c(u,v)=q(u,v)+c(u,v)_{L^2}$ spełnia
$q_c[u]\ge a_0\|\nabla u\|_2^2+\|u\|_2^2$.
Ponadto z symetrii $A$ i ograniczenia jej wartości własnych
wynika $|A\xi|\le a_1|\xi|$, więc
\[
 |q_c(u,v)|\le a_1\|\nabla u\|_2\|\nabla v\|_2
 +(\|V\|_\infty+c)\|u\|_2\|v\|_2.
\]
Forma jest zatem ciągła i koercywna na $H_0^1(\Omega)$.
Norma $q_c[u]^{1/2}$ jest równoważna normie $H_0^1$;
w szczególności $q_c$ jest iloczynem skalarnym
czyniącym tę przestrzeń zupełną.
Dla każdego $f\in L^2$ twierdzenie Laxa--Milgrama
wyznacza jednoznacznie $u=T f\in H_0^1$ takie, że
$q_c(u,v)=(f,v)_{L^2}$ dla wszystkich $v\in H_0^1$.
Otrzymujemy ciągły operator $T:L^2\to H_0^1\hookrightarrow L^2$.
Ponieważ zanurzenie $H_0^1(\Omega)\hookrightarrow L^2(\Omega)$
jest zwarte, operator $T$ działający w $L^2$ jest zwarty.

Z symetrii formy wynika
$(Tf,g)_{L^2}=q_c(Tf,Tg)=(f,Tg)_{L^2}$;
ponadto $(Tf,f)_{L^2}=q_c[Tf]\ge0$. Gdy $Tf=0$,
to $(f,v)_{L^2}=q_c(0,v)=0$ dla wszystkich
$v\in H_0^1(\Omega)$; z gęstości wynika $f=0$.
Zatem $\ker T=\{0\}$.
Twierdzenie spektralne dla zwartych operatorów samosprzężonych
(rozdział~\ref{chap:pde12}) daje bazę ortonormalną
$(\varphi_k)$ z $T\varphi_k=\mu_k\varphi_k$,
$\mu_k>0$, $\mu_k\downarrow0$. Z definicji $T$ mamy $q(Tf,v)=(f-cTf,v)_{L^2}$,
więc $Tf\in D(L_{A,V})$ oraz $(L_{A,V}+cI)Tf=f$.
Odwrotnie, dla $u\in D(L_{A,V})$ kładziemy
$f=L_{A,V}u+cu$. Wtedy $q_c(u,v)=(f,v)_{L^2}$,
a jednoznaczność rozwiązania daje $u=Tf$.
Zatem $D(L_{A,V})=\operatorname{ran}T$ oraz
$L_{A,V}=T^{-1}-cI$.

Współrzędne w bazie $(\varphi_k)$ dają dokładny opis
dziedziny i działania operatora. Dla $u_k=(u,\varphi_k)_{L^2}$
mamy
\[
 D(L_{A,V})=
 \left\{u\in L^2:\sum_{k\ge1}\mu_k^{-2}|u_k|^2<\infty\right\},
 \qquad
 L_{A,V}u=\sum_{k\ge1}(\mu_k^{-1}-c)u_k\varphi_k.
\]
Istotnie, $u=Tf$ wtedy i tylko wtedy, gdy współrzędne
$f_k=\mu_k^{-1}u_k$ tworzą ciąg z $\ell^2$.
W szczególności dziedzina zawiera wszystkie skończone
kombinacje funkcji własnych, jest więc gęsta.
Oznaczmy $\lambda_k=\mu_k^{-1}-c$. Ponieważ $c$ jest
stałe, warunek dziedziny jest równoważny
$\sum_k\lambda_k^2|u_k|^2<\infty$.
Operator diagonalny jest symetryczny. Jeżeli $g$ należy
do dziedziny jego sprzężenia, testowanie definicji
sprzężenia na $\varphi_k$ daje współrzędne
$(L_{A,V}^*g,\varphi_k)=\lambda_k(g,\varphi_k)$.
Ich całkowalność kwadratowa wymusza
$g\in D(L_{A,V})$, co dowodzi samosprzężoności.
Identyczny argument działa po zespoleniu.

Mamy $\lambda_k\to+\infty$. Dla każdego
$z\in\C\setminus\{\lambda_k:k\ge1\}$ odległość $z$
od zbioru wartości własnych jest dodatnia, a ciągi
$(1/(\lambda_k-z))$ oraz $(\lambda_k/(\lambda_k-z))$
są ograniczone. Operator diagonalny o mnożnikach
$1/(\lambda_k-z)$ jest więc ograniczoną odwrotnością
$L_{A,V}-zI$, przyjmującą wartości w jego dziedzinie.
To dowodzi, że nie ma innych punktów widma.
Ponadto $(L_{A,V}+cI)^{-1}=T$ jest zwarty;
z tożsamości rezolwenty wynika zwartość każdej rezolwenty.
Ograniczenie z dołu wynika także bezpośrednio
z $q[u]\ge-M\|u\|_2^2$.
\end{proof}

\section{Zasada min--max i jej konsekwencje}
\label{sec:pde18-minmax}
\begin{twierdzenie}[Zasada min--max Couranta--Fischera, wersja formowa]
\label{thm:pde18-minmax}
Niech $L_{A,V}$ będzie operatorem z
twierdzenia~\ref{thm:pde18-compact}, a $q$ jego formą.
Dla każdego $k\ge1$ zachodzi
\begin{equation}\label{eq:pde18-minmax}
 \lambda_k=\min_{\substack{E\subset H_0^1(\Omega)\\\dim E=k}}
 \ \max_{u\in E\setminus\{0\}}\frac{q[u]}{\|u\|_2^2}
 =\max_{\substack{F\subset L^2(\Omega)\\\dim F=k-1}}
 \ \inf_{\substack{u\in H_0^1(\Omega)\setminus\{0\}\\u\perp F}}
 \frac{q[u]}{\|u\|_2^2}.
\end{equation}
W drugim wyrażeniu dla $k=1$ przyjmujemy $F=\{0\}$.
\end{twierdzenie}
\begin{proof}
Niech $(\varphi_j)$ będzie ortonormalną bazą funkcji
własnych, a wartości własne niech będą uporządkowane rosnąco.
Niech $c$ będzie przesunięciem użytym w dowodzie
twierdzenia~\ref{thm:pde18-compact}. Dla każdego
$u\in H_0^1(\Omega)$ z definicji operatora i symetrii formy
mamy $q_c(u,\varphi_j)=(\lambda_j+c)(u,\varphi_j)_{L^2}$.
W konsekwencji układ
$\psi_j=(\lambda_j+c)^{-1/2}\varphi_j$ jest ortonormalny
w przestrzeni Hilberta $(H_0^1(\Omega),q_c)$.
Jest też zupełny: jeśli $q_c(u,\psi_j)=0$ dla wszystkich
$j$, to $(u,\varphi_j)_{L^2}=0$ dla wszystkich $j$,
a zupełność bazy w $L^2$ daje $u=0$.
Z tożsamości Parsevala w przestrzeni formowej wynika
dla $c_j=(u,\varphi_j)_{L^2}$, że
\[
 q_c[u]=\sum_{j\ge1}|q_c(u,\psi_j)|^2
       =\sum_{j\ge1}(\lambda_j+c)|c_j|^2.
\]
Ponadto sumy częściowe $\sum_{j=1}^n c_j\varphi_j$
zbiegają do $u$ w normie formy. Odejmując
$c\|u\|_2^2=c\sum_j|c_j|^2$, otrzymujemy
$q[u]=\sum_j\lambda_j|c_j|^2$; szereg jest bezwzględnie
zbieżny, ponieważ $|\lambda_j|\le\lambda_j+c+c$.
Maksimum ilorazu Rayleigha na każdej skończenie wymiarowej
podprzestrzeni istnieje: wystarczy ograniczyć ciągłą
funkcję $q$ do zwartej sfery jednostkowej tej podprzestrzeni
w normie $L^2$.

Jeśli $E$ ma wymiar $k$, odwzorowanie
$u\mapsto((u,\varphi_1),\ldots,(u,\varphi_{k-1}))$
z $E$ do $\R^{k-1}$ ma niezerowe jądro. Istnieje więc
$0\ne u\in E$ ortogonalne do $\varphi_1,\ldots,
\varphi_{k-1}$, a wtedy $q[u]\ge\lambda_k\|u\|_2^2$.
Przy $E_k=\operatorname{span}\{\varphi_1,\ldots,\varphi_k\}$
każdy niezerowy $u$ ma iloraz Rayleigha najwyżej
$\lambda_k$. To dowodzi pierwszej równości.

Dla dowolnego $F$ o wymiarze $k-1$ istnieje niezerowy
$u\in E_k\cap F^\perp$. Stąd
$\inf_{u\perp F}q[u]/\|u\|_2^2\le\lambda_k$.
Jeśli $F=\operatorname{span}\{\varphi_1,\ldots,
\varphi_{k-1}\}$, wspomniana suma spektralna daje
odwrotną nierówność; równość osiąga $\varphi_k$.
Otrzymujemy drugą równość, a oba zewnętrzne ekstrema
są osiągane na wskazanych podprzestrzeniach.
\end{proof}

\begin{wniosek}[Porównanie współczynników i potencjałów]
\label{cor:pde18-compare}
Niech $L_1,L_2$ będą operatorami z twierdzenia
~\ref{thm:pde18-compact}, mającymi wspólną dziedzinę formową
$H_0^1(\Omega)$, a ich symetryczne formy spełniają
$q_1[u]\le q_2[u]$ dla każdego $u\in H_0^1(\Omega)$.
Wtedy $\lambda_k(L_1)\le\lambda_k(L_2)$ dla każdego $k$.
W szczególności, gdy $A$ spełnia warunki z
\ref{sec:pde18-form}, a $v_-\le V\le v_+$ prawie wszędzie,
to
\begin{equation}\label{eq:pde18-compare}
 a_0\lambda_k(-\Delta_D)+v_-
 \le\lambda_k(L_{A,V})\le
 a_1\lambda_k(-\Delta_D)+v_+.
\end{equation}
\end{wniosek}
\begin{proof}
Dla dowolnej $k$-wymiarowej podprzestrzeni $E$ maksimum
ilorazu Rayleigha $q_1$ nie przewyższa odpowiedniego maksimum
$q_2$. Bierzemy minimum po $E$ i stosujemy
\eqref{eq:pde18-minmax}. Nierówności
$a_0\|\nabla u\|_2^2+v_-\|u\|_2^2\le q[u]
\le a_1\|\nabla u\|_2^2+v_+\|u\|_2^2$
prowadzą do~\eqref{eq:pde18-compare}.
\end{proof}

\begin{wniosek}[Monotoniczność Dirichleta względem obszaru]
\label{cor:pde18-domain}
Jeżeli $\Omega_1\subset\Omega_2$ są niepustymi ograniczonymi
zbiorami otwartymi,
a $-\Delta_{D,\Omega_i}$ oznacza laplasjan Dirichleta
w $L^2(\Omega_i)$, to
$\lambda_k(-\Delta_{D,\Omega_2})\le
\lambda_k(-\Delta_{D,\Omega_1})$ dla każdego $k$.
\end{wniosek}
\begin{proof}
Dla $u\in C_c^\infty(\Omega_1)$ przedłużenie przez zero
należy do $C_c^\infty(\Omega_2)$ i zachowuje
$\|u\|_2$ oraz $\|\nabla u\|_2$. Przez przejście
do granicy w normie $H^1$ otrzymujemy izometryczne
zanurzenie $H_0^1(\Omega_1)$ w $H_0^1(\Omega_2)$
z tymi samymi własnościami.
Minimum w~\eqref{eq:pde18-minmax} dla $\Omega_2$
biegnie zatem po większej rodzinie podprzestrzeni.
\end{proof}

\begin{przyklad}[Odcinek i prostokąt]
Na przedziale $(0,L)$, gdzie $L>0$, rozwiązujemy
$-u''=\lambda u$, $u(0)=u(L)=0$.
Z tożsamości $\lambda\|u\|_2^2=\|u'\|_2^2$ wynika
$\lambda>0$ dla niezerowej funkcji własnej:
przy zerowej energii $u$ byłaby stała i znikałaby na brzegu.
Ogólne rozwiązanie dla $\lambda>0$ ma postać
$u(x)=a\cos(\sqrt\lambda x)+b\sin(\sqrt\lambda x)$.
Pierwszy warunek daje $a=0$, a drugi
$\sqrt\lambda L=n\pi$, $n\ge1$. Stąd
$\lambda_n=(\pi n/L)^2$, a unormowane funkcje własne
$\varphi_n(x)=\sqrt{2/L}\sin(n\pi x/L)$.
Normalizacja wynika z całki $\int_0^L\sin^2(n\pi x/L)
\,\mathrm dx=L/2$. Zupełność układu sinusów w $L^2(0,L)$
wynika z zupełności układu trygonometrycznego na $(-L,L)$
po nieparzystym przedłużeniu funkcji.

Na prostokącie $\Omega=(0,L_1)\times(0,L_2)$ produkty
$\varphi_m^{(L_1)}(x_1)\varphi_n^{(L_2)}(x_2)$
spełniają warunek Dirichleta i są funkcjami własnymi
o wartościach $\lambda_{m,n}=\pi^2(m^2/L_1^2+n^2/L_2^2)$,
$m,n\ge1$. Ortonormalność wynika z twierdzenia Fubiniego,
a zupełność z zupełności układów jednowymiarowych:
jeżeli funkcja $f\in L^2(\Omega)$ jest ortogonalna do
wszystkich produktów, to dla każdego $n$ funkcja
$x_1\mapsto\int_0^{L_2}f(x_1,x_2)\varphi_n^{(L_2)}(x_2)
\,\mathrm dx_2$ należy do $L^2(0,L_1)$ i znika
prawie wszędzie. Po odrzuceniu przeliczalnej sumy
zbiorów miary zero, zupełność układu w drugim przedziale
daje $f=0$ prawie wszędzie. Gdy $L_1=L_2$ i $m\ne n$,
pary $(m,n)$ i $(n,m)$ dają tę samą wartość własną
oraz ortogonalne funkcje własne. Widmo nie musi być proste.
\end{przyklad}


\begin{stwierdzenie}[Skalowanie widma przy dylatacji obszaru]
\label{prop:pde18-scale}
Niech $\Omega\subset\R^d$ będzie ograniczonym obszarem,
$a>0$, a $a\Omega=\{ax:x\in\Omega\}$.
Dla każdego $k\ge1$ zachodzi
$\lambda_k(-\Delta_{D,a\Omega})=a^{-2}\lambda_k(-\Delta_{D,\Omega})$.
W konsekwencji
$N_{D,a\Omega}(\Lambda)=N_{D,\Omega}(a^2\Lambda)$.
\end{stwierdzenie}
\begin{proof}
Dla $v\in H_0^1(\Omega)$ określamy
$Uv(x)=a^{-d/2}v(x/a)$ dla $x\in a\Omega$.
Podstawienie $x=ay$ daje
\[
 \|Uv\|_{L^2(a\Omega)}^2
 =a^{-d}\int_{a\Omega}|v(x/a)|^2\,\mathrm d\lambda_d(x)
 =\|v\|_{L^2(\Omega)}^2.
\]
Odwrotnością jest $U^{-1}w(y)=a^{d/2}w(ay)$,
więc $U$ jest izometrycznym izomorfizmem (unitarnym
po zespoleniu). Dla funkcji gładkich o zwartym nośniku
$\nabla Uv(x)=a^{-d/2-1}\nabla v(x/a)$, skąd
$\|\nabla Uv\|_{L^2(a\Omega)}^2
=a^{-2}\|\nabla v\|_{L^2(\Omega)}^2$.
Przez gęstość i zastosowanie odwrotności wnioskujemy,
że $U$ przeprowadza $H_0^1(\Omega)$ na $H_0^1(a\Omega)$
oraz że tożsamość dla gradientu obowiązuje na całej
dziedzinie formowej.
Zasada min--max zastosowana do $U(E)$ i $U^{-1}(E)$
daje pierwszą równość, a definicja funkcji zliczającej drugą.
\end{proof}

\section{Funkcja zliczająca i prawo Weyla}
\label{sec:pde18-weyl}
\begin{definicja}[Funkcja zliczająca]
Dla operatora o dyskretnym widmie ograniczonym z dołu
i wartościach własnych $\lambda_k(L)\to+\infty$ określamy
$N_L(\Lambda)=\#\{k\ge1:\lambda_k(L)\le\Lambda\}$,
licząc krotności wartości własnych. Symbol $\omega_d$
oznacza objętość jednostkowej kuli w $\R^d$.
\end{definicja}

\begin{lemat}[Liczba modów na sześcianie]
\label{lem:pde18-cube}
Niech $Q=(0,h)^d$, $h>0$. Dla laplasjanu z warunkiem
Dirichleta albo Neumanna zachodzi
\begin{equation}\label{eq:pde18-cube}
 N_{D,Q}(\Lambda)=\frac{\omega_dh^d}{(2\pi)^d}
 \Lambda^{d/2}+O_{d,h}(1+\Lambda^{(d-1)/2}),
 \qquad \Lambda\to\infty,
\end{equation}
a dla $N_{N,Q}(\Lambda)$ obowiązuje ten sam człon główny
i oszacowanie reszty tego samego rzędu.
\end{lemat}
\begin{proof}
W przypadku Dirichleta bierzemy produkty unormowanych
sinusów z poprzedniego przykładu. Dla Neumanna na $(0,h)$
bierzemy $h^{-1/2}$ oraz $\sqrt{2/h}\cos(n\pi x/h)$,
$n\ge1$. Spełniają one zerowy warunek na pochodną
na końcach, a zupełność wynika z parzystego przedłużenia
na $(-h,h)$ i zupełności układu trygonometrycznego.
Produkty w $d$ wymiarach są zupełne przez powtarzanie
argumentu Fubiniego z przykładu prostokąta. Całkowanie
przez części w każdej zmiennej sprawdza równanie formowe
dla testów gładkich (w przypadku Neumanna aż do brzegu);
przez gęstość równanie obowiązuje dla wszystkich testów
z odpowiedniej dziedziny formowej. Zatem wszystkie
wartości własne mają postać
$\pi^2h^{-2}\sum_{j=1}^d n_j^2$, przy czym
$n\in\{1,2,\ldots\}^d$ dla Dirichleta i
$n\in\{0,1,\ldots\}^d$ dla Neumanna. Każdy indeks
wielokrotny liczymy oddzielnie, także gdy odpowiadające
mu wartości własne się pokrywają.

Niech $R=h\sqrt\Lambda/\pi$ dla $\Lambda\ge0$ oraz
$K_R=\{x\in[0,\infty)^d:|x|\le R\}$.
Dla każdego $n\in\{1,2,\ldots\}^d$ o $|n|\le R$
bierzemy półotwartą kostkę $\prod_{i=1}^d[n_i-1,n_i)$.
Kostki te są rozłączne i ich suma $S_R$ ma miarę
$N_{D,Q}(\Lambda)$. Mamy
$K_{(R-\sqrt d)_+}\subset S_R\subset K_R$
z dokładnością do zbiorów miary zero: dla punktu
w mniejszym ortancie wybieramy $n_i=\lfloor x_i\rfloor+1$,
wtedy $|n|\le|x|+\sqrt d\le R$ (dla $R<\sqrt d$
mniejszy ortant ma miarę zero). Drugie zawieranie
wynika z $0\le x_i\le n_i$.
Zatem
\[
 2^{-d}\omega_d(R-\sqrt d)_+^d
 \le N_{D,Q}(\Lambda)\le 2^{-d}\omega_dR^d.
\]
Różnica skrajnych wyrażeń jest
$O_d(1+R^{d-1})$: dla $R\ge\sqrt d$ stosujemy
twierdzenie o wartości średniej do funkcji $t^d$,
a dla mniejszych $R$ różnica jest ograniczona.
Indeksy Neumanna z co najmniej jedną zerową współrzędną
mają liczbę nie większą niż $d(\lfloor R\rfloor+1)^{d-1}$,
gdyż wybieramy zerową współrzędną, a pozostałe
należą do $\{0,\ldots,\lfloor R\rfloor\}$.
Po podstawieniu $R=h\sqrt\Lambda/\pi$ otrzymujemy
oba oszacowania, także dla $d=1$.
\end{proof}

\begin{lemat}[Rozdzielanie na sześciany i ograniczenie liczby modów]
\label{lem:pde18-bracket}
Niech $Q_1,\ldots,Q_m$ będą otwartymi sześcianami
pochodzącymi z tej samej siatki, o parami rozłącznych
wnętrzach, a $P=\operatorname{int}(\bigcup_{j=1}^m
\overline{Q_j})$, gdzie $m\ge1$. Dla każdego $\Lambda\in\R$
\begin{equation}\label{eq:pde18-bracket}
 \sum_{j=1}^m N_{D,Q_j}(\Lambda)
 \le N_{D,P}(\Lambda)
 \le\sum_{j=1}^m N_{N,Q_j}(\Lambda).
\end{equation}
\end{lemat}
\begin{proof}
Przestrzeń $L^2(P)$ utożsamiamy izometrycznie z
$\bigoplus_{j=1}^m L^2(Q_j)$ za pomocą ograniczeń,
gdyż $P\setminus\bigcup_j Q_j$ leży w skończonej sumie
hiperpłaszczyzn i ma miarę zero. Przy tym utożsamieniu
mamy zawierania dziedzin formowych
\[
 \bigoplus_{j=1}^m H_0^1(Q_j)
 \subset H_0^1(P)\subset
 \bigoplus_{j=1}^m H^1(Q_j).
\]
Pierwsze wynika z przedłużenia przez zero z każdej
komórki i skończonego sumowania, drugie z ograniczenia
funkcji Sobolewa na otwarte podzbiory. Energia na każdej
z tych przestrzeni jest dana przez sumę
$\sum_j\int_{Q_j}|\nabla u_j|^2\,\mathrm d\lambda_d$,
a kwadrat normy $L^2$ przez sumę kwadratów norm.

Dla form na skrajnych przestrzeniach bazę własną tworzy
suma układów własnych poszczególnych sześcianów,
przedłużonych przez zero jako funkcje $L^2(P)$.
W przypadku Neumanna konstrukcja z twierdzenia
~\ref{thm:pde18-compact} działa na każdej przestrzeni
$H^1(Q_j)$: po dodaniu normy $L^2$ forma jest
koercywna, a zanurzenie jest zwarte z twierdzenia
~\ref{thm:pde10-Rellich-p}. Dowód zasady min--max
pozostaje ten sam dla tych form i ich skończonych sum prostych.
Jeżeli $\lambda_k^-$, $\lambda_k^+$ oznaczają uporządkowane
wartości własne sum odpowiednio Dirichleta i Neumanna,
to zawierania dziedzin dają
$\lambda_k^+\le\lambda_k(-\Delta_{D,P})\le\lambda_k^-$.
Dla dowolnego progu $\Lambda$ oznacza to
$N^-(\Lambda)\le N_{D,P}(\Lambda)\le N^+(\Lambda)$,
gdzie $N^\pm$ są sumami funkcji zliczających komórek,
co dowodzi~\eqref{eq:pde18-bracket}.
\end{proof}

\begin{twierdzenie}[Prawo Weyla dla laplasjanu Dirichleta]
\label{thm:pde18-weyl}
Niech $\Omega\subset\R^d$ będzie ograniczonym obszarem
o brzegu miary Lebesgue'a zero, tj.
$\lambda_d(\partial\Omega)=0$. Niech $-\Delta_D$
oznacza laplasjan Dirichleta w $L^2(\Omega)$.
Wówczas
\begin{equation}\label{eq:pde18-weyl}
 \lim_{\Lambda\to\infty}
 \frac{N_{-\Delta_D}(\Lambda)}{\Lambda^{d/2}}
 =\frac{\omega_d}{(2\pi)^d}|\Omega|.
\end{equation}
Równoważnie, gdy $k\to\infty$,
\begin{equation}\label{eq:pde18-lambda-weyl}
 \lambda_k(-\Delta_D)\sim
 (2\pi)^2\left(\frac{k}{\omega_d|\Omega|}\right)^{2/d}.
\end{equation}
\end{twierdzenie}
\begin{proof}
Ustalmy siatkę sześcianów o boku $h>0$ w $\R^d$.
Ze względu na ograniczoność $\Omega$ tylko skończenie
wiele komórek siatki przecina $\Omega$. Przy dostatecznie
małym $h$ istnieje przynajmniej jedna domknięta komórka
zawarta w $\Omega$ (wybieramy kulę wewnętrzną
o promieniu większym niż $\sqrt d\,h$).
Przez $P_h^-$ oznaczmy wnętrze sumy wszystkich domkniętych
komórek zawartych w $\Omega$, a przez $P_h^+$ wnętrze sumy
wszystkich domkniętych komórek, których wnętrza przecinają
$\Omega$. Mamy $P_h^-\subset\Omega\subset P_h^+$.
Istotnie, pierwszy zbiór jest zawarty w $\Omega$
z definicji. Dla każdego $x\in\Omega$ istnieje kula
$B_\rho(x)\subset\Omega$. Każdy punkt tej kuli należy do
domknięcia komórki, której wnętrze przecina $\Omega$
(wnętrze komórki jest gęste w jej domknięciu,
a $\Omega$ jest otwarty). Cała kula znajduje się więc
w sumie domkniętych komórek wyznaczających $P_h^+$,
czyli $x$ należy do wnętrza tej sumy.

Z monotoniczności Dirichleta względem obszaru wynika
$N_{D,P_h^-}(\Lambda)\le N_{D,\Omega}(\Lambda)
\le N_{D,P_h^+}(\Lambda)$.
Stosujemy nierówności~\eqref{eq:pde18-bracket} do obu
wielościanów, a następnie~\eqref{eq:pde18-cube} do ich
skończenie wielu komórek. Dla każdego ustalonego $h$
otrzymujemy po podzieleniu przez $\Lambda^{d/2}$ i
przejściu $\Lambda\to\infty$:
\[
 \frac{\omega_d}{(2\pi)^d}|P_h^-|
 \le\liminf_{\Lambda\to\infty}\frac{N_{D,\Omega}(\Lambda)}{\Lambda^{d/2}}
 \le\limsup_{\Lambda\to\infty}\frac{N_{D,\Omega}(\Lambda)}{\Lambda^{d/2}}
 \le\frac{\omega_d}{(2\pi)^d}|P_h^+|.
\]
Z dokładnością do zbioru miary zero różnica
$P_h^+\setminus P_h^-$ jest sumą wnętrz komórek,
których wnętrza przecinają $\Omega$, ale których
domknięcia nie są zawarte w $\Omega$. Każde takie
domknięcie zawiera punkt z $\Omega$ i punkt z jej
dopełnienia; łączący je odcinek w domkniętej komórce
przecina $\partial\Omega$. Stąd wszystkie punkty
tej komórki leżą w odległości co najwyżej $\sqrt d\,h$
od $\partial\Omega$. Zatem
\[
 |P_h^+\setminus P_h^-|
 \le\lambda_d(\{x:\operatorname{dist}(x,\partial\Omega)
 \le\sqrt d\,h\}).
\]
Ponieważ $\partial\Omega$ jest zwartym zbiorem miary zero,
prawa strona dąży do zera: domknięte otoczenia maleją
do $\partial\Omega$, a pierwsze z nich ma skończoną miarę,
więc stosujemy ciągłość miary z góry.
Zatem $|P_h^-|\to|\Omega|$ oraz
$|P_h^+|\to|\Omega|$, co daje~\eqref{eq:pde18-weyl}.

Aby uzyskać~\eqref{eq:pde18-lambda-weyl}, oznaczmy
$C_\Omega=\omega_d|\Omega|/(2\pi)^d>0$. Dla
$0<\eps<C_\Omega$ i dużych $\Lambda$ mamy
$(C_\Omega-\eps)\Lambda^{d/2}\le N(\Lambda)
\le(C_\Omega+\eps)\Lambda^{d/2}$.
Dla dostatecznie dużego $k$ stosujemy dolne oszacowanie
do $\Lambda_+=(k/(C_\Omega-\eps))^{2/d}$, co daje
$N(\Lambda_+)\ge k$, a zatem $\lambda_k\le\Lambda_+$.
Górne oszacowanie zastosowane do $\Lambda=\lambda_k$
daje $k\le N(\lambda_k)\le(C_\Omega+\eps)\lambda_k^{d/2}$,
czyli
\[
 \left(\frac{k}{C_\Omega+\eps}\right)^{2/d}
 \le\lambda_k\le
 \left(\frac{k}{C_\Omega-\eps}\right)^{2/d}.
\]
Dzielimy przez $(k/C_\Omega)^{2/d}$ i przechodzimy
do granic dolnej i górnej, a następnie $\eps\downarrow0$.
Odwrotnie, jeśli $\lambda_k\sim(k/C_\Omega)^{2/d}$,
to dla dużych $k$ i dowolnego $0<\eps<1$ mamy
$(1-\eps)(k/C_\Omega)^{2/d}\le\lambda_k
\le(1+\eps)(k/C_\Omega)^{2/d}$. Dla dużego $\Lambda$
kładziemy $k=N(\Lambda)$; wtedy
$\lambda_k\le\Lambda<\lambda_{k+1}$, więc
\[
 C_\Omega\left(\frac{\Lambda}{1+\eps}\right)^{d/2}-1
 <N(\Lambda)\le
 C_\Omega\left(\frac{\Lambda}{1-\eps}\right)^{d/2}.
\]
Po podzieleniu przez $\Lambda^{d/2}$ i przejściu
$\Lambda\to\infty$, a następnie $\eps\downarrow0$,
odzyskujemy~\eqref{eq:pde18-weyl}. Oba argumenty
uwzględniają krotności, także w punktach skoku funkcji
zliczającej.
\end{proof}

\begin{uwaga}[Co pozostaje trudne w ogólnej teorii Weyla?]
Powyższy dowód wykorzystuje wyłącznie geometrię siatki,
min--max i jawne widmo sześcianu. Dla operatorów z macierzą
$A(x)$ oczekiwany człon główny dla odpowiednio regularnych
współczynników to
$(2\pi)^{-d}\omega_d\Lambda^{d/2}
\int_\Omega(\det A(x))^{-1/2}\,\mathrm d\lambda_d(x)$.
Nie wyprowadzamy tu tej mocniejszej asymptotyki, ponieważ
wymaga ona dodatkowej lokalizacji w przestrzeni fazowej.
Natomiast porównanie~\eqref{eq:pde18-compare} natychmiast
daje ilościowe dwustronne ograniczenia $\lambda_k(L_{A,V})
\asymp k^{2/d}$ dla dostatecznie dużych $k$.
\end{uwaga}

\section{Nierówność Li--Yau: dolne oszacowanie sumy wartości własnych}
\label{sec:pde18-liyau}
\begin{lemat}[Zasada wypełniania najniższych częstotliwości]
\label{lem:pde18-bathtub}
Niech $F:\R^d\to[0,M]$ będzie mierzalna,
$0<\int_{\R^d}F\,\mathrm d\lambda_d=k<\infty$,
$M>0$, i niech $\rho=(k/(M\omega_d))^{1/d}$.
Wtedy
\begin{equation}\label{eq:pde18-bathtub}
 \int_{\R^d}|\xi|^2F(\xi)\,\mathrm d\lambda_d(\xi)
 \ge\frac{d}{d+2}k\rho^2.
\end{equation}
\end{lemat}
\begin{proof}
Jeżeli moment $\int_{\R^d}|\xi|^2F\,\mathrm d\lambda_d$
jest nieskończony, nierówność jest natychmiastowa.
Zakładamy dalej, że jest skończony.
Oznaczmy $F_0=M\chi_{B_\rho(0)}$.
Na $B_\rho$ mamy $(|\xi|^2-\rho^2)(F-F_0)\ge0$,
ponieważ oba czynniki są niedodatnie; poza kulą oba są
nieujemne. Całkując i korzystając z równości
$\int_{\R^d}F\,\mathrm d\lambda_d=\int_{\R^d}F_0\,\mathrm d\lambda_d=k$, dostajemy
$\int_{\R^d}|\xi|^2F(\xi)\,\mathrm d\lambda_d(\xi)\ge\int_{\R^d}|\xi|^2F_0(\xi)\,\mathrm d\lambda_d(\xi)$.
Całkowanie radialne daje
\[
 \int_{\R^d}|\xi|^2F_0(\xi)\,\mathrm d\lambda_d(\xi)
 =Md\omega_d\int_0^\rho r^{d+1}\,\mathrm dr
 =\frac{d}{d+2}k\rho^2.
\]
Równość zachodzi wtedy i tylko wtedy, gdy $F=F_0$
prawie wszędzie: nieujemna funkcja
$(|\xi|^2-\rho^2)(F-F_0)$ musi mieć całkę zero,
a pierwszy czynnik znika tylko na sferze miary zero.
\end{proof}

\begin{twierdzenie}[Nierówność Li--Yau dla laplasjanu Dirichleta]
\label{thm:pde18-liyau}
Niech $\Omega\subset\R^d$ będzie ograniczonym obszarem,
a $0<\lambda_1\le\lambda_2\le\cdots$ wartościami własnymi
laplasjanu Dirichleta. Dla każdego $k\ge1$ zachodzi
\begin{equation}\label{eq:pde18-liyau}
 \sum_{j=1}^k\lambda_j\ge
 \frac{d}{d+2}(2\pi)^2
 \left(\frac{1}{\omega_d|\Omega|}\right)^{2/d}
 k^{1+2/d}.
\end{equation}
W szczególności
$\lambda_k\ge\frac d{d+2}(2\pi)^2
(k/(\omega_d|\Omega|))^{2/d}$.
\end{twierdzenie}
\begin{proof}
Niech $\varphi_1,\ldots,\varphi_k$ będą ortonormalnymi
funkcjami własnymi w $L^2(\Omega)$ i przedłużmy je przez
zero do $\R^d$. Ponieważ należą do $H_0^1(\Omega)$,
ich przedłużenia należą do $H^1(\R^d)$
(twierdzenie~\ref{thm:pde9-zero-extension}), a ich słabe
gradienty są przedłużeniami przez zero gradientów na $\Omega$.
Ponadto przedłużone funkcje należą do $L^1(\R^d)$, gdyż
$\|\varphi_j\|_1\le|\Omega|^{1/2}\|\varphi_j\|_2$.
Używamy unitarnej transformaty Fouriera
$\widehat f(\xi)=(2\pi)^{-d/2}
\int_{\R^d}f(x)e^{-ix\cdot\xi}\,\mathrm d\lambda_d(x)$.
Definiujemy $F(\xi)=\sum_{j=1}^k|\widehat\varphi_j(\xi)|^2$.

Z twierdzenia Plancherela wynika
$\int_{\R^d}F(\xi)\,\mathrm d\lambda_d(\xi)=k$.
Dla $f\in H^1(\R^d)$ zachodzi
$\widehat{\partial_\ell f}=i\xi_\ell\widehat f$
w $L^2$: najpierw całkujemy przez części dla
$f\in C_c^\infty(\R^d)$, następnie przybliżamy w $H^1$
i stosujemy Plancherela. Identyfikację granicy
$i\xi_\ell\widehat f$ uzasadnia zbieżność w $L^2$
na każdej ograniczonej kuli w przestrzeni częstotliwości,
gdzie mnożnik $\xi_\ell$ jest ograniczony.
Testując równanie własne funkcją $\varphi_j$, mamy
$\|\nabla\varphi_j\|_2^2=\lambda_j$.
Stąd
\begin{equation}\label{eq:pde18-fft-energy}
 \int_{\R^d}|\xi|^2F(\xi)\,\mathrm d\lambda_d(\xi)
 =\sum_{j=1}^k\int_\Omega|\nabla\varphi_j|^2
 \,\mathrm d\lambda_d(x)=\sum_{j=1}^k\lambda_j.
\end{equation}
Dla każdego ustalonego $\xi$ funkcje
$\varphi_1,\ldots,\varphi_k$ są ortonormalne także
w zespolonej przestrzeni $L^2(\Omega;\C)$.
Przy konwencji $(f,g)=\int_\Omega f\overline g\,\mathrm d\lambda_d$
mamy $\widehat\varphi_j(\xi)=(\varphi_j,g_\xi)$,
gdzie $g_\xi(x)=(2\pi)^{-d/2}e^{ix\cdot\xi}$.
Nierówność Bessela daje
\[
 F(\xi)=\sum_{j=1}^k|(\varphi_j,g_\xi)|^2
 \le\|g_\xi\|_2^2=(2\pi)^{-d}|\Omega|=:M.
\]
Stosujemy lemat~\ref{lem:pde18-bathtub} z tym $M$,
$\rho^d=k/(M\omega_d)=(2\pi)^dk/(\omega_d|\Omega|)$.
Wstawiając oszacowanie do~\eqref{eq:pde18-fft-energy},
otrzymujemy~\eqref{eq:pde18-liyau}. Ponieważ
$k\lambda_k\ge\sum_{j=1}^k\lambda_j$, otrzymujemy
podany dolny szacunek pojedynczej wartości własnej.
\end{proof}

\begin{uwaga}[Dlaczego współczynnik jest naturalny?]
Prawa strona~\eqref{eq:pde18-liyau} jest dokładnie całką
średniego kosztu $|\xi|^2$ przy równomiernym wypełnianiu
najniższych częstotliwości w przestrzeni fazowej. Jest to wariacyjna interpretacja problemu
optymalizacyjnego dla funkcji $F$ o zadanej masie i ograniczeniu
punktowym; dowód nie wymaga znajomości prawa Weyla.
\end{uwaga}


\begin{wniosek}[Oszacowanie funkcji zliczającej bez prawa Weyla]
\label{cor:pde18-count-liyau}
Niech $d\ge1$, a $\Omega\subset\R^d$ będzie ograniczonym
obszarem. Dla $\Lambda>0$ funkcja zliczająca laplasjanu
Dirichleta spełnia
\[
 N_{-\Delta_D}(\Lambda)\le
 \frac{\omega_d|\Omega|}{(2\pi)^d}
 \left(\frac{d+2}{d}\right)^{d/2}\Lambda^{d/2}.
\]
Wniosek nie wymaga założenia $\lambda_d(\partial\Omega)=0$.
\end{wniosek}
\begin{proof}
Niech $k=N_{-\Delta_D}(\Lambda)$. Gdy $k=0$, nie ma
czego dowodzić. W przeciwnym razie $\lambda_k\le\Lambda$.
Na mocy nierówności~\eqref{eq:pde18-liyau} mamy
$\Lambda\ge\lambda_k\ge\frac d{d+2}(2\pi)^2
(k/(\omega_d|\Omega|))^{2/d}$. Przekształcenie tego
oszacowania daje tezę. Zwróćmy uwagę, że jest to
\emph{nierówność dla każdej} $\Lambda>0$, a nie
asymptotyka z prawa Weyla.
\end{proof}

\section{Pierwsza wartość własna, dodatniość i prostota}
\label{sec:pde18-ground}
\begin{twierdzenie}[Prostota pierwszej wartości własnej laplasjanu]
\label{thm:pde18-simple}
Niech $\Omega\subset\R^d$ będzie ograniczonym obszarem
(nie zakładamy regularności brzegu). Dla laplasjanu Dirichleta
$\lambda_1>0$ jest prosta, a odpowiadającą jej funkcję
własną można wybrać ściśle dodatnią we wnętrzu $\Omega$.
\end{twierdzenie}
\begin{proof}
Wybieramy kostkę $Q$ zawierającą $\overline\Omega$.
Przedłużając przez zero do $Q$ i stosując wniosek
~\ref{cor:pde9-Poincare-zero}, otrzymujemy
$\|w\|_2\le C_Q\|\nabla w\|_2$ dla $w\in H_0^1(\Omega)$.
Z zasady min--max wynika więc
$\lambda_1=\min_{w\in H_0^1\setminus\{0\}}
\|\nabla w\|_2^2/\|w\|_2^2\ge C_Q^{-2}>0$.
Niech $u$ będzie rzeczywistą niezerową funkcją własną
dla $\lambda_1$. Uzasadnimy, że $|u|\in H_0^1(\Omega)$
bez odwoływania się do śladu na brzegu.
Wybieramy $u_n\in C_c^\infty(\Omega)$ zbieżne do $u$
w $H^1$. Funkcje $|u_n|$ należą do $H_0^1(\Omega)$,
gdyż mają zwarty nośnik i można je przybliżyć splotami
w $H^1$. Z reguł dla obcięć (wniosek~\ref{cor:pde9-truncation})
wynika $\|\nabla|u_n|\|_2\le\|\nabla u_n\|_2$.
Ciąg jest ograniczony w $H_0^1$, a w $L^2$ zbiega
do $|u|$ na mocy nierówności
$||u_n|-|u||\le|u_n-u|$. Słabo zbieżny podciąg
w przestrzeni Hilberta $H_0^1$ musi zatem mieć
granicę $|u|$, co daje żądane przynależenie.
Ponadto $|\nabla|u||\le|\nabla u|$ prawie wszędzie,
więc iloraz Rayleigha $w=|u|$ jest nie większy od
$\lambda_1$ i z minimalności jest równy $\lambda_1$.

Dla każdego $v\in H_0^1(\Omega)$ i $s\in\R$ mamy
$\|\nabla(w+sv)\|_2^2-\lambda_1\|w+sv\|_2^2\ge0$,
z równością w $s=0$. Pochodna tego wielomianu
w zerze znika, skąd
\[
 \int_\Omega\nabla w\cdot\nabla v\,\mathrm d\lambda_d
 =\lambda_1\int_\Omega wv\,\mathrm d\lambda_d.
\]
Zatem $w$ jest również funkcją własną.

Gładkość wewnętrzną zarówno $u$, jak i $w$ uzyskamy
z lematu Weyla~\ref{thm:pde7-weyl}, bez założeń o brzegu
i bez osobnego argumentu dla $d=1$.
Jeżeli $f$ jest dowolną z tych funkcji, określamy
$H(x,t)=f(x)\cosh(\sqrt{\lambda_1}t)$ na
$\Omega\times\R$. Funkcja jest lokalnie całkowalna
i dystrybucyjnie harmoniczna:
dla $\Phi\in C_c^\infty(\Omega\times\R)$ całkujemy
najpierw w $x$, używając
$\Delta_x f=-\lambda_1 f$, a dla członu
$\partial_t^2\Phi$ całkujemy dwukrotnie przez części w $t$.
Człony $-\lambda_1H$ i $\lambda_1H$ się znoszą.
Lemat Weyla w wymiarze $d+1\ge2$ daje gładkiego
reprezentanta $H$. Z twierdzenia Fubiniego dla pewnego
$t_0\in\R$ zachodzi
$H(\cdot,t_0)=f(\cdot)\cosh(\sqrt{\lambda_1}t_0)$
prawie wszędzie. Dzieląc przez dodatni czynnik,
otrzymujemy gładkiego reprezentanta $f$.
Ciągłość wymusza, że jego równanie dystrybucyjne
jest równaniem punktowym.

Dla $f=w$ reprezentant $H$ jest nieujemny wszędzie
(z ciągłości i nieujemności prawie wszędzie).
Gdyby $w$ miała zero w $\Omega$, to ciągłe funkcje
$H(x,t)$ i $w(x)\cosh(\sqrt{\lambda_1}t)$, równe
prawie wszędzie, byłyby równe wszędzie i $H$ osiągałaby
minimum zero w spójnym obszarze $\Omega\times\R$.
Mocna zasada maksimum dla funkcji harmonicznych
~\ref{thm:pde14-strong} dawałaby $H=0$, sprzecznie
z $u\ne0$. Zatem $w>0$ w $\Omega$.
Ponieważ gładki reprezentant $u$ spełnia $|u|=w$
wszędzie, nie ma zer; spójność $\Omega$ wymusza
stały znak $u$.

Przypuśćmy, że przestrzeń własna ma wymiar co najmniej dwa.
Można wtedy znaleźć niezerowe, rzeczywiste, ortogonalne
funkcje własne $u,v$ dla $\lambda_1$.
Po ewentualnej zmianie znaków obie byłyby dodatnie w $\Omega$,
co przeczyłoby $(u,v)_{L^2}=0$: ich iloczyn jest
dodatni na całym obszarze i całkowalny z nierówności
Cauchy'ego--Schwarza. Pierwsza wartość własna jest
więc prosta w rzeczywistej przestrzeni Hilberta,
a po zespoleniu także w zespolonej, gdyż części
rzeczywista i urojona funkcji własnej są rozwiązaniami
tego samego rzeczywistego równania.
\end{proof}

\section{Warunek Neumanna i problem zerowej energii}
\label{sec:pde18-neumann}
\begin{twierdzenie}[Widmo Neumanna i pierwsza dodatnia wartość]
\label{thm:pde18-neumann}
Niech $\Omega\subset\R^d$ będzie ograniczonym, spójnym
obszarem Lipschitzowskim, a $-\Delta_N$ laplasjanem
związanym z formą $\int_\Omega|\nabla u|^2\,
\mathrm d\lambda_d$ o dziedzinie formowej $H^1(\Omega)$.
Jego widmo jest dyskretne i nieujemne:
$0=\mu_1<\mu_2\le\mu_3\le\cdots\to\infty$.
Pierwsza przestrzeń własna składa się z funkcji stałych.
Ponadto
\begin{equation}\label{eq:pde18-neumann}
 \mu_2=\min_{\substack{u\in H^1(\Omega)\setminus\{0\}\\
 \int_\Omega u\,\mathrm d\lambda_d=0}}
 \frac{\int_\Omega|\nabla u|^2\,\mathrm d\lambda_d}
 {\int_\Omega|u|^2\,\mathrm d\lambda_d}>0.
\end{equation}
\end{twierdzenie}
\begin{proof}
Z twierdzenia~\ref{thm:pde10-Rellich-p} zanurzenie
$H^1(\Omega)\hookrightarrow L^2(\Omega)$ jest zwarte.
Forma przesunięta o iloczyn skalarny $L^2$ jest
dokładnie iloczynem skalarnym przestrzeni $H^1$.
Przestrzeń ta jest gęsta w $L^2$, ponieważ zawiera
$C_c^\infty(\Omega)$. Wszystkie kroki konstrukcji
z dowodu twierdzenia~\ref{thm:pde18-compact}
działają z $H^1$ zamiast $H_0^1$ oraz przesunięciem
$c=1$. Otrzymujemy operator samosprzężony o zwartej
rezolwencie, bazę własną oraz zasadę min--max
z dowodu twierdzenia~\ref{thm:pde18-minmax}.

Funkcje stałe należą do $H^1(\Omega)$ i spełniają
równanie formowe z prawą stroną zero. Odwrotnie,
jeżeli $u$ jest funkcją własną dla zera, testowanie
funkcją $u$ daje $\|\nabla u\|_2^2=0$.
Wniosek~\ref{cor:pde10-Poincare} daje
$\|u-u_\Omega\|_2\le C_\Omega\|\nabla u\|_2=0$,
gdzie $u_\Omega=|\Omega|^{-1}\int_\Omega u\,\mathrm d\lambda_d$.
Zatem zero jest prostą wartością własną.
Wszystkie wartości własne są nieujemne, ponieważ
energia jest nieujemna. W drugiej postaci zasady min--max
dla $k=2$ wybieramy osiągającą maksimum przestrzeń
$F=\operatorname{span}\{1\}$, czyli pierwszą przestrzeń
własną. Ortogonalność do niej jest równoważna zerowej
średniej. Minimum osiąga druga funkcja własna;
nierówność Poincarégo--Wirtingera daje dolną granicę
$C_\Omega^{-2}>0$ dla jej ilorazu Rayleigha.
Otrzymujemy~\eqref{eq:pde18-neumann}.
\end{proof}

\begin{uwaga}[Pułapka z monotonicznością]
Monotoniczność wartości własnych względem zawierania
obszarów udowodniliśmy dla warunku Dirichleta.
Nie wolno bez dodatkowych hipotez przenosić jej na
laplasjan Neumanna: przedłużenie przez zero nie działa
wtedy jako zanurzenie $H^1(\Omega_1)\to H^1(\Omega_2)$.
\end{uwaga}

\section{Ślad półgrupy cieplnej i odczytywanie wymiaru}
\label{sec:pde18-heat}
\begin{twierdzenie}[Wiodący człon śladu cieplnego]
\label{thm:pde18-heat}
Niech $\Omega$ spełnia założenia
Twierdzenia~\ref{thm:pde18-weyl}. Dla każdego $t>0$
operator $e^{t\Delta_D}$ jest klasy śladowej w $L^2(\Omega)$,
a przy $t\downarrow0$ zachodzi
\begin{equation}\label{eq:pde18-heat}
 \operatorname{Tr}(e^{t\Delta_D})=
 \sum_{k=1}^\infty e^{-t\lambda_k}
 \sim\frac{|\Omega|}{(4\pi t)^{d/2}}.
\end{equation}
\end{twierdzenie}
\begin{proof}
Prawo Weyla daje $N(\Lambda)\le C(1+\Lambda^{d/2})$.
Dla $t>0$ suma $\sum_ke^{-t\lambda_k}$ jest więc skończona.
Ponieważ $N(\Lambda)=\sum_k\chi_{[\lambda_k,\infty)}(\Lambda)$,
zasada Tonellego prowadzi do tożsamości
\begin{equation}\label{eq:pde18-heat-count}
 \sum_ke^{-t\lambda_k}=t\int_0^\infty e^{-t\Lambda}
 N(\Lambda)\,\mathrm d\Lambda
 =\int_0^\infty e^{-s}N(s/t)\,\mathrm ds.
\end{equation}
Oznaczmy $C_\Omega=\omega_d|\Omega|/(2\pi)^d$.
Z~\eqref{eq:pde18-weyl} wynika, że dla ustalonego $s>0$
$t^{d/2}N(s/t)\to C_\Omega s^{d/2}$.
Wskutek oszacowania liczby modów funkcje pod całką są
zdominowane przez $C e^{-s}(1+s^{d/2})$ dla $0<t\le1$.
Stosujemy zbieżność zdominowaną:
\[
 \lim_{t\downarrow0}t^{d/2}\sum_ke^{-t\lambda_k}
 =C_\Omega\int_0^\infty e^{-s}s^{d/2}\,\mathrm ds
 =C_\Omega\Gamma(1+d/2).
\]
Równość $\omega_d=\pi^{d/2}/\Gamma(1+d/2)$ daje
$C_\Omega\Gamma(1+d/2)=|\Omega|/(4\pi)^{d/2}$.
To dowodzi wzoru~\eqref{eq:pde18-heat}.
\end{proof}

\begin{uwaga}[Powiązanie z geometrią spektralną]
Z zachowania $\operatorname{Tr}(e^{t\Delta_D})\asymp t^{-d/2}$
można odczytać wymiar $d$, a z wiodącego współczynnika
objętość obszaru. Wyższe człony rozwinięcia dla gładkiego
brzegu zawierają jego pole i dane krzywiznowe, ale ich
wyznaczenie wymaga osobnych oszacowań jądra ciepła.
Nie wnioskujemy tu automatycznie z prawa Weyla o pełnym
rozwinięciu asymptotycznym w potęgach $t^{1/2}$.
\end{uwaga}


\begin{lemat}[Przejście abelowe od liczby modów do śladu]
\label{lem:pde18-abel}
Niech $(\nu_j)_{j\ge1}$ będzie niemalejącym ciągiem
liczb dodatnich, $\nu_j\to\infty$, oraz
$N(\Lambda)=\#\{j:\nu_j\le\Lambda\}$. Załóżmy, że dla
pewnych $a>0$, $C>0$, $0\le\gamma<a$ i $K>0$ zachodzi
$|N(\Lambda)-C\Lambda^a|\le K(1+\Lambda^\gamma)$
dla wszystkich $\Lambda\ge0$. Wtedy, gdy $t\downarrow0$,
\[
 \sum_{j\ge1}e^{-t\nu_j}
 =C\Gamma(a+1)t^{-a}+O(1+t^{-\gamma}).
\]
\end{lemat}
\begin{proof}
Ponieważ $N(\Lambda)\le C'(1+\Lambda^a)$, szereg
po lewej jest zbieżny dla $t>0$. Z twierdzenia Tonellego
oraz $e^{-t\nu_j}=t\int_{\nu_j}^\infty e^{-t\Lambda}
\,\mathrm d\Lambda$ otrzymujemy
$\sum_j e^{-t\nu_j}=t\int_0^\infty e^{-t\Lambda}
N(\Lambda)\,\mathrm d\Lambda$.
Po podstawieniu $s=t\Lambda$ wyraz główny ma wartość
$C t^{-a}\int_0^\infty e^{-s}s^a\,\mathrm ds
=C\Gamma(a+1)t^{-a}$. Moduł reszty nie przekracza
$K\int_0^\infty e^{-s}(1+t^{-\gamma}s^\gamma)
\,\mathrm ds=K(1+\Gamma(\gamma+1)t^{-\gamma})$.
\end{proof}

\begin{uwaga}[Różnica między wiodącym wyrazem a rozwinięciem]
Prawo Weyla~\eqref{eq:pde18-weyl} daje tylko $o(t^{-d/2})$
w~\eqref{eq:pde18-heat}. Oszacowanie liczby modów z jawną
resztą prowadzi, przez poprzedni lemat, do oszacowania
reszty śladu. Nie wolno jednak wyprowadzać z samego prawa
Weyla ani dwuczłonowej asymptotyki śladu, ani współczynnika
zależnego od pola brzegu.
\end{uwaga}

\section{Perturbacje widma i wzór Hellmanna--Feynmana}
\label{sec:pde18-perturb}
\begin{wniosek}[Stabilność względem potencjału]
Niech $V_1,V_2\in L^\infty(\Omega;\R)$, a macierz $A$
niech będzie ustalona, symetryczna i jednostajnie eliptyczna.
Wtedy
$|\lambda_k(L_{A,V_1})-\lambda_k(L_{A,V_2})|
\le\|V_1-V_2\|_{L^\infty(\Omega)}$ dla każdego $k\ge1$.
\end{wniosek}
\begin{proof}
Dla każdego $u$ zachodzi
$|q_{V_1}[u]-q_{V_2}[u]|
\le\|V_1-V_2\|_\infty\|u\|_2^2$.
Stosujemy dwukrotnie zasadę min--max.
\end{proof}

\begin{stwierdzenie}[Wzór Hellmanna--Feynmana przy różniczkowalnej gałęzi]
\label{prop:pde18-hf}
Niech $L_t=L_0+tW$, gdzie $L_0$ jest samosprzężonym
operatorem z dyskretnym widmem, a $W$ ograniczonym,
samosprzężonym operatorem na $L^2(\Omega)$.
Przyjmijmy, że na otwartym przedziale $I\subset\R$
można wybrać różniczkowalne w normie grafowej funkcje
$t\mapsto\lambda(t)\in\R$ i
$t\mapsto\varphi(t)\in D(L_0)$ takie, że
$L_t\varphi(t)=\lambda(t)\varphi(t)$ oraz
$\|\varphi(t)\|_2=1$. Wtedy
\begin{equation}\label{eq:pde18-hf}
 \lambda'(t)=(W\varphi(t),\varphi(t))_{L^2}.
\end{equation}
Jeżeli $W$ jest operatorem mnożenia przez funkcję $W(x)$,
to prawa strona ma postać
\[
\int_\Omega W(x)|\varphi(t,x)|^2\,\mathrm d\lambda_d(x).
\]
\end{stwierdzenie}
\begin{proof}
Różniczkujemy równość
$L_t\varphi(t)=\lambda(t)\varphi(t)$ w $L^2$.
Bierzemy iloczyn skalarny z $\varphi(t)$.
Samosprzężoność $L_t$ powoduje skasowanie składników
z $\varphi'(t)$, natomiast z warunku normalizacji
wynika $\|\varphi(t)\|_2=1$. Pozostaje~\eqref{eq:pde18-hf}.
\end{proof}

\begin{uwaga}[Dlaczego założenie różniczkowalności jest jawne?]
Wzór nie mówi sam przez się, że dowolna numerowana gałąź
$\lambda_k(t)$ jest różniczkowalna w punkcie zderzenia
wartości własnych. Dla prostych, izolowanych wartości własnych
regularność gałęzi wynika z teorii perturbacji, której pełnego
dowodu nie wprowadzamy w tej części skryptu.
\end{uwaga}

\section{Podsumowanie i mapa dalszych związków}
\label{sec:pde18-summary}
Z \eqref{eq:pde18-minmax}, \eqref{eq:pde18-weyl} i
\eqref{eq:pde18-liyau} otrzymujemy trzy wzajemnie
uzupełniające się opisy: porządek geometryczny wartości
własnych, asymptotyczną gęstość wysokich częstotliwości
i globalne, ilościowe dolne ograniczenie sumy wartości
własnych. Wzór~\eqref{eq:pde18-heat} pokazuje, że to samo
widmo kontroluje krótkoczasowe zachowanie półgrupy cieplnej.
W następnym rozdziale przejdziemy do pojemności
Sobolewa, osobliwości usuwalnych i własności subtelnych
zbiorów brzegowych.
% Rozdział 19: włączany przed \appendix do PDE_tom_I_wyklad.tex.

\chapter[Pojemność Sobolewa i osobliwości]{Pojemność Sobolewa, osobliwości usuwalne i regularność brzegu}
\label{chap:pde19}

\section{Dlaczego miara Lebesgue'a nie wystarcza?}
\label{sec:pde19-motivation}
W teorii słabych rozwiązań PDE rozróżniamy zbiory pomijalne w sensie
miary Lebesgue'a od zbiorów pomijalnych dla energii Dirichleta.
Przykładowo pojedynczy punkt w $\R^2$ ma miarę zero, a funkcja
$\log|x|$ jest harmoniczna poza tym punktem, gdyż
dla $r=|x|>0$ mamy $\Delta\log r=-r^{-2}+r^{-2}=0$.
Jej gradient nie jest jednak całkowalny z kwadratem
w otoczeniu zera:
\[
 \int_{B_R(0)}|\nabla\log|x||^2\,\mathrm d\lambda_2(x)
 =2\pi\int_0^R\frac{1}{r}\,\mathrm dr=+\infty.
\] Tymczasem funkcja
harmoniczna, która \emph{należy do $H^1$ także w pobliżu osobliwości},
może nie odczuwać usunięcia punktu. Sformułujemy tę zasadę przy
pomocy pojemności.

Zawsze, gdy piszemy $\Omega\subset\R^d$ jako \emph{obszar},
zakładamy, że jest on otwarty i spójny. W definicjach lokalnych
pojemności nie wymagamy spójności zbioru otwartego.
Rozpatrujemy funkcje rzeczywiste; argumenty dotyczące
równań z rzeczywistymi współczynnikami dla funkcji zespolonych
prowadzi się osobno dla części rzeczywistej i urojonej.
Przyjmujemy $|E|=\lambda_d(E)$ dla zbiorów mierzalnych,
$\omega_d=|B_1(0)|$ i $\|\cdot\|_p=\|\cdot\|_{L^p(\Omega)}$,
o ile zbiór, na którym liczymy normę, nie został wskazany inaczej.

\begin{uwaga}[Równość prawie wszędzie a quasi-ciągłość]
\label{rem:pde-v43-capacity-ae}
Element $H^1(\Omega)$ jest klasą funkcji równych prawie
wszędzie; w tym sensie nie ma zadanego punktowego śladu
na dowolnym zbiorze miary zero. Pojemność oparta na energii
pozwala rozróżniać zbiory, na których reprezentant
quasi-ciągły może mieć znaczenie. Nie wynika z tego, że
każdy zbiór miary Lebesgue'a zero ma zerową pojemność.
Relację z operatorem śladu na brzegu Lipschitzowskim
omówiono w rozdziale~\ref{chap:pde-embeddings}; są to
różne konstrukcje, wymagające odmiennych założeń.
\end{uwaga}

\section{Definicja pojemności i jej podstawowe własności}
\label{sec:pde19-capacity}
\begin{definicja}[Względna pojemność Dirichleta]
\label{def:pde19-cap}
Niech $\Omega\subset\R^d$ będzie \emph{ograniczonym zbiorem otwartym},
a $K\Subset\Omega$ zbiorem zwartym. Oznaczmy przez $\mathcal A_\Omega(K)$ zbiór tych $v\in H_0^1(\Omega)$,
dla których istnieje otwarty $U$ spełniający $K\subset U\Subset\Omega$
oraz $v\ge1$ prawie wszędzie na $U$. Definiujemy
\begin{equation}\label{eq:pde19-cap}
 \operatorname{cap}_{2,\Omega}(K)=
 \inf_{v\in\mathcal A_\Omega(K)}
 \int_\Omega|\nabla v|^2\,\mathrm d\lambda_d(x).
\end{equation}
Pojemność zbioru pustego wynosi zero. Warunek $v\ge1$ oznacza,
że rozważamy rzeczywistą nierówność prawie wszędzie, a nie
niezdefiniowane wartości punktowe funkcji Sobolewa.
\end{definicja}

\begin{uwaga}[Pojemność i norma $H^1$]
Z nierówności Poincarégo na kostce i przedłużenia
przez zero (jak w sekcji~\ref{sec:pde18-intro}) wynika
$\|v\|_2\le C_\Omega\|\nabla v\|_2$ dla $v\in H_0^1(\Omega)$,
także bez regularności brzegu. Zatem seminorma
$\|\nabla v\|_2$ jest normą równoważną normie $H^1$
na $H_0^1(\Omega)$. Zatem zerowanie się pojemności z
\eqref{eq:pde19-cap} jest równoważne zerowaniu się pojemności
zdefiniowanej przez pełną normę $H^1$. Wartości liczbowe tych dwóch
pojemności są jednak różne. Wzory na pojemności kul poniżej dotyczą
\emph{energii gradientu}, a nie pełnej normy $H^1$.
\end{uwaga}

\begin{lemat}[Próby gładkie i obcinanie]
\label{lem:pde19-smooth}
W definicji~\ref{def:pde19-cap} można zastąpić $H_0^1(\Omega)$
przestrzenią $C_c^\infty(\Omega)$ bez zmiany wartości infimum.
Można ponadto wymagać $0\le v\le1$ prawie wszędzie, używając
wówczas funkcji Sobolewa.
\end{lemat}
\begin{proof}
Niech $v\in\mathcal A_\Omega(K)$. Obcięcie
$T(v)=\min\{1,\max\{v,0\}\}$ spełnia $T(v)=1$ prawie
wszędzie na pewnym otoczeniu $K$. Z wniosku
~\ref{cor:pde9-truncation} (stosowanego do części dodatniej
i obcięcia na poziomie $1$) wynika
$|\nabla T(v)|\le|\nabla v|$ prawie wszędzie.
Uzasadnijmy również $T(v)\in H_0^1(\Omega)$.
Dla $v_n\in C_c^\infty(\Omega)$ zbieżnych do $v$ w $H^1$
funkcje $T(v_n)$ należą do $H^1$ i mają zwarty nośnik
w $\Omega$, więc po przedłużeniu przez zero można je
wygładzić splotem o dostatecznie małym nośniku.
Należą zatem do $H_0^1(\Omega)$.
Ciąg $(T(v_n))$ jest ograniczony w $H^1$ i zbiega
do $T(v)$ w $L^2$, ponieważ $T$ jest lipschitzowska
ze stałą $1$. Słabo zbieżny podciąg w przestrzeni Hilberta
$H_0^1(\Omega)$ ma więc granicę $T(v)$.
W ten sposób obcięcie pozostaje dopuszczalne
i nie zwiększa energii.

Możemy teraz zakładać, że $0\le v\le1$ i $v=1$
prawie wszędzie na otwartym $U$ takim, że
$K\subset U\Subset\Omega$. Wybieramy
$\chi\in C_c^\infty(U)$ równą $1$ na otoczeniu $K$,
oraz $v_n\in C_c^\infty(\Omega)$ zbieżne do $v$ w $H^1$.
Kładziemy $w_n=\chi+(1-\chi)v_n$. Są to funkcje gładkie
o zwartym nośniku, równe $1$ na otoczeniu $K$.
Ponieważ $\chi(1-v)=0$ prawie wszędzie, mamy
$w_n-v=(1-\chi)(v_n-v)$ w $H^1(\Omega)$.
Reguła Leibniza i ograniczoność $\chi,\nabla\chi$
dają $w_n\to v$ w $H^1$.
W szczególności $\|\nabla w_n\|_2^2\to\|\nabla v\|_2^2$.
Infimum po funkcjach gładkich jest nie mniejsze niż
infimum po funkcjach Sobolewa, a z powyższej konstrukcji
jest nie większe od energii dowolnej dopuszczalnej funkcji
Sobolewa. Oba infima są więc równe.
\end{proof}

\begin{twierdzenie}[Monotoniczność i (skończona) podaddytywność]
\label{thm:pde19-basic}
Niech $\Omega\subset\R^d$ będzie ograniczonym zbiorem otwartym,
a $K,L\Subset\Omega$ zbiorami zwartymi. Wówczas:
\begin{enumerate}[label=\textup{(\roman*)}]
\item $0\le\operatorname{cap}_{2,\Omega}(K)<\infty$;
\item $K\subset L$ implikuje
$\operatorname{cap}_{2,\Omega}(K)\le\operatorname{cap}_{2,\Omega}(L)$;
\item $\operatorname{cap}_{2,\Omega}(K\cup L)\le
\operatorname{cap}_{2,\Omega}(K)+\operatorname{cap}_{2,\Omega}(L)$;
\item jeśli $K\Subset\Omega_1\subset\Omega_2$ i oba zbiory są
ograniczone i otwarte, to
$\operatorname{cap}_{2,\Omega_2}(K)\le
\operatorname{cap}_{2,\Omega_1}(K)$.
\end{enumerate}
Ponadto $\operatorname{cap}_{2,\Omega}(K)=0$ implikuje
$\lambda_d(K)=0$.
\end{twierdzenie}
\begin{proof}
Dla (i) wybieramy funkcję $\eta\in C_c^\infty(\Omega)$ równą
$1$ w otoczeniu $K$. Monotoniczność (ii) wynika z zawierania
zbiorów funkcji dopuszczalnych.

Niech $v,w\in H_0^1(\Omega)$ będą nieujemnymi funkcjami
dopuszczalnymi odpowiednio dla $K$ i $L$, z energiami nie większymi
od odpowiednich infimów plus $\eps$. Funkcja
$z=\max\{v,w\}=w+(v-w)^+$ należy do $H_0^1(\Omega)$:
dla części dodatniej stosujemy ten sam argument aproksymacyjny
co w lemacie~\ref{lem:pde19-smooth}. Funkcja $z$ jest
dopuszczalna dla $K\cup L$, bo jest co najmniej równa $1$
prawie wszędzie na sumie obu otoczeń dopuszczalności.
Z reguł obcięcia i zanikania gradientu na zbiorze
$\{v=w\}$ dla funkcji $v-w$ wynika, że prawie wszędzie
\[
 |\nabla\max(v,w)|^2+|\nabla\min(v,w)|^2
 =|\nabla v|^2+|\nabla w|^2,
\]
energia $z$ nie przewyższa sumy obu energii. Granica
$\eps\downarrow0$ dowodzi (iii). W (iv) przedłużamy funkcję
$H_0^1(\Omega_1)$ przez zero do $\Omega_2$, co nie zmienia
energii ani warunku dopuszczalności.

Jeżeli $v$ jest dopuszczalna dla $K$, to
$\lambda_d(K)\le\int_\Omega|v|^2\,\mathrm d\lambda_d$.
Nierówność Poincarégo (obowiązująca dla $H_0^1$ na
ograniczonym zbiorze otwartym) daje
$\lambda_d(K)\le C_\Omega\int_\Omega|\nabla v|^2$.
Biorąc infimum, otrzymujemy ostatnią tezę.
\end{proof}

\begin{uwaga}[Miara zero nie wystarcza]
Implikacji z ostatniego twierdzenia nie można odwrócić.
W wymiarze jeden pojedynczy punkt ma dodatnią pojemność,
choć oczywiście ma miarę zero. W większych wymiarach
podobnie mogą zachowywać się odpowiednio duże fragmenty
hiperpłaszczyzn (widać to z twierdzenia o śladzie).
\end{uwaga}

\section{Dokładna pojemność kondensatora kulistego}
\label{sec:pde19-ball}
\begin{twierdzenie}[Pojemność współśrodkowych kul]
\label{thm:pde19-balls}
Niech $d\ge1$, $0<r<R$ i $K=\overline{B_r(0)}$.
Dla względnej pojemności w $B_R(0)$ zachodzi
\begin{equation}\label{eq:pde19-ball-exact}
 \operatorname{cap}_{2,B_R}(K)=
 \begin{cases}
 \displaystyle\frac{2}{R-r},&d=1,\\[5pt]
 \displaystyle\frac{2\pi}{\log(R/r)},&d=2,\\[6pt]
 \displaystyle\frac{(d-2)|S^{d-1}|}{r^{2-d}-R^{2-d}},&d\ge3,
 \end{cases}
\end{equation}
gdzie $|S^{d-1}|=d\omega_d$ jest polem sfery jednostkowej.
\end{twierdzenie}
\begin{proof}
Dzięki lematowi~\ref{lem:pde19-smooth} dla dolnego
oszacowania wystarczy rozpatrzyć dopuszczalne funkcje
$v\in C_c^\infty(B_R)$. Przedłużamy je przez zero poza
kulę. Dla każdego $\vartheta\in S^{d-1}$ funkcja
$s\mapsto v(s\vartheta)$ jest gładka na $[r,R]$,
$v(r\vartheta)\ge1$ (nierówność prawie wszędzie na
otoczeniu przechodzi na punktową przez ciągłość)
i $v(R\vartheta)=0$. Z twierdzenia Newtona--Leibniza
i nierówności Cauchy'ego--Schwarza otrzymujemy
\[
 1\le\int_r^R|\partial_sv(s\vartheta)|\,\mathrm ds
 \le\left(\int_r^R|\partial_sv|^2s^{d-1}\,\mathrm ds\right)^{1/2}
 \left(\int_r^R s^{1-d}\,\mathrm ds\right)^{1/2}.
\]
Integrujemy po sferze i wykorzystujemy
$|\nabla v|\ge|\partial_sv|$. Energia całkowita jest
więc nie mniejsza od
$|S^{d-1}|/I_{r,R}$, gdzie
$I_{r,R}=\int_r^R s^{1-d}\,\mathrm ds$.

Aby uzyskać odwrotne oszacowanie, definiujemy $v=1$
na $\overline{B_r}$, $v=0$ poza $B_R$ i dla $r<s<R$
przyjmujemy funkcję radialną
\[
 v(s)=\frac{\displaystyle\int_s^R t^{1-d}\,\mathrm dt}
 {\displaystyle\int_r^R t^{1-d}\,\mathrm dt}.
\]
Jest to funkcja lipschitzowska: jej profil radialny
jest ciągły i kawałkami gładki, a pochodna
$-s^{1-d}/I_{r,R}$ na $[r,R]$ jest ograniczona.
Ma zerowy ślad na brzegu kuli, więc należy do $H_0^1(B_R)$
z twierdzenia~\ref{thm:pde10-trace-zero}.
Z rachunku radialnego otrzymujemy
\[
 \int_{B_R}|\nabla v|^2\,\mathrm d\lambda_d
 =\frac{|S^{d-1}|}{I_{r,R}^2}\int_r^R s^{1-d}\,\mathrm ds
 =\frac{|S^{d-1}|}{I_{r,R}}.
\]
Funkcja ta nie jest dopuszczalna w pierwotnej definicji,
gdyż nie jest równa co najmniej $1$ na otwartym otoczeniu
$\overline{B_r}$. Dla $0<\eps<R-r$ stosujemy jednak
ten sam wzór z $r+\eps$ zamiast $r$. Otrzymujemy
funkcję dopuszczalną o energii $|S^{d-1}|/I_{r+\eps,R}$.
Przejście $\eps\downarrow0$ daje odwrotne oszacowanie.
Co więcej, infimum w tej definicji nie jest osiągane:
otoczenie dopuszczalności każdej funkcji zawiera
$\overline{B_{r+\eps}}$ dla pewnego $\eps>0$, więc jej
energia jest co najmniej $|S^{d-1}|/I_{r+\eps,R}$,
co jest ściśle większe od $|S^{d-1}|/I_{r,R}$.
Dla $d=1$ rachunek wykonujemy na dwóch odcinkach,
a $|S^0|=2$. Mamy kolejno
$I_{r,R}=R-r$ dla $d=1$, $I_{r,R}=\log(R/r)$ dla $d=2$
oraz $I_{r,R}=(r^{2-d}-R^{2-d})/(d-2)$ dla $d\ge3$.
Daje to~\eqref{eq:pde19-ball-exact}.
\end{proof}

\begin{wniosek}[Pojemność punktu]
\label{cor:pde19-point}
Jeżeli $\Omega\subset\R^d$ jest ograniczonym zbiorem otwartym,
a $x_0\in\Omega$, to
\begin{equation}\label{eq:pde19-point}
 \operatorname{cap}_{2,\Omega}(\{x_0\})=0\quad(d\ge2),
 \qquad
 \operatorname{cap}_{2,\Omega}(\{x_0\})>0\quad(d=1).
\end{equation}
\end{wniosek}
\begin{proof}
Wybierzmy $B_R(x_0)\Subset\Omega$ i $0<r<R$.
Wzór~\eqref{eq:pde19-ball-exact} obowiązuje również dla
kul o środku $x_0$, ponieważ translacja zachowuje
dziedzinę formową, dopuszczalność i energię gradientu.
Z monotoniczności
względem zbioru i obszaru otrzymujemy
$\operatorname{cap}_{2,\Omega}(\{x_0\})\le
\operatorname{cap}_{2,B_R}(\overline{B_r(x_0)})$.
Prawa strona dąży do zera przy $r\downarrow0$ dla $d\ge2$;
w $d=2$ tempo jest logarytmiczne, a w $d\ge3$ rzędu
$r^{d-2}$. W jednym wymiarze nierówność
$|v(x_0)|\le\sqrt{|\Omega|}\,\|v'\|_{L^2(\Omega)}$
wynika następująco: dla składowej $(a,b)$ zawierającej
$x_0$ ograniczenie $v$ należy do $H_0^1(a,b)$ przez
ograniczenie ciągu aproksymujących funkcji testowych.
Jego absolutnie ciągły reprezentant spełnia $v(a)=0$,
więc $v(x_0)=\int_a^{x_0}v'(s)\,\mathrm d\lambda_1(s)$.
Stosujemy nierówność Cauchy'ego--Schwarza i $b-a\le|\Omega|$.
Dopuszczalność i ciągłość dają $v(x_0)\ge1$, skąd
$\operatorname{cap}_{2,\Omega}(\{x_0\})\ge|\Omega|^{-1}>0$.
\end{proof}

\section[Pojemność Sobolewa i wymiar krytyczny]{Pojemność $W^{1,p}$ i wymiar krytyczny}
\label{sec:pde19-p}
\begin{definicja}[Pojemność $p$-energetyczna]
Niech $1\le p<\infty$, a $K\Subset\Omega$ będzie zbiorem
zwartym, gdzie $\Omega$ jest ograniczonym zbiorem otwartym. Definiujemy
\[
 \operatorname{cap}_{p,\Omega}(K)=\inf\left\{
 \int_\Omega|\nabla v|^p\,\mathrm d\lambda_d(x):
 v\in W_0^{1,p}(\Omega),\ v\ge1\ \text{p.w. na otoczeniu }K\right\}.
\]
Otoczenie w warunku dopuszczalności można zawsze
wybrać o zwartym domknięciu w $\Omega$, dzięki zwartości
$K$. Przyjmujemy pojemność zbioru pustego równą zero.
Dla $p=2$ jest to poprzednia definicja.
\end{definicja}
\begin{twierdzenie}[Punkt a wykładnik Sobolewa]
\label{thm:pde19-point-p}
Niech $d\ge2$, $1\le p<\infty$ i $x_0$ należy do ograniczonego
zbioru otwartego $\Omega\subset\R^d$. Wtedy
$\operatorname{cap}_{p,\Omega}(\{x_0\})=0$ dla $p\le d$,
natomiast $\operatorname{cap}_{p,\Omega}(\{x_0\})>0$
dla $p>d$. W przypadku $d=1$ pojemność punktu jest dodatnia
dla każdego $p\ge1$.
\end{twierdzenie}
\begin{proof}
Po translacji przyjmijmy $x_0=0$, wybierając
$B_R\Subset\Omega$. Dla $p<d$ bierzemy
$0<r<R/2$ i próbę $v_r(x)=\eta(x/r)$, gdzie
$\eta\in C_c^\infty(B_2)$ i $\eta=1$ na $B_1$.
Mamy $v_r\in C_c^\infty(\Omega)$ i $v_r=1$ na $B_r$,
a więc funkcja jest dopuszczalna dla punktu.
Ponieważ $\nabla v_r(x)=r^{-1}\nabla\eta(x/r)$,
podstawienie $x=ry$ daje
\[
 \int_\Omega|\nabla v_r|^p\,\mathrm d\lambda_d(x)
 =r^{d-p}\int_{\R^d}|\nabla\eta(y)|^p\,\mathrm d\lambda_d(y)
 \longrightarrow0.
\]

Dla $p=d>1$ wybieramy $0<r<R$ i określamy
\[
 v_r(x)=
 \begin{cases}
 1,& |x|\le r,\\
 \log(R/|x|)/\log(R/r),&r<|x|<R,\\
 0,&|x|\ge R.
 \end{cases}
\]
Profil jest ciągły i lipschitzowski dla ustalonych $r,R$.
Funkcja ma zwarty nośnik w $\Omega$, należy do
$W^{1,d}(\R^d)$ i przez wygładzanie splotem należy
do $W_0^{1,d}(\Omega)$. Jest dopuszczalna dla punktu.
Na pierścieniu $|\nabla v_r|=1/(|x|\log(R/r))$,
poza nim gradient znika prawie wszędzie.
Jej energia $d$-ta wynosi
\[
 |S^{d-1}|\,\frac{\displaystyle\int_r^R s^{-1}\,\mathrm ds}
 {\bigl(\log(R/r)\bigr)^d}
 =\frac{|S^{d-1}|}{(\log(R/r))^{d-1}}\longrightarrow0.
\]
Dla $p>d$ przedłużamy $v\in W_0^{1,p}(\Omega)$
przez zero do $\R^d$, korzystając z twierdzenia
~\ref{thm:pde9-zero-extension}. Wybieramy ustalony
punkt $y\notin\overline\Omega$. Nierówność Morreya
~\ref{thm:pde10-Morrey} dla ciągłego reprezentanta
przedłużenia $E_0v$ daje
\[
 |v(x_0)|=|E_0v(x_0)-E_0v(y)|
 \le C_{d,p}|x_0-y|^{1-d/p}\|\nabla v\|_{L^p(\Omega)}.
\]
Równość $E_0v(y)=0$ wynika z ciągłości i zerowania się
prawie wszędzie na otoczeniu $y$.
Dopuszczalność oraz ciągłość dają $v(x_0)\ge1$.
Podnosząc oszacowanie do potęgi $p$, otrzymujemy
dolną granicę energii niezależną od $v$, ściśle dodatnią.

W $d=1$ dla każdego $1\le p<\infty$ stosujemy
absolutną ciągłość na składowej $(a,b)$, jak w dowodzie
wniosku~\ref{cor:pde19-point}. Dla $p>1$ nierówność
Höldera daje
$1\le |v(x_0)|\le (b-a)^{1-1/p}\|v'\|_{L^p(a,b)}$.
Dla $p=1$ stosujemy bezpośrednio
$1\le\int_a^b|v'|\,\mathrm d\lambda_1$.
W obu przypadkach energia jest co najmniej
$(b-a)^{1-p}>0$.
\end{proof}

\begin{uwaga}[Wymiar, miara i pojemność]
Wynik nie orzeka, że każdy zbiór wymiaru Hausdorffa $d-p$
ma zerową pojemność: przypadki graniczne są subtelne.
Zależności z miarą Hausdorffa wymagają osobnych twierdzeń,
a dla $p=d$ pojawiają się poprawki logarytmiczne.
\end{uwaga}

\section{Usuwalność osobliwości dla słabych rozwiązań}
\label{sec:pde19-removable}
\begin{twierdzenie}[Zbiór zerowej pojemności jest niewidoczny dla rozwiązań energetycznych]
\label{thm:pde19-remove}
Niech $\Omega\subset\R^d$ będzie ograniczonym zbiorem otwartym,
a $K\Subset\Omega$ zbiorem zwartym takim, że
$\operatorname{cap}_{2,\Omega}(K)=0$.
Niech $A\in L^\infty(\Omega;\R^{d\times d})$;
nie zakładamy tu ani symetrii, ani eliptyczności.
Jeżeli $u\in H^1(\Omega\setminus K)$ oraz
\begin{equation}\label{eq:pde19-off-K}
 \int_{\Omega\setminus K}A(x)\nabla u\cdot\nabla\varphi
 \,\mathrm d\lambda_d(x)=0
 \quad\text{dla każdego }\varphi\in C_c^\infty(\Omega\setminus K),
\end{equation}
to $u$ ma jedyne (z dokładnością do równości prawie
wszędzie) przedłużenie do $H^1(\Omega)$, a przedłużenie
jest słabym rozwiązaniem $-\operatorname{div}(A\nabla u)=0$
na całym $\Omega$.
\end{twierdzenie}
\begin{proof}
Z lematu~\ref{lem:pde19-smooth} i zerowej pojemności
istnieją $\eta_j\in C_c^\infty(\Omega)$, $\eta_j\ge1$
na otoczeniu $K$, takie że $\|\nabla\eta_j\|_2\to0$.
Wybieramy gładką funkcję $\Theta\colon\R\to[0,1]$,
równą $0$ na $(-\infty,0]$ i $1$ na $[1,\infty)$,
z ograniczoną pochodną (np. całkę z nieujemnej
funkcji testowej na $(0,1)$ o całce $1$). Wtedy
$\psi_j=\Theta(\eta_j)\in C_c^\infty(\Omega)$,
$\psi_j=1$ w otoczeniu $K$ i
$\|\psi_j\|_{H^1(\Omega)}\to0$: ostatni wniosek wynika
z nierówności Poincarégo i oszacowań
$|\Theta(s)|\le C|s|$, $|\Theta'(s)|\le C$.

Z twierdzenia~\ref{thm:pde19-basic} mamy $\lambda_d(K)=0$.
Przedłużamy $u$ oraz każdą jego słabą pochodną
$g_\ell=\partial_\ell u$ przez zero na $K$ jako funkcje
$L^2(\Omega)$. Aby wykazać, że przedłużone $g_\ell$
są słabymi pochodnymi przedłużonego $u$, dla
$\varphi\in C_c^\infty(\Omega)$ stosujemy równanie
definiujące słabą pochodną na $\Omega\setminus K$
do testu $(1-\psi_j)\varphi$. Po rozwinięciu otrzymujemy
\[
 \int_\Omega u\,\partial_\ell\varphi\,\mathrm d\lambda_d
 +\int_\Omega g_\ell\varphi\,\mathrm d\lambda_d
 =
 \int_\Omega (u\psi_j\partial_\ell\varphi+
 u\varphi\partial_\ell\psi_j+g_\ell\psi_j\varphi)
 \,\mathrm d\lambda_d.
\]
Moduł prawej strony jest nie większy od
\[
 (\|u\|_2\|\partial_\ell\varphi\|_\infty+
 \|g_\ell\|_2\|\varphi\|_\infty)\|\psi_j\|_2+
 \|u\|_2\|\varphi\|_\infty\|\partial_\ell\psi_j\|_2,
\]
co dąży do zera. Zatem $u\in H^1(\Omega)$ po
przedłużeniu i jego gradient zgadza się prawie wszędzie
z gradientem na $\Omega\setminus K$.
Jednoznaczność wynika z $\lambda_d(K)=0$.

Weźmy teraz dowolną $\varphi\in C_c^\infty(\Omega)$.
Funkcja $\varphi_j=(1-\psi_j)\varphi$ jest gładka,
ma zwarty nośnik w $\Omega\setminus K$, więc z
\eqref{eq:pde19-off-K} mamy
\[
 0=\int_\Omega A\nabla u\cdot\nabla\varphi_j
 \mathrm d\lambda_d
 =\int_\Omega A\nabla u\cdot\nabla\varphi
 \mathrm d\lambda_d
 -\int_\Omega A\nabla u\cdot
 (\psi_j\nabla\varphi+\varphi\nabla\psi_j)\mathrm d\lambda_d.
\]
Dwa ostatnie składniki dążą do zera dzięki
nierówności Cauchy'ego--Schwarza, ograniczoności
$A,\varphi,\nabla\varphi$ i zbieżności
$\psi_j\to0$ w $H^1$. Otrzymujemy tożsamość słabą
na całym $\Omega$.
\end{proof}

\begin{wniosek}[Równość przestrzeni Dirichleta po usunięciu zbioru]
\label{cor:pde19-H0}
Niech $\Omega$ będzie ograniczonym zbiorem otwartym,
a $K\Subset\Omega$ zbiorem zwartym o
$\operatorname{cap}_{2,\Omega}(K)=0$. Wtedy zachodzi
\[
 H_0^1(\Omega\setminus K)=H_0^1(\Omega)
\]
po naturalnym utożsamieniu funkcji przez przedłużenie zerem
na $K$ (który ma miarę zero).
\end{wniosek}
\begin{proof}
Z twierdzenia~\ref{thm:pde19-basic} zbiór $K$ ma miarę
Lebesgue'a zero. Przedłużenie przez zero funkcji
testowych z $\Omega\setminus K$ do $\Omega$ zachowuje
normę $H^1$. Przez gęstość daje ono izometryczne
zanurzenie $H_0^1(\Omega\setminus K)\to H_0^1(\Omega)$,
którego obraz jest domknięty dzięki zupełności.
Dla odwrotnego zawierania wystarczy więc wykazać,
że obraz zawiera $C_c^\infty(\Omega)$.
Jeżeli $\varphi\in C_c^\infty(\Omega)$ i $\psi_j$ są jak
w poprzednim dowodzie, to
$\varphi_j=(1-\psi_j)\varphi\in C_c^\infty(\Omega\setminus K)$
i $\varphi_j\to\varphi$ w $H^1(\Omega)$, ponieważ
$\|\varphi\psi_j\|_2\le\|\varphi\|_\infty\|\psi_j\|_2$
oraz
$\|\nabla(\varphi\psi_j)\|_2\le
\|\nabla\varphi\|_\infty\|\psi_j\|_2+
\|\varphi\|_\infty\|\nabla\psi_j\|_2$.
Z domkniętości obrazu i gęstości $C_c^\infty(\Omega)$
wynika żądana równość przestrzeni.
\end{proof}

\begin{kontrprzyklad}[Zerowa pojemność nie wystarcza bez kontroli energii]
\label{ex:pde19-fundamental}
Dla $d\ge3$ funkcja $u(x)=|x|^{2-d}$, a dla $d=2$
$u(x)=\log|x|$, jest harmoniczna na $B_1\setminus\{0\}$,
ale nie należy do $H^1(B_1\setminus\{0\})$ ani nie
ma przedłużenia do $H^1(B_1)$. Rzeczywiście,
$|\nabla u(x)|\asymp|x|^{1-d}$ (dla $d\ge2$), więc
\[
 \int_{B_\eps\setminus\{0\}}|\nabla u|^2
 \,\mathrm d\lambda_d(x)
 \asymp\int_0^\eps r^{1-d}\,\mathrm dr=+\infty.
\]
Funkcja $u$ jest lokalnie całkowalna, gdyż dla $d\ge3$
całka radialna modułu w pobliżu zera sprowadza się
do $\int_0^\eps r\,\mathrm dr$, a dla $d=2$
do $\int_0^\eps r|\log r|\,\mathrm dr$.
W sensie dystrybucyjnym, przy normalizacji rozwiązania
fundamentalnego z~\eqref{eq:pde7-laplace-fundamental}, mamy
\[
 -\Delta u=
 \begin{cases}
 (d-2)|S^{d-1}|\delta_0,&d\ge3,\\
 -2\pi\delta_0,&d=2.
 \end{cases}
\]
W szczególności przedłużenie lokalnie całkowalne
nie jest harmoniczne w sensie dystrybucji. Zerowa pojemność punktu nie
usuwa więc osobliwości przy nieograniczonej energii.
\end{kontrprzyklad}

\section{Funkcje harmoniczne ograniczone na obszarze z przekłuciem}
\label{sec:pde19-puncture}
\begin{twierdzenie}[Usuwalność ograniczonej osobliwości izolowanej]
\label{thm:pde19-bounded}
Niech $d\ge2$, $B_R(x_0)\subset\R^d$ będzie kulą,
a $u\in C^2(B_R(x_0)\setminus\{x_0\})$ funkcją harmoniczną.
Jeżeli $u$ jest ograniczona w otoczeniu $x_0$, to istnieje
jedyna funkcja harmoniczna $\widetilde u\in C^\infty(B_R(x_0))$
równa $u$ poza $x_0$.
\end{twierdzenie}
\begin{proof}
Po translacji przyjmijmy $x_0=0$. Wybierzmy
$\chi\in C_c^\infty(B_R)$, $0\le\chi\le1$, równą $1$
na $B_{R/2}$. Ustalmy promienie $R_j\downarrow0$,
$R_j<R/2$. Z wniosku~\ref{cor:pde19-point} pojemność
punktu w każdej kuli $B_{R_j}$ wynosi zero. Z lematu
~\ref{lem:pde19-smooth} oraz gładkiego obcięcia
z dowodu twierdzenia~\ref{thm:pde19-remove} wybieramy
$\psi_j\in C_c^\infty(B_{R_j})$ równe $1$ na
otoczeniu zera, $0\le\psi_j\le1$, tak aby
$\|\nabla\psi_j\|_{L^2(\R^d)}\le1/j$.
Po przedłużeniu przez zero mamy $\psi_j(x)\to0$
dla każdego $x\ne0$.

Funkcja $u$ jest ograniczona na
$\operatorname{supp}\chi\setminus\{0\}$:
w pobliżu zera z założenia, a poza mniejszą kulą
z ciągłości na zbiorze zwartym. Niech $M_0$ będzie
kresem $|u|$ na tym zbiorze. Kładziemy
$\zeta_j=\chi(1-\psi_j)$. Funkcja $\zeta_j^2u$
ma zwarty nośnik poza zerem i należy do $H^1$;
przez lokalne wygładzanie jest dopuszczalnym testem
w równaniu harmonicznym. Otrzymujemy
\[
 \int\zeta_j^2|\nabla u|^2\,\mathrm d\lambda_d
 =-2\int \zeta_j u\,\nabla u\cdot\nabla\zeta_j\,
 \mathrm d\lambda_d.
\]
Całki liczymy na $B_R\setminus\{0\}$.
Z nierówności Cauchy'ego--Schwarza wynika
$\int\zeta_j^2|\nabla u|^2\,\mathrm d\lambda_d
\le4\int |u|^2|\nabla\zeta_j|^2\,\mathrm d\lambda_d$
(po podzieleniu przez pierwiastek z lewej strony,
gdy jest dodatni; przypadek zerowy jest natychmiastowy).
Ponieważ $\nabla\zeta_j=(1-\psi_j)\nabla\chi-\chi\nabla\psi_j$,
otrzymujemy
\[
 \int_{B_R}\chi^2(1-\psi_j)^2|\nabla u|^2
 \,\mathrm d\lambda_d
 \le 8 M_0^2\int_{B_R}
 (|\nabla\chi|^2+\chi^2|\nabla\psi_j|^2)
 \,\mathrm d\lambda_d,
\]
Prawa strona jest jednostajnie ograniczona względem $j$.
Lemat Fatou i $\psi_j\to0$ poza zerem dają
$\nabla u\in L^2(B_{R/2}\setminus\{0\})$.
Ponadto $u\in L^2(B_{R/2}\setminus\{0\})$ z ograniczoności.
Ponieważ $u$ jest klasycznie gładka poza zerem (wystarcza
tu $C^2$), jej klasyczne pochodne są tam słabymi
pochodnymi. Zatem $u\in H^1(B_{R/2}\setminus\{0\})$.
Twierdzenie~\ref{thm:pde19-remove}, zastosowane z $A=I$,
daje przedłużenie do $H^1(B_{R/2})$ harmoniczne
dystrybucyjnie. Lemat Weyla~\ref{thm:pde7-weyl}
daje jego gładkiego harmonicznego reprezentanta.
Na przekłutej kuli $B_{R/2}\setminus\{0\}$ zgadza się on
z pierwotną funkcją wszędzie, bo obie są ciągłe i równe
prawie wszędzie. Możemy więc skleić go z $u$ na
$B_R\setminus\{0\}$, otrzymując funkcję harmoniczną
na $B_R$. Jednoznaczność wynika z ciągłości
i gęstości $B_R\setminus\{0\}$ w $B_R$.
\end{proof}

\begin{przyklad}[Punkt brzegowy, w którym dane Dirichleta giną]
\label{ex:pde19-irregular}
Niech $d\ge2$ i $\Omega=B_1(0)\setminus\{0\}$.
Na brzegu $\partial\Omega=S^{d-1}\cup\{0\}$ zadajmy
funkcję ciągłą $g$, równą $0$ na $S^{d-1}$ oraz $1$ w zerze.
Nie istnieje ograniczona funkcja harmoniczna na $\Omega$,
ciągła na $\overline\Omega$ i przyjmująca wartości $g$.
W istocie taka funkcja na mocy poprzedniego twierdzenia
przedłużyłaby się harmonicznie na $B_1$. Z zasady maksimum~\ref{thm:pde14-weak-max},
stosowanej do przedłużenia oraz do jego przeciwności
przy zerowych danych na $S^{d-1}$, byłaby wszędzie zerowa.
Przeczy to wartości $1$ narzuconej w zerze.
To przykład \emph{nieregularnego punktu brzegu}.
\end{przyklad}

\section{Pojemność, reprezentanci i zbiory wyjątkowe}
\label{sec:pde19-quasi}
\begin{definicja}[Pojemność zbioru otwartego i quasi-ciągłość]
Niech $\Omega$ będzie ograniczonym zbiorem otwartym.
Dla otwartego $U\subset\Omega$ definiujemy
\[
 \operatorname{cap}_{2,\Omega}(U)=\inf\left\{
 \|\nabla v\|_2^2:v\in H_0^1(\Omega),\
 v\ge1\ \text{p.w. na }U\right\}.
\]
Nie wymagamy $U\Subset\Omega$ i przyjmujemy
$\inf\varnothing=+\infty$.
Dla dowolnego $E\subset\Omega$ kładziemy
\[
 \operatorname{cap}_{2,\Omega}(E)
 =\inf\{\operatorname{cap}_{2,\Omega}(U):E\subset U\subset\Omega,
 \ U\text{ otwarty}\}.
\]
Rozszerzenie jest zgodne z definicją dla zbiorów
otwartych dzięki monotoniczności warunku dopuszczalności,
a dla zbiorów zwartych $K\Subset\Omega$ dzięki możliwości
zmniejszenia otwartego otoczenia do otoczenia
o zwartym domknięciu w $\Omega$.
Funkcję $\widetilde u:\Omega\to\R$ nazywamy
\emph{quasi-ciągłą}, jeśli dla każdego $\eps>0$
istnieje otwarty $G\subset\Omega$,
$\operatorname{cap}_{2,\Omega}(G)<\eps$, taki że
$\widetilde u|_{\Omega\setminus G}$ jest ciągła.
Własność zachodzącą poza zbiorem pojemności zero nazywamy
własnością zachodzącą \emph{quasi-wszędzie}.
\end{definicja}

\begin{uwaga}
Rozszerzona pojemność jest monotoniczna. Dla otwartego
$U$ nierówność Poincarégo daje
$\lambda_d(U)\le C_\Omega^2\operatorname{cap}_{2,\Omega}(U)$
(przez oszacowanie każdej próby, również gdy infimum
wynosi $+\infty$). Dla dowolnego $E\subset\Omega$ stąd
$\lambda_d^*(E)\le C_\Omega^2\operatorname{cap}_{2,\Omega}(E)$,
gdzie $\lambda_d^*$ jest miarą zewnętrzną Lebesgue'a.
W szczególności zbiory pojemności zero mają miarę zero.
\end{uwaga}

\begin{lemat}[Podaddytywność przeliczalna pojemności]
\label{lem:pde19-countable}
Dla dowolnych $E_j\subset\Omega$, $j\ge1$, zachodzi
$\operatorname{cap}_{2,\Omega}(\bigcup_jE_j)
 \le\sum_j\operatorname{cap}_{2,\Omega}(E_j)$.
\end{lemat}
\begin{proof}
Gdy suma pojemności po prawej stronie nierówności
jest nieskończona, teza jest natychmiastowa.
Najpierw rozważmy zbiory otwarte $U_j$ o skończonej sumie
pojemności i ustalmy $\eps>0$. Wybieramy $v_j\in H_0^1(\Omega)$, $v_j\ge1$
prawie wszędzie na $U_j$, $0\le v_j\le1$, z
$\int_\Omega|\nabla v_j|^2\,\mathrm d\lambda_d(x)\le\operatorname{cap}_{2,\Omega}(U_j)
 +\eps 2^{-j}$. Dla $m\ge1$ kładziemy
$w_m=\max(v_1,\ldots,v_m)$. Z tożsamości gradientowej
użytej w dowodzie Twierdzenia~\ref{thm:pde19-basic}
wynika $\|\nabla w_m\|_2^2\le\sum_{j=1}^m
\|\nabla v_j\|_2^2$. Normy $H_0^1$ są więc jednostajnie
ograniczone. Ponadto $0\le w_m\le1$ i $w_m$ jest
niemalejący prawie wszędzie. Po odrzuceniu przeliczalnej
sumy zbiorów miary zero wybieramy punktową granicę $w$.
Z twierdzenia o zbieżności zmajoryzowanej zastosowanego
do $|w_m-w|^2\le1$ na ograniczonym zbiorze $\Omega$
wynika $w_m\to w$ silnie w $L^2$.
Z ograniczoności w przestrzeni Hilberta $H_0^1$
wybieramy słabo zbieżny podciąg; jego granica musi być
równa $w$ przez ciągłe zanurzenie do $L^2$.
Słaba dolna półciągłość normy gradientu daje
$w\in H_0^1$ oraz
$\|\nabla w\|_2^2\le\sum_j\|\nabla v_j\|_2^2$.
Ponieważ $w\ge1$ prawie wszędzie na $\bigcup_jU_j$,
uzyskujemy
$\operatorname{cap}_{2,\Omega}(\bigcup_jU_j)
\le\sum_j\operatorname{cap}_{2,\Omega}(U_j)+\eps$.
Przejście $\eps\downarrow0$ dowodzi nierówności
dla zbiorów otwartych.
Z definicji zewnętrznej pojemności wybieramy
$U_j\supset E_j$ z
$\operatorname{cap}_{2,\Omega}(U_j)
\le\operatorname{cap}_{2,\Omega}(E_j)+\eps2^{-j}$.
Stosujemy monotoniczność i już udowodnioną nierówność
dla $\bigcup_jU_j$, otrzymując górną granicę
$\sum_j\operatorname{cap}_{2,\Omega}(E_j)+\eps$.
Przechodzimy z $\eps\downarrow0$.
\end{proof}

\begin{twierdzenie}[Quasi-ciągły reprezentant funkcji Sobolewa]
\label{thm:pde19-quasi}
Niech $\Omega\subset\R^d$ będzie ograniczonym zbiorem otwartym.
Każda rzeczywista funkcja $u\in H_0^1(\Omega)$ ma
reprezentanta quasi-ciągłego $\widetilde u$, równego $u$
prawie wszędzie. Jest on jednoznaczny z dokładnością
do zbioru pojemności zero.
\end{twierdzenie}
\begin{proof}
Wybieramy $u_n\in C_c^\infty(\Omega)$ tak szybko zbieżne
do $u$ w $H_0^1(\Omega)$, aby
$\|u_n-u\|_{H^1}\le2^{-2n-1}$ dla $n\ge1$.
Z nierówności trójkąta wynika wtedy
$\|u_{n+1}-u_n\|_{H^1}^2\le2^{-4n}$.
Niech $G_n=\{x\in\Omega:|u_{n+1}(x)-u_n(x)|>2^{-n}\}$.
Ponieważ $u_{n+1}-u_n\in C_c^\infty(\Omega)$,
funkcja $2^n|u_{n+1}-u_n|$ należy do $H_0^1$ i jest
co najmniej $1$ na $G_n$. Zatem
\[
 \operatorname{cap}_{2,\Omega}(G_n)
 \le2^{2n}\|\nabla(u_{n+1}-u_n)\|_2^2\le2^{-2n}.
\]
Dla $F_m=\bigcup_{n\ge m}G_n$ z
lematu~\ref{lem:pde19-countable} mamy
$\operatorname{cap}_{2,\Omega}(F_m)
\le\sum_{n\ge m}2^{-2n}\to0$.
Na $\Omega\setminus F_m$ szereg
$\sum_{n\ge m}(u_{n+1}-u_n)$ jest jednostajnie i bezwzględnie
zbieżny, ponieważ każdy składnik ma tam moduł nie większy niż
$2^{-n}$. Granica $\widetilde u$ jest na tym zbiorze ciągła
jako jednostajna granica funkcji ciągłych.

Zbiór $F=\bigcap_m F_m$ ma pojemność zero, ponieważ
$F\subset F_m$ dla wszystkich $m$; poza $F$
ciąg $u_n$ zbiega punktowo. Ponieważ $F$ ma miarę zero,
a ze zbieżności w $L^2$ można wybrać podciąg zbieżny
prawie wszędzie do $u$, otrzymujemy $\widetilde u=u$
prawie wszędzie. Wartości na $F$ można zdefiniować
dowolnie. Dla każdego $\eps>0$ wybór dużego $m$ sprawia,
że $\operatorname{cap}_{2,\Omega}(F_m)<\eps$,
a ograniczenie $\widetilde u$ do dopełnienia $F_m$
jest ciągłe.

Dla jednoznaczności niech $v,w$ będą dwoma
quasi-ciągłymi reprezentantami tej samej klasy.
Dla $\eps>0$ wybieramy otwarte zbiory wyjątkowe
dla obu funkcji, każdy o pojemności mniejszej niż
$\eps/2$. Ich suma $G$ ma pojemność mniejszą niż
$\eps$ i obie funkcje są ciągłe na $\Omega\setminus G$.
Niech $D=\{x:v(x)\ne w(x)\}$. Jest to zbiór zawarty
w zbiorze miary zero. Z względnej ciągłości wynika,
że $D\setminus G$ jest otwarty w $\Omega\setminus G$,
czyli istnieje otwarty $V\subset\Omega$ taki, że
$D\setminus G=V\setminus G$. Zatem $D\subset V\cup G$,
a $\lambda_d(V\setminus G)=0$. Każda próba dopuszczalna
dla $G$ jest dopuszczalna także dla otwartego $V\cup G$,
gdyż warunek obowiązuje prawie wszędzie.
Stąd $\operatorname{cap}_{2,\Omega}(V\cup G)
=\operatorname{cap}_{2,\Omega}(G)<\eps$.
Monotoniczność daje $\operatorname{cap}_{2,\Omega}(D)<\eps$,
a dowolność $\eps$ dowodzi jednoznaczności.
\end{proof}

\begin{uwaga}[Po co quasi-ciągłość?]
Funkcję Sobolewa określamy zwykle jedynie prawie wszędzie,
więc zapis $u(x_0)$ nie ma sensu dla dowolnie wybranego
punktu. Twierdzenie~\ref{thm:pde19-quasi} zastępuje
nieistniejącą ciągłość klasyczną ciągłością poza zbiorem
\emph{małym energetycznie}. Jest to podstawa subtelnej
interpretacji warunków brzegowych oraz przeszkód w zagadnieniach
wariacyjnych. Nie utożsamiamy przy tym automatycznie
pojęcia quasi-ciągłości z regularnym śladem na
brzegu Lipschitzowskim.
\end{uwaga}

\section{Bariera zewnętrznej kuli i kryterium Wienera}
\label{sec:pde19-wiener}
\begin{definicja}[Bariera harmoniczna]
Niech $\Omega\subset\R^d$ będzie ograniczonym obszarem,
a $x_0\in\partial\Omega$. Lokalną barierą harmoniczną
w $x_0$ nazywamy nieujemną funkcję harmoniczną $b$
w $\Omega\cap B_\rho(x_0)$, dodatnią w jego wnętrzu,
która spełnia $b(x)\to0$ przy $x\to x_0$ od strony $\Omega$.
Do kryterium regularności używamy bariery spełniającej
dodatkowo warunek oddzielenia od zera poza każdym
otoczeniem $x_0$: dla $0<\delta<\rho$ infimum $b$
na $\Omega\cap B_\rho(x_0)\setminus B_\delta(x_0)$
jest dodatnie, jeśli ten zbiór jest niepusty.
W szczególności nie wystarcza sama dodatniość
we wnętrzu.
\end{definicja}

\begin{stwierdzenie}[Jawna bariera przy zewnętrznej kuli]
\label{prop:pde19-exterior-ball}
Niech $d\ge2$, $\Omega\subset\R^d$ będzie ograniczonym
obszarem, a $x_0\in\partial\Omega$. Załóżmy, że istnieje
kula $B_r(y)\subset\R^d\setminus\overline\Omega$
taka, że $|x_0-y|=r$. Wtedy funkcja
\[
 b(x)=\begin{cases}
 r^{2-d}-|x-y|^{2-d},&d\ge3,\\[2pt]
 \log(|x-y|/r),&d=2
 \end{cases}
\]
jest dodatnia i harmoniczna w $\Omega$, a
$b(x)\to0$ przy $x\to x_0$ w $\Omega$.
\end{stwierdzenie}
\begin{proof}
Ponieważ kula $B_r(y)$ jest zawarta w dopełnieniu
$\overline\Omega$, każdy $x\in\Omega$ spełnia
$|x-y|>r$. Funkcja radialna $s\mapsto s^{2-d}$
(dla $d\ge3$) albo $s\mapsto\log s$ (dla $d=2$)
jest harmoniczna poza środkiem $y$. Wskazane kombinacje
są dodatnie dla $s>r$ i zerują się dla $s=r$.
Ciągłość funkcji radialnej daje granicę zero w $x_0$.
Jeśli kula styka się z brzegiem także w innych punktach,
sama ta funkcja może nie spełniać warunku oddzielenia
od zera. Wybierzmy $0<r'<r$ oraz
$y'=x_0+(r'/r)(y-x_0)$. Z nierówności trójkąta
$\overline{B_{r'}(y')}\subset\overline{B_r(y)}$.
Równość w nierówności trójkąta pokazuje, że sfery
obu kul mają tylko punkt $x_0$ wspólny: dla
$|x-y'|\le r'$ i $|x-y|=r$ konieczne jest
$x-y'$ współliniowe z $y'-y$ i zgodnie skierowane,
o długości $r'$, co daje $x=x_0$.
W konsekwencji mniejsza domknięta kula przecina
$\overline\Omega$ wyłącznie w $x_0$. Wstawiając
$r',y'$ do wzoru na $b$, otrzymujemy funkcję ciągłą
na $\overline\Omega$, dodatnią na
$\overline\Omega\setminus\{x_0\}$ i zerową w $x_0$.
Dla każdego $\delta>0$ jej minimum na niepustym
zwartym zbiorze $\overline\Omega\setminus B_\delta(x_0)$
jest ściśle dodatnie. Jest to bariera z wymaganym
warunkiem oddzielenia.
\end{proof}

\begin{uwaga}[Warunek geometryczny a rozwiązywalność]
Istnienie bariery zewnętrznej kuli prowadzi do regularności
$x_0$ dla zagadnienia Dirichleta w sensie Perrona, po
zastosowaniu zasady porównania. Korzystamy tu z klasycznego twierdzenia o barierach
dla rozwiązania Perrona (wynik przywołany, nieudowodniony
w tym rozdziale). Pełna konstrukcja Perrona
nie jest tutaj dowodzona. Warunek zewnętrznej kuli jest
wystarczający, ale niekonieczny.
\end{uwaga}

\begin{definicja}[Bezwzględna pojemność Newtonowska]
Dla $d\ge3$ i zwartego $K\subset\R^d$ określamy
\[
 \operatorname{Cap}_2(K)=\inf\left\{
 \int_{\R^d}|\nabla\varphi|^2\,\mathrm d\lambda_d(x):
 \varphi\in C_c^\infty(\R^d),\ \varphi\ge1
 \text{ na otoczeniu }K\right\}.
\]
W odróżnieniu od~\eqref{eq:pde19-cap} nie występuje tu
ustalony obszar zewnętrzny ani wyraz $L^2$.
\end{definicja}

\begin{stwierdzenie}[Skalowanie i pojemność bezwzględna kuli]
Dla $d\ge3$, zwartego $K\subset\R^d$, $a>0$ i
$z\in\R^d$ zachodzi
\[
 \operatorname{Cap}_2(z+aK)=a^{d-2}\operatorname{Cap}_2(K),
 \qquad
 \operatorname{Cap}_2(\overline{B_\rho(z)})
 =(d-2)|S^{d-1}|\rho^{d-2}\quad(\rho>0).
\]
Pojemność jest monotoniczna względem zawierania zbiorów.
\end{stwierdzenie}
\begin{proof}
Monotoniczność wynika z odwrotnego zawierania klas
dopuszczalnych. Przyporządkowanie
$\varphi\mapsto\varphi((\cdot-z)/a)$ jest bijekcją
między klasami prób dla $K$ i $z+aK$. Zmiana
zmiennych w energii gradientu daje czynnik $a^{d-2}$.
Dla kuli o środku zero każda próba gładka o zwartym
nośniku należy do $C_c^\infty(B_R)$ dla pewnego
$R>\rho$. Odwrotnie, każda taka próba jest również
próbą w całej przestrzeni po przedłużeniu przez zero.
Z lematu~\ref{lem:pde19-smooth} wynika zatem
\[
 \operatorname{Cap}_2(\overline{B_\rho})
 =\inf_{R>\rho}\operatorname{cap}_{2,B_R}(\overline{B_\rho})
 =\lim_{R\to\infty}
 \frac{(d-2)|S^{d-1}|}{\rho^{2-d}-R^{2-d}}.
\]
Ostatnią równość daje monotoniczność względem obszaru
i wzór~\eqref{eq:pde19-ball-exact}.
Granica wynosi $(d-2)|S^{d-1}|\rho^{d-2}$.
Przesunięcie środka nie zmienia pojemności.
\end{proof}


\begin{lemat}[Porównanie pojemności względnej i bezwzględnej]
\label{lem:pde19-cap-comparison}
Niech $d\ge3$, $\Omega\subset\R^d$ będzie ograniczonym
zbiorem otwartym, a $D\Subset\Omega$ zwartym zbiorem.
Istnieje stała $C=C(d,D,\Omega)$ taka, że dla każdego
zwartego $K\subset D$
\[
 \operatorname{Cap}_2(K)\le\operatorname{cap}_{2,\Omega}(K)
 \le C\operatorname{Cap}_2(K).
\]
W szczególności zerowość tych dwóch pojemności jest
równoważna na ustalonym zwartym podzbiorze wnętrza $\Omega$.
\end{lemat}
\begin{proof}
W pierwszej nierówności stosujemy przedłużenie przez
zero funkcji dopuszczalnych. Dzięki
Lematowi~\ref{lem:pde19-smooth} możemy ograniczyć się
do prób gładkich. Dla drugiej nierówności wybierzmy
$\chi\in C_c^\infty(\Omega)$, $\chi=1$ na otoczeniu $D$.
Jeżeli $\varphi\in C_c^\infty(\R^d)$ jest dopuszczalna
dla $\operatorname{Cap}_2(K)$, to $\chi\varphi$ jest
dopuszczalna dla $\operatorname{cap}_{2,\Omega}(K)$.
Mamy $|\nabla(\chi\varphi)|^2\le
2\chi^2|\nabla\varphi|^2+2|\nabla\chi|^2|\varphi|^2$.
Niech $S=\operatorname{supp}\nabla\chi$ oraz
$2^*=2d/(d-2)$. Z nierówności Höldera
i twierdzenia Sobolewa~\ref{thm:pde10-Sobp} dla $p=2$
otrzymujemy
\[
 \int_S|\varphi|^2\,\mathrm d\lambda_d
 \le |S|^{2/d}\|\varphi\|_{L^{2^*}(\R^d)}^2
 \le C_d^2|S|^{2/d}\|\nabla\varphi\|_{L^2(\R^d)}^2.
\]
Po uwzględnieniu ograniczoności $\chi$ i $\nabla\chi$
wynika
\[
 \int_\Omega|\nabla(\chi\varphi)|^2
 \,\mathrm d\lambda_d(x)
 \le C\int_{\R^d}|\nabla\varphi|^2
 \,\mathrm d\lambda_d(x).
\]
Bierzemy infimum po $\varphi$.
\end{proof}

\begin{twierdzenie}[Kryterium Wienera, $d\ge3$ — wynik przywołany]
\label{thm:pde19-wiener}
Niech $d\ge3$, $\Omega\subset\R^d$ będzie ograniczonym
obszarem i $x_0\in\partial\Omega$. Dla $k$ dostatecznie
dużych niech
\[
 E_k=(\R^d\setminus\Omega)\cap
 \{x:2^{-k-1}\le|x-x_0|\le2^{-k}\}.
\]
Punkt $x_0$ jest regularny dla klasycznego zagadnienia
Dirichleta w sensie Perrona wtedy i tylko wtedy, gdy
\begin{equation}\label{eq:pde19-wiener}
 \sum_{k=k_0}^\infty
 2^{k(d-2)}\operatorname{Cap}_2(E_k)=+\infty.
\end{equation}
Tutaj regularność oznacza, że rozwiązanie Perrona
przyjmuje w granicy od strony $\Omega$ wartość
$f(x_0)$ dla każdej ciągłej funkcji brzegowej
$f\in C(\partial\Omega)$.
\end{twierdzenie}

\begin{uwaga}[Zakres dowodu kryterium]
Kryterium Wienera jest głębokim twierdzeniem teorii
potencjału: \emph{nie dowodzimy tu jego pełnej implikacji
w obu kierunkach}. Z poprzednich twierdzeń wykazaliśmy
natomiast jawnie wzór pojemności kul i nieregularność
izolowanego punktu brzegowego. Dla przekłutej kuli
z~\ref{ex:pde19-irregular} wszystkie dostatecznie małe
$E_k$ są puste, zatem szereg jest zbieżny. Dla dopełnienia
zawierającego stały stożek z wierzchołkiem w $x_0$
pojemności kolejnych fragmentów są porównywalne z
$2^{-k(d-2)}$, co daje rozbieżność. Kryterium w wymiarze
dwa przyjmuje postać logarytmiczną i nie jest tu
formułowane.
\end{uwaga}


\begin{stwierdzenie}[Zewnętrzny stożek a szereg Wienera]
\label{prop:pde19-cone-wiener}
Niech $d\ge3$, $\Omega\subset\R^d$ będzie ograniczonym obszarem
oraz $x_0\in\partial\Omega$. Załóżmy, że istnieją $r_0>0$
i niepusty, otwarty zbiór kierunków $\Sigma\subset S^{d-1}$
takie, że $x_0+r\theta\notin\Omega$ dla każdego
$0<r<r_0$ i $\theta\in\Sigma$.
Wtedy dla zbiorów $E_k$ z Twierdzenia~\ref{thm:pde19-wiener}
istnieje $c>0$ i $k_1$ takie, że
\[
  2^{k(d-2)}\operatorname{Cap}_2(E_k)\ge c
  \qquad(k\ge k_1).
\]
W szczególności szereg Wienera jest rozbieżny, a
$x_0$ jest regularnym punktem brzegowym, jeśli stosujemy
przywołane Twierdzenie~\ref{thm:pde19-wiener}.
\end{stwierdzenie}
\begin{proof}
Wybierzmy $\theta_0\in\Sigma$ i $\delta>0$ tak małe,
aby kierunki z kuli sferycznej wokół $\theta_0$ o promieniu
$2\delta$ należały do $\Sigma$, przy czym $\delta<1/16$.
Dla $h=2^{-k}<r_0$ weźmy
$z_k=x_0+3h\theta_0/4$. Możemy dodatkowo zmniejszyć
$\delta$, tak aby każdy kierunek o odległości euklidesowej
od $\theta_0$ nie większej niż $3\delta$ należał do $\Sigma$.
Jeśli $x\in\overline{B_{\delta h}(z_k)}$, to
$(3/4-\delta)h\le|x-x_0|\le(3/4+\delta)h$,
więc $h/2<|x-x_0|<h$. Dla $w=(x-x_0)/h$ mamy
$|w-(3/4)\theta_0|\le\delta$, a z nierówności trójkąta
\[
 \left|\frac{w}{|w|}-\theta_0\right|
 \le\frac{\bigl||w|-3/4\bigr|+|w-(3/4)\theta_0|}{3/4}
 \le\frac83\delta<3\delta.
\]
Kierunek należy zatem do $\Sigma$, a promień jest
mniejszy niż $r_0$. To dowodzi zawierania
$\overline{B_{\delta h}(z_k)}\subset E_k$.
Monotoniczność pojemności bezwzględnej daje
$\operatorname{Cap}_2(E_k)\ge
\operatorname{Cap}_2(\overline{B_{\delta h}(z_k)})$.
Przesunięcie nie zmienia pojemności, a z obliczenia energii
radialnej na pierścieniu i przejścia z jego promieniem
zewnętrznym do nieskończoności (dla $d\ge3$) mamy
$\operatorname{Cap}_2(\overline{B_\rho})
=(d-2)|S^{d-1}|\rho^{d-2}$.
Stąd dolne ograniczenie jest równe
$(d-2)|S^{d-1}|\delta^{d-2}h^{d-2}$, co po
pomnożeniu przez $h^{-(d-2)}$ dowodzi tezy.
\end{proof}

\begin{przyklad}[Przekłuta kula a kryterium Wienera]
\label{ex:pde19-wiener-puncture}
Dla $d\ge3$ oraz $\Omega=B_1(0)\setminus\{0\}$
w punkcie $x_0=0$ wszystkie dostatecznie małe pierścienie
$E_k$ są puste, gdyż dopełnienie obszaru zawiera
wewnątrz $B_1$ jedynie punkt $0$, niewchodzący do
żadnego z tych pierścieni. Suma z kryterium Wienera wynosi
zero, zatem $x_0$ jest nieregularny. Wynik zgadza się z
Przykładem~\ref{ex:pde19-irregular}.
\end{przyklad}



\begin{lemat}[Usuwalność ograniczonej osobliwości harmonicznej w wymiarze dwa]
\label{lem:pde-v5-removable-bounded-2d}
Niech $R>0$ i niech $h$ będzie rzeczywistą funkcją harmoniczną
na $B_R(0)\setminus\{0\}\subset\R^2$. Jeżeli $h$ jest ograniczona
w otoczeniu zera, to ma jednoznaczne przedłużenie harmoniczne
na całą kulę $B_R(0)$.
\end{lemat}
\begin{proof}
Wybierzmy $0<r_0<R$ i zapiszmy $x=(r\cos\theta,r\sin\theta)$.
Dla $0<r<r_0$ oznaczmy przez
\[
 c_n(r)=\frac1{2\pi}\int_{0}^{2\pi}
 h(r\cos\theta,r\sin\theta)e^{-in\theta}
 \,\mathrm d\lambda_1(\theta),\qquad n\in\mathbb Z,
\]
zespolone współczynniki Fouriera ograniczonej funkcji
$\theta\mapsto h(r\cos\theta,r\sin\theta)$.
Na każdym zwartym podpierścieniu możemy różniczkować
dwukrotnie pod całką, ponieważ pochodne $h$ są ciągłe.
Mnożąc równanie
$h_{rr}+r^{-1}h_r+r^{-2}h_{\theta\theta}=0$
przez $e^{-in\theta}$ i całkując po kącie,
dwukrotnie całkujemy przez części (wyrazy brzegowe
znikają przez okresowość). Dostajemy
$r^2c_n''+rc_n'-n^2c_n=0$. Podstawienie $r=e^t$
sprowadza je do równania o stałych współczynnikach
$\frac{\mathrm d^2}{\mathrm dt^2}c_n(e^t)-n^2c_n(e^t)=0$. Dla $n=0$ zatem
$c_0(r)=A_0+B_0\log r$, a dla $n\ne0$
$c_n(r)=A_nr^{|n|}+B_nr^{-|n|}$.
Ograniczoność $h$ oznacza $|c_n(r)|\le\|h\|_\infty$
dla dostatecznie małych $r$, więc wszystkie współczynniki
$B_n$ muszą znikać. Niech $H$ będzie harmonicznym
przedłużeniem Poissona funkcji $h|_{\partial B_{r_0}}$
do $B_{r_0}(0)$ (sekcja~\ref{sec:pde14-poisson}).
Wzór Poissona na dysku wykorzystuje jądro
$(1-q^2)/(1-2q\cos\alpha+q^2)$, gdzie $q=r/r_0<1$.
Z dwóch szeregów geometrycznych wynika rozwinięcie
tego jądra jako $\sum_{n\in\mathbb Z}q^{|n|}e^{in\alpha}$,
jednostajnie i bezwzględnie zbieżne dla ustalonego $q$.
Całkując szereg przeciwko danym brzegowym, otrzymujemy
współczynniki $(r/r_0)^{|n|}c_n(r_0)$; dla $n=0$
są stałe.
Są zatem identyczne z $c_n(r)$ dla $0<r<r_0$.
Z jednoznaczności szeregów Fouriera dla funkcji ciągłych
na okręgu wynika $H=h$ w przekłutej kuli $B_{r_0}$.
Funkcja otrzymana przez dołączenie $H(0)$ jest więc
harmonicznym przedłużeniem $h$ na całą $B_R$.
Jednoznaczność wynika z ciągłości.
\end{proof}

\begin{stwierdzenie}[Nieregularność izolowanego punktu brzegu w wymiarze dwa]
\label{prop:pde-v5-punctured-disk}
Niech $\Omega=B_1(0)\setminus\{0\}\subset\R^2$.
Punkt $0\in\partial\Omega$ nie jest regularny dla
klasycznego zagadnienia Dirichleta w sensie Perrona.
Dokładniej, dla danych brzegowych
$f(0)=1$ oraz $f|_{\partial B_1}=0$ rozwiązanie
Perrona jest identycznie równe zeru w $\Omega$.
\end{stwierdzenie}
\begin{proof}
Funkcja $f$ jest ciągła na rozłącznej sumie
$\partial\Omega=\{0\}\cup\partial B_1$.
Z zasady porównania ze stałymi wynika, że rozwiązanie
Perrona $u$ jest ograniczone przez $0\le u\le1$;
standardowa konstrukcja Perrona zapewnia jego harmoniczność
wewnątrz $\Omega$.
Z lematu~\ref{lem:pde-v5-removable-bounded-2d} funkcja $u$
przedłuża się harmonicznie przez zero. Każdy punkt
zewnętrznej sfery $\partial B_1$ ma barierę pochodzącą
od kuli zewnętrznej, więc $u$ dąży tam do wartości $0$
(przywołujemy tu standardową własność rozwiązania Perrona
w punktach z barierą).
Zasada maksimum dla funkcji harmonicznych w $B_1$
implikuje $u\equiv0$. Stąd granica $u(x)$ przy $x\to0$
wynosi $0$, a nie $f(0)=1$, co dowodzi nieregularności.
Argument nie wymaga pełnego logarytmicznego kryterium Wienera.
\end{proof}

\begin{przyklad}[Rodzina punktów regularnych i nieregularnych
sterowana szeregiem]
\label{ex:pde19-wiener-balls}
Niech $d\ge3$, $e_1=(1,0,\ldots,0)$, $h_k=2^{-k}$,
$k\ge2$, oraz $0<\delta_k\le1/16$. Określmy
$z_k=\tfrac34h_ke_1$, $K_k=\overline{B_{\delta_kh_k}(z_k)}$
i
\[
 \Omega=B_1(0)\setminus\biggl(\{0\}\cup
 \bigcup_{k=2}^\infty K_k\biggr).
\]
Wówczas $\Omega$ jest ograniczonym, otwartym i spójnym
obszarem, a $0\in\partial\Omega$. Kryterium
Wienera~\ref{thm:pde19-wiener} daje dokładnie
\[
 0\text{ jest regularny}\quad\Longleftrightarrow\quad
 \sum_{k=2}^\infty\delta_k^{d-2}=+\infty.
\]
W szczególności, dla $\delta_k=\frac1{16}k^{-a}$,
$a\ge0$, punkt $0$ jest regularny wtedy i tylko wtedy,
gdy $a(d-2)\le1$.
\end{przyklad}
\begin{proof}
Kule $K_k$ są parami rozłączne. Ich środki leżą na
promieniach $3h_k/4$, a promienie nie przekraczają
$h_k/16$. Dlatego $K_k$ zawiera się w otwartym
pierścieniu $h_k/2<|x|<h_k$. Pozostałe kule nie przecinają
tego pierścienia; zbiór $E_k$ z definicji kryterium Wienera
jest więc równy $K_k$. Usunięty zbiór jest domknięty
w $B_1$, ponieważ jego jedynym punktem skupienia poza
kulami jest $0$; dołączono ten punkt. Ponieważ punkty $-te_1$ dla $0<t<1$ należą do
$\Omega$, punkt $0$ jest rzeczywiście punktem brzegu.
Wykażemy łukową spójność. Dla dowolnych $x,y\in\Omega$
wybieramy $m$ tak duże, by $h_m<\min(|x|,|y|)$.
Kule $K_k$ dla $k\ge m$ leżą wewnątrz $B_{h_m}$,
a kule dla $k<m$ leżą poza $\overline{B_{h_m}}$.
Wszystkie domknięte kule z tej skończonej rodziny
(wraz z $\overline{B_{h_m}}$) są parami rozłączne,
leżą wewnątrz $B_1$ i nie zawierają $x,y$.
Możemy nieznacznie zwiększyć każdy z ich promieni,
zachowując te własności. Odcinek od $x$ do $y$ leży
w $B_1$. Każdy jego fragment przechodzący przez jedną
z powiększonych kul zastępujemy łukiem na jej sferze
(która jest łukowo spójna dla $d\ge2$).
Sfery są rozłączne z pozostałymi powiększonymi kulami,
więc po skończenie wielu zastąpieniach otrzymujemy
drogę w $\Omega$. To dowodzi spójności.

Dla $d\ge3$ bezwzględna pojemność kuli ma wzór
$\operatorname{Cap}_2(K_k)=(d-2)|S^{d-1}|
(\delta_kh_k)^{d-2}$ (Twierdzenie~\ref{thm:pde19-balls}
z promieniem zewnętrznym dążącym do nieskończoności).
Stąd
\[
 2^{k(d-2)}\operatorname{Cap}_2(E_k)
 =(d-2)|S^{d-1}|\delta_k^{d-2}.
\]
Szereg Wienera jest zatem rozbieżny dokładnie wtedy,
gdy rozbieżny jest wskazany szereg liczbowy. Ostatnia
równoważność wynika z kryterium zbieżności szeregu
$\sum_{k\ge2}k^{-a(d-2)}$.
\end{proof}

\begin{uwaga}[Granica samodzielności argumentu]
Geometria kul, obliczenie pojemności i redukcja do szeregu
liczbowego zostały wyprowadzone powyżej. Równoważność
między rozbieżnością szeregu a regularnością punktu brzegu
jest zastosowaniem \emph{przywołanego}, a nie udowodnionego
tutaj kryterium Wienera.
\end{uwaga}

\section{Związki z dalszą regularnością i podsumowanie}
\label{sec:pde19-summary}
Pojemność rozróżnia zbiory, na których w naturalny sposób
może koncentrować się energia rozwiązań, od tych, które
są dla przestrzeni Sobolewa niewidoczne. Wykazaliśmy
trzy konkretne mechanizmy: znikanie pojemności punktów
w wymiarach $d\ge2$, przedłużanie równań dywergencyjnych
przez zbiory pojemności zero dla rozwiązań $H^1$ oraz
istnienie quasi-ciągłych reprezentantów funkcji Sobolewa.
Kryterium Wienera opisuje natomiast regularność warunku
brzegowego za pomocą rozłożenia pojemności na skalach.

W późniejszej teorii jedyności kontynuacji i analizie
mikrolokalnej pojawią się inne sposoby mierzenia wielkości
zbiorów wyjątkowych; nie są one równoważne pojemności
w dowolnych klasach współczynników.


% Rozdzial 20 -- samodzielny fragment do wlaczenia przed \appendix.

\chapter[Paraboliczne równania ewolucyjne]{Paraboliczne równania ewolucyjne: rozwiązania słabe, metoda Galerkina i regularność}
\label{chap:pde20}

\section{Od równania ciepła do współczynników mierzalnych}
W rozdziale~4 rozwiązywaliśmy równanie ciepła w przestrzeni całej
przy użyciu jądra Gaussa. Obecnie rozważymy znacznie ogólniejsze
zagadnienie na ograniczonym obszarze $\Omega\subset\R^d$:
\begin{equation}\label{eq:pde20-eq}
 \partial_t u-\operatorname{div}(A(t,x)\nabla u)=f
 \quad\text{w }(0,T)\times\Omega,\qquad
 u|_{\partial\Omega}=0,\qquad u(0)=u_0.
\end{equation}
Macierz $A$ może być tylko mierzalna zarówno w czasie, jak i w przestrzeni.
W takich warunkach nie należy oczekiwać klasycznych pochodnych drugiego
rzędu. W przeciwieństwie do zagadnienia eliptycznego rozwiązanie zależy
od czasu, ale nie zakładamy istnienia pochodnej $u_t$ o wartościach w $L^2(\Omega)$.

\begin{definicja}[Dane i eliptyczność paraboliczna]\label{def:pde20-data}
Niech $\Omega\subset\R^d$ będzie ograniczonym obszarem Lipschitzowskim,
$0<T<\infty$, $H=L^2(\Omega)$, $V=H_0^1(\Omega)$, $V^*=H^{-1}(\Omega)$.
W całym rozdziale rozpatrujemy rzeczywiste przestrzenie
Hilberta, a $V^*$ oznacza przestrzeń ciągłych
funkcjonałów liniowych na $V$.
Rozpatrujemy rzeczywistą macierz $A=(a_{ij})\in
L^\infty((0,T)\times\Omega;\R^{d\times d})$ taką, że dla pewnych
$0<\lambda\le\Lambda<\infty$, dla prawie wszystkich $(t,x)$ i każdego
$\xi\in\R^d$ zachodzi
\[
 \lambda|\xi|^2\le A(t,x)\xi\cdot\xi,
 \qquad |A(t,x)\xi|\le\Lambda|\xi|.
\]
Nie wymagamy symetrii macierzy $A$. Dla $v,w\in V$ określamy
$a(t;v,w)=\int_\Omega A(t,x)\nabla v(x)\cdot\nabla w(x)
\,\mathrm d\lambda_d(x)$. Przyjmujemy $u_0\in H$ oraz
$f\in L^2(0,T;V^*)$; dualność $V^*,V$ oznaczamy przez
$\langle\cdot,\cdot\rangle$.
\end{definicja}


\begin{center}
\begin{tikzpicture}[scale=0.9,>=Latex]
  \fill[blue!4] (0,0) rectangle (6,3.4);
  \draw[very thick,blue!45!black] (0,0) -- (0,3.4);
  \draw[very thick,blue!45!black] (6,3.4) -- (6,0);
  \draw[dashed,gray] (0,3.4) -- (6,3.4);
  \draw[very thick,teal!65!black] (0,0)--(6,0);
  \draw[->] (-0.4,0)--(6.8,0) node[right] {$x$};
  \draw[->] (0,-0.25)--(0,4.1) node[above] {$t$};
  \node[below] at (3,-0.07) {$u(0,x)=u_0(x)$};
  \node[rotate=90,above] at (0,1.8) {$u(t,0)=0$};
  \node[rotate=90,below] at (6,1.8) {$u(t,1)=0$};
  \node[above] at (3,3.4) {$t=T$ (wartość nie jest zadawana)};
  \node at (3,1.75) {$Q_T=(0,T)\times(0,1)$};
  \node[below] at (0,-0.4) {$0$};
  \node[below] at (6,-0.4) {$1$};
\end{tikzpicture}

\small Rysunek: jednowymiarowy przykład ($\Omega=(0,1)$). Podstawa i brzeg boczny walca parabolicznego.
Dane początkowe i warunek Dirichleta określa się na różnych częściach brzegu.
\end{center}

Normę w $V$ ustalamy jako $\|v\|_V=\|\nabla v\|_{L^2(\Omega)}$.
Dzięki nierówności Poincarégo (por. sekcję~\ref{sec:pde18-intro})
jest ona równoważna zwykłej normie $H^1$. Oznaczamy przez
$C_P$ stałą w nierówności $\|v\|_H\le C_P\|v\|_V$.
Mamy trójkę Gelfanda
\[
 V\hookrightarrow H\cong H^*\hookrightarrow V^*,
 \qquad h\longmapsto[v\mapsto(h,v)_H].
\]
Pierwsze zanurzenie jest ciągłe i ma gęsty obraz,
bo zawiera $C_c^\infty(\Omega)$. Drugie jest iniekcją
dzięki tej gęstości, a nierówność Cauchy'ego--Schwarza
daje $\|h\|_{V^*}\le C_P\|h\|_H$.
Jego obraz jest gęsty w $V^*$: gdyby nie był, istniałby
niezerowy ciągły funkcjonał na $V^*$ zerujący się na $H$.
Z refleksywności przestrzeni Hilberta $V$ taki funkcjonał
miałby postać $g\mapsto\langle g,v\rangle$ dla pewnego
$v\in V$. Warunek $(h,v)_H=0$ dla wszystkich $h\in H$
dawałby jednak $v=0$.

Dla prawie każdego $t$ forma $a(t;\cdot,\cdot)$ spełnia
\[
 |a(t;v,w)|\le\Lambda\|v\|_V\|w\|_V,\qquad
 a(t;v,v)\ge\lambda\|v\|_V^2.
\]
Pierwsze oszacowanie wynika z nierówności
Cauchy'ego--Schwarza i ograniczenia macierzy, drugie
z eliptyczności. Niech $\mathcal A(t):V\to V^*$ będzie
operatorem określonym przez
$\langle\mathcal A(t)v,w\rangle=a(t;v,w)$.
Dla ustalonych $v,w\in V$ mierzalność tej funkcji czasu
wynika z twierdzenia Fubiniego, ponieważ całkowany
iloczyn jest zmajoryzowany przez
$\Lambda|\nabla v||\nabla w|\in L^1(\Omega)$.
Przestrzeń $V$ jest ośrodkową przestrzenią Hilberta
(jako domknięta podprzestrzeń przestrzeni izometrycznie
zanurzonej w $L^2(\Omega)^{d+1}$); również $V^*$
jest ośrodkowa. Mierzalność skalarna $\mathcal A(t)v$
pociąga więc silną mierzalność, na przykład przez
rozwinięcie w przeliczalnej bazie ortonormalnej $V^*$.
Aproksymacja $u(t)$ funkcjami prostymi i jednolite
ograniczenie $\|\mathcal A(t)\|\le\Lambda$ pokazują,
że $\mathcal A(\cdot)u(\cdot)$ jest silnie mierzalna
dla $u\in L^2(0,T;V)$ oraz
$\|\mathcal A(\cdot)u(\cdot)\|_{L^2(0,T;V^*)}
\le\Lambda\|u\|_{L^2(0,T;V)}$. Wszystkie przestrzenie
$L^2(0,T;X)$ rozumiemy w sensie Bochnera.

\begin{uwaga}[Trzy odrębne narzędzia: forma, baza i zwartość]
\label{rem:pde-v43-parabolic-basis}
W dowodzie istnienia wykorzystamy koercywność formy
(znaną z rozdziału~\ref{chap:pde-variational}), bazę własną
pomocniczego $-\Delta_D$ z rozdziału~\ref{chap:pde18}
oraz zwarte zanurzenie z rozdziału~\ref{chap:pde-embeddings}.
Baza laplasjanu nie diagonalizuje na ogół macierzy
$A(t,x)$; służy do zbudowania stabilnych rzutów
$P_n$ w $V$, $H$ i $V^*$. Dlatego ani niesymetryczność,
ani mierzalna zależność $A$ od czasu nie podważają
przedstawionej konstrukcji. Własności samosprzężonej
półgrupy $e^{-tL}$ w dalszej części rozdziału
odnoszą się natomiast do wyraźnie wskazanego przypadku
$L=-\Delta_D$ (lub innego stałego operatora dodatniego
samosprzężonego), nie do ogólnego $A(t,x)$.
\end{uwaga}

\begin{definicja}[Rozwiązanie słabe w przestrzeni energii]
\label{def:pde20-weak}
Rozwiązaniem słabym zagadnienia~\eqref{eq:pde20-eq} nazywamy funkcję
\[
 u\in L^2(0,T;V),\qquad u'\in L^2(0,T;V^*),
\]
której pochodną rozumiemy dystrybucyjnie w czasie:
dla każdej $\eta\in C_c^\infty(0,T)$ zachodzi tożsamość
w $V^*$
\[
 \int_0^T u(t)\eta'(t)\,\mathrm d\lambda_1(t)
 =-\int_0^T u'(t)\eta(t)\,\mathrm d\lambda_1(t).
\]
Tutaj $u$ po lewej stronie utożsamiamy z jego obrazem
w $V^*$. Ponadto rozwiązanie spełnia
$\langle u'(t),v\rangle+a(t;u(t),v)=\langle f(t),v\rangle$
dla każdego $v\in V$ i prawie każdego $t\in(0,T)$, przy czym zbiór
wyjątkowy można wybrać niezależnie od $v$.
Równoważnie mamy $u'(t)+\mathcal A(t)u(t)=f(t)$
w $V^*$ dla prawie każdego $t$. Gdyby równanie
skalarne sformułowano z osobnym zbiorem wyjątkowym
dla każdego $v$, wystarczyłoby wziąć przeliczalny
gęsty podzbiór $V$, zsumować jego zbiory wyjątkowe
i zastosować ciągłość obu stron względem $v$.
Warunek $u(0)=u_0$ dotyczy ciągłego reprezentanta $H$-wartościowego,
którego istnienie wykażemy poniżej.
\end{definicja}

\section{Ciągłość czasowa i tożsamość energii}
\begin{lemat}[Tożsamość energii w trójce Gelfanda]\label{lem:pde20-energy}
Niech $V=H_0^1(\Omega)$, $H=L^2(\Omega)$, a
$u\in L^2(0,T;V)$ ma pochodną $u'\in L^2(0,T;V^*)$.
Wówczas $u$ ma jednoznacznego reprezentanta należącego do
$C([0,T];H)$ i dla wszystkich $0\le s\le t\le T$ zachodzi
\begin{equation}\label{eq:pde20-energy-identity}
 \frac12\|u(t)\|_H^2-\frac12\|u(s)\|_H^2
 =\int_s^t\langle u'(r),u(r)\rangle\,\mathrm dr.
\end{equation}
\end{lemat}
\begin{proof}
Niech $\{e_j\}_{j\ge1}$ będzie ortonormalną bazą $H$ złożoną
z funkcji własnych dodatniego laplasjanu Dirichleta,
$-\Delta_D e_j=\mu_j e_j$, $\mu_j>0$.
Z rozwinięcia formy energii w dowodzie twierdzenia
~\ref{thm:pde18-minmax} wynika, że
$\|v\|_V^2=\sum_{j\ge1}\mu_j|(v,e_j)_H|^2$ dla $v\in V$.
Układ $(\mu_j^{-1/2}e_j)$ jest bazą ortonormalną $V$.
Rzut ortogonalny $P_n$ w $H$ na
$\operatorname{span}\{e_1,\ldots,e_n\}$ jest więc
również rzutem ortogonalnym w $V$.
Na $V^*$ definiujemy zgodne rozszerzenie przez
$\langle P_ng,v\rangle=\langle g,P_nv\rangle$,
czyli $P_ng=\sum_{j=1}^n\langle g,e_j\rangle e_j$
przy utożsamieniu elementów $H$ z funkcjonałami na $V$.
Z twierdzenia Riesza i tożsamości Parsevala w $V$
wynika
\[
 \|g\|_{V^*}^2=\sum_{j\ge1}\mu_j^{-1}
 |\langle g,e_j\rangle|^2.
\]
Zatem $P_n$ jest kontrakcją we wszystkich trzech
przestrzeniach i zbiega silnie do identyczności:
kwadraty norm różnic są ogonami odpowiednich szeregów.
Twierdzenie o zbieżności zmajoryzowanej zastosowane
do kwadratów norm różnic daje
$P_nu\to u$ w $L^2(0,T;V)$ oraz
$P_nu'\to u'$ w $L^2(0,T;V^*)$.

Współczynniki $t\mapsto (u(t),e_j)_H$ należą do $H^1(0,T)$,
ponieważ testując definicję pochodnej czasowej
funkcją $\eta$ i funkcjonałem ewaluacji na $e_j$,
otrzymujemy ich słabe pochodne
$\langle u'(t),e_j\rangle\in L^2(0,T)$.
Wybieramy ich absolutnie ciągłe reprezentanty
(twierdzenie~\ref{thm:pde9-ac}); wtedy $P_nu$
jest ciągłe na $[0,T]$ i jego pochodna to $P_nu'$.
Dla skończonego $n$ zwykła reguła różniczkowania daje
\[
 \|P_nu(t)\|_H^2-\|P_nu(s)\|_H^2
 =2\int_s^t\langle P_nu'(r),P_nu(r)\rangle\,\mathrm dr.
\]
Wykażemy najpierw zbieżność $P_nu$ w $C([0,T];H)$.
Dla $w=P_nu-P_mu$ wybieramy punkt $s\in[0,T]$ taki, że
$\|w(s)\|_H^2\le T^{-1}\|w\|_{L^2(0,T;H)}^2$;
istnieje on, gdyż średnia nie może być mniejsza od wszystkich wartości.
Stosując poprzednią tożsamość na przedziale o końcach $s,t$
i nierówność Cauchy'ego--Schwarza, otrzymujemy
\begin{equation}\label{eq:pde20-cont-est}
 \sup_{0\le t\le T}\|w(t)\|_H^2
 \le T^{-1}\|w\|_{L^2(0,T;H)}^2
 +2\|w'\|_{L^2(0,T;V^*)}\|w\|_{L^2(0,T;V)}.
\end{equation}
Oba składniki po prawej dążą do zera dla $m,n\to\infty$.
Granica jest więc $H$-wartościowo ciągła i równa $u$
prawie wszędzie. Przechodząc do granicy w skończenie wymiarowej
tożsamości energii, korzystamy z jednostajnej zbieżności
wartości końcowych oraz oszacowania
\[
 \begin{split}
 \int_0^T|\langle P_nu',P_nu\rangle-\langle u',u\rangle|
 \,\mathrm d\lambda_1
 &\le \|P_nu'-u'\|_{L^2(0,T;V^*)}\|P_nu\|_{L^2(0,T;V)}\\
 &\quad+\|u'\|_{L^2(0,T;V^*)}\|P_nu-u\|_{L^2(0,T;V)}.
 \end{split}
\]
Prawa strona dąży do zera. To daje~\eqref{eq:pde20-energy-identity}.
Jednoznaczność reprezentanta wynika z ciągłości w $H$:
dwa ciągłe reprezentanty równe prawie wszędzie muszą
być równe w każdym punkcie $[0,T]$.
W szczególności funkcja $t\mapsto\frac12\|u(t)\|_H^2$
jest absolutnie ciągła, a jej pochodną prawie wszędzie
jest $\langle u'(t),u(t)\rangle$, która jest całkowalna
z nierówności Cauchy'ego--Schwarza.
\end{proof}

\section{Istnienie i jednoznaczność: pełna konstrukcja Galerkina}
\begin{twierdzenie}[Paraboliczne zagadnienie Dirichleta]\label{thm:pde20-existence}
Niech $\Omega,T,A,f,u_0$ spełniają założenia
definicji~\ref{def:pde20-data}. Wtedy istnieje dokładnie jedno rozwiązanie
słabe zagadnienia~\eqref{eq:pde20-eq} w sensie
definicji~\ref{def:pde20-weak}. Spełnia ono oszacowanie
\begin{equation}\label{eq:pde20-apriori}
 \sup_{0\le t\le T}\|u(t)\|_H^2
 +\lambda\int_0^T\|u(t)\|_V^2\,\mathrm dt
 \le 2\|u_0\|_H^2
 +2\lambda^{-1}\int_0^T\|f(t)\|_{V^*}^2\,\mathrm dt.
\end{equation}
Stała w nierówności nie zależy od gładkości współczynników.
\end{twierdzenie}
\begin{proof}
\emph{Krok 1: równanie skończenie wymiarowe.}
Niech $V_n=\operatorname{span}\{e_1,\ldots,e_n\}$, gdzie $e_j$ są
funkcjami własnymi użytymi w lemacie~\ref{lem:pde20-energy}.
Szukamy $u_n(t)=\sum_{j=1}^n c_j(t)e_j$, które spełnia
\[
 (u_n'(t),e_i)_H+a(t;u_n(t),e_i)
 =\langle f(t),e_i\rangle,\quad i=1,\ldots,n,
 \qquad u_n(0)=P_nu_0.
\]
Jest to liniowy układ $c'(t)+B_n(t)c(t)=F_n(t)$,
gdzie $(B_n(t))_{ij}=a(t;e_j,e_i)$ oraz
$(F_n(t))_i=\langle f(t),e_i\rangle$.
Mierzalność i oszacowania formy z sekcji pierwszej dają
$|(B_n(t))_{ij}|\le\Lambda\sqrt{\mu_i\mu_j}$ oraz
$|(F_n(t))_i|\le\sqrt{\mu_i}\|f(t)\|_{V^*}$.
W konsekwencji $M_n=\operatorname*{ess\,sup}_{t\in(0,T)}
\|B_n(t)\|_{\mathcal L(\R^n)}<\infty$ i
$F_n\in L^2(0,T;\R^n)$.
Na przedziale $[a,b]\subset[0,T]$ o długości
$(b-a)M_n<1$ i dla ustalonego $c_a$ definiujemy
\[
 (\mathcal Tc)(t)=c_a+\int_a^t
 (F_n(r)-B_n(r)c(r))\,\mathrm d\lambda_1(r).
\]
Jest to odwzorowanie zupełnej przestrzeni
$C([a,b];\R^n)$ w siebie: całkowana funkcja należy
do $L^1$, a całka nieoznaczona jest absolutnie ciągła.
Ponadto $\|\mathcal Tc-\mathcal T\widetilde c\|_\infty
\le(b-a)M_n\|c-\widetilde c\|_\infty$.
Twierdzenie Banacha o punkcie stałym daje dokładnie
jedno rozwiązanie równania całkowego.
Dzieląc $[0,T]$ na skończenie wiele takich przedziałów,
za każdym razem bierzemy wartość końcową poprzedniego
rozwiązania jako dane początkowe następnego.
Otrzymujemy jednoznaczne absolutnie ciągłe $c$
na $[0,T]$, z $c_j(0)=(u_0,e_j)_H$.
Ponieważ $c$ jest ograniczona na tym przedziale,
$c'=F_n-B_nc\in L^2(0,T;\R^n)$, czyli $c\in H^1(0,T;\R^n)$.

\emph{Krok 2: energia.}
Mnożymy $i$-te równanie przez $c_i(t)$ i sumujemy po $i$.
Otrzymujemy dla prawie każdego $t$
\[
 \frac12\frac{\mathrm d}{\mathrm dt}\|u_n\|_H^2
 +\lambda\|u_n\|_V^2
 \le\|f\|_{V^*}\|u_n\|_V
 \le\frac{\lambda}{2}\|u_n\|_V^2
 +\frac1{2\lambda}\|f\|_{V^*}^2.
\]
Całkując od $0$ do $t$ i korzystając z
$\|P_nu_0\|_H\le\|u_0\|_H$, dostajemy
\begin{equation}\label{eq:pde20-energy-n}
 \|u_n(t)\|_H^2+\lambda\int_0^t\|u_n(r)\|_V^2\,\mathrm dr
 \le\|u_0\|_H^2+\lambda^{-1}\|f\|_{L^2(0,T;V^*)}^2.
\end{equation}
To jednocześnie zapobiega utracie kontroli nad rozwiązaniami przy $n\to\infty$.

\emph{Krok 3: kontrola pochodnej czasowej.}
Dla każdego $v\in V$ zachodzi
$\langle u_n',v\rangle=\langle f,P_nv\rangle-a(t;u_n,P_nv)$,
ponieważ $u_n'\in V_n$ i $P_n$ jest ortogonalny w $H$.
Korzystamy z $\|P_nv\|_V\le\|v\|_V$ oraz ograniczoności formy:
\[
 \|u_n'(t)\|_{V^*}\le\|f(t)\|_{V^*}+\Lambda\|u_n(t)\|_V.
\]
Zatem $\{u_n'\}$ jest ograniczony w $L^2(0,T;V^*)$.

\emph{Krok 4: granica.}
Po wybraniu podciągu, dzięki refleksywności przestrzeni Hilberta,
$u_n\rightharpoonup u$ w $L^2(0,T;V)$,
$u_n'\rightharpoonup g$ w $L^2(0,T;V^*)$.
Ponadto $u_n$ jest ograniczony w $L^\infty(0,T;H)$.
Dla dowolnego $v\in V$ i $\eta\in C_c^\infty(0,T)$ mamy
\[
 \int_0^T\langle u_n'(t),v\rangle\eta(t)\,\mathrm dt
 =-\int_0^T(u_n(t),v)_H\eta'(t)\,\mathrm dt.
\]
Obie całki są wartościami ciągłych funkcjonałów
liniowych odpowiednio na $L^2(0,T;V^*)$ i $L^2(0,T;V)$.
Przechodząc do granicy dla każdego $v$ i $\eta$,
stwierdzamy, że $g=u'$ w sensie dystrybucyjnym
o wartościach w $V^*$.
Ustalmy $v=e_j$. Dla $n\ge j$ równanie Galerkina
mnożymy przez $\eta\in L^2(0,T)$ i całkujemy po czasie.
Funkcjonał
\[
 w\longmapsto\int_0^T a(t;w(t),e_j)\eta(t)
 \,\mathrm d\lambda_1(t)
\]
jest dobrze określony dzięki mierzalności formy
i ma normę nie większą niż
$\Lambda\|e_j\|_V\|\eta\|_{L^2(0,T)}$.
Słaba zbieżność daje więc
\[
 \int_0^T
 (\langle u',e_j\rangle+a(t;u,e_j)-\langle f,e_j\rangle)
 \eta\,\mathrm d\lambda_1=0.
\]
Funkcja w nawiasie należy do $L^2(0,T)$; wybór jej
samej jako $\eta$ dowodzi, że znika prawie wszędzie.
Po odrzuceniu przeliczalnej sumy zbiorów wyjątkowych
równość zachodzi jednocześnie dla wszystkich $e_j$.
Przez liniowość, gęstość ich rozpiętości w $V$
i ciągłość funkcjonałów otrzymujemy równanie w $V^*$
dla prawie każdego $t$.

\emph{Krok 5: warunek początkowy.}
Dla $v\in V_n$ mnożymy równanie Galerkina przez
dowolną $\eta\in C^1([0,T])$ spełniającą $\eta(T)=0$,
całkujemy od $0$ do $T$ i całkujemy pierwszy składnik
przez części. Otrzymujemy
\[
 -\int_0^T(u_n(t),v)_H\eta'(t)\,\mathrm dt
 +\int_0^T a(t;u_n(t),v)\eta(t)\,\mathrm dt
 =\int_0^T\langle f(t),v\rangle\eta(t)\,\mathrm dt
 +(P_nu_0,v)_H\eta(0).
\]
Dla każdego ustalonego $v=e_j$ przechodzimy do granicy jak wyżej.
Z lematu~\ref{lem:pde20-energy} funkcja
$t\mapsto(u(t),e_j)_H$ ma absolutnie ciągły reprezentant
o pochodnej $\langle u'(t),e_j\rangle$.
Całkowanie równania granicznego przez części
dla tego współczynnika daje tę samą tożsamość
z wyrazem $(u(0),e_j)_H\eta(0)$ zamiast
$(u_0,e_j)_H\eta(0)$.
Odejmując obie tożsamości i biorąc np.
$\eta(t)=1-t/T$, otrzymujemy
$(u(0),e_j)_H=(u_0,e_j)_H$ dla każdego $j$.
Zatem $u(0)=u_0$ w $H$.

\emph{Krok 6: oszacowanie i jednoznaczność.}
Wstawiamy $v=u(t)$ do równania słabego (jest to dopuszczalne
w równości w $V^*$ dla prawie każdego $t$) i używamy
lematu~\ref{lem:pde20-energy}. Z koercywności i nierówności
Younga, tak jak w kroku 2, dla każdego $t$ dostajemy
\[
 \|u(t)\|_H^2+\lambda\int_0^t\|u(r)\|_V^2\,\mathrm d\lambda_1(r)
 \le M,\qquad
 M=\|u_0\|_H^2+\lambda^{-1}\|f\|_{L^2(0,T;V^*)}^2.
\]
Pierwszy składnik jest ograniczony przez $M$ jednostajnie
w czasie, a drugi przez $M$ dla $t=T$. Dodanie tych
dwóch oszacowań daje~\eqref{eq:pde20-apriori}.
Z równania w $V^*$ otrzymujemy również
$\|u'\|_{L^2(0,T;V^*)}\le
\|f\|_{L^2(0,T;V^*)}+\Lambda\|u\|_{L^2(0,T;V)}$.
Jeżeli $u,w$ są rozwiązaniami z tymi samymi danymi,
to $z=u-w$ spełnia równanie z zerową prawą stroną
i zerową wartością początkową. Tożsamość energii daje
$\|z(t)\|_H^2+2\lambda\int_0^t\|z\|_V^2\,\mathrm dr\le0$.
Stąd $z=0$.
\end{proof}

\begin{uwaga}[Zmienne w czasie współczynniki i brak symetrii]
W dowodzie nie różniczkowaliśmy $A$ względem czasu
ani nie odwoływaliśmy się do symetrii $a(t;\cdot,\cdot)$.
Jest to istotna przewaga słabego sformułowania nad próbami
obliczania $\Delta u$ punkt po punkcie.
\end{uwaga}

\section{Zależność od danych i paraboliczna zasada maksimum}
\begin{wniosek}[Stabilność rozwiązania]\label{cor:pde20-stability}
Jeżeli $u_i$ odpowiadają danym $(u_{0,i},f_i)$, $i=1,2$,
przy tej samej macierzy $A$, to dla $z=u_1-u_2$ mamy
\[
 \sup_{t\in[0,T]}\|z(t)\|_H^2
 +\lambda\|z\|_{L^2(0,T;V)}^2
 \le2\|u_{0,1}-u_{0,2}\|_H^2
 +2\lambda^{-1}\|f_1-f_2\|_{L^2(0,T;V^*)}^2.
\]
\end{wniosek}
\begin{proof}
Równanie na $z$ ma tę samą formę i odpowiadającą jej prawą
stronę $f_1-f_2$. Stosujemy twierdzenie~\ref{thm:pde20-existence}.
\end{proof}

\begin{lemat}[Reguła obcięcia w czasie]\label{lem:pde20-trunc}
Jeżeli $u\in L^2(0,T;V)$ i $u'\in L^2(0,T;V^*)$, to
$u^-:=\max\{-u,0\}\in L^2(0,T;V)$ i funkcja
$t\mapsto\|u^-(t)\|_H^2$ jest absolutnie ciągła, przy czym
\[
 \frac12\frac{\mathrm d}{\mathrm dt}\|u^-(t)\|_H^2
 =-\langle u'(t),u^-(t)\rangle
 \quad\text{dla prawie wszystkich }t.
\]
\end{lemat}
\begin{proof}
Niech $N(v)=v^-$. Z reguł obcięcia
~\ref{cor:pde9-truncation} oraz zanikania gradientu
na zbiorze poziomicowym~\ref{lem:pde9-levelsets}
mamy $\nabla N(v)=-\chi_{\{v<0\}}\nabla v$
prawie wszędzie. Dla $v\in V$ obcięcie należy do $V$:
przybliżamy $v$ funkcjami z $C_c^\infty(\Omega)$,
bierzemy ich części ujemne (funkcje Sobolewa o zwartym
nośniku, należące do $V$ przez wygładzanie splotem).
Są ograniczone w $V$ i zbiegają do $N(v)$ w $H$,
ponieważ $N$ jest lipschitzowskie ze stałą $1$.
Słabo zbieżny podciąg w $V$ ma więc granicę $N(v)$.
Ponadto $\|N(v)\|_V\le\|v\|_V$.

Jeśli $v_n\to v$ w $L^2(0,T;V)$, to
$N(v_n)\to N(v)$ w tej samej przestrzeni. Istotnie,
\[
 \begin{split}
 \nabla N(v_n)-\nabla N(v)
 &=-\chi_{\{v_n<0\}}(\nabla v_n-\nabla v)\\
 &\quad+(\chi_{\{v<0\}}-\chi_{\{v_n<0\}})\nabla v.
 \end{split}
\]
Pierwszy składnik dąży do zera w $L^2$ na
$(0,T)\times\Omega$. Z każdego podciągu wybieramy
dalszy podciąg, dla którego $v_n\to v$ prawie wszędzie.
Poza $\{v=0\}$ funkcje charakterystyczne w drugim
składniku zbiegają, a na $\{v=0\}$ gradient $v$
znika prawie wszędzie (stosujemy lemat o poziomicach
dla prawie każdego czasu i twierdzenie Fubiniego).
Z twierdzenia o zbieżności zmajoryzowanej, z majorantą
$|\nabla v|^2$ dla kwadratu modułu, drugi składnik
dąży do zera w $L^2$ wzdłuż tego dalszego podciągu.
Argument podciągowy daje zbieżność całego ciągu.
Analogicznie dla pojedynczej funkcji czasu otrzymujemy
ciągłość $N:V\to V$, a stąd silną mierzalność $N(u)$
oraz $N(u)\in L^2(0,T;V)$.

Na $H$ określamy funkcjonał
$J(v)=\frac12\|v^-\|_H^2$.
Funkcja skalarna $\Phi(s)=\frac12(s^-)^2$ jest klasy
$C^1$ i $\Phi'(s)=-s^-$ jest lipschitzowska
ze stałą $1$. Z twierdzenia Newtona--Leibniza wynika
$|\Phi(a+b)-\Phi(a)-\Phi'(a)b|\le\frac12|b|^2$.
Całkując, dostajemy
\[
 |J(v+h)-J(v)+(v^-,h)_H|\le\tfrac12\|h\|_H^2.
\]
Zatem $J$ jest różniczkowalny w sensie Frécheta,
$DJ(v)h=-(v^-,h)_H$, a jego różniczka jest ciągła
dzięki lipschitzowskości $N$ w $H$.

Stosujemy teraz rzuty $P_n$ z dowodu lematu
~\ref{lem:pde20-energy}. Funkcje $u_n=P_nu$ są
absolutnie ciągłe o wartościach w skończenie
wymiarowej przestrzeni $V_n$ i $u_n'=P_nu'$.
Reguła łańcuchowa dla funkcji absolutnie ciągłych
i funkcjonału $J$ daje dla $0\le s\le t\le T$
\[
 J(u_n(t))-J(u_n(s))
 =-\int_s^t\langle P_nu'(r),N(u_n(r))\rangle
 \,\mathrm d\lambda_1(r).
\]
Iloczyn skalarny w $H$ jest tu równy dualności,
ponieważ $u_n'\in V_n$ prawie wszędzie.
Z lematu~\ref{lem:pde20-energy} mamy zbieżność
$u_n\to u$ w $C([0,T];H)$ i $L^2(0,T;V)$
oraz $P_nu'\to u'$ w $L^2(0,T;V^*)$.
Wykazana ciągłość obcięcia daje
$N(u_n)\to N(u)$ w $L^2(0,T;V)$.
Stąd przez nierówność Cauchy'ego--Schwarza
całki w dualności zbiegają do całki
$-\int_s^t\langle u',N(u)\rangle\,\mathrm d\lambda_1$,
a wartości końcowe $J(u_n(s)),J(u_n(t))$ zbiegają
dzięki ciągłości $J$ na $H$.
Otrzymujemy tożsamość całkową dla wszystkich $s,t$.
Ponieważ $\langle u',N(u)\rangle\in L^1(0,T)$,
wynika z niej absolutna ciągłość $J(u(\cdot))$
i żądana równość po zróżniczkowaniu prawie wszędzie.
\end{proof}

\begin{twierdzenie}[Zachowanie dodatniości]\label{thm:pde20-positivity}
Niech $u$ będzie rozwiązaniem z twierdzenia~\ref{thm:pde20-existence}.
Jeżeli $u_0\ge0$ prawie wszędzie i $\langle f(t),v\rangle\ge0$
dla każdej nieujemnej $v\in V$ i prawie każdego $t$, to
$u(t,x)\ge0$ dla prawie każdego $x\in\Omega$ i każdego $t\in[0,T]$.
\end{twierdzenie}
\begin{proof}
Stożek nieujemnych elementów $V$ ma przeliczalny
gęsty podzbiór, ponieważ $V$ jest ośrodkową
przestrzenią metryczną. Odrzucając sumę zbiorów
wyjątkowych dla tego podzbioru i korzystając
z ciągłości $f(t)$ na $V$, możemy przyjąć, że
$\langle f(t),v\rangle\ge0$ dla wszystkich
nieujemnych $v\in V$ poza jednym zbiorem miary zero.
Testujemy równanie w $V^*$ przez $u^-(t)$
dla prawie każdego $t$. Z lematu~\ref{lem:pde20-trunc}
$\langle u',u^-\rangle=-\frac12\mathrm d\|u^-\|_H^2/\mathrm dt$.
Na zbiorze $\{u<0\}$ mamy $\nabla u=-\nabla u^-$,
a na dopełnieniu $\nabla u^-=0$ prawie wszędzie. Zatem
$a(t;u,u^-)=-a(t;u^-,u^-)\le-\lambda\|u^-\|_V^2$.
W rezultacie
\[
 \frac12\frac{\mathrm d}{\mathrm dt}\|u^-(t)\|_H^2
 +\lambda\|u^-(t)\|_V^2\le0.
\]
Ponieważ $u^-(0)=0$, lewa strona po scałkowaniu daje
$\|u^-(t)\|_H^2=0$ dla każdego $t$.
\end{proof}

\begin{wniosek}[Zasada porównania]
Dla tej samej macierzy $A$ niech rozwiązania $u,v$ mają
odpowiednio dane $(u_0,f)$ i $(v_0,g)$.
Jeżeli $u_0\le v_0$ prawie wszędzie i
$\langle g(t)-f(t),\phi\rangle\ge0$ dla wszystkich
nieujemnych $\phi\in V$ i prawie wszystkich $t$,
to $u(t)\le v(t)$ prawie wszędzie dla każdego $t$.
\end{wniosek}
\begin{proof}
Stosujemy twierdzenie~\ref{thm:pde20-positivity}
do różnicy $v-u$.
\end{proof}

\section{Wygładzanie i zanik energii dla laplasjanu Dirichleta}
Niech $L=-\Delta_D$ będzie dodatnim laplasjanem Dirichleta
w $H$. Jego wartości własne oznaczamy przez
$0<\mu_1\le\mu_2\le\cdots\to\infty$ (z krotnościami),
a ortonormalne funkcje własne przez $e_j$
(twierdzenie~\ref{thm:pde18-compact}).
Dla $s\ge0$ określamy potęgę spektralną przez
\[
 D(L^s)=\left\{v\in H:\sum_{j\ge1}\mu_j^{2s}|(v,e_j)_H|^2<\infty\right\},
 \qquad L^sv=\sum_{j\ge1}\mu_j^s(v,e_j)_He_j.
\]
W szczególności $D(L^{1/2})=V$ i
$\|L^{1/2}v\|_H=\|v\|_V$ z rozwinięcia formy
w dowodzie twierdzenia~\ref{thm:pde18-minmax}.
Odróżniamy operator $L:D(L)\subset H\to H$
od jego rozszerzenia formowego $\mathcal L:V\to V^*$,
określonego przez
$\langle\mathcal Lv,w\rangle=(\nabla v,\nabla w)_{L^2}$.
Mamy $\|\mathcal Lv\|_{V^*}=\|v\|_V$, co wynika
z twierdzenia Riesza w $V$.
W tej sekcji zakładamy
$A=I$ i $f=0$, chyba że wyraźnie zaznaczymy inaczej.

\begin{twierdzenie}[Półgrupa cieplna i wygładzanie]\label{thm:pde20-semigroup}
Dla każdego $u_0\in H$ funkcja
$u(t)=e^{-tL}u_0=\sum_{j\ge1}e^{-\mu_jt}(u_0,e_j)_He_j$
jest jedynym rozwiązaniem słabym zagadnienia jednorodnego.
Dla każdego $s\ge0$ i $t>0$ zachodzi
\[
 \|L^se^{-tL}u_0\|_H
 \le C_s t^{-s}\|u_0\|_H,
 \qquad C_s=(s/\mathrm e)^s\ (s>0),\quad C_0=1.
\]
Ponadto $\|u(t)\|_H\le e^{-\mu_1t}\|u_0\|_H$.
\end{twierdzenie}
\begin{proof}
Dla skończonych sum własnych obliczamy
$u'(t)+Lu(t)=0$, a tożsamość energii jest bezpośrednia.
Niech $b_j=(u_0,e_j)_H$. Dla ogólnego $u_0$
kwadrat normy ogona szeregu w $H$ nie przewyższa
$\sum_{j>N}|b_j|^2$, jednostajnie dla $t\ge0$.
Szereg określa więc ciągłą funkcję na $[0,\infty)$,
przy czym $u(0)=u_0$.
Mnożenie współczynników przez $e^{-\mu_jt}$ określa
ograniczony operator $S(t)$ na $H$. Tożsamość
$e^{-\mu_j(t+r)}=e^{-\mu_jt}e^{-\mu_jr}$ daje
$S(t+r)=S(t)S(r)$ oraz $S(0)=I$.
Jest to silnie ciągła półgrupa, a zapis $e^{-tL}$
oznacza właśnie tę konstrukcję spektralną.
Z twierdzenia Tonellego dla nieujemnych składników mamy
\[
 \int_0^T\|L^{1/2}e^{-tL}u_0\|_H^2\,\mathrm dt
 =\sum_j\mu_j|(u_0,e_j)_H|^2
 \int_0^T e^{-2\mu_jt}\,\mathrm dt
 \le\tfrac12\|u_0\|_H^2.
\]
To samo oszacowanie dla ogonów pokazuje, że
skończone sumy $u_N$ zbiegają w $L^2(0,T;V)$.
Ich pochodne $u_N'=-\mathcal Lu_N$ zbiegają
w $L^2(0,T;V^*)$ do $-\mathcal Lu$ dzięki izometrii
$\mathcal L$. Przechodząc do granicy w definicji
pochodnej dystrybucyjnej, otrzymujemy
$u'=-\mathcal Lu$. Razem z warunkiem początkowym
dowodzi to, że granica spełnia definicję rozwiązania słabego;
jednoznaczność wynika z twierdzenia~\ref{thm:pde20-existence}.
Dla $s>0$ supremum funkcji $y\mapsto y^se^{-ty}$ na
$[0,\infty)$ wynosi $(s/(\mathrm et))^s$: pochodna
dla $y>0$ ma znak $s-ty$, a funkcja dąży do zera
na obu krańcach. Dla $s=0$ używamy oszacowania
$e^{-ty}\le1$. Zatem
\[
 \sum_j\mu_j^{2s}e^{-2\mu_jt}|b_j|^2
 \le C_s^2t^{-2s}\sum_j|b_j|^2.
\]
Dowodzi to jednocześnie $u(t)\in D(L^s)$ i nierówności
wygładzającej. Ponieważ $\mu_j\ge\mu_1$,
z tożsamości Parsevala otrzymujemy
$\|u(t)\|_H^2\le e^{-2\mu_1t}\|u_0\|_H^2$.
Ponadto dla $0\le r\le t$ dokładne całkowanie
każdego składnika i twierdzenie Tonellego dają
\[
 \|u(t)\|_H^2+2\int_r^t\|u(\sigma)\|_V^2
 \,\mathrm d\lambda_1(\sigma)=\|u(r)\|_H^2.
\]
\end{proof}


\begin{twierdzenie}[Wzór Duhamela dla wymuszenia w $L^2(H)$]
\label{thm:pde20-duhamel}
Niech $\Omega\subset\R^d$ będzie ograniczonym obszarem
Lipschitzowskim, $H=L^2(\Omega)$, $L=-\Delta_D$,
$u_0\in H$, $f\in L^2(0,T;H)$.
Jedynym rozwiązaniem słabym zagadnienia $u'+Lu=f$,
$u(0)=u_0$, jest funkcja
\[
 u(t)=e^{-tL}u_0+\int_0^t e^{-(t-s)L}f(s)\,\mathrm ds,
 \qquad 0\le t\le T,
\]
gdzie całka jest całką Bochnera w $H$.
\end{twierdzenie}
\begin{proof}
Operator $S(t)=e^{-tL}$ jest kontrakcją w $H$.
Dla każdego ustalonego $t\in[0,T]$ funkcja
$s\mapsto S(t-s)f(s)$ na $[0,t]$ jest silnie
mierzalna: dla funkcji prostych $f$ wynika to
z silnej ciągłości $S$, a dla ogólnej $f$
z aproksymacji punktowej prawie wszędzie.
Jej norma jest ograniczona przez $\|f(s)\|_H$,
która należy do $L^1(0,T)$, gdyż
$\|f\|_{L^1(0,T;H)}\le\sqrt T\|f\|_{L^2(0,T;H)}$.
Całka Bochnera istnieje więc dla każdego $t$.
Jeżeli $t_n\to t$, przedłużamy integrandy zerem
na $(t_n,T)$; dla prawie każdego $s\ne t$
zbiegają silnie w $H$ do odpowiedniego integrandu
dla $t$. Twierdzenie o zbieżności zmajoryzowanej
dla całek Bochnera daje ciągłość składnika całkowego.

Niech $u_n$ będzie rozwiązaniem w $V_n$ z danymi
$P_nu_0$ i $f_n=P_nf$. Dla $1\le j\le n$
jego współczynniki spełniają
$r_j'+\mu_jr_j=(f,e_j)_H$,
$r_j(0)=(u_0,e_j)_H$. Mnożymy przez $e^{\mu_jt}$
i całkujemy, uzyskując
$r_j(t)=e^{-\mu_jt}b_j+
\int_0^t e^{-\mu_j(t-s)}(f(s),e_j)_H
\,\mathrm d\lambda_1(s)$.
Sumując skończenie wiele współczynników, dostajemy
wzór Duhamela z danymi $P_nu_0,P_nf$. Dla różnic dwóch skończonych przybliżeń
stosujemy oszacowanie energii z
Twierdzenia~\ref{thm:pde20-existence}:
\[
 \sup_{0\le t\le T}\|u_n(t)-u_m(t)\|_H^2
 +\int_0^T\|u_n-u_m\|_V^2\,\mathrm dt
 \le C\bigl(\|P_nu_0-P_mu_0\|_H^2
 +\|P_nf-P_mf\|_{L^2(0,T;H)}^2\bigr),
\]
przy czym w prawej stronie korzystamy z ciągłego
zanurzenia $H\hookrightarrow V^*$. Prawa strona dąży do
zera. Równanie dla różnic daje również
$\|u_n'-u_m'\|_{L^2(0,T;V^*)}\le
 C_P\|P_nf-P_mf\|_{L^2(0,T;H)}
 +\|u_n-u_m\|_{L^2(0,T;V)}$.
Z zupełności i przejścia granicznego w definicji
pochodnej wynika, że granica spełnia równanie
$u'+\mathcal Lu=f$ w $L^2(0,T;V^*)$.
Warunek początkowy przechodzi dzięki zbieżności
w $C([0,T];H)$. Jednocześnie różnica między wzorem
całkowym dla $u_n$ a wzorem z tezy ma normę w
$C([0,T];H)$ nie większą niż
$\|P_nu_0-u_0\|_H+\sqrt T\|P_nf-f\|_{L^2(0,T;H)}$.
Obie granice są więc identyczne; jednoznaczność
daje twierdzenie~\ref{thm:pde20-existence}.
\end{proof}

\begin{twierdzenie}[Wyższa regularność energii]\label{thm:pde20-max-reg}
Niech $L=-\Delta_D$, $u_0\in V=D(L^{1/2})$,
$f\in L^2(0,T;H)$. Wtedy rozwiązanie słabe równania
$u'+Lu=f$ spełnia
$u\in H^1(0,T;H)\cap L^2(0,T;D(L))$ oraz
\[
 \int_0^T\|u'(t)\|_H^2\,\mathrm dt
 +\sup_{t\in[0,T]}\|L^{1/2}u(t)\|_H^2
 \le2\|L^{1/2}u_0\|_H^2
 +2\int_0^T\|f(t)\|_H^2\,\mathrm dt.
\]
Dla obszarów klasy $C^{1,1}$ można ponadto utożsamić
$D(L)=H^2(\Omega)\cap H_0^1(\Omega)$; bez regularności brzegu
takiego utożsamienia nie zakładamy.
\end{twierdzenie}
\begin{proof}
Dla $u_n=P_nu$ (równoważnie skończonych rozwiązań
z dowodu twierdzenia~\ref{thm:pde20-duhamel})
korzystamy z równania $u_n'+Lu_n=P_nf$
z początkową daną $P_nu_0$.
Iloczyn skalarny z $u_n'$ daje
$\|u_n'\|_H^2+\frac12\frac{\mathrm d}{\mathrm dt}
\|L^{1/2}u_n\|_H^2=(P_nf,u_n')_H$.
Stosujemy nierówność
$|(P_nf,u_n')|\le(\|f\|_H^2+\|u_n'\|_H^2)/2$,
po czym całkujemy:
\[
 \int_0^t\|u_n'(r)\|_H^2\,\mathrm d\lambda_1(r)
 +\|L^{1/2}u_n(t)\|_H^2
 \le \|P_nu_0\|_V^2+\int_0^t\|P_nf(r)\|_H^2\,\mathrm d\lambda_1(r).
\]
To samo oszacowanie stosujemy do $u_n-u_m$, z danymi
$(P_n-P_m)u_0$ i $(P_n-P_m)f$. Prawa strona
dąży do zera z silnej zbieżności rzutów w $V$
oraz $L^2(0,T;H)$. Otrzymujemy zatem ciąg Cauchy'ego
w $C([0,T];V)$ oraz ciąg pochodnych Cauchy'ego
w $L^2(0,T;H)$. Równanie pokazuje ponadto, że
$Lu_n$ jest ciągiem Cauchy'ego w $L^2(0,T;H)$.
Przy normie grafowej
$\|v\|_{D(L)}^2=\|v\|_H^2+\|Lv\|_H^2$
przestrzeń $D(L)$ jest zupełna dzięki domkniętości $L$.
W konsekwencji $u_n$ zbiega również w
$L^2(0,T;D(L))$. Wszystkie granice funkcji $u_n$
są równe $u$, gdyż zanurzenia do $L^2(0,T;H)$
są ciągłe. Przechodząc do granicy w tożsamości
dystrybucyjnej dla pochodnych, otrzymujemy
$u\in H^1(0,T;H)$ i $Lu=f-u'$ w $L^2(0,T;H)$.
Mamy także $u\in C([0,T];V)$.
Oznaczmy $M=\|u_0\|_V^2+\|f\|_{L^2(0,T;H)}^2$.
Przejście do granicy w oszacowaniu daje
$\|u'\|_{L^2(0,T;H)}^2\le M$ oraz
$\sup_t\|u(t)\|_V^2\le M$; dodanie obu nierówności
daje tezę. Dodatkowo
$\|Lu\|_{L^2(0,T;H)}\le\|f\|_{L^2(0,T;H)}
+\|u'\|_{L^2(0,T;H)}$.
\end{proof}

\section{Zwartość czasowo-przestrzenna: lemat Aubina--Lionsa w naszym modelu}
\begin{twierdzenie}[Zwartość dla parabolicznych ciągów Galerkina]\label{thm:pde20-compact}
Niech $\Omega\subset\R^d$ będzie ograniczonym obszarem
Lipschitzowskim. Jeżeli ciąg $\{u_n\}$ jest ograniczony
w $L^2(0,T;V)$ i ciąg $\{u_n'\}$ jest ograniczony
w $L^2(0,T;V^*)$, to $\{u_n\}$ jest względnie zwarty
w $L^2(0,T;H)$.
\end{twierdzenie}
\begin{proof}
Niech $P_N$ oznacza stały rzut na pierwsze $N$ funkcji
własnych laplasjanu Dirichleta. Z rozwinięcia w bazie
$\{e_j\}$ dostajemy dla prawie każdego $t$
\[
 \|(I-P_N)u_n(t)\|_H^2
 \le\mu_{N+1}^{-1}\|u_n(t)\|_V^2.
\]
Całkowanie po czasie i wspólne ograniczenie norm $V$ dają
\begin{equation}\label{eq:pde20-tail}
 \sup_n\|(I-P_N)u_n\|_{L^2(0,T;H)}^2
 \le C\mu_{N+1}^{-1}\xrightarrow[N\to\infty]{}0.
\end{equation}
Dla ustalonego $N$ współczynniki
$c_{n,j}(t)=(u_n(t),e_j)_H$, $j\le N$, są ograniczone
w $L^2(0,T)$, a ich pochodne spełniają
$|c_{n,j}'(t)|=|\langle u_n'(t),e_j\rangle|
\le\|u_n'(t)\|_{V^*}\|e_j\|_V$.
Zatem dla każdego $j\le N$ ciąg $\{c_{n,j}\}_n$ jest
ograniczony w $H^1(0,T)$, a więc względnie zwarty
w $C([0,T])$. Dla absolutnie ciągłego reprezentanta
nierówność Cauchy'ego--Schwarza daje
$|c(t)-c(s)|\le\|c'\|_{L^2}|t-s|^{1/2}$.
Wybierając $s$ z $|c(s)|\le T^{-1/2}\|c\|_{L^2}$,
otrzymujemy również
$\|c\|_\infty\le T^{-1/2}\|c\|_{L^2}
+\sqrt T\|c'\|_{L^2}$.
Mamy zatem jednostajne ograniczenie i równociągłość,
więc stosujemy twierdzenie Arzelà--Ascoliego.
Dla każdego $N$ wybieramy kolejno podciąg poprzedniego
podciągu, na którym pierwsze $N$ współczynników
zbiegają w $C([0,T])$. Podciąg przekątny $(u_{n_k})$
ma tę własność dla każdego ustalonego $N$,
a zatem $P_Nu_{n_k}$ zbiega w $L^2(0,T;H)$.
Dla $\eps>0$ wybieramy $N$ tak, by z~\eqref{eq:pde20-tail}
$\sup_n\|(I-P_N)u_n\|_{L^2(0,T;H)}<\eps/3$.
Następnie dla dostatecznie dużych $k,l$
$\|P_N(u_{n_k}-u_{n_l})\|_{L^2(0,T;H)}<\eps/3$.
Nierówność trójkąta daje
$\|u_{n_k}-u_{n_l}\|_{L^2(0,T;H)}<\eps$.
Z zupełności tej przestrzeni podciąg jest zbieżny,
co dowodzi względnej zwartości.
\end{proof}


\begin{wniosek}[Szacunek przesunięć czasowych w normie $H$]
\label{cor:pde20-translations}
Niech $\Omega\subset\R^d$ będzie ograniczonym obszarem
Lipschitzowskim, $V=H_0^1(\Omega)$, $H=L^2(\Omega)$,
a $u\in L^2(0,T;V)$, $u'\in L^2(0,T;V^*)$.
Dla $0<h<T$ zachodzi
\[
 \int_0^{T-h}\|u(t+h)-u(t)\|_H^2\,\mathrm dt
 \le 2h\,\|u\|_{L^2(0,T;V)}
 \|u'\|_{L^2(0,T;V^*)}.
\]
W szczególności rodzina funkcji o jednostajnie
ograniczonych obu normach ma jednostajny moduł ciągłości
względem translacji czasu w $L^2(H)$.
\end{wniosek}
\begin{proof}
Używamy ciągłego reprezentanta z lematu
~\ref{lem:pde20-energy}. Dla skończonych rzutów
mamy tożsamość
$P_Nu(t+h)-P_Nu(t)=\int_t^{t+h}P_Nu'(s)
\,\mathrm d\lambda_1(s)$ w $V^*$.
Przechodzimy do granicy: lewa strona zbiega dzięki
jednostajnej zbieżności w $H$ i ciągłemu zanurzeniu
do $V^*$, a prawa z nierówności
$\|\int_t^{t+h}(P_Nu'-u')\,\mathrm d\lambda_1\|_{V^*}
\le\sqrt h\|P_Nu'-u'\|_{L^2(0,T;V^*)}$.
Dla wszystkich $t\in[0,T-h]$ otrzymujemy w $V^*$
$u(t+h)-u(t)=\int_t^{t+h}u'(s)\,\mathrm d\lambda_1(s)$.
Dzięki nierówności Minkowskiego dla przesunięć oraz
przedłużeniu $u'$ przez zero poza $(0,T)$ mamy
\[
 \|u(\cdot+h)-u\|_{L^2(0,T-h;V^*)}
 \le h\|u'\|_{L^2(0,T;V^*)}.
\]
Z drugiej strony nierówność trójkąta daje
$\|u(\cdot+h)-u\|_{L^2(0,T-h;V)}
\le2\|u\|_{L^2(0,T;V)}$.
Z osadzeń $V\hookrightarrow H\hookrightarrow V^*$,
przy identyfikacji $H=H^*$, otrzymujemy
$\|w\|_H^2\le\|w\|_V\|w\|_{V^*}$.
Całkujemy tę nierówność dla $w(t)=u(t+h)-u(t)$,
a następnie stosujemy Cauchy'ego--Schwarza w czasie.
\end{proof}

\begin{uwaga}[Zastosowanie do nieliniowości]
Zbieżność słaba w $L^2(0,T;V)$ na ogół nie pozwala
przechodzić do granicy w wyrażeniu nieliniowym $F(u_n)$.
Zbieżność silna w $L^2(0,T;H)$ otrzymana powyżej
jest często brakującym elementem dowodów istnienia
w nieliniowych PDE. Ogólne twierdzenie Aubina--Lionsa
obejmuje więcej par przestrzeni i inne wykładniki;
tutaj udowodniliśmy samodzielnie wystarczający przypadek.
\end{uwaga}

\section{Nieliniowy kontrast: możliwość eksplozji}
\begin{twierdzenie}[Warunkowa eksplozja dla równania reakcji--dyfuzji]
\label{thm:pde20-blowup}
Niech $p>1$, $\Omega\subset\R^d$ będzie ograniczonym,
gładkim obszarem, a $\phi_1>0$ wewnątrz $\Omega$
będzie pierwszą funkcją własną laplasjanu Dirichleta,
$-\Delta\phi_1=\mu_1\phi_1$, z normalizacją
$\int_\Omega\phi_1\,\mathrm d\lambda_d(x)=1$.
Niech $u\ge0$ będzie dostatecznie regularnym rozwiązaniem
$u_t-\Delta u=u^p$ na $[0,T_*)\times\Omega$
z jednorodnym warunkiem Dirichleta. Jeżeli
$F(0):=\int_\Omega u(0,x)\phi_1(x)\,\mathrm d\lambda_d(x)
>\mu_1^{1/(p-1)}$, to nie istnieje globalne, nieujemne rozwiązanie klasyczne
na $[0,\infty)$ z funkcjonałem $F(t)$ skończonym w każdej chwili.
Dokładniej, jeśli $T_y<\infty$ jest czasem eksplozji rozwiązania
$y'=-\mu_1y+y^p$, $y(0)=F(0)$, to rozwiązanie $u$ nie może
pozostawać klasyczne na żadnym przedziale $[0,T]$ z $T\ge T_y$.
Jeśli istnieje aż do $T_y$, wówczas $F(t)\to\infty$ przy $t\uparrow T_y$.
\end{twierdzenie}
\begin{proof}
Mnożąc równanie przez $\phi_1$ i całkując po $\Omega$,
a następnie korzystając z całkowania przez części
i z warunków Dirichleta, otrzymujemy
\[
 F'(t)=-\mu_1F(t)+\int_\Omega u(t,x)^p\phi_1(x)
 \,\mathrm d\lambda_d(x).
\]
Miara $\phi_1(x)\,\mathrm d\lambda_d(x)$ jest probabilistyczna,
więc nierówność Jensena daje
$\int_\Omega u^p\phi_1\,\mathrm d\lambda_d\ge F(t)^p$.
Stąd $F'\ge-\mu_1F+F^p$.
Dla $F(0)>\mu_1^{1/(p-1)}$ rozwiązanie skalarnego ODE
$y'=-\mu_1y+y^p$ rośnie i osiąga nieskończoność
w skończonym czasie, ponieważ dla dużych $y$
mamy $y^p-\mu_1y\ge y^p/2$, a
$\int_{y_0}^\infty 2s^{-p}\,\mathrm ds<\infty$.
Z zasady porównania dla skalarnych ODE
$F(t)\ge y(t)$, dopóki oba wyrażenia są skończone.
To wyklucza globalne ograniczone rozwiązanie wskazanego typu.
\end{proof}

\section{Stabilność wstecznego Eulera i rozdzielenie czasu od przestrzeni}
Dla dyskretyzacji laplasjanu Dirichleta na siatce jednowymiarowej
niech $L_h$ będzie symetryczną dodatnio określoną macierzą
różnicową. Przy kroku czasowym $\tau>0$ schemat
wstecznego Eulera ma postać
\begin{equation}\label{eq:pde20-implicit}
 \frac{U^{m+1}-U^m}{\tau}+L_hU^{m+1}=F^{m+1}.
\end{equation}
\begin{stwierdzenie}[Bezwarunkowa stabilność energetyczna]
\label{prop:pde20-be}
Jeżeli $F^{m+1}=0$ dla wszystkich $m$, to rozwiązanie
schematu~\eqref{eq:pde20-implicit} spełnia
$\|U^{m+1}\|_h\le\|U^m\|_h$ niezależnie od $\tau/h^2$,
gdzie $\|U\|_h^2=h\sum_j|U_j|^2$.
\end{stwierdzenie}
\begin{proof}
Bierzemy dyskretny iloczyn skalarny równania przez $U^{m+1}$.
Tożsamość $2(a-b,a)=\|a\|^2-\|b\|^2+\|a-b\|^2$
daje
\[
 \|U^{m+1}\|_h^2-\|U^m\|_h^2
 +\|U^{m+1}-U^m\|_h^2
 +2\tau(L_hU^{m+1},U^{m+1})_h=0.
\]
Ostatnie dwa składniki są nieujemne; stąd teza.
Jest to twierdzenie o stabilności, nie samoistny dowód
zbieżności: potrzebna jest jeszcze kontrola błędu aproksymacji.
\end{proof}

\begin{uwaga}[Granice bieżącego rozdziału]
Dla współczynników zmiennych w czasie udowodniliśmy
istnienie i jednoznaczność rozwiązania w naturalnej klasie
$L^2(V)\cap H^1(V^*)$ bez gładkości $A$. Nie dowodziliśmy
tu pełnej parabolicznej teorii De Giorgiego--Nasha--Mosera,
teorii oszacowań Schaudera ani maksymalnej regularności
w klasach $L^p$ dla dowolnych operatorów niesymetrycznych.
Odpowiednie zagadnienia wymagają odrębnych założeń i metod.
\end{uwaga}


\chapter[Układy hiperboliczne i propagacja]{Równania hiperboliczne: układy symetryczno-hiperboliczne, energia i propagacja informacji}
\label{chap:pde21}

\section{Dlaczego potrzebujemy układów pierwszego rzędu?}
Rozdział~\ref{chap:pde-wave} dotyczył klasycznego równania falowego, przede wszystkim
jego jawnych formuł. Obecnie rozważamy układy o zmiennych współczynnikach,
bez dostępnej formuły rozwiązania. Ich centralną własnością będzie możliwość
kontrolowania energii bez tracenia pochodnych. Pozwoli to oddzielić
\emph{hiperboliczność symbolu} od \emph{poprawności zagadnienia Cauchy'ego}.

\begin{uwaga}[Dwie zasady energetyczne: dyfuzja i propagacja]
\label{rem:pde-v43-energy-types}
Oszacowanie paraboliczne z
twierdzenia~\ref{thm:pde20-existence} kontroluje także
całkę w czasie z normy przestrzennej $H_0^1$ dzięki
koercywności operatora dyfuzji. W układzie
symetryczno-hiperbolicznym energia kontroluje przede
wszystkim normę $L^2$ w chwili $t$, natomiast jej
lokalny bilans daje stożki zależności. Nie wolno
wywnioskować wygładzania przestrzennego z samej
nierówności energetycznej dla układu hiperbolicznego.
\end{uwaga}

Niech $d,m$ będą dodatnimi liczbami całkowitymi i $0<T<\infty$.
Dla kolumny $u=(u_1,\ldots,u_m)^\mathsf T$
rozważmy na $[0,T]\times\R^d$ równanie
\begin{equation}\label{eq:pde21-main}
 \partial_tu+\sum_{j=1}^d A_j(t,x)\partial_{x_j}u+B(t,x)u=f(t,x),
 \qquad u(0,x)=u_0(x).
\end{equation}
Macierze $A_j,B$ mają rozmiar $m\times m$ i rzeczywiste współczynniki.
Pracujemy nad ciałem liczb rzeczywistych. Przez $H$
oznaczamy przestrzeń Hilberta $L^2(\R^d;\R^m)$
z iloczynem skalarnym
$(v,w)_H=\int_{\R^d}v(x)^\mathsf Tw(x)\,\mathrm d\lambda_d(x)$
i normą $\|v\|_2^2=\int_{\R^d}|v(x)|^2\,\mathrm d\lambda_d(x)$.
Norma $|\cdot|$ w $\R^m$ jest normą euklidesową,
a $\|\cdot\|_{\mathrm{op}}$ jest indukowaną przez nią
normą operatorową macierzy. Przestrzenie funkcji czasu
o wartościach w przestrzeniach Sobolewa i w $H$
rozumiemy w sensie Bochnera.

\begin{definicja}[Układ symetryczno-hiperboliczny]\label{def:pde21-sym}
Układ~\eqref{eq:pde21-main} nazywamy \emph{symetryczno-hiperbolicznym},
jeżeli $A_j(t,x)=A_j(t,x)^\mathsf T$ dla wszystkich $j$ i wszystkich
$(t,x)$ (lub prawie wszędzie, gdy pracujemy ze współczynnikami mierzalnymi).
Jego symbolem głównym (z pominięciem wspólnego
czynnika $\mathrm i$ w konwencji Fouriera) jest
$P(t,x;\tau,\xi)=\tau I_m+\sum_{j=1}^d\xi_jA_j(t,x)$.
Dla współczynników mierzalnych rozpatrujemy symbol
dla prawie każdego $(t,x)$.
Kowektor $(\tau,\xi)\ne0$ jest charakterystyczny, jeżeli
$\det P(t,x;\tau,\xi)=0$.
\end{definicja}

\begin{stwierdzenie}[Rzeczywiste prędkości charakterystyczne]
\label{prop:pde21-real}
W każdym punkcie $(t,x)$, w którym wszystkie macierze
$A_j$ są symetryczne, dla każdego
$\xi\in\R^d\setminus\{0\}$ wszystkie pierwiastki wielomianu
$\tau\mapsto\det(\tau I_m+\sum_j\xi_jA_j)$ są rzeczywiste.
Płaszczyzna $t=\mathrm{const}$ jest niecharakterystyczna
niezależnie od symetrii $A_j$.
\end{stwierdzenie}
\begin{proof}
Macierz $A(\xi)=\sum_j\xi_jA_j$ jest rzeczywista i symetryczna.
Z twierdzenia spektralnego istnieje baza ortonormalna jej wektorów
własnych i rzeczywiste liczby $\lambda_1(\xi),\ldots,\lambda_m(\xi)$.
W tej bazie $\det(\tau I_m+A(\xi))=\prod_k(\tau+\lambda_k(\xi))$.
Dla konormalnej do płaszczyzny $t=\mathrm{const}$ mamy $(\tau,\xi)=(1,0)$,
a zatem $P=I_m$ jest odwracalna.
\end{proof}

\begin{uwaga}
Symetryczna hiperboliczność nie wymaga, aby wartości
własne macierzy $A(\xi)=\sum_j\xi_jA_j$ były różne.
Ścisła hiperboliczność względem czasu oznacza, że
dla każdego $\xi\ne0$ wielomian charakterystyczny
w zmiennej $\tau$ ma $m$ różnych pierwiastków
rzeczywistych (w każdym rozpatrywanym punkcie).
Nie należy utożsamiać tych dwóch własności. Co więcej,
sama rzeczywistość pierwiastków wielomianu charakterystycznego nie
zapewnia oszacowania energetycznego; zobaczymy to w
kontrprzykładzie~\ref{ex:pde21-jordan}.
\end{uwaga}

\section{Tożsamość energii i zależność od danych}
\begin{lemat}[Lokalny bilans energii]\label{lem:pde21-local}
Niech $A_j\in C^1([0,T]\times\R^d;\R^{m\times m})$ będą symetryczne,
$B$ ciągła, a $u\in C^1([0,T]\times\R^d;\R^m)$ rozwiązuje
\eqref{eq:pde21-main}. Wówczas punktowo zachodzi
\begin{equation}\label{eq:pde21-flux}
 \partial_t|u|^2+\sum_{j=1}^d\partial_{x_j}(u^\mathsf TA_ju)
 =u^\mathsf T\!\left(\sum_{j=1}^d\partial_{x_j}A_j-B-B^\mathsf T\right)u
 +2u^\mathsf Tf.
\end{equation}
\end{lemat}
\begin{proof}
Mnożymy równanie przez $2u^\mathsf T$. Otrzymujemy
$2u^\mathsf Tu_t+2\sum_j u^\mathsf TA_ju_{x_j}+2u^\mathsf TBu=2u^\mathsf Tf$.
Pierwszy składnik jest $\partial_t|u|^2$. Symetria $A_j$ daje
$\partial_{x_j}(u^\mathsf TA_ju)=2u^\mathsf TA_ju_{x_j}
+u^\mathsf T(\partial_{x_j}A_j)u$.
Wreszcie $2u^\mathsf TBu=u^\mathsf T(B+B^\mathsf T)u$.
Przenosząc odpowiednie wyrazy, otrzymujemy~\eqref{eq:pde21-flux}.
\end{proof}

\begin{twierdzenie}[Oszacowanie energii na całej przestrzeni]\label{thm:pde21-energy}
Niech $A_j$ będą wspólnie mierzalnymi macierzami
symetrycznymi prawie wszędzie, o jednostajnie ograniczonych
normach w $W^{1,\infty}(\R^d;\R^{m\times m})$
dla prawie każdego $t$ (nie wymagamy pochodnej czasowej),
$B\in L^\infty((0,T)\times\R^d;\R^{m\times m})$,
$f\in L^1(0,T;H)$ oraz $u_0\in H$.
Niech $u\in C([0,T];H)\cap L^1(0,T;H^1(\R^d;\R^m))$
rozwiązuje~\eqref{eq:pde21-main} w sensie dystrybucyjnym
i spełnia $u(0)=u_0$ w $H$.
Oznaczmy
\[
 D=\sum_{j=1}^d\partial_{x_j}A_j-B-B^\mathsf T,\qquad
 M=\operatorname*{ess\,sup}_{(t,x)}
       \|D(t,x)\|_{\mathrm{op}}.
\]
Pochodne współczynników rozumiemy w sensie słabym.
Wtedy dla każdego $t\in[0,T]$
\begin{equation}\label{eq:pde21-L2}
 \|u(t)\|_2\le e^{Mt/2}\left(\|u_0\|_2+
 \int_0^t\|f(r)\|_2\,\mathrm d\lambda_1(r)\right).
\end{equation}
To samo oszacowanie zachodzi dla różnicy dwóch rozwiązań, z prawą
stroną równą różnicy wymuszeń.
\end{twierdzenie}
\begin{proof}
Z równania dystrybucyjnego i ograniczoności
współczynników wynika
\[
 u'=f-\sum_j A_j\partial_{x_j}u-Bu\in L^1(0,T;H).
\]
Zatem $u$ należy do $W^{1,1}(0,T;H)$,
a jego absolutnie ciągły reprezentant zgadza się
z przyjętym ciągłym reprezentantem.
Dla $E(t)=\|u(t)\|_2^2$ reguła łańcuchowa w przestrzeni
Hilberta daje $E'(t)=2(u'(t),u(t))_H$ prawie wszędzie.
Dla prawie każdego $t$ mamy $u(t)\in H^1$,
a iloczyn $u^\mathsf TA_ju$ należy do $W^{1,1}(\R^d)$.
Istotnie, aproksymacja $u(t)$ przez funkcje z
$C_c^\infty$ w $H^1$, ograniczoność $A_j$ i jego
pochodnych oraz nierówność Cauchy'ego--Schwarza
dają zbieżność iloczynów w $W^{1,1}$, z pochodną
\[
 \partial_{x_j}(u^\mathsf TA_ju)
 =2u^\mathsf TA_j\partial_{x_j}u+
 u^\mathsf T(\partial_{x_j}A_j)u.
\]
Wybierzmy $\eta\in C_c^\infty(\R^d)$, $0\le\eta\le1$,
równą $1$ w otoczeniu zera, i połóżmy
$\eta_R(x)=\eta(x/R)$. Całkowanie przez części daje
\[
 \int_{\R^d}\eta_R\partial_{x_j}(u^\mathsf TA_ju)
 \,\mathrm d\lambda_d(x)
 =-\int_{\R^d}(\partial_{x_j}\eta_R)u^\mathsf TA_ju
 \,\mathrm d\lambda_d(x).
\]
Moduł prawej strony jest nie większy niż
$R^{-1}\|\nabla\eta\|_\infty\|A_j\|_\infty\|u(t)\|_2^2$,
a lewa zbiega do całki po całej przestrzeni
z twierdzenia o zbieżności zmajoryzowanej.
W konsekwencji całka pochodnej jest zerowa.
Po podstawieniu równania do pochodnej $E$
otrzymujemy globalny odpowiednik bilansu
~\eqref{eq:pde21-flux}:
\[
 E'(t)=\int_{\R^d}u(t,x)^\mathsf TD(t,x)u(t,x)
 \,\mathrm d\lambda_d(x)+2(u(t),f(t))_H.
\]
Wszystkie składniki są całkowalne w czasie,
ponieważ $u$ jest ograniczona w $H$, $D$ jest
ograniczona, a $f\in L^1(0,T;H)$.
Stąd $E'\le ME+2\sqrt E\|f\|_2$ prawie wszędzie.
Aby nie dzielić przez zero, definiujemy
$y_\eps(t)=\sqrt{E(t)+\eps^2}$ dla $\eps>0$.
Z absolutnej ciągłości i reguły łańcuchowej mamy
\[
 y_\eps'=\frac{E'}{2y_\eps}
 \le\frac{M}{2}y_\eps+\|f\|_2.
\]
Mnożymy przez $e^{-Mt/2}$ i całkujemy od $0$ do $t$:
\[
 y_\eps(t)\le e^{Mt/2}y_\eps(0)
 +\int_0^t e^{M(t-r)/2}\|f(r)\|_2\,\mathrm d\lambda_1(r).
\]
Przechodzimy z $\eps\downarrow0$. Ponieważ $M\ge0$,
dostajemy~\eqref{eq:pde21-L2}. Dla różnicy dwóch
rozwiązań z tej samej klasy stosujemy ten sam dowód
z różnicą danych i wymuszeń.
\end{proof}

\begin{wniosek}[Jednoznaczność]
W klasie rozwiązań z twierdzenia~\ref{thm:pde21-energy}
zagadnienie Cauchy'ego dla~\eqref{eq:pde21-main}
ma co najwyżej jedno rozwiązanie.
\end{wniosek}
\begin{proof}
Dla różnicy dwóch rozwiązań z tymi samymi danymi mamy zerowe
dane początkowe i zerową prawą stronę. Stosujemy~\eqref{eq:pde21-L2}.
\end{proof}

\section{Stożki zależności i skończona prędkość propagacji}
\begin{definicja}[Maksymalna prędkość symbolu]
Dla symetrycznych $A_j$ definiujemy
$c_*:=\operatorname*{ess\,sup}_{(t,x)\in(0,T)\times\R^d}
\sup_{|\nu|=1}\|\sum_j\nu_jA_j(t,x)\|_{\mathrm{op}}$.
Supremum po $\nu$ jest mierzalną funkcją $(t,x)$,
gdyż ciągłość względem $\nu$ pozwala zastąpić sferę
jednostkową jej przeliczalnym gęstym podzbiorem.
Dla symetrycznej macierzy $A(\nu)=\sum_j\nu_jA_j$
norma operatorowa jest maksimum modułów jej wartości
własnych. W szczególności
$c_*\le\sum_j\|A_j\|_{L^\infty;\mathrm{op}}$.
W dalszych rozważaniach zakładamy $c_*<\infty$.
Jeżeli współczynniki są ciągłe na $[0,T]\times\R^d$,
to nierówność $\|A(\nu)\|_{\mathrm{op}}\le c_*$
zachodzi w każdym punkcie i dla każdego $|\nu|=1$:
przekroczenie tej granicy dla ustalonego $\nu$
utrzymywałoby się na względnie otwartym zbiorze
dodatniej miary w walcu.
\end{definicja}

\begin{twierdzenie}[Stożkowa zasada zależności]\label{thm:pde21-cone}
Niech $A_j\in C^1_b([0,T]\times\R^d)$ będą symetryczne,
$B\in C_b([0,T]\times\R^d)$ i niech $u$ będzie klasycznym
rozwiązaniem~\eqref{eq:pde21-main}. Ustalmy $x_0\in\R^d$,
$t_0\in(0,T]$, $R\ge0$ i $c\ge c_*$. Jeśli
$u(0,x)=0$ dla $|x-x_0|\le R+ct_0$ oraz
$f(t,x)=0$ dla $0\le t\le t_0$, $|x-x_0|\le R+c(t_0-t)$,
to $u(t_0,x)=0$ dla $|x-x_0|\le R$.
Dla $R=0$ wniosek dotyczy punktu $x_0$.
\end{twierdzenie}
\begin{proof}
Najpierw załóżmy $R>0$. Niech $r(t)=R+c(t_0-t)$
i $D(t)=B_{r(t)}(x_0)$. Funkcję
$E_D(t)=\int_{D(t)}|u(t,x)|^2\,\mathrm d\lambda_d(x)$
zapisujemy na ustalonej kuli przez podstawienie
$x=x_0+r(t)y$:
\[
 E_D(t)=r(t)^d\int_{B_1(0)}
 g(t,x_0+r(t)y)\,\mathrm d\lambda_d(y),\qquad
 g(t,x)=|u(t,x)|^2.
\]
Promień jest dodatni, a $g$ i jej pierwsze pochodne
są ciągłe i ograniczone na rozpatrywanym zwartym
walcu. Różniczkowanie pod znakiem całki i powrót
do zmiennej $x$ dają
\[
 \begin{split}
 E_D'(t)
 &=\int_{D(t)}\partial_tg(t,x)\,\mathrm d\lambda_d(x)\\
 &\quad+\frac{r'(t)}{r(t)}
 \int_{D(t)}\operatorname{div}_x((x-x_0)g(t,x))
 \,\mathrm d\lambda_d(x).
 \end{split}
\]
Na brzegu kuli $(x-x_0)\cdot\nu=r(t)$, gdzie
$\nu$ jest normalną jednostkową zewnętrzną. Z twierdzenia
Gaussa otrzymujemy więc wzór transportu Reynoldsa
w tym szczególnym przypadku:
\[
 E_D'(t)=\int_{D(t)}\partial_tg(t,x)\,\mathrm d\lambda_d(x)
 +r'(t)\int_{\partial D(t)}g(t,x)\,\mathrm d\mathcal H^{d-1}(x).
\]
Dla $d=1$ to samo obliczenie jest regułą różniczkowania
całki o ruchomych końcach, a $\mathcal H^0$ jest miarą
liczącą na dwóch końcach przedziału.

Oznaczmy $D_A=\sum_j\partial_{x_j}A_j-B-B^\mathsf T$
i $M=\|D_A\|_{L^\infty;\mathrm{op}}$.
Wstawiając lokalny bilans~\eqref{eq:pde21-flux}
oraz $r'=-c$ i stosując twierdzenie Gaussa do strumienia,
otrzymujemy
\[
 \begin{split}
 E_D'(t)
 &=\int_{D(t)}(u^\mathsf TD_Au+2u^\mathsf Tf)
       \,\mathrm d\lambda_d(x)\\
 &\quad-\int_{\partial D(t)}
       (u^\mathsf TA_\nu u+c|u|^2)\,\mathrm d\mathcal H^{d-1}(x),
 \end{split}
\]
gdzie $A_\nu=\sum_j\nu_jA_j$. Z punktowego ograniczenia
$\|A_\nu\|_{\mathrm{op}}\le c_*\le c$ wynika
$u^\mathsf TA_\nu u+c|u|^2\ge0$. Całka brzegowa
w bilansie ma więc znak niedodatni. Ponieważ
$f=0$ w rozważanym stożku, otrzymujemy
$E_D'(t)\le M E_D(t)$. Dane początkowe dają $E_D(0)=0$.
Funkcja $e^{-Mt}E_D(t)$ jest nieujemna i nierosnąca,
zatem $E_D(t)=0$ dla $0\le t\le t_0$.
W otwartej kuli $D(t)$ funkcja $u(t,\cdot)$ znika
punktowo: niezerowa wartość przez ciągłość dawałaby
otoczenie dodatniej miary z dodatnią energią.
Przez ciągłość znika również na brzegu kuli. Dla $R=0$ i $c>0$ przeprowadzamy ten sam rachunek na $[0,t_0)$:
promień $r(t)=c(t_0-t)$ jest tam dodatni, a zerowa energia daje
$u(t,x_0)=0$ dla wszystkich $t<t_0$. Ciągłość w czasie pozwala
przejść do $t=t_0$. Jeżeli $c=0$, to również $c_*=0$, a więc
$A_j=0$ i równanie w każdym punkcie jest zwykłym równaniem
różniczkowym. Istotnie, podstawienie $\nu=e_j$
daje $A_j=0$ w każdym punkcie. Na prostej $x=x_0$
mamy $u_t=-Bu$, gdyż $f(t,x_0)=0$, i $u(0,x_0)=0$.
Nierówność $(|u|^2)'\le2\|B\|_\infty |u|^2$
daje zerowanie także w tym przypadku.
\end{proof}

\begin{wniosek}[Lokalna zależność od danych i propagacja nośnika]
Dwa klasyczne rozwiązania tego samego układu
(z różnymi danymi początkowymi i wymuszeniami)
mają równe wartości na $\{x:|x-x_0|\le R\}$ w chwili
$t_0$, jeśli ich dane początkowe i wymuszenia są równe
w zbiorach wskazanych w twierdzeniu~\ref{thm:pde21-cone}.
W szczególności przy $f=0$ i zwartym nośniku danych
początkowych $K$ zachodzi
\[
 \operatorname{supp}u(t,\cdot)
 \subset \{x\in\R^d:\operatorname{dist}(x,K)\le ct\},
 \qquad 0\le t\le T,\quad c\ge c_*.
\]
\end{wniosek}
\begin{proof}
Pierwszą część otrzymujemy, stosując twierdzenie
~\ref{thm:pde21-cone} do różnicy rozwiązań.
Dla $t>0$ i punktu $x$ o $\operatorname{dist}(x,K)>ct$
zbiór $\{y:|y-x|\le ct\}$ nie przecina $K$,
więc twierdzenie z $R=0$ daje $u(t,x)=0$.
Zbiór takich punktów jest otwarty (funkcja odległości
od niepustego zbioru jest lipschitzowska ze stałą $1$
na mocy nierówności trójkąta), więc nie przecina nośnika.
Dla $t=0$ wniosek wynika z definicji $K$.
Jeżeli dane są identycznie zerowe, twierdzenie zastosowane
w każdym punkcie daje $u=0$; przyjmujemy wtedy nośnik
pusty i odległość od niego równą $+\infty$.
\end{proof}



\begin{center}
\begin{tikzpicture}[scale=.88,>=Latex]
\fill[blue!8] (0.8,0) -- (6.0,0) -- (3.4,3.55) -- cycle;
\draw[->] (0,0) -- (6.8,0) node[right] {$x$};
\draw[->] (0,0) -- (0,4.2) node[above] {$t$};
\draw[thick,blue!65!black] (.8,0)--(3.4,3.55)--(6,0);
\draw[dashed] (1.6789,1.2)--(5.1211,1.2);
\draw[gray,dotted] (0,3.55)--(3.4,3.55);
\draw[gray,dotted] (3.4,0)--(3.4,3.55);
\draw[<->,blue!65!black] (.8,-.25)--(3.4,-.25)
 node[midway,below] {$ct_0$};
\draw[<->,blue!65!black] (3.4,-.25)--(6,-.25)
 node[midway,below] {$ct_0$};
\node[left] at (0,3.55) {$t_0$};
\node[below] at (3.4,-.65) {$x_0$};
\node[above right] at (3.4,3.55) {$(t_0,x_0)$};
\node[align=center,font=\small] at (3.4,.55) {dane i wymuszenie\\w stożku};
\node[right] at (5.1,1.2) {$t=t_1$};
\end{tikzpicture}

\small Rysunek: wsteczny stożek zależności w jednym wymiarze.
Gdy $R=0$, podstawa ma promień $ct_0$.
\end{center}

\section{Równanie falowe jako układ symetryczny}
Niech $c>0$ i niech $u\in C^2([0,T]\times\R)$
spełnia $u_{tt}-c^2u_{xx}=0$. Określamy $v=u_t$,
$w=cu_x$ i $U=(v,w)^\mathsf T$.
Z równania falowego mamy $v_t=cw_x$, a z równości
pochodnych mieszanych $w_t=cv_x$. Wówczas
\begin{equation}\label{eq:pde21-wavefirst}
 U_t+\begin{pmatrix}0&-c\\-c&0\end{pmatrix}U_x=0,
 \qquad |U|^2=u_t^2+c^2u_x^2.
\end{equation}
Macierz główna jest symetryczna; wektory
$2^{-1/2}(1,-1)^\mathsf T$ i $2^{-1/2}(1,1)^\mathsf T$
są jej wektorami własnymi odpowiednio dla $c$ i $-c$.
Są to prędkości przenoszenia dwóch niezależnych modów.

Odwrotnie, niech $(v,w)\in C^1([0,T]\times\R;\R^2)$
spełnia~\eqref{eq:pde21-wavefirst}.
Jeżeli $a\in C^2(\R)$ i $w(0,x)=ca'(x)$, to definiujemy
\[
 u(t,x)=a(x)+\int_0^t v(r,x)\,\mathrm d\lambda_1(r).
\]
Z twierdzenia Newtona--Leibniza mamy $u_t=v$.
Różniczkując pod znakiem całki względem $x$ i używając
$w_t=cv_x$, otrzymujemy
\[
 cu_x(t,x)=ca'(x)+\int_0^t w_t(r,x)\,\mathrm d\lambda_1(r)
 =w(t,x).
\]
Obie pierwsze pochodne $u$ są klasy $C^1$, więc
$u\in C^2$ i $u_{tt}=v_t=cw_x=c^2u_{xx}$.
Dane początkowe to $u(0,x)=a(x)$ i $u_t(0,x)=v(0,x)$.
Wybór $a$ jest jednoznaczny z dokładnością do
stałej addytywnej, jeśli znamy tylko $w(0,\cdot)$:
można wziąć $a(x)=a(0)+c^{-1}\int_0^x w(0,y)\,
\mathrm d\lambda_1(y)$. Układ opisuje więc pochodne
rozwiązania falowego, ale sam nie wyznacza jego
stałej addytywnej.

\begin{stwierdzenie}[Rozkład na fale podróżujące]
Niech $v,w$ będą gładkimi funkcjami rozwiązującymi
\eqref{eq:pde21-wavefirst} na $[0,T]\times\R$. Funkcje
$r_+=v-w$ i $r_-=v+w$ spełniają odpowiednio
$(\partial_t+c\partial_x)r_+=0$ oraz
$(\partial_t-c\partial_x)r_-=0$.
W konsekwencji $r_+(t,x)=r_+(0,x-ct)$ i
$r_-(t,x)=r_-(0,x+ct)$.
\end{stwierdzenie}
\begin{proof}
Równania składowych to $v_t=cw_x$ i $w_t=cv_x$.
Po odjęciu otrzymujemy $(v-w)_t=-c(v-w)_x$,
a po dodaniu $(v+w)_t=c(v+w)_x$.
Dla ustalonego $(t,x)$ funkcja
$r_+(s,x-c(t-s))$ ma pochodną
$(\partial_t+c\partial_x)r_+$ w punkcie
$(s,x-c(t-s))$, czyli zero. Porównując jej wartości
dla $s=0$ i $s=t$, otrzymujemy pierwszą formułę.
Analogicznie funkcja $r_-(s,x+c(t-s))$ ma pochodną
$(\partial_t-c\partial_x)r_-=0$, co daje drugą.
Jeżeli $u(0,x)=a(x)$ i $u_t(0,x)=b(x)$, to
\[
 r_+(0,x)=b(x)-ca'(x),\qquad
 r_-(0,x)=b(x)+ca'(x).
\]
Z $v=(r_++r_-)/2$ i $w=(r_--r_+)/2$
odzyskujemy obie pochodne rozwiązania falowego.
Całkując $v$ względem czasu z danymi $a$, dostajemy
\[
 u(t,x)=\frac{a(x-ct)+a(x+ct)}{2}
 +\frac1{2c}\int_{x-ct}^{x+ct}b(y)\,\mathrm d\lambda_1(y),
\]
gdyż podstawienia $y=x\mp cr$ w całkach z $b(x\mp cr)$
dają drugi składnik, a całki z $a'$ dają pierwszy
po dodaniu $a(x)$. Jest to formuła d'Alemberta
z rozdziału~\ref{chap:pde-wave}.
\end{proof}

\section{Oszacowania wyższych pochodnych bez ich utraty}
\begin{lemat}[Komutator z pochodną wielokrotną]\label{lem:pde21-comm}
Niech $s\ge1$ będzie liczbą całkowitą,
$A\in W^{s,\infty}(\R^d;\R^{m\times m})$ oraz
$u\in H^s(\R^d;\R^m)$. Dla każdego wielowskaźnika
$\alpha$ z $|\alpha|\le s$ zachodzi
\[
 \|[\partial^\alpha,A]\partial_j u\|_2
 \le C_{d,s}\|A\|_{W^{s,\infty}}\|u\|_{H^s},
\]
gdzie komutator oznacza
$[\partial^\alpha,A]v=\partial^\alpha(Av)-A\partial^\alpha v$.
Dla $v=\partial_j u$ różnicę rozumiemy początkowo
w $H^{-1}(\R^d;\R^m)$: mnożenie przez $A$ na $H^{-1}$
definiujemy przez dualność, używając ograniczonego
operatora mnożenia przez $A^\mathsf T$ na $H^1$.
Teza stwierdza, że różnica należy faktycznie do $H$.
\end{lemat}
\begin{proof}
Najpierw dla $u\in C_c^\infty$ stosujemy słabą
regułę Leibniza (przez iterację reguły iloczynu
funkcji Sobolewa z funkcją gładką):
\[
 [\partial^\alpha,A]\partial_ju
 =\sum_{0<\beta\le\alpha}\binom{\alpha}{\beta}
 (\partial^\beta A)\partial_j\partial^{\alpha-\beta}u,
 \qquad
 \binom{\alpha}{\beta}=\prod_{k=1}^d
 \binom{\alpha_k}{\beta_k}.
\]
Nierówność $\beta\le\alpha$ rozumiemy współrzędnie.
Dla $\alpha=0$ suma jest pusta. Dla każdego
pozostałego składnika mamy
$1+|\alpha|-|\beta|\le|\alpha|\le s$.
Norma $L^2$ tego składnika nie przekracza
$\binom{\alpha}{\beta}\|\partial^\beta A\|_\infty
\|\partial_j\partial^{\alpha-\beta}u\|_2$.
Sumowanie daje żądane oszacowanie.
Dla ogólnego $u\in H^s$ wybieramy $u_k\in C_c^\infty$
zbieżne w $H^s$. Prawa strona wzoru zbiega w $H$
na mocy powyższych oszacowań. Jednocześnie
$A\partial_j u_k\to A\partial_j u$ w $H^{s-1}$:
słaba reguła Leibniza i ograniczoność pochodnych $A$
dają ciągłość mnożenia na $H^{s-1}$.
Zatem pierwszy składnik komutatora zbiega w $H^{-1}$.
Drugi zbiega w $H^{-1}$ dzięki ciągłości mnożenia
na tej przestrzeni. Ponieważ zanurzenie $H\hookrightarrow H^{-1}$
jest ciągłe i iniektywne, granica w $H$ jest właśnie
żądanym komutatorem. Nie przybliżamy $A$ w normie
$W^{s,\infty}$.
\end{proof}

\begin{twierdzenie}[Energia w $H^s$]\label{thm:pde21-Hs}
Niech $s\ge1$ będzie liczbą całkowitą,
$A_j,B\in C([0,T];W^{s,\infty}(\R^d;\R^{m\times m}))$,
a macierze $A_j$ niech będą symetryczne.
Przyjmijmy normę
$\|u\|_{H^s}^2=\sum_{|\alpha|\le s}\|\partial^\alpha u\|_2^2$.
Niech $u_0\in H^s$, $f\in L^1(0,T;H^s)$
i niech $u\in C([0,T];H^s)\cap L^1(0,T;H^{s+1})$
będzie rozwiązaniem~\eqref{eq:pde21-main} w sensie
dystrybucyjnym, z danymi początkowymi $u_0$.
Wówczas
\begin{equation}\label{eq:pde21-Hs}
 \|u(t)\|_{H^s}\le e^{Ct}
 \left(\|u_0\|_{H^s}+\int_0^t\|f(r)\|_{H^s}\,\mathrm d\lambda_1(r)\right),
\end{equation}
gdzie $C$ zależy od $d,m,s$ i jednostajnych norm
$W^{s,\infty}$ współczynników, nie zaś od $u$ i $f$.
\end{twierdzenie}
\begin{proof}
Niech $|\alpha|\le s$. Po zastosowaniu $\partial^\alpha$
do~\eqref{eq:pde21-main} otrzymujemy
\[
 \partial_t(\partial^\alpha u)+\sum_jA_j\partial_j(\partial^\alpha u)
 +B\partial^\alpha u
 =\partial^\alpha f-\sum_j[\partial^\alpha,A_j]\partial_ju
 -[\partial^\alpha,B]u.
\]
Oznaczmy $v_\alpha=\partial^\alpha u$ i
$R_\alpha=-\sum_j[\partial^\alpha,A_j]\partial_ju
-[\partial^\alpha,B]u$.
Z lematu~\ref{lem:pde21-comm} oraz słabej reguły
Leibniza dla $B$ (tu każdy składnik zawiera co
najwyżej $s-1$ pochodnych $u$) wynika
\[
 \left(\sum_{|\alpha|\le s}\|R_\alpha(t)\|_2^2\right)^{1/2}
 \le C_0\|u(t)\|_{H^s},
\]
gdzie $C_0$ zależy wyłącznie od podanych norm
współczynników i liczby wielowskaźników.
Każde $v_\alpha$ należy do $C([0,T];H)\cap L^1(0,T;H^1)$,
a wymuszenie $\partial^\alpha f+R_\alpha$ do
$L^1(0,T;H)$. Stosujemy więc globalną tożsamość
energii z dowodu twierdzenia~\ref{thm:pde21-energy}.
Po zsumowaniu i zastosowaniu nierówności
Cauchy'ego--Schwarza do skończonej sumy iloczynów
skalarnych otrzymujemy dla
$E_s=\sum_{|\alpha|\le s}\|v_\alpha\|_2^2$
\[
 E_s'\le M E_s+2\sqrt{E_s}
   \bigl(\|f\|_{H^s}+C_0\sqrt{E_s}\bigr).
\]
Zatem $E_s'\le(M+2C_0)E_s+2\sqrt{E_s}\|f\|_{H^s}$.
Argument z $\sqrt{E_s+\eps^2}$ i czynnikiem całkującym
z dowodu twierdzenia~\ref{thm:pde21-energy}
daje~\eqref{eq:pde21-Hs} z $C=M/2+C_0$.
Poniższy lemat pozwoli usunąć pomocniczy warunek
całkowalności w $H^{s+1}$.
\end{proof}


\begin{lemat}[Przedłużenie tożsamości energii na rozwiązania słabe]
\label{lem:pde21-friedrichs-comm}
Niech $A_j,B$ spełniają założenia dotyczące
współczynników z twierdzenia~\ref{thm:pde21-energy}.
Jeżeli $v\in L^\infty(0,T;H)$ rozwiązuje dystrybucyjnie
$v_t+\sum_j A_j\partial_jv+Bv=F$ z $F\in L^1(0,T;H)$,
to ma jednoznacznego reprezentanta w $C([0,T];H)$.
Dla każdego $0\le a\le b\le T$ zachodzi
\[
 \begin{split}
 \|v(b)\|_2^2-\|v(a)\|_2^2
 &=\int_a^b\int_{\R^d}v(t,x)^\mathsf TD_A(t,x)v(t,x)
       \,\mathrm d\lambda_d(x)\,\mathrm d\lambda_1(t)\\
 &\quad+2\int_a^b(v(t),F(t))_H\,\mathrm d\lambda_1(t),
 \end{split}
\]
gdzie $D_A=\sum_j\partial_{x_j}A_j-B-B^\mathsf T$.
W szczególności zachodzi oszacowanie~\eqref{eq:pde21-L2}
z dowolnej chwili początkowej.
\end{lemat}
\begin{proof}
\emph{Krok 1: słabo ciągły reprezentant.}
Iloczyn z pochodną dystrybucyjną określamy przez
$A_j\partial_j v=\partial_j(A_jv)-(\partial_jA_j)v$,
co należy do $H^{-1}$ dla prawie każdego $t$.
Z równania wynika
$v'\in L^1(0,T;H^{-1})$, ponieważ $\partial_j:H\to H^{-1}$
jest ograniczony, współczynniki mają jednostajne
ograniczenia, a $F\in L^1(0,T;H)$.
Wybieramy absolutnie ciągły reprezentant
$w:[0,T]\to H^{-1}$ klasy $v\in W^{1,1}(0,T;H^{-1})$.
Niech $K=\|v\|_{L^\infty(0,T;H)}$ i niech $G$
będzie zbiorem czasów pełnej miary, dla których
$w(t)=v(t)$ w $H^{-1}$ oraz $\|v(t)\|_2\le K$.
Dla każdego $t$ wybieramy $t_n\in G$ z $t_n\to t$.
Ciąg $v(t_n)$ ma słabo zbieżny w $H$ podciąg
z granicą $h_t$, $\|h_t\|_2\le K$. Z ciągłości
zanurzenia $H\hookrightarrow H^{-1}$ jego słaba
granica w $H^{-1}$ jest obrazem $h_t$, ale
$w(t_n)\to w(t)$ silnie w $H^{-1}$.
Stąd $w(t)=h_t$ po utożsamieniu przez zanurzenie.
Jego iniektywność zapewnia jednoznaczność $h_t$.
Przyjmijmy $v(t)=h_t$ dla wszystkich $t$.
Dla $\phi\in H^1$ funkcja $t\mapsto(v(t),\phi)_H$
jest ciągła, bo jest wartością dualności $H^{-1},H^1$.
Aproksymacja dowolnego $\phi\in H$ funkcjami z $H^1$
i jednolite ograniczenie $\|v(t)\|_2\le K$ dają
słabą ciągłość $v:[0,T]\to H$.

\emph{Krok 2: komutator wygładzania.}
Wybierzmy $\rho\in C_c^\infty(\R^d)$ nieujemną,
o całce $1$, i połóżmy
$\rho_\eps(z)=\eps^{-d}\rho(z/\eps)$,
$J_\eps v=\rho_\eps*v$, gdzie splot jest tylko
względem zmiennych przestrzennych.
Operator $J_\eps$ jest kontrakcją w $H$ oraz
$J_\eps g\to g$ w $H$ dla każdego $g\in H$.
Funkcja $v_\eps=J_\eps v$ spełnia
\[
 \partial_t v_\eps+\sum_j A_j\partial_j v_\eps+Bv_\eps
 =J_\eps F+R_\eps,\qquad
 R_\eps=\sum_j[A_j,J_\eps]\partial_jv+[B,J_\eps]v,
\]
gdzie $[A,J_\eps]=AJ_\eps-J_\eps A$.
Dla ustalonego $t$ i początkowo gładkiej $v$
\[
 [J_\eps,A]\partial_jv(x)=
 \int_{\R^d}\rho_\eps(z)(A(x-z)-A(x))\partial_jv(x-z)
 \,\mathrm d\lambda_d(z).
\]
Całkując przez części względem $z_j$ (pochodna
$v(x-z)$ względem $z_j$ ma przeciwny znak),
otrzymujemy
\[
 \int_{\R^d}
 \bigl(\partial_j\rho_\eps(z)(A(x-z)-A(x))
 -\rho_\eps(z)(\partial_j A)(x-z)\bigr)v(x-z)
 \,\mathrm d\lambda_d(z).
\]
Całka samego współczynnika przy $v(x-z)$ jest
zerowa: jest całką pochodnej
$\partial_{z_j}(\rho_\eps(z)(A(x-z)-A(x)))$.
Możemy więc zastąpić $v(x-z)$ przez $v(x-z)-v(x)$.
Dla $A\in W^{1,\infty}$ używamy lipschitzowskiego
reprezentanta; powyższe całkowanie przez części
jest dopuszczalne w sensie słabym.
Przy $\operatorname{supp}\rho\subset\overline{B_{R_\rho}(0)}$
nierówność Minkowskiego dla całki daje
\[
 \|[J_\eps,A]\partial_jv\|_2
 \le C_\rho\|\nabla A\|_\infty
       \sup_{|z|\le\eps R_\rho}\|v(\cdot-z)-v\|_2,
\]
gdzie $C_\rho$ można wybrać jako
$\int |z|\,|\partial_j\rho(z)|\,\mathrm d\lambda_d(z)
+\|\rho\|_{L^1}$, z odpowiednią normą pochodnych
macierzy. Stała nie zależy od $\eps$, gdyż podstawienie
$z=\eps y$ znosi czynniki skalowania w pierwszej całce.
Wzór i oszacowanie rozszerzają się na $v\in H$ przez
aproksymację w $H$: prawa strona określa ograniczony
operator na $H$, a lewa zbiega dystrybucyjnie dzięki
definicji $A\partial_j v$ i ciągłości splotu.
Ciągłość translacji w $H$ daje zbieżność do zera
dla prawie każdego $t$, a oszacowanie przez
$2C_\rho\|\nabla A\|_\infty K$ i twierdzenie
o zbieżności zmajoryzowanej dają zbieżność
w $L^1(0,T;H)$. Dla członu bez pochodnej piszemy
\[
 [B,J_\eps]v=B(J_\eps v-v)+(Bv-J_\eps(Bv)).
\]
Oba składniki zbiegają do zera w $L^1(0,T;H)$
z silnej zbieżności splotu, ograniczoności $B$
i tego samego twierdzenia granicznego.
W konsekwencji $\|R_\eps\|_{L^1(0,T;H)}\to0$.
Również $J_\eps F\to F$ w $L^1(0,T;H)$.

\emph{Krok 3: bilans dla wszystkich chwil.}
Dla ustalonego $\eps$ operator $J_\eps$ jest
ograniczony z $H^{-1}$ do $H^k$ dla każdego
całkowitego $k\ge0$ (wynika to np. z szybkiego zaniku
transformaty Fouriera $\rho_\eps$).
Stąd $v_\eps=J_\eps w$ jest ciągły w $H^1$
na całym $[0,T]$ i należy do $W^{1,1}(0,T;H)$.
Jego wartości w każdej chwili są dokładnie
$J_\eps v(t)$, a nie tylko prawie wszędzie.
Stosujemy tożsamość z dowodu twierdzenia
~\ref{thm:pde21-energy} do równania wygładzonego,
z wymuszeniem $J_\eps F+R_\eps$, i całkujemy
od dowolnego $a$ do dowolnego $b$.
Dla każdego ustalonego $t\in[0,T]$ mamy
$J_\eps v(t)\to v(t)$ silnie w $H$ na mocy
własności aproksymacji jedynki, także w chwilach
spoza zbioru $G$. W szczególności obie wartości
końcowe w bilansie zbiegają do właściwych norm.
Ponadto $\|v_\eps(t)\|_2\le K$, a różnica wyrazów
kwadratowych z $D_A$ jest ograniczona przez
$2MK\|v_\eps(t)-v(t)\|_2$.
Twierdzenie o zbieżności zmajoryzowanej daje
zbieżność tych całek. Dla członu wymuszenia
rozkładamy różnicę na składniki z
$J_\eps F-F+R_\eps$ oraz $v_\eps-v$.
Pierwsze całki mają moduł nie większy niż
$K\|J_\eps F-F+R_\eps\|_{L^1(0,T;H)}$,
druga zbiega do zera przez majorantę $2K\|F(t)\|_2$.
Otrzymujemy zatem tożsamość z tezy dla \emph{każdej}
pary $a,b$, nie tylko prawie wszystkich końców przedziału.

W szczególności funkcja $t\mapsto\|v(t)\|_2^2$
jest absolutnie ciągła, gdyż prawa strona tożsamości
dla $a=0$ jest całką funkcji z $L^1(0,T)$.
Dla $t_n\to t$ słaba ciągłość i ciągłość normy dają
\[
 \|v(t_n)-v(t)\|_2^2
 =\|v(t_n)\|_2^2+\|v(t)\|_2^2-2(v(t_n),v(t))_H\to0.
\]
To dowodzi silnej ciągłości; jednoznaczność
reprezentanta wynika z gęstości zbioru pełnej miary.
Oszacowanie~\eqref{eq:pde21-L2} wynika z tej samej
nierówności różniczkowej dla normy, co w twierdzeniu
~\ref{thm:pde21-energy}, z początkiem całkowania w $a$.
\end{proof}

\section{Istnienie przez regularizację Friedricha}
Warto odróżnić metodę zastosowaną poniżej od rozdziału~20:
w równaniach parabolicznych otrzymywaliśmy zwartość z całki po czasie
normy $H^1$. Tutaj podstawowym narzędziem jest kontrola \emph{wszystkich}
pochodnych przestrzennych dzięki symetrii macierzy głównych.

\begin{twierdzenie}[Zagadnienie Cauchy'ego dla gładkich współczynników]
\label{thm:pde21-existence}
Niech $s>d/2+2$ będzie liczbą całkowitą. Załóżmy, że
$A_j=A_j^\mathsf T$ oraz
$A_j,B\in C([0,T];W^{s+1,\infty}(\R^d;\R^{m\times m}))$,
zaś $u_0\in H^s(\R^d;\R^m)$ i
$f\in C([0,T];H^s(\R^d;\R^m))$.
Wtedy istnieje dokładnie jedno rozwiązanie zagadnienia
\eqref{eq:pde21-main}, dla którego
\[
 u\in C([0,T];H^s(\R^d;\R^m))\cap
 C^1([0,T];H^{s-1}(\R^d;\R^m)).
\]
Spełnia ono~\eqref{eq:pde21-Hs}.
\end{twierdzenie}
\begin{proof}
\emph{Krok 1: regularizacja.}
Niech $P_N$ będzie ortogonalnym rzutem Fouriera na częstotliwości
$|\xi|\le N$. Operator $P_N$ jest kontrakcją w każdym $H^k$
i komutuje z pochodnymi. Szukamy rozwiązania
\begin{equation}\label{eq:pde21-friedrichs}
 u_N'+P_N\left(\sum_j A_j\partial_j u_N+Bu_N\right)=P_Nf,
 \quad u_N(0)=P_Nu_0,
 \quad P_Nu_N=u_N.
\end{equation}
Na podprzestrzeni $P_NH^s$ operator
$v\mapsto P_N(A_j\partial_jv+Bv)$ jest ograniczony, ponieważ
$\|\partial_jv\|_{H^s}\le N\|v\|_{H^s}$.
Równanie~\eqref{eq:pde21-friedrichs} jest więc zwykłym
liniowym równaniem różniczkowym w przestrzeni Banacha, o ciągłych
w czasie współczynnikach; metoda kolejnych przybliżeń daje
rozwiązanie na krótkim przedziale, a liniowość pozwala je
przedłużyć do $T$.

\emph{Krok 2: jednostajne oszacowania.}
Ponieważ $P_N$ jest samosprzężony i $P_N\partial^\alpha u_N=\partial^\alpha u_N$,
rachunek energii po zastosowaniu $\partial^\alpha$ wygląda tak
samo jak w dowodzie~\eqref{eq:pde21-Hs}. Istotnie
$\langle\partial^\alpha u_N,P_N(A_j\partial_j\partial^\alpha u_N)\rangle
=\langle\partial^\alpha u_N,A_j\partial_j\partial^\alpha u_N\rangle$.
Komutatory są oszacowane przez Lemat~\ref{lem:pde21-comm}.
Stąd $\sup_{N,t}\|u_N(t)\|_{H^s}\le C$.
Z~\eqref{eq:pde21-friedrichs} i ograniczoności mnożenia
przez $W^{s,\infty}$ dostajemy także
$\sup_{N,t}\|u_N'(t)\|_{H^{s-1}}\le C$.

\emph{Krok 3: zwartość lokalna.}
Na każdej kuli $B_R$ zanurzenie $H^s(B_R)\hookrightarrow
H^{s-1}(B_R)$ jest zwarte. Ponadto
$\|u_N(t)-u_N(r)\|_{H^{s-1}(B_R)}\le C|t-r|$.
Z twierdzenia Arzelà--Ascoliego w metrycznej przestrzeni
$H^{s-1}(B_R)$ otrzymujemy podciąg zbieżny w
$C([0,T];H^{s-1}(B_R))$. Wybierając podciąg przekątny dla
$R=1,2,\ldots$, dostajemy funkcję $u$, która jest granicą lokalną
oraz słabą-* granicą w $L^\infty(0,T;H^s(\R^d))$.

\emph{Krok 4: równanie i dane.}
Dla $\varphi\in C_c^\infty((0,T)\times\R^d;\R^m)$
mnożymy~\eqref{eq:pde21-friedrichs} przez $\varphi$
i całkujemy. Przenosimy $P_N$ na funkcję testową
$P_N\varphi$; ponieważ $P_N\varphi\to\varphi$ w każdym
$H^k$, przechodzimy do granicy w wyrazach liniowych,
korzystając z lokalnej silnej zbieżności oraz jednostajnych
oszacowań na całej przestrzeni. Otrzymujemy~\eqref{eq:pde21-main}
w sensie dystrybucyjnym. Zbieżność
$u_N(0)=P_Nu_0\to u_0$ w $H^{s-1}$ i lokalna ciągłość implikują
$u(0)=u_0$ jako element $H^{s-1}(\R^d)$.

\emph{Krok 5: energia, jednoznaczność i ciągłość silna.}
Dla każdego wielowskaźnika $|\alpha|\le s$ funkcja
$v=\partial^\alpha u$ należy do $L^\infty(0,T;L^2)$.
Po zróżniczkowaniu równania jej prawa strona jest sumą
$\partial^\alpha f$ i komutatorów z
Lematu~\ref{lem:pde21-comm}, a więc należy do $L^1(0,T;L^2)$.
Lemat~\ref{lem:pde21-friedrichs-comm} stosuje się zatem
osobno do wszystkich $v=\partial^\alpha u$; daje
$v\in C([0,T];L^2)$ i po zsumowaniu oszacowanie
~\eqref{eq:pde21-Hs}. Dzięki temu $u\in C([0,T];H^s)$
bez konieczności zakładania z góry ciągłości najwyższych pochodnych.
Równanie ma prawą stronę ciągłą z wartościami w $H^{s-1}$,
ponieważ mnożenie przez przyjęte współczynniki jest ciągłe
na $H^{s-1}$, zatem $u_t\in C([0,T];H^{s-1})$.
Ostatecznie $u\in C^1([0,T];H^{s-1})$.
Jednoznaczność wynika z Twierdzenia~\ref{thm:pde21-energy}
zastosowanego do różnicy dwóch rozwiązań.
\end{proof}

\begin{uwaga}[Wersja o mniejszej regularności]
Następnie można przez aproksymację danych w $L^2$
i oszacowanie~\eqref{eq:pde21-L2} otrzymać rozwiązania energetyczne
z ciągłością $C([0,T];L^2)$, ale wymaga to osobnego uzasadnienia
reguły energii dla granicznych rozwiązań. Twierdzenie powyżej
celowo sformułowaliśmy z mocniejszymi założeniami i bez ukrytej
utraty pochodnych.
\end{uwaga}

\section{Symetryzatory: gdy macierze nie są symetryczne}
\begin{definicja}[Układ symetryzowalny]\label{def:pde21-symmetrizer}
Niech $A_j$ będą rzeczywistymi macierzami $m\times m$.
Mówimy, że układ~\eqref{eq:pde21-main} ma \emph{symetryzator} $S(t,x)$,
jeżeli $S=S^\mathsf T$, istnieją stałe $0<\sigma\le\Sigma$ takie,
że $\sigma I_m\le S\le\Sigma I_m$, a macierze $SA_j$ są symetryczne.
Dla oszacowań energii zakładamy dodatkowo
$S\in W^{1,\infty}$ i $A_j,B\in W^{1,\infty}$.
\end{definicja}

\begin{twierdzenie}[Energia ważona symetryzatorem]
\label{thm:pde21-symmetrizer}
Niech założenia Definicji~\ref{def:pde21-symmetrizer} zachodzą
na $[0,T]\times\R^d$, a $u$ jest gładkim, zanikającym
w nieskończoności rozwiązaniem~\eqref{eq:pde21-main}.
Wówczas dla $E_S(t)=\int_{\R^d}u^\mathsf TS u\,\mathrm d\lambda_d$
 istnieje stała $C$, zależna tylko od
$\sigma,\Sigma$ i norm odpowiednich współczynników,
taka że
\[
 \sqrt{E_S(t)}\le e^{Ct}\left(\sqrt{E_S(0)}+
 \sqrt\Sigma\int_0^t\|f(r)\|_2\,\mathrm dr\right).
\]
\end{twierdzenie}
\begin{proof}
Różniczkując, otrzymujemy
$\partial_t(u^\mathsf TS u)=2u^\mathsf TS u_t+u^\mathsf TS_tu$.
Po wstawieniu równania i użyciu symetrii $SA_j$:
\[
 \partial_t(u^\mathsf TS u)+\sum_j\partial_j(u^\mathsf TSA_ju)
 =u^\mathsf T\Bigl(S_t+\sum_j\partial_j(SA_j)-SB-B^\mathsf TS\Bigr)u
 +2u^\mathsf TSf.
\]
Całkujemy po przestrzeni. Z $E_S\ge\sigma\|u\|_2^2$
wynika ograniczenie wyrazu kwadratowego przez $C E_S$,
a z nierówności Cauchy'ego--Schwarza w iloczynie skalarnym
wyznaczonym przez $S$ dostajemy
$\left|\int_{\R^d}u(t,x)^\mathsf TS(t,x)f(t,x)\,\mathrm d\lambda_d(x)\right|\le\sqrt{E_S(t)}\sqrt\Sigma\|f(t)\|_2$.
Argument z $\sqrt{E_S+\eps^2}$ i Gronwallem kończy dowód.
\end{proof}

\begin{kontrprzyklad}[Rzeczywiste charakterystyki bez stabilności]\label{ex:pde21-jordan}
Niech $U=(u,v)^\mathsf T$ i rozważmy na
$[0,T]\times\R$ układ $u_t=v_x$, $v_t=0$.
Macierz główna w zapisie $U_t+AU_x=0$ jest równa
$A=\begin{pmatrix}0&-1\\0&0\end{pmatrix}$.
Jej jedyną wartością własną jest $0$, lecz nie ma bazy
wektorów własnych. Wybierzmy niezerową funkcję
$\phi\in C_c^\infty(\R)$ i dane
$u(0,x)=0$, $v(0,x)=n^{-1}\phi(x)\sin(nx)$.
Równanie ma jawne rozwiązanie
\[
 v(t,x)=n^{-1}\phi(x)\sin(nx),\qquad
 u(t,x)=t\bigl(\phi(x)\cos(nx)+n^{-1}\phi'(x)\sin(nx)\bigr).
\]
Norma początkowa $\|(u(0),v(0))\|_{L^2(\R;\R^2)}$
nie przekracza $n^{-1}\|\phi\|_2$ i dąży do zera.
Z kolei dla każdego ustalonego $t>0$,
na mocy lematu Riemanna--Lebesgue'a,
$\|u(t)\|_2^2\to t^2\|\phi\|_2^2/2>0$:
całka $\int_{\R}|\phi(x)|^2\cos^2(nx)\,
\mathrm d\lambda_1(x)$ dąży do $\|\phi\|_2^2/2$,
a składniki z $n^{-1}\phi'$ znikają w $L^2$.
Nie istnieje zatem ograniczony operator ewolucji
$L^2(\R;\R^2)\to L^2(\R;\R^2)$ w żadnej chwili $t>0$.
Przykład na torusie powstaje prościej, biorąc $\phi\equiv1$,
lecz funkcji okresowych nie wolno traktować jako elementów
$L^2(\R)$.
\end{kontrprzyklad}

\section{Zmiennowspółczynnikowe równanie falowe}
Niech $G(x)=G(x)^\mathsf T$ będzie macierzą $d\times d$,
$G\in W^{1,\infty}(\R^d)$ oraz
$\lambda|\xi|^2\le\xi^\mathsf TG(x)\xi\le\Lambda|\xi|^2$.
Rozważamy
\begin{equation}\label{eq:pde21-wavevariable}
 u_{tt}-\operatorname{div}(G(x)\nabla u)=f.
\end{equation}
\begin{twierdzenie}[Energia równania falowego]\label{thm:pde21-wave-energy}
Dla dostatecznie gładkiego rozwiązania~\eqref{eq:pde21-wavevariable},
które wraz z pochodnymi dostatecznie szybko zanika, energia
\[
 \mathcal E(t)=\frac12\int_{\R^d}
 (|u_t|^2+\nabla u^\mathsf TG\nabla u)\,\mathrm d\lambda_d
\]
spełnia $\mathcal E'(t)=\int_{\R^d}f u_t\,\mathrm d\lambda_d$.
Jeżeli $f=0$, energia jest stała. Jeśli $f$ znika we wstecznym
stożku o prędkości $c\ge\sqrt\Lambda$, a dane $u_0,u_1$ znikają
na jego podstawie, to $u$ znika w całym wnętrzu stożka.
\end{twierdzenie}
\begin{proof}
Mnożymy równanie przez $u_t$ i całkujemy.
Ponieważ $G$ jest niezależna od $t$, całkowanie przez części daje
$\int_{\R^d}u_{tt}u_t\,\mathrm d\lambda_d(x)
=\frac12\frac{\mathrm d}{\mathrm dt}\|u_t\|_2^2$
oraz
\[
-\int_{\R^d}\operatorname{div}(G\nabla u)u_t\,\mathrm d\lambda_d(x)
=\int_{\R^d}G\nabla u\cdot\nabla u_t\,\mathrm d\lambda_d(x)
=\frac12\frac{\mathrm d}{\mathrm dt}
\int_{\R^d}\nabla u^\mathsf TG\nabla u\,\mathrm d\lambda_d(x).
\]
To daje tożsamość energii.

Dla skończonej prędkości propagacji używamy lokalnej gęstości
$e=(|u_t|^2+\nabla u^\mathsf TG\nabla u)/2$
i strumienia $J=-u_tG\nabla u$. Z równania otrzymujemy
$\partial_te+\operatorname{div}J=f u_t$.
Na brzegu ruchomej kuli mamy
$|J\cdot\nu|\le |u_t|\,|G\nabla u\cdot\nu|
\le\sqrt\Lambda |u_t|\sqrt{\nabla u^\mathsf TG\nabla u}
\le\sqrt\Lambda\,e$.
Zatem przy ruchu brzegu z szybkością $-c$ całka energii
w malejącej kuli jest nierosnąca, jeśli $f=0$ wewnątrz.
Zerowe dane $u_1$ i $\nabla u_0$ oznaczają zerową energię
na podstawie. Stąd $u_t=\nabla u=0$ we wnętrzu stożka;
ponieważ $u_0=0$ na podstawie, otrzymujemy tam $u=0$.
\end{proof}

\section{Warunki brzegowe: dlaczego nie można narzucić dowolnych danych?}
Na obszarze ograniczonym energia układu symetrycznego spełnia
prawo~\eqref{eq:pde21-flux}, ale po całkowaniu pojawia się
strumień przez brzeg $\partial\Omega$:
$-\int_{\partial\Omega}u^\mathsf TA_\nu u\,\mathrm d\mathcal H^{d-1}$.
W szczególności dla zagadnienia Cauchy'ego na półprostej
$x>0$ z macierzą diagonalną $A=\operatorname{diag}(a_1,\ldots,a_m)$
warunki brzegowe podaje się tylko na składowych z $a_k>0$,
które wchodzą do obszaru przez $x=0$; dla $a_k<0$
wartości brzegowe są wyznaczane przez wnętrze.

\begin{przyklad}[Dwa kierunki transportu na półprostej]
Dla $a>0$ rozwiązanie $u_t+au_x=0$ na $x>0$ wymaga
jednej wartości brzegowej $u(t,0)=g(t)$: charakterystyki
wychodzące z brzegu wchodzą do obszaru. Natomiast dla $a<0$
charakterystyki biegną ku brzegowi i wartość $u(t,0)$
wynika już z danej początkowej. Dodatkowe narzucenie
niezależnego $g$ może prowadzić do braku rozwiązania.
\end{przyklad}


\begin{twierdzenie}[Dysypatywne dane brzegowe na półprostej]
\label{thm:pde21-halfline-energy}
Niech $A=A^{\mathsf T}\in\R^{m\times m}$ będzie macierzą stałą,
$T>0$, $U\in C^1([0,T]\times[0,\infty);\R^m)$ i
$F\in C([0,T]\times[0,\infty);\R^m)$ spełniają
$U_t+A U_x=F$ na $(0,T)\times(0,\infty)$.
Załóżmy, że $U$, $F$ i ich pochodne mają zanik pozwalający
na całkowanie przez części przy $x\to\infty$ oraz że
$U(t,\cdot),F(t,\cdot)\in L^2(0,\infty)$.
Niech $P_+$ będzie rzutem ortogonalnym na sumę przestrzeni
własnych $A$ odpowiadających dodatnim wartościom własnym,
a $\lambda_+=\max(0,\lambda_{\max}(A))$. Jeżeli
$P_+U(t,0)=g(t)\in\operatorname{Ran}P_+$, to
\[
 \|U(t)\|_{L^2(0,\infty)}\le
 \|U(0)\|_{L^2(0,\infty)}+
 \sqrt{\lambda_+}\,\|g\|_{L^2(0,t)}+
 \int_0^t\|F(r)\|_{L^2(0,\infty)}\,\mathrm dr.
\]
W szczególności dla $g=F=0$ energia nie rośnie.
Nie narzuca się niezależnych danych na ujemnej podprzestrzeni
spektralnej $A$; warunek dla zerowej wartości własnej
wynika z równania zwyczajnego w każdym punkcie przestrzeni.
\end{twierdzenie}
\begin{proof}
Tożsamość lokalna dla macierzy stałej ma postać
$\partial_t|U|^2+\partial_x(U^{\mathsf T}AU)=2U^{\mathsf T}F$.
Po scałkowaniu po półprostej otrzymujemy
\[
 \frac{\mathrm d}{\mathrm dt}\|U(t)\|_2^2
 =U(t,0)^{\mathsf T}AU(t,0)+
 2\langle U(t),F(t)\rangle.
\]
Rozkład spektralny $A=A_+-A_-$, gdzie $A_\pm\ge0$
i ich nośniki spektralne są ortogonalne, daje
$U^{\mathsf T}AU\le\lambda_+|P_+U|^2
=\lambda_+|g|^2$ na $x=0$. Oznaczmy
$E(t)=\|U(t)\|_2^2$, $b(t)=\|F(t)\|_2$ oraz
\[
 a(t)=E(0)+\lambda_+\int_0^t|g(r)|^2\,\mathrm dr.
\]
Z nierówności Cauchy'ego--Schwarza dostajemy
$E'(t)\le a'(t)+2\sqrt{E(t)}\,b(t)$.
Dla $\delta>0$ wprowadzamy
\[
 z_\delta(t)=\left(a(t)+\delta^2+
 2\int_0^t\sqrt{E(r)}b(r)\,\mathrm dr\right)^{1/2}.
\]
Wtedy $\sqrt{E(t)}\le z_\delta(t)$ i
$z_\delta(t)\ge\sqrt{a(t)+\delta^2}$. Dlatego
\[
 z_\delta'(t)=\frac{a'(t)+2\sqrt{E(t)}b(t)}
 {2z_\delta(t)}
 \le\frac{a'(t)}{2\sqrt{a(t)+\delta^2}}+b(t).
\]
Całkując od zera do $t$, wykorzystujemy tożsamość
$\frac{\mathrm d}{\mathrm dt}\sqrt{a(t)+\delta^2}
=a'(t)/(2\sqrt{a(t)+\delta^2})$.
Otrzymujemy
$\sqrt{E(t)}\le\sqrt{a(t)+\delta^2}+
\int_0^t b(r)\,\mathrm dr$. Granica $\delta\downarrow0$
oraz nierówność $\sqrt{r+s}\le\sqrt r+\sqrt s$
kończą dowód. Wartości własne dodatnie odpowiadają
charakterystykom wchodzącym do $x>0$ przez $x=0$,
a ujemne --- wychodzącym, na których nie można niezależnie
zadawać śladu. Dla zerowej wartości własnej projekcja
równania ma postać $(P_0U)_t=P_0F$.
\end{proof}


\section{Podsumowanie}
Symetria części głównej daje tożsamość energii, z której wynikają
jednoznaczność, stabilność i skończona prędkość propagacji.
Po zróżniczkowaniu równania komutatory nie zawierają pochodnych
wyższego rzędu niż $s$; dzięki temu otrzymujemy poprawność
w $H^s$ bez utraty pochodnych. Nie każda macierz o rzeczywistych
wartościach własnych ma te własności; pokazuje to blok Jordana.
W następnym rozdziale porównamy tę metodę z innym
mechanizmem ewolucji: unitarnością dla równania Schrödingera.
Nie jest ono ogólnym układem symetryczno-hiperbolicznym
pierwszego rzędu w przestrzeni, toteż stożkowa zasada
propagacji nie przechodzi na nie bezpośrednio.


\chapter[Równanie Schrödingera]{Równanie Schrödingera: hamiltoniany samosprzężone, ewolucja unitarna i dyspersja}
\label{chap:pde22}

\section{Od równania falowego do ewolucji unitarnej}
Niech $d$ będzie dodatnią liczbą całkowitą. W tym
rozdziale niewiadoma jest zespolona:
$u\colon\R\times\R^d\to\C$. Przyjmujemy jednostki, w których
$\hbar=1$ oraz masa cząstki jest równa $1/2$. Równanie
Schrödingera przyjmuje wówczas postać
\begin{equation}\label{eq:pde22-schrod}
 i\partial_tu=Hu,\qquad H=-\Delta+V(x),\qquad u(0)=u_0.
\end{equation}
Potencjał $V$ ma być rzeczywisty, chyba że wyraźnie podamy inne
założenia. Symbolem $H$ oznaczamy \emph{operator wraz z jego
dziedziną}, a nie tylko formalne wyrażenie różniczkowe.
Przez $\mathcal H$ oznaczamy przestrzeń Hilberta
$L^2(\R^d;\C)$. Przyjmujemy w niej iloczyn skalarny
$\langle f,g\rangle=\int_{\R^d}f(x)\overline{g(x)}\,\mathrm d\lambda_d(x)$,
liniowy w pierwszym argumencie.

Istotna różnica względem równania ciepła jest już widoczna w
transformacie Fouriera. Dla swobodnego równania $i u_t=-\Delta u$
mnożnik wynosi $e^{-it|\xi|^2}$: jego moduł jest stale równy $1$.
Nie ma więc tłumienia wysokich częstotliwości ani analogicznego
do ciepła wygładzania w skali Sobolewa. Za to oscylacja fazy
prowadzi do \emph{dyspersji} przestrzennej.

\begin{uwaga}[Porównanie trzech ewolucji]
\label{rem:pde-v43-evolution-contrast}
Równanie ciepła (rozdziały~\ref{chap:pde18} i~\ref{chap:pde20})
prowadzi dla nieujemnego operatora samosprzężonego
do mnożnika spektralnego $e^{-t\lambda}$, $t>0$,
który tłumi wysokie częstotliwości spektralne
(nie tylko wartości własne w przypadku widma ciągłego). Układ hiperboliczny z rozdziału~\ref{chap:pde21}
ma skończoną prędkość propagacji przy odpowiednich
założeniach. Swobodne równanie Schrödingera zachowuje
normę $L^2$, ale nie ma ogólnej skończonej prędkości
propagacji; samo zachowanie normy nie implikuje
zwartości ani wygładzania. Wszystkie trzy stwierdzenia
odnoszą się do \emph{różnych} operatorów i klas rozwiązań.
\end{uwaga}

\section{Samosprzężoność swobodnego hamiltonianu}
Używamy unitarnej normalizacji transformaty Fouriera z
rozdziału~\ref{chap:pde-fourier-Rd}:
\[
 \widehat f(\xi)=(2\pi)^{-d/2}
 \int_{\R^d}e^{-\mathrm i x\cdot\xi}f(x)\,\mathrm d\lambda_d(x).
\]
Wzór ten przyjmujemy początkowo dla funkcji
$f\in\mathcal S(\R^d)$ z przestrzeni Schwartza.
Twierdzenie Plancherela~\ref{thm:pde8-plancherel}
daje unitarne przedłużenie $\mathcal F$ na $\mathcal H$;
przez dualność przedłużamy je również na przestrzeń
dystrybucji temperowanych $\mathcal S'(\R^d)$.
Dla $s\in\R$ przyjmujemy przestrzeń Sobolewa
\[
 H^s(\R^d;\C)=
 \{f\in\mathcal S'(\R^d):
 (1+|\xi|^2)^{s/2}\widehat f\in L^2(\R^d;\C)\},
\]
z normą
$\|f\|_{H^s}^2=\int_{\R^d}(1+|\xi|^2)^s
|\widehat f(\xi)|^2\,\mathrm d\lambda_d(\xi)$.
Warunek oznacza, że ważona transformata jest
reprezentowana przez funkcję z $L^2$.
Dla całkowitego $s\ge0$ jest to przestrzeń słabo
różniczkowalnych funkcji z pochodnymi do rzędu $s$
w $L^2$, z równoważną normą: wynika to z
$\widehat{\partial^\alpha f}=(\mathrm i\xi)^\alpha\widehat f$
i porównania wielomianu
$(1+|\xi|^2)^s$ z sumą $\sum_{|\alpha|\le s}|\xi^\alpha|^2$.
Norma Fouriera będzie w tym rozdziale normą domyślną.

\begin{twierdzenie}[Swobodny hamiltonian]\label{thm:pde22-free-sa}
Niech $H_0\colon D(H_0)\subset L^2(\R^d;\C)\to L^2(\R^d;\C)$
będzie operatorem $H_0u=-\Delta u$ z dziedziną
$D(H_0)=H^2(\R^d;\C)$. Wtedy $H_0$ jest samosprzężony,
nieujemny i domknięty, a $\mathcal F H_0\mathcal F^{-1}$
jest operatorem mnożenia przez $|\xi|^2$ na dziedzinie
$\{g\in L^2:|\xi|^2g\in L^2\}$.
\end{twierdzenie}
\begin{proof}
Równość $\widehat{-\Delta u}=|\xi|^2\widehat u$
wynika z całkowania przez części dla $u\in\mathcal S$
i obowiązuje także dystrybucyjnie.
Dla $g\in L^2$ warunek $|\xi|^2g\in L^2$
jest równoważny warunkowi $(1+|\xi|^2)g\in L^2$:
implikacja w jedną stronę wynika z nierówności
trójkąta, a w drugą z $|\xi|^2\le1+|\xi|^2$.
Zatem $\mathcal F$ przenosi $H^2$ dokładnie na
dziedzinę operatora mnożenia
$Mg=mg$, $m(\xi)=|\xi|^2$.
Dziedzina ta jest gęsta w $L^2$, gdyż zawiera
wszystkie funkcje z $L^2$ o ograniczonym nośniku,
a $g\chi_{\{|\xi|\le n\}}\to g$ w $L^2$.
Operator jest symetryczny, ponieważ $m$ jest
rzeczywista i dla $g,h\in D(M)$ mamy
$\langle Mg,h\rangle=\langle g,Mh\rangle$.
W szczególności $M\subset M^*$.
Jeżeli $h\in D(M^*)$, to istnieje $k\in L^2$
takie, że $\langle Mg,h\rangle=\langle g,k\rangle$
dla wszystkich $g\in D(M)$. Dla każdego $R>0$
funkcja $g=\chi_{\{|\xi|\le R\}}(mh-k)$ należy
do $D(M)$, bo $m$ jest ograniczona na tej kuli.
Po podstawieniu otrzymujemy
\[
 \int_{\{|\xi|\le R\}}|mh-k|^2\,\mathrm d\lambda_d(\xi)=0.
\]
Biorąc kolejno $R\in\N$, dostajemy $mh=k$ prawie
wszędzie na $\R^d$. Zatem $h\in D(M)$ i $M^*h=Mh$,
co dowodzi przeciwnej inkluzji i samosprzężoności.
Nieujemność wynika z
\[
 \langle Mg,g\rangle
 =\int_{\R^d}m(\xi)|g(\xi)|^2\,\mathrm d\lambda_d(\xi)\ge0.
\]
Domkniętość można zobaczyć bezpośrednio: ze
zbieżności $g_n\to g$ i $mg_n\to h$ w $L^2$
wybieramy wspólny podciąg zbieżny prawie wszędzie
w obu przypadkach. W granicy otrzymujemy $h=mg$
prawie wszędzie, więc $g\in D(M)$ i $Mg=h$.
Unitarność $\mathcal F$ przenosi wszystkie te
własności na $H_0$.
\end{proof}

\begin{uwaga}[Dziedzina jest częścią twierdzenia]
Formalny laplasjan na $C_c^\infty(\Omega)$ dla obszaru
z brzegiem jest zwykle jedynie symetryczny, a nie samosprzężony.
Realizacje Dirichleta i Neumanna mogą dawać różne ewolucje.
Na $\R^d$ wybór $H^2$ jest kanoniczną samosprzężoną
realizacją swobodnego hamiltonianu. Jest ona
domknięciem operatora $-\Delta$ z dziedziną
$C_c^\infty(\R^d)$: jego norma grafowa
$(\|u\|_2^2+\|\Delta u\|_2^2)^{1/2}$ jest
równoważna normie $H^2$, ponieważ
$1+|\xi|^4\le(1+|\xi|^2)^2\le2(1+|\xi|^4)$,
a $C_c^\infty$ jest gęste w $H^2$.
W szczególności ta realizacja nie wymaga wyboru
dodatkowych warunków brzegowych.
\end{uwaga}

\section{Ewolucja unitarna i poprawność zagadnienia Cauchy'ego}
\begin{twierdzenie}[Grupa swobodnego Schrödingera]\label{thm:pde22-group}
Dla $t\in\R$ i $f\in L^2(\R^d;\C)$ niech
$U_0(t)f=\mathcal F^{-1}(e^{-it|\xi|^2}\widehat f)$.
Wtedy $(U_0(t))_{t\in\R}$ jest silnie ciągłą grupą
unitarną. Dla każdego $s\in\R$ ten sam wzór definiuje
silnie ciągłą grupę izometrii na $H^s$ (także
dla $s<0$, gdy $f$ nie musi należeć do $L^2$).
Dla każdego $f\in H^s$ zachodzi
$\|U_0(t)f\|_{H^s}=\|f\|_{H^s}$ oraz
$U_0(t)H^s=H^s$. Jej generator na $\mathcal H$
jest równy $-\mathrm iH_0$ z dziedziną $H^2$. Jeśli $f\in H^2$, to
$u(t)=U_0(t)f$ należy do
$C(\R;H^2)\cap C^1(\R;L^2)$ i spełnia
$i u_t=-\Delta u$. Jest to jedyne rozwiązanie w tej klasie
z daną $u(0)=f$.
\end{twierdzenie}
\begin{proof}
Ponieważ $|e^{-it|\xi|^2}|=1$, twierdzenie Plancherela
daje zachowanie iloczynu skalarnego w $L^2$,
a definicja normy Fouriera daje zachowanie normy
$H^s$ dla każdego $s\in\R$. Dla $s<0$ mnożnik
działa również na dystrybucje temperowane,
ponieważ jest gładki, a każda jego pochodna
ma co najwyżej wzrost wielomianowy.
Tożsamość mnożników
$e^{-i(t+r)|\xi|^2}=e^{-it|\xi|^2}e^{-ir|\xi|^2}$
daje prawo grupowe. Ponieważ $U_0(0)=I$,
operator $U_0(-t)$ jest odwrotnością $U_0(t)$,
więc izometrie są surjektywne.
Ponadto
$\|(U_0(t)-I)f\|_2^2=
\int_{\R^d}|e^{-it|\xi|^2}-1|^2|\widehat f(\xi)|^2\,\mathrm d\lambda_d(\xi)$.
Funkcja podcałkowa dąży punktowo do zera i jest ograniczona
przez $4|\widehat f|^2$; twierdzenie o zbieżności zmajoryzowanej
zapewnia silną ciągłość w zerze, a prawo grupowe w każdym $t$.
Tę samą argumentację można wykonać z wagą
$(1+|\xi|^2)^s$.
Jeśli $f\in H^2$, iloraz różnicowy w dziedzinie Fouriera
ma postać
$\frac{e^{-i(t+h)|\xi|^2}-e^{-it|\xi|^2}}h\widehat f$.
Dąży do $-i|\xi|^2e^{-it|\xi|^2}\widehat f$ i ma moduł
nie większy niż $|\xi|^2|\widehat f|$,
na mocy $|e^{\mathrm i y}-1|\le |y|$ dla $y\in\R$.
Kwadrat modułu różnicy ilorazu i jego granicy
ma całkowalną majorantę $4|\xi|^4|\widehat f|^2$.
Twierdzenie o zbieżności zmajoryzowanej daje
pochodną czasową $u_t=-\mathrm iH_0u$ w $L^2$.
Pochodna jest ciągła, gdyż $H_0:H^2\to L^2$
jest ograniczony i $u$ jest ciągła w $H^2$.
Odwrotnie, jeśli ilorazy
$(U_0(h)f-f)/h$ mają granicę w $L^2$ przy $h\to0$,
to ich transformaty wzdłuż pewnego ciągu $h_n\to0$
mają podciąg zbieżny prawie wszędzie do transformaty
tej granicy. Granica punktowa jest koniecznie
$-\mathrm i|\xi|^2\widehat f$, więc
$|\xi|^2\widehat f\in L^2$ i $f\in H^2$.
To identyfikuje dokładną dziedzinę generatora.
Dla dwóch rozwiązań
różnica $w$ spełnia to samo równanie z zerowymi danymi.
Różniczkując normę, uzyskujemy
$\frac{\mathrm d}{\mathrm dt}\|w\|_2^2=
2\operatorname{Re}\langle-iH_0w,w\rangle=0$,
ponieważ $\langle H_0w,w\rangle$ jest rzeczywiste.
Zatem $w=0$.
\end{proof}

\begin{stwierdzenie}[Granica regularności]
Dla każdej liczby $s\in\R$ i $t\in\R$ grupa $U_0(t)$
jest izometrią na $H^s(\R^d)$, ale dla $t\ne0$ nie
działa w ogólności jako operator wygładzający
$H^s\to H^{s+\eps}$ dla żadnego $\eps>0$.
\end{stwierdzenie}
\begin{proof}
Gdyby $U_0(t)f\in H^{s+\eps}$ dla każdego $f\in H^s$,
to ponieważ mnożnik odwrotny $U_0(-t)$ zachowuje
$H^{s+\eps}$, mielibyśmy $f=U_0(-t)U_0(t)f\in H^{s+\eps}$
dla wszystkich $f\in H^s$, co jest fałszem.
Dla ścisłości wybieramy dystrybucję temperowaną $f$
o transformacie danej funkcją
\[
 \widehat f(\xi)=
 \frac{\chi_{\{|\xi|\ge2\}}(\xi)}
 {(1+|\xi|^2)^{s/2+d/4}\log(2+|\xi|)}.
\]
Funkcja ta jest lokalnie całkowalna i ma co
najwyżej wzrost wielomianowy, więc określa
dystrybucję temperowaną. We współrzędnych
biegunowych kwadrat normy $H^s$ jest skończony,
gdyż całka radialna jest porównywalna z
$\int_2^\infty(r(\log r)^2)^{-1}\,\mathrm d\lambda_1(r)$,
która jest zbieżna przez podstawienie $y=\log r$.
Dla normy $H^{s+\eps}$ całka radialna jest
porównywalna z
$\int_2^\infty r^{2\eps-1}(\log r)^{-2}\,\mathrm d\lambda_1(r)$.
Jest rozbieżna, ponieważ dla dostatecznie dużego $r$
mamy $(\log r)^2\le r^\eps$, a całka
z $r^{\eps-1}$ jest rozbieżna.
Stąd $f\in H^s\setminus H^{s+\eps}$, a moduł
mnożnika $U_0(t)$ równy $1$ pokazuje również
bezpośrednio, że $U_0(t)f$ nie należy do
$H^{s+\eps}$.
\end{proof}

\begin{twierdzenie}[Zasada Duhamela]\label{thm:pde22-duhamel}
Niech $0<T<\infty$, $f\in L^1(0,T;\mathcal H)$ i
$u_0\in\mathcal H$. Funkcja
\[
 u(t)=U_0(t)u_0-i\int_0^tU_0(t-r)f(r)\,\mathrm d\lambda_1(r)
\]
należy do $C([0,T];\mathcal H)$ i jest jedynym
rozwiązaniem łagodnym $\mathrm i u_t=H_0u+f$
z warunkiem $u(0)=u_0$ (rozwiązanie łagodne
definiujemy jako funkcję ciągłą w $\mathcal H$
spełniającą powyższe równanie całkowe).
Spełnia również równanie w sensie dystrybucji
o wartościach w $H^{-2}$.
Spełnia oszacowanie
$\|u(t)\|_2\le\|u_0\|_2+\int_0^t\|f(r)\|_2\,\mathrm d\lambda_1(r)$.
Jeżeli ponadto $u_0\in H^2$ i
$f\in C([0,T];H^2)$, to rozwiązanie należy do
$C([0,T];H^2)\cap C^1([0,T];L^2)$ i spełnia równanie w $L^2$.
\end{twierdzenie}
\begin{proof}
Funkcja $r\mapsto U_0(-r)f(r)$ jest silnie
mierzalna w $\mathcal H$: dla funkcji prostych
wynika to z silnej ciągłości grupy, a dla ogólnej
$f$ przechodzimy do granicy prawie wszędzie
przez aproksymację funkcjami prostymi i izometryczność.
Jej norma jest równa $\|f(r)\|_2$,
więc jest całkowalna w sensie Bochnera.
Korzystając z prawa grupowego oraz ciągłości
operatora $U_0(t)$, zapisujemy rozwiązanie jako
\[
 u(t)=U_0(t)w(t),\qquad
 w(t)=u_0-\mathrm i\int_0^t U_0(-r)f(r)\,\mathrm d\lambda_1(r).
\]
Funkcja $w$ jest absolutnie ciągła w $\mathcal H$,
a $w'=-\mathrm i U_0(-t)f(t)$ prawie wszędzie.
Dla $t_n\to t$ mamy
\[
 \|U_0(t_n)w(t_n)-U_0(t)w(t)\|_2
 \le\|w(t_n)-w(t)\|_2+
       \|(U_0(t_n)-U_0(t))w(t)\|_2\to0.
\]
To daje ciągłość na całym $[0,T]$, warunek początkowy
i, przez nierówność dla normy całki Bochnera,
oszacowanie z tezy.

Operator $-\Delta$ ma ograniczone przedłużenie
$H_0:\mathcal H\to H^{-2}$, ponieważ
$|\xi|^2/(1+|\xi|^2)\le1$. Dla każdego $g\in\mathcal H$
ilorazy różnicowe $U_0(t)g$ zbiegają w $H^{-2}$
do $-\mathrm i H_0U_0(t)g$ (ten sam argument
zmajoryzowanej zbieżności z wagą $(1+|\xi|^2)^{-2}$).
Ponadto
$\|(U_0(t+h)-U_0(t))g\|_{H^{-2}}\le |h|\|g\|_2$.
Rozkładając różnicę wartości $u$ na składnik
z różnicą $w$ i składnik z różnicą $U_0$,
otrzymujemy absolutną ciągłość w $H^{-2}$:
pierwszy składnik jest kontrolowany przez
$\|w(t+h)-w(t)\|_2$, drugi przez $|h|\sup_t\|w(t)\|_2$.
W chwilach różniczkowalności $w$ w $\mathcal H$
reguła iloczynu, wynikająca z tego samego rozkładu,
daje w $H^{-2}$
\[
 u'=-\mathrm iH_0u+U_0(t)w'
    =-\mathrm i(H_0u+f)
\]
prawie wszędzie. Stąd równanie dystrybucyjne.

Przy mocniejszych danych $w$ jest klasy $C^1$
w $H^2$, dzięki silnej ciągłości grupy na $H^2$
i ciągłości $f$. Prawo grupowe daje ciągłość
$u$ w $H^2$. Rozkład ilorazu różnicowego teraz
w $\mathcal H$ i identyfikacja generatora dają
$u'=-\mathrm i(H_0u+f)$ w $\mathcal H$ we wszystkich
chwilach (na końcach jednostronnie).
Prawa strona jest ciągła w $\mathcal H$, co
dowodzi żądanej regularności.
Jednoznaczność łagodna wynika z definicji, a
jednoznaczność w klasie silnej z zachowania normy
dla różnicy dwóch rozwiązań.
\end{proof}

\section{Potencjały rzeczywiste i zachowanie energii}
\begin{twierdzenie}[Ograniczony potencjał rzeczywisty]\label{thm:pde22-potential}
Niech $V\in L^\infty(\R^d;\R)$ i niech $V$ działa jako
operator mnożenia na $L^2(\R^d;\C)$. Wtedy
$H=H_0+V$, $D(H)=H^2(\R^d)$, jest samosprzężony.
Grupa $U_V(t)=e^{-itH}$ jest unitarna i silnie ciągła.
Dla $u_0\in H^2$ rozwiązanie $u(t)=U_V(t)u_0$
zachowuje zarówno normę $L^2$, jak i energię
$E(u)=\langle Hu,u\rangle=
\int_{\R^d}(|\nabla u|^2+V|u|^2)\,\mathrm d\lambda_d$.
\end{twierdzenie}
\begin{proof}
Operator $V$ jest ograniczony i samosprzężony.
Suma $H_0+V$ jest symetryczna na $H^2$.
Pokażemy jej samosprzężoność bez przyjmowania
ogólnego twierdzenia o perturbacjach.
Dla $a>\|V\|_\infty$ mamy
$\|(H_0\pm ia)^{-1}\|_{L^2\to L^2}\le a^{-1}$,
co wprost wynika z mnożnika $(|\xi|^2\pm ia)^{-1}$.
Faktoryzacja
$H\pm ia=(I+V(H_0\pm ia)^{-1})(H_0\pm ia)$
pokazuje, że $H\pm ia$ są surjektywne, gdyż
$\|V(H_0\pm ia)^{-1}\|<1$ i odwrotność pierwszego
czynnika otrzymujemy z szeregu Neumanna.
Symetryczny operator o obu surjektywnych operatorach
$H\pm ia$ jest samosprzężony: istotnie dla
$z\in D(H^*)$ i $h=(H^*+ia)z$ dobieramy
$y\in D(H)$ takie, że $(H+ia)y=h$.
Wtedy $(H^*+ia)(z-y)=0$, a
$\ker(H^*+ia)=\operatorname{Ran}(H-ia)^\perp=\{0\}$;
zatem $z=y\in D(H)$. To dowodzi $H=H^*$.
Z twierdzenia spektralnego dla operatorów samosprzężonych
wynika istnienie silnie ciągłej grupy $e^{-itH}$ oraz
$U_V(t)D(H)=D(H)$, $HU_V(t)=U_V(t)H$.
Unitarność daje zachowanie normy, a ostatnia komutacja
prowadzi do
$\langle HU_V(t)u_0,U_V(t)u_0\rangle=
\langle Hu_0,u_0\rangle$. Całkowanie przez części
jest uzasadnione dla $u_0\in H^2$ i $V$ ograniczone.
\end{proof}

\begin{stwierdzenie}[Lokalny bilans prawdopodobieństwa]
Niech $V\in L^\infty(\R^d;\R)$ oraz
$u_0\in H^2(\R^d)$. Dla rozwiązania z
Twierdzenia~\ref{thm:pde22-potential} zachodzi w sensie
dystrybucji na $\R\times\R^d$
\[
 \partial_t|u|^2+\operatorname{div}j=0,
 \qquad j=2\operatorname{Im}(\overline u\,\nabla u).
\]
\end{stwierdzenie}
\begin{proof}
Ponieważ $u_t=-iHu=i\Delta u-iVu$,
$\partial_t|u|^2=2\operatorname{Re}(\overline u\,u_t)
=-2\operatorname{Im}(\overline u\,\Delta u)$.
Wyraz $-iV|u|^2$ nie wnosi wkładu do części rzeczywistej,
bo $V$ jest rzeczywisty. Z kolei
$\operatorname{Im}(\nabla\overline u\cdot\nabla u)=0$,
więc $\operatorname{Im}(\overline u\,\Delta u)
=\operatorname{div}\operatorname{Im}(\overline u\,\nabla u)$.
Tożsamość jest poprawna dla gładkich funkcji, a następnie
w sensie dystrybucji przez aproksymację w $H^2$.
\end{proof}


\begin{twierdzenie}[Zależny od czasu, ograniczony potencjał rzeczywisty]
\label{thm:pde22-time-potential}
Niech $T>0$, $V\in L^\infty((0,T)\times\R^d;\R)$,
$F\in L^1(0,T;L^2(\R^d;\C))$ oraz
$u_0\in L^2(\R^d;\C)$. Istnieje dokładnie jedna funkcja
$u\in C([0,T];L^2(\R^d;\C))$ spełniająca formułę Duhamela
\[
 u(t)=U_0(t)u_0-i\int_0^tU_0(t-s)
       \bigl(V(s,\cdot)u(s)+F(s)\bigr)\,\mathrm ds.
\]
Zachodzi oszacowanie
$\|u(t)\|_2\le\|u_0\|_2+\int_0^t\|F(s)\|_2\,\mathrm ds$.
Jeżeli $F=0$, to $\|u(t)\|_2=\|u_0\|_2$ dla każdego
$t\in[0,T]$, mimo że operator $H(t)=-\Delta+V(t,\cdot)$
może zależeć od czasu. Nie twierdzimy, że funkcja ta jest
rozwiązaniem silnym o wartościach w $H^2$.
\end{twierdzenie}
\begin{proof}
Stosujemy obraz oddziaływania $w(t)=U_0(-t)u(t)$.
Szukana równość jest równoważna równaniu całkowemu
\[
 w(t)=u_0-i\int_0^t\bigl(K(s)w(s)+G(s)\bigr)\,\mathrm ds,
 \quad K(s)=U_0(-s)M_{V(s)}U_0(s),\quad
 G(s)=U_0(-s)F(s).
\]
Operator $M_{V(s)}$ jest mnożeniem przez rzeczywistą
funkcję $V(s,\cdot)$. Dla prawie każdego $s$ operator $K(s)$
jest ograniczony, samosprzężony i
$\|K(s)\|\le\|V\|_\infty$. Funkcja
$s\mapsto K(s)z$ jest silnie mierzalna dla każdego $z\in L^2$;
wynika to z mierzalności $V$, separowalności $L^2$ i silnej
ciągłości $U_0(s)$.

Iteracje Picarda równania całkowego zbiegają jednostajnie
w $C([0,T];L^2)$: różnica kolejnych iteracji jest
ograniczona przez stałą pomnożoną przez
$(\|V\|_\infty T)^n/n!$. Dla dwóch rozwiązań nierówność
Gronwalla daje jednoznaczność. Ponadto
$w\in W^{1,1}(0,T;L^2)$ i prawie wszędzie
$w'=-iKw-iG$. Dla naszego iloczynu skalarnego, liniowego
w pierwszym argumencie, samosprzężoność $K$ daje
\[
 \frac{\mathrm d}{\mathrm dt}\|w(t)\|_2^2
 =2\operatorname{Re}\langle-iK(t)w(t)-iG(t),w(t)\rangle
 =2\operatorname{Re}\langle-iG(t),w(t)\rangle.
\]
Przy $G=0$ norma jest stała. Ogólnie, regularizując
pierwiastek $\sqrt{\|w\|_2^2+\delta^2}$ i przechodząc
$\delta\downarrow0$, otrzymujemy
$\|w(t)\|_2\le\|u_0\|_2+\int_0^t\|G(s)\|_2\,\mathrm ds$.
Unitarność $U_0(t)$ oraz $\|G(s)\|_2=\|F(s)\|_2$
kończy dowód.
\end{proof}


\begin{uwaga}[Potencjały nieograniczone]
Potencjały Coulomba i oscylatora harmonicznego wymagają
innych twierdzeń o samosprzężoności lub konstrukcji
przez formy kwadratowe. Nie wynikają one z powyższego
twierdzenia, ponieważ nie należą do $L^\infty$.
\end{uwaga}

\section{Jądro swobodnego propagatora i oszacowanie dyspersyjne}
\begin{twierdzenie}[Wzór na propagator i zanik dyspersyjny]
\label{thm:pde22-disp}
Dla $t\in\R\setminus\{0\}$ i $f\in\mathcal S(\R^d)$ zachodzi
\begin{equation}\label{eq:pde22-kernel}
 U_0(t)f(x)=(4\pi it)^{-d/2}
 \int_{\R^d}\exp\!\left(\frac{i|x-y|^2}{4t}\right)
 f(y)\,\mathrm d\lambda_d(y),
\end{equation}
gdzie potęgę zespoloną wybiera się zgodnie z ciągłą
granicą regularizacji gaussowskiej. W szczególności
\begin{equation}\label{eq:pde22-L1Linf}
 \|U_0(t)f\|_\infty\le(4\pi|t|)^{-d/2}\|f\|_1.
\end{equation}
Dla $2\le r\le\infty$ i $r'$ sprzężonego do $r$
zachodzi ponadto
\[
 \|U_0(t)f\|_{L^r}\le C_{d,r}
 |t|^{-d(1/2-1/r)}\|f\|_{L^{r'}}
\]
dla $f\in L^{r'}\cap L^2$, z przedłużeniem do wskazanych
klas przez gęstość.
\end{twierdzenie}
\begin{proof}
Wstawiamy czynnik $e^{-\eps|\xi|^2}$ do mnożnika
Fouriera i otrzymujemy, dla $\eps>0$, splot z funkcją
\[
 K_{t,\eps}(x)=[4\pi(\eps+it)]^{-d/2}
 \exp\!\left(-\frac{|x|^2}{4(\eps+it)}\right).
\]
Formuła wynika przez obliczenie transformaty Gaussa,
początkowo dla rzeczywistego parametru dodatniego,
a następnie przez analityczną kontynuację do półpłaszczyzny
$\operatorname{Re}(\eps+it)>0$. W dziedzinie Fouriera
mamy $e^{-(\eps+it)|\xi|^2}\widehat f\to
 e^{-it|\xi|^2}\widehat f$ w $L^2$ na mocy dominacji.
Dla $f\in\mathcal S$ i ustalonego $t\ne0$ dominacja
w całce przestrzennej daje także zbieżność punktową
$K_{t,\eps}*f\to K_t*f$, gdzie
$K_t(x)=(4\pi it)^{-d/2}\exp(i|x|^2/(4t))$.
Zbieżności muszą opisywać tę samą dystrybucję,
co dowodzi wzoru~\eqref{eq:pde22-kernel}.
Ponieważ $|K_t(x)|=(4\pi|t|)^{-d/2}$,
nierówność~\eqref{eq:pde22-L1Linf} jest natychmiastowa.
Drugie oszacowanie wynika przez interpolację
Riesza--Thorina między $L^1\to L^\infty$
a unitarnością $L^2\to L^2$. Przy parametrze
$\theta\in[0,1]$ mamy $1/r=\theta/2$
i $1/r'=1-\theta/2$; potęga czasu to
$-d(1-\theta)/2=-d(1/2-1/r)$.
\end{proof}

\begin{przyklad}[Pakiet gaussowski i brak stożka propagacji]
Dla $u_0(x)=e^{-|x|^2/2}$ otrzymujemy
\[
 U_0(t)u_0(x)=(1+2it)^{-d/2}
 \exp\!\left(-\frac{|x|^2}{2(1+2it)}\right).
\]
Dla każdego $t\ne0$ funkcja ta jest niezerowa we wszystkich
punktach $\R^d$. Zatem równanie Schrödingera, w przeciwieństwie
do równania falowego z rozdziału 21, nie ma skończonej prędkości
propagacji w sensie zachowywania zwartego nośnika.
Ściśle rzecz biorąc, pakiet Gaussa ma już w chwili początkowej
niezwarty nośnik. Poprawny kontrprzykład ze zwartymi danymi
otrzymujemy jednak ze wzoru~\eqref{eq:pde22-kernel}:
jeżeli $f\in C_c^\infty(\R^d)$ jest niezerowa, to dla $t\ne0$
$U_0(t)f$ jest, po wyłączeniu czynnika
$e^{i|x|^2/(4t)}$, transformatą Fouriera funkcji
o zwartym nośniku. Jest więc rzeczywiście analityczna
jako funkcja zmiennej rzeczywistej i nie może znikać na
niepustym zbiorze otwartym bez identycznego zerowania.
Unitarność wyklucza identyczne zerowanie. W szczególności
$U_0(t)f$ nie może mieć zwartego nośnika.
\end{przyklad}

\section[Oszacowania Strichartza w dwóch wymiarach]{Oszacowania Strichartza w modelowym przypadku $d=2$}
Oszacowanie dyspersyjne dotyczy ustalonego czasu. Pokażemy,
jak z niego uzyskać normę czasowo-przestrzenną.
Wybrany przypadek $(q,r)=(4,4)$ jest wystarczająco bogaty,
aby zobaczyć całą ideę metody $TT^*$, a zarazem pozwala
uniknąć najtrudniejszych zagadnień końcowych.

\begin{lemat}[Oszacowanie całkowe w czasie]\label{lem:pde22-hls}
Dla każdego $g\in L^{4/3}(\R)$ zachodzi
\begin{equation}\label{eq:pde22-hls}
 \left\|\int_{\R}|t-s|^{-1/2}|g(s)|\,\mathrm ds
 \right\|_{L^4_t}\le C\|g\|_{L^{4/3}_t}.
\end{equation}
\end{lemat}
\begin{proof}
Niech $Mg(t)=\sup_{R>0}(2R)^{-1}\int_{t-R}^{t+R}|g(s)|\,\mathrm ds$
będzie operatorem maksymalnym Hardy'ego--Littlewooda.
Najpierw przypomnimy wraz z uzasadnieniem jego ograniczoność
na $L^p(\R)$ dla $p>1$. Dla skończenie wielu przedziałów
wybranych z rodziny $\{I:|I|^{-1}\int_I|g|>\lambda\}$
można wybrać rozłączne $I_j$ tak, aby ich trzykrotne
powiększenia pokrywały całą rozważaną skończoną rodzinę.
Wynika stąd oszacowanie słabego typu
$|\{Mg>\lambda\}|\le C\|g\|_1/\lambda$
przez aproksymację otwartego zbioru poziomicowego.
Rozdzielmy teraz $g=g_1+g_2$,
$g_1=g\,\chi_{\{|g|>\lambda/2\}}$ oraz
$g_2=g\,\chi_{\{|g|\le\lambda/2\}}$.
Wtedy $Mg_2\le\lambda/2$, więc
$|\{Mg>\lambda\}|\le C\lambda^{-1}
\int_{\{|g|>\lambda/2\}}|g(s)|\,\mathrm ds$.
Po pomnożeniu przez $p\lambda^{p-1}$, całkowaniu
w $\lambda$ i zastosowaniu twierdzenia Tonellego
dostajemy $\|Mg\|_p^p\le C_p\|g\|_p^p$.

Dla ustalonego $R>0$ rozbijamy całkę w~\eqref{eq:pde22-hls}
na $|t-s|\le R$ i $|t-s|>R$. W pierwszej części
rozkład na pierścienie dyadyczne daje oszacowanie
$C R^{1/2}Mg(t)$. W drugiej nierówność Höldera z
wykładnikami $4$ i $4/3$ daje
$C R^{-1/4}\|g\|_{4/3}$, gdyż
$\int_{|r|>R}|r|^{-2}\,\mathrm dr=2/R$.
Wybieramy $R>0$ tak, aby
$R^{3/4}=\|g\|_{4/3}/Mg(t)$, jeśli oba człony
są dodatnie (przypadki zerowe otrzymujemy granicznie).
W rezultacie lewa strona całki punktowej jest ograniczona
przez $C(Mg(t))^{1/3}\|g\|_{4/3}^{2/3}$.
Podnosząc do czwartej potęgi i całkując, otrzymujemy
\[
 \|I_{1/2}|g|\|_4^4\le C\|Mg\|_{4/3}^{4/3}
 \|g\|_{4/3}^{8/3}\le C\|g\|_{4/3}^4.
\]
To dowodzi tezy.
\end{proof}

\begin{twierdzenie}[Oszacowanie Strichartza $L^4_{t,x}$ w $\R^2$]
\label{thm:pde22-strichartz}
Niech $d=2$, $u_0\in L^2(\R^2)$ i
$F\in L^{4/3}(\R\times\R^2)$. Wtedy
\begin{align}
 \|U_0(t)u_0\|_{L^4(\R\times\R^2)}&\le C\|u_0\|_2,
 \label{eq:pde22-str-hom}\\
 \left\|\int_{s<t}U_0(t-s)F(s)\,\mathrm ds
 \right\|_{L^4(\R\times\R^2)}&\le
 C\|F\|_{L^{4/3}(\R\times\R^2)}.
 \label{eq:pde22-str-inhom}
\end{align}
Drugą całkę definiujemy początkowo dla gładkich danych
o zwartym nośniku, a następnie przez domknięcie.
\end{twierdzenie}
\begin{proof}
Oznaczmy $T\varphi(t)=U_0(t)\varphi$ dla
$\varphi\in\mathcal S(\R^2)$. Formalnym operatorem
sprzężonym jest
$T^*G=\int_{\R}U_0(-s)G(s)\,\mathrm ds$
(dla $G\in C_c^\infty(\R\times\R^2)$).
Stąd
$TT^*G(t)=\int_{\R}U_0(t-s)G(s)\,\mathrm ds$.
Interpolacja oszacowania dyspersyjnego dla $d=2$
między $L^1\to L^\infty$ i $L^2\to L^2$ daje
$\|U_0(t-s)G(s)\|_{L^4_x}\le
 C|t-s|^{-1/2}\|G(s)\|_{L^{4/3}_x}$.
Z nierówności Minkowskiego oraz
Lematu~\ref{lem:pde22-hls} wynika zatem
\begin{equation}\label{eq:pde22-ttstar}
 \|TT^*G\|_{L^4_{t,x}}\le C\|G\|_{L^{4/3}_{t,x}}.
\end{equation}
Wykorzystujemy teraz dodatniość formy
$\langle TT^*G,G\rangle_{t,x}=\|T^*G\|_{L^2_x}^2$.
Nierówność Höldera wraz z~\eqref{eq:pde22-ttstar}
daje $\|T^*G\|_2^2\le C\|G\|_{4/3}^2$.
Dualność pomiędzy $L^{4/3}_{t,x}$ i $L^4_{t,x}$
implikuje zatem
$\|T\varphi\|_4\le C^{1/2}\|\varphi\|_2$.
Przez gęstość otrzymujemy~\eqref{eq:pde22-str-hom}.

Dla operatora opóźnionego zastępujemy całkę
$\int_{\R}$ przez $\int_{s<t}$. To samo oszacowanie
punktowe pozostaje prawdziwe, ponieważ
$\int_{s<t}|t-s|^{-1/2}g(s)\,\mathrm ds\le
\int_{\R}|t-s|^{-1/2}g(s)\,\mathrm ds$
dla nieujemnego $g$; Lemat~\ref{lem:pde22-hls}
prowadzi wprost do~\eqref{eq:pde22-str-inhom}.
\end{proof}

\begin{uwaga}[Admisybilność i zakres twierdzenia]
W wymiarze $d$ para wykładników $2\le q,r\le\infty$
jest formalnie zgodna ze skalowaniem, gdy
$2/q+d/r=d/2$. Powyżej \emph{nie} dowodziliśmy
wszystkich takich oszacowań ani przypadków końcowych.
W szczególności przypadek $(q,r)=(2,\infty)$ w wymiarze
$d=2$ nie spełnia standardowego twierdzenia Strichartza.
Oszacowania dla operatorów z potencjałem wymagają dodatkowych
założeń spektralnych; nie wynikają z samej samosprzężoności.
\end{uwaga}

\section{Stan związany i brak dyspersji dla pewnych hamiltonianów}
\begin{przyklad}[Oscylator harmoniczny]
W przestrzeni $L^2(\R^d)$ rozważmy nieujemny operator
$H_{\mathrm{osc}}=-\Delta+|x|^2$ w jego standardowej
samosprzężonej realizacji. Funkcja
$\psi_0(x)=e^{-|x|^2/2}$ spełnia
$H_{\mathrm{osc}}\psi_0=d\psi_0$.
W szczególności $e^{-itH_{\mathrm{osc}}}\psi_0
=e^{-itd}\psi_0$: jej gęstość przestrzenna
$|\psi_0(x)|^2$ nie zmienia się w czasie.
Nie może więc istnieć uniwersalne oszacowanie
$\|e^{-itH_{\mathrm{osc}}}f\|_\infty\le C|t|^{-d/2}\|f\|_1$
dla wszystkich $t\ne0$ i wszystkich $f\in L^1\cap L^2$,
bo zastosowanie go do $\psi_0$ i przejście $t\to\infty$
prowadzi do sprzeczności.
\end{przyklad}

\section{Zależność od warunków brzegowych i dalsze pytania}
Na ograniczonym obszarze gładkim laplasjan Dirichleta
ma dyskretne widmo. Jego ewolucja Schrödingera jest nadal
unitarna, lecz na funkcji własnej działa tylko przez zmianę
fazy, a zatem nie występuje globalny zanik dyspersyjny
analogiczny do~\eqref{eq:pde22-L1Linf}. Jest to kolejny
powód, dla którego należy wyraźnie oddzielać twierdzenia
na $\R^d$ od zagadnień z brzegiem.

Pozostają ważne kierunki dalszego rozwoju:
bezpośrednia konstrukcja hamiltonianów z osobliwymi potencjałami,
lokalne i globalne oszacowania Strichartza dla nieliniowych
równań Schrödingera, zjawiska rozpraszania oraz analiza
rezonansów spektralnych. Zagadnienia te wykraczają poza
podstawowy model liniowy niniejszego rozdziału.



\chapter[Równania Hamiltona--Jacobiego]{Równania Hamiltona--Jacobiego: charakterystyki, rozwiązania lepkościowe i wzór Hopfa--Laxa}
\label{chap:pde23}

\section{Zagadnienie Cauchy'ego i geometryczne znaczenie hamiltonianu}
Niech $0<T<\infty$, $d\in\N$, $d\ge1$,
$g\colon\R^d\to\R$, a $H\colon\R^d\times\R^d\to\R$
będzie hamiltonianem. Rozważamy równanie
\begin{equation}\label{eq:pde23-hj}
 \partial_tu(t,x)+H(x,\nabla_xu(t,x))=0,
 \qquad u(0,x)=g(x).
\end{equation}
W pierwszej części zakładamy gładkość danych, potem będziemy
rozpatrywać funkcje ciągłe, niekoniecznie różniczkowalne.
Gradienty $D_xH$ i $D_pH$ traktujemy jako wektory kolumnowe;
drugi z nich jest gradientem względem zmiennej pędowej $p$.
Macierze Jacobiego działają na wektory kolumnowe.

Równanie~\eqref{eq:pde23-hj} występuje w mechanice
hamiltonowskiej, optymalnym sterowaniu i prawach zachowania.
Nawet jeśli $g$ i $H$ są gładkie, rozwiązanie klasyczne
może utracić regularność po skończonym czasie.
Po przecięciu charakterystyk ich klasyczna konstrukcja
może przestać wyznaczać jednowartościowe rozwiązanie.
Potrzebujemy wówczas pojęcia rozwiązania i zasady
jednoznaczności, które nie zakładają różniczkowalności.
Lokalną metodę charakterystyk poznaliśmy w rozdziale
~\ref{chap:pde-characteristics}; poniżej wyprowadzimy
jej postać właściwą równaniu Hamiltona--Jacobiego.

\begin{definicja}[Rozwiązanie klasyczne]
Niech $H$ i $g$ będą ciągłe. Funkcja
$u\in C([0,T]\times\R^d)\cap C^1((0,T)\times\R^d)$
jest rozwiązaniem klasycznym~\eqref{eq:pde23-hj}, jeżeli
spełnia równanie punktowo dla $0<t<T$ oraz $u(0,x)=g(x)$.
W konstrukcji charakterystyk z danych początkowych
będziemy zakładać dodatkową regularność aż do chwili
$t=0$; sama powyższa definicja jej nie wymaga.
\end{definicja}

\section{Charakterystyki układu Hamiltona i utrata gładkości}
\begin{twierdzenie}[Układ charakterystyczny]
\label{thm:pde23-char}
Niech $H\in C^2(\R^d_x\times\R^d_p)$ i niech
$u\in C^2((0,T)\times\R^d)$ spełnia
$u_t+H(x,\nabla_xu)=0$. Dla $(t_0,x_0)\in(0,T)\times\R^d$
rozważmy rozwiązanie układu
\begin{equation}\label{eq:pde23-hamilton}
 \dot X=D_pH(X,P),\qquad \dot P=-D_xH(X,P),
 \qquad X(t_0)=x_0,\quad P(t_0)=\nabla_xu(t_0,x_0).
\end{equation}
Na całym jego przedziale istnienia zawartym w $(0,T)$
zachodzi $P(t)=\nabla_xu(t,X(t))$ i
\begin{equation}\label{eq:pde23-action}
 \frac{\mathrm d}{\mathrm dt}u(t,X(t))
 =P(t)\cdot D_pH(X(t),P(t))-H(X(t),P(t)).
\end{equation}
Ponadto $H(X(t),P(t))$ jest stały.
\end{twierdzenie}
\begin{proof}
Pole wektorowe w układzie~\eqref{eq:pde23-hamilton}
jest klasy $C^1$, więc lokalne istnienie i jednoznaczność
wynikają z twierdzenia Picarda--Lindelöfa.
Różniczkujemy równanie Hamiltona--Jacobiego względem $x_j$:
\[
 \partial_tu_{x_j}+\sum_{k=1}^d H_{p_k}(x,\nabla_xu)u_{x_jx_k}
 +H_{x_j}(x,\nabla_xu)=0.
\]
Równanie dla gradientu przyjmuje postać
\[
 \partial_t\nabla_xu+
 D_x^2u\,D_pH(x,\nabla_xu)=-D_xH(x,\nabla_xu).
\]
Pole $(t,x)\mapsto D_pH(x,\nabla_xu(t,x))$
jest klasy $C^1$. Rozwiązujemy lokalnie
$\dot Y=D_pH(Y,\nabla_xu(t,Y))$, $Y(t_0)=x_0$,
i kładziemy $Q(t)=\nabla_xu(t,Y(t))$.
Reguła łańcuchowa i powyższa równość dają
\[
 \dot Q=\partial_t\nabla_xu(t,Y)+D_x^2u(t,Y)\dot Y
       =-D_xH(Y,Q).
\]
Para $(Y,Q)$ spełnia zatem układ Hamiltona
z tymi samymi danymi co $(X,P)$.
Jednoznaczność daje lokalnie $Y=X$, $Q=P$.
Zbiór chwil, w których $P(t)=\nabla_xu(t,X(t))$,
jest względnie domknięty przez ciągłość i względnie
otwarty przez powtórzenie powyższego argumentu
z dowolnej takiej chwili. Ponieważ zawiera $t_0$,
jest całym rozpatrywanym przedziałem.
Reguła łańcuchowa daje teraz
\[
 \frac{\mathrm d}{\mathrm dt}u(t,X(t))
 =u_t(t,X(t))+\nabla_xu(t,X(t))\cdot\dot X(t)
 =-H(X,P)+P\cdot D_pH(X,P).
\]
Wreszcie
\[
 \frac{\mathrm d}{\mathrm dt}H(X,P)
 =D_xH(X,P)\cdot D_pH(X,P)
 -D_pH(X,P)\cdot D_xH(X,P)=0.
\]
\end{proof}

\begin{stwierdzenie}[Lokalna rekonstrukcja rozwiązania z charakterystyk]
Niech $H\in C^2(\R^{2d})$, $g\in C^2(\R^d)$
i $y_0\in\R^d$. Istnieje otoczenie punktu $(0,y_0)$,
w którym charakterystyki z danymi
$X(0,y)=y$, $P(0,y)=\nabla g(y)$ wyznaczają
rozwiązanie klasy $C^2$ równania Hamiltona--Jacobiego
z warunkiem początkowym $g$ (dla nieujemnych czasów).
Konstrukcja jest ważna lokalnie wszędzie tam,
gdzie odwzorowanie $(t,y)\mapsto(t,X(t,y))$
jest lokalnym dyfeomorfizmem.
\end{stwierdzenie}
\begin{proof}
Twierdzenie o lokalnym przepływie pola klasy $C^1$
daje funkcje $(X,P)$ klasy $C^1$ zmiennych $(t,y)$
w otoczeniu $(0,y_0)$; można tu rozważać także
ujemne czasy. Definiujemy
\[
 Z(t,y)=g(y)+\int_0^t
 \bigl(P(r,y)\cdot\dot X(r,y)-H(X(r,y),P(r,y))\bigr)
 \,\mathrm d\lambda_1(r).
\]
Funkcja podcałkowa jest klasy $C^1$ względem $(r,y)$,
ponieważ $\dot X=D_pH(X,P)$; stąd $Z$ jest klasy $C^1$.
Kluczowa jest tożsamość
\[
 D_yZ(t,y)=P(t,y)^{\mathsf T}D_yX(t,y),
\]
gdzie pochodna funkcji skalarnej jest wierszem.
Dla $j=1,\ldots,d$ różniczkujemy definicję $Z$:
\[
 \begin{split}
 \partial_{y_j}\dot Z
 &=\partial_{y_j}P\cdot\dot X+
 P\cdot\partial_{y_j}\dot X\\
 &\quad-D_xH\cdot\partial_{y_j}X
          -D_pH\cdot\partial_{y_j}P\\
 &=P\cdot\partial_{y_j}\dot X+
          \dot P\cdot\partial_{y_j}X\\
 &=\partial_t(P\cdot\partial_{y_j}X).
 \end{split}
\]
Pochodna mieszana $\partial_t\partial_{y_j}X$
istnieje i jest równa $\partial_{y_j}\dot X$:
równanie całkowe przepływu można różniczkować
względem $y$, ponieważ pole i przepływ są klasy $C^1$.
W chwili $0$ mamy $\partial_{y_j}Z=g_{y_j}$
oraz $P(0,y)\cdot\partial_{y_j}X(0,y)=g_{y_j}$.
Całkowanie daje żądaną tożsamość.

W chwili $0$ zachodzi $D_yX(0,y_0)=I$.
Twierdzenie o funkcji odwrotnej zastosowane do
$(t,y)\mapsto(t,X(t,y))$ daje lokalną odwrotność
$(t,x)\mapsto(t,Y(t,x))$ klasy $C^1$.
Definiujemy $u(t,x)=Z(t,Y(t,x))$, czyli
$u(t,X(t,y))=Z(t,y)$.
Różniczkując względem $y$, otrzymujemy
\[
 (\nabla_xu(t,X(t,y)))^{\mathsf T}D_yX(t,y)
 =P(t,y)^{\mathsf T}D_yX(t,y).
\]
Odwracalność $D_yX$ daje $\nabla_xu(t,X(t,y))=P(t,y)$.
Różniczkując względem $t$ przy ustalonym $y$,
otrzymujemy $u_t+P\cdot\dot X=\dot Z$,
czyli $u_t=-H(X,P)$.
Po podstawieniu $y=Y(t,x)$ daje to równanie.
Obie pierwsze pochodne $u$ są klasy $C^1$,
ponieważ $P,Y,H$ mają taką regularność.
Wobec tego $u$ jest klasy $C^2$.
Warunek początkowy wynika z $X(0,y)=y$ i $Z(0,y)=g(y)$.
Dowód działa też na późniejszej lokalnej gałęzi
odwrotności przepływu, dopóki charakterystyki
istnieją i macierz $D_yX$ jest odwracalna.
\end{proof}

Dla $H=H(p)$ druga część układu daje $P(t)=P(0)$.
Jeśli $g\in C^2(\R^d)$, charakterystyka wychodząca
z $y$ ma zatem postać
\begin{equation}\label{eq:pde23-explicit-char}
 X(t,y)=y+tD_pH(\nabla g(y)),\qquad P(t,y)=\nabla g(y).
\end{equation}
Wartość rozwiązania na odwracalnej gałęzi wynosi
\[
 Z(t,y)=g(y)+t\bigl[
 \nabla g(y)\cdot D_pH(\nabla g(y))-H(\nabla g(y))\bigr].
\]
Warunek lokalnej odwracalności jest kontrolowany przez
\begin{equation}\label{eq:pde23-jacobian}
 D_yX(t,y)=I+tD_{pp}^2H(\nabla g(y))D^2g(y).
\end{equation}
Równość wynika z reguły łańcuchowej.
Zerowanie się wyznacznika oznacza utratę lokalnej
odwracalności w klasie dyfeomorfizmów $C^1$;
nie oznacza samo przez się utraty różnowartościowości.

\begin{przyklad}[Powstanie osobliwości z gładkich danych]
\label{ex:pde23-cos}
Niech $d=1$, $H(p)=p^2/2$ i $g(y)=\cos y$. Wtedy
$X(t,y)=y-t\sin y$, $P(t,y)=-\sin y$,
$Z(t,y)=\cos y+(t/2)\sin^2 y$,
a $\partial_yX=1-t\cos y$.
Dla $0\le t<1$ mamy $\partial_yX\ge1-t>0$.
Ponadto $|X(t,y)-y|\le t$, więc odwzorowanie
$y\mapsto X(t,y)$ jest różnowartościowe i na $\R$.
Jego odwrotność jest klasy $C^1$ i daje globalne
w przestrzeni rozwiązanie klasy $C^2$ dla $t<1$.
Różniczkując $u_x(t,X(t,y))=-\sin y$
względem $y$, dostajemy
\[
 u_{xx}(t,X(t,y))=-\frac{\cos y}{1-t\cos y}.
\]
W szczególności $u_{xx}(t,0)=-1/(1-t)$;
rozwiązania nie można przedłużyć w klasie $C^2$
przez punkt $(1,0)$.

W chwili $t=1$ pochodna $X(1,\cdot)$ zeruje się
w $2\pi\mathbb Z$, ale odwzorowanie pozostaje ściśle
rosnące: całka nieujemnej pochodnej po każdym
przedziale dodatniej długości jest dodatnia.
Dla $t>1$ funkcja $y-t\sin y$ ma ujemną pochodną
w zerze i jest dodatnia w $\pi$. Twierdzenie
o wartości pośredniej daje $y_*\in(0,\pi)$
z $X(t,y_*)=X(t,0)=0$.
Pędy na tych dwóch charakterystykach są różne,
gdyż $P(t,0)=0$ i $P(t,y_*)=-\sin y_*\ne0$.
Nie można zatem zachować wszystkich charakterystyk
wychodzących z danych początkowych jako wykresu
jednego gradientu. Nie jest to jeszcze dowód
nieistnienia dowolnego rozwiązania w klasie $C^1$
z samej definicji. Poniżej skonstruujemy rozwiązanie
lepkościowe, które istnieje również dla $t>1$.
\end{przyklad}

\section{Rozwiązania lepkościowe: definicja bez pochodnych rozwiązania}
\begin{definicja}[Rozwiązanie lepkościowe]
\label{def:pde23-visc}
Niech $Q=(0,T)\times\Omega$, gdzie $\Omega\subset\R^d$
jest obszarem, a $H\colon\Omega\times\R^d\to\R$ jest
ciągły. Funkcja $u\in C(Q)$ jest \emph{podrozwiązaniem
lepkościowym} równania $u_t+H(x,\nabla_xu)=0$, jeśli dla
każdego $(t_0,x_0)\in Q$ i każdej funkcji
$\varphi$ klasy $C^1$ w otoczeniu tego punktu, dla której
$u-\varphi$ ma tam lokalne maksimum, zachodzi
\begin{equation}\label{eq:pde23-sub}
 \varphi_t(t_0,x_0)+H(x_0,\nabla_x\varphi(t_0,x_0))\le0.
\end{equation}
\emph{Nadrozwiązanie lepkościowe} to funkcja ciągła, dla której
w każdym punkcie lokalnego minimum $u-\varphi$ zachodzi
odwrotna nierówność $\ge0$ (dla wszystkich takich funkcji
testowych). \emph{Rozwiązanie lepkościowe}
jest jednocześnie podrozwiązaniem i nadrozwiązaniem.
Warunek początkowy $u(0,x)=g(x)$ oznacza, że
$u$ ma ciągłe przedłużenie do chwili $0$ o tej wartości.
Na płaskim torusie stosujemy tę samą definicję
w lokalnych współrzędnych otrzymanych z ilorazu
$\R^d/(2\pi\mathbb Z)^d$; zmiana takiej reprezentacji
jest lokalnie translacją, więc nie zmienia pochodnych.
\end{definicja}

\begin{uwaga}
Można równoważnie stosować funkcje testowe z $C^1(Q)$:
lokalną funkcję mnożymy przez gładką funkcję odcinającą,
równą $1$ w mniejszym otoczeniu punktu i mającą nośnik
w otoczeniu, na którym funkcja jest określona,
a następnie przedłużamy przez zero.
Nie zmienia to zachowania $u-\varphi$ przy punkcie
ani pochodnych funkcji testowej w tym punkcie.
Dodanie do $\varphi$ stałej pozwala założyć
$u(t_0,x_0)=\varphi(t_0,x_0)$ bez zmiany nierówności.
Ekstrema w definicji nie muszą być ścisłe.
\end{uwaga}

\begin{uwaga}
Słowo „lepkościowe” pochodzi z przybliżenia równań pierwszego
rzędu równaniami $u_t+H(x,\nabla_xu)=\eps\Delta u$,
$\eps\downarrow0$. Definicja nie wymaga istnienia takiego
przybliżenia ani pochodnych $u$. Korzysta wyłącznie
z pochodnych funkcji testowej.
Nie utożsamiamy jej automatycznie z rozwiązaniem słabym
w sensie dystrybucyjnym.
\end{uwaga}

\begin{stwierdzenie}[Zgodność z rozwiązaniem klasycznym]
\label{prop:pde23-smooth-visc}
Jeżeli $u\in C^1(Q)$, to $u$ jest rozwiązaniem
lepkościowym równania $u_t+H(x,\nabla_xu)=0$
wtedy i tylko wtedy, gdy spełnia je punktowo.
\end{stwierdzenie}
\begin{proof}
Jeżeli $u$ spełnia równanie punktowo i $u-\varphi$ ma
lokalne ekstremum, to na mocy warunku koniecznego
ekstremum wewnętrznego pochodne $u-\varphi$ znikają:
$\varphi_t=u_t$ i $\nabla_x\varphi=\nabla_xu$.
Dostajemy właściwą nierówność, a nawet równość.
W drugą stronę w ustalonym punkcie $(t_0,x_0)$
bierzemy $\varphi=u$. Różnica $u-\varphi\equiv0$
ma zarazem maksimum i minimum.
Obie nierówności wymuszają równość.
\end{proof}

\section{Zasada porównania na torusie}
Przez $\T^d=\R^d/(2\pi\mathbb Z)^d$ rozumiemy płaski torus,
a przez $\lambda_d$ na $\T^d$ miarę indukowaną przez miarę
Lebesgue'a na $[0,2\pi)^d$, bez normalizacji do jedynki;
w szczególności $\lambda_d(\T^d)=(2\pi)^d$.
Najważniejszym argumentem zapewniającym jednoznaczność
jest \emph{podwajanie zmiennych}. Podajemy pełny dowód
modelowego twierdzenia, w którym brak brzegu pozwala
oddzielić istotę argumentu od technicznych barier.

\begin{twierdzenie}[Porównanie rozwiązań lepkościowych]
\label{thm:pde23-comparison}
Niech $T>0$, a $H\colon\T^d\times\R^d\to\R$ będzie
ciągły i niech istnieje $L>0$ takie, że dla wszystkich
$x,y\in\T^d$, $p\in\R^d$ i dla dostatecznie bliskich
$x,y$, przy użyciu tej samej lokalnej reprezentacji
współrzędnych,
\begin{equation}\label{eq:pde23-HxLip}
 |H(x,p)-H(y,p)|\le L|x-y|(1+|p|).
\end{equation}
Niech $u,v\in C([0,T]\times\T^d)$ będą odpowiednio
podrozwiązaniem i nadrozwiązaniem lepkościowym
$w_t+H(x,\nabla w)=0$ w $(0,T)\times\T^d$.
Jeśli $u(0,x)\le v(0,x)$ dla każdego $x$, to
$u(t,x)\le v(t,x)$ dla wszystkich $(t,x)\in[0,T]\times\T^d$.
\end{twierdzenie}
\begin{proof}
Przypuśćmy, że nierówność jest fałszywa. Wybieramy
$0<\tau<T$ oraz $(t_*,x_*)\in(0,\tau)\times\T^d$
takie, że $u(t_*,x_*)-v(t_*,x_*)>0$.
Wybieramy małe $\delta,\kappa>0$ spełniające
$u(t_*,x_*)-v(t_*,x_*)-\delta t_*-\kappa/(\tau-t_*)>0$.
Dla $\eps,\alpha>0$ maksymalizujemy dla
$t\in[0,\tau)$, $s\in[0,\tau]$ i $x,y\in\T^d$ funkcję
\begin{equation}\label{eq:pde23-doubling}
 \Phi_{\eps,\alpha}(t,x,s,y)=u(t,x)-v(s,y)
 -\frac{d_{\T}(x,y)^2}{2\eps}
 -\frac{(t-s)^2}{2\alpha}
 -\delta t-\frac{\kappa}{\tau-t}.
\end{equation}
Z powodu ostatniego składnika maksimum jest osiągane
przy $t<\tau$; przy ustalonym $t$ wybór $s\in[0,\tau]$
nie wymaga osobnej bariery. Maksimum jest dodatnie,
bo można wstawić $t=s=t_*$ i $x=y=x_*$.
Niech $C_0=2\max(\|u\|_\infty,\|v\|_\infty)+1$.
W punkcie dodatniego maksimum wszystkie składniki kary
są nieujemne, toteż $\kappa/(\tau-t_\eps)\le C_0$.
A zatem $t_\eps\le\tau-\kappa/C_0$ niezależnie od
$\eps,\alpha$. Po wykazaniu, że
$|t_\eps-s_\alpha|\to0$, dostajemy również $s_\alpha<\tau$
dla dostatecznie małych parametrów. Odrębnie poniżej
wykluczamy ich zerowanie się przy $t=0$.

Niech $(t_\eps,x_\eps,s_\alpha,y_\eps)$ będzie
punktem maksimum; dla czytelności pomijamy w indeksach
część parametrów. Porównanie z punktem
$(t_\eps,x_\eps,s_\alpha,x_\eps)$ oraz jednostajna
ciągłość $v$ dają
$\frac{d_{\T}(x_\eps,y_\eps)^2}{2\eps}
\le \omega_v(d_{\T}(x_\eps,y_\eps))$.
Z ograniczoności $u,v$ wynika
$d_{\T}(x_\eps,y_\eps)\to0$, a z powyższej
nierówności
\begin{equation}\label{eq:pde23-pen-small}
 \frac{d_{\T}(x_\eps,y_\eps)^2}{\eps}\longrightarrow0.
\end{equation}
Analogicznie porównanie z $s=t$ i jednostajna ciągłość
$v$ w czasie dają $|t_\eps-s_\alpha|\to0$ i
$(t_\eps-s_\alpha)^2/\alpha\to0$. Wybieramy
kolejność granic $\eps\downarrow0$, $\alpha\downarrow0$
lub wspólną kolejność zbieżną do zera.

Dla dostatecznie małych parametrów zarówno $t_\eps$,
jak i $s_\alpha$ należą do $(0,\tau)$.
Gdyby na przykład $t_\eps=0$ dla pewnego ciągu parametrów,
to $s_\alpha\to0$, $d_{\T}(x_\eps,y_\eps)\to0$;
jednostajna ciągłość i warunek początkowy dawałyby
graniczną wartość maksimum $\le0$, w sprzeczności
z ustaloną dodatnią wartością w punkcie $(t_*,x_*)$.
Tak samo wykluczamy $s_\alpha=0$.

Ponieważ punkty przestrzenne są bliskie, kwadrat
odległości ma w ich otoczeniu gładką postać euklidesową.
Oznaczmy $p=(x_\eps-y_\eps)/\eps$ w lokalnych
współrzędnych. Z definicji maksimum skonstruowane
przez zamrożenie $(s,y)$ funkcje testowe dotykają $u$
z góry, a przez zamrożenie $(t,x)$ dotykają $v$ z dołu.
Otrzymujemy zatem
\begin{align*}
 \frac{t_\eps-s_\alpha}{\alpha}+\delta
  +\frac{\kappa}{(\tau-t_\eps)^2}+H(x_\eps,p)&\le0,\\
 \frac{t_\eps-s_\alpha}{\alpha}+H(y_\eps,p)&\ge0.
\end{align*}
Po odjęciu dostajemy
$\delta\le H(y_\eps,p)-H(x_\eps,p)$.
Założenie~\eqref{eq:pde23-HxLip} pozwala ograniczyć
prawą stronę przez
$L|x_\eps-y_\eps|+L|x_\eps-y_\eps|^2/\eps$,
co dąży do zera na mocy~\eqref{eq:pde23-pen-small}.
Sprzeczność dowodzi twierdzenia dla $t<\tau$,
a dowolność $\tau$ i ciągłość dają tezę do czasu $T$.
\end{proof}

\begin{wniosek}[Jednoznaczność i kontrakcja]
\label{cor:pde23-contract}
Przy założeniach Twierdzenia~\ref{thm:pde23-comparison}
rozwiązanie lepkościowe z danymi $g\in C(\T^d)$ jest
co najwyżej jedno. Jeżeli $u,v$ są dwoma rozwiązaniami
dla danych $g,h\in C(\T^d)$, to
$\|u(t)-v(t)\|_{L^\infty(\T^d)}\le
\|g-h\|_{L^\infty(\T^d)}$ dla $0\le t\le T$.
\end{wniosek}
\begin{proof}
Równanie nie zależy od wartości $u$, zatem $v+c$
jest rozwiązaniem, ilekroć $v$ nim jest. Stosujemy
porównanie do pary $(u,v+c)$, gdzie
$c=\|g-h\|_\infty$, a następnie zamieniamy rolami $u,v$.
\end{proof}

\section{Transformata Legendre'a i wzór Hopfa--Laxa}
\begin{definicja}[Funkcja sprzężona Legendre'a]
Niech $H\colon\R^d\to\R$ będzie funkcją wypukłą.
Jej funkcję sprzężoną definiujemy wzorem
$L(v)=\sup_{p\in\R^d}(p\cdot v-H(p))$.
Funkcję $H$ nazywamy \emph{superliniową}, gdy
$H(p)/|p|\to\infty$ przy $|p|\to\infty$.
\end{definicja}

\begin{lemat}[Dualność i kontrola prędkości minimalizatorów]
\label{lem:pde23-legendre}
Niech $H\in C^1(\R^d)$ będzie wypukła i superliniowa.
Wtedy $L$ jest skończona, ciągła, wypukła i superliniowa,
a dla każdego $p\in\R^d$ zachodzi
\begin{equation}\label{eq:pde23-duality}
 H(p)=\sup_{v\in\R^d}(p\cdot v-L(v)).
\end{equation}
Jeżeli $g$ jest globalnie $K$-lipschitzowska, to
$y\mapsto g(y)+tL((x-y)/t)$ osiąga minimum, a każda
jego minimalizująca wartość $v=(x-y)/t$ spełnia
$L(v)-K|v|\le L(0)$. W szczególności prędkości
minimalizatorów są ograniczone liczbą zależną tylko od
$H$ i $K$.
\end{lemat}
\begin{proof}
Ponieważ $H$ jest superliniowa, funkcja
$p\mapsto p\cdot v-H(p)$ dąży do $-\infty$ dla
$|p|\to\infty$, więc supremum definiujące $L(v)$
jest skończone i osiągane. Jako supremum funkcji
afinicznych $L$ jest wypukła i dolnie półciągła,
a skończona funkcja wypukła na całej $\R^d$ jest ciągła.
Dla dowolnego $R>0$ i $v\ne0$ wybieramy
$p=R v/|v|$, zatem $L(v)\ge R|v|-\max_{|p|=R}H(p)$.
Po podzieleniu przez $|v|$ i przejściu $|v|\to\infty$
wnioskujemy o superliniowości $L$, ponieważ $R$ jest
dowolne.

Z definicji $L$ mamy $p\cdot v-L(v)\le H(p)$.
Jeśli jednak $v=D H(p)$, wypukłość daje dla każdego $q$
$H(q)\ge H(p)+v\cdot(q-p)$, stąd
$L(v)=p\cdot v-H(p)$; powstaje równość w~\eqref{eq:pde23-duality}.

Dla ustalonych $x,t>0$ i $y=x-tv$ warunek Lipschitza
daje $g(x-tv)\ge g(x)-Kt|v|$. Ponieważ $L$ jest
superliniowa, funkcja $g(x-tv)+tL(v)$ dąży do
$+\infty$ przy $|v|\to\infty$; ciągłość zapewnia
osiągnięcie minimum. Porównując je z próbą $v=0$,
dostajemy
$g(x-tv)+tL(v)\le g(x)+tL(0)$.
Wykorzystanie nierówności Lipschitza daje ostatnią tezę.
\end{proof}

\begin{twierdzenie}[Wzór Hopfa--Laxa]
\label{thm:pde23-hopflax}
Niech $H\in C^1(\R^d)$ będzie wypukła i superliniowa,
a $g\colon\R^d\to\R$ niech będzie funkcją ograniczoną
i globalnie $K$-lipschitzowską. Dla $t>0$ zdefiniujmy
\begin{equation}\label{eq:pde23-hopflax}
 u(t,x)=\min_{y\in\R^d}\left\{
 g(y)+tL\!\left(\frac{x-y}{t}\right)\right\},
 \qquad u(0,x)=g(x).
\end{equation}
Wtedy $u$ jest ograniczona i globalnie Lipschitzowska
na $[0,T]\times\R^d$ dla każdego $T>0$ oraz jest
rozwiązaniem lepkościowym równania
$u_t+H(\nabla u)=0$. Ponadto
\begin{equation}\label{eq:pde23-dpp}
 u(t,x)=\min_{z\in\R^d}
 \left\{u(s,z)+(t-s)L\!\left(\frac{x-z}{t-s}\right)\right\}
 \quad (0\le s<t).
\end{equation}
\end{twierdzenie}
\begin{proof}
\textbf{Krok 1: ograniczenie minimalizatorów i ciągłość.}
Lemat~\ref{lem:pde23-legendre} daje osiąganie minimum
oraz wspólne ograniczenie $|x-y|/t\le R$.
Wybór dowolnego $v$ w równoważnej reprezentacji
$u(t,x)=\min_v\{g(x-tv)+tL(v)\}$ daje
$u(t,x)\le g(x)+tL(0)$, a nierówność Lipschitza
$g(x-tv)\ge g(x)-Kt|v|$ daje
$u(t,x)\ge g(x)+t\min_v(L(v)-K|v|)$.
Wynika stąd $\|u(t)-g\|_\infty\le Ct$.
Przy przesunięciu $x\mapsto x+h$ każda próba $v$
zmienia składnik $g(x-tv)$ o co najwyżej $K|h|$;
po wzięciu minimów otrzymujemy
$|u(t,x+h)-u(t,x)|\le K|h|$.

\textbf{Krok 2: zasada programowania dynamicznego.}
Oznaczmy $a=t-s$, $b=s$ i użyjmy wypukłości $L$.
Dla dowolnych $y,z\in\R^d$ mamy
\[
 aL\!\left(\frac{x-z}{a}\right)
 +bL\!\left(\frac{z-y}{b}\right)
 \ge (a+b)L\!\left(\frac{x-y}{a+b}\right).
\]
Minimalizacja kolejno po $y,z$ daje kierunek „$\ge$”
w~\eqref{eq:pde23-dpp}. Kierunek przeciwny otrzymujemy
przez wybór $z=y+\frac{s}{t}(x-y)$ na prostym odcinku
łączącym $y$ z $x$, kiedy obie prędkości są równe
$(x-y)/t$ i nierówność wypukłości staje się równością.
Przypadek $s=0$ wynika bezpośrednio z definicji.

Stosując wzór dynamiczny na przedziale długości $h>0$
i argument kroku pierwszego do funkcji $u(s,\cdot)$,
dostajemy $|u(s+h,x)-u(s,x)|\le Ch$ ze stałą
zależną tylko od $K,H$. Razem z oszacowaniem
przestrzennym dowodzi to globalnej ciągłości Lipschitza.

\textbf{Krok 3: nierówność dla funkcji dotykającej z góry.}
Niech $u-\varphi$ ma lokalne maksimum równe zero
w $(t,x)$, $t>0$. Ustalmy dowolne $v\in\R^d$.
Z~\eqref{eq:pde23-dpp} dla małych $h>0$ dostajemy
$u(t,x)\le u(t-h,x-hv)+hL(v)$.
Ponieważ $u(t-h,x-hv)\le\varphi(t-h,x-hv)$,
mamy $\varphi(t,x)\le\varphi(t-h,x-hv)+hL(v)$.
Po podzieleniu przez $h$ i przejściu do granicy
$\varphi_t+v\cdot\nabla\varphi\le L(v)$.
Warunek zachodzi dla każdego $v$; dzięki
\eqref{eq:pde23-duality} otrzymujemy
$\varphi_t+H(\nabla\varphi)\le0$.

\textbf{Krok 4: nierówność dla funkcji dotykającej z dołu.}
Niech $u-\varphi$ ma lokalne minimum równe zero
w $(t,x)$. Wybieramy minimalizator $z_h=x-hv_h$
w~\eqref{eq:pde23-dpp}. Wspólne oszacowanie Lipschitza
funkcji $u(t-h,\cdot)$ i Lemat~\ref{lem:pde23-legendre}
dają ograniczenie $|v_h|\le R$.
Z minimalności i dotykania z dołu wynika
$\varphi(t,x)\ge\varphi(t-h,x-hv_h)+hL(v_h)$.
Taylor z resztą $o(h)$ jednostajną dla $|v_h|\le R$
daje
$\varphi_t+v_h\cdot\nabla\varphi\ge L(v_h)+o(1)$.
Ale $L(v_h)-v_h\cdot\nabla\varphi\ge-H(\nabla\varphi)$,
więc $\varphi_t+H(\nabla\varphi)\ge o(1)$.
Przechodzimy do $h\downarrow0$.
Obie nierówności dowodzą rozwiązania lepkościowego.
\end{proof}

\begin{uwaga}[O jednoznaczności na całej przestrzeni]
Dla autonomicznego $H=H(p)$ powyższe rozwiązanie jest
jedynym ograniczonym, globalnie Lipschitzowskim
rozwiązaniem lepkościowym o danych $g$.
Dowód porównania na $\R^d$ otrzymuje się z dowodu
Twierdzenia~\ref{thm:pde23-comparison}, dodając do
funkcji podwajanej zmiennej składnik
$-\beta\sqrt{1+|x|^2}$, który zapewnia osiąganie
maksimum. Dla funkcji Lipschitzowskich gradienty funkcji
testowych w punktach styczności są jednostajnie
ograniczone. Nierówność dla $H$ wynika więc z jego
lokalnej lipschitzowskości w zmiennej pędowej, a dodatkowy
błąd ma rozmiar $O(\beta)$. Najpierw przechodzimy
$\eps,\alpha\downarrow0$, później $\beta\downarrow0$.
Jest to ta sama zasada porównania, z dodatkową lokalizacją
maksimum; nie korzysta z wypukłości $H$.
\end{uwaga}


\begin{wniosek}[Skończony obszar zależności we wzorze Hopfa--Laxa]
\label{cor:pde23-hopflax-local}
Niech $H$ spełnia założenia Twierdzenia~\ref{thm:pde23-hopflax},
a $g,h\colon\R^d\to\R$ będą ograniczone i globalnie
$K$-lipschitzowskie. Niech $u_g,u_h$ oznaczają funkcje
zdefiniowane przez wzór Hopfa--Laxa odpowiednio dla $g$ i $h$.
Istnieje $R_{H,K}<\infty$, zależne tylko od $H,K$, takie, że
jeśli $g=h$ na domkniętej kuli $\overline B_R(x_0)$, to
\[
 u_g(t,x)=u_h(t,x)
 \quad\text{dla }t>0,\qquad |x-x_0|+tR_{H,K}\le R.
\]
Jeżeli $R_{H,K}=0$, nierówność interpretujemy bez
ograniczenia czasu. Analogiczna własność zachodzi w chwili $t=0$.
\end{wniosek}
\begin{proof}
Z Lematu~\ref{lem:pde23-legendre} wszyscy minimalizatorzy
dla obu danych są postaci $y=x-tv$, gdzie
$L(v)-K|v|\le L(0)$. Ponieważ $L$ jest superliniowa,
ten podpoziom jest zwarty; wybieramy $R_{H,K}$ nie mniejsze
niż maksimum $|v|$ na nim. Gdy
$|x-x_0|+tR_{H,K}\le R$, każdy minimalizator dla $g$
i każdy minimalizator dla $h$ leży w $\overline B_R(x_0)$.
Na tej kuli obie funkcje minimalizowane są identyczne.
Minimum globalne każdej jest osiągane wewnątrz kuli,
więc obie wartości minimalne są równe.
\end{proof}


\section{Rozwiązanie na przeciętych charakterystykach i półwklęsłość}
\begin{twierdzenie}[Półwklęsłość w przypadku kwadratowym]
\label{thm:pde23-semiconcave}
Niech $H(p)=|p|^2/2$, $g$ będzie ograniczona i
Lipschitzowska, a $u$ będzie dane przez
\eqref{eq:pde23-hopflax}. Dla każdego $t>0$
funkcja $x\mapsto u(t,x)-|x|^2/(2t)$ jest wklęsła.
W szczególności $u(t,\cdot)$ jest lokalnie Lipschitzowska
i dwukrotnie różniczkowalna prawie wszędzie w sensie
Alexandrowa (ostatni fakt jest klasycznym twierdzeniem
o funkcjach półwklęsłych, którego tutaj nie dowodzimy).
\end{twierdzenie}
\begin{proof}
Ponieważ $L(v)=|v|^2/2$, mamy
\[
 u(t,x)-\frac{|x|^2}{2t}
 =\inf_{y\in\R^d}\left\{
 g(y)+\frac{|y|^2}{2t}-\frac{x\cdot y}{t}\right\}.
\]
Dla ustalonego $y$ prawa strona pod nawiasem jest
funkcją afiniczną względem $x$. Infimum funkcji
afinicznych jest funkcją wklęsłą. Lokalna
lipschitzowskość wynika tu również bezpośrednio
z Twierdzenia~\ref{thm:pde23-hopflax}.
\end{proof}

\begin{przyklad}[Utrata różniczkowalności i wybór lepkościowy]
\label{ex:pde23-shock}
Dla $g(x)=\cos x$ i $H(p)=p^2/2$ otrzymujemy
$u(t,x)=\min_y\{\cos y+(x-y)^2/(2t)\}$.
Jeśli $0<t<1$, mapa charakterystyk
$y\mapsto y-t\sin y$ jest ściśle rosnąca.
Gdy $t>1$ i $x=0$, funkcja minimalizowana jest parzysta;
punkt $y=0$ jest lokalnym maksimum, ponieważ druga pochodna
w $0$ wynosi $-1+1/t<0$. Dla $t$ dostatecznie bliskiego $1$
ma ona dwa symetryczne globalne minima $\pm a(t)$,
$a(t)>0$, spełniające $a(t)=t\sin a(t)$.
Dla $x$ dodatnich bliskich zera minimalizatorem jest
gałąź $+a(t)$, a dla ujemnych gałąź $-a(t)$.
Z twierdzenia o różniczkowaniu wartości minimalnej
(lub bezpośrednio z różnic ilorazowych) pochodne jednostronne
$u_x(t,0+)$ i $u_x(t,0-)$ wynoszą odpowiednio
$-a(t)/t$ i $a(t)/t$. Funkcja jest więc ciągła, lecz
nie jest różniczkowalna w $x=0$. Rozwiązanie lepkościowe
istnieje mimo przecięcia charakterystyk i jest jednoznaczne
w klasie ograniczonych rozwiązań Lipschitzowskich.
\end{przyklad}

\section{Związek z równaniem Burgersa i warunkiem entropii}
Jeżeli $u$ jest gładka i spełnia
$u_t+(u_x)^2/2=0$, to $w=u_x$ spełnia
$w_t+(w^2/2)_x=0$. Po powstaniu osobliwości
różniczkowanie można interpretować dystrybucyjnie.
Pochodna rozwiązania lepkościowego odpowiada
rozwiązaniu entropijnemu równania Burgersa dla
odpowiednich danych Lipschitzowskich.
Warunek entropii jest więc odpowiednikiem wyboru
lepkościowego na poziomie prawa zachowania.
Nie wyprowadzamy tutaj pełnego twierdzenia Krużkowa;
zostanie ono udowodnione w rozdziale o prawach zachowania.

\begin{stwierdzenie}[Różniczkowanie lepkiego równania Hamiltona--Jacobiego]
\label{prop:pde-v43-hj-viscous}
Niech $f\in C^2(\R)$, $\varepsilon>0$ oraz
$U\in C^{1,3}((0,T)\times\R;\R)$ spełnia
$U_t+f(U_x)=\varepsilon U_{xx}$ punktowo.
Wówczas $w=U_x\in C^{1,2}$ spełnia
$w_t+\partial_x f(w)=\varepsilon w_{xx}$.
Ponadto dla każdej wypukłej $\eta\in C^2(\R)$ i
strumienia $q$ określonego przez $q'=\eta' f'$ zachodzi
\[
 \partial_t\eta(w)+\partial_x q(w)
 =\varepsilon\partial_{xx}\eta(w)
 -\varepsilon\eta''(w)|w_x|^2
 \le\varepsilon\partial_{xx}\eta(w).
\]
\end{stwierdzenie}
\begin{proof}
Różniczkujemy równanie dla $U$ względem $x$:
$w_t+f'(w)w_x=\varepsilon w_{xx}$.
Po pomnożeniu przez $\eta'(w)$ otrzymujemy
$\partial_t\eta(w)+\partial_xq(w)
=\varepsilon\eta'(w)w_{xx}$.
Reguła łańcuchowa daje
$\eta'(w)w_{xx}=\partial_{xx}\eta(w)-\eta''(w)|w_x|^2$.
Ostatni wyraz jest niedodatni przez wypukłość $\eta$.
\end{proof}

\begin{uwaga}[Różniczkowanie nie jest jeszcze przejściem do granicy]
Powyższa tożsamość obowiązuje dla gładkich rozwiązań
\emph{z tą samą lepkością} $\varepsilon$ i łączy
rachunek w rozdziałach~\ref{chap:pde23} i~\ref{chap:pde24}.
Aby stwierdzić, że $\partial_xU$ jest rozwiązaniem
entropijnym dla granicznego $U$, trzeba ponadto
uzasadnić zbieżność pochodnych przy
$\varepsilon\downarrow0$ i zgodność danych początkowych;
sama zbieżność jednostajna $U^\varepsilon\to U$
nie zapewnia silnej zbieżności $U_x^\varepsilon$.
\end{uwaga}


\begin{uwaga}[Zakres wyników]
Twierdzenie o porównaniu udowodniliśmy na torusie
w założeniach~\eqref{eq:pde23-HxLip}; wzór Hopfa--Laxa
udowodniliśmy dla autonomicznego hamiltonianu wypukłego
i superliniowego. Ogólna teoria dla hamiltonianów
niewypukłych, zależnych od $x$ i sterowania optymalnego
wymaga osobnych hipotez i metod. Samo stosowanie
wzoru~\eqref{eq:pde23-hopflax} poza tymi założeniami
nie jest uzasadnione.
\end{uwaga}

\section{Wskazówki do dalszej lektury rozdziału}
W następnym rozdziale przejdziemy od potencjału $u$
do skalarnego prawa zachowania dla $w=u_x$, gdzie
pojawią się rozwiązania słabe, fale uderzeniowe,
warunki Rankine'a--Hugoniota i selekcja entropijna.
Charakterystyki i geometria ich przecinania
pozostaną podstawowym źródłem przykładów.


\chapter[Skalarne prawa zachowania]{Skalarne prawa zachowania: fale uderzeniowe, entropia i twierdzenie Krużkowa}
\label{chap:pde24}

\section{Od charakterystyk do rozwiązań słabych}
Niech $T>0$, $f\in C^2(\R)$ i $u_0\colon\R\to\R$ będzie
funkcją mierzalną. Rozważamy zagadnienie Cauchy'ego
\begin{equation}\label{eq:pde24-cl}
 u_t+\partial_x f(u)=0\quad\text{w }(0,T)\times\R,
 \qquad u(0,x)=u_0(x).
\end{equation}
Funkcję $f$ nazywamy \emph{strumieniem}. Dla gładkiego rozwiązania
równanie ma postać $u_t+f'(u)u_x=0$: wartości $u$ są stałe
wzdłuż charakterystyk o prędkości $f'(u)$. Po przecięciu
charakterystyk nie można na ogół kontynuować rozwiązania
w klasie $C^1$; musimy dopuścić nieciągłości.

Poniższe sformułowanie całkowe jest szczególnym
przypadkiem definicji~\ref{def:pde2-weak}.
Przyjmujemy dodatkowo lokalną ograniczoność rozwiązania,
odpowiednią dla teorii entropijnej.

\begin{uwaga}[Potencjał a prawo zachowania]
\label{rem:pde-v43-hj-entropy-bridge}
Dla $f(w)=w^2/2$ związek z równaniem
Hamiltona--Jacobiego opisano w
rozdziale~\ref{chap:pde23}. Pełna rachunkowa zgodność
\emph{lepkich} równań jest udowodniona w
Stwierdzeniu~\ref{prop:pde-v43-hj-viscous}.
Selekcja rozwiązania entropijnego po zaniku lepkości
wymaga dodatkowej zwartości i jest odrębnym etapem,
którego nie zastępuje formalne różniczkowanie
rozwinięcia charakterystyk.
\end{uwaga}

\begin{definicja}[Rozwiązanie słabe]
Niech $u_0\in L^1_{\mathrm{loc}}(\R)$ i
$u\in L^\infty_{\mathrm{loc}}([0,T)\times\R)$.
Funkcja $u$ jest słabym rozwiązaniem~\eqref{eq:pde24-cl}, jeżeli
\begin{equation}\label{eq:pde24-weak}
 \int_0^T\!\int_\R
  \bigl(u\varphi_t+f(u)\varphi_x\bigr)\,\mathrm d\lambda_1(x)\,\mathrm d\lambda_1(t)
 +\int_\R u_0(x)\varphi(0,x)\,\mathrm d\lambda_1(x)=0
\end{equation}
dla każdej $\varphi\in C_c^\infty([0,T)\times\R)$.
W globalnej teorii kontrakcji będziemy dodatkowo wymagać
$u\in C([0,T];L^1(\R))\cap L^\infty((0,T)\times\R)$.
W teorii lokalnej zastąpimy globalną ciągłość w $L^1$
ciągłością w $L^1_{\mathrm{loc}}$.
\end{definicja}

\begin{stwierdzenie}[Bilans na ruchomym przedziale]
\label{prop:pde24-balance}
Niech $u\in C^1$ spełnia~\eqref{eq:pde24-cl} i niech
$a,b\in C^1([0,T])$ oraz $a(t)<b(t)$. Wówczas
\[
 \frac{\mathrm d}{\mathrm dt}\int_{a(t)}^{b(t)}u(t,x)\,\mathrm d\lambda_1(x)
 =f(u(t,a(t)))-f(u(t,b(t)))
 +b'(t)u(t,b(t))-a'(t)u(t,a(t)).
\]
\end{stwierdzenie}
\begin{proof}
Reguła Leibniza daje dwa składniki pochodzące z przesuwających
się końców przedziału. W pozostałej całce podstawiamy
$u_t=-\partial_xf(u)$ i stosujemy twierdzenie Newtona--Leibniza.
\end{proof}

\begin{definicja}[Wariacja całkowita i przestrzeń $BV$]
Niech $w\in L^1(\R)$ będzie funkcją o wartościach
rzeczywistych. Definiujemy wariację całkowitą
\[
 \operatorname{TV}(w)=\sup\left\{
 \int_\R w(x)\phi'(x)\,\mathrm d\lambda_1(x):
 \phi\in C_c^1(\R),\ \|\phi\|_\infty\le1\right\}.
\]
Piszemy $w\in BV(\R)$, jeżeli $\operatorname{TV}(w)<\infty$.
Równoważnie pochodna dystrybucyjna $Dw$ jest skończoną
miarą ze znakiem, której wariacja całkowita jest równa
$\operatorname{TV}(w)$.
Normę miary o skończonej wariacji oznaczamy przez
$\|\mu\|_{\mathcal M}=|\mu|(\R)$.
Na $BV(\R)$ przyjmujemy normę
$\|w\|_{BV}=\|w\|_1+\operatorname{TV}(w)$.
\end{definicja}

\begin{uwaga}
$BV(\R)$ jest przestrzenią Banacha w tej normie.
Istotnie wariacja jest seminormą (wynika to z symetrii
zbioru testów i nierówności trójkąta).
Jeżeli $(w_n)$ jest ciągiem Cauchy'ego w normie $BV$,
to $w_n\to w$ w $L^1(\R)$. Dla każdego testu $\phi$
z definicji wariacji i każdego $n$ mamy
\[
 \left|\int_\R(w_n-w)\phi'\,\mathrm d\lambda_1\right|
 \le\liminf_{m\to\infty}\operatorname{TV}(w_n-w_m).
\]
Przechodząc do supremum po $\phi$, dostajemy to samo
oszacowanie dla $\operatorname{TV}(w_n-w)$.
Dla dużego $n$ prawa strona jest dowolnie mała.
W szczególności $w_n-w\in BV$, więc $w\in BV$,
a $\|w_n-w\|_{BV}\to0$.
\end{uwaga}

\begin{lemat}[Oszacowanie translacji w przestrzeni $BV$]
\label{lem:pde24-bv-translation}
Jeżeli $w\in BV(\R)$, to dla każdego $h\in\R$ zachodzi
\[
 \|w(\cdot+h)-w\|_1\le |h|\operatorname{TV}(w).
\]
\end{lemat}
\begin{proof}
Najpierw uzasadniamy równoważność z definicji przestrzeni
$BV$. Funkcjonał $\phi\mapsto-\int_\R w\phi'\,
\mathrm d\lambda_1$ na $C_c^1(\R)$ jest ograniczony
w normie supremum, z normą $\operatorname{TV}(w)$
(zbiór testów jest symetryczny). Rozszerza się
jednoznacznie na $C_0(\R)$ przez gęstość. Twierdzenie
Riesza o reprezentacji funkcjonałów na $C_0(\R)$ daje
skończoną miarę ze znakiem $Dw$ o tej samej normie
wariacji. Odwrotnie, jeżeli $Dw$ jest taką miarą,
to z definicji pochodnej dystrybucyjnej otrzymujemy
to samo ograniczenie, a gęstość testów daje równość norm.

Niech $\rho_\delta$ będzie nieujemnym gładkim jądrem
aproksymującym jedność i $w_\delta=\rho_\delta*w$.
Wtedy $w_\delta\to w$ w $L^1(\R)$ i
$w_\delta'=\rho_\delta*Dw$. Z twierdzenia Fubiniego
względem miary wariacji wynika
$\|w_\delta'\|_1\le |Dw|(\R)$.
Dla $h>0$ twierdzenie Newtona--Leibniza daje
\[
 \|w_\delta(\cdot+h)-w_\delta\|_1
 \le\int_0^h\int_\R|w_\delta'(x+s)|
 \,\mathrm d\lambda_1(x)\,\mathrm d\lambda_1(s)
 \le h\operatorname{TV}(w).
\]
Dla $h<0$ korzystamy z niezmienniczości normy względem
translacji, a dla $h=0$ nierówność jest trywialna.
Przechodzimy do granicy w $L^1$.
\end{proof}

\section{Skoki i warunek Rankine'a--Hugoniota}
\begin{twierdzenie}[Warunek skoku]
\label{thm:pde24-RH}
Niech $s\in C^1((0,T))$ i niech $u_-,u_+$ będą funkcjami
klasy $C^1$ odpowiednio po lewej i prawej stronie krzywej
$x=s(t)$, przy czym obie funkcje wraz z pierwszymi
pochodnymi przedłużają się ciągle do krzywej od swojej
strony. Oznaczamy ich ślady przez $u_L(t),u_R(t)$.
Załóżmy, że każda gałąź spełnia klasycznie
$u_t+\partial_xf(u)=0$. Funkcja sklejona z tych gałęzi
jest rozwiązaniem dystrybucyjnym w otoczeniu krzywej wtedy
i tylko wtedy, gdy dla każdego $t$ zachodzi
\begin{equation}\label{eq:pde24-RH}
 s'(t)\bigl(u_R(t)-u_L(t)\bigr)
 =f(u_R(t))-f(u_L(t)).
\end{equation}
\end{twierdzenie}
\begin{proof}
Niech $\varphi\in C_c^\infty((0,T)\times\R)$
ma nośnik w rozważanym otoczeniu. Z reguły Leibniza
dla całki po $x<s(t)$ i z równania klasycznego wynika
\[
 \int_0^T\!\int_{x<s(t)}(u_-\varphi_t+
 f(u_-)\varphi_x)\,\mathrm d\lambda_1(x)\,
 \mathrm d\lambda_1(t)
 =\int_0^T(f(u_L)-s'u_L)\varphi(t,s(t))\,
 \mathrm d\lambda_1(t).
\]
Istotnie człon czasowy daje po całkowaniu przez części
$-s'u_L\varphi$ na ruchomym brzegu, a człon przestrzenny
daje $f(u_L)\varphi$. Całki wewnętrzne znoszą się
na mocy równania klasycznego. Na obszarze $x>s(t)$
orientacja brzegu jest przeciwna, zatem otrzymujemy
analogicznie $(s'u_R-f(u_R))\varphi(t,s(t))$.
Po zsumowaniu i zmianie znaku, zgodnie z definicją
pochodnej dystrybucyjnej, dostajemy
\[
 \langle u_t+\partial_x f(u),\varphi\rangle
 =\int_0^T([f(u)]-s'[u])\varphi(t,s(t))\,
 \mathrm d\lambda_1(t),
 \qquad [h]=h_R-h_L.
\]
Prawą stronę zapisujemy także jako
$([f(u)]-s'[u])\delta_{x=s(t)}$, gdzie delta oznacza
miarę zadaną przez całkowanie względem czasu wzdłuż
wykresu, a nie względem jego długości łuku.
Współczynnik jest ciągły. Gdyby był niezerowy
w pewnym punkcie, miałby stały ścisły znak na pewnym
przedziale czasu; nieujemny test o nośniku
w małym otoczeniu odpowiedniego fragmentu wykresu
dałby niezerową całkę. Stąd znikanie dystrybucji
jest równoważne~\eqref{eq:pde24-RH}.
\end{proof}

Dla strumienia Burgersa $f(u)=u^2/2$, przy $u_L\ne u_R$,
prędkość skoku jest równa $(u_L+u_R)/2$. Sam warunek
Rankine'a--Hugoniota nie rozstrzyga jednak, czy taka fala
jest zgodna z fizycznym kierunkiem przepływu informacji.

\begin{kontrprzyklad}[Słabe rozwiązanie nie jest jednoznaczne]
\label{ex:pde24-expansion}
Dla $f(u)=u^2/2$ i danych Riemanna $u_0(x)=0$ dla $x<0$
oraz $u_0(x)=1$ dla $x>0$ funkcja
$u(t,x)=0$ dla $x<t/2$ i $u(t,x)=1$ dla $x>t/2$
spełnia warunek Rankine'a--Hugoniota, więc jest rozwiązaniem
słabym. Innym rozwiązaniem słabym tych samych danych jest
ciągły wachlarz $u(t,x)=0$ dla $x\le0$,
$u(t,x)=x/t$ dla $0<x<t$ oraz $u(t,x)=1$ dla $x\ge t$.
Na wnętrzu wachlarza $u_t=-x/t^2$ oraz
$\partial_x(u^2/2)=x/t^2$; na pozostałych
obszarach funkcja jest stała. Na prostych $x=0$ i
$x=t$ ślady obu gałęzi są równe, więc warunek
~\eqref{eq:pde24-RH} jest spełniony. Dla obu
rozwiązań różnica od $u_0$ jest ograniczona przez
$1$ i znika poza $[0,t]$, zatem norma tej różnicy
w $L^1_{\mathrm{loc}}$ dąży do zera przy $t\downarrow0$.
Całkowanie przez części od dodatniej chwili i przejście
z tą chwilą do zera daje pełny wzór
~\eqref{eq:pde24-weak}. Dodatkowa zasada selekcji
jest konieczna.
\end{kontrprzyklad}

\section{Entropie i poprawny kierunek fal uderzeniowych}
\begin{definicja}[Para entropijna]
Niech $f\in C^1(\R)$. Parą entropijną nazywamy funkcje
$\eta\in C^2(\R)$ oraz $q\in C^1(\R)$ takie, że $\eta$ jest
wypukła i $q'(z)=\eta'(z)f'(z)$. Rozwiązanie spełnia
nierówność entropijną, jeśli
$\partial_t\eta(u)+\partial_x q(u)\le0$ w sensie dystrybucji.
\end{definicja}

\begin{lemat}[Produkcja entropii przez lepkość]
\label{lem:pde24-viscentropy}
Jeżeli $u^\eps$ jest rozwiązaniem klasy $C^{1,2}$
(jedna ciągła pochodna czasowa i dwie przestrzenne)
równania $u_t^\eps+\partial_xf(u^\eps)=\eps u_{xx}^\eps$,
$\eps>0$, to dla każdej pary entropijnej zachodzi tożsamość
\begin{equation}\label{eq:pde24-entropyvisc}
 \partial_t\eta(u^\eps)+\partial_xq(u^\eps)
 =\eps\partial_{xx}\eta(u^\eps)
 -\eps\eta''(u^\eps)|u_x^\eps|^2
 \le\eps\partial_{xx}\eta(u^\eps).
\end{equation}
\end{lemat}
\begin{proof}
Mnożymy równanie przez $\eta'(u^\eps)$ i używamy reguły
łańcuchowej. Z równości
$\eta'(u)u_{xx}=\partial_{xx}\eta(u)-\eta''(u)u_x^2$
wynika wzór; znak ostatniego członu wynika z wypukłości.
\end{proof}

\begin{definicja}[Rozwiązanie entropijne Krużkowa]
\label{def:pde24-kru}
Niech $u_0\in L^1(\R)\cap L^\infty(\R)$.
Funkcja $u\in C([0,T];L^1(\R))\cap L^\infty((0,T)\times\R)$
jest \emph{rozwiązaniem entropijnym}~\eqref{eq:pde24-cl}, jeżeli
$u(0)=u_0$ w $L^1$ i dla każdego $k\in\R$ zachodzi
\begin{equation}\label{eq:pde24-kruzineq}
 \partial_t|u-k|+\partial_x\left(
   \operatorname{sgn}(u-k)(f(u)-f(k))\right)\le0
 \quad\text{w }\mathcal D'((0,T)\times\R),
\end{equation}
gdzie $\operatorname{sgn}(0)=0$.
Oznaczając $q_k(z)=\operatorname{sgn}(z-k)(f(z)-f(k))$,
zapisujemy tę nierówność równoważnie jako
\[
 \int_0^T\!\int_\R
 \bigl(|u-k|\varphi_t+q_k(u)\varphi_x\bigr)
 \,\mathrm d\lambda_1(x)\,\mathrm d\lambda_1(t)\ge0
\]
dla każdej nieujemnej funkcji
$\varphi\in C_c^\infty((0,T)\times\R)$.
\end{definicja}

\begin{uwaga}
Z warunków dla wszystkich $k$
wynika również słabe równanie: jeśli $k$ jest większe od
istotnego supremum $u$, w~\eqref{eq:pde24-kruzineq}
uzyskujemy jedną nierówność liniową, a jeśli jest mniejsze
od istotnego infimum, nierówność przeciwną. Obie nierówności dają równość dystrybucyjną.
Aby odzyskać również~\eqref{eq:pde24-weak},
testujemy równanie przez $\varphi(t,x)\theta_\delta(t)$,
gdzie $\theta_\delta$ rośnie gładko od $0$ do $1$
w $(\delta,2\delta)$ i jest stała poza tym przedziałem.
Wyraz zawierający $\theta_\delta'$ dąży do
$\int_\R u_0(x)\varphi(0,x)\,\mathrm d\lambda_1(x)$:
pochodna ma masę $1$, a całka z $u(t,x)\varphi(t,x)$
jest ciągła w chwili $0$. Pozostałe wyrazy przechodzą
do granicy przez zbieżność zmajoryzowaną na nośniku
testu. Tak samo postępujemy przy lokalnym śladzie
w $L^1_{\mathrm{loc}}$; przy lokalnej ograniczoności
stałe $k$ dobieramy stosownie do nośnika testu.
\end{uwaga}

\begin{lemat}[Równoważność warunków entropijnych]
\label{lem:pde24-entropy-equivalence}
Niech $u$ będzie lokalnie ograniczonym rozwiązaniem
dystrybucyjnym równania zachowania na $(0,T)\times\R$.
Nierówności Krużkowa~\eqref{eq:pde24-kruzineq} dla
wszystkich $k\in\R$ są równoważne nierównościom
$\partial_t\eta(u)+\partial_xq(u)\le0$ dla wszystkich
gładkich par entropijnych w sensie podanej wyżej definicji.
\end{lemat}
\begin{proof}
Ustalmy nieujemny test $\varphi$ i przedział $[A,B]$
zawierający istotny zakres wartości $u$ na jego nośniku.
Jeżeli zachodzą nierówności Krużkowa, to dla wypukłej
$\eta\in C^2(\R)$ na $[A,B]$ mamy
\[
 \eta(z)=a+bz+\frac12\int_A^B |z-k|\eta''(k)\,
                         \mathrm d\lambda_1(k),
\]
gdzie $a,b$ są stałymi. Istotnie pochodna rzędu
drugiego prawej strony na $(A,B)$ jest równa
$\eta''(z)$, więc różnica jest afiniczna.
Dla $q'=\eta'f'$ analogicznie
\[
 q(z)=c+bf(z)+\frac12\int_A^B q_k(z)\eta''(k)\,
                         \mathrm d\lambda_1(k).
\]
Druga tożsamość wynika z różniczkowania względem
$z$: dla $k\ne z$ pochodna $q_k(z)$ wynosi
$\operatorname{sgn}(z-k)f'(z)$, a ilorazy różnicowe
są jednostajnie ograniczone przez supremum $|f'|$
na przedziale. Twierdzenie o zbieżności zmajoryzowanej
pozwala różniczkować pod całką; stała $c$ uwzględnia
dowolną normalizację strumienia.
Części afiniczne dają zerowy wkład do nierówności,
ponieważ $u_t+\partial_x f(u)=0$ dystrybucyjnie.
Pozostały wkład jest całką z nierówności Krużkowa
z nieujemną wagą $\eta''/2$. Twierdzenie Fubiniego
stosujemy na zwartym nośniku testu i ograniczonym
przedziale parametrów $k$. Otrzymujemy żądaną
nierówność dla $\eta,q$.

Odwrotnie, dla ustalonego $k$ bierzemy
$\eta_\delta(z)=\sqrt{(z-k)^2+\delta^2}$ oraz
$q_\delta(z)=\int_k^z\eta_\delta'(r)f'(r)\,
\mathrm d\lambda_1(r)$. Są to gładkie pary entropijne
w wymaganym sensie. Entropie zbiegają jednostajnie
do $|z-k|$, a strumienie do $q_k$ jednostajnie na
$[A,B]$: błąd strumienia nie przekracza całki
z $|\eta_\delta'(r)-\operatorname{sgn}(r-k)|\,|f'(r)|$
po przedziale zawierającym $[A,B]$ i $k$.
Ta całka dąży do zera przez zbieżność zmajoryzowaną.
Przechodząc do granicy w całkach z pochodnymi testu,
otrzymujemy nierówność Krużkowa.
\end{proof}

\begin{twierdzenie}[Warunek entropijny na skoku]
\label{thm:pde24-entropyshock}
Niech $u_L\ne u_R$, $f\in C^2(\R)$ i niech skok między
stałymi stanami $u_L,u_R$ porusza się z prędkością
$\sigma=(f(u_R)-f(u_L))/(u_R-u_L)$.
Jest on entropijny w sensie lokalnych nierówności
Krużkowa wtedy i tylko wtedy, gdy dla każdego $k$ pomiędzy
$u_L$ i $u_R$ zachodzi
\begin{equation}\label{eq:pde24-chord}
 \begin{cases}
 f(k)\le \ell(k),&u_L>u_R,\\
 f(k)\ge \ell(k),&u_L<u_R,
 \end{cases}
 \qquad
 \ell(k)=f(u_L)+\sigma(k-u_L).
\end{equation}
W szczególności gdy $f$ jest ściśle wypukła, skok z
$u_L>u_R$ jest entropijny, natomiast skok z $u_L<u_R$
nie jest entropijny.
\end{twierdzenie}
\begin{proof}
Przyjmujemy $[h]=h_R-h_L$. Z rachunku dystrybucyjnego
nierówność dla entropii $\eta$ jest równoważna
$[q]-\sigma[\eta]\le0$. Jeżeli $k$ leży poza odcinkiem
wyznaczonym przez stany, wyrażenie to znika, ponieważ
$[q]=\pm[f]$ i $[\eta]=\pm[u]$.
Gdy $u_L>u_R$ i $u_R<k<u_L$, mamy
$[q_k]=2f(k)-f(u_R)-f(u_L)$ oraz
$[|u-k|]=2k-u_R-u_L$. Stąd
\[
 \begin{aligned}
 [q_k]-\sigma[|u-k|]
 &=2f(k)-f(u_R)-f(u_L)\\
 &\quad-\sigma(2k-u_R-u_L)\\
 &=2(f(k)-\ell(k)),
 \end{aligned}
\]
gdzie w ostatniej równości korzystamy z
$f(u_R)=f(u_L)+\sigma(u_R-u_L)$.
Przy $u_L<u_R$ i $u_L<k<u_R$ oba wyrażenia
mają przeciwne znaki, więc wynik to
$2(\ell(k)-f(k))$. Dla $k$ równego jednemu ze
stanów współczynnik jest zerowy przez ciągłość.
Te warunki są konieczne i wystarczające. Ścisła wypukłość
umieszcza wykres $f$ ściśle poniżej cięciwy łączącej
odpowiednie punkty, co kończy dowód.
\end{proof}

\section{Zagadnienie Riemanna dla ściśle wypukłego strumienia}
\begin{twierdzenie}[Fala uderzeniowa i rozrzedzenia]
\label{thm:pde24-riemann}
Niech $f\in C^2(\R)$ i $f''(z)>0$ na domkniętym
przedziale o końcach $u_L,u_R$, jeśli $u_L\ne u_R$. Rozważmy dane $u_0=u_L$ dla $x<0$ oraz
$u_0=u_R$ dla $x>0$ (nie muszą należeć do $L^1$).
Dla $t>0$ istnieje samopodobne lokalnie ograniczone
rozwiązanie entropijne $u(t,x)=U(x/t)$:
\begin{enumerate}
\item Jeśli $u_L>u_R$, jest to skok przemieszczający się
z prędkością $\sigma=(f(u_L)-f(u_R))/(u_L-u_R)$.
\item Jeśli $u_L<u_R$, to
\[
 U(\xi)=\begin{cases}
 u_L,&\xi\le f'(u_L),\\
 (f')^{-1}(\xi),&f'(u_L)<\xi<f'(u_R),\\
 u_R,&\xi\ge f'(u_R).
 \end{cases}
\]
\end{enumerate}
Jeżeli $u_L=u_R$, rozwiązaniem jest funkcja stała $u(t,x)=u_L$.
Rozwiązania te są jednoznaczne w klasie
$C([0,T];L^1_{\mathrm{loc}}(\R))\cap L^\infty((0,T)\times\R)$
rozwiązań lokalnych nierówności entropijnych z zadanym
śladem początkowym. Dowód wykorzystuje skończoną prędkość
propagacji (Wniosek~\ref{cor:pde24-finite}).
Nie zakładamy całkowalności na całej prostej, lecz
wymagamy globalnego ograniczenia wartości rozwiązania.
\end{twierdzenie}
\begin{proof}
W pierwszym przypadku warunek Rankine'a--Hugoniota
sprawia, że równanie jest spełnione słabo, a
Twierdzenie~\ref{thm:pde24-entropyshock} zapewnia
nierówności entropijne. W drugim przypadku $U$ jest
ciągła, a na wachlarzu spełnia $f'(U(\xi))=\xi$.
Obliczamy $u_t=-\xi U'(\xi)/t$ oraz
$\partial_x f(u)=f'(U)U'(\xi)/t$; składniki się znoszą.
Poza wachlarzem rozwiązanie jest stałe. Ponieważ funkcja
jest ciągła, w przejściu przez brzegi wachlarza nie
powstają wyrazy typu delta. Dla każdej gładkiej pary entropijnej
na wnętrzu wachlarza zachodzi równość
$\partial_t\eta(u)+\partial_xq(u)=0$.
Po uzasadnieniu braku składników osobliwych poniżej
stosujemy Lemat~\ref{lem:pde24-entropy-equivalence}.
W obu przypadkach $u(t,\cdot)$ jest globalnie ograniczona
przez $\max(|u_L|,|u_R|)$. Różnica od danych początkowych
może wystąpić tylko w $[-Ct,Ct]$, gdzie
$C=\max_{z\in[\min(u_L,u_R),\max(u_L,u_R)]}|f'(z)|$
(zaś $|\sigma|\le C$ przez twierdzenie o wartości średniej).
Jej norma $L^1$ na dowolnym ograniczonym przedziale
dąży więc do zera. Dla $t_n\to t>0$ profile zbiegają
punktowo poza ewentualnym punktem skoku, więc zbieżność
zmajoryzowana daje ciągłość w $L^1_{\mathrm{loc}}$.
W przypadku wachlarza założenie o dodatniości $f''$
na domkniętym przedziale gwarantuje ograniczoność
pochodnej $(f')^{-1}$; na każdym zwartym przedziale
dodatnich czasów profil jest kawałkami klasy $C^1$
i ciągły na granicach kawałków. Całkowanie przez części
jak w dowodzie warunku skoku pokazuje, że nie ma
składników osobliwych także dla $\eta(u),q(u)$.
Jednoznaczność wynika z Wniosku~\ref{cor:pde24-finite};
dla równych stanów stosujemy rozwiązanie stałe.
\end{proof}

\begin{center}
\begin{tikzpicture}[x=1.05cm,y=1.55cm,>=Latex]
 \begin{scope}
  \draw[->] (-2.15,0)--(2.15,0) node[right] {$x$};
  \draw[->] (0,-.08)--(0,1.9) node[above] {$t$};
  \foreach \a in {-1.8,-1.2,-.6}{
   \pgfmathsetmacro{\ttt}{min(1.65,-2*\a)}
   \draw[blue!70!black] (\a,0)--({\a+\ttt},\ttt);
  }
  \foreach \a in {.3,.6,.9,1.2,1.5,1.8}{
   \pgfmathsetmacro{\ttt}{min(1.65,2*\a)}
   \draw[teal!70!black] (\a,0)--(\a,\ttt);
  }
  \draw[red!75!black,very thick] (0,0)--(.825,1.65)
     node[above right,font=\small] {$x=t/2$};
  \node[font=\small] at (-1.25,1.73) {$u_L=1$};
  \node[font=\small] at (1.55,1.73) {$u_R=0$};
  \node[font=\small] at (0,-.35) {fala uderzeniowa};
 \end{scope}
 \begin{scope}[xshift=5.1cm]
  \draw[->] (-2.15,0)--(2.15,0) node[right] {$x$};
  \draw[->] (0,-.08)--(0,1.9) node[above] {$t$};
  \foreach \a in {-1.8,-1.2,-.6}
   \draw[blue!70!black] (\a,0)--(\a,1.65);
  \foreach \a in {.3,.6,.9,1.2,1.5,1.8}{
   \pgfmathsetmacro{\ttt}{min(1.65,2.1-\a)}
   \draw[teal!70!black] (\a,0)--({\a+\ttt},\ttt);
  }
  \foreach \v in {.2,.4,.6,.8}
   \draw[orange!80!black] (0,0)--({1.65*\v},1.65);
  \draw[orange!80!black,very thick] (0,0)--(1.65,1.65)
      node[above right,font=\small] {$x=t$};
  \node[font=\small] at (-1.25,1.73) {$u_L=0$};
  \node[font=\small,anchor=west] at (.7,-.14) {$u_R=1$};
  \node[font=\small] at (0,-.35) {fala rozrzedzenia};
 \end{scope}
\end{tikzpicture}

\small Rysunek: charakterystyki równania Burgersa dla
danych $(u_L,u_R)=(1,0)$ (po lewej) oraz $(0,1)$
(po prawej). W pierwszym przypadku charakterystyki
dochodzą do skoku o prędkości $1/2$; w drugim
wachlarz promieni $x=vt$, $0<v<1$, wypełnia
obszar między charakterystykami $x=0$ i $x=t$.
\end{center}

\section{Nierówność Kato dla dwóch rozwiązań entropijnych}
\begin{twierdzenie}[Nierówność Kato--Krużkowa]
\label{thm:pde24-kato}
Niech $f\in C^1(\R)$ oraz niech
$u,v\in C([0,T];L^1(\R))\cap L^\infty((0,T)\times\R)$
będą rozwiązaniami entropijnymi z danymi $u_0,v_0$.
Wówczas w sensie dystrybucji na $(0,T)\times\R$
\begin{equation}\label{eq:pde24-kato}
 \partial_t|u-v|+\partial_x\bigl[
 \operatorname{sgn}(u-v)(f(u)-f(v))\bigr]\le0.
\end{equation}
\end{twierdzenie}
\begin{proof}
\textbf{Krok 1: podwojenie zmiennych.}
Dla ustalonej pary $(s,y)$ w nierówności entropijnej
dla $u(t,x)$ wybieramy $k=v(s,y)$, a dla ustalonej pary
$(t,x)$ w nierówności dla $v(s,y)$ wybieramy $k=u(t,x)$.
Są to dopuszczalne stałe, ponieważ w każdej nierówności
całkujemy tylko po zmiennych jej rozwiązania. Po
scalaniu nierówności twierdzeniem Fubiniego otrzymujemy
sumę dwóch nierówności na $(t,x,s,y)$.

\textbf{Krok 2: funkcja testowa.}
Niech $\rho\in C_c^\infty(\R)$ jest nieujemna, parzysta
i ma całkę równą $1$. Dla $\delta>0$ kładziemy
$\rho_\delta(z)=\delta^{-1}\rho(z/\delta)$.
Bierzemy $\psi\in C_c^\infty((0,T)\times\R)$, $\psi\ge0$,
i testujemy obie nierówności funkcją
\[
 \Phi_\delta(t,x,s,y)=
 \rho_\delta(t-s)\rho_\delta(x-y)\,
 \psi\!\left(\frac{t+s}{2},\frac{x+y}{2}\right).
\]
Dla małego $\delta$ nośnik leży w zwartym podzbiorze
$((0,T)\times\R)^2$: średnie zmiennych należą do
nośnika $\psi$, a ich różnice do nośników jąder.
Zatem dla każdej ustalonej pary zmiennych drugi
przekrój jest dopuszczalnym testem. Wybieramy mierzalne
reprezentanty $u,v$; wszystkie funkcje podcałkowe są
mierzalne, ograniczone na tym zwartym nośniku
i całkowalne, więc stosowanie twierdzenia Fubiniego
jest uzasadnione. Warunek entropijny obowiązuje dla
każdej stałej, nie tylko dla prawie każdej.
Suma pochodnych spełnia dokładnie
\[
 (\partial_t+\partial_s)\Phi_\delta
 =\rho_\delta(t-s)\rho_\delta(x-y)\psi_t,
 \qquad
 (\partial_x+\partial_y)\Phi_\delta
 =\rho_\delta(t-s)\rho_\delta(x-y)\psi_x,
\]
gdzie pochodne $\psi$ obliczamy w zmiennych średnich.
Po dodaniu i zastosowaniu twierdzenia Fubiniego dostajemy
\[
 \int_{((0,T)\times\R)^2}
 \left(
 |u(t,x)-v(s,y)|(\partial_t+\partial_s)\Phi_\delta
 +Q(u(t,x),v(s,y))(\partial_x+\partial_y)\Phi_\delta
 \right)
 \,\mathrm d\lambda_4(t,x,s,y)\ge0,
\]
gdzie $Q(a,b)=\operatorname{sgn}(a-b)(f(a)-f(b))$.
W obu nierównościach występują te same wartości
$|u(t,x)-v(s,y)|$ i ten sam strumień
$\operatorname{sgn}(u(t,x)-v(s,y))
(f(u(t,x))-f(v(s,y)))$; w nierówności dla $v$
oba znaki zmieniają się jednocześnie.

\textbf{Krok 3: granica diagonalna.}
Podstawiamy $h=t-s$, $k=x-y$; wartość bezwzględna
jakobianu tego podstawienia wynosi $1$. Wówczas
$v(s,y)=v(t-h,x-k)$, a pochodne $\psi$ obliczamy
w $(t-h/2,x-k/2)$. Na ustalonym zwartym powiększeniu
nośnika $\psi$ ciągłość translacji w $L^1_{\mathrm{loc}}$
daje
\[
 \int |v(t-h,x-k)-v(t,x)|\,\mathrm d\lambda_2(t,x)
 \longrightarrow0
 \qquad ((h,k)\to(0,0)),
\]
gdzie całka jest po tym powiększeniu. Różnicę wartości
$|u-v|$ ograniczamy przez $|v(t-h,x-k)-v(t,x)|$.
Jeżeli $M$ ogranicza $|f'|$ na wspólnym przedziale wartości
$u,v$, to dla ustalonego $a$ funkcja $b\mapsto Q(a,b)$
jest $M$-lipschitzowska: na każdej stronie punktu $a$
wynika to z twierdzenia o wartości średniej, a dla
argumentów po przeciwnych stronach stosujemy nierówność
trójkąta i $Q(a,a)=0$. Tak samo postępujemy w pierwszym
argumencie. Błąd w strumieniu ma więc analogiczne
oszacowanie. Całkowanie względem nieujemnego jądra
$\rho_\delta(h)\rho_\delta(k)$ o masie $1$ nie zwiększa
supremum tych błędów po jego nośniku. Pochodne $\psi$
w przesuniętych argumentach zbiegają jednostajnie
do pochodnych w $(t,x)$, a pozostałe czynniki są
ograniczone na wspólnym zwartym nośniku. Ostatecznie
\[
 \int_0^T\!\int_\R
 \bigl(|u-v|\psi_t+
 \operatorname{sgn}(u-v)(f(u)-f(v))\psi_x\bigr)
 \,\mathrm d\lambda_1(x)\,\mathrm d\lambda_1(t)\ge0.
\]
Jest to~\eqref{eq:pde24-kato} w sensie dystrybucyjnym.
\end{proof}

\begin{twierdzenie}[Kontrakcja w $L^1$ i jednoznaczność]
\label{thm:pde24-L1}
Przy założeniach Twierdzenia~\ref{thm:pde24-kato} zachodzi
\begin{equation}\label{eq:pde24-L1}
 \|u(t)-v(t)\|_{L^1(\R)}
 \le\|u_0-v_0\|_{L^1(\R)},\qquad 0\le t\le T.
\end{equation}
W szczególności dane początkowe wyznaczają co najwyżej
jedno rozwiązanie entropijne w powyższej klasie.
\end{twierdzenie}
\begin{proof}
Oznaczmy $w=|u-v|\in C([0,T];L^1)$ i
$q=\operatorname{sgn}(u-v)(f(u)-f(v))$. Jeśli
$M=\max_{|z|\le B}|f'(z)|$, gdzie $B$ ogranicza
$|u|,|v|$, to $|q|\le Mw$. W nierówności Kato
testujemy nieujemną funkcją aproksymującą
$\chi_{[s,t]}(\tau)\zeta_R(x)$, gdzie $\zeta_R$ jest
równa $1$ na $[-R,R]$, równa $0$ poza $[-2R,2R]$
i spełnia $|\zeta_R'|\le C/R$. Przechodząc do granicy
w aproksymacji czasowej, otrzymujemy
\[
 \int_\R w(t,x)\zeta_R(x)\,\mathrm d\lambda_1(x)
 \le \int_\R w(s,x)\zeta_R(x)\,\mathrm d\lambda_1(x)
 +\frac{CM}{R}\int_s^t\!\int_\R w(\tau,x)
 \,\mathrm d\lambda_1(x)\,\mathrm d\tau.
\]
Wybieramy $0\le\zeta_R\le1$. Dla $R\to\infty$
całki po obu stronach przechodzą do norm $L^1$
na mocy twierdzenia o zbieżności zmajoryzowanej.
Ostatni człon znika, ponieważ $w$ jest ciągła
w $L^1$ na zwartym przedziale czasowym.
Otrzymujemy kontrakcję dla $0<s<t<T$,
a przypadki $s=0$ i $t=T$ wynikają z ciągłości
w $L^1$.
\end{proof}

\begin{lemat}[Lokalna postać nierówności Kato]
\label{lem:pde24-local-kato}
Niech $f\in C^1(\R)$, a funkcje $u,v$ niech należą do
\[
 C([0,T];L^1_{\mathrm{loc}}(\R))\cap
 L^\infty_{\mathrm{loc}}([0,T]\times\R).
\]
Załóżmy, że spełniają lokalnie wszystkie nierówności entropijne
Krużkowa i mają lokalne ślady początkowe $u_0,v_0$.
Wówczas nierówność~\eqref{eq:pde24-kato} zachodzi
w $\mathcal D'((0,T)\times\R)$. Do wyprowadzenia
lokalnej kontrakcji na zwężających się przedziałach
wystarcza ograniczenie $|f'|$ na zakresie wartości
$u,v$ w rozważanym stożku.
\end{lemat}
\begin{proof}
Każda funkcja testowa $\psi\in C_c^\infty((0,T)\times\R)$
ma nośnik wewnątrz pewnego zwartego walca.
W dowodzie Twierdzenia~\ref{thm:pde24-kato}
całkujemy cztery zmienne tylko w niewielkim powiększeniu
tego walca, bo funkcje
$\rho_\delta(t-s)\rho_\delta(x-y)$ mają małe nośniki.
Zatem potrzebujemy jedynie ograniczeń $L^\infty$
i ciągłości translacji w $L^1$ na tym walcu,
a nie całkowalności $u,v$ na całej prostej.
Przejście diagonalne daje nierówność Kato lokalnie.
W testowaniu przez przybliżenie wskaźnika stożka
stosujemy analogiczne odcięcie przestrzenne; używane
granice czasu wynikają z założonej ciągłości w
$L^1_{\mathrm{loc}}$. Warunek $|f'|\le M$ na
wspólnym zakresie wartości zapewnia $|q|\le M|u-v|$.
\end{proof}

\begin{wniosek}[Skończona prędkość propagacji]
\label{cor:pde24-finite}
Niech $u,v$ spełniają założenia
Lematu~\ref{lem:pde24-local-kato}. Ustalmy $a<b$,
$t\in[0,T]$ oraz $M\ge0$. Załóżmy, że $|f'|\le M$
na przedziale zawierającym istotne zakresy wartości
$u,v$ w zbiorze
\[
 \{(\tau,x):0\le\tau\le t,\ 
 a-M(t-\tau)\le x\le b+M(t-\tau)\}
\]
i jego pewnym otoczeniu względnym w $[0,T]\times\R$.
Jeżeli $u,v$ są globalnie ograniczone przez $B$,
to $M=\max_{|z|\le B}|f'(z)|$ działa dla wszystkich
$a,b,t$. Przy podanych założeniach zachodzi
\begin{equation}\label{eq:pde24-cone}
 \int_a^b|u(t,x)-v(t,x)|\,\mathrm d\lambda_1(x)
 \le\int_{a-Mt}^{b+Mt}|u_0(x)-v_0(x)|\,\mathrm d\lambda_1(x).
\end{equation}
\end{wniosek}
\begin{proof}
Na mocy Lematu~\ref{lem:pde24-local-kato}
możemy zastosować nierówność Kato lokalnie. Używamy
funkcji testowych przybliżających wskaźnik zwężającego się przedziału
$[a-M(t-\tau),b+M(t-\tau)]$ dla $0\le\tau\le t$.
Jego lewy koniec $L(\tau)=a-M(t-\tau)$ porusza się
z prędkością $M$, a prawy $R(\tau)=b+M(t-\tau)$
z prędkością $-M$. Niech $\rho_\delta$ będzie nieujemnym gładkim jądrem
o całce $1$ i nośniku w $[-\delta,\delta]$. Zamiast
podstawiać niegładką funkcję charakterystyczną lub
wymagać istnienia śladów $w=|u-v|$ i $q$ na brzegach
stożka, definiujemy
\[
 \zeta_\delta(\tau,x)=
 \int_{L(\tau)}^{R(\tau)}\rho_\delta(x-z)\,
 \mathrm d\lambda_1(z).
\]
Z reguły Leibniza oraz całkowania po $z$ wynika
\[
 \partial_\tau\zeta_\delta
 =-M\bigl(\rho_\delta(x-L(\tau))+
             \rho_\delta(x-R(\tau))\bigr),
 \qquad
 \partial_x\zeta_\delta
 =\rho_\delta(x-L(\tau))-\rho_\delta(x-R(\tau)).
\]
Stąd prawie wszędzie
\[
 w\partial_\tau\zeta_\delta+q\partial_x\zeta_\delta
 =(q-Mw)\rho_\delta(x-L(\tau))
   +(-q-Mw)\rho_\delta(x-R(\tau))\le0.
\]
Mnożymy $\zeta_\delta$ przez nieujemne gładkie
odcięcie czasowe przybliżające $\chi_{[s,t]}$,
gdzie $0<s<t<T$, i stosujemy nierówność Kato.
Po przejściu do granicy w odcięciu czasowym otrzymujemy
\[
 \int_\R w(t,x)\zeta_\delta(t,x)\,\mathrm d\lambda_1(x)
 \le
 \int_\R w(s,x)\zeta_\delta(s,x)\,\mathrm d\lambda_1(x).
\]
Ciągłość $w$ w $L^1_{\mathrm{loc}}$ pozwala
przejść z $s$ do $0$, a z $t$ do $T$, gdy jest to
potrzebne. Ponieważ $0\le\zeta_\delta\le1$,
nośniki pozostają w jednym ograniczonym przedziale,
a $\zeta_\delta$ zbiega prawie wszędzie do funkcji
charakterystycznej odpowiedniego przedziału,
twierdzenie o zbieżności zmajoryzowanej przy
$\delta\downarrow0$ daje~\eqref{eq:pde24-cone}.
Dla $t=0$ nierówność jest równością.
\end{proof}

\section{Istnienie przez zanikającą lepkość}
\begin{twierdzenie}[Istnienie dla danych $BV$]
\label{thm:pde24-existence}
Niech $f\in C^2(\R)$ oraz
$u_0\in BV(\R)\cap L^1(\R)\cap L^\infty(\R)$.
Wtedy dla każdego $T>0$ istnieje dokładnie jedno
rozwiązanie entropijne $u\in C([0,T];L^1(\R))\cap L^\infty((0,T)\times\R)$
z danymi $u_0$. Dla każdego $t\in[0,T]$ spełnia ono
\begin{equation}\label{eq:pde24-BVbound}
 \|u(t)\|_\infty\le\|u_0\|_\infty,
 \qquad \operatorname{TV}(u(t))\le\operatorname{TV}(u_0),
 \qquad \|u(t)\|_1\le\|u_0\|_1.
\end{equation}
\end{twierdzenie}
\begin{proof}
\textbf{Krok 1: aproksymacja lepkościowa.}
Zastępujemy $f$ przez $f-f(0)$, co nie zmienia równania,
i oznaczamy nadal tę funkcję przez $f$.
Niech $u_0^\eps=G_{\eps^2}*u_0$, gdzie $G_{\eps^2}$ jest
nieujemnym gładkim jądrem o masie $1$ i szerokości
$O(\eps)$. Rozwiązujemy
\begin{equation}\label{eq:pde24-visc}
 u_t^\eps+\partial_x f(u^\eps)=\eps u_{xx}^\eps,
 \qquad u^\eps(0)=u_0^\eps.
\end{equation}
Dla ustalonego $\eps>0$ rozwiązanie otrzymujemy
lokalnie z wzoru Duhamela
\[
 u^\eps(t)=G_{\eps t}*u_0^\eps
 -\int_0^t\partial_xG_{\eps(t-s)}*f(u^\eps(s))\,\mathrm ds,
\]
gdzie $G_t(x)=(4\pi t)^{-1/2}e^{-x^2/(4t)}$.
Uzasadnimy istnienie i całkowalność jeszcze przed
użyciem kontrakcji lepkościowej. Ustalmy $\eps>0$ i
$B=\|u_0\|_\infty$. Zastępujemy tymczasowo $f$ przez
funkcję $\widetilde f\in C^2(\R)$ o ograniczonych
pierwszej i drugiej pochodnej, równą $f$ na $[-B,B]$,
na przykład mnożąc $f$ przez gładkie odcięcie równe
$1$ na otoczeniu tego przedziału. Niech
$L=\|\widetilde f'\|_\infty$.
W przestrzeni Banacha
$X=L^1(\R)\cap BUC(\R)$, z normą
$\|w\|_X=\|w\|_1+\|w\|_\infty$, przez $BUC$ rozumiemy
ograniczone funkcje jednostajnie ciągłe. Zupełność
$X$ wynika z zupełności obu przestrzeni: granice ciągu
Cauchy'ego w $L^1$ oraz w normie supremum są równe
prawie wszędzie, co widać po wyborze podciągu
zbieżnego prawie wszędzie w $L^1$. Jądro Gaussa
jest silnie ciągłą półgrupą kontrakcji na $X$:
w normie $L^1$ wynika to z ciągłości translacji,
a w normie supremum z jednostajnej ciągłości funkcji.
Ponieważ $\widetilde f(0)=0$, superpozycja
$w\mapsto\widetilde f(w)$ jest $L$-lipschitzowska na $X$.
Dla prawej strony wzoru Duhamela, oznaczonej przez
$\mathcal P$, nierówność Younga i oszacowanie
$\|\partial_xG_{\eps t}\|_1=C(\eps t)^{-1/2}$ dają
\[
 \|\mathcal P a-\mathcal P b\|_{C([0,S];X)}
 \le 2CL\sqrt{S/\eps}\,
     \|a-b\|_{C([0,S];X)}.
\]
Całka definiuje funkcję ciągłą w czasie: osobno
rozpatrujemy część oddaloną od $s=t$, gdzie jądro
zależy ciągle od czasu w $L^1$, oraz część bliską
$s=t$, ograniczoną przez całkę z $(t-s)^{-1/2}$.
Dla małego $S$ stała kontrakcji jest mniejsza od $1$.
Twierdzenie Banacha o punkcie stałym i powtarzanie
konstrukcji na przedziałach tej samej długości daje
rozwiązanie w $C([0,T];X)$. Ten sam argument
w samej przestrzeni $BUC(\R)$ daje jednoznaczność
rozwiązania całkowego także w $C([0,T];BUC(\R))$:
stała kontrakcji w normie supremum jest taka sama.
Będziemy korzystać z tej jednoznaczności przy
uzasadnianiu dodatkowej regularności.

Dla gładkich danych Gaussa rozwiązanie ma także
regularność klasyczną. Można to uzasadnić lokalnie
tym samym argumentem w przestrzeni funkcji klasy
$C^1$, dla których funkcja i jej pochodna są
ograniczone oraz jednostajnie ciągłe, z normą
$\|w\|_\infty+\|w_x\|_\infty$. Jest to przestrzeń
Banacha: jednostajne granice funkcji i pochodnych
są związane twierdzeniem Newtona--Leibniza.
Stosujemy odwzorowanie Picarda na domkniętej kuli
wokół ewolucji danych przez półgrupę Gaussa;
dla małego czasu pozostaje ono w tej kuli i jest
kontrakcją, gdyż całka z normy pochodnej jądra
jest rzędu $\sqrt{S/\eps}$. Dokładniej, po zróżniczkowaniu wzoru Duhamela
przenosimy jedną pochodną na strumień, czyli używamy
$\partial_x\widetilde f(w)=\widetilde f'(w)w_x$.
Dla funkcji $a,b$ z ograniczonej kuli
$\|a_x\|_\infty,\|b_x\|_\infty\le K$ mamy
\[
 \begin{aligned}
 \|(\widetilde f(a)-\widetilde f(b))_x\|_\infty
 &\le L\|a_x-b_x\|_\infty\\
 &\quad+K\|\widetilde f''\|_\infty\|a-b\|_\infty .
 \end{aligned}
\]
Wraz z oszacowaniem w normie supremum daje to
lokalną lipschitzowskość superpozycji w tej normie.
Otrzymany punkt stały jest równy wcześniej
skonstruowanemu rozwiązaniu dzięki jednoznaczności
w $C([0,S];BUC(\R))$. Dla $z(t)=\|u_x^\eps(t)\|_\infty$
otrzymujemy na każdym przedziale zaczynającym się
w chwili $r$
\[
 z(t)\le z(r)+CL\eps^{-1/2}
 \int_r^t(t-s)^{-1/2}z(s)\,\mathrm d\lambda_1(s).
\]
Wybór długości przedziału takiej, by
$2CL\sqrt{(t-r)/\eps}\le1/2$, daje
$\sup z(t)\le2z(r)$ na tym przedziale.
Iteracja wyklucza utratę ograniczoności pierwszej
pochodnej w skończonym czasie. Ponieważ norma supremum
rozwiązania jest również ograniczona (z konstrukcji
w $C([0,T];X)$), lokalne twierdzenie o punkcie stałym
można ponawiać z czasem istnienia zależnym tylko od
tych dwóch ograniczeń, $\eps$ i strumienia.
Uzasadnia to kontynuację regularnego rozwiązania
na cały przedział $[0,T]$. Aby uzyskać drugą
pochodną, zapisujemy wzór Duhamela w postaci
z wymuszeniem $F=-\widetilde f'(u^\eps)u_x^\eps$.
Dla $0<\alpha<1$ oszacowania jąder
$\|\partial_xG_\tau\|_1\le C\tau^{-1/2}$ oraz
$\|\partial_{xx}G_\tau\|_1\le C\tau^{-1}$ dają
\[
 \|\partial_xG_\tau(\cdot+h)-\partial_xG_\tau\|_1
 \le C_\alpha |h|^\alpha\tau^{-(1+\alpha)/2}.
\]
Istotnie bierzemy minimum oszacowań przez
$2\|\partial_xG_\tau\|_1$ i
$|h|\|\partial_{xx}G_\tau\|_1$.
Wykładnik $(1+\alpha)/2$ jest mniejszy od $1$,
zatem $u_x^\eps$, a w konsekwencji $F$, są
jednostajnie hölderowskie względem $x$ na skończonym
przedziale czasu (dane Gaussa mają ograniczone
pochodne). Ponieważ $\int_\R\partial_{xx}G_\tau=0$,
przy obliczaniu drugiej pochodnej całki Duhamela
odejmujemy $F(s,x)$ pod całką przestrzenną.
Otrzymujemy całkowalne w czasie oszacowanie
$C_{\eps,\alpha}(t-s)^{-1+\alpha/2}$ dzięki
skalowaniu jądra Gaussa. Aby uzasadnić różniczkowanie, najpierw obcinamy
całkę Duhamela do $0\le s\le t-\eta$, $\eta>0$.
Pochodne rzędu drugiego tych całek zbiegają
jednostajnie (na każdym zwartym przedziale dodatnich
czasów), ponieważ brakujący fragment jest ograniczony
przez $C_{\eps,\alpha}\int_0^\eta
r^{-1+\alpha/2}\,\mathrm d\lambda_1(r)$.
Całki i ich pierwsze pochodne również zbiegają
jednostajnie, z oszacowań jąder rzędu $0$ i $1$.
Zatem granica ma ciągłą drugą pochodną przestrzenną.
Wzór Duhamela spełnia równanie z wymuszeniem $F$
w sensie dystrybucji (testowanie i Fubini, następnie
równanie ciepła dla jądra i jego ślad w chwili $0$).
Prawa strona $H=\eps u_{xx}^\eps+F$ jest ciągła.
Dla ustalonego $r>0$ funkcja
$u^\eps(t,x)-\int_r^t H(s,x)\,\mathrm d\lambda_1(s)$
ma zerową pochodną dystrybucyjną względem $t$.
Testując oddzielnie względem $x$, a następnie
korzystając z ciągłości, wnioskujemy, że jest równa
$u^\eps(r,x)$. Stąd $\partial_t u^\eps=H$
punktowo i rozwiązanie należy do klasy $C^{1,2}$.

Zasada maksimum dla
$u_t^\eps-\eps u_{xx}^\eps+
\widetilde f'(u^\eps)u_x^\eps=0$ daje
$|u^\eps|\le B$. Na całej prostej można ją zastosować
z barierą $\delta e^{At}(1+x^2)$, gdzie
$A>L+2\eps+1$: operator paraboliczny tej bariery
jest ściśle dodatni, bo $2L|x|\le L(1+x^2)$.
Funkcja $u^\eps-B$ pomniejszona o barierę jest
ujemna początkowo i dla dużych $|x|$, jednostajnie
na skończonym przedziale czasu. Gdyby osiągała
dodatnie maksimum na takim walcu, to w punkcie
maksimum (także na jego górnej podstawie) pochodna
czasowa od lewej byłaby nieujemna, pochodna
przestrzenna zerowa, a druga pochodna przestrzenna
niedodatnia. Operator paraboliczny miałby zatem
wartość nieujemną, wbrew jego ścisłej ujemności.
Analogicznie rozpatrujemy $-u^\eps-B$, po czym
przechodzimy z $\delta$ do $0$. Rozwiązanie przyjmuje
więc wartości w przedziale, na którym
$\widetilde f=f$, i rozwiązuje pierwotne równanie.
W szczególności $u^\eps\in C([0,T];L^1(\R))$.

\textbf{Krok 2: kontrakcja lepkościowa.}
Dla dwóch rozwiązań $a,b$ klasy $C^{1,2}$ równania lepkościowego
ich różnica $w=a-b$ spełnia
$w_t+\partial_x(f(a)-f(b))=\eps w_{xx}$.
Przyjmijmy $\eta_\delta(w)=\sqrt{w^2+\delta^2}-\delta$.
Po pomnożeniu przez $\eta_\delta'(w)$ i całkowaniu przez
części, korzystając z $\eta_\delta''\ge0$, otrzymujemy
w granicy $\delta\downarrow0$ nierówność Kato
\[
 \partial_t|a-b|+\partial_x\bigl[
 \operatorname{sgn}(a-b)(f(a)-f(b))\bigr]
 \le\eps\partial_{xx}|a-b|.
\]
Aby uzasadnić przejście do granicy, piszemy
$f(a)-f(b)=A(a,b)(a-b)$, gdzie
$A(a,b)=\int_0^1f'(b+\theta(a-b))\,\mathrm d\theta$.
Dokładna tożsamość ma postać
\[
 \partial_t\eta_\delta(w)+\partial_x(A\eta_\delta(w))
 =\eps\partial_{xx}\eta_\delta(w)
 -\eps\eta_\delta''(w)|w_x|^2
 -\bigl(w\eta_\delta'(w)-\eta_\delta(w)\bigr)\partial_x A.
\]
Ponieważ $w\eta_\delta'(w)-\eta_\delta(w)
=\delta-\delta^2/\sqrt{w^2+\delta^2}\in[0,\delta]$,
ostatni wyraz dąży do zera lokalnie w $L^1$.
Pozostałe wyrazy przechodzą do granicy po testowaniu,
z pominięciem niedodatniego składnika dyssypacyjnego.
Jeśli $a-b\in C([0,T];L^1(\R))$ i oba rozwiązania
są ograniczone (jak rozwiązania z kroku 1), to
$|A|\le M$ dla pewnego $M$. Testujemy otrzymaną
nierówność odcięciem $\chi_{[s,t]}$ wygładzonym
w czasie oraz przestrzenną funkcją $\zeta_R$ równą
$1$ na $[-R,R]$, o nośniku w $[-2R,2R]$,
$0\le\zeta_R\le1$,
$|\zeta_R'|\le C/R$, $|\zeta_R''|\le C/R^2$.
Po granicy czasowej
\[
 \int_\R |w(t,x)|\zeta_R(x)\,\mathrm d\lambda_1(x)
 \le \int_\R |w(s,x)|\zeta_R(x)\,\mathrm d\lambda_1(x)
 +C\left(\frac M R+\frac{\eps}{R^2}\right)
   \int_s^t\|w(\tau)\|_1\,\mathrm d\lambda_1(\tau).
\]
Całka po czasie jest skończona dzięki wcześniejszej
konstrukcji w $C([0,T];L^1)$. Zbieżność
zmajoryzowana dla $R\to\infty$ i ciągłość przy $s=0$
dają $\|a(t)-b(t)\|_1\le\|a(0)-b(0)\|_1$.
Stosując to do $b\equiv0$, otrzymujemy ograniczenie
normy $L^1$. Stosując do przesunięć $u^\eps(t,x+h)$
i $u^\eps(t,x)$, dostajemy
\[
 \|u^\eps(t,\cdot+h)-u^\eps(t,\cdot)\|_1
 \le\|u_0^\eps(\cdot+h)-u_0^\eps\|_1
 \le |h|\operatorname{TV}(u_0).
\]
Aby otrzymać oszacowanie wariacji bez dodatkowego
założenia o całkowalności pochodnej, dla
$\phi\in C_c^1(\R)$, $\|\phi\|_\infty\le1$,
zapisujemy
\[
 \int_\R u^\eps(t,x)\phi'(x)\,\mathrm d\lambda_1(x)
 =\lim_{h\to0}\int_\R
 \frac{u^\eps(t,x-h)-u^\eps(t,x)}{h}\phi(x)
 \,\mathrm d\lambda_1(x)
 \le\operatorname{TV}(u_0).
\]
Równość wynika z podstawienia zmiennej i zbieżności
ilorazów różnicowych $\phi$; oszacowanie z poprzedniej
nierówności. Supremum po $\phi$ daje
$\operatorname{TV}(u^\eps(t))\le\operatorname{TV}(u_0)$.

\textbf{Krok 3: jednostajny moduł ciągłości w czasie.}
Z tożsamości Duhamela stosowanej od czasu $t$ do $t+h$
i z $\|\partial_x f(u^\eps(t))\|_{\mathcal M}\le
M\operatorname{TV}(u_0)$, gdzie
$M=\max_{|z|\le\|u_0\|_\infty}|f'(z)|$, wynika
\[
 \|u^\eps(t+h)-u^\eps(t)\|_1
 \le C\operatorname{TV}(u_0)\sqrt{\eps h}
 +Mh\operatorname{TV}(u_0),\qquad 0<\eps\le1.
\]
Dla $w=u^\eps(t)$ twierdzenie Fubiniego i Lemat
~\ref{lem:pde24-bv-translation} dają
\[
 \begin{aligned}
 \|G_{\eps h}*w-w\|_1
 &\le \int_\R G_{\eps h}(z)\|w(\cdot-z)-w\|_1\,
       \mathrm d\lambda_1(z)\\
 &\le \operatorname{TV}(w)
       \int_\R |z|G_{\eps h}(z)\,\mathrm d\lambda_1(z)\\
 &=2\operatorname{TV}(w)\sqrt{\eps h/\pi}.
 \end{aligned}
\]
Ostatnia równość wynika z podstawienia
$r=z^2/(4\eps h)$ na dodatniej półprostej.
Dla $s>0$ reguła łańcuchowa i ograniczenie
$|f'|\le M$ dają
$\|\partial_x f(u^\eps(s))\|_1
\le M\|u_x^\eps(s)\|_1
=M\operatorname{TV}(u^\eps(s))$;
ostatnia równość wynika z reprezentacji $Dw$ jako
miary ze znakiem. Zatem z nierówności Younga
\[
 \left\|\int_t^{t+h}G_{\eps(t+h-s)}*
 \partial_x f(u^\eps(s))\,\mathrm d\lambda_1(s)\right\|_1
 \le Mh\operatorname{TV}(u_0).
\]
Oba oszacowania razem dają żądany moduł ciągłości.

\textbf{Krok 4: zwartość i granica.}
Uzasadnimy względną zwartość lokalną bez dodatkowego
twierdzenia o zwartości w przestrzeni $BV$.
Dzielimy $[-R,R]$ na $N$ równych przedziałów $I_j$
o długości $\ell=2R/N$ i oznaczamy przez $P_Nw$
funkcję równą średniej $w$ na każdym $I_j$.
Korzystając z nierówności trójkąta, Fubiniego
i Lematu~\ref{lem:pde24-bv-translation}, otrzymujemy
\[
 \begin{aligned}
 \|w-P_Nw\|_{L^1([-R,R])}
 &\le \frac1\ell\sum_{j=1}^N
 \int_{I_j}\!\int_{I_j}|w(x)-w(y)|\,
 \mathrm d\lambda_1(y)\,\mathrm d\lambda_1(x)\\
 &\le\frac1\ell\int_{-\ell}^{\ell}
    \|w(\cdot+h)-w\|_1\,\mathrm d\lambda_1(h)\\
 &\le\ell\operatorname{TV}(w).
 \end{aligned}
\]
W drugiej nierówności podstawiamy $h=y-x$;
dla ustalonego $h$ przedziały całkowania po $x$
są rozłączne i zawarte w $\R$.
Zbiór funkcji $P_Nu^\eps(t)$ jest ograniczony
w skończenie wymiarowej przestrzeni funkcji
stałych na $I_j$ (norma $L^1$ nie wzrasta przy $P_N$).
Dla dowolnej zadanej dokładności najpierw wybieramy
$N$ duże, a potem skończoną sieć w tym ograniczonym
zbiorze. Otrzymujemy całkowitą ograniczoność rodziny
$u^\eps(t)|_{[-R,R]}$, a z zupełności $L^1([-R,R])$
wynika jej względna zwartość. Wspólny moduł ciągłości w
czasie oraz diagonalny argument Arzeli--Ascoliego dają
zbieżność podciągu w $C([0,T];L^1_{\mathrm{loc}}(\R))$.
Aby przejść do zbieżności globalnej, wybieramy
$\zeta_R\in C^\infty(\R)$ równą $0$ na $[-R,R]$ i
równą $1$ poza $[-2R,2R]$, przy
$|\zeta_R'|\le C/R$, $|\zeta_R''|\le C/R^2$.
Nierówność Kato dla $a=u^\eps$, $b=0$ daje
\[
 \int_\R|u^\eps(t,x)|\zeta_R(x)\,\mathrm d\lambda_1(x)
 \le\int_\R|u_0^\eps(x)|\zeta_R(x)\,\mathrm d\lambda_1(x)
 +CT\|u_0\|_1\left(\frac M R+\frac1{R^2}\right).
\]
Testy niegładkie względem czasu i o niezwartym
nośniku względem zmiennej przestrzennej uzasadniamy przez wygładzenie czasu i dodatkowe
odcięcie na $[-N,N]$, po czym $N\to\infty$; pochodne
dodatkowego odcięcia mają rząd $N^{-1}$ i $N^{-2}$,
a całki błędów znikają dzięki oszacowaniu normy $L^1$.
Dla danych początkowych twierdzenie Fubiniego
i nierówność trójkąta dają
\[
 \int_{|x|>R}|u_0^\eps(x)|\,\mathrm d\lambda_1(x)
 \le\int_{|y|>R/2}|u_0(y)|\,\mathrm d\lambda_1(y)
 +\|u_0\|_1\int_{|z|>R/2}G_{\eps^2}(z)\,
      \mathrm d\lambda_1(z).
\]
Ostatnią całkę ograniczamy przez
$4R^{-2}\int_\R z^2G_{\eps^2}(z)\,
\mathrm d\lambda_1(z)=8\eps^2/R^2\le8/R^2$.
Zatem ogony rozwiązań znikają jednostajnie
w $\eps\in(0,1]$ i $t\in[0,T]$. Zbieżny lokalnie
podciąg jest więc ciągiem Cauchy'ego w
$C([0,T];L^1(\R))$: na $[-2R,2R]$ używamy zbieżności
lokalnej, a poza nim sumy dwóch oszacowań ogonów.
Zupełność tej przestrzeni daje zbieżność globalną.

\textbf{Krok 5: równanie i nierówności entropijne.}
Wzór słaby dla $u^\eps$ przechodzi do granicy, gdyż
$f$ jest lipschitzowska na wspólnym zakresie wartości,
a składnik lepkościowy
\[
 \eps\int_0^T\!\int_\R u^\eps\varphi_{xx}\,
 \mathrm d\lambda_1(x)\,\mathrm d\lambda_1(t)
\]
dąży do zera. Dla gładkiej wypukłej entropii
korzystamy z~\eqref{eq:pde24-entropyvisc}.
Po przetestowaniu nieujemną $\varphi$ otrzymujemy
\[
 \int_0^T\!\int_\R
 \bigl(\eta(u^\eps)\varphi_t+q(u^\eps)\varphi_x+
       \eps\eta(u^\eps)\varphi_{xx}\bigr)
 \,\mathrm d\lambda_1(x)\,\mathrm d\lambda_1(t)\ge0.
\]
Ostatni składnik dąży do zera dzięki wspólnemu
ograniczeniu $u^\eps$ i zwartemu nośnikowi testu.
Funkcje $\eta,q$ są lipschitzowskie na wspólnym
przedziale wartości, więc silna zbieżność w
$L^1_{\mathrm{loc}}$ pozwala przejść do granicy
w pozostałych składnikach. Lemat
~\ref{lem:pde24-entropy-equivalence} daje warunki
Krużkowa. Zbieżność w $C([0,T];L^1)$ oraz
$G_{\eps^2}*u_0\to u_0$ w $L^1$ gwarantują $u(0)=u_0$. Ograniczenia~\eqref{eq:pde24-BVbound}
przechodzą do granicy w każdej chwili $t$:
norma $L^1$ jest ciągła, dla ograniczenia $L^\infty$
wybieramy podciąg zbieżny prawie wszędzie w tej
ustalonej chwili, a wariacja jest dolnie półciągła,
gdyż jest supremum ciągłych funkcjonałów
$w\mapsto\int_\R w\phi'\,\mathrm d\lambda_1$.
Jednoznaczność wynika z Twierdzenia~\ref{thm:pde24-L1}.
Zbiega cała rodzina: w przeciwnym razie istniałby
ciąg $\eps_n\downarrow0$ o odległości od $u$
w $C([0,T];L^1)$ ograniczonej od dołu przez dodatnią
stałą. Poprzedni argument zwartości dałby podciąg
zbieżny do rozwiązania entropijnego, które z
jednoznaczności musi być równe $u$, sprzeczność.
\end{proof}


\begin{twierdzenie}[Istnienie dla danych $L^1\cap L^\infty$ bez założenia $BV$]
\label{thm:pde24-existence-L1Linf}
Niech $f\in C^2(\R)$, $u_0\in L^1(\R)\cap L^\infty(\R)$
i $T>0$. Wtedy istnieje dokładnie jedno rozwiązanie entropijne
$u\in C([0,T];L^1(\R))\cap L^\infty((0,T)\times\R)$
z warunkiem $u(0)=u_0$. Spełnia ono
\[
 \|u(t)\|_\infty\le\|u_0\|_\infty,\qquad
 \|u(t)\|_1\le\|u_0\|_1
 \quad(0\le t\le T).
\]
Jeżeli dodatkowo $u_0\in BV(\R)$, to
$\operatorname{TV}(u(t))\le\operatorname{TV}(u_0)$.
\end{twierdzenie}
\begin{proof}
Wybieramy ciąg $u_0^n\in C_c^\infty(\R)$ taki, że
$u_0^n\to u_0$ w $L^1(\R)$ oraz
$\|u_0^n\|_\infty\le\|u_0\|_\infty$: wystarczy
obciąć $u_0$ przez funkcje z $[0,1]$, a następnie
zastosować splot z nieujemnym jądrem aproksymującym jedność.
Każde $u_0^n$ należy do $BV\cap L^1\cap L^\infty$,
a z Twierdzenia~\ref{thm:pde24-existence} istnieje
odpowiadające mu rozwiązanie entropijne $u^n$.
Oszacowanie~\eqref{eq:pde24-L1} daje
\[
 \sup_{0\le t\le T}\|u^n(t)-u^m(t)\|_1
 \le\|u_0^n-u_0^m\|_1.
\]
Zupełność $C([0,T];L^1(\R))$ zapewnia więc istnienie
$u$ z $u^n\to u$ w tej przestrzeni. Dla każdej chwili $t$ zbieżność w $L^1$ daje
\[
 \int_\R (|u(t,x)|-B)_+\,\mathrm d\lambda_1(x)
 \le\|u(t)-u^n(t)\|_1\longrightarrow0,
 \qquad B=\|u_0\|_\infty,
\]
gdzie $r_+=\max(r,0)$. Zatem $\|u(t)\|_\infty\le B$
dla każdego $t$, a z Fubiniego otrzymujemy
$\|u\|_{L^\infty((0,T)\times\R)}\le B$.
Ponieważ $f$ i strumienie Krużkowa
$q_k(z)=\operatorname{sgn}(z-k)(f(z)-f(k))$
są lipschitzowskie na wspólnym ograniczonym zakresie
wartości (dla każdego ustalonego $k$), możemy przejść
do granicy we wszystkich nierównościach entropijnych
na testach o zwartym nośniku. Ciągłość w $L^1$
zapewnia ślad $u(0)=u_0$. Ograniczenie normy $L^1$
przechodzi do granicy; jednoznaczność jest konsekwencją
Twierdzenia~\ref{thm:pde24-L1}. Jeżeli $u_0\in BV$,
możemy też skorzystać wprost z Twierdzenia
\ref{thm:pde24-existence}, które daje ograniczenie wariacji.
\end{proof}


\begin{twierdzenie}[Istnienie dla danych ograniczonych]
\label{thm:pde24-existence-Linf}
Niech $f\in C^2(\R)$, $u_0\in L^\infty(\R)$ oraz $T>0$.
Istnieje dokładnie jedna funkcja
\[
 u\in C([0,T];L^1_{\mathrm{loc}}(\R))
 \cap L^\infty((0,T)\times\R)
\]
spełniająca lokalne nierówności Krużkowa
~\eqref{eq:pde24-kruzineq} oraz $u(0)=u_0$
w $L^1_{\mathrm{loc}}(\R)$. Ponadto
$\|u(t)\|_\infty\le\|u_0\|_\infty$ dla każdego $t$.
\end{twierdzenie}
\begin{proof}
Niech $B=\|u_0\|_\infty$,
$M=\max_{|z|\le B}|f'(z)|$ i
$u_0^n=u_0\chi_{[-n,n]}$ dla $n\ge1$.
Z Twierdzenia~\ref{thm:pde24-existence-L1Linf}
otrzymujemy rozwiązania $u^n$ o globalnym
ograniczeniu przez $B$.
Ustalmy $R>0$ i wybierzmy $n,m>R+MT$.
Dla każdego $t\in[0,T]$ Wniosek~\ref{cor:pde24-finite}
daje
\[
 \int_{-R}^R |u^n(t,x)-u^m(t,x)|\,\mathrm d\lambda_1(x)
 \le\int_{-R-Mt}^{R+Mt}|u_0^n(x)-u_0^m(x)|
 \,\mathrm d\lambda_1(x)=0.
\]
A zatem ograniczenia $u^n$ do $[-R,R]$ stabilizują się
w $C([0,T];L^1([-R,R]))$. Dla $R=1,2,\ldots$
otrzymane granice są zgodne przy ograniczaniu
do mniejszych przedziałów. Sklejamy je w funkcję
$u\in C([0,T];L^1_{\mathrm{loc}}(\R))$.
Możemy wybrać wspólnie mierzalnego reprezentanta,
sklejając reprezentanty $u^n$ na kolejnych
pierścieniach przestrzennych; ograniczenie przez $B$
przechodzi na każdą chwilę oraz prawie wszędzie
w czasoprzestrzeni. Każdy test entropijny ma nośnik
w pewnym walcu $(0,T)\times[-R,R]$, gdzie $u=u^n$
prawie wszędzie dla dostatecznie dużego $n$.
Stąd $u$ spełnia wszystkie lokalne nierówności.
Zgodność w chwili $0$ na każdym takim przedziale
daje ślad początkowy. Dwa dowolne rozwiązania
z podanej klasy mają wspólne globalne ograniczenie
(a więc wspólną skończoną prędkość w oszacowaniu
stożkowym), toteż Wniosek~\ref{cor:pde24-finite}
zapewnia ich jednoznaczność.
\end{proof}

\section{Związek z równaniem Hamiltona--Jacobiego i dalsze problemy}
Dla $f(w)=w^2/2$ i funkcji $g\in W^{1,\infty}(\R)$
rozwiązanie lepkościowe równania
$U_t+\tfrac12(U_x)^2=0$ z danymi $U(0,\cdot)=g$
ma pochodną
$w=U_x$, która jest rozwiązaniem entropijnym
równania Burgersa, w odpowiednim sensie dystrybucyjnym.
Dla kwadratowego hamiltonianu jest to już udowodnione
w rozdziale~\ref{chap:pde23}, w sekcji o związku
z równaniem Burgersa i warunkiem entropii, metodą
półwklęsłości rozwiązania Hopfa--Laxa. Otrzymujemy
tam także $w\in C([0,T];L^1_{\mathrm{loc}}(\R))$,
$|w|\le\|g'\|_\infty$ i ślad $w(0)=g'$.
Wniosek~\ref{cor:pde24-finite} zapewnia jednoznaczność
w klasie globalnie ograniczonych rozwiązań lokalnych.
Jeżeli dodatkowo $g'\in L^1(\R)$, to twierdzenie
~\ref{thm:pde24-existence-L1Linf} oraz ta sama lokalna
jednoznaczność identyfikują $w$ z rozwiązaniem
w $C([0,T];L^1(\R))$. Dla dowolnej globalnie
lipschitzowskiej funkcji $g$ (niekoniecznie ograniczonej)
istnienie i jednoznaczność lokalnego rozwiązania
entropijnego wynikają też z Twierdzenia
~\ref{thm:pde24-existence-Linf}, stosowanego do $g'$.
Nie utożsamiamy natomiast
automatycznie granic pochodnych przybliżeń
lepkościowych na podstawie samej zbieżności potencjałów.

\begin{przyklad}[Dyssypacja energii w fali Burgersa]
Dla $f(w)=w^2/2$ wybieramy $\eta(w)=w^2/2$, wtedy
$q(w)=w^3/3$. Przy skoku z $u_L=1$ do $u_R=0$
prędkość wynosi $\sigma=1/2$, a współczynnik
entropijnej delty jest równy
$[q]-\sigma[\eta]=-1/3+1/4=-1/12<0$.
Dla odwróconego skoku współczynnik jest dodatni,
więc warunek entropii zostaje naruszony. Na mocy
Lematu~\ref{lem:pde24-entropy-equivalence} skok ten
nie spełnia również wszystkich nierówności Krużkowa.
\end{przyklad}

\begin{uwaga}[Zakres twierdzeń]
Globalna kontrakcja $L^1$ dotyczy rozwiązań
o całkowalnych i ograniczonych danych; dla oszacowania
wariacji dodatkowo zakładamy $BV$. Twierdzenie
~\ref{thm:pde24-existence-Linf} obejmuje także dane
jedynie ograniczone, w tym dane Riemanna. W tej klasie
stosujemy oszacowanie w stożku zamiast globalnej normy
$L^1$. Wyniki dla układów praw zachowania,
strumieni zależnych od $x$, wyższych wymiarów lub
nie-wypukłych zadań Riemanna wymagają dodatkowych metod.
\end{uwaga}



% ROZDZIAL 25: nowe rownania pololiniowe

\chapter[Równania półliniowe i eksplozja]{Równania półliniowe: lokalne istnienie, przedłużanie, eksplozja i energia}
\label{chap:pde25}

\section{Model, pojęcia i konieczność doboru przestrzeni}
Równanie jest \emph{półliniowe}, gdy najwyższe pochodne niewiadomej
występują liniowo, a ich współczynniki nie zależą od niewiadomej.
Przykładami są $u_t-\Delta u=f(u)$ oraz
$u_{tt}-\Delta u=f(u)$; natomiast $u_t-\diver(a(u)\nabla u)=0$
jest na ogół równaniem quasi-liniowym.
Dla zwięzłości pełny dowód lokalnego istnienia przeprowadzimy
najpierw na płaskim torusie $\T^d=\R^d/(2\pi\mathbb Z)^d$.
Przez $C(\T^d)$ rozumiemy rzeczywistą przestrzeń Banacha funkcji
ciągłych z normą supremum. W przeciwieństwie do znormalizowanej
normy $L^2(\T)$ z rozdziału~\ref{chap:pde-fourier},
tu używamy nieznormalizowanej miary objętości. Miara $\lambda_d$ na torusie jest
miarą Lebesgue'a przeniesioną z kostki $[0,2\pi)^d$;
$|\T^d|=(2\pi)^d$. Na torusie nie ma warunków brzegowych,
więc nieliniowość i mechanizmy rozwoju osobliwości nie są zasłaniane
kwestią regularności brzegu.

\begin{definicja}[Rozwiązanie łagodne]
\label{def:pde25-mild}
Niech $u_0\in C(\T^d)$, $f\in C^1(\R)$ i $T>0$.
Funkcję $u\in C([0,T];C(\T^d))$ nazywamy
\emph{rozwiązaniem łagodnym} równania $u_t-\Delta u=f(u)$
z warunkiem $u(0)=u_0$, jeżeli
\begin{equation}\label{eq:pde25-duhamel}
 u(t)=S(t)u_0+\int_0^tS(t-s)f(u(s))\,\mathrm d\lambda_1(s),
 \qquad 0\le t\le T,
\end{equation}
gdzie $S(t)=e^{t\Delta}$ jest półgrupą cieplną na $\T^d$,
a $f(u(s))$ oznacza złożenie funkcji $f$ z $u(s,\cdot)$.
Całka w~\eqref{eq:pde25-duhamel} jest całką Bochnera
w przestrzeni Banacha $C(\T^d)$.
\end{definicja}

\begin{lemat}[Półgrupa ciepła na torusie]
\label{lem:pde25-heat}
Dla każdego $d\ge1$ operatory $S(t)$, $t\ge0$, tworzą silnie
ciągłą półgrupę na $C(\T^d)$. Są dodatnie, zachowują stałe
i spełniają $\|S(t)v\|_\infty\le\|v\|_\infty$.
Ponadto $\int_{\T^d}S(t)v\,\mathrm d\lambda_d
=\int_{\T^d}v\,\mathrm d\lambda_d$.
\end{lemat}
\begin{proof}
Dla $t>0$ periodyzujemy jądro Gaussa:
$K_t(z)=\sum_{k\in\mathbb Z^d}(4\pi t)^{-d/2}
\exp(-|z+2\pi k|^2/(4t))$.
Na kostce fundamentalnej $|z+2\pi k|\ge 2\pi|k|-2\pi\sqrt d$.
Każda pochodna przestrzenna składnika szeregu jest
wielomianem w $z+2\pi k$ pomnożonym przez funkcję Gaussa.
Na zwartym przedziale dodatnich czasów szeregi tych
pochodnych mają zatem sumowalną majorantę, niezależną
od $z$ i $t$. To samo dotyczy pochodnych czasowych.
Mamy $K_t\ge0$ i
$\int_{\T^d}K_t(z)\,\mathrm d\lambda_d(z)=1$,
ponieważ kostki przesunięte o $2\pi k$ pokrywają
$\R^d$ z dokładnością do zbiorów miary zero, a całka
jądra Gaussa na $\R^d$ wynosi $1$;
stąd dodatniość, kontrakcja i zachowanie stałych dla
$S(t)v=K_t*v$. Splot jąder Gaussa na $\R^d$ oraz
periodyzacja dają $S(t)S(s)=S(t+s)$. Dla $v$ ciągłej
i jednostajnie ciągłej na torusie rozdzielamy całkę splotową
na $|z|<\delta$ i $|z|\ge\delta$ (odległość liczona na torusie).
Pierwszą część oszacowujemy przez moduł ciągłości $v$,
a druga jest ograniczona przez $2\|v\|_\infty$
pomnożone przez masę jądra poza kulą torusową
o promieniu $\delta$. Ta masa nie przekracza
$\int_{|z|\ge\delta}(4\pi t)^{-d/2}
e^{-|z|^2/(4t)}\,\mathrm d\lambda_d(z)$,
gdyż odległość torusowa od zera nie przekracza
odległości euklidesowej żadnego reprezentanta.
Po podstawieniu $z=\sqrt t\,w$ ostatnia całka
dąży do zera przez całkowalność funkcji Gaussa.
Zatem $\|S(t)v-v\|_\infty\to0$. Tożsamość całek wynika
z twierdzenia Fubiniego i jednostkowej masy jądra.
\end{proof}

\begin{uwaga}[Dlaczego nie zaczynamy od $L^2$?]
Dla wielomianowej nieliniowości $f(u)=u^3$ operator
Nemyckiego $u\mapsto f(u)$ nie jest lokalnie
lipschitzowski z $L^2(\T^d)$ do $L^2(\T^d)$.
Istotnie wybierzmy mierzalne zbiory $E_n$ o mierze
$n^{-4}$ i połóżmy $v_n=n\chi_{E_n}$.
Wówczas $\|v_n\|_2=n^{-1}\to0$, ale
$\|v_n^3\|_2=n\to\infty$.
Ponadto na małej kuli w kostce fundamentalnej
funkcja $v(x)=|x|^{-d/4}$, przedłużona przez zero
poza kulą, należy do $L^2$, lecz jej sześcian
nie należy do $L^2$ (całki radialne mają wykładniki
odpowiednio $d/2-1$ oraz $-d/2-1$). Twierdzenie Picarda w tej przestrzeni
nie byłoby więc uzasadnione. Przestrzeń $C(\T^d)$ pozwala
przeprowadzić argument bez założeń o wykładnikach Sobolewa.
\end{uwaga}

\section{Twierdzenie Picarda i kryterium przedłużania}
\begin{twierdzenie}[Lokalne istnienie i jednoznaczność]
\label{thm:pde25-local}
Niech $d\ge1$, $u_0\in C(\T^d)$ i $f\in C^1(\R)$.
Istnieje $T_0>0$ i dokładnie jedno rozwiązanie łagodne
$u\in C([0,T_0];C(\T^d))$ równania~\eqref{eq:pde25-duhamel}.
Rozwiązanie posiada jednoznaczne maksymalne przedłużenie
na przedział $[0,T_{\max})$, gdzie $T_{\max}\in(0,\infty]$.
Jeżeli $T_{\max}<\infty$, to
\begin{equation}\label{eq:pde25-blowalt}
 \limsup_{t\uparrow T_{\max}}\|u(t)\|_\infty=\infty.
\end{equation}
W rzeczywistości w takim przypadku norma dąży do
nieskończoności w sensie granicy.
\end{twierdzenie}
\begin{proof}
\textbf{Krok 1: wybór kuli.} Ustalmy $R>\|u_0\|_\infty+1$.
Na $[-R,R]$ przyjmijmy
$L_R=\max_{|z|\le R}|f'(z)|$ i
$M_R=\max_{|z|\le R}|f(z)|$.
W zupełnej przestrzeni $Z_T=C([0,T];C(\T^d))$
rozważmy domkniętą kulę
$B_R=\{v:\|v\|_{Z_T}\le R\}$.
Definiujemy $\Phi(v)$ jako prawą stronę~\eqref{eq:pde25-duhamel}.
Dla $v\in Z_T$ superpozycja $s\mapsto f(v(s))$ jest
ciągła w normie supremum, gdyż $f$ jest lipschitzowska
na ograniczonym zakresie wartości $v$.
Silna ciągłość i kontraktywność półgrupy gwarantują
ciągłość całki Bochnera względem górnej granicy:
dla $t_n\to t$ rozpatrujemy całki na wspólnym
przedziale i stosujemy zbieżność zmajoryzowaną w normie,
a całkę po różnicy przedziałów ograniczamy przez
$|t_n-t|\sup_{s\in[0,T]}\|f(v(s))\|_\infty$
(dla $v\in B_R$ jest to co najwyżej $|t_n-t|M_R$). Zatem $\Phi$ rzeczywiście przyjmuje
wartości w $Z_T$.
Lemat~\ref{lem:pde25-heat} daje
$\|\Phi(v)(t)\|_\infty\le\|u_0\|_\infty+TM_R$.
Wybierając $T$ tak, by $TM_R\le R-\|u_0\|_\infty$,
otrzymujemy $\Phi(B_R)\subseteq B_R$.

\textbf{Krok 2: kontrakcja.} Twierdzenie o wartości średniej
i kontrakcja półgrupy implikują
$\|\Phi(v)-\Phi(w)\|_{Z_T}\le TL_R\|v-w\|_{Z_T}$.
Jeśli $TL_R<1$, zasada Banacha daje jedyny punkt stały
w $B_R$. Jednoznaczność nie zależy od wybranej kuli:
dla dowolnych dwóch rozwiązań na wspólnym przedziale
ich zakresy są ograniczone przez pewne $R$, a nierówność
$\|u(t)-v(t)\|_\infty\le
L_R\int_0^t\|u(s)-v(s)\|_\infty\,\mathrm d\lambda_1(s)$
i lemat Gronwalla wymuszają równość.

\textbf{Krok 3: sklejanie.} Półgrupowość pozwala przepisać
wzór Duhamela od dowolnego czasu $t_0$:
$u(t_0+t)=S(t)u(t_0)+\int_0^tS(t-s)f(u(t_0+s))\,\mathrm d\lambda_1(s)$.
Tożsamość otrzymujemy, rozdzielając całkę w chwili
$t_0$ i stosując własność półgrupy do pierwszej części.
Odwrotnie, rozwiązanie rozpoczynające się od $u(t_0)$
skleja się z poprzednim, gdyż po podstawieniu wzoru
na $u(t_0)$ spełnia pierwotny wzór Duhamela.
Niech $T_{\max}$ będzie supremum czasów istnienia.
Jednoznaczność na wspólnych przedziałach pozwala
zdefiniować rozwiązanie na $[0,T_{\max})$ przez
zgodne ograniczenia; każde jego przedłużenie poza
ten przedział przeczyłoby definicji supremum.

\textbf{Krok 4: alternatywa eksplozji.} Jeśli $T_{\max}<\infty$
i $\|u(t)\|_\infty\le R$ dla wszystkich $t$ bliskich
$T_{\max}$, to w konstrukcji z kroku pierwszego można
zastosować stale tę samą większą kulę oraz ten sam czas
lokalnego istnienia $\tau_R>0$, niezależny od chwili restartu.
Wybierając $t_0\in(T_{\max}-\tau_R/2,T_{\max})$
przedłużamy rozwiązanie poza $T_{\max}$, sprzeczność.
Ten sam argument pokazuje, że nie może istnieć ciąg
$t_n\uparrow T_{\max}$ z $\|u(t_n)\|_\infty\le R$,
ponieważ restart w $t_n$ daje przedłużenie o stały czas.
Zatem norma rzeczywiście dąży do nieskończoności.
\end{proof}

\begin{twierdzenie}[Stabilność i ciągła zależność od danych]
\label{thm:pde25-stability}
Niech $u,v$ będą rozwiązaniami łagodnymi na $[0,T]$,
z początkami $u_0,v_0\in C(\T^d)$, dla tej samej
$f\in C^1(\R)$. Jeśli
$\max\{\|u\|_{Z_T},\|v\|_{Z_T}\}\le R$, to
$\|u(t)-v(t)\|_\infty\le e^{L_Rt}\|u_0-v_0\|_\infty$
dla każdego $0\le t\le T$.
\end{twierdzenie}
\begin{proof}
Odejmujemy wzory Duhamela, używamy oszacowania
$|f(a)-f(b)|\le L_R|a-b|$ dla $|a|,|b|\le R$,
a następnie lematu Gronwalla w postaci całkowej.
\end{proof}


\begin{uwaga}[Przenoszenie konstrukcji między rozdziałami]
\label{rem:pde-v44-picard-scaling}
Dowód kontrakcji z twierdzenia~\ref{thm:pde25-local} opiera się na
silnej ciągłości i kontraktywności półgrupy w $C(\T^d)$ oraz na
lokalnej lipschitzowskości operatora Nemyckiego w tej normie.
W rozdziale~\ref{chap:pde27} użyjemy dokładnie tego samego
mechanizmu na torusie o okresie $L$ i z dyfuzją $D>0$;
zmiana okresu i czasu zostanie tam uzasadniona wzorem
\eqref{eq:pde-v44-torus-scale}. Nie jest natomiast poprawne
przeniesienie dowodu bezpośrednio na $L^2$ dla dowolnej
nieliniowości wielomianowej: odwzorowanie $v\mapsto v^3$
nie musi przyjmować wartości w $L^2$ ani być lokalnie
lipschitzowskie w tej przestrzeni.
\end{uwaga}


\section{Zasada porównania, dodatniość i istnienie globalne}
\begin{twierdzenie}[Porównanie rozwiązań półliniowych]
\label{thm:pde25-comparison}
Niech $f\in C^1(\R)$, $T>0$, a $u,v\in Z_T$
będą rozwiązaniami łagodnymi z danymi $u_0,v_0$.
Jeżeli $u_0(x)\le v_0(x)$ dla każdego $x\in\T^d$,
to $u(t,x)\le v(t,x)$ dla wszystkich
$(t,x)\in[0,T]\times\T^d$.
Nie zakładamy monotoniczności $f$.
\end{twierdzenie}
\begin{proof}
Oznaczmy $w=v-u$ i
$a(t,x)=\int_0^1f'(u+\theta(v-u))\,\mathrm d\lambda_1(\theta)$.
Wtedy $f(v)-f(u)=aw$ i $|a|\le L$ dla pewnej
stałej $L$ zależnej od zakresów rozwiązań.
Dla $z(t,x)=e^{Lt}w(t,x)$ otrzymujemy równanie
\[
 z(t)=e^{Lt}S(t)w(0)+
 \int_0^t e^{L(t-s)}S(t-s)a(s)z(s)\,\mathrm d\lambda_1(s).
\]
Aby zmienić je na wzór z półgrupą $S$, używamy
tożsamości
\[
 e^{Lt}S(t)=S(t)+L\int_0^t S(t-r)e^{Lr}S(r)\,
 \mathrm d\lambda_1(r),
\]
wynikającej z własności półgrupy i całki z $Le^{Lr}$.
Podstawiając ją do obu składników wzoru dla $z$
oraz zamieniając kolejność całkowania w trójkącie
$0\le s\le r\le t$, otrzymujemy
\[
 z(t)=S(t)w(0)+
 \int_0^t S(t-s)((a(s)+L)z(s))\,\mathrm d\lambda_1(s).
\]
Zamiana kolejności jest dopuszczalna dla całek
Bochnera dzięki ciągłości i ograniczoności
funkcji podcałkowych na zwartym trójkącie.
Ten sam rachunek działa dla dodatkowego źródła
$b(t)$ w równaniu dla $w$: w ostatnim wzorze
pojawia się dodatkowo $e^{Ls}b(s)$ pod całką.
Tu $a+L\ge0$. Kolejne iteracje tego liniowego równania
Volterry, zaczynając od $z_0(t)=S(t)w(0)\ge0$, pozostają
nieujemne, gdyż półgrupa jest dodatnia. Iteracje zbiegają
jednostajnie na dostatecznie krótkim przedziale przez
kontrakcję ($T\|a+L\|_\infty<1$). Na pozostałych
częściach $[0,T]$ powtarzamy ten sam argument.
Otrzymujemy $w\ge0$.
\end{proof}

\begin{wniosek}[Globalność przy wzroście liniowym]
\label{cor:pde25-global}
Jeśli $f\in C^1(\R)$ i dla pewnych $A,B\ge0$ zachodzi
$|f(s)|\le A+B|s|$ dla każdego $s\in\R$, to każde
rozwiązanie łagodne na $\T^d$ istnieje dla wszystkich
$t\ge0$. Ponadto
\[
 \|u(t)\|_\infty\le
 e^{Bt}\|u_0\|_\infty+\frac AB(e^{Bt}-1)
\]
dla $B>0$; przy $B=0$ ostatni składnik należy zastąpić
przez $At$.
\end{wniosek}
\begin{proof}
Z kontrakcji półgrupy i wzoru Duhamela wynika
$\|u(t)\|_\infty\le\|u_0\|_\infty+
\int_0^t(A+B\|u(s)\|_\infty)\,\mathrm d\lambda_1(s)$.
Gronwall daje żądane ograniczenie na każdym skończonym
przedziale. Alternatywa~\eqref{eq:pde25-blowalt}
wyklucza skończony czas maksymalnego istnienia.
\end{proof}

\begin{przyklad}[Równanie logistyczne]
\label{ex:pde25-logistic}
Dla $u_t=\Delta u+u(1-u)$ i
$0\le u_0\in C(\T^d)$ otrzymujemy
$0\le u(t,x)\le M=\max\{1,\|u_0\|_\infty\}$.
Funkcja stała $0$ jest rozwiązaniem, a $M$ jest
supersolucją, bo $M(1-M)\le0$. Porównanie z funkcją
stałą $M$ można uzasadnić bez dodatkowej teorii
supersolucji: niech $m(t)$ rozwiązuje
$m'=m(1-m)$, $m(0)=M$; wtedy $m(t)\le M$.
Twierdzenie~\ref{thm:pde25-comparison} zastosowane
do rozwiązań $u$ i $m$ daje ograniczenie i globalność.
Choć $f(s)=s(1-s)$ rośnie kwadratowo w module,
nie występuje eksplozja dla takich danych nieujemnych.
\end{przyklad}


\begin{twierdzenie}[Oszacowanie przez równanie skalarne; warunek Osgooda]
\label{thm:pde25-osgood}
Niech $d\ge1$, $u_0\in C(\T^d)$ i $f\in C^1(\R)$.
Załóżmy, że funkcja $h:[0,\infty)\to(0,\infty)$ jest
niemalejąca, lokalnie lipschitzowska i spełnia
$|f(s)|\le h(|s|)$ dla wszystkich $s\in\R$.
Niech $M=\|u_0\|_\infty$, a $y:[0,\tau_M)\to[0,\infty)$
będzie maksymalnym rozwiązaniem $y'=h(y)$, $y(0)=M$.
Wówczas maksymalne rozwiązanie łagodne równania
$u_t=\Delta u+f(u)$ spełnia
$T_{\max}\ge\tau_M$ oraz
$\|u(t)\|_\infty\le y(t)$ dla $0\le t<\tau_M$.
Ponadto
\begin{equation}\label{eq:pde25-osgood}
 \tau_M=\int_M^\infty\frac{\mathrm d\lambda_1(s)}{h(s)},
\end{equation}
gdzie dopuszczamy wartość $+\infty$.
W szczególności rozbieżność całki w~\eqref{eq:pde25-osgood}
gwarantuje globalne istnienie rozwiązania $u$.
\end{twierdzenie}
\begin{proof}
Funkcja $H(r)=\int_M^r h(s)^{-1}\,\mathrm d\lambda_1(s)$,
$r\ge M$, jest klasy $C^1$ i ma ściśle dodatnią
pochodną. Jej obrazem jest przedział $[0,\tau_M)$,
gdzie $\tau_M$ jest prawą stroną~\eqref{eq:pde25-osgood}.
Odwrotność $y=H^{-1}$ spełnia $y(0)=M$ oraz
$y'=1/H'(y)=h(y)$. Każde rozwiązanie ma tę postać
przez regułę łańcuchową stosowaną do $H(y(t))$.
Jeżeli $\tau_M<\infty$, to $y(t)\to\infty$ przy
$t\uparrow\tau_M$, więc przedłużenie nie istnieje.
Jeżeli $\tau_M=\infty$, odwrotność jest określona
dla wszystkich nieujemnych czasów.

Ustalmy $T<\min\{T_{\max},\tau_M\}$. Funkcje stałe
w przestrzeni $\overline y(t,x)=y(t)$ oraz
$\underline y(t,x)=-y(t)$ spełniają odpowiednio
$\overline y_t-\Delta\overline y-f(\overline y)\ge0$
i $\underline y_t-\Delta\underline y-f(\underline y)\le0$.
Dla uzasadnienia porównania w klasie łagodnej zauważmy,
że $w=\overline y-u$ spełnia wzór Duhamela z nieujemnym
danym początkowym $M-u_0$ oraz dodatkowym nieujemnym
źródłem $h(y)-f(y)$. Pozostała część źródła jest równa
$a(t,x)w$, gdzie
$a=\int_0^1f'(u+\theta(\overline y-u))\,\mathrm d\lambda_1(\theta)$
jest ograniczona. Mnożnik $e^{Lt}$, $L\ge\|a\|_\infty$,
przekształca to równanie w równanie Volterry z dodatnim
jądrem i nieujemnym współczynnikiem. Iteracja jak w dowodzie
twierdzenia~\ref{thm:pde25-comparison} daje $w\ge0$.
Analogicznie $u-\underline y\ge0$ (dodatkowe źródło wynosi
$f(-y)+h(y)\ge0$). Zatem $|u(t,x)|\le y(t)$.

Gdyby $T_{\max}<\tau_M$, funkcja $y$ byłaby ograniczona
na $[0,T_{\max}]$. Uzyskane oszacowanie przeczyłoby
alternatywie eksplozji z twierdzenia~\ref{thm:pde25-local}.
Wobec tego $T_{\max}\ge\tau_M$.
\end{proof}

\begin{przyklad}[Wzrost szybszy niż liniowy bez eksplozji z tego kryterium]
Dla $h(s)=(1+s)\log(e+s)$ całka
$\int_M^\infty h(s)^{-1}\,\mathrm d\lambda_1(s)$ jest rozbieżna,
gdyż dla $s\ge e$ zachodzi $1+s\le2s$ i
$\log(e+s)\le\log(2s)\le2\log s$, a całka
$\int_e^\infty (s\log s)^{-1}\,\mathrm d\lambda_1(s)$
jest rozbieżna po podstawieniu $r=\log s$.
Każda nieliniowość $f\in C^1(\R)$ spełniająca
$|f(s)|\le(1+|s|)\log(e+|s|)$ daje więc globalne
rozwiązania na torusie dla dowolnych danych ciągłych.
Pokazuje to, że warunek wzrostu liniowego z
wniosku~\ref{cor:pde25-global} nie jest konieczny.
\end{przyklad}


\section{Eksplozja w skończonym czasie: zasada średniej}
\begin{twierdzenie}[Wybuch dla dodatniego źródła na torusie]
\label{thm:pde25-blowup}
Niech $d\ge1$, $p>1$ oraz
$u_0\in C(\T^d)$, $u_0\ge0$, $u_0\not\equiv0$.
Niech $u$ będzie maksymalnym rozwiązaniem łagodnym
równania $u_t=\Delta u+(u_+)^p$, gdzie $u_+=\max\{u,0\}$.
Oznaczmy $m_0=|\T^d|^{-1}\int_{\T^d}u_0(x)
\,\mathrm d\lambda_d(x)>0$. Wówczas $u\ge0$ i
\begin{equation}\label{eq:pde25-mass-blowup}
 T_{\max}\le\frac{m_0^{1-p}}{p-1}<\infty.
\end{equation}
\end{twierdzenie}
\begin{proof}
Funkcja $s\mapsto(s_+)^p$ należy do $C^1(\R)$,
więc twierdzenie lokalne obowiązuje. Porównanie z
rozwiązaniem zerowym daje $u\ge0$. Z zachowania masy
przez półgrupę i~\eqref{eq:pde25-duhamel} wynika
\[
 m(t)=m_0+\int_0^t
 \frac1{|\T^d|}\int_{\T^d}u(s,x)^p\,
 \mathrm d\lambda_d(x)\,\mathrm d\lambda_1(s).
\]
Funkcja $m$ jest klasy $C^1$ i, na mocy nierówności
Jensena dla wypukłej funkcji $s\mapsto s^p$ na
$[0,\infty)$, spełnia $m'(t)\ge m(t)^p$.
Ponieważ $m(t)>0$, możemy podzielić przez $m(t)^p$:
$\frac{\mathrm d}{\mathrm dt}m(t)^{1-p}
=(1-p)m(t)^{-p}m'(t)\le-(p-1)$.
Stąd $m(t)^{1-p}\le m_0^{1-p}-(p-1)t$.
Lewa strona jest dodatnia na całym przedziale istnienia,
co wymusza~\eqref{eq:pde25-mass-blowup}.
\end{proof}

\begin{uwaga}[Źródło eksplozji]
Wzrost normy wynika tu z nierówności różniczkowej dla
średniej, a nie z nieciągłości współczynników lub
braku regularności danych. Co więcej, nawet niewielkie,
niezerowe dane nieujemne prowadzą do eksplozji na torusie.
To twierdzenie nie przenosi się bez zmian na $\R^d$,
gdzie rozpraszanie masy może zapobiegać wybuchowi.
Wymagane są tam inne argumenty i krytyczne wykładniki.
\end{uwaga}

\begin{twierdzenie}[Wybuch dla dużych danych Dirichleta]
\label{thm:pde25-kaplan}
Niech $\Omega\subset\R^d$ będzie ograniczonym obszarem
klasy $C^2$, $p>1$ i $u$ nieujemnym rozwiązaniem klasycznym
na maksymalnym przedziale $[0,T_{\max})$
z regularnością klasyczną aż do brzegu zagadnienia
$u_t-\Delta u=u^p$, $u|_{\partial\Omega}=0$.
Niech $\lambda_1>0$ będzie pierwszą wartością własną
$-\Delta$ z warunkiem Dirichleta, a $\phi_1>0$
odpowiadającą jej funkcją własną
(Twierdzenie~\ref{thm:pde18-simple}),
$\int_\Omega\phi_1\,\mathrm d\lambda_d=1$.
Jeśli
$y_0=\int_\Omega u_0(x)\phi_1(x)\,\mathrm d\lambda_d(x)
>\lambda_1^{1/(p-1)}$,
to $T_{\max}<\infty$ i
$\|u(t)\|_\infty\to\infty$ dla $t\uparrow T_{\max}$,
przy czym wniosek o granicy normy wymaga założenia,
że w rozważanej klasie istnieje jednostajny dodatni
czas lokalnego przedłużenia dla danych o ograniczonej
normie supremum. Sam wniosek $T_{\max}<\infty$
nie wymaga tego dodatkowego założenia.
\end{twierdzenie}
\begin{proof}
Dla $y(t)=\int_\Omega u(t,x)\phi_1(x)
\,\mathrm d\lambda_d(x)$ korzystamy ze słabego równania
własnego dla $\phi_1\in H_0^1(\Omega)$ i z warunku
$u(t)\in H_0^1(\Omega)$:
\[
 \int_\Omega\Delta u\,\phi_1\,\mathrm d\lambda_d
 =-\int_\Omega\nabla u\cdot\nabla\phi_1\,\mathrm d\lambda_d
 =-\lambda_1\int_\Omega u\phi_1\,\mathrm d\lambda_d .
\]
Różniczkowanie pod całką i równanie paraboliczne dają
$y'=-\lambda_1y+\int_\Omega u^p\phi_1\,
\mathrm d\lambda_d$.
Miara $\phi_1\,\mathrm d\lambda_d$ ma masę $1$,
więc Jensen zapewnia $y'\ge-\lambda_1y+y^p$.
Ponieważ $y_0>\lambda_1^{1/(p-1)}$, prawa strona
jest dodatnia w otoczeniu $y_0$. Funkcja $y$ nie może
spaść poniżej $y_0$: przy pierwszym powrocie do
tej wartości z góry jej pochodna byłaby niedodatnia,
podczas gdy nierówność różniczkowa daje pochodną
ściśle dodatnią. Stąd $y$ jest ściśle rosnąca.
Ustalmy $K>\max\{y_0,(2\lambda_1)^{1/(p-1)}\}$.
Jeżeli rozwiązanie przestaje istnieć przed osiągnięciem
$K$, jego maksymalny czas już jest skończony.
W przeciwnym razie, dopóki $y\in[y_0,K]$, pochodna
jest ograniczona od dołu przez dodatnią stałą
$c=\min_{s\in[y_0,K]}(s^p-\lambda_1s)$.
Zatem przy założeniu globalnego istnienia $y$ musi
osiągnąć $K$ najpóźniej w chwili $(K-y_0)/c$.
Od tej chwili $y'\ge\frac12 y^p$, co po całkowaniu
nierówności dla $y^{1-p}$ prowadzi do sprzeczności
najpóźniej po kolejnych $2K^{1-p}/(p-1)$ jednostkach
czasu. Nie istnieje więc globalne rozwiązanie,
nawet nieograniczone w czasie.
Dokładniej, ponieważ $y'\ge y^p-\lambda_1y>0$,
całkowanie daje
\[
 t\le\int_{y_0}^{y(t)}\frac{\mathrm d\lambda_1(s)}
 {s^p-\lambda_1s}
 <\int_{y_0}^{\infty}\frac{\mathrm d\lambda_1(s)}
 {s^p-\lambda_1s}<\infty.
\]
Ostatnia całka jest skończona, gdyż mianownik
jest dodatni na $[y_0,\infty)$ i dla dużych $s$
jest co najmniej równy $s^p/2$.
Jeżeli obowiązuje założone kryterium przedłużania
z jednostajnym czasem lokalnego istnienia dla
ograniczonych danych, to argument restartu
z Twierdzenia~\ref{thm:pde25-local} daje
$\|u(t)\|_\infty\to\infty$ przy $t\uparrow T_{\max}$.
\end{proof}

\begin{twierdzenie}[Regularność rozwiązania łagodnego dla dodatnich czasów]
\label{thm:pde25-parabolic-regularity}
Niech $f\in C^1(\R)$ i niech $u$ będzie rozwiązaniem
łagodnym z Twierdzenia~\ref{thm:pde25-local}.
Dla $0<a<T<T_{\max}$ zachodzi
\[
 u\in C([a,T];C^2(\T^d))\cap C^1([a,T];C(\T^d)),
 \qquad u_t-\Delta u=f(u).
\]
Jeżeli ponadto $u_0\in H^1(\T^d)$, to
$u\in C([0,T];H^1(\T^d))$.
\end{twierdzenie}
\begin{proof}
Dla periodyzowanego jądra cieplnego z Lematu
~\ref{lem:pde25-heat} zachodzi dla $t>0$
\[
 \|\partial_j K_t\|_{L^1(\T^d)}
 \le\|\partial_j G_t\|_{L^1(\R^d)}
 =c_d t^{-1/2},
\]
gdzie $G_t$ jest jądrem Gaussa na $\R^d$.
Nierówność wynika z nierówności trójkąta i sumowania
po przesuniętych kostkach; równość z podstawienia
$z=\sqrt t\,w$ w całce z modułu pochodnej Gaussa.
Stąd $\|\partial_j S(t)v\|_\infty
\le c_dt^{-1/2}\|v\|_\infty$.

Ustalmy $0<r<a<T<T_{\max}$. Ze wzoru Duhamela
i ograniczoności $f(u)$ otrzymujemy na $[r,T]$
\[
 \partial_j u(t)=\partial_j S(t)u_0+
 \int_0^t\partial_j S(t-s)f(u(s))\,
 \mathrm d\lambda_1(s).
\]
Aby uzasadnić różniczkowanie, najpierw obcinamy
całkę do $[0,t-\delta]$. Zarówno całki, jak i ich
pochodne przestrzenne zbiegają jednostajnie na
$[r,T]\times\T^d$, gdyż ogon całki z pochodną ma
normę co najwyżej $2c_d\sqrt\delta
\sup_s\|f(u(s))\|_\infty$. Całki obcięte są ciągłe
w czasie i przestrzeni. Twierdzenie o jednostajnej
zbieżności funkcji i pochodnych daje zatem
$u\in C([r,T];C^1(\T^d))$.

Teraz rozpoczynamy wzór Duhamela w chwili $r$.
Ponieważ $\partial_j f(u)=f'(u)\partial_j u$
jest ciągła w czasie i przestrzeni na $[r,T]$,
możemy przenieść jedną pochodną na funkcję
pod splotem (przez periodyczne całkowanie przez części).
Dostajemy dla $t\ge a$
\[
 \partial_i\partial_j u(t)=\partial_i\partial_j S(t-r)u(r)
 +\int_r^t\partial_i S(t-s)
       \bigl(f'(u(s))\partial_j u(s)\bigr)\,
       \mathrm d\lambda_1(s).
\]
To samo oszacowanie ogona i jednostajna zbieżność
pochodnych uzasadniają wzór oraz ciągłość drugich
pochodnych. W szczególności $\Delta u+f(u)$ jest
ciągła na $[a,T]\times\T^d$.

Wzór Duhamela spełnia równanie dystrybucyjne:
w każdym modzie Fouriera $k\in\mathbb Z^d$ otrzymujemy
\[
 \widehat u_k'(t)=-|k|^2\widehat u_k(t)
                  +\widehat{f(u)}_k(t).
\]
Wynika to z mnożnika $e^{-t|k|^2}$ półgrupy
(otrzymanego przez całkowanie periodyzowanego jądra)
i twierdzenia Newtona--Leibniza. Całkując ten wzór
od $a$ do $t$, otrzymujemy równość wszystkich
współczynników Fouriera obu stron tożsamości
\[
 u(t)=u(a)+\int_a^t(\Delta u(s)+f(u(s)))\,
                      \mathrm d\lambda_1(s).
\]
Jednoznaczność współczynników Fouriera dla funkcji
ciągłych wynika z gęstości wielomianów trygonometrycznych.
Tożsamość w $C(\T^d)$ daje żądaną ciągłą pochodną
czasową.

Na koniec, dla $u_0\in H^1$, półgrupa jest silnie
ciągła w $H^1$ (komutuje ze słabymi pochodnymi
i jest silnie ciągła w $L^2$, co wynika z ciągłości
translacji i aproksymacji Gaussa).
Nierówność Younga i oszacowanie pochodnych jądra
dają
\[
 \left\|\int_0^t S(t-s)f(u(s))\,
              \mathrm d\lambda_1(s)\right\|_{H^1}
 \le C_d(t+\sqrt t)\sup_{0\le s\le t}\|f(u(s))\|_2.
\]
Całka Bochnera w $H^1$ istnieje dzięki tej całkowalnej
majorancie. Prawa strona dąży do zera przy $t\downarrow0$,
co wraz z ciągłością dla dodatnich czasów dowodzi
ciągłości w $H^1$ na $[0,T]$.
\end{proof}

\section{Funkcjonał energii i równanie Allena--Cahna}
\begin{twierdzenie}[Dyssypacja energii]
\label{thm:pde25-energy}
Niech $V\in C^2(\R)$ i niech $u$ będzie rozwiązaniem
łagodnym $u_t=\Delta u-V'(u)$ na $[0,T]$, z danymi
$u_0\in C(\T^d)$. Dla dodatnich czasów, zgodnie
z Twierdzeniem~\ref{thm:pde25-parabolic-regularity},
rozwiązanie jest klasyczne.
Definiujemy
\[
 E(u)=\int_{\T^d}\left(
 \frac12|\nabla u(x)|^2+V(u(x))\right)
 \,\mathrm d\lambda_d(x).
\]
Wtedy $t\mapsto E(u(t))$ jest różniczkowalna dla
$t>0$ i
\begin{equation}\label{eq:pde25-energy}
 \frac{\mathrm d}{\mathrm dt}E(u(t))
 =-\int_{\T^d}|u_t(t,x)|^2\,\mathrm d\lambda_d(x)\le0.
\end{equation}
Jeżeli dodatkowo $u_0\in H^1(\T^d)$, to dla
$0\le t\le T$ zachodzi
\[
 E(u(t))+\int_0^t\|u_t(s)\|_2^2\,
                  \mathrm d\lambda_1(s)=E(u_0).
\]
\end{twierdzenie}
\begin{proof}
Nie zakładamy istnienia mieszanej pochodnej
$\partial_t\nabla u$. Dla $t,t+h>0$ z tożsamości
różnicy kwadratów i periodycznego całkowania przez
części otrzymujemy
\[
 \begin{aligned}
 &\frac{1}{2h}\bigl(\|\nabla u(t+h)\|_2^2-
                       \|\nabla u(t)\|_2^2\bigr)\\
 &\quad=-\frac12\int_{\T^d}
    (\Delta u(t+h)+\Delta u(t))
    \frac{u(t+h)-u(t)}h\,\mathrm d\lambda_d .
 \end{aligned}
\]
Ciągłość $\Delta u$ i różniczkowalność $u$ w normie
supremum dają granicę $-\int_{\T^d}\Delta u\,u_t\,
\mathrm d\lambda_d$. Dla części potencjalnej
twierdzenie o wartości średniej w postaci całkowej
i jednostajna ciągłość $V'$ na wspólnym ograniczonym
zakresie wartości dają pochodną
$\int_{\T^d}V'(u)u_t\,\mathrm d\lambda_d$.
Z równania wynika $-\Delta u+V'(u)=-u_t$,
co dowodzi~\eqref{eq:pde25-energy}.
Pochodna jest ciągła dla dodatnich czasów.
Całkujemy od $a>0$ do $t$. Jeśli $u_0\in H^1$,
to $u(a)\to u_0$ w $H^1$ oraz w normie supremum
przy $a\downarrow0$, zatem $E(u(a))\to E(u_0)$.
Twierdzenie o zbieżności monotonicznej dla całki
z nieujemnej funkcji $\|u_t\|_2^2$ daje tożsamość
od chwili $0$, a także skończoność tej całki.
\end{proof}

\begin{przyklad}[Równanie Allena--Cahna]
Dla $V(s)=\frac14(s^2-1)^2$ równanie
$u_t=\Delta u+u-u^3$ ma stabilne, stałe stany $u\equiv\pm1$
i niestabilny stan $u\equiv0$ w sensie liniaryzacji
jednorodnych perturbacji. Istotnie dla perturbacji jednorodnej liniaryzacja
wokół stanu $a\in\{-1,0,1\}$ daje równanie
$h'=(1-3a^2)h$, a więc zanik dla $a=\pm1$
i wzrost dla $a=0$. Niech
$M\ge\max\{1,\|u_0\|_\infty\}$.
Rozwiązanie skalarne $m'=m-m^3$, $m(0)=M$,
pozostaje w $[1,M]$: pochodna jest niedodatnia
dla $m\ge1$, a jednoznaczność wyklucza przecięcie
rozwiązania stałego $1$. Jest zatem globalne.
Funkcje $m(t)$ i $-m(t)$, stałe względem $x$,
są rzeczywistymi rozwiązaniami równania półliniowego
(nie utożsamiamy ich ze stałymi $\pm M$).
Twierdzenie~\ref{thm:pde25-comparison} daje
$-M\le-m(t)\le u(t,x)\le m(t)\le M$,
a alternatywa eksplozji zapewnia globalność.
Dla dodatnich czasów energia maleje zgodnie
z~\eqref{eq:pde25-energy}; dla danych należących
również do $H^1$ zachodzi tożsamość energetyczna
od chwili początkowej.
Sama monotoniczność energii nie oznacza jednak,
że rozwiązanie musi zbiegać do konkretnego stanu
stałego bez dalszych argumentów zwartości.
\end{przyklad}

\begin{center}
\begin{tikzpicture}[x=1.65cm,y=3.0cm,>=Latex]
 \draw[->] (-1.65,0)--(1.7,0) node[right] {$s$};
 \draw[->] (0,-.045)--(0,.66) node[above] {$V(s)$};
 \draw[blue!75!black,very thick,domain=-1.52:1.52,samples=121]
  plot (\x,{.25*(\x*\x-1)^2});
 \foreach \a in {-1,1}{
  \fill[blue!75!black] (\a,0) circle (1.7pt);
  \node[below] at (\a,0) {$\a$};
 }
 \fill[red!75!black] (0,.25) circle (1.7pt);
 \draw[densely dashed,gray] (-.95,.25)--(0,.25);
 \node[left,font=\small] at (-.95,.25) {$1/4$};
 \node[below left] at (0,0) {$0$};
 \begin{scope}[xshift=6.9cm,y=1cm]
  \draw[->] (-1.65,0)--(1.7,0) node[right] {$s$};
  \foreach \a in {-1,0,1}{
   \fill (\a,0) circle (1.7pt);
   \node[below] at (\a,0) {$\a$};
  }
  \draw[->,thick,blue!75!black] (-1.55,.18)--(-1.12,.18);
  \draw[->,thick,blue!75!black] (-.15,.18)--(-.85,.18);
  \draw[->,thick,blue!75!black] (.15,.18)--(.85,.18);
  \draw[->,thick,blue!75!black] (1.55,.18)--(1.12,.18);
  \node[font=\small] at (0,.65) {$s'=s-s^3$};
 \end{scope}
\end{tikzpicture}

\small Rysunek: potencjał Allena--Cahna
$V(s)=(s^2-1)^2/4$ i kierunki ruchu dla jednorodnego
przestrzennie równania $s'=-V'(s)$. Minima potencjału
$s=\pm1$ są przyciągające, a maksimum lokalne $s=0$
jest odpychające dla tego równania skalarnego.
\end{center}

\section{Półliniowe równanie falowe w jednym wymiarze}
Dotychczas wykorzystywaliśmy wygładzającą półgrupę ciepła.
Dla równania falowego nieliniowość nie jest wygładzana:
wymagana jest przestrzeń fazowa, w której składanie
z $g$ jest lokalnie lipschitzowskim odwzorowaniem
z przestrzeni pierwszej składowej do przestrzeni
drugiej składowej. Pracujemy na $\T^1$.
Niech $A=1-\partial_x^2$ z dziedziną $H^2(\T^1)$
w przestrzeni Hilberta $L^2(\T^1)$, a
$Y=H^1(\T^1)\times L^2(\T^1)$ z normą
$\|(u,v)\|_Y^2=\|A^{1/2}u\|_2^2+\|v\|_2^2$.

\begin{twierdzenie}[Lokalne istnienie dla równania falowego]
\label{thm:pde25-wave}
Niech $g\in C^1(\R)$, $(u_0,u_1)\in Y$
i $T>0$ będzie wystarczająco mały, zależnie od normy
danych i $g$. Wtedy zagadnienie
\[
 u_{tt}-u_{xx}+u=g(u),\quad
 u(0)=u_0,\quad u_t(0)=u_1
\]
ma dokładnie jedno rozwiązanie łagodne
$(u,u_t)\in C([0,T];Y)$. Rozwiązanie zależy
lokalnie lipschitzowsko od danych w normie $Y$.
Ma jednoznaczne maksymalne przedłużenie na
$[0,T_{\max})$. Jeżeli $T_{\max}<\infty$, to
\[
 \lim_{t\uparrow T_{\max}}\|(u(t),u_t(t))\|_Y=\infty.
\]
\end{twierdzenie}
\begin{proof}
\textbf{Krok 1: grupa liniowa.} W bazie Fouriera operator
$A$ ma wartości własne $1+k^2$, $k\in\mathbb Z$.
Definiujemy $C(t)=\cos(tA^{1/2})$ i
$D(t)=A^{-1/2}\sin(tA^{1/2})$ przez mnożniki Fouriera.
Równanie liniowe ma postać
$u(t)=C(t)u_0+D(t)u_1$,
$u_t(t)=-A^{1/2}\sin(tA^{1/2})u_0+C(t)u_1$.
Dla $k\in\mathbb Z$ oznaczmy $\omega_k=\sqrt{1+k^2}$.
W zmiennych $(\omega_k\widehat u(k),\widehat v(k))$
operator przejścia $U(t)$ działa przez macierz
\[
 \begin{pmatrix}
 \cos(t\omega_k)&\sin(t\omega_k)\\
 -\sin(t\omega_k)&\cos(t\omega_k)
 \end{pmatrix}.
\]
Macierz ta jest unitarna, a wzory dodawania argumentów
dają własność grupy. Z tożsamości Parsevala
~\ref{thm:pde6-parseval} (z uwzględnieniem czynnika
$2\pi$ dla nieznormalizowanej miary) wynika,
że $U(t)$ jest izometrią na $Y$. Silna ciągłość
wynika najpierw dla skończonych sum Fouriera,
a następnie dla wszystkich elementów $Y$ przez
gęstość i jednostajne ograniczenie normy operatorów.

\textbf{Krok 2: lokalna lipschitzowskość.} Wymiar jeden
zapewnia zanurzenie $H^1(\T^1)\hookrightarrow C(\T^1)$.
Można je otrzymać bezpośrednio z nierówności
Cauchy'ego--Schwarza:
\[
 \sum_{k\in\mathbb Z}|\widehat w(k)|
 \le\left(\sum_{k\in\mathbb Z}\frac1{1+k^2}\right)^{1/2}
 \left(\sum_{k\in\mathbb Z}(1+k^2)|\widehat w(k)|^2\right)^{1/2}.
\]
Pierwszy szereg jest zbieżny, a drugi odpowiada
normie $H^1$ przez Parsevala. Szereg Fouriera
zbiega więc jednostajnie do ciągłego reprezentanta
$w$ i $\|w\|_\infty\le C\|w\|_{H^1}$.
Jeśli $\|v\|_{H^1},\|w\|_{H^1}\le R$, to
$\|g(v)-g(w)\|_2\le L_{CR}\|v-w\|_2
\le L_{CR}\|v-w\|_{H^1}$, gdzie
$L_{CR}=\max_{|s|\le CR}|g'(s)|$.

\textbf{Krok 3: punkt stały.} Dla $z=(u,v)$ definiujemy
\[
 \Psi(z)(t)=U(t)(u_0,u_1)
 +\int_0^tU(t-s)(0,g(u(s)))\,\mathrm d\lambda_1(s).
\]
Niech $R>\|(u_0,u_1)\|_Y+1$ i
$M_R=\sqrt{2\pi}\max_{|s|\le CR}|g(s)|$.
Silna ciągłość grupy i ciągłość superpozycji
$u\mapsto g(u)$ z $H^1$ do $L^2$ zapewniają
istnienie całki Bochnera i ciągłość $\Psi(z)$.
Na kuli $\|z\|_{C([0,T];Y)}\le R$ mamy
$\|\Psi(z)\|_{C([0,T];Y)}\le\|(u_0,u_1)\|_Y+TM_R$,
ponieważ $\|U(t-s)(0,h)\|_Y=\|h\|_2$.
Ponadto
\[
 \|\Psi(z)-\Psi(w)\|_{C([0,T];Y)}
 \le TL_{CR}\|z-w\|_{C([0,T];Y)}.
\]
Wybieramy $T$ tak, by $TM_R\le R-\|(u_0,u_1)\|_Y$
oraz $TL_{CR}<1$. Zasada Banacha daje punkt stały.
Pierwsza składowa wzoru ma postać
\[
 u(t)=C(t)u_0+D(t)u_1+
 \int_0^tD(t-s)g(u(s))\,\mathrm d\lambda_1(s).
\]
Mnożniki Fouriera pokazują, że
$D'(t)=C(t)$ w sensie silnym na $L^2$
oraz $C'(t)u_0=-A^{1/2}\sin(tA^{1/2})u_0$ w $L^2$.
Ponieważ $D(0)=0$, różniczkowanie całki w $L^2$
daje dokładnie drugą składową punktu stałego.
Jest ona ciągła w $L^2$. Równanie falowe zachodzi
w sensie dystrybucyjnym: każdy współczynnik Fouriera
spełnia wymuszone równanie skalarne
$\widehat u_k''+\omega_k^2\widehat u_k=\widehat{g(u)}_k$.
Przejście od skończonych sum testów Fouriera do
dowolnych testów gładkich jest uzasadnione ich
zbieżnością wraz z pochodnymi.
Dla dwóch rozwiązań o wspólnym ograniczeniu $R$
różnica wzorów całkowych i lemat Gronwalla dają
\[
 \|z(t)-\widetilde z(t)\|_Y
 \le e^{L_{CR}t}\|z(0)-\widetilde z(0)\|_Y .
\]
Dowodzi to jednoznaczności również poza pierwotną
kulą i lokalnej lipschitzowskiej zależności od danych;
czas istnienia jest wspólny dla danych z ustalonej
ograniczonej kuli.

Własność grupy $U(t)$ pozwala rozdzielić całkę
w dowolnej chwili $r$ i zapisać ten sam wzór
z danymi $z(r)$ oraz czasem liczonym od $r$.
Odwrotnie, rozwiązanie skonstruowane od tych danych
skleja się z poprzednim, gdyż po podstawieniu
wzoru na $z(r)$ spełnia pierwotne równanie całkowe.
Jednoznaczność na wspólnych przedziałach umożliwia
sklejenie rozwiązań do supremum czasów istnienia.
Gdyby dla skończonego $T_{\max}$ istniał ciąg
$t_n\uparrow T_{\max}$ z $\|z(t_n)\|_Y\le B$,
konstrukcja z ustaloną kulą o promieniu $B+2$
dawałaby rozwiązanie od każdej chwili $t_n$
na przedziale o tej samej dodatniej długości.
Dla dużego $n$ byłoby to przedłużenie poza
$T_{\max}$, sprzeczność. Dowodzi to podanej granicy.
\end{proof}

\begin{stwierdzenie}[Energia półliniowego równania falowego]
\label{prop:pde25-waveenergy}
Niech $g\in C^1(\R)$, $G(s)=\int_0^s g(r)\,\mathrm d\lambda_1(r)$,
a $(u,u_t)\in C([0,T];Y)$ będzie rzeczywistym
rozwiązaniem łagodnym z Twierdzenia~\ref{thm:pde25-wave}.
Wtedy
\[
 \mathcal E(t)=\int_{\T^1}\left(
 \tfrac12u_t^2+\tfrac12u_x^2+\tfrac12u^2-G(u)\right)
 \,\mathrm d\lambda_1(x)
\]
jest stała w czasie.
\end{stwierdzenie}
\begin{proof}
Dla rozwiązań o regularności
$u\in C([0,T];H^2)\cap C^1([0,T];H^1)
\cap C^2([0,T];L^2)$
mnożymy równanie przez $u_t$ i całkujemy po torusie.
Składnik $-u_{xx}u_t$ po całkowaniu przez części
jest pochodną $\tfrac12\|u_x\|_2^2$; pozostałe
składniki są pochodnymi $\tfrac12\|u_t\|_2^2$,
$\tfrac12\|u\|_2^2$ i
$-\int_{\T^1}G(u)\,\mathrm d\lambda_1$.
Suma pochodnych wynosi zero.

Dla ogólnych danych energetycznych nie zakładamy
istnienia $u_{tt}$ w $L^2$. Niech $P_N$ oznacza
rzut ortogonalny na mody Fouriera $|k|\le N$,
a $u_N=P_Nu$ i $v_N=P_Nu_t$.
Rzuty komutują z $A$; ze skalarnych równań
dla współczynników Fouriera otrzymujemy
\[
 u_N'=v_N,\qquad v_N'=-Au_N+P_Ng(u).
\]
Obie funkcje są klasy $C^1$ w skończenie wymiarowej
przestrzeni funkcji trygonometrycznych. Różniczkując
energię $\mathcal E_N$ zdefiniowaną tym samym
wzorem dla pary $(u_N,v_N)$, dostajemy
\[
 \mathcal E_N'(t)=
 \int_{\T^1}v_N\bigl(P_Ng(u)-g(u_N)\bigr)
 \,\mathrm d\lambda_1 .
\]
Nie jest to na ogół zero, ponieważ rzut rozwiązania
nie musi być rozwiązaniem równania nieliniowego.

Rzuty są kontrakcjami zarówno w $H^1$, jak i $L^2$
oraz zbiegają silnie do identyczności. Zbieżność jest
jednostajna na zwartych podzbiorach tych przestrzeni:
wystarczy przybliżyć zbiór skończoną siecią i użyć
jednostajnego ograniczenia norm rzutów.
Stosując to do zwartych obrazów $[0,T]$ przez
$u,u_t$ i $g(u)$, otrzymujemy jednostajną w czasie
zbieżność $(u_N,v_N)\to(u,u_t)$ w $Y$ oraz
$P_Ng(u)\to g(u)$ w $L^2$.
Z lokalnej lipschitzowskości $g:H^1\to L^2$
wynika również $g(u_N)\to g(u)$ jednostajnie w czasie
w $L^2$. Stąd
\[
 \sup_{0\le t\le T}|\mathcal E_N'(t)|
 \le \sup_t\|v_N(t)\|_2
 \sup_t\|P_Ng(u(t))-g(u_N(t))\|_2\longrightarrow0.
\]
Z zanurzenia $H^1\hookrightarrow C$ wynika jednostajna
zbieżność $u_N$ do $u$ także w zmiennych $(t,x)$.
Funkcja $G$ jest lipschitzowska na wspólnym ograniczonym
zakresie tych funkcji, więc $\mathcal E_N\to\mathcal E$
jednostajnie w czasie. Całkując tożsamość dla
$\mathcal E_N'$ od $0$ do $t$ i przechodząc do granicy,
otrzymujemy $\mathcal E(t)=\mathcal E(0)$.
\end{proof}

\begin{uwaga}[Energia nie zawsze kontroluje normę]
Jeżeli $G$ rośnie szybciej niż kwadratowo i z dodatnim
znakiem, składnik $-G(u)$ może sprawić, że energia
nie jest ograniczona od dołu. Jej zachowanie samo w sobie
nie daje globalnego istnienia. Dla nieliniowości
$g(u)=-u^3$ energia ma natomiast dodatni człon
$\int_{\T^1}u(t,x)^4/4\,\mathrm d\lambda_1(x)$
i kontroluje normę energetyczną.
\end{uwaga}

\begin{wniosek}[Globalność z koercywności energii]
\label{cor:pde25-wave-global}
Niech $g\in C^1(\R)$ i $G(s)=\int_0^s g(r)\,
\mathrm d\lambda_1(r)$. Jeśli dla pewnych
$0\le a<1/2$ i $b\ge0$ zachodzi
$G(s)\le a s^2+b$ dla każdego $s\in\R$,
to rozwiązanie łagodne z dowolnymi danymi w $Y$
istnieje dla wszystkich $t\ge0$.
W szczególności dotyczy to $g(s)=-s^3$.
\end{wniosek}
\begin{proof}
Niech $c=1/2-a>0$. Ze Stwierdzenia
~\ref{prop:pde25-waveenergy} oraz nierówności
dla $G$ wynika
\[
 \mathcal E(0)+2\pi b
 \ge \tfrac12\|u_t(t)\|_2^2+
 \tfrac12\|u_x(t)\|_2^2+c\|u(t)\|_2^2
 \ge c\|(u(t),u_t(t))\|_Y^2 .
\]
Użyliśmy $c\le1/2$ i tożsamości
$\|A^{1/2}u\|_2^2=\|u_x\|_2^2+\|u\|_2^2$
wynikającej z Parsevala.
Alternatywa eksplozji z Twierdzenia
~\ref{thm:pde25-wave} wyklucza skończony
maksymalny czas istnienia. Dla $g(s)=-s^3$
mamy $G(s)=-s^4/4\le0$, więc wystarczy $a=b=0$.
\end{proof}

\section{Eliptyczne równanie półliniowe i tożsamość Pohozaeva}
Rozważmy $-\Delta u=u^p$ w ograniczonym obszarze
$\Omega\subset\R^d$, $d\ge3$, z zerowym warunkiem
Dirichleta i $u>0$ wewnątrz. W przeciwieństwie do teorii
liniowej rozwiązanie może nie istnieć z przyczyn
wynikających z geometrii obszaru i wykładnika $p$.

\begin{lemat}[Tożsamość Pohozaeva]
\label{lem:pde25-pohozaev}
Niech $d\ge2$, $\Omega\subset\R^d$ będzie ograniczonym
obszarem klasy $C^2$, $F\in C^1(\R)$, $F(0)=0$,
a $u\in C^2(\overline\Omega)$ spełnia
$-\Delta u=F'(u)$ w $\Omega$ i $u=0$ na $\partial\Omega$.
Jeśli $\nu$ jest jednostkową normalną zewnętrzną, to
\begin{equation}\label{eq:pde25-pohozaev}
 \frac{d-2}{2}\int_\Omega|\nabla u|^2\,\mathrm d\lambda_d
 -d\int_\Omega F(u)\,\mathrm d\lambda_d
 =-\frac12\int_{\partial\Omega}(x\cdot\nu)
 (\partial_\nu u)^2\,\mathrm d\sigma(x),
\end{equation}
gdzie $\sigma$ jest miarą powierzchniową na brzegu.
\end{lemat}
\begin{proof}
Mnożymy równanie przez $x\cdot\nabla u$ i całkujemy.
Z jednej strony, ponieważ $\nabla F(u)=F'(u)\nabla u$,
$\int_\Omega F'(u)(x\cdot\nabla u)\,\mathrm d\lambda_d
=-d\int_\Omega F(u)\,\mathrm d\lambda_d$;
wyraz brzegowy $\int_{\partial\Omega}(x\cdot\nu)F(u)
\,\mathrm d\sigma$ znika z uwagi na $F(0)=0$.
Z drugiej strony całkowanie przez części daje
\[
 \begin{aligned}
 \int_\Omega(-\Delta u)(x\cdot\nabla u)\,\mathrm d\lambda_d
 &= -\int_{\partial\Omega}\partial_\nu u
 (x\cdot\nabla u)\,\mathrm d\sigma
 +\int_\Omega\nabla u\cdot\nabla(x\cdot\nabla u)
 \,\mathrm d\lambda_d\\
 &= -\int_{\partial\Omega}(x\cdot\nu)(\partial_\nu u)^2
 \,\mathrm d\sigma
 +\int_\Omega\bigl(|\nabla u|^2
 +\tfrac12 x\cdot\nabla|\nabla u|^2\bigr)
 \,\mathrm d\lambda_d.
 \end{aligned}
\]
Ponieważ $u$ ma stały ślad na brzegu, jego pochodne
styczne znikają tam, zatem $\nabla u=(\partial_\nu u)\nu$.
Twierdzenie o dywergencji zastosowane do
$x|\nabla u|^2$ zamienia ostatni składnik objętościowy
na $-\frac d2\int_\Omega|\nabla u|^2\,
\mathrm d\lambda_d+\frac12\int_{\partial\Omega}
(x\cdot\nu)(\partial_\nu u)^2\,\mathrm d\sigma$.
Po porównaniu z pierwszą stroną otrzymujemy
\eqref{eq:pde25-pohozaev}.
\end{proof}

\begin{twierdzenie}[Nieistnienie dodatnich rozwiązań przy wykładniku krytycznym]
\label{thm:pde25-pohozaev-nonexist}
Niech $d\ge3$, $p\ge(d+2)/(d-2)$ oraz
$\Omega\subset\R^d$ będzie ograniczonym obszarem klasy
$C^2$, \emph{ściśle gwiaździstym względem zera}:
$x\cdot\nu(x)>0$ dla każdego $x\in\partial\Omega$.
Nie istnieje funkcja $u\in C^2(\overline\Omega)$
taka, że $u>0$ w $\Omega$, $u=0$ na $\partial\Omega$
i $-\Delta u=u^p$ w $\Omega$.
\end{twierdzenie}
\begin{proof}
Mnożenie równania przez $u$ i całkowanie przez części
daje $\int_\Omega|\nabla u|^2\,\mathrm d\lambda_d
=\int_\Omega u^{p+1}\,\mathrm d\lambda_d$.
Stosujemy Lemat~\ref{lem:pde25-pohozaev}
do $F(s)=s^{p+1}/(p+1)$ dla $s\ge0$,
z dowolnym przedłużeniem $C^1$ na $\R$ i $F(0)=0$.
Otrzymujemy
\[
 \left(\frac{d-2}{2}-\frac{d}{p+1}\right)
 \int_\Omega|\nabla u|^2\,\mathrm d\lambda_d
 =-\frac12\int_{\partial\Omega}(x\cdot\nu)
 (\partial_\nu u)^2\,\mathrm d\sigma.
\]
Dla $p>(d+2)/(d-2)$ lewy współczynnik jest dodatni,
a prawa strona niedodatnia, więc $u\equiv0$,
sprzeczność. W przypadku krytycznym współczynnik
znika; ścisła gwiaździstość wymusza
$\partial_\nu u\equiv0$ na brzegu. Tymczasem $-\Delta u=u^p>0$
i $u>0$ wewnątrz, więc lemat Hopfa (Wniosek~\ref{cor:pde14-hopf-super})
dla obszaru klasy $C^2$ daje $\partial_\nu u<0$ na brzegu.
To sprzeczność kończąca dowód.
\end{proof}

\begin{uwaga}[Wykładnik krytyczny i ograniczenia twierdzenia]
Liczba $p_*=(d+2)/(d-2)$ odpowiada krytycznemu
wykładnikowi zanurzenia Sobolewa, ponieważ
$p_*+1=2d/(d-2)$, a to właśnie potęga $p+1$
występuje w całce z potencjału w tożsamości Pohozaeva.
Zanurzenie ma postać $H_0^1\hookrightarrow L^{2d/(d-2)}$.
Nie twierdzimy, że nie istnieją rozwiązania zmieniające znak
w przypadku krytycznym, ani że zachodzi ten sam wynik
na dowolnych obszarach: użyliśmy dodatniości oraz
ściśle gwiaździstej geometrii w twierdzeniu Pohozaeva.
\end{uwaga}

\section{Porównanie trzech mechanizmów i dalszy plan}
Dla równania ciepła z nieliniowością odpowiedni dobór
przestrzeni i kontrakcja dają lokalną jednoznaczność,
ale wzrost reakcji może wymusić eksplozję. Dla równania
falowego nieliniowość działa w układzie bez wygładzania,
co wymaga lokalnej lipschitzowskości nieliniowości
jako odwzorowania z $H^1$ do $L^2$. Dla zagadnienia
eliptycznego nawet bardzo gładkie dane i brzeg nie
zapewniają istnienia z powodu tożsamości całkowej.
W następnym rozdziale wykorzystamy operatory monotoniczne
i $p$-laplasjan, które prowadzą do innego mechanizmu istnienia
rozwiązań nieliniowych, bez kontrakcji Picarda.


% ROZDZIAL 26: monotonicznosc i plaplasjan

\chapter[Operatory monotoniczne i p-laplasjan]{Operatory monotoniczne, twierdzenie Minty'ego--Browdera i nieliniowe równania eliptyczne}
\label{chap:pde26}

\section{Dlaczego metoda Laxa--Milgrama przestaje wystarczać?}
W rozdziale~\ref{chap:pde25} oraz we wcześniejszych zagadnieniach
eliptycznych występowały albo operatory liniowe, albo nieliniowości
niższego rzędu. Teraz najwyższa pochodna występuje nieliniowo:
\begin{equation}\label{eq:pde26-main}
 -\diver\bigl(|\nabla u|^{p-2}\nabla u\bigr)=f,\qquad p>1.
\end{equation}
Dla $p\ne2$ nie istnieje forma dwuliniowa pozwalająca zastosować
wprost twierdzenie Laxa--Milgrama. Zastępujemy dwuliniowość
\emph{monotonicznością}, która w teorii nieliniowej odgrywa rolę
dodatniości formy kwadratowej.

W całym rozdziale $\Omega\subset\R^d$ oznacza ograniczony obszar
Lipschitzowski, o ile przy konkretnym wyniku nie zaznaczono inaczej.
Ustalamy $1<p<\infty$, $p'=p/(p-1)$ oraz rzeczywistą
przestrzeń $V=W_0^{1,p}(\Omega)$ wyposażoną w normę
\[
 \|u\|_V=\left(\int_\Omega|\nabla u(x)|^p\,
                         \mathrm d\lambda_d(x)\right)^{1/p},
\]
gdzie moduł gradientu jest normą euklidesową.
Nierówność Poincarégo zapewnia jej równoważność ze zwykłą
normą Sobolewa: przedłużamy funkcje testowe przez zero
do kostki $Q$ zawierającej $\overline\Omega$ i stosujemy
Wniosek~\ref{cor:pde9-Poincare-zero}, po czym przechodzimy
do granicy w $W^{1,p}$. Odwzorowanie
$u\mapsto(u,\partial_1u,\ldots,\partial_du)$ identyfikuje
$W^{1,p}$ z domkniętą podprzestrzenią skończonego
iloczynu przestrzeni $L^p$ (domkniętość wynika
z definicji słabych pochodnych). Ponieważ $1<p<\infty$,
przestrzeń ta jest separowalna i refleksywna; te własności
dziedziczy domknięta podprzestrzeń $V$. Jej przestrzeń
sprzężoną oznaczamy przez $V^*=W^{-1,p'}(\Omega)$,
zawsze z normą funkcjonału względem normy $\|\cdot\|_V$.
Dla innej równoważnej normy Sobolewa w oszacowaniach
liczbowych może pojawić się stała równoważności.
Przez $\langle f,v\rangle$ oznaczamy wartość funkcjonału
$f\in V^*$ na $v\in V$; w przypadku $f\in L^{p'}(\Omega)$
jest to całka $\int_\Omega f(x)v(x)\,\mathrm d\lambda_d(x)$.
Dla funkcji $v\in V$ oznaczamy $v_+=\max\{v,0\}$ i
$v_-=\max\{-v,0\}$; obie należą do $V$, a
$\nabla v_+=\chi_{\{v>0\}}\nabla v$ prawie wszędzie
(Wniosek~\ref{cor:pde9-truncation}). Aby uzasadnić
zerowy warunek brzegowy, wybieramy $v_n\in C_c^\infty(\Omega)$
zbieżne do $v$ w $W^{1,p}$. Funkcje $(v_n)_+$
mają zwarty nośnik w $\Omega$ i należą do $W^{1,p}$,
więc po przedłużeniu przez zero można je przybliżyć
splotami o nośnikach nadal zawartych w $\Omega$.
Tym samym należą do $V$. Są ograniczone w $V$
dzięki wzorowi na gradient, a w $L^p$ zbiegają do $v_+$
z nierówności $|a_+-b_+|\le|a-b|$.
Z refleksywności $V$ wybieramy podciąg słabo zbieżny
w $V$ do pewnego $w\in V$. Ciągłe zanurzenie
$V\hookrightarrow L^p$ daje jego słabą zbieżność
do $w$ również w $L^p$. Jednoznaczność granicy
słabej wymusza $w=v_+$ jako klasy funkcji.
Tak samo postępujemy dla $v_-=(-v)_+$.

\begin{definicja}[Monotoniczność i hemiciągłość]\label{def:pde26-monotone}
Niech $X$ będzie rzeczywistą przestrzenią Banacha, a
$A:X\to X^*$ odwzorowaniem określonym na całym $X$.
Mówimy, że $A$ jest
\emph{monotoniczne}, jeżeli
$\langle Au-Av,u-v\rangle\ge0$ dla wszystkich $u,v\in X$,
zaś \emph{ściśle monotoniczne}, jeżeli nierówność jest ostra
przy $u\ne v$. Operator jest \emph{hemiciągły}, jeśli dla
wszystkich $u,v,w\in X$ odwzorowanie
$t\mapsto\langle A(u+tv),w\rangle$ jest ciągłe na $\R$.
Operator $A$ jest \emph{ograniczony na zbiorach ograniczonych},
jeśli dla każdego $R>0$ zachodzi
$\sup_{\|u\|_X\le R}\|Au\|_{X^*}<\infty$.
Nazywamy go \emph{koercywnym}, jeśli
$\langle Au,u\rangle/\|u\|_X\to+\infty$
przy $\|u\|_X\to\infty$.
\end{definicja}

\begin{uwaga}[Warianty słowa koercywność]
Spotyka się również warunek silniejszy
$\langle Au,u\rangle\ge c\|u\|_X^p-C$ dla
$c>0$, $C\ge0$ i $p>1$.
Jest on wystarczający do naszej definicji, ale nie jest z nią
równoważny: dla $X=\R$ i $A(s)=\operatorname{sgn}(s)
\log(1+|s|)$ iloraz koercywności wynosi $\log(1+|s|)$,
natomiast $\langle A(s),s\rangle=|s|\log(1+|s|)$
rośnie wolniej niż $|s|^p$ dla każdego $p>1$.
Koercywność nie oznacza monotoniczności:
$A(s)=s^3-s$ jest koercywny, lecz
$(A(1/2)-A(0))(1/2-0)<0$.
Z kolei operator zerowy jest monotoniczny, ale nie
jest koercywny na $\R$. Obie własności są zatem
niezależnymi hipotezami twierdzenia Minty'ego--Browdera.
\end{uwaga}

\section{Nierówności monotoniczności dla wektorów}
\begin{lemat}[Monotoniczność strumienia potęgowego]\label{lem:pde26-vector}
Dla $p>1$ odwzorowanie $a_p:\R^d\to\R^d$,
$a_p(\xi)=|\xi|^{p-2}\xi$ dla $\xi\ne0$ i $a_p(0)=0$,
jest ciągłe i ściśle monotoniczne. Co więcej,
dla wszystkich $\xi,\eta\in\R^d$ zachodzą:
\begin{align}
 (a_p(\xi)-a_p(\eta))\cdot(\xi-\eta)
 &\ge \frac{2^{2-p}}{p-1}|\xi-\eta|^p &&(p\ge2),\label{eq:pde26-strong}\\
\intertext{a dla $1<p<2$ i $\xi\ne\eta$}
 (a_p(\xi)-a_p(\eta))\cdot(\xi-\eta)
 &\ge(p-1)(|\xi|+|\eta|)^{p-2}|\xi-\eta|^2.\label{eq:pde26-subquad}
\end{align}
W obu przypadkach prawa strona jest dodatnia przy $\xi\ne\eta$.
\end{lemat}
\begin{proof}
Funkcja $\Phi(\xi)=|\xi|^p/p$ jest ściśle wypukła,
a jej gradientem jest $a_p$. Wypukłość daje monotoniczność;
wyjaśniamy dodatkowo jej ilościowy charakter.
Dla $\zeta\ne0$ macierz pochodnej $a_p$ wynosi
$Da_p(\zeta)=|\zeta|^{p-2}I+(p-2)|\zeta|^{p-4}\zeta\otimes\zeta$.
W kierunku $\zeta$ macierz ma wartość własną
$(p-1)|\zeta|^{p-2}$, a w kierunkach prostopadłych
$|\zeta|^{p-2}$, zatem jest ograniczona od dołu
w sensie form kwadratowych przez
$\min\{1,p-1\}|\zeta|^{p-2}I$ (także dla $d=1$).
Dla $\xi=\eta$ nierówność monotoniczności jest trywialna.
Niech $h=\xi-\eta\ne0$. Jeżeli odcinek
$\eta+th$ zawiera zero dla $t=t_*$, to
$|\eta+th|=|h||t-t_*|$. Pochodna funkcji
$t\mapsto a_p(\eta+th)$ jest wtedy ograniczona
w normie przez stałą razy $|t-t_*|^{p-2}$,
co jest całkowalne dla $p>1$.
Stosując twierdzenie Newtona--Leibniza oddzielnie
po obu stronach $t_*$ i przechodząc do granicy
dzięki ciągłości $a_p$ w zerze, uzyskujemy ten sam
wzór całkowy co dla odcinka omijającego zero.
Otrzymujemy
\[
 (a_p(\xi)-a_p(\eta))\cdot(\xi-\eta)
 \ge \min\{1,p-1\}|\xi-\eta|^2
 \int_0^1|\eta+t(\xi-\eta)|^{p-2}\,\mathrm d\lambda_1(t).
\]
Jeśli $1<p<2$, na całym odcinku
$|\eta+t(\xi-\eta)|\le|\xi|+|\eta|$, a wykładnik
$p-2$ jest ujemny. Otrzymujemy~\eqref{eq:pde26-subquad}.
Dla $p=2$ mamy dokładnie $a_2(\xi)=\xi$.
Dla $p>2$ korzystamy z elementarnej nierówności
\[
 \int_0^1|\eta+t(\xi-\eta)|^{p-2}\,\mathrm d\lambda_1(t)
 \ge \frac{2^{2-p}}{p-1}|\xi-\eta|^{p-2}.
\]
Dla przejrzystości można ją uzasadnić po rozłożeniu
$\eta$ na składową równoległą i prostopadłą do
$h=\xi-\eta$: usunięcie składowej prostopadłej
nie zwiększa lewej strony, a funkcja
$b\mapsto\int_0^1|b+t|^{p-2}\,\mathrm d\lambda_1(t)$
ma pochodną $|b+1|^{p-2}-|b|^{p-2}$,
na mocy twierdzenia Newtona--Leibniza.
Jest ona ujemna dla $b<-1/2$ i dodatnia dla
$b>-1/2$, ponieważ $p-2>0$. Minimum osiągamy
więc w $b=-1/2$ i wynosi ono
$2\int_0^{1/2}t^{p-2}\,\mathrm d\lambda_1(t)
=2^{2-p}/(p-1)$.
To prowadzi do stałej $2^{2-p}/(p-1)$ w
\eqref{eq:pde26-strong}. Nie jest to najlepsza możliwa
stała, ale wystarcza do dalszych zastosowań.
Ścisłość wynika z dodatniości uzyskanych prawych stron.
\end{proof}

\begin{przyklad}[Dlaczego $p=2$ jest szczególny]
Dla $p=2$ mamy $a_2(\xi)=\xi$ i
$(a_2(\xi)-a_2(\eta))\cdot(\xi-\eta)=|\xi-\eta|^2$.
Dla $p>2$ różniczka $Da_p(0)$ jest zerowa, zatem
jednostajna eliptyczność linearyzacji zanika przy
$\nabla u=0$. Dla $1<p<2$ pochodna $a_p$ ma
osobliwość przy zerze. Żadne z tych zjawisk nie
niszczy jednak monotoniczności nieliniowego operatora.
\end{przyklad}

\section{Twierdzenie Minty'ego--Browdera z dowodem}
\begin{lemat}[Skończeniewymiarowa zasada istnienia]
\label{lem:pde26-brouwer}
Jeśli $F:\R^n\to\R^n$ jest ciągłe, a dla pewnego
$R>0$ zachodzi $F(x)\cdot x>0$ na sferze $|x|=R$,
to istnieje $x\in\R^n$ takie, że $|x|<R$ i $F(x)=0$.
\end{lemat}
\begin{proof}
Załóżmy przeciwnie, że $F$ nie zeruje się na
$\overline B(0,R)$. Definiujemy na tej kuli
odwzorowanie ciągłe
\[
 T(x)=-\frac{RF(x)}{|F(x)|}.
\]
Jego wartości należą do sfery o promieniu $R$,
a więc do rozważanej kuli. Twierdzenie Brouwera
o punkcie stałym daje $x=-RF(x)/|F(x)|$.
Wtedy $|x|=R$ oraz $F(x)\cdot x=-R|F(x)|<0$,
co jest sprzeczne z hipotezą. Wniosek jest prawdziwy.
\end{proof}

\begin{twierdzenie}[Twierdzenie Minty'ego--Browdera; wersja z pełną dziedziną]
\label{thm:pde26-minty}
Niech $X$ będzie rzeczywistą, separowalną, refleksywną
przestrzenią Banacha. Załóżmy, że $A:X\to X^*$ jest
monotoniczny, hemiciągły, ograniczony na zbiorach ograniczonych
i koercywny w sensie definicji~\ref{def:pde26-monotone}.
Wówczas dla każdego $f\in X^*$ istnieje $u\in X$ takie, że
$Au=f$. Jeśli ponadto $A$ jest ściśle monotoniczny,
to rozwiązanie jest jedyne.
\end{twierdzenie}
\begin{proof}
\textbf{Krok 1: przybliżenia Galerkina.}
Jeżeli $X=\{0\}$, teza jest natychmiastowa.
W przypadku nieskończenie wymiarowym wybieramy
przeliczalny gęsty podzbiór $X$ i usuwamy kolejno
elementy należące do powłoki liniowej poprzednich.
Otrzymujemy liniowo niezależny ciąg $(e_j)$
o gęstej powłoce liniowej. Kładziemy
$X_n=\operatorname{span}\{e_1,\ldots,e_n\}$.
Dla skończenie wymiarowego $X$ wybieramy bazę
i od chwili osiągnięcia jej wymiaru przyjmujemy
$X_n=X$; w poniższym układzie używamy bazy $X_n$.
Szukamy $u_n\in X_n$ spełniającego
\begin{equation}\label{eq:pde26-galerkin}
 \langle Au_n,v\rangle=\langle f,v\rangle
 \quad\text{dla każdego }v\in X_n.
\end{equation}
Oznaczmy $m_n=\dim X_n$. W bazie
$(e_1,\ldots,e_{m_n})$ tworzymy układ
$F:\R^{m_n}\to\R^{m_n}$, którego współrzędne to
\[
 F_j(c)=\left\langle A\left(\sum_{k=1}^{m_n}c_ke_k\right)
                      -f,e_j\right\rangle.
\]
Równość $F(c)=0$ jest równoważna
~\eqref{eq:pde26-galerkin} dla $u_n=\sum_{k=1}^{m_n}c_ke_k$,
gdyż funkcjonały są liniowe względem argumentu testowego.
Hemiciągłość wraz z monotonicznością i lokalną
ograniczonością implikuje ciągłość $A$ od normy $X$
do słabej topologii $X^*$: istotnie, jeśli $z_m\to z$,
to każda słaba granica $Az_m$ (istniejąca na podciągu)
spełnia, po przejściu do granicy w nierówności
$\langle Az_m-Av,z_m-v\rangle\ge0$,
relację $\langle g-Av,z-v\rangle\ge0$ dla wszystkich $v$.
Dokładniej, rozkładamy iloczyn na
$\langle Az_m-Av,z-v\rangle+
\langle Az_m-Av,z_m-z\rangle$.
Pierwszy składnik zbiega do podanej wartości,
a moduł drugiego jest ograniczony przez
$(M+\|Av\|_{X^*})\|z_m-z\|_X\to0$.
Teraz,
podstawiając $v=z-tw$ i przechodząc z $t\downarrow0$
przez hemiciągłość, otrzymujemy
$\langle g-Az,w\rangle\ge0$. Podstawienie $-w$
daje równość, a więc $g=Az$.
Ponieważ $(Az_m)$ jest ograniczony i $X^*$ refleksywna,
każdy podciąg ma słabo zbieżny dalszy podciąg.
Gdyby dla pewnego $w\in X$ ciąg
$\langle Az_m,w\rangle$ nie zbiegał do
$\langle Az,w\rangle$, można by wybrać podciąg
oddalony od tej wartości o dodatnią stałą,
co przeczyłoby właśnie wykazanej identyfikacji.
Stąd wszystkie współrzędne $F_j$ są ciągłe,
a więc $F$ jest ciągłe.

Koercywność i nierówność
$\langle Au-f,u\rangle\ge\langle Au,u\rangle
-\|f\|_{X^*}\|u\|_X$ pokazują, że
$\langle Au-f,u\rangle>0$ dla wszystkich
$\|u\|_X\ge R$ przy dostatecznie dużym $R$,
niezależnym od $n$.
Dla każdego $n$ równoważność norm na $X_n$ daje
stałą $b_n>0$ taką, że
$\|\sum_{j=1}^{m_n} c_je_j\|_X\ge b_n|c|$.
Wybieramy $r_n=R/b_n$, więc norma elementu $X_n$
jest co najmniej $R$ dla $|c|=r_n$.
Wówczas $F(c)\cdot c=\langle Au-f,u\rangle>0$.
Lemat~\ref{lem:pde26-brouwer} daje rozwiązanie
\eqref{eq:pde26-galerkin}.

\textbf{Krok 2: oszacowania i granice słabe.}
Testując~\eqref{eq:pde26-galerkin} przez $u_n$, mamy
$\langle Au_n,u_n\rangle=\langle f,u_n\rangle$.
Koercywność wymusza jednostajne ograniczenie
$\|u_n\|_X\le R$. Lokalna ograniczoność $A$ daje
$\|Au_n\|_{X^*}\le M_R$. Refleksywność $X$ oraz $X^*$
oraz twierdzenie Eberleina--Šmuliana pozwalają
wybrać rosnący ciąg indeksów $(n_j)$, dla którego
$u_{n_j}\rightharpoonup u$ w $X$ oraz
$Au_{n_j}\rightharpoonup g$ w $X^*$.
W dalszych granicach zapisujemy nadal $n$,
rozumiejąc przez niego indeks z tego podciągu,
a $X_n$ oznacza odpowiadającą mu pierwotną podprzestrzeń.
Dla dowolnego $j$ i $n\ge j$ mamy
$\langle Au_n,e_j\rangle=\langle f,e_j\rangle$.
Po przejściu do granicy i wykorzystaniu gęstości
otrzymujemy $g=f$.

\textbf{Krok 3: kluczowa nierówność graniczna.}
Wybieramy $v_n\in X_n$ takie, że $v_n\to u$ silnie
w $X$. Istnienie takiego ciągu wynika z gęstości
rosnących podprzestrzeni $X_n$: dla każdego $k$ wybieramy
$z_k\in X_{n_k}$ z $\|z_k-u\|_X<1/k$ i
$n_k\uparrow\infty$; kładziemy $v_n=z_k$ dla
$n_k\le n<n_{k+1}$. Dzięki równaniu Galerkina
$\langle Au_n,u_n-v_n\rangle=\langle f,u_n-v_n\rangle$.
Ponieważ $u_n\rightharpoonup u$ i $v_n\to u$,
mamy $\langle f,u_n-v_n\rangle\to0$.
Ponadto
$|\langle Au_n,v_n-u\rangle|\le M_R\|v_n-u\|_X\to0$.
Dodając obie równości, otrzymujemy
\begin{equation}\label{eq:pde26-limsup}
 \lim_{n\to\infty}\langle Au_n,u_n-u\rangle=0.
\end{equation}
Dla ustalonego $v\in X$ monotoniczność daje
$\langle Au_n-Av,u_n-v\rangle\ge0$.
Rozkładamy pierwszy człon jako
$\langle Au_n,u_n-u\rangle+\langle Au_n,u-v\rangle$
i korzystamy z~\eqref{eq:pde26-limsup}. W granicy
otrzymujemy $\langle f-Av,u-v\rangle\ge0$.

\textbf{Krok 4: identyfikacja operatora.}
Dla dowolnego $w\in X$ i $t>0$ podstawiamy $v=u-tw$.
Wynika stąd $\langle f-A(u-tw),w\rangle\ge0$.
Przy $t\downarrow0$ hemiciągłość daje
$\langle f-Au,w\rangle\ge0$. Stosując nierówność także
do $-w$, otrzymujemy $Au=f$.
Jeśli $Au=Av=f$, ścisła monotoniczność wymusza $u=v$.
\end{proof}

\begin{uwaga}[Dokładny zakres twierdzenia]
Nie wymagamy liniowości $A$ ani istnienia potencjału
wariacyjnego. Warunek ograniczoności na zbiorach
ograniczonych jest tu przyjęty jawnie, ponieważ używamy
go w dowodzie do wydzielenia słabo zbieżnego podciągu $Au_n$.
Twierdzenie nie obejmuje automatycznie operatorów
określonych jedynie na właściwej części $X$ ani równań,
w których zawodzi monotoniczność.
\end{uwaga}

\section{p-laplasjan jako operator na przestrzeni Sobolewa}
\begin{definicja}[Rozwiązanie słabe $p$-laplasjanu]\label{def:pde26-weak}
Niech $f\in V^*$. Funkcję $u\in V=W_0^{1,p}(\Omega)$
nazywamy słabym rozwiązaniem zagadnienia
\eqref{eq:pde26-main} z zerowym warunkiem Dirichleta,
jeśli dla każdego $v\in V$ zachodzi
\begin{equation}\label{eq:pde26-weak}
 \int_\Omega|\nabla u|^{p-2}\nabla u\cdot\nabla v
 \,\mathrm d\lambda_d(x)=\langle f,v\rangle.
\end{equation}
Operator $A_p:V\to V^*$ definiujemy przez lewą stronę
tego wzoru; zapis $-\Delta_pu=f$ oznacza równość
$A_pu=f$ w $V^*$.
\end{definicja}

\begin{twierdzenie}[Istnienie i jednoznaczność]\label{thm:pde26-plap}
Niech $1<p<\infty$ i niech $\Omega\subset\R^d$ będzie
ograniczonym obszarem Lipschitzowskim. Dla każdego
$f\in W^{-1,p'}(\Omega)$ istnieje dokładnie jedno
$u\in W_0^{1,p}(\Omega)$ spełniające
\eqref{eq:pde26-weak}. Ponadto
\begin{equation}\label{eq:pde26-apriori}
 \|\nabla u\|_{L^p(\Omega)}^{p-1}
 \le\|f\|_{W^{-1,p'}(\Omega)}.
\end{equation}
\end{twierdzenie}
\begin{proof}
Sprawdzamy każdą hipotezę twierdzenia
\ref{thm:pde26-minty} dla $A_p$.
Nierówność Höldera daje
$|\langle A_pu,v\rangle|\le
\|\nabla u\|_p^{p-1}\|\nabla v\|_p$, więc
$\|A_pu\|_{V^*}\le\|u\|_V^{p-1}$.
Dla $u\ne0$ test $v=u/\|u\|_V$ osiąga to ograniczenie:
$\langle A_pu,v\rangle=\|u\|_V^{p-1}$.
Dla $u=0$ obie strony są zerowe, zatem
$\|A_pu\|_{V^*}=\|u\|_V^{p-1}$ dla każdego $u\in V$.
Operator jest określony na całym $V$ i ograniczony
na zbiorach ograniczonych. Ciągłość strumienia $a_p$ oraz
oszacowanie
$|a_p(\nabla u+t\nabla w)|\le
C_p(|\nabla u|+|\nabla w|)^{p-1}$ dla $|t|\le1$
po pomnożeniu przez $|\nabla v|$ dają całkowalną
majorantę: $(p-1)p'=p$, więc stosujemy nierówność
Höldera. Twierdzenie o zbieżności zmajoryzowanej
zapewnia ciągłość $\langle A_p(u+tw),v\rangle$
w $t=0$, a po zastąpieniu $u$ przez $u+t_0w$
również w dowolnym $t_0\in\R$.
Monotoniczność wynika z Lematu~\ref{lem:pde26-vector}.
Jeżeli $u\ne v$ w $V$, to z nierówności Poincarégo
gradienty różnią się na zbiorze dodatniej miary,
na którym całkowany iloczyn jest ściśle dodatni.
Stąd operator jest ściśle monotoniczny.
Koercywność jest dokładna:
$\langle A_pu,u\rangle=\|\nabla u\|_p^p$,
a więc iloraz przez $\|u\|_V$ dąży do nieskończoności.
Twierdzenie Minty'ego--Browdera zapewnia jedyne rozwiązanie.
Testując jego równanie przez $u$, otrzymujemy
$\|u\|_V^p=\langle f,u\rangle\le
\|f\|_{V^*}\|u\|_V$. Dzielenie przez $\|u\|_V$
(dla $u=0$ teza jest oczywista) kończy dowód.
\end{proof}

\section{Metoda bezpośrednia rachunku wariacyjnego}
\begin{twierdzenie}[Charakteryzacja minimizacyjna]\label{thm:pde26-minimizer}
Niech $p>1$, $\Omega$ będzie ograniczonym obszarem
Lipschitzowskim, a $f\in V^*$. Funkcjonał
\[
 J_f(v)=\frac1p\int_\Omega|\nabla v|^p
 \,\mathrm d\lambda_d(x)-\langle f,v\rangle,\qquad v\in V,
\]
ma dokładnie jeden punkt minimum. Jest nim jedyne
rozwiązanie zagadnienia~\eqref{eq:pde26-weak}.
\end{twierdzenie}
\begin{proof}
Z definicji normy funkcjonału otrzymujemy
$J_f(v)\ge\|v\|_V^p/p-\|f\|_{V^*}\|v\|_V$.
Ponieważ $p>1$, prawa strona dąży do nieskończoności
przy $\|v\|_V\to\infty$ i jest ograniczona od dołu
jako ciągła funkcja nieujemnej zmiennej $\|v\|_V$.
Ciąg minimalizujący jest ograniczony w refleksywnej
przestrzeni $V$; wybrany podciąg słabo zbiega do $u$.
Ciągły operator liniowy $\nabla:V\to L^p(\Omega;\R^d)$
zachowuje zbieżność słabą. Norma w przestrzeni $L^p$
jest słabo półciągła z dołu (jako supremum modułów
ciągłych funkcjonałów liniowych o normie co najwyżej $1$),
a funkcjonał liniowy $v\mapsto\langle f,v\rangle$
jest słabo ciągły. Stąd $J_f(u)\le\liminf J_f(v_n)$,
czyli $u$ jest minimizatorem. Ścisła wypukłość
$z\mapsto|z|^p$ daje ścisłą wypukłość $J_f$ na $V$:
jeżeli $u\ne v$ w $V$, nierówność Poincarégo
zapewnia, że gradienty różnią się na zbiorze dodatniej
miary. Dla $0<\theta<1$ całkowanie punktowej nierówności
ścisłej wypukłości daje zatem
$J_f((1-\theta)u+\theta v)<(1-\theta)J_f(u)+\theta J_f(v)$.
Stąd minimum jest jedyne.

Dla dowolnych $u,v\in V$ i $|t|\le1$ iloraz różnicowy
funkcji $|\nabla u+t\nabla v|^p/p$ jest ograniczony
w module przez $(|\nabla u|+|\nabla v|)^{p-1}|\nabla v|$,
co wynika z twierdzenia o wartości średniej.
Majoranta jest całkowalna dzięki nierówności Höldera.
Z twierdzenia o zbieżności zmajoryzowanej wynika
\[
 \left.\frac{\mathrm d}{\mathrm dt}J_f(u+tv)\right|_{t=0}
 =\langle A_pu-f,v\rangle.
\]
Poniżej uzasadnimy również, że $J_f$ jest klasy $C^1$
w sensie różniczkowania Frécheta.
Dla minimizatora pochodna ta jest zerowa w każdym
kierunku, więc zachodzi~\eqref{eq:pde26-weak}.
Odwrotnie, jeżeli $A_pu=f$, to dla dowolnego $w\in V$
wypukła funkcja $t\mapsto J_f(u+t(w-u))$ spełnia
$J_f(w)-J_f(u)\ge\langle A_pu-f,w-u\rangle=0$.
Dowodzi to minimalności.
\end{proof}

\begin{uwaga}[Dlaczego przedstawiliśmy dwa dowody?]
Metoda minimizacji wykorzystuje szczególną strukturę
potencjałową: pochodna kierunkowa $J_f$ w $u$
jest dana przez funkcjonał $A_pu-f$. Twierdzenie Minty'ego--Browdera
działa również dla monotonicznych operatorów bez potencjału.
Przykładem jest strumień $a(x,\xi)$ monotoniczny w $\xi$,
który nie musi być gradientem funkcji wypukłej względem $\xi$.
\end{uwaga}


\begin{twierdzenie}[Strumienie typu Leraya--Lionsa]
\label{thm:pde26-leraylions}
Niech $\Omega\subset\R^d$ będzie ograniczonym obszarem
Lipschitzowskim, $1<p<\infty$ i $V=W_0^{1,p}(\Omega)$.
Niech $a:\Omega\times\R^d\to\R^d$ będzie funkcją
Carathéodory'ego: $x\mapsto a(x,\xi)$ jest mierzalna
przy ustalonym $\xi$, a $\xi\mapsto a(x,\xi)$ jest
ciągła dla prawie każdego $x$. Załóżmy, że istnieją
$C,\alpha>0$, $h\in L^{p'}(\Omega)$, $h\ge0$,
i $\beta\in L^1(\Omega)$, $\beta\ge0$, takie że
prawie wszędzie w $x$ i dla wszystkich $\xi,\eta\in\R^d$
zachodzą
\begin{align}
 |a(x,\xi)|&\le h(x)+C|\xi|^{p-1},\label{eq:pde26-ll-growth}\\
 a(x,\xi)\cdot\xi&\ge\alpha|\xi|^p-\beta(x),\label{eq:pde26-ll-coerc}\\
 (a(x,\xi)-a(x,\eta))\cdot(\xi-\eta)&\ge0.\label{eq:pde26-ll-mono}
\end{align}
Wówczas dla każdego $f\in V^*$ istnieje $u\in V$
spełniające
\begin{equation}\label{eq:pde26-ll-weak}
 \int_\Omega a(x,\nabla u(x))\cdot\nabla v(x)
 \,\mathrm d\lambda_d(x)=\langle f,v\rangle
 \qquad(v\in V).
\end{equation}
Jeżeli nierówność w~\eqref{eq:pde26-ll-mono} jest ostra
przy $\xi\ne\eta$ dla prawie każdego $x$, rozwiązanie
jest jednoznaczne.
\end{twierdzenie}
\begin{proof}
Definiujemy $A:V\to V^*$ lewą stroną
\eqref{eq:pde26-ll-weak}. Najpierw zauważmy, że dla
dowolnej mierzalnej funkcji $\zeta:\Omega\to\R^d$
złożenie $x\mapsto a(x,\zeta(x))$ jest mierzalne
w sensie Lebesgue'a: przybliżamy $\zeta$ punktowo
mierzalnymi funkcjami prostymi $\zeta_n$ o skończonym
zbiorze wartości. Funkcje $a(x,\zeta_n(x))$ są mierzalne
na mocy mierzalności przekrojów, a poza wspólnym
zbiorem wyjątkowym ciągłość względem $\xi$ daje
ich punktową zbieżność do $a(x,\zeta(x))$.
Stosujemy to do $\zeta=\nabla u$.
Zmiana reprezentanta gradientu na zbiorze miary zero
nie zmienia klasy otrzymanego strumienia.
Warunek wzrostu i nierówność Höldera zapewniają
$\|Au\|_{V^*}\le\|h\|_{p'}+C\|u\|_V^{p-1}$;
odwzorowanie jest więc dobrze określone i ograniczone
na zbiorach ograniczonych. Monotoniczność wynika
z całkowania~\eqref{eq:pde26-ll-mono}.
Dla $u,v,w\in V$ oraz $|t|\le1$ zachodzi prawie wszędzie
\[
 |a(x,\nabla u+t\nabla v)\cdot\nabla w|
 \le\bigl(h(x)+C(|\nabla u|+|\nabla v|)^{p-1}\bigr)
 |\nabla w|.
\]
Prawa strona jest całkowalna przez nierówność Höldera.
Ciągłość w zmiennej $\xi$ i twierdzenie o zbieżności
zmajoryzowanej dowodzą ciągłości w $t=0$.
Aby uzyskać ciągłość w dowolnym $t_0\in\R$,
zastępujemy $u$ przez $u+t_0v$.

W rzeczywistości $A:V\to V^*$ jest ciągły w normie.
Niech $u_n\to u$ w $V$. Z każdego podciągu wybieramy
dalszy podciąg $(u_{n_j})$ tak, by
$\|\nabla u_{n_j}-\nabla u\|_p\le 2^{-j}$.
Funkcja
\[
 H=|\nabla u|+\sum_{j=1}^\infty
                         |\nabla u_{n_j}-\nabla u|
\]
należy do $L^p$ na mocy nierówności Minkowskiego
dla sum częściowych i lematu Fatou.
W szczególności szereg jest skończony prawie wszędzie,
więc $\nabla u_{n_j}\to\nabla u$ prawie wszędzie.
Ponadto $|\nabla u_{n_j}|\le H$. Z warunku wzrostu
wynika oszacowanie
\[
 |a(x,\nabla u_{n_j})-a(x,\nabla u)|^{p'}
 \le (2h(x)+2CH(x)^{p-1})^{p'}.
\]
Prawa strona jest całkowalna, ponieważ $(p-1)p'=p$.
Ciągłość strumienia względem $\xi$ i twierdzenie
o zbieżności zmajoryzowanej dają zbieżność strumieni
w $L^{p'}(\Omega;\R^d)$ wzdłuż dalszego podciągu.
Z nierówności Höldera wynika
$\|Au_{n_j}-Au\|_{V^*}\to0$.
Gdyby cały ciąg nie zbiegał do $Au$, można by
wybrać podciąg oddalony w normie od $Au$ o dodatnią
stałą, wbrew właśnie udowodnionej własności.
Z~\eqref{eq:pde26-ll-coerc} otrzymujemy
$\langle Au,u\rangle\ge\alpha\|u\|_V^p-\|\beta\|_1$,
a zatem $A$ jest koercywny.
Wszystkie hipotezy twierdzenia~\ref{thm:pde26-minty}
są spełnione, więc istnieje rozwiązanie.
Dla dwóch rozwiązań $u_1,u_2$ testujemy różnicę
równań przez $u_1-u_2$. Jeśli monotoniczność jest
ostra, nieujemna całkowana funkcja może mieć całkę
zero tylko gdy $\nabla u_1=\nabla u_2$ prawie wszędzie.
Nierówność Poincarégo daje $u_1=u_2$.
\end{proof}

\begin{uwaga}[Różniczka funkcjonału energii]
Strumień $a(x,\xi)=a_p(\xi)$ spełnia hipotezy
Twierdzenia~\ref{thm:pde26-leraylions} z $h=\beta=0$,
$C=\alpha=1$. Z dowodu wynika zatem ciągłość
$A_p:V\to V^*$ w normie.
Dla $w\in V$ stosujemy twierdzenie Newtona--Leibniza
do funkcji $t\mapsto J_f(u+tw)$, której pochodną
obliczyliśmy w dowodzie Twierdzenia~\ref{thm:pde26-minimizer}.
Otrzymujemy
\[
 \begin{aligned}
 &|J_f(u+w)-J_f(u)-\langle A_pu-f,w\rangle|\\
 &\quad\le\|w\|_V\sup_{0\le t\le1}
                         \|A_p(u+tw)-A_pu\|_{V^*}.
 \end{aligned}
\]
Supremum dąży do zera przy $\|w\|_V\to0$, ponieważ
$\|tw\|_V\le\|w\|_V$ i $A_p$ jest ciągły w $u$.
W konsekwencji $J_f$ jest różniczkowalny w sensie
Frécheta, $DJ_f(u)=A_pu-f$, a różniczka jest ciągła.
\end{uwaga}

\begin{przyklad}[Monotoniczny strumień bez potencjału]
\label{ex:pde26-skewflux}
Niech $d=2$, $p=2$ i
$J=\begin{psmallmatrix}0&-1\\1&0\end{psmallmatrix}$.
Dla niekonstantnej funkcji $\kappa\in W^{1,\infty}(\Omega)$
kładziemy $a(x,\xi)=(I+\kappa(x)J)\xi$.
Wówczas $(a(x,\xi)-a(x,\eta))\cdot(\xi-\eta)
=|\xi-\eta|^2$ oraz $a(x,\xi)\cdot\xi=|\xi|^2$.
Twierdzenie~\ref{thm:pde26-leraylions} zapewnia jednoznaczne
rozwiązanie dla każdego $f\in H^{-1}(\Omega)$.
Na zbiorze $\{\kappa\ne0\}$ macierz $I+\kappa J$
nie jest symetryczna, zatem strumień $a(x,\cdot)$
nie jest tam gradientem funkcji skalarnej względem
$\xi$: gdyby był, ta funkcja byłaby klasy $C^2$
i jej hesjan byłby symetryczny.
Dla gładkich funkcji testowych rachunek daje
\[
 -\diver(\kappa J\nabla u)
 =-\nabla\kappa\cdot J\nabla u,
\]
gdyż $\diver(J\nabla u)=-\partial_1\partial_2 u
+\partial_2\partial_1 u=0$.
Dla przykładu $\kappa(x)=x_1$ otrzymujemy dodatkowy
człon $\partial_2 u$. Rachunek obowiązuje także
dystrybucyjnie dla $u\in H_0^1(\Omega)$, przez
aproksymację funkcjami testowymi w $H^1$.
Dla stałego $\kappa$ człon ten byłby zerowy mimo
niesymetryczności strumienia. W przykładzie
$\kappa=x_1$ sam operator $A=-\Delta+\partial_2$
jest niesymetryczny: dla $u,v\in C_c^\infty(\Omega)$
różnica $\langle Au,v\rangle-\langle Av,u\rangle$
wynosi $2\int_\Omega(\partial_2u)v\,\mathrm d\lambda_2$.
Wybór $u$ z $\partial_2u\not\equiv0$ i
$v=\partial_2u$ daje ściśle dodatnią różnicę.
Nie istnieje więc różniczkowalny w sensie Frécheta
potencjał o różniczce $Au$: taka różniczka byłaby
ciągła i liniowa względem $u$, a druga różniczka
potencjału byłaby symetryczna.
W przypadku $\kappa(x)=x_1$ minimalizacja samej energii
$\frac12\int_\Omega|\nabla u|^2\,\mathrm d\lambda_2
-\langle f,u\rangle$ nie prowadzi do tego równania.
\end{przyklad}



\begin{uwaga}[Monotoniczność, energia i ograniczenia jednostronne]
\label{rem:pde-v44-minty-obstacle}
Twierdzenie~\ref{thm:pde26-leraylions} nie wymaga istnienia
potencjału: operator monotoniczny może być niesymetryczny,
jak w przykładzie~\ref{ex:pde26-skewflux}. Nie wolno więc
wnioskować, że każde takie równanie jest warunkiem stacjonarności
funkcjonału energii. W rozdziale~\ref{chap:pde29} korzystamy
z innej, dodatkowej struktury: \emph{symetryczna} forma
koercywna określa iloczyn skalarny na $H_0^1(\Omega)$,
a ograniczenie $u\ge\psi$ prowadzi do nierówności wariacyjnej.
Reprezentację jej rozwiązania jako rzutu w normie energetycznej
podamy w stwierdzeniu~\ref{prop:pde-v44-obstacle-projection}.
\end{uwaga}


\section{Porównanie, dodatniość i stabilność względem prawej strony}
\begin{twierdzenie}[Zasada porównania]\label{thm:pde26-comparison}
Niech $f,g\in V^*$, a $u,v\in V$ będą odpowiednio
jedynymi rozwiązaniami $A_pu=f$ i $A_pv=g$.
Jeśli $\langle f-g,\varphi\rangle\le0$ dla wszystkich
nieujemnych $\varphi\in V$, to $u\le v$ prawie wszędzie
w $\Omega$. W szczególności z $f\ge0$ w powyższym
porządku w przestrzeni sprzężonej wynika $u\ge0$ prawie wszędzie.
\end{twierdzenie}
\begin{proof}
Kładziemy $w=(u-v)_+\in V$ i odejmujemy równania słabe:
\[
 \int_{\{u>v\}}(a_p(\nabla u)-a_p(\nabla v))\cdot
 (\nabla u-\nabla v)\,\mathrm d\lambda_d(x)
 =\langle f-g,w\rangle\le0.
\]
Całkowana funkcja jest nieujemna, a przy nierównych
gradientach ściśle dodatnia, na mocy lematu~\ref{lem:pde26-vector}.
Stąd $\nabla w=0$ prawie wszędzie, a nierówność Poincarégo
daje $w=0$. Dla drugiego twierdzenia porównujemy z
zerowym rozwiązaniem równania $A_p0=0$.
\end{proof}

\begin{twierdzenie}[Ilościowa stabilność]\label{thm:pde26-stability}
Niech $u,v$ będą rozwiązaniami z Twierdzenia
~\ref{thm:pde26-plap} dla dowolnych $f,g\in V^*$.
Nie zakładamy uporządkowania $f$ i $g$.
Jeśli $p\ge2$, to
\begin{equation}\label{eq:pde26-holder-data}
 \|u-v\|_V\le
 \bigl((p-1)2^{p-2}\|f-g\|_{V^*}\bigr)^{1/(p-1)}.
\end{equation}
Jeżeli $1<p<2$ i $\max\{\|u\|_V,\|v\|_V\}\le R$,
to istnieje $C_p>0$ takie, że
\begin{equation}\label{eq:pde26-lipschitz-data}
 \|u-v\|_V\le C_p(2R)^{2-p}\|f-g\|_{V^*}.
\end{equation}
W każdym przypadku odwzorowanie $f\mapsto u$
jest ciągłe z $V^*$ do $V$.
\end{twierdzenie}
\begin{proof}
Odejmując równania i testując przez $u-v$ otrzymujemy
\[
 \langle A_pu-A_pv,u-v\rangle
 =\langle f-g,u-v\rangle
 \le\|f-g\|_{V^*}\|u-v\|_V.
\]
Jeżeli $u=v$, oba oszacowania zachodzą natychmiast.
Dalej zakładamy $u\ne v$, więc $\|u-v\|_V>0$.
Dla $p\ge2$ całkujemy~\eqref{eq:pde26-strong},
co daje $2^{2-p}\|u-v\|_V^p/(p-1)$ po lewej stronie,
a zatem~\eqref{eq:pde26-holder-data}.
Dla $1<p<2$ wprowadzamy
$W(x)=|\nabla u(x)|+|\nabla v(x)|$ i
$h=|\nabla u-\nabla v|$. Na zbiorze $W>0$ rozkładamy
$h^p=(W^{p-2}h^2)^{p/2}W^{p(2-p)/2}$
i stosujemy nierówność Höldera z wykładnikami
$2/p$ i $2/(2-p)$:
\[
 \int_\Omega h^p\,\mathrm d\lambda_d
 \le \left(\int_\Omega W^{p-2}h^2\,\mathrm d\lambda_d\right)^{p/2}
 \left(\int_\Omega W^p\,\mathrm d\lambda_d\right)^{(2-p)/2}.
\]
Na zbiorze $W=0$ przyjmujemy $W^{p-2}h^2=0$.
Ponieważ $0\le h\le W$, tak określona funkcja jest
całkowalna, ograniczona przez $W^p$.
Po podniesieniu do potęgi $2/p$ mamy
\[
 \|u-v\|_V^2
 \le \left(\int_\Omega W^{p-2}h^2\,\mathrm d\lambda_d\right)
       \|W\|_p^{2-p}.
\]
Nierówność Minkowskiego daje $\|W\|_p\le2R$.
Z założenia $u\ne v$ wynika $\|W\|_p>0$, więc
$R>0$ i można dzielić przez $\|W\|_p^{2-p}$.
Wykorzystując~\eqref{eq:pde26-subquad}, dostajemy
$\langle A_pu-A_pv,u-v\rangle
\ge(p-1)(2R)^{p-2}\|u-v\|_V^2$.
Dzielimy przez $\|u-v\|_V$ i otrzymujemy
\eqref{eq:pde26-lipschitz-data} z $C_p=1/(p-1)$.
Wreszcie oszacowanie~\eqref{eq:pde26-apriori} zapewnia
wspólne ograniczenie norm rozwiązań, gdy $f,g$
pozostają w ograniczonym zbiorze $V^*$.
\end{proof}

\section{Jawne rozwiązania: kule i pierścienie}
\begin{przyklad}[Rozwiązanie skrętnościowe w kuli]\label{ex:pde26-torsion}
Niech $B_R=\{x\in\R^d:|x|<R\}$, $R>0$, oraz
$\alpha=p/(p-1)$. Definiujemy
\begin{equation}\label{eq:pde26-torsion}
 u(x)=\frac{p-1}{p}\,d^{-1/(p-1)}
 \bigl(R^\alpha-|x|^\alpha\bigr),\qquad x\in\overline B_R.
\end{equation}
Ponieważ $\alpha>1$, funkcja $u$ jest klasy $C^1$
na domkniętej kuli, z gradientem równym zero w środku,
i znika na brzegu. Należy do $W_0^{1,p}(B_R)$:
mnożymy ją przez gładkie odcięcie równe $1$ w odległości
co najmniej $2\varepsilon$ od brzegu i równe $0$
w odległości mniejszej niż $\varepsilon$, o gradiencie
ograniczonym przez $C/\varepsilon$. W warstwie
brzegowej $|u|\le C\varepsilon$, więc błąd w normie
$W^{1,p}$ dąży do zera wraz z miarą warstwy.
Obcięte funkcje przybliżamy splotami o zwartym nośniku.
Dla $x\ne0$ mamy
$\nabla u=-d^{-1/(p-1)}|x|^{\alpha-2}x$, a zatem,
ponieważ $(\alpha-1)(p-1)=1$,
$|\nabla u|^{p-2}\nabla u=-x/d$.
Równość strumieni zachodzi również w zerze przez
ciągłość. Dywergencja pola $-x/d$ wynosi $-1$.
Całkowanie przez części dla funkcji testowych i
gęstość $C_c^\infty(B_R)$ w $W_0^{1,p}(B_R)$
dają $-\Delta_pu=1$ w sensie słabym.
Dla $p>2$ mamy $1<\alpha<2$:
rozwiązanie jest różniczkowalne z ciągłym gradientem,
ale nie należy do $C^2$ w środku kuli:
na dowolnej prostej przez zero, sparametryzowanej
wektorem jednostkowym, drugi symetryczny iloraz
różnicowy w zerze wynosi
$-2d^{-1/(p-1)}|t|^{\alpha-2}/\alpha$,
co dąży do $-\infty$ przy $t\to0$.
To przykład pokazujący, dlaczego nie można automatycznie
przenosić regularności $C^2$ z przypadku liniowego $p=2$.
\end{przyklad}

\begin{przyklad}[Rozwiązanie radialne w pierścieniu]\label{ex:pde26-annulus}
Niech $0<r_0<r_1<\infty$, $d\ge2$ i $u(x)=U(|x|)$
jest radialnym rozwiązaniem słabym $\Delta_pu=0$
w pierścieniu, z $U\in W^{1,p}(r_0,r_1)$.
Współrzędne biegunowe i testy radialne
$v(x)=\psi(|x|)$, $\psi\in C_c^\infty(r_0,r_1)$,
dają po skróceniu przez pole sfery jednostkowej
\[
 \int_{r_0}^{r_1}r^{d-1}|U'|^{p-2}U'\psi'
                 \,\mathrm d\lambda_1(r)=0.
\]
Zatem $r^{d-1}|U'|^{p-2}U'$ jest stałą prawie
wszędzie, gdyż jego pochodna dystrybucyjna znika
na przedziale.
Stąd $U'(r)=c\,r^{-(d-1)/(p-1)}$ z odpowiednio
zdefiniowaną stałą $c\in\R$. Po całkowaniu
\[
 U(r)=\begin{cases}
 a+b\,r^{(p-d)/(p-1)},&p\ne d,\\
 a+b\log r,&p=d.
 \end{cases}
\]
Dla zadanych wartości $U(r_0)=A$, $U(r_1)=B$
i $\gamma=(p-d)/(p-1)\ne0$ otrzymujemy
\[
 U(r)=A+(B-A)\frac{r^\gamma-r_0^\gamma}
                         {r_1^\gamma-r_0^\gamma}.
\]
Dla $p=d$ ułamek zastępujemy przez
$\log(r/r_0)/\log(r_1/r_0)$.
Jest to też granica powyższego ułamka przy
$\gamma\to0$, po podzieleniu licznika i mianownika
przez $\gamma$ i użyciu pochodnej wykładniczej.
Odwrotnie, otrzymane profile mają strumień
$C x/|x|^d$ o zerowej dywergencji w pierścieniu,
więc spełniają równanie dla wszystkich testów,
nie tylko radialnych. Jeżeli $C\ne0$, pole to
jest lokalnie całkowalne również w całej kuli,
lecz jego dywergencja dystrybucyjna wynosi
$C|\mathbb S^{d-1}|\delta_0$: całkowanie przez
części poza kulą promienia $\varepsilon$ daje
wyraz brzegowy $C\int_{\mathbb S^{d-1}}
\varphi(\varepsilon\omega)\,\mathrm d\sigma(\omega)$,
który dąży do $C|\mathbb S^{d-1}|\varphi(0)$.
Dla $C\ne0$ całka z potęgi $p$ gradientu w kuli
promienia $\rho$ jest dodatnią stałą razy
\[
 \int_0^\rho r^{-(d-1)/(p-1)}\,\mathrm d\lambda_1(r).
\]
Jest skończona dokładnie wtedy, gdy $p>d$.
Wtedy profil ma skończoną granicę w zerze, jest
ograniczony i należy do $W^{1,p}$ w kuli:
całkowanie przez części poza kulą promienia
$\varepsilon$ daje dla pierwszych słabych pochodnych
wyraz brzegowy ograniczony przez stałą razy
$\varepsilon^{d-1}$, który znika dla $d\ge2$.
Dla $p\le d$ nie może należeć do $W^{1,p}$,
gdyż jego klasyczny gradient poza zerem nie
jest w $L^p$. Nawet w przypadku $p>d$ profil
nie jest rozwiązaniem jednorodnego równania
w całej kuli, ze względu na obliczony człon Diraca.
\end{przyklad}

\begin{przyklad}[Jednowymiarowe równanie z prawą stroną]\label{ex:pde26-1d}
Niech $\Omega=(0,1)$ i $f\equiv1$. Jednoznaczne rozwiązanie
$-(|u'|^{p-2}u')'=1$ z $u(0)=u(1)=0$ wynosi
\[
 u(x)=\frac{p-1}{p}\left(2^{-p/(p-1)}
 -|x-\tfrac12|^{p/(p-1)}\right).
\]
W istocie strumień $q=|u'|^{p-2}u'$ jest dokładnie
$q(x)=\tfrac12-x$, więc $-q'=1$; warunki brzegowe
wynikają z symetrii profilu. Dla $p=2$ odzyskujemy
$u(x)=x(1-x)/2$. Przynależność $u$ do
$W_0^{1,p}(0,1)$ wynika również z konstrukcji
w Przykładzie~\ref{ex:pde26-torsion} dla $d=1$,
po przesunięciu środka i przyjęciu $R=1/2$.
\end{przyklad}

\begin{center}
\begin{tikzpicture}[x=8.5cm,y=4cm,>=Latex]
 \draw[->] (-.025,0)--(1.05,0) node[right] {$x$};
 \draw[->] (0,-.025)--(0,.66) node[above] {$u(x)$};
 \foreach \x/\lab in {0/0,.5/{1/2},1/1}
  \draw (\x,.009)--(\x,-.009) node[below] {$\lab$};
 \draw[densely dashed,gray] (.5,0)--(.5,.53);
 \draw[blue!75!black,thick,domain=0:1,samples=121]
  plot (\x,{(1/3)*(.5^3-abs(\x-.5)^3)});
 \draw[teal!75!black,thick,domain=0:1,samples=121]
  plot (\x,{.5*(.5^2-abs(\x-.5)^2)});
 \draw[red!75!black,thick,domain=0:1,samples=121]
  plot (\x,{.75*(.5^(4/3)-abs(\x-.5)^(4/3))});
 \draw[blue!75!black,thick] (.7,.58)--(.77,.58)
  node[right,black,font=\small] {$p=3/2$};
 \draw[teal!75!black,thick] (.7,.49)--(.77,.49)
  node[right,black,font=\small] {$p=2$};
 \draw[red!75!black,thick] (.7,.40)--(.77,.40)
  node[right,black,font=\small] {$p=4$};
\end{tikzpicture}

\small Rysunek: profile rozwiązania jednowymiarowego
dla $p=3/2,2,4$ przy prawej stronie równej $1$.
Dla $p>2$ gradient jest ciągły, lecz druga pochodna
w środku przedziału nie istnieje jako skończona liczba.
\end{center}

\section{Pierwsza wartość własna nieliniowego operatora}
\begin{definicja}[Wartość własna $p$-laplasjanu]
\label{def:pde26-eigen}
Liczbę $\lambda\in\R$ nazywamy wartością własną
Dirichleta $p$-laplasjanu na $\Omega$, jeśli istnieje
niezerowa funkcja $u\in W_0^{1,p}(\Omega)$ spełniająca
$-\Delta_pu=\lambda|u|^{p-2}u$ w $V^*$.
Definiujemy też liczbę
\[
 \lambda_{1,p}(\Omega)=
 \inf_{v\in V\setminus\{0\}}
 \frac{\int_\Omega|\nabla v|^p\,\mathrm d\lambda_d(x)}
 {\int_\Omega|v|^p\,\mathrm d\lambda_d(x)}.
\]
\end{definicja}
\begin{twierdzenie}[Osiągalność minimum i równanie własne]
\label{thm:pde26-eigen}
Niech $\Omega\subset\R^d$ będzie ograniczonym obszarem
Lipschitzowskim i $1<p<\infty$.
Wówczas $0<\lambda_{1,p}(\Omega)<\infty$ oraz
istnieje $u\in W_0^{1,p}(\Omega)$ z
$\|u\|_{L^p(\Omega)}=1$, $u\ge0$ prawie wszędzie,
takie, że
\[
 \int_\Omega|\nabla u|^{p-2}\nabla u\cdot\nabla v
 \,\mathrm d\lambda_d(x)=
 \lambda_{1,p}(\Omega)
 \int_\Omega |u|^{p-2}uv\,\mathrm d\lambda_d(x)
 \quad(v\in V).
\]
Dla $r>0$ zachodzi ponadto prawo skalowania
$\lambda_{1,p}(r\Omega)=r^{-p}\lambda_{1,p}(\Omega)$.
\end{twierdzenie}
\begin{proof}
Z nierówności Poincarégo mamy
$\|v\|_p\le C\|\nabla v\|_p$,
więc iloraz Rayleigha jest ograniczony od dołu przez
$C^{-p}>0$. Skończoność otrzymujemy po podstawieniu
dowolnej niezerowej funkcji testowej z $C_c^\infty(\Omega)$.
Bierzemy ciąg minimalizujący z $\|v_n\|_p=1$.
Jego gradienty są ograniczone w $L^p$. Dzięki
refleksywności, po przejściu do podciągu,
$v_n\rightharpoonup u$ w $V$. Zwarte zanurzenie
$W_0^{1,p}(\Omega)\hookrightarrow L^p(\Omega)$
(Twierdzenie~\ref{thm:pde10-Rellich-p})
daje dalszy podciąg silnie zbieżny w $L^p$.
Jego granicą musi być $u$, ponieważ ciągłe
zanurzenie $V\hookrightarrow L^p$ zachowuje
zbieżność słabą. Stąd $\|u\|_p=1$.
Słaba półciągłość normy gradientu dowodzi osiągalności.
Ponieważ $|u|\in V$, $\||u|\|_p=1$ oraz
$\|\nabla|u|\|_p=\|\nabla u\|_p$, możemy wybrać
minimizator nieujemny (wzory na moduł wynikają
z Wniosku~\ref{cor:pde9-truncation}).
Dla ustalonego $v\in V$ kładziemy
$P(t)=\|\nabla(u+tv)\|_p^p$ i $Q(t)=\|u+tv\|_p^p$.
Mamy $Q(0)=1$, zatem $Q(t)>0$ w otoczeniu zera.
Różniczkowanie pod całką uzasadniamy majorantami
z dowodu Twierdzenia~\ref{thm:pde26-minimizer},
odpowiednio dla gradientów i dla samych funkcji.
Otrzymujemy
\[
 \begin{aligned}
 P'(0)&=p\langle A_pu,v\rangle,\\
 Q'(0)&=p\int_\Omega |u|^{p-2}uv\,\mathrm d\lambda_d.
 \end{aligned}
\]
Ponieważ $P/Q$ osiąga minimum w zerze,
$0=P'(0)-\lambda_{1,p}Q'(0)$.
To daje wypisaną tożsamość słabą.
Każda wartość własna $\lambda$ spełnia
$\lambda\ge\lambda_{1,p}$, gdyż testowanie równania
własnego przez niezerową funkcję własną daje
$\lambda=\|\nabla u\|_p^p/\|u\|_p^p$.
Liczba $\lambda_{1,p}$ jest więc rzeczywiście
najmniejszą wartością własną.
Wreszcie podstawienie $v(x)=w(x/r)$
w definicji wartości własnej na $r\Omega$ i
zmiana zmiennych $x=ry$ mnoży licznik przez
$r^{d-p}$, a mianownik przez $r^d$. Skalowanie
jest bijekcją między $W_0^{1,p}(\Omega)$ i
$W_0^{1,p}(r\Omega)$: najpierw sprawdzamy to
dla funkcji testowych, potem stosujemy gęstość
i ciągłość skalowania oraz jego odwrotności. Stąd
prawo skalowania. Nie twierdzimy tutaj, że
nieujemny minimizator jest punktowo ściśle dodatni;
taki wniosek wymaga dodatkowej teorii regularności.
\end{proof}

\section{Monotoniczny schemat różnicowy w jednym wymiarze}
Dzięki jednowymiarowej postaci strumienia możemy
zbudować schemat numeryczny, którego jednoznaczność
wynika z tej samej monotoniczności co twierdzenie
\ref{thm:pde26-plap}. Co istotne, dla stałej prawej
strony nie musimy korzystać z iteracyjnego solvera
nieliniowego.

\begin{twierdzenie}[Dokładna konstrukcja dyskretnego strumienia]
\label{thm:pde26-discrete}
Niech $p>1$ i $N\ge2$ będzie parzystą liczbą całkowitą.
Ustalmy $h=1/N$ oraz $x_i=ih$ dla $0\le i\le N$.
Szukamy liczb $U_0,\ldots,U_N$ spełniających
$U_0=U_N=0$ oraz dla $1\le i\le N-1$ układ
\begin{equation}\label{eq:pde26-discrete}
 -\frac{Q_{i+1/2}-Q_{i-1/2}}h=1,\qquad
 Q_{i+1/2}=a_p\left(\frac{U_{i+1}-U_i}h\right).
\end{equation}
Strumienie są zdefiniowane dla $0\le i\le N-1$,
a w równaniu bilansu $1\le i\le N-1$.
Używamy jednowymiarowej funkcji $a_p$ z Lematu
~\ref{lem:pde26-vector}, w szczególności $a_p(0)=0$.
Układ ma dokładnie jedno rozwiązanie, określone przez
\begin{equation}\label{eq:pde26-flux}
 Q_{i+1/2}=\frac12-\left(i+\frac12\right)h,\qquad
 U_i=h\sum_{j=0}^{i-1}\operatorname{sgn}(Q_{j+1/2})
 |Q_{j+1/2}|^{1/(p-1)}.
\end{equation}
Jeśli $u$ jest rozwiązaniem jawnym z
przykładu~\ref{ex:pde26-1d}, to przy
$\beta=1/(p-1)$ zachodzi
\[
 \max_{0\le i\le N}|U_i-u(x_i)|
 \le C_p h^{\min\{2,1+\beta\}},
\]
gdzie $C_p$ nie zależy od $N$ (dla $p=2$ błąd
w węzłach jest dokładnie zerowy).
\end{twierdzenie}
\begin{proof}
\textbf{Krok 1: bilans strumienia.} Pierwsze równanie
\eqref{eq:pde26-discrete} daje
$Q_{i+1/2}=C-(i+1/2)h$. Odwrotnością funkcji
$a_p(s)=|s|^{p-2}s$ jest
$b_p(q)=\operatorname{sgn}(q)|q|^{1/(p-1)}$.
Zatem $U_{i+1}-U_i=h\,b_p(Q_{i+1/2})$.
Warunek $U_N=U_0$ wymaga
$\sum_{i=0}^{N-1}b_p(C-(i+1/2)h)=0$.
Lewa strona jest ściśle rosnącą i ciągłą funkcją
$C$, a dla $C=1/2$ wyrazy sumy znoszą się parami
przez symetrię względem środka odcinka. Stąd jedyne
$C$ jest równe $1/2$. Formuła~\eqref{eq:pde26-flux}
spełnia więc oba warunki brzegowe i wszystkie równania.

\textbf{Krok 2: niezależny dowód jednoznaczności.}
Dla dwóch rozwiązań $U,V$ pomnóżmy różnicę
równań przez $h(U_i-V_i)$ i zsumujmy po węzłach
wewnętrznych. Oznaczmy $w_i=U_i-V_i$ i $R_i=a_p(D_iU)-a_p(D_iV)$,
gdzie $D_iU=(U_{i+1}-U_i)/h$.
Dyskretne całkowanie przez części ma postać
\[
 -\sum_{i=1}^{N-1}(R_i-R_{i-1})w_i
 =\sum_{i=0}^{N-1}R_i(w_{i+1}-w_i),
\]
co wynika z przesunięcia indeksu i $w_0=w_N=0$.
Stąd
\[
 \sum_{i=0}^{N-1}h\bigl(a_p(D_iU)-a_p(D_iV)\bigr)
 (D_iU-D_iV)=0,
\]
gdzie $D_iU=(U_{i+1}-U_i)/h$. Ścisła monotoniczność
$a_p$ implikuje $D_iU=D_iV$ dla wszystkich $i$,
a zatem $U=V$.

\textbf{Krok 3: zbieżność.} Wzór~\eqref{eq:pde26-flux}
wyraża $U_i$ jako sumę środków przedziałów dla
całki $\int_0^{x_i}b_p(1/2-x)\,\mathrm d\lambda_1(x)$,
która jest równa $u(x_i)$.
Na przedziałach odległych od $1/2$ druga pochodna
funkcji $g(x)=b_p(1/2-x)$ spełnia
$|g''(x)|\le C_p|x-1/2|^{\beta-2}$.
Na przedziale $I$ długości $h$ ze środkiem $m$
rozwinięcie Taylora daje
$|g(x)-g(m)-g'(m)(x-m)|\le
\frac12\sup_I|g''|\,|x-m|^2$.
Po scałkowaniu składnik liniowy znika, a błąd
jest nie większy niż $h^3\sup_I|g''|/24$.
Ponieważ $N$ jest parzyste, środek $1/2$ jest węzłem.
Dla przedziałów o odległości od środka należącej
do $[jh,(j+1)h]$, $j\ge1$, supremum pochodnej
jest nie większe niż $C_p(jh)^{\beta-2}$,
także dla $\beta>2$, bo $(j+1)/j\le2$.
Na każdym poziomie $j$ są co najwyżej dwa przedziały.
Suma błędów jest zatem ograniczona przez
$C_ph^{1+\beta}\sum_{j=1}^{N/2}j^{\beta-2}$.
Dla każdego z dwóch przedziałów przylegających
do środka błąd jest nie większy niż
$h(h/2)^\beta+\int_0^h s^\beta\,\mathrm d\lambda_1(s)$,
czyli co najwyżej $C_ph^{1+\beta}$.
Oszacowanie sumy modułów błędów obowiązuje również
dla dowolnej sumy częściowej, a więc dla każdego węzła.
Dla $\beta<1$ wspomniany szereg jest jednostajnie
ograniczony; dla $\beta>1$ jego suma jest
$O(N^{\beta-1})$, co daje $O(h^2)$.
Gdy $\beta=1$, funkcja $g$ jest afiniczna i metoda
punktu środkowego jest dokładna. Dowód jest zakończony.
\end{proof}

\section{Granice twierdzeń i przypadki niemonotoniczne}
\begin{kontrprzyklad}[Zanik monotoniczności po odjęciu potencjału]
\label{ex:pde26-nonmonotone}
Dla $\lambda>0$ rozważmy
$Bu=A_pu-\lambda|u|^{p-2}u$ w $V^*$,
przy czym drugi składnik rozumiemy przez
$\langle |u|^{p-2}u,v\rangle=
\int_\Omega|u|^{p-2}uv\,\mathrm d\lambda_d(x)$.
Jeśli $\lambda>\lambda_{1,p}(\Omega)$, bierzemy
minimizator $\varphi$ z twierdzenia~\ref{thm:pde26-eigen}.
Wtedy
$\langle B\varphi-B0,\varphi-0\rangle
=\lambda_{1,p}-\lambda<0$, więc $B$ nie jest monotoniczny.
Dla $\lambda=\lambda_{1,p}$ równanie $Bu=0$ ma
co najmniej dwa różne rozwiązania: $u=0$ i $u=\varphi$.
Nie wolno więc automatycznie stosować twierdzenia
Minty'ego--Browdera po dodaniu wyrazu o niekorzystnym znaku.
\end{kontrprzyklad}

\begin{uwaga}[Co zostaje na później?]
Nie udowodniliśmy regularności $C^{1,\alpha}$
ani nierówności Harnacka dla $p$-harmonicznych funkcji;
nie wynikają one z samej monotoniczności. Również
nierówności wariacyjne z przeszkodą wymagają
modyfikacji relacji $Au=f$ na nierówność w zbiorze wypukłym;
wrócimy do nich w rozdziale 29.
Zaproponowana tu metoda Galerkina obejmuje również
strumienie typu Leraya--Lionsa przy właściwych
hipotezach wzrostu, monotoniczności i koercywności,
ale pełnego twierdzenia z zależnością od $x$ i $u$
nie ogłaszamy bez osobnej weryfikacji warunków.
\end{uwaga}


% ROZDZIAŁ 27: reakcja i dyfuzja

\chapter[Reakcja--dyfuzja i fale biegnące]{Równania reakcji--dyfuzji, fale biegnące, dynamika długoczasowa i niestabilność Turinga}
\label{chap:pde27}

\section{Mechanizmy reakcji i dyfuzji: zakres rozdziału}
Rozważamy skalarne równanie
\begin{equation}\label{eq:pde27-main}
 u_t=D\Delta u+f(u),\qquad D>0,
\end{equation}
gdzie dyfuzja wygładza rozkład przestrzenny, a reakcja opisuje
lokalną dynamikę. Badamy zarówno torus
$\T_L^d=(\R/L\mathbb Z)^d$, jak i prostą $\R$.
Na torusie stosujemy periodyczne warunki brzegowe;
na prostej badamy fronty $u(t,x)=U(x-ct)$.
Jeżeli $f$ jest tylko lokalnie lipschitzowska, rozwiązanie
nie musi istnieć dla wszystkich czasów (por. rozdział~\ref{chap:pde25}).
W szczególności nie wolno wyprowadzać zachowania długoczasowego
z samego twierdzenia o lokalnym istnieniu.

W obliczeniach dla modelu Fishera--KPP użyjemy
$f(s)=rs(1-s)$, $r>0$. Rozwiązanie z zakresu $0\le u\le1$
pozostaje w tym przedziale, mimo że $f$ nie jest monotoniczna
na całej osi. Jest to zastosowanie zasady porównania,
nie monotoniczności funkcji $f$.

\section{Istnienie, ciągłość względem danych i zasada porównania}
\begin{twierdzenie}[Istnienie lokalne i alternatywa eksplozji]\label{thm:pde27-local}
Niech $L,D>0$, $d\ge1$, $f\in C^1(\R)$ oraz
$u_0\in C(\T_L^d)$. Istnieje dokładnie jedno maksymalne
rozwiązanie łagodne $u\in C([0,T_{\max});C(\T_L^d))$
równania~\eqref{eq:pde27-main} z $u(0)=u_0$.
Dla każdego $T<T_{\max}$ zachodzi wzór Duhamela
\[
 u(t)=S_D(t)u_0+\int_0^tS_D(t-s)f(u(s))\,\mathrm d\lambda_1(s),
 \qquad S_D(t)=e^{tD\Delta}.
\]
Jeżeli $T_{\max}<\infty$, to
$\|u(t)\|_{\infty}\to\infty$ przy $t\uparrow T_{\max}$.
Jeśli $u,v$ są rozwiązaniami o wartościach w $[-M,M]$
na $[0,T]$, to dla $K=\max_{|s|\le M}|f'(s)|$ mamy
\[
 \|u(t)-v(t)\|_{\infty}\le e^{Kt}\|u_0-v_0\|_{\infty},
 \qquad 0\le t\le T.
\]
\end{twierdzenie}
\begin{proof}
Na przestrzeni Banacha $X_T=C([0,T];C(\T_L^d))$
z normą $\|w\|_{X_T}=\sup_{0\le t\le T}\|w(t)\|_\infty$
rozważamy operator $\mathcal F$ prawej strony wzoru Duhamela,
w którym całka jest całką Bochnera w $C(\T_L^d)$.
Półgrupa cieplna jest silnie ciągła, dodatnia
i kontraktywna w normie maksimum; wynika to
z Lematu~\ref{lem:pde25-heat} przez skalowanie
w Stwierdzeniu~\ref{prop:pde-v44-torus-scaling}.
Dla $w\in X_T$ funkcja $s\mapsto f(w(s))$ jest ciągła
w normie maksimum, ponieważ zakres $w$ jest ograniczony.
Silna ciągłość półgrupy i jej kontraktywność dają
ciągłość całki względem $t$ (przez zbieżność
zmajoryzowaną na wspólnej części przedziałów całkowania
i oszacowanie całki na ich różnicy).
Zatem $\mathcal F$ przyjmuje wartości w $X_T$.

Ustalmy $M>\|u_0\|_\infty+1$ i oznaczmy
$F_M=\max_{|s|\le M}|f(s)|$, $K_M=\max_{|s|\le M}|f'(s)|$.
Dla $w,z$ z domkniętej kuli o promieniu $M$ mamy
\[
 \|\mathcal Fw\|_{X_T}\le\|u_0\|_\infty+TF_M,\qquad
 \|\mathcal Fw-\mathcal Fz\|_{X_T}
 \le TK_M\|w-z\|_{X_T}.
\]
Wybieramy $T>0$ tak, by $TF_M\le M-\|u_0\|_\infty$
i $TK_M<1$. Kula jest zupełna, więc twierdzenie
Banacha o punkcie stałym daje rozwiązanie.
Jednoznaczność w całej klasie rozwiązań łagodnych
wynika z oszacowania Gronwalla poniżej, zastosowanego
na dowolnym wspólnym zwartym przedziale istnienia.
Zgodne na przecięciach przedłużenia wyznaczają
jedno rozwiązanie na sumie wszystkich przedziałów
istnienia, czyli na $[0,T_{\max})$.
Wzór Duhamela można rozpocząć w dowolnej chwili
$t_0<T_{\max}$:
\[
 u(t_0+t)=S_D(t)u(t_0)+
 \int_0^t S_D(t-s)f(u(t_0+s))\,\mathrm d\lambda_1(s).
\]
Otrzymujemy go przez rozdzielenie pierwotnej całki
w chwili $t_0$ i użycie własności półgrupy.
Czas lokalnego istnienia można wybrać jednostajnie
dla danych o ograniczonej normie: dla
$\|u(t_0)\|_\infty\le B$ przyjmujemy $M=B+2$
i stosujemy powyższe warunki na $T$.
Nie zależy on od momentu rozpoczęcia ewolucji. Gdyby przy skończonym $T_{\max}$
norma rozwiązania była ograniczona wzdłuż pewnego ciągu
tendującego do $T_{\max}$, można by uruchamiać kontrakcję
od tych chwil na wspólnym przedziale o dodatniej długości,
sprzecznie z maksymalnością. To dowodzi alternatywy.
Dla różnicy $w=u-v$ wzór Duhamela i twierdzenie o wartości
średniej dają
$\|w(t)\|_\infty\le\|w(0)\|_\infty+
K\int_0^t\|w(s)\|_\infty\,\mathrm d\lambda_1(s)$.
Lemat Gronwalla kończy dowód stabilności.
\end{proof}


\begin{stwierdzenie}[Skalowanie półgrupy cieplnej na torusie]
\label{prop:pde-v44-torus-scaling}
Niech $L,D>0$ i $d\ge1$. Dla $v\in C(\T_L^d)$ połóżmy
$(Rv)(\theta)=v(L\theta/(2\pi))$, $\theta\in\T^d$.
Wówczas $R:C(\T_L^d)\to C(\T^d)$ jest liniową
izometrią na całą przestrzeń i
\begin{equation}\label{eq:pde-v44-torus-scale}
 R\,e^{tD\Delta_x}\,R^{-1}
 =e^{\,tD(2\pi/L)^2\Delta_\theta},\qquad t\ge0.
\end{equation}
Po podstawieniu $\tau=D(2\pi/L)^2t$ oraz
$w(\tau,\theta)=u(t,L\theta/(2\pi))$ równanie
$u_t=D\Delta_xu+f(u)$ przyjmuje postać
$w_\tau=\Delta_\theta w+\bigl[D(2\pi/L)^2\bigr]^{-1}f(w)$.
W konsekwencji konstrukcja lokalna i alternatywa eksplozji
z rozdziału~\ref{chap:pde25} stosują się do twierdzenia
\ref{thm:pde27-local} po tej zmianie zmiennych.
\end{stwierdzenie}
\begin{proof}
Skalowanie argumentu jest dobrze określone na klasach
modulo odpowiedni okres, a odwrotność ma postać
$(R^{-1}w)(x)=w(2\pi x/L)$. Obie normy maksimum
są równe, ponieważ zmiana zmiennych jest bijekcją torusów.
Z reguły łańcuchowej dla funkcji gładkiej
$R(\Delta_xv)=(2\pi/L)^2\Delta_\theta(Rv)$.
Dla wielomianów trygonometrycznych tożsamość półgrup
wynika także z działania na każdej częstotliwości
$k\in\mathbb Z^d$: obie strony mnożą współczynnik
Fouriera przez $\exp[-tD(2\pi/L)^2|k|^2]$.
Gęstość wielomianów trygonometrycznych w $C(\T^d)$
i kontraktywność półgrup rozszerzają równość na każdą
funkcję ciągłą. Zmiana czasu daje ostatni wzór.
Zmiana czasu we wzorze Duhamela daje tę samą
transformację także dla rozwiązań łagodnych.
Ponieważ zmodyfikowana nieliniowość nadal należy do
$C^1(\R)$, stosujemy Twierdzenie~\ref{thm:pde25-local}.
Ponadto Twierdzenie~\ref{thm:pde25-parabolic-regularity}
daje dla $0<a<T<T_{\max}$
\[
 u\in C([a,T];C^2(\T_L^d))\cap C^1([a,T];C(\T_L^d)),
\]
a równanie~\eqref{eq:pde27-main} jest spełnione
klasycznie dla dodatnich czasów.
\end{proof}


\begin{twierdzenie}[Zasada porównania i przedziały niezmiennicze]
\label{thm:pde27-comparison}
Niech $f\in C^1(\R)$, a $u,v$ będą rozwiązaniami z
poprzedniego twierdzenia, istniejącymi na $[0,T]$.
Jeżeli $u_0\le v_0$, to $u(t,\cdot)\le v(t,\cdot)$
dla każdego $0\le t\le T$. Jeśli $a\le b$,
$f(a)\ge0$ i $f(b)\le0$, to dane $a\le u_0\le b$
implikują $a\le u(t,\cdot)\le b$ na całym
przedziale istnienia; wówczas $T_{\max}=\infty$.
\end{twierdzenie}
\begin{proof}
Dzięki regularności dla dodatnich czasów, wykazanej
po Stwierdzeniu~\ref{prop:pde-v44-torus-scaling},
nie potrzebujemy gładkich danych początkowych.
Różnica $w=u-v$ spełnia
$w_t-D\Delta w=q(t,x)w$, gdzie
\[
 q(t,x)=\int_0^1 f'\bigl(v(t,x)+\theta(u(t,x)-v(t,x))\bigr)
 \,\mathrm d\lambda_1(\theta).
\]
Współczynnik $q$ jest ciągły i ograniczony na
$[0,T]\times\T_L^d$, gdyż $u,v$ są ciągłe na tym zbiorze.
Wybierzmy $K>\sup q$ i połóżmy $z=e^{-Kt}w$.
Wtedy $z_t-D\Delta z=(q-K)z$ dla $t>0$.
Gdyby $z$ przyjmowała wartość dodatnią, na zwartym
walcu $[0,T]\times\T_L^d$ osiągałaby dodatnie maksimum
w $(t_*,x_*)$ z $t_*>0$, ponieważ $z(0,\cdot)\le0$.
W tym punkcie $z_t\ge0$ (pochodna lewostronna,
jeśli $t_*=T$) i $\Delta z\le0$, natomiast
$(q-K)z<0$. Sprzeczność daje $z\le0$, czyli $u\le v$.

Ten sam argument działa, gdy $u$ jest podrozwiązaniem,
a $v$ nadrozwiązaniem, ciągłymi do chwili $0$
i klasycznymi dla $t>0$: równanie dla $z$
zastępujemy nierównością $z_t-D\Delta z\le(q-K)z$.
Stałe $a,b$ są odpowiednio pod- i nadrozwiązaniami,
gdyż $0\le f(a)$ i $0\ge f(b)$. Porównujemy je z $u$
na każdym zwartym przedziale jego istnienia.
Przedział $[a,b]$ jest niezmienniczy, a ograniczenie normy
wyklucza skończony czas maksymalny.
\end{proof}

\begin{kontrprzyklad}[Porównanie nie wymaga $f'\ge0$]
Dla $f(s)=s(1-s)$ pochodna $f'(s)=1-2s$
jest ujemna na $(1/2,\infty)$.
Mimo to zasada porównania działa, ponieważ różnica
$f(u)-f(v)$ ma postać $a(t,x)(u-v)$ z ograniczonym
współczynnikiem na każdym ograniczonym przedziale czasu.
Dla układów reakcji--dyfuzji takie porównanie
składowa po składowej na ogół nie zachodzi bez
warunku kooperatywności.
\end{kontrprzyklad}

\section{Równanie Fishera--KPP: zbieżność do stanu nasycenia na torusie}
\begin{twierdzenie}[Zbieżność do stanu nasycenia]\label{thm:pde27-kpp-torus}
Niech $u_0\in C(\T_L^d)$, $0\le u_0\le1$ i
$u_0\not\equiv0$. Dla $D,r>0$ jedyne rozwiązanie
zagadnienia
$u_t=D\Delta u+ru(1-u)$ z danymi $u_0$
istnieje globalnie i
\[
 \lim_{t\to\infty}\|u(t,\cdot)-1\|_{C(\T_L^d)}=0.
\]
Jeśli $u_0\equiv0$, to $u(t,\cdot)\equiv0$.
\end{twierdzenie}
\begin{proof}
Porównanie ze stałymi rozwiązaniami $0$ i $1$
daje $0\le u\le1$ i istnienie globalne.
Dla $t>0$ półgrupa $S_D(t)$ ma dodatnie jądro
periodyczne względem miary Lebesgue'a:
\[
 K_{D,t}^L(z)=\sum_{k\in\mathbb Z^d}
 (4\pi Dt)^{-d/2}\exp\left(-\frac{|z+Lk|^2}{4Dt}\right).
\]
Wzór wynika z Lematu~\ref{lem:pde25-heat}
przez zmianę skali, łącznie z jakobianem miary.
Szereg jest jednostajnie zbieżny dla ustalonego $t>0$
na kostce fundamentalnej i wszystkie składniki są
dodatnie. Ponieważ $u_0$ jest nieujemna, ciągła
i niezerowa, jest dodatnia na zbiorze dodatniej miary.
Stąd $S_D(t)u_0>0$ punktowo.
Wzór Duhamela z nieujemnym źródłem $ru(1-u)$ daje
$u(t,\cdot)\ge S_D(t)u_0$, więc dla każdego $t_0>0$
funkcja $u(t_0,\cdot)$ jest ściśle dodatnia. Zwarta przestrzeń
$\T_L^d$ zapewnia minimum $m_0>0$.
Niech $y'=ry(1-y)$ i $z'=rz(1-z)$,
$y(t_0)=m_0$, $z(t_0)=\max u(t_0,\cdot)\le1$.
Porównanie daje $y(t)\le u(t,x)\le z(t)$.
Z jawnego rozwiązania równania logistycznego
\[
 y(t)=\frac{1}{1+(m_0^{-1}-1)e^{-r(t-t_0)}}
\]
wynika $y(t)\uparrow1$ (dla $m_0=1$ funkcja
jest stale równa $1$), podczas gdy $z(t)\le1$.
Wzór na $y$ sprawdzamy przez różniczkowanie
i podstawienie warunku początkowego; analogiczny
wzór obowiązuje dla $z$. Porównanie jest dopuszczalne
na każdym skończonym przedziale po chwili $t_0$,
gdyż oba rozwiązania skalarne są globalne.
Mamy $\|1-u(t)\|_\infty\le1-y(t)\to0$.
Dla $u_0\equiv0$ teza wynika z jednoznaczności.
\end{proof}

\begin{uwaga}[Dlaczego torus, a nie cała prosta?]
Na $\R$ rozwiązanie wychodzące z danych o zwartym nośniku
nie zbiega jednostajnie do $1$ na całej prostej
przy $t\to\infty$: dla każdego skończonego $t>0$
rozwiązanie dąży do zera przy $|x|\to\infty$, co
wynika z oszacowania gaussowskiego w dowodzie
Twierdzenia~\ref{thm:pde27-spreading-upper}
(dla $0\le u_0\le1$). W szczególności
$\|1-u(t)\|_\infty=1$. Właściwą asymptotyką jest
\emph{propagacja frontu}, z którą zapoznamy się dalej.
\end{uwaga}


\begin{wniosek}[Ilościowa zbieżność rozwiązania Fishera--KPP na torusie]
\label{cor:pde27-kpp-rate}
Przy hipotezach twierdzenia~\ref{thm:pde27-kpp-torus}
niech $t_0>0$ i
$m_{t_0}=\min_{x\in\T_L^d}(S_D(t_0)u_0)(x)>0$.
Wówczas dla wszystkich $t\ge t_0$
\begin{equation}\label{eq:pde27-kpp-rate}
 \|1-u(t,\cdot)\|_\infty\le
 \frac{1-m_{t_0}}{m_{t_0}}e^{-r(t-t_0)}.
\end{equation}
Ponadto jeśli $K_{D,t_0}^L$ oznacza dodatnie
periodyczne jądro cieplne na $\T_L^d$, to
\[
 m_{t_0}\ge\left(\min_{z\in\T_L^d}K_{D,t_0}^L(z)\right)
 \int_{\T_L^d}u_0(x)\,\mathrm d\lambda_d(x)>0.
\]
\end{wniosek}
\begin{proof}
Nieujemne źródło $ru(1-u)$ i wzór Duhamela dają
$u(t_0)\ge S_D(t_0)u_0$, stąd
$\min u(t_0)\ge m_{t_0}>0$.
Rozwiązanie skalarnego równania
$y'=ry(1-y)$ z $y(t_0)=m_{t_0}$ wynosi
$y(t)=(1+(m_{t_0}^{-1}-1)e^{-r(t-t_0)})^{-1}$.
Porównanie ze stałym w przestrzeni rozwiązaniem $y$
i nadrozwiązaniem $1$ daje
$0\le1-u(t,x)\le1-y(t)$.
Ponieważ $1-(1+a)^{-1}=a/(1+a)\le a$
dla $a\ge0$, otrzymujemy~\eqref{eq:pde27-kpp-rate}.
Ostatnia nierówność wynika bezpośrednio z dodatniości
i ciągłości jądra na zwartym torusie.
\end{proof}


\section{Fale biegnące: od PDE do układu równań zwyczajnych}
Niech $d=1$ i rozważmy równanie
$u_t=Du_{xx}+ru(1-u)$ na $\R$.
Dla postaci fali biegnącej $u(t,x)=U(z)$, $z=x-ct$,
mamy $u_t=-cU'$, $u_x=U'$ oraz $u_{xx}=U''$,
a zatem otrzymujemy równanie profilu
\begin{equation}\label{eq:pde27-wave}
 DU''+cU'+rU(1-U)=0.
\end{equation}
Szukamy funkcji $U\in C^2(\R)$, dla której
$U(-\infty)=1$, $U(+\infty)=0$ i $U'<0$.
Prędkość $c$ jest prędkością przemieszczania się
profilu w prawo. Sama zamiana zmiennych
nie dowodzi jeszcze istnienia takiego profilu.

\begin{lemat}[Obszar niezmienniczy dla układu profilowego]
\label{lem:pde27-wedge}
Niech $D,r>0$, $c\ge2\sqrt{Dr}$ oraz
\[
 k=\frac{c-\sqrt{c^2-4Dr}}{2D}>0.
\]
Wtedy obszar
$\mathcal W=\{(U,q):0<U<1,\ 0<q<kU\}$
jest dodatnio niezmienniczy dla układu
\begin{equation}\label{eq:pde27-phase}
 U'=-q,\qquad q'=-\frac cD q+\frac rD U(1-U).
\end{equation}
Oznacza to, że trajektorie rozpoczynające się w $\mathcal W$
pozostają w nim dla dodatnich czasów istnienia.
Ponadto rozwiązanie wychodzące z punktu
$U(0)=1-\delta$, $q(0)=0$, gdzie $0<\delta<1/2$,
zachowuje $0<U<1$, $0<q<kU$ dla $z>0$ i
spełnia $U(z)\to0$, $q(z)\to0$ przy $z\to\infty$.
\end{lemat}
\begin{proof}
Na brzegu $q=0$, gdy $0<U<1$, mamy
$q'=rU(1-U)/D>0$. Na brzegu $q=kU$ otrzymujemy
\[
 (q-kU)'=\bigl(k^2-\tfrac cDk+\tfrac rD(1-U)\bigr)U
 =-\tfrac rD U^2<0,
\]
ponieważ $k^2-(c/D)k+r/D=0$.
Uzasadnijmy brak pierwszego wyjścia z obszaru.
Dopóki trajektoria pozostaje w $\mathcal W$, mamy
$U'<0$ oraz $U'+kU>0$. Stąd
$U(0)e^{-kz}\le U(z)\le U(0)<1$.
Jeżeli pierwszy kontakt z brzegiem nastąpiłby
w skończonej chwili $z_*>0$, to powyższe oszacowanie
wykluczałoby brzegi $U=0$ i $U=1$.
Przy pierwszym kontakcie z $q=0$ od strony $q>0$
musiałoby zachodzić $q'(z_*)\le0$, wbrew
$q'(z_*)=rU(z_*)(1-U(z_*))/D>0$.
Przy pierwszym kontakcie z $q=kU$ od strony $q-kU<0$
musiałoby zachodzić $(q-kU)'(z_*)\ge0$, wbrew
obliczonej ujemnej wartości.
Trajektoria pozostaje zatem w $\mathcal W$.
Pole wektorowe jest wielomianowe, więc jest lokalnie
lipschitzowskie; ograniczoność trajektorii w zwartym
domknięciu $\mathcal W$ wyklucza skończony maksymalny
czas istnienia (lokalny czas istnienia można wybrać
jednostajnie dla danych z tego zwartego zbioru).
Dla danych $U(0)=1-\delta$, $q(0)=0$ mamy
$q'(0)>0$ i $q(0)-kU(0)<0$. Ciągłość daje
$0<U<1$ i $0<q<kU$ dla dostatecznie małych $z>0$,
po czym stosujemy poprzedni argument.
Funkcja $U$ jest malejąca, więc ma granicę
$\ell\in[0,1-\delta)$. Ponieważ $q=-U'\ge0$,
$\int_0^\infty q(z)\,\mathrm d\lambda_1(z)
=U(0)-\ell<\infty$, przez twierdzenie Newtona--Leibniza
na przedziałach skończonych i przejście do granicy.
Pochodna $q'$ jest ograniczona w $\mathcal W$,
więc $q$ jest jednostajnie ciągła.
Gdyby $q$ nie dążyła do zera, istniałby ciąg
$z_j\to\infty$ z $q(z_j)\ge\varepsilon>0$.
Jednostajna ciągłość dawałaby wspólny promień $\rho>0$
taki, że $q\ge\varepsilon/2$ na przedziałach
$[z_j-\rho,z_j+\rho]$. Wybierając rozłączny
podciąg tych przedziałów, otrzymalibyśmy nieskończoną
całkę z $q$. Zatem $q(z)\to0$. Gdyby $\ell>0$, drugie
równanie~\eqref{eq:pde27-phase} prowadziłoby do
$q'(z)\to r\ell(1-\ell)/D>0$, co jest sprzeczne
z $q(z)\to0$. W konsekwencji $\ell=0$.
\end{proof}

\begin{twierdzenie}[Prędkości fal biegnących równania Fishera--KPP]
\label{thm:pde27-wave}
Niech $D,r>0$. Dla każdej $c\ge2\sqrt{Dr}$
istnieje $U\in C^2(\R)$ spełniająca
\eqref{eq:pde27-wave}, $0<U<1$, $U'<0$ oraz
$U(-\infty)=1$, $U(+\infty)=0$.
Natomiast dla żadnej $c<2\sqrt{Dr}$
taki profil nie istnieje. W szczególności
minimalna prędkość monotonicznej fali wynosi
$c_*=2\sqrt{Dr}$.
\end{twierdzenie}
\begin{proof}
\textbf{Krok 1: skonstruowanie pełnej trajektorii.}
Ustalmy $c\ge2\sqrt{Dr}$ i $k$ z
lematu~\ref{lem:pde27-wedge}. Dla $\delta\downarrow0$
rozważmy rozwiązania $(U_\delta,q_\delta)$
rozpoczynające się w $(1-\delta,0)$.
Każda trajektoria osiąga poziom $U=1/2$ dokładnie raz
na mocy ciągłości, ścisłego malenia i granicy zero;
niech $t_\delta>0$ oznacza odpowiedni czas.
Pokażemy, że $t_\delta\to\infty$.
Kładziemy $v=1-U$ i korzystamy z równań
$v'=q$, $q'=-(c/D)q+(r/D)(1-v)v\le(r/D)v$.
Dla $v,q\ge0$ dostajemy
$(v+q)'\le C(v+q)$ z $C=\max\{1,r/D\}$,
a zatem $v(t)+q(t)\le\delta e^{Ct}$.
Gdy $U(t_\delta)=1/2$, mamy $v(t_\delta)=1/2$,
więc $t_\delta\ge C^{-1}\log(1/(2\delta))\to\infty$.

Przesuwamy trajektorie: $\widetilde U_\delta(z)
=U_\delta(t_\delta+z)$, analogicznie dla $q$.
Są one określone na $(-t_\delta,\infty)$,
mają wartości w zwartym domknięciu $\mathcal W$
i jednostajnie ograniczone pochodne.
Oznaczmy przez $F$ pole wektorowe układu
~\eqref{eq:pde27-phase}. Jest ono klasy $C^\infty$,
więc na zwartym domknięciu $\mathcal W$ zarówno $F$,
jak i $DF$ są ograniczone. Dla przesuniętych trajektorii
$Y_\delta=(\widetilde U_\delta,\widetilde q_\delta)$
mamy $Y_\delta'=F(Y_\delta)$ i
$Y_\delta''=DF(Y_\delta)F(Y_\delta)$.
Zatem na każdym $[-m,m]$ dla dostatecznie małych $\delta$
funkcje i ich pierwsze pochodne są jednostajnie
ograniczone i jednakowo ciągłe.
Wybierając najpierw ciąg $\delta_j\downarrow0$,
stosujemy twierdzenie Arzeli--Ascolego do par
$(Y_{\delta_j},Y_{\delta_j}')$ na $[-m,m]$
i wydzielamy podciąg diagonalny po $m\in\mathbb N$.
Jednostajna zbieżność funkcji i pochodnych
oraz twierdzenie Newtona--Leibniza dają zbieżność
w $C^1_{\mathrm{loc}}(\R)$. Przechodząc do granicy
w $Y_\delta'=F(Y_\delta)$, otrzymujemy pełną
trajektorię $(U,q)$ układu~\eqref{eq:pde27-phase}
o wartościach w $\overline{\mathcal W}$, z $U(0)=1/2$. Ponieważ $q\ge0$ i $U'=-q$,
$U$ jest nierosnąca i ma granice $\ell_\pm$
przy $z\to\pm\infty$. Analogicznie jak w lemacie,
całkowalność $q$ na obu półprostych i jego
jednostajna ciągłość implikują $q(z)\to0$
na obu krańcach. Drugie równanie układu daje
$q'(z)\to (r/D)\ell_\pm(1-\ell_\pm)$
przy $z\to\pm\infty$. Gdyby któraś z tych granic
była dodatnia, to $q'$ byłaby ograniczona od dołu
przez dodatnią stałą na odpowiedniej półprostej.
Całkowanie nierówności na dowolnie długim odcinku
przeczyłoby ograniczoności $q$. Stąd
$\ell_\pm(1-\ell_\pm)=0$.
Z $U(0)=1/2$ i monotoniczności wynika
$\ell_-=1$, $\ell_+=0$.
Gdyby $U(z_0)=0$ lub $U(z_0)=1$, to $U$
osiągałaby odpowiednio minimum lub maksimum
w punkcie wewnętrznym prostej, a zatem $q(z_0)=0$.
Oba punkty $(0,0)$ i $(1,0)$ są równowagami
układu. Jednoznaczność rozwiązania zagadnienia
początkowego w obu kierunkach czasu wymuszałaby
trajektorię stałą, wbrew $U(0)=1/2$.
Zatem $0<U<1$ na całej prostej.
Gdyby $q(z_0)=0$ przy $0<U(z_0)<1$,
to $q'(z_0)>0$, co jest niemożliwe przy $q\ge0$
na całej prostej. Stąd $q>0$ i $U'<0$.

\textbf{Krok 2: wykluczenie prędkości mniejszych od krytycznej.}
Jeżeli $c\le0$, po podstawieniu $q=-U'>0$
do~\eqref{eq:pde27-phase} otrzymujemy $q'>0$,
a to przeczy jednocześnie całkowalności
$\int_\R q(z)\,\mathrm d\lambda_1(z)=1$
(wynika to z granic $U$ i $q=-U'$):
dla dowolnego $a\in\R$ mamy wówczas
$q(z)\ge q(a)>0$ dla $z\ge a$, co przeczy
skończoności całki.
Niech teraz $0<c<2\sqrt{Dr}$. Dla $z$ dostatecznie dużych
$U(z)$ jest tak małe, że
$\mu(z)=\frac rD(1-U(z))-\frac{c^2}{4D^2}\ge\kappa^2>0$.
Po podstawieniu $y(z)=e^{cz/(2D)}U(z)>0$
równanie~\eqref{eq:pde27-wave} przyjmuje postać
$y''+\mu(z)y=0$. Istotnie, pisząc
$U=e^{-cz/(2D)}y$, dostajemy
\[
 DU''+cU'=e^{-cz/(2D)}
 \left(Dy''-\frac{c^2}{4D}y\right),
\]
a składnik reakcji po podzieleniu przez
$De^{-cz/(2D)}$ ma postać $(r/D)(1-U)y$.
Możemy wybrać na przykład
$\kappa^2=\frac12(r/D-c^2/(4D^2))>0$,
a następnie $a$ tak duże, by
$(r/D)U(z)\le\kappa^2$ dla $z\ge a$.
Na przedziale $[a,a+\pi/\kappa]$, gdzie $a$ jest
dostatecznie duże, mnożymy równanie przez
$\sin(\kappa(z-a))$ i całkujemy względem $z$.
Po dwukrotnym całkowaniu przez części dostajemy
\[
 \int_a^{a+\pi/\kappa}(y''+\kappa^2y)
 \sin(\kappa(z-a))\,\mathrm d\lambda_1(z)
 =\kappa\bigl(y(a)+y(a+\pi/\kappa)\bigr)>0.
\]
Tymczasem $y''+\kappa^2y=-(\mu-\kappa^2)y\le0$,
a sinus jest nieujemny na całym przedziale.
Sprzeczność kończy dowód.
\end{proof}

\begin{przyklad}[Jawna fala o prędkości $5/\sqrt6$]
\label{ex:pde27-explicit}
Dla $D=r=1$ i $c=5/\sqrt6$ profil
\[
 U(z)=\frac{1}{(1+e^{z/\sqrt6})^2}
\]
rozwiązuje~\eqref{eq:pde27-wave} i spełnia
$U(-\infty)=1$, $U(+\infty)=0$, $U'<0$.
Podstawiając $w=e^{z/\sqrt6}$ mamy
$U'=-2w/(\sqrt6(1+w)^3)$ oraz
$U''=-w/(3(1+w)^3)+w^2/(1+w)^4$;
po sprowadzeniu do wspólnego mianownika $(1+w)^4$
licznik wyrażenia $U''+(5/\sqrt6)U'+U(1-U)$
wynosi
$-\frac13 w(1+w)+w^2-\frac53w(1+w)+2w+w^2=0$.
Znak $U'$ i granice wynikają bezpośrednio ze wzoru.
Dla ogólnych $D,r>0$ piszemy
$U(z)=V(\sqrt{r/D}\,z)$, gdzie $V$ jest powyższym
profilem dla $D=r=1$. Po podstawieniu i podzieleniu
przez $r$ równanie profilu ma postać
$V''+(c/\sqrt{Dr})V'+V(1-V)=0$.
Wybór $c=(5/\sqrt6)\sqrt{Dr}$ daje zatem
rozwiązanie
\[
 U(z)=\left(1+\exp\left(\sqrt{\frac r{6D}}\,z\right)\right)^{-2}.
\]
\end{przyklad}

\begin{center}
\begin{tikzpicture}[>=Latex]
 \begin{scope}[x=4.8cm,y=4.0cm]
  \fill[blue!9] (0,0)--(1,0)--(1,.8164966)--cycle;
  \draw[->] (-.04,0)--(1.13,0) node[right] {$U$};
  \draw[->] (0,-.035)--(0,1.02) node[above] {$q$};
  \draw[blue!70!black,thick] (0,0)--(1,.8164966)
    node[above left] {$q=kU$};
  \draw[dashed] (1,0)--(1,.8164966);
  \draw[->,thick] (.5,0)--(.5,.12);
  \draw[->,thick] (.68,.5552177)--(.58,.42);
  \draw[->,thick] (1,.45)--(.88,.40);
  \draw[red!75!black,very thick,domain=0:1,samples=101]
    plot (\x,{2/sqrt(6)*\x*(1-sqrt(\x))});
  \draw[red!75!black,->,thick] (.56,{2/sqrt(6)*.56*(1-sqrt(.56))})
    --(.46,{2/sqrt(6)*.46*(1-sqrt(.46))});
  \fill (0,0) circle (1.3pt) node[below left] {$0$};
  \fill (1,0) circle (1.3pt) node[below] {$1$};
  \node at (.52,.96) {Obszar fazowy};
 \end{scope}
 \begin{scope}[xshift=6.9cm,x=.43cm,y=3.0cm]
  \draw[->] (-6.5,0)--(6.7,0) node[right] {$z$};
  \draw[->] (0,-.05)--(0,1.19) node[above] {$U$};
  \draw[dashed,gray] (-6.3,1)--(6.3,1);
  \node[left] at (0,1) {$1$};
  \draw[red!75!black,very thick,domain=-6.3:6.3,samples=101]
    plot (\x,{1/(1+exp(\x/sqrt(6)))^2});
  \node[below] at (-6,0) {$-6$};
  \node[below] at (6,0) {$6$};
  \node at (0,1.36) {Profil fali};
 \end{scope}
\end{tikzpicture}

\small Rysunek: obszar niezmienniczy $\mathcal W$ i jawna
trajektoria $q=(2/\sqrt6)U(1-\sqrt U)$ (czerwona)
dla $D=r=1$, $c=5/\sqrt6$, $k=2/\sqrt6$.
Po prawej odpowiadający jej profil $U(z)$.
\end{center}

\section{Granice propagacji na prostej}
\begin{twierdzenie}[Górne oszacowanie prędkości ekspansji]
\label{thm:pde27-spreading-upper}
Niech $D,r>0$, a $u_0\in C_c(\R)$ spełnia
$0\le u_0\le1$. W przestrzeni $\operatorname{BUC}(\R)$ rzeczywistych
ograniczonych funkcji jednostajnie ciągłych istnieje
dokładnie jedno globalne rozwiązanie łagodne
$u\in C([0,\infty);\operatorname{BUC}(\R))$
równania $u_t=Du_{xx}+ru(1-u)$ z danymi $u_0$.
Spełnia ono $0\le u\le1$.
Wówczas dla każdej $c>2\sqrt{Dr}$ istnieją
stałe $C,\gamma>0$, zależne od $c,u_0,D,r$, takie, że
\[
 \sup_{|x|\ge ct}u(t,x)\le C e^{-\gamma t}
 \quad\text{dla wszystkich }t\ge0.
\]
\end{twierdzenie}
\begin{proof}
\textbf{Istnienie i porównanie na prostej.}
Przestrzeń $\operatorname{BUC}(\R)$ jest domkniętą
podprzestrzenią przestrzeni funkcji ograniczonych
z normą supremum, a więc jest przestrzenią Banacha.
Dla $t>0$ oznaczmy przez $S_D(t)$ splot z jądrem
$G_{Dt}(x)=(4\pi Dt)^{-1/2}\exp(-x^2/(4Dt))$,
a $S_D(0)=I$. Jądro jest nieujemne i ma całkę $1$.
Z tożsamości splotowej jąder Gaussa wynika własność
półgrupy; całkowanie daje kontraktywność, dodatniość
i zachowanie stałych. Splot zachowuje jednostajną
ciągłość, ponieważ
$\|S_D(t)v(\cdot+h)-S_D(t)v\|_\infty
\le\|v(\cdot+h)-v\|_\infty$.
Silna ciągłość przy $t=0$ wynika z rozdzielenia
całki z $G_{Dt}(y)(v(x-y)-v(x))$ na $|y|<\eta$
i $|y|\ge\eta$: na pierwszym zbiorze stosujemy
jednostajną ciągłość $v$, a na drugim zanik masy
ogona Gaussa przy $t\downarrow0$.
Silna ciągłość w pozostałych chwilach wynika
z własności półgrupy i kontraktywności.

Niech $\pi(s)=\min\{1,\max\{0,s\}\}$ oraz
$g(s)=r\pi(s)(1-\pi(s))$. Funkcja $g$ jest globalnie
lipschitzowska ze stałą $r$. Dla każdego $T>0$
na $C([0,T];\operatorname{BUC}(\R))$ rozważamy
operator
\[
 (\mathcal Pv)(t)=e^{-rt}S_D(t)u_0+
 \int_0^t e^{-r(t-s)}S_D(t-s)
       (g(v(s))+rv(s))\,\mathrm d\lambda_1(s).
\]
Ciągłość całki i kontrakcję uzasadniamy tak samo
jak w dowodzie Twierdzenia~\ref{thm:pde27-local}.
Dla $2rT<1$ operator jest kontrakcją na całej
przestrzeni, ponieważ $g(s)+rs$ ma stałą
Lipschitza co najwyżej $2r$. Funkcja $g(s)+rs$
jest niemalejąca z nierówności Lipschitza.
Stąd $\mathcal P$ zachowuje porządek punktowy.
Dla funkcji stale równych $0$ i $1$ mamy
$\mathcal P0\ge0$ oraz $\mathcal P1\le1$,
ponieważ $0\le u_0\le1$, $g(0)=g(1)=0$
i $S_D(t)$ zachowuje stałe.
Iteracje Picarda rozpoczęte od zera pozostają
w przedziale $[0,1]$; ich granica również.
Powtarzając konstrukcję na kolejnych przedziałach
tej samej długości, otrzymujemy rozwiązanie globalne.

W przedziale $[0,1]$ funkcja $g$ jest równa
$f(s)=rs(1-s)$. Wzór z $\mathcal P$ jest równoważny
zwykłemu wzorowi Duhamela z półgrupą $S_D$.
Istotnie, dla dowolnej rzeczywistej stałej $a$
mamy
\[
 e^{at}S_D(t)=S_D(t)+a\int_0^t
 S_D(t-s)e^{as}S_D(s)\,\mathrm d\lambda_1(s).
\]
Podstawienie tej tożsamości do wzoru Duhamela
i zastosowanie twierdzenia Fubiniego do całek
Bochnera (na skończonych przedziałach całkowane
funkcje są ograniczone w normie) pokazują,
że przejściu od $S_D$ do $e^{at}S_D(t)$
towarzyszy zastąpienie $f(u)$ przez $f(u)-au$.
Odwrotna transformacja odpowiada stałej $-a$.
Jednoznaczność w całej klasie rozwiązań łagodnych
dla $f$ wynika z oszacowania Gronwalla na każdym
skończonym przedziale, na którym oba rozwiązania
mają ograniczone normy supremum.

\textbf{Oszacowanie gaussowskie.}
Stosujemy powyższą transformację z $a=r$.
Ponieważ $f(u)-ru=-ru^2$, otrzymujemy
\[
 u(t)=e^{rt}S_D(t)u_0-
 r\int_0^t e^{r(t-s)}S_D(t-s)u(s)^2\,
                       \mathrm d\lambda_1(s).
\]
Dodatniość półgrupy i wcześniej wykazane
$0\le u\le1$ dają
\[
 0\le u(t,x)\le e^{rt}\int_{\R}
 \frac{e^{-(x-y)^2/(4Dt)}}{\sqrt{4\pi Dt}}u_0(y)\,\mathrm d\lambda_1(y),
 \qquad t>0.
\]
Wybierzmy $R>0$ tak, aby $\supp u_0\subset[-R,R]$,
oraz liczbę $c_0$ z $2\sqrt{Dr}<c_0<c$.
Niech $t_0=\max\{1,R/(c-c_0)\}$. Dla $t\ge t_0$
i $|x|\ge ct$ zachodzi $|x-y|\ge c_0t$
dla każdego $y\in\supp u_0$. Stąd
\[
 u(t,x)\le\frac{\|u_0\|_{L^1(\R)}}{\sqrt{4\pi Dt}}
 \exp\!\left(-\left(\frac{c_0^2}{4D}-r\right)t\right).
\]
Liczba $\gamma=c_0^2/(4D)-r$ jest dodatnia,
a czynnik $t^{-1/2}\le1$ dla $t\ge t_0$.
Na odcinku $0\le t<t_0$ korzystamy z $u\le1$.
Wystarczy więc przyjąć
$C=\max\{e^{\gamma t_0},\|u_0\|_{L^1}/\sqrt{4\pi D}\}$,
co daje żądane oszacowanie dla wszystkich $t\ge0$.
\end{proof}

\begin{uwaga}[Co udowodniono, a czego nie]
Twierdzenie~\ref{thm:pde27-spreading-upper} daje wyłącznie
górne oszacowanie szybkości rozprzestrzeniania się.
Dopasowana dolna granica prędkości i pełne twierdzenie
o asymptotyce rozwiązania na całej prostej wymagają
osobnej konstrukcji podrozwiązań. Nie wynikają
same z istnienia fal ani z powyższej bariery.
\end{uwaga}

\section{Funkcjonały Lapunowa i równanie Allena--Cahna}
\begin{twierdzenie}[Dyssypacja energii i stany graniczne]
\label{thm:pde27-allen}
Niech $D,L>0$, $d\ge1$, a
$u_0\in C(\T_L^d)\cap H^1(\T_L^d)$.
Istnieje dokładnie jedno globalne rozwiązanie łagodne
$u$ równania $u_t=D\Delta u+u-u^3$ z danymi $u_0$.
Jest ono klasyczne dla dodatnich czasów.
Dla funkcjonału
\[
 E(u)=\int_{\T_L^d}\left(\frac D2|\nabla u|^2
 +\frac14(u^2-1)^2\right)\,\mathrm d\lambda_d(x)
\]
zachodzi dla $t\ge0$ tożsamość
$E(u(t))+\int_0^t\|u_s(s)\|_2^2\,\mathrm d\lambda_1(s)=E(u_0)$.
Ponadto istnieje ciąg $t_n\to\infty$ i funkcja
$u_\infty\in H^1(\T_L^d)$ takie, że
$u(t_n)\to u_\infty$ silnie w $L^2$ i
$-D\Delta u_\infty=u_\infty-u_\infty^3$
w sensie słabym. Nie twierdzimy bez dodatkowych
założeń o zbieżności całej trajektorii do jednego
stanu stacjonarnego.
\end{twierdzenie}
\begin{proof}
Wybieramy $M\ge\max\{1,\|u_0\|_\infty\}$.
Stałe $-M$ i $M$ są odpowiednio pod- i
nadrozwiązaniem, gdyż $f(-M)=-M+M^3\ge0$
i $f(M)=M-M^3\le0$. Z zasady porównania
otrzymujemy $|u(t,x)|\le M$ i globalność.
Aby uzasadnić tożsamość energii bez dodatkowych
założeń o mieszanych pochodnych, stosujemy
Stwierdzenie~\ref{prop:pde-v44-torus-scaling}.
Oznaczmy $a=D(2\pi/L)^2$, $b=L/(2\pi)$,
$\tau=at$ i $w(\tau,\theta)=u(t,b\theta)$.
Wówczas $w_\tau=\Delta_\theta w-\widetilde V'(w)$,
gdzie $\widetilde V(s)=(s^2-1)^2/(4a)$.
Dane $w(0)$ należą do $C(\T^d)\cap H^1(\T^d)$,
ponieważ liniowe skalowanie zachowuje te przestrzenie.
Stosujemy Twierdzenie~\ref{thm:pde25-energy},
z funkcjonałem
$\widetilde E(w)=\int_{\T^d}
(\frac12|\nabla_\theta w|^2+\widetilde V(w))
\,\mathrm d\lambda_d(\theta)$.
Zmiana zmiennych i reguła łańcuchowa dają
\[
 E(u(t))=ab^d\widetilde E(w(at)),\qquad
 \|u_t(t)\|_2^2=a^2b^d\|w_\tau(at)\|_2^2.
\]
Mnożąc tożsamość energii dla $w$ przez $ab^d$
i podstawiając $\sigma=as$ w całce czasowej,
otrzymujemy dokładnie żądaną tożsamość dla $u$,
również od chwili $0$.
W szczególności
$\|\nabla u(t)\|_2^2\le2E(u_0)/D$ oraz
$\|u(t)\|_2^2\le M^2L^d$. Zatem
$\sup_{t\ge0}\|u(t)\|_{H^1}<\infty$.
Ponieważ energia jest nieujemna,
$\int_0^\infty\|u_t(t)\|_2^2\,\mathrm d\lambda_1(t)
\le E(u_0)$ na mocy zbieżności monotonicznej.
Dla każdego $n\ge1$ wybieramy $t_n\in[n,n+1]$
takie, że
$\|u_t(t_n)\|_2^2\le\int_n^{n+1}\|u_t(s)\|_2^2
\,\mathrm d\lambda_1(s)$.
Jest to możliwe dzięki ciągłości funkcji podcałkowej
dla dodatnich czasów. Prawa strona dąży do zera
z całkowalności na półprostej.
Słaba zwartość ograniczonych ciągów w przestrzeni
Hilberta $H^1$ oraz zwarte zanurzenie
$H^1(\T_L^d)\hookrightarrow L^2(\T_L^d)$
pozwalają po przejściu do podciągu założyć
$u(t_n)\rightharpoonup u_\infty$ w $H^1$ i
$u(t_n)\to u_\infty$ w $L^2$.
Z silnej zbieżności w $L^2$ możemy wybrać dalszy
podciąg zbieżny prawie wszędzie; stąd $|u_\infty|\le M$
prawie wszędzie. Funkcja $f(s)=s-s^3$ jest
lipschitzowska na $[-M,M]$ ze stałą co najwyżej
$1+3M^2$, więc
$\|f(u(t_n))-f(u_\infty)\|_2
\le(1+3M^2)\|u(t_n)-u_\infty\|_2\to0$.
Dla każdej $\varphi\in H^1(\T_L^d)$ równanie
w chwili $t_n>0$ daje
\[
 D\int_{\T_L^d}\nabla u(t_n)\cdot\nabla\varphi\,
                         \mathrm d\lambda_d(x)
 =\int_{\T_L^d}(f(u(t_n))-u_t(t_n))\varphi\,
                         \mathrm d\lambda_d(x).
\]
Lewa strona zbiega dzięki słabej zbieżności
gradientów, prawa dzięki silnej zbieżności w $L^2$
i zanikaniu $\|u_t(t_n)\|_2$.
Otrzymujemy słabą postać równania stacjonarnego
dla $u_\infty$.
\end{proof}

\section{Układy i niestabilność Turinga}
Dla wektorowej funkcji $U=(u_1,u_2)^{\mathsf T}$
rozważamy $U_t=\mathcal D\Delta U+F(U)$,
gdzie $F\in C^1(\R^2;\R^2)$,
$\mathcal D=\operatorname{diag}(D_1,D_2)$, $D_i>0$.
Stan jednorodny $U_*$ spełnia $F(U_*)=0$.
Po podstawieniu $U=U_*+\varepsilon V$
i podzieleniu równania przez $\varepsilon$
składnik reakcji przyjmuje postać
$\varepsilon^{-1}(F(U_*+\varepsilon V)-F(U_*))$.
Z różniczkowalności $F$ jego granicą jest $JV$,
gdzie $J=DF(U_*)$. Linearyzacja ma więc postać
$V_t=\mathcal D\Delta V+JV$.
W odróżnieniu od przypadku skalarnego dyfuzja może
wywołać niestabilność stanu stabilnego dla samego ODE.


\begin{twierdzenie}[Kryterium liniowej niestabilności Turinga]
\label{thm:pde27-turing-general}
Niech $J=\begin{psmallmatrix}a&b\\c&d\end{psmallmatrix}$
będzie rzeczywistą macierzą $2\times2$ oraz
$\mathcal D=\operatorname{diag}(D_1,D_2)$, $D_1,D_2>0$.
Załóżmy $\operatorname{tr}J<0$ i $\det J>0$,
tak że jednorodny układ liniowy $V_t=JV$ jest
asymptotycznie stabilny. Dla $\mu>0$ macierz
$J-\mu\mathcal D$ ma wartość własną o dodatniej
części rzeczywistej wtedy i tylko wtedy, gdy
\begin{equation}\label{eq:pde27-turing-poly}
 D_1D_2\mu^2-(D_2a+D_1d)\mu+\det J<0.
\end{equation}
Takie $\mu>0$ istnieje wtedy i tylko wtedy, gdy
\[
 S:=D_2a+D_1d>0,\qquad S^2>4D_1D_2\det J.
\]
Na torusie lub ograniczonym obszarze Lipschitzowskim
z jednorodnym warunkiem Neumanna występowanie
wykładniczo rosnącego modu jest równoważne temu,
że co najmniej jedna dodatnia wartość własna
$\mu_k$ operatora $-\Delta$ spełnia
\eqref{eq:pde27-turing-poly}; warunki na $S$ i wyróżnik
nie wystarczają bez tego warunku spektralnego.
\end{twierdzenie}
\begin{proof}
Jeżeli $-\Delta\phi=\mu\phi$ i
$V(t,x)=v(t)\phi(x)$, gdzie $v(t)\in\R^2$,
to podstawienie do równania liniowego daje
$v'=(J-\mu\mathcal D)v$. Oznaczamy tę macierz
przez $M_\mu$.
Z hipotez wynika
$\operatorname{tr}M_\mu=\operatorname{tr}J-
\mu(D_1+D_2)<0$. Dla rzeczywistej macierzy $2\times2$
o ujemnym śladzie dodatnią część rzeczywistą
wartości własnej otrzymujemy dokładnie przy
$\det M_\mu<0$: wartości własne są wtedy
rzeczywiste i mają przeciwne znaki.
Jeśli wyznacznik jest dodatni, to rzeczywiste
pierwiastki mają ten sam znak, a ich ujemna suma
wymusza oba znaki ujemne; dla pary zespolonej
sprzężonej część rzeczywista jest połową śladu.
Przy wyznaczniku zerowym pierwiastki to $0$
i $\operatorname{tr}M_\mu<0$. Są różne, więc
mod jest ograniczony, lecz nie asymptotycznie stabilny.
Zatem dodatni wykładnik wzrostu zachodzi dokładnie
dla ujemnego wyznacznika. Rachunek
\[
 \det M_\mu=(a-D_1\mu)(d-D_2\mu)-bc
 =D_1D_2\mu^2-(D_2a+D_1d)\mu+\det J
\]
daje wielomian z~\eqref{eq:pde27-turing-poly}.
Jeżeli $S\le0$, wszystkie składniki trójmianu
są nieujemne dla $\mu>0$, a wyraz wolny jest dodatni.
Jeżeli $S>0$, minimum na dodatniej półosi
jest osiągane w $\mu_0=S/(2D_1D_2)$ i wynosi
$\det J-S^2/(4D_1D_2)$. Jest ujemne dokładnie
przy dodatnim wyróżniku. Wtedy oba pierwiastki
są dodatnie, bo mają dodatnią sumę i iloczyn.

Na torusie $\T_L^d$ funkcje
$\exp(2\pi i k\cdot x/L)$, $k\in\mathbb Z^d$,
tworzą po normalizacji bazę ortonormalną w $L^2$
i odpowiadają wartościom
$\mu_k=(2\pi/L)^2|k|^2$. Dla warunku Neumanna
na ograniczonym obszarze Lipschitzowskim analogiczny
rozkład daje baza funkcji własnych realizacji $-\Delta$
z formą $\int_\Omega\nabla v\cdot\nabla w\,
\mathrm d\lambda_d(x)$ na $H^1(\Omega)$.
Każdy współczynnik ewoluuje niezależnie według
macierzy $M_{\mu_k}$, zatem rosnący wykładniczo
mod istnieje dokładnie przy ujemnym wyznaczniku
dla co najmniej jednej dopuszczalnej wartości $\mu_k$.
Dla dużych $\mu$ zachodzi ponadto
\[
 \tfrac12\frac{\mathrm d}{\mathrm dt}|v|^2
 \le(\|J\|-\mu\min\{D_1,D_2\})|v|^2,
\]
więc wszystkie dostatecznie wysokie mody
zanikają wykładniczo z jednostajnym oszacowaniem.
Pozostało tylko skończenie wiele wartości własnych
(z krotnościami); gdy żaden mod nie rośnie,
ich ewolucje są ograniczone, a jeśli wszystkie
wyznaczniki są dodatnie, zanikają wykładniczo.
Przez tożsamość Parsevala daje to odpowiednio
ograniczoność lub wykładniczy zanik w $L^2$.
\end{proof}

\begin{uwaga}[Przedział niestabilnych częstości przestrzennych]
Jeżeli warunki z poprzedniego twierdzenia są spełnione,
niestałe mody rosną dokładnie dla
$\mu\in(\mu_-,\mu_+)$, gdzie
$\mu_\pm=(S\pm\sqrt{S^2-4D_1D_2\det J})/(2D_1D_2)$.
Na torusie należy więc sprawdzić, czy
$\mu_-<(2\pi/L)^2|k|^2<\mu_+$ dla pewnego $k\ne0$.
Dla jednakowych współczynników dyfuzji $D_1=D_2=D$
mamy $M_\mu=J-\mu DI$, a jego wartości własne
powstają z wartości własnych $J$ przez odjęcie
$\mu D$; destabilizacja dyfuzyjna nie jest możliwa.
W przykładzie poniżej $\mu_-=1/5$ i $\mu_+=1/2$;
wybór $\mu_1=3/10$ rzeczywiście trafia do tego przedziału.
\end{uwaga}


\begin{stwierdzenie}[Kryterium niestabilności konkretnego modu]
\label{prop:pde27-turing}
Niech $\Omega=(0,L)$, na końcach niech obowiązują
jednorodne warunki Neumanna, oraz
\[
 J=\begin{pmatrix}1&-2\\2&-3\end{pmatrix},\qquad
 \mathcal D=\begin{pmatrix}1&0\\0&10\end{pmatrix}.
\]
Macierz $J$ jest stabilna w sensie liniowym dla ODE
$\dot V=JV$: obie jej wartości własne wynoszą $-1$.
Jeżeli $L=\pi/\sqrt{3/10}$, to odpowiedni
jednomodowy układ reakcji--dyfuzji ma dodatnią
rzeczywistą wartość własną. W szczególności
zero jest liniowo niestabilne dla PDE, mimo
stabilności przestrzennie jednorodnej dynamiki.
\end{stwierdzenie}
\begin{proof}
Funkcje $\cos(n\pi x/L)$ są funkcjami własnymi
laplasjanu Neumanna z wartościami $-\mu_n$,
$\mu_n=(n\pi/L)^2$. Na takim modzie ewolucja
jest wyznaczona przez macierz $J-\mu_n\mathcal D$.
Dla $n=0$ mamy
$\operatorname{tr}J=-2$, $\det J=1$, więc
wielomian charakterystyczny jest $(\lambda+1)^2$.
Macierz $J+I$ jest niezerowa i $(J+I)^2=0$,
więc $e^{tJ}=e^{-t}(I+t(J+I))\to0$, także
mimo braku diagonalizowalności.
Dla $n=1$, $\mu_1=3/10$, otrzymujemy
\[
 \operatorname{tr}(J-\mu_1\mathcal D)=-\frac{53}{10},
 \qquad \det(J-\mu_1\mathcal D)
 =1-7\frac3{10}+10\frac9{100}=-\frac15.
\]
Ujemny wyznacznik oznacza jedną dodatnią i jedną
ujemną rzeczywistą wartość własną. Dodatnia wartość własna wynosi
$\lambda_+=(-53+\sqrt{2889})/20>0$.
Jeżeli $v\ne0$ jest odpowiadającym jej rzeczywistym
wektorem własnym, to
$V(t,x)=e^{\lambda_+t}v\cos(\pi x/L)$
jest rzeczywistym rozwiązaniem liniowego PDE
z warunkiem Neumanna, rosnącym wykładniczo.
Ponieważ $\mu_n=(3/10)n^2$, jedynie $n=1$
należy do przedziału $(1/5,1/2)$.
To dowodzi twierdzenia.
\end{proof}

\begin{center}
\begin{tikzpicture}[x=11cm,y=2.1cm,>=Latex]
 \fill[red!10] (.2,-.33) rectangle (.5,0);
 \draw[->] (-.025,0)--(.74,0) node[right] {$\mu$};
 \draw[->] (0,-.36)--(0,1.35) node[above] {$\det M_\mu$};
 \draw[blue!75!black,thick,domain=0:.69,samples=101]
  plot (\x,{10*\x*\x-7*\x+1});
 \draw[densely dashed,gray] (.2,-.30)--(.2,0);
 \draw[densely dashed,gray] (.5,-.30)--(.5,0);
 \fill (.2,0) circle (1.5pt) node[above left] {$\mu_-$};
 \fill (.5,0) circle (1.5pt) node[above right] {$\mu_+$};
 \draw[red!75!black,dashed] (.3,0)--(.3,-.2);
 \fill[red!75!black] (.3,-.2) circle (2pt);
 \node[below,red!75!black] at (.3,-.34) {$\mu_1=3/10$};
 \node[left] at (0,1) {$1$};
\end{tikzpicture}

\small Rysunek: trójmian $\det M_\mu=10\mu^2-7\mu+1$.
Zacieniowany przedział odpowiada ujemnemu wyznacznikowi;
jedyny dodatni mod Neumanna w tym przedziale to $\mu_1$.
\end{center}

\begin{uwaga}[Stabilność liniowa a nieliniowa]
Powyższy rachunek dowodzi niestabilności
\emph{zlinearyzowanego} PDE. Aby wnioskować o
nieliniowej niestabilności układu
$U_t=\mathcal D\Delta U+F(U)$, trzeba dodatkowo
sprecyzować regularność $F$ oraz zastosować
twierdzenie o niestabilności ze spektralnej wartości
własnej o dodatniej części rzeczywistej.
Takiego twierdzenia nie ukrywamy w argumencie liniowym.
\end{uwaga}

\section{Podsumowanie pojęciowe}
Monotoniczność operatorów z rozdziału~\ref{chap:pde26}
prowadziła do istnienia nieliniowych rozwiązań eliptycznych.
W obecnym rozdziale kluczową rolę odgrywają
\emph{porównanie} i \emph{dynamika półgrupy}.
W skalarnej teorii KPP dyfuzja i reakcja prowadzą
do frontu o minimalnej prędkości $2\sqrt{Dr}$;
w układach z różnymi współczynnikami dyfuzji
mogą natomiast powstawać niestabilne mody przestrzenne.
Tych dwóch mechanizmów nie należy utożsamiać.


\chapter[Równania Naviera--Stokesa]{Równania Naviera--Stokesa: rozkład Helmholtza, rozwiązania Leraya--Hopfa i struktura wirów}
\label{chap:pde28}

\section{Równanie nieściśliwego płynu i wybór przestrzeni}
Niech $d\in\{2,3\}$, $\nu>0$, $T>0$ i $\T^d=(\R/2\pi\mathbb Z)^d$.
Szukamy prędkości $u:(0,T)\times\T^d\to\R^d$ i ciśnienia
$p:(0,T)\times\T^d\to\R$ spełniających
\begin{equation}\label{eq:pde28-ns}
 \partial_tu+(u\cdot\nabla)u-\nu\Delta u+\nabla p=0,
 \qquad \diver u=0,
 \qquad u(0)=u_0.
\end{equation}
Współczynnik $\nu$ jest lepkością kinematyczną, $u$ jest \emph{polem
prędkości}, natomiast ciśnienie jest mnożnikiem wymuszającym warunek
nieściśliwości. Przyjęcie torusa pozwala wyprowadzić zasadnicze
argumenty bez dodatkowych założeń o regularności brzegu. Dla
pól wektorowych symbol $\nabla u$ oznacza macierz
$(\partial_j u_i)_{i,j}$, a $\diver(u\otimes u)$ jest wektorem
o współrzędnych $\sum_j\partial_j(u_i u_j)$.

Zakładamy $\int_{\T^d}u_0\,\mathrm d\lambda_d(x)=0$.
Dla gładkich rozwiązań średnia prędkości zachowuje się
w czasie: z $\diver u=0$ wynika
$(u\cdot\nabla)u=\diver(u\otimes u)$, a całki
z pochodnych periodycznych znikają. Całkując
równanie~\eqref{eq:pde28-ns} w przestrzeni, otrzymujemy
$\frac{\mathrm d}{\mathrm dt}\int_{\T^d}u\,\mathrm d\lambda_d=0$.
Dla rozwiązań dystrybucyjnych analogiczny rachunek
przeprowadzamy, testując równanie stałym wektorem
w przestrzeni pomnożonym przez gładką funkcję czasu
o zwartym nośniku.
Jeżeli średnia wynosi $m=(2\pi)^{-d}\int_{\T^d}u_0\,
\mathrm d\lambda_d$, to zmiana układu odniesienia
\[
 v(t,x)=u(t,x+tm)-m,\qquad
 \widetilde p(t,x)=p(t,x+tm)
\]
daje równanie tej samej postaci dla $v$ o średniej zero.
Istotnie, dodatkowy składnik $m\cdot\nabla u$
w pochodnej czasowej znosi się ze składnikiem
$-m\cdot\nabla u$ w członie konwekcyjnym;
laplasjan i dywergencja komutują z translacją.
Dla pól o średniej zero używamy przestrzeni
\[
 \begin{aligned}
 H&=\Bigl\{v\in L^2(\T^d;\R^d):\
 \diver v=0\text{ w }\mathcal D'(\T^d),\
 \int_{\T^d}v\,\mathrm d\lambda_d=0\Bigr\},\\
 V&=H\cap H^1(\T^d;\R^d).
 \end{aligned}
\]
Przestrzeń $H$ jest domknięta w $L^2$: średnia
jest ciągłym odwzorowaniem liniowym, a warunek
dystrybucyjny jest przecięciem jąder ciągłych
funkcjonałów $v\mapsto\int_{\T^d}v\cdot\nabla\psi\,
\mathrm d\lambda_d$, $\psi\in C^\infty(\T^d)$.
Analogicznie $V$ jest domknięta w $H^1$.
Przyjmujemy współczynniki Fouriera unormowane przez
$\widehat v(k)=(2\pi)^{-d}\int_{\T^d}v(x)e^{-ik\cdot x}
\,\mathrm d\lambda_d(x)$, natomiast normy $L^2$
obliczamy względem pełnej miary Lebesgue'a.
Z twierdzenia Parsevala~\ref{thm:pde6-parseval}
i zerowej średniej dla $v\in V$ mamy
\[
 \|v\|_2^2=(2\pi)^d\sum_{k\ne0}|\widehat v(k)|^2
 \le (2\pi)^d\sum_{k\ne0}|k|^2|\widehat v(k)|^2
 =\|\nabla v\|_2^2.
\]
Na $V$ przyjmujemy iloczyn skalarny
$(v,w)_V=(\nabla v,\nabla w)_{L^2}$ i normę
$\|v\|_V=\|\nabla v\|_2$, równoważną normie $H^1$,
zatem $V$ jest przestrzenią Hilberta.
Skończone sumy Fouriera pól z $H$ zachowują zerową
średnią i warunek $k\cdot\widehat v(k)=0$,
a przy symetrycznym obcięciu częstości również
rzeczywistość pola. Są gęste w $H$, a dla $v\in V$
zbiegają do $v$ również w normie $V$.
Zatem zanurzenie $V\hookrightarrow H$ jest gęste
i ciągłe. Oznaczamy przez $V'$ przestrzeń dualną $V$
z normą $\|F\|_{V'}=\sup_{\|\varphi\|_V\le1}
|\langle F,\varphi\rangle|$.
Zanurzenie $H\hookrightarrow V'$ przyporządkowuje
$h$ funkcjonałowi $\varphi\mapsto(h,\varphi)_H$;
jest ciągłe i injektywne dzięki gęstości $V$.
Jego obraz jest gęsty w $V'$: gdyby nie był, istniałby
niezerowy ciągły funkcjonał na $V'$ zerujący się
na tym obrazie. Refleksywność $V$ pozwala zapisać go
jako ewaluację w pewnym $\varphi\in V\setminus\{0\}$,
co dawałoby $(h,\varphi)_H=0$ dla wszystkich $h\in H$.
Otrzymujemy trójkę przestrzeni Hilberta
$V\subset H\simeq H'\subset V'$ z gęstymi,
ciągłymi zanurzeniami.

\section{Rozkład Helmholtza i rzut Leraya}
\begin{twierdzenie}[Rozkład Helmholtza na torusie]\label{thm:pde28-helmholtz}
Niech $d\ge2$ (przestrzeń $H$ definiujemy jak powyżej
również dla tych wymiarów) i niech
$F\in L^2(\T^d;\R^d)$ ma średnią zero.
Istnieje dokładnie jedna para $(v,\phi)$, gdzie
$v\in H$, $\phi\in H^1(\T^d)$ ma średnią zero, taka że
$F=v+\nabla\phi$. Składniki są ortogonalne w $L^2$.
Przyporządkowanie $PF=v$ jest ortogonalnym rzutem
$L^2_0(\T^d;\R^d)\to H$, gdzie indeks $0$ oznacza
podprzestrzeń pól o średniej zero. Dla
$k\in\mathbb Z^d\setminus\{0\}$
\begin{equation}\label{eq:pde28-P}
 \widehat{PF}(k)=\left(I-\frac{k\otimes k}{|k|^2}\right)\widehat F(k),
 \qquad \widehat\phi(k)=-\frac{i k\cdot\widehat F(k)}{|k|^2}.
\end{equation}
\end{twierdzenie}
\begin{proof}
Przyjmujemy konwencję $\widehat F(k)=(2\pi)^{-d}
\int_{\T^d}F(x)e^{-ik\cdot x}\,\mathrm d\lambda_d(x)$.
Tak jak dla szeregów Fouriera w rozdziale~\ref{chap:pde-fourier}
współczynniki są uśredniane; normy $L^2$ użyte przy oszacowaniach
energii płynu odnoszą się jednak do pełnej miary $\lambda_d$,
bez dzielenia przez $(2\pi)^d$.
Dla każdego $k\ne0$ macierz w~\eqref{eq:pde28-P} jest rzutem
ortogonalnym na podprzestrzeń prostopadłą do $k$.
Ustalamy $\widehat{PF}(0)=\widehat\phi(0)=0$.
Ponieważ norma operatorowa macierzy nie przekracza $1$,
twierdzenie Parsevala daje ciąg współczynników
o skończonej sumie kwadratów, a zupełność bazy Fouriera
wyznacza $PF\in L^2$. Dla rzeczywistego $F$
mamy $\widehat F(-k)=\overline{\widehat F(k)}$.
Macierz rzutu jest rzeczywista i parzysta względem $k$,
więc współczynniki $PF$ mają tę samą symetrię.
Współczynniki $\phi$ również spełniają
$\widehat\phi(-k)=\overline{\widehat\phi(k)}$.
Obie funkcje są więc rzeczywiste.
Mamy $k\cdot\widehat{PF}(k)=0$ dla każdego $k$,
co jest równoważne $\diver PF=0$ w sensie dystrybucji
(przez różniczkowanie szeregów w sensie dystrybucyjnym).
Średnia $PF$ jest zerowa.
Z drugiej strony $ik\widehat\phi(k)=\widehat F(k)-\widehat{PF}(k)$.
Nierówność $|k||\widehat\phi(k)|\le|\widehat F(k)|$
oraz $|k|\ge1$ dla $k\ne0$ dają
\[
 \sum_k(1+|k|^2)|\widehat\phi(k)|^2
 \le 2\sum_k|\widehat F(k)|^2<\infty.
\]
Stąd $\phi\in H^1$ i równość $F=PF+\nabla\phi$
wynika z równości wszystkich współczynników Fouriera.
Dla $v\in H$ i $\phi\in H^1$ wybieramy gładkie
periodyczne wielomiany trygonometryczne
$\phi_n\to\phi$ w $H^1$. Definicja dywergencji
dystrybucyjnej daje $\int_{\T^d}v\cdot\nabla\phi_n\,
\mathrm d\lambda_d=0$. Nierówność Cauchy'ego--Schwarza
pozwala przejść do granicy, więc
$\int_{\T^d}v\cdot\nabla\phi\,\mathrm d\lambda_d=0$.
Jest to ortogonalność składników. Jeżeli $F=v_1+\nabla\phi_1
=v_2+\nabla\phi_2$, to różnica obu rozkładów leży jednocześnie
w $H$ i w jego ortogonalnym dopełnieniu, zatem znika.
Warunek zerowej średniej wyznacza jednoznacznie $\phi$:
jeśli $\nabla(\phi_1-\phi_2)=0$, to nierówność
Poincarégo daje $\phi_1=\phi_2$.
Ponieważ $PF\in H$ oraz $F-PF\perp H$, jest to
ortogonalny rzut na $H$, a więc $P^2=P$ i $P^*=P$.
Wzór na współczynniki definiuje również $P$ na
periodycznych przestrzeniach $H^s$ o średniej zero
dla $s\in\R$; mnożnik ma normę co najwyżej $1$
dla każdego $k$, więc jest kontrakcją w normie
Fouriera z wagą $(1+|k|^2)^s$.
\end{proof}

\begin{uwaga}[Operator Stokesa]
Na $H$ określamy $A=-P\Delta$ z dziedziną
$D(A)=H^2(\T^d;\R^d)\cap H$. Ponieważ $P$ komutuje
z $\Delta$ na torusie, mamy $Av=-\Delta v$, a
$(Av,v)_H=\|\nabla v\|_2^2$. Dokładniej, dla każdej
pary częstości $\{k,-k\}$, $k\ne0$, wybieramy bazę
ortonormalną $k^\perp$ w $\R^d$ i mnożymy jej wektory
przez odpowiednio unormowane $\cos(k\cdot x)$
i $\sin(k\cdot x)$. Otrzymujemy rzeczywistą bazę
ortonormalną $(w_j)$ przestrzeni $H$, złożoną
z funkcji własnych $A$ o wartościach własnych
$\lambda_j=|k_j|^2$. Numerujemy je niemalejąco,
z krotnościami, więc $\lambda_j\ge1$ i
$\lambda_j\to\infty$.
Dziedzina $D(A)$ odpowiada warunkowi
$\sum_j\lambda_j^2|(v,w_j)_H|^2<\infty$,
równoważnemu warunkowi $H^2$ na polach o średniej zero.
Operator jest zatem unitarnie równoważny rzeczywistemu
operatorowi mnożenia przez ciąg $(\lambda_j)$
na jego maksymalnej dziedzinie w $\ell^2$.
Jest samosprzężony: jeśli ciąg należy do dziedziny
operatora sprzężonego, testowanie przez wektory
jednostkowe wymusza, by ciąg pomnożony przez
$(\lambda_j)$ należał do $\ell^2$, co jest dokładnie
warunkiem dziedziny $A$.
Mamy $A\ge I$ i $A^{-1}$ jest ograniczony na $H$.
Obcięcie mnożnika $\lambda_j^{-1}$ do pierwszych $m$
współrzędnych daje operator skończonego rzędu,
którego odległość od $A^{-1}$ w normie operatorowej
wynosi $\sup_{j>m}\lambda_j^{-1}\to0$.
Stąd $A^{-1}$ jest zwarty. Ponadto
$\|v\|_V^2=\sum_j\lambda_j|(v,w_j)_H|^2$,
więc
$\sum_{j>m}|(v,w_j)_H|^2\le
\lambda_{m+1}^{-1}\|v\|_V^2$.
Ten jednostajny szacunek ogonów i skończenie wymiarowe
projekcje dowodzą zwartości zanurzenia $V\hookrightarrow H$.
Rzut $P$ nie jest mnożeniem przez stałą macierz w przestrzeni
fizycznej, lecz nielokalnym operatorem Fouriera.
\end{uwaga}

\section{Tożsamość energii i słaba postać równania}
Dla $a,b,c\in V$ określamy formę trójliniową
\[
 b(a,b,c)=\int_{\T^d}(a\cdot\nabla)b\cdot c\,\mathrm d\lambda_d(x).
\]
W wymiarach $2$ i $3$ zanurzenie Sobolewa
$H^1(\T^d)\hookrightarrow L^4(\T^d)$ jest ciągłe.
Stąd nierówność Höldera z wykładnikami $4,2,4$ daje
\[
 |b(a,b,c)|\le\|a\|_4\|\nabla b\|_2\|c\|_4
 \le C\|a\|_V\|b\|_V\|c\|_V.
\]
Forma jest zatem dobrze określona i ciągła na $V^3$.
Dla gładkich pól periodycznych, gdy $\diver a=0$,
\[
 b(a,b,c)+b(a,c,b)
 =\int_{\T^d}a\cdot\nabla(b\cdot c)\,
 \mathrm d\lambda_d
 =-\int_{\T^d}(\diver a)(b\cdot c)\,\mathrm d\lambda_d=0.
\]
Aproksymujemy trzy pola ich skończonymi sumami Fouriera
w $V$ i przechodzimy do granicy dzięki ciągłości formy.
Otrzymujemy $b(a,b,c)=-b(a,c,b)$, w szczególności
$b(a,b,b)=0$. Oznacza to, że konwekcja nie wnosi
wkładu do zmiany energii kinetycznej; za dyssypację
odpowiada człon lepkościowy.
Dla gładkiego rozwiązania testowanie równania przez $u$
i periodyczne całkowanie przez części dają bowiem
\[
 \tfrac12\frac{\mathrm d}{\mathrm dt}\|u\|_2^2
 +\nu\|\nabla u\|_2^2=0,
\]
gdyż $\int_{\T^d}\nabla p\cdot u\,\mathrm d\lambda_d
=-\int_{\T^d}p\,\diver u\,\mathrm d\lambda_d=0$.
Testowanie przez stałe w czasie $\varphi\in V$
i całkowanie względem czasu prowadzi do słabej
postaci równania podanej niżej.

\begin{definicja}[Rozwiązanie Leraya--Hopfa na torusie]\label{def:pde28-leray}
Niech $u_0\in H$, $\nu>0$ i $T>0$. Pole $u$ nazywamy
rozwiązaniem Leraya--Hopfa zagadnienia~\eqref{eq:pde28-ns} na $[0,T]$,
jeżeli $u\in L^\infty(0,T;H)\cap L^2(0,T;V)$,
$u\in C([0,T];H_{\mathrm w})$, $u(0)=u_0$ oraz dla każdego
$\varphi\in V$ i $t\in[0,T]$ spełnia tożsamość
\begin{equation}\label{eq:pde28-weak}
 (u(t),\varphi)_H-(u_0,\varphi)_H
 +\nu\int_0^t(\nabla u(s),\nabla\varphi)_{L^2}\,\mathrm d\lambda_1(s)
 +\int_0^t b(u(s),u(s),\varphi)\,\mathrm d\lambda_1(s)=0.
\end{equation}
Zapis $C([0,T];H_{\mathrm w})$ oznacza ciągłość wszystkich
funkcji skalarnych $t\mapsto(u(t),h)_H$, $h\in H$;
wybieramy ten słabo ciągły reprezentant klasy
funkcji określonych prawie wszędzie.
Wymagamy ponadto nierówności energii
\begin{equation}\label{eq:pde28-energy}
 \frac12\|u(t)\|_2^2+\nu\int_0^t\|\nabla u(s)\|_2^2\,\mathrm d\lambda_1(s)
 \le\frac12\|u_0\|_2^2\qquad (0\le t\le T).
\end{equation}
Wymagamy również silnej nierówności energii:
dla prawie każdego $s\in(0,T)$ i każdego $t\in[s,T]$
\begin{equation}\label{eq:pde28-energy-shift}
 \frac12\|u(t)\|_2^2+
 \nu\int_s^t\|\nabla u(\tau)\|_2^2\,\mathrm d\lambda_1(\tau)
 \le\frac12\|u(s)\|_2^2.
\end{equation}
Zbiór dopuszczalnych chwil $s$ można wybrać
niezależnie od $t$. Nierówność~\eqref{eq:pde28-energy}
odpowiada dodatkowo chwili $s=0$.
\end{definicja}

\begin{lemat}[Oszacowanie nieliniowości w przestrzeni dualnej]
\label{lem:pde28-B-est}
Dla $u,v\in V$ niech $B(u,v)\in V'$ będzie określony przez
$\langle B(u,v),\varphi\rangle=b(u,v,\varphi)$.
Dla $u\in V$ zachodzą oszacowania
\[
 \|B(u,u)\|_{V'}\le
 \begin{cases}
 C\|u\|_2\|\nabla u\|_2,&d=2,\\
 C\|u\|_2^{1/2}\|\nabla u\|_2^{3/2},&d=3.
 \end{cases}
\]
Stała $C$ zależy tylko od wymiaru i wielkości torusa.
\end{lemat}
\begin{proof}
Wykorzystujemy antysymetrię w dwóch ostatnich zmiennych:
$|b(u,u,\varphi)|=|b(u,\varphi,u)|
\le\|u\|_4^2\|\nabla\varphi\|_2$.
W wymiarze $2$ stosujemy nierówność Ladyżenskiej
na torusie:
$\|u\|_4^2\le C\|u\|_2\|u\|_{H^1}$,
a następnie równoważność normy $H^1$ i normy $V$.
W wymiarze $3$ interpolujemy między $L^2$ i $L^6$:
$\|u\|_4\le\|u\|_2^{1/4}\|u\|_6^{3/4}$, ponieważ
$1/4=(1/4)/2+(3/4)/6$.
Z zanurzenia Sobolewa $H^1\hookrightarrow L^6$
i nierówności Poincarégo mamy
$\|u\|_6\le C\|\nabla u\|_2$.
Po podniesieniu nierówności do kwadratu
otrzymujemy obie żądane potęgi.
Po podzieleniu przez $\|\varphi\|_V=\|\nabla\varphi\|_2$
dla $\varphi\ne0$ i wzięciu supremum otrzymujemy tezę.
\end{proof}

\begin{uwaga}[Pochodna czasowa rozwiązania słabego]
Niech $\mathcal A:V\to V'$ będzie operatorem
określonym przez
$\langle\mathcal Au,\varphi\rangle=(u,\varphi)_V$.
Dla każdego rozwiązania z Definicji~\ref{def:pde28-leray}
mamy $\mathcal Au\in L^2(0,T;V')$, a z Lematu
~\ref{lem:pde28-B-est} wynika $B(u,u)\in L^q(0,T;V')$,
gdzie $q=2$ dla $d=2$ i $q=4/3$ dla $d=3$.
Silna mierzalność wynika z ciągłości odwzorowań
$\mathcal A$ i $u\mapsto B(u,u)$ na $V$.
Różniczkując~\eqref{eq:pde28-weak} w sensie
dystrybucyjnym względem czasu dla każdego $\varphi\in V$,
otrzymujemy w $V'$
\[
 u_t=-\nu\mathcal Au-B(u,u)\in L^q(0,T;V').
\]
Własność ta dotyczy wszystkich rozwiązań z definicji,
nie tylko skonstruowanych metodą Galerkina.
\end{uwaga}

\section{Globalne istnienie rozwiązania Leraya--Hopfa}
\begin{twierdzenie}[Istnienie słabego rozwiązania w wymiarach $2$ i $3$]
\label{thm:pde28-existence}
Niech $d\in\{2,3\}$, $\nu>0$, $T>0$ i $u_0\in H$.
Istnieje rozwiązanie Leraya--Hopfa zagadnienia~\eqref{eq:pde28-ns}
na $[0,T]$. Każde takie rozwiązanie spełnia
$\partial_tu\in L^{4/3}(0,T;V')$ dla $d=3$,
a $\partial_tu\in L^2(0,T;V')$ dla $d=2$.
Ponadto $u(t)\to u_0$ silnie w $H$ przy $t\downarrow0$.
\end{twierdzenie}
\begin{proof}
\textbf{Krok 1: układy Galerkina.}
Niech $(w_j)_{j\ge1}$ będzie ortonormalną bazą $H$
złożoną z gładkich, rzeczywistych, bezdywergencyjnych
funkcji własnych operatora Stokesa, z uwzględnieniem krotności.
Oznaczmy $H_m=\operatorname{span}\{w_1,\dots,w_m\}$
i niech $P_m$ będzie rzutem ortogonalnym na $H_m$.
Szukamy $u_m:[0,T]\to H_m$ o współrzędnych
spełniających dla $1\le j\le m$
\begin{equation}\label{eq:pde28-galerkin}
 (u_m',w_j)_H+\nu(\nabla u_m,\nabla w_j)_{L^2}
 +b(u_m,u_m,w_j)=0,\quad u_m(0)=P_m u_0.
\end{equation}
Pisząc $u_m=\sum_{i=1}^m a_iw_i$, dostajemy
układ współrzędnych
\[
 a_j'=-\nu\lambda_j a_j-
 \sum_{i,\ell=1}^m a_i a_\ell b(w_i,w_\ell,w_j),
 \qquad a_j(0)=(u_0,w_j)_H.
\]
Prawa strona jest wielomianem w $(a_1,\ldots,a_m)$,
a więc jest lokalnie lipschitzowska. Twierdzenie
Picarda--Lindelöfa daje jednoznaczne rozwiązanie
na maksymalnym przedziale istnienia.

\textbf{Krok 2: energia i istnienie globalne przybliżeń.}
Mnożymy $j$-te równanie przez współrzędną $u_m$ przy $w_j$
i sumujemy. Wykorzystując $b(u_m,u_m,u_m)=0$, otrzymujemy
\begin{equation}\label{eq:pde28-gal-energy}
 \frac12\|u_m(t)\|_2^2+\nu\int_0^t\|\nabla u_m(s)\|_2^2
 \,\mathrm d\lambda_1(s)=\frac12\|P_m u_0\|_2^2\le\frac12\|u_0\|_2^2.
\end{equation}
Norma $u_m$ nie może eksplodować w skończonym czasie,
a skończony układ ODE można przedłużyć na dowolny $[0,T]$.
Oszacowania są jednostajne względem $m$.

\textbf{Krok 3: pochodne czasowe.}
Rozszerzamy $P_m$ z $H$ na $V'$ przez
$\langle P_mF,\varphi\rangle=\langle F,P_m\varphi\rangle$.
Rozszerzenie jest zgodne z rzutem na $H$, ponieważ
rzut jest samosprzężony względem iloczynu skalarnego $H$.
Rzut $P_m$ jest kontrakcją na $V$ i $V'$:
bazę wybrano z funkcji własnych operatora Stokesa,
które są ortogonalne również dla iloczynu skalarnego
$(\nabla v,\nabla w)_{L^2}$. Wobec tego
$\|P_m\varphi\|_V\le\|\varphi\|_V$;
następnie z definicji normy dualnej
$\|P_m F\|_{V'}=\sup_{\|\varphi\|_V\le1}
|\langle F,P_m\varphi\rangle|\le\|F\|_{V'}$.
Z równania Galerkina oraz lematu~\ref{lem:pde28-B-est} mamy
\[
 \|u_m'\|_{V'}\le\nu\|\nabla u_m\|_2+
 \begin{cases}
 C\|u_m\|_2\|\nabla u_m\|_2,&d=2,\\
 C\|u_m\|_2^{1/2}\|\nabla u_m\|_2^{3/2},&d=3.
 \end{cases}
\]
Prawa strona jest ograniczona odpowiednio w
$L^2(0,T)$ i $L^{4/3}(0,T)$; dla pierwszego składnika
w drugim przypadku korzystamy z $T<\infty$.

\textbf{Krok 4: zwartość w przestrzeni i czasie.}
Z oszacowań i refleksywności odpowiednich
przestrzeni $L^2(0,T;V)$ i $L^q(0,T;V')$
oraz twierdzenia Banacha--Alaoglu (z ośrodkowością
predualu $L^1(0,T;H)$) wydzielamy podciąg taki, że
$u_m\stackrel{*}{\rightharpoonup}u$ w $L^\infty(0,T;H)$,
$u_m\rightharpoonup u$ w $L^2(0,T;V)$ i
$u_m'\rightharpoonup g$ w $L^q(0,T;V')$.
Obie granice $u_m$ są równe przez testowanie
funkcjami czasu i wektorami z $H$.
Testowanie pochodnych przez $\eta(t)\varphi$,
$\eta\in C_c^\infty(0,T)$, $\varphi\in V$,
i całkowanie przez części w czasie daje $g=u_t$
w sensie dystrybucyjnym.

Uzasadnimy zwartość czasoprzestrzenną wprost,
modyfikując argument z Twierdzenia
~\ref{thm:pde20-compact} dla wykładnika $q>1$,
w szczególności $q=4/3$.
Dla $j$ ustalonego współczynniki
$c_{m,j}=(u_m,w_j)_H$ są ograniczone jednostajnie
w czasie i spełniają
\[
 |c_{m,j}(t)-c_{m,j}(s)|
 \le \|u_m'\|_{L^q(0,T;V')}\|w_j\|_V
 |t-s|^{1-1/q}.
\]
Twierdzenie Arzeli--Ascolego i wybór podciągu
diagonalnego zapewniają jednostajną zbieżność
każdego współczynnika na $[0,T]$.
Dla ustalonego $N$ rzuty $P_Nu_m$ zbiegają więc
w $L^2(0,T;H)$. Jednocześnie
\[
 \sup_m\|(I-P_N)u_m\|_{L^2(0,T;H)}^2
 \le\lambda_{N+1}^{-1}
 \sup_m\|u_m\|_{L^2(0,T;V)}^2\longrightarrow0.
\]
Dla dowolnej dokładności wybieramy najpierw $N$,
by oba ogony były małe, a następnie indeksy
podciągu, by różnica rzutów była mała.
Podciąg jest Cauchy'ego w $L^2(0,T;H)$ i jego
granica musi być równa $u$ z poprzednich zbieżności.
Otrzymujemy zatem
\begin{equation}\label{eq:pde28-strong}
 u_m\longrightarrow u\quad\text{silnie w }L^2(0,T;H).
\end{equation}
W szczególności $u_m\otimes u_m\to u\otimes u$ w
$L^1((0,T)\times\T^d)$, gdyż
$\|a\otimes a-b\otimes b\|_1
\le(\|a\|_2+\|b\|_2)\|a-b\|_2$ na czasoprzestrzeni.

\textbf{Krok 5: przejście do granicy w równaniu.}
Niech $\varphi\in\bigcup_{j\ge1}H_j$ i
$\eta\in C_c^\infty(0,T)$. Testujemy~\eqref{eq:pde28-galerkin}
przez $\eta(t)\varphi$ i całkujemy w czasie.
Człon liniowy przechodzi do granicy słabo, a człon nieliniowy
po całkowaniu przez części przestrzennym ma postać
\[
 -\int_0^T\eta(t)\int_{\T^d}
 (u_m\otimes u_m):\nabla\varphi\,\mathrm d\lambda_d(x)
 \,\mathrm d\lambda_1(t).
\]
Ponieważ $\eta\nabla\varphi$ jest ograniczona,
silna zbieżność iloczynów w $L^1$ pozwala przejść
do granicy. Dostajemy
\[
 \int_0^T\bigl(- (u,\varphi)_H\eta'
  +\nu(u,\varphi)_V\eta+b(u,u,\varphi)\eta\bigr)
 \,\mathrm d\lambda_1(t)=0.
\]
Gęstość sumy $H_j$ w $V$ i oszacowanie
$B(u,u)\in L^q(0,T;V')$ wynikające z Lematu
~\ref{lem:pde28-B-est} rozszerzają tę tożsamość
na wszystkie $\varphi\in V$ (każdy ze składników
jest ciągłym funkcjonałem względem $\varphi$
w normie $V$). W kroku następnym uzasadnimy
postać całkową z warunkiem początkowym.

\textbf{Krok 6: warunek początkowy i nierówność energii.}
Dla każdego $\varphi\in V$ funkcje
$t\mapsto(u_m(t),\varphi)_H$ są jednostajnie ograniczone
i jednostajnie Hölderowskie w czasie: z kroku 3,
$|(u_m(t)-u_m(s),\varphi)|\le
C\|\varphi\|_V|t-s|^{1-1/q}$, gdzie $q=4/3$ lub $2$.
Na podciągu z kroku 4 dla każdego $j$
$c_{m,j}\to c_j$ jednostajnie na $[0,T]$.
Z oszacowania energii dla każdego $N$ i $t$
$\sum_{j=1}^N|c_j(t)|^2\le\|u_0\|_2^2$.
Szereg $\sum_j c_j(t)w_j$ określa więc wektor
$U(t)\in H$ o normie co najwyżej $\|u_0\|_2$.
Aproksymując dowolny $h\in H$ przez $P_Nh$,
otrzymujemy jednostajną w czasie zbieżność
$(u_m(t),h)_H\to(U(t),h)_H$: błąd ogona
jest nie większy niż $2\|u_0\|_2\|(I-P_N)h\|_2$.
W szczególności $U\in C([0,T];H_{\mathrm w})$
i $u_m(t)\rightharpoonup U(t)$ w $H$ dla każdego $t$.
Silna zbieżność w $L^2(0,T;H)$ identyfikuje $U=u$
prawie wszędzie; odtąd wybieramy reprezentant $U$.
Ponieważ $u_m(0)=P_m u_0\to u_0$ w $H$,
warunek początkowy jest zachowany.
Z kroku 5 dla każdego $\varphi\in V$
dystrybucyjna pochodna ciągłej funkcji
$t\mapsto(u(t),\varphi)_H$ należy do $L^q(0,T)$
i jest równa $-\nu(u,\varphi)_V-b(u,u,\varphi)$.
Jej reprezentant jest zatem absolutnie ciągły,
a całkowanie od $0$ do dowolnego $t$ daje
~\eqref{eq:pde28-weak}.
Dla ustalonego $t$ stosujemy słabą półciągłość dolną
normy $\|u_m(t)\|_2$ i normy
$\|\nabla u_m\|_{L^2(0,t;L^2)}$ do~\eqref{eq:pde28-gal-energy},
otrzymując~\eqref{eq:pde28-energy} dla każdego $t$.
Ze zbieżności~\eqref{eq:pde28-strong} wydzielamy
dalszy podciąg taki, że $u_m(s)\to u(s)$ w $H$
dla prawie każdego $s\in(0,T)$; na przykład
wybieramy podciąg z sumowalnymi kwadratami
norm błędów w $L^2(0,T;H)$ i stosujemy twierdzenie
Tonellego. Dla każdej takiej chwili $s$ odejmujemy
tożsamość energii Galerkina w chwili $s$ od
tożsamości w dowolnej chwili $t\ge s$.
Prawa strona $\frac12\|u_m(s)\|_2^2$ zbiega
do $\frac12\|u(s)\|_2^2$, natomiast po lewej
stosujemy słabą półciągłość dolną normy $u_m(t)$
oraz normy gradientu ograniczonego do $(s,t)$.
Otrzymujemy~\eqref{eq:pde28-energy-shift}.
Jeden podciąg zapewnia słabą zbieżność
we wszystkich chwilach $t$, więc zbiór
dopuszczalnych $s$ nie zależy od $t$.
Wreszcie $u(t)\rightharpoonup u_0$ w $H$ przy $t\downarrow0$,
a nierówność energii daje
$\limsup_{t\downarrow0}\|u(t)\|_2\le\|u_0\|_2$.
Słaba zbieżność wraz ze zbieżnością norm dowodzi
$u(t)\to u_0$ silnie w $H$.
\end{proof}


\begin{uwaga}[Zwartość paraboliczna a nieliniowość Naviera--Stokesa]
\label{rem:pde-v44-ns-compactness}
Zwartość z kroku czwartego dowodu twierdzenia
\ref{thm:pde28-existence} ma dokładnie tę samą konstrukcję
funkcjonalną co w rozdziale~\ref{chap:pde20}:
$V\hookrightarrow H$ zwarte i $H\hookrightarrow V'$ ciągłe,
$u_m$ ograniczone w $L^2(0,T;V)$ oraz $\partial_tu_m$
ograniczone w $L^q(0,T;V')$, gdzie $q=2$ dla $d=2$
i $q=4/3$ dla $d=3$. Sama słaba zbieżność w przestrzeni
energii nie wystarczałaby do przejścia do granicy
w $u_m\otimes u_m$. Silna zbieżność w $L^2(0,T;H)$
zapewnia bowiem bezpośrednio
\[
 \|u_m\otimes u_m-u\otimes u\|_{L^1((0,T)\times\T^d)}
 \le (\|u_m\|_{L^2}+\|u\|_{L^2})\|u_m-u\|_{L^2}\longrightarrow0.
\]
Nie wnioskujemy stąd o silnej zbieżności gradientów ani
o równości energii w trzech wymiarach.
\end{uwaga}


\begin{uwaga}[Ciśnienie rozwiązania słabego]
Dla gładkiego rozwiązania warunek $\diver u=0$
i równanie~\eqref{eq:pde28-ns} dają
$-\Delta p=\sum_{i,j}\partial_i\partial_j(u_i u_j)$.
Na torusie wybieramy normalizację
$\int_{\T^d}p(t,x)\,\mathrm d\lambda_d(x)=0$.
Dla rozwiązania Leraya--Hopfa wzór ten określa $p$
w sensie dystrybucji przestrzennych i czasowych;
nie należy bez dodatkowych założeń traktować go
jako tożsamości klasycznych funkcji. Dla ustalonego czasu w sensie dystrybucji
przestrzennych rozwiązujemy równanie Poissona dla
$p$ o średniej zero; w bazie Fouriera, dla $k\ne0$,
$\widehat p(k)=-|k|^{-2}\sum_{i,j}k_i k_j
\widehat{(u_i u_j)}(k)$.
Współczynniki iloczynów $u_i u_j$ są dobrze określone,
gdy $u(t)\in L^2$. Ustalamy $\widehat p(0)=0$.
Dla $k\ne0$ możemy zapisać wzór jako
$\widehat p(t,k)=-(2\pi)^{-d}\int_{\T^d}
 |k\cdot u(t,x)|^2|k|^{-2} e^{-ik\cdot x}
 \,\mathrm d\lambda_d(x)$.
W szczególności $|\widehat p(t,k)|\le
(2\pi)^{-d}\|u(t)\|_2^2$ jednostajnie w $k$ i $t$.
Współczynniki są mierzalne względem czasu,
a współczynniki gładkich funkcji testowych
maleją szybciej niż dowolna potęga $|k|$,
jednostajnie na zwartych przedziałach czasu.
Szereg definiuje więc dystrybucję na
$(0,T)\times\T^d$, o zerowym współczynniku dla $k=0$.
Bez normalizacji ciśnienie jest wyznaczone
z dokładnością do dystrybucji zależnej tylko od czasu.

Pozostaje wykazać pełne równanie pędu.
Niech $R=u_t-\nu\Delta u+\diver(u\otimes u)$
w sensie dystrybucji. Dla $k\ne0$ słaba postać
~\eqref{eq:pde28-weak}, testowana sinusami i cosinusami
z amplitudami w $k^\perp$, daje $P_k\widehat R(k)=0$,
gdzie $P_k=I-k\otimes k/|k|^2$.
Z warunku nieściśliwości również
$P_k\widehat u_t(k)=\widehat u_t(k)$ i
$P_k\widehat{\Delta u}(k)=\widehat{\Delta u}(k)$.
Wobec tego
$\widehat R(k)=(I-P_k)\widehat{\diver(u\otimes u)}(k)$.
Wzór na $\widehat p(k)$ daje dokładnie
$ik\widehat p(k)=-\widehat R(k)$.
Dla $k=0$ oba współczynniki znikają na mocy
zerowej średniej $u$ i periodyczności.
Warto zauważyć dodatkową regularność ciśnienia:
z oszacowań $L^4$ w Lemacie~\ref{lem:pde28-B-est}
wynika $u\otimes u\in L^q(0,T;L^2)$
z $q=2$ dla $d=2$ i $q=4/3$ dla $d=3$.
Dla macierzy $G=(u_i u_j)_{i,j}$ mnożnik
$-\sum_{i,j}k_i k_j\widehat G_{ij}(k)/|k|^2$
ma normę co najwyżej $1$ względem normy
Frobeniusa, przez nierówność Cauchy'ego--Schwarza.
Twierdzenie Parsevala daje więc
$\|p(t)\|_2\le\|u(t)\otimes u(t)\|_2=\|u(t)\|_4^2$
dla prawie każdego $t$.
Mnożnik jest ograniczonym operatorem liniowym,
stąd $p\in L^q(0,T;L^2)$ i ma średnią zero.
Zatem $R+\nabla p=0$ w sensie dystrybucji;
otrzymaliśmy pełne równanie~\eqref{eq:pde28-ns}.
\end{uwaga}


\begin{twierdzenie}[Równość energii i ciągłość silna w wymiarze dwa]
\label{thm:pde28-energy-equality-2d}
Niech $\nu>0$, $T>0$, $d=2$ i $u_0\in H$.
Każde rozwiązanie Leraya--Hopfa $u$ na $[0,T]$
należy do $C([0,T];H)$ i spełnia dla wszystkich
$t\in[0,T]$ dokładną równość
\begin{equation}\label{eq:pde28-energy-equality-2d}
 \frac12\|u(t)\|_2^2+
 \nu\int_0^t\|\nabla u(s)\|_2^2\,\mathrm d\lambda_1(s)
 =\frac12\|u_0\|_2^2.
\end{equation}
Ponadto, ponieważ na torusie $2\pi$ funkcje w $H$
mają zerową średnią,
$\|u(t)\|_2^2\le e^{-2\nu t}\|u_0\|_2^2$.
\end{twierdzenie}
\begin{proof}
Z lematu~\ref{lem:pde28-B-est} i warunków
$u\in L^\infty(0,T;H)\cap L^2(0,T;V)$ wynika
$\|B(u,u)\|_{V'}\le C\|u\|_2\|\nabla u\|_2\in
L^2(0,T)$. Niech $\mathcal A:V\to V'$ oznacza rozszerzenie
operatora Stokesa określone przez
$\langle\mathcal Au,\varphi\rangle=
(\nabla u,\nabla\varphi)_{L^2}$.
Wówczas $\|\mathcal Au\|_{V'}\le\|u\|_V$,
a równość słaba daje
$u_t=-\nu\mathcal Au-B(u,u)\in L^2(0,T;V')$.
W trójce Hilberta $V\subset H\subset V'$ mamy
zatem $u\in L^2(0,T;V)$ oraz $u_t\in L^2(0,T;V')$.
Argument z Lematu~\ref{lem:pde20-energy} przenosi się
na obecną trójkę: układ $(w_j)$ jest ortonormalny
w $H$, a $(\lambda_j^{-1/2}w_j)$ w $V$.
W przestrzeni $V'$ mamy, na mocy twierdzenia Riesza,
$\|F\|_{V'}^2=\sum_j\lambda_j^{-1}
|\langle F,w_j\rangle|^2$. Rzuty $P_m$ są więc
kontrakcjami i zbiegają silnie do identyczności
w każdej z trzech przestrzeni.
Zbieżność zmajoryzowana daje $P_mu\to u$ w $L^2(V)$
oraz $P_mu_t\to u_t$ w $L^2(V')$, gdzie pomijamy
w skrótach wspólny przedział czasu $(0,T)$.
Dla $z=P_mu-P_nu$ skończenie wymiarowa reguła
łańcuchowa i wybór chwili o normie nie większej
od średniej kwadratowej dają
\[
 \sup_{0\le t\le T}\|z(t)\|_2^2
 \le T^{-1}\|z\|_{L^2(H)}^2+
 2\|z_t\|_{L^2(V')}\|z\|_{L^2(V)}.
\]
Prawa strona dąży do zera, więc rzuty zbiegają
w $C([0,T];H)$ do reprezentanta $u$.
Jest to ten sam reprezentant co w definicji
rozwiązania Leraya--Hopfa: oba są słabo ciągłe
i równe prawie wszędzie.
Przechodząc do granicy w całkowej regule
łańcuchowej dla $P_mu$ za pomocą nierówności
Höldera w czasie, otrzymujemy absolutną ciągłość
$t\mapsto\|u(t)\|_2^2$ i
$\frac12\frac{\mathrm d}{\mathrm dt}\|u(t)\|_2^2
=\langle u_t(t),u(t)\rangle$ prawie wszędzie.
Równanie zachodzi w $V'$ dla prawie każdego $t$:
korzystamy z tożsamości dystrybucyjnej i ośrodkowości
$V$, aby wybrać wspólny zbiór pełnej miary dla
wszystkich testów. W tych chwilach, w których
również $u(t)\in V$, parujemy równanie z $u(t)$.
Dla takich chwil $b(u,u,u)=0$ na mocy
antysymetrii i gęstości gładkich pól bezdywergencyjnych.
Otrzymujemy pochodną
$\frac12(\|u\|_2^2)'=-\nu\|\nabla u\|_2^2$;
całkowanie w czasie daje~\eqref{eq:pde28-energy-equality-2d}.

Na torusie $\T^2$ najmniejsza dodatnia wartość
własna $-\Delta$ wynosi $1$, więc nierówność
Poincarégo ma postać $\|u\|_2\le\|\nabla u\|_2$.
Dla $E(t)=\|u(t)\|_2^2$ otrzymujemy
$E'(t)\le-2\nu E(t)$ prawie wszędzie.
Całkowanie tej nierówności daje wskazany zanik.
\end{proof}

\begin{wniosek}[Wykładniczy zanik energii rozwiązania Galerkina w 3D]
\label{cor:pde28-decay-3d}
Niech $\nu>0$, $T>0$, $u_0\in H$ i $d=3$.
Rozwiązanie Leraya--Hopfa otrzymane w dowodzie
twierdzenia~\ref{thm:pde28-existence} można wybrać tak, że
\[
 \|u(t)\|_2^2\le e^{-2\nu t}\|u_0\|_2^2,
 \qquad 0\le t\le T.
\]
Nie jest tu zakładana gładkość ani jednoznaczność
rozwiązań trójwymiarowych.
\end{wniosek}
\begin{proof}
Dla przybliżeń Galerkina z~\eqref{eq:pde28-gal-energy}
różniczkujemy energię i używamy
$\|\nabla u_m\|_2^2\ge\|u_m\|_2^2$:
$\frac{\mathrm d}{\mathrm dt}\|u_m(t)\|_2^2
\le-2\nu\|u_m(t)\|_2^2$.
Stąd $\|u_m(t)\|_2^2\le
 e^{-2\nu t}\|P_mu_0\|_2^2\le e^{-2\nu t}\|u_0\|_2^2$.
W kroku 6 dowodu~\ref{thm:pde28-existence} wybrano
jeden podciąg zbieżny słabo w $H$ dla każdej chwili
$t\in[0,T]$. Słaba półciągłość dolna normy przy
ustalonym $t$ daje pożądane oszacowanie dla tej granicy.
W szczególności w tym wniosku nie korzystamy z
nieudowodnionej równości energii dla wszystkich
rozwiązań Leraya--Hopfa w wymiarze trzy.
\end{proof}


\section{Wymiar dwa: jednoznaczność i wirowość}
\begin{twierdzenie}[Jednoznaczność słabych rozwiązań w wymiarze dwa]
\label{thm:pde28-unique2d}
Niech $\nu>0$, $T>0$ i $u_0\in H$ dla $d=2$.
Rozwiązanie Leraya--Hopfa na $[0,T]$ jest jednoznaczne.
Jeśli $u,v$ odpowiadają danym $u_0,v_0\in H$, to
\[
 \|u(t)-v(t)\|_2^2\le\|u_0-v_0\|_2^2
 \exp\left(\frac C\nu\int_0^t\|\nabla v(s)\|_2^2\,\mathrm d\lambda_1(s)\right),
 \qquad 0\le t\le T.
\]
\end{twierdzenie}
\begin{proof}
Z oszacowania w Lemacie~\ref{lem:pde28-B-est}
i słabej postaci równań wynika
$u_t,v_t\in L^2(0,T;V')$ dla wszystkich rozwiązań
Leraya--Hopfa w wymiarze dwa.
Stosujemy regułę łańcuchową w trójce Hilberta,
uzasadnioną w dowodzie Twierdzenia
~\ref{thm:pde28-energy-equality-2d}, do $w=u-v$.
Odejmujemy równania w $V'$ i parujemy z $w(t)$
dla prawie każdego $t$.
Z trójliniowości mamy
$b(u,u,w)-b(v,v,w)=b(u,w,w)+b(w,v,w)$;
pierwszy składnik znika przez antysymetrię.
Otrzymujemy
\[
 \frac12\frac{\mathrm d}{\mathrm dt}\|w\|_2^2
 +\nu\|\nabla w\|_2^2=-b(w,v,w).
\]
Wymiar dwa pozwala użyć nierówności Ladyżenskiej:
$|b(w,v,w)|\le\|w\|_4^2\|\nabla v\|_2
\le C\|w\|_2\|\nabla w\|_2\|\nabla v\|_2$.
Nierówność Younga $ab\le\nu a^2/2+b^2/(2\nu)$
dla $a=\|\nabla w\|_2$,
$b=C\|w\|_2\|\nabla v\|_2$ daje
\[
 \frac{\mathrm d}{\mathrm dt}\|w\|_2^2
 +\nu\|\nabla w\|_2^2
 \le (C^2/\nu)\|\nabla v\|_2^2\|w\|_2^2.
\]
Po pominięciu składnika nieujemnego i zmianie
stałej otrzymujemy żądaną nierówność różniczkową.
Ponieważ $\nabla v\in L^2(0,T;L^2)$,
lemat Gronwalla daje podane oszacowanie, a dla równych
danych otrzymujemy $u=v$.
\end{proof}

\begin{stwierdzenie}[Równanie wirowości i zanik enstrofii]
\label{prop:pde28-vorticity}
Niech $d=2$, a $(u,p)$ będzie gładkim rozwiązaniem
\eqref{eq:pde28-ns} na $[0,T]\times\T^2$.
Przyjmijmy $\omega=\partial_1u_2-\partial_2u_1$
oraz $\mathcal E_\omega(t)=\frac12\|\omega(t)\|_2^2$.
Wówczas
\[
 \partial_t\omega+u\cdot\nabla\omega=\nu\Delta\omega,
 \qquad \mathcal E_\omega(t)+\nu\int_0^t\|\nabla\omega(s)\|_2^2
 \,\mathrm d\lambda_1(s)=\mathcal E_\omega(0).
\]
Dodatkowo dla bezdywergencyjnych pól o zerowej średniej
$\|\nabla u(t)\|_2=\|\omega(t)\|_2$.
\end{stwierdzenie}
\begin{proof}
Stosując operator $\partial_1$ do drugiego równania
wektorowego i odejmując $\partial_2$ pierwszego,
eliminujemy gradient ciśnienia, ponieważ pochodne
mieszane są równe. Dla członu konwekcyjnego
stosujemy regułę Leibniza:
\[
 \begin{aligned}
 &\partial_1((u\cdot\nabla)u_2)
 -\partial_2((u\cdot\nabla)u_1)\\
 &\quad=u\cdot\nabla\omega+
 (\partial_1u_1+\partial_2u_2)
 (\partial_1u_2-\partial_2u_1).
 \end{aligned}
\]
Ostatni iloczyn znika na mocy nieściśliwości.
Operator rotacji komutuje z laplasjanem,
co daje równanie wirowości.
Mnożymy równanie przez $\omega$ i całkujemy po torusie.
Składnik konwekcyjny znika:
$\int_{\T^2}u\cdot\nabla(\omega^2/2)\,\mathrm d\lambda_2(x)=0$.
Człon dyfuzyjny daje $-\nu\|\nabla\omega\|_2^2$
po periodycznym całkowaniu przez części.
Całkowanie tożsamości różniczkowej w czasie daje
równość enstrofii.
Równość norm można wykazać na współczynnikach
Fouriera: dla $k\ne0$
\[
 |k|^2|\widehat u(k)|^2
 =|k\cdot\widehat u(k)|^2+
 |k_1\widehat u_2(k)-k_2\widehat u_1(k)|^2.
\]
Pierwszy składnik po prawej znika, a drugi to
$|\widehat\omega(k)|^2$.
Sumowanie i twierdzenie Parsevala dają
$\|\nabla u\|_2=\|\omega\|_2$.
Po pomnożeniu równości przez $|k|^2$ otrzymujemy
także $\|\Delta u\|_2=\|\nabla\omega\|_2$
dla pól w $H^2$.
\end{proof}

\begin{twierdzenie}[Globalne oszacowanie $H^1$ w wymiarze dwa]
\label{thm:pde28-H1-2d}
Niech $d=2$, $\nu>0$, $T>0$ oraz $u_0\in V$.
Jedyne rozwiązanie Leraya--Hopfa $u$ na $[0,T]$ spełnia
$u\in C([0,T];V)\cap L^2(0,T;D(A))$
i $u_t\in L^2(0,T;H)$
oraz dla każdego $t\in[0,T]$
\[
 \|\nabla u(t)\|_2^2+2\nu\int_0^t\|\Delta u(s)\|_2^2
 \,\mathrm d\lambda_1(s)\le\|\nabla u_0\|_2^2.
\]
\end{twierdzenie}
\begin{proof}
Wybieramy w dowodzie twierdzenia~\ref{thm:pde28-existence}
przestrzenie Galerkina złożone ze wszystkich bezdywergencyjnych
modów Fouriera $0<|k|\le N$ (są to podprzestrzenie
odpowiadające pełnym grupom wartości własnych).
Oznaczmy przez $Q_N$ obcięcie do tych częstości,
działające składowo na polach wektorowych.
Rzut Galerkina jest równy $Q_NP$ dla pól o średniej zero.
Równanie przybliżone ma postać
$\partial_t u_N-\nu\Delta u_N+
Q_NP((u_N\cdot\nabla)u_N)=0$.
Rotacja komutuje z $Q_N$ i $\Delta$, a usuwa
część gradientową $I-P$ (na każdej częstości
jest ona równoległa do $k$).
Dla $\omega_N=\partial_1u_{N,2}
-\partial_2u_{N,1}$ z obliczenia w Stwierdzeniu
~\ref{prop:pde28-vorticity} równanie rzutowane ma postać
$\partial_t\omega_N+Q_N(u_N\cdot\nabla\omega_N)
=\nu\Delta\omega_N$, gdzie $Q_N$ jest ortogonalnym
rzutem na skalarne mody z $0<|k|\le N$.
Ponieważ $Q_N\omega_N=\omega_N$, mnożenie przez
$\omega_N$ daje
\[
 \frac12\frac{\mathrm d}{\mathrm dt}\|\omega_N\|_2^2
 +\nu\|\nabla\omega_N\|_2^2=0.
\]
Przy tym człon transportowy znika przez $\diver u_N=0$.
Dla bezdywergencyjnego pola o zerowej średniej zachodzą
$\|\omega_N\|_2=\|\nabla u_N\|_2$ i
$\|\nabla\omega_N\|_2=\|\Delta u_N\|_2$,
co można sprawdzić oddzielnie na każdym modzie Fouriera.
Otrzymujemy zatem jednostajne ograniczenia w
$L^\infty(0,T;V)$ i $L^2(0,T;H^2)$.
Dokładniej, $\|\nabla u_N(0)\|_2\le\|\nabla u_0\|_2$,
a $\|\Delta u_N\|_2$ kontroluje normę $H^2$
dzięki zerowej średniej i oszacowaniu
$(1+|k|^2)^2\le4|k|^4$ dla $k\ne0$.
Wydzielamy podciąg słabo zbieżny w $L^2(H^2)$
i słabo-* zbieżny w $L^\infty(V)$.
Obie granice są równe granicy z dowodu istnienia
dzięki silnej zbieżności w $L^2(H)$.
Dla każdej chwili $t$ zachodzi
$u_N(t)\rightharpoonup u(t)$ w $V$:
istotnie, ciąg jest ograniczony w $V$,
a jego słaba granica w $H$ została już ustalona
dla każdej chwili; każda słabo zbieżna
w $V$ podsekwencja ma zatem tę samą granicę.
W szczególności $u(t)\in V$ dla każdego $t$.
Stosując słabą półciągłość dolną normy $V$
oraz normy $\Delta u_N$ w $L^2(0,t;L^2)$
do tożsamości enstrofii, otrzymujemy żądane
oszacowanie nawet dla wszystkich $t\in[0,T]$.
Na mocy jednoznaczności z Twierdzenia
~\ref{thm:pde28-unique2d} jest to jedyne
rozwiązanie Leraya--Hopfa.

Możemy również uzasadnić silną ciągłość w $V$.
Dla prawie każdego $t$ pole $u(t)\in D(A)$,
a nierówność Ladyżenskiej zastosowana do
pochodnych pierwszego rzędu (które mają średnią zero)
daje
\[
 \|\nabla u\|_4^2\le C\|\nabla u\|_2\|\Delta u\|_2.
\]
Łącząc to z $\|u\|_4^2\le C\|u\|_2\|\nabla u\|_2$,
otrzymujemy
\[
 \|P((u\cdot\nabla)u)\|_2^2
 \le C\|u\|_2\|\nabla u\|_2^2\|\Delta u\|_2.
\]
Prawa strona jest całkowalna po czasie,
gdyż $u$ jest ograniczona w $V$ i
$\Delta u\in L^2(0,T;H)$. Z równania rzutowanego
wynika więc $u_t\in L^2(0,T;H)$.
Niech $A^{1/2}$ będzie mnożnikiem
$\lambda_j^{1/2}$ w bazie Stokesa.
Dla $z=A^{1/2}u$ mamy $z\in L^2(0,T;V)$,
gdyż $\|z\|_V^2=\|Au\|_2^2$, oraz
$z_t=A^{1/2}u_t\in L^2(0,T;V')$, ponieważ
$\|A^{1/2}h\|_{V'}=\|h\|_2$ dla $h\in H$.
Równość pochodnych sprawdzamy na skończonych
rzutach spektralnych, a następnie przechodzimy
do granicy w obu przestrzeniach Bochnera.
Reguła łańcuchowa z dowodu Twierdzenia
~\ref{thm:pde28-energy-equality-2d} daje
$z\in C([0,T];H)$. Ponieważ $A^{-1/2}:H\to V$
jest izometrią, otrzymujemy $u\in C([0,T];V)$;
reprezentant ten jest zgodny z wcześniejszym
ciągłym reprezentantem w $H$.
\end{proof}

\begin{uwaga}[Dlaczego wymiar trzy jest inny]
W wymiarze $3$ wirowość $\boldsymbol\omega=\nabla\times u$
spełnia dla gładkiego rozwiązania
$\partial_t\boldsymbol\omega+(u\cdot\nabla)\boldsymbol\omega
=(\boldsymbol\omega\cdot\nabla)u+\nu\Delta\boldsymbol\omega$.
Pojawia się człon \emph{rozciągania wirów}, który nie ma
odpowiednika w wymiarze $2$. Z tożsamości energii $L^2$
nie otrzymuje się wprost ograniczenia normy $\|\nabla u\|_2$
we wszystkich czasach w wymiarze $3$.
\end{uwaga}

\section{Wymiar trzy: kryterium warunkowej jednoznaczności}
\begin{twierdzenie}[Jednoznaczność przy dodatkowej całkowalności]
\label{thm:pde28-serrin}
Niech $d=3$, $\nu>0$, $T>0$ i $u_0\in H$.
Niech $u,v$ będą rozwiązaniami Leraya--Hopfa na $[0,T]$
o tych samych danych początkowych. Załóżmy dodatkowo, że
$v\in L^\infty(0,T;V)\cap L^2(0,T;D(A))$,
$\partial_tv\in L^2(0,T;H)$ oraz
$v\in L^q(0,T;L^p(\T^3))$, gdzie $p>3$ i
$2/q+3/p\le1$. Wówczas $u=v$ prawie wszędzie na
$(0,T)\times\T^3$.
\end{twierdzenie}
\begin{proof}
Z Lematu~\ref{lem:pde28-B-est} i słabej postaci
równania wiemy, że $u_t\in L^{4/3}(0,T;V')$.
Ponieważ $v\in L^\infty(0,T;V)$ i
$v_t\in L^2(0,T;H)$, reguła łańcuchowa uzasadniona
w dowodzie Twierdzenia~\ref{thm:pde28-energy-equality-2d}
daje $v\in C([0,T];H)$ oraz równość energii dla $v$:
parujemy równanie dla $v$ z $v$ w $V'$ i stosujemy
$b(v,v,v)=0$. Nie stosujemy tej reguły bezpośrednio
do $u$, gdyż jego pochodna czasowa nie musi należeć
do $L^2(0,T;V')$.

Aby uzyskać tożsamość mieszaną, niech $P_n$ oznacza
rzut na pierwsze $n$ wektorów bazy Stokesa.
Współczynniki $P_nu$ i $P_nv$ są absolutnie ciągłe:
ich pochodne są odpowiednio współczynnikami $u_t$
i $v_t$. Z reguły Leibniza w skończonym wymiarze
otrzymujemy dla każdego $t\in[0,T]$
\[
 \begin{aligned}
 &(P_nu(t),P_nv(t))_H-(P_nu_0,P_nv(0))_H\\
 &\quad=\int_0^t
 \bigl(\langle u_t,P_nv\rangle+(P_nu,v_t)_H\bigr)
 \,\mathrm d\lambda_1(s).
 \end{aligned}
\]
Po lewej stronie przechodzimy do granicy w $H$
dla każdej ustalonej chwili, korzystając z
wybranych słabo ciągłych reprezentantów.
Po prawej stronie $P_nv\to v$ w $L^4(0,T;V)$
oraz $P_nu\to u$ w $L^2(0,T;H)$ na mocy
silnej zbieżności rzutów i zbieżności zmajoryzowanej.
Nierówności Höldera dają zbieżność obu całek,
jednostajną względem górnej granicy całkowania.
Stąd
\[
 (u(t),v(t))_H-(u_0,v(0))_H
 =\int_0^t
 (\langle u_t,v\rangle+(u,v_t)_H)\,\mathrm d\lambda_1(s).
\]
Oba równania zachodzą w $V'$ prawie wszędzie.
Wstawiając je do ostatniej tożsamości, dostajemy
\[
 \begin{aligned}
 (u(t),v(t))_H={}&(u_0,v(0))_H\\
 &-2\nu\int_0^t(u,v)_V\,\mathrm d\lambda_1(s)\\
 &-\int_0^t
 [b(u,u,v)+b(v,v,u)]\,\mathrm d\lambda_1(s).
 \end{aligned}
\]
Całkowalność pierwszego członu nieliniowego
wynika z $B(u,u)\in L^{4/3}(V')$ i $v\in L^4(V)$;
drugiego z $B(v,v)\in L^\infty(V')$
oraz $u\in L^2(V)$. W tych skrótach przedziałem
całkowania jest $(0,T)$.
Łącząc otrzymaną tożsamość mieszaną z nierównością
energii dla $u$ i równością energii dla $v$,
uzyskujemy dla $w=u-v$ nierówność
\[
 \frac12\|w(t)\|_2^2+\nu\int_0^t\|\nabla w\|_2^2\,\mathrm d\lambda_1(s)
 \le-\int_0^t b(w(s),v(s),w(s))\,\mathrm d\lambda_1(s).
\]
Istotnie, sumujemy nierówność energii dla $u$
i równość energii dla $v$, po czym odejmujemy
tożsamość mieszaną. Składniki lepkościowe łączą się
w $\nu\|\nabla(u-v)\|_2^2$, a składnik początkowy
jest równy $\frac12\|u_0-v(0)\|_2^2=0$.
Po prawej stronie pozostaje
$b(u,u,v)+b(v,v,u)=-b(w,v,w)$:
wstawiamy $u=v+w$ i stosujemy antysymetrię,
w szczególności $b(a,v,v)=0$ oraz
$b(w,w,v)=-b(w,v,w)$.
Otrzymujemy zatem nierówność względnej energii.

Niech najpierw $3<p<\infty$.
Z nierówności Höldera i interpolacji pomiędzy
$L^2$ i $L^6$, przy parametrze $\theta=3/p$
spełniającym $1/(2p/(p-2))=(1-\theta)/2+\theta/6$,
mamy
\[
 |b(w,w,v)|\le\|v\|_p\|w\|_{2p/(p-2)}\|\nabla w\|_2
 \le C\|v\|_p\|w\|_2^{1-3/p}
 \|\nabla w\|_2^{1+3/p}.
\]
Stosujemy nierówność Younga z wykładnikami
$2/(1+\theta)$ i $2/(1-\theta)$,
dobierając współczynnik przy pierwszym składniku
tak, aby pochłonąć połowę dyssypacji.
Dla $r=2/(1-\theta)=2p/(p-3)$ dostajemy
\[
 |b(w,w,v)|\le\frac{\nu}{2}\|\nabla w\|_2^2+
 C_{\nu,p}\|v\|_p^r\|w\|_2^2.
\]
Dla $p=\infty$ to samo oszacowanie obowiązuje
z $r=2$: bez interpolacji stosujemy
$|b(w,w,v)|\le\|v\|_\infty\|w\|_2\|\nabla w\|_2$
i nierówność Younga.
Warunek $2/q+3/p\le1$ zapewnia $q\ge r$,
więc $\|v\|_p^r$ jest całkowalna w czasie.
Z antysymetrii i nierówności względnej energii
wynika zatem
\[
 \|w(t)\|_2^2\le
 2C_{\nu,p}\int_0^t\|v(s)\|_p^r\|w(s)\|_2^2
 \,\mathrm d\lambda_1(s).
\]
Funkcja po lewej stronie jest mierzalna
i ograniczona. Całkowa postać lematu Gronwalla
daje $\|w(t)\|_2=0$ dla wszystkich $t\in[0,T]$.
\end{proof}

\begin{uwaga}[Co pozostaje otwarte]
Dla $d=3$ twierdzenie~\ref{thm:pde28-existence} gwarantuje
istnienie globalnych rozwiązań słabych, ale nie dowodzi
ich jednoznaczności ani gładkości dla dowolnych danych w $H$.
Dowód twierdzenia~\ref{thm:pde28-serrin} jest dowodem
jednoznaczności słabego rozwiązania w obecności rozwiązania
o silniejszej regularności. Przy podanych założeniach
warunek całkowalności Prodiego--Serrina jest w istocie
automatyczny dla pewnej pary wykładników:
$v\in L^\infty(0,T;V)$ i zanurzenie Sobolewa dają
$v\in L^\infty(0,T;L^6)$, a przy $q=\infty$ mamy $2/q+3/6=1/2<1$.
Zatem do jednoznaczności wystarczają tu same
podane założenia o silniejszej regularności.
Zachowaliśmy rachunek dla ogólnych wykładników
$p,q$, by wyjaśnić znaczenie warunku
$2/q+3/p\le1$. Usunięcie tych dodatkowych założeń
i pozostawienie jedynie warunku $v\in L^q_tL^p_x$
daje znacznie silniejsze twierdzenie Prodiego--Serrina,
którego pełnego dowodu tutaj nie podajemy.
Argument wymiaru dwa wykorzystuje całkowalność
$\|\nabla v\|_2^2$ i nierówność Ladyżenskiej, której
odpowiednik w wymiarze trzy daje zbyt wysoką potęgę
$\|\nabla v\|_2$. Przejście od warunkowej jednoznaczności
do pełnej wymagałoby wykazania dodatkowej regularności;
nie należy tego zakładać bez dowodu.
\end{uwaga}

\section{Lokalne oszacowanie silnych rozwiązań i skala problemu}
\begin{stwierdzenie}[Oszacowanie normy $H^1$ w wymiarze trzy]
\label{prop:pde28-H1}
Niech $d=3$, $\nu>0$ i $(u,p)$ będzie gładkim rozwiązaniem
\eqref{eq:pde28-ns} na $[0,T]\times\T^3$, o średniej $u$ równej zero.
Istnieje $C>0$, zależna tylko od torusa, taka że
\[
 \frac{\mathrm d}{\mathrm dt}\|\nabla u(t)\|_2^2
 +\nu\|\Delta u(t)\|_2^2
 \le C\nu^{-3}\|\nabla u(t)\|_2^6.
\]
\end{stwierdzenie}
\begin{proof}
Bierzemy iloczyn skalarny równania z $-\Delta u$.
Ponieważ $\diver\Delta u=\Delta\diver u=0$,
człon ciśnienia znika przez całkowanie przez części.
Człon czasowy daje
$(u_t,-\Delta u)_H=(\nabla u_t,\nabla u)_{L^2}
=\frac12(\|\nabla u\|_2^2)'$, a lepkościowy
$\nu\|\Delta u\|_2^2$. Stosując nierówność Höldera
z wykładnikami $6,3,2$ do członu konwekcyjnego,
otrzymujemy
\[
 \frac12\frac{\mathrm d}{\mathrm dt}\|\nabla u\|_2^2
 +\nu\|\Delta u\|_2^2
 \le\|u\|_6\|\nabla u\|_3\|\Delta u\|_2.
\]
Z zanurzenia Sobolewa i zerowej średniej mamy
$\|u\|_6\le C\|\nabla u\|_2$.
Każda pochodna przestrzenna $u$ również ma
średnią zero, więc $\|\nabla u\|_6\le C\|D^2u\|_2$.
W bazie Fouriera suma kwadratów wszystkich
pochodnych drugiego rzędu (z uporządkowanymi
indeksami) ma mnożnik
$\sum_{i,j}k_i^2 k_j^2=|k|^4$, stąd
$\|D^2u\|_2=\|\Delta u\|_2$.
Interpolacja $\|\nabla u\|_3\le
\|\nabla u\|_2^{1/2}\|\nabla u\|_6^{1/2}$
daje więc
$\|\nabla u\|_3\le C\|\nabla u\|_2^{1/2}
\|\Delta u\|_2^{1/2}$.
Stąd prawa strona jest ograniczona przez
$C\|\nabla u\|_2^{3/2}\|\Delta u\|_2^{3/2}$.
Nierówność Younga z wykładnikami $4/3$ i $4$
pozwala pochłonąć połowę składnika
$\nu\|\Delta u\|_2^2$ i pozostawia
$C\nu^{-3}\|\nabla u\|_2^6$; pomnożenie przez $2$
i zmianę stałej kończy dowód.
\end{proof}

\begin{uwaga}[Czas obowiązywania oszacowania a priori]
Jeśli $y(t)=\|\nabla u(t)\|_2^2$ oraz $y_0=y(0)>0$,
to Stwierdzenie~\ref{prop:pde28-H1} daje
$y'\le C\nu^{-3}y^3$. Z porównania ze skalarnym
równaniem $z'=C\nu^{-3}z^3$, $z(0)=y_0$,
wynika
\[
 y(t)\le
 \frac{y_0}{\sqrt{1-2C\nu^{-3}y_0^2t}},
 \qquad 0\le t<\frac{\nu^3}{2Cy_0^2},
\]
o ile gładkie rozwiązanie istnieje w rozważanej chwili.
Wzór na $z$ wynika z różniczkowania $z^{-2}$;
porównanie stosujemy na każdym zwartym przedziale
wspólnego istnienia, gdzie sześcian jest lipschitzowski
na wspólnym ograniczonym zakresie.
Dla $y_0=0$ zerowa średnia daje $u_0=0$,
a tożsamość energii wymusza $u=0$.
Powyższy rachunek jest oszacowaniem a priori,
a nie samodzielnym dowodem istnienia silnego
rozwiązania ani dowodem eksplozji w chwili,
w której mianownik oszacowania znika.
\end{uwaga}

\begin{kontrprzyklad}[Dlaczego sama energia $L^2$ nie wystarcza]
Funkcje $u_n(x)=\sin(nx_1)e_2$ na $\T^3$, gdzie
$e_2=(0,1,0)$, mają zerową dywergencję i średnią.
Ich normy $\|u_n\|_2$ są jednakowo ograniczone,
podczas gdy $\|\nabla u_n\|_2\to\infty$.
Wobec tego kontrola samej energii $L^2$ nie daje
jednostajnego oszacowania $H^1$ w pojedynczych chwilach.
Nie jest to kontrprzykład dla regularności rozwiązań
równań Naviera--Stokesa, lecz przykład braku
odpowiedniego wniosku funkcjonalno-analitycznego.
\end{kontrprzyklad}

\begin{uwaga}[Skalowanie na całej przestrzeni i na torusie]
Dla $\ell>0$ połóżmy
\[
 u_\ell(t,x)=\ell u(\ell^2t,\ell x),\qquad
 p_\ell(t,x)=\ell^2p(\ell^2t,\ell x).
\]
Piszemy $s=\ell^2t$, $y=\ell x$. Z reguły
łańcuchowej otrzymujemy
$\partial_t u_\ell=\ell^3\partial_s u$,
$\nabla_xu_\ell=\ell^2\nabla_yu$,
$\Delta_xu_\ell=\ell^3\Delta_yu$ oraz
$\nabla_xp_\ell=\ell^3\nabla_yp$,
gdzie prawe strony obliczamy w $(s,y)$.
W szczególności $(u_\ell\cdot\nabla_x)u_\ell
=\ell^3(u\cdot\nabla_y)u$,
a $\diver_x u_\ell=\ell^2\diver_y u$.
Dlatego przeskalowane pole spełnia równanie
Naviera--Stokesa z tą samą lepkością $\nu$.
Na $\R^d$ zmiana zmiennych daje
\[
 \|u_\ell(t)\|_{L^p(\R^d)}
 =\ell^{1-d/p}\|u(\ell^2t)\|_{L^p(\R^d)}
 \quad(1\le p\le\infty),
\]
gdzie $d/\infty=0$. Dla normy mieszanej na
odpowiednio przeskalowanym przedziale czasu
czynnik wynosi $\ell^{1-d/p-2/q}$, również
dla $q=\infty$ z konwencją $2/\infty=0$.
W wymiarze trzy warunek $2/q+3/p=1$
jest więc warunkiem niezmienniczości tej normy.
Norma $L^2$ skaluje się z czynnikiem
$\ell^{1-d/2}$: jest niezmiennicza w wymiarze dwa,
ale nie w wymiarze trzy.
Na torusie o okresie $2\pi$ powyższa transformacja
prowadzi naturalnie do torusa o okresie $2\pi/\ell$;
nie jest dowolnym ciągłym skalowaniem
zagadnienia na ustalonym torusie.
\end{uwaga}

\section{Dokładne rozwiązanie: wir Taylora--Greena}
\begin{przyklad}[Wir Taylora--Greena na torusie dwuwymiarowym]
\label{ex:pde28-TG}
Niech $\nu>0$ i $(x,y)\in\T^2$. Funkcje
\begin{align*}
 u(t,x,y)&=e^{-2\nu t}(\sin x\cos y,-\cos x\sin y),\\
 p(t,x,y)&=\tfrac14e^{-4\nu t}(\cos 2x+\cos 2y)
\end{align*}
rozwiązują~\eqref{eq:pde28-ns} z danymi
$u_0=(\sin x\cos y,-\cos x\sin y)$.
Oznaczmy $a=e^{-2\nu t}$. Bezpośrednio
$\partial_xu_1=a\cos x\cos y=-\partial_yu_2$,
$\Delta u=-2u$ i $u_t=-2\nu u$.
Rozpisując konwekcję, otrzymujemy
\[
 \begin{aligned}
 ((u\cdot\nabla)u)_1
 &=a^2\sin x\cos x(\cos^2y+\sin^2y),\\
 ((u\cdot\nabla)u)_2
 &=a^2\sin y\cos y(\sin^2x+\cos^2x).
 \end{aligned}
\]
Jest to $\frac12a^2(\sin2x,\sin2y)=-\nabla p$.
Średnie $u$ i $p$ są zerowe. Ponieważ
$\int_0^{2\pi}\sin^2x\,\mathrm d\lambda_1(x)
=\int_0^{2\pi}\cos^2x\,\mathrm d\lambda_1(x)=\pi$,
mamy $\|u\|_2^2=2\pi^2a^2$.
Z całkowania przez części
$\|\nabla u\|_2^2=(u,-\Delta u)_H=2\|u\|_2^2$,
a różniczkowanie daje $\omega=2a\sin x\sin y$.
Ponadto
\[
 \frac12\|u(t)\|_2^2=\pi^2e^{-4\nu t},\qquad
 \|\nabla u(t)\|_2^2=4\pi^2e^{-4\nu t},\qquad
 \omega(t,x,y)=2e^{-2\nu t}\sin x\sin y.
\]
Różniczkowanie energii daje
$-4\nu\pi^2e^{-4\nu t}=-\nu\|\nabla u(t)\|_2^2$.
Jawny wzór pozwala porównać obliczenia numeryczne
z rozwiązaniem dokładnym na dowolnej siatce
i badać błąd oraz rząd zbieżności schematu.

\begin{center}
\begin{tikzpicture}[x=1.5cm,y=1.5cm,>=Stealth,font=\footnotesize]
 \foreach \i in {0,...,15}{
  \foreach \j in {0,...,15}{
   \pgfmathsetmacro{\xx}{360*(\i+.5)/16}
   \pgfmathsetmacro{\yy}{360*(\j+.5)/16}
   \pgfmathsetmacro{\ww}{sin(\xx)*sin(\yy)}
   \pgfmathsetmacro{\opa}{.65*abs(\ww)}
   \ifdim\ww pt>0pt
    \fill[red,opacity=\opa] ({\i/4},{\j/4})
       rectangle ({(\i+1)/4},{(\j+1)/4});
   \else
    \fill[blue,opacity=\opa] ({\i/4},{\j/4})
       rectangle ({(\i+1)/4},{(\j+1)/4});
   \fi
  }
 }
 \foreach \i in {0,...,7}{
  \foreach \j in {0,...,7}{
   \pgfmathsetmacro{\xx}{360*(\i+.5)/8}
   \pgfmathsetmacro{\yy}{360*(\j+.5)/8}
   \draw[-{Stealth[length=2.5pt]},black!85]
   ({(\i+.5)/2},{(\j+.5)/2})
   --++({.23*sin(\xx)*cos(\yy)},{-.23*cos(\xx)*sin(\yy)});
  }
 }
 \draw[->] (0,0)--(4.3,0) node[right] {$x$};
 \draw[->] (0,0)--(0,4.3) node[above] {$y$};
 \draw[black!60] (0,0) rectangle (4,4);
 \foreach \p/\lab in {0/0,2/\pi,4/{2\pi}}{
  \node[below] at (\p,0) {$\lab$};
  \node[left] at (0,\p) {$\lab$};
 }
\end{tikzpicture}

{\small Rysunek: pole prędkości i wirowość w chwili $t=0$.
Czerwony kolor oznacza dodatnią, niebieski ujemną wirowość;
strzałki są obliczone z jawnego wzoru.}
\end{center}
\end{przyklad}


\begin{przyklad}[Trójwymiarowy wir Taylora--Greena: dane początkowe]
\label{ex:pde28-TG3D}
Na torusie $\T^3=([0,2\pi]/\sim)^3$ rozważmy
bezdywergencyjne pole początkowe
\[
 u_0(x,y,z)=(\sin x\cos y\cos z,\,-\cos x\sin y\cos z,\,0).
\]
Istotnie, $\partial_x u_{0,1}+\partial_y u_{0,2}=0$,
a $\partial_z u_{0,3}=0$. Wszystkie składowe
mają średnią zero, więc $u_0\in H\cap C^\infty$.
Jego wirowość jest \emph{wektorem}
\[
 \nabla\times u_0=(-\cos x\sin y\sin z,\,
 -\sin x\cos y\sin z,\,2\sin x\sin y\cos z).
\]
W szczególności trzecia składowa wirowości wynosi
$\omega_{0,z}=2\sin x\sin y\cos z$. Energia początkowa
względem miary Lebesgue'a na sześcianie okresowym wynosi
\[
 E(0)=\frac12\int_{[0,2\pi]^3}|u_0(x,y,z)|^2
 \,\mathrm d\lambda_3(x,y,z)=\pi^3.
\]
Istotnie, całka z kwadratu każdej z dwóch
niezerowych składowych jest równa $\pi^3$.
Ponadto $\Delta u_0=-3u_0$ i
$\|\nabla u_0\|_2^2=6\pi^3$.
Dla gładkiej ewolucji z tymi danymi
początkowa pochodna energii wynosi więc $-6\nu\pi^3$.

Ponieważ $u_{0,3}=0$, \emph{linie prądu pola początkowego}
leżą w płaszczyznach $z=\mathrm{const}$. Dla ustalonego
$z$ regularne linie poziomicowe funkcji prądu
$\psi_z(x,y)=\sin x\sin y\cos z$ są liniami prądu,
ponieważ $(u_{0,1},u_{0,2})=(\partial_y\psi_z,
-\partial_x\psi_z)$, skąd $u_0\cdot\nabla\psi_z=0$
przy ustalonym $z$. W punktach stagnacji
trajektorie są stałe; dla $\cos z=0$ całe pole zanika.

Dane początkowe nie są jawnym rozwiązaniem
trójwymiarowego zagadnienia w czasie.
W istocie bezpośrednio
\[
 (u_0\cdot\nabla)u_0
 =\tfrac12\cos^2z(\sin2x,\sin2y,0).
\]
Równanie Poissona na ciśnienie w chwili początkowej
ma prawą stronę
$(\cos2x+\cos2y)\cos^2z$.
Jedynym rozwiązaniem o średniej zero jest
\[
 p_0=\tfrac1{16}(\cos2x+\cos2y)(2+\cos2z),
\]
co sprawdzamy, stosując $-\Delta$ do obu składników
i używając $2\cos^2z=1+\cos2z$.
W konsekwencji dla gładkiego rozwiązania
$\partial_t u_3(0)= -\partial_zp_0
=\tfrac18(\cos2x+\cos2y)\sin2z$,
co nie znika tożsamościowo. Trzecia składowa
prędkości nie pozostaje zatem zerowa.

\begin{center}
\begin{tikzpicture}[x=.80cm,y=.80cm,>=Stealth,font=\footnotesize]
\foreach \shift/\az/\zlabel in
 {0/0/{0},5.25/90/{\pi/2},10.5/180/{\pi}} {
  \begin{scope}[shift={(\shift,0)}]
  \fill[gray!5] (0,0) rectangle (4.6,4.6);
  \foreach \i in {0,...,7} {
   \foreach \j in {0,...,7} {
    \pgfmathsetmacro{\xx}{360*(\i+.5)/8}
    \pgfmathsetmacro{\yy}{360*(\j+.5)/8}
    \pgfmathsetmacro{\ww}{2*sin(\xx)*sin(\yy)*cos(\az)}
    \pgfmathsetmacro{\opa}{min(.68,abs(\ww)*.43)}
    \pgfmathsetmacro{\xlo}{4.6*\i/8}
    \pgfmathsetmacro{\ylo}{4.6*\j/8}
    \pgfmathsetmacro{\xhi}{4.6*(\i+1)/8}
    \pgfmathsetmacro{\yhi}{4.6*(\j+1)/8}
    \ifdim\ww pt>0pt
      \fill[red,opacity=\opa] (\xlo,\ylo) rectangle (\xhi,\yhi);
    \else
      \fill[blue,opacity=\opa] (\xlo,\ylo) rectangle (\xhi,\yhi);
    \fi
   }
  }
  \foreach \i in {1,...,8} {
   \foreach \j in {1,...,8} {
    \pgfmathsetmacro{\xx}{360*(\i-.5)/8}
    \pgfmathsetmacro{\yy}{360*(\j-.5)/8}
    \pgfmathsetmacro{\ux}{sin(\xx)*cos(\yy)*cos(\az)}
    \pgfmathsetmacro{\uy}{-cos(\xx)*sin(\yy)*cos(\az)}
    \pgfmathsetmacro{\px}{4.6*(\i-.5)/8}
    \pgfmathsetmacro{\py}{4.6*(\j-.5)/8}
    \draw[-{Stealth[length=2pt]},black!85,line width=.4pt]
      (\px,\py)--++({.36*\ux},{.36*\uy});
   }
  }
  \draw[black!60] (0,0) rectangle (4.6,4.6);
  \node[anchor=south] at (2.3,4.68) {$z=\zlabel$};
  \node[anchor=north] at (0,0) {$0$};
  \node[anchor=north] at (4.6,0) {$2\pi$};
  \end{scope}
}
\node[anchor=west,font=\scriptsize] at (0,-.9)
  {\textcolor{red}{\rule{8pt}{6pt}}\ $\omega_{0,z}>0$\qquad
   \textcolor{blue}{\rule{8pt}{6pt}}\ $\omega_{0,z}<0$\qquad
   strzałki: $(u_{0,1},u_{0,2})$};
\end{tikzpicture}

{\small\emph{Rysunek:} Trzy przekroje pola początkowego wiru
Taylora--Greena, obliczone wprost ze wzorów na $u_0$ i
$\omega_{0,z}$. Strzałki przedstawiają rzuty prędkości na
płaszczyznę $xy$, a czerwień i błękit oznaczają znak pionowej
składowej wirowości (nasycenie barwy rośnie z jej modułem).
Rysunek powstaje w TikZ bez zewnętrznych obrazów i skryptów.}
\end{center}

\begin{center}
\begin{tikzpicture}[x={(.82cm,0cm)},y={(.33cm,.28cm)},
 z={(0cm,1.65cm)},>=Stealth,font=\footnotesize]
 \foreach \h/\az/\lab in {0/0/0,1/90/{\pi/2},2/180/\pi}{
  \fill[gray!5] (0,0,\h)--(6,0,\h)--(6,6,\h)--(0,6,\h)--cycle;
  \foreach \i in {0,...,7}{
   \foreach \j in {0,...,7}{
    \pgfmathsetmacro{\xx}{360*(\i+.5)/8}
    \pgfmathsetmacro{\yy}{360*(\j+.5)/8}
    \pgfmathsetmacro{\ww}{sin(\xx)*sin(\yy)*cos(\az)}
    \pgfmathsetmacro{\opa}{.68*abs(\ww)}
    \pgfmathsetmacro{\lo}{.75*\i}
    \pgfmathsetmacro{\hi}{.75*(\i+1)}
    \pgfmathsetmacro{\bot}{.75*\j}
    \pgfmathsetmacro{\topc}{.75*(\j+1)}
    \ifdim\ww pt>0pt
      \fill[red,opacity=\opa] (\lo,\bot,\h)--(\hi,\bot,\h)
        --(\hi,\topc,\h)--(\lo,\topc,\h)--cycle;
    \else
      \fill[blue,opacity=\opa] (\lo,\bot,\h)--(\hi,\bot,\h)
        --(\hi,\topc,\h)--(\lo,\topc,\h)--cycle;
    \fi
   }
  }
  \foreach \i in {0,...,5}{
   \foreach \j in {0,...,5}{
    \pgfmathsetmacro{\xx}{360*(\i+.5)/6}
    \pgfmathsetmacro{\yy}{360*(\j+.5)/6}
    \pgfmathsetmacro{\ux}{sin(\xx)*cos(\yy)*cos(\az)}
    \pgfmathsetmacro{\uy}{-cos(\xx)*sin(\yy)*cos(\az)}
    \ifnum\az=90\relax
    \else
      \draw[-{Stealth[length=2pt]},black!85,line width=.4pt]
        ({\i+.5},{\j+.5},\h)--++({.42*\ux},{.42*\uy},0);
    \fi
   }
  }
  \draw[black!65] (0,0,\h)--(6,0,\h)--(6,6,\h)--(0,6,\h)--cycle;
  \node[anchor=west] at (6,6,\h) {$z=\lab$};
 }
 \draw[->] (0,0,0)--(6.6,0,0) node[right] {$x$};
 \draw[->] (0,0,0)--(0,6.6,0) node[above] {$y$};
 \draw[->] (0,0,0)--(0,0,2.4) node[above] {$z$};
 \node[below] at (6,0,0) {$2\pi$};
 \node[left] at (0,6,0) {$2\pi$};
\end{tikzpicture}

{\small\emph{Rysunek:} Perspektywiczny widok trzech przekrojów
pola początkowego. Barwy i strzałki wynikają z rzeczywistych
wzorów na $\omega_{0,z}$ i prędkość, a nie ze schematycznych
okręgów. Na $z=\pi/2$ pole zanika; na $z=\pi$ zwroty
prędkości i znaki pionowej składowej wirowości odwracają się
względem $z=0$. Nie jest to wizualizacja ewolucji w czasie.}
\end{center}
\end{przyklad}

\section{Porównanie zagadnień brzegowych i dalsze kierunki}
Na ograniczonym obszarze Lipschitzowskim $\Omega\subset\R^d$
z warunkiem przylegania $u|_{\partial\Omega}=0$
przestrzenie $H,V$ tworzy się przez domknięcie
gładkich bezdywergencyjnych pól o zwartym nośniku
odpowiednio w $L^2(\Omega;\R^d)$ i
$H^1_0(\Omega;\R^d)$ (w drugim przypadku
równoważnie w normie gradientu).
Odzyskiwanie ciśnienia wymaga wówczas analizy warunków
brzegowych, a operator Stokesa nie sprowadza się do
prostego mnożnika Fouriera. Dlatego powyższe dowody
są sformułowane na torusie; ich bezpośrednie przeniesienie
na dowolny obszar bez dodatkowych twierdzeń byłoby nieuzasadnione.

Podstawowa różnica między wymiarami jest następująca:
energia $L^2$ i globalne istnienie rozwiązań Leraya--Hopfa
działają w $d=2,3$; w wymiarze dwa mamy dodatkowo
jednoznaczność i kontrolę wirowości, podczas gdy w wymiarze trzy
pojawia się rozciąganie wirów i warunkowe kryteria regularności.
W kolejnych rozdziałach będziemy badać inne mechanizmy
nieliniowe: swobodne brzegi i homogenizację.


\chapter[Nierówności wariacyjne i przeszkoda]{Nierówności wariacyjne, zagadnienie przeszkody i swobodny brzeg}
\label{chap:pde29}

\section{Motywacja: energia z ograniczeniem jednostronnym}
W poprzednich rozdziałach minimalizowaliśmy funkcjonały energii na całych
przestrzeniach Sobolewa. W wielu zagadnieniach mechanicznych funkcja
$u$ nie może jednak przekroczyć przeszkody w ustalonym kierunku.
Na przykład odkształcenie membrany nie może być mniejsze niż wysokość
sztywnej podpory $\psi$. W takim przypadku warunek Eulera--Lagrange'a
nie musi być równaniem w całym obszarze. Na zbiorze kontaktu pojawia
się dodatkowa reakcja, a sam zbiór kontaktu nie jest znany z góry.

Niech $\Omega\subset\R^d$ będzie ograniczonym obszarem Lipschitzowskim.
Pracujemy na rzeczywistych przestrzeniach funkcyjnych.
Oznaczamy $V=H_0^1(\Omega)$, z iloczynem skalarnym
$(v,w)_V=\int_\Omega\nabla v\cdot\nabla w\,\mathrm d\lambda_d(x)$
i normą $\|v\|_V=\|\nabla v\|_{L^2(\Omega)}$.
Nierówność Poincarégo~\ref{cor:pde10-Poincare}
daje równoważność tej normy z normą $H^1$, więc
$V$ jest przestrzenią Hilberta.
Przez $V'=H^{-1}(\Omega)$ oznaczamy jej przestrzeń
dualną z normą względem $\|\cdot\|_V$.
Wszystkie relacje $v\ge w$ rozumiemy prawie wszędzie.
Przyjmujemy $\psi\in H^1(\Omega)$ i $\psi^+\in V$,
gdzie $\psi^+=\max(\psi,0)$ należy do $H^1$
na mocy Wniosku~\ref{cor:pde9-truncation}.
Ostatni warunek gwarantuje istnienie funkcji dopuszczalnej.
Definiujemy
\[
 K_\psi=\{v\in V:v\geq\psi\ \text{prawie wszędzie w }\Omega\}.
\]
Zbiór ten jest niepusty, bo $\psi^+\in K_\psi$, oraz wypukły.
Przez $f\in V'$ oznaczamy obciążenie, a przez
$\langle f,v\rangle$ wartość funkcjonału $f$ na $v$.

\begin{lemat}[Domkniętość zbioru funkcji dopuszczalnych]
\label{lem:pde29-K}
Dla ograniczonego obszaru Lipschitzowskiego $\Omega\subset\R^d$
i dowolnej funkcji $\psi\in H^1(\Omega)$ spełniającej
$\psi^+\in H_0^1(\Omega)$ zbiór $K_\psi$ jest niepusty,
wypukły i słabo domknięty w $H_0^1(\Omega)$.
\end{lemat}
\begin{proof}
Niepustość wynika z $\psi^+\in K_\psi$.
Dla $v,w\in K_\psi$ i $\theta\in[0,1]$ mamy
$\theta v+(1-\theta)w\ge\psi$ prawie wszędzie.
Dla $v\in V$ zachodzi równoważność
\[
 v\in K_\psi\quad\Longleftrightarrow\quad
 \int_\Omega(v-\psi)g\,\mathrm d\lambda_d\ge0
 \quad\text{dla każdego }g\in L^2(\Omega),\ g\ge0.
\]
Implikacja w prawo wynika z nieujemności czynników.
W drugą stronę wybieramy $g=(\psi-v)^+$:
całka wynosi $-\|(\psi-v)^+\|_2^2$, więc $v\ge\psi$.
Dla ustalonego $g$ funkcjonał
$v\mapsto\int_\Omega vg\,\mathrm d\lambda_d$
jest ciągły w $V$, na mocy nierówności Poincarégo
i nierówności Cauchy'ego--Schwarza.
Zatem $K_\psi$ jest przecięciem słabo domkniętych
półprzestrzeni, a więc jest słabo domknięty,
nie tylko słabo ciągowo domknięty.
\end{proof}

\section{Rozwiązanie wariacyjne i jego jednoznaczność}
Niech $A(x)=(a_{ij}(x))_{i,j=1}^d$ będzie mierzalną macierzą
symetryczną, $A\in L^\infty(\Omega;\R^{d\times d})$, oraz
$\lambda,\Lambda>0$ takie, że dla prawie każdego $x\in\Omega$
i każdego $\xi\in\R^d$ zachodzi
$\lambda|\xi|^2\le\xi^{\mathsf T}A(x)\xi\le\Lambda|\xi|^2$.
Wprowadzamy formę $a(u,v)=\int_\Omega A\nabla u\cdot\nabla v\,
\mathrm d\lambda_d(x)$ i funkcjonał
$J_f(v)=\frac12a(v,v)-\langle f,v\rangle$.
Z symetrii $A(x)$ i oszacowań jej wartości własnych
wynika $\|A(x)\|_{\mathrm{op}}\le\Lambda$ prawie wszędzie.
Nierówność Cauchy'ego--Schwarza daje
$|a(u,v)|\le\Lambda\|u\|_V\|v\|_V$,
a eliptyczność $a(v,v)\ge\lambda\|v\|_V^2$.
Forma $a$ jest iloczynem skalarnym, którego norma
$\|v\|_a=\sqrt{a(v,v)}$ spełnia
\[
 \sqrt\lambda\|v\|_V\le\|v\|_a\le
 \sqrt\Lambda\|v\|_V.
\]
Jest więc zupełna i wyznacza tę samą topologię
słabą, gdyż obie normy mają te same ciągłe
funkcjonały liniowe.

\begin{twierdzenie}[Istnienie i jednoznaczność w zagadnieniu przeszkody]
\label{thm:pde29-existence}
Niech $\Omega$, $A$, $\psi$ spełniają powyższe założenia,
a $f\in H^{-1}(\Omega)$. Istnieje dokładnie jedna funkcja
$u\in K_\psi$, dla której
\begin{equation}\label{eq:pde29-VI}
 a(u,v-u)\ge\langle f,v-u\rangle
 \qquad\text{dla każdego }v\in K_\psi.
\end{equation}
Jest ona jedynym punktem, w którym $J_f$ osiąga minimum na $K_\psi$.
\end{twierdzenie}
\begin{proof}
Z eliptyczności, definicji normy dualnej
i nierówności Younga mamy
\[
 J_f(v)\ge\tfrac\lambda2\|v\|_V^2-\|f\|_{V'}\|v\|_V
 \ge\tfrac\lambda4\|v\|_V^2-\lambda^{-1}\|f\|_{V'}^2.
\]
Niech $m=\inf_{K_\psi}J_f$. Z niepustości zbioru
i powyższego oszacowania $-\infty<m\le J_f(\psi^+)<\infty$.
Wybieramy $u_n\in K_\psi$ z $J_f(u_n)\le m+1/n$.
Ciąg jest ograniczony w $V$; refleksywność i twierdzenie
Eberleina--Šmuliana pozwalają wydzielić podciąg
słabo zbieżny do $u\in V$.
Z Lematu~\ref{lem:pde29-K} mamy $u\in K_\psi$.
Norma $\|\cdot\|_a$ jest słabo półciągła z dołu,
bo $\|u\|_a=\sup_{\|w\|_a\le1}|a(u,w)|$
i każdy funkcjonał $a(\cdot,w)$ jest słabo ciągły.
Składnik liniowy jest słabo ciągły, więc
$m\le J_f(u)\le\liminf_n J_f(u_n)=m$.
Otrzymujemy minimizator.

Dla $v\in K_\psi$ odcinek $u+t(v-u)$ należy do $K_\psi$
dla $t\in[0,1]$. Minimalność daje
$0\le[J_f(u+t(v-u))-J_f(u)]/t
 =a(u,v-u)-\langle f,v-u\rangle+
 \frac t2a(v-u,v-u)$. Przechodząc z $t$ do zera od strony dodatniej,
otrzymujemy~\eqref{eq:pde29-VI}. Odwrotnie, dla funkcji
spełniającej tę nierówność
$J_f(v)-J_f(u)=a(u,v-u)-\langle f,v-u\rangle
+\frac12a(v-u,v-u)\ge0$.
Dokładniej, nierówność wariacyjna i eliptyczność dają
\[
 J_f(v)-J_f(u)\ge\tfrac12\|v-u\|_a^2
 \ge\tfrac\lambda2\|v-u\|_V^2.
\]
Każdy ciąg minimalizujący zbiega więc silnie do $u$
w normie $V$.
Jeżeli $u_1,u_2$ spełniają nierówność, testujemy pierwszą przez
$v=u_2$, drugą przez $v=u_1$ i dodajemy. Otrzymujemy
$a(u_1-u_2,u_1-u_2)\le0$, skąd $u_1=u_2$.
\end{proof}


\begin{stwierdzenie}[Zagadnienie przeszkody jako rzut energetyczny]
\label{prop:pde-v44-obstacle-projection}
Niech $V=H_0^1(\Omega)$, $a$ będzie symetryczną,
ciągłą i koercywną formą z twierdzenia~\ref{thm:pde29-existence},
a $K_\psi\subset V$ niepustym, domkniętym i wypukłym
zbiorem dopuszczalnym. Niech $z\in V$ będzie jedynym
rozwiązaniem równania $a(z,v)=\langle f,v\rangle$ dla
każdego $v\in V$. Wtedy rozwiązanie przeszkody $u$
jest rzutem metrycznym $z$ na $K_\psi$ w normie
$\|v\|_a=a(v,v)^{1/2}$. Co więcej, dla ustalonego
$K_\psi$ i dwóch danych $f_i\in V'$ zachodzi
\begin{equation}\label{eq:pde-v44-projection-nonexpansive}
 \|u_1-u_2\|_a\le \|z_1-z_2\|_a
 \le\lambda^{-1/2}\|f_1-f_2\|_{V'},
\end{equation}
gdzie $z_i$ rozwiązuje równanie liniowe dla $f_i$,
a $u_i$ odpowiednie zagadnienie przeszkody.
\end{stwierdzenie}
\begin{proof}
Istnienie i jednoznaczność $z$ wynikają z twierdzenia Laxa--Milgrama;
forma $a$ jest iloczynem skalarnym równoważnym
standardowemu na $V$. Równanie dla $z$ daje
\[
 J_f(v)=\tfrac12\|v-z\|_a^2-\tfrac12\|z\|_a^2.
\]
Minimalizacja energii jest zatem dokładnie
minimalizacją odległości od $z$. Z warunku
$a(u,v-u)\ge\langle f,v-u\rangle=a(z,v-u)$
dla $v\in K_\psi$ otrzymujemy
$a(u-z,v-u)\ge0$. Jest to dokładnie
charakteryzacja rzutu na domknięty zbiór wypukły
w przestrzeni Hilberta $(V,a)$.
Dla dwóch rzutów wykorzystujemy w nierównościach
wariacyjnych testy $v=u_2$ i $v=u_1$, a następnie
dodajemy. Dla $w=u_1-u_2$ dostajemy
$a(w,w)\le a(z_1-z_2,w)$; nierówność Cauchy'ego--Schwarza
daje pierwsze oszacowanie, po podzieleniu
przez $\|w\|_a$ dla $w\ne0$.
Dla $\delta z=z_1-z_2$ testujemy równanie
liniowe przez $\delta z$:
\[
 \|\delta z\|_a^2=\langle f_1-f_2,\delta z\rangle
 \le\lambda^{-1/2}\|f_1-f_2\|_{V'}\|\delta z\|_a.
\]
Dla $\delta z\ne0$ dzielimy przez normę, a dla
$\delta z=0$ oszacowanie jest natychmiastowe.
\end{proof}


\begin{stwierdzenie}[Stabilność względem obciążenia]
\label{prop:pde29-stable}
Przy ustalonych $\Omega$, $A$, $\psi$ z
 twierdzenia~\ref{thm:pde29-existence}, jeżeli $u_i$ odpowiadają
$f_i\in V'$ ($i=1,2$), to
$\|u_1-u_2\|_V\le\lambda^{-1}\|f_1-f_2\|_{V'}$.
\end{stwierdzenie}
\begin{proof}
W obu nierównościach wariacyjnych wybieramy jako test funkcję
będącą drugim rozwiązaniem. Po dodaniu i zastosowaniu
eliptyczności otrzymujemy
$\lambda\|u_1-u_2\|_V^2\le
\langle f_1-f_2,u_1-u_2\rangle$.
Dalej stosujemy nierówność dualną; przypadek $u_1=u_2$
nie wymaga dzielenia przez normę.
\end{proof}

\begin{twierdzenie}[Zasada porównania dla przeszkód]
\label{thm:pde29-order}
Niech $A$, $\Omega$ spełniają powyższe założenia,
a $\psi_i\in H^1(\Omega)$, $\psi_i^+\in V$ oraz
$f_i\in L^2(\Omega)$, $i=1,2$. Jeśli
$\psi_1\le\psi_2$ i $f_1\le f_2$ prawie wszędzie,
to odpowiadające rozwiązania spełniają $u_1\le u_2$
prawie wszędzie w $\Omega$.
\end{twierdzenie}
\begin{proof}
Niech $w=(u_1-u_2)^+$.
Z Wniosku~\ref{cor:pde9-truncation} mamy $w\in H^1$
oraz $\nabla w=\chi_{\{u_1>u_2\}}\nabla(u_1-u_2)$.
Aby wykazać $w\in V$, przybliżamy różnicę
$u_1-u_2$ przez funkcje $g_n\in C_c^\infty(\Omega)$
w normie $H^1$. Funkcje $g_n^+$ mają zwarty
nośnik w $\Omega$ i należą do $V$ przez wygładzanie.
Ich normy $H^1$ są ograniczone, a
$\|g_n^+-w\|_2\le\|g_n-(u_1-u_2)\|_2\to0$.
Słaba domkniętość $V$ w $H^1$ identyfikuje
granicę jako element $V$.
Ponieważ
$\min(u_1,u_2)\ge\psi_1$ i
$\max(u_1,u_2)\ge\psi_2$,
w nierówności dla $u_1$ wolno wziąć $v=u_1-w$,
a w nierówności dla $u_2$ funkcję $v=u_2+w$.
Daje to $a(u_1,w)\le\int_\Omega f_1w\,\mathrm d\lambda_d(x)$
oraz $a(u_2,w)\ge\int_\Omega f_2w\,\mathrm d\lambda_d(x)$.
Po odjęciu i wykorzystaniu $\nabla w=\chi_{\{u_1>u_2\}}
\nabla(u_1-u_2)$ otrzymujemy
$\lambda\|w\|_V^2\le\int_\Omega(f_1-f_2)w\,\mathrm d\lambda_d(x)
\le0$. Stąd $w=0$.
\end{proof}


\begin{twierdzenie}[Ciągła zależność od przeszkody w normie energetycznej]
\label{thm:pde29-obstacle-stability}
Niech $\Omega\subset\R^d$ będzie ograniczonym obszarem
Lipschitzowskim, $V=H_0^1(\Omega)$, a $A$ niech będzie
symetryczna, mierzalna i niech spełnia warunki
$\lambda|\xi|^2\le\xi^{\mathsf T}A(x)\xi\le\Lambda|\xi|^2$
dla $\lambda,\Lambda>0$. Niech $f\in V'$ będzie ustalone.
Dla $j=1,2$ wybierzmy przeszkodę $\psi_j\in V$
i oznaczmy przez $u_j\in K_{\psi_j}$ rozwiązanie
nierówności wariacyjnej~\eqref{eq:pde29-VI}.
Wtedy
\begin{equation}\label{eq:pde29-obstacle-stability}
 \|u_1-u_2\|_V\le\sqrt{\frac{\Lambda}{\lambda}}
 \|\psi_1-\psi_2\|_V.
\end{equation}
\end{twierdzenie}
\begin{proof}
Niech $h=\psi_1-\psi_2\in V$.
Ponieważ $u_2\ge\psi_2$, mamy $u_2+h\ge\psi_1$;
podobnie $u_1-h\ge\psi_2$. Są to zatem poprawne
funkcje testowe w obu nierównościach wariacyjnych.
Pierwszą testujemy przez $u_2+h$, a drugą przez $u_1-h$.
Po dodaniu, przy $w=u_1-u_2$, otrzymujemy
\[
 a(w,w-h)\le\langle f,w-h\rangle
                 -\langle f,w-h\rangle=0.
\]
Stąd $a(w,w)\le a(w,h)$.
Nierówność Cauchy'ego--Schwarza w iloczynie
skalarnym $a$ daje $\|w\|_a^2\le\|w\|_a\|h\|_a$,
czyli $\|w\|_a\le\|h\|_a$ (dla $w=0$ nie dzielimy).
Z równoważności norm dostajemy
\[
 \|w\|_V\le\lambda^{-1/2}\|w\|_a
 \le\lambda^{-1/2}\|h\|_a
 \le\sqrt{\Lambda/\lambda}\|h\|_V.
\]
Jest to~\eqref{eq:pde29-obstacle-stability}.
\end{proof}

\begin{uwaga}[Ograniczenie założenia o przeszkodach]
Twierdzenie dotyczy przeszkód $\psi_j\in H_0^1(\Omega)$,
a nie tylko warunku $\psi_j^+\in H_0^1(\Omega)$.
Jest to istotne: przesunięte funkcje testowe
$u_2+\psi_1-\psi_2$ i $u_1-\psi_1+\psi_2$
muszą należeć do $H_0^1(\Omega)$.
\end{uwaga}

\section{Reakcja przeszkody i warunki komplementarności}
W tej sekcji rozważamy $A=I$. Oznaczamy przez $-\Delta u$
dystrybucję określoną wzorem
$\langle-\Delta u,\varphi\rangle
=\int_\Omega\nabla u\cdot\nabla\varphi\,\mathrm d\lambda_d(x)$.

\begin{stwierdzenie}[Dodatnia reakcja jako miara]
\label{prop:pde29-lambda-measure}
Niech $u\in K_\psi$ rozwiązuje~\eqref{eq:pde29-VI}
dla $A=I$ i $f\in V'$. Dystrybucja
$\mu=-\Delta u-f$ jest nieujemna, czyli
$\langle\mu,\varphi\rangle\ge0$ dla każdego
$\varphi\in C_c^\infty(\Omega)$, $\varphi\ge0$.
Wobec tego $\mu$ jest nieujemną miarą Radona w $\Omega$.
\end{stwierdzenie}
\begin{proof}
Jeśli $\varphi\in V$ i $\varphi\ge0$ prawie wszędzie,
to $u+\varphi\in K_\psi$,
więc~\eqref{eq:pde29-VI} daje
$\langle-\Delta u-f,\varphi\rangle\ge0$.
W szczególności $\mu$ jest dodatnią dystrybucją
i ciągłym funkcjonałem na $V$.
Dla dowolnego zwartego $K\Subset\Omega$ wybierzmy
$\eta\in C_c^\infty(\Omega)$, $\eta\ge0$, $\eta=1$
w otoczeniu $K$. Jeżeli $\varphi\in C_c^\infty(\Omega)$ i
$\operatorname{supp}\varphi\subset K$,
to $-\|\varphi\|_\infty\eta\le\varphi\le
\|\varphi\|_\infty\eta$. Dodatniość daje
$|\langle\mu,\varphi\rangle|\le
\|\varphi\|_\infty\langle\mu,\eta\rangle$.
Dla $g\in C_c(\Omega)$ wybieramy gładkie
przybliżenia $g_n$ z nośnikami w jednym zwartym
podzbiorze $\Omega$, zbiegające jednostajnie do $g$.
Można je otrzymać przez przedłużenie przez zero
i splot z jądrem wygładzającym o małym nośniku.
Powyższe oszacowanie zapewnia istnienie granicy
$L(g)=\lim_n\langle\mu,g_n\rangle$ oraz jej
niezależność od przybliżeń (dla dwóch ciągów
stosujemy oszacowanie na sumie ich nośników).
Funkcjonał $L$ jest liniowy i dodatni:
dla $g\ge0$ wybieramy nieujemne jądro splotowe.
Twierdzenie Riesza o reprezentacji dodatniego
funkcjonału na $C_c(\Omega)$ daje jedyną
nieujemną miarę Radona reprezentującą $L$.
Jej działanie na funkcjach testowych jest
równe działaniu pierwotnej dystrybucji.
Nie wnioskujemy, że miara ma skończoną masę
na całej $\Omega$. Nie utożsamiamy też parowania
w $V',V$ z całką dowolnego reprezentanta
funkcji Sobolewa względem $\mu$; reprezentanci
równi prawie wszędzie względem $\lambda_d$
mogą różnić się na zbiorze o dodatniej mierze $\mu$.
\end{proof}

\begin{uwaga}[Komplementarność w parowaniu dualnym]
Jeżeli dodatkowo $\psi\in V$, to można testować
nierówność wariacyjną przez $v=\psi$.
Otrzymujemy $\langle\mu,u-\psi\rangle\le0$.
Ponieważ $u-\psi\ge0$ i należy do $V$, dodatniość
funkcjonału daje nierówność przeciwną, a zatem
$\langle\mu,u-\psi\rangle=0$.
Odwrotnie, jeśli $u\in K_\psi$, $\mu=-\Delta u-f$
jest dodatnia na stożku nieujemnych elementów $V$
i zachodzi ta równość, to dla każdego $v\in K_\psi$
mamy $\langle\mu,v-u\rangle
=\langle\mu,v-\psi\rangle\ge0$, czyli~\eqref{eq:pde29-VI}.
Przy przeszkodzie spoza $V$ takie parowanie
nie jest na ogół określone.
\end{uwaga}

\begin{lemat}[Pochodne funkcji Sobolewa na zbiorze jej zer]
\label{lem:pde29-zero-level}
Niech $\Omega\subset\R^d$ będzie zbiorem otwartym.
Jeżeli $z\in H^1_{\mathrm{loc}}(\Omega)$ i
$E=\{x\in\Omega:z(x)=0\}$, to $\nabla z=0$
prawie wszędzie na $E$. Jeśli ponadto $z\in H^2_{\mathrm{loc}}(\Omega)$,
to $D^2z=0$ prawie wszędzie na $E$.
\end{lemat}
\begin{proof}
Pierwszą część można otrzymać z Lematu
~\ref{lem:pde9-levelsets}, stosując go do ograniczeń
$z$ na relatywnie zwartych kostkach zawartych w $\Omega$.
Przeliczalna rodzina takich kostek pokrywa $\Omega$,
więc suma wyjątkowych zbiorów miary zero również
ma miarę zero.
Jeśli $z\in H^2_{\mathrm{loc}}$, to dla każdego $j$
funkcja $\partial_jz\in H^1_{\mathrm{loc}}$.
Na mocy pierwszej części $E$ jest, z dokładnością
do zbioru miary zero, zawarty w zbiorze
$\{\partial_jz=0\}$. Ponownie stosując pierwszą część,
tym razem do $\partial_jz$, otrzymujemy
$\partial_k\partial_jz=0$ prawie wszędzie na $E$.
Biorąc skończoną sumę wyjątkowych zbiorów
dla wszystkich $j,k$, otrzymujemy tezę o $D^2z$.
\end{proof}

\begin{twierdzenie}[Regularność $H^2$ i nierówność Lewy'ego--Stampacchii]
\label{thm:pde29-LS}
Niech $\Omega\subset\R^d$ będzie ograniczonym obszarem
klasy $C^{1,1}$, $f\in L^2(\Omega)$ i
$\psi\in H^2(\Omega)$, przy czym ślad $\psi$ na
$\partial\Omega$ jest niedodatni. Niech $u\in H_0^1(\Omega)$
będzie rozwiązaniem zagadnienia przeszkody dla $A=I$.
Wówczas $u\in H^2(\Omega)$, a $\mu=-\Delta u-f$ jest funkcją
z $L^2(\Omega)$ i prawie wszędzie zachodzi
\begin{equation}\label{eq:pde29-LS}
 0\le\mu\le(-\Delta\psi-f)^+,
 \qquad f\le-\Delta u\le\max\{f,-\Delta\psi\}.
\end{equation}
Ponadto $(u-\psi)\mu=0$ prawie wszędzie w $\Omega$.
\end{twierdzenie}
\begin{proof}
\textbf{1. Zadanie z karą.}
Dla $\eps>0$ minimalizujemy na $V=H_0^1(\Omega)$
funkcjonał
\[
 J_\eps(v)=\frac12\|\nabla v\|_2^2
 -\int_\Omega fv\,\mathrm d\lambda_d(x)
 +\frac1{2\eps}\int_\Omega((\psi-v)^+)^2\,\mathrm d\lambda_d(x).
\]
Nieujemny składnik kary zachowuje koercywność
funkcjonału $J_f$. Funkcja $s\mapsto(s^+)^2$
jest wypukła (jej pochodna $2s^+$ jest niemalejąca),
więc całka kary jest wypukła. Energia Dirichleta
jest ściśle wypukła na $V$.
Dla ograniczonego ciągu minimalizującego wybieramy
podciąg słabo zbieżny w $V$ i silnie zbieżny w $L^2$
na mocy twierdzenia Rellicha~\ref{thm:pde10-Rellich-p}.
Ponieważ $s\mapsto s^+$ jest $1$-lipschitzowska,
$(\psi-v_n)^+\to(\psi-v)^+$ w $L^2$, więc kara
zbiega, a energia jest słabo półciągła z dołu.
Otrzymujemy jedyny minimizator $u_\eps$.

Dla $\varphi\in V$ różniczkujemy
$J_\eps(u_\eps+t\varphi)$ w $t=0$.
Funkcja $F(s)=(s^+)^2/2$ ma pochodną $F'(s)=s^+$
o stałej Lipschitza $1$. Z twierdzenia Newtona--Leibniza
wynika $|F(s+h)-F(s)-F'(s)h|\le |h|^2/2$.
Stosując to punktowo do $s=\psi-u_\eps$ i $h=-t\varphi$,
po całkowaniu otrzymujemy resztę ograniczoną przez
$t^2\|\varphi\|_2^2/(2\eps)$ dla składnika kary.
Zatem warunek stacjonarności daje dla każdego $\varphi\in V$
\[
 \int_\Omega\nabla u_\eps\cdot\nabla\varphi\,
 \mathrm d\lambda_d=\int_\Omega(f+\mu_\eps)\varphi\,
 \mathrm d\lambda_d,\qquad
 \mu_\eps=\eps^{-1}(\psi-u_\eps)^+\ge0.
\]
Jest to równanie $-\Delta u_\eps=f+\mu_\eps$ w $V'$.

\textbf{2. Jednostajna kontrola reakcji.}
Niech $w_\eps=(\psi-u_\eps)^+$.
Operator śladu komutuje z częścią dodatnią:
dla gładkich przybliżeń w $H^1$ ich części dodatnie
są ograniczone w $H^1$ i zbiegają w $L^2$,
a więc zbiegają słabo w $H^1$.
Z ciągłości śladu i jego zgodności z ograniczeniem
funkcji ciągłych (Twierdzenie~\ref{thm:pde10-trace})
otrzymujemy słabą zbieżność śladów w $L^2$
na brzegu; jednocześnie części dodatnie śladów
zbiegają silnie w $L^2$ dzięki własności Lipschitza.
Stąd ślad $w_\eps$ jest równy części dodatniej
śladu $\psi-u_\eps$, czyli zeru.
Na mocy Twierdzenia~\ref{thm:pde10-trace-zero}
mamy $w_\eps\in V$. Ten sam argument dla $\psi^+$
uzasadnia niepustość $K_\psi$ w hipotezach twierdzenia.
Wprowadzamy $h=-\Delta\psi-f\in L^2(\Omega)$.
Odejmujemy równanie $-\Delta\psi=f+h$ od
równania dla $u_\eps$ i testujemy różnicę przez
$w_\eps$. Reguła obcięcia
~\ref{cor:pde9-truncation} daje
$\nabla w_\eps=\chi_{\{\psi>u_\eps\}}
\nabla(\psi-u_\eps)$, więc lewa strona równania
dla różnicy wynosi $-\|\nabla w_\eps\|_2^2$.
Prawa strona jest równa
$\int_\Omega(\mu_\eps-h)w_\eps\,\mathrm d\lambda_d$.
Podstawiając $w_\eps=\eps\mu_\eps$ i przestawiając
składniki, dostajemy
\begin{equation}\label{eq:pde29-penalty}
 \|\nabla w_\eps\|_2^2+\eps\|\mu_\eps\|_2^2
 =\eps\int_\Omega h\mu_\eps\,\mathrm d\lambda_d(x)
 \le\eps\|h^+\|_2\|\mu_\eps\|_2.
\end{equation}
W szczególności $\|\mu_\eps\|_2\le\|h^+\|_2$,
$\|w_\eps\|_2\le\eps\|h^+\|_2$ oraz
$\|\nabla w_\eps\|_2^2\le
\eps\|h^+\|_2^2\to0$.
Testujemy równanie $-\Delta u_\eps=f+\mu_\eps$ przez $u_\eps$.
Z nierówności Poincarégo i oszacowania reakcji otrzymujemy
\[
 \|\nabla u_\eps\|_2^2
 =\int_\Omega(f+\mu_\eps)u_\eps\,\mathrm d\lambda_d(x)
 \le C_P(\|f\|_2+\|h^+\|_2)\|\nabla u_\eps\|_2.
\]
Zatem $\|u_\eps\|_V\le C_P(\|f\|_2+\|h^+\|_2)$
niezależnie od $\eps$, co uzasadnia wybranie słabo zbieżnego
podciągu w następnym kroku.

\textbf{3. Granica i warunek komplementarności.}
Dla dowolnego ciągu $\eps_n\downarrow0$
z jednostajnych oszacowań, refleksywności i twierdzenia
Rellicha wydzielamy podciąg (zachowujemy zapis $\eps$)
taki, że $u_\eps\rightharpoonup u$ w $V$,
$u_\eps\to u$ w $L^2(\Omega)$ i
$\mu_\eps\rightharpoonup\mu$ w $L^2(\Omega)$.
Nierówność Lipschitza daje
$(\psi-u_\eps)^+\to(\psi-u)^+$ w $L^2$,
a norma pierwszej funkcji dąży do zera, więc $u\ge\psi$.
Dla ustalonego $\varphi\in V$ przechodzimy
do granicy w równaniu: gradient zbiega słabo w $L^2$,
a reakcja słabo w $L^2$, gdzie $\varphi\in L^2$.
Otrzymujemy $-\Delta u=f+\mu$ w $V'$.
Dodatniość $\mu$ wynika z testowania słabej
zbieżności przez $\mu^-\in L^2$: mamy
$0\le\int_\Omega\mu\mu^-\,\mathrm d\lambda_d
=-\|\mu^-\|_2^2$.
Dodatkowo
\[
 \int_\Omega(u_\eps-\psi)\mu_\eps\,\mathrm d\lambda_d(x)
 =-\eps\|\mu_\eps\|_2^2\longrightarrow0.
\]
Różnicę całek rozbijamy na
$\int_\Omega(u_\eps-u)\mu_\eps\,\mathrm d\lambda_d$
oraz $\int_\Omega(u-\psi)(\mu_\eps-\mu)\,\mathrm d\lambda_d$.
Pierwsza dąży do zera z nierówności
Cauchy'ego--Schwarza i jednostajnego ograniczenia
$\|\mu_\eps\|_2$, a druga ze słabej zbieżności reakcji. Stąd
$\int_\Omega(u-\psi)\mu\,\mathrm d\lambda_d(x)=0$;
nieujemność czynników daje $(u-\psi)\mu=0$ prawie wszędzie.
Dla $v\in K_\psi$ zachodzi więc
$\int_\Omega\mu(v-u)\,\mathrm d\lambda_d(x)
 =\int_\Omega\mu(v-\psi)\,\mathrm d\lambda_d(x)\ge0$.
Wynika stąd~\eqref{eq:pde29-VI}, a jednoznaczność z
twierdzenia~\ref{thm:pde29-existence} identyfikuje $u$.
W istocie zachodzi silna zbieżność całej rodziny
w $V$, z oszacowaniem
\[
 \|u_\eps-u\|_V^2\le\tfrac{\eps}{4}\|\mu\|_2^2.
\]
Odejmujemy równania dla $u_\eps$ i $u$ oraz testujemy
przez $u_\eps-u$. Ponieważ $\mu_\eps\ge0$, $u\ge\psi$,
mamy $\mu_\eps(u_\eps-u)\le-\eps\mu_\eps^2$.
Z kolei komplementarność daje
$-\mu(u_\eps-u)=\mu(\psi-u_\eps)\le\eps\mu\mu_\eps$.
Po całkowaniu
$\|u_\eps-u\|_V^2\le
\eps(-\|\mu_\eps\|_2^2+(\mu,\mu_\eps)_{L^2})
\le\eps\|\mu\|_2^2/4$,
przez dopełnienie kwadratu w $L^2$.
Argument dotyczy każdego $\eps>0$, nie tylko podciągu.

\textbf{4. Regularność i nierówność punktowa.}
Z równania $-\Delta u=f+\mu\in L^2(\Omega)$ oraz
globalnego oszacowania regularności Dirichleta na obszarze
klasy $C^{1,1}$ (Twierdzenie~\ref{thm:pde15-boundary}
w rozdziale~\ref{chap:pde15}) otrzymujemy
$u\in H^2(\Omega)$ oraz
$\|u\|_{H^2}\le C_\Omega(\|f\|_2+\|h^+\|_2)$,
gdyż $\|\mu\|_2\le\|h^+\|_2$ ze słabej
półciągłości normy. Niech $z=u-\psi\in H^2(\Omega)$.
Z lematu~\ref{lem:pde29-zero-level} wynika
$D^2z=0$ prawie wszędzie na zbiorze $\{z=0\}$.
Na zbiorze kontaktu $\{u=\psi\}$ mamy więc
$-\Delta u=-\Delta\psi$ prawie wszędzie;
na jego dopełnieniu $\mu=0$ na mocy komplementarności.
Zatem $\mu=h$ prawie wszędzie na $\{u=\psi\}$,
gdzie dodatniość $\mu$ wymusza $h\ge0$,
i $\mu=0$ prawie wszędzie na $\{u>\psi\}$.
Otrzymujemy $0\le\mu\le h^+$ i po dodaniu $f$
drugi zapis~\eqref{eq:pde29-LS}.
\end{proof}

\begin{uwaga}[Granica regularności]
Oszacowanie $H^2$ zależy od regularności brzegu, a
nierówność~\eqref{eq:pde29-LS} od drugich pochodnych przeszkody.
Nie wolno ich przypisywać rozwiązaniu dla dowolnej
przeszkody $H^1$ ani dowolnego obszaru Lipschitzowskiego.
Na ogólnych obszarach $\mu$ jest na początku jedynie miarą
Radona, a swobodny brzeg może być nieregularny.
\end{uwaga}

\section{Zbiór kontaktu i swobodny brzeg}
Przyjmijmy założenia Twierdzenia~\ref{thm:pde29-LS}
oraz $d\le3$. Wówczas $u,\psi\in H^2(\Omega)$
mają jednoznaczne ciągłe reprezentanty na $\overline\Omega$.
Istotnie, dla $d=3$ stosujemy Twierdzenie
~\ref{thm:pde10-domain-embed} do funkcji i jej pochodnych
pierwszego rzędu: $H^2\hookrightarrow W^{1,6}$,
a następnie $W^{1,6}\hookrightarrow C^{0,1/2}$.
Dla $d=2$ analogicznie wybieramy wykładnik $4$,
otrzymując $H^2\hookrightarrow W^{1,4}\hookrightarrow C^{0,1/2}$.
Dla $d=1$ wystarcza jednowymiarowa charakterystyka
$H^1$ z Twierdzenia~\ref{thm:pde9-ac}.
Ciągły reprezentant różnicy $u-\psi$ jest nieujemny
w każdym punkcie $\Omega$: gdyby był ujemny w pewnym
punkcie, byłby ujemny w całej kuli dodatniej miary,
wbrew warunkowi dopuszczalności.
Definiujemy wtedy
$\mathcal C=\{x\in\Omega:u(x)=\psi(x)\}$ (zbiór kontaktu)
oraz $\mathcal N=\{x\in\Omega:u(x)>\psi(x)\}$ (zbiór braku kontaktu).
Zbiory $\mathcal C$ i $\mathcal N$ są odpowiednio
względnie domknięty i otwarty w $\Omega$, a
$\Omega=\mathcal C\mathbin{\dot\cup}\mathcal N$.
Zbiorem swobodnego brzegu nazywamy
$\Gamma=\partial_\Omega\mathcal N
=\mathcal C\cap\overline{\mathcal N}^{\,\Omega}$,
gdzie domknięcie i brzeg rozumiemy w topologii względnej.
Z komplementarności wynika $\mu=0$ prawie wszędzie
na $\mathcal N$, więc $-\Delta u=f$ w sensie
dystrybucji na tym zbiorze. Można to też wykazać
bez korzystania z regularności reakcji: dla
$\varphi\in C_c^\infty(\mathcal N)$ różnica
$u-\psi$ ma dodatnie minimum na nośniku $\varphi$
(o ile $\varphi\ne0$), zatem oba testy
$u\pm t\varphi$ są dopuszczalne dla dostatecznie małego
$t>0$. Nierówność wariacyjna dla obu znaków
daje $\int_\Omega\nabla u\cdot\nabla\varphi\,
\mathrm d\lambda_d=\int_\Omega f\varphi\,\mathrm d\lambda_d$.
Na $\mathcal C$ z Lematu~\ref{lem:pde29-zero-level}
wynika $-\Delta u=-\Delta\psi$ prawie wszędzie,
a zatem $\mu=(-\Delta\psi-f)\chi_{\mathcal C}$
prawie wszędzie. Nie utożsamiamy zbioru kontaktu
z nośnikiem reakcji: reakcja może znikać na części
zbioru kontaktu (na przykład dla $\psi=f=0$). Zbiór $\Gamma$ nie jest dany
z góry i jego regularność wymaga dodatkowej analizy.

\begin{przyklad}[Swobodny brzeg i skok drugiej pochodnej w wymiarze jeden]
\label{ex:pde29-explicit}
Niech $\Omega=(-1,1)$, $f=0$ i
$\psi(x)=(1-x^2)^2$.
Szukamy parzystego rozwiązania, które na $[0,a]$
pokrywa się z przeszkodą, a na $[a,1]$ jest afiniczne
i znika w punkcie $1$, gdzie $0<a<1$.
Ciągłość wymusza wzór $u(x)=m(1-x)$ na $[a,1]$,
z $m=\psi(a)/(1-a)=(1-a)(1+a)^2$.
Ciągłość pochodnej w $a$ wymaga
$m=-\psi'(a)=4a(1-a^2)$.
Po podzieleniu przez $(1-a)(1+a)>0$
otrzymujemy $1+a=4a$, czyli $a=1/3$ i $m=32/27$.
Jest to sposób znalezienia kandydata, a nie
założenie o postaci rozwiązania.
Zweryfikujemy, że jedynym rozwiązaniem jest funkcja
\begin{equation}\label{eq:pde29-explicit}
 u(x)=\begin{cases}
 (1-x^2)^2,& |x|\le\frac13,\\
 \frac{32}{27}(1-|x|),&\frac13<|x|\le1.
 \end{cases}
\end{equation}
Istotnie, $u(\pm1)=0$. W punkcie $1/3$
obie wartości wynoszą $64/81$, a obie pochodne
$-32/27$; w punkcie $-1/3$ korzystamy z parzystości.
Funkcja jest więc klasy $C^1$ na $[-1,1]$,
a jej pochodna jest odcinkami gładka z ograniczoną
drugą pochodną. Stąd $u'\in W^{1,\infty}(-1,1)$,
gdyż całkowanie przez części na trzech podprzedziałach
nie daje wyrazów skokowych w punktach sklejenia.
W szczególności $u\in H^2(-1,1)\cap H_0^1(-1,1)$. Dla $x\in[1/3,1]$ zachodzi
\[
 u(x)-\psi(x)=\frac{(1-x)(3x-1)^2(3x+5)}{27}\ge0;
\]
na lewej połowie korzystamy z parzystości.
Dla $1/3<x<1$ wszystkie czynniki są dodatnie,
więc zbiór kontaktu wewnątrz $\Omega$ jest dokładnie
$[-1/3,1/3]$.
Ponadto $-u''=(4-12x^2)\chi_{\{|x|<1/3\}}\ge0$
prawie wszędzie. Jest to również dystrybucyjna
druga pochodna ze zmienionym znakiem,
bez mas Diraca w punktach sklejenia.
Reakcja jest zerowa poza zbiorem kontaktu,
a na nim $v-u=v-\psi\ge0$ dla $v\in K_\psi$.
Całkowanie przez części (z gęstości funkcji
testowych w $H_0^1$) daje więc
$\int_{-1}^1u'(v'-u')\,\mathrm d\lambda_1(x)
=\int_{-1}^1(-u'')(v-u)\,\mathrm d\lambda_1(x)\ge0$.
Z jednoznaczności wynika wzór~\eqref{eq:pde29-explicit}.
Swobodny brzeg składa się z dwóch punktów
$\Gamma=\{-1/3,1/3\}$. Funkcja $u$ należy do
$C^{1,1}([-1,1])$, ale nie do $C^2([-1,1])$,
ponieważ jednostronne granice $u''$ w punkcie $1/3$
wynoszą $-8/3$ od lewej i $0$ od prawej.
W punkcie $-1/3$ kolejność jest odwrotna.
Całkowita masa reakcji wynosi
\[
 \int_{-1}^1\mu\,\mathrm d\lambda_1
 =2\int_0^{1/3}(4-12x^2)\,\mathrm d\lambda_1(x)
 =\tfrac{64}{27}.
\]
Jest ona równa $u'(-1)-u'(1)$, zgodnie
z całkowaniem równania $-u''=\mu$.
\end{przyklad}

\begin{center}
\begin{tikzpicture}[x=4.4cm,y=2.5cm]
 \draw[-{Latex}] (-1.06,0)--(1.08,0) node[right] {$x$};
 \draw[-{Latex}] (-1,-.12)--(-1,1.15) node[above] {$u,\psi$};
 \draw[dashed,gray] (-.3333,0)--(-.3333,1.07);
 \draw[dashed,gray] (.3333,0)--(.3333,1.07);
 \fill[blue!6] ({-1/3},0) rectangle ({1/3},1);
 \draw[blue!70!black,thick,samples=181,domain=-1:1]
   plot (\x,{(1-\x*\x)^2});
 \draw[red!75!black,very thick,samples=81,domain=-1/3:1/3]
   plot (\x,{(1-\x*\x)^2});
 \draw[red!75!black,very thick] (-1,0)--({-1/3},{64/81})
   ({1/3},{64/81})--(1,0);
 \node[below] at (-1,0) {$-1$};
 \node[below] at (-.3333,0) {$-\tfrac13$};
 \node[below] at (.3333,0) {$\tfrac13$};
 \node[below] at (1,0) {$1$};
 \node[blue!65!black,right] at (.58,.26) {$\psi$};
 \node[red!70!black,right] at (.62,.54) {$u$};
 \node[above] at (0,1.02) {zbiór kontaktu};
\end{tikzpicture}

{\small Rysunek: przeszkoda (niebieska) i rozwiązanie (czerwone)
z Przykładu~\ref{ex:pde29-explicit}. W środku zachodzi $u=\psi$;
poza tym przedziałem $-u''=0$.}
\end{center}

\begin{center}
\begin{tikzpicture}[x=4.4cm,y=.66cm,>=Latex,font=\small]
 \fill[orange!20] ({-1/3},0)
  -- plot[domain=-1/3:1/3,samples=81] (\x,{4-12*\x*\x})
  -- ({1/3},0)--cycle;
 \draw[->] (-1.06,0)--(1.08,0) node[right] {$x$};
 \draw[->] (-1,-.12)--(-1,4.7) node[above] {$\mu$};
 \draw[orange!85!black,very thick] (-1,0)--({-1/3},0)
    ({1/3},0)--(1,0);
 \draw[orange!85!black,very thick]
    plot[domain=-1/3:1/3,samples=81] (\x,{4-12*\x*\x});
 \foreach \a in {-1/3,1/3}{
   \draw[gray,dashed] ({\a},0)--({\a},{8/3});
   \draw[orange!85!black,fill=white] ({\a},0) circle[radius=1.5pt];
   \draw[orange!85!black,fill=white] ({\a},{8/3}) circle[radius=1.5pt];
 }
 \foreach \a/\lab in {-1/-1,-1/3/{-\frac13},1/3/{\frac13},1/1}
   \node[below] at ({\a},0) {$\lab$};
 \draw (-1,4)--++(-.025,0) node[left] {$4$};
 \draw (-1,{8/3})--++(-.025,0) node[left] {$8/3$};
 \node[above] at (0,4.15) {reakcja na zbiorze kontaktu};
\end{tikzpicture}

{\small Rysunek: gęstość reakcji $\mu=-u''$.
Przerywane odcinki zaznaczają skoki; wartości w pojedynczych
punktach swobodnego brzegu nie mają znaczenia dla klasy $L^2$.}
\end{center}

\section{Projekcja na zbiór wypukły i mnożniki ograniczeń}
Dla funkcjonału $f\in V'$ niech $w\in V$ będzie jedynym
rozwiązaniem $a(w,v)=\langle f,v\rangle$ dla $v\in V$
(twierdzenie Laxa--Milgrama). Wówczas
$J_f(v)=\frac12a(v-w,v-w)-\frac12a(w,w)$.
Rozwiązanie nierówności wariacyjnej jest zatem rzutem $w$
na $K_\psi$ w metryce określonej przez iloczyn skalarny $a$.
Nie jest to na ogół rzut punktowy $\max\{w,\psi\}$:
energia kontroluje pochodne, więc zmiana w jednym miejscu
wpływa na rozwiązanie w całym obszarze.

Dla domkniętego wypukłego zbioru $K\subset V$ definiujemy
stożek normalny
$N_K(u)=\{\eta\in V':\langle\eta,v-u\rangle\le0
\text{ dla każdego }v\in K\}$, jeżeli $u\in K$,
i $N_K(u)=\varnothing$ w przeciwnym przypadku.
Jeżeli $\mathcal A:V\to V'$ jest operatorem odpowiadającym $a$,
to nierówność~\eqref{eq:pde29-VI} jest równoważna
$0\in\mathcal Au-f+N_{K_\psi}(u)$.
Ta postać pozwala przenieść argumenty na inne problemy
jednostronne, także na przeszkody ruchome oraz przepływy
gradientowe z ograniczeniami.

\section{Schemat numeryczny i warunki punktowe}
W wymiarze jeden, dla $\Omega=(-1,1)$ i równomiernej siatki
$x_i=-1+ih$, $h=2/N$, $i=0,\ldots,N$, oznaczamy
$U_i\approx u(x_i)$ i $\Psi_i=\psi(x_i)$.
Dla $f=0$ dyskretna energia jest równa
$E_h(U)=\frac1{2h}\sum_{i=0}^{N-1}(U_{i+1}-U_i)^2$
przy $U_0=U_N=0$. Minimalizujemy ją na zbiorze
$U_i\ge\Psi_i$ dla $i=1,\dots,N-1$.
W węzłach wewnętrznych warunki optymalności mają postać
\begin{equation}\label{eq:pde29-discrete}
 U_i\ge\Psi_i,\qquad
 (2U_i-U_{i-1}-U_{i+1})/h^2\ge0,\qquad
 (U_i-\Psi_i)(2U_i-U_{i-1}-U_{i+1})=0.
\end{equation}
Dla przeszkody z przykładu~\ref{ex:pde29-explicit}
można zastosować metodę Gaussa--Seidela z projekcją,
zastępując każdy nowy węzeł przez większą z wartości
$\Psi_i$ i średniej sąsiednich wartości. Sprawdzanie
nierówności~\eqref{eq:pde29-discrete} jest ważniejsze niż
sam obraz krzywej: wykres może wyglądać poprawnie również
wtedy, gdy warunek komplementarności jest naruszony.

\section{Dalsze pytania i granice wykazanej teorii}
Udowodniliśmy istnienie i jednoznaczność nierówności
wariacyjnej w $H_0^1$, a przy mocniejszych założeniach
regularność $H^2$ i nierówność Lewy'ego--Stampacchii.
Nie udowodniliśmy ogólnej regularności $C^{1,1}$
przy dowolnej gładkiej przeszkodzie ani regularności
swobodnego brzegu. Analiza jego punktów regularnych
i osobliwych wymaga dodatkowych oszacowań wzrostu,
monotoniczności i argumentów geometrycznych;
nie wynika z samego sformułowania wariacyjnego.



\chapter[Homogenizacja okresowa równań eliptycznych]{Homogenizacja równań eliptycznych: zadanie komórkowe, macierz efektywna i granica oscylacyjna}
\label{chap:pde30}

\section{Dlaczego współczynników nie można po prostu uśrednić?}
Rozważmy ograniczony obszar Lipschitzowski $\Omega\subset\R^d$, $d\ge1$,
liczbę $\eps>0$ i macierz $A(y)$, która zmienia się w skali
pojedynczej komórki $Y=(0,1)^d$. W zagadnieniu Dirichleta
\begin{equation}\label{eq:pde30-micro}
 -\diver\bigl(A(x/\eps)\nabla u_\eps\bigr)=f\quad\text{w }\Omega,
 \qquad u_\eps\in H_0^1(\Omega),
\end{equation}
rozwiązanie odczuwa coraz drobniejsze niejednorodności. Pytamy,
czy istnieje granica $u_0$ i czy spełnia równanie tej samej postaci,
z pewną stałą macierzą $A^{\mathrm{hom}}$. Samo zastąpienie $A$
jego średnią arytmetyczną jest na ogół błędne: w kierunku
prostopadłym do warstw o przepływie decyduje średnia harmoniczna.

Przyjmujemy okres komórki równy $1$ w każdym kierunku.
Jest to inna normalizacja geometryczna niż dla torusa
$\T^d=\R^d/(2\pi\mathbb Z)^d$ z rozdziału~\ref{chap:pde25};
przy porównywaniu wzorów trzeba przeskalować zmienne.
Pracujemy na rzeczywistych przestrzeniach funkcyjnych.
Przez $H^{-1}(\Omega)$ rozumiemy przestrzeń dualną
do $H_0^1(\Omega)$ z normą względem normy gradientu
$\|v\|_{H_0^1}=\|\nabla v\|_2$,
równoważnej normie Sobolewa na mocy
Wniosku~\ref{cor:pde10-Poincare}.
Przez $H^1_{\mathrm{per}}(Y)$ rozumiemy przestrzeń ograniczeń do
$Y$ funkcji z $H^1_{\mathrm{loc}}(\R^d)$, których przedłużenie jest
$\mathbb Z^d$-okresowe. Jej podprzestrzenią o średniej zero jest
$H^1_{\mathrm{per},0}(Y)=\{v\in H^1_{\mathrm{per}}(Y):
\int_Yv\,\mathrm d\lambda_d(y)=0\}$.
Niech $A\in L^\infty_{\mathrm{per}}(Y;\R^{d\times d})$ będzie
\emph{symetryczna} oraz niech $0<\lambda\le\Lambda$ spełniają
\begin{equation}\label{eq:pde30-ellipticity}
 \lambda|\xi|^2\le\xi^{\mathsf T}A(y)\xi\le\Lambda|\xi|^2
 \quad (\xi\in\R^d),\quad\text{dla prawie każdego }y\in Y.
\end{equation}
Macierz przedłużamy okresowo na $\R^d$.
Przeskalowanie $x\mapsto x/\eps$ zachowuje zbiory
miary zero w tym sensie, że ich przeciwobrazy
również mają miarę zero. Z symetrii i warunku
~\eqref{eq:pde30-ellipticity} wynika
$\|A(y)\|_{\mathrm{op}}\le\Lambda$ prawie wszędzie.
Definiujemy na $H_0^1(\Omega)$ formę
\[
 a_\eps(u,v)=\int_\Omega A(x/\eps)\nabla u\cdot\nabla v\,
 \mathrm d\lambda_d(x).
\]
Mamy $|a_\eps(u,v)|\le\Lambda\|\nabla u\|_2\|\nabla v\|_2$
oraz $a_\eps(v,v)\ge\lambda\|\nabla v\|_2^2$.
Z twierdzenia Laxa--Milgrama dla każdego $f\in H^{-1}(\Omega)$
istnieje dokładnie jedno $u_\eps\in H_0^1(\Omega)$
spełniające $a_\eps(u_\eps,v)=\langle f,v\rangle$
dla każdego $v\in H_0^1(\Omega)$.
Jest to równoważne równaniu~\eqref{eq:pde30-micro}
w sensie dystrybucji przez gęstość funkcji testowych.
Testując przez $v=u_\eps$, dostajemy
$\lambda\|\nabla u_\eps\|_2^2\le
\|f\|_{H^{-1}}\|\nabla u_\eps\|_2$,
a stąd $\|\nabla u_\eps\|_2\le\lambda^{-1}\|f\|_{H^{-1}}$
(również gdy gradient znika).
Zapis $\diver$ oznacza dywergencję, a $\mathrm d\lambda_d$ miarę Lebesgue'a.

\section{Przypadek jednowymiarowy: dokładny wzór i tempo zbieżności}
Niech $\Omega=(0,L)$, $L>0$, $a\in L^\infty_{\mathrm{per}}(0,1)$
oraz $0<\lambda\le a\le\Lambda$ prawie wszędzie. Definiujemy
$b=1/a$ i $\bar b=\int_0^1b(y)\,\mathrm d\lambda_1(y)$, a także
$a_{\mathrm{hom}}=\bar b^{-1}$. Dla $f\in L^2(0,L)$ połóżmy
$F(x)=\int_0^xf(t)\,\mathrm d\lambda_1(t)$.
Mamy $1/\Lambda\le b\le1/\lambda$, więc
$\lambda\le a_{\mathrm{hom}}\le\Lambda$.
Funkcja $F$ należy do $H^1(0,L)$ i ma ciągły,
absolutnie ciągły reprezentant, $F'=f$ prawie wszędzie,
oraz $\|F\|_\infty\le\sqrt L\|f\|_2$.

\begin{twierdzenie}[Jednowymiarowe równanie uśrednione]\label{thm:pde30-1d}
Dla każdego $\eps>0$ istnieje dokładnie jedno $u_\eps\in H_0^1(0,L)$
spełniające $-(a(x/\eps)u_\eps')'=f$ w sensie dystrybucji.
Niech $u_0$ będzie jedynym rozwiązaniem
$-a_{\mathrm{hom}}u_0''=f$ w $H_0^1(0,L)$.
Wtedy, gdy $\eps\downarrow0$,
$u_\eps\to u_0$ w $L^\infty(0,L)$ oraz
$u_\eps\rightharpoonup u_0$ w $H_0^1(0,L)$, a strumienie
$a(x/\eps)u_\eps'$ zbiegają silnie w $L^2(0,L)$ do
$a_{\mathrm{hom}}u_0'$. Istnieje stała $C$, zależna tylko od
$L,\lambda,\Lambda$, taka że dla $0<\eps\le1$
\begin{equation}\label{eq:pde30-1d-rate}
 \|u_\eps-u_0\|_{L^\infty(0,L)}
 +\|a(x/\eps)u_\eps'-a_{\mathrm{hom}}u_0'\|_{L^2(0,L)}
 \le C\eps\|f\|_{L^2(0,L)}.
\end{equation}
\end{twierdzenie}
\begin{proof}
Istnienie i jednoznaczność wynikają z poprzedniej sekcji.
Ponieważ $q_\eps=a(x/\eps)u_\eps'\in L^2(0,L)$
i $q_\eps'=-f\in L^2(0,L)$ w sensie dystrybucji,
mamy $q_\eps\in H^1(0,L)$. Z jednowymiarowej
charakterystyki Sobolewa (Twierdzenie~\ref{thm:pde9-ac})
otrzymujemy absolutnie ciągły reprezentant
$q_\eps(x)=C_\eps-F(x)$. Całkując $u_\eps'$ oraz
wykorzystując $u_\eps(L)=u_\eps(0)=0$, otrzymujemy
\begin{equation}\label{eq:pde30-explicit-1d}
 C_\eps=\frac{\int_0^L b(t/\eps)F(t)\,\mathrm d\lambda_1(t)}
 {\int_0^L b(t/\eps)\,\mathrm d\lambda_1(t)},\qquad
 u_\eps(x)=\int_0^xb(t/\eps)(C_\eps-F(t))\,\mathrm d\lambda_1(t).
\end{equation}
Odwrotnie, wzór~\eqref{eq:pde30-explicit-1d}
określa funkcję absolutnie ciągłą z pochodną
$b(x/\eps)(C_\eps-F(x))\in L^2$, zerowymi wartościami
na końcach i strumieniem $C_\eps-F$; jest zatem
rozwiązaniem zagadnienia. W rozwiązaniu granicznym
analogicznie $u_0'=\bar b(C_0-F)$,
gdzie $C_0=L^{-1}\int_0^LF(t)\,\mathrm d\lambda_1(t)$.
W szczególności $u_0\in H^2(0,L)$ i
$u_0''=-\bar b f$.

Wyjaśnijmy ilościowo uśrednianie. Jeśli $g\in L^\infty_{\mathrm{per}}(0,1)$
ma średnią zero, to jej funkcja pierwotna
$G(s)=\int_0^sg(y)\,\mathrm d\lambda_1(y)$ jest ograniczona
i okresowa, ponieważ
$G(s+1)-G(s)=\int_s^{s+1}g\,\mathrm d\lambda_1=0$.
Ponadto $\|G\|_\infty\le\|g\|_{L^1(0,1)}$.
Dla $h\in W^{1,1}(0,L)$ stosujemy regułę iloczynu
dla absolutnie ciągłych reprezentantów oraz
$(\eps G(x/\eps))'=g(x/\eps)$ prawie wszędzie:
\[
 \int_0^Lg(x/\eps)h(x)\,\mathrm d\lambda_1(x)
 =\eps[G(x/\eps)h(x)]_0^L
 -\eps\int_0^LG(x/\eps)h'(x)\,\mathrm d\lambda_1(x).
\]
Stąd
\begin{equation}\label{eq:pde30-average-rate}
 \left|\int_0^Lg(x/\eps)h(x)\,\mathrm d\lambda_1(x)\right|
 \le\eps\|G\|_\infty
 \bigl(|h(0)|+|h(L)|+\|h'\|_{L^1(0,L)}\bigr).
\end{equation}
Połóżmy $g=b-\bar b$ oraz $h=F-C_0$.
Mamy $\int_0^L h\,\mathrm d\lambda_1=0$,
$|C_0|\le\|F\|_\infty$ i
$\|h\|_\infty\le2\sqrt L\|f\|_2$,
$\|h'\|_1\le\sqrt L\|f\|_2$.
Odejmując $C_0$ od ilorazu
~\eqref{eq:pde30-explicit-1d}, dostajemy
\[
 C_\eps-C_0=
 \frac{\int_0^L g(t/\eps)h(t)\,\mathrm d\lambda_1(t)}
 {\int_0^Lb(t/\eps)\,\mathrm d\lambda_1(t)}.
\]
Mianownik jest nie mniejszy od $L/\Lambda$,
więc~\eqref{eq:pde30-average-rate} daje
$|C_\eps-C_0|\le C\eps\|f\|_2$.
Stała zależy wyłącznie od $L,\lambda,\Lambda$,
gdyż $\|G\|_\infty\le\|b-\bar b\|_1\le2/\lambda$.
Następnie odejmujemy wzory na pochodne i całkujemy:
\[
 \begin{aligned}
 u_\eps(x)-u_0(x)
 &=(C_\eps-C_0)\int_0^x b(t/\eps)\,\mathrm d\lambda_1(t)\\
 &\quad+\int_0^x g(t/\eps)(C_0-F(t))\,\mathrm d\lambda_1(t).
 \end{aligned}
\]
Pierwszy składnik jest ograniczony przez
$(L/\lambda)|C_\eps-C_0|$, a drugi przez
$\eps\|G\|_\infty(2\|h\|_\infty+\|h'\|_1)$
na mocy~\eqref{eq:pde30-average-rate} na $(0,x)$.
Otrzymujemy oszacowanie jednostajne na $[0,L]$.
Ponieważ $q_\eps-q_0=C_\eps-C_0$, norma
$L^2$ różnicy strumieni jest równa
$\sqrt L|C_\eps-C_0|$; w istocie zachodzi
także zbieżność strumieni w normie jednostajnej.
To dowodzi~\eqref{eq:pde30-1d-rate}.
Jednostajne ograniczenie gradientów i refleksywność
$H_0^1$ pozwalają z dowolnego ciągu $\eps\to0$
wydzielić słabo zbieżny podciąg. Ciągłe zanurzenie
w $L^2$ i wykazana zbieżność jednostajna identyfikują
jego słabą granicę z $u_0$. Ponieważ każdy ciąg
ma taki podciąg, cała rodzina zbiega słabo do $u_0$
(testując przez dowolny ciągły funkcjonał liniowy).
\end{proof}

\begin{stwierdzenie}[Poprawka do gradientu w wymiarze jeden]\label{prop:pde30-corrector1d}
Przy założeniach twierdzenia~\ref{thm:pde30-1d} istnieje jedyna
funkcja $\chi\in W^{1,\infty}_{\mathrm{per}}(0,1)$ o średniej zero,
której pochodna spełnia $\chi'(y)=a_{\mathrm{hom}}/a(y)-1$.
Ponadto
\[
\left\|u_\eps'-\bigl(1+\chi'(x/\eps)\bigr)u_0'\right\|_{L^2(0,L)}
\le C\eps\|f\|_2.
\]
Jeśli $a$ nie jest prawie wszędzie stała i $f\not\equiv0$,
to nie zachodzi silna zbieżność $u_\eps'\to u_0'$.
\end{stwierdzenie}
\begin{proof}
Funkcja $g_1=a_{\mathrm{hom}}b-1$ należy do
$L^\infty_{\mathrm{per}}(0,1)$ i ma średnią zero,
bo $a_{\mathrm{hom}}\bar b=1$.
Jej funkcja pierwotna $G_1(y)=\int_0^y g_1\,\mathrm d\lambda_1$
jest okresowa i lipschitzowska. Definiujemy
$\chi=G_1-\int_0^1G_1\,\mathrm d\lambda_1$.
Dwie funkcje o wymaganej pochodnej różnią się stałą,
a średnia zero wymusza jej znikanie.
Z tożsamości $u_0'=\bar b(C_0-F)$ wynika
$(1+\chi'(x/\eps))u_0'=b(x/\eps)(C_0-F)$.
Różnica z $u_\eps'$ to $b(x/\eps)(C_\eps-C_0)$,
więc twierdzenie~\ref{thm:pde30-1d} daje oszacowanie.
Jeśli $a$ nie jest stała, to średnia z
$(a_{\mathrm{hom}}/a-1)^2$ jest dodatnia.
Przy $f\not\equiv0$ mamy $u_0'\not\equiv0$,
bo $u_0''=-\bar b f$.
Ponieważ $u_0'\in H^1(0,L)\subset L^\infty(0,L)$,
waga $|u_0'|^2$ należy do $W^{1,1}(0,L)$
z pochodną $2u_0'u_0''$.
Stosujemy~\eqref{eq:pde30-average-rate} do
$g_1^2-\int_0^1g_1^2\,\mathrm d\lambda_1$
oraz tej wagi. Różnica
$u_\eps'-u_0'-g_1(x/\eps)u_0'$ dąży do zera w $L^2$
z już wykazanego oszacowania poprawki,
a normy $g_1(x/\eps)u_0'$ są jednostajnie ograniczone.
Zatem różnica kwadratów norm również dąży do zera,
co daje
\[
\lim_{\eps\downarrow0}\|u_\eps'-u_0'\|_2^2=
\left(\int_0^1|a_{\mathrm{hom}}/a(y)-1|^2
\,\mathrm d\lambda_1(y)\right)\|u_0'\|_2^2>0.
\]
Istotne są więc poprawki oscylacyjne, a nie tylko granica $u_0$.
\end{proof}

\begin{przyklad}[Dlaczego średnia arytmetyczna jest błędna]
Niech na połowie komórki $a=1$, a na drugiej połowie $a=4$.
Wówczas średnia arytmetyczna wynosi $5/2$, a efektywny
współczynnik to $a_{\mathrm{hom}}=(\tfrac12+\tfrac18)^{-1}=8/5$.
Dla $f=1$ na $(0,1)$ rozwiązanie graniczne ma postać
$u_0(x)=x(1-x)/(2a_{\mathrm{hom}})$.
Zastąpienie $a_{\mathrm{hom}}$ przez $5/2$ daje
$u_{\mathrm{ar}}(x)=x(1-x)/5$ zamiast
$u_0(x)=5x(1-x)/16$. Różnica jest równa
$9x(1-x)/80$ i ma normę jednostajną $9/320$;
nie znika mimo dowolnie drobnych warstw.
Dla tych dwóch faz $g_1$ przyjmuje wartości
$3/5$ i $-3/5$, więc
$\int_0^1g_1^2\,\mathrm d\lambda_1=9/25$.
Ponieważ $u_0'(x)=5(1-2x)/16$,
$\|u_0'\|_2^2=25/768$. Ze Stwierdzenia
~\ref{prop:pde30-corrector1d} otrzymujemy
$\lim_{\eps\downarrow0}\|u_\eps'-u_0'\|_2^2=3/256$.
\end{przyklad}

\section{Uśrednianie funkcji okresowych}
Do przejścia w wielu wymiarach potrzebujemy prostego lematu.
\begin{lemat}[Słaba granica funkcji okresowych]\label{lem:pde30-average}
Niech $D\subset\R^d$ będzie ograniczonym zbiorem otwartym,
a $g\in L^2_{\mathrm{per}}(Y)$ (czyli $g$ jest okresowym
przedłużeniem funkcji z $L^2(Y)$). Wtedy
$g(x/\eps)\rightharpoonup\bar g:=\int_Yg(y)\,\mathrm d\lambda_d(y)$
słabo w $L^2(D)$, gdy $\eps\downarrow0$.
Ponadto $\eps g(x/\eps)\to0$ silnie w $L^2(D)$.
\end{lemat}
\begin{proof}
Wybierzmy $R\ge1$ takie, że $D\subset(-R,R)^d$.
Dla $0<\eps\le1$ co najwyżej $(2R/\eps+2)^d$
komórek $\eps(k+Y)$, $k\in\mathbb Z^d$, przecina
ten sześcian. Ich brzegi mają miarę zero.
Przez podstawienie $x=\eps(k+y)$ i okresowość
całka na każdej pełnej komórce wynosi
$\eps^d\|g\|_{L^2(Y)}^2$. W szczególności
\[
 \|g(\cdot/\eps)\|_{L^2(D)}
 \le (2R+2)^{d/2}\|g\|_{L^2(Y)}.
\]
To daje również silną zbieżność $\eps g(\cdot/\eps)$
do zera w $L^2(D)$.

Wykażemy słabą zbieżność bez założeń o regularności brzegu.
Niech $h=g-\bar g$ i $\varphi\in C_c^\infty(D)$,
przedłużona przez zero na $\R^d$. Z zerowej średniej
$h$ na każdej przeskalowanej komórce wynika
\[
 \begin{aligned}
 \int_D h(x/\eps)\varphi(x)\,\mathrm d\lambda_d(x)
 =\sum_{k\in\mathbb Z^d}\eps^d\int_Y h(y)
 \bigl(\varphi(\eps(k+y))-\varphi(\eps k)\bigr)
 \,\mathrm d\lambda_d(y).
 \end{aligned}
\]
Wystarczy sumować po komórkach, których domknięcia
przecinają nośnik $\varphi$; pozostałe składniki są zerowe.
Ich liczba jest nie większa niż $(2R/\eps+2)^d$.
Różnica wartości funkcji testowej jest ograniczona
przez $\eps\sqrt d\|\nabla\varphi\|_\infty$.
Moduł całej sumy jest więc nie większy niż
\[
 (2R+2)^d\eps\sqrt d\|\nabla\varphi\|_\infty
 \|h\|_{L^1(Y)},
\]
co dąży do zera. Dla dowolnego testu $\varphi\in L^2(D)$
wybieramy $\varphi_n\in C_c^\infty(D)$ zbieżne
do niego w $L^2$. Błąd zastąpienia testu przez
$\varphi_n$ jest ograniczony przez
$(2R+2)^{d/2}\|h\|_{L^2(Y)}\|\varphi-\varphi_n\|_2$,
jednostajnie względem $\eps$. Najpierw przechodzimy
do granicy z $\eps$, potem z $n$.
Otrzymujemy tezę o słabej zbieżności.
\end{proof}

\section{Zadanie komórkowe i macierz efektywna}

\begin{lemat}[Okresowa przestrzeń Sobolewa]
\label{lem:pde30-periodic-H1}
Przestrzeń $H^1_{\mathrm{per}}(Y)$ jest domkniętą
podprzestrzenią $H^1(Y)$, a gładkie funkcje okresowe
są w niej gęste. Dla każdego $v\in H^1_{\mathrm{per}}(Y)$
zachodzi $\int_Y\nabla v\,\mathrm d\lambda_d=0$.
Podprzestrzeń $H^1_{\mathrm{per},0}(Y)$ jest przestrzenią
Hilberta z normą $\|\nabla v\|_{L^2(Y)}$.
\end{lemat}
\begin{proof}
Niech $v_n\in H^1_{\mathrm{per}}(Y)$ oraz
$v_n\to v$ w $H^1(Y)$.
Oznaczmy okresowe przedłużenia przez $\widetilde v_n$.
Dla dodatniej liczby całkowitej $m$ sześcian
$B_m=(-m,m)^d$ składa się, z dokładnością
do zbioru miary zero, z $(2m)^d$ komórek jednostkowych.
Okresowość funkcji i ich słabych pochodnych daje
\[
 \|\widetilde v_n-\widetilde v_l\|_{H^1(B_m)}^2
 =(2m)^d\|v_n-v_l\|_{H^1(Y)}^2.
\]
Ciąg jest Cauchy'ego w każdej z przestrzeni $H^1(B_m)$.
Z ich zupełności otrzymujemy granice zgodne na
przecięciach (przez ciągłość ograniczenia).
Definiują one funkcję $\widetilde v\in H^1_{\mathrm{loc}}(\R^d)$,
która ogranicza się do $v$ na $Y$.
Dla każdego $k\in\mathbb Z^d$ mamy
$\widetilde v(\cdot+k)=\widetilde v$ prawie wszędzie:
na dowolnym ograniczonym sześcianie obie strony
są granicami w $L^2$ odpowiednich przesunięć
$\widetilde v_n$. Zatem $v\in H^1_{\mathrm{per}}(Y)$.

Dla ustalonego $v\in H^1_{\mathrm{per}}(Y)$
wygładzamy $\widetilde v$ przez splot z gładkim
jądrem o zwartym nośniku dążącym do zera.
Sploty są okresowe, ponieważ przesunięcie o
$k\in\mathbb Z^d$ nie zmienia funkcji pod całką.
Zbieżność w $H^1(Y)$ wynika z lokalnego
twierdzenia o aproksymacji przez splot:
można pomnożyć $\widetilde v$ przez gładką
funkcję o zwartym nośniku równą $1$ w otoczeniu
$\overline Y$, gdyż dla małych jąder nie zmienia
to splotu na $Y$. Dla gładkich funkcji okresowych
całka każdej pochodnej po $Y$ znika, dzięki
twierdzeniu Newtona--Leibniza w odpowiedniej
współrzędnej i równości wartości na przeciwległych
ścianach. Przechodząc do granicy w $H^1(Y)$,
otrzymujemy $\int_Y\nabla v\,\mathrm d\lambda_d=0$.

Warunek średniej zero definiuje domkniętą
podprzestrzeń, gdyż $\lambda_d(Y)=1$ i całka
jest ciągła w $L^2(Y)$. Z nierówności
Poincarégo--Wirtingera na kostce
(Twierdzenie~\ref{thm:pde9-cube-Poincare})
mamy $\|v\|_2\le\sqrt d\|\nabla v\|_2$ dla
$v$ o średniej zero. Norma gradientu jest więc
równoważna normie $H^1$ i pochodzi od iloczynu
skalarnego $(v,w)\mapsto\int_Y\nabla v\cdot\nabla w\,
\mathrm d\lambda_d$. Domkniętość zapewnia jej zupełność.
\end{proof}

\begin{twierdzenie}[Istnienie i jednoznaczność poprawek komórkowych]\label{thm:pde30-cell}
Niech $A$ spełnia~\eqref{eq:pde30-ellipticity}.
Dla każdego $j\in\{1,\ldots,d\}$ istnieje dokładnie jedna
funkcja $\chi_j\in H^1_{\mathrm{per},0}(Y)$ taka, że
\begin{equation}\label{eq:pde30-cell}
\int_Y A(y)(e_j+\nabla_y\chi_j(y))\cdot\nabla_y\varphi(y)
\,\mathrm d\lambda_d(y)=0
\quad\text{dla każdego }\varphi\in H^1_{\mathrm{per},0}(Y).
\end{equation}
Definiujemy wektory strumienia $Q_j(y)=A(y)(e_j+\nabla\chi_j(y))$
i macierz stałą
\begin{equation}\label{eq:pde30-ahom}
 A^{\mathrm{hom}}e_j=\int_YQ_j(y)\,\mathrm d\lambda_d(y).
\end{equation}
Wtedy $\diver_yQ_j=0$ w sensie dystrybucji na torusie,
$A^{\mathrm{hom}}$ jest symetryczna i
$\lambda|\xi|^2\le\xi^{\mathsf T}A^{\mathrm{hom}}\xi
\le\Lambda|\xi|^2$ dla każdego $\xi\in\R^d$.
\end{twierdzenie}
\begin{proof}
Z Lematu~\ref{lem:pde30-periodic-H1}
przestrzeń $H^1_{\mathrm{per},0}(Y)$ z normą
gradientu jest przestrzenią Hilberta.
Dla $v,\varphi$ z tej przestrzeni forma
$b(v,\varphi)=\int_Y A\nabla v\cdot\nabla\varphi\,
\mathrm d\lambda_d$ spełnia
$|b(v,\varphi)|\le\Lambda\|\nabla v\|_2\|\nabla\varphi\|_2$
oraz $b(v,v)\ge\lambda\|\nabla v\|_2^2$.
Funkcjonał
$\varphi\mapsto-\int_YAe_j\cdot\nabla\varphi\,
\mathrm d\lambda_d$ ma normę co najwyżej $\Lambda$,
bo $\lambda_d(Y)=1$. Twierdzenie Laxa--Milgrama
daje jedyne $\chi_j$, a test $\chi_j$ daje
$\|\nabla\chi_j\|_2\le\Lambda/\lambda$.
W szczególności $Q_j\in L^2_{\mathrm{per}}(Y;\R^d)$.

Równanie~\eqref{eq:pde30-cell} zachodzi dla wszystkich
testów z $H^1_{\mathrm{per}}(Y)$, gdyż odjęcie średniej
nie zmienia gradientu. Stąd $\diver_yQ_j=0$
w sensie dystrybucji na torusie jednostkowym.
Zachodzi to również na $\R^d$ dla okresowego
przedłużenia: dla $\varphi\in C_c^\infty(\R^d)$
funkcja $\Phi(y)=\sum_{k\in\mathbb Z^d}\varphi(y+k)$
jest gładka i okresowa, a suma jest lokalnie skończona.
Z okresowości i rozbicia $\R^d$ na komórki wynika
\[
 \int_{\R^d}Q_j\cdot\nabla\varphi\,\mathrm d\lambda_d
 =\int_YQ_j\cdot\nabla\Phi\,\mathrm d\lambda_d=0.
\]

Dla $\xi\in\R^d$ niech $\chi_\xi=\sum_j\xi_j\chi_j$.
Z liniowości równań komórkowych
$\int_Y A(\xi+\nabla\chi_\xi)\cdot\nabla\varphi\,
\mathrm d\lambda_d=0$ dla każdego testu okresowego.
Z definicji macierzy efektywnej mamy
$A^{\mathrm{hom}}\xi=\int_Y A(\xi+\nabla\chi_\xi)\,
\mathrm d\lambda_d$.
Dodając do iloczynu skalarnego z $\xi$ zerową
całkę odpowiadającą testowi $\varphi=\chi_\xi$, otrzymujemy
\begin{equation}\label{eq:pde30-energycell}
 \xi^{\mathsf T}A^{\mathrm{hom}}\xi
 =\int_Y A(y)(\xi+\nabla\chi_\xi)\cdot
             (\xi+\nabla\chi_\xi)\,\mathrm d\lambda_d(y)
 =\min_{v\in H^1_{\mathrm{per},0}(Y)}
 \int_Y (\xi+\nabla v)^{\mathsf T}A(y)(\xi+\nabla v)\,\mathrm d\lambda_d(y).
\end{equation}
Aby uzasadnić drugą równość, oznaczmy energię
komórkową przez $\mathcal E_\xi(v)$.
Dla $v\in H^1_{\mathrm{per},0}(Y)$ rozwijamy kwadrat:
\[
 \mathcal E_\xi(v)-\mathcal E_\xi(\chi_\xi)
 =\int_Y A\nabla(v-\chi_\xi)\cdot\nabla(v-\chi_\xi)
 \,\mathrm d\lambda_d\ge0,
\]
gdyż wyraz mieszany znika na mocy
~\eqref{eq:pde30-cell} z testem $v-\chi_\xi$.
Równość wymusza $\nabla(v-\chi_\xi)=0$,
a warunek średniej zero i nierówność Poincarégo
dają $v=\chi_\xi$.
Ponieważ $\int_Y\nabla v\,\mathrm d\lambda_d(y)=0$,
mamy $\int_Y|\xi+\nabla v|^2\,\mathrm d\lambda_d(y)
=|\xi|^2+\int_Y|\nabla v|^2\,\mathrm d\lambda_d(y)\ge|\xi|^2$.
Stosując tę tożsamość dla $v=\chi_\xi$ i dolną
granicę eliptyczności $A$, otrzymujemy
$\xi^{\mathsf T}A^{\mathrm{hom}}\xi\ge\lambda|\xi|^2$.
Dla górnej granicy wybieramy $v=0$ w problemie
minimalizacyjnym, co daje
$\xi^{\mathsf T}A^{\mathrm{hom}}\xi\le
\int_Y\xi^{\mathsf T}A(y)\xi\,\mathrm d\lambda_d
\le\Lambda|\xi|^2$. Symetrię otrzymujemy z tożsamości
$e_i\cdot A^{\mathrm{hom}}e_j
=\int_Y(e_i+\nabla\chi_i)\cdot A(e_j+\nabla\chi_j)
\,\mathrm d\lambda_d(y)$, wynikającej z testu $\chi_i$
w zadaniu komórkowym. Prawa strona jest symetryczna
względem $i,j$, gdyż $A=A^{\mathsf T}$.
\end{proof}

\begin{wniosek}[Ograniczenia harmoniczne i arytmetyczne]
\label{cor:pde30-harmonic-arithmetic}
W założeniach Twierdzenia~\ref{thm:pde30-cell}
dla każdej $\xi\in\R^d$ zachodzi
\[
 \xi^{\mathsf T}\left(\int_Y A(y)^{-1}
 \,\mathrm d\lambda_d(y)\right)^{-1}\xi
 \le\xi^{\mathsf T}A^{\mathrm{hom}}\xi
 \le\xi^{\mathsf T}\left(\int_Y A(y)
 \,\mathrm d\lambda_d(y)\right)\xi .
\]
\end{wniosek}
\begin{proof}
Górna nierówność wynika z wyboru $v=0$
w~\eqref{eq:pde30-energycell}.
Niech $B=\int_Y A^{-1}\,\mathrm d\lambda_d$.
Macierz $A^{-1}$ jest mierzalna (odwracanie
jest ciągłe na zbiorze macierzy dodatnio określonych)
i spełnia oszacowania $\Lambda^{-1}I\le A^{-1}\le\lambda^{-1}I$
w sensie form kwadratowych. Zatem $B$ jest
symetryczna i dodatnio określona.
Dla dowolnego $v\in H^1_{\mathrm{per},0}(Y)$,
$z=\xi+\nabla v$ oraz stałego $\eta\in\R^d$
nieujemność formy daje
\[
 0\le\int_Y (z-A^{-1}\eta)^{\mathsf T}
 A(z-A^{-1}\eta)\,\mathrm d\lambda_d
 =\mathcal E_\xi(v)-2\eta\cdot\xi+\eta^{\mathsf T}B\eta.
\]
Użyliśmy $\int_Y z\,\mathrm d\lambda_d=\xi$.
Wybieramy $\eta=B^{-1}\xi$, a następnie
$v=\chi_\xi$. Z~\eqref{eq:pde30-energycell}
otrzymujemy dolną nierówność.
Dla $d=1$ dolna granica jest osiągana dzięki
poprawce ze Stwierdzenia~\ref{prop:pde30-corrector1d}.
\end{proof}

\begin{uwaga}[Co jest mikroskopowe, a co makroskopowe?]
Funkcja $\chi_j$ jest okresową poprawką komórkową i zależy
od zmiennej szybkiej $y$. Macierz $A^{\mathrm{hom}}$ jest stała
w zmiennej wolnej $x$. Współczynniki $A(x/\eps)$ nie muszą
zbiegać punktowo ani silnie do $A^{\mathrm{hom}}$; zbiega
rozwiązanie i odpowiednio rozumiany strumień. Nawet słaba
granica $A(x/\eps)$ jest na ogół inna niż $A^{\mathrm{hom}}$.
\end{uwaga}

\section{Twierdzenie o homogenizacji: pełny dowód przez oscylujące testy}
\begin{twierdzenie}[Homogenizacja okresowa przy mierzalnych współczynnikach]
\label{thm:pde30-homogenization}
Niech $d\ge1$, $\Omega\subset\R^d$ będzie ograniczonym
obszarem Lipschitzowskim, a $A$ spełnia~\eqref{eq:pde30-ellipticity}.
Dla $f\in H^{-1}(\Omega)$ niech $u_\eps\in H_0^1(\Omega)$
będzie jedynym słabym rozwiązaniem~\eqref{eq:pde30-micro}.
Niech $u_0\in H_0^1(\Omega)$ będzie jedynym słabym
rozwiązaniem zagadnienia
\begin{equation}\label{eq:pde30-macro}
 -\diver(A^{\mathrm{hom}}\nabla u_0)=f
 \quad\text{w }\Omega,\qquad u_0|_{\partial\Omega}=0.
\end{equation}
Wówczas dla $\eps\downarrow0$ zachodzą zbieżności
\[
 u_\eps\rightharpoonup u_0\text{ w }H_0^1(\Omega),\quad
 u_\eps\to u_0\text{ w }L^2(\Omega),\quad
 A(x/\eps)\nabla u_\eps\rightharpoonup
 A^{\mathrm{hom}}\nabla u_0\text{ w }L^2(\Omega;\R^d).
\]
W szczególności twierdzenie obejmuje obciążenia
$f\in L^2(\Omega)$, utożsamiane z funkcjonałami
$v\mapsto\int_\Omega fv\,\mathrm d\lambda_d$.
\end{twierdzenie}
\begin{proof}
\textbf{Krok 1: zwartość i strumień.}
Z oszacowania energetycznego
$\|\nabla u_\eps\|_2\le\lambda^{-1}\|f\|_{H^{-1}}$
oraz nierówności Poincarégo rodzina $(u_\eps)$
jest ograniczona w $H_0^1(\Omega)$.
Dla dowolnego ciągu dodatnich $\eps_k\to0$
z refleksywności przestrzeni Hilberta wybieramy
słabo zbieżny podciąg w $H_0^1$.
Z twierdzenia Rellicha~\ref{thm:pde10-Rellich-p}
po dalszym wydzieleniu $u_{\eps_k}\to u$ silnie
w $L^2$. Granice są zgodne, ponieważ zanurzenie
$H_0^1\hookrightarrow L^2$ jest ciągłe.
Strumienie $q_\eps=A(x/\eps)\nabla u_\eps$ spełniają
$\|q_\eps\|_2\le(\Lambda/\lambda)\|f\|_{H^{-1}}$;
po dalszym wydzieleniu $q_\eps\rightharpoonup q$ w $L^2$.
Dla każdego $\varphi\in H_0^1(\Omega)$
mamy $\int_\Omega q_\eps\cdot\nabla\varphi\,
\mathrm d\lambda_d=\langle f,\varphi\rangle$.
Przechodząc do granicy, otrzymujemy tę samą
tożsamość dla $q$, czyli $\diver q=-f$.
W dalszym ciągu piszemy $\eps$ zamiast indeksów podciągu.

\textbf{Krok 2: oscylujące funkcje testowe.}
Dla $i\in\{1,\ldots,d\}$ zdefiniujmy
$w_i^\eps(x)=x_i+\eps\chi_i(x/\eps)$.
Reguła łańcuchowa dla przeskalowania w $H^1_{\mathrm{loc}}$
daje $\nabla w_i^\eps=e_i+(\nabla\chi_i)(x/\eps)$.
Można ją uzasadnić przez gładkie okresowe przybliżenia
z Lematu~\ref{lem:pde30-periodic-H1}: przy ustalonym
$\eps$ podstawienie zmiennych i pokrycie komórkami
zapewniają zbieżność funkcji oraz gradientów w $L^2(\Omega)$.
W szczególności $w_i^\eps\in H^1(\Omega)$.
Mamy
$Q_i^\eps(x)=Q_i(x/\eps)=A(x/\eps)\nabla w_i^\eps(x)$.
Z Lematu~\ref{lem:pde30-average} oraz zerowej średniej
$\nabla\chi_i$ otrzymujemy
\[
 w_i^\eps\to x_i\ \text{w }L^2(\Omega),
 \qquad \nabla w_i^\eps\rightharpoonup e_i\
 \text{w }L^2(\Omega;\R^d).
\]
Zatem $w_i^\eps\rightharpoonup x_i$ w $H^1(\Omega)$.
Ten sam lemat daje
$Q_i^\eps\rightharpoonup A^{\mathrm{hom}}e_i$ w $L^2$.
Ponadto $\diver Q_i^\eps=0$ w sensie dystrybucji:
z dowodu Twierdzenia~\ref{thm:pde30-cell} wiemy,
że $\diver_yQ_i=0$ na całej $\R^d$.
Dla $\varphi\in C_c^\infty(\Omega)$ przedłużonej
przez zero podstawiamy $x=\eps y$:
\[
 \int_\Omega Q_i^\eps\cdot\nabla_x\varphi\,
 \mathrm d\lambda_d(x)
 =\eps^{d-1}\int_{\R^d} Q_i(y)\cdot
 \nabla_y(\varphi(\eps y))\,\mathrm d\lambda_d(y)=0.
\]
Nie wykorzystujemy klasycznych pochodnych $A$
ani ograniczoności $\chi_i$.

\textbf{Krok 3: identyfikacja granicy iloczynu.}
Dla $\varphi\in C_c^\infty(\Omega)$ iloczyn
$\varphi w_i^\eps$ należy do $H_0^1(\Omega)$:
ma zwarty nośnik, należy do $H^1$ na mocy reguły
iloczynu i można go wygładzić wewnątrz $\Omega$.
Ponadto mnożenie przez $\varphi$ jest ograniczonym
operatorem $H^1(\Omega)\to H_0^1(\Omega)$,
więc $\varphi w_i^\eps\rightharpoonup\varphi x_i$
w $H_0^1$. Testując równanie przez $\varphi w_i^\eps$
i rozwijając gradient iloczynu, otrzymujemy
\begin{align*}
\int_\Omega\varphi\,q_\eps\cdot\nabla w_i^\eps
\,\mathrm d\lambda_d(x)
&=\langle f,\varphi w_i^\eps\rangle
 -\int_\Omega w_i^\eps q_\eps\cdot\nabla\varphi
 \,\mathrm d\lambda_d(x)\\
&\longrightarrow \langle f,\varphi x_i\rangle
 -\int_\Omega x_i q\cdot\nabla\varphi\,\mathrm d\lambda_d(x)
 =\int_\Omega\varphi q_i\,\mathrm d\lambda_d(x).
\end{align*}
Pierwszy składnik zbiega dzięki słabej zbieżności
$\varphi w_i^\eps$ w $H_0^1$; w drugim odejmujemy
$\int_\Omega x_iq_\eps\cdot\nabla\varphi\,\mathrm d\lambda_d$,
a błąd jest ograniczony przez
$\|\nabla\varphi\|_\infty\|w_i^\eps-x_i\|_2\|q_\eps\|_2$.
Pozostała całka zbiega ze słabej zbieżności $q_\eps$.
Ostatnia równość wynika z testu $\varphi x_i$
w równaniu dla $q$.
Z drugiej strony, zerową dywergencję $Q_i^\eps$
możemy testować przez $\varphi u_\eps\in H_0^1$
dzięki gęstości funkcji testowych i $Q_i^\eps\in L^2$.
Otrzymujemy
\begin{align*}
\int_\Omega\varphi\,Q_i^\eps\cdot\nabla u_\eps
\,\mathrm d\lambda_d(x)
&=-\int_\Omega u_\eps Q_i^\eps\cdot\nabla\varphi
\,\mathrm d\lambda_d(x)\\
&\longrightarrow
 -\int_\Omega u(A^{\mathrm{hom}}e_i)\cdot\nabla\varphi
 \,\mathrm d\lambda_d(x)
 =\int_\Omega\varphi(A^{\mathrm{hom}}e_i)\cdot\nabla u
 \,\mathrm d\lambda_d(x).
\end{align*}
Tutaj błąd zastąpienia $u_\eps$ przez $u$ jest
ograniczony przez
$\|\nabla\varphi\|_\infty\|u_\eps-u\|_2\|Q_i^\eps\|_2$,
a pozostała całka zbiega ze słabej zbieżności
$Q_i^\eps$. Ostatnia równość to całkowanie przez
części dla stałego wektora $A^{\mathrm{hom}}e_i$.
Kluczowa jest dokładna tożsamość
$q_\eps\cdot\nabla w_i^\eps
=Q_i^\eps\cdot\nabla u_\eps$, wynikająca z symetrii $A$.
Porównując granice, dostajemy
$q_i=(A^{\mathrm{hom}}e_i)\cdot\nabla u$ dla każdego $i$.
Ponieważ $A^{\mathrm{hom}}$ jest symetryczna, oznacza to
$q=A^{\mathrm{hom}}\nabla u$.

\textbf{Krok 4: identyfikacja rozwiązania i cały ciąg.}
Równość $\diver q=-f$ i poprzednia identyfikacja dają
$-\diver(A^{\mathrm{hom}}\nabla u)=f$, $u\in H_0^1(\Omega)$.
Istnienie i jednoznaczność rozwiązania granicznego
wynikają z Twierdzenia~\ref{thm:pde30-cell}
i twierdzenia Laxa--Milgrama, zatem $u=u_0$.
Dla każdego ciągu $\eps_k\to0$ znaleźliśmy
podciąg realizujący wszystkie trzy zbieżności
do tych samych granic. Gdyby nie zachodziła
silna zbieżność w $L^2$, istniałby ciąg z normą
błędu ograniczoną od dołu przez dodatnią liczbę.
Analogicznie dla każdej z dwóch słabych zbieżności
testujemy przez ustalony ciągły funkcjonał liniowy.
W obu przypadkach otrzymalibyśmy sprzeczność
z własnością podciągu. Dowodzi to zbieżności całej
rodziny. Argument od początku obejmuje
$f\in H^{-1}(\Omega)$, bez aproksymowania $f$
funkcjami z $L^2$.
\end{proof}


\begin{twierdzenie}[Zbieżność energii i stabilność względem danych]
\label{thm:pde30-energy-convergence}
Niech $\Omega\subset\R^d$ będzie ograniczonym obszarem
Lipschitzowskim, a $A$ będzie symetryczna, okresowa,
mierzalna i jednostajnie eliptyczna jak w
\eqref{eq:pde30-ellipticity}. Dla $f\in H^{-1}(\Omega)$
niech $u_\eps$ i $u_0$ będą rozwiązaniami zagadnień
\eqref{eq:pde30-micro} i~\eqref{eq:pde30-macro}.
Wówczas
\begin{equation}\label{eq:pde30-energy-convergence}
 \lim_{\eps\downarrow0}\int_\Omega
 A(x/\eps)\nabla u_\eps\cdot\nabla u_\eps
 \,\mathrm d\lambda_d(x)
 =\int_\Omega A^{\mathrm{hom}}\nabla u_0\cdot\nabla u_0
 \,\mathrm d\lambda_d(x).
\end{equation}
Co więcej, jeśli $f_\eps\to f$ silnie w $H^{-1}(\Omega)$,
a $v_\eps\in H_0^1(\Omega)$ rozwiązuje
$-\diver(A(x/\eps)\nabla v_\eps)=f_\eps$,
to $v_\eps\rightharpoonup u_0$ w $H_0^1$,
$v_\eps\to u_0$ w $L^2$, a odpowiednie strumienie
zbiegają słabo w $L^2$. Zachodzi także
\eqref{eq:pde30-energy-convergence} po zastąpieniu
$u_\eps$ przez $v_\eps$.
\end{twierdzenie}
\begin{proof}
Testując równanie dla $u_\eps$ funkcją $u_\eps$,
otrzymujemy, przy oznaczeniu parowania dualnego,
\[
 \int_\Omega A(x/\eps)\nabla u_\eps\cdot\nabla u_\eps
 \,\mathrm d\lambda_d(x)=\langle f,u_\eps\rangle.
\]
Słaba zbieżność w $H_0^1$ z
\ref{thm:pde30-homogenization} oznacza zbieżność
prawej strony do $\langle f,u_0\rangle$.
Test równania granicznego funkcją $u_0$ identyfikuje
tę liczbę z prawą stroną~\eqref{eq:pde30-energy-convergence}.

Dla zmiennego obciążenia niech $u_\eps$ nadal oznacza
rozwiązanie przy ustalonym $f$. Różnica
$z_\eps=v_\eps-u_\eps$ spełnia równanie wariacyjne
$a_\eps(z_\eps,\varphi)=\langle f_\eps-f,\varphi\rangle$.
Test $\varphi=z_\eps$ i koercywność dają
$\lambda\|z_\eps\|_{H_0^1}^2\le
\|f_\eps-f\|_{H^{-1}}\|z_\eps\|_{H_0^1}$, czyli
\[
 \lambda\|z_\eps\|_{H_0^1}\le
 \|f_\eps-f\|_{H^{-1}},
\]
gdzie w normie $H_0^1$ bierzemy normę gradientu.
W szczególności $z_\eps\to0$ silnie w $H_0^1$,
a $\|A(x/\eps)\nabla z_\eps\|_{L^2}\to0$.
Wszystkie wymienione zbieżności dziedziczymy więc
z twierdzenia~\ref{thm:pde30-homogenization}.
Tożsamość energii dla $v_\eps$ ma postać
$\int_\Omega A(x/\eps)\nabla v_\eps\cdot\nabla v_\eps\,\mathrm d\lambda_d(x)
=\langle f_\eps,v_\eps\rangle$.
Z silnej zbieżności $f_\eps$ w $H^{-1}$,
ograniczoności $v_\eps$ w $H_0^1$ i jego słabej
zbieżności do $u_0$ wynika
$\langle f_\eps,v_\eps\rangle\to\langle f,u_0\rangle$:
różnica jest sumą
$\langle f_\eps-f,v_\eps\rangle+\langle f,v_\eps-u_0\rangle$.
Pierwszy składnik dąży do zera przez ograniczoność
norm $v_\eps$, drugi przez słabą zbieżność.
To kończy dowód.
\end{proof}

\begin{wniosek}[Lokalna zbieżność gęstości energii]
\label{cor:pde30-local-energy}
W założeniach Twierdzenia~\ref{thm:pde30-energy-convergence}
dla każdego $\varphi\in C_c(\Omega)$ zachodzi
\[
 \begin{aligned}
 &\lim_{\eps\downarrow0}\int_\Omega
 \varphi A(x/\eps)\nabla v_\eps\cdot\nabla v_\eps
 \,\mathrm d\lambda_d\\
 &\qquad=\int_\Omega
 \varphi A^{\mathrm{hom}}\nabla u_0\cdot\nabla u_0
 \,\mathrm d\lambda_d .
 \end{aligned}
\]
Obejmuje to przypadek $v_\eps=u_\eps$, $f_\eps=f$.
\end{wniosek}
\begin{proof}
Najpierw niech $\varphi\in C_c^\infty(\Omega)$.
Testując przez $\varphi v_\eps\in H_0^1$, mamy
\[
 \int_\Omega\varphi A(x/\eps)\nabla v_\eps\cdot\nabla v_\eps
 \,\mathrm d\lambda_d
 =\langle f_\eps,\varphi v_\eps\rangle
 -\int_\Omega v_\eps A(x/\eps)\nabla v_\eps\cdot\nabla\varphi
 \,\mathrm d\lambda_d .
\]
Mnożenie przez $\varphi$ jest ograniczone w $H_0^1$,
zatem pierwszy składnik zbiega do
$\langle f,\varphi u_0\rangle$ dzięki silnej zbieżności
$f_\eps$ i słabej zbieżności $\varphi v_\eps$.
Drugi składnik zbiega do
$-\int_\Omega u_0 A^{\mathrm{hom}}\nabla u_0\cdot\nabla\varphi
\,\mathrm d\lambda_d$ dzięki silnej zbieżności $v_\eps$
w $L^2$ i słabej zbieżności strumieni.
Stosując równanie graniczne z testem $\varphi u_0$,
identyfikujemy granicę z żądaną całką.
Dla $\varphi\in C_c(\Omega)$ wybieramy gładkie
przybliżenia zbieżne jednostajnie, z nośnikami
w ustalonym zbiorze zwartym wewnątrz $\Omega$.
Błąd aproksymacji jest ograniczony przez normę
jednostajną różnicy testów pomnożoną przez
$\Lambda\|\nabla v_\eps\|_2^2$, jednostajnie
ograniczoną. Analogicznie dla granicy.
Otrzymujemy tezę.
\end{proof}

\begin{uwaga}[Zbieżność energii nie oznacza silnej zbieżności gradientów]
W~\eqref{eq:pde30-energy-convergence} występują
\emph{zmienne formy kwadratowe}; nie jest to
zbieżność $\|\nabla u_\eps-\nabla u_0\|_2\to0$.
Przykład jednowymiarowy z
\ref{prop:pde30-corrector1d} pokazuje, że ta ostatnia
zbieżność może być fałszywa.
\end{uwaga}


\begin{uwaga}[Dlaczego w homogenizacji występuje energia?]
\label{rem:pde-v44-homogenization-monotone}
Dla każdego $\eps>0$ macierz $A(x/\eps)$
określa liniowy operator monotoniczny
$V\to V'$ z rozdziału~\ref{chap:pde26}. Istnienie
rozwiązania wynika z koercywności; nie wymaga ono
własnej konstrukcji energii homogenizacyjnej.
Jednak w twierdzeniu~\ref{thm:pde30-energy-convergence}
\emph{symetria} $A$ pozwala interpretować
$\frac12\int_\Omega A(x/\eps)\nabla v\cdot\nabla v\,
\mathrm d\lambda_d(x)-\langle f,v\rangle$ jako funkcjonał,
którego punkt minimalny rozwiązuje równanie. Dla
macierzy niesymetrycznych forma pozostaje koercywna,
lecz jej część antysymetryczna nie jest widoczna
w funkcjonale kwadratowym. Oszacowanie energii z testu
$v=u_\eps$ nadal ma sens; nie daje samo przez się
charakteryzacji minimizacyjnej całego operatora.
\end{uwaga}


\section{Warstwy w wymiarze dwa: anizotropia z izotropowych faz}
\begin{przyklad}[Efektywna macierz materiału warstwowego]\label{ex:pde30-layer}
Niech $d=2$, $A(y)=a(y_1)I_2$, gdzie $a$ jest okresowa,
$0<\lambda\le a\le\Lambda$. Dla $j=1$ poprawka zależy
tylko od $y_1$ i spełnia
$1+\chi_1'(y_1)=a_{\mathrm{hom}}/a(y_1)$.
Dla $j=2$ funkcja $\chi_2=0$ rozwiązuje zadanie komórkowe,
ponieważ $\diver_y(Ae_2)=\partial_{y_2}a(y_1)=0$.
Stąd
\begin{equation}\label{eq:pde30-layers}
 A^{\mathrm{hom}}=
 \begin{pmatrix}
  \bigl(\int_0^1 a(y)^{-1}\,\mathrm d\lambda_1(y)\bigr)^{-1}&0\\
  0&\int_0^1a(y)\,\mathrm d\lambda_1(y)
 \end{pmatrix}.
\end{equation}
Jeśli $a$ przyjmuje wartości $1$ i $4$ na dwóch równych
połówkach komórki, to $A^{\mathrm{hom}}=\operatorname{diag}(8/5,5/2)$.
W każdej fazie materiał jest lokalnie izotropowy,
ale jego zachowanie efektywne jest anizotropowe.
\end{przyklad}

\begin{stwierdzenie}[Granice średniej harmonicznej i arytmetycznej]
\label{prop:pde30-bounds}
Dla $d=1$ i współczynnika $a$ jak wyżej zachodzi
$\lambda\le a_{\mathrm{hom}}\le\int_0^1 a(y)
\,\mathrm d\lambda_1(y)\le\Lambda$,
a równość między średnią harmoniczną i arytmetyczną
zachodzi wtedy i tylko wtedy, gdy $a$ jest stała prawie wszędzie.
\end{stwierdzenie}
\begin{proof}
Pierwsza i ostatnia nierówność wynikają z oszacowań $a$.
Nierówność Cauchy'ego--Schwarza zastosowana do
$\sqrt a$ i $1/\sqrt a$ daje
$1\le(\int_0^1a)(\int_0^1a^{-1})$.
Równość zachodzi dokładnie, gdy $\sqrt a$ jest
proporcjonalna do $1/\sqrt a$ prawie wszędzie.
\end{proof}

\section{Pułapki zbieżności i poprawki pierwszego rzędu}
\begin{kontrprzyklad}[Słaba zbieżność współczynników nie określa równania granicznego]
Na $(0,1)$ niech $a(y)$ będzie niejednorodna i okresowa.
Wtedy $a(x/\eps)\rightharpoonup\int_0^1 a$ w $L^2(0,1)$,
lecz rozwiązania równań $-(a(x/\eps)u_\eps')'=f$
zbiegną do rozwiązania równania z $a_{\mathrm{hom}}
=(\int_0^1a^{-1})^{-1}$. Obie liczby różnią się,
gdy $a$ nie jest stała. Niedozwolone jest więc przejście
do granicy osobno w iloczynie dwóch słabo zbieżnych funkcji
$a(x/\eps)$ i $u_\eps'$.
\end{kontrprzyklad}

\begin{uwaga}[Granica zakresu twierdzenia]
Twierdzenie~\ref{thm:pde30-homogenization} dowodzi zbieżności
bez oszacowania jej szybkości w wymiarach $d\ge2$. Nie
implikuje samo z siebie globalnego szacowania
$\|\nabla u_\eps-(I+\nabla\chi(x/\eps))\nabla u_0\|_2\le C\eps$:
takie oszacowanie wymaga dodatkowej regularności danych,
współczynników i brzegu oraz analizy poprawek brzegowych.
Nie wystarcza formalne podstawienie ansatzu
$u_\eps\approx u_0+\eps\sum_j\chi_j(x/\eps)\partial_j u_0$.
W wymiarze jeden otrzymaliśmy natomiast
oszacowanie gradientowe z poprawką w
stwierdzeniu~\ref{prop:pde30-corrector1d} bez dodatkowej
gładkości współczynnika $a$.
\end{uwaga}

\section{Przykład numeryczny i interpretacja strumienia}
Dla $\Omega=(0,1)$ wybieramy $a(y)=2+\sin(2\pi y)$ i $f=1$.
Wtedy $\bar a=2$, podczas gdy
$a_{\mathrm{hom}}=\bigl(\int_0^1(2+\sin(2\pi y))^{-1}
\,\mathrm d\lambda_1(y)\bigr)^{-1}=\sqrt3$.
Rozwiązanie graniczne jest równe
$u_0(x)=x(1-x)/(2\sqrt3)$, a dokładne $u_\eps$ daje
wzór~\eqref{eq:pde30-explicit-1d}. W każdym mikrośrodowisku
strumień $q_\eps=a(x/\eps)u_\eps'$ jest funkcją afiniczną
$x$, mimo że sam gradient $u_\eps'$ szybko oscyluje.

W dodatku numerycznym~\ref{sec:pde30-code} przedstawiamy
kod Plotly, który wyznacza rodzinę rozwiązań $u_\eps$,
porównuje średnie arytmetyczną i harmoniczną, ilustruje
oscylacje gradientów oraz bada zbieżność numeryczną.




\chapter[Analiza mikrolokalna]{Analiza mikrolokalna: zbiór falowy, symbole i operatory pseudoróżniczkowe}
\label{chap:pde31}

\section{Motywacja i konwencja transformaty Fouriera}
Nośnik osobliwy dystrybucji mówi, \emph{gdzie} dystrybucja
nie jest lokalnie reprezentowana przez funkcję gładką. Nie mówi jednak, w \emph{jakich kierunkach} jej
oscylacje wysokoczęstotliwościowe nie zanikają. Tę brakującą informację
koduje \emph{zbiór falowy} (ang. \emph{wavefront set}).
W analizie mikrolokalnej każdemu punktowi $x$ przyporządkowujemy
kowektor $\xi\ne0$; para $(x,\xi)$ należy do zbioru falowego, jeżeli
lokalna transformata Fouriera nie maleje szybko w stożku wokół $\xi$.

W tym rozdziale $d\ge1$, a $U\subset\R^d$ jest otwarty. Dla
$h\in\mathcal S(\R^d)$ ustalamy normalizację
\begin{equation}\label{eq:pde31-fourier}
 \widehat h(\xi)=\int_{\R^d}e^{-ix\cdot\xi}h(x)
 \,\mathrm d\lambda_d(x),\qquad
 h(x)=(2\pi)^{-d}\int_{\R^d}e^{ix\cdot\xi}\widehat h(\xi)
 \,\mathrm d\lambda_d(\xi).
\end{equation}
Przedłużamy ją na dystrybucje temperowane przez dualność,
zgodnie z definicją w rozdziale~\ref{chap:pde-fourier-Rd}.
Konwencja jest dokładnie tą z równania~\eqref{eq:pde8-convention};
nie zmieniamy znaku fazy ani czynników $(2\pi)^{-d}$.
Przez $\langle\xi\rangle=(1+|\xi|^2)^{1/2}$ oznaczamy standardową
wagę Sobolewa. \emph{Stożek otwarty} $\Gamma\subset\R^d\setminus\{0\}$
jest otwarty w $\R^d\setminus\{0\}$ i niezmienniczy
względem mnożenia przez dodatnie liczby.
Dystrybucje mogą być zespolone; stosujemy parowanie
liniowe względem funkcji testowej, zgodnie z konwencją
rozdziału~\ref{chap:pde-fourier-Rd}.

\begin{lemat}[Transformata dystrybucji o zwartym nośniku]
\label{lem:pde31-compact-fourier}
Niech $T\in\mathcal E'(\R^d)$ i niech
$\theta\in C_c^\infty(\R^d)$ będzie równa $1$
w otoczeniu $\operatorname{supp}T$. Wtedy
$\widehat T$ jest reprezentowana przez funkcję
\[
 F_T(\xi)=\langle T(x),\theta(x)e^{-ix\cdot\xi}\rangle.
\]
Funkcja ta nie zależy od wyboru $\theta$, jest gładka
i istnieje $M\in\N_0$ takie, że dla każdego
wielowskaźnika $\alpha$ istnieje $C_\alpha>0$ z
\[
 |\partial_\xi^\alpha F_T(\xi)|\le C_\alpha
 \langle\xi\rangle^M.
\]
\end{lemat}
\begin{proof}
Różnica dwóch dopuszczalnych funkcji $\theta$
znika w otoczeniu nośnika $T$, więc nie zmienia
wartości parowania. Z oszacowania lokalnego rzędu
dystrybucji (Twierdzenie~\ref{thm:pde7-local-order})
istnieją $C,M$ takie, że dla testów o nośniku
w ustalonym zwartym otoczeniu $\operatorname{supp}\theta$
wartość $|T(\psi)|$ nie przekracza
$C\max_{|\beta|\le M}\|\partial^\beta\psi\|_\infty$.
Stąd również $T$ przedłuża się na $\mathcal S$
przez $T(\psi)=T(\theta\psi)$, ciągle względem
półnorm Schwartza.
Różniczkując funkcję $\theta(x)e^{-ix\cdot\xi}$
względem $\xi$ w topologii funkcji testowych,
otrzymujemy
\[
 \partial_\xi^\alpha F_T(\xi)
 =\langle T(x),(-ix)^\alpha\theta(x)e^{-ix\cdot\xi}\rangle.
\]
Różniczkowanie jest dopuszczalne: na ustalonym zwartym
nośniku różnicowe ilorazy i wszystkie ich pochodne
względem $x$ zbiegają jednostajnie, z rozwinięcia
Taylora funkcji wykładniczej.
Reguła Leibniza i ograniczoność $x$ na nośniku
$\theta$ dają żądane oszacowanie, z tym samym $M$.

Dla $\psi\in\mathcal S$ całka z
$\theta(x)e^{-ix\cdot\xi}\psi(\xi)$ względem $\xi$
zbiega w każdej półnormie funkcji testowych
o ustalonym nośniku: pochodne względem $x$
są majorowane przez stałą razy
$\langle\xi\rangle^{|\beta|}|\psi(\xi)|$.
Można więc przenieść $T$ pod całkę
(przez obcięcie całki i przybliżenie sumami Riemanna).
Otrzymujemy
\[
 \int_{\R^d}F_T(\xi)\psi(\xi)\,\mathrm d\lambda_d(\xi)
 =\langle T,\theta\widehat\psi\rangle
 =\langle T,\widehat\psi\rangle
 =\langle\widehat T,\psi\rangle.
\]
To dowodzi reprezentacji transformaty.
\end{proof}

\begin{definicja}[Regularność mikrolokalna]\label{def:pde31-wf}
Niech $u\in\mathcal D'(U)$, $x_0\in U$ i
$\xi_0\in\R^d\setminus\{0\}$. Mówimy, że
$(x_0,\xi_0)\notin\operatorname{WF}(u)$, jeżeli istnieją
$\chi\in C_c^\infty(U)$ z $\chi=1$ w otoczeniu $x_0$
(iloczyn $\chi u$ przedłużamy przez zero na $\R^d$)
i otwarty stożek $\Gamma$ zawierający $\xi_0$, takie że dla każdego
$N\in\N$ istnieje $C_N>0$ spełniająca
\begin{equation}\label{eq:pde31-wf-decay}
 |\widehat{\chi u}(\xi)|\le C_N\langle\xi\rangle^{-N}
 \qquad(\xi\in\Gamma).
\end{equation}
Dopełnienie zbioru tak określonych par nazywamy zbiorem falowym $u$.
Jest to domknięty podzbiór $U\times(\R^d\setminus\{0\})$,
niezmienniczy względem dodatniego skalowania drugiej współrzędnej.
\end{definicja}
Istotnie, jeśli para jest regularna mikrolokalnie,
to ta sama funkcja $\chi$ i ten sam stożek świadczą
o regularności wszystkich par z iloczynu otoczenia,
na którym $\chi=1$, i stożka $\Gamma$.
Dopełnienie zbioru falowego jest więc otwarte.
Niezmienniczość wynika bezpośrednio ze stożkowości
$\Gamma$.

\begin{lemat}[Niezależność od dalszego obcięcia]\label{lem:pde31-cutoff}
Załóżmy, że $u\in\mathcal D'(U)$, $\chi\in C_c^\infty(U)$,
$\chi=1$ w otoczeniu $x_0$, i że
$\widehat{\chi u}$ szybko maleje w stożku $\Gamma$.
Jeżeli $\eta\in C_c^\infty(U)$ ma nośnik w tym otoczeniu,
a $\Gamma'$ jest stożkiem, którego domknięcie przecięte ze sferą
jednostkową leży wewnątrz $\Gamma$, to
$\widehat{\eta u}$ szybko maleje na $\Gamma'$.
\end{lemat}
\begin{proof}
Połóżmy $T=\chi u$, przedłużone przez zero na $\R^d$.
Ponieważ $\eta u=\eta T$, mamy
\[
 \widehat{\eta u}(\xi)=(2\pi)^{-d}\int_{\R^d}
 \widehat\eta(\xi-\zeta)\widehat T(\zeta)
 \,\mathrm d\lambda_d(\zeta).
\]
Istotnie, stosując odwrotną transformatę Fouriera
do $\eta$ i mnożąc przez $e^{-ix\cdot\xi}$,
otrzymujemy całkę funkcji wykładniczych z wagą
$\widehat\eta$. Po pomnożeniu przez funkcję
obcinającą równą $1$ blisko nośnika $T$,
całka zbiega w półnormach funkcji testowych,
jak w Lemacie~\ref{lem:pde31-compact-fourier}.
Możemy więc przenieść $T$ pod całkę i podstawić
$\zeta=\xi-\omega$, gdzie $\omega$ jest zmienną
całkowania we wzorze odwrotnym dla $\eta$.
Całka jest bezwzględnie zbieżna: z tego samego
lematu mamy $|\widehat T(\zeta)|\le C\langle\zeta\rangle^M$,
a $\widehat\eta$ należy do $\mathcal S$.

Z założenia o $\Gamma'$ istnieje $0<\delta<1/2$
takie, że dla $\xi\in\Gamma'$ i
$|\xi-\zeta|\le\delta|\xi|$ zachodzi $\zeta\in\Gamma$.
Wynika to z dodatniej odległości zwartego zbioru
$\overline{\Gamma'}\cap S^{d-1}$ od dopełnienia
$\Gamma\cap S^{d-1}$ na sferze: jednostkowe
kierunki $\xi$ i $\zeta$ różnią się o co najwyżej
$2|\xi-\zeta|/|\xi|\le2\delta$.
Jeśli dopełnienie na sferze jest puste,
wystarcza dowolne $0<\delta<1/2$.
W tym obszarze $|\zeta|\ge(1-\delta)|\xi|$,
więc dla dowolnego $N\in\N$
\[
 \int_{|\xi-\zeta|\le\delta|\xi|}
 |\widehat\eta(\xi-\zeta)\widehat T(\zeta)|
 \,\mathrm d\lambda_d(\zeta)
 \le C_N'\langle\xi\rangle^{-N}\|\widehat\eta\|_1 .
\]
W pozostałym obszarze kładziemy $w=\xi-\zeta$.
Mamy $|w|>\delta|\xi|$, a zatem
$\langle\xi\rangle\le\delta^{-1}\langle w\rangle$
oraz $\langle\xi-w\rangle\le(1+\delta^{-1})\langle w\rangle$.
Z wielomianowego wzrostu $\widehat T$ i szybkiego
zaniku $\widehat\eta$, wybierając wykładnik
$L=M+N+d+1$, dostajemy
\[
 \begin{aligned}
 &\int_{|w|>\delta|\xi|}
 |\widehat\eta(w)\widehat T(\xi-w)|
 \,\mathrm d\lambda_d(w)\\
 &\qquad\le C_L\int_{|w|>\delta|\xi|}
 \langle w\rangle^{-N-d-1}\,\mathrm d\lambda_d(w)\\
 &\qquad\le C_{L,\delta}\langle\xi\rangle^{-N}
 \int_{\R^d}\langle w\rangle^{-d-1}\,\mathrm d\lambda_d(w).
 \end{aligned}
\]
Ostatnia całka jest skończona, co widać
współrzędnymi biegunowymi dla $|w|\ge1$
i z ograniczoności funkcji na kuli jednostkowej.
Oba oszacowania dają szybki zanik w $\Gamma'$.
\end{proof}

\begin{stwierdzenie}[Rzut zbioru falowego]\label{prop:pde31-singular-support}
Dla $u\in\mathcal D'(U)$ zachodzi
$\operatorname{sing\,supp}(u)=\pi_x(\operatorname{WF}(u))$,
gdzie $\pi_x(x,\xi)=x$, a $\operatorname{sing\,supp}(u)$
oznacza zbiór punktów, w których $u$ nie jest funkcją klasy $C^\infty$
lokalnie.
\end{stwierdzenie}
\begin{proof}
Jeżeli $u$ jest gładka w otoczeniu $x_0$, to dla
$\chi\in C_c^\infty$ równej $1$ blisko $x_0$
i o nośniku w tym otoczeniu mamy
$\chi u\in C_c^\infty$; jej transformata maleje szybko we wszystkich
kierunkach, zatem $x_0$ nie leży w rzucie zbioru falowego.

Przeciwnie, załóżmy, że wszystkie pary $(x_0,\xi)$, $|\xi|=1$,
nie należą do zbioru falowego. Dla każdego kierunku $\xi$ wybieramy stożek
regularności $\Gamma_\xi$, funkcję obcinającą $\chi_\xi$
oraz mniejszy otwarty stożek $\Gamma_\xi'$
zawierający $\xi$, którego domknięcie na sferze
mieści się w $\Gamma_\xi$.
Ze zwartości sfery wybieramy skończone pokrycie
przez $\Gamma_{\xi_j}'$. W przecięciu otoczeń,
na których wszystkie $\chi_{\xi_j}=1$,
wybieramy $\eta\in C_c^\infty(U)$, równą $1$
blisko $x_0$. Z Lematu~\ref{lem:pde31-cutoff}
$\widehat{\eta u}$ szybko maleje w każdym
$\Gamma_{\xi_j}'$, a biorąc maksimum skończenie
wielu stałych, otrzymujemy szybki zanik na całej
$\R^d$ (również w zerze, przez ciągłość transformaty).
Definiujemy
\[
 v(x)=(2\pi)^{-d}\int_{\R^d}
 e^{ix\cdot\xi}\widehat{\eta u}(\xi)\,\mathrm d\lambda_d(\xi).
\]
Dla każdego wielowskaźnika $\alpha$ funkcja
$|\xi|^{|\alpha|}|\widehat{\eta u}(\xi)|$ jest całkowalna,
więc twierdzenie o zbieżności zmajoryzowanej
uzasadnia różniczkowanie pod całką i ciągłość
wszystkich pochodnych $v$.
Twierdzenie~\ref{thm:pde8-dist-rules} o odwracalności
transformaty w $\mathcal S'$ i twierdzenie Fubiniego
dla parowania z funkcjami Schwartza dają $v=\eta u$
w sensie dystrybucji.
Ponieważ $\eta=1$ w otoczeniu $x_0$, dystrybucja
$u$ jest tam reprezentowana przez funkcję gładką $v$.
\end{proof}

\section{Przykłady kierunkowych osobliwości}
\begin{przyklad}[Delta Diraca i skok]\label{ex:pde31-delta}
Dla $\delta_0\in\mathcal D'(\R^d)$ z Lematu
~\ref{lem:pde31-compact-fourier} mamy
$\widehat{\chi\delta_0}(\xi)=\chi(0)$.
Jeśli $\chi=1$ blisko zera, transformata jest
identycznie równa $1$ i nie maleje szybko w żadnym
niepustym stożku. Poza zerem dystrybucja jest
reprezentowana lokalnie przez funkcję zerową.
Zatem $\operatorname{WF}(\delta_0)=\{(0,\xi):\xi\ne0\}$.

Dla funkcji Heaviside'a $H(x)=\chi_{(0,\infty)}(x)$
w $\R$ zbiór falowy jest taki sam jak dla delty
w wymiarze jeden. Istotnie, dla dowolnej
$\chi\in C_c^\infty(\R)$ równej $1$ blisko zera
oraz $\xi\ne0$, całkowanie przez części daje
\[
 \widehat{\chi H}(\xi)
 =\int_0^\infty e^{-ix\xi}\chi(x)\,\mathrm d\lambda_1(x)
 =\frac{1}{i\xi}+
 \frac{1}{i\xi}\int_0^\infty e^{-ix\xi}\chi'(x)
 \,\mathrm d\lambda_1(x).
\]
Funkcja $\chi'\chi_{(0,\infty)}$ jest gładka
i ma zwarty nośnik, ponieważ $\chi'=0$
w otoczeniu zera. Jej transformata należy do
$\mathcal S(\R)$, więc ostatni składnik jest
$O(|\xi|^{-N})$ dla każdego $N$ przy $|\xi|\to\infty$.
W szczególności $|\xi\widehat{\chi H}(\xi)|\to1$
w obu kierunkach, co wyklucza szybki zanik
dla każdej dopuszczalnej lokalizacji.
Poza zerem $H$ jest lokalnie stała.
\end{przyklad}

\begin{przyklad}[Osobliwość hiperpłaszczyzny]\label{ex:pde31-conormal}
Niech $d\ge2$, $x=(x_1,x')$ i
$u=\delta(x_1)\otimes1(x')$, tj.
$\langle u,\varphi\rangle=\int_{\R^{d-1}}\varphi(0,x')
\,\mathrm d\lambda_{d-1}(x')$ dla $\varphi\in C_c^\infty(\R^d)$.
Wtedy
\[
 \operatorname{WF}(u)=
 \{((0,x'),(\xi_1,0)):\ x'\in\R^{d-1},\ \xi_1\ne0\}.
\]
Poza hiperpłaszczyzną $x_1=0$ dystrybucja znika
lokalnie. Ustalmy punkt $(0,x_0')$ na niej
i dowolną funkcję obcinającą $\chi$ równą $1$
blisko tego punktu. Z definicji dystrybucji
\[
 \widehat{\chi u}(\xi_1,\xi')
 =\int_{\R^{d-1}} e^{-ix'\cdot\xi'}\chi(0,x')
 \,\mathrm d\lambda_{d-1}(x')
 =\widehat h(\xi'),\qquad h(x')=\chi(0,x').
\]
Mamy $h\in C_c^\infty(\R^{d-1})$ i $h=1$
blisko $x_0'$. Jeśli kierunek $\xi_0=(\xi_{0,1},\xi_0')$
spełnia $\xi_0'\ne0$, to w odpowiednim otwartym
stożku wokół niego zachodzi $|\xi'|\ge c|\xi|$
z $c>0$. Szybki zanik $\widehat h$ daje więc
regularność mikrolokalną w tym stożku.
Jeśli natomiast $\xi_0'=0$, to z $h\not\equiv0$
i odwracalności transformaty Fouriera wynika,
że istnieje $\eta'\in\R^{d-1}$ z
$\widehat h(\eta')\ne0$. Dla $r\to+\infty$
wektory $(r\xi_{0,1},\eta')$ ostatecznie należą
do każdego otwartego stożka zawierającego $\xi_0$,
a transformata ma na nich tę samą niezerową wartość.
To wyklucza szybki zanik dla dowolnej $\chi$.
W szczególności nie wystarcza badanie samej półprostej
$\xi'=0$, gdyż całka z obcięcia $h$ może być zerowa.
\end{przyklad}

\begin{center}
\begin{tikzpicture}[scale=0.83,>=Latex]
  \draw[->,gray!70] (-2.8,0)--(2.8,0) node[right] {$x_1$};
  \draw[->,gray!70] (0,-2.5)--(0,2.6) node[above] {$x_2$};
  \draw[blue!65!black,very thick] (0,-2.25)--(0,2.25);
  \foreach \y in {-1.5,0,1.5}{
    \draw[-{Latex[length=3mm]},red!65!black,thick] (0,\y)--(1.25,\y);
    \draw[-{Latex[length=3mm]},red!65!black,thick] (0,\y)--(-1.25,\y);
  }
  \node[blue!65!black,below right] at (0,-2.25) {nośnik $u$};
  \node[red!65!black,above right] at (1.15,1.5) {kowektory normalne};
\end{tikzpicture}

\small Rysunek: nośnik $\delta(x_1)$ i kierunki zbioru falowego
w wymiarze dwa (schemat).
\end{center}

\begin{uwaga}[Zbiór falowy a nośnik]
Dystrybucja $\delta(x_1)$ ma nośnik będący hiperpłaszczyzną,
lecz zbiór falowy składa się z \emph{kowektorów normalnych} do niej.
Każda funkcja gładka, nawet o niezwartym nośniku,
ma pusty zbiór falowy; przykładem jest funkcja $1$.
Nie należy utożsamiać nośnika, nośnika osobliwego
ani zbioru falowego.
\end{uwaga}


\begin{twierdzenie}[Dokładna propagacja zbioru falowego przy translacji]
\label{thm:pde31-transport-wf}
Niech $u_0\in\mathcal D'(\R^d)$, $c\in\R^d$,
a dla $t\in\R$ niech $u(t)=\tau_{ct}u_0$,
gdzie $\tau_a u_0$ oznacza translację dystrybucji,
która dla funkcji ma postać $(\tau_a u_0)(x)=u_0(x-a)$.
Wtedy $t\mapsto u(t)$ jest gładkim odwzorowaniem
$\R\to\mathcal D'(\R^d)$ i rozwiązuje
$\partial_tu+c\cdot\nabla_xu=0$ przy $u(0)=u_0$.
Dla każdego $t\in\R$ zachodzi dokładna tożsamość
\begin{equation}\label{eq:pde31-transport-wf}
 \operatorname{WF}(u(t))=
 \{(x+ct,\xi):(x,\xi)\in\operatorname{WF}(u_0)\}.
\end{equation}
\end{twierdzenie}
\begin{proof}
Definiujemy translację przez
$\langle\tau_a u_0,\varphi\rangle
=\langle u_0,\varphi(\cdot+a)\rangle$.
Dla ustalonej funkcji testowej $\varphi$
i $t$ z przedziału zwartego $I$ funkcje
$\varphi(\cdot+ct)$ mają nośniki w jednym
zbiorze zwartym $K$. Taylorowski wzór z resztą
całkową pokazuje, że ilorazy różnicowe
zbiegają w każdej półnormie funkcji testowych
o nośniku w $K$. Dla pochodnych kolejnych rzędów
stosujemy ten sam argument do
$(c\cdot\nabla)^k\varphi(\cdot+ct)$.
Z ciągłości $u_0$ otrzymujemy
\[
 \frac{\mathrm d^k}{\mathrm dt^k}
 \langle u(t),\varphi\rangle
 =\langle u_0,(c\cdot\nabla)^k\varphi(\cdot+ct)\rangle
 =\langle(-c\cdot\nabla_x)^k u(t),\varphi\rangle .
\]
Zbieżność ilorazów różnicowych zachodzi także
w silnej topologii dystrybucji: dla ograniczonego
zbioru testów ich nośniki leżą w jednym zbiorze
zwartym i wszystkie półnormy są jednostajnie
ograniczone. Przesunięcia dla $t\in I$
zachowują te własności. Oszacowanie lokalnego
rzędu $u_0$ ogranicza resztę Taylora
po sparowaniu przez stałą razy moduł przyrostu,
jednostajnie na tym zbiorze testów.
Stąd $t\mapsto u(t)$ jest gładkie i
$\partial_tu=-c\cdot\nabla_xu$, a $u(0)=u_0$.

Dla dowolnej funkcji odcinającej
$\chi\in C_c^\infty(\R^d)$ definiujemy
$\eta(y)=\chi(y+a)$. Wtedy
$\chi\tau_a u_0=\tau_a(\eta u_0)$: dla testu
$\varphi$ obie strony mają wartość
$\langle u_0,\chi(\cdot+a)\varphi(\cdot+a)\rangle$.
Ponieważ $\eta u_0$ ma zwarty nośnik, wzór
z Lematu~\ref{lem:pde31-compact-fourier} i definicja
translacji dają
\[
 \widehat{\chi\tau_a u_0}(\xi)
 =e^{-ia\cdot\xi}\widehat{\eta u_0}(\xi).
\]
Współczynnik wykładniczy ma moduł jeden, zatem
szybki zanik w dowolnym stożku jest równoważny
po obu stronach. Funkcja $\chi$ jest równa jeden
w otoczeniu $x_0+a$ dokładnie wtedy, gdy $\eta$
jest równa jeden blisko $x_0$.
Stąd $(x_0+a,\xi_0)\notin\operatorname{WF}(\tau_a u_0)$
wtedy i tylko wtedy, gdy
$(x_0,\xi_0)\notin\operatorname{WF}(u_0)$.
Podstawiamy $a=ct$.
\end{proof}

\begin{przyklad}[Transport punktowej osobliwości]
\label{ex:pde31-moving-delta}
W jednym wymiarze dla $u_0=\delta_0$ otrzymujemy
$u(t,x)=\delta(x-ct)$. Z
\eqref{eq:pde31-transport-wf} wynika
$\operatorname{WF}(u(t))=\{(ct,\xi):\xi\ne0\}$.
Na czasoprzestrzeni $\R^2$ definiujemy $u$ przez
$\langle u,\varphi\rangle=\int_\R\varphi(t,ct)
\,\mathrm d\lambda_1(t)$.
Jest to dystrybucja rzędu zero: dla testów o nośniku
w ustalonym zbiorze zwartym całkowanie odbywa się
po ograniczonym przedziale czasu, więc moduł
parowania ogranicza stała razy $\|\varphi\|_\infty$.
Poza prostą $x=ct$ dystrybucja lokalnie znika.
Obliczymy
\[
 \operatorname{WF}(u)=
 \{((t,ct),(-c\xi,\xi)):t\in\R,\ \xi\ne0\}.
\]
Aby to sprawdzić, dla $\rho\in C_c^\infty(\R^2)$
obliczamy
$\widehat{\rho u}(\tau,\xi)=
\int_\R e^{-it(\tau+c\xi)}\rho(t,ct)\,
\mathrm d\lambda_1(t)$.
Jest to jednowymiarowa transformata funkcji
$h(t)=\rho(t,ct)\in C_c^\infty(\R)$ w punkcie
$\tau+c\xi$. W stożku wokół kierunku
$(\tau_0,\xi_0)$ z $\tau_0+c\xi_0\ne0$
mamy $|\tau+c\xi|\ge c_0|(\tau,\xi)|$
dla pewnego $c_0>0$, więc transformata
maleje szybko.
Ustalmy teraz niezerowy kierunek normalny
$(-c\xi_0,\xi_0)$ oraz dowolne obcięcie $\rho$
równe $1$ w otoczeniu $(t_0,ct_0)$.
Wtedy $h=1$ blisko $t_0$, więc z odwracalności
transformaty istnieje $\sigma\in\R$ z
$\widehat h(\sigma)\ne0$.
Dla $r\to+\infty$ pary
$(\sigma-rc\xi_0,r\xi_0)$ należą ostatecznie
do każdego otwartego stożka wokół kierunku
normalnego, ale transformata na nich jest
równa $\widehat h(\sigma)$. Wyklucza to
regularność mikrolokalną w tym kierunku.
Nie dowodzi to ogólnego twierdzenia o propagacji
bicharakterystyk dla operatorów o zmiennych współczynnikach.
\end{przyklad}


\begin{stwierdzenie}[Wygładzanie cieplne usuwa zbiór falowy]
\label{prop:pde-v44-heat-wavefront}
Niech $T\in\mathcal E'(\R^d)$ będzie dystrybucją
o zwartym nośniku i $t>0$. Dla jądra ciepła
$G_t(x)=(4\pi t)^{-d/2}e^{-|x|^2/(4t)}$
funkcja $G_t*T$ należy do $\mathcal S(\R^d)$,
a więc $\operatorname{WF}(G_t*T)=\varnothing$.
W szczególności dla $T=\delta_0$ zbiór falowy
w chwili początkowej jest niepusty we wszystkich
kierunkach, ale znika dla każdego $t>0$.
\end{stwierdzenie}
\begin{proof}
Z Lematu~\ref{lem:pde31-compact-fourier}
$\widehat T$ jest funkcją gładką na $\R^d$,
a ona sama i wszystkie jej pochodne mają
co najwyżej wzrost wielomianowy.
Jądro $G_t$ należy do $\mathcal S$, a jego
transformata wynosi $e^{-t|\xi|^2}$
(por. Lemat~\ref{lem:pde8-gauss}).
Twierdzenie~\ref{thm:pde8-convol-temp} daje
$\widehat{G_t*T}(\xi)=e^{-t|\xi|^2}\widehat T(\xi)$.
Każda pochodna mnożnika gaussowskiego jest
iloczynem wielomianu i tego samego mnożnika.
Z reguły Leibniza dla dowolnych $\alpha,N$
funkcja $\langle\xi\rangle^N
\partial_\xi^\alpha(e^{-t|\xi|^2}\widehat T(\xi))$
jest ograniczona, gdyż wykładniczy zanik
przeważa dowolną potęgę.
Iloczyn należy więc do przestrzeni Schwartza,
a jego odwrotna transformata również.
Z odwracalności transformaty w $\mathcal S'$
jest ona równa $G_t*T$ w sensie dystrybucji,
a ponieważ obie funkcje są gładkie,
równość zachodzi punktowo.
Z definicji~\ref{def:pde31-wf} każda funkcja gładka
ma pusty zbiór falowy.
\end{proof}
\begin{uwaga}[Nieodwracalność informacji wysokoczęstotliwościowej]
Powyższy fakt opisuje kierunki osobliwości, ale nie
jest oszacowaniem stabilności problemu odwrotnego.
Na obszarze ograniczonym mnożnik cieplny
$e^{-t\lambda_k}$ na funkcjach własnych laplasjanu
powoduje silną niestabilność odtwarzania warunku
początkowego; por. stwierdzenie
\ref{prop:pde-v44-heat-inverse} w rozdziale~\ref{chap:pde32}.
\end{uwaga}


\section{Symbole i kwantyzacja Kohna--Nirenberga}
\begin{definicja}[Klasa symboli $S^m_{1,0}$]\label{def:pde31-symbol}
Dla $m\in\R$ klasą $S^m_{1,0}(\R^d\times\R^d)$
nazywamy zbiór funkcji $a\in C^\infty(\R^d\times\R^d)$,
dla których dla każdych wielowskaźników $\alpha,\beta$
istnieje $C_{\alpha\beta}$ takie, że
\begin{equation}\label{eq:pde31-symbol}
 |\partial_x^\alpha\partial_\xi^\beta a(x,\xi)|
 \le C_{\alpha\beta}\langle\xi\rangle^{m-|\beta|}
 \quad(x,\xi\in\R^d).
\end{equation}
Dla $u\in\mathcal S(\R^d)$ operator pseudoróżniczkowy
w kwantyzacji Kohna--Nirenberga określamy wzorem
\begin{equation}\label{eq:pde31-op}
 \operatorname{Op}(a)u(x)=(2\pi)^{-d}
 \int_{\R^d}e^{ix\cdot\xi}a(x,\xi)\widehat u(\xi)
 \,\mathrm d\lambda_d(\xi).
\end{equation}
Liczba $m$ jest górnym ograniczeniem rzędu symbolu:
dla $m_1\le m_2$ mamy $S^{m_1}_{1,0}\subset S^{m_2}_{1,0}$.
Nie zakładamy jednorodności symbolu w $\xi$.
Całka~\eqref{eq:pde31-op} jest bezwzględnie zbieżna,
ponieważ $|a(x,\xi)|\le C_{00}\langle\xi\rangle^m$
i $\widehat u\in\mathcal S$.
\end{definicja}

\begin{lemat}[Ciągłość na przestrzeni Schwartza]
\label{lem:pde31-schwartz-op}
Jeżeli $a\in S^m_{1,0}$, to $\operatorname{Op}(a)$
jest ciągłym operatorem liniowym
$\mathcal S(\R^d)\to\mathcal S(\R^d)$.
\end{lemat}
\begin{proof}
Dla wielowskaźnika $\alpha$ różniczkujemy pod całką:
\[
 \partial_x^\alpha\operatorname{Op}(a)u(x)
 =(2\pi)^{-d}\int_{\R^d}e^{ix\cdot\xi}
 b_\alpha(x,\xi)\widehat u(\xi)\,\mathrm d\lambda_d(\xi),
 \qquad
 b_\alpha(x,\xi)=
 \sum_{\gamma\le\alpha}\binom{\alpha}{\gamma}
 (i\xi)^{\alpha-\gamma}\partial_x^\gamma a(x,\xi).
\]
Z reguły Leibniza i oszacowań symbolu wynika
$|\partial_\xi^\nu b_\alpha(x,\xi)|
\le C_{\alpha\nu}\langle\xi\rangle^{m+|\alpha|-|\nu|}$.
Każdy iloczyn $b_\alpha\widehat u$ jest więc
jednostajnie względem $x$ szybko malejący wraz
ze wszystkimi pochodnymi po $\xi$.
W szczególności funkcje pod całkami otrzymanymi
przez kolejne różniczkowania względem $x$
mają całkowalne majoranty niezależne od $x$.
Uzasadnia to różniczkowanie przez twierdzenie
o zbieżności zmajoryzowanej.

Dla wielowskaźnika $\beta$ wykorzystujemy
$x^\beta e^{ix\cdot\xi}=i^{-|\beta|}
\partial_\xi^\beta e^{ix\cdot\xi}$.
Całkowanie przez części daje
\[
 x^\beta\partial_x^\alpha\operatorname{Op}(a)u(x)
 =(2\pi)^{-d}i^{|\beta|}
 \int_{\R^d}e^{ix\cdot\xi}
 \partial_\xi^\beta(b_\alpha(x,\xi)\widehat u(\xi))
 \,\mathrm d\lambda_d(\xi).
\]
Nie ma wyrazów brzegowych dzięki szybkiemu
zanikowi iloczynu i jego pochodnych.
Rozwijając pochodną przez regułę Leibniza,
dla liczby całkowitej $L> d+\max(0,m+|\alpha|)$
otrzymujemy
\[
 \sup_x|x^\beta\partial_x^\alpha\operatorname{Op}(a)u(x)|
 \le C_{\alpha\beta L}
 \sum_{\nu\le\beta}\sup_\xi
 \langle\xi\rangle^L|\partial_\xi^\nu\widehat u(\xi)|.
\]
Istotnie, po wyłączeniu wskazanych półnorm
pozostaje całkowalna majoranta
$\langle\xi\rangle^{m+|\alpha|-L}$.
Prawa strona jest ciągłą półnormą funkcji $u$
z ciągłości transformaty Fouriera na przestrzeni
Schwartza (rozdział~\ref{chap:pde-fourier-Rd}).
Półnormy po lewej wyznaczają jej topologię,
co dowodzi tezy.
\end{proof}

\begin{przyklad}[Mnożenie, pochodne, potencjały Bessela]
Jeżeli $b\in C_b^\infty(\R^d)$ (wszystkie pochodne są ograniczone),
to $a(x,\xi)=b(x)$ należy do $S^0_{1,0}$ i
$\operatorname{Op}(a)u=bu$.
Dla $a(\xi)=i\xi_j\in S^1_{1,0}$ dostajemy
$\partial_{x_j}$ przez różniczkowanie wzoru
odwrotnego Fouriera. Ogólniej, dla
$P=\sum_{|\alpha|\le k}b_\alpha(x)\partial^\alpha$
z $b_\alpha\in C_b^\infty$ symbolem jest
$p(x,\xi)=\sum_{|\alpha|\le k}b_\alpha(x)(i\xi)^\alpha
\in S^k_{1,0}$; kolejność mnożenia i różniczkowania
jest dokładnie taka jak we wzorze na $P$.

Symbol $\langle\xi\rangle^m$ definiuje operator
$(I-\Delta)^{m/2}$ (potencjał Bessela dla $m<0$).
Przynależność do $S^m_{1,0}$ wynika z następującego
oszacowania: każda pochodna rzędu $|\beta|$
jest sumą wyrazów postaci
$p(\xi)(1+|\xi|^2)^{m/2-j}$, gdzie $p$
jest wielomianem stopnia co najwyżej $2j-|\beta|$.
Dowodzimy tego indukcyjnie: różniczkowanie $p$
zmniejsza stopień o jeden, a różniczkowanie potęgi
zwiększa $j$ o jeden i mnoży $p$ przez współrzędną.
Każdy taki wyraz jest ograniczony przez stałą
razy $\langle\xi\rangle^{m-|\beta|}$.
Wzór na operator wynika z definicji mnożnika
Fouriera; por. Twierdzenie~\ref{thm:pde8-Hs}.
\end{przyklad}

\begin{twierdzenie}[Oszacowanie Sobolewa dla symboli o zwartym nośniku w $x$]
\label{thm:pde31-compact-symbol}
Niech $m,s\in\R$ i $a\in S^m_{1,0}$ ma nośnik w zmiennej $x$
zawarty w ustalonym zwartym zbiorze $K\subset\R^d$.
Wtedy $\operatorname{Op}(a)$ przedłuża się jednoznacznie do
ograniczonego operatora $H^s(\R^d)\to H^{s-m}(\R^d)$.
Jego norma zależy od $m,s,d,K$ i skończonej liczby
stałych z~\eqref{eq:pde31-symbol}.
\end{twierdzenie}
\begin{proof}
Najpierw niech $u\in\mathcal S(\R^d)$.
Niech $\widehat a_x(\eta,\xi)$ oznacza transformatę
Fouriera $a$ po zmiennej $x$. Dla każdej liczby
całkowitej $k\ge0$ mamy
\[
 \langle\eta\rangle^{2k}\widehat a_x(\eta,\xi)
 =\int_{\R^d}e^{-ix\cdot\eta}(I-\Delta_x)^ka(x,\xi)
 \,\mathrm d\lambda_d(x).
\]
Całkowanie przez części nie daje wyrazów brzegowych,
bo $a(\cdot,\xi)$ ma nośnik w $K$.
Moduł prawej strony jest nie większy niż
$C_k\lambda_d(K)\langle\xi\rangle^m$,
gdzie $C_k$ zależy tylko od stałych symbolu
dla pochodnych po $x$ rzędu co najwyżej $2k$.
W szczególności dla dowolnego $N\ge0$,
wybierając $2k\ge N$, otrzymujemy
$|\widehat a_x(\eta,\xi)|\le
C_N\langle\eta\rangle^{-N}\langle\xi\rangle^m$.

Z twierdzenia Fubiniego
\[
 \widehat{\operatorname{Op}(a)u}(\zeta)
 =(2\pi)^{-d}\int_{\R^d}
 \widehat a_x(\zeta-\xi,\xi)\widehat u(\xi)
 \,\mathrm d\lambda_d(\xi).
\]
Istotnie, całka z modułu funkcji pod podwójną
całką nie przekracza
$C_{00}\lambda_d(K)\int_{\R^d}
\langle\xi\rangle^m|\widehat u(\xi)|\,\mathrm d\lambda_d(\xi)$,
co jest skończone.

Dla $s\in\R$ kładziemy
$v(\xi)=\langle\xi\rangle^s\widehat u(\xi)$.
Jądro operatora przekształcającego $v$ w
$\langle\zeta\rangle^{s-m}\widehat{\operatorname{Op}(a)u}(\zeta)$
wynosi
\[
 J(\zeta,\xi)=(2\pi)^{-d}\langle\zeta\rangle^{s-m}
 \widehat a_x(\zeta-\xi,\xi)\langle\xi\rangle^{-s}.
\]
Nierówność $\langle x+y\rangle\le
\sqrt2\langle x\rangle\langle y\rangle$
wynika z $|x+y|^2\le2|x|^2+2|y|^2$.
Stosując ją do $\zeta=\xi+(\zeta-\xi)$
dla $s-m\ge0$ albo do
$\xi=\zeta+(\xi-\zeta)$ dla $s-m<0$,
dostajemy nierówność Peetre'a
\[
 \langle\zeta\rangle^{s-m}\langle\xi\rangle^{m-s}
 \le2^{|s-m|/2}\langle\zeta-\xi\rangle^{|s-m|}.
\]
Stąd $|J(\zeta,\xi)|\le
C_N\langle\zeta-\xi\rangle^{-N+|s-m|}$.
Wybieramy $N>d+|s-m|$ i oznaczamy
\[
 A_1=\sup_\zeta\int_{\R^d}|J(\zeta,\xi)|
 \,\mathrm d\lambda_d(\xi),\qquad
 A_2=\sup_\xi\int_{\R^d}|J(\zeta,\xi)|
 \,\mathrm d\lambda_d(\zeta).
\]
Obie liczby są skończone przez podstawienie
$\eta=\zeta-\xi$ i całkowalność
$\langle\eta\rangle^{-N+|s-m|}$.
Nierówność Schura wynika tutaj bezpośrednio
z nierówności Cauchy'ego--Schwarza:
\[
 \left|\int_{\R^d}J(\zeta,\xi)v(\xi)
 \,\mathrm d\lambda_d(\xi)\right|^2
 \le A_1\int_{\R^d}|J(\zeta,\xi)|\,|v(\xi)|^2
 \,\mathrm d\lambda_d(\xi).
\]
Całkując po $\zeta$ i stosując twierdzenie Tonellego,
otrzymujemy oszacowanie kwadratu normy $L^2$
przez $A_1A_2\|v\|_2^2$.
W normach Sobolewa po obu stronach występuje
ten sam czynnik $(2\pi)^{-d/2}$, zgodny z
normalizacją Plancherela. W konsekwencji
\[
 \|\operatorname{Op}(a)u\|_{H^{s-m}}
 \le\sqrt{A_1A_2}\|u\|_{H^s}.
\]
Z gęstości $\mathcal S$ w $H^s$ i zupełności
$H^{s-m}$ (Twierdzenie~\ref{thm:pde8-Hs})
wynika istnienie i jednoznaczność ciągłego
przedłużenia. Podane oszacowania pokazują
zależność jego normy tylko od wymienionych danych.
\end{proof}

\begin{uwaga}[Zasięg twierdzenia]
Zwartego nośnika symbolu w zmiennej $x$ nie wymaga pełne
twierdzenie o ograniczoności operatorów pseudoróżniczkowych,
lecz nasze założenie umożliwia samodzielny dowód przez jądro Fouriera.
Dla symboli ogólnych potrzebna jest dodatkowa lokalizacja
lub pełna teoria ograniczoności operatorów pseudoróżniczkowych;
nie utożsamiamy tego skrótu z jej dowodem.
\end{uwaga}

\begin{wniosek}[Lokalne oszacowanie operatora]
\label{cor:pde31-local-sobolev}
Dla $a\in S^m_{1,0}$, $\rho\in C_c^\infty(\R^d)$
i $s\in\R$ operator $\rho\operatorname{Op}(a)$
ma jednoznaczne ograniczone przedłużenie
$H^s(\R^d)\to H^{s-m}(\R^d)$.
\end{wniosek}
\begin{proof}
Na przestrzeni Schwartza
$\rho\operatorname{Op}(a)=\operatorname{Op}(\rho a)$
bez żadnego wyrazu resztowego, bezpośrednio
ze wzoru~\eqref{eq:pde31-op}.
Z reguły Leibniza $\rho a\in S^m_{1,0}$:
pochodne $\rho$ są ograniczone, a pochodne
po $\xi$ działają jedynie na $a$.
Nośnik w zmiennej $x$ mieści się w
$\operatorname{supp}\rho$. Stosujemy
Twierdzenie~\ref{thm:pde31-compact-symbol}.
\end{proof}

\section{Jądro poza przekątną i lokalność osobliwości}
\begin{stwierdzenie}[Gładkość jądra poza przekątną]
\label{prop:pde31-kernel}
Niech $a\in S^m_{1,0}$. Jądro dystrybucyjne
$K_a(x,y)$ operatora~\eqref{eq:pde31-op} jest funkcją gładką
na zbiorze $\{(x,y):x\ne y\}$.
\end{stwierdzenie}
\begin{proof}
Formalnie
$K_a(x,y)=(2\pi)^{-d}\int_{\R^d}
 e^{i(x-y)\cdot\xi}a(x,\xi)\,\mathrm d\lambda_d(\xi)$.
Najpierw wstawiamy obcięcie częstotliwości
$\rho(\xi/R)\in C_c^\infty$ równe jeden dla $|\xi|\le R$.
Dla $z=x-y\ne0$ zachodzi
$e^{iz\cdot\xi}=(i|z|^2)^{-1}z\cdot\nabla_\xi
 e^{iz\cdot\xi}$. Po $N$ całkowaniach przez części
pochodne symbolu maleją jak $\langle\xi\rangle^{m-N}$.
Na każdym zwartym podzbiorze $\{|x-y|>0\}$ wybieramy
$N>m+d+k$ dla pochodnych do rzędu $k$, co zapewnia
jednostajną zbieżność całek i ich pochodnych przy $R\to\infty$.
Ponieważ $k$ jest dowolne, jądro jest gładkie poza przekątną.
\end{proof}


\begin{lemat}[Przedłużenie działania na dystrybucje o zwartym nośniku]
\label{lem:pde31-distribution-op}
Niech $a\in S^m_{1,0}(\R^d\times\R^d)$, $m\in\R$.
Operator $\operatorname{Op}(a)$ określony wzorem
\eqref{eq:pde31-op} dla funkcji Schwartza ma
kanoniczne przedłużenie
$\mathcal E'(\R^d)\to\mathcal D'(\R^d)$.
Dla $u\in\mathcal E'$ i $\varphi\in C_c^\infty$ jest ono
zdefiniowane przez
\[
\langle\operatorname{Op}(a)u,\varphi\rangle
=\langle u,\operatorname{Op}(a)^t\varphi\rangle,
\]
gdzie $t$ oznacza transpozycję względem biliniowego
parowania dystrybucji z funkcją testową.
\end{lemat}
\begin{proof}
Dla $\varphi\in C_c^\infty(\R^d)$ formalna transpozycja ma
postać
\[
 (\operatorname{Op}(a)^t\varphi)(y)
 =(2\pi)^{-d}\int_{\R^d}e^{-iy\cdot\xi}
 \left(\int_{\R^d}e^{ix\cdot\xi}a(x,\xi)\varphi(x)
 \,\mathrm d\lambda_d(x)\right)\mathrm d\lambda_d(\xi).
\]
Całka wewnętrzna jest szybko malejąca jako funkcja $\xi$.
W istocie, na zwartym nośniku $\varphi$ pochodne w $x$
symbolu $a$ są oszacowane przez $C\langle\xi\rangle^m$.
Wielokrotne całkowanie przez części względem $x$ daje dla
każdego $N$ oszacowanie $C_N\langle\xi\rangle^{m-N}$.
Ten sam argument działa po różniczkowaniu w $\xi$,
a zatem całka wewnętrzna należy do $\mathcal S(\R^d)$.
Jej odwrotna transformata jest funkcją gładką w $y$,
łącznie ze wszystkimi pochodnymi. Ponieważ $u$ ma zwarty
nośnik, działa w sposób ciągły na funkcjach gładkich
w otoczeniu swojego nośnika. Powyższy wzór definiuje
zatem dystrybucję $\operatorname{Op}(a)u$; ciągłość względem
$\varphi$ wynika z oszacowań na skończoną liczbę
półnorm funkcji testowych. Dla $u\in C_c^\infty$ dostajemy
pierwotną całkę przez twierdzenie Fubiniego.
\end{proof}

\begin{wniosek}[Własność pseudolokalna]
\label{cor:pde31-pseudolocal}
Jeśli $a\in S^m_{1,0}$, a $u\in\mathcal E'(\R^d)$
jest dystrybucją o zwartym nośniku, przy czym
$\operatorname{Op}(a)u$ rozumiemy w sensie
lematu~\ref{lem:pde31-distribution-op}, to
$\operatorname{sing\,supp}(\operatorname{Op}(a)u)
\subset\operatorname{sing\,supp}(u)$.
\end{wniosek}
\begin{proof}
Niech $x_0$ leży poza zbiorem osobliwości $u$.
Rozkładamy $u=\rho u+(1-\rho)u$, gdzie $\rho=1$ blisko $x_0$
i jest wsparta w otoczeniu gładkości $u$.
Funkcja $\rho u$ jest gładka i ma zwarty nośnik,
więc jej transformata jest szybko malejąca, a
\eqref{eq:pde31-op} można różniczkować dowolnie wiele razy.
Drugie wyrażenie daje funkcję gładką przy $x_0$ dzięki
stwierdzeniu~\ref{prop:pde31-kernel}, gdyż nośnik
$(1-\rho)u$ nie przecina otoczenia $x_0$.
\end{proof}

\section{Parametryks i regularność eliptyczna}
\begin{twierdzenie}[Parametryks operatora o stałych współczynnikach]
\label{thm:pde31-parametrix}
Niech $P(D)=\sum_{|\alpha|\le m}c_\alpha\partial^\alpha$
będzie operatorem o stałych współczynnikach o rzędzie $m\ge1$.
Załóżmy, że jego jednorodny symbol główny
$P_m(i\xi)=\sum_{|\alpha|=m}c_\alpha(i\xi)^\alpha$
nie znika dla $\xi\ne0$. Wtedy istnieje
$q\in S^{-m}_{1,0}$ i funkcja $r\in C_c^\infty(\R^d)$
takie, że
\begin{equation}\label{eq:pde31-param}
 q(\xi)P(i\xi)=1-r(\xi).
\end{equation}
Dla każdej dystrybucji temperowanej $u$ zachodzi
$u=q(D)P(D)u+r(D)u$. Operator $r(D)$ jest wygładzający,
tj. przekształca $\mathcal S'$ w $C^\infty$.
Jeśli $P(D)u\in H^s(\R^d)$ oraz $u\in H^t(\R^d)$ dla pewnego $t\in\R$,
to $u\in H^{s+m}(\R^d)$.
\end{twierdzenie}
\begin{proof}
Z eliptyczności symbolu głównego i zwartości sfery wynika
$|P_m(i\xi)|\ge c|\xi|^m$. Ponieważ symbole niższych
rzędów są $O(|\xi|^{m-1})$, mamy
$|P(i\xi)|\ge c'\langle\xi\rangle^m$ przy $|\xi|\ge R$.
Wybieramy $\rho\in C^\infty$ równą zero dla $|\xi|\le R$
i jeden dla $|\xi|\ge2R$. Kładziemy
$q(\xi)=\rho(\xi)/P(i\xi)$ i $r=1-\rho$.
Wielokrotne różniczkowanie ilorazu pokazuje
$|\partial_\xi^\beta q(\xi)|\le C_\beta
\langle\xi\rangle^{-m-|\beta|}$, a więc $q\in S^{-m}_{1,0}$.
Otrzymujemy~\eqref{eq:pde31-param}.
Ponieważ $r$ ma zwarty nośnik w częstotliwościach,
$\partial^\alpha r(D)u$ jest funkcją gładką dla każdego
wielowskaźnika $\alpha$, co wynika z reprezentacji Fouriera
dystrybucji o zwartym nośniku w przestrzeni częstotliwości.
Wreszcie $q(D):H^s\to H^{s+m}$ jest ograniczony
z definicji norm Sobolewa, a $r(D):H^t\to H^{s+m}$
dla dowolnych $s,t$, gdyż mnożnik ma zwarty nośnik i jest gładki.
\end{proof}

\begin{uwaga}[Eliptyczność mikrolokalna]
Dla gładkich operatorów o zmiennych współczynnikach analogiczny
parametryks konstruuje się iteracyjnie w rachunku symboli.
Standardowa konsekwencja to
$\operatorname{WF}(Pu)\subset\operatorname{WF}(u)$ oraz
$\operatorname{WF}(u)\subset\operatorname{WF}(Pu)\cup\operatorname{Char}(P)$,
gdzie $\operatorname{Char}(P)$ jest zbiorem zer symbolu głównego
poza przekrojem zerowym. Pełny dowód ogólnego rachunku składania
symboli wykracza poza przedstawiony tu samodzielny dowód
parametryksu dla stałych współczynników; nie stosujemy tego
wyniku do operatorów z jedynie mierzalnymi współczynnikami.
\end{uwaga}

\section{Zbiór falowy Sobolewa i powierzchnie charakterystyczne}
\begin{definicja}[Zbiór falowy Sobolewa rzędu $s$]
\label{def:pde31-wfs}
Dla $s\in\R$, $u\in\mathcal D'(U)$, para $(x_0,\xi_0)$
nie należy do $\operatorname{WF}^s(u)$, jeśli istnieje
$\chi\in C_c^\infty(U)$ równa jeden blisko $x_0$ i stożek
$\Gamma\ni\xi_0$ taki, że
$\int_\Gamma\langle\xi\rangle^{2s}|\widehat{\chi u}(\xi)|^2
\,\mathrm d\lambda_d(\xi)<\infty$.
\end{definicja}

\begin{przyklad}[Próg Sobolewa delty]
Dla $\delta_0$ lokalna transformata jest stała, więc
$\delta_0\in H^s_{\mathrm{loc}}(\R^d)$ dokładnie wtedy, gdy
$s<-d/2$. W istocie w współrzędnych biegunowych
zbieżność całki jest równoważna warunkowi
$\int_1^\infty r^{2s+d-1}\,\mathrm d\lambda_1(r)<\infty$.
Dla $s\ge-d/2$ cały stożek kierunków w punkcie zero
należy do $\operatorname{WF}^s(\delta_0)$.
\end{przyklad}

\begin{przyklad}[Charakterystyki równania falowego]
Symbol główny $\partial_t^2-c^2\Delta_x$ wynosi
$-\tau^2+c^2|\xi|^2$. Jego zbiór charakterystyczny jest stożkiem
$\tau^2=c^2|\xi|^2$, $(\tau,\xi)\ne0$.
W odróżnieniu od laplasjanu, który nie ma rzeczywistych
niezerowych kowektorów charakterystycznych, operator falowy
może przenosić osobliwości wzdłuż charakterystyk.
Twierdzenie o propagacji zbioru falowego wzdłuż bicharakterystyk
jest głębsze niż uzyskana tutaj własność pseudolokalna:
nie przedstawiamy jego pełnego dowodu bez rachunku operatorów
całkowych Fouriera.
\end{przyklad}

\section{Przybliżenie optyki geometrycznej}
\begin{stwierdzenie}[Równanie eikonalne jako warunek wysokich częstotliwości]
\label{prop:pde31-eikonal}
Niech $c>0$, a $a,\varphi$ będą rzeczywistymi funkcjami klasy
$C^2$ na otwartym podzbiorze $\R^{1+d}$.
Dla $u_\eps=a\exp(i\varphi/\eps)$ zachodzi
\[
 (\partial_t^2-c^2\Delta_x)u_\eps
 =e^{i\varphi/\eps}
 \Bigl[\eps^{-2}a(c^2|\nabla_x\varphi|^2-\varphi_t^2)
 +i\eps^{-1}\mathcal T(a,\varphi)
 +(\partial_t^2-c^2\Delta_x)a\Bigr],
\]
gdzie
$\mathcal T(a,\varphi)=2\varphi_ta_t+a\varphi_{tt}
 -2c^2\nabla_x\varphi\cdot\nabla_xa-c^2a\Delta_x\varphi$.
Zanik wyrazu rzędu $\eps^{-2}$ w punktach, gdzie $a\ne0$,
wymaga równania eikonalnego $\varphi_t^2=c^2|\nabla_x\varphi|^2$.
\end{stwierdzenie}
\begin{proof}
Różniczkujemy dwukrotnie iloczyn $ae^{i\varphi/\eps}$:
każda pochodna wykładnika wnosi $i\partial\varphi/\eps$.
Po zebraniu współczynników przy $\eps^{-2}$ otrzymujemy
$a(c^2|\nabla_x\varphi|^2-\varphi_t^2)$,
a przy $i\eps^{-1}$ dokładnie $\mathcal T(a,\varphi)$.
Pozostałe wyrazy odpowiadają operatorowi falowemu
zastosowanemu do amplitudy $a$.
\end{proof}

\section{Jak stosować analizę mikrolokalną w PDE?}
Praktyczny schemat polega na identyfikacji symbolu głównego,
wyznaczeniu zbioru charakterystycznego i zbadaniu,
czy rozważane pary $(x,\xi)$ są eliptyczne.
Tam, gdzie symbol nie znika, parametryks usuwa możliwość
występowania nowych osobliwości. Na zbiorze charakterystycznym
istotna staje się geometria przepływu hamiltonowskiego symbolu.
Przed użyciem tego schematu trzeba określić \emph{gładkość
współczynników}; metody z tego rozdziału nie zastępują teorii
De Giorgiego--Nasha--Mosera dla macierzy mierzalnych.




\chapter[Problemy odwrotne i sterowanie PDE]{Problemy odwrotne i sterowanie optymalne dla równań różniczkowych cząstkowych}
\label{chap:pde32}

\section{Problem prosty, odwrotny i sens stabilności}
Zagadnienie proste polega na wyznaczeniu rozwiązania $y$ PDE
z danych, np. źródła $f$ i współczynników operatora.
W zagadnieniu odwrotnym obserwujemy częściowo $y$ i próbujemy
odzyskać nieznane $f$, współczynniki albo warunki początkowe.
Problem sterowania polega natomiast na doborze danych wejściowych
$u$, aby wymusić pożądaną własność rozwiązania.

\begin{definicja}[Stabilność zagadnienia odwrotnego]
Niech $X,Y$ będą przestrzeniami unormowanymi, $K:X\to Y$
odwzorowaniem obserwacji (niekoniecznie liniowym),
$M\subset X$ klasą dopuszczalnych danych, a $y=Kf$
danymi bez szumu dla $f\in M$.
Zagadnienie odtworzenia $f$ nazywamy stabilnym
w sensie Lipschitza w tych normach, jeśli istnieje
stała $C>0$ spełniająca
$\|f_1-f_2\|_X\le C\|Kf_1-Kf_2\|_Y$ dla wszystkich
$f_1,f_2\in M$.
Warunek ten implikuje różnowartościowość $K|_M$,
a więc istnienie jednoznacznej odwrotności na $K(M)$.
Sama różnowartościowość nie gwarantuje stabilności.
Jeżeli $K$ jest liniowy i $M=X$, jednoznaczność
jest równoważna warunkowi $\ker K=\{0\}$.
Ciągłość odwrotności jest warunkiem słabszym niż
stabilność w sensie Lipschitza; należy odróżniać
te pojęcia, zwłaszcza dla problemów nieliniowych.
\end{definicja}

\begin{przyklad}[Odtwarzanie źródła równania eliptycznego]
\label{ex:pde32-elliptic-inverse}
Niech $\Omega\subset\R^d$ będzie ograniczonym obszarem
Lipschitzowskim, a $A=-\Delta+I$ z jednorodnym warunkiem
Dirichleta. Dla $f\in L^2(\Omega)$ niech $y=Kf=A^{-1}f$.
Twierdzenie Laxa--Milgrama zastosowane do formy
\[
 b(y,v)=\int_\Omega(\nabla y\cdot\nabla v+yv)
 \,\mathrm d\lambda_d(x)
\]
na rzeczywistej przestrzeni $H_0^1(\Omega)$ daje
istnienie i jednoznaczność rozwiązania słabego.
Testując równanie przez $v=y$, otrzymujemy
\[
 \|\nabla y\|_2^2+\|y\|_2^2
 =\int_\Omega fy\,\mathrm d\lambda_d(x)
 \le\|f\|_2\|y\|_2.
\]
Stąd $\|y\|_{H^1}\le\|f\|_2$ (dla $y=0$
nierówność zachodzi bez dzielenia).
Operator $K:L^2(\Omega)\to L^2(\Omega)$ jest zwarty,
ponieważ odwzorowuje zbiory ograniczone w $L^2$
na zbiory ograniczone w $H_0^1(\Omega)$, a zanurzenie
$H_0^1\hookrightarrow L^2$ jest zwarte.
Jeśli $e_k$ jest znormalizowaną funkcją własną laplasjanu
Dirichleta z wartością własną $\lambda_k\to\infty$, to
$Ke_k=(1+\lambda_k)^{-1}e_k$. Mamy
$\|e_k\|_2=1$ i $\|Ke_k\|_2\to0$.
Operator jest różnowartościowy: jeśli $Kf=0$,
to równanie słabe daje $\int_\Omega fv\,\mathrm d\lambda_d(x)=0$
dla wszystkich $v\in H_0^1(\Omega)$, skąd $f=0$
z gęstości $C_c^\infty(\Omega)$ w $L^2(\Omega)$.
Odwrotność $K^{-1}$ na obrazie $K$ nie jest jednak
ciągła nawet w zerze, ponieważ $Ke_k\to0$,
lecz $K^{-1}(Ke_k)=e_k$ nie zbiega do
$K^{-1}(0)=0$.
\end{przyklad}

\begin{uwaga}[Odtwarzanie warunku początkowego równania ciepła]
Dla operatora $-\Delta_D$ obserwacja rozwiązania po czasie
$T>0$ jest dana przez $K_T=\exp(T\Delta_D)$.
Na funkcjach własnych $K_Te_k=e^{-T\lambda_k}e_k$.
Współczynniki wysokoczęstotliwościowe są tłumione
wykładniczo; ich odzyskanie wzmacnia nawet niewielki szum.
Jest to przykład \emph{silnej niestabilności} zagadnienia odwrotnego.
\end{uwaga}


\begin{stwierdzenie}[Zwartość obserwacji cieplnej i brak stabilnego odwrócenia]
\label{prop:pde-v44-heat-inverse}
Niech $\Omega\subset\R^d$ będzie niepustym ograniczonym
obszarem Lipschitzowskim, $H=L^2(\Omega)$ i $T>0$.
Dla laplasjanu Dirichleta niech
$K_T=e^{T\Delta_D}:H\to H$. Wtedy $K_T$ jest
zwartym, samosprzężonym, dodatnim i różnowartościowym
operatorem. Jego odwrotność $K_T^{-1}:\operatorname{Ran}(K_T)\to H$
nie jest ograniczona, przy czym $\operatorname{Ran}(K_T)$
jest gęsty, lecz nie jest domknięty w $H$.
\end{stwierdzenie}
\begin{proof}
Z teorii spektralnej w rozdziale~\ref{chap:pde18}
istnieje baza ortonormalna $(e_k)$ przestrzeni $H$
i wartości własne $0<\lambda_k\to\infty$ operatora
$-\Delta_D$. Mamy $K_Te_k=e^{-T\lambda_k}e_k$.
Obcięcia skończonego rzędu
$K_T^{(N)}f=\sum_{k=1}^Ne^{-T\lambda_k}(f,e_k)e_k$
spełniają
$\|K_T-K_T^{(N)}\|=\sup_{k>N}e^{-T\lambda_k}\to0$.
Istotnie, tożsamość Parsevala daje ograniczenie normy
przez ten kres górny, a testowanie na $e_k$ dla $k>N$
daje nierówność przeciwną. Zatem $K_T$ jest zwarty. Dodatnie rzeczywiste
mnożniki dają samosprzężoność i dodatniość;
ponieważ żaden nie znika, jądro jest zerowe.
Dokładniej, dla $g=\sum_k g_ke_k\in H$ zachodzi
\[
 g\in\operatorname{Ran}(K_T)
 \quad\Longleftrightarrow\quad
 \sum_k e^{2T\lambda_k}|g_k|^2<\infty.
\]
Jeżeli $g=K_Tf$, to $g_k=e^{-T\lambda_k}f_k$,
a warunek wynika z tożsamości Parsevala.
Odwrotnie, gdy szereg jest zbieżny, funkcja
$f=\sum_k e^{T\lambda_k}g_ke_k$ należy do $H$
i spełnia $K_Tf=g$.
W szczególności obraz zawiera wszystkie skończone
kombinacje funkcji własnych, jest więc gęsty.
Nie jest jednak domknięty: wybierając podciąg
$\lambda_{k_j}$ taki, że $e^{T\lambda_{k_j}}\ge j$,
oraz $g_{k_j}=1/j$, a pozostałe współczynniki zerowe,
otrzymujemy $g\in H$ niespełniające warunku
przynależności do obrazu.
Gdyby odwrotność była ograniczona, mielibyśmy
$1=\|e_k\|_H\le C\|K_Te_k\|_H
=C e^{-T\lambda_k}$ dla wszystkich $k$, sprzeczność.
\end{proof}
\begin{uwaga}[Wygładzanie a identyfikowalność]
Zwartość operatora obserwacji wiąże się z wygładzaniem
parabolicznym z rozdziału~\ref{chap:pde20}; stwierdzenie
\ref{prop:pde-v44-heat-wavefront} pokazuje analogiczny
mechanizm z perspektywy zbioru falowego.
Jednoznaczność danych początkowych nie zapewnia ich
stabilnej rekonstrukcji. Twierdzenie~\ref{thm:pde32-tikhonov}
dodaje regularyzację, nie odtwarza utraconej
stabilności nieuregulowanej odwrotności.
\end{uwaga}


\section{Regularyzacja Tichonowa: istnienie i równanie normalne}
\begin{twierdzenie}[Regularyzacja kwadratowa]\label{thm:pde32-tikhonov}
Niech $X,Y$ będą rzeczywistymi przestrzeniami Hilberta,
$K\in\mathcal L(X,Y)$, $y^\delta\in Y$ oraz $\alpha>0$.
Funkcjonał
\begin{equation}\label{eq:pde32-tikhonov}
 J_\alpha(f)=\frac12\|Kf-y^\delta\|_Y^2
 +\frac\alpha2\|f\|_X^2
\end{equation}
ma dokładnie jeden minimizator $f_\alpha^\delta\in X$.
Jest on jedynym rozwiązaniem równania normalnego
\begin{equation}\label{eq:pde32-normal}
 (K^*K+\alpha I)f_\alpha^\delta=K^*y^\delta.
\end{equation}
Dla ustalonego $\alpha>0$ odwzorowanie danych
w minimizator jest liniowe i stabilne w sensie
Lipschitza, ze stałą nie większą niż $1/(2\sqrt\alpha)$.
Jeżeli $\|y^\delta-y\|_Y\le\delta$ i $y=Kf^\dagger$,
gdzie $f^\dagger$ jest rozwiązaniem równania $Kf=y$
o najmniejszej normie, to przy $\alpha\downarrow0$
i $\delta/\sqrt\alpha\to0$ zachodzi
$\|f_\alpha^\delta-f^\dagger\|_X\to0$.
\end{twierdzenie}
\begin{proof}
Funkcjonał jest ściśle wypukły, ciągły i koercywny.
Każdy ciąg minimalizujący jest ograniczony; z refleksywności
przestrzeni Hilberta wybieramy słabo zbieżny podciąg.
Normy po zastosowaniu ograniczonego operatora $K$ są słabo
dolnie półciągłe, więc granica minimalizuje $J_\alpha$.
Ścisła wypukłość składnika $\alpha\|f\|^2/2$ daje
jednoznaczność. Dla dowolnego $f,h\in X$ rozwinięcie kwadratów daje
\[
 \begin{aligned}
 J_\alpha(f+h)-J_\alpha(f)
 &= (Kf-y^\delta,Kh)_Y+\alpha(f,h)_X\\
 &\quad+\tfrac12\|Kh\|_Y^2+\tfrac\alpha2\|h\|_X^2.
 \end{aligned}
\]
Reszta jest ograniczona przez
$(\|K\|^2+\alpha)\|h\|_X^2/2$, więc część liniowa
jest pochodną Frécheta. W minimum pochodna
znika dla każdego $h$, co jest równoważne
~\eqref{eq:pde32-normal}. Odwrotnie, jeśli
równanie normalne zachodzi, to powyższa różnica
jest nieujemna dla każdego $h$, zatem $f$
jest minimizatorem.
Ponadto operator $K^*K+\alpha I$ jest odwracalny
na $X$: forma $(Kf,Kh)_Y+\alpha(f,h)_X$
jest ograniczona i koercywna ze stałą $\alpha$.
Twierdzenie Laxa--Milgrama daje dla dowolnego
$g\in X$ dokładnie jedno rozwiązanie równania
$(K^*K+\alpha I)f=g$. Testując przez $f$,
otrzymujemy $\|f\|_X\le\alpha^{-1}\|g\|_X$,
co dowodzi ograniczoności odwrotności.

Niech $A_0=K^*K\ge0$ i $f_\alpha^0=(A_0+\alpha I)^{-1}A_0f^\dagger$.
Z równania normalnego
\[
 f_\alpha^\delta-f^\dagger
 =-\alpha(A_0+\alpha I)^{-1}f^\dagger
 +(A_0+\alpha I)^{-1}K^*(y^\delta-y).
\]
Ponieważ $(A_0 f,f)_X=\|Kf\|_Y^2$, mamy
$\ker A_0=\ker K$. Zbiór rozwiązań równania $Kf=y$
jest domkniętą przestrzenią afiniczną $f^\dagger+\ker K$,
a rozwiązanie o najmniejszej normie jest jej rzutem
metrycznym z zera. W szczególności
$f^\dagger\perp\ker A_0$.
Korzystamy z twierdzenia spektralnego dla ograniczonego
operatora samosprzężonego (nie zakładamy zwartości $K$).
W przypadku rzeczywistej przestrzeni Hilberta można
stosować rachunek spektralny na jej zespoleniu:
rzeczywiste funkcje borelowskie operatora zachowują
podprzestrzeń rzeczywistą.
Niech $E$ będzie miarą spektralną operatora $A_0$
i $\mu(B)=(E(B)f^\dagger,f^\dagger)_X$ dla zbiorów
borelowskich $B$. Wówczas $\mu(\{0\})=0$, ponieważ
$E(\{0\})$ jest rzutem na $\ker A_0$.
Z rachunku funkcyjnego otrzymujemy
\[
 \|\alpha(A_0+\alpha I)^{-1}f^\dagger\|_X^2
 =\int_{[0,\|A_0\|]}
 \left(\frac{\alpha}{t+\alpha}\right)^2\,\mathrm d\mu(t).
\]
Miara $\mu$ jest skończona, funkcja podcałkowa jest
ograniczona przez $1$ i zbiega do zera $\mu$-prawie
wszędzie. Twierdzenie o zbieżności zmajoryzowanej
daje silną zbieżność pierwszego składnika do zera.
Dla drugiego składnika połóżmy
$B_\alpha=(A_0+\alpha I)^{-1}K^*$.
Ponieważ $(A_0+\alpha I)^{-1}$ jest samosprzężony,
\[
 B_\alpha B_\alpha^*
 =(A_0+\alpha I)^{-1}A_0(A_0+\alpha I)^{-1}
 =A_0(A_0+\alpha I)^{-2}.
\]
Z tożsamości $\|B_\alpha\|^2=\|B_\alpha B_\alpha^*\|$
i rachunku spektralnego wynika
\[
 \|B_\alpha\|^2
 \le\sup_{t\ge0}\frac{t}{(t+\alpha)^2}
 =\frac1{4\alpha}.
\]
Ostatnią równość daje nierówność
$(t+\alpha)^2\ge4t\alpha$, z równością dla $t=\alpha$.
Dla dowolnych danych $y_1,y_2\in Y$ ich minimizatory
$f_{\alpha,1},f_{\alpha,2}$ spełniają
$f_{\alpha,1}-f_{\alpha,2}=B_\alpha(y_1-y_2)$,
więc
\[
 \|f_{\alpha,1}-f_{\alpha,2}\|_X
 \le\frac{\|y_1-y_2\|_Y}{2\sqrt\alpha}.
\]
W szczególności błąd od szumu nie przekracza
$\delta/(2\sqrt\alpha)$. Razem z zanikiem
składnika obciążenia przy $\alpha\downarrow0$
oraz założeniem $\delta/\sqrt\alpha\to0$
daje to zbieżność z tezy.
\end{proof}

\begin{stwierdzenie}[Oszacowanie ze źródłowym warunkiem regularności]
\label{prop:pde32-source}
Przy założeniach twierdzenia~\ref{thm:pde32-tikhonov}, jeśli
$f^\dagger=K^*w$ dla pewnego $w\in Y$, to
\begin{equation}\label{eq:pde32-rate}
 \|f_\alpha^\delta-f^\dagger\|_X
 \le \frac{\sqrt\alpha}{2}\|w\|_Y
 +\frac\delta{2\sqrt\alpha}.
\end{equation}
Wybór $\alpha\asymp\delta$ daje błąd $O(\sqrt\delta)$,
przy ustalonym $\|w\|_Y$.
\end{stwierdzenie}
\begin{proof}
Z rozkładu błędu w dowodzie twierdzenia~\ref{thm:pde32-tikhonov} i warunku $f^\dagger=K^*w$
wynika, że składnik obciążenia jest równy
$-\alpha(A_0+\alpha I)^{-1}K^*w=-\alpha B_\alpha w$.
Wykazane tam oszacowanie $\|B_\alpha\|\le1/(2\sqrt\alpha)$
daje
\[
 \|\alpha B_\alpha w\|_X
 \le\frac{\sqrt\alpha}{2}\|w\|_Y .
\]
Część odpowiadającą szumowi oszacowano w poprzednim dowodzie.
\end{proof}


\begin{twierdzenie}[Ułamkowy warunek źródłowy i tempo Tichonowa]
\label{thm:pde32-fractional-source}
Niech $X,Y$ będą rzeczywistymi przestrzeniami Hilberta,
$K\in\mathcal L(X,Y)$, $A_0=K^*K$,
$0<\nu\le1$, $w\in X$ oraz
$f^\dagger=A_0^\nu w$. Przyjmijmy
$y=Kf^\dagger$ i $\|y^\delta-y\|_Y\le\delta$.
Wtedy $f^\dagger$ jest rozwiązaniem $Kf=y$
o najmniejszej normie, a minimizator
\eqref{eq:pde32-tikhonov} spełnia
\begin{equation}\label{eq:pde32-fractional-rate}
 \|f_\alpha^\delta-f^\dagger\|_X
 \le C_\nu\alpha^\nu\|w\|_X+
 \frac{\delta}{2\sqrt\alpha},\qquad
 C_\nu=\sup_{r\ge0}\frac{r^\nu}{1+r}<\infty.
\end{equation}
Jeżeli $\delta\downarrow0$ i $\|w\|_X$ pozostaje
ograniczona, to wybór
$\alpha\asymp\delta^{2/(2\nu+1)}$
daje błąd $O(\delta^{2\nu/(2\nu+1)})$.
\end{twierdzenie}
\begin{proof}
Z twierdzenia spektralnego dla dodatniego ograniczonego
operatora samosprzężonego $A_0$ definiujemy
$A_0^\nu$. Dla $\nu>0$ i $v\in\ker A_0=\ker K$
rachunek funkcyjny daje $A_0^\nu v=0$.
Ponieważ $A_0^\nu$ jest samosprzężony,
$(f^\dagger,v)_X=(w,A_0^\nu v)_X=0$.
Każde rozwiązanie równania $Kf=y$ ma postać
$f=f^\dagger+v$, gdzie $v\in\ker K$, a wtedy
$\|f\|_X^2=\|f^\dagger\|_X^2+\|v\|_X^2$.
Stąd $f^\dagger$ jest jedynym rozwiązaniem
o najmniejszej normie.
Rozkład błędu w dowodzie
\ref{thm:pde32-tikhonov} ma postać
\[
 f_\alpha^\delta-f^\dagger
 =-\alpha(A_0+\alpha I)^{-1}A_0^\nu w
 +(A_0+\alpha I)^{-1}K^*(y^\delta-y).
\]
Twierdzenie spektralne daje
\[
 \|\alpha(A_0+\alpha I)^{-1}A_0^\nu\|
 \le\sup_{t\ge0}\frac{\alpha t^\nu}{t+\alpha}
 =\alpha^\nu\sup_{r\ge0}\frac{r^\nu}{1+r}
 =C_\nu\alpha^\nu.
\]
Dla $0<\nu<1$ pochodna funkcji
$g_\nu(r)=r^\nu/(1+r)$ dla $r>0$ wynosi
\[
 g_\nu'(r)=
 \frac{r^{\nu-1}(\nu-(1-\nu)r)}{(1+r)^2}.
\]
Funkcja rośnie do $r=\nu/(1-\nu)$, a następnie
maleje; ponadto dąży do zera przy $r\downarrow0$
oraz przy $r\to\infty$.
Podstawienie punktu maksimum daje
\[
 C_\nu=\nu^\nu(1-\nu)^{1-\nu}
 \quad(0<\nu<1).
\]
Dla $\nu=1$ funkcja $r/(1+r)$ rośnie do $1$,
zatem $C_1=1$.
Część pochodząca od szumu ma normę nie większą
niż $\delta/(2\sqrt\alpha)$, jak wykazano w
\ref{thm:pde32-tikhonov}. To dowodzi oszacowania.
Aby oba człony oszacowania miały ten sam rząd,
przyjmujemy $\alpha^{\nu+1/2}\asymp\delta$.
Dokładniej, jeśli dla dostatecznie małych $\delta>0$
zachodzi $c_1\delta^{2/(2\nu+1)}\le\alpha
\le c_2\delta^{2/(2\nu+1)}$ z $c_1,c_2>0$, to
\[
 \alpha^\nu\le c_2^\nu\delta^{2\nu/(2\nu+1)},
 \qquad
 \frac{\delta}{\sqrt\alpha}\le
 c_1^{-1/2}\delta^{2\nu/(2\nu+1)}.
\]
To uzasadnia deklarowane tempo przy jednostajnym
ograniczeniu normy $w$. Nie stwierdzamy, że ten wybór
parametru minimalizuje rzeczywisty błąd dla każdych danych.
\end{proof}

\begin{uwaga}[Zakres tempa]
Jest to oszacowanie warunkowe: bez założenia
$f^\dagger\in\operatorname{Ran}(A_0^\nu)$
nie można przypisać wszystkim danym tempa
$\delta^{2\nu/(2\nu+1)}$. Ponadto w standardowej
regularyzacji Tichonowa wykładnik $\nu>1$ nie
prowadzi do analogicznego polepszenia rzędu
członu obciążenia. Istotnie, dla
$f^\dagger=A_0^\nu w$ i danych dokładnych
$y=Kf^\dagger$ mamy
\[
 \|f_\alpha^0-f^\dagger\|_X
 \le\alpha\|A_0\|^{\nu-1}\|w\|_X,
\]
gdyż $t^\nu/(t+\alpha)\le t^{\nu-1}$
dla $t>0$. Co więcej, z rachunku spektralnego
i twierdzenia o zbieżności zmajoryzowanej wynika
\[
 \frac{f_\alpha^0-f^\dagger}{\alpha}
 =-(A_0+\alpha I)^{-1}A_0^\nu w
 \longrightarrow -A_0^{\nu-1}w
 \quad\text{w }X .
\]
Rzeczywiście, mnożniki $t^\nu/(t+\alpha)$
zbiegają punktowo na widmie do $t^{\nu-1}$
(również w zerze) i są jednostajnie ograniczone
przez $\|A_0\|^{\nu-1}$.
Jeżeli $f^\dagger\ne0$, to $A_0^{\nu-1}w\ne0$,
ponieważ $A_0 A_0^{\nu-1}w=f^\dagger$.
Wówczas błąd dla danych dokładnych jest dokładnie
rzędu $\alpha$, nie zaś $o(\alpha)$.
Jest to nasycenie rzędu obciążenia; nie jest to
jeszcze twierdzenie o optymalnym tempie rekonstrukcji
dla dowolnej reguły doboru parametru i klasy szumów.
\end{uwaga}

\begin{przyklad}[Zbyt mała regularyzacja]
Dla $K_T=\exp(T\Delta_D)$ mamy
$\sigma_k=e^{-T\lambda_k}$. Dla dokładnego stanu początkowego $f^\dagger=e_k$
i obserwacji $y^\delta=(\sigma_k+\delta)e_k$
nieuregulowana odwrotność daje
$(1+\delta/\sigma_k)e_k$, zatem mnoży szum przez
$1/\sigma_k=e^{T\lambda_k}$.
Z równania normalnego i samosprzężoności $K_T$
wynika natomiast
\[
 f_\alpha^\delta
 =\frac{\sigma_k(\sigma_k+\delta)}
 {\sigma_k^2+\alpha}e_k,\qquad
 f_\alpha^\delta-f^\dagger
 =\frac{-\alpha+\sigma_k\delta}
 {\sigma_k^2+\alpha}e_k .
\]
Mnożnik szumu $\sigma_k/(\sigma_k^2+\alpha)$
jest ograniczony przez $1/(2\sqrt\alpha)$.
Dla ustalonej $\alpha>0$ i $\delta=0$ błąd ma normę
$\alpha/(\sigma_k^2+\alpha)$, która dąży do $1$
przy $k\to\infty$. Regularyzacja tłumi więc
również część sygnału, nie tylko szum. W celu uzyskania zbieżności do rozwiązania dokładnego przy
zanikającym szumie trzeba dobierać $\alpha=\alpha(\delta)$;
parametr stały ogranicza wpływ szumu, ale dla
niezerowego dokładnego rozwiązania o najmniejszej
normie pozostawia niezanikający błąd regularyzacji:
$\alpha(A_0+\alpha I)^{-1}f^\dagger\ne0$
dla $\alpha>0$, ponieważ $(A_0+\alpha I)^{-1}$
jest różnowartościowy.
\end{przyklad}

\section{Sterowanie rozłożone w obszarze: operator stanu}
Niech $\Omega\subset\R^d$ będzie ograniczonym obszarem
Lipschitzowskim, $V=H_0^1(\Omega)$ i $H=L^2(\Omega)$
nad ciałem liczb rzeczywistych. W dalszej części
przyjmujemy normę $\|v\|_V=\|\nabla v\|_2$,
równoważną standardowej normie $H^1$ na $V$
na mocy nierówności Poincarégo
$\|v\|_H\le C_P\|v\|_V$.
Normę na przestrzeni funkcjonałów liniowych i ciągłych
$V'=H^{-1}(\Omega)$
wyznaczamy względem tej normy na $V$.
Przyjmujemy mierzalną symetryczną macierz
$A\in L^\infty(\Omega;\R^{d\times d})$
i $\xi^{\mathsf T}A(x)\xi\ge\lambda|\xi|^2$ dla
każdego $\xi$ i prawie każdego $x$, z $\lambda>0$.
Niech $c\in L^\infty(\Omega)$, $c\ge0$.
Forma
$\mathfrak a(y,v)=\int_\Omega
(A\nabla y\cdot\nabla v+cyv)\,\mathrm d\lambda_d(x)$
jest symetryczna, ograniczona i koercywna na $V$.
Dokładniej, z nierówności Cauchy'ego--Schwarza
i nierówności Poincarégo mamy
\[
 \begin{aligned}
 |\mathfrak a(y,v)|
 &\le(\|A\|_{L^\infty}+C_P^2\|c\|_\infty)
 \|y\|_V\|v\|_V,\\
 \mathfrak a(v,v)&\ge\lambda\|v\|_V^2,
 \end{aligned}
\]
gdzie normę macierzy w $L^\infty$ rozumiemy
jako istotny kres górny normy operatorowej macierzy.
Niech $f\in V'$ oraz $u\in H$. Funkcjonał
$v\mapsto\langle f,v\rangle_{V',V}+(u,v)_H$
jest liniowy i ograniczony przez
$(\|f\|_{V'}+C_P\|u\|_H)\|v\|_V$.
Z twierdzenia Laxa--Milgrama istnieje zatem
dokładnie jedno $y(u)\in V$ spełniające równanie
\begin{equation}\label{eq:pde32-state}
 \mathfrak a(y(u),v)=\langle f,v\rangle_{V',V}
 +\int_\Omega uv\,\mathrm d\lambda_d(x)
 \quad\text{dla każdego }v\in V.
\end{equation}
Testując równanie przez $v=y(u)$, otrzymujemy
\[
 \|y(u)\|_V\le\lambda^{-1}
 (\|f\|_{V'}+C_P\|u\|_H).
\]
Kładziemy $S(u)=y(u)-y(0)$. Odejmując równania
dla $y(u)$ oraz $y(0)$, mamy
$\mathfrak a(S(u),v)=(u,v)_H$ dla każdego $v\in V$.
Liniowość prawej strony oraz jednoznaczność
rozwiązania dają liniowość $S$.
Z koercywności, po podstawieniu $v=S(u)$,
wynika $\|S(u)\|_V\le C_P\lambda^{-1}\|u\|_H$.
W szczególności $u\mapsto y(u)=y(0)+S(u)$
jest afiniczne i ciągłe jako odwzorowanie $H\to V$.
Traktowany jako operator $H\to H$, operator $S$
jest także zwarty dzięki zwartemu zanurzeniu $V\hookrightarrow H$.

\begin{twierdzenie}[Istnienie sterowania optymalnego]\label{thm:pde32-control}
Niech $y_d\in H$, $\alpha>0$ i $U_{\mathrm{ad}}\subset H$
będzie niepustym, domkniętym i wypukłym podzbiorem.
Wtedy problem
\begin{equation}\label{eq:pde32-opt}
 \min_{u\in U_{\mathrm{ad}}}
 \mathcal J(u),\qquad
 \mathcal J(u)=\frac12\|y(u)-y_d\|_H^2
 +\frac\alpha2\|u\|_H^2
\end{equation}
ma dokładnie jedno rozwiązanie optymalne $\bar u$.
\end{twierdzenie}
\begin{proof}
Dla dowolnego ciągu minimalizującego $(u_n)$ nierówność
$\alpha\|u_n\|_H^2/2\le\mathcal J(u_n)$ daje ograniczoność.
Istnieje słabo zbieżny podciąg $u_n\rightharpoonup\bar u$ w $H$.
Każdy domknięty wypukły podzbiór przestrzeni Hilberta jest
słabo domknięty, więc $\bar u\in U_{\mathrm{ad}}$.
Ciągłość liniowa $S$ daje słabą zbieżność
$y(u_n)\rightharpoonup y(\bar u)$ w $H$.
Słaba dolna półciągłość norm daje
$\mathcal J(\bar u)\le\liminf\mathcal J(u_n)$.
Dla $u,v\in H$ i $0\le\theta\le1$ z afiniczności
stanu i tożsamości dla kwadratu normy w przestrzeni
Hilberta otrzymujemy
\[
 \begin{aligned}
 \mathcal J((1-\theta)u+\theta v)
 &=(1-\theta)\mathcal J(u)+\theta\mathcal J(v)\\
 &\quad-\frac{\theta(1-\theta)}2
 \bigl(\|S(u-v)\|_H^2+\alpha\|u-v\|_H^2\bigr).
 \end{aligned}
\]
Dla $u\ne v$ i $0<\theta<1$ ostatni człon jest
ściśle ujemny, co daje ścisłą wypukłość.
Dwa różne minimizatory w $U_{\mathrm{ad}}$
miałyby dopuszczalny środek o mniejszej wartości
funkcjonału, co jest sprzecznością.
Ponadto każdy ciąg minimalizujący zbiega silnie
do $\bar u$: podstawiając $u=\bar u$, $v=u_n$,
$\theta=1/2$ i korzystając z minimalności
$\bar u$, otrzymujemy
\[
 \frac{\alpha}{8}\|u_n-\bar u\|_H^2
 \le \frac12\bigl(\mathcal J(u_n)-\mathcal J(\bar u)\bigr)
 \longrightarrow0.
\]
\end{proof}

\section{Równanie sprzężone i warunki optymalności}
\begin{twierdzenie}[Warunek pierwszego rzędu i rzut na sterowania dopuszczalne]
\label{thm:pde32-adjoint}
Niech $\bar u\in U_{\mathrm{ad}}$ będzie dowolnym
sterowaniem dopuszczalnym, $\bar y=y(\bar u)$,
a $\bar p\in V$ niech będzie jedynym rozwiązaniem
równania sprzężonego
\begin{equation}\label{eq:pde32-adj}
 \mathfrak a(v,\bar p)
 =\int_\Omega(\bar y-y_d)v\,\mathrm d\lambda_d(x)
 \qquad(v\in V).
\end{equation}
Wówczas $\bar u$ jest rozwiązaniem
~\eqref{eq:pde32-opt} wtedy i tylko wtedy, gdy
spełnia nierówność wariacyjną
\begin{equation}\label{eq:pde32-variational-control}
 \int_\Omega(\bar p+\alpha\bar u)(v-\bar u)
 \,\mathrm d\lambda_d(x)\ge0
 \quad\text{dla każdego }v\in U_{\mathrm{ad}},
\end{equation}
a równoważnie
$\bar u=P_{U_{\mathrm{ad}}}(-\bar p/\alpha)$,
gdzie $P_{U_{\mathrm{ad}}}$ jest rzutem metrycznym na
niepusty domknięty wypukły podzbiór przestrzeni Hilberta.
Jeżeli $U_{\mathrm{ad}}=\{v\in L^2(\Omega):a\le v\le b
\text{ prawie wszędzie}\}$ dla liczb $a<b$,
to powyższe warunki są równoważne warunkowi
$\bar u(x)=\min\{b,\max\{a,-\bar p(x)/\alpha\}\}$
dla prawie każdego $x\in\Omega$.
\end{twierdzenie}
\begin{proof}
Forma $\mathfrak a$ jest ograniczona i koercywna,
a funkcjonał $v\mapsto(\bar y-y_d,v)_H$ jest
ograniczony na $V$ przez nierówność Poincarégo.
Twierdzenie Laxa--Milgrama zastosowane do formy
$(p,v)\mapsto\mathfrak a(v,p)$ daje istnienie
i jednoznaczność $\bar p$.
Ponieważ $y(\bar u+h)=\bar y+S h$, pochodna
Frécheta stanu w kierunku $h\in H$ jest równa
$z=S h$ i spełnia
$\mathfrak a(z,v)=\int_\Omega hv\,\mathrm d\lambda_d(x)$.
Rozwijając kwadraty norm w funkcjonale, otrzymujemy
\[
 \begin{aligned}
 \mathcal J(\bar u+h)-\mathcal J(\bar u)
 &= (\bar y-y_d,S h)_H+\alpha(\bar u,h)_H\\
 &\quad+\tfrac12\|S h\|_H^2+\tfrac\alpha2\|h\|_H^2.
 \end{aligned}
\]
Pierwsze dwa składniki wyznaczają ograniczony
funkcjonał liniowy względem $h$, a suma pozostałych
jest $O(\|h\|_H^2)$ z ograniczoności $S$.
Zatem $\mathcal J'(\bar u)h=(\bar y-y_d,z)_H+
\alpha(\bar u,h)_H$. Z równań dla $z$ i $\bar p$
wynika
$(\bar y-y_d,z)_H=\mathfrak a(z,\bar p)=
(h,\bar p)_H$. Stąd $\mathcal J'(\bar u)h=
(\bar p+\alpha\bar u,h)_H$.
Załóżmy teraz, że $\bar u$ jest minimizatorem.
Dla $v\in U_{\mathrm{ad}}$ z wypukłości zbioru
dopuszczalnego wynika
$\bar u+t(v-\bar u)\in U_{\mathrm{ad}}$ dla $t\in[0,1]$.
Funkcja $t\mapsto\mathcal J(\bar u+t(v-\bar u))$
ma zatem minimum w $t=0$ na $[0,1]$.
Jej pochodna prawostronna w zerze jest nieujemna
i równa $\mathcal J'(\bar u)(v-\bar u)$,
co daje~\eqref{eq:pde32-variational-control}.
Odwrotnie, jeśli nierówność wariacyjna zachodzi,
to dla $h=v-\bar u$, $v\in U_{\mathrm{ad}}$,
powyższe rozwinięcie funkcjonału daje
\[
 \mathcal J(v)-\mathcal J(\bar u)
 \ge\tfrac\alpha2\|v-\bar u\|_H^2\ge0.
\]
Dowodzi to optymalności i jednoznaczności.
Dla niepustego domkniętego wypukłego zbioru
$C\subset H$ rzut metryczny $z=P_Cw$ istnieje
i jest jednoznaczny z twierdzenia o rzucie
w przestrzeni Hilberta. Jego charakterystyka
\[
 z=P_Cw\quad\Longleftrightarrow\quad
 z\in C,\quad (z-w,v-z)_H\ge0\quad(v\in C)
\]
wynika w jedną stronę z różniczkowania
$\|z+t(v-z)-w\|_H^2$ w $t=0$ od prawej,
a w drugą z tożsamości
$\|v-w\|_H^2=\|z-w\|_H^2+
2(z-w,v-z)_H+\|v-z\|_H^2$.
Podstawiając $C=U_{\mathrm{ad}}$ i $w=-\bar p/\alpha$
oraz dzieląc nierówność wariacyjną przez
$\alpha>0$, otrzymujemy wzór na rzut.

W przypadku ograniczeń punktowych zbiór
$U_{\mathrm{ad}}$ jest niepusty (zawiera funkcję
stałą $a$), wypukły i domknięty w $L^2$:
z ciągu zbieżnego w $L^2$ można wybrać
podciąg zbieżny prawie wszędzie, więc ograniczenia
$a\le v\le b$ przechodzą na granicę.
Dla $w\in H$ określamy
$q(x)=\min\{b,\max\{a,w(x)\}\}$.
Funkcja $q$ jest mierzalna i ograniczona, należy
do $U_{\mathrm{ad}}$, ponieważ $\lambda_d(\Omega)<\infty$.
Dla każdej liczby $r\in[a,b]$ zachodzi
$(q(x)-w(x))(r-q(x))\ge0$: gdy $w(x)<a$,
oba czynniki są nieujemne, gdy $w(x)>b$,
oba są niedodatnie, a w pozostałym przypadku
pierwszy czynnik jest zerowy.
Podstawiając $r=v(x)$ i całkując, dostajemy
$(q-w,v-q)_H\ge0$ dla wszystkich $v\in U_{\mathrm{ad}}$.
Z charakterystyki rzutu $q=P_{U_{\mathrm{ad}}}w$,
co daje podany wzór punktowy.
\end{proof}

\begin{uwaga}[Dlaczego równanie sprzężone jest użyteczne?]
Bez zmiennej sprzężonej obliczenie gradientu kosztu
wymagałoby rozwiązania nowego PDE dla każdego kierunku $h$.
Równanie~\eqref{eq:pde32-adj} wymaga jednego dodatkowego
rozwiązania PDE i daje gradient jako funkcję $\alpha u+p$.
To podstawowy mechanizm metod numerycznych sterowania.
\end{uwaga}

\section{Sterowanie równaniem ciepła i czas wsteczny}
\begin{stwierdzenie}[Równanie sprzężone dla parabolicznego sterowania rozłożonego]
\label{prop:pde32-heat-adjoint}
Niech $T>0$, $\Omega$ będzie ograniczonym obszarem
Lipschitzowskim, $u_0\in L^2(\Omega)$, $y_d\in
L^2((0,T)\times\Omega)$ oraz $\alpha>0$.
Dla $v\in L^2((0,T)\times\Omega)$ niech $y(v)$ będzie
słabym rozwiązaniem
$y_t-\Delta y=v$ w $(0,T)\times\Omega$,
$y=0$ na $(0,T)\times\partial\Omega$ i $y(0)=u_0$.
Koszt
$\mathcal J_T(v)=\frac12\|y(v)-y_d\|_{L^2((0,T)\times\Omega)}^2
+\frac\alpha2\|v\|_{L^2((0,T)\times\Omega)}^2$
ma pochodną
$\mathcal J_T'(v)h=
\int_0^T\!\int_\Omega(\alpha v+p)h
\,\mathrm d\lambda_d(x)\,\mathrm d\lambda_1(t)$,
gdzie $p$ jest rozwiązaniem równania sprzężonego
\[
 -p_t-\Delta p=y(v)-y_d,
 \qquad p(T)=0,\qquad p|_{\partial\Omega}=0.
\]
\end{stwierdzenie}
\begin{proof}
Piszemy $Q=(0,T)\times\Omega$, $V=H_0^1(\Omega)$,
$H=L^2(\Omega)$ oraz $V'=H^{-1}(\Omega)$.
Z Twierdzenia~\ref{thm:pde20-existence}
w rozdziale~\ref{chap:pde20}
dla każdego $v\in L^2(Q)$ istnieje dokładnie jedno
rozwiązanie słabe w przestrzeni
\[
 W=\{w\in L^2(0,T;V):w_t\in L^2(0,T;V')\},
\]
zanurzonej w $C([0,T];H)$. Wartości początkowe
i końcowe rozumiemy w tym ciągłym reprezentancie.
Liniowość równania i jednoznaczność dają
$y(v+h)-y(v)=z_h$, gdzie $z_h\in W$ spełnia
\[
 \langle(z_h)_t,\phi\rangle_{V',V}
 +(\nabla z_h,\nabla\phi)_{L^2(\Omega)}
 =(h,\phi)_H,\qquad z_h(0)=0.
\]
Równość zachodzi w $V'$ dla prawie każdego $t$.
Testując przez $z_h(t)$ i stosując nierówność
Poincarégo oraz nierówność Younga, otrzymujemy
\[
 \tfrac12\frac{\mathrm d}{\mathrm dt}\|z_h\|_H^2
 +\tfrac12\|z_h\|_V^2
 \le\tfrac{C_P^2}{2}\|h\|_H^2
\]
dla prawie każdego $t$; pochodna normy jest
uzasadniona formułą energii w trójce Gelfanda.
Po całkowaniu i użyciu $z_h(0)=0$ dostajemy
$\|z_h\|_{L^2(0,T;V)}\le C_P\|h\|_{L^2(Q)}$,
a więc $\|z_h\|_{L^2(Q)}\le C_P^2\|h\|_{L^2(Q)}$.
Operator $h\mapsto z_h$ jest liniowy i ograniczony,
stąd jest pochodną Frécheta stanu do $L^2(Q)$.

Aby skonstruować $p$, rozwiązujemy zwykłe
zagadnienie paraboliczne
$q_s-\Delta q=(y(v)-y_d)(T-s)$ z $q(0)=0$
i jednorodnym warunkiem brzegowym.
Prawa strona należy do $L^2(Q)$.
Podstawienie $p(t)=q(T-t)$ daje $p\in W$
spełniające równanie sprzężone z $p(T)=0$;
to samo odwrócenie czasu dowodzi jednoznaczności.
Nie jest to niestabilne odwracanie ewolucji
cieplnej z zadanym stanem obserwowanym,
lecz zwykłe zagadnienie początkowe po zmianie czasu.

Testujemy słabe równanie $z_h$ przez $p(t)$,
a słabe równanie $p$ przez $z_h(t)$ i całkujemy.
Wszystkie parowania są całkowalne z nierówności
Cauchy'ego--Schwarza w czasie. Odejmując równości,
otrzymujemy
\[
 \begin{aligned}
 &\int_0^T\bigl((h,p)_H-(y(v)-y_d,z_h)_H\bigr)
 \,\mathrm d\lambda_1(t)\\
 &\qquad=\int_0^T
 \bigl(\langle(z_h)_t,p\rangle_{V',V}
 +\langle p_t,z_h\rangle_{V',V}\bigr)
 \,\mathrm d\lambda_1(t)\\
 &\qquad=(z_h(T),p(T))_H-(z_h(0),p(0))_H=0.
 \end{aligned}
\]
W ostatnim kroku użyliśmy formuły całkowania
przez części w trójce Gelfanda, wynikającej
z Lematu~\ref{lem:pde20-energy} przez polaryzację:
dla $z_h+p$ oraz $z_h-p$ odejmujemy tożsamości
energii i dzielimy przez dwa, uzyskując tożsamość
dla przyrostu iloczynu skalarnego $(z_h,p)_H$.
Odsyłamy do rozdziału~\ref{chap:pde20}.
Wyrazy przestrzenne znoszą się
dzięki symetrii formy Dirichleta.
Rozwinięcie kwadratów w funkcjonale daje
\[
 \begin{aligned}
 \mathcal J_T(v+h)-\mathcal J_T(v)
 &= (y(v)-y_d,z_h)_{L^2(Q)}
   +\alpha(v,h)_{L^2(Q)}\\
 &\quad+\tfrac12\|z_h\|_{L^2(Q)}^2+
 \tfrac\alpha2\|h\|_{L^2(Q)}^2 .
 \end{aligned}
\]
Reszta jest $O(\|h\|_{L^2(Q)}^2)$ z oszacowania
operatora stanu. Podstawiając wykazaną tożsamość
do części liniowej, otrzymujemy pochodną Frécheta
z tezy.
\end{proof}

\section{Odtwarzanie współczynnika i granice identyfikowalności}
\begin{stwierdzenie}[Pochodna względem współczynnika dyfuzji]
\label{prop:pde32-coeff-gradient}
Niech $\Omega$ będzie ograniczonym obszarem Lipschitzowskim,
$0<\lambda\le a\le\Lambda$ prawie wszędzie,
$a\in L^\infty(\Omega)$, $f\in H^{-1}(\Omega)$.
Niech $y(a)\in H_0^1(\Omega)$ będzie jedynym rozwiązaniem
$-\diver(a\nabla y(a))=f$.
Dla $h\in L^\infty(\Omega)$ i $|\eps|$ dostatecznie małych
niech $a_\eps=a+\eps h$. Wtedy istnieje pochodna kierunkowa
$z=\lim_{\eps\to0}(y(a_\eps)-y(a))/\eps$ w $H_0^1(\Omega)$
i spełnia
\begin{equation}\label{eq:pde32-coeff-linear}
 \int_\Omega a\nabla z\cdot\nabla v
 \,\mathrm d\lambda_d(x)
 =-\int_\Omega h\nabla y(a)\cdot\nabla v
 \,\mathrm d\lambda_d(x)\qquad(v\in H_0^1(\Omega)).
\end{equation}
Jeżeli $y_d\in L^2(\Omega)$,
$\mathcal F(a)=\frac12\|y(a)-y_d\|_2^2$,
a $p\in H_0^1(\Omega)$ jest jedynym rozwiązaniem
$-\diver(a\nabla p)=y(a)-y_d$ w $H^{-1}$,
z jednorodnym warunkiem Dirichleta, to
$\mathcal F'(a)h=-\int_\Omega
h\nabla y(a)\cdot\nabla p\,\mathrm d\lambda_d(x)$.
\end{stwierdzenie}
\begin{proof}
Pracujemy z normami $\|v\|_V=\|\nabla v\|_2$
oraz odpowiadającą jej normą na $V'=H^{-1}(\Omega)$.
Dla $|\eps|\|h\|_\infty\le\lambda/2$ mamy
$a_\eps\ge\lambda/2$ prawie wszędzie.
Istnienie i jednoznaczność $y(a_\eps)$ wynikają
z twierdzenia Laxa--Milgrama.
Testując równanie przez $y(a_\eps)$, dostajemy
$\|y(a_\eps)\|_V\le2\lambda^{-1}\|f\|_{V'}$.
Oznaczmy $d_\eps=y(a_\eps)-y(a)$.
Po odjęciu równań stanów otrzymujemy
\[
 \int_\Omega a\nabla d_\eps\cdot\nabla v
 \,\mathrm d\lambda_d(x)
 =-\eps\int_\Omega h\nabla y(a_\eps)\cdot\nabla v
 \,\mathrm d\lambda_d(x).
\]
Dla $v=d_\eps$ koercywność i nierówność
Cauchy'ego--Schwarza dają
\[
 \|d_\eps\|_V
 \le 2\lambda^{-2}|\eps|\|h\|_\infty\|f\|_{V'}.
\]
Gdy $d_\eps=0$, oszacowanie zachodzi bez dzielenia.
Iloraz różnicowy $z_\eps$ spełnia
$\int_\Omega a\nabla z_\eps\cdot\nabla v\,\mathrm d\lambda_d(x)
=-\int_\Omega h\nabla y(a_\eps)\cdot\nabla v\,\mathrm d\lambda_d(x)$.
Prawa strona równania~\eqref{eq:pde32-coeff-linear}
jest ograniczonym funkcjonałem na $V$, zatem
twierdzenie Laxa--Milgrama daje dokładnie jedno $z$.
Odejmując równanie dla $z$ od równania dla
$z_\eps$ i testując przez $z_\eps-z$, dostajemy
\[
 \|z_\eps-z\|_V
 \le\lambda^{-1}\|h\|_\infty\|d_\eps\|_V
 \le2\lambda^{-3}|\eps|\|h\|_\infty^2\|f\|_{V'}.
\]
W szczególności $z_\eps\to z$ silnie w $V$.
Równanie dla $p$ ma dokładnie jedno rozwiązanie
w $V$ na mocy twierdzenia Laxa--Milgrama, ponieważ
$y(a)-y_d\in L^2(\Omega)\hookrightarrow V'$.
Ponadto
\[
 \frac{\mathcal F(a_\eps)-\mathcal F(a)}{\eps}
 =(y(a)-y_d,z_\eps)_2+\frac{\eps}{2}\|z_\eps\|_2^2.
\]
Z silnej zbieżności $z_\eps$ w $V$, a więc w $L^2$,
wynika istnienie pochodnej kierunkowej kosztu.
Z definicji $p$ i~\eqref{eq:pde32-coeff-linear}
$\mathcal F'(a)h=(y(a)-y_d,z)_2
=\int_\Omega a\nabla p\cdot\nabla z\,\mathrm d\lambda_d(x)
=-\int_\Omega h\nabla y(a)\cdot\nabla p\,\mathrm d\lambda_d(x)$.
\end{proof}


\begin{wniosek}[Różniczkowalność w normie $L^\infty$]
\label{cor:pde32-frechet}
Przy założeniach stwierdzenia~\ref{prop:pde32-coeff-gradient}
operator rozwiązania $a\mapsto y(a)$ jest różniczkowalny
w sensie Frécheta jako odwzorowanie z otwartego zbioru
$\{a\in L^\infty(\Omega):\operatorname*{ess\,inf}a>0\}$
do $H_0^1(\Omega)$. Jego pochodna w kierunku $h\in L^\infty$
jest jedynym rozwiązaniem~\eqref{eq:pde32-coeff-linear}.
W poniższym oszacowaniu normą na $H_0^1(\Omega)$
jest $\|\nabla\cdot\|_2$, a na $H^{-1}(\Omega)$
odpowiadająca jej norma funkcjonałów liniowych i ciągłych.
Dla $\|h\|_\infty\le\lambda/2$ istnieje stała $C$
zależna wyłącznie od $\lambda$, dla której
\[
 \|y(a+h)-y(a)-z_h\|_{H_0^1}
 \le C\|h\|_\infty^2\|f\|_{H^{-1}}.
\]
\end{wniosek}
\begin{proof}
Oznaczmy $y_h=y(a+h)$, $y=y(a)$ i $\delta y=y_h-y$.
Odejmowanie równań stanów daje
$\int_\Omega a\nabla\delta y\cdot\nabla v\,\mathrm d\lambda_d(x)
=-\int_\Omega h\nabla y_h\cdot\nabla v\,\mathrm d\lambda_d(x)$.
Ponieważ $a+h\ge\lambda/2$, z oszacowań energetycznych
wynikają $\|y_h\|_{H_0^1}\le2\lambda^{-1}\|f\|_{H^{-1}}$
i $\|\delta y\|_{H_0^1}\le
2\lambda^{-2}\|h\|_\infty\|f\|_{H^{-1}}$.
Kandydat na wartość pochodnej $z_h$ spełnia
$\int_\Omega a\nabla z_h\cdot\nabla v\,\mathrm d\lambda_d(x)
=-\int_\Omega h\nabla y\cdot\nabla v\,\mathrm d\lambda_d(x)$.
Dla reszty $r_h=\delta y-z_h$ otrzymujemy więc
\[
 \int_\Omega a\nabla r_h\cdot\nabla v\,\mathrm d\lambda_d(x)
 =-\int_\Omega h\nabla\delta y\cdot\nabla v\,\mathrm d\lambda_d(x).
\]
Testując przez $v=r_h$ i wykorzystując koercywność,
otrzymujemy $\|r_h\|_{H_0^1}\le
2\lambda^{-3}\|h\|_\infty^2\|f\|_{H^{-1}}$.
Liniowość $h\mapsto z_h$ wynika z liniowości
prawej strony równania~\eqref{eq:pde32-coeff-linear}
i jednoznaczności rozwiązania. Testując przez
$z_h$, otrzymujemy
$\|z_h\|_V\le\lambda^{-2}\|h\|_\infty\|f\|_{V'}$.
Zbiór współczynników z dodatnim istotnym kresem
dolnym jest otwarty w rzeczywistej przestrzeni
$L^\infty$, gdyż kula o promieniu $\lambda/2$
wokół $a$ jest w nim zawarta.
Oszacowanie reszty dowodzi różniczkowalności
Frécheta w każdym punkcie tego zbioru.
Z reguły łańcuchowej i ciągłego zanurzenia $V$
w $L^2$ wynika także różniczkowalność Frécheta
kosztu $\mathcal F$. Jego pochodna jest reprezentowana
przez funkcję $-\nabla y(a)\cdot\nabla p\in L^1(\Omega)$,
gdyż oba gradienty należą do $L^2$.
Nie wynika stąd jednak istnienie gradientu
w przestrzeni Hilberta $L^2(\Omega)$: iloczyn
dwóch funkcji z $L^2$ nie musi należeć do $L^2$.
\end{proof}

\begin{kontrprzyklad}[Brak identyfikowalności z jednej obserwacji]
Dla $f=0$ i jednorodnych warunków Dirichleta rozwiązaniem
$-\diver(a\nabla y)=0$ jest $y=0$ dla każdego
$a\in L^\infty(\Omega)$ o dodatnim istotnym kresie
dolnym. Istotnie, testowanie równania przez $y$
daje $\lambda\|\nabla y\|_2^2\le0$, gdzie
$\lambda=\operatorname*{ess\,inf}_\Omega a>0$,
a nierówność Poincarégo implikuje $y=0$.
Zatem obserwacja całej
funkcji $y$ nie wyznacza $a$ bez odpowiedniego pobudzenia.
To nie jest wada regularyzacji, lecz strukturalny brak informacji.
\end{kontrprzyklad}

\section{Sterowanie modalne i osiągalność stanów równania ciepła}
\begin{twierdzenie}[Dokładny koszt sterowania pojedynczego modu]
\label{thm:pde32-gramian}
Niech $\lambda>0$, $T>0$ i $a\in\R$. Spośród wszystkich
$h\in L^2(0,T)$, dla których rozwiązanie
$z'(t)+\lambda z(t)=h(t)$, $z(0)=0$, spełnia $z(T)=a$,
minimalną normę ma
\[
 h_*(s)=\frac{a\,e^{-\lambda(T-s)}}{W_\lambda(T)},
 \qquad W_\lambda(T)=\frac{1-e^{-2\lambda T}}{2\lambda}.
\]
Ponadto $\int_0^T|h_*(s)|^2\,\mathrm d\lambda_1(s)
=a^2/W_\lambda(T)$.
\end{twierdzenie}
\begin{proof}
Ponieważ $h\in L^2(0,T)\subset L^1(0,T)$,
funkcja
\[
 z(t)=e^{-\lambda t}\int_0^t e^{\lambda s}h(s)
 \,\mathrm d\lambda_1(s)
\]
jest bezwzględnie ciągła, spełnia równanie
prawie wszędzie i ma wartość $0$ w chwili $0$.
Jednoznaczność wynika z tego, że po pomnożeniu
równania jednorodnego przez $e^{\lambda t}$
pochodna iloczynu znika prawie wszędzie.
Niech $g(s)=e^{-\lambda(T-s)}$. Warunek końcowy
ma postać $(h,g)_{L^2}=a$, a
\[
 \|g\|_{L^2}^2
 =\int_0^T e^{-2\lambda(T-s)}\,\mathrm d\lambda_1(s)
 =W_\lambda(T)>0.
\]
Nierówność Cauchy'ego--Schwarza daje
$|a|^2\le W_\lambda(T)\|h\|_{L^2}^2$.
Funkcja $h_*=ag/W_\lambda(T)$ spełnia warunek
końcowy i osiąga dolne ograniczenie.
Dla każdego innego dopuszczalnego $h$ mamy
$(h-h_*,g)_{L^2}=0$, a stąd
\[
 \|h\|_{L^2}^2
 =\|h_*\|_{L^2}^2+\|h-h_*\|_{L^2}^2.
\]
Dowodzi to jednoznaczności minimalnego sterowania
również dla $a=0$.
\end{proof}

\begin{twierdzenie}[Przestrzeń osiągalnych stanów końcowych]
\label{thm:pde32-reachable}
Niech $\Omega$ będzie niepustym ograniczonym obszarem
Lipschitzowskim, $(e_k)$ bazą ortonormalną $L^2(\Omega)$
złożoną z funkcji własnych laplasjanu Dirichleta,
a $0<\lambda_1\le\lambda_2\le\cdots$ odpowiadającymi
wartościami własnymi liczonymi z krotnościami
(Twierdzenie~\ref{thm:pde18-compact}).
Pracujemy w rzeczywistych przestrzeniach Hilberta.
Rozważmy równanie ciepła $y_t-\Delta y=h$ z zerowym
stanem początkowym i jednorodnym warunkiem Dirichleta,
sterowane na całym obszarze przez $h\in L^2((0,T)\times\Omega)$.
Stan $a=\sum_k a_ke_k\in L^2(\Omega)$ jest osiągalny
w chwili $T>0$ wtedy i tylko wtedy, gdy
\[
 \sum_k\frac{|a_k|^2}{W_{\lambda_k}(T)}<\infty .
\]
Minimalny kwadrat normy sterowania jest równy tej
sumie i jest osiągany przez dokładnie jedno sterowanie.
\end{twierdzenie}
\begin{proof}
Dla $h_k(t)=(h(t),e_k)_{L^2(\Omega)}$ tożsamość
Parsevala i twierdzenie Tonellego dają
\[
 \|h\|_{L^2((0,T)\times\Omega)}^2
 =\sum_k\|h_k\|_{L^2(0,T)}^2.
\]
Z równania słabego testowanego przez $e_k$
wynika $y_k'+\lambda_k y_k=h_k$ w sensie
dystrybucji na $(0,T)$, gdzie $y_k(t)=(y(t),e_k)$.
Współczynniki $y_k$ są bezwzględnie ciągłe:
ich pochodne należą do $L^2(0,T)$, a warunek
$y_k(0)=0$ wynika z ciągłości stanu w $L^2$.
Jeżeli $y(T)=a$, to Twierdzenie~\ref{thm:pde32-gramian}
daje $\|h_k\|_{L^2(0,T)}^2\ge |a_k|^2/W_{\lambda_k}(T)$.
Po zsumowaniu otrzymujemy konieczność warunku
i dolne ograniczenie kosztu.

Odwrotnie, jeśli suma jest skończona, określamy
$h_k^*(s)=a_ke^{-\lambda_k(T-s)}/W_{\lambda_k}(T)$.
Skończone sumy $\sum_{k=1}^N h_k^*(s)e_k(x)$
tworzą ciąg Cauchy'ego w $L^2((0,T)\times\Omega)$,
ponieważ kwadraty norm różnic są ogonami
zbieżnego szeregu kosztów. Ich granica $h^*$
ma współczynniki $h_k^*$; istnieje dla niej
słabe rozwiązanie paraboliczne.
Każdy współczynnik końcowy tego rozwiązania
jest równy $a_k$ z twierdzenia o pojedynczym
modzie. Zupełność bazy daje $y(T)=a$ w $L^2$.
Norma $h^*$ osiąga dolne ograniczenie.
Jednoznaczność wynika z tożsamości
\[
 \|h\|_{L^2((0,T)\times\Omega)}^2
 =\|h^*\|_{L^2((0,T)\times\Omega)}^2+
 \|h-h^*\|_{L^2((0,T)\times\Omega)}^2
\]
dla każdego sterowania osiągającego $a$,
otrzymanej przez zsumowanie tożsamości
dla poszczególnych modów.
\end{proof}
\begin{uwaga}[Zasięg twierdzenia o sterowalności]
Ponieważ $\lambda_k\ge\lambda_1>0$, mamy
\[
 2\lambda_k\le W_{\lambda_k}(T)^{-1}
 \le\frac{2\lambda_k}{1-e^{-2\lambda_1 T}}.
\]
Warunek osiągalności jest więc równoważny
$\sum_k\lambda_k|a_k|^2<\infty$.
Z charakterystyki dziedziny formy Dirichleta
w dowodzie Twierdzenia~\ref{thm:pde18-minmax}
wynika, że zbiór stanów osiągalnych jest dokładnie
przestrzenią $H_0^1(\Omega)$: dla $a$ z tej
przestrzeni mamy
$\|\nabla a\|_2^2=\sum_k\lambda_k|a_k|^2$,
a skończoność szeregu implikuje przynależność do
$H_0^1$ przez zbieżność sum częściowych w normie
formy. W szczególności minimalny kwadrat normy
sterowania, oznaczony przez $M_T(a)$, spełnia
\[
 2\|\nabla a\|_2^2\le M_T(a)
 \le\frac{2}{1-e^{-2\lambda_1T}}\|\nabla a\|_2^2 .
\]
Każdy skończony układ modów jest osiągalny,
a zbiór stanów osiągalnych jest gęsty w $L^2$.
Nie jest on jednak całym $L^2$: wybierając
podciąg $\lambda_{k_j}\ge j^2$ i współczynniki
$a_{k_j}=1/j$, a pozostałe równe zero, dostajemy
$a\in L^2$ oraz $\sum_k\lambda_k|a_k|^2=\infty$.
Dowodzi to sterowalności przybliżonej, ale nie
dokładnej w $L^2$ dla opisanego operatora sterowania.
Nie rozstrzyga to sterowalności przy sterowaniu
ograniczonym do właściwego podobszaru.
\end{uwaga}

\begin{wniosek}[Sterowalność do zera przy sterowaniu na całym obszarze]
\label{cor:pde32-null-control}
Przy założeniach twierdzenia~\ref{thm:pde32-reachable}
dla każdego stanu początkowego $u_0\in L^2(\Omega)$
i każdego $T>0$ istnieje sterowanie
$h\in L^2((0,T)\times\Omega)$ prowadzące do $y(T)=0$.
Istnieje dokładnie jedno takie sterowanie o najmniejszej
normie i spełnia ono
\[
 \|h_*\|_{L^2((0,T)\times\Omega)}^2
 \le T^{-1}\|u_0\|_2^2.
\]
\end{wniosek}
\begin{proof}
Niech $u_{0,k}=(u_0,e_k)_{L^2(\Omega)}$. Z liniowości
równania stan końcowy jest sumą ewolucji swobodnej
$e^{T\Delta_D}u_0$ i stanu osiąganego z zera
przez sterowanie $h$. Aby uzyskać $y(T)=0$,
wystarczy więc osiągnąć stan o współczynnikach
$a_k=-e^{-T\lambda_k}u_{0,k}$ z zerowego stanu
początkowego. Koszt modalny spełnia
\[
 \frac{|a_k|^2}{W_{\lambda_k}(T)}
 =\frac{2\lambda_k}{e^{2T\lambda_k}-1}|u_{0,k}|^2
 \le\frac{|u_{0,k}|^2}{T},
\]
gdzie użyliśmy nierówności $e^r-1\ge r$ dla $r\ge0$.
Sumowanie i tożsamość Parsevala dają skończony
koszt nie większy niż $T^{-1}\|u_0\|_2^2$.
Istnienie i jednoznaczność sterowania o najmniejszej
normie wynikają teraz z twierdzenia
~\ref{thm:pde32-reachable}. Sterowalność do zera
nie przeczy brakowi dokładnej sterowalności w $L^2$:
swobodna ewolucja cieplna wygładza stan początkowy,
a do zera trzeba skompensować jedynie ten wygładzony
stan.
\end{proof}

\section{Podsumowanie metod i ograniczeń}
Regularyzacja stabilizuje problem odwrotny, lecz nie przywraca
informacji utraconej przez operator obserwacji.
W rozważanym kwadratowym zadaniu sterowania optymalnego
istnienie wynika z koercywności, słabej dolnej
półciągłości kosztu i słabej domkniętości zbioru
dopuszczalnego, a jednoznaczność ze ścisłej wypukłości.
Równanie sprzężone pozwala obliczyć pochodną
kosztu bez rozwiązywania osobnego PDE dla każdego kierunku.
Ogólna teoria identyfikacji współczynników z danych brzegowych
(np. problem Calderóna), nierówności obserwowalności,
sterowalność przy sterowaniu na właściwym podobszarze
oraz problemy optymalizacji nieliniowych układów PDE
wymagają dodatkowych metod. Wynik o osiągalności
stanów końcowych dotyczy tutaj równania ciepła
ze sterowaniem rozłożonym na całym obszarze,
a nie ogólnego zagadnienia sterowalności.


\appendix

\chapter{Eksperymenty numeryczne do rozdziałów 2--32}
\label{chap:pde-code}

Każdy poniższy listing jest samodzielnym programem w języku Python,
korzystającym z bibliotek \texttt{numpy} i \texttt{plotly}.
Wizualizacje są interaktywne i zapisują się do plików \texttt{.html};
pliki zawierają również lokalną kopię biblioteki JavaScript Plotly,
więc otwierają się bez połączenia z internetem.
Programy uruchamia się przez \texttt{python nazwa\_pliku.py}.
W załączonym projekcie wszystkie 29 listingów zapisano także jako
osobne pliki w katalogu \texttt{scripts/listings/}. Aby uruchomić
przykład, należy wykonać odpowiadający mu plik \texttt{.py};
kompletną tabelę nazw umieszczono w pliku
\texttt{INSTRUKCJA\_URUCHOMIENIA.md}. Eksperymenty korzystają
z NumPy w wersji co najmniej 2.0 oraz Plotly. Część przykładów,
w szczególności wizualizacja 3D wiru Taylora--Greena, wymaga ponadto
SciPy i Matplotlib. Wykonanie programu testuje jego działanie
numeryczne, ale samo nie stanowi dowodu zbieżności schematu.

\section{Charakterystyki, równanie Poissona i ciepło}
W kodzie wykorzystujemy rozwiązanie charakterystyk
równania Burgersa dla danych $-\tanh x$, tłumienie
harmonicznych modów przez $r^k$ oraz jawny profil Gaussa.

\begin{lstlisting}[style=pythonpde,caption={Burgers, funkcje harmoniczne i dyfuzja: Plotly},label={code:pde-old}]
"""Interaktywne wykresy PDE: charakterystyki, profile, dyfuzja i Poisson.
Wymagania: numpy, plotly. Zapisuje samodzielne pliki HTML (offline).
"""
from pathlib import Path
import numpy as np
import plotly.graph_objects as go

OUT = Path(__file__).resolve().parent


def save(fig: go.Figure, name: str) -> None:
    fig.update_layout(template="plotly_white")
    fig.write_html(OUT / name, include_plotlyjs=True, full_html=True)


def data(x):
    return -np.tanh(x)


def position(t, xi):
    return xi + t * data(xi)


def burgers(t, x):
    xi = np.linspace(-8.0, 8.0, 100001)
    X = position(t, xi)
    if not np.all(np.diff(X) > 0):
        raise ValueError("Brak rozwiazania klasycznego dla tego czasu")
    return data(np.interp(x, X, xi))


def characteristic_plot():
    fig = go.Figure()
    ts = np.linspace(0, 0.99, 250)
    for xi in np.linspace(-2.4, 2.4, 13):
        fig.add_trace(go.Scatter(x=position(ts, xi), y=ts,
                                 mode="lines", name=f"xi={xi:.1f}"))
    fig.update_layout(title="Charakterystyki Burgersa", xaxis_title="x",
                      yaxis_title="t")
    save(fig, "burgers_characteristics.html")


def profile_plot():
    x = np.linspace(-3.0, 3.0, 1001)
    fig = go.Figure()
    for t in (0.0, 0.5, 0.9, 0.99):
        fig.add_trace(go.Scatter(x=x, y=burgers(t, x),
                                 mode="lines", name=f"t={t:.2f}"))
    fig.update_layout(title="Burgers: profile przed utrata regularnosci",
                      xaxis_title="x", yaxis_title="u(t,x)")
    save(fig, "burgers_profiles.html")


def poisson_modes():
    k = np.arange(0, 30)
    fig = go.Figure()
    for r in (0.2, 0.5, 0.8):
        fig.add_trace(go.Scatter(x=k, y=r**k, mode="lines+markers",
                                 name=f"r={r}"))
    fig.update_layout(title="Wygaszanie modow harmonicznych",
                      xaxis_title="czestosc k", yaxis_title="amplituda r^k",
                      yaxis_type="log")
    save(fig, "poisson_modes.html")


def heat_gaussian_plot():
    x = np.linspace(-4, 4, 1200)
    fig = go.Figure()
    for t in (0.0, 0.10, 0.40, 1.20):
        u = np.exp(-x**2 / (1 + 4*t)) / np.sqrt(1 + 4*t)
        fig.add_trace(go.Scatter(x=x, y=u, mode="lines", name=f"t={t:.2f}"))
    fig.update_layout(title="Dokladne rozwiazanie Gaussowskie (kappa=1)",
                      xaxis_title="x", yaxis_title="u(t,x)")
    save(fig, "heat_gaussian_profiles.html")


if __name__ == "__main__":
    characteristic_plot()
    profile_plot()
    poisson_modes()
    heat_gaussian_plot()
    for t in (0.0, 0.5, 0.9, 0.99):
        print(f"t={t:.2f}, pochodna w x=0: {-1/(1-t):.6f}")
\end{lstlisting}

\section{Schemat falowy i kontrola dokładności}
\begin{lstlisting}[style=pythonpde,caption={Schemat różnicowy równania falowego: Plotly i test błędu},label={code:pde-wave}]
"""Jawny schemat falowy: interaktywny wykres Plotly i blad max.
Wymagania: numpy, plotly. Zapisuje plik HTML do katalogu skryptu.
"""
from pathlib import Path
import numpy as np
import plotly.graph_objects as go

c, L, N = 1.0, 1.0, 300
dx = L / N
lam = 0.8
dt = lam * dx / c
x = np.linspace(0.0, L, N + 1)
T = 0.8
steps = int(round(T / dt))


def exact(t):
    return (np.cos(2*np.pi*c*t/L)*np.sin(2*np.pi*x/L)
            + 0.4*np.cos(7*np.pi*c*t/L)*np.sin(7*np.pi*x/L))


u_prev = exact(0.0)
u = u_prev.copy()
u[1:-1] += 0.5*lam**2*(u_prev[2:] - 2*u_prev[1:-1] + u_prev[:-2])
u[0] = u[-1] = 0.0
snapshots = {0: u_prev.copy()}
for m in range(1, steps):
    u_next = np.zeros_like(u)
    u_next[1:-1] = (2*u[1:-1] - u_prev[1:-1]
                    + lam**2*(u[2:] - 2*u[1:-1] + u[:-2]))
    u_prev, u = u, u_next
    if m + 1 in (steps//3, 2*steps//3, steps):
        snapshots[m + 1] = u.copy()

fig = go.Figure()
for m, values in snapshots.items():
    fig.add_trace(go.Scatter(x=x, y=values, mode="lines", name=f"t={m*dt:.3f}"))
fig.update_layout(template="plotly_white", title="Struna: lambda=0.8",
                  xaxis_title="x", yaxis_title="u(t,x)")
fig.write_html(Path(__file__).with_name("wave_leapfrog.html"),
               include_plotlyjs=True, full_html=True)
print("CFL:", lam, "blad max:", np.max(np.abs(u - exact(steps*dt))))
\end{lstlisting}

\section{Szeregi Fouriera i ewolucje modalne}
\begin{lstlisting}[style=pythonpde,caption={Średnie Fejéra, Gibbs oraz powierzchnie ewolucji: Plotly},label={code:pde-fourier}]
"""Szeregi Fouriera i PDE na odcinku, wykresy interaktywne Plotly.
Wymagania: numpy, plotly. Zapisuje samodzielne pliki HTML.
"""
from pathlib import Path
import numpy as np
import plotly.graph_objects as go

OUT = Path(__file__).resolve().parent


def save(fig, name):
    fig.update_layout(template="plotly_white")
    fig.write_html(OUT / name, include_plotlyjs=True, full_html=True)


def fourier_vs_fejer():
    # f(x) = sign(x), x in (-pi, pi). W wartosciach skoku: f=0.
    x = np.linspace(-np.pi, np.pi, 1400)
    f = np.sign(x)
    fig = go.Figure([go.Scatter(x=x, y=f, name="dane skokowe")])
    for N in (5, 15, 45):
        ks = np.arange(1, N+1, 2)
        terms = (4/np.pi) * np.sin(np.outer(ks, x))/ks[:, None]
        fourier = terms.sum(axis=0)
        fejer = ((1-ks/(N+1))[:, None] * terms).sum(axis=0)
        fig.add_trace(go.Scatter(x=x, y=fourier, name=f"S_{N}"))
        fig.add_trace(go.Scatter(x=x, y=fejer, name=f"sigma_{N}"))
    fig.update_layout(title="Zjawisko Gibbsa i srednie Fejera",
                      xaxis_title="x", yaxis_title="wartosc")
    save(fig, "fourier_vs_fejer.html")


def heat_modes():
    L, kappa = 1.0, 0.06
    x = np.linspace(0.0, L, 500)
    t = np.linspace(0.0, 0.6, 100)
    X, T = np.meshgrid(x, t)
    U = (np.exp(-kappa*np.pi**2*T)*np.sin(np.pi*X)
         + 0.45*np.exp(-kappa*(5*np.pi)**2*T)*np.sin(5*np.pi*X))
    fig = go.Figure(data=go.Surface(x=X, y=T, z=U, colorscale="Viridis"))
    fig.update_layout(title="Rownanie ciepla: tlumienie wysokich modow",
                      scene=dict(xaxis_title="x", yaxis_title="t",
                                 zaxis_title="u(t,x)"))
    save(fig, "heat_modes.html")


def wave_modes():
    L, c = 1.0, 1.0
    x = np.linspace(0.0, L, 500)
    t = np.linspace(0.0, 2.0, 140)
    X, T = np.meshgrid(x, t)
    U = (np.cos(np.pi*c*T/L)*np.sin(np.pi*X/L)
         + 0.45*np.cos(5*np.pi*c*T/L)*np.sin(5*np.pi*X/L))
    fig = go.Figure(data=go.Surface(x=X, y=T, z=U, colorscale="Viridis"))
    fig.update_layout(title="Rownanie falowe: utrzymywanie sie modow",
                      scene=dict(xaxis_title="x", yaxis_title="t",
                                 zaxis_title="u(t,x)"))
    save(fig, "wave_modes.html")


if __name__ == "__main__":
    fourier_vs_fejer()
    heat_modes()
    wave_modes()
    print("Wykresy Plotly zapisane w", OUT)
\end{lstlisting}

\section{Przybliżone jedynki, skok i jądro cieplne}
\begin{lstlisting}[style=pythonpde,caption={Dystrybucje: przybliżone jedynki, funkcja Heaviside'a i jądro cieplne (Plotly)},label={code:pde-distributions}]
"""Przyblizone jedynki, pochodne slabego skoku i jadro ciepla.
Wymagania: numpy, plotly. Eksportuje interaktywne HTML.
"""
from pathlib import Path
import numpy as np
import plotly.graph_objects as go

OUT = Path(__file__).resolve().parent
x = np.linspace(-2.0, 2.0, 10001)
dx = x[1] - x[0]


def bump(z):
    """Gladki profil o zwartym nosniku [-1,1], bez normalizacji."""
    z = np.asarray(z)
    y = np.zeros_like(z, dtype=float)
    mask = np.abs(z) < 1.0
    y[mask] = np.exp(-1.0 / (1.0 - z[mask] ** 2))
    return y


zq = np.linspace(-1.0, 1.0, 5001)
MASS = np.trapezoid(bump(zq), zq)


def mollifier(epsilon):
    return bump(x / epsilon) / (epsilon * MASS)


def save(fig, name):
    fig.update_layout(template="plotly_white", hovermode="x unified")
    fig.write_html(OUT / name, include_plotlyjs=True, full_html=True)


def delta_and_heaviside():
    fig_delta, fig_h = go.Figure(), go.Figure()
    eps_values = (0.5, 0.25, 0.10)
    for eps in eps_values:
        rho = mollifier(eps)
        fig_delta.add_trace(go.Scatter(x=x, y=rho, mode="lines",
                                       name=f"epsilon={eps}"))
        # H * rho = integral od -nieskonczonosci do x z rho(s) ds
        h = np.r_[0.0, np.cumsum((rho[:-1] + rho[1:]) * dx / 2)]
        fig_h.add_trace(go.Scatter(x=x, y=h, mode="lines",
                                   name=f"epsilon={eps}"))
        # Test dystrybucyjny: int rho_epsilon(x) * phi(x) dx -> phi(0)=1
        phi = np.exp(-x*x)
        pairing = np.trapezoid(rho * phi, x)
        print(f"epsilon={eps:5.2f}, masa={np.trapezoid(rho,x):.7f}, "
              f"<rho_epsilon,phi>={pairing:.7f}, phi(0)=1")
    fig_h.add_trace(go.Scatter(x=x, y=(x > 0).astype(float),
                               mode="lines", name="H (skok)",
                               line=dict(dash="dot", color="gray")))
    fig_delta.update_layout(title="Gladkie aproksymacje delty Diraca",
                            xaxis_title="x", yaxis_title="rho_epsilon(x)")
    fig_h.update_layout(title="Wygladzona funkcja Heaviside'a",
                        xaxis_title="x", yaxis_title="(H * rho_epsilon)(x)")
    save(fig_delta, "pde7_delta_plotly.html")
    save(fig_h, "pde7_heaviside_plotly.html")


def heat_kernel():
    kappa = 0.15
    fig = go.Figure()
    for t in (0.01, 0.04, 0.2):
        gamma = np.exp(-x*x/(4*kappa*t)) / np.sqrt(4*np.pi*kappa*t)
        fig.add_trace(go.Scatter(x=x, y=gamma, name=f"t={t}"))
        print(f"t={t:0.3f}, masa na [-2,2]={np.trapezoid(gamma,x):.7f}")
    fig.update_layout(title="Jadro cieplne w 1D jako aproksymacja delty",
                      xaxis_title="x", yaxis_title="Gamma(t,x)")
    save(fig, "pde7_heat_kernel_plotly.html")


if __name__ == "__main__":
    delta_and_heaviside()
    heat_kernel()
\end{lstlisting}


\section[Transformata, Parseval i mnożniki]{Transformata na $\R^d$, tożsamość Parsevala i mnożniki}
\begin{lstlisting}[style=pythonpde,
caption={Transformata Gaussa, dyskretna tożsamość Parsevala, ewolucje spektralne oraz normy $H^s$ (Plotly)},
label={code:pde8-plotly}]
"""PDE, rozdzial 8: FFT, Plancherel, ewolucja ciepla i fal, Plotly.
Wymagania: numpy>=2.0, plotly>=5. Zapisuje samodzielne pliki HTML.
Transformata ciagla: F f(xi)=int exp(-i x xi) f(x) dx.
Siatka FFT jest okresowa: L=N*dx, xi=2*pi*fftfreq(N,d=dx).
"""
from pathlib import Path
import numpy as np
import plotly.graph_objects as go
OUT = Path(__file__).resolve().parent
OUT.mkdir(parents=True, exist_ok=True)

def save_plot(fig, filename):
    fig.update_layout(template="plotly_white", hovermode="x unified",
                      margin=dict(l=60, r=25, t=65, b=60))
    fig.write_html(OUT / filename, include_plotlyjs=True, full_html=True)

def gauss_fft():
    n, length, sigma = 2048, 24.0, 1.1
    dx = length / n
    x = (np.arange(n) - n//2) * dx
    f = np.exp(-x*x/(2*sigma*sigma))
    xi = np.fft.fftshift(np.fft.fftfreq(n, d=dx))*2*np.pi
    fhat = dx * np.fft.fftshift(np.fft.fft(np.fft.ifftshift(f)))
    exact = np.sqrt(2*np.pi)*sigma*np.exp(-sigma*sigma*xi*xi/2)
    mask = np.abs(xi) <= 8.0
    err = np.max(np.abs(fhat[mask]-exact[mask]))
    parseval_x = dx * np.sum(np.abs(f)**2)
    dxi = 2*np.pi/length
    parseval_k = (dxi/(2*np.pi))*np.sum(np.abs(fhat)**2)
    print("FFT vs Gauss: blad max na |xi|<=8:", f"{err:.3e}")
    print("Parseval: przestrzen", f"{parseval_x:.10f}",
          "czestotliwosc", f"{parseval_k:.10f}")
    fig = go.Figure()
    fig.add_trace(go.Scatter(x=xi[mask], y=exact[mask], mode="lines",
                             name="transformata dokladna"))
    fig.add_trace(go.Scatter(x=xi[mask], y=np.real(fhat[mask]), mode="markers",
                             marker=dict(size=4), name="FFT * dx"))
    fig.update_layout(title="Transformata Gaussa: wzor dokladny i FFT",
                      xaxis_title="czestotliwosc xi",
                      yaxis_title="F[f](xi)")
    save_plot(fig, "pde8_gauss_fft.html")

def modes_heat_wave():
    n, length = 1024, 2*np.pi
    dx = length/n
    x = np.arange(n)*dx
    frequencies = 2*np.pi*np.fft.fftfreq(n, d=dx)
    # Dane na torusie: te same dla obu rownan, predkosc poczatkowa fali 0.
    f = np.sin(2*x) + 0.55*np.cos(9*x) + 0.25*np.sin(23*x)
    fhat = np.fft.fft(f)
    kappa, c = 0.09, 1.0
    times = (0.0, 0.08, 0.3, 0.85)
    fig_heat, fig_wave = go.Figure(), go.Figure()
    for t in times:
        heat = np.fft.ifft(np.exp(-kappa*t*frequencies**2)*fhat).real
        wave = np.fft.ifft(np.cos(c*t*np.abs(frequencies))*fhat).real
        fig_heat.add_trace(go.Scatter(x=x, y=heat, name=f"t={t:g}"))
        fig_wave.add_trace(go.Scatter(x=x, y=wave, name=f"t={t:g}"))
    fig_heat.update_layout(title="Rownanie ciepla: tlumienie wysokich modow",
                           xaxis_title="x", yaxis_title="u(t,x)")
    fig_wave.update_layout(title="Rownanie falowe: mody oscyluja",
                           xaxis_title="x", yaxis_title="u(t,x)")
    save_plot(fig_heat, "pde8_heat_fft.html")
    save_plot(fig_wave, "pde8_wave_fft.html")

def h_s_delta():
    # Obciecie czestotliwosci: dla s<-1/2 calka stabilizuje sie.
    s_values = (-0.8, -0.5, -0.2)
    cutoffs = np.geomspace(2, 250, 60)
    # Dla d=1, ||delta||_{H^s}^2=(2*pi)^-1 * int(1+xi^2)^s dxi.
    xi = np.linspace(0, 250, 20001)
    fig = go.Figure()
    for s in s_values:
        vals = (1+xi*xi)**s
        increments = 0.5*(vals[1:]+vals[:-1])*np.diff(xi)
        cumulative = np.r_[0.0, np.cumsum(increments)] / np.pi
        approx = np.interp(cutoffs, xi, cumulative)
        fig.add_trace(go.Scatter(x=cutoffs, y=approx, name=f"s={s}"))
    fig.update_layout(title="Delta w H^s(R): normy z obcieciem czestotliwosci",
                      xaxis_title="odciecie |xi| <= R", xaxis_type="log",
                      yaxis_title="(2*pi)^-1 calka do R")
    save_plot(fig, "pde8_delta_Hs.html")
if __name__ == "__main__":
    gauss_fft()
    modes_heat_wave()
    h_s_delta()
    print("Wyniki zapisane w", OUT)
\end{lstlisting}

\section{Wygładzanie Sobolewa i warunki brzegowe}
Poniższy eksperyment porównuje zbieżność splotu w normie $H^1$
ze zbieżnością obcięć brzegowych jedynie w $L^2$.
\begin{lstlisting}[style=pythonpde,caption={Wygładzanie w $H^1$, pochodna słaba i obcięcia brzegowe -- Plotly},label={code:pde-sobolev9}]
"""Iteracja 9: Sobolev W^{1,2}, periodyczny splot i warunki brzegowe.

Wymagania: numpy, plotly. Zapisuje trzy samodzielne pliki HTML.
"""
from pathlib import Path
import numpy as np
import plotly.graph_objects as go

HERE = Path(__file__).resolve().parent


def save(fig, filename):
    fig.update_layout(template="plotly_white", hovermode="x unified",
                      font=dict(size=13), legend=dict(orientation="h", y=1.11))
    fig.write_html(HERE / filename, include_plotlyjs=True, full_html=True)


def periodic_mollification():
    n = 4096
    x = -np.pi + 2 * np.pi * np.arange(n) / n
    h = 2 * np.pi / n
    u = np.abs(np.sin(x))
    du = np.sign(np.sin(x)) * np.cos(x)  # pochodna slaba a.e.
    k = np.fft.fftfreq(n, d=h) * (2 * np.pi)
    uhat = np.fft.fft(u)
    epsilons = [0.40, 0.20, 0.10, 0.05, 0.025, 0.0125]
    fig = go.Figure()
    fig.add_trace(go.Scatter(x=x, y=u, name="u=|sin(x)|", line=dict(width=3)))
    errors = []
    for eps in epsilons:
        mult = np.exp(-0.5 * eps**2 * k**2)
        smooth = np.fft.ifft(mult * uhat).real
        d_smooth = np.fft.ifft(1j * k * mult * uhat).real
        err_l2 = np.sqrt(h * np.sum((smooth - u)**2))
        err_grad = np.sqrt(h * np.sum((d_smooth - du)**2))
        errors.append((eps, np.hypot(err_l2, err_grad)))
        if eps in [0.40, 0.10, 0.025]:
            fig.add_trace(go.Scatter(x=x, y=smooth, name=f"eps={eps:.3f}"))
    fig.update_layout(title="Wygladzanie w H^1(T): u(x)=|sin x|",
                      xaxis_title="x", yaxis_title="u_epsilon(x)")
    save(fig, "pde9_wygladzanie_plotly.html")

    fig2 = go.Figure()
    fig2.add_trace(go.Scatter(x=x, y=du, name="pochodna slaba", line=dict(width=2)))
    for eps in [0.40, 0.10, 0.025]:
        derivative = np.fft.ifft(1j * k * np.exp(-0.5*eps**2*k**2) * uhat).real
        fig2.add_trace(go.Scatter(x=x, y=derivative, name=f"eps={eps:.3f}"))
    fig2.update_layout(title="Wygladzanie pochodnej slabej",
                       xaxis_title="x", yaxis_title="d u_epsilon / dx")
    save(fig2, "pde9_pochodne_plotly.html")
    print("Spadek bledu H1:", errors)
    assert errors[-1][1] < errors[0][1], "Blad powinien malec przy eps->0"


def smoothstep(z):
    v = np.clip(z, 0.0, 1.0)
    return 3*v*v - 2*v*v*v


def smoothstep_deriv(z):
    v = np.clip(z, 0.0, 1.0)
    return 6*v*(1-v)


def boundary_cutoffs():
    x = (np.arange(20001) + 0.5) / 20001
    vals = []
    fig = go.Figure()
    for n in [4, 8, 16, 32, 64]:
        a, b = smoothstep(n*x), smoothstep(n*(1-x))
        d_a = n*smoothstep_deriv(n*x)
        d_b = -n*smoothstep_deriv(n*(1-x))
        chi = a*b
        d_chi = d_a*b + a*d_b
        err = np.sqrt(np.mean((chi-1)**2))
        grad = np.sqrt(np.mean(d_chi**2))
        vals.append((n, err, grad))
        fig.add_trace(go.Scatter(x=x, y=chi, name=f"n={n}"))
    fig.update_layout(title="Obciecia daza do 1 w L^2, ale nie w H^1",
                      xaxis_title="x w (0,1)", yaxis_title="chi_n(x)")
    save(fig, "pde9_slady_obciecia_plotly.html")
    print("n, ||chi-1||L2, ||chi'||L2:", vals)
    assert vals[-1][1] < vals[0][1]
    assert vals[-1][2] > vals[0][2]


if __name__ == "__main__":
    periodic_mollification()
    boundary_cutoffs()
\end{lstlisting}


\section{Koncentracja krytyczna i ciąg Mosera — Plotly}
Poniższy program wizualizuje koncentrację przy wykładniku $p^*=6$ dla $d=3$, $p=2$, porównuje skalowanie norm oraz ciąg Mosera w wymiarze $2$. Wyniki są zapisywane w trzech interaktywnych plikach HTML.
\begin{lstlisting}[style=pythonpde,caption={Koncentracja i ostrosc zanurzen Sobolewa - Plotly},label={code:pde-embed10}]
"""PDE, rozdzial 10. Interaktywne wizualizacje zanurzen Sobolewa.
Uruchom: python demo_zanurzenia_plotly.py
Pliki HTML zapisuja lokalna kopie Plotly.js i dzialaja offline.
"""
from pathlib import Path
import numpy as np
import plotly.graph_objects as go

OUT = Path(__file__).resolve().parent
EPS = [0.8, 0.4, 0.2, 0.1, 0.05]


def bump(r):
    r = np.asarray(r, dtype=float)
    v = np.zeros_like(r)
    idx = np.abs(r) < 1
    v[idx] = np.exp(-1 / (1 - r[idx] ** 2))
    return v


def bump_prime(r):
    r = np.asarray(r, dtype=float)
    v = np.zeros_like(r)
    idx = np.abs(r) < 1
    v[idx] = bump(r[idx]) * (-2 * r[idx] / (1 - r[idx] ** 2) ** 2)
    return v


# d=3, p=2. Wykladnik krytyczny p*=6, skala eps^{-1/2}.
x = np.linspace(-1, 1, 1601)
fig = go.Figure()
for eps in EPS:
    fig.add_trace(go.Scatter(
        x=x, y=eps ** (-0.5) * bump(x / eps),
        name=f"eps={eps:g}", mode="lines"))
fig.update_layout(title="Koncentracja: W1,2 w wymiarze 3 (przekroj radialny)",
                  xaxis_title="x", yaxis_title="u_eps(x)", template="plotly_white")
fig.write_html(OUT / "pde10_koncentracja.html", include_plotlyjs=True)

# Calkowanie sferyczne: 4*pi int_0^1 |v(r)|^q r2 dr.
r = np.linspace(0, 1, 50001)
q_values = [2, 4, 6, 8]
baseline = {q: (4 * np.pi * np.trapezoid(bump(r) ** q * r ** 2, r)) ** (1 / q)
            for q in q_values}
grad = np.sqrt(4 * np.pi * np.trapezoid(bump_prime(r) ** 2 * r ** 2, r))
fig = go.Figure()
for q in q_values:
    vals = [e ** (-0.5 + 3 / q) * baseline[q] for e in EPS]
    fig.add_trace(go.Scatter(x=EPS, y=vals, mode="lines+markers", name=f"L^{q}"))
fig.add_trace(go.Scatter(x=EPS, y=[grad] * len(EPS),
                         mode="lines+markers", name="norma gradientu L2"))
fig.update_layout(title="Skalowanie norm: krytyczne L6 nie zalezy od eps",
                  xaxis_title="eps", yaxis_title="norma",
                  xaxis_type="log", yaxis_type="log", template="plotly_white")
fig.write_html(OUT / "pde10_skalowanie_norm.html", include_plotlyjs=True)

# Ciag Mosera w wymiarze 2: sup rosnie, norma gradientu stala.
r = np.geomspace(1e-5, 1, 900)
fig = go.Figure()
for eps in [0.1, 0.01, 0.001, 0.0001]:
    L = np.log(1 / eps)
    val = np.minimum(np.log(1 / r), L) / np.sqrt(L)
    fig.add_trace(go.Scatter(x=r, y=val, mode="lines", name=f"eps={eps:g}"))
fig.update_layout(title="Ciag Mosera: wzrost maksimum przy stalej normie gradu w L2",
                  xaxis_title="promien r", yaxis_title="u_eps(r)",
                  xaxis_type="log", template="plotly_white")
fig.write_html(OUT / "pde10_moser.html", include_plotlyjs=True)

for q in q_values:
    assert np.isfinite(baseline[q]) and baseline[q] > 0
assert np.isfinite(grad) and grad > 0
print(f"Kontrola norm: L6(eps=0.8)={baseline[6]:.8f}; "
      f"L6(eps=0.05)={baseline[6]:.8f}; grad L2={grad:.8f}")
print("Wygenerowano: pde10_koncentracja.html, pde10_skalowanie_norm.html, pde10_moser.html")
\end{lstlisting}


\section{Galerkin, zbieżność i energia — Plotly}
\label{sec:pde11-code}
Poniższy program tworzy macierz sztywności metody elementów liniowych,
porównuje aproksymację z rozwiązaniem dokładnym, weryfikuje wzory
na błędy w $L^2$ i seminormie $H^1$ oraz wizualizuje funkcjonał energii.
Wynikami są trzy interaktywne pliki HTML, otwierane bez połączenia
z Internetem.
\begin{lstlisting}[style=pythonpde,caption={Galerkin, błędy i funkcjonał energii (Plotly)},label={code:pde11-galerkin}]
"""PDE - rozdzial 11: metoda Galerkina dla -u''=1, u(0)=u(1)=0.

Instalacja: pip install numpy plotly
Uruchom: python demo_wariacyjne_plotly.py
Wyniki: trzy samodzielne pliki HTML z Plotly (offline).
"""
from pathlib import Path

import numpy as np
import plotly.graph_objects as go

OUT = Path(__file__).resolve().parent


def exact(x):
    return x * (1.0 - x) / 2.0


def fem_solution(N):
    """Zwraca wezly i wartosci liniowych elementow skonczonych."""
    h = 1.0 / N
    diagonal = np.full(N - 1, 2.0 / h)
    K = np.diag(diagonal)
    if N > 2:
        offdiag = np.full(N - 2, -1.0 / h)
        K += np.diag(offdiag, 1) + np.diag(offdiag, -1)
    b = np.full(N - 1, h)
    U = np.linalg.solve(K, b)
    nodes = np.linspace(0.0, 1.0, N + 1)
    values = np.r_[0.0, U, 0.0]
    assert np.allclose(values, exact(nodes), atol=1e-12)
    assert np.linalg.eigvalsh(K)[0] > 0
    return nodes, values


def quadrature_errors(nodes, values):
    """Bledy L2 i H1 liczone calkowaniem Gaussa na kazdej komorce."""
    z, weights = np.polynomial.legendre.leggauss(4)
    sq_l2 = sq_h1 = 0.0
    for i in range(len(nodes) - 1):
        x0, x1 = nodes[i : i + 2]
        h = x1 - x0
        x = (x0 + x1) / 2.0 + h * z / 2.0
        slope = (values[i + 1] - values[i]) / h
        approx = values[i] + slope * (x - x0)
        sq_l2 += h / 2.0 * np.sum(weights * (exact(x) - approx) ** 2)
        sq_h1 += h / 2.0 * np.sum(weights * ((0.5 - x) - slope) ** 2)
    return np.sqrt(sq_l2), np.sqrt(sq_h1)


def save(fig, name):
    fig.update_layout(template="plotly_white", font=dict(size=13))
    fig.write_html(OUT / name, include_plotlyjs=True, full_html=True)


def main():
    Ns = [4, 8, 16, 32, 64]
    h_values, errors_l2, errors_h1 = [], [], []
    fine = np.linspace(0.0, 1.0, 1200)

    profile = go.Figure()
    profile.add_trace(go.Scatter(x=fine, y=exact(fine), name="Rozwiazanie dokladne"))

    for N in Ns:
        nodes, values = fem_solution(N)
        err_l2, err_h1 = quadrature_errors(nodes, values)
        h = 1.0 / N
        expected_l2 = h**2 / np.sqrt(120.0)
        expected_h1 = h / np.sqrt(12.0)
        assert np.isclose(err_l2, expected_l2, rtol=1e-9, atol=1e-12)
        assert np.isclose(err_h1, expected_h1, rtol=1e-9, atol=1e-12)
        h_values.append(h)
        errors_l2.append(err_l2)
        errors_h1.append(err_h1)
        if N <= 16:
            profile.add_trace(go.Scatter(x=nodes, y=values, name=f"FEM N={N}",
                                         mode="lines+markers"))
        print(f"N={N:2d}, h={h:.5f}, blad L2={err_l2:.8e}, blad H1={err_h1:.8e}")

    profile.update_layout(title="Galerkin: rozwiazanie dokladne i aproksymacje liniowe",
                          xaxis_title="x", yaxis_title="u(x)")
    save(profile, "pde11_fem_rozwiazania.html")

    convergence = go.Figure()
    convergence.add_trace(go.Scatter(x=h_values, y=errors_l2,
                                     name="||u-u_h||_L2 (rzad 2)", mode="lines+markers"))
    convergence.add_trace(go.Scatter(x=h_values, y=errors_h1,
                                     name="||u'-u_h'||_L2 (rzad 1)", mode="lines+markers"))
    convergence.update_layout(title="Zbieznosc FEM: dokladne bledy dla -u''=1",
                              xaxis_title="h", yaxis_title="norma bledu",
                              xaxis_type="log", yaxis_type="log")
    save(convergence, "pde11_fem_zbieznosc.html")

    # Dla u(x)=x(1-x)/2 mamy int_0^1 |u'|^2 dx = int_0^1 u dx = 1/12.
    t = np.linspace(-0.25, 2.25, 300)
    energy = (0.5 * t * t - t) / 12.0
    energy_plot = go.Figure()
    energy_plot.add_trace(go.Scatter(x=t, y=energy, name="J(tu)"))
    energy_plot.add_trace(go.Scatter(x=[1], y=[-1.0 / 24.0], mode="markers",
                                     marker=dict(size=11), name="Minimum t=1"))
    energy_plot.update_layout(title="Zasada Dirichleta: energia na prostej tu",
                              xaxis_title="t", yaxis_title="J(tu)")
    save(energy_plot, "pde11_energia.html")

    orders = np.polyfit(np.log(h_values), np.log(errors_l2), 1)[0], \
             np.polyfit(np.log(h_values), np.log(errors_h1), 1)[0]
    assert np.allclose(orders, (2.0, 1.0), atol=1e-9)
    print(f"Zweryfikowane rzedy: L2={orders[0]:.12f}, H1={orders[1]:.12f}")


if __name__ == "__main__":
    main()
\end{lstlisting}

\section{Spektrum, półgrupy i rezonans — Plotly}
\label{sec:pde12-code}
Poniższy program porównuje dokładne wartości własne operatora
$-\frac{\mathrm d^2}{\mathrm dx^2}$ na $(0,1)$ z widmem macierzy
różnicowej, wizualizuje ewolucję cieplną i falową tych samych
danych początkowych oraz wzrost normy rezolwenty w pobliżu
pierwszej wartości własnej. Sprawdza również rząd zbieżności
wartości własnych, równy $2$ dla zastosowanego schematu.
\begin{lstlisting}[style=pythonpde,caption={Widmo operatora Dirichleta, ewolucje i rezonans (Plotly)},label={code:pde12-operators}]
"""PDE, rozdzial 12: widmo laplasjanu Dirichleta i ewolucje.

Instalacja: pip install numpy plotly
Uruchom: python demo_operatory_plotly.py
Wyniki: trzy samodzielne pliki HTML, dzialajace offline.
"""
from pathlib import Path
import numpy as np
import plotly.graph_objects as go

OUT = Path(__file__).resolve().parent
PI = np.pi


def discrete_spectrum(n_intervals):
    h = 1.0 / n_intervals
    matrix = np.diag(np.full(n_intervals - 1, 2.0 / h**2))
    matrix += np.diag(np.full(n_intervals - 2, -1.0 / h**2), 1)
    matrix += np.diag(np.full(n_intervals - 2, -1.0 / h**2), -1)
    values = np.linalg.eigvalsh(matrix)
    j = np.arange(1, n_intervals, dtype=float)
    formula = 4.0 / h**2 * np.sin(j * PI * h / 2.0) ** 2
    assert np.allclose(values, formula, rtol=2e-12, atol=2e-11)
    return values


def save(fig, name):
    fig.update_layout(template="plotly_white", font=dict(size=13))
    fig.write_html(OUT / name, include_plotlyjs=True, full_html=True)


def main():
    sizes = [12, 24, 48, 96, 192]
    nplot = 6
    exact = (PI * np.arange(1, nplot + 1)) ** 2
    fig = go.Figure()
    fig.add_trace(go.Scatter(x=np.arange(1, nplot + 1), y=exact,
                             name="Wartosci dokladne", mode="lines+markers"))
    errors = []
    for n in sizes:
        eigen = discrete_spectrum(n)
        fig.add_trace(go.Scatter(x=np.arange(1, nplot + 1), y=eigen[:nplot],
                                 name=f"FD, N={n}", mode="lines+markers"))
        errors.append(abs(eigen[0] - PI**2))
    fig.update_layout(title="Laplasjan Dirichleta: widmo ciagle vs. macierz FD",
                      xaxis_title="Numer modu j", yaxis_title="lambda_j")
    save(fig, "pde12_widmo_plotly.html")
    h = 1.0 / np.array(sizes, dtype=float)
    order = np.polyfit(np.log(h), np.log(errors), 1)[0]
    print(f"Rzad bledu pierwszej wartosci wlasnej = {order:.4f}")
    assert 1.95 < order < 2.05

    x = np.linspace(0.0, 1.0, 1000)
    modes = np.array([1, 3, 6], dtype=float)
    coeffs = np.array([1.0, 0.35, -0.2])
    basis = np.sin(np.pi * modes[:, None] * x[None, :])
    lambdas = (np.pi * modes) ** 2
    fig2 = go.Figure()
    for t in (0.0, 0.01, 0.05, 0.12):
        y = np.sum((coeffs * np.exp(-t * lambdas))[:, None] * basis, axis=0)
        fig2.add_trace(go.Scatter(x=x, y=y, name=f"Cieplo t={t:g}"))
    for t in (0.05, 0.12):
        y = np.sum((coeffs * np.cos(t * np.sqrt(lambdas)))[:, None] * basis,
                   axis=0)
        fig2.add_trace(go.Scatter(x=x, y=y, name=f"Fala t={t:g}",
                                  line=dict(dash="dash")))
    fig2.update_layout(title="Te same dane: wygladzanie ciepla i oscylacje fali",
                       xaxis_title="x", yaxis_title="u(t,x)")
    save(fig2, "pde12_cieplo_fala_plotly.html")

    lambda1 = PI**2
    delta = np.geomspace(0.001, 4.0, 160)
    fig3 = go.Figure()
    fig3.add_trace(go.Scatter(x=delta, y=1.0 / delta,
                              name="|| (A-(lambda_1+delta)I)^-1 phi_1 ||"))
    fig3.update_layout(title="Alternatywa Fredholma: rezonans blisko lambda_1",
                       xaxis_title="delta = |zeta-lambda_1|",
                       yaxis_title="Norma rozwiazania", xaxis_type="log",
                       yaxis_type="log")
    save(fig3, "pde12_rezolwenta_plotly.html")
    print("Zweryfikowano formula FD, blad O(h^2) i trzy wizualizacje Plotly.")


if __name__ == "__main__":
    main()
\end{lstlisting}


\section{Eliptyczne zagadnienie transmisyjne — Plotly}
\label{sec:pde13-code}
Program rozwiązuje metodą elementów skończonych zadanie
$-(a u')'=1$ na $(-1,1)$, z warunkami Dirichleta i skokiem
współczynnika $a$ w zerze. Porównuje rozwiązanie numeryczne ze
wzorem~\eqref{eq:pde13-exact}, przedstawia strumień
$q=-a u'$ oraz bada błędy w normie $L^2$ i seminormie $H^1$.
Sprawdza także dokładne wzory na błędy interpolacji:
$\|u-u_h\|_{L^2}=h^2\sqrt{17/1920}$ i
$\|(u-u_h)'\|_{L^2}=h\sqrt{17/192}$.
\begin{lstlisting}[style=pythonpde,caption={Transmisja i zbieżność metody Galerkina (Plotly)},label={code:pde13-transmission}]
"""PDE, rozdzial 13: dywergencyjne zagadnienie transmisyjne.

Instalacja: pip install numpy plotly
Uruchom: python demo_eliptyczne_plotly.py
Wyniki: samodzielne pliki HTML (dzialaja offline).
"""
from pathlib import Path
import numpy as np
import plotly.graph_objects as go

OUT = Path(__file__).resolve().parent


def exact(x):
    x = np.asarray(x, dtype=float)
    return np.where(x <= 0, -x*x/2 - 0.3*x + 0.2,
                    -x*x/8 - 0.075*x + 0.2)


def exact_deriv(x):
    x = np.asarray(x, dtype=float)
    return np.where(x < 0, -x - 0.3, -x/4 - 0.075)


def solve_tridiagonal(subdiag, diag, rhs):
    """Algorytm Thomasa: macierz symetryczna trojdiagonalna."""
    b, d, c = diag.copy(), rhs.copy(), subdiag.copy()
    for i in range(1, len(b)):
        fac = c[i - 1] / b[i - 1]
        b[i] -= fac * c[i - 1]
        d[i] -= fac * d[i - 1]
    x = np.empty_like(d)
    x[-1] = d[-1] / b[-1]
    for i in range(len(x)-2, -1, -1):
        x[i] = (d[i] - c[i] * x[i+1]) / b[i]
    return x


def fem_solution(n_per_half):
    """Liniowe FEM dla -(a u')'=1, u(-1)=u(1)=0.

    Punkt 0 jest wezelm siatki. a=1 po lewej i a=4 po prawej.
    """
    nodes = np.linspace(-1., 1., 2 * n_per_half + 1)
    h = nodes[1] - nodes[0]
    a_elem = np.where((nodes[:-1] + nodes[1:]) / 2 < 0, 1., 4.)
    diag = (a_elem[:-1] + a_elem[1:]) / h
    subdiag = -a_elem[1:-1] / h
    rhs = np.full(len(nodes)-2, h)
    sol = np.r_[0., solve_tridiagonal(subdiag, diag, rhs), 0.]
    assert np.max(np.abs(sol - exact(nodes))) < 5e-10
    return nodes, sol, a_elem


def errors(nodes, sol):
    z, w = np.polynomial.legendre.leggauss(5)
    l2sq = h1sq = 0.
    for j in range(len(nodes) - 1):
        x0, x1 = nodes[j:j+2]
        h = x1 - x0
        xx = (x0+x1)/2 + h*z/2
        approx = sol[j] + (sol[j+1]-sol[j])*(xx-x0)/h
        slope = (sol[j+1]-sol[j])/h
        l2sq += (h/2) * np.dot(w, (approx-exact(xx))**2)
        h1sq += (h/2) * np.dot(w, (slope-exact_deriv(xx))**2)
    return np.sqrt(l2sq), np.sqrt(h1sq)


def save(fig, name):
    fig.update_layout(template="plotly_white", font=dict(size=13),
                      hovermode="x unified")
    fig.write_html(OUT / name, include_plotlyjs=True, full_html=True)


def main():
    hs, l2errors, h1errors = [], [], []
    fig_sol = go.Figure()
    x = np.linspace(-1, 1, 1201)
    fig_sol.add_trace(go.Scatter(x=x, y=exact(x), mode="lines",
                                 name="Rozwiazanie dokladne"))
    fig_flux = go.Figure()
    fig_flux.add_trace(go.Scatter(x=x, y=x + 0.3, mode="lines",
                                  name="Dokladny strumien q=-a u'"))
    for N in [4, 8, 16, 32, 64, 128]:
        nodes, vals, a_elem = fem_solution(N)
        h = nodes[1]-nodes[0]
        l2, h1 = errors(nodes, vals)
        assert np.isclose(l2, h*h * np.sqrt(17/1920), rtol=5e-8, atol=5e-11)
        assert np.isclose(h1, h * np.sqrt(17/192), rtol=5e-8, atol=5e-11)
        hs.append(h)
        l2errors.append(l2)
        h1errors.append(h1)
        if N in (4, 16):
            fig_sol.add_trace(go.Scatter(x=nodes, y=vals, mode="lines+markers",
                                         name=f"FEM, N={N} na polowe"))
        if N in (4, 16):
            mid = (nodes[:-1]+nodes[1:])/2
            qfem = -a_elem*np.diff(vals)/h
            fig_flux.add_trace(go.Scatter(x=mid, y=qfem, mode="markers",
                                          name=f"q_h w srodkach, N={N}"))
        print(f"N={N:3d} h={h:.5f} L2={l2:.9e} H1-semi={h1:.9e}")
    fig_sol.add_vline(x=0, line_dash="dash", annotation_text="Interfejs")
    fig_sol.update_layout(title="Transmisja przez dwa osrodki: rozwiazanie slabe",
                          xaxis_title="x", yaxis_title="u(x)")
    save(fig_sol, "pde13_transmisja_rozwiazanie.html")
    fig_flux.add_vline(x=0, line_dash="dash", annotation_text="Interfejs")
    fig_flux.update_layout(title="Strumien ciagly w rozwiazaniu dokladnym",
                           xaxis_title="x", yaxis_title="q=-a u'(x)")
    save(fig_flux, "pde13_transmisja_strumien.html")
    fig_error = go.Figure()
    fig_error.add_trace(go.Scatter(x=hs, y=l2errors, mode="lines+markers",
                                   name="Blad L2, rzad 2"))
    fig_error.add_trace(go.Scatter(x=hs, y=h1errors, mode="lines+markers",
                                   name="Blad seminormy H1, rzad 1"))
    fig_error.update_layout(title="Zbieznosc FEM dla rownania eliptycznego",
                            xaxis_title="h", yaxis_title="Blad",
                            xaxis_type="log", yaxis_type="log")
    save(fig_error, "pde13_transmisja_zbieznosc.html")
    p2 = np.polyfit(np.log(hs), np.log(l2errors), 1)[0]
    p1 = np.polyfit(np.log(hs), np.log(h1errors), 1)[0]
    assert np.isclose(p2, 2., atol=1e-6) and np.isclose(p1, 1., atol=1e-6)
    print(f"Rzedy: L2={p2:.6f}, H1-semi={p1:.6f}; testy zaliczone.")


if __name__ == "__main__":
    main()
\end{lstlisting}

\section{Zasady maksimum, nierówność Harnacka i lemat Hopfa — Plotly}
Kod oblicza harmoniczne przedłużenie danych brzegowych w dysku,
weryfikuje nierówność Harnacka oraz pochodną normalną w przykładzie
ilustrującym lemat Hopfa. Porównanie jądra Poissona z rozwiązaniem
analitycznym jest jednocześnie testem poprawności normalizacji jądra.
\begin{lstlisting}[style=pythonpde,caption={Jądro Poissona, nierówność Harnacka i lemat Hopfa (Plotly)},label={code:pde14-plotly}]
"""PDE, rozdzial 14: jadro Poissona, Harnack i Hopf.

Instalacja: pip install numpy plotly
Uruchomienie: python demo_maksimum_harnack_plotly.py
Wykresy HTML zawieraja Plotly offline.
"""
from pathlib import Path
import numpy as np
import plotly.graph_objects as go

OUT = Path(__file__).resolve().parent
TWO_PI = 2.0 * np.pi


def save(fig, filename):
    fig.update_layout(template="plotly_white", font=dict(size=13))
    fig.write_html(OUT / filename, include_plotlyjs=True, full_html=True)


def boundary(phi):
    return 0.5 + np.cos(2.0 * phi) + 0.3 * np.sin(3.0 * phi)


def poisson_disk(r, theta, phi):
    # Jadrowa gestosc wzgledem miary dlugosci na okregu jednostkowym.
    return (1.0 - r * r) / (TWO_PI * (1.0 + r * r - 2.0 * r * np.cos(theta - phi)))


def harmonic_exact(r, theta):
    return 0.5 + r**2 * np.cos(2.0 * theta) + 0.3 * r**3 * np.sin(3.0 * theta)


def main():
    phi = np.linspace(0.0, TWO_PI, 2048, endpoint=False)
    theta = np.linspace(0.0, TWO_PI, 160)
    g = boundary(phi)
    fig = go.Figure()
    max_error = 0.0
    for r in [0.0, 0.3, 0.6, 0.85]:
        p = poisson_disk(r, theta[:, None], phi[None, :])
        vals = TWO_PI * np.mean(p * g[None, :], axis=1)
        exact = harmonic_exact(r, theta)
        error = np.max(np.abs(vals - exact))
        max_error = max(max_error, error)
        fig.add_trace(go.Scatter(x=theta, y=vals, mode="lines", name=f"r={r:g}"))
        assert error < 1e-9, (r, error)
    fig.update_layout(title="Przedluzenie harmoniczne przez jadro Poissona",
                      xaxis_title="kat theta", yaxis_title="u(r, theta)")
    save(fig, "pde14_poisson_plotly.html")
    print(f"Poisson: maksymalny blad dyskretyzacji = {max_error:.3e}")

    # Jawna funkcja harmoniczna: u(x,y)=1+0.6*x w B_1(0).
    r = 0.5
    t = np.linspace(0.0, TWO_PI, 500)
    vals = 1.0 + 0.6 * r * np.cos(t)
    ratio = np.max(vals) / np.min(vals)
    lower = (1.0 - r**2) / (1.0 + r)**2
    upper = 1.0 / (1.0 - r)**2
    assert np.min(vals) >= lower and np.max(vals) <= upper
    fig = go.Figure()
    fig.add_trace(go.Scatter(x=t, y=vals, name="u(0.5 cos(theta), 0.5 sin(theta))"))
    fig.add_hline(y=lower, line_dash="dash", annotation_text="dolne Harnack")
    fig.add_hline(y=upper, line_dash="dash", annotation_text="gorne Harnack")
    fig.update_layout(title=f"Harnack: rzeczywisty sup/inf = {ratio:.4f}",
                      xaxis_title="kat theta", yaxis_title="u")
    save(fig, "pde14_harnack_plotly.html")
    print(f"Harnack: sup/inf = {ratio:.6f}, zakres teoretyczny = [{lower:.3f}, {upper:.3f}]")

    # Dodatnia funkcja harmoniczna u(x,y)=1-x; w x0=(1,0) minimum 0.
    x = np.linspace(0.75, 1.0, 150)
    u = 1.0 - x
    fig = go.Figure()
    fig.add_trace(go.Scatter(x=x, y=u, name="u(x,0)=1-x"))
    fig.add_trace(go.Scatter(x=x, y=(1.0-x), mode="markers", name="styczna w x0", marker=dict(size=3)))
    fig.update_layout(title="Lemat Hopfa: u(1,0)=0, pochodna normalna = -1",
                      xaxis_title="x (promien do brzegu)", yaxis_title="u(x,0)")
    save(fig, "pde14_hopf_plotly.html")
    assert np.allclose(np.gradient(u, x), -1.0, atol=1e-10)
    print("Hopf: pochodna zewnetrzna = -1 (blad < 1e-10)")


if __name__ == "__main__":
    main()
\end{lstlisting}

\section{Ilorazy różnicowe i skok pochodnej — Plotly}
Przykład pokazuje, jak skok pochodnej prowadzi do nieograniczonej normy dyskretnych drugich pochodnych.
\begin{lstlisting}[style=pythonpde,caption={Brak regularnosci H2 dla transmisji (Plotly)},label={code:pde15-plotly}]
"""PDE, rozdzial 15: H2-regularnosc i transmisja (Plotly).
Instalacja: pip install numpy plotly
"""
from pathlib import Path
import numpy as np
import plotly.graph_objects as go

OUT = Path(__file__).resolve().parent

def save(fig, name):
    fig.update_layout(template="plotly_white", font=dict(size=13))
    fig.write_html(OUT / name, include_plotlyjs=True, full_html=True)

def transmission(x):
    return np.where(x <= 0.5, x, 0.5 + 0.5*(x-0.5))

def main():
    x = np.linspace(0., 1., 1001)
    u = transmission(x)
    fig = go.Figure()
    fig.add_trace(go.Scatter(x=x, y=u, name="u: ciagla, z zalamaniem"))
    fig.add_trace(go.Scatter(x=x, y=x*(1-x), name="v: gladka"))
    fig.add_vline(x=0.5, line_dash="dot", annotation_text="granica osrodkow")
    fig.update_layout(title="Regularnosc rozwiazan: gladkie i transmisyjne",
                      xaxis_title="x", yaxis_title="wartosc")
    save(fig, "pde15_regularnosc_plotly.html")

    Ns = np.array([20, 40, 80, 160, 320, 640])
    hs = 1.0/Ns
    norms = []
    for N in Ns:
        xgrid = np.linspace(0, 1, N+1)
        uh = transmission(xgrid)
        second = (uh[2:] - 2*uh[1:-1] + uh[:-2]) / (1/N)**2
        norms.append(np.sqrt(np.sum(second**2) / N))
    norms = np.array(norms)
    reference = 0.5 / np.sqrt(hs)
    assert np.allclose(norms, reference, rtol=1e-10, atol=1e-10)
    slope = np.polyfit(np.log(hs), np.log(norms), 1)[0]
    assert abs(slope + .5) < 1e-10
    fig = go.Figure()
    fig.add_trace(go.Scatter(x=hs, y=norms, mode="lines+markers", name="dyskretna norma H2"))
    fig.add_trace(go.Scatter(x=hs, y=reference, mode="lines", name="0.5 / sqrt(h)"))
    fig.update_layout(title=f"Brak H2: norma drugiej roznicy ~ h^({slope:.2f})",
                      xaxis_type="log", yaxis_type="log",
                      xaxis_title="h", yaxis_title="norma L2 drugiej roznicy")
    save(fig, "pde15_skok_h2_plotly.html")
    print(f"Nachylenie log-log: {slope:.8f} (oczekiwane -0.5)")

if __name__ == "__main__":
    main()
\end{lstlisting}

\section[Drugie pochodne i ostrość oszacowań Schaudera — Plotly]{Drugie pochodne, brak oszacowania w $L^\infty$ i ostrość oszacowań Schaudera — Plotly}
Poniższy eksperyment na torusie dwuwymiarowym ilustruje kontrast
między normą średniokwadratową a normą supremum Hessianu.
Jest to eksperyment periodyczny, nie dowód twierdzenia dla $\R^d$.
Na końcu kod przedstawia dokładny przykład ostrości wykładnika Höldera.
\begin{lstlisting}[style=pythonpde,caption={Rachunek spektralny dla Hessianu i wykresy ilustrujące oszacowania Schaudera (Plotly)},label={code:pde16-plotly}]
"""Rozdzial 16: regularnosc W2p i oszacowania Schaudera; interaktywne wykresy Plotly.
Instalacja: pip install numpy plotly
Uruchomienie: python demo_calderon_schauder_plotly.py
"""
from pathlib import Path
import numpy as np
import plotly.graph_objects as go

OUT = Path(__file__).resolve().parent


def save(fig, name):
    fig.update_layout(template='plotly_white', font=dict(size=13))
    fig.write_html(OUT / name, include_plotlyjs=True, full_html=True)


def mixed_hessian(N):
    # Kwadrat w torusie [0,2pi)^2; z jego funkcji wskaznikowej
    # usuwamy srednia, by rozwiazanie -Delta u=f istnialo okresowo.
    x = 2 * np.pi * np.arange(N) / N
    square = ((x[:, None] < np.pi) & (x[None, :] < np.pi)).astype(float)
    f = square - square.mean()
    k = np.fft.fftfreq(N, d=1.0/N)  # calkowite czestosci
    K1, K2 = np.meshgrid(k, k, indexing='ij')
    k2 = K1*K1 + K2*K2
    m = np.zeros_like(k2, dtype=float)
    np.divide(-K1*K2, k2, out=m, where=k2 != 0)
    d12u = np.fft.ifft2(np.fft.fft2(f)*m).real
    return x, f, d12u


def demo_fft():
    N = 256
    x, f, d12u = mixed_hessian(N)
    fig = go.Figure(data=go.Heatmap(x=x, y=x, z=d12u.T,
                                     colorbar=dict(title='u_xy')))
    fig.update_layout(title='Drugie pochodne rozwiazania -Delta u=f, dane naroznikowe',
                      xaxis_title='x', yaxis_title='y', yaxis_scaleanchor='x')
    save(fig, 'pde16_fft_hessjan.html')

    Ns = [32, 64, 128, 256, 512]
    peaks = []
    rms = []
    for n in Ns:
        _, fn, hn = mixed_hessian(n)
        peaks.append(float(np.max(np.abs(hn))))
        rms.append(float(np.mean(hn**2)**0.5))
        # Mnoznik Hessianu ma modul <= 1 w normie L2:
        assert rms[-1] <= float(np.mean(fn**2)**0.5) + 1e-11
    fig = go.Figure()
    fig.add_trace(go.Scatter(x=Ns, y=peaks, mode='lines+markers', name='maksimum |u_xy|'))
    fig.add_trace(go.Scatter(x=Ns, y=rms, mode='lines+markers', name='norma sredniokwadratowa'))
    fig.update_layout(title='Wzrost maksimum przy zachowanym oszacowaniu L2',
                      xaxis_title='liczba wezlow N', xaxis_type='log',
                      yaxis_title='wartosc normy')
    save(fig, 'pde16_fft_normy.html')
    print('N / maksimum / RMS')
    for n, peak, r in zip(Ns, peaks, rms):
        print(f'{n:4d} / {peak:.5f} / {r:.5f}')
    assert peaks[-1] > peaks[0], 'Oczekiwano wzrostu maksimum przy narozniku'


def demo_holder():
    x = np.linspace(-1, 1, 1501)
    alpha = 0.6
    u = -np.abs(x)**(alpha+2)/((alpha+1)*(alpha+2))
    f = np.abs(x)**alpha
    fig = go.Figure()
    fig.add_trace(go.Scatter(x=x, y=u, name='u: rozwiazanie -u\'\'=f'))
    fig.add_trace(go.Scatter(x=x, y=f, name='f=|x|^alpha = -u\'\''))
    fig.update_layout(title=f'Ostrosc oszacowan Schaudera, alpha={alpha}',
                      xaxis_title='x', yaxis_title='wartosc')
    save(fig, 'pde16_holder_ostrosc.html')
    h = np.logspace(-4, -1, 30)
    beta = 0.8
    qa = h**alpha / h**alpha
    qb = h**alpha / h**beta
    fig = go.Figure()
    fig.add_trace(go.Scatter(x=h, y=qa, name='iloraz Holdera dla beta=alpha'))
    fig.add_trace(go.Scatter(x=h, y=qb, name='iloraz Holdera dla beta>alpha'))
    fig.update_layout(title='Wyzszy wykladnik Holdera jest nieosiagalny',
                      xaxis_type='log', yaxis_type='log',
                      xaxis_title='odleglosc h', yaxis_title='iloraz')
    save(fig, 'pde16_holder_ilorazy.html')
    assert np.allclose(qa, 1)
    assert qb[0] > qb[-1]


if __name__ == '__main__':
    demo_fft()
    demo_holder()
\end{lstlisting}

\section{Regularność De Giorgiego--Nasha--Mosera: przykład radialny — Plotly}
\label{sec:pde17-code}
Program wizualizuje słabe rozwiązanie o optymalnej regularności
$C^{0,\gamma}$, macierz współczynników oraz ilorazy Höldera.
Dodatkowo sprawdza wartości maksimum i minimum przykładowej
funkcji harmonicznej. Jest to eksperyment ilustracyjny,
nie dowód twierdzenia o regularności przy mierzalnych współczynnikach.
\begin{lstlisting}[style=pythonpde,caption={Mierzalne współczynniki, regularność Höldera i Harnack (Plotly)},label={code:pde17-plotly}]
"""PDE rozdzial 17: ostrosc Holdera przy mierzalnych wspolczynnikach.

Zaleznosci: pip install numpy plotly
Uruchomienie: python demo_degiorgi_moser_plotly.py
Wyniki: pde17_profil_holder.html, pde17_macierz.html,
        pde17_ilorazy_holdera.html, pde17_harnack_affine.html
"""
from pathlib import Path
import numpy as np
import plotly.graph_objects as go

OUT = Path(__file__).resolve().parent


def store(fig: go.Figure, filename: str) -> None:
    fig.update_layout(template="plotly_white", font=dict(size=13))
    fig.write_html(OUT / filename, include_plotlyjs=True, full_html=True)


def radial_example(gamma: float, n: int = 301):
    xx = np.linspace(-1.0, 1.0, n)
    x, y = np.meshgrid(xx, xx, indexing="xy")
    radius = np.hypot(x, y)
    theta_cos = np.divide(x, radius, out=np.zeros_like(x), where=radius > 0)
    u = radius**gamma * theta_cos
    u[radius > 1] = np.nan
    # A = gamma^2 I + (1-gamma^2) e_r \otimes e_r
    e1 = np.divide(x, radius, out=np.zeros_like(x), where=radius > 0)
    e2 = np.divide(y, radius, out=np.zeros_like(y), where=radius > 0)
    a11 = gamma**2 + (1.0-gamma**2) * e1**2
    a22 = gamma**2 + (1.0-gamma**2) * e2**2
    a12 = (1.0-gamma**2) * e1 * e2
    return xx, x, y, radius, u, a11, a12, a22


def main():
    gamma = 0.55
    xx, x, y, radius, u, a11, a12, a22 = radial_example(gamma)
    fig = go.Figure(go.Heatmap(x=xx, y=xx, z=u, colorscale="RdBu",
                               zmid=0, colorbar=dict(title="u")))
    fig.update_layout(title=f"Rozwiazanie r^gamma cos(theta), gamma={gamma}",
                      xaxis_title="x", yaxis_title="y", yaxis_scaleanchor="x")
    store(fig, "pde17_profil_holder.html")

    fig = go.Figure()
    fig.add_trace(go.Heatmap(x=xx, y=xx, z=np.where(radius <= 1, a12, np.nan),
                             colorscale="RdBu", zmid=0, colorbar=dict(title="a12")))
    fig.update_layout(title="Mierzalna macierz A: wspolczynnik a12",
                      xaxis_title="x", yaxis_title="y", yaxis_scaleanchor="x")
    store(fig, "pde17_macierz.html")

    r = np.logspace(-4, -1, 90)
    fig = go.Figure()
    for beta in [0.35, gamma, 0.75]:
        quotient = r**(gamma-beta)
        fig.add_trace(go.Scatter(x=r, y=quotient, mode="lines",
                                 name=f"beta={beta}"))
    fig.update_layout(title="Ilorazy Holdera w zerze: r^(gamma-beta)",
                      xaxis_title="r", yaxis_title="iloraz", xaxis_type="log",
                      yaxis_type="log")
    store(fig, "pde17_ilorazy_holdera.html")
    assert np.all(np.diff(r**(gamma-0.75)) < 0)

    rr = np.linspace(0.0, 0.95, 100)
    # harmonic positive v(x,y)=2+0.4x on B_1
    sup = 2.0 + 0.4 * rr
    inf = 2.0 - 0.4 * rr
    ratio = sup/inf
    assert np.max(ratio) < 1.5
    fig = go.Figure()
    fig.add_trace(go.Scatter(x=rr, y=sup, name="sup na B_r"))
    fig.add_trace(go.Scatter(x=rr, y=inf, name="inf na B_r"))
    fig.add_trace(go.Scatter(x=rr, y=ratio, name="sup/inf"))
    fig.update_layout(title="Przykladowa dodatnia funkcja harmoniczna i Harnack",
                      xaxis_title="promien r", yaxis_title="wartosc")
    store(fig, "pde17_harnack_affine.html")
    print("Zweryfikowano cztery wykresy i prosty test Harnacka.")


if __name__ == "__main__":
    main()
\end{lstlisting}


\section{Widmo eliptyczne: wartości własne, Weyl i ślad cieplny — Plotly}
\label{sec:pde18-code}
Poniższy program oblicza dyskretne wartości własne jednowymiarowego
laplasjanu Dirichleta, porównuje funkcję zliczającą kwadratowego
laplasjanu różnicowego z wyrazem głównym prawa Weyla oraz
wizualizuje ślad półgrupy cieplnej. Ostatnie porównanie wykorzystuje
jawne wartości własne ciągłego laplasjanu na kwadracie.
Przy ustalonej siatce nie należy interpretować końca widma
macierzy jako asymptotyki wysokich częstotliwości operatora ciągłego.
\begin{lstlisting}[style=pythonpde,caption={Wartości własne, prawo Weyla i ślad cieplny (Plotly)},label={code:pde18-plotly}]
"""PDE, rozdzial 18: min-max, Weyl, slad cieplny.

Zaleznosci: pip install numpy scipy plotly
Uruchom: python demo_widmo_weyl_plotly.py
Wyniki: pde18_wartosci_wlasne.html, pde18_funkcja_zliczajaca.html,
        pde18_slad_cieplny.html
"""
from pathlib import Path
import numpy as np
import plotly.graph_objects as go
from scipy.linalg import eigh_tridiagonal

OUT = Path(__file__).resolve().parent


def save(fig, filename):
    fig.update_layout(template='plotly_white', font=dict(size=13),
                      margin=dict(l=60, r=35, t=65, b=65))
    fig.write_html(OUT / filename, include_plotlyjs=True, full_html=True)


def dirichlet_interval():
    ns = np.array([12, 20, 32, 50, 80, 128, 200])
    errors = []
    for n in ns:
        h = 1.0 / (n + 1)
        vals = eigh_tridiagonal(np.full(n, 2 / h**2),
                                np.full(n - 1, -1 / h**2),
                                select='i', select_range=(0, 0),
                                eigvals_only=True)
        errors.append(abs(vals[0] - np.pi**2))
    err = np.asarray(errors)
    h = 1 / (ns + 1)
    fig = go.Figure()
    fig.add_trace(go.Scatter(x=h, y=err, mode='lines+markers',
                             name='Blad pierwszej wartosci wlasnej'))
    fig.add_trace(go.Scatter(x=h, y=err[-1] * (h/h[-1])**2,
                             mode='lines', name='Linia odniesienia: h^2'))
    fig.update_layout(title='Zbieznosc wartosci wlasnej laplasjanu (Dirichlet)',
                      xaxis_type='log', yaxis_type='log',
                      xaxis_title='Krok h', yaxis_title='Blad bezwzgledny')
    save(fig, 'pde18_wartosci_wlasne.html')
    slope = np.polyfit(np.log(h), np.log(err), 1)[0]
    assert 1.95 < slope < 2.05, f'Oczekiwany rzad 2, otrzymano {slope}'
    print('Rzad zbieznosci wartosci wlasnej:', round(slope, 6))


def weyl_square():
    m = 100
    h = 1 / (m + 1)
    n = np.arange(1, m+1)
    one_d = 4 * np.sin(np.pi * n/(2*(m+1)))**2 / h**2
    two_d = np.sort((one_d[:, None] + one_d[None, :]).ravel())
    cutoffs = np.linspace(0, 2400, 450)
    counts = np.searchsorted(two_d, cutoffs, side='right')
    leading = cutoffs / (4*np.pi)  # |Omega|=1, d=2
    fig = go.Figure()
    fig.add_trace(go.Scatter(x=cutoffs, y=counts, mode='lines',
                             line_shape='hv', name='Dyskretne zliczanie modow'))
    fig.add_trace(go.Scatter(x=cutoffs, y=leading, mode='lines',
                             name='Wyraz glowny Weyla: Lambda/(4 pi)'))
    fig.update_layout(title='Funkcja zliczajaca na kwadracie jednostkowym',
                      xaxis_title='Prog Lambda', yaxis_title='Liczba modow')
    save(fig, 'pde18_funkcja_zliczajaca.html')
    # m=100 makes low modes accurate; at high frequencies the grid
    # is fixed and hence Weyl is only illustrative.
    print('N_dyskretne(1000)=', np.searchsorted(two_d, 1000, side='right'))


def heat_trace():
    j = np.arange(1, 180)
    eigs = np.pi**2 * (j[:, None]**2 + j[None, :]**2)
    t = np.logspace(-3.3, -1.2, 170)
    exact = np.array([np.exp(-tau*eigs).sum() for tau in t])
    leading = 1 / (4*np.pi*t)  # |Omega|=1, d=2
    fig = go.Figure()
    fig.add_trace(go.Scatter(x=t, y=exact, mode='lines',
                             name='Slad cieplny (suma modow)'))
    fig.add_trace(go.Scatter(x=t, y=leading, mode='lines',
                             name='Asymptotyka 1/(4 pi t)'))
    fig.update_layout(title='Slad cieplny laplasjanu Dirichleta na kwadracie',
                      xaxis_type='log', yaxis_type='log',
                      xaxis_title='Czas t', yaxis_title='Tr exp(t Delta_D)')
    save(fig, 'pde18_slad_cieplny.html')
    ratio = exact[0] / leading[0]
    assert 0.8 < ratio < 1.0
    print('Stosunek sladu do wyrazu glownego dla malego t:', round(ratio, 6))


if __name__ == '__main__':
    dirichlet_interval()
    weyl_square()
    heat_trace()
\end{lstlisting}



\section{Pojemność kondensatora kulistego — Plotly}
\label{sec:pde19-code}
Poniższy program porównuje pojemności względne
$\operatorname{cap}_{2,B_1}(\overline B_r)$ w wymiarach $2$ i $3$.
Każda krzywa jest obliczana z dokładnego wzoru
\eqref{eq:pde19-ball-exact}; program nie zastępuje dowodu
minimalności energii ani nie numeryzuje równania Laplace'a.
Wykresy zapisywane są jako interaktywne pliki HTML.
\begin{lstlisting}[style=pythonpde,caption={Pojemność Sobolewa i tempo zaniku pojemności punktu (Plotly)},label={code:pde19-plotly}]
"""PDE, rozdzial 19: pojemnosc kondensatora kulistego.
Wymaga: numpy, plotly. Wyniki: interaktywne pliki HTML.
"""
from pathlib import Path
import numpy as np
import plotly.graph_objects as go

OUT = Path(__file__).resolve().parent
R = 1.0
r = np.geomspace(1e-6, 0.9, 240)
cap2 = 2 * np.pi / np.log(R / r)
cap3 = 4 * np.pi / (1 / r - 1 / R)


def save(fig, name):
    fig.update_layout(template="plotly_white", font=dict(size=13))
    fig.write_html(OUT / name, include_plotlyjs=True, full_html=True)


def condenser_plot(cap, d, name):
    fig = go.Figure()
    fig.add_trace(go.Scatter(x=r, y=cap, mode="lines",
                             name=f"pojemnosc d={d}"))
    fig.update_layout(title=f"Pojemnosc wzgledna kuli w B_1, d={d}",
                      xaxis_title="promien r", yaxis_title="energia minimalna")
    save(fig, name)


def comparison():
    fig = go.Figure()
    fig.add_trace(go.Scatter(x=r, y=cap2, name="d=2 (logarytm)"))
    fig.add_trace(go.Scatter(x=r, y=cap3, name="d=3 (potega)"))
    fig.update_layout(title="Zanik pojemnosci punktu, rozne tempa",
                      xaxis_type="log", yaxis_type="log",
                      xaxis_title="promien r", yaxis_title="pojemnosc")
    save(fig, "pde19_pojemnosc_porownanie.html")


def test_exact_values():
    assert np.all(np.diff(cap2) > 0)
    assert np.all(np.diff(cap3) > 0)
    assert np.isclose(2 * np.pi / np.log(2), 9.064720283654388)
    assert np.isclose(4 * np.pi / (2 - 1), 4 * np.pi)
    print("d=2, r=1/2:", 2 * np.pi / np.log(2))
    print("d=3, r=1/2:", 4 * np.pi)
    print("d=2, r=1e-6:", cap2[0])
    print("d=3, r=1e-6:", cap3[0])


if __name__ == "__main__":
    test_exact_values()
    condenser_plot(cap2, 2, "pde19_pojemnosc_2d.html")
    condenser_plot(cap3, 3, "pde19_pojemnosc_3d.html")
    comparison()
\end{lstlisting}


\section{Paraboliczne równania ewolucyjne: wsteczny Euler i stabilność — Plotly}
\label{sec:pde20-code}
Eksperyment porównuje niejawny schemat różnicowy z rozwiązaniem dokładnym
na odcinku, testuje rząd zbieżności przy $\tau\asymp h^2$ oraz
ilustruje tłumienie poszczególnych modów. W algorytmie stosujemy
dyskretną transformatę sinusową, która diagonalizuje laplasjan
z jednorodnymi warunkami Dirichleta. Jest to implementacja
wstecznego Eulera, a nie schematu dokładnej ewolucji macierzy.
\begin{lstlisting}[style=pythonpde,caption={Niejawny schemat cieplny, zbieżność i tłumienie modów (Plotly)},label={code:pde20-plotly}]
"""Iteracja 20, rownanie ciepa na (0,1) z Dirichletem.

Wsteczny Euler i centralny laplasjan roznicowy,
diagonalizacja macierzy przez dyskretna transformate sinusowa.
Instalacja: python -m pip install numpy scipy plotly
Uruchomienie: python demo_paraboliczne_plotly.py
"""
from pathlib import Path
import numpy as np
import plotly.graph_objects as go
from scipy.fft import dst, idst

OUT = Path(__file__).resolve().parent
T = 0.035


def initial(x):
    return np.sin(np.pi * x) + 0.5 * np.sin(3 * np.pi * x)


def exact(x, t):
    return np.exp(-np.pi**2 * t) * np.sin(np.pi*x) + 0.5 * np.exp(-9*np.pi**2*t) * np.sin(3*np.pi*x)


def backward_euler(n, at=T, cfl=0.32):
    """Rozwiazanie po nstep krokach: kazdy krok jest niejawny."""
    h = 1 / (n + 1)
    nstep = max(1, int(np.ceil(at / (cfl*h*h))))
    tau = at / nstep
    x = np.arange(1, n+1) * h
    j = np.arange(1, n+1)
    eig = 4 / h**2 * np.sin(j*np.pi/(2*(n+1)))**2
    coeffs = dst(initial(x), type=1, norm='ortho')
    ampl = (1 + tau*eig)**(-nstep)
    uh = idst(coeffs*ampl, type=1, norm='ortho')
    return x, uh, tau, eig


def save(fig, file):
    fig.update_layout(template='plotly_white', font=dict(size=13),
                      margin=dict(l=70, r=35, b=60, t=75))
    fig.write_html(OUT/file, include_plotlyjs=True, full_html=True)


def profiles():
    n=80
    h = 1/(n+1)
    x = np.arange(1,n+1)*h
    fig = go.Figure()
    for t in [0.0, 0.004, 0.012, T]:
        if t == 0:
            u=initial(x)
        else:
            _, u, _, _ = backward_euler(n, at=t)
        fig.add_trace(go.Scatter(x=x, y=u, mode='lines', name=f't={t:g}'))
    fig.update_layout(title='Rownanie ciepa: wygadzanie i zanikanie modow',
                      xaxis_title='x',yaxis_title='u(t,x)')
    save(fig, 'pde20_profile_plotly.html')
    # Zanik nie zwieksza normy w l2 dla danych bez wymuszenia.
    assert np.linalg.norm(backward_euler(n)[1]) <= np.linalg.norm(initial(x))


def convergence():
    sizes=np.array([12,18,27,40,60,90,135])
    errs=[]
    for n in sizes:
        x, u, _, _ = backward_euler(int(n))
        h = 1/(n+1)
        errs.append(np.sqrt(h*np.sum((u-exact(x,T))**2)))
    hs=1/(sizes+1)
    errs=np.array(errs)
    slope=np.polyfit(np.log(hs),np.log(errs),1)[0]
    print('Wsteczny Euler / laplasjan roznicowy: rzad bedu',round(slope,5))
    assert 1.85 < slope < 2.12, f'Nieoczekiwany rzad: {slope}'
    fig=go.Figure()
    fig.add_trace(go.Scatter(x=hs,y=errs,mode='lines+markers',name='Bad w L2_h'))
    fig.add_trace(go.Scatter(x=hs,y=errs[-1]*(hs/hs[-1])**2,name='Odniesienie: h2'))
    fig.update_layout(title=f'Zbieznosc dla 0.32h2, rzad {slope:.3f}',
                      xaxis_type='log',yaxis_type='log',
                      xaxis_title='h',yaxis_title='||U-u||_(L2_h)')
    save(fig, 'pde20_zbieznosc_plotly.html')


def damping():
    z=np.logspace(-3,2,500)
    fig=go.Figure()
    fig.add_trace(go.Scatter(x=z,y=np.exp(-z),name='Dokadny: exp(-z)'))
    fig.add_trace(go.Scatter(x=z,y=1/(1+z),name='Wsteczny Euler: 1/(1+z)'))
    fig.add_trace(go.Scatter(x=z,y=np.abs(1-z),name='Jawny Euler: |1-z|'))
    fig.update_layout(title='Tumienie jednego modu; z=',
                      xaxis_title='z',yaxis_title='wspoczynnik amplitudy',
                      xaxis_type='log',yaxis_type='log')
    save(fig,'pde20_tlumienie_plotly.html')
    assert np.all(1/(1+z)<=1)
    assert np.any(np.abs(1-z)>1)


if __name__ == '__main__':
    profiles()
    convergence()
    damping()
\end{lstlisting}


\section{Układ symetryczno-hiperboliczny: schemat pod prąd — Plotly}
\label{sec:pde21-code}
Program oblicza ewolucję dwóch modów o różnych prędkościach,
porównuje rozwiązanie numeryczne z dokładnym,
sprawdza nierosnącą energię dyskretną i wyznacza rząd
zbieżności jawnego schematu pod prąd. Rozdzielenie macierzy
$A=A_++A_-$ odbywa się przez diagonalizację ortogonalną,
a warunek CFL ma postać $\tau\max_k|a_k|/h\le1$.
\begin{lstlisting}[style=pythonpde,caption={Hiperboliczny układ, energia i zbieżność: Plotly},label={code:pde21-plotly}]
"""Rozdzial 21: dwa przemieszczajace sie mody ukadu symetrycznego.
Zaleznosci: numpy, plotly; uruchomienie: python demo_hiperboliczne_plotly.py
Schemat pod prad z rozdzieleniem macierzy A=A_++A_-.
"""
from pathlib import Path
import numpy as np
import plotly.graph_objects as go

OUT = Path(__file__).resolve().parent
THETA = 0.48
Q = np.array([[np.cos(THETA), -np.sin(THETA)],
              [np.sin(THETA), np.cos(THETA)]])
SPEEDS = np.array([-1.0, 0.6])
A = Q @ np.diag(SPEEDS) @ Q.T
A_PLUS = Q @ np.diag(np.maximum(SPEEDS, 0.0)) @ Q.T
A_MINUS = Q @ np.diag(np.minimum(SPEEDS, 0.0)) @ Q.T


def wrapped_gauss(x, center, width):
    y = (x - center + 0.5) % 1.0 - 0.5
    return np.exp(-(y / width) ** 2)


def initial(x):
    w = np.vstack([wrapped_gauss(x, 0.25, 0.065),
                   -0.8 * wrapped_gauss(x, 0.71, 0.075)])
    return Q @ w


def exact(x, t):
    w = np.vstack([wrapped_gauss(x - SPEEDS[0] * t, 0.25, 0.065),
                   -0.8 * wrapped_gauss(x - SPEEDS[1] * t, 0.71, 0.075)])
    return Q @ w


def upwind(n, T=0.35, cfl=0.8, record=False):
    h = 1.0 / n
    x = np.arange(n) * h
    steps = int(np.ceil(T * np.max(abs(SPEEDS)) / (cfl * h)))
    dt = T / steps
    U = initial(x)
    en = [h * np.sum(U * U)]
    for _ in range(steps):
        left = U - np.roll(U, 1, axis=1)
        right = np.roll(U, -1, axis=1) - U
        U = U - dt / h * (A_PLUS @ left + A_MINUS @ right)
        if record:
            en.append(h * np.sum(U * U))
    ref = exact(x, T)
    err = np.sqrt(h * np.sum((U - ref) ** 2))
    return x, U, ref, err, np.array(en), dt, h


def export(fig, name):
    fig.update_layout(template='plotly_white', font=dict(size=13),
                      margin=dict(l=75, r=35, b=55, t=75))
    fig.write_html(OUT / name, include_plotlyjs=True, full_html=True)


def main():
    n = 240
    x, U, target, err, energy, dt, h = upwind(n, record=True)
    fig = go.Figure()
    for k in range(2):
        fig.add_trace(go.Scatter(x=x, y=U[k], mode='lines',
                     name=f'numeryczne u{k+1}'))
        fig.add_trace(go.Scatter(x=x, y=target[k], mode='lines',
                     line=dict(dash='dash'), name=f'dokladne u{k+1}'))
    fig.update_layout(title=f'Ukad symetryczny: mody o predkosciach -1 i 0.6; blad={err:.3g}',
                      xaxis_title='x', yaxis_title='u(t,x)')
    export(fig, 'pde21_profile_plotly.html')

    fig = go.Figure()
    fig.add_trace(go.Scatter(x=np.arange(energy.size) * dt, y=energy,
                             mode='lines', name='Dyskretna energia l2'))
    fig.update_layout(title='Kontrakcja energii: schemat pod prad, CFL=0.8',
                      xaxis_title='czas', yaxis_title='h sum |U_j|^2')
    export(fig, 'pde21_energia_plotly.html')
    assert np.max(np.diff(energy)) <= 1e-12, 'Energia powinna byc nierosnaca'
    assert np.allclose(A, A.T)

    sizes = np.array([40, 60, 90, 135, 200, 300, 450])
    errors = np.array([upwind(int(k))[3] for k in sizes])
    hs = 1.0 / sizes
    slope = np.polyfit(np.log(hs[-4:]), np.log(errors[-4:]), 1)[0]
    print('Rzad zbieznosci, cztery ostatnie siatki:', round(slope, 4))
    print('Blad l2 dla N=240:', round(err, 7))
    assert 0.65 < slope < 1.15, f'Nieoczekiwany rzad {slope}'
    fig = go.Figure()
    fig.add_trace(go.Scatter(x=hs, y=errors, mode='lines+markers', name='Blad l2'))
    fig.add_trace(go.Scatter(x=hs, y=errors[-1] * hs / hs[-1], mode='lines',
                             line=dict(dash='dash'), name='Skala O(h)'))
    fig.update_layout(title=f'Zbieznosc schematu pod prad: rzad {slope:.3f}',
                      xaxis_title='h', yaxis_title='blad l2',
                      xaxis_type='log', yaxis_type='log')
    export(fig, 'pde21_zbieznosc_plotly.html')


if __name__ == '__main__':
    main()
\end{lstlisting}

\section{Swobodna ewolucja Schrödingera na torusie — Plotly}
\label{sec:pde22-code}
Kod stosuje reprezentację spektralną na torusie, a nie swobodne
jądro propagatora na $\R^d$. Sprawdza zachowanie normy
$L^2$, ruch pakietu i wzór Parsevala na błąd obcięcia.

\begin{lstlisting}[style=pythonpde,caption={Ewolucja Schrodingera: pakiet, faza i błąd widmowy (Plotly)},label={code:pde22-plotly}]
"""Rozdzial 22: swobodny propagator Schrodingera na torusie.

Zaleznosci: numpy, plotly
Uruchomienie: python demo_schrodinger_plotly.py
Wyniki: pde22_pakiet_plotly.html, pde22_faza_plotly.html,
        pde22_obciecie_plotly.html
"""
from pathlib import Path
import numpy as np
import plotly.graph_objects as go

OUT = Path(__file__).resolve().parent
N = 512
x = -np.pi + (2 * np.pi / N) * np.arange(N)
k = np.arange(-100, 101)
k0, sigma, x0 = 12.0, 5.0, -1.0
coef = np.exp(-0.5 * ((k - k0) / sigma)**2) * np.exp(-1j * k * x0)

def psi(t, cutoff=None):
    mask = np.ones(k.shape, dtype=bool) if cutoff is None else np.abs(k) <= cutoff
    z = (coef[mask] * np.exp(-1j * t * k[mask]**2))[:, None]
    return np.sum(z * np.exp(1j * k[mask, None] * x[None, :]), axis=0) / np.sqrt(2 * np.pi)

def store(fig, name):
    fig.update_layout(template="plotly_white", font=dict(size=13),
                      margin=dict(l=70, r=30, t=70, b=60))
    fig.write_html(OUT / name, include_plotlyjs=True, full_html=True)

def main():
    times = [0.0, 0.015, 0.03, 0.06, 0.09]
    fig = go.Figure()
    for t in times:
        p = psi(t)
        fig.add_trace(go.Scatter(x=x, y=np.abs(p)**2, mode='lines',
                                 name=f't={t:.3f}'))
    fig.update_layout(title='Pakiet swobodny na torusie: propagacja i dyspersja',
                      xaxis_title='x', yaxis_title='|u(t,x)|^2')
    store(fig, 'pde22_pakiet_plotly.html')

    # Przestrzenna siatka nie zawiera prawego konca okresu:
    # okresowa kwadratura trapezowa jest dokladna dla pasmowych iloczynow.
    norms_periodic = [(2*np.pi/N)*np.sum(np.abs(psi(t))**2) for t in times]
    parseval = np.sum(np.abs(coef)**2)
    print('Maksymalna zmiana normy (kwadrat):',
          max(abs(v - parseval) for v in norms_periodic))
    assert max(abs(v-parseval) for v in norms_periodic) < 1e-10

    t = 0.04
    fig = go.Figure()
    fig.add_trace(go.Scatter(x=k, y=np.abs(coef), name='|c_k|', mode='lines'))
    fig.add_trace(go.Scatter(x=k, y=np.angle(np.exp(-1j*t*k*k)),
                             name='faza propagatora', yaxis='y2', mode='markers',
                             marker=dict(size=3)))
    fig.update_layout(title='Zachowany modul i zmieniajaca sie faza',
                      xaxis_title='k', yaxis_title='|c_k|',
                      yaxis2=dict(title='faza [rad]', overlaying='y', side='right'))
    store(fig, 'pde22_faza_plotly.html')

    full = psi(0.05)
    cutoffs = [6, 10, 14, 18, 24, 32, 48, 64]
    errors = [(2*np.pi/N*np.sum(np.abs(full-psi(0.05,K))**2))**0.5
              for K in cutoffs]
    exact_tails = [np.sqrt(np.sum(np.abs(coef[np.abs(k)>K])**2))
                   for K in cutoffs]
    assert np.allclose(errors, exact_tails, rtol=1e-7, atol=1e-12)
    fig = go.Figure()
    fig.add_trace(go.Scatter(x=cutoffs, y=errors, name='blad numeryczny',
                             mode='lines+markers'))
    fig.add_trace(go.Scatter(x=cutoffs, y=exact_tails, name='ogon Parsevala',
                             mode='lines', line=dict(dash='dash')))
    fig.update_layout(title='Obciecie szeregu i tozsamosc Parsevala',
                      xaxis_title='maksymalne |k|', yaxis_title='blad L2',
                      yaxis_type='log')
    store(fig, 'pde22_obciecie_plotly.html')

if __name__ == '__main__':
    main()
\end{lstlisting}



\section{Równanie Hamiltona--Jacobiego: charakterystyki i rozwiązania lepkościowe (Plotly)}
Kod oblicza wartości wzoru Hopfa--Laxa przez minimalizację na siatce,
rysuje charakterystyki i pochodną rozwiązania po powstaniu osobliwości.
Wykresy ilustrują teorię; nie zastępują dowodów z Rozdziału~\ref{chap:pde23}.
\begin{lstlisting}[style=pythonpde,caption={Wzór Hopfa--Laxa i powstawanie narożnika (Plotly)},label={code:pde23-plotly}]
from pathlib import Path
import numpy as np
import plotly.graph_objects as go
OUT = Path(__file__).resolve().parent
def hopf_lax(t: float, x: np.ndarray, y: np.ndarray):
    """Minimalizacja dla g(y)=cos(y) i H(p)=p^2/2."""
    cost = np.cos(y)[None, :] + (x[:, None] - y[None, :])**2 / (2 * t)
    argmin = np.argmin(cost, axis=1)
    return cost[np.arange(len(x)), argmin], y[argmin]
def save(fig, filename):
    fig.update_layout(template="plotly_white", font=dict(size=13))
    fig.write_html(OUT / filename, include_plotlyjs=True, full_html=True)
def main():
    x = np.linspace(-np.pi, np.pi, 601)
    y = np.linspace(-np.pi - 2.2, np.pi + 2.2, 1801)
    fig = go.Figure()
    for t in (0.2, 0.8, 1.02, 1.35):
        u, minimizer = hopf_lax(t, x, y)
        fig.add_scatter(x=x, y=u, mode="lines", name=f"t={t:.2f}")
    fig.add_scatter(x=x, y=np.cos(x), mode="lines", name="t=0", line=dict(dash="dash"))
    fig.update_layout(title="Rozwiazanie lepkosciowe: g(x)=cos x, H(p)=p^2/2",
                      xaxis_title="x", yaxis_title="u(t,x)")
    save(fig, "pde23_profile_plotly.html")
    yy = np.linspace(-np.pi, np.pi, 35)
    tt = np.linspace(0, 1.6, 125)
    fig = go.Figure()
    for a in yy:
        fig.add_scatter(x=a - tt * np.sin(a), y=tt, mode="lines",
                        showlegend=False, line=dict(width=1.2))
    fig.add_hline(y=1, line_dash="dash", annotation_text="pierwsza degeneracja: t=1")
    fig.update_layout(title="Charakterystyki: x(t,y)=y-t sin y",
                      xaxis_title="poozenie x", yaxis_title="czas t")
    save(fig, "pde23_charakterystyki_plotly.html")
    t = 1.35
    dx = x[1] - x[0]
    u, minimizer = hopf_lax(t, x, y)
    grad = np.gradient(u, dx)
    fig = go.Figure()
    fig.add_scatter(x=x, y=grad, mode="lines", name="gradient numeryczny")
    fig.add_scatter(x=x, y=(x - minimizer) / t, mode="lines",
                    name="(x-y_min)/t", line=dict(dash="dash"))
    fig.update_layout(title="Skok gradientu rozwiazania po t=1",
                      xaxis_title="x", yaxis_title="pochodna przestrzenna")
    save(fig, "pde23_skok_gradientu_plotly.html")
    k = np.argmin(np.abs(x))
    xx = x[k]
    costs = np.cos(y) + (xx - y)**2 / (2*t)
    near = np.where(costs <= costs.min() + 1e-4)[0]
    assert np.any(y[near] < -0.4) and np.any(y[near] > 0.4)
    assert np.all(1 - 0.8*np.cos(yy) > 0)
    print("Trzy wizualizacje zapisane. Test symetrycznych minimow: OK.")
if __name__ == "__main__":
    main()
\end{lstlisting}


\section{Skalarne prawa zachowania: rozwiązania Riemanna — Plotly}
\label{sec:pde24-code}
Program implementuje monotoniczny schemat Godunowa dla strumienia
$f(u)=u^2/2$, porównuje go ze ścisłymi rozwiązaniami
entropijnymi zagadnienia Riemanna i wyznacza błąd w normie
$L^1$ na skończonym przedziale obliczeniowym.
Dane Riemanna nie są całkowalne na całej prostej; eksperyment
wykonujemy w skończonym oknie przestrzennym.
\begin{lstlisting}[style=pythonpde,caption={Fala uderzeniowa, fala rozrzedzenia i metoda Godunowa — Plotly},label={code:pde24-plotly}]
"""Skalarne prawa zachowania: fala uderzeniowa, rozrzedzenie, zbieznosc.

Uruchomienie: python demo_prawa_zachowania_plotly.py
Wymagania: numpy, plotly
"""
from pathlib import Path
import numpy as np
import plotly.graph_objects as go

OUT = Path(__file__).resolve().parent


def flux(u):
    return 0.5 * np.asarray(u) ** 2


def godunov(left, right):
    """Strumien Godunowa dla wypukego f(u)=u^2/2."""
    left = np.asarray(left)
    right = np.asarray(right)
    rarefaction = np.where(left >= 0, flux(left),
                          np.where(right <= 0, flux(right), 0.0))
    shock = np.maximum(flux(left), flux(right))
    return np.where(left <= right, rarefaction, shock)


def riemann_exact(x, t, left, right):
    if left > right:
        speed = (left + right) / 2
        return np.where(x < speed*t, left, right)
    if t == 0:
        return np.where(x < 0, left, right)
    return np.where(x <= left*t, left,
                    np.where(x >= right*t, right, x/t))


def numerical(n, tfinal, left, right):
    dx = 6.0/n
    x = -3.0 + (np.arange(n) + 0.5) * dx
    u = np.where(x < 0, left, right).astype(float)
    t = 0.0
    while t < tfinal - 1e-13:
        max_speed = max(1.0, np.max(np.abs(u)))
        dt = min(0.43*dx/max_speed, tfinal-t)
        padded = np.concatenate(([left], u, [right]))
        numerical_flux = godunov(padded[:-1], padded[1:])
        u -= dt/dx * (numerical_flux[1:] - numerical_flux[:-1])
        t += dt
    return x, u, dx


def graph(filename, title, traces, xaxis='x', yaxis='u'):
    fig = go.Figure()
    for x, y, name, style in traces:
        fig.add_trace(go.Scatter(x=x, y=y, name=name, mode='lines',
                                 line=dict(dash=style)))
    fig.update_layout(template='plotly_white', title=title,
                      xaxis_title=xaxis, yaxis_title=yaxis)
    if filename.endswith('zbieznosc.html'):
        fig.update_xaxes(type='log')
        fig.update_yaxes(type='log')
    fig.write_html(OUT / filename, include_plotlyjs=True, full_html=True)


def main():
    t = 0.8
    for left, right, name in [(1.0, 0.0, 'uderzenie'),
                               (0.0, 1.0, 'rozrzedzenie')]:
        x, u, dx = numerical(600, t, left, right)
        exact = riemann_exact(x, t, left, right)
        graph(f'pde24_{name}.html', f'Burgers: {name} (t={t})',
              [(x, exact, 'Rozwiazanie entropijne', 'solid'),
               (x, u, 'Schemat Godunowa', 'dash')])
        print(name, 'L1 error:', np.sum(abs(u-exact))*dx)

    ns = np.array([100, 200, 400, 800, 1600])
    trace_data = []
    for left, right, name, dash in [(1.0, 0.0, 'Uderzenie', 'solid'),
                                   (0.0, 1.0, 'Rozrzedzenie', 'dash')]:
        errors = []
        for n in ns:
            x, u, dx = numerical(int(n), t, left, right)
            errors.append(np.sum(abs(u-riemann_exact(x,t,left,right)))*dx)
        errors = np.asarray(errors)
        rate = np.polyfit(np.log(6.0/ns), np.log(errors), 1)[0]
        trace_data.append((6.0/ns, errors, f'{name}, nachylenie {rate:.2f}', dash))
        print(f'{name}: measured rate {rate:.4f}')
    graph('pde24_zbieznosc.html', 'Schemat Godunowa: bad L1',
          trace_data, xaxis='Krok siatki h', yaxis='Bad L1')

    # Kontrola zachowania ograniczenia maksimum przez monotoniczny schemat.
    x, u, _ = numerical(800, t, 1.0, 0.0)
    assert np.all(u >= -1e-12) and np.all(u <= 1+1e-12)


if __name__ == '__main__':
    main()
\end{lstlisting}


\section{Półliniowe równania ewolucyjne: energia i wybuch — Plotly}
\label{sec:pde25-code}
Program ilustruje spadek energii w równaniu Allena--Cahna
oraz dolną barierę na średnią rozwiązania równania
$u_t=\Delta u+u^2$ na torusie. Schemat Stranga jest
eksperymentem numerycznym, a nie dowodem twierdzeń.
\begin{lstlisting}[style=pythonpde,caption={Równania półliniowe: energia i eksplozja (Plotly)},label={code:pde25-plotly}]
"""PDE, rozdzial 25: reakcja, energia i eksplozja na torusie.

Uruchomienie: python demo_pololiniowe_plotly.py
Zaleznosci: numpy, plotly
Zapisuje trzy samodzielne wykresy interaktywne do katalogu skryptu.
"""
from pathlib import Path
import numpy as np
import plotly.graph_objects as go

OUT = Path(__file__).resolve().parent
N = 512
x = 2 * np.pi * np.arange(N) / N
k = np.fft.fftfreq(N, d=1 / N)
dx = 2 * np.pi / N

def diffuse(u, dt):
    return np.fft.ifft(np.fft.fft(u) * np.exp(-(k * k) * dt)).real

def save(fig, name):
    fig.update_layout(template="plotly_white", font=dict(size=13))
    fig.write_html(OUT / name, include_plotlyjs=True, full_html=True)

def allen_cahn():
    # Splitting Stranga: dyfuzja, dokladny lokalny przeplyw u'=u-u^3,
    # dyfuzja. Dokladny przeplyw reakcji zachowuje znak.
    u = 0.17 + 0.72 * np.cos(x) + 0.14 * np.cos(3 * x)
    dt = 0.001
    snapshots = {0.0: u.copy()}
    ts, energies = [], []
    def energy(v):
        vx = np.fft.ifft(1j * k * np.fft.fft(v)).real
        return dx * np.sum(0.5 * vx**2 + 0.25 * (v*v - 1)**2)
    for j in range(1601):
        if j % 20 == 0:
            ts.append(j * dt)
            energies.append(energy(u))
        if j in (400, 800, 1600):
            snapshots[round(j * dt, 2)] = u.copy()
        if j == 1600:
            break
        u = diffuse(u, dt / 2)
        e2 = np.exp(2 * dt)
        u = e2**0.5 * u / np.sqrt(1 + u*u * (e2 - 1))
        u = diffuse(u, dt / 2)
    # Wybieramy dostatecznie dokladny krok, by energia byla nierosnaca.
    assert np.max(np.diff(energies)) < 1e-6
    fig = go.Figure()
    for t, v in snapshots.items():
        fig.add_trace(go.Scatter(x=x, y=v, name=f"t={t:g}"))
    fig.update_layout(title="Rownanie Allena-Cahna: wygadzanie i reakcja", xaxis_title="x", yaxis_title="u(t,x)")
    save(fig, "pde25_allen_cahn_profile.html")
    fig2 = go.Figure(go.Scatter(x=ts, y=energies, name="E(u(t))"))
    fig2.update_layout(title="Rownanie Allena-Cahna: dyssypacja energii", xaxis_title="czas t", yaxis_title="energia")
    save(fig2, "pde25_allen_cahn_energia.html")
    print("Rownanie Allena-Cahna: energia poczatkowa i koncowa:", energies[0], energies[-1])

def blowup_mean():
    # u_t=Delta u+u^2, u0=0.35+0.15 cos(x), dodatnia srednia 0.35.
    u = 0.35 + 0.15 * np.cos(x)
    m0 = np.mean(u)
    bound = 1.0 / m0
    dt = 0.0005
    tmax = 2.2
    ts, means, mins, maxs = [], [], [], []
    for j in range(int(round(tmax/dt)) + 1):
        if j % 40 == 0:
            ts.append(j * dt)
            means.append(float(np.mean(u)))
            mins.append(float(np.min(u)))
            maxs.append(float(np.max(u)))
        if j == int(round(tmax/dt)):
            break
        u = diffuse(u, dt / 2)
        assert dt * np.max(u) < 1
        u = u / (1 - dt * u)  # scisly przeplyw reakcji u'=u^2
        u = diffuse(u, dt / 2)
    ts = np.asarray(ts)
    barrier = m0 / (1 - m0 * ts)
    # Blad splittigu i kwadratury dopuszczamy ponizej 1e-3.
    assert np.min(np.asarray(means) - barrier) > -1e-3
    fig = go.Figure()
    fig.add_trace(go.Scatter(x=ts, y=means, name="srednia numeryczna"))
    fig.add_trace(go.Scatter(x=ts, y=barrier, name="bariera Jensena / ODE", line=dict(dash="dash")))
    fig.add_trace(go.Scatter(x=ts, y=maxs, name="maksimum"))
    fig.update_layout(title=f"u_t=Delta u + u^2: srednia i dolna bariera, T*  {bound:.3f}", xaxis_title="czas t", yaxis_title="wartosc")
    save(fig, "pde25_eksplozja_srednia.html")
    print("Bariera T* <=", bound, "roznica sredniej i bariery w koncu", means[-1] - barrier[-1])

if __name__ == "__main__":
    allen_cahn()
    blowup_mean()

\end{lstlisting}



\section{Nieliniowy strumień p-laplasjanu: przykłady — Plotly}
\label{sec:pde26-code}
Program rozwiązujący jednowymiarowe zagadnienie $-\Delta_pu=1$
z zerowymi warunkami Dirichleta konstruuje dyskretny strumień
z równania bilansu i odtwarza węzłowe wartości rozwiązania.
Osobno mierzymy błąd i wzrost przybliżonej drugiej pochodnej
w maksimum. Przy $p>2$ rozwiązanie jest klasy $C^1$,
lecz nie musi należeć do $C^2$. Obliczenia stanowią
ilustrację, a nie dowód twierdzenia Minty'ego--Browdera.
\begin{lstlisting}[style=pythonpde,caption={Nieliniowy strumień, błąd aproksymacji i krzywizna (Plotly)},label={code:pde26-plotly}]
"""PDE 26: p-Laplacian Dirichlet, finite-volume flux and Plotly.
Requires: numpy, plotly. Run: python demo_plaplasjan_plotly.py
"""
from pathlib import Path
import numpy as np
import plotly.graph_objects as go

OUT = Path(__file__).resolve().parent
PS = (1.5, 2.0, 3.0, 5.0)


def exact(x, p):
    exponent = p / (p - 1)
    return ((p - 1) / p) * (0.5**exponent - np.abs(x - 0.5)**exponent)


def discrete(p, n):
    """Solve -D_h(|D_h u|^(p-2) D_h u)=1 with zero end values.
    In 1D the numerical flux is q_{i+1/2}=1/2-(i+1/2)h.
    """
    x = np.linspace(0, 1, n + 1)
    h = 1.0 / n
    xm = (x[1:] + x[:-1]) / 2
    flux = 0.5 - xm
    grad = np.sign(flux) * np.abs(flux)**(1 / (p - 1))
    values = np.r_[0.0, np.cumsum(h * grad)]
    assert abs(values[-1]) < 5e-13
    q_recovered = np.sign(np.diff(values) / h) * np.abs(np.diff(values) / h)**(p - 1)
    residual = -(np.diff(q_recovered) / h) - 1
    assert np.max(np.abs(residual)) < 1e-8
    return x, values


def save(fig, filename):
    fig.update_layout(template='plotly_white', font=dict(size=13))
    fig.write_html(OUT / filename, include_plotlyjs=True, full_html=True)


def make_profiles():
    x_dense = np.linspace(0, 1, 1001)
    fig = go.Figure()
    for p in PS:
        xd, ud = discrete(p, 16)
        fig.add_trace(go.Scatter(x=x_dense, y=exact(x_dense, p),
                                 mode='lines', name=f'p={p:g}, exact'))
        fig.add_trace(go.Scatter(x=xd, y=ud, mode='markers',
                                 name=f'p={p:g}, grid n=16'))
    fig.update_layout(title='One-dimensional p-Laplacian: -Delta_p u = 1',
                      xaxis_title='x', yaxis_title='u(x)')
    save(fig, 'pde26_profile.html')


def make_convergence():
    ns = np.array((20, 40, 80, 160, 320, 640))
    fig = go.Figure()
    for p in (1.5, 3.0, 5.0):
        errors = []
        for n in ns:
            x, u = discrete(p, int(n))
            errors.append(np.max(np.abs(u - exact(x, p))))
        errors = np.asarray(errors)
        rate = -np.polyfit(np.log(ns), np.log(errors), 1)[0]
        print('p=', p, 'convergence rate=', float(rate))
        assert errors[-1] < errors[0]
        fig.add_trace(go.Scatter(x=1/ns, y=errors, mode='lines+markers',
                                 name=f'p={p:g}, slope={rate:.2f}'))
    fig.update_layout(title='Error of nonlinear flux discretization',
                      xaxis_title='grid step h', yaxis_title='max nodal error',
                      xaxis_type='log', yaxis_type='log')
    save(fig, 'pde26_zbieznosc.html')


def make_curvature():
    ns = np.array((20, 40, 80, 160, 320, 640))
    fig = go.Figure()
    for p in (2.0, 3.0, 5.0):
        alpha = p / (p - 1)
        c = (p - 1) / p
        hs = 1 / ns
        # Exact centered second difference of the continuum solution.
        curvature = 2 * c * hs**(alpha - 2)
        assert curvature[-1] >= curvature[0] - 1e-12
        fig.add_trace(go.Scatter(x=hs, y=curvature, mode='lines+markers',
                                 name=f'p={p:g}, power={alpha-2:.3f}'))
    fig.update_layout(title='Growth of the centered second difference at x=1/2',
                      xaxis_title='h', yaxis_title='-D_h^2 u(1/2)',
                      xaxis_type='log', yaxis_type='log')
    save(fig, 'pde26_krzywizna.html')


if __name__ == '__main__':
    make_profiles()
    make_convergence()
    make_curvature()
    print('Discrete flux residuals and endpoint conditions verified.')
\end{lstlisting}


\section{Reakcja--dyfuzja: fale biegnące, energia i niestabilność Turinga (Plotly)}
\label{sec:pde27-code}
Program symuluje fale biegnące dla równania Fishera--KPP,
rozpad energii w równaniu Allena--Cahna i wzrost modów
w niestabilności Turinga. Wizualizacje ilustrują wyniki analityczne;
nie zastępują ich dowodów.
\begin{lstlisting}[style=pythonpde,caption={Fale reakcji--dyfuzji, energia i niestabilność Turinga (Plotly)},label={code:pde27-plotly}]
"""PDE 27: Fisher--KPP, fala biegnaca, Allen--Cahn i Turing.

Uruchomienie: pip install numpy plotly; python demo_reakcja_dyfuzja_plotly.py
Wyniki: pde27_fronty.html, pde27_fala_jawna.html,
        pde27_energia.html, pde27_turing.html
"""
from pathlib import Path
import numpy as np
import plotly.graph_objects as go

OUT = Path(__file__).resolve().parent


def save(fig, filename):
    fig.update_layout(template="plotly_white", font=dict(size=12),
                      legend=dict(orientation="h", y=-0.22))
    fig.write_html(OUT / filename, include_plotlyjs=True, full_html=True)


def reaction_kpp(u, dt, r=1.0):
    e = np.exp(r * dt)
    return u * e / (1 + u * (e - 1))


def split_step_kpp():
    n, L, dt, tmax = 1024, 120.0, 0.04, 16.0
    x = np.linspace(-L/2, L/2, n, endpoint=False)
    k = 2 * np.pi * np.fft.fftfreq(n, d=L/n)
    diffusion_half = np.exp(-k*k*dt/2)
    u = .8 * np.exp(-x*x/6)
    fig = go.Figure()
    fig.add_scatter(x=x, y=u, name="t=0")
    for j in range(1, round(tmax/dt)+1):
        u = np.fft.ifft(np.fft.fft(u) * diffusion_half).real
        u = reaction_kpp(np.clip(u, 0, 1), dt)
        u = np.fft.ifft(np.fft.fft(u) * diffusion_half).real
        u = np.clip(u, 0, 1)
        if j % 100 == 0:
            fig.add_scatter(x=x, y=u, name=f"t={j*dt:g}")
    assert np.min(u) >= 0 and np.max(u) <= 1
    fig.update_layout(title="Rownanie Fishera-KPP: rozprzestrzenianie danych skupionych",
                      xaxis_title="x", yaxis_title="u(t,x)")
    save(fig, 'pde27_fronty.html')


def explicit_front():
    z = np.linspace(-28, 28, 1601)
    w = np.exp(z / np.sqrt(6))
    u = (1+w)**-2
    up = -2*w/(np.sqrt(6)*(1+w)**3)
    upp = -w/(3*(1+w)**3)+w*w/(1+w)**4
    c = 5/np.sqrt(6)
    residual = upp+c*up+u*(1-u)
    assert np.max(np.abs(residual)) < 5e-15
    fig = go.Figure()
    for t in (0, 3, 6, 9):
        x = z+c*t
        fig.add_scatter(x=x, y=u, name=f"t={t}")
    fig.update_layout(title=f"Fala rownania Fishera-KPP: c=5/sqrt6~{c:.4f}, residuum < 5e-15",
                      xaxis_title="x", yaxis_title="U(x-ct)")
    save(fig, 'pde27_fala_jawna.html')
    return float(np.max(np.abs(residual)))


def reaction_allen(u, dt):
    e = np.exp(dt)
    return u*e/np.sqrt(1+u*u*(e*e-1))


def allen_cahn_energy():
    n, L, D, dt = 512, 20.0, .3, .02
    x = np.linspace(0, L, n, endpoint=False)
    k = 2*np.pi*np.fft.fftfreq(n, d=L/n)
    half = np.exp(-D*k*k*dt/2)
    u = .6*np.cos(2*np.pi*x/L)+.15*np.cos(6*np.pi*x/L)
    energy, time = [], []
    for j in range(601):
        if j % 5 == 0:
            ux = np.fft.ifft(1j*k*np.fft.fft(u)).real
            E = (L/n)*np.sum(D*ux**2/2+(u*u-1)**2/4)
            energy.append(E)
            time.append(j*dt)
        if j < 600:
            u = np.fft.ifft(np.fft.fft(u)*half).real
            u = reaction_allen(u, dt)
            u = np.fft.ifft(np.fft.fft(u)*half).real
    energy = np.asarray(energy)
    assert energy[-1] < energy[0] and np.all(np.isfinite(energy))
    fig = go.Figure(go.Scatter(x=time, y=energy, mode="lines+markers", name="E(u)"))
    fig.update_layout(title="Rownanie Allena-Cahna: dyssypacja energii (rozszczepienie Stranga)",
                      xaxis_title="t", yaxis_title="energia")
    save(fig, 'pde27_energia.html')
    return float(np.max(np.diff(energy)))


def turing_mode():
    mu = np.linspace(0, 1.1, 400)
    lam = []
    for m in mu:
        A = np.array([[1-m,-2],[2,-3-10*m]])
        lam.append(np.max(np.linalg.eigvals(A).real))
    lam = np.asarray(lam)
    at = np.linalg.eigvals(np.array([[.7,-2],[2,-6]])).real
    assert max(at) > 0
    assert np.all(lam[(mu > .2) & (mu < .5)] > 0)
    fig = go.Figure(go.Scatter(x=mu, y=lam, name="max Re lambda(J-muD)"))
    fig.add_hline(y=0, line_dash="dash")
    fig.add_vrect(x0=.2, x1=.5, fillcolor="lightgray", opacity=.18,
                  line_width=0, annotation_text="niestabilne mody")
    fig.update_layout(title="Niestabilnosc Turinga: stabilna reakcja, niestabilna dyfuzja",
                      xaxis_title="mu (wartosc wasna -Delta)", yaxis_title="tempo wzrostu")
    save(fig, 'pde27_turing.html')
    return float(max(at))


if __name__ == '__main__':
    split_step_kpp()
    print('maks. residuum jawnej fali:', explicit_front())
    print('maks. przyrost energii:', allen_cahn_energy())
    print('dodatnia wartosc wlasna Turinga:', turing_mode())
    print('Zapisano 4 wykresy Plotly w:', OUT)

\end{lstlisting}


\section{Równania Naviera--Stokesa: wir Taylora--Greena i rzut Leraya — Plotly}
\label{sec:pde28-code}
Poniższy program sprawdza dokładną energię wiru Taylora--Greena,
przedstawia wirowość i ciśnienie oraz oblicza rzut Leraya
w przestrzeni Fouriera. W modelowym polu dywergencja po rzutowaniu
jest bliska zeru z dokładnością maszynową. To eksperyment numeryczny,
nie zamiennik dowodu istnienia rozwiązania Leraya--Hopfa.
\begin{lstlisting}[style=pythonpde,caption={Wir Taylora--Greena i rzut Leraya (Plotly)},label={code:pde28-plotly}]
"""Rownania Naviera-Stokesa: wir Taylora-Greena i rzut Leraya.

pip install numpy plotly
python demo_navier_stokes_plotly.py
Wyniki: pde28_wiry.html, pde28_energia.html,
        pde28_cisnienie.html, pde28_rzut_helmholtza.html
"""
from pathlib import Path
import numpy as np
import plotly.graph_objects as go

OUT = Path(__file__).resolve().parent
N, NU = 128, 0.08
x = np.linspace(0, 2*np.pi, N, endpoint=False)
X, Y = np.meshgrid(x, x, indexing='ij')
k = np.fft.fftfreq(N, d=1/N)  # Fourier integers, not physical cycles
KX, KY = np.meshgrid(k, k, indexing='ij')


def store(fig, name):
    fig.update_layout(template='plotly_white', font=dict(size=13))
    fig.write_html(OUT / name, include_plotlyjs=True, full_html=True)


def taylor_green(t):
    decay = np.exp(-2*NU*t)
    ux = decay*np.sin(X)*np.cos(Y)
    uy = -decay*np.cos(X)*np.sin(Y)
    p = 0.25*decay**2*(np.cos(2*X)+np.cos(2*Y))
    omega = 2*decay*np.sin(X)*np.sin(Y)
    return ux, uy, p, omega


def divergence_fft(vx, vy):
    return np.fft.ifft2(1j*KX*np.fft.fft2(vx)
                          +1j*KY*np.fft.fft2(vy)).real


def project(vx, vy):
    hx, hy = np.fft.fft2(vx), np.fft.fft2(vy)
    den = KX*KX+KY*KY
    safe = np.where(den == 0, 1, den)
    dot = KX*hx+KY*hy
    hx = hx-KX*dot/safe
    hy = hy-KY*dot/safe
    hx[0,0]=hy[0,0]=0
    return np.fft.ifft2(hx).real, np.fft.ifft2(hy).real


def plots():
    u1,u2,p,omega = taylor_green(0.5)
    fig = go.Figure(go.Contour(x=x,y=x,z=omega.T, colorscale='RdBu',
                 contours=dict(showlabels=True)))
    fig.update_layout(title='Wir Taylora-Greena: wirowosc w chwili t=0.5',
                      xaxis_title='x',yaxis_title='y', yaxis_scaleanchor='x')
    store(fig, 'pde28_wiry.html')

    times=np.linspace(0,10,150)
    e_exact=np.pi**2*np.exp(-4*NU*times)
    e_numeric=np.array([0.5*(taylor_green(t)[0]**2+
                           taylor_green(t)[1]**2).mean()*(2*np.pi)**2
                        for t in times])
    assert np.max(np.abs(e_exact-e_numeric))<2e-12
    fig=go.Figure()
    fig.add_trace(go.Scatter(x=times,y=e_numeric,name='Energia z siatki',mode='markers',
                             marker=dict(size=4)))
    fig.add_trace(go.Scatter(x=times,y=e_exact,name='Energia dokladna',mode='lines'))
    fig.update_layout(title='Tozsamosc energii: E(t)=pi^2 exp(-4 nu t)',
                      xaxis_title='czas t',yaxis_title='E(t)')
    store(fig,'pde28_energia.html')

    fig=go.Figure(go.Heatmap(x=x,y=x,z=p.T,colorscale='RdBu'))
    fig.update_layout(title='Cisnienie wiru Taylora-Greena w chwili t=0.5',
                      xaxis_title='x',yaxis_title='y',yaxis_scaleanchor='x')
    store(fig,'pde28_cisnienie.html')

    fx=np.sin(Y)+np.cos(X)
    fy=np.sin(X)+np.cos(Y)
    px,py=project(fx,fy)
    assert np.max(np.abs(px-np.sin(Y)))<2e-12
    assert np.max(np.abs(py-np.sin(X)))<2e-12
    div_before=np.max(np.abs(divergence_fft(fx,fy)))
    div_after=np.max(np.abs(divergence_fft(px,py)))
    assert div_before>1.9 and div_after<1e-10
    assert np.max(np.abs(divergence_fft(u1,u2)))<1e-10
    fig=go.Figure()
    step=9
    for vx,vy,name,dx in [(fx,fy,'F przed rzutem',0), (px,py,'PF po rzucie',2*np.pi+1)]:
        first=True
        for i in range(0,N,step):
            for j in range(0,N,step):
                xi,yi=x[i]+dx,x[j]
                fig.add_trace(go.Scatter(x=[xi,xi+0.27*vx[i,j],None],
                      y=[yi,yi+0.27*vy[i,j],None],mode='lines',
                      line=dict(width=1.1),name=name,showlegend=first))
                first=False
    fig.update_layout(title=f'Rzut Leraya: max div przed={div_before:.2f}, po={div_after:.1e}',
                      xaxis_title='dwa pola obok siebie',yaxis_title='y',
                      yaxis_scaleanchor='x')
    store(fig,'pde28_rzut_helmholtza.html')
    print('Energia: maksymalny blad', np.max(np.abs(e_exact-e_numeric)))
    print('Dywergencja przed i po rzucie',div_before,div_after)


if __name__ == '__main__':
    plots()
\end{lstlisting}


\section{Zagadnienie przeszkody i reakcja kontaktowa — Plotly}
\label{sec:pde29-code}
Program implementuje metodę Gaussa--Seidela z projekcją
na zbiór funkcji spełniających przeszkodę. Porównuje wyniki
z dokładnym rozwiązaniem z przykładu~\ref{ex:pde29-explicit}
i kontroluje nieujemność reakcji oraz warunek komplementarności.
Listing jest obramowany; samodzielny plik \texttt{demo\_przeszkoda\_plotly.py}
znajduje się również w pakiecie projektu.
\begin{lstlisting}[style=pythonpde,caption={Przeszkoda, reakcja i zbieżność schematu (Plotly)},label={code:pde29-plotly}]
"""Rozdzial 29 PDE: metoda Gaussa-Seidela z projekcja dla przeszkody.

Wymagania: pip install numpy plotly
Wykonanie: python demo_przeszkoda_plotly.py
Wyniki: pde29_przeszkoda.html, pde29_reakcja.html,
        pde29_zbieznosc.html
"""
from pathlib import Path
import numpy as np
import plotly.graph_objects as go

OUT = Path(__file__).resolve().parent


def obstacle(x: np.ndarray) -> np.ndarray:
    return (1.0 - x * x) ** 2


def exact(x: np.ndarray) -> np.ndarray:
    a = np.abs(x)
    return np.where(a <= 1 / 3, (1.0 - x * x) ** 2,
                    32 / 27 * (1.0 - a))


def projected_gauss_seidel(n: int, omega=1.75, tol=2e-12,
                            max_sweeps=18000):
    h = 2.0 / n
    x = np.linspace(-1.0, 1.0, n + 1)
    psi = obstacle(x)
    u = psi.copy()
    u[0] = u[-1] = 0
    history = []
    for k in range(max_sweeps):
        change = 0.0
        for i in range(1, n):
            new = max(psi[i], u[i] + omega *
                      (0.5 * (u[i - 1] + u[i + 1]) - u[i]))
            change = max(change, abs(new - u[i]))
            u[i] = new
        if k % 25 == 0:
            energy = np.sum(np.diff(u) ** 2) / (2 * h)
            history.append((k, energy, change))
        if change < tol:
            break
    else:
        raise RuntimeError(f'Projected GS did not converge for n={n}')
    reaction = (2 * u[1:-1] - u[:-2] - u[2:]) / (h * h)
    slack = u[1:-1] - psi[1:-1]
    return x, u, psi, reaction, slack, np.array(history), k + 1


def save(fig, name):
    fig.update_layout(template='plotly_white', font=dict(size=13))
    fig.write_html(OUT / name, include_plotlyjs=True, full_html=True)


def main():
    ns = [40, 80, 160]
    results = [projected_gauss_seidel(n) for n in ns]
    errors = []
    for n, (x, u, psi, reaction, slack, hist, sweeps) in zip(ns, results):
        err = np.max(np.abs(u - exact(x)))
        errors.append(err)
        assert np.min(slack) > -1e-11, (n, np.min(slack))
        assert np.min(reaction) > -2e-7, (n, np.min(reaction))
        assert np.max(np.abs(reaction * slack)) < 5e-7
        print(f'N={n}, sweeps={sweeps}, max error={err:.4e}, '
              f'min reaction={reaction.min():.2e}, '
              f'max complementarity={np.max(np.abs(reaction * slack)):.2e}')

    x, u, psi, reaction, slack, hist, _ = results[-1]
    xf = np.linspace(-1, 1, 901)
    fig = go.Figure()
    fig.add_trace(go.Scatter(x=xf, y=obstacle(xf),
                             name='Przeszkoda psi', mode='lines'))
    fig.add_trace(go.Scatter(x=xf, y=exact(xf),
                             name='Rozwi\u0105zanie dok\u0142adne', mode='lines'))
    fig.add_trace(go.Scatter(x=x, y=u,
                             name='Wynik numeryczny', mode='markers',
                             marker=dict(size=5)))
    for point in (-1 / 3, 1 / 3):
        fig.add_vline(x=point, line_dash='dash', line_color='gray')
    fig.update_layout(title='Zagadnienie przeszkody: kontakt i swobodny brzeg',
                      xaxis_title='x', yaxis_title='Wysoko\u015b\u0107')
    save(fig, 'pde29_przeszkoda.html')

    xnodes = x[1:-1]
    reaction_exact = np.where(np.abs(xnodes) < 1 / 3,
                              4 - 12 * xnodes * xnodes, 0.0)
    fig = go.Figure()
    fig.add_trace(go.Scatter(x=xnodes, y=reaction,
                             mode='lines+markers',
                             name='Reakcja dyskretna',
                             marker=dict(size=4)))
    fig.add_trace(go.Scatter(x=xnodes, y=reaction_exact,
                             name='Reakcja dok\u0142adna', mode='lines'))
    fig.update_layout(title='Reakcja przeszkody: -u" >= 0',
                      xaxis_title='x', yaxis_title='Mno\u017cnik ograniczenia')
    save(fig, 'pde29_reakcja.html')

    fig = go.Figure()
    fig.add_trace(go.Scatter(x=np.array(ns), y=errors, mode='lines+markers',
                             name='Maksymalny b\u0142\u0105d w\u0119z\u0142owy'))
    fig.update_layout(title='Zbie\u017cno\u015b\u0107 metody z projekcj\u0105',
                      xaxis_type='log', yaxis_type='log',
                      xaxis_title='Liczba przedzia\u0142\xf3w N',
                      yaxis_title='B\u0142\u0105d maksimum')
    save(fig, 'pde29_zbieznosc.html')
    if all(err > 0 for err in errors):
        orders = np.log(np.array(errors[:-1]) / np.array(errors[1:])) / np.log(2)
        print('Obserwowane rz\u0119dy zbie\u017cno\u015bci:', orders)


if __name__ == '__main__':
    main()
\end{lstlisting}

\section{Homogenizacja periodyczna: rozwiązanie i strumień — Plotly}
\label{sec:pde30-code}
Program porównuje dokładne rozwiązania jednowymiarowe z równaniem
uśrednionym, pokazuje oscylacje gradientu oraz zbieżność błędu.
Ostatni wykres przedstawia anizotropię macierzy efektywnej
materiału dwuwarstwowego; jest ilustracją prawa
~\eqref{eq:pde30-layers}, a nie obliczeniem metodą elementów skończonych.
\begin{lstlisting}[style=pythonpde,caption={Homogenizacja okresowa i macierz efektywna (Plotly)},label={code:pde30-plotly}]
"""PDE, rozdzial 30: 1D homogenizacja okresowa i srednie efektywne.

Wymagania: pip install numpy plotly
Uruchomienie: python demo_homogenizacja_plotly.py
Pliki wynikowe: pde30_profile.html, pde30_gradient_strumien.html,
                pde30_zbieznosc.html, pde30_anizotropia.html
"""
from pathlib import Path
import numpy as np
import plotly.graph_objects as go

DIR = Path(__file__).resolve().parent
A_HOM = np.sqrt(3.0)


def cumulative_trapezoid(v, h):
    return np.concatenate(([0.0], np.cumsum((v[1:] + v[:-1]) * (h / 2.0))))


def micro_solution(periods, samples_per_period=240):
    eps = 1.0 / periods
    x = np.linspace(0.0, 1.0, periods * samples_per_period + 1)
    h = x[1] - x[0]
    a = 2.0 + np.sin(2.0 * np.pi * x / eps)
    b = 1.0 / a
    ib = cumulative_trapezoid(b, h)
    ibx = cumulative_trapezoid(b * x, h)
    c = ibx[-1] / ib[-1]
    u = c * ib - ibx
    gradient = b * (c - x)
    flux = c - x
    u0 = x * (1.0 - x) / (2.0 * A_HOM)
    grad0 = (0.5 - x) / A_HOM
    err = np.sqrt(np.trapezoid((u - u0) ** 2, x))
    return x, u, gradient, flux, u0, grad0, err


def export(fig, name):
    fig.update_layout(template="plotly_white", font=dict(size=13))
    fig.write_html(DIR / name, full_html=True, include_plotlyjs=True)


def main():
    first = go.Figure()
    for n in [4, 8, 16, 32]:
        x, u, du, q, u0, du0, err = micro_solution(n)
        skip = max(1, len(x) // 2000)
        first.add_trace(go.Scatter(x=x[::skip], y=u[::skip], name=f"eps=1/{n}", mode="lines"))
    first.add_trace(go.Scatter(x=x, y=u0, name="rozwiazanie graniczne", line=dict(color="black", dash="dash")))
    first.update_layout(title="Homogenizacja: rozwiazania mikroskopowe i graniczne",
                        xaxis_title="x", yaxis_title="u(x)")
    export(first, "pde30_profile.html")

    x, u, grad, flux, u0, grad0, err = micro_solution(24)
    fig = go.Figure()
    fig.add_trace(go.Scatter(x=x, y=grad, name="gradient u'eps"))
    fig.add_trace(go.Scatter(x=x, y=grad0, name="gradient u'0", line=dict(dash="dot")))
    fig.add_trace(go.Scatter(x=x, y=flux, name="strumien aepsu'eps", line=dict(width=3)))
    fig.add_trace(go.Scatter(x=x, y=0.5-x, name="strumien graniczny", line=dict(dash="dash")))
    fig.update_layout(title="Silna zbieznosc strumienia, oscylacje gradientu",
                      xaxis_title="x", yaxis_title="wartosc")
    export(fig, "pde30_gradient_strumien.html")

    ns = np.array([4, 8, 16, 32, 64], dtype=int)
    errors = np.array([micro_solution(int(n))[6] for n in ns])
    order = float(np.polyfit(np.log(1.0/ns), np.log(errors), 1)[0])
    fig = go.Figure()
    fig.add_trace(go.Scatter(x=1.0/ns, y=errors, mode="lines+markers", name="blad L^2"))
    fig.add_trace(go.Scatter(x=1.0/ns, y=errors[0]*(ns[0]/ns),
                             mode="lines", name="wzorzec O(eps)", line=dict(dash="dash")))
    fig.update_layout(title=f"Zbieznosc do rozwiazania granicznego; rzad ~{order:.3f}",
                      xaxis_type="log", yaxis_type="log",
                      xaxis_title="eps", yaxis_title="||ueps-u0|| L^2")
    export(fig, "pde30_zbieznosc.html")

    t = np.linspace(0, 2*np.pi, 360)
    fig = go.Figure()
    fig.add_trace(go.Scatter(x=np.sqrt(8.0/5.0)*np.cos(t),
                             y=np.sqrt(5.0/2.0)*np.sin(t),
                             name="warstwy: diag(8/5,5/2)"))
    fig.add_trace(go.Scatter(x=np.sqrt(2.5)*np.cos(t),
                             y=np.sqrt(2.5)*np.sin(t),
                             name="bledna srednia arytmetyczna", line=dict(dash="dash")))
    fig.update_yaxes(scaleanchor="x", scaleratio=1)
    fig.update_layout(title="Izotropowe fazy -> anizotropowa macierz efektywna",
                      xaxis_title="kierunek prostopadly do warstw (y1)",
                      yaxis_title="kierunek wzdluz warstw (y2)")
    export(fig, "pde30_anizotropia.html")
    assert abs(A_HOM - 1.7320508075688772) < 1e-12
    assert all(np.diff(errors) < 0), errors
    assert 0.8 < order < 1.2, order
    assert abs(flux[-1] - flux[0] + 1.0) < 1e-12
    assert abs(u[-1]) < 5e-12
    print("Efektywny wspolczynnik:", A_HOM)
    print("Rzad zbieznosci L2:", order)
    print("Bledy L2:", errors)


if __name__ == "__main__":
    main()
\end{lstlisting}

\section{Analiza mikrolokalna i symbole: eksperymenty Plotly}
\label{sec:pde31-code}
Rozkład częstości dla wąskiej warstwy, stożek charakterystyczny i parametryks.
\begin{lstlisting}[style=pythonpde,caption={Eksperymenty rozdziału 31 (Python i Plotly)},label={code:pde31-plotly}]
"""Rozdzial 31: kierunki osobliwosci i symbole Fourierowskie.
Wymaga numpy, plotly; zapisuje trzy interaktywne wykresy HTML.
Przyblizenie delta(x1) jest gladkie dla skonczonej szerokosci;
wykres ilustruje, a nie dowodzi twierdzenia o zbiorze falowym.
"""
from pathlib import Path
import numpy as np
import plotly.graph_objects as go

OUT = Path(__file__).resolve().parent


def save(fig, filename):
    fig.update_layout(template="plotly_white", font=dict(size=13))
    fig.write_html(OUT / filename, include_plotlyjs=True, full_html=True)


def directional_singularity():
    n = 256
    x = np.linspace(-np.pi, np.pi, n, endpoint=False)
    dx = x[1] - x[0]
    xx, yy = np.meshgrid(x, x, indexing="ij")
    angle = np.pi / 5
    normal = xx * np.cos(angle) + yy * np.sin(angle)
    # Okresowy profil waskiego pasa; zwezone przyblizenie delty.
    profile = np.exp(-0.5 * (np.sin(normal / 2) / 0.075) ** 2)
    window = np.exp(-0.5 * ((xx / 2.2) ** 8 + (yy / 2.2) ** 8))
    img = profile * window
    freq = np.fft.fftshift(np.fft.fftfreq(n, d=dx)) * 2 * np.pi
    amp = np.log1p(np.abs(np.fft.fftshift(np.fft.fft2(img))))
    fig = go.Figure(go.Heatmap(x=freq, y=freq, z=amp.T, colorscale="Viridis"))
    fig.update_layout(title="Kierunki wysokich czestotliwosci: waski pas",
                      xaxis_title="xi_1", yaxis_title="xi_2")
    fig.update_xaxes(range=[-45, 45]); fig.update_yaxes(range=[-45, 45])
    save(fig, "pde31_kierunki_fouriera.html")
    assert np.isfinite(amp).all()


def characteristic_cone():
    xi = np.linspace(-5, 5, 801)
    fig = go.Figure()
    fig.add_trace(go.Scatter(x=xi, y=2 * np.abs(xi), name="tau = +2|xi|"))
    fig.add_trace(go.Scatter(x=xi, y=-2 * np.abs(xi), name="tau = -2|xi|"))
    fig.update_layout(title="Zbior charakterystyczny rownania falowego, c=2",
                      xaxis_title="xi", yaxis_title="tau")
    save(fig, "pde31_charakterystyki.html")
    assert np.max(np.abs((2 * np.abs(xi)) ** 2 - 4 * xi ** 2)) < 1e-10


def parametrix_filter():
    xi = np.linspace(-20, 20, 1201)
    p = 1 + xi ** 2
    inverse = 1 / p
    fig = go.Figure()
    fig.add_trace(go.Scatter(x=xi, y=p, name="P(i xi)=1+xi^2"))
    fig.add_trace(go.Scatter(x=xi, y=inverse, name="q(xi)=1/(1+xi^2)"))
    fig.update_layout(title="Symbol eliptyczny i odwrotny mnoznik",
                      xaxis_title="xi", yaxis_title="wartosc")
    save(fig, "pde31_parametryks.html")
    assert np.allclose(p * inverse, 1)


if __name__ == "__main__":
    directional_singularity()
    characteristic_cone()
    parametrix_filter()
    print("Mikrolokalna: zapisano 3 wykresy Plotly.")
\end{lstlisting}


\section{Problemy odwrotne i sterowanie optymalne: eksperymenty Plotly}
\label{sec:pde32-code}
Zanik wartości singularnych, regularyzacja Tichonowa i sterowanie.
\begin{lstlisting}[style=pythonpde,caption={Eksperymenty rozdziału 32 (Python i Plotly)},label={code:pde32-plotly}]
"""Rozdzial 32: niestabilnosc, regularyzacja Tichonowa i sterowanie PDE.
Wymaga numpy, plotly. Zawiera testy formul dla modelu spektralnego.
"""
from pathlib import Path
import numpy as np
import plotly.graph_objects as go

OUT = Path(__file__).resolve().parent


def save(fig, filename):
    fig.update_layout(template="plotly_white", font=dict(size=13))
    fig.write_html(OUT / filename, include_plotlyjs=True, full_html=True)


def singular_values():
    n = np.arange(1, 41)
    sig_ell = 1 / (1 + n ** 2)
    sig_heat = np.exp(-0.01 * n ** 2)
    fig = go.Figure()
    fig.add_trace(go.Scatter(x=n, y=sig_ell, mode="lines+markers", name="(-Delta+I)^(-1)"))
    fig.add_trace(go.Scatter(x=n, y=sig_heat, mode="lines+markers", name="exp(T Delta), T=0.01"))
    fig.update_layout(title="Male wartosci singularne dwoch operatorow PDE",
                      xaxis_title="numer modu", yaxis_title="wartosc singularna",
                      yaxis_type="log")
    save(fig, "pde32_wartosci_singularne.html")


def inverse_heat():
    rng = np.random.default_rng(713)
    n = np.arange(1, 51)
    T = 0.007
    sigma = np.exp(-T * n ** 2)
    truth = np.exp(-n / 7) * np.sin(n / 3)
    delta = 0.002
    noise = rng.normal(0, 1, n.size)
    noise = delta * noise / np.linalg.norm(noise)
    data = sigma * truth + noise
    alpha = np.logspace(-10, -2, 130)
    recon = sigma[:, None] * data[:, None] / (sigma[:, None] ** 2 + alpha[None, :])
    error = np.linalg.norm(recon - truth[:, None], axis=0)
    idx = int(np.argmin(error)); alpha_best = alpha[idx]
    fig = go.Figure()
    fig.add_trace(go.Scatter(x=alpha, y=error, name="blad rekonstrukcji"))
    fig.add_vline(x=alpha_best, line_dash="dash",
                  annotation_text="minimum dla danych testowych")
    fig.update_layout(title="Regularyzacja Tichonowa: kompromis szum-gladzenie",
                      xaxis_title="alpha", xaxis_type="log",
                      yaxis_title="blad wspolczynnikow")
    save(fig, "pde32_regularyzacja.html")
    fig2 = go.Figure()
    fig2.add_trace(go.Scatter(x=n, y=truth, name="sygnal prawdziwy"))
    fig2.add_trace(go.Scatter(x=n, y=recon[:, idx], name="rekonstrukcja"))
    fig2.update_layout(title=f"Odtwarzanie modow, alpha={alpha_best:.2e}",
                       xaxis_title="numer modu", yaxis_title="wspolczynnik")
    save(fig2, "pde32_rekonstrukcja.html")
    assert error[idx] <= error[0] and error[idx] <= error[-1]
    return float(alpha_best), float(error[idx])


def elliptic_control():
    x = np.linspace(0, np.pi, 601)
    alpha = 0.15
    n = np.array([1.0, 3.0])
    desired_coeff = np.array([1.0, 0.35])
    sigma = 1 / (n ** 2)
    control_coeff = sigma * desired_coeff / (sigma ** 2 + alpha)
    observed_coeff = sigma * control_coeff
    sin_modes = np.sin(n[:, None] * x)
    y_desired = desired_coeff @ sin_modes
    y_state = observed_coeff @ sin_modes
    control = control_coeff @ sin_modes
    fig = go.Figure()
    for values, name in ((y_desired, "stan pozadany"),
                         (y_state, "stan optymalny"), (control, "sterowanie")):
        fig.add_trace(go.Scatter(x=x, y=values, name=name))
    fig.update_layout(title="Sterowanie rozlozone dla -y''=u",
                      xaxis_title="x", yaxis_title="wartosc")
    save(fig, "pde32_sterowanie.html")
    assert np.allclose((sigma ** 2 + alpha) * control_coeff,
                       sigma * desired_coeff)


if __name__ == "__main__":
    singular_values()
    a, err = inverse_heat()
    elliptic_control()
    print(f"Optymalne testowe alpha={a:.3e}, blad={err:.4f}; zapisano 4 wykresy.")
\end{lstlisting}



\end{document}
